\documentclass[10pt, ukrainian]{article}
%\setlength\textwidth{190mm} \setlength\textheight{277mm}
%\setlength\voffset{-38mm} \setlength\hoffset{-38mm}
%\setlength\parindent{0mm}
\pagestyle{empty}
\usepackage[T2A]{fontenc}
\usepackage[cp1251]{inputenc}
\usepackage[ukrainian]{babel}
\usepackage{graphicx}
\usepackage{caption}
\DeclareCaptionFormat{nocaption}{}
\captionsetup{format=nocaption,aboveskip=0pt,belowskip=0pt}
\usepackage[Export]{adjustbox} 
\adjustboxset{max size={0.9\linewidth}{0.9\paperheight}}
\begin{document}
%begin
{{\begin {small}\par
{ \center  Київський національний університет імені Тараса Шевченка \par}
{\parbox {5 cm}{Освітня програма \hfill}{\parbox {7 cm}{Прикладна математика, Системний аналіз \hfill}}\parbox {6 cm}{\hfill Семестр 2}}\par
{\parbox {5 cm} {Навчальний предмет \hfill}\parbox {7 cm}{Програмування \hfill}\parbox {6cm}{\hfill}\par}\vspace{-7pt} \end{small}}{\center Екзаменаційний білет \No \textbf 
{ 1 }\par}}
{
\begin {large}
1. Що буде виведено за виконання?
code
try:
    res = 8/8*3*9
    if isinstance(res, bool):
        print(f'bool;{res}')
    elif isinstance(res, int):
        print(f'int;{res}')
    elif isinstance(res, float):
        print(f'float;{res:.2f}')
    else:
        print('???')
except:
    print('ERR')
code

2. Що буде виведено за виконання?
code
a = 1
b = 9
c = 5

def f(b):
    a = 3
    b = 5
    c = 4
    return a + b + c

a = 4
b = 3
c = 7

try:
    print(f(b), a, b, c, sep=';')
except:
    print('ERR', a, b, c, sep=';')
code

3. Що буде виведено за виконання?
code
try:
    m = 3
    while m<=6:
        m += 1
    print(m, end=';')
except:
    print('ERR')
code

4. Що буде виведено за виконання?
code
try:
    for b in range(6, -15, 3):
        if b<=5:
            continue
        print(b, end=';');b = 6
    else:
        print(b,end=';')
    print(b)
except:
    print('ERR')
code

5. Що буде виведено за виконання?\par (Дійсні значення, якщо наявні, в цьому завданні будуть виводитись у форматі з фіксованою крапкою та, можливо, з доволі великою кількістю десяткових розрядів. У відповіді для дійсних значень у ЦЬОМУ завданні достатньо вказати один десятковий розряд після крапки, відкинувши решту розрядів без округлення. Наприклад, замість 13.66666666666667 достатньо вказати 13.6, замість 25.23 достатньо вказати 25.2)
code
try:
    print(4, end=';')
    print(int('3'), end=';')
    print(9, end=';')
except ValueError:
    print(8, end=';')
except Exception:
    print(1, end=';')
except:
    print('ERR', end=';')
else:
    print(5, end=';')
finally:
    print(7)
code

6. Що буде виведено за виконання?
code
def f(n):
    if n<=9:
        print(n%10, end='')
    else:
        f(n//10)

f(123456)
code

7. Що буде виведено за виконання?
code
class A:
    b = 1

    def _g1(self): print(2, end=';')
    def _g2(self): print(3, end=';')

    def main(self):
        try: print(self.b, end=';')
        except: print('ERR', end=';')

        try: self._g1()
        except: print('ERR', end=';')

        try: self._g2()
        except: print('ERR', end=';')


class B(A):
    b = 4

    def _g1(self): print(5, end=';')
    def _g2(self): print(6, end=';')

try: B().main()
except: print('ERR', end=';')
code

8. Що буде виведено за виконання?
code
def f(arg, n):
    arg[78] = 0
    n = 7
    arg = set(arg.keys())

try:
    n = 0
    d = {28:4, 78:5, 28:7}
    f(d, n)
    print(n, end=';')
    for x in d.items():
        print(x, sep=';', end=';')
except:
    print('ERR')
code

9. Що буде виведено за виконання?
code
class A:
    b = 1
    c = 2

    def __init__(self):
        b = 3

class B(A):
    b = 4

    def __init__(self):
        super().__init__()
        self.a = 7

obj = B()
obj.c = obj.b + 10
B.c = 13

try: print(obj.c, end= ';')
except: print('ERR', end=';')

try: print(A.c, end= ';')
except: print('ERR', end=';')

try: print(B.c, end= ';')
except: print('ERR', end=';')
code

10. 
Для графа на рисунку з вершини 2 запущено алгоритм пошуку в глибину, що будує кістякове дерево. Списки суміжності вершин переглядаються алгоритмом за зростанням міток вершин.\par
\begin{center}
	\adjustimage{max size={0.9\linewidth}{0.9\paperheight}}
	{../10/Traversals/73.png}
\end{center}
\par Які з наведених ребер входять до складу отриманого кістякового дерева?\par
1) (0, 6)\par
2) (3, 6)\par
3) (4, 6)\par
4) (2, 4)\par
5) (1, 2)\par
6) (0, 1)\par
{\vskip 15pt} У формі вибрати номери відповідних ребер.\par
%answer:2 3 5 6
\end {large}}
%end
{\smallskip \begin {small} Затверджено на засіданні кафедри теоретичної кібернетики, протокол \No 12 від   19.05.2020 р.\par\smallskip\smallskip
{\parbox {6.5 cm}{Зав. кафедри \dotfill Ю.В.~Крак}}\hspace {0.7cm}{\hbox to 6.5 cm{ Екзаменатор \dotfill Т.О.~Карнаух}} \end{small}}
%begin
{{\begin {small}\par
{ \center  Київський національний університет імені Тараса Шевченка \par}
{\parbox {5 cm}{Освітня програма \hfill}{\parbox {7 cm}{Прикладна математика, Системний аналіз \hfill}}\parbox {6 cm}{\hfill Семестр 2}}\par
{\parbox {5 cm} {Навчальний предмет \hfill}\parbox {7 cm}{Програмування \hfill}\parbox {6cm}{\hfill}\par}\vspace{-7pt} \end{small}}{\center Екзаменаційний білет \No \textbf 
{ 2 }\par}}
{
\begin {large}
1. Що буде виведено за виконання?
code
try:
    res = 4**False+2
    if isinstance(res, bool):
        print(f'bool;{res}')
    elif isinstance(res, int):
        print(f'int;{res}')
    elif isinstance(res, float):
        print(f'float;{res:.2f}')
    else:
        print('???')
except:
    print('ERR')
code

2. Що буде виведено за виконання?
code
a = 8
b = 2
c = 0

def h(a):
    global c
    a = 1
    b = 2
    c = 4
    return a + b + c

a = 6
b = 5
c = 0

try:
    print(h(b), a, b, c, sep=';')
except:
    print('ERR', a, b, c, sep=';')
code

3. Що буде виведено за виконання?
code
try:
    k = 4
    while k<8:
        k += 2
    print(k, end=';')
except:
    print('ERR')
code

4. Що буде виведено за виконання?
code
try:
    for a in range(-13, 1, -2):
        if a<8:
            break
        print(a, end=';')
        if a>5:
            break
    else:
        print(a,end=';')
    print(a)
except:
    print('ERR')
code

5. Що буде виведено за виконання?\par (Дійсні значення, якщо наявні, в цьому завданні будуть виводитись у форматі з фіксованою крапкою та, можливо, з доволі великою кількістю десяткових розрядів. У відповіді для дійсних значень у ЦЬОМУ завданні достатньо вказати один десятковий розряд після крапки, відкинувши решту розрядів без округлення. Наприклад, замість 13.66666666666667 достатньо вказати 13.6, замість 25.23 достатньо вказати 25.2)
code
try:
    print(2, end=';')
    print(0%0.0, end=';')
    print(6, end=';')
except TypeError:
    print(4, end=';')
except ZeroDivisionError:
    print(8, end=';')
except:
    print('ERR', end=';')
else:
    print(6, end=';')
finally:
    print(1)
code

6. Що буде виведено за виконання?
code
def f(s):
    for i in range(len(s)):
        s[i] += 1
        return
 
s = [1, 2, 3]
f(s)
print(s)
code

7. Що буде виведено за виконання?
code
class A:
    a = 1

    def _f1(self): print(2, end=';')
    def __f2(self): print(3, end=';')

    def main(self):
        try: print(self.a, end=';')
        except: print('ERR', end=';')

        try: self._f1()
        except: print('ERR', end=';')

        try: self.__f2()
        except: print('ERR', end=';')


class B(A):
    a = 4

    def _f1(self): print(5, end=';')
    def __f2(self): print(6, end=';')

try: B().main()
except: print('ERR', end=';')
code

8. Що буде виведено за виконання?
code
def f(arg, n):
    arg[37] = 8
    n = 9
    arg = tuple(arg.items())

try:
    n = 2
    t = {37:3, 16:4, 57:9, 16:7}
    f(t, n)
    print(n, end=';')
    for x in t.keys():
        print(x, sep=';', end=';')
except:
    print('ERR')
code

9. Що буде виведено за виконання?
code
class A:
    a = 1
    b = 2

    def __init__(self):
        b = 3

class B(A):
    b = 4

    def __init__(self):
        super().__init__()
        self.c = 7

obj = B()
obj.a = obj.a + 10
B.a = 13

try: print(obj.a, end= ';')
except: print('ERR', end=';')

try: print(A.a, end= ';')
except: print('ERR', end=';')

try: print(B.a, end= ';')
except: print('ERR', end=';')
code

10. 
Для графа на рисунку з вершини 0 запущено алгоритм пошуку в глибину, що будує кістякове дерево. Списки суміжності вершин переглядаються алгоритмом за зростанням міток вершин.\par
\begin{center}
	\adjustimage{max size={0.9\linewidth}{0.9\paperheight}}
	{../10/Traversals/79.png}
\end{center}
\par Які з наведених ребер входять до складу отриманого кістякового дерева?\par
1) (4, 5)\par
2) (1, 5)\par
3) (3, 4)\par
4) (0, 2)\par
5) (0, 1)\par
6) (1, 4)\par
{\vskip 15pt} У формі вибрати номери відповідних ребер.\par
%answer:3 5 6
\end {large}}
%end
{\smallskip \begin {small} Затверджено на засіданні кафедри теоретичної кібернетики, протокол \No 12 від   19.05.2020 р.\par\smallskip\smallskip
{\parbox {6.5 cm}{Зав. кафедри \dotfill Ю.В.~Крак}}\hspace {0.7cm}{\hbox to 6.5 cm{ Екзаменатор \dotfill Т.О.~Карнаух}} \end{small}}
%begin
{{\begin {small}\par
{ \center  Київський національний університет імені Тараса Шевченка \par}
{\parbox {5 cm}{Освітня програма \hfill}{\parbox {7 cm}{Прикладна математика, Системний аналіз \hfill}}\parbox {6 cm}{\hfill Семестр 2}}\par
{\parbox {5 cm} {Навчальний предмет \hfill}\parbox {7 cm}{Програмування \hfill}\parbox {6cm}{\hfill}\par}\vspace{-7pt} \end{small}}{\center Екзаменаційний білет \No \textbf 
{ 3 }\par}}
{
\begin {large}
1. Що буде виведено за виконання?
code
try:
    res = -3**2--1
    if isinstance(res, bool):
        print(f'bool;{res}')
    elif isinstance(res, int):
        print(f'int;{res}')
    elif isinstance(res, float):
        print(f'float;{res:.2f}')
    else:
        print('???')
except:
    print('ERR')
code

2. Що буде виведено за виконання?
code
a = 4
b = 2
c = 0

def f(b):
    global c
    a *= 1
    b = 4
    c = 2
    return a + b + c

a = 6
b = 3
c = 7

try:
    print(f(a), a, b, c, sep=';')
except:
    print('ERR', a, b, c, sep=';')
code

3. Що буде виведено за виконання?
code
try:
    j = 6
    while j<9:
        j += 1
    print(j, end=';')
except:
    print('ERR')
code

4. Що буде виведено за виконання?
code
try:
    for c in range(-9, 11, -1):
        if c>=4:
            continue
        print(c, end=';')
    else:
        print(c,end=';')
    print(c)
except:
    print('ERR')
code

5. Що буде виведено за виконання?\par (Дійсні значення, якщо наявні, в цьому завданні будуть виводитись у форматі з фіксованою крапкою та, можливо, з доволі великою кількістю десяткових розрядів. У відповіді для дійсних значень у ЦЬОМУ завданні достатньо вказати один десятковий розряд після крапки, відкинувши решту розрядів без округлення. Наприклад, замість 13.66666666666667 достатньо вказати 13.6, замість 25.23 достатньо вказати 25.2)
code
try:
    print(3, end=';')
    print(7/1, end=';')
    print(2, end=';')
except KeyboardInterrupt:
    print(0, end=';')
except ZeroDivisionError:
    print(5, end=';')
except:
    print('ERR', end=';')
else:
    print(9, end=';')
finally:
    print(4)
code

6. Що буде виведено за виконання?
code
def f(n):
    print(n%10,end='')
    if n>9:
        f(n//100)
 
f(123456)
code

7. Що буде виведено за виконання?
code
class A:
    d = 1

    def __h1(self): print(2, end=';')
    def _h2(self): print(3, end=';')

    def main(self):
        try: print(self.d, end=';')
        except: print('ERR', end=';')

        try: self.__h1()
        except: print('ERR', end=';')

        try: self._h2()
        except: print('ERR', end=';')


class B(A):
    d = 4

    def __h1(self): print(5, end=';')
    def _h2(self): print(6, end=';')

try: B().main()
except: print('ERR', end=';')
code

8. Що буде виведено за виконання?
code
def f(arg, n):
    arg[42] = 3
    n *= -3

try:
    n = 5
    s = {83:6, 26:1, 26:2}
    f(s, n)
    print(n, end=';')
    for x in s :
        print(x, sep=';', end=';')
except:
    print('ERR')
code

9. Що буде виведено за виконання?
code
class A:
    a = 1
    c = 2

    def __init__(self):
        c = 3

class B(A):
    a = 4

    def __init__(self):
        super().__init__()
        self.b = 7

obj = B()
obj.c = obj.b + 10
B.c = 13

try: print(obj.c, end= ';')
except: print('ERR', end=';')

try: print(A.c, end= ';')
except: print('ERR', end=';')

try: print(B.c, end= ';')
except: print('ERR', end=';')
code

10. 
Для графа на рисунку з вершини 1 запущено алгоритм пошуку в глибину, що будує кістякове дерево. Списки суміжності вершин переглядаються алгоритмом за зростанням міток вершин.\par
\begin{center}
	\adjustimage{max size={0.9\linewidth}{0.9\paperheight}}
	{../10/Traversals/69.png}
\end{center}
\par Які з наведених ребер входять до складу отриманого кістякового дерева?\par
1) (1, 6)\par
2) (3, 4)\par
3) (3, 6)\par
4) (0, 2)\par
5) (2, 6)\par
6) (0, 1)\par
{\vskip 15pt} У формі вибрати номери відповідних ребер.\par
%answer:2 3 4 5 6
\end {large}}
%end
{\smallskip \begin {small} Затверджено на засіданні кафедри теоретичної кібернетики, протокол \No 12 від   19.05.2020 р.\par\smallskip\smallskip
{\parbox {6.5 cm}{Зав. кафедри \dotfill Ю.В.~Крак}}\hspace {0.7cm}{\hbox to 6.5 cm{ Екзаменатор \dotfill Т.О.~Карнаух}} \end{small}}
%begin
{{\begin {small}\par
{ \center  Київський національний університет імені Тараса Шевченка \par}
{\parbox {5 cm}{Освітня програма \hfill}{\parbox {7 cm}{Прикладна математика, Системний аналіз \hfill}}\parbox {6 cm}{\hfill Семестр 2}}\par
{\parbox {5 cm} {Навчальний предмет \hfill}\parbox {7 cm}{Програмування \hfill}\parbox {6cm}{\hfill}\par}\vspace{-7pt} \end{small}}{\center Екзаменаційний білет \No \textbf 
{ 4 }\par}}
{
\begin {large}
1. Що буде виведено за виконання?
code
try:
    res = 2.0>="True"
    if isinstance(res, bool):
        print(f'bool;{res}')
    elif isinstance(res, int):
        print(f'int;{res}')
    elif isinstance(res, float):
        print(f'float;{res:.2f}')
    else:
        print('???')
except:
    print('ERR')
code

2. Що буде виведено за виконання?
code
a = 1
b = 9
c = 8

def h(a):
    a -= 3
    b = 5
    c = 5
    return a + b + c

a = 5
b = 3
c = 5

try:
    print(h(b), a, b, c, sep=';')
except:
    print('ERR', a, b, c, sep=';')
code

3. Що буде виведено за виконання?
code
try:
    n = 1
    while n<=11:
        n += 2
    print(n, end=';')
except:
    print('ERR')
code

4. Що буде виведено за виконання?
code
try:
    for e in range(4, -11, 2):
        if e>6:
            break
        print(e, end=';')
    else:
        print(e,end=';')
    print(e)
except:
    print('ERR')
code

5. Що буде виведено за виконання?\par (Дійсні значення, якщо наявні, в цьому завданні будуть виводитись у форматі з фіксованою крапкою та, можливо, з доволі великою кількістю десяткових розрядів. У відповіді для дійсних значень у ЦЬОМУ завданні достатньо вказати один десятковий розряд після крапки, відкинувши решту розрядів без округлення. Наприклад, замість 13.66666666666667 достатньо вказати 13.6, замість 25.23 достатньо вказати 25.2)
code
try:
    print(1, end=';')
    print(3j<8j, end=';')
    print(9, end=';')
except TypeError:
    print(2, end=';')
except ValueError:
    print(5, end=';')
except:
    print('ERR', end=';')
else:
    print(7, end=';')
finally:
    print(6)
code

6. Що буде виведено за виконання?
code
def f(s1):
    s = s1[:]
    for i in range(len(s)):
        s[i] += 1
 
s = [1, 2, 3]
f(s)
print(s)
code

7. Що буде виведено за виконання?
code
class A:
    c = 1

    def f1(self): print(2, end=';')
    def _f2(self): print(3, end=';')

    def main(self):
        try: print(self.c, end=';')
        except: print('ERR', end=';')

        try: self.f1()
        except: print('ERR', end=';')

        try: self._f2()
        except: print('ERR', end=';')


class B(A):
    c = 4

    def f1(self): print(5, end=';')
    def _f2(self): print(6, end=';')

try: B().main()
except: print('ERR', end=';')
code

8. Що буде виведено за виконання?
code
def f(arg, n):
    n = 0
    tmp = arg
    tmp[8:-5] = [1,0]
    return len(tmp)>3

try:
    f = [10,11,12,13,14]
    n = 1
    f(f, n)
    print(n, end=';')
    for el in f:
        print(el, end=';')
except:
    print('ERR')
code

9. Що буде виведено за виконання?
code
class A:
    c = 1
    b = 2

    def __init__(self):
        c = 3

class B(A):
    b = 4

    def __init__(self):
        super().__init__()
        self.a = 7

obj = B()
obj.c = obj.b + 10
B.c = 13

try: print(obj.c, end= ';')
except: print('ERR', end=';')

try: print(A.c, end= ';')
except: print('ERR', end=';')

try: print(B.c, end= ';')
except: print('ERR', end=';')
code

10. 
Для графа на рисунку з вершини 6 запущено алгоритм пошуку в глибину, що будує кістякове дерево. Списки суміжності вершин переглядаються алгоритмом за зростанням міток вершин.\par
\begin{center}
	\adjustimage{max size={0.9\linewidth}{0.9\paperheight}}
	{../10/Traversals/280.png}
\end{center}
\par Які з наведених ребер входять до складу отриманого кістякового дерева?\par
1) (0, 6)\par
2) (2, 5)\par
3) (1, 4)\par
4) (0, 1)\par
5) (1, 3)\par
6) (0, 3)\par
{\vskip 15pt} У формі вибрати номери відповідних ребер.\par
%answer:1 2 3 4 5
\end {large}}
%end
{\smallskip \begin {small} Затверджено на засіданні кафедри теоретичної кібернетики, протокол \No 12 від   19.05.2020 р.\par\smallskip\smallskip
{\parbox {6.5 cm}{Зав. кафедри \dotfill Ю.В.~Крак}}\hspace {0.7cm}{\hbox to 6.5 cm{ Екзаменатор \dotfill Т.О.~Карнаух}} \end{small}}
%begin
{{\begin {small}\par
{ \center  Київський національний університет імені Тараса Шевченка \par}
{\parbox {5 cm}{Освітня програма \hfill}{\parbox {7 cm}{Прикладна математика, Системний аналіз \hfill}}\parbox {6 cm}{\hfill Семестр 2}}\par
{\parbox {5 cm} {Навчальний предмет \hfill}\parbox {7 cm}{Програмування \hfill}\parbox {6cm}{\hfill}\par}\vspace{-7pt} \end{small}}{\center Екзаменаційний білет \No \textbf 
{ 5 }\par}}
{
\begin {large}
1. Що буде виведено за виконання?
code
try:
    res = False<"4j"
    if isinstance(res, bool):
        print(f'bool;{res}')
    elif isinstance(res, int):
        print(f'int;{res}')
    elif isinstance(res, float):
        print(f'float;{res:.2f}')
    else:
        print('???')
except:
    print('ERR')
code

2. Що буде виведено за виконання?
code
a = 2
b = 7
c = 6

def g(a):
    global c
    a = 3
    b += 4
    c = 2
    return a + b + c

a = 9
b = 8
c = 0

try:
    print(g(a), a, b, c, sep=';')
except:
    print('ERR', a, b, c, sep=';')
code

3. Що буде виведено за виконання?
code
try:
    t = 5
    while t<5:
        t += 2
    print(t, end=';')
except:
    print('ERR')
code

4. Що буде виведено за виконання?
code
try:
    for d in range(-7, -14, -3):
        if d>=7:
            continue
        print(d, end=';')
    else:
        print('end',end=';')
    print(d)
except:
    print('ERR')
code

5. Що буде виведено за виконання?\par (Дійсні значення, якщо наявні, в цьому завданні будуть виводитись у форматі з фіксованою крапкою та, можливо, з доволі великою кількістю десяткових розрядів. У відповіді для дійсних значень у ЦЬОМУ завданні достатньо вказати один десятковий розряд після крапки, відкинувши решту розрядів без округлення. Наприклад, замість 13.66666666666667 достатньо вказати 13.6, замість 25.23 достатньо вказати 25.2)
code
try:
    print(0, end=';')
    print(int('a8'), end=';')
    print(9, end=';')
except KeyboardInterrupt:
    print(4, end=';')
except Exception:
    print(1, end=';')
except:
    print('ERR', end=';')
else:
    print(0, end=';')
finally:
    print(7)
code

6. Що буде виведено за виконання?
code
def f(s):
    for i in range(len(s)):
        s[i] += 1
    return
 
s = [1, 2, 3]
f(s)
print(s)
code

7. Що буде виведено за виконання?
code
class A:
    a = 1

    def _g1(self): print(2, end=';')
    def g2(self): print(3, end=';')

    def main(self):
        try: print(self.a, end=';')
        except: print('ERR', end=';')

        try: self._g1()
        except: print('ERR', end=';')

        try: self.g2()
        except: print('ERR', end=';')


class B(A):
    a = 4

    def _g1(self): print(5, end=';')
    def g2(self): print(6, end=';')

try: B().main()
except: print('ERR', end=';')
code

8. Що буде виведено за виконання?
code
def f(arg, n):
    arg[16] = 9
    n += 3

try:
    n = 4
    t = {41:7, 12:2, 31:8}
    f(t, n)
    print(n, end=';')
    for x in t :
        print(x, sep=';', end=';')
except:
    print('ERR')
code

9. Що буде виведено за виконання?
code
class A:
    b = 1
    a = 2

    def __init__(self):
        a = 3

class B(A):
    b = 4

    def __init__(self):
        super().__init__()
        self.c = 7

obj = B()
obj.a = obj.b + 10
B.a = 13

try: print(obj.a, end= ';')
except: print('ERR', end=';')

try: print(A.a, end= ';')
except: print('ERR', end=';')

try: print(B.a, end= ';')
except: print('ERR', end=';')
code

10. 
Для графа на рисунку з вершини 6 запущено алгоритм пошуку в глибину, що будує кістякове дерево. Списки суміжності вершин переглядаються алгоритмом за зростанням міток вершин.\par
\begin{center}
	\adjustimage{max size={0.9\linewidth}{0.9\paperheight}}
	{../10/Traversals/173.png}
\end{center}
\par Які з наведених ребер входять до складу отриманого кістякового дерева?\par
1) (3, 4)\par
2) (5, 6)\par
3) (0, 4)\par
4) (1, 2)\par
5) (2, 3)\par
6) (1, 3)\par
{\vskip 15pt} У формі вибрати номери відповідних ребер.\par
%answer:1 4 6
\end {large}}
%end
{\smallskip \begin {small} Затверджено на засіданні кафедри теоретичної кібернетики, протокол \No 12 від   19.05.2020 р.\par\smallskip\smallskip
{\parbox {6.5 cm}{Зав. кафедри \dotfill Ю.В.~Крак}}\hspace {0.7cm}{\hbox to 6.5 cm{ Екзаменатор \dotfill Т.О.~Карнаух}} \end{small}}
%begin
{{\begin {small}\par
{ \center  Київський національний університет імені Тараса Шевченка \par}
{\parbox {5 cm}{Освітня програма \hfill}{\parbox {7 cm}{Прикладна математика, Системний аналіз \hfill}}\parbox {6 cm}{\hfill Семестр 2}}\par
{\parbox {5 cm} {Навчальний предмет \hfill}\parbox {7 cm}{Програмування \hfill}\parbox {6cm}{\hfill}\par}\vspace{-7pt} \end{small}}{\center Екзаменаційний білет \No \textbf 
{ 6 }\par}}
{
\begin {large}
1. Що буде виведено за виконання?
code
try:
    res = 5//5/8*8
    if isinstance(res, bool):
        print(f'bool;{res}')
    elif isinstance(res, int):
        print(f'int;{res}')
    elif isinstance(res, float):
        print(f'float;{res:.2f}')
    else:
        print('???')
except:
    print('ERR')
code

2. Що буде виведено за виконання?
code
a = 1
b = 3
c = 9

def g(a):
    a = 5
    b = 2
    c = 1
    return a + b + c

a = 7
b = 6
c = 4

try:
    print(g(a), a, b, c, sep=';')
except:
    print('ERR', a, b, c, sep=';')
code

3. Що буде виведено за виконання?
code
try:
    j = 8
    while j>1:
        j += -2
    print(j, end=';')
except:
    print('ERR')
code

4. Що буде виведено за виконання?
code
try:
    for f in range(-5, 6, -2):
        if f<3:
            continue
        print(f, end=';');f = -10
    else:
        print(f,end=';')
    print(f)
except:
    print('ERR')
code

5. Що буде виведено за виконання?\par (Дійсні значення, якщо наявні, в цьому завданні будуть виводитись у форматі з фіксованою крапкою та, можливо, з доволі великою кількістю десяткових розрядів. У відповіді для дійсних значень у ЦЬОМУ завданні достатньо вказати один десятковий розряд після крапки, відкинувши решту розрядів без округлення. Наприклад, замість 13.66666666666667 достатньо вказати 13.6, замість 25.23 достатньо вказати 25.2)
code
try:
    print(3, end=';')
    print(2//0, end=';')
    print(6, end=';')
except ValueError:
    print(8, end=';')
except ZeroDivisionError:
    print(5, end=';')
except:
    print('ERR', end=';')
else:
    print(6, end=';')
finally:
    print(7)
code

6. Що буде виведено за виконання?
code
def f(n):
    if n>9: f(n//10)
    else:
        print(n%10, end='')
 
f(123456)
code

7. Що буде виведено за виконання?
code
class A:
    c = 1

    def _h1(self): print(2, end=';')
    def _h2(self): print(3, end=';')

    def main(self):
        try: print(self.c, end=';')
        except: print('ERR', end=';')

        try: self._h1()
        except: print('ERR', end=';')

        try: self._h2()
        except: print('ERR', end=';')


class B(A):
    c = 4

    def _h1(self): print(5, end=';')
    def _h2(self): print(6, end=';')

try: B().main()
except: print('ERR', end=';')
code

8. Що буде виведено за виконання?
code
def f(arg, n):
    n = 6
    tmp = arg
    tmp[-3:4] = [2,3]
    return len(tmp)>3

try:
    e = ['0','1','2','3','4']
    n = 2
    f(e, n)
    print(n, end=';')
    for el in e:
        print(el, end=';')
except:
    print('ERR')
code

9. Що буде виведено за виконання?
code
class A:
    c = 1
    a = 2

    def __init__(self):
        c = 3

class B(A):
    a = 4

    def __init__(self):
        super().__init__()
        self.b = 7

obj = B()
obj.a = obj.c + 10
B.a = 13

try: print(obj.a, end= ';')
except: print('ERR', end=';')

try: print(A.a, end= ';')
except: print('ERR', end=';')

try: print(B.a, end= ';')
except: print('ERR', end=';')
code

10. 
Для графа на рисунку з вершини 0 запущено алгоритм пошуку в глибину, що будує кістякове дерево. Списки суміжності вершин переглядаються алгоритмом за зростанням міток вершин.\par
\begin{center}
	\adjustimage{max size={0.9\linewidth}{0.9\paperheight}}
	{../10/Traversals/32.png}
\end{center}
\par Які з наведених ребер входять до складу отриманого кістякового дерева?\par
1) (0, 1)\par
2) (2, 6)\par
3) (0, 3)\par
4) (2, 5)\par
5) (3, 4)\par
6) (0, 6)\par
{\vskip 15pt} У формі вибрати номери відповідних ребер.\par
%answer:1 2 4 5
\end {large}}
%end
{\smallskip \begin {small} Затверджено на засіданні кафедри теоретичної кібернетики, протокол \No 12 від   19.05.2020 р.\par\smallskip\smallskip
{\parbox {6.5 cm}{Зав. кафедри \dotfill Ю.В.~Крак}}\hspace {0.7cm}{\hbox to 6.5 cm{ Екзаменатор \dotfill Т.О.~Карнаух}} \end{small}}
%begin
{{\begin {small}\par
{ \center  Київський національний університет імені Тараса Шевченка \par}
{\parbox {5 cm}{Освітня програма \hfill}{\parbox {7 cm}{Прикладна математика, Системний аналіз \hfill}}\parbox {6 cm}{\hfill Семестр 2}}\par
{\parbox {5 cm} {Навчальний предмет \hfill}\parbox {7 cm}{Програмування \hfill}\parbox {6cm}{\hfill}\par}\vspace{-7pt} \end{small}}{\center Екзаменаційний білет \No \textbf 
{ 7 }\par}}
{
\begin {large}
1. Що буде виведено за виконання?
code
try:
    res = "5.0j"==9
    if isinstance(res, bool):
        print(f'bool;{res}')
    elif isinstance(res, int):
        print(f'int;{res}')
    elif isinstance(res, float):
        print(f'float;{res:.2f}')
    else:
        print('???')
except:
    print('ERR')
code

2. Що буде виведено за виконання?
code
a = 8
b = 2
c = 1

def g(b):
    global c
    a = 3
    b = 2
    c = 4
    return a + b + c

a = 6
b = 2
c = 8

try:
    print(g(b), a, b, c, sep=';')
except:
    print('ERR', a, b, c, sep=';')
code

3. Що буде виведено за виконання?
code
try:
    i = 0
    while i<=13:
        i += 1
    print(i, end=';')
except:
    print('ERR')
code

4. Що буде виведено за виконання?
code
try:
    for d in range(-10, 4, -1):
        if d<=6:
            break
        print(d, end=';')
    else:
        print(d,end=';')
    print(d)
except:
    print('ERR')
code

5. Що буде виведено за виконання?\par (Дійсні значення, якщо наявні, в цьому завданні будуть виводитись у форматі з фіксованою крапкою та, можливо, з доволі великою кількістю десяткових розрядів. У відповіді для дійсних значень у ЦЬОМУ завданні достатньо вказати один десятковий розряд після крапки, відкинувши решту розрядів без округлення. Наприклад, замість 13.66666666666667 достатньо вказати 13.6, замість 25.23 достатньо вказати 25.2)
code
try:
    print(9, end=';')
    print(int('3'), end=';')
    print(1, end=';')
except Exception:
    print(0, end=';')
except TypeError:
    print(8, end=';')
except:
    print('ERR', end=';')
else:
    print(2, end=';')
finally:
    print(4)
code

6. Що буде виведено за виконання?
code
def f(s):
    for el in s:
        el += 1
    return s
 
s = [1, 2, 3]
s = f(s)
print(s)
code

7. Що буде виведено за виконання?
code
class A:
    d = 1

    def f1(self): print(2, end=';')
    def _f2(self): print(3, end=';')

    def main(self):
        try: print(self.d, end=';')
        except: print('ERR', end=';')

        try: self.f1()
        except: print('ERR', end=';')

        try: self._f2()
        except: print('ERR', end=';')


class B(A):
    d = 4

    def f1(self): print(5, end=';')
    def _f2(self): print(6, end=';')

try: B().main()
except: print('ERR', end=';')
code

8. Що буде виведено за виконання?
code
def f(arg, n):
    n = 9
    tmp = arg
    del tmp[:6]
    return len(tmp)>3

try:
    h = [10,11,12,13,14]
    n = 4
    f(h, n)
    print(n, end=';')
    for el in h:
        print(el, end=';')
except:
    print('ERR')
code

9. Що буде виведено за виконання?
code
class A:
    a = 1
    c = 2

    def __init__(self):
        c = 3

class B(A):
    c = 4

    def __init__(self):
        super().__init__()
        self.b = 7

obj = B()
obj.a = obj.b + 10
B.a = 13

try: print(obj.a, end= ';')
except: print('ERR', end=';')

try: print(A.a, end= ';')
except: print('ERR', end=';')

try: print(B.a, end= ';')
except: print('ERR', end=';')
code

10. 
Для графа на рисунку з вершини 2 запущено алгоритм пошуку в глибину, що будує кістякове дерево. Списки суміжності вершин переглядаються алгоритмом за зростанням міток вершин.\par
\begin{center}
	\adjustimage{max size={0.9\linewidth}{0.9\paperheight}}
	{../10/Traversals/105.png}
\end{center}
\par Які з наведених ребер входять до складу отриманого кістякового дерева?\par
1) (0, 1)\par
2) (5, 6)\par
3) (4, 5)\par
4) (3, 6)\par
5) (1, 4)\par
6) (4, 6)\par
{\vskip 15pt} У формі вибрати номери відповідних ребер.\par
%answer:1 2
\end {large}}
%end
{\smallskip \begin {small} Затверджено на засіданні кафедри теоретичної кібернетики, протокол \No 12 від   19.05.2020 р.\par\smallskip\smallskip
{\parbox {6.5 cm}{Зав. кафедри \dotfill Ю.В.~Крак}}\hspace {0.7cm}{\hbox to 6.5 cm{ Екзаменатор \dotfill Т.О.~Карнаух}} \end{small}}
%begin
{{\begin {small}\par
{ \center  Київський національний університет імені Тараса Шевченка \par}
{\parbox {5 cm}{Освітня програма \hfill}{\parbox {7 cm}{Прикладна математика, Системний аналіз \hfill}}\parbox {6 cm}{\hfill Семестр 2}}\par
{\parbox {5 cm} {Навчальний предмет \hfill}\parbox {7 cm}{Програмування \hfill}\parbox {6cm}{\hfill}\par}\vspace{-7pt} \end{small}}{\center Екзаменаційний білет \No \textbf 
{ 8 }\par}}
{
\begin {large}
1. Що буде виведено за виконання?
code
try:
    res = "False"!=7.0j
    if isinstance(res, bool):
        print(f'bool;{res}')
    elif isinstance(res, int):
        print(f'int;{res}')
    elif isinstance(res, float):
        print(f'float;{res:.2f}')
    else:
        print('???')
except:
    print('ERR')
code

2. Що буде виведено за виконання?
code
a = 3
b = 9
c = 5

def h(a):
    a = 5
    b = 1
    c = 3
    return a + b + c

a = 0
b = 7
c = 4

try:
    print(h(a), a, b, c, sep=';')
except:
    print('ERR', a, b, c, sep=';')
code

3. Що буде виведено за виконання?
code
try:
    t = 9
    while t>=2:
        t += -3
    print(t, end=';')
except:
    print('ERR')
code

4. Що буде виведено за виконання?
code
try:
    for f in range(5, -1, 2):
        if f>4:
            continue
        print(f, end=';')
    else:
        print(f,end=';')
    print(f)
except:
    print('ERR')
code

5. Що буде виведено за виконання?\par (Дійсні значення, якщо наявні, в цьому завданні будуть виводитись у форматі з фіксованою крапкою та, можливо, з доволі великою кількістю десяткових розрядів. У відповіді для дійсних значень у ЦЬОМУ завданні достатньо вказати один десятковий розряд після крапки, відкинувши решту розрядів без округлення. Наприклад, замість 13.66666666666667 достатньо вказати 13.6, замість 25.23 достатньо вказати 25.2)
code
try:
    print(5, end=';')
    print(7j!=3j, end=';')
    print(9, end=';')
except KeyboardInterrupt:
    print(4, end=';')
except Exception:
    print(8, end=';')
except:
    print('ERR', end=';')
else:
    print(0, end=';')
finally:
    print(1)
code

6. Що буде виведено за виконання?
code
def f(s):
    for el in s:
        el += 1
        return s
 
s = [1, 2, 3]
s = f(s)
print(s)
code

7. Що буде виведено за виконання?
code
class A:
    b = 1

    def _h1(self): print(2, end=';')
    def h2(self): print(3, end=';')

    def main(self):
        try: print(self.b, end=';')
        except: print('ERR', end=';')

        try: self._h1()
        except: print('ERR', end=';')

        try: self.h2()
        except: print('ERR', end=';')


class B(A):
    b = 4

    def _h1(self): print(5, end=';')
    def h2(self): print(6, end=';')

try: B().main()
except: print('ERR', end=';')
code

8. Що буде виведено за виконання?
code
def f(arg, n):
    arg[46] = 2
    n -= -2

try:
    n = 7
    d = {10:0, 24:6, 37:5, 10:8}
    f(d, n)
    print(n, end=';')
    for x in d.items():
        print(x, sep=';', end=';')
except:
    print('ERR')
code

9. Що буде виведено за виконання?
code
class A:
    c = 1
    a = 2

    def __init__(self):
        c = 3

class B(A):
    c = 4

    def __init__(self):
        super().__init__()
        self.b = 7

obj = B()
obj.c = obj.a + 10
B.c = 13

try: print(obj.c, end= ';')
except: print('ERR', end=';')

try: print(A.c, end= ';')
except: print('ERR', end=';')

try: print(B.c, end= ';')
except: print('ERR', end=';')
code

10. 
Для графа на рисунку з вершини 4 запущено алгоритм пошуку в глибину, що будує кістякове дерево. Списки суміжності вершин переглядаються алгоритмом за зростанням міток вершин.\par
\begin{center}
	\adjustimage{max size={0.9\linewidth}{0.9\paperheight}}
	{../10/Traversals/44.png}
\end{center}
\par Які з наведених ребер входять до складу отриманого кістякового дерева?\par
1) (1, 5)\par
2) (4, 6)\par
3) (1, 4)\par
4) (3, 6)\par
5) (0, 6)\par
6) (1, 2)\par
{\vskip 15pt} У формі вибрати номери відповідних ребер.\par
%answer:3 5
\end {large}}
%end
{\smallskip \begin {small} Затверджено на засіданні кафедри теоретичної кібернетики, протокол \No 12 від   19.05.2020 р.\par\smallskip\smallskip
{\parbox {6.5 cm}{Зав. кафедри \dotfill Ю.В.~Крак}}\hspace {0.7cm}{\hbox to 6.5 cm{ Екзаменатор \dotfill Т.О.~Карнаух}} \end{small}}
%begin
{{\begin {small}\par
{ \center  Київський національний університет імені Тараса Шевченка \par}
{\parbox {5 cm}{Освітня програма \hfill}{\parbox {7 cm}{Прикладна математика, Системний аналіз \hfill}}\parbox {6 cm}{\hfill Семестр 2}}\par
{\parbox {5 cm} {Навчальний предмет \hfill}\parbox {7 cm}{Програмування \hfill}\parbox {6cm}{\hfill}\par}\vspace{-7pt} \end{small}}{\center Екзаменаційний білет \No \textbf 
{ 9 }\par}}
{
\begin {large}
1. Що буде виведено за виконання?
code
try:
    res = True>5j
    if isinstance(res, bool):
        print(f'bool;{res}')
    elif isinstance(res, int):
        print(f'int;{res}')
    elif isinstance(res, float):
        print(f'float;{res:.2f}')
    else:
        print('???')
except:
    print('ERR')
code

2. Що буде виведено за виконання?
code
a = 4
b = 7
c = 1

def f(b):
    global c
    a = 4
    b = 3
    c = 1
    return a + b + c

a = 2
b = 6
c = 3

try:
    print(f(a), a, b, c, sep=';')
except:
    print('ERR', a, b, c, sep=';')
code

3. Що буде виведено за виконання?
code
try:
    k = 14
    while k>8:
        k += -2
    print(k, end=';')
except:
    print('ERR')
code

4. Що буде виведено за виконання?
code
try:
    for c in range(-12, -4, -3):
        if c<3:
            continue
        print(c, end=';');c = 0
        if c>=7:
            break
    else:
        print('end',end=';')
    print(c)
except:
    print('ERR')
code

5. Що буде виведено за виконання?\par (Дійсні значення, якщо наявні, в цьому завданні будуть виводитись у форматі з фіксованою крапкою та, можливо, з доволі великою кількістю десяткових розрядів. У відповіді для дійсних значень у ЦЬОМУ завданні достатньо вказати один десятковий розряд після крапки, відкинувши решту розрядів без округлення. Наприклад, замість 13.66666666666667 достатньо вказати 13.6, замість 25.23 достатньо вказати 25.2)
code
try:
    print(5, end=';')
    print(3/3, end=';')
    print(2, end=';')
except ZeroDivisionError:
    print(6, end=';')
except ValueError:
    print(8, end=';')
except:
    print('ERR', end=';')
else:
    print(5, end=';')
finally:
    print(3)
code

6. Що буде виведено за виконання?
code
def f(s):
    s1 = s
    for i in range(len(s)):
        s1[i] += 1
    return s1

s = [1, 2, 3]
f(s)
print(s)
code

7. Що буде виведено за виконання?
code
class A:
    c = 1

    def _g1(self): print(2, end=';')
    def __g2(self): print(3, end=';')

    def main(self):
        try: print(self.c, end=';')
        except: print('ERR', end=';')

        try: self._g1()
        except: print('ERR', end=';')

        try: self.__g2()
        except: print('ERR', end=';')


class B(A):
    c = 4

    def _g1(self): print(5, end=';')
    def __g2(self): print(6, end=';')

try: B().main()
except: print('ERR', end=';')
code

8. Що буде виведено за виконання?
code
def f(arg, n):
    arg[90] = 7
    n = 5
    arg = list(arg.values())

try:
    n = 1
    s = {68:8, 34:5, 34:5}
    f(s, n)
    print(n, end=';')
    for x in s.values():
        print(x, sep=';', end=';')
except:
    print('ERR')
code

9. Що буде виведено за виконання?
code
class A:
    b = 1
    c = 2

    def __init__(self):
        c = 3

class B(A):
    b = 4

    def __init__(self):
        super().__init__()
        self.a = 7

obj = B()
obj.c = obj.b + 10
B.c = 13

try: print(obj.c, end= ';')
except: print('ERR', end=';')

try: print(A.c, end= ';')
except: print('ERR', end=';')

try: print(B.c, end= ';')
except: print('ERR', end=';')
code

10. 
Для графа на рисунку з вершини 5 запущено алгоритм пошуку в глибину, що будує кістякове дерево. Списки суміжності вершин переглядаються алгоритмом за зростанням міток вершин.\par
\begin{center}
	\adjustimage{max size={0.9\linewidth}{0.9\paperheight}}
	{../10/Traversals/36.png}
\end{center}
\par Які з наведених ребер входять до складу отриманого кістякового дерева?\par
1) (0, 2)\par
2) (2, 5)\par
3) (1, 4)\par
4) (0, 3)\par
5) (2, 6)\par
6) (1, 5)\par
{\vskip 15pt} У формі вибрати номери відповідних ребер.\par
%answer:1 5
\end {large}}
%end
{\smallskip \begin {small} Затверджено на засіданні кафедри теоретичної кібернетики, протокол \No 12 від   19.05.2020 р.\par\smallskip\smallskip
{\parbox {6.5 cm}{Зав. кафедри \dotfill Ю.В.~Крак}}\hspace {0.7cm}{\hbox to 6.5 cm{ Екзаменатор \dotfill Т.О.~Карнаух}} \end{small}}
%begin
{{\begin {small}\par
{ \center  Київський національний університет імені Тараса Шевченка \par}
{\parbox {5 cm}{Освітня програма \hfill}{\parbox {7 cm}{Прикладна математика, Системний аналіз \hfill}}\parbox {6 cm}{\hfill Семестр 2}}\par
{\parbox {5 cm} {Навчальний предмет \hfill}\parbox {7 cm}{Програмування \hfill}\parbox {6cm}{\hfill}\par}\vspace{-7pt} \end{small}}{\center Екзаменаційний білет \No \textbf 
{ 10 }\par}}
{
\begin {large}
1. Що буде виведено за виконання?
code
try:
    res = 1**0.5*1
    if isinstance(res, bool):
        print(f'bool;{res}')
    elif isinstance(res, int):
        print(f'int;{res}')
    elif isinstance(res, float):
        print(f'float;{res:.2f}')
    else:
        print('???')
except:
    print('ERR')
code

2. Що буде виведено за виконання?
code
a = 8
b = 9
c = 0

def f(a):
    a = 2
    b = 5
    c = 3
    return a + b + c

a = 5
b = 0
c = 3

try:
    print(f(b), a, b, c, sep=';')
except:
    print('ERR', a, b, c, sep=';')
code

3. Що буде виведено за виконання?
code
try:
    k = 7
    while k<=12:
        k += 2
    print(k, end=';')
except:
    print('ERR')
code

4. Що буде виведено за виконання?
code
try:
    for a in range(-14, -9, 3):
        if a>=8:
            break
        print(a, end=';')
    else:
        print(a,end=';')
    print(a)
except:
    print('ERR')
code

5. Що буде виведено за виконання?\par (Дійсні значення, якщо наявні, в цьому завданні будуть виводитись у форматі з фіксованою крапкою та, можливо, з доволі великою кількістю десяткових розрядів. У відповіді для дійсних значень у ЦЬОМУ завданні достатньо вказати один десятковий розряд після крапки, відкинувши решту розрядів без округлення. Наприклад, замість 13.66666666666667 достатньо вказати 13.6, замість 25.23 достатньо вказати 25.2)
code
try:
    print(9, end=';')
    print(int('b4'), end=';')
    print(0, end=';')
except TypeError:
    print(1, end=';')
except KeyboardInterrupt:
    print(6, end=';')
except:
    print('ERR', end=';')
else:
    print(2, end=';')
finally:
    print(7)
code

6. Що буде виведено за виконання?
code
def f(n):
    if n>9:
        f(n//10)
    else:
        print(n%10, end='')

f(543216)
code

7. Що буде виведено за виконання?
code
class A:
    d = 1

    def __g1(self): print(2, end=';')
    def _g2(self): print(3, end=';')

    def main(self):
        try: print(self.d, end=';')
        except: print('ERR', end=';')

        try: self.__g1()
        except: print('ERR', end=';')

        try: self._g2()
        except: print('ERR', end=';')


class B(A):
    d = 4

    def __g1(self): print(5, end=';')
    def _g2(self): print(6, end=';')

try: B().main()
except: print('ERR', end=';')
code

8. Що буде виведено за виконання?
code
def f(arg, n):
    arg[22] = 6
    n *= 2

try:
    n = 3
    d = {71:6, 36:9, 36:7}
    f(d, n)
    print(n, end=';')
    for x in d.keys():
        print(x, sep=';', end=';')
except:
    print('ERR')
code

9. Що буде виведено за виконання?
code
class A:
    b = 1
    a = 2

    def __init__(self):
        b = 3

class B(A):
    a = 4

    def __init__(self):
        super().__init__()
        self.c = 7

obj = B()
obj.a = obj.c + 10
B.a = 13

try: print(obj.a, end= ';')
except: print('ERR', end=';')

try: print(A.a, end= ';')
except: print('ERR', end=';')

try: print(B.a, end= ';')
except: print('ERR', end=';')
code

10. 
Для графа на рисунку з вершини 6 запущено алгоритм пошуку в глибину, що будує кістякове дерево. Списки суміжності вершин переглядаються алгоритмом за зростанням міток вершин.\par
\begin{center}
	\adjustimage{max size={0.9\linewidth}{0.9\paperheight}}
	{../10/Traversals/153.png}
\end{center}
\par Які з наведених ребер входять до складу отриманого кістякового дерева?\par
1) (1, 6)\par
2) (1, 2)\par
3) (0, 2)\par
4) (4, 5)\par
5) (0, 4)\par
6) (3, 5)\par
{\vskip 15pt} У формі вибрати номери відповідних ребер.\par
%answer:1 2 3 5 6
\end {large}}
%end
{\smallskip \begin {small} Затверджено на засіданні кафедри теоретичної кібернетики, протокол \No 12 від   19.05.2020 р.\par\smallskip\smallskip
{\parbox {6.5 cm}{Зав. кафедри \dotfill Ю.В.~Крак}}\hspace {0.7cm}{\hbox to 6.5 cm{ Екзаменатор \dotfill Т.О.~Карнаух}} \end{small}}
%begin
{{\begin {small}\par
{ \center  Київський національний університет імені Тараса Шевченка \par}
{\parbox {5 cm}{Освітня програма \hfill}{\parbox {7 cm}{Прикладна математика, Системний аналіз \hfill}}\parbox {6 cm}{\hfill Семестр 2}}\par
{\parbox {5 cm} {Навчальний предмет \hfill}\parbox {7 cm}{Програмування \hfill}\parbox {6cm}{\hfill}\par}\vspace{-7pt} \end{small}}{\center Екзаменаційний білет \No \textbf 
{ 11 }\par}}
{
\begin {large}
1. Що буде виведено за виконання?
code
try:
    res = "7"<="4.0"
    if isinstance(res, bool):
        print(f'bool;{res}')
    elif isinstance(res, int):
        print(f'int;{res}')
    elif isinstance(res, float):
        print(f'float;{res:.2f}')
    else:
        print('???')
except:
    print('ERR')
code

2. Що буде виведено за виконання?
code
a = 1
b = 4
c = 9

def h(b):
    a -= 1
    b = 1
    c = 2
    return a + b + c

a = 3
b = 0
c = 4

try:
    print(h(a), a, b, c, sep=';')
except:
    print('ERR', a, b, c, sep=';')
code

3. Що буде виведено за виконання?
code
try:
    m = 2
    while m<10:
        m += 1
    print(m, end=';')
except:
    print('ERR')
code

4. Що буде виведено за виконання?
code
try:
    for e in range(3, 7, -2):
        if e<=7:
            continue
        print(e, end=';');e = 8
    else:
        print(e,end=';')
    print(e)
except:
    print('ERR')
code

5. Що буде виведено за виконання?\par (Дійсні значення, якщо наявні, в цьому завданні будуть виводитись у форматі з фіксованою крапкою та, можливо, з доволі великою кількістю десяткових розрядів. У відповіді для дійсних значень у ЦЬОМУ завданні достатньо вказати один десятковий розряд після крапки, відкинувши решту розрядів без округлення. Наприклад, замість 13.66666666666667 достатньо вказати 13.6, замість 25.23 достатньо вказати 25.2)
code
try:
    print(8, end=';')
    print(int('d7'), end=';')
    print(4, end=';')
except KeyboardInterrupt:
    print(3, end=';')
except Exception:
    print(5, end=';')
except:
    print('ERR', end=';')
else:
    print(9, end=';')
finally:
    print(1)
code

6. Що буде виведено за виконання?
code
def f(n):
    if n>9:
        f(n//100)
    else:
        print(n%10,end='')

f(987654)
code

7. Що буде виведено за виконання?
code
class A:
    a = 1

    def _h1(self): print(2, end=';')
    def __h2(self): print(3, end=';')

    def main(self):
        try: print(self.a, end=';')
        except: print('ERR', end=';')

        try: self._h1()
        except: print('ERR', end=';')

        try: self.__h2()
        except: print('ERR', end=';')


class B(A):
    a = 4

    def _h1(self): print(5, end=';')
    def __h2(self): print(6, end=';')

try: B().main()
except: print('ERR', end=';')
code

8. Що буде виведено за виконання?
code
def f(arg, n):
    n = 7
    tmp = arg
    tmp[:-2] = (14)
    return len(tmp)>3

try:
    a = ['0','1','2','3','4']
    n = 5
    f(a, n)
    print(n, end=';')
    for el in a:
        print(el, end=';')
except:
    print('ERR')
code

9. Що буде виведено за виконання?
code
class A:
    c = 1
    b = 2

    def __init__(self):
        c = 3

class B(A):
    c = 4

    def __init__(self):
        super().__init__()
        self.a = 7

obj = B()
obj.b = obj.a + 10
B.b = 13

try: print(obj.b, end= ';')
except: print('ERR', end=';')

try: print(A.b, end= ';')
except: print('ERR', end=';')

try: print(B.b, end= ';')
except: print('ERR', end=';')
code

10. 
Для графа на рисунку з вершини 2 запущено алгоритм пошуку в глибину, що будує кістякове дерево. Списки суміжності вершин переглядаються алгоритмом за зростанням міток вершин.\par
\begin{center}
	\adjustimage{max size={0.9\linewidth}{0.9\paperheight}}
	{../10/Traversals/135.png}
\end{center}
\par Які з наведених ребер входять до складу отриманого кістякового дерева?\par
1) (5, 6)\par
2) (4, 5)\par
3) (1, 5)\par
4) (3, 4)\par
5) (3, 6)\par
6) (1, 3)\par
{\vskip 15pt} У формі вибрати номери відповідних ребер.\par
%answer:3 5 6
\end {large}}
%end
{\smallskip \begin {small} Затверджено на засіданні кафедри теоретичної кібернетики, протокол \No 12 від   19.05.2020 р.\par\smallskip\smallskip
{\parbox {6.5 cm}{Зав. кафедри \dotfill Ю.В.~Крак}}\hspace {0.7cm}{\hbox to 6.5 cm{ Екзаменатор \dotfill Т.О.~Карнаух}} \end{small}}
%begin
{{\begin {small}\par
{ \center  Київський національний університет імені Тараса Шевченка \par}
{\parbox {5 cm}{Освітня програма \hfill}{\parbox {7 cm}{Прикладна математика, Системний аналіз \hfill}}\parbox {6 cm}{\hfill Семестр 2}}\par
{\parbox {5 cm} {Навчальний предмет \hfill}\parbox {7 cm}{Програмування \hfill}\parbox {6cm}{\hfill}\par}\vspace{-7pt} \end{small}}{\center Екзаменаційний білет \No \textbf 
{ 12 }\par}}
{
\begin {large}
1. Що буде виведено за виконання?
code
try:
    res = 9/7%7*3
    if isinstance(res, bool):
        print(f'bool;{res}')
    elif isinstance(res, int):
        print(f'int;{res}')
    elif isinstance(res, float):
        print(f'float;{res:.2f}')
    else:
        print('???')
except:
    print('ERR')
code

2. Що буде виведено за виконання?
code
a = 5
b = 9
c = 8

def g(b):
    global c
    a = 5
    b = 4
    c = 1
    return a + b + c

a = 2
b = 7
c = 6

try:
    print(g(b), a, b, c, sep=';')
except:
    print('ERR', a, b, c, sep=';')
code

3. Що буде виведено за виконання?
code
try:
    i = 6
    while i>=0:
        i += -3
    print(i, end=';')
except:
    print('ERR')
code

4. Що буде виведено за виконання?
code
try:
    for b in range(1, 15, 3):
        if b>5:
            continue
        print(b, end=';')
    else:
        print(b,end=';')
    print(b)
except:
    print('ERR')
code

5. Що буде виведено за виконання?\par (Дійсні значення, якщо наявні, в цьому завданні будуть виводитись у форматі з фіксованою крапкою та, можливо, з доволі великою кількістю десяткових розрядів. У відповіді для дійсних значень у ЦЬОМУ завданні достатньо вказати один десятковий розряд після крапки, відкинувши решту розрядів без округлення. Наприклад, замість 13.66666666666667 достатньо вказати 13.6, замість 25.23 достатньо вказати 25.2)
code
try:
    print(0, end=';')
    print(2//2, end=';')
    print(6, end=';')
except ValueError:
    print(7, end=';')
except ZeroDivisionError:
    print(1, end=';')
except:
    print('ERR', end=';')
else:
    print(5, end=';')
finally:
    print(8)
code

6. Що буде виведено за виконання?
code
def f(n):
    if n>9:
        f(n//10)
    print(n%10, end='')
 
f(123456)
code

7. Що буде виведено за виконання?
code
class A:
    b = 1

    def _f1(self): print(2, end=';')
    def f2(self): print(3, end=';')

    def main(self):
        try: print(self.b, end=';')
        except: print('ERR', end=';')

        try: self._f1()
        except: print('ERR', end=';')

        try: self.f2()
        except: print('ERR', end=';')


class B(A):
    b = 4

    def _f1(self): print(5, end=';')
    def f2(self): print(6, end=';')

try: B().main()
except: print('ERR', end=';')
code

8. Що буде виведено за виконання?
code
def f(arg, n):
    arg[41] = 2
    n = 6
    arg = list(arg.items())

try:
    n = 4
    d = {44:0, 79:1, 79:1}
    f(d, n)
    print(n, end=';')
    for x in d.values():
        print(x, sep=';', end=';')
except:
    print('ERR')
code

9. Що буде виведено за виконання?
code
class A:
    a = 1
    b = 2

    def __init__(self):
        b = 3

class B(A):
    b = 4

    def __init__(self):
        super().__init__()
        self.c = 7

obj = B()
obj.b = obj.c + 10
B.b = 13

try: print(obj.b, end= ';')
except: print('ERR', end=';')

try: print(A.b, end= ';')
except: print('ERR', end=';')

try: print(B.b, end= ';')
except: print('ERR', end=';')
code

10. 
Для графа на рисунку з вершини 6 запущено алгоритм пошуку в глибину, що будує кістякове дерево. Списки суміжності вершин переглядаються алгоритмом за зростанням міток вершин.\par
\begin{center}
	\adjustimage{max size={0.9\linewidth}{0.9\paperheight}}
	{../10/Traversals/129.png}
\end{center}
\par Які з наведених ребер входять до складу отриманого кістякового дерева?\par
1) (2, 3)\par
2) (1, 2)\par
3) (2, 6)\par
4) (3, 6)\par
5) (4, 5)\par
6) (5, 6)\par
{\vskip 15pt} У формі вибрати номери відповідних ребер.\par
%answer:1 3 5
\end {large}}
%end
{\smallskip \begin {small} Затверджено на засіданні кафедри теоретичної кібернетики, протокол \No 12 від   19.05.2020 р.\par\smallskip\smallskip
{\parbox {6.5 cm}{Зав. кафедри \dotfill Ю.В.~Крак}}\hspace {0.7cm}{\hbox to 6.5 cm{ Екзаменатор \dotfill Т.О.~Карнаух}} \end{small}}
%begin
{{\begin {small}\par
{ \center  Київський національний університет імені Тараса Шевченка \par}
{\parbox {5 cm}{Освітня програма \hfill}{\parbox {7 cm}{Прикладна математика, Системний аналіз \hfill}}\parbox {6 cm}{\hfill Семестр 2}}\par
{\parbox {5 cm} {Навчальний предмет \hfill}\parbox {7 cm}{Програмування \hfill}\parbox {6cm}{\hfill}\par}\vspace{-7pt} \end{small}}{\center Екзаменаційний білет \No \textbf 
{ 13 }\par}}
{
\begin {large}
1. Що буде виведено за виконання?
code
def f(n):
    print('f', end='')
    return n==1

try:
    print(f(7) and f(2))
except:
    print('ERR')
code

2. Що буде виведено за виконання?
code
a = 7
b = 6
c = 2

def f(a):
    a += 4
    b = 5
    c = 3
    return a + b + c

a = 5
b = 8
c = 1

try:
    print(f(b), a, b, c, sep=';')
except:
    print('ERR', a, b, c, sep=';')
code

3. Що буде виведено за виконання?
code
try:
    i = 8
    while i<7:
        i += 2
    print(i, end=';')
except:
    print('ERR')
code

4. Що буде виведено за виконання?
code
try:
    for c in range(-15, -6, -1):
        if c>3:
            continue
        print(c, end=';')
    else:
        print(c,end=';')
    print(c)
except:
    print('ERR')
code

5. Що буде виведено за виконання?\par (Дійсні значення, якщо наявні, в цьому завданні будуть виводитись у форматі з фіксованою крапкою та, можливо, з доволі великою кількістю десяткових розрядів. У відповіді для дійсних значень у ЦЬОМУ завданні достатньо вказати один десятковий розряд після крапки, відкинувши решту розрядів без округлення. Наприклад, замість 13.66666666666667 достатньо вказати 13.6, замість 25.23 достатньо вказати 25.2)
code
try:
    print(6, end=';')
    print(9j>9j, end=';')
    print(2, end=';')
except TypeError:
    print(4, end=';')
except KeyboardInterrupt:
    print(0, end=';')
except:
    print('ERR', end=';')
else:
    print(3, end=';')
finally:
    print(3)
code

6. Що буде виведено за виконання?
code
def f(n):
    print(n%10, end='')
    if n>9:
        f(n//10)
 
f(543216)
code

7. Що буде виведено за виконання?
code
class A:
    b = 1

    def _f1(self): print(2, end=';')
    def _f2(self): print(3, end=';')

    def main(self):
        try: print(self.b, end=';')
        except: print('ERR', end=';')

        try: self._f1()
        except: print('ERR', end=';')

        try: self._f2()
        except: print('ERR', end=';')


class B(A):
    b = 4

    def _f1(self): print(5, end=';')
    def _f2(self): print(6, end=';')

try: B().main()
except: print('ERR', end=';')
code

8. Що буде виведено за виконання?
code
def f(arg, n):
    arg[11] = 0
    n += -3

try:
    n = 6
    t = {61:8, 14:0, 14:3}
    f(t, n)
    print(n, end=';')
    for x in t.values():
        print(x, sep=';', end=';')
except:
    print('ERR')
code

9. Що буде виведено за виконання?
code
class A:
    b = 1
    a = 2

    def __init__(self):
        a = 3

class B(A):
    a = 4

    def __init__(self):
        super().__init__()
        self.c = 7

obj = B()
obj.a = obj.c + 10
B.a = 13

try: print(obj.a, end= ';')
except: print('ERR', end=';')

try: print(A.a, end= ';')
except: print('ERR', end=';')

try: print(B.a, end= ';')
except: print('ERR', end=';')
code

10. 
Для графа на рисунку з вершини 6 запущено алгоритм пошуку в глибину, що будує кістякове дерево. Списки суміжності вершин переглядаються алгоритмом за зростанням міток вершин.\par
\begin{center}
	\adjustimage{max size={0.9\linewidth}{0.9\paperheight}}
	{../10/Traversals/9.png}
\end{center}
\par Які з наведених ребер входять до складу отриманого кістякового дерева?\par
1) (2, 6)\par
2) (0, 2)\par
3) (5, 6)\par
4) (3, 4)\par
5) (0, 5)\par
6) (3, 5)\par
{\vskip 15pt} У формі вибрати номери відповідних ребер.\par
%answer:1 2 6
\end {large}}
%end
{\smallskip \begin {small} Затверджено на засіданні кафедри теоретичної кібернетики, протокол \No 12 від   19.05.2020 р.\par\smallskip\smallskip
{\parbox {6.5 cm}{Зав. кафедри \dotfill Ю.В.~Крак}}\hspace {0.7cm}{\hbox to 6.5 cm{ Екзаменатор \dotfill Т.О.~Карнаух}} \end{small}}
%begin
{{\begin {small}\par
{ \center  Київський національний університет імені Тараса Шевченка \par}
{\parbox {5 cm}{Освітня програма \hfill}{\parbox {7 cm}{Прикладна математика, Системний аналіз \hfill}}\parbox {6 cm}{\hfill Семестр 2}}\par
{\parbox {5 cm} {Навчальний предмет \hfill}\parbox {7 cm}{Програмування \hfill}\parbox {6cm}{\hfill}\par}\vspace{-7pt} \end{small}}{\center Екзаменаційний білет \No \textbf 
{ 14 }\par}}
{
\begin {large}
1. Що буде виведено за виконання?
code
try:
    res = 1<="1.0j"
    if isinstance(res, bool):
        print(f'bool;{res}')
    elif isinstance(res, int):
        print(f'int;{res}')
    elif isinstance(res, float):
        print(f'float;{res:.2f}')
    else:
        print('???')
except:
    print('ERR')
code

2. Що буде виведено за виконання?
code
a = 7
b = 1
c = 0

def g(b):
    global c
    a = 2
    b *= 5
    c = 4
    return a + b + c

a = 2
b = 9
c = 5

try:
    print(g(b), a, b, c, sep=';')
except:
    print('ERR', a, b, c, sep=';')
code

3. Що буде виведено за виконання?
code
try:
    j = 9
    while j<=13:
        j += 1
    print(j, end=';')
except:
    print('ERR')
code

4. Що буде виведено за виконання?
code
try:
    for f in range(13, 0, -3):
        if f>=6:
            continue
        print(f, end=';');f = 7
    else:
        print(f,end=';')
    print(f)
except:
    print('ERR')
code

5. Що буде виведено за виконання?\par (Дійсні значення, якщо наявні, в цьому завданні будуть виводитись у форматі з фіксованою крапкою та, можливо, з доволі великою кількістю десяткових розрядів. У відповіді для дійсних значень у ЦЬОМУ завданні достатньо вказати один десятковий розряд після крапки, відкинувши решту розрядів без округлення. Наприклад, замість 13.66666666666667 достатньо вказати 13.6, замість 25.23 достатньо вказати 25.2)
code
try:
    print(1, end=';')
    print(7%0, end=';')
    print(2, end=';')
except Exception:
    print(9, end=';')
except TypeError:
    print(5, end=';')
except:
    print('ERR', end=';')
else:
    print(6, end=';')
finally:
    print(0)
code

6. Що буде виведено за виконання?
code
def f(n):
    print(n%2, end='')
    if n>2:
        f(n//2)
 
f(16)
code

7. Що буде виведено за виконання?
code
class A:
    a = 1

    def h1(self): print(2, end=';')
    def _h2(self): print(3, end=';')

    def main(self):
        try: print(self.a, end=';')
        except: print('ERR', end=';')

        try: self.h1()
        except: print('ERR', end=';')

        try: self._h2()
        except: print('ERR', end=';')


class B(A):
    a = 4

    def h1(self): print(5, end=';')
    def _h2(self): print(6, end=';')

try: B().main()
except: print('ERR', end=';')
code

8. Що буде виведено за виконання?
code
def f(arg, n):
    arg[71] = 1
    n -= -2

try:
    n = 4
    s = {27:8, 37:9, 16:9, 27:6}
    f(s, n)
    print(n, end=';')
    for x in s.keys():
        print(x, sep=';', end=';')
except:
    print('ERR')
code

9. Що буде виведено за виконання?
code
class A:
    a = 1
    b = 2

    def __init__(self):
        a = 3

class B(A):
    a = 4

    def __init__(self):
        super().__init__()
        self.c = 7

obj = B()
obj.b = obj.c + 10
B.b = 13

try: print(obj.b, end= ';')
except: print('ERR', end=';')

try: print(A.b, end= ';')
except: print('ERR', end=';')

try: print(B.b, end= ';')
except: print('ERR', end=';')
code

10. 
Для графа на рисунку з вершини 3 запущено алгоритм пошуку в глибину, що будує кістякове дерево. Списки суміжності вершин переглядаються алгоритмом за зростанням міток вершин.\par
\begin{center}
	\adjustimage{max size={0.9\linewidth}{0.9\paperheight}}
	{../10/Traversals/22.png}
\end{center}
\par Які з наведених ребер входять до складу отриманого кістякового дерева?\par
1) (3, 6)\par
2) (2, 4)\par
3) (4, 6)\par
4) (1, 5)\par
5) (3, 5)\par
6) (1, 3)\par
{\vskip 15pt} У формі вибрати номери відповідних ребер.\par
%answer:2 3 4 6
\end {large}}
%end
{\smallskip \begin {small} Затверджено на засіданні кафедри теоретичної кібернетики, протокол \No 12 від   19.05.2020 р.\par\smallskip\smallskip
{\parbox {6.5 cm}{Зав. кафедри \dotfill Ю.В.~Крак}}\hspace {0.7cm}{\hbox to 6.5 cm{ Екзаменатор \dotfill Т.О.~Карнаух}} \end{small}}
%begin
{{\begin {small}\par
{ \center  Київський національний університет імені Тараса Шевченка \par}
{\parbox {5 cm}{Освітня програма \hfill}{\parbox {7 cm}{Прикладна математика, Системний аналіз \hfill}}\parbox {6 cm}{\hfill Семестр 2}}\par
{\parbox {5 cm} {Навчальний предмет \hfill}\parbox {7 cm}{Програмування \hfill}\parbox {6cm}{\hfill}\par}\vspace{-7pt} \end{small}}{\center Екзаменаційний білет \No \textbf 
{ 15 }\par}}
{
\begin {large}
1. Що буде виведено за виконання?
code
try:
    res = -2**-4-False
    if isinstance(res, bool):
        print(f'bool;{res}')
    elif isinstance(res, int):
        print(f'int;{res}')
    elif isinstance(res, float):
        print(f'float;{res:.2f}')
    else:
        print('???')
except:
    print('ERR')
code

2. Що буде виведено за виконання?
code
a = 4
b = 1
c = 7

def h(a):
    a = 2
    b = 2
    c = 1
    return a + b + c

a = 1
b = 3
c = 0

try:
    print(h(a), a, b, c, sep=';')
except:
    print('ERR', a, b, c, sep=';')
code

3. Що буде виведено за виконання?
code
try:
    m = 13
    while m>7:
        m += -3
    print(m, end=';')
except:
    print('ERR')
code

4. Що буде виведено за виконання?
code
try:
    for b in range(9, -7, 2):
        if b<4:
            break
        print(b, end=';')
    else:
        print(b,end=';')
    print(b)
except:
    print('ERR')
code

5. Що буде виведено за виконання?\par (Дійсні значення, якщо наявні, в цьому завданні будуть виводитись у форматі з фіксованою крапкою та, можливо, з доволі великою кількістю десяткових розрядів. У відповіді для дійсних значень у ЦЬОМУ завданні достатньо вказати один десятковий розряд після крапки, відкинувши решту розрядів без округлення. Наприклад, замість 13.66666666666667 достатньо вказати 13.6, замість 25.23 достатньо вказати 25.2)
code
try:
    print(8, end=';')
    print(int('4'), end=';')
    print(5, end=';')
except ValueError:
    print(6, end=';')
except ZeroDivisionError:
    print(8, end=';')
except:
    print('ERR', end=';')
else:
    print(3, end=';')
finally:
    print(1)
code

6. Що буде виведено за виконання?
code
def f(n):
    print(n%10, end='')
    if n>9:
        f(n//100)
 
f(987654)
code

7. Що буде виведено за виконання?
code
class A:
    d = 1

    def __g1(self): print(2, end=';')
    def _g2(self): print(3, end=';')

    def main(self):
        try: print(self.d, end=';')
        except: print('ERR', end=';')

        try: self.__g1()
        except: print('ERR', end=';')

        try: self._g2()
        except: print('ERR', end=';')


class B(A):
    d = 4

    def __g1(self): print(5, end=';')
    def _g2(self): print(6, end=';')

try: B().main()
except: print('ERR', end=';')
code

8. Що буде виведено за виконання?
code
def f(arg, n):
    arg[17] = 6
    n = 8
    arg = set(arg.keys())

try:
    n = 3
    t = {49:5, 22:7, 73:1}
    f(t, n)
    print(n, end=';')
    for x in t.items():
        print(*x, sep=';', end=';')
except:
    print('ERR')
code

9. Що буде виведено за виконання?
code
class A:
    b = 1
    c = 2

    def __init__(self):
        b = 3

class B(A):
    b = 4

    def __init__(self):
        super().__init__()
        self.a = 7

obj = B()
obj.a = obj.b + 10
B.a = 13

try: print(obj.a, end= ';')
except: print('ERR', end=';')

try: print(A.a, end= ';')
except: print('ERR', end=';')

try: print(B.a, end= ';')
except: print('ERR', end=';')
code

10. 
Для графа на рисунку з вершини 3 запущено алгоритм пошуку в глибину, що будує кістякове дерево. Списки суміжності вершин переглядаються алгоритмом за зростанням міток вершин.\par
\begin{center}
	\adjustimage{max size={0.9\linewidth}{0.9\paperheight}}
	{../10/Traversals/182.png}
\end{center}
\par Які з наведених ребер входять до складу отриманого кістякового дерева?\par
1) (1, 4)\par
2) (3, 4)\par
3) (5, 6)\par
4) (2, 3)\par
5) (1, 6)\par
6) (1, 3)\par
{\vskip 15pt} У формі вибрати номери відповідних ребер.\par
%answer:3 6
\end {large}}
%end
{\smallskip \begin {small} Затверджено на засіданні кафедри теоретичної кібернетики, протокол \No 12 від   19.05.2020 р.\par\smallskip\smallskip
{\parbox {6.5 cm}{Зав. кафедри \dotfill Ю.В.~Крак}}\hspace {0.7cm}{\hbox to 6.5 cm{ Екзаменатор \dotfill Т.О.~Карнаух}} \end{small}}
%begin
{{\begin {small}\par
{ \center  Київський національний університет імені Тараса Шевченка \par}
{\parbox {5 cm}{Освітня програма \hfill}{\parbox {7 cm}{Прикладна математика, Системний аналіз \hfill}}\parbox {6 cm}{\hfill Семестр 2}}\par
{\parbox {5 cm} {Навчальний предмет \hfill}\parbox {7 cm}{Програмування \hfill}\parbox {6cm}{\hfill}\par}\vspace{-7pt} \end{small}}{\center Екзаменаційний білет \No \textbf 
{ 16 }\par}}
{
\begin {large}
1. Що буде виведено за виконання?
code
try:
    res = 6//6//9*7
    if isinstance(res, bool):
        print(f'bool;{res}')
    elif isinstance(res, int):
        print(f'int;{res}')
    elif isinstance(res, float):
        print(f'float;{res:.2f}')
    else:
        print('???')
except:
    print('ERR')
code

2. Що буде виведено за виконання?
code
a = 6
b = 8
c = 4

def h(b):
    global c
    a -= 3
    b = 1
    c = 5
    return a + b + c

a = 3
b = 4
c = 8

try:
    print(h(a), a, b, c, sep=';')
except:
    print('ERR', a, b, c, sep=';')
code

3. Що буде виведено за виконання?
code
try:
    n = 5
    while n>=4:
        n += -2
    print(n, end=';')
except:
    print('ERR')
code

4. Що буде виведено за виконання?
code
try:
    for a in range(-6, 10, 2):
        if a<=8:
            continue
        print(a, end=';')
    else:
        print(a,end=';')
    print(a)
except:
    print('ERR')
code

5. Що буде виведено за виконання?\par (Дійсні значення, якщо наявні, в цьому завданні будуть виводитись у форматі з фіксованою крапкою та, можливо, з доволі великою кількістю десяткових розрядів. У відповіді для дійсних значень у ЦЬОМУ завданні достатньо вказати один десятковий розряд після крапки, відкинувши решту розрядів без округлення. Наприклад, замість 13.66666666666667 достатньо вказати 13.6, замість 25.23 достатньо вказати 25.2)
code
try:
    print(9, end=';')
    print(int('0'), end=';')
    print(4, end=';')
except ZeroDivisionError:
    print(2, end=';')
except TypeError:
    print(7, end=';')
except:
    print('ERR', end=';')
else:
    print(9, end=';')
finally:
    print(0)
code

6. Що буде виведено за виконання?
code
def f(n):
    if n>9:
        return f(n//10)+1
    else:
        return 1

print(f(123456))
code

7. Що буде виведено за виконання?
code
class A:
    c = 1

    def _h1(self): print(2, end=';')
    def __h2(self): print(3, end=';')

    def main(self):
        try: print(self.c, end=';')
        except: print('ERR', end=';')

        try: self._h1()
        except: print('ERR', end=';')

        try: self.__h2()
        except: print('ERR', end=';')


class B(A):
    c = 4

    def _h1(self): print(5, end=';')
    def __h2(self): print(6, end=';')

try: B().main()
except: print('ERR', end=';')
code

8. Що буде виведено за виконання?
code
def f(arg, n):
    arg[88] = 8
    n = 7
    arg = tuple(arg.values())

try:
    n = 4
    s = {37:8, 87:2, 37:8}
    f(s, n)
    print(n, end=';')
    for x, y in s.items():
        print(x, y, sep=';', end=';')
except:
    print('ERR')
code

9. Що буде виведено за виконання?
code
class A:
    c = 1
    a = 2

    def __init__(self):
        a = 3

class B(A):
    a = 4

    def __init__(self):
        super().__init__()
        self.b = 7

obj = B()
obj.b = obj.c + 10
B.b = 13

try: print(obj.b, end= ';')
except: print('ERR', end=';')

try: print(A.b, end= ';')
except: print('ERR', end=';')

try: print(B.b, end= ';')
except: print('ERR', end=';')
code

10. 
Для графа на рисунку з вершини 5 запущено алгоритм пошуку в глибину, що будує кістякове дерево. Списки суміжності вершин переглядаються алгоритмом за зростанням міток вершин.\par
\begin{center}
	\adjustimage{max size={0.9\linewidth}{0.9\paperheight}}
	{../10/Traversals/49.png}
\end{center}
\par Які з наведених ребер входять до складу отриманого кістякового дерева?\par
1) (4, 5)\par
2) (4, 6)\par
3) (2, 6)\par
4) (0, 1)\par
5) (1, 5)\par
6) (3, 6)\par
{\vskip 15pt} У формі вибрати номери відповідних ребер.\par
%answer:2 3 4 5
\end {large}}
%end
{\smallskip \begin {small} Затверджено на засіданні кафедри теоретичної кібернетики, протокол \No 12 від   19.05.2020 р.\par\smallskip\smallskip
{\parbox {6.5 cm}{Зав. кафедри \dotfill Ю.В.~Крак}}\hspace {0.7cm}{\hbox to 6.5 cm{ Екзаменатор \dotfill Т.О.~Карнаух}} \end{small}}
%begin
{{\begin {small}\par
{ \center  Київський національний університет імені Тараса Шевченка \par}
{\parbox {5 cm}{Освітня програма \hfill}{\parbox {7 cm}{Прикладна математика, Системний аналіз \hfill}}\parbox {6 cm}{\hfill Семестр 2}}\par
{\parbox {5 cm} {Навчальний предмет \hfill}\parbox {7 cm}{Програмування \hfill}\parbox {6cm}{\hfill}\par}\vspace{-7pt} \end{small}}{\center Екзаменаційний білет \No \textbf 
{ 17 }\par}}
{
\begin {large}
1. Що буде виведено за виконання?
code
try:
    res = 3/4/6*5
    if isinstance(res, bool):
        print(f'bool;{res}')
    elif isinstance(res, int):
        print(f'int;{res}')
    elif isinstance(res, float):
        print(f'float;{res:.2f}')
    else:
        print('???')
except:
    print('ERR')
code

2. Що буде виведено за виконання?
code
a = 3
b = 2
c = 6

def g(a):
    a *= 3
    b = 1
    c = 4
    return a + b + c

a = 5
b = 0
c = 7

try:
    print(g(b), a, b, c, sep=';')
except:
    print('ERR', a, b, c, sep=';')
code

3. Що буде виведено за виконання?
code
try:
    i = 7
    while i>3:
        i += -2
    print(i, end=';')
except:
    print('ERR')
code

4. Що буде виведено за виконання?
code
try:
    for e in range(-2, 12, -2):
        if e>7:
            continue
        print(e, end=';')
    else:
        print(e,end=';')
    print(e)
except:
    print('ERR')
code

5. Що буде виведено за виконання?\par (Дійсні значення, якщо наявні, в цьому завданні будуть виводитись у форматі з фіксованою крапкою та, можливо, з доволі великою кількістю десяткових розрядів. У відповіді для дійсних значень у ЦЬОМУ завданні достатньо вказати один десятковий розряд після крапки, відкинувши решту розрядів без округлення. Наприклад, замість 13.66666666666667 достатньо вказати 13.6, замість 25.23 достатньо вказати 25.2)
code
try:
    print(8, end=';')
    print(5j>=4j, end=';')
    print(6, end=';')
except ValueError:
    print(7, end=';')
except Exception:
    print(4, end=';')
except:
    print('ERR', end=';')
else:
    print(2, end=';')
finally:
    print(3)
code

6. Що буде виведено за виконання?
code
def f(n):
    if n>9:
        f(n//100)
    else:
        print(n%10,end='')

f(123456)
code

7. Що буде виведено за виконання?
code
class A:
    a = 1

    def f1(self): print(2, end=';')
    def _f2(self): print(3, end=';')

    def main(self):
        try: print(self.a, end=';')
        except: print('ERR', end=';')

        try: self.f1()
        except: print('ERR', end=';')

        try: self._f2()
        except: print('ERR', end=';')


class B(A):
    a = 4

    def f1(self): print(5, end=';')
    def _f2(self): print(6, end=';')

try: B().main()
except: print('ERR', end=';')
code

8. Що буде виведено за виконання?
code
def f(arg, n):
    n = 3
    tmp = arg
    del tmp[-7:3]
    return len(tmp)>3

try:
    g = [10,11,12,13,14]
    n = 8
    f(g, n)
    print(n, end=';')
    for el in g:
        print(el, end=';')
except:
    print('ERR')
code

9. Що буде виведено за виконання?
code
class A:
    c = 1
    b = 2

    def __init__(self):
        c = 3

class B(A):
    c = 4

    def __init__(self):
        super().__init__()
        self.a = 7

obj = B()
obj.a = obj.c + 10
B.a = 13

try: print(obj.a, end= ';')
except: print('ERR', end=';')

try: print(A.a, end= ';')
except: print('ERR', end=';')

try: print(B.a, end= ';')
except: print('ERR', end=';')
code

10. 
Для графа на рисунку з вершини 1 запущено алгоритм пошуку в глибину, що будує кістякове дерево. Списки суміжності вершин переглядаються алгоритмом за зростанням міток вершин.\par
\begin{center}
	\adjustimage{max size={0.9\linewidth}{0.9\paperheight}}
	{../10/Traversals/51.png}
\end{center}
\par Які з наведених ребер входять до складу отриманого кістякового дерева?\par
1) (2, 5)\par
2) (0, 2)\par
3) (1, 4)\par
4) (0, 5)\par
5) (1, 3)\par
6) (0, 6)\par
{\vskip 15pt} У формі вибрати номери відповідних ребер.\par
%answer:2 3 4
\end {large}}
%end
{\smallskip \begin {small} Затверджено на засіданні кафедри теоретичної кібернетики, протокол \No 12 від   19.05.2020 р.\par\smallskip\smallskip
{\parbox {6.5 cm}{Зав. кафедри \dotfill Ю.В.~Крак}}\hspace {0.7cm}{\hbox to 6.5 cm{ Екзаменатор \dotfill Т.О.~Карнаух}} \end{small}}
%begin
{{\begin {small}\par
{ \center  Київський національний університет імені Тараса Шевченка \par}
{\parbox {5 cm}{Освітня програма \hfill}{\parbox {7 cm}{Прикладна математика, Системний аналіз \hfill}}\parbox {6 cm}{\hfill Семестр 2}}\par
{\parbox {5 cm} {Навчальний предмет \hfill}\parbox {7 cm}{Програмування \hfill}\parbox {6cm}{\hfill}\par}\vspace{-7pt} \end{small}}{\center Екзаменаційний білет \No \textbf 
{ 18 }\par}}
{
\begin {large}
1. Що буде виведено за виконання?
code
try:
    res = 4**-1.0*True
    if isinstance(res, bool):
        print(f'bool;{res}')
    elif isinstance(res, int):
        print(f'int;{res}')
    elif isinstance(res, float):
        print(f'float;{res:.2f}')
    else:
        print('???')
except:
    print('ERR')
code

2. Що буде виведено за виконання?
code
a = 9
b = 1
c = 0

def f(a):
    a = 2
    b += 2
    c = 1
    return a + b + c

a = 1
b = 9
c = 5

try:
    print(f(a), a, b, c, sep=';')
except:
    print('ERR', a, b, c, sep=';')
code

3. Що буде виведено за виконання?
code
try:
    n = 8
    while n<=10:
        n += 2
    print(n, end=';')
except:
    print('ERR')
code

4. Що буде виведено за виконання?
code
try:
    for d in range(7, 13, 3):
        if d<5:
            break
        print(d, end=';');d = -2
    else:
        print('end',end=';')
    print(d)
except:
    print('ERR')
code

5. Що буде виведено за виконання?\par (Дійсні значення, якщо наявні, в цьому завданні будуть виводитись у форматі з фіксованою крапкою та, можливо, з доволі великою кількістю десяткових розрядів. У відповіді для дійсних значень у ЦЬОМУ завданні достатньо вказати один десятковий розряд після крапки, відкинувши решту розрядів без округлення. Наприклад, замість 13.66666666666667 достатньо вказати 13.6, замість 25.23 достатньо вказати 25.2)
code
try:
    print(1, end=';')
    print(int('c3'), end=';')
    print(4, end=';')
except KeyboardInterrupt:
    print(7, end=';')
except KeyboardInterrupt:
    print(8, end=';')
except:
    print('ERR', end=';')
else:
    print(0, end=';')
finally:
    print(9)
code

6. Що буде виведено за виконання?
code
def f(n):
    if n<=9:
        print(n%10, end='')
    else:
        f(n//10)

f(987651)
code

7. Що буде виведено за виконання?
code
class A:
    c = 1

    def __g1(self): print(2, end=';')
    def _g2(self): print(3, end=';')

    def main(self):
        try: print(self.c, end=';')
        except: print('ERR', end=';')

        try: self.__g1()
        except: print('ERR', end=';')

        try: self._g2()
        except: print('ERR', end=';')


class B(A):
    c = 4

    def __g1(self): print(5, end=';')
    def _g2(self): print(6, end=';')

try: B().main()
except: print('ERR', end=';')
code

8. Що буде виведено за виконання?
code
def f(arg, n):
    n = 1
    tmp = arg
    tmp[0:-3] = 'yz'
    return len(tmp)>3

try:
    c = ['0','1','2','3','4']
    n = 4
    f(c, n)
    print(n, end=';')
    for el in c:
        print(el, end=';')
except:
    print('ERR')
code

9. Що буде виведено за виконання?
code
class A:
    a = 1
    c = 2

    def __init__(self):
        c = 3

class B(A):
    c = 4

    def __init__(self):
        super().__init__()
        self.b = 7

obj = B()
obj.a = obj.b + 10
B.a = 13

try: print(obj.a, end= ';')
except: print('ERR', end=';')

try: print(A.a, end= ';')
except: print('ERR', end=';')

try: print(B.a, end= ';')
except: print('ERR', end=';')
code

10. 
Для графа на рисунку з вершини 5 запущено алгоритм пошуку в глибину, що будує кістякове дерево. Списки суміжності вершин переглядаються алгоритмом за зростанням міток вершин.\par
\begin{center}
	\adjustimage{max size={0.9\linewidth}{0.9\paperheight}}
	{../10/Traversals/21.png}
\end{center}
\par Які з наведених ребер входять до складу отриманого кістякового дерева?\par
1) (1, 5)\par
2) (4, 6)\par
3) (3, 4)\par
4) (1, 3)\par
5) (5, 6)\par
6) (0, 3)\par
{\vskip 15pt} У формі вибрати номери відповідних ребер.\par
%answer:2 4 6
\end {large}}
%end
{\smallskip \begin {small} Затверджено на засіданні кафедри теоретичної кібернетики, протокол \No 12 від   19.05.2020 р.\par\smallskip\smallskip
{\parbox {6.5 cm}{Зав. кафедри \dotfill Ю.В.~Крак}}\hspace {0.7cm}{\hbox to 6.5 cm{ Екзаменатор \dotfill Т.О.~Карнаух}} \end{small}}
%begin
{{\begin {small}\par
{ \center  Київський національний університет імені Тараса Шевченка \par}
{\parbox {5 cm}{Освітня програма \hfill}{\parbox {7 cm}{Прикладна математика, Системний аналіз \hfill}}\parbox {6 cm}{\hfill Семестр 2}}\par
{\parbox {5 cm} {Навчальний предмет \hfill}\parbox {7 cm}{Програмування \hfill}\parbox {6cm}{\hfill}\par}\vspace{-7pt} \end{small}}{\center Екзаменаційний білет \No \textbf 
{ 19 }\par}}
{
\begin {large}
1. Що буде виведено за виконання?
code
try:
    res = 7//3%4*6
    if isinstance(res, bool):
        print(f'bool;{res}')
    elif isinstance(res, int):
        print(f'int;{res}')
    elif isinstance(res, float):
        print(f'float;{res:.2f}')
    else:
        print('???')
except:
    print('ERR')
code

2. Що буде виведено за виконання?
code
a = 6
b = 5
c = 2

def f(b):
    global c
    a = 5
    b = 3
    c = 4
    return a + b + c

a = 4
b = 8
c = 9

try:
    print(f(a), a, b, c, sep=';')
except:
    print('ERR', a, b, c, sep=';')
code

3. Що буде виведено за виконання?
code
try:
    k = 15
    while k>=9:
        k += -3
    print(k, end=';')
except:
    print('ERR')
code

4. Що буде виведено за виконання?
code
try:
    for a in range(15, 3, -1):
        if a<=8:
            continue
        print(a, end=';')
    else:
        print(a,end=';')
    print(a)
except:
    print('ERR')
code

5. Що буде виведено за виконання?\par (Дійсні значення, якщо наявні, в цьому завданні будуть виводитись у форматі з фіксованою крапкою та, можливо, з доволі великою кількістю десяткових розрядів. У відповіді для дійсних значень у ЦЬОМУ завданні достатньо вказати один десятковий розряд після крапки, відкинувши решту розрядів без округлення. Наприклад, замість 13.66666666666667 достатньо вказати 13.6, замість 25.23 достатньо вказати 25.2)
code
try:
    print(6, end=';')
    print(1//1, end=';')
    print(5, end=';')
except ZeroDivisionError:
    print(2, end=';')
except Exception:
    print(4, end=';')
except:
    print('ERR', end=';')
else:
    print(2, end=';')
finally:
    print(1)
code

6. Що буде виведено за виконання?
code
def f(n):
    if n>9:
        return f(n//100)+1
    else:
        return 1

print(f(123456))
code

7. Що буде виведено за виконання?
code
class A:
    b = 1

    def _g1(self): print(2, end=';')
    def _g2(self): print(3, end=';')

    def main(self):
        try: print(self.b, end=';')
        except: print('ERR', end=';')

        try: self._g1()
        except: print('ERR', end=';')

        try: self._g2()
        except: print('ERR', end=';')


class B(A):
    b = 4

    def _g1(self): print(5, end=';')
    def _g2(self): print(6, end=';')

try: B().main()
except: print('ERR', end=';')
code

8. Що буде виведено за виконання?
code
def f(arg, n):
    arg[90] = 4
    n += 2

try:
    n = 7
    d = {90:6, 88:2, 90:3}
    f(d, n)
    print(n, end=';')
    for x in d.values():
        print(x, sep=';', end=';')
except:
    print('ERR')
code

9. Що буде виведено за виконання?
code
class A:
    c = 1
    b = 2

    def __init__(self):
        b = 3

class B(A):
    c = 4

    def __init__(self):
        super().__init__()
        self.a = 7

obj = B()
obj.b = obj.c + 10
B.b = 13

try: print(obj.b, end= ';')
except: print('ERR', end=';')

try: print(A.b, end= ';')
except: print('ERR', end=';')

try: print(B.b, end= ';')
except: print('ERR', end=';')
code

10. 
Для графа на рисунку з вершини 0 запущено алгоритм пошуку в глибину, що будує кістякове дерево. Списки суміжності вершин переглядаються алгоритмом за зростанням міток вершин.\par
\begin{center}
	\adjustimage{max size={0.9\linewidth}{0.9\paperheight}}
	{../10/Traversals/273.png}
\end{center}
\par Які з наведених ребер входять до складу отриманого кістякового дерева?\par
1) (1, 3)\par
2) (2, 3)\par
3) (1, 6)\par
4) (2, 6)\par
5) (1, 2)\par
6) (2, 5)\par
{\vskip 15pt} У формі вибрати номери відповідних ребер.\par
%answer:1 5
\end {large}}
%end
{\smallskip \begin {small} Затверджено на засіданні кафедри теоретичної кібернетики, протокол \No 12 від   19.05.2020 р.\par\smallskip\smallskip
{\parbox {6.5 cm}{Зав. кафедри \dotfill Ю.В.~Крак}}\hspace {0.7cm}{\hbox to 6.5 cm{ Екзаменатор \dotfill Т.О.~Карнаух}} \end{small}}
%begin
{{\begin {small}\par
{ \center  Київський національний університет імені Тараса Шевченка \par}
{\parbox {5 cm}{Освітня програма \hfill}{\parbox {7 cm}{Прикладна математика, Системний аналіз \hfill}}\parbox {6 cm}{\hfill Семестр 2}}\par
{\parbox {5 cm} {Навчальний предмет \hfill}\parbox {7 cm}{Програмування \hfill}\parbox {6cm}{\hfill}\par}\vspace{-7pt} \end{small}}{\center Екзаменаційний білет \No \textbf 
{ 20 }\par}}
{
\begin {large}
1. Що буде виведено за виконання?
code
def f(n):
    print('f', end='')
    return n

try:
    print(f(-7) or f(-1))
except:
    print('ERR')
code

2. Що буде виведено за виконання?
code
a = 2
b = 6
c = 3

def g(b):
    global c
    a = 4
    b -= 3
    c = 5
    return a + b + c

a = 7
b = 8
c = 4

try:
    print(g(a), a, b, c, sep=';')
except:
    print('ERR', a, b, c, sep=';')
code

3. Що буде виведено за виконання?
code
try:
    t = 1
    while t<12:
        t += 1
    print(t, end=';')
except:
    print('ERR')
code

4. Що буде виведено за виконання?
code
try:
    for b in range(8, -2, -3):
        if b>=3:
            break
        print(b, end=';')
    else:
        print('end',end=';')
    print(b)
except:
    print('ERR')
code

5. Що буде виведено за виконання?\par (Дійсні значення, якщо наявні, в цьому завданні будуть виводитись у форматі з фіксованою крапкою та, можливо, з доволі великою кількістю десяткових розрядів. У відповіді для дійсних значень у ЦЬОМУ завданні достатньо вказати один десятковий розряд після крапки, відкинувши решту розрядів без округлення. Наприклад, замість 13.66666666666667 достатньо вказати 13.6, замість 25.23 достатньо вказати 25.2)
code
try:
    print(7, end=';')
    print(6/0.0, end=';')
    print(8, end=';')
except TypeError:
    print(3, end=';')
except ValueError:
    print(0, end=';')
except:
    print('ERR', end=';')
else:
    print(9, end=';')
finally:
    print(5)
code

6. Що буде виведено за виконання?
code
def f(n):
    print(n%10, end='')
    if n>9:
        f(n//10)
 
f(123456)
code

7. Що буде виведено за виконання?
code
class A:
    d = 1

    def _h1(self): print(2, end=';')
    def h2(self): print(3, end=';')

    def main(self):
        try: print(self.d, end=';')
        except: print('ERR', end=';')

        try: self._h1()
        except: print('ERR', end=';')

        try: self.h2()
        except: print('ERR', end=';')


class B(A):
    d = 4

    def _h1(self): print(5, end=';')
    def h2(self): print(6, end=';')

try: B().main()
except: print('ERR', end=';')
code

8. Що буде виведено за виконання?
code
def f(arg, n):
    n = 7
    tmp = arg
    del tmp[7:5]
    return len(tmp)>3

try:
    d = ['0','1','2','3','4']
    n = 8
    f(d, n)
    print(n, end=';')
    for el in d:
        print(el, end=';')
except:
    print('ERR')
code

9. Що буде виведено за виконання?
code
class A:
    a = 1
    c = 2

    def __init__(self):
        a = 3

class B(A):
    c = 4

    def __init__(self):
        super().__init__()
        self.b = 7

obj = B()
obj.a = obj.b + 10
B.a = 13

try: print(obj.a, end= ';')
except: print('ERR', end=';')

try: print(A.a, end= ';')
except: print('ERR', end=';')

try: print(B.a, end= ';')
except: print('ERR', end=';')
code

10. 
Для графа на рисунку з вершини 4 запущено алгоритм пошуку в глибину, що будує кістякове дерево. Списки суміжності вершин переглядаються алгоритмом за зростанням міток вершин.\par
\begin{center}
	\adjustimage{max size={0.9\linewidth}{0.9\paperheight}}
	{../10/Traversals/133.png}
\end{center}
\par Які з наведених ребер входять до складу отриманого кістякового дерева?\par
1) (2, 6)\par
2) (2, 3)\par
3) (1, 6)\par
4) (2, 4)\par
5) (5, 6)\par
6) (1, 4)\par
{\vskip 15pt} У формі вибрати номери відповідних ребер.\par
%answer:1 2 5 6
\end {large}}
%end
{\smallskip \begin {small} Затверджено на засіданні кафедри теоретичної кібернетики, протокол \No 12 від   19.05.2020 р.\par\smallskip\smallskip
{\parbox {6.5 cm}{Зав. кафедри \dotfill Ю.В.~Крак}}\hspace {0.7cm}{\hbox to 6.5 cm{ Екзаменатор \dotfill Т.О.~Карнаух}} \end{small}}
%begin
{{\begin {small}\par
{ \center  Київський національний університет імені Тараса Шевченка \par}
{\parbox {5 cm}{Освітня програма \hfill}{\parbox {7 cm}{Прикладна математика, Системний аналіз \hfill}}\parbox {6 cm}{\hfill Семестр 2}}\par
{\parbox {5 cm} {Навчальний предмет \hfill}\parbox {7 cm}{Програмування \hfill}\parbox {6cm}{\hfill}\par}\vspace{-7pt} \end{small}}{\center Екзаменаційний білет \No \textbf 
{ 21 }\par}}
{
\begin {large}
1. Що буде виведено за виконання?
code
try:
    res = 4/9//5*4
    if isinstance(res, bool):
        print(f'bool;{res}')
    elif isinstance(res, int):
        print(f'int;{res}')
    elif isinstance(res, float):
        print(f'float;{res:.2f}')
    else:
        print('???')
except:
    print('ERR')
code

2. Що буде виведено за виконання?
code
a = 5
b = 9
c = 7

def h(a):
    global c
    a *= 2
    b = 5
    c = 3
    return a + b + c

a = 6
b = 1
c = 4

try:
    print(h(b), a, b, c, sep=';')
except:
    print('ERR', a, b, c, sep=';')
code

3. Що буде виведено за виконання?
code
try:
    t = 0
    while t<9:
        t += 2
    print(t, end=';')
except:
    print('ERR')
code

4. Що буде виведено за виконання?
code
try:
    for d in range(10, -8, -3):
        if d<=7:
            continue
        print(d, end=';');d = 2
    else:
        print(d,end=';')
    print(d)
except:
    print('ERR')
code

5. Що буде виведено за виконання?\par (Дійсні значення, якщо наявні, в цьому завданні будуть виводитись у форматі з фіксованою крапкою та, можливо, з доволі великою кількістю десяткових розрядів. У відповіді для дійсних значень у ЦЬОМУ завданні достатньо вказати один десятковий розряд після крапки, відкинувши решту розрядів без округлення. Наприклад, замість 13.66666666666667 достатньо вказати 13.6, замість 25.23 достатньо вказати 25.2)
code
try:
    print(2, end=';')
    print(9%0, end=';')
    print(1, end=';')
except TypeError:
    print(5, end=';')
except ZeroDivisionError:
    print(4, end=';')
except:
    print('ERR', end=';')
else:
    print(0, end=';')
finally:
    print(7)
code

6. Що буде виведено за виконання?
code
def f(n):
    if n>2:
        f(n//2)
    print(n%2, end='')
 
f(16)
code

7. Що буде виведено за виконання?
code
class A:
    a = 1

    def _f1(self): print(2, end=';')
    def _f2(self): print(3, end=';')

    def main(self):
        try: print(self.a, end=';')
        except: print('ERR', end=';')

        try: self._f1()
        except: print('ERR', end=';')

        try: self._f2()
        except: print('ERR', end=';')


class B(A):
    a = 4

    def _f1(self): print(5, end=';')
    def _f2(self): print(6, end=';')

try: B().main()
except: print('ERR', end=';')
code

8. Що буде виведено за виконання?
code
def f(arg, n):
    arg[49] = 4
    n = 5
    arg = list(arg.values())

try:
    n = 3
    d = {87:4, 69:8, 87:5}
    f(d, n)
    print(n, end=';')
    for x in d.keys():
        print(x, sep=';', end=';')
except:
    print('ERR')
code

9. Що буде виведено за виконання?
code
class A:
    c = 1
    a = 2

    def __init__(self):
        a = 3

class B(A):
    a = 4

    def __init__(self):
        super().__init__()
        self.b = 7

obj = B()
obj.c = obj.a + 10
B.c = 13

try: print(obj.c, end= ';')
except: print('ERR', end=';')

try: print(A.c, end= ';')
except: print('ERR', end=';')

try: print(B.c, end= ';')
except: print('ERR', end=';')
code

10. 
Для графа на рисунку з вершини 1 запущено алгоритм пошуку в глибину, що будує кістякове дерево. Списки суміжності вершин переглядаються алгоритмом за зростанням міток вершин.\par
\begin{center}
	\adjustimage{max size={0.9\linewidth}{0.9\paperheight}}
	{../10/Traversals/55.png}
\end{center}
\par Які з наведених ребер входять до складу отриманого кістякового дерева?\par
1) (2, 6)\par
2) (0, 2)\par
3) (4, 6)\par
4) (3, 6)\par
5) (2, 5)\par
6) (2, 4)\par
{\vskip 15pt} У формі вибрати номери відповідних ребер.\par
%answer:2 3 6
\end {large}}
%end
{\smallskip \begin {small} Затверджено на засіданні кафедри теоретичної кібернетики, протокол \No 12 від   19.05.2020 р.\par\smallskip\smallskip
{\parbox {6.5 cm}{Зав. кафедри \dotfill Ю.В.~Крак}}\hspace {0.7cm}{\hbox to 6.5 cm{ Екзаменатор \dotfill Т.О.~Карнаух}} \end{small}}
%begin
{{\begin {small}\par
{ \center  Київський національний університет імені Тараса Шевченка \par}
{\parbox {5 cm}{Освітня програма \hfill}{\parbox {7 cm}{Прикладна математика, Системний аналіз \hfill}}\parbox {6 cm}{\hfill Семестр 2}}\par
{\parbox {5 cm} {Навчальний предмет \hfill}\parbox {7 cm}{Програмування \hfill}\parbox {6cm}{\hfill}\par}\vspace{-7pt} \end{small}}{\center Екзаменаційний білет \No \textbf 
{ 22 }\par}}
{
\begin {large}
1. Що буде виведено за виконання?
code
def f(n):
    print('f', end='')
    return n

try:
    print(f(4) and f(-2))
except:
    print('ERR')
code

2. Що буде виведено за виконання?
code
a = 3
b = 1
c = 8

def h(a):
    a = 3
    b = 5
    c = 1
    return a + b + c

a = 2
b = 4
c = 0

try:
    print(h(b), a, b, c, sep=';')
except:
    print('ERR', a, b, c, sep=';')
code

3. Що буде виведено за виконання?
code
try:
    j = 10
    while j>=5:
        j += -2
    print(j, end=';')
except:
    print('ERR')
code

4. Що буде виведено за виконання?
code
try:
    for f in range(-11, -3, 3):
        if f>=4:
            break
        print(f, end=';')
    else:
        print(f,end=';')
    print(f)
except:
    print('ERR')
code

5. Що буде виведено за виконання?\par (Дійсні значення, якщо наявні, в цьому завданні будуть виводитись у форматі з фіксованою крапкою та, можливо, з доволі великою кількістю десяткових розрядів. У відповіді для дійсних значень у ЦЬОМУ завданні достатньо вказати один десятковий розряд після крапки, відкинувши решту розрядів без округлення. Наприклад, замість 13.66666666666667 достатньо вказати 13.6, замість 25.23 достатньо вказати 25.2)
code
try:
    print(6, end=';')
    print(int('d3'), end=';')
    print(8, end=';')
except KeyboardInterrupt:
    print(9, end=';')
except ValueError:
    print(6, end=';')
except:
    print('ERR', end=';')
else:
    print(1, end=';')
finally:
    print(4)
code

6. Що буде виведено за виконання?
code
def f(s):
    s1 = s
    for i in range(len(s)):
        s1[i] += 1
        return s1

s = [1, 2, 3]
f(s)
print(s)
code

7. Що буде виведено за виконання?
code
class A:
    b = 1

    def _g1(self): print(2, end=';')
    def __g2(self): print(3, end=';')

    def main(self):
        try: print(self.b, end=';')
        except: print('ERR', end=';')

        try: self._g1()
        except: print('ERR', end=';')

        try: self.__g2()
        except: print('ERR', end=';')


class B(A):
    b = 4

    def _g1(self): print(5, end=';')
    def __g2(self): print(6, end=';')

try: B().main()
except: print('ERR', end=';')
code

8. Що буде виведено за виконання?
code
def f(arg, n):
    n = 5
    tmp = arg
    tmp[5:] = 'bd'
    return len(tmp)>3

try:
    b = [10,11,12,13,14]
    n = 2
    f(b, n)
    print(n, end=';')
    for el in b:
        print(el, end=';')
except:
    print('ERR')
code

9. Що буде виведено за виконання?
code
class A:
    b = 1
    c = 2

    def __init__(self):
        b = 3

class B(A):
    b = 4

    def __init__(self):
        super().__init__()
        self.a = 7

obj = B()
obj.c = obj.b + 10
B.c = 13

try: print(obj.c, end= ';')
except: print('ERR', end=';')

try: print(A.c, end= ';')
except: print('ERR', end=';')

try: print(B.c, end= ';')
except: print('ERR', end=';')
code

10. 
Для графа на рисунку з вершини 4 запущено алгоритм пошуку в глибину, що будує кістякове дерево. Списки суміжності вершин переглядаються алгоритмом за зростанням міток вершин.\par
\begin{center}
	\adjustimage{max size={0.9\linewidth}{0.9\paperheight}}
	{../10/Traversals/72.png}
\end{center}
\par Які з наведених ребер входять до складу отриманого кістякового дерева?\par
1) (1, 4)\par
2) (1, 6)\par
3) (4, 5)\par
4) (1, 3)\par
5) (0, 4)\par
6) (3, 4)\par
{\vskip 15pt} У формі вибрати номери відповідних ребер.\par
%answer:2 4 5
\end {large}}
%end
{\smallskip \begin {small} Затверджено на засіданні кафедри теоретичної кібернетики, протокол \No 12 від   19.05.2020 р.\par\smallskip\smallskip
{\parbox {6.5 cm}{Зав. кафедри \dotfill Ю.В.~Крак}}\hspace {0.7cm}{\hbox to 6.5 cm{ Екзаменатор \dotfill Т.О.~Карнаух}} \end{small}}
%begin
{{\begin {small}\par
{ \center  Київський національний університет імені Тараса Шевченка \par}
{\parbox {5 cm}{Освітня програма \hfill}{\parbox {7 cm}{Прикладна математика, Системний аналіз \hfill}}\parbox {6 cm}{\hfill Семестр 2}}\par
{\parbox {5 cm} {Навчальний предмет \hfill}\parbox {7 cm}{Програмування \hfill}\parbox {6cm}{\hfill}\par}\vspace{-7pt} \end{small}}{\center Екзаменаційний білет \No \textbf 
{ 23 }\par}}
{
\begin {large}
1. Що буде виведено за виконання?
code
try:
    res = False>="6j"
    if isinstance(res, bool):
        print(f'bool;{res}')
    elif isinstance(res, int):
        print(f'int;{res}')
    elif isinstance(res, float):
        print(f'float;{res:.2f}')
    else:
        print('???')
except:
    print('ERR')
code

2. Що буде виведено за виконання?
code
a = 7
b = 6
c = 9

def g(b):
    global c
    a = 4
    b = 2
    c = 4
    return a + b + c

a = 5
b = 2
c = 7

try:
    print(g(a), a, b, c, sep=';')
except:
    print('ERR', a, b, c, sep=';')
code

3. Що буде виведено за виконання?
code
try:
    m = 11
    while m>6:
        m += -3
    print(m, end=';')
except:
    print('ERR')
code

4. Що буде виведено за виконання?
code
try:
    for e in range(2, -5, -1):
        if e<5:
            continue
        print(e, end=';');e = 9
    else:
        print(e,end=';')
    print(e)
except:
    print('ERR')
code

5. Що буде виведено за виконання?\par (Дійсні значення, якщо наявні, в цьому завданні будуть виводитись у форматі з фіксованою крапкою та, можливо, з доволі великою кількістю десяткових розрядів. У відповіді для дійсних значень у ЦЬОМУ завданні достатньо вказати один десятковий розряд після крапки, відкинувши решту розрядів без округлення. Наприклад, замість 13.66666666666667 достатньо вказати 13.6, замість 25.23 достатньо вказати 25.2)
code
try:
    print(8, end=';')
    print(int('0'), end=';')
    print(5, end=';')
except Exception:
    print(3, end=';')
except ValueError:
    print(7, end=';')
except:
    print('ERR', end=';')
else:
    print(2, end=';')
finally:
    print(0)
code

6. Що буде виведено за виконання?
code
def f(n):
    if n>9:
        f(n//10)
    print(n%10, end='')
 
f(543216)
code

7. Що буде виведено за виконання?
code
class A:
    c = 1

    def __f1(self): print(2, end=';')
    def _f2(self): print(3, end=';')

    def main(self):
        try: print(self.c, end=';')
        except: print('ERR', end=';')

        try: self.__f1()
        except: print('ERR', end=';')

        try: self._f2()
        except: print('ERR', end=';')


class B(A):
    c = 4

    def __f1(self): print(5, end=';')
    def _f2(self): print(6, end=';')

try: B().main()
except: print('ERR', end=';')
code

8. Що буде виведено за виконання?
code
def f(arg, n):
    arg[77] = 4
    n = 9
    arg = tuple(arg.items())

try:
    n = 0
    s = {30:0, 70:6, 50:7, 50:2}
    f(s, n)
    print(n, end=';')
    for x, y in s.items():
        print(x, y, sep=';', end=';')
except:
    print('ERR')
code

9. Що буде виведено за виконання?
code
class A:
    a = 1
    b = 2

    def __init__(self):
        a = 3

class B(A):
    a = 4

    def __init__(self):
        super().__init__()
        self.c = 7

obj = B()
obj.a = obj.b + 10
B.a = 13

try: print(obj.a, end= ';')
except: print('ERR', end=';')

try: print(A.a, end= ';')
except: print('ERR', end=';')

try: print(B.a, end= ';')
except: print('ERR', end=';')
code

10. 
Для графа на рисунку з вершини 4 запущено алгоритм пошуку в глибину, що будує кістякове дерево. Списки суміжності вершин переглядаються алгоритмом за зростанням міток вершин.\par
\begin{center}
	\adjustimage{max size={0.9\linewidth}{0.9\paperheight}}
	{../10/Traversals/58.png}
\end{center}
\par Які з наведених ребер входять до складу отриманого кістякового дерева?\par
1) (2, 4)\par
2) (0, 6)\par
3) (0, 1)\par
4) (4, 5)\par
5) (1, 6)\par
6) (0, 3)\par
{\vskip 15pt} У формі вибрати номери відповідних ребер.\par
%answer:3 6
\end {large}}
%end
{\smallskip \begin {small} Затверджено на засіданні кафедри теоретичної кібернетики, протокол \No 12 від   19.05.2020 р.\par\smallskip\smallskip
{\parbox {6.5 cm}{Зав. кафедри \dotfill Ю.В.~Крак}}\hspace {0.7cm}{\hbox to 6.5 cm{ Екзаменатор \dotfill Т.О.~Карнаух}} \end{small}}
%begin
{{\begin {small}\par
{ \center  Київський національний університет імені Тараса Шевченка \par}
{\parbox {5 cm}{Освітня програма \hfill}{\parbox {7 cm}{Прикладна математика, Системний аналіз \hfill}}\parbox {6 cm}{\hfill Семестр 2}}\par
{\parbox {5 cm} {Навчальний предмет \hfill}\parbox {7 cm}{Програмування \hfill}\parbox {6cm}{\hfill}\par}\vspace{-7pt} \end{small}}{\center Екзаменаційний білет \No \textbf 
{ 24 }\par}}
{
\begin {large}
1. Що буде виведено за виконання?
code
try:
    res = 9.0=="True"
    if isinstance(res, bool):
        print(f'bool;{res}')
    elif isinstance(res, int):
        print(f'int;{res}')
    elif isinstance(res, float):
        print(f'float;{res:.2f}')
    else:
        print('???')
except:
    print('ERR')
code

2. Що буде виведено за виконання?
code
a = 2
b = 3
c = 0

def f(b):
    a = 1
    b += 4
    c = 1
    return a + b + c

a = 8
b = 8
c = 0

try:
    print(f(b), a, b, c, sep=';')
except:
    print('ERR', a, b, c, sep=';')
code

3. Що буде виведено за виконання?
code
try:
    n = 12
    while n>=3:
        n += -2
    print(n, end=';')
except:
    print('ERR')
code

4. Що буде виведено за виконання?
code
try:
    for c in range(11, -10, -2):
        if c>6:
            continue
        print(c, end=';')
        if c<=6:
            break
    else:
        print(c,end=';')
    print(c)
except:
    print('ERR')
code

5. Що буде виведено за виконання?\par (Дійсні значення, якщо наявні, в цьому завданні будуть виводитись у форматі з фіксованою крапкою та, можливо, з доволі великою кількістю десяткових розрядів. У відповіді для дійсних значень у ЦЬОМУ завданні достатньо вказати один десятковий розряд після крапки, відкинувши решту розрядів без округлення. Наприклад, замість 13.66666666666667 достатньо вказати 13.6, замість 25.23 достатньо вказати 25.2)
code
try:
    print(9, end=';')
    print(6/2, end=';')
    print(8, end=';')
except Exception:
    print(5, end=';')
except KeyboardInterrupt:
    print(1, end=';')
except:
    print('ERR', end=';')
else:
    print(7, end=';')
finally:
    print(3)
code

6. Що буде виведено за виконання?
code
def f(n):
    if n<=9:
        print(n%10,end='')
    else:
        f(n//10)

f(543216)
code

7. Що буде виведено за виконання?
code
class A:
    d = 1

    def _h1(self): print(2, end=';')
    def h2(self): print(3, end=';')

    def main(self):
        try: print(self.d, end=';')
        except: print('ERR', end=';')

        try: self._h1()
        except: print('ERR', end=';')

        try: self.h2()
        except: print('ERR', end=';')


class B(A):
    d = 4

    def _h1(self): print(5, end=';')
    def h2(self): print(6, end=';')

try: B().main()
except: print('ERR', end=';')
code

8. Що буде виведено за виконання?
code
def f(arg, n):
    arg[15] = 4
    n = 6
    arg = set(arg.keys())

try:
    n = 1
    t = {37:9, 51:4, 15:2, 54:8}
    f(t, n)
    print(n, end=';')
    for x in t.values():
        print(x, sep=';', end=';')
except:
    print('ERR')
code

9. Що буде виведено за виконання?
code
class A:
    b = 1
    a = 2

    def __init__(self):
        a = 3

class B(A):
    a = 4

    def __init__(self):
        super().__init__()
        self.c = 7

obj = B()
obj.c = obj.a + 10
B.c = 13

try: print(obj.c, end= ';')
except: print('ERR', end=';')

try: print(A.c, end= ';')
except: print('ERR', end=';')

try: print(B.c, end= ';')
except: print('ERR', end=';')
code

10. 
Для графа на рисунку з вершини 2 запущено алгоритм пошуку в глибину, що будує кістякове дерево. Списки суміжності вершин переглядаються алгоритмом за зростанням міток вершин.\par
\begin{center}
	\adjustimage{max size={0.9\linewidth}{0.9\paperheight}}
	{../10/Traversals/96.png}
\end{center}
\par Які з наведених ребер входять до складу отриманого кістякового дерева?\par
1) (1, 6)\par
2) (0, 5)\par
3) (1, 2)\par
4) (0, 1)\par
5) (2, 3)\par
6) (3, 5)\par
{\vskip 15pt} У формі вибрати номери відповідних ребер.\par
%answer:1 2 3 4 6
\end {large}}
%end
{\smallskip \begin {small} Затверджено на засіданні кафедри теоретичної кібернетики, протокол \No 12 від   19.05.2020 р.\par\smallskip\smallskip
{\parbox {6.5 cm}{Зав. кафедри \dotfill Ю.В.~Крак}}\hspace {0.7cm}{\hbox to 6.5 cm{ Екзаменатор \dotfill Т.О.~Карнаух}} \end{small}}
%begin
{{\begin {small}\par
{ \center  Київський національний університет імені Тараса Шевченка \par}
{\parbox {5 cm}{Освітня програма \hfill}{\parbox {7 cm}{Прикладна математика, Системний аналіз \hfill}}\parbox {6 cm}{\hfill Семестр 2}}\par
{\parbox {5 cm} {Навчальний предмет \hfill}\parbox {7 cm}{Програмування \hfill}\parbox {6cm}{\hfill}\par}\vspace{-7pt} \end{small}}{\center Екзаменаційний білет \No \textbf 
{ 25 }\par}}
{
\begin {large}
1. Що буде виведено за виконання?
code
try:
    res = 6.0>"False"
    if isinstance(res, bool):
        print(f'bool;{res}')
    elif isinstance(res, int):
        print(f'int;{res}')
    elif isinstance(res, float):
        print(f'float;{res:.2f}')
    else:
        print('???')
except:
    print('ERR')
code

2. Що буде виведено за виконання?
code
a = 7
b = 4
c = 5

def h(a):
    a += 2
    b = 3
    c = 4
    return a + b + c

a = 6
b = 3
c = 2

try:
    print(h(a), a, b, c, sep=';')
except:
    print('ERR', a, b, c, sep=';')
code

3. Що буде виведено за виконання?
code
try:
    k = 6
    while k<=6:
        k += 1
    print(k, end=';')
except:
    print('ERR')
code

4. Що буде виведено за виконання?
code
try:
    for f in range(0, 14, 2):
        if f>=8:
            break
        print(f, end=';');f = 4
    else:
        print(f,end=';')
    print(f)
except:
    print('ERR')
code

5. Що буде виведено за виконання?\par (Дійсні значення, якщо наявні, в цьому завданні будуть виводитись у форматі з фіксованою крапкою та, можливо, з доволі великою кількістю десяткових розрядів. У відповіді для дійсних значень у ЦЬОМУ завданні достатньо вказати один десятковий розряд після крапки, відкинувши решту розрядів без округлення. Наприклад, замість 13.66666666666667 достатньо вказати 13.6, замість 25.23 достатньо вказати 25.2)
code
try:
    print(2, end=';')
    print(4j==5j, end=';')
    print(8, end=';')
except TypeError:
    print(9, end=';')
except ZeroDivisionError:
    print(2, end=';')
except:
    print('ERR', end=';')
else:
    print(7, end=';')
finally:
    print(3)
code

6. Що буде виведено за виконання?
code
def f(n):
    if n>9:
        f(n//100)
    print(n%10,end='')
 
f(987654)
code

7. Що буде виведено за виконання?
code
class A:
    a = 1

    def f1(self): print(2, end=';')
    def _f2(self): print(3, end=';')

    def main(self):
        try: print(self.a, end=';')
        except: print('ERR', end=';')

        try: self.f1()
        except: print('ERR', end=';')

        try: self._f2()
        except: print('ERR', end=';')


class B(A):
    a = 4

    def f1(self): print(5, end=';')
    def _f2(self): print(6, end=';')

try: B().main()
except: print('ERR', end=';')
code

8. Що буде виведено за виконання?
code
def f(arg, n):
    arg[75] = 7
    n = 8
    arg = list(arg.keys())

try:
    n = 2
    d = {78:9, 77:7, 14:1, 14:9}
    f(d, n)
    print(n, end=';')
    for x in d.items():
        print(*x, sep=';', end=';')
except:
    print('ERR')
code

9. Що буде виведено за виконання?
code
class A:
    c = 1
    a = 2

    def __init__(self):
        c = 3

class B(A):
    c = 4

    def __init__(self):
        super().__init__()
        self.b = 7

obj = B()
obj.a = obj.b + 10
B.a = 13

try: print(obj.a, end= ';')
except: print('ERR', end=';')

try: print(A.a, end= ';')
except: print('ERR', end=';')

try: print(B.a, end= ';')
except: print('ERR', end=';')
code

10. 
Для графа на рисунку з вершини 2 запущено алгоритм пошуку в глибину, що будує кістякове дерево. Списки суміжності вершин переглядаються алгоритмом за зростанням міток вершин.\par
\begin{center}
	\adjustimage{max size={0.9\linewidth}{0.9\paperheight}}
	{../10/Traversals/291.png}
\end{center}
\par Які з наведених ребер входять до складу отриманого кістякового дерева?\par
1) (3, 6)\par
2) (2, 4)\par
3) (0, 2)\par
4) (2, 3)\par
5) (4, 5)\par
6) (3, 4)\par
{\vskip 15pt} У формі вибрати номери відповідних ребер.\par
%answer:3 4 6
\end {large}}
%end
{\smallskip \begin {small} Затверджено на засіданні кафедри теоретичної кібернетики, протокол \No 12 від   19.05.2020 р.\par\smallskip\smallskip
{\parbox {6.5 cm}{Зав. кафедри \dotfill Ю.В.~Крак}}\hspace {0.7cm}{\hbox to 6.5 cm{ Екзаменатор \dotfill Т.О.~Карнаух}} \end{small}}
%begin
{{\begin {small}\par
{ \center  Київський національний університет імені Тараса Шевченка \par}
{\parbox {5 cm}{Освітня програма \hfill}{\parbox {7 cm}{Прикладна математика, Системний аналіз \hfill}}\parbox {6 cm}{\hfill Семестр 2}}\par
{\parbox {5 cm} {Навчальний предмет \hfill}\parbox {7 cm}{Програмування \hfill}\parbox {6cm}{\hfill}\par}\vspace{-7pt} \end{small}}{\center Екзаменаційний білет \No \textbf 
{ 26 }\par}}
{
\begin {large}
1. Що буде виведено за виконання?
code
try:
    res = 7//6*8*9
    if isinstance(res, bool):
        print(f'bool;{res}')
    elif isinstance(res, int):
        print(f'int;{res}')
    elif isinstance(res, float):
        print(f'float;{res:.2f}')
    else:
        print('???')
except:
    print('ERR')
code

2. Що буде виведено за виконання?
code
a = 8
b = 1
c = 4

def h(b):
    global c
    a = 1
    b = 5
    c = 2
    return a + b + c

a = 0
b = 9
c = 6

try:
    print(h(a), a, b, c, sep=';')
except:
    print('ERR', a, b, c, sep=';')
code

3. Що буде виведено за виконання?
code
try:
    t = 12
    while t>5:
        t += -3
    print(t, end=';')
except:
    print('ERR')
code

4. Що буде виведено за виконання?
code
try:
    for d in range(14, 8, -1):
        if d>3:
            continue
        print(d, end=';')
    else:
        print(d,end=';')
    print(d)
except:
    print('ERR')
code

5. Що буде виведено за виконання?\par (Дійсні значення, якщо наявні, в цьому завданні будуть виводитись у форматі з фіксованою крапкою та, можливо, з доволі великою кількістю десяткових розрядів. У відповіді для дійсних значень у ЦЬОМУ завданні достатньо вказати один десятковий розряд після крапки, відкинувши решту розрядів без округлення. Наприклад, замість 13.66666666666667 достатньо вказати 13.6, замість 25.23 достатньо вказати 25.2)
code
try:
    print(5, end=';')
    print(0%0.0, end=';')
    print(6, end=';')
except KeyboardInterrupt:
    print(4, end=';')
except Exception:
    print(1, end=';')
except:
    print('ERR', end=';')
else:
    print(4, end=';')
finally:
    print(6)
code

6. Що буде виведено за виконання?
code
def f(n):
    if n>2:
        f(n//2)
    else:
        print(n%2, end='')

f(16)
code

7. Що буде виведено за виконання?
code
class A:
    d = 1

    def _h1(self): print(2, end=';')
    def __h2(self): print(3, end=';')

    def main(self):
        try: print(self.d, end=';')
        except: print('ERR', end=';')

        try: self._h1()
        except: print('ERR', end=';')

        try: self.__h2()
        except: print('ERR', end=';')


class B(A):
    d = 4

    def _h1(self): print(5, end=';')
    def __h2(self): print(6, end=';')

try: B().main()
except: print('ERR', end=';')
code

8. Що буде виведено за виконання?
code
def f(arg, n):
    arg[69] = 8
    n *= 3

try:
    n = 3
    s = {33:9, 30:1, 35:7, 35:8}
    f(s, n)
    print(n, end=';')
    for x in s.items():
        print(*x, sep=';', end=';')
except:
    print('ERR')
code

9. Що буде виведено за виконання?
code
class A:
    a = 1
    c = 2

    def __init__(self):
        c = 3

class B(A):
    c = 4

    def __init__(self):
        super().__init__()
        self.b = 7

obj = B()
obj.c = obj.c + 10
B.c = 13

try: print(obj.c, end= ';')
except: print('ERR', end=';')

try: print(A.c, end= ';')
except: print('ERR', end=';')

try: print(B.c, end= ';')
except: print('ERR', end=';')
code

10. 
Для графа на рисунку з вершини 3 запущено алгоритм пошуку в глибину, що будує кістякове дерево. Списки суміжності вершин переглядаються алгоритмом за зростанням міток вершин.\par
\begin{center}
	\adjustimage{max size={0.9\linewidth}{0.9\paperheight}}
	{../10/Traversals/231.png}
\end{center}
\par Які з наведених ребер входять до складу отриманого кістякового дерева?\par
1) (1, 4)\par
2) (3, 4)\par
3) (4, 5)\par
4) (1, 5)\par
5) (0, 6)\par
6) (4, 6)\par
{\vskip 15pt} У формі вибрати номери відповідних ребер.\par
%answer:2 4
\end {large}}
%end
{\smallskip \begin {small} Затверджено на засіданні кафедри теоретичної кібернетики, протокол \No 12 від   19.05.2020 р.\par\smallskip\smallskip
{\parbox {6.5 cm}{Зав. кафедри \dotfill Ю.В.~Крак}}\hspace {0.7cm}{\hbox to 6.5 cm{ Екзаменатор \dotfill Т.О.~Карнаух}} \end{small}}
%begin
{{\begin {small}\par
{ \center  Київський національний університет імені Тараса Шевченка \par}
{\parbox {5 cm}{Освітня програма \hfill}{\parbox {7 cm}{Прикладна математика, Системний аналіз \hfill}}\parbox {6 cm}{\hfill Семестр 2}}\par
{\parbox {5 cm} {Навчальний предмет \hfill}\parbox {7 cm}{Програмування \hfill}\parbox {6cm}{\hfill}\par}\vspace{-7pt} \end{small}}{\center Екзаменаційний білет \No \textbf 
{ 27 }\par}}
{
\begin {large}
1. Що буде виведено за виконання?
code
try:
    res = 5<=True or not 2.0>7 and 8.0!=8
    if isinstance(res, bool):
        print(f'bool;{res}')
    elif isinstance(res, int):
        print(f'int;{res}')
    elif isinstance(res, float):
        print(f'float;{res:.2f}')
    else:
        print('???')
except:
    print('ERR')
code

2. Що буде виведено за виконання?
code
a = 5
b = 3
c = 9

def g(a):
    a = 3
    b = 4
    c = 2
    return a + b + c

a = 4
b = 8
c = 3

try:
    print(g(a), a, b, c, sep=';')
except:
    print('ERR', a, b, c, sep=';')
code

3. Що буде виведено за виконання?
code
try:
    j = 14
    while j>=0:
        j += -2
    print(j, end=';')
except:
    print('ERR')
code

4. Що буде виведено за виконання?
code
try:
    for b in range(-8, 9, 2):
        if b<5:
            continue
        print(b, end=';')
    else:
        print('end',end=';')
    print(b)
except:
    print('ERR')
code

5. Що буде виведено за виконання?\par (Дійсні значення, якщо наявні, в цьому завданні будуть виводитись у форматі з фіксованою крапкою та, можливо, з доволі великою кількістю десяткових розрядів. У відповіді для дійсних значень у ЦЬОМУ завданні достатньо вказати один десятковий розряд після крапки, відкинувши решту розрядів без округлення. Наприклад, замість 13.66666666666667 достатньо вказати 13.6, замість 25.23 достатньо вказати 25.2)
code
try:
    print(1, end=';')
    print(int('7'), end=';')
    print(3, end=';')
except TypeError:
    print(9, end=';')
except ValueError:
    print(0, end=';')
except:
    print('ERR', end=';')
else:
    print(5, end=';')
finally:
    print(8)
code

6. Що буде виведено за виконання?
code
def f(n):
    if n>9:
        f(n//100)
        print(n%10, end='')
 
f(123456)
code

7. Що буде виведено за виконання?
code
class A:
    c = 1

    def __g1(self): print(2, end=';')
    def _g2(self): print(3, end=';')

    def main(self):
        try: print(self.c, end=';')
        except: print('ERR', end=';')

        try: self.__g1()
        except: print('ERR', end=';')

        try: self._g2()
        except: print('ERR', end=';')


class B(A):
    c = 4

    def __g1(self): print(5, end=';')
    def _g2(self): print(6, end=';')

try: B().main()
except: print('ERR', end=';')
code

8. Що буде виведено за виконання?
code
def f(arg, n):
    arg[61] = 6
    n = 5
    arg = tuple(arg.values())

try:
    n = 0
    s = {49:8, 31:6, 49:6}
    f(s, n)
    print(n, end=';')
    for x in s.values():
        print(x, sep=';', end=';')
except:
    print('ERR')
code

9. Що буде виведено за виконання?
code
class A:
    b = 1
    c = 2

    def __init__(self):
        c = 3

class B(A):
    b = 4

    def __init__(self):
        super().__init__()
        self.a = 7

obj = B()
obj.a = obj.b + 10
B.a = 13

try: print(obj.a, end= ';')
except: print('ERR', end=';')

try: print(A.a, end= ';')
except: print('ERR', end=';')

try: print(B.a, end= ';')
except: print('ERR', end=';')
code

10. 
Для графа на рисунку з вершини 5 запущено алгоритм пошуку в глибину, що будує кістякове дерево. Списки суміжності вершин переглядаються алгоритмом за зростанням міток вершин.\par
\begin{center}
	\adjustimage{max size={0.9\linewidth}{0.9\paperheight}}
	{../10/Traversals/158.png}
\end{center}
\par Які з наведених ребер входять до складу отриманого кістякового дерева?\par
1) (3, 4)\par
2) (2, 6)\par
3) (2, 4)\par
4) (2, 5)\par
5) (5, 6)\par
6) (0, 4)\par
{\vskip 15pt} У формі вибрати номери відповідних ребер.\par
%answer:4 6
\end {large}}
%end
{\smallskip \begin {small} Затверджено на засіданні кафедри теоретичної кібернетики, протокол \No 12 від   19.05.2020 р.\par\smallskip\smallskip
{\parbox {6.5 cm}{Зав. кафедри \dotfill Ю.В.~Крак}}\hspace {0.7cm}{\hbox to 6.5 cm{ Екзаменатор \dotfill Т.О.~Карнаух}} \end{small}}
%begin
{{\begin {small}\par
{ \center  Київський національний університет імені Тараса Шевченка \par}
{\parbox {5 cm}{Освітня програма \hfill}{\parbox {7 cm}{Прикладна математика, Системний аналіз \hfill}}\parbox {6 cm}{\hfill Семестр 2}}\par
{\parbox {5 cm} {Навчальний предмет \hfill}\parbox {7 cm}{Програмування \hfill}\parbox {6cm}{\hfill}\par}\vspace{-7pt} \end{small}}{\center Екзаменаційний білет \No \textbf 
{ 28 }\par}}
{
\begin {large}
1. Що буде виведено за виконання?
code
try:
    res = 9**1.0*-2
    if isinstance(res, bool):
        print(f'bool;{res}')
    elif isinstance(res, int):
        print(f'int;{res}')
    elif isinstance(res, float):
        print(f'float;{res:.2f}')
    else:
        print('???')
except:
    print('ERR')
code

2. Що буде виведено за виконання?
code
a = 5
b = 0
c = 7

def f(b):
    global c
    a = 5
    b = 3
    c = 1
    return a + b + c

a = 2
b = 6
c = 1

try:
    print(f(b), a, b, c, sep=';')
except:
    print('ERR', a, b, c, sep=';')
code

3. Що буде виведено за виконання?
code
try:
    m = 7
    while m>1:
        m += -3
    print(m, end=';')
except:
    print('ERR')
code

4. Що буде виведено за виконання?
code
try:
    for e in range(-1, -12, -2):
        if e<=6:
            continue
        print(e, end=';')
    else:
        print(e,end=';')
    print(e)
except:
    print('ERR')
code

5. Що буде виведено за виконання?\par (Дійсні значення, якщо наявні, в цьому завданні будуть виводитись у форматі з фіксованою крапкою та, можливо, з доволі великою кількістю десяткових розрядів. У відповіді для дійсних значень у ЦЬОМУ завданні достатньо вказати один десятковий розряд після крапки, відкинувши решту розрядів без округлення. Наприклад, замість 13.66666666666667 достатньо вказати 13.6, замість 25.23 достатньо вказати 25.2)
code
try:
    print(2, end=';')
    print(2j<=7j, end=';')
    print(5, end=';')
except ZeroDivisionError:
    print(0, end=';')
except ZeroDivisionError:
    print(3, end=';')
except:
    print('ERR', end=';')
else:
    print(8, end=';')
finally:
    print(6)
code

6. Що буде виведено за виконання?
code
def f(n):
    if n>9:
        f(n//100)
    print(n%10, end='')
 
f(123456)
code

7. Що буде виведено за виконання?
code
class A:
    b = 1

    def f1(self): print(2, end=';')
    def _f2(self): print(3, end=';')

    def main(self):
        try: print(self.b, end=';')
        except: print('ERR', end=';')

        try: self.f1()
        except: print('ERR', end=';')

        try: self._f2()
        except: print('ERR', end=';')


class B(A):
    b = 4

    def f1(self): print(5, end=';')
    def _f2(self): print(6, end=';')

try: B().main()
except: print('ERR', end=';')
code

8. Що буде виведено за виконання?
code
def f(arg, n):
    arg[32] = 9
    n -= 3

try:
    n = 6
    t = {58:4, 83:1, 83:6}
    f(t, n)
    print(n, end=';')
    for x, y in t.items():
        print(x, y, sep=';', end=';')
except:
    print('ERR')
code

9. Що буде виведено за виконання?
code
class A:
    b = 1
    a = 2

    def __init__(self):
        b = 3

class B(A):
    a = 4

    def __init__(self):
        super().__init__()
        self.c = 7

obj = B()
obj.b = obj.a + 10
B.b = 13

try: print(obj.b, end= ';')
except: print('ERR', end=';')

try: print(A.b, end= ';')
except: print('ERR', end=';')

try: print(B.b, end= ';')
except: print('ERR', end=';')
code

10. 
Для графа на рисунку з вершини 3 запущено алгоритм пошуку в глибину, що будує кістякове дерево. Списки суміжності вершин переглядаються алгоритмом за зростанням міток вершин.\par
\begin{center}
	\adjustimage{max size={0.9\linewidth}{0.9\paperheight}}
	{../10/Traversals/274.png}
\end{center}
\par Які з наведених ребер входять до складу отриманого кістякового дерева?\par
1) (2, 6)\par
2) (1, 4)\par
3) (0, 2)\par
4) (1, 5)\par
5) (0, 3)\par
6) (3, 6)\par
{\vskip 15pt} У формі вибрати номери відповідних ребер.\par
%answer:1 2 3 4 5
\end {large}}
%end
{\smallskip \begin {small} Затверджено на засіданні кафедри теоретичної кібернетики, протокол \No 12 від   19.05.2020 р.\par\smallskip\smallskip
{\parbox {6.5 cm}{Зав. кафедри \dotfill Ю.В.~Крак}}\hspace {0.7cm}{\hbox to 6.5 cm{ Екзаменатор \dotfill Т.О.~Карнаух}} \end{small}}
%begin
{{\begin {small}\par
{ \center  Київський національний університет імені Тараса Шевченка \par}
{\parbox {5 cm}{Освітня програма \hfill}{\parbox {7 cm}{Прикладна математика, Системний аналіз \hfill}}\parbox {6 cm}{\hfill Семестр 2}}\par
{\parbox {5 cm} {Навчальний предмет \hfill}\parbox {7 cm}{Програмування \hfill}\parbox {6cm}{\hfill}\par}\vspace{-7pt} \end{small}}{\center Екзаменаційний білет \No \textbf 
{ 29 }\par}}
{
\begin {large}
1. Що буде виведено за виконання?
code
try:
    res = 2<False and not 4.0>=3 or 5.0==9
    if isinstance(res, bool):
        print(f'bool;{res}')
    elif isinstance(res, int):
        print(f'int;{res}')
    elif isinstance(res, float):
        print(f'float;{res:.2f}')
    else:
        print('???')
except:
    print('ERR')
code

2. Що буде виведено за виконання?
code
a = 7
b = 6
c = 4

def g(a):
    a = 1
    b = 5
    c = 2
    return a + b + c

a = 8
b = 0
c = 2

try:
    print(g(b), a, b, c, sep=';')
except:
    print('ERR', a, b, c, sep=';')
code

3. Що буде виведено за виконання?
code
try:
    i = 7
    while i<=7:
        i += 2
    print(i, end=';')
except:
    print('ERR')
code

4. Що буде виведено за виконання?
code
try:
    for a in range(-4, 2, -3):
        if a>7:
            continue
        print(a, end=';')
        if a<3:
            break
    else:
        print(a,end=';')
    print(a)
except:
    print('ERR')
code

5. Що буде виведено за виконання?\par (Дійсні значення, якщо наявні, в цьому завданні будуть виводитись у форматі з фіксованою крапкою та, можливо, з доволі великою кількістю десяткових розрядів. У відповіді для дійсних значень у ЦЬОМУ завданні достатньо вказати один десятковий розряд після крапки, відкинувши решту розрядів без округлення. Наприклад, замість 13.66666666666667 достатньо вказати 13.6, замість 25.23 достатньо вказати 25.2)
code
try:
    print(4, end=';')
    print(int('c9'), end=';')
    print(1, end=';')
except ValueError:
    print(7, end=';')
except KeyboardInterrupt:
    print(2, end=';')
except:
    print('ERR', end=';')
else:
    print(6, end=';')
finally:
    print(8)
code

6. Що буде виведено за виконання?
code
def f(s):
    s1 = s
    for i in range(len(s)):
        s1[i] += 1
    return s1

s = [1, 2, 3]
f(s)
print(s)
code

7. Що буде виведено за виконання?
code
class A:
    c = 1

    def _h1(self): print(2, end=';')
    def h2(self): print(3, end=';')

    def main(self):
        try: print(self.c, end=';')
        except: print('ERR', end=';')

        try: self._h1()
        except: print('ERR', end=';')

        try: self.h2()
        except: print('ERR', end=';')


class B(A):
    c = 4

    def _h1(self): print(5, end=';')
    def h2(self): print(6, end=';')

try: B().main()
except: print('ERR', end=';')
code

8. Що буде виведено за виконання?
code
def f(arg, n):
    n = 9
    tmp = arg
    tmp[-9:] = ()
    return len(tmp)>3

try:
    f = [10,11,12,13,14]
    n = 6
    f(f, n)
    print(n, end=';')
    for el in f:
        print(el, end=';')
except:
    print('ERR')
code

9. Що буде виведено за виконання?
code
class A:
    c = 1
    b = 2

    def __init__(self):
        b = 3

class B(A):
    c = 4

    def __init__(self):
        super().__init__()
        self.a = 7

obj = B()
obj.c = obj.a + 10
B.c = 13

try: print(obj.c, end= ';')
except: print('ERR', end=';')

try: print(A.c, end= ';')
except: print('ERR', end=';')

try: print(B.c, end= ';')
except: print('ERR', end=';')
code

10. 
Для графа на рисунку з вершини 2 запущено алгоритм пошуку в глибину, що будує кістякове дерево. Списки суміжності вершин переглядаються алгоритмом за зростанням міток вершин.\par
\begin{center}
	\adjustimage{max size={0.9\linewidth}{0.9\paperheight}}
	{../10/Traversals/67.png}
\end{center}
\par Які з наведених ребер входять до складу отриманого кістякового дерева?\par
1) (2, 5)\par
2) (3, 4)\par
3) (0, 1)\par
4) (5, 6)\par
5) (1, 2)\par
6) (0, 2)\par
{\vskip 15pt} У формі вибрати номери відповідних ребер.\par
%answer:2 3 4 6
\end {large}}
%end
{\smallskip \begin {small} Затверджено на засіданні кафедри теоретичної кібернетики, протокол \No 12 від   19.05.2020 р.\par\smallskip\smallskip
{\parbox {6.5 cm}{Зав. кафедри \dotfill Ю.В.~Крак}}\hspace {0.7cm}{\hbox to 6.5 cm{ Екзаменатор \dotfill Т.О.~Карнаух}} \end{small}}
%begin
{{\begin {small}\par
{ \center  Київський національний університет імені Тараса Шевченка \par}
{\parbox {5 cm}{Освітня програма \hfill}{\parbox {7 cm}{Прикладна математика, Системний аналіз \hfill}}\parbox {6 cm}{\hfill Семестр 2}}\par
{\parbox {5 cm} {Навчальний предмет \hfill}\parbox {7 cm}{Програмування \hfill}\parbox {6cm}{\hfill}\par}\vspace{-7pt} \end{small}}{\center Екзаменаційний білет \No \textbf 
{ 30 }\par}}
{
\begin {large}
1. Що буде виведено за виконання?
code
def f(n):
    print('f', end='')
    return n<4

try:
    print(f(1) or f(5))
except:
    print('ERR')
code

2. Що буде виведено за виконання?
code
a = 1
b = 9
c = 9

def f(b):
    global c
    a = 5
    b *= 1
    c = 4
    return a + b + c

a = 1
b = 5
c = 0

try:
    print(f(b), a, b, c, sep=';')
except:
    print('ERR', a, b, c, sep=';')
code

3. Що буде виведено за виконання?
code
try:
    j = 5
    while j<5:
        j += 1
    print(j, end=';')
except:
    print('ERR')
code

4. Що буде виведено за виконання?
code
try:
    for c in range(12, -13, 3):
        if c<=4:
            break
        print(c, end=';');c = -1
    else:
        print(c,end=';')
    print(c)
except:
    print('ERR')
code

5. Що буде виведено за виконання?\par (Дійсні значення, якщо наявні, в цьому завданні будуть виводитись у форматі з фіксованою крапкою та, можливо, з доволі великою кількістю десяткових розрядів. У відповіді для дійсних значень у ЦЬОМУ завданні достатньо вказати один десятковий розряд після крапки, відкинувши решту розрядів без округлення. Наприклад, замість 13.66666666666667 достатньо вказати 13.6, замість 25.23 достатньо вказати 25.2)
code
try:
    print(9, end=';')
    print(7//3, end=';')
    print(1, end=';')
except Exception:
    print(0, end=';')
except TypeError:
    print(3, end=';')
except:
    print('ERR', end=';')
else:
    print(5, end=';')
finally:
    print(4)
code

6. Що буде виведено за виконання?
code
def f(s):
    for el in s:
        el += 1
    return s
 
s = [1, 2, 3]
s = f(s)
print(s)
code

7. Що буде виведено за виконання?
code
class A:
    a = 1

    def _g1(self): print(2, end=';')
    def _g2(self): print(3, end=';')

    def main(self):
        try: print(self.a, end=';')
        except: print('ERR', end=';')

        try: self._g1()
        except: print('ERR', end=';')

        try: self._g2()
        except: print('ERR', end=';')


class B(A):
    a = 4

    def _g1(self): print(5, end=';')
    def _g2(self): print(6, end=';')

try: B().main()
except: print('ERR', end=';')
code

8. Що буде виведено за виконання?
code
def f(arg, n):
    arg[69] = 0
    n += -3

try:
    n = 5
    d = {69:8, 84:3, 26:6, 26:4}
    f(d, n)
    print(n, end=';')
    for x in d :
        print(x, sep=';', end=';')
except:
    print('ERR')
code

9. Що буде виведено за виконання?
code
class A:
    a = 1
    b = 2

    def __init__(self):
        a = 3

class B(A):
    b = 4

    def __init__(self):
        super().__init__()
        self.c = 7

obj = B()
obj.b = obj.c + 10
B.b = 13

try: print(obj.b, end= ';')
except: print('ERR', end=';')

try: print(A.b, end= ';')
except: print('ERR', end=';')

try: print(B.b, end= ';')
except: print('ERR', end=';')
code

10. 
Для графа на рисунку з вершини 2 запущено алгоритм пошуку в глибину, що будує кістякове дерево. Списки суміжності вершин переглядаються алгоритмом за зростанням міток вершин.\par
\begin{center}
	\adjustimage{max size={0.9\linewidth}{0.9\paperheight}}
	{../10/Traversals/70.png}
\end{center}
\par Які з наведених ребер входять до складу отриманого кістякового дерева?\par
1) (1, 6)\par
2) (2, 6)\par
3) (0, 2)\par
4) (0, 1)\par
5) (0, 3)\par
6) (0, 4)\par
{\vskip 15pt} У формі вибрати номери відповідних ребер.\par
%answer:3 4
\end {large}}
%end
{\smallskip \begin {small} Затверджено на засіданні кафедри теоретичної кібернетики, протокол \No 12 від   19.05.2020 р.\par\smallskip\smallskip
{\parbox {6.5 cm}{Зав. кафедри \dotfill Ю.В.~Крак}}\hspace {0.7cm}{\hbox to 6.5 cm{ Екзаменатор \dotfill Т.О.~Карнаух}} \end{small}}
%begin
{{\begin {small}\par
{ \center  Київський національний університет імені Тараса Шевченка \par}
{\parbox {5 cm}{Освітня програма \hfill}{\parbox {7 cm}{Прикладна математика, Системний аналіз \hfill}}\parbox {6 cm}{\hfill Семестр 2}}\par
{\parbox {5 cm} {Навчальний предмет \hfill}\parbox {7 cm}{Програмування \hfill}\parbox {6cm}{\hfill}\par}\vspace{-7pt} \end{small}}{\center Екзаменаційний білет \No \textbf 
{ 31 }\par}}
{
\begin {large}
1. Що буде виведено за виконання?
code
try:
    res = True==3.0 and not 4<=6 or 2>6.0
    if isinstance(res, bool):
        print(f'bool;{res}')
    elif isinstance(res, int):
        print(f'int;{res}')
    elif isinstance(res, float):
        print(f'float;{res:.2f}')
    else:
        print('???')
except:
    print('ERR')
code

2. Що буде виведено за виконання?
code
a = 2
b = 7
c = 6

def g(a):
    global c
    a -= 3
    b = 5
    c = 2
    return a + b + c

a = 3
b = 8
c = 4

try:
    print(g(a), a, b, c, sep=';')
except:
    print('ERR', a, b, c, sep=';')
code

3. Що буде виведено за виконання?
code
try:
    k = 13
    while k>7:
        k += -2
    print(k, end=';')
except:
    print('ERR')
code

4. Що буде виведено за виконання?
code
try:
    for e in range(-3, 5, -1):
        if e<6:
            continue
        print(e, end=';')
    else:
        print('end',end=';')
    print(e)
except:
    print('ERR')
code

5. Що буде виведено за виконання?\par (Дійсні значення, якщо наявні, в цьому завданні будуть виводитись у форматі з фіксованою крапкою та, можливо, з доволі великою кількістю десяткових розрядів. У відповіді для дійсних значень у ЦЬОМУ завданні достатньо вказати один десятковий розряд після крапки, відкинувши решту розрядів без округлення. Наприклад, замість 13.66666666666667 достатньо вказати 13.6, замість 25.23 достатньо вказати 25.2)
code
try:
    print(6, end=';')
    print(int('b3'), end=';')
    print(2, end=';')
except TypeError:
    print(5, end=';')
except ZeroDivisionError:
    print(7, end=';')
except:
    print('ERR', end=';')
else:
    print(4, end=';')
finally:
    print(8)
code

6. Що буде виведено за виконання?
code
def f(n):
    if n>9:
        return f(n//10)+1
    else:
        return 1

print(f(123456))
code

7. Що буде виведено за виконання?
code
class A:
    d = 1

    def f1(self): print(2, end=';')
    def _f2(self): print(3, end=';')

    def main(self):
        try: print(self.d, end=';')
        except: print('ERR', end=';')

        try: self.f1()
        except: print('ERR', end=';')

        try: self._f2()
        except: print('ERR', end=';')


class B(A):
    d = 4

    def f1(self): print(5, end=';')
    def _f2(self): print(6, end=';')

try: B().main()
except: print('ERR', end=';')
code

8. Що буде виведено за виконання?
code
def f(arg, n):
    n = 0
    tmp = arg
    del tmp[4:]
    return len(tmp)>3

try:
    h = ['0','1','2','3','4']
    n = 3
    f(h, n)
    print(n, end=';')
    for el in h:
        print(el, end=';')
except:
    print('ERR')
code

9. Що буде виведено за виконання?
code
class A:
    a = 1
    c = 2

    def __init__(self):
        c = 3

class B(A):
    a = 4

    def __init__(self):
        super().__init__()
        self.b = 7

obj = B()
obj.c = obj.a + 10
B.c = 13

try: print(obj.c, end= ';')
except: print('ERR', end=';')

try: print(A.c, end= ';')
except: print('ERR', end=';')

try: print(B.c, end= ';')
except: print('ERR', end=';')
code

10. 
Для графа на рисунку з вершини 6 запущено алгоритм пошуку в глибину, що будує кістякове дерево. Списки суміжності вершин переглядаються алгоритмом за зростанням міток вершин.\par
\begin{center}
	\adjustimage{max size={0.9\linewidth}{0.9\paperheight}}
	{../10/Traversals/278.png}
\end{center}
\par Які з наведених ребер входять до складу отриманого кістякового дерева?\par
1) (4, 5)\par
2) (0, 5)\par
3) (1, 6)\par
4) (1, 5)\par
5) (1, 2)\par
6) (2, 6)\par
{\vskip 15pt} У формі вибрати номери відповідних ребер.\par
%answer:1 2 3 5
\end {large}}
%end
{\smallskip \begin {small} Затверджено на засіданні кафедри теоретичної кібернетики, протокол \No 12 від   19.05.2020 р.\par\smallskip\smallskip
{\parbox {6.5 cm}{Зав. кафедри \dotfill Ю.В.~Крак}}\hspace {0.7cm}{\hbox to 6.5 cm{ Екзаменатор \dotfill Т.О.~Карнаух}} \end{small}}
%begin
{{\begin {small}\par
{ \center  Київський національний університет імені Тараса Шевченка \par}
{\parbox {5 cm}{Освітня програма \hfill}{\parbox {7 cm}{Прикладна математика, Системний аналіз \hfill}}\parbox {6 cm}{\hfill Семестр 2}}\par
{\parbox {5 cm} {Навчальний предмет \hfill}\parbox {7 cm}{Програмування \hfill}\parbox {6cm}{\hfill}\par}\vspace{-7pt} \end{small}}{\center Екзаменаційний білет \No \textbf 
{ 32 }\par}}
{
\begin {large}
1. Що буде виведено за виконання?
code
try:
    res = 1**-0.5*True
    if isinstance(res, bool):
        print(f'bool;{res}')
    elif isinstance(res, int):
        print(f'int;{res}')
    elif isinstance(res, float):
        print(f'float;{res:.2f}')
    else:
        print('???')
except:
    print('ERR')
code

2. Що буде виведено за виконання?
code
a = 3
b = 2
c = 1

def g(b):
    a = 1
    b *= 4
    c = 5
    return a + b + c

a = 9
b = 0
c = 5

try:
    print(g(b), a, b, c, sep=';')
except:
    print('ERR', a, b, c, sep=';')
code

3. Що буде виведено за виконання?
code
try:
    i = 9
    while i>=2:
        i += -3
    print(i, end=';')
except:
    print('ERR')
code

4. Що буде виведено за виконання?
code
try:
    for d in range(-3, -12, 2):
        if d>=5:
            break
        print(d, end=';');d = -8
    else:
        print(d,end=';')
    print(d)
except:
    print('ERR')
code

5. Що буде виведено за виконання?\par (Дійсні значення, якщо наявні, в цьому завданні будуть виводитись у форматі з фіксованою крапкою та, можливо, з доволі великою кількістю десяткових розрядів. У відповіді для дійсних значень у ЦЬОМУ завданні достатньо вказати один десятковий розряд після крапки, відкинувши решту розрядів без округлення. Наприклад, замість 13.66666666666667 достатньо вказати 13.6, замість 25.23 достатньо вказати 25.2)
code
try:
    print(1, end=';')
    print(9%2, end=';')
    print(0, end=';')
except ValueError:
    print(1, end=';')
except KeyboardInterrupt:
    print(5, end=';')
except:
    print('ERR', end=';')
else:
    print(9, end=';')
finally:
    print(3)
code

6. Що буде виведено за виконання?
code
def f(s):
    s1 = s
    for i in range(len(s)):
        s1[i] += 1
        return s1

s = [1, 2, 3]
f(s)
print(s)
code

7. Що буде виведено за виконання?
code
class A:
    b = 1

    def __h1(self): print(2, end=';')
    def _h2(self): print(3, end=';')

    def main(self):
        try: print(self.b, end=';')
        except: print('ERR', end=';')

        try: self.__h1()
        except: print('ERR', end=';')

        try: self._h2()
        except: print('ERR', end=';')


class B(A):
    b = 4

    def __h1(self): print(5, end=';')
    def _h2(self): print(6, end=';')

try: B().main()
except: print('ERR', end=';')
code

8. Що буде виведено за виконання?
code
def f(arg, n):
    arg[48] = 4
    n = 7
    arg = set(arg.items())

try:
    n = 3
    t = {79:9, 71:5, 79:7}
    f(t, n)
    print(n, end=';')
    for x in t.keys():
        print(x, sep=';', end=';')
except:
    print('ERR')
code

9. Що буде виведено за виконання?
code
class A:
    b = 1
    c = 2

    def __init__(self):
        b = 3

class B(A):
    c = 4

    def __init__(self):
        super().__init__()
        self.a = 7

obj = B()
obj.b = obj.b + 10
B.b = 13

try: print(obj.b, end= ';')
except: print('ERR', end=';')

try: print(A.b, end= ';')
except: print('ERR', end=';')

try: print(B.b, end= ';')
except: print('ERR', end=';')
code

10. 
Для графа на рисунку з вершини 2 запущено алгоритм пошуку в глибину, що будує кістякове дерево. Списки суміжності вершин переглядаються алгоритмом за зростанням міток вершин.\par
\begin{center}
	\adjustimage{max size={0.9\linewidth}{0.9\paperheight}}
	{../10/Traversals/267.png}
\end{center}
\par Які з наведених ребер входять до складу отриманого кістякового дерева?\par
1) (3, 4)\par
2) (0, 3)\par
3) (0, 1)\par
4) (0, 5)\par
5) (0, 4)\par
6) (1, 3)\par
{\vskip 15pt} У формі вибрати номери відповідних ребер.\par
%answer:1 3 4 6
\end {large}}
%end
{\smallskip \begin {small} Затверджено на засіданні кафедри теоретичної кібернетики, протокол \No 12 від   19.05.2020 р.\par\smallskip\smallskip
{\parbox {6.5 cm}{Зав. кафедри \dotfill Ю.В.~Крак}}\hspace {0.7cm}{\hbox to 6.5 cm{ Екзаменатор \dotfill Т.О.~Карнаух}} \end{small}}
%begin
{{\begin {small}\par
{ \center  Київський національний університет імені Тараса Шевченка \par}
{\parbox {5 cm}{Освітня програма \hfill}{\parbox {7 cm}{Прикладна математика, Системний аналіз \hfill}}\parbox {6 cm}{\hfill Семестр 2}}\par
{\parbox {5 cm} {Навчальний предмет \hfill}\parbox {7 cm}{Програмування \hfill}\parbox {6cm}{\hfill}\par}\vspace{-7pt} \end{small}}{\center Екзаменаційний білет \No \textbf 
{ 33 }\par}}
{
\begin {large}
1. Що буде виведено за виконання?
code
try:
    res = 3<True
    if isinstance(res, bool):
        print(f'bool;{res}')
    elif isinstance(res, int):
        print(f'int;{res}')
    elif isinstance(res, float):
        print(f'float;{res:.2f}')
    else:
        print('???')
except:
    print('ERR')
code

2. Що буде виведено за виконання?
code
a = 1
b = 9
c = 3

def f(a):
    global c
    a = 3
    b = 4
    c = 2
    return a + b + c

a = 5
b = 0
c = 2

try:
    print(f(a), a, b, c, sep=';')
except:
    print('ERR', a, b, c, sep=';')
code

3. Що буде виведено за виконання?
code
try:
    n = 11
    while n>8:
        n += -2
    print(n, end=';')
except:
    print('ERR')
code

4. Що буде виведено за виконання?
code
try:
    for f in range(-15, -6, -2):
        if f<=4:
            continue
        print(f, end=';')
    else:
        print(f,end=';')
    print(f)
except:
    print('ERR')
code

5. Що буде виведено за виконання?\par (Дійсні значення, якщо наявні, в цьому завданні будуть виводитись у форматі з фіксованою крапкою та, можливо, з доволі великою кількістю десяткових розрядів. У відповіді для дійсних значень у ЦЬОМУ завданні достатньо вказати один десятковий розряд після крапки, відкинувши решту розрядів без округлення. Наприклад, замість 13.66666666666667 достатньо вказати 13.6, замість 25.23 достатньо вказати 25.2)
code
try:
    print(7, end=';')
    print(8j>6j, end=';')
    print(6, end=';')
except Exception:
    print(4, end=';')
except ZeroDivisionError:
    print(0, end=';')
except:
    print('ERR', end=';')
else:
    print(2, end=';')
finally:
    print(8)
code

6. Що буде виведено за виконання?
code
def f(s1):
    s = s1[:]
    for i in range(len(s)):
        s[i] += 1
 
s = [1, 2, 3]
f(s)
print(s)
code

7. Що буде виведено за виконання?
code
class A:
    d = 1

    def _g1(self): print(2, end=';')
    def g2(self): print(3, end=';')

    def main(self):
        try: print(self.d, end=';')
        except: print('ERR', end=';')

        try: self._g1()
        except: print('ERR', end=';')

        try: self.g2()
        except: print('ERR', end=';')


class B(A):
    d = 4

    def _g1(self): print(5, end=';')
    def g2(self): print(6, end=';')

try: B().main()
except: print('ERR', end=';')
code

8. Що буде виведено за виконання?
code
def f(arg, n):
    n = 9
    tmp = arg
    tmp[2:2] = [4,5]
    return len(tmp)>3

try:
    b = [10,11,12,13,14]
    n = 6
    f(b, n)
    print(n, end=';')
    for el in b:
        print(el, end=';')
except:
    print('ERR')
code

9. Що буде виведено за виконання?
code
class A:
    c = 1
    b = 2

    def __init__(self):
        b = 3

class B(A):
    c = 4

    def __init__(self):
        super().__init__()
        self.a = 7

obj = B()
obj.a = obj.c + 10
B.a = 13

try: print(obj.a, end= ';')
except: print('ERR', end=';')

try: print(A.a, end= ';')
except: print('ERR', end=';')

try: print(B.a, end= ';')
except: print('ERR', end=';')
code

10. 
Для графа на рисунку з вершини 3 запущено алгоритм пошуку в глибину, що будує кістякове дерево. Списки суміжності вершин переглядаються алгоритмом за зростанням міток вершин.\par
\begin{center}
	\adjustimage{max size={0.9\linewidth}{0.9\paperheight}}
	{../10/Traversals/17.png}
\end{center}
\par Які з наведених ребер входять до складу отриманого кістякового дерева?\par
1) (0, 4)\par
2) (1, 6)\par
3) (0, 1)\par
4) (2, 4)\par
5) (0, 6)\par
6) (1, 4)\par
{\vskip 15pt} У формі вибрати номери відповідних ребер.\par
%answer:1 3 4 5
\end {large}}
%end
{\smallskip \begin {small} Затверджено на засіданні кафедри теоретичної кібернетики, протокол \No 12 від   19.05.2020 р.\par\smallskip\smallskip
{\parbox {6.5 cm}{Зав. кафедри \dotfill Ю.В.~Крак}}\hspace {0.7cm}{\hbox to 6.5 cm{ Екзаменатор \dotfill Т.О.~Карнаух}} \end{small}}
%begin
{{\begin {small}\par
{ \center  Київський національний університет імені Тараса Шевченка \par}
{\parbox {5 cm}{Освітня програма \hfill}{\parbox {7 cm}{Прикладна математика, Системний аналіз \hfill}}\parbox {6 cm}{\hfill Семестр 2}}\par
{\parbox {5 cm} {Навчальний предмет \hfill}\parbox {7 cm}{Програмування \hfill}\parbox {6cm}{\hfill}\par}\vspace{-7pt} \end{small}}{\center Екзаменаційний білет \No \textbf 
{ 34 }\par}}
{
\begin {large}
1. Що буде виведено за виконання?
code
try:
    res = not 5<7.0 or 9.0>=3 and False!=9
    if isinstance(res, bool):
        print(f'bool;{res}')
    elif isinstance(res, int):
        print(f'int;{res}')
    elif isinstance(res, float):
        print(f'float;{res:.2f}')
    else:
        print('???')
except:
    print('ERR')
code

2. Що буде виведено за виконання?
code
a = 8
b = 7
c = 6

def h(a):
    global c
    a = 3
    b += 2
    c = 1
    return a + b + c

a = 4
b = 4
c = 2

try:
    print(h(b), a, b, c, sep=';')
except:
    print('ERR', a, b, c, sep=';')
code

3. Що буде виведено за виконання?
code
try:
    t = 6
    while t>=6:
        t += -3
    print(t, end=';')
except:
    print('ERR')
code

4. Що буде виведено за виконання?
code
try:
    for a in range(-2, -1, 3):
        if a<8:
            continue
        print(a, end=';')
    else:
        print(a,end=';')
    print(a)
except:
    print('ERR')
code

5. Що буде виведено за виконання?\par (Дійсні значення, якщо наявні, в цьому завданні будуть виводитись у форматі з фіксованою крапкою та, можливо, з доволі великою кількістю десяткових розрядів. У відповіді для дійсних значень у ЦЬОМУ завданні достатньо вказати один десятковий розряд після крапки, відкинувши решту розрядів без округлення. Наприклад, замість 13.66666666666667 достатньо вказати 13.6, замість 25.23 достатньо вказати 25.2)
code
try:
    print(3, end=';')
    print(2//0, end=';')
    print(6, end=';')
except KeyboardInterrupt:
    print(1, end=';')
except TypeError:
    print(7, end=';')
except:
    print('ERR', end=';')
else:
    print(0, end=';')
finally:
    print(5)
code

6. Що буде виведено за виконання?
code
def f(s):
    for i in range(len(s)):
        s[i] += 1
        return
 
s = [1, 2, 3]
f(s)
print(s)
code

7. Що буде виведено за виконання?
code
class A:
    c = 1

    def _g1(self): print(2, end=';')
    def _g2(self): print(3, end=';')

    def main(self):
        try: print(self.c, end=';')
        except: print('ERR', end=';')

        try: self._g1()
        except: print('ERR', end=';')

        try: self._g2()
        except: print('ERR', end=';')


class B(A):
    c = 4

    def _g1(self): print(5, end=';')
    def _g2(self): print(6, end=';')

try: B().main()
except: print('ERR', end=';')
code

8. Що буде виведено за виконання?
code
def f(arg, n):
    arg[28] = 8
    n -= 2

try:
    n = 5
    t = {28:2, 51:1, 51:9}
    f(t, n)
    print(n, end=';')
    for x in t.values():
        print(x, sep=';', end=';')
except:
    print('ERR')
code

9. Що буде виведено за виконання?
code
class A:
    a = 1
    b = 2

    def __init__(self):
        a = 3

class B(A):
    b = 4

    def __init__(self):
        super().__init__()
        self.c = 7

obj = B()
obj.b = obj.a + 10
B.b = 13

try: print(obj.b, end= ';')
except: print('ERR', end=';')

try: print(A.b, end= ';')
except: print('ERR', end=';')

try: print(B.b, end= ';')
except: print('ERR', end=';')
code

10. 
Для графа на рисунку з вершини 5 запущено алгоритм пошуку в глибину, що будує кістякове дерево. Списки суміжності вершин переглядаються алгоритмом за зростанням міток вершин.\par
\begin{center}
	\adjustimage{max size={0.9\linewidth}{0.9\paperheight}}
	{../10/Traversals/208.png}
\end{center}
\par Які з наведених ребер входять до складу отриманого кістякового дерева?\par
1) (3, 5)\par
2) (3, 6)\par
3) (0, 3)\par
4) (0, 4)\par
5) (1, 5)\par
6) (3, 4)\par
{\vskip 15pt} У формі вибрати номери відповідних ребер.\par
%answer:2 3 4 5
\end {large}}
%end
{\smallskip \begin {small} Затверджено на засіданні кафедри теоретичної кібернетики, протокол \No 12 від   19.05.2020 р.\par\smallskip\smallskip
{\parbox {6.5 cm}{Зав. кафедри \dotfill Ю.В.~Крак}}\hspace {0.7cm}{\hbox to 6.5 cm{ Екзаменатор \dotfill Т.О.~Карнаух}} \end{small}}
%begin
{{\begin {small}\par
{ \center  Київський національний університет імені Тараса Шевченка \par}
{\parbox {5 cm}{Освітня програма \hfill}{\parbox {7 cm}{Прикладна математика, Системний аналіз \hfill}}\parbox {6 cm}{\hfill Семестр 2}}\par
{\parbox {5 cm} {Навчальний предмет \hfill}\parbox {7 cm}{Програмування \hfill}\parbox {6cm}{\hfill}\par}\vspace{-7pt} \end{small}}{\center Екзаменаційний білет \No \textbf 
{ 35 }\par}}
{
\begin {large}
1. Що буде виведено за виконання?
code
try:
    res = 3**-2+-2
    if isinstance(res, bool):
        print(f'bool;{res}')
    elif isinstance(res, int):
        print(f'int;{res}')
    elif isinstance(res, float):
        print(f'float;{res:.2f}')
    else:
        print('???')
except:
    print('ERR')
code

2. Що буде виведено за виконання?
code
a = 9
b = 7
c = 0

def f(b):
    a = 3
    b -= 4
    c = 2
    return a + b + c

a = 1
b = 3
c = 6

try:
    print(f(a), a, b, c, sep=';')
except:
    print('ERR', a, b, c, sep=';')
code

3. Що буде виведено за виконання?
code
try:
    m = 2
    while m<=11:
        m += 1
    print(m, end=';')
except:
    print('ERR')
code

4. Що буде виведено за виконання?
code
try:
    for b in range(7, 15, -3):
        if b>7:
            continue
        print(b, end=';');b = 3
    else:
        print(b,end=';')
    print(b)
except:
    print('ERR')
code

5. Що буде виведено за виконання?\par (Дійсні значення, якщо наявні, в цьому завданні будуть виводитись у форматі з фіксованою крапкою та, можливо, з доволі великою кількістю десяткових розрядів. У відповіді для дійсних значень у ЦЬОМУ завданні достатньо вказати один десятковий розряд після крапки, відкинувши решту розрядів без округлення. Наприклад, замість 13.66666666666667 достатньо вказати 13.6, замість 25.23 достатньо вказати 25.2)
code
try:
    print(9, end=';')
    print(int('4'), end=';')
    print(8, end=';')
except Exception:
    print(9, end=';')
except ValueError:
    print(7, end=';')
except:
    print('ERR', end=';')
else:
    print(1, end=';')
finally:
    print(4)
code

6. Що буде виведено за виконання?
code
def f(s):
    for i in range(len(s)):
        s[i] += 1
    return
 
s = [1, 2, 3]
f(s)
print(s)
code

7. Що буде виведено за виконання?
code
class A:
    a = 1

    def _f1(self): print(2, end=';')
    def __f2(self): print(3, end=';')

    def main(self):
        try: print(self.a, end=';')
        except: print('ERR', end=';')

        try: self._f1()
        except: print('ERR', end=';')

        try: self.__f2()
        except: print('ERR', end=';')


class B(A):
    a = 4

    def _f1(self): print(5, end=';')
    def __f2(self): print(6, end=';')

try: B().main()
except: print('ERR', end=';')
code

8. Що буде виведено за виконання?
code
def f(arg, n):
    n = 7
    tmp = arg
    tmp[-6:-1] = 'ac'
    return len(tmp)>3

try:
    c = ['0','1','2','3','4']
    n = 1
    f(c, n)
    print(n, end=';')
    for el in c:
        print(el, end=';')
except:
    print('ERR')
code

9. Що буде виведено за виконання?
code
class A:
    b = 1
    a = 2

    def __init__(self):
        a = 3

class B(A):
    b = 4

    def __init__(self):
        super().__init__()
        self.c = 7

obj = B()
obj.c = obj.a + 10
B.c = 13

try: print(obj.c, end= ';')
except: print('ERR', end=';')

try: print(A.c, end= ';')
except: print('ERR', end=';')

try: print(B.c, end= ';')
except: print('ERR', end=';')
code

10. 
Для графа на рисунку з вершини 1 запущено алгоритм пошуку в глибину, що будує кістякове дерево. Списки суміжності вершин переглядаються алгоритмом за зростанням міток вершин.\par
\begin{center}
	\adjustimage{max size={0.9\linewidth}{0.9\paperheight}}
	{../10/Traversals/185.png}
\end{center}
\par Які з наведених ребер входять до складу отриманого кістякового дерева?\par
1) (0, 4)\par
2) (4, 6)\par
3) (1, 2)\par
4) (1, 5)\par
5) (0, 3)\par
6) (1, 6)\par
{\vskip 15pt} У формі вибрати номери відповідних ребер.\par
%answer:2 3 5
\end {large}}
%end
{\smallskip \begin {small} Затверджено на засіданні кафедри теоретичної кібернетики, протокол \No 12 від   19.05.2020 р.\par\smallskip\smallskip
{\parbox {6.5 cm}{Зав. кафедри \dotfill Ю.В.~Крак}}\hspace {0.7cm}{\hbox to 6.5 cm{ Екзаменатор \dotfill Т.О.~Карнаух}} \end{small}}
%begin
{{\begin {small}\par
{ \center  Київський національний університет імені Тараса Шевченка \par}
{\parbox {5 cm}{Освітня програма \hfill}{\parbox {7 cm}{Прикладна математика, Системний аналіз \hfill}}\parbox {6 cm}{\hfill Семестр 2}}\par
{\parbox {5 cm} {Навчальний предмет \hfill}\parbox {7 cm}{Програмування \hfill}\parbox {6cm}{\hfill}\par}\vspace{-7pt} \end{small}}{\center Екзаменаційний білет \No \textbf 
{ 36 }\par}}
{
\begin {large}
1. Що буде виведено за виконання?
code
try:
    res = -3**-2+False
    if isinstance(res, bool):
        print(f'bool;{res}')
    elif isinstance(res, int):
        print(f'int;{res}')
    elif isinstance(res, float):
        print(f'float;{res:.2f}')
    else:
        print('???')
except:
    print('ERR')
code

2. Що буде виведено за виконання?
code
a = 6
b = 4
c = 5

def h(b):
    a = 3
    b = 5
    c = 4
    return a + b + c

a = 3
b = 7
c = 1

try:
    print(h(b), a, b, c, sep=';')
except:
    print('ERR', a, b, c, sep=';')
code

3. Що буде виведено за виконання?
code
try:
    m = 8
    while m>=9:
        m += -2
    print(m, end=';')
except:
    print('ERR')
code

4. Що буде виведено за виконання?
code
try:
    for c in range(-4, 6, -1):
        if c>=3:
            break
        print(c, end=';')
        if c>=4:
            break
    else:
        print(c,end=';')
    print(c)
except:
    print('ERR')
code

5. Що буде виведено за виконання?\par (Дійсні значення, якщо наявні, в цьому завданні будуть виводитись у форматі з фіксованою крапкою та, можливо, з доволі великою кількістю десяткових розрядів. У відповіді для дійсних значень у ЦЬОМУ завданні достатньо вказати один десятковий розряд після крапки, відкинувши решту розрядів без округлення. Наприклад, замість 13.66666666666667 достатньо вказати 13.6, замість 25.23 достатньо вказати 25.2)
code
try:
    print(2, end=';')
    print(3j!=0j, end=';')
    print(0, end=';')
except ValueError:
    print(5, end=';')
except TypeError:
    print(6, end=';')
except:
    print('ERR', end=';')
else:
    print(8, end=';')
finally:
    print(0)
code

6. Що буде виведено за виконання?
code
def f(n):
    if n>9:
        f(n//10)
    else:
        print(n%10, end='')

f(543216)
code

7. Що буде виведено за виконання?
code
class A:
    b = 1

    def _h1(self): print(2, end=';')
    def _h2(self): print(3, end=';')

    def main(self):
        try: print(self.b, end=';')
        except: print('ERR', end=';')

        try: self._h1()
        except: print('ERR', end=';')

        try: self._h2()
        except: print('ERR', end=';')


class B(A):
    b = 4

    def _h1(self): print(5, end=';')
    def _h2(self): print(6, end=';')

try: B().main()
except: print('ERR', end=';')
code

8. Що буде виведено за виконання?
code
def f(arg, n):
    arg[63] = 7
    n *= -2

try:
    n = 7
    s = {21:7, 63:7, 63:7}
    f(s, n)
    print(n, end=';')
    for x in s.values():
        print(x, sep=';', end=';')
except:
    print('ERR')
code

9. Що буде виведено за виконання?
code
class A:
    c = 1
    a = 2

    def __init__(self):
        c = 3

class B(A):
    a = 4

    def __init__(self):
        super().__init__()
        self.b = 7

obj = B()
obj.b = obj.c + 10
B.b = 13

try: print(obj.b, end= ';')
except: print('ERR', end=';')

try: print(A.b, end= ';')
except: print('ERR', end=';')

try: print(B.b, end= ';')
except: print('ERR', end=';')
code

10. 
Для графа на рисунку з вершини 0 запущено алгоритм пошуку в глибину, що будує кістякове дерево. Списки суміжності вершин переглядаються алгоритмом за зростанням міток вершин.\par
\begin{center}
	\adjustimage{max size={0.9\linewidth}{0.9\paperheight}}
	{../10/Traversals/148.png}
\end{center}
\par Які з наведених ребер входять до складу отриманого кістякового дерева?\par
1) (0, 3)\par
2) (0, 5)\par
3) (5, 6)\par
4) (4, 6)\par
5) (0, 4)\par
6) (2, 6)\par
{\vskip 15pt} У формі вибрати номери відповідних ребер.\par
%answer:4 6
\end {large}}
%end
{\smallskip \begin {small} Затверджено на засіданні кафедри теоретичної кібернетики, протокол \No 12 від   19.05.2020 р.\par\smallskip\smallskip
{\parbox {6.5 cm}{Зав. кафедри \dotfill Ю.В.~Крак}}\hspace {0.7cm}{\hbox to 6.5 cm{ Екзаменатор \dotfill Т.О.~Карнаух}} \end{small}}
%begin
{{\begin {small}\par
{ \center  Київський національний університет імені Тараса Шевченка \par}
{\parbox {5 cm}{Освітня програма \hfill}{\parbox {7 cm}{Прикладна математика, Системний аналіз \hfill}}\parbox {6 cm}{\hfill Семестр 2}}\par
{\parbox {5 cm} {Навчальний предмет \hfill}\parbox {7 cm}{Програмування \hfill}\parbox {6cm}{\hfill}\par}\vspace{-7pt} \end{small}}{\center Екзаменаційний білет \No \textbf 
{ 37 }\par}}
{
\begin {large}
1. Що буде виведено за виконання?
code
def f(n):
    print('f', end='')
    return n>=-1

try:
    print(f(8) and f(-5))
except:
    print('ERR')
code

2. Що буде виведено за виконання?
code
a = 9
b = 8
c = 6

def g(a):
    a = 1
    b = 5
    c = 4
    return a + b + c

a = 9
b = 7
c = 0

try:
    print(g(a), a, b, c, sep=';')
except:
    print('ERR', a, b, c, sep=';')
code

3. Що буде виведено за виконання?
code
try:
    n = 9
    while n<8:
        n += 2
    print(n, end=';')
except:
    print('ERR')
code

4. Що буде виведено за виконання?
code
try:
    for e in range(-8, 2, 3):
        if e>=5:
            continue
        print(e, end=';')
    else:
        print(e,end=';')
    print(e)
except:
    print('ERR')
code

5. Що буде виведено за виконання?\par (Дійсні значення, якщо наявні, в цьому завданні будуть виводитись у форматі з фіксованою крапкою та, можливо, з доволі великою кількістю десяткових розрядів. У відповіді для дійсних значень у ЦЬОМУ завданні достатньо вказати один десятковий розряд після крапки, відкинувши решту розрядів без округлення. Наприклад, замість 13.66666666666667 достатньо вказати 13.6, замість 25.23 достатньо вказати 25.2)
code
try:
    print(1, end=';')
    print(7/3, end=';')
    print(3, end=';')
except KeyboardInterrupt:
    print(9, end=';')
except ZeroDivisionError:
    print(4, end=';')
except:
    print('ERR', end=';')
else:
    print(2, end=';')
finally:
    print(5)
code

6. Що буде виведено за виконання?
code
def f(n):
    if n>9:
        f(n//10)
    print(n%10, end='')
 
f(123456)
code

7. Що буде виведено за виконання?
code
class A:
    a = 1

    def _g1(self): print(2, end=';')
    def __g2(self): print(3, end=';')

    def main(self):
        try: print(self.a, end=';')
        except: print('ERR', end=';')

        try: self._g1()
        except: print('ERR', end=';')

        try: self.__g2()
        except: print('ERR', end=';')


class B(A):
    a = 4

    def _g1(self): print(5, end=';')
    def __g2(self): print(6, end=';')

try: B().main()
except: print('ERR', end=';')
code

8. Що буде виведено за виконання?
code
def f(arg, n):
    arg[22] = 5
    n += -2

try:
    n = 6
    s = {58:4, 90:5, 58:8}
    f(s, n)
    print(n, end=';')
    for x in s.keys():
        print(x, sep=';', end=';')
except:
    print('ERR')
code

9. Що буде виведено за виконання?
code
class A:
    b = 1
    c = 2

    def __init__(self):
        b = 3

class B(A):
    c = 4

    def __init__(self):
        super().__init__()
        self.a = 7

obj = B()
obj.a = obj.c + 10
B.a = 13

try: print(obj.a, end= ';')
except: print('ERR', end=';')

try: print(A.a, end= ';')
except: print('ERR', end=';')

try: print(B.a, end= ';')
except: print('ERR', end=';')
code

10. 
Для графа на рисунку з вершини 5 запущено алгоритм пошуку в глибину, що будує кістякове дерево. Списки суміжності вершин переглядаються алгоритмом за зростанням міток вершин.\par
\begin{center}
	\adjustimage{max size={0.9\linewidth}{0.9\paperheight}}
	{../10/Traversals/174.png}
\end{center}
\par Які з наведених ребер входять до складу отриманого кістякового дерева?\par
1) (5, 6)\par
2) (1, 3)\par
3) (0, 4)\par
4) (3, 5)\par
5) (2, 5)\par
6) (1, 6)\par
{\vskip 15pt} У формі вибрати номери відповідних ребер.\par
%answer:2 5 6
\end {large}}
%end
{\smallskip \begin {small} Затверджено на засіданні кафедри теоретичної кібернетики, протокол \No 12 від   19.05.2020 р.\par\smallskip\smallskip
{\parbox {6.5 cm}{Зав. кафедри \dotfill Ю.В.~Крак}}\hspace {0.7cm}{\hbox to 6.5 cm{ Екзаменатор \dotfill Т.О.~Карнаух}} \end{small}}
%begin
{{\begin {small}\par
{ \center  Київський національний університет імені Тараса Шевченка \par}
{\parbox {5 cm}{Освітня програма \hfill}{\parbox {7 cm}{Прикладна математика, Системний аналіз \hfill}}\parbox {6 cm}{\hfill Семестр 2}}\par
{\parbox {5 cm} {Навчальний предмет \hfill}\parbox {7 cm}{Програмування \hfill}\parbox {6cm}{\hfill}\par}\vspace{-7pt} \end{small}}{\center Екзаменаційний білет \No \textbf 
{ 38 }\par}}
{
\begin {large}
1. Що буде виведено за виконання?
code
try:
    res = 9//7%6*4
    if isinstance(res, bool):
        print(f'bool;{res}')
    elif isinstance(res, int):
        print(f'int;{res}')
    elif isinstance(res, float):
        print(f'float;{res:.2f}')
    else:
        print('???')
except:
    print('ERR')
code

2. Що буде виведено за виконання?
code
a = 8
b = 5
c = 7

def g(a):
    global c
    a -= 5
    b = 5
    c = 4
    return a + b + c

a = 9
b = 8
c = 6

try:
    print(g(b), a, b, c, sep=';')
except:
    print('ERR', a, b, c, sep=';')
code

3. Що буде виведено за виконання?
code
try:
    j = 4
    while j<12:
        j += 1
    print(j, end=';')
except:
    print('ERR')
code

4. Що буде виведено за виконання?
code
try:
    for c in range(10, 13, -3):
        if c>8:
            break
        print(c, end=';')
    else:
        print(c,end=';')
    print(c)
except:
    print('ERR')
code

5. Що буде виведено за виконання?\par (Дійсні значення, якщо наявні, в цьому завданні будуть виводитись у форматі з фіксованою крапкою та, можливо, з доволі великою кількістю десяткових розрядів. У відповіді для дійсних значень у ЦЬОМУ завданні достатньо вказати один десятковий розряд після крапки, відкинувши решту розрядів без округлення. Наприклад, замість 13.66666666666667 достатньо вказати 13.6, замість 25.23 достатньо вказати 25.2)
code
try:
    print(6, end=';')
    print(int('a8'), end=';')
    print(4, end=';')
except Exception:
    print(9, end=';')
except Exception:
    print(5, end=';')
except:
    print('ERR', end=';')
else:
    print(1, end=';')
finally:
    print(3)
code

6. Що буде виведено за виконання?
code
def f(n):
    if n<=9:
        print(n%10,end='')
    else:
        f(n//10)

f(543216)
code

7. Що буде виведено за виконання?
code
class A:
    d = 1

    def __h1(self): print(2, end=';')
    def _h2(self): print(3, end=';')

    def main(self):
        try: print(self.d, end=';')
        except: print('ERR', end=';')

        try: self.__h1()
        except: print('ERR', end=';')

        try: self._h2()
        except: print('ERR', end=';')


class B(A):
    d = 4

    def __h1(self): print(5, end=';')
    def _h2(self): print(6, end=';')

try: B().main()
except: print('ERR', end=';')
code

8. Що буде виведено за виконання?
code
def f(arg, n):
    arg[67] = 3
    n = 9
    arg = tuple(arg.items())

try:
    n = 1
    d = {44:2, 54:0, 54:5}
    f(d, n)
    print(n, end=';')
    for x in d.items():
        print(x, sep=';', end=';')
except:
    print('ERR')
code

9. Що буде виведено за виконання?
code
class A:
    c = 1
    a = 2

    def __init__(self):
        a = 3

class B(A):
    c = 4

    def __init__(self):
        super().__init__()
        self.b = 7

obj = B()
obj.b = obj.c + 10
B.b = 13

try: print(obj.b, end= ';')
except: print('ERR', end=';')

try: print(A.b, end= ';')
except: print('ERR', end=';')

try: print(B.b, end= ';')
except: print('ERR', end=';')
code

10. 
Для графа на рисунку з вершини 0 запущено алгоритм пошуку в глибину, що будує кістякове дерево. Списки суміжності вершин переглядаються алгоритмом за зростанням міток вершин.\par
\begin{center}
	\adjustimage{max size={0.9\linewidth}{0.9\paperheight}}
	{../10/Traversals/0.png}
\end{center}
\par Які з наведених ребер входять до складу отриманого кістякового дерева?\par
1) (0, 5)\par
2) (1, 5)\par
3) (1, 4)\par
4) (4, 5)\par
5) (0, 6)\par
6) (1, 3)\par
{\vskip 15pt} У формі вибрати номери відповідних ребер.\par
%answer:4 6
\end {large}}
%end
{\smallskip \begin {small} Затверджено на засіданні кафедри теоретичної кібернетики, протокол \No 12 від   19.05.2020 р.\par\smallskip\smallskip
{\parbox {6.5 cm}{Зав. кафедри \dotfill Ю.В.~Крак}}\hspace {0.7cm}{\hbox to 6.5 cm{ Екзаменатор \dotfill Т.О.~Карнаух}} \end{small}}
%begin
{{\begin {small}\par
{ \center  Київський національний університет імені Тараса Шевченка \par}
{\parbox {5 cm}{Освітня програма \hfill}{\parbox {7 cm}{Прикладна математика, Системний аналіз \hfill}}\parbox {6 cm}{\hfill Семестр 2}}\par
{\parbox {5 cm} {Навчальний предмет \hfill}\parbox {7 cm}{Програмування \hfill}\parbox {6cm}{\hfill}\par}\vspace{-7pt} \end{small}}{\center Екзаменаційний білет \No \textbf 
{ 39 }\par}}
{
\begin {large}
1. Що буде виведено за виконання?
code
def f(n):
    print('f', end='')
    return n

try:
    print(f(0) or f(-9))
except:
    print('ERR')
code

2. Що буде виведено за виконання?
code
a = 4
b = 2
c = 5

def h(b):
    global c
    a = 2
    b = 1
    c = 3
    return a + b + c

a = 8
b = 1
c = 3

try:
    print(h(b), a, b, c, sep=';')
except:
    print('ERR', a, b, c, sep=';')
code

3. Що буде виведено за виконання?
code
try:
    k = 3
    while k<=11:
        k += 2
    print(k, end=';')
except:
    print('ERR')
code

4. Що буде виведено за виконання?
code
try:
    for a in range(-13, 10, -2):
        if a<=4:
            continue
        print(a, end=';');a = -3
    else:
        print(a,end=';')
    print(a)
except:
    print('ERR')
code

5. Що буде виведено за виконання?\par (Дійсні значення, якщо наявні, в цьому завданні будуть виводитись у форматі з фіксованою крапкою та, можливо, з доволі великою кількістю десяткових розрядів. У відповіді для дійсних значень у ЦЬОМУ завданні достатньо вказати один десятковий розряд після крапки, відкинувши решту розрядів без округлення. Наприклад, замість 13.66666666666667 достатньо вказати 13.6, замість 25.23 достатньо вказати 25.2)
code
try:
    print(7, end=';')
    print(2/0.0, end=';')
    print(6, end=';')
except ZeroDivisionError:
    print(0, end=';')
except ValueError:
    print(8, end=';')
except:
    print('ERR', end=';')
else:
    print(1, end=';')
finally:
    print(2)
code

6. Що буде виведено за виконання?
code
def f(n):
    if n<=9:
        print(n%10, end='')
    else:
        f(n//10)

f(123456)
code

7. Що буде виведено за виконання?
code
class A:
    b = 1

    def _f1(self): print(2, end=';')
    def f2(self): print(3, end=';')

    def main(self):
        try: print(self.b, end=';')
        except: print('ERR', end=';')

        try: self._f1()
        except: print('ERR', end=';')

        try: self.f2()
        except: print('ERR', end=';')


class B(A):
    b = 4

    def _f1(self): print(5, end=';')
    def f2(self): print(6, end=';')

try: B().main()
except: print('ERR', end=';')
code

8. Що буде виведено за виконання?
code
def f(arg, n):
    arg[23] = 5
    n *= 2

try:
    n = 3
    t = {51:4, 50:0, 60:4, 48:6}
    f(t, n)
    print(n, end=';')
    for x in t.items():
        print(x, sep=';', end=';')
except:
    print('ERR')
code

9. Що буде виведено за виконання?
code
class A:
    c = 1
    b = 2

    def __init__(self):
        b = 3

class B(A):
    b = 4

    def __init__(self):
        super().__init__()
        self.a = 7

obj = B()
obj.a = obj.b + 10
B.a = 13

try: print(obj.a, end= ';')
except: print('ERR', end=';')

try: print(A.a, end= ';')
except: print('ERR', end=';')

try: print(B.a, end= ';')
except: print('ERR', end=';')
code

10. 
Для графа на рисунку з вершини 4 запущено алгоритм пошуку в глибину, що будує кістякове дерево. Списки суміжності вершин переглядаються алгоритмом за зростанням міток вершин.\par
\begin{center}
	\adjustimage{max size={0.9\linewidth}{0.9\paperheight}}
	{../10/Traversals/201.png}
\end{center}
\par Які з наведених ребер входять до складу отриманого кістякового дерева?\par
1) (1, 4)\par
2) (1, 5)\par
3) (3, 5)\par
4) (0, 6)\par
5) (0, 2)\par
6) (3, 4)\par
{\vskip 15pt} У формі вибрати номери відповідних ребер.\par
%answer:1 2 3 4 5
\end {large}}
%end
{\smallskip \begin {small} Затверджено на засіданні кафедри теоретичної кібернетики, протокол \No 12 від   19.05.2020 р.\par\smallskip\smallskip
{\parbox {6.5 cm}{Зав. кафедри \dotfill Ю.В.~Крак}}\hspace {0.7cm}{\hbox to 6.5 cm{ Екзаменатор \dotfill Т.О.~Карнаух}} \end{small}}
%begin
{{\begin {small}\par
{ \center  Київський національний університет імені Тараса Шевченка \par}
{\parbox {5 cm}{Освітня програма \hfill}{\parbox {7 cm}{Прикладна математика, Системний аналіз \hfill}}\parbox {6 cm}{\hfill Семестр 2}}\par
{\parbox {5 cm} {Навчальний предмет \hfill}\parbox {7 cm}{Програмування \hfill}\parbox {6cm}{\hfill}\par}\vspace{-7pt} \end{small}}{\center Екзаменаційний білет \No \textbf 
{ 40 }\par}}
{
\begin {large}
1. Що буде виведено за виконання?
code
try:
    res = 2.0!=4 or not False<6 or 8.0>8
    if isinstance(res, bool):
        print(f'bool;{res}')
    elif isinstance(res, int):
        print(f'int;{res}')
    elif isinstance(res, float):
        print(f'float;{res:.2f}')
    else:
        print('???')
except:
    print('ERR')
code

2. Що буде виведено за виконання?
code
a = 2
b = 1
c = 0

def h(b):
    a *= 1
    b = 3
    c = 2
    return a + b + c

a = 3
b = 4
c = 5

try:
    print(h(a), a, b, c, sep=';')
except:
    print('ERR', a, b, c, sep=';')
code

3. Що буде виведено за виконання?
code
try:
    t = 9
    while t<=7:
        t += 1
    print(t, end=';')
except:
    print('ERR')
code

4. Що буде виведено за виконання?
code
try:
    for f in range(4, 9, 2):
        if f<7:
            continue
        print(f, end=';')
    else:
        print(f,end=';')
    print(f)
except:
    print('ERR')
code

5. Що буде виведено за виконання?\par (Дійсні значення, якщо наявні, в цьому завданні будуть виводитись у форматі з фіксованою крапкою та, можливо, з доволі великою кількістю десяткових розрядів. У відповіді для дійсних значень у ЦЬОМУ завданні достатньо вказати один десятковий розряд після крапки, відкинувши решту розрядів без округлення. Наприклад, замість 13.66666666666667 достатньо вказати 13.6, замість 25.23 достатньо вказати 25.2)
code
try:
    print(7, end=';')
    print(int('6'), end=';')
    print(4, end=';')
except TypeError:
    print(5, end=';')
except KeyboardInterrupt:
    print(3, end=';')
except:
    print('ERR', end=';')
else:
    print(9, end=';')
finally:
    print(0)
code

6. Що буде виведено за виконання?
code
def f(n):
    if n>9:
        f(n//100)
    print(n%10,end='')
 
f(987654)
code

7. Що буде виведено за виконання?
code
class A:
    c = 1

    def h1(self): print(2, end=';')
    def _h2(self): print(3, end=';')

    def main(self):
        try: print(self.c, end=';')
        except: print('ERR', end=';')

        try: self.h1()
        except: print('ERR', end=';')

        try: self._h2()
        except: print('ERR', end=';')


class B(A):
    c = 4

    def h1(self): print(5, end=';')
    def _h2(self): print(6, end=';')

try: B().main()
except: print('ERR', end=';')
code

8. Що буде виведено за виконання?
code
def f(arg, n):
    arg[84] = 4
    n -= -3

try:
    n = 4
    d = {61:1, 17:1, 61:9}
    f(d, n)
    print(n, end=';')
    for x in d :
        print(x, sep=';', end=';')
except:
    print('ERR')
code

9. Що буде виведено за виконання?
code
class A:
    b = 1
    a = 2

    def __init__(self):
        b = 3

class B(A):
    b = 4

    def __init__(self):
        super().__init__()
        self.c = 7

obj = B()
obj.a = obj.b + 10
B.a = 13

try: print(obj.a, end= ';')
except: print('ERR', end=';')

try: print(A.a, end= ';')
except: print('ERR', end=';')

try: print(B.a, end= ';')
except: print('ERR', end=';')
code

10. 
Для графа на рисунку з вершини 4 запущено алгоритм пошуку в глибину, що будує кістякове дерево. Списки суміжності вершин переглядаються алгоритмом за зростанням міток вершин.\par
\begin{center}
	\adjustimage{max size={0.9\linewidth}{0.9\paperheight}}
	{../10/Traversals/75.png}
\end{center}
\par Які з наведених ребер входять до складу отриманого кістякового дерева?\par
1) (2, 5)\par
2) (0, 3)\par
3) (0, 1)\par
4) (5, 6)\par
5) (1, 4)\par
6) (0, 4)\par
{\vskip 15pt} У формі вибрати номери відповідних ребер.\par
%answer:1 2 3 4 6
\end {large}}
%end
{\smallskip \begin {small} Затверджено на засіданні кафедри теоретичної кібернетики, протокол \No 12 від   19.05.2020 р.\par\smallskip\smallskip
{\parbox {6.5 cm}{Зав. кафедри \dotfill Ю.В.~Крак}}\hspace {0.7cm}{\hbox to 6.5 cm{ Екзаменатор \dotfill Т.О.~Карнаух}} \end{small}}
%begin
{{\begin {small}\par
{ \center  Київський національний університет імені Тараса Шевченка \par}
{\parbox {5 cm}{Освітня програма \hfill}{\parbox {7 cm}{Прикладна математика, Системний аналіз \hfill}}\parbox {6 cm}{\hfill Семестр 2}}\par
{\parbox {5 cm} {Навчальний предмет \hfill}\parbox {7 cm}{Програмування \hfill}\parbox {6cm}{\hfill}\par}\vspace{-7pt} \end{small}}{\center Екзаменаційний білет \No \textbf 
{ 41 }\par}}
{
\begin {large}
1. Що буде виведено за виконання?
code
def f(n):
    print('f', end='')
    return n

try:
    print(f(-3) or f(6))
except:
    print('ERR')
code

2. Що буде виведено за виконання?
code
a = 8
b = 4
c = 6

def f(a):
    global c
    a = 1
    b = 2
    c = 4
    return a + b + c

a = 5
b = 9
c = 2

try:
    print(f(b), a, b, c, sep=';')
except:
    print('ERR', a, b, c, sep=';')
code

3. Що буде виведено за виконання?
code
try:
    n = 15
    while n>4:
        n += -3
    print(n, end=';')
except:
    print('ERR')
code

4. Що буде виведено за виконання?
code
try:
    for b in range(-10, -11, -2):
        if b>3:
            continue
        print(b, end=';');b = 5
    else:
        print(b,end=';')
    print(b)
except:
    print('ERR')
code

5. Що буде виведено за виконання?\par (Дійсні значення, якщо наявні, в цьому завданні будуть виводитись у форматі з фіксованою крапкою та, можливо, з доволі великою кількістю десяткових розрядів. У відповіді для дійсних значень у ЦЬОМУ завданні достатньо вказати один десятковий розряд після крапки, відкинувши решту розрядів без округлення. Наприклад, замість 13.66666666666667 достатньо вказати 13.6, замість 25.23 достатньо вказати 25.2)
code
try:
    print(5, end=';')
    print(2//1, end=';')
    print(9, end=';')
except KeyboardInterrupt:
    print(8, end=';')
except ValueError:
    print(4, end=';')
except:
    print('ERR', end=';')
else:
    print(0, end=';')
finally:
    print(7)
code

6. Що буде виведено за виконання?
code
def f(n):
    print(n%2, end='')
    if n>2:
        f(n//2)
 
f(16)
code

7. Що буде виведено за виконання?
code
class A:
    b = 1

    def _f1(self): print(2, end=';')
    def _f2(self): print(3, end=';')

    def main(self):
        try: print(self.b, end=';')
        except: print('ERR', end=';')

        try: self._f1()
        except: print('ERR', end=';')

        try: self._f2()
        except: print('ERR', end=';')


class B(A):
    b = 4

    def _f1(self): print(5, end=';')
    def _f2(self): print(6, end=';')

try: B().main()
except: print('ERR', end=';')
code

8. Що буде виведено за виконання?
code
def f(arg, n):
    arg[89] = 5
    n = 6
    arg = list(arg.keys())

try:
    n = 2
    t = {89:4, 32:0, 80:5, 32:4}
    f(t, n)
    print(n, end=';')
    for x in t.keys():
        print(x, sep=';', end=';')
except:
    print('ERR')
code

9. Що буде виведено за виконання?
code
class A:
    a = 1
    b = 2

    def __init__(self):
        b = 3

class B(A):
    a = 4

    def __init__(self):
        super().__init__()
        self.c = 7

obj = B()
obj.c = obj.b + 10
B.c = 13

try: print(obj.c, end= ';')
except: print('ERR', end=';')

try: print(A.c, end= ';')
except: print('ERR', end=';')

try: print(B.c, end= ';')
except: print('ERR', end=';')
code

10. 
Для графа на рисунку з вершини 5 запущено алгоритм пошуку в глибину, що будує кістякове дерево. Списки суміжності вершин переглядаються алгоритмом за зростанням міток вершин.\par
\begin{center}
	\adjustimage{max size={0.9\linewidth}{0.9\paperheight}}
	{../10/Traversals/29.png}
\end{center}
\par Які з наведених ребер входять до складу отриманого кістякового дерева?\par
1) (1, 4)\par
2) (1, 6)\par
3) (4, 6)\par
4) (1, 2)\par
5) (2, 3)\par
6) (5, 6)\par
{\vskip 15pt} У формі вибрати номери відповідних ребер.\par
%answer:3 4 5
\end {large}}
%end
{\smallskip \begin {small} Затверджено на засіданні кафедри теоретичної кібернетики, протокол \No 12 від   19.05.2020 р.\par\smallskip\smallskip
{\parbox {6.5 cm}{Зав. кафедри \dotfill Ю.В.~Крак}}\hspace {0.7cm}{\hbox to 6.5 cm{ Екзаменатор \dotfill Т.О.~Карнаух}} \end{small}}
%begin
{{\begin {small}\par
{ \center  Київський національний університет імені Тараса Шевченка \par}
{\parbox {5 cm}{Освітня програма \hfill}{\parbox {7 cm}{Прикладна математика, Системний аналіз \hfill}}\parbox {6 cm}{\hfill Семестр 2}}\par
{\parbox {5 cm} {Навчальний предмет \hfill}\parbox {7 cm}{Програмування \hfill}\parbox {6cm}{\hfill}\par}\vspace{-7pt} \end{small}}{\center Екзаменаційний білет \No \textbf 
{ 42 }\par}}
{
\begin {large}
1. Що буде виведено за виконання?
code
try:
    res = 7==4.0 and not 2<=5.0 and 6>=True
    if isinstance(res, bool):
        print(f'bool;{res}')
    elif isinstance(res, int):
        print(f'int;{res}')
    elif isinstance(res, float):
        print(f'float;{res:.2f}')
    else:
        print('???')
except:
    print('ERR')
code

2. Що буде виведено за виконання?
code
a = 2
b = 3
c = 5

def f(a):
    global c
    a += 5
    b = 1
    c = 3
    return a + b + c

a = 7
b = 4
c = 6

try:
    print(f(b), a, b, c, sep=';')
except:
    print('ERR', a, b, c, sep=';')
code

3. Що буде виведено за виконання?
code
try:
    j = 5
    while j>9:
        j += -2
    print(j, end=';')
except:
    print('ERR')
code

4. Що буде виведено за виконання?
code
try:
    for d in range(5, -7, -1):
        if d<=6:
            break
        print(d, end=';')
    else:
        print('end',end=';')
    print(d)
except:
    print('ERR')
code

5. Що буде виведено за виконання?\par (Дійсні значення, якщо наявні, в цьому завданні будуть виводитись у форматі з фіксованою крапкою та, можливо, з доволі великою кількістю десяткових розрядів. У відповіді для дійсних значень у ЦЬОМУ завданні достатньо вказати один десятковий розряд після крапки, відкинувши решту розрядів без округлення. Наприклад, замість 13.66666666666667 достатньо вказати 13.6, замість 25.23 достатньо вказати 25.2)
code
try:
    print(1, end=';')
    print(3j<1j, end=';')
    print(6, end=';')
except TypeError:
    print(7, end=';')
except ZeroDivisionError:
    print(9, end=';')
except:
    print('ERR', end=';')
else:
    print(5, end=';')
finally:
    print(3)
code

6. Що буде виведено за виконання?
code
def f(n):
    if n<=9:
        print(n%10, end='')
    else:
        f(n//10)

f(987651)
code

7. Що буде виведено за виконання?
code
class A:
    c = 1

    def __g1(self): print(2, end=';')
    def _g2(self): print(3, end=';')

    def main(self):
        try: print(self.c, end=';')
        except: print('ERR', end=';')

        try: self.__g1()
        except: print('ERR', end=';')

        try: self._g2()
        except: print('ERR', end=';')


class B(A):
    c = 4

    def __g1(self): print(5, end=';')
    def _g2(self): print(6, end=';')

try: B().main()
except: print('ERR', end=';')
code

8. Що буде виведено за виконання?
code
def f(arg, n):
    arg[23] = 8
    n = 8
    arg = set(arg.values())

try:
    n = 4
    s = {89:4, 23:3, 89:2}
    f(s, n)
    print(n, end=';')
    for x in s.items():
        print(x, sep=';', end=';')
except:
    print('ERR')
code

9. Що буде виведено за виконання?
code
class A:
    a = 1
    c = 2

    def __init__(self):
        a = 3

class B(A):
    c = 4

    def __init__(self):
        super().__init__()
        self.b = 7

obj = B()
obj.c = obj.a + 10
B.c = 13

try: print(obj.c, end= ';')
except: print('ERR', end=';')

try: print(A.c, end= ';')
except: print('ERR', end=';')

try: print(B.c, end= ';')
except: print('ERR', end=';')
code

10. 
Для графа на рисунку з вершини 5 запущено алгоритм пошуку в глибину, що будує кістякове дерево. Списки суміжності вершин переглядаються алгоритмом за зростанням міток вершин.\par
\begin{center}
	\adjustimage{max size={0.9\linewidth}{0.9\paperheight}}
	{../10/Traversals/250.png}
\end{center}
\par Які з наведених ребер входять до складу отриманого кістякового дерева?\par
1) (1, 2)\par
2) (4, 5)\par
3) (2, 5)\par
4) (0, 3)\par
5) (5, 6)\par
6) (0, 5)\par
{\vskip 15pt} У формі вибрати номери відповідних ребер.\par
%answer:1 4 6
\end {large}}
%end
{\smallskip \begin {small} Затверджено на засіданні кафедри теоретичної кібернетики, протокол \No 12 від   19.05.2020 р.\par\smallskip\smallskip
{\parbox {6.5 cm}{Зав. кафедри \dotfill Ю.В.~Крак}}\hspace {0.7cm}{\hbox to 6.5 cm{ Екзаменатор \dotfill Т.О.~Карнаух}} \end{small}}
%begin
{{\begin {small}\par
{ \center  Київський національний університет імені Тараса Шевченка \par}
{\parbox {5 cm}{Освітня програма \hfill}{\parbox {7 cm}{Прикладна математика, Системний аналіз \hfill}}\parbox {6 cm}{\hfill Семестр 2}}\par
{\parbox {5 cm} {Навчальний предмет \hfill}\parbox {7 cm}{Програмування \hfill}\parbox {6cm}{\hfill}\par}\vspace{-7pt} \end{small}}{\center Екзаменаційний білет \No \textbf 
{ 43 }\par}}
{
\begin {large}
1. Що буде виведено за виконання?
code
try:
    res = not 8!=9 or 9.0>=6.0 and False==3
    if isinstance(res, bool):
        print(f'bool;{res}')
    elif isinstance(res, int):
        print(f'int;{res}')
    elif isinstance(res, float):
        print(f'float;{res:.2f}')
    else:
        print('???')
except:
    print('ERR')
code

2. Що буде виведено за виконання?
code
a = 1
b = 8
c = 9

def h(b):
    a -= 2
    b = 4
    c = 2
    return a + b + c

a = 0
b = 5
c = 0

try:
    print(h(a), a, b, c, sep=';')
except:
    print('ERR', a, b, c, sep=';')
code

3. Що буде виведено за виконання?
code
try:
    i = 10
    while i>=0:
        i += -3
    print(i, end=';')
except:
    print('ERR')
code

4. Що буде виведено за виконання?
code
try:
    for a in range(0, 11, 2):
        if a<5:
            continue
        print(a, end=';')
    else:
        print(a,end=';')
    print(a)
except:
    print('ERR')
code

5. Що буде виведено за виконання?\par (Дійсні значення, якщо наявні, в цьому завданні будуть виводитись у форматі з фіксованою крапкою та, можливо, з доволі великою кількістю десяткових розрядів. У відповіді для дійсних значень у ЦЬОМУ завданні достатньо вказати один десятковий розряд після крапки, відкинувши решту розрядів без округлення. Наприклад, замість 13.66666666666667 достатньо вказати 13.6, замість 25.23 достатньо вказати 25.2)
code
try:
    print(4, end=';')
    print(6%0, end=';')
    print(2, end=';')
except Exception:
    print(8, end=';')
except KeyboardInterrupt:
    print(0, end=';')
except:
    print('ERR', end=';')
else:
    print(1, end=';')
finally:
    print(6)
code

6. Що буде виведено за виконання?
code
def f(n):
    print(n%10, end='')
    if n>9:
        f(n//10)
 
f(543216)
code

7. Що буде виведено за виконання?
code
class A:
    a = 1

    def _g1(self): print(2, end=';')
    def g2(self): print(3, end=';')

    def main(self):
        try: print(self.a, end=';')
        except: print('ERR', end=';')

        try: self._g1()
        except: print('ERR', end=';')

        try: self.g2()
        except: print('ERR', end=';')


class B(A):
    a = 4

    def _g1(self): print(5, end=';')
    def g2(self): print(6, end=';')

try: B().main()
except: print('ERR', end=';')
code

8. Що буде виведено за виконання?
code
def f(arg, n):
    arg[22] = 0
    n = 7
    arg = tuple(arg.values())

try:
    n = 4
    s = {45:4, 16:4, 22:8, 87:2}
    f(s, n)
    print(n, end=';')
    for x, y in s.items():
        print(x, y, sep=';', end=';')
except:
    print('ERR')
code

9. Що буде виведено за виконання?
code
class A:
    c = 1
    a = 2

    def __init__(self):
        c = 3

class B(A):
    a = 4

    def __init__(self):
        super().__init__()
        self.b = 7

obj = B()
obj.b = obj.c + 10
B.b = 13

try: print(obj.b, end= ';')
except: print('ERR', end=';')

try: print(A.b, end= ';')
except: print('ERR', end=';')

try: print(B.b, end= ';')
except: print('ERR', end=';')
code

10. 
Для графа на рисунку з вершини 2 запущено алгоритм пошуку в глибину, що будує кістякове дерево. Списки суміжності вершин переглядаються алгоритмом за зростанням міток вершин.\par
\begin{center}
	\adjustimage{max size={0.9\linewidth}{0.9\paperheight}}
	{../10/Traversals/128.png}
\end{center}
\par Які з наведених ребер входять до складу отриманого кістякового дерева?\par
1) (2, 5)\par
2) (0, 2)\par
3) (3, 6)\par
4) (3, 5)\par
5) (0, 4)\par
6) (1, 3)\par
{\vskip 15pt} У формі вибрати номери відповідних ребер.\par
%answer:2 3 4 5 6
\end {large}}
%end
{\smallskip \begin {small} Затверджено на засіданні кафедри теоретичної кібернетики, протокол \No 12 від   19.05.2020 р.\par\smallskip\smallskip
{\parbox {6.5 cm}{Зав. кафедри \dotfill Ю.В.~Крак}}\hspace {0.7cm}{\hbox to 6.5 cm{ Екзаменатор \dotfill Т.О.~Карнаух}} \end{small}}
%begin
{{\begin {small}\par
{ \center  Київський національний університет імені Тараса Шевченка \par}
{\parbox {5 cm}{Освітня програма \hfill}{\parbox {7 cm}{Прикладна математика, Системний аналіз \hfill}}\parbox {6 cm}{\hfill Семестр 2}}\par
{\parbox {5 cm} {Навчальний предмет \hfill}\parbox {7 cm}{Програмування \hfill}\parbox {6cm}{\hfill}\par}\vspace{-7pt} \end{small}}{\center Екзаменаційний білет \No \textbf 
{ 44 }\par}}
{
\begin {large}
1. Що буде виведено за виконання?
code
def f(n):
    print('f', end='')
    return n<=3

try:
    print(f(-8) and f(9))
except:
    print('ERR')
code

2. Що буде виведено за виконання?
code
a = 1
b = 3
c = 7

def g(b):
    a = 5
    b = 3
    c = 2
    return a + b + c

a = 0
b = 1
c = 6

try:
    print(g(a), a, b, c, sep=';')
except:
    print('ERR', a, b, c, sep=';')
code

3. Що буде виведено за виконання?
code
try:
    t = 15
    while t>6:
        t += -3
    print(t, end=';')
except:
    print('ERR')
code

4. Що буде виведено за виконання?
code
try:
    for b in range(-6, 0, 3):
        if b>=6:
            break
        print(b, end=';')
    else:
        print(b,end=';')
    print(b)
except:
    print('ERR')
code

5. Що буде виведено за виконання?\par (Дійсні значення, якщо наявні, в цьому завданні будуть виводитись у форматі з фіксованою крапкою та, можливо, з доволі великою кількістю десяткових розрядів. У відповіді для дійсних значень у ЦЬОМУ завданні достатньо вказати один десятковий розряд після крапки, відкинувши решту розрядів без округлення. Наприклад, замість 13.66666666666667 достатньо вказати 13.6, замість 25.23 достатньо вказати 25.2)
code
try:
    print(3, end=';')
    print(int('9'), end=';')
    print(0, end=';')
except TypeError:
    print(1, end=';')
except ValueError:
    print(7, end=';')
except:
    print('ERR', end=';')
else:
    print(8, end=';')
finally:
    print(5)
code

6. Що буде виведено за виконання?
code
def f(n):
    if n>2:
        f(n//2)
    else:
        print(n%2, end='')

f(16)
code

7. Що буде виведено за виконання?
code
class A:
    d = 1

    def _f1(self): print(2, end=';')
    def __f2(self): print(3, end=';')

    def main(self):
        try: print(self.d, end=';')
        except: print('ERR', end=';')

        try: self._f1()
        except: print('ERR', end=';')

        try: self.__f2()
        except: print('ERR', end=';')


class B(A):
    d = 4

    def _f1(self): print(5, end=';')
    def __f2(self): print(6, end=';')

try: B().main()
except: print('ERR', end=';')
code

8. Що буде виведено за виконання?
code
def f(arg, n):
    n = 4
    tmp = arg
    tmp[6:8] = [1,2]
    return len(tmp)>3

try:
    e = [10,11,12,13,14]
    n = 0
    f(e, n)
    print(n, end=';')
    for el in e:
        print(el, end=';')
except:
    print('ERR')
code

9. Що буде виведено за виконання?
code
class A:
    a = 1
    b = 2

    def __init__(self):
        b = 3

class B(A):
    a = 4

    def __init__(self):
        super().__init__()
        self.c = 7

obj = B()
obj.a = obj.c + 10
B.a = 13

try: print(obj.a, end= ';')
except: print('ERR', end=';')

try: print(A.a, end= ';')
except: print('ERR', end=';')

try: print(B.a, end= ';')
except: print('ERR', end=';')
code

10. 
Для графа на рисунку з вершини 2 запущено алгоритм пошуку в глибину, що будує кістякове дерево. Списки суміжності вершин переглядаються алгоритмом за зростанням міток вершин.\par
\begin{center}
	\adjustimage{max size={0.9\linewidth}{0.9\paperheight}}
	{../10/Traversals/160.png}
\end{center}
\par Які з наведених ребер входять до складу отриманого кістякового дерева?\par
1) (4, 6)\par
2) (0, 5)\par
3) (1, 2)\par
4) (0, 3)\par
5) (1, 5)\par
6) (2, 4)\par
{\vskip 15pt} У формі вибрати номери відповідних ребер.\par
%answer:2 3 4
\end {large}}
%end
{\smallskip \begin {small} Затверджено на засіданні кафедри теоретичної кібернетики, протокол \No 12 від   19.05.2020 р.\par\smallskip\smallskip
{\parbox {6.5 cm}{Зав. кафедри \dotfill Ю.В.~Крак}}\hspace {0.7cm}{\hbox to 6.5 cm{ Екзаменатор \dotfill Т.О.~Карнаух}} \end{small}}
%begin
{{\begin {small}\par
{ \center  Київський національний університет імені Тараса Шевченка \par}
{\parbox {5 cm}{Освітня програма \hfill}{\parbox {7 cm}{Прикладна математика, Системний аналіз \hfill}}\parbox {6 cm}{\hfill Семестр 2}}\par
{\parbox {5 cm} {Навчальний предмет \hfill}\parbox {7 cm}{Програмування \hfill}\parbox {6cm}{\hfill}\par}\vspace{-7pt} \end{small}}{\center Екзаменаційний білет \No \textbf 
{ 45 }\par}}
{
\begin {large}
1. Що буде виведено за виконання?
code
def f(n):
    print('f', end='')
    return n!=-2

try:
    print(f(3) and f(-6))
except:
    print('ERR')
code

2. Що буде виведено за виконання?
code
a = 5
b = 3
c = 8

def h(b):
    global c
    a = 5
    b = 3
    c = 4
    return a + b + c

a = 0
b = 7
c = 2

try:
    print(h(a), a, b, c, sep=';')
except:
    print('ERR', a, b, c, sep=';')
code

3. Що буде виведено за виконання?
code
try:
    k = 14
    while k>=5:
        k += -2
    print(k, end=';')
except:
    print('ERR')
code

4. Що буде виведено за виконання?
code
try:
    for c in range(-7, 5, -3):
        if c<=3:
            continue
        print(c, end=';');c = -5
    else:
        print(c,end=';')
    print(c)
except:
    print('ERR')
code

5. Що буде виведено за виконання?\par (Дійсні значення, якщо наявні, в цьому завданні будуть виводитись у форматі з фіксованою крапкою та, можливо, з доволі великою кількістю десяткових розрядів. У відповіді для дійсних значень у ЦЬОМУ завданні достатньо вказати один десятковий розряд після крапки, відкинувши решту розрядів без округлення. Наприклад, замість 13.66666666666667 достатньо вказати 13.6, замість 25.23 достатньо вказати 25.2)
code
try:
    print(2, end=';')
    print(int('c4'), end=';')
    print(0, end=';')
except Exception:
    print(1, end=';')
except ZeroDivisionError:
    print(5, end=';')
except:
    print('ERR', end=';')
else:
    print(8, end=';')
finally:
    print(9)
code

6. Що буде виведено за виконання?
code
def f(n):
    print(n%10, end='')
    if n>9:
        f(n//10)
 
f(123456)
code

7. Що буде виведено за виконання?
code
class A:
    c = 1

    def h1(self): print(2, end=';')
    def _h2(self): print(3, end=';')

    def main(self):
        try: print(self.c, end=';')
        except: print('ERR', end=';')

        try: self.h1()
        except: print('ERR', end=';')

        try: self._h2()
        except: print('ERR', end=';')


class B(A):
    c = 4

    def h1(self): print(5, end=';')
    def _h2(self): print(6, end=';')

try: B().main()
except: print('ERR', end=';')
code

8. Що буде виведено за виконання?
code
def f(arg, n):
    arg[50] = 1
    n = 6
    arg = list(arg.keys())

try:
    n = 0
    d = {53:2, 50:3, 38:2, 50:2}
    f(d, n)
    print(n, end=';')
    for x in d.keys():
        print(x, sep=';', end=';')
except:
    print('ERR')
code

9. Що буде виведено за виконання?
code
class A:
    c = 1
    b = 2

    def __init__(self):
        c = 3

class B(A):
    b = 4

    def __init__(self):
        super().__init__()
        self.a = 7

obj = B()
obj.b = obj.a + 10
B.b = 13

try: print(obj.b, end= ';')
except: print('ERR', end=';')

try: print(A.b, end= ';')
except: print('ERR', end=';')

try: print(B.b, end= ';')
except: print('ERR', end=';')
code

10. 
Для графа на рисунку з вершини 2 запущено алгоритм пошуку в глибину, що будує кістякове дерево. Списки суміжності вершин переглядаються алгоритмом за зростанням міток вершин.\par
\begin{center}
	\adjustimage{max size={0.9\linewidth}{0.9\paperheight}}
	{../10/Traversals/229.png}
\end{center}
\par Які з наведених ребер входять до складу отриманого кістякового дерева?\par
1) (0, 3)\par
2) (2, 3)\par
3) (1, 3)\par
4) (0, 2)\par
5) (1, 5)\par
6) (3, 4)\par
{\vskip 15pt} У формі вибрати номери відповідних ребер.\par
%answer:1 3 4 5 6
\end {large}}
%end
{\smallskip \begin {small} Затверджено на засіданні кафедри теоретичної кібернетики, протокол \No 12 від   19.05.2020 р.\par\smallskip\smallskip
{\parbox {6.5 cm}{Зав. кафедри \dotfill Ю.В.~Крак}}\hspace {0.7cm}{\hbox to 6.5 cm{ Екзаменатор \dotfill Т.О.~Карнаух}} \end{small}}
%begin
{{\begin {small}\par
{ \center  Київський національний університет імені Тараса Шевченка \par}
{\parbox {5 cm}{Освітня програма \hfill}{\parbox {7 cm}{Прикладна математика, Системний аналіз \hfill}}\parbox {6 cm}{\hfill Семестр 2}}\par
{\parbox {5 cm} {Навчальний предмет \hfill}\parbox {7 cm}{Програмування \hfill}\parbox {6cm}{\hfill}\par}\vspace{-7pt} \end{small}}{\center Екзаменаційний білет \No \textbf 
{ 46 }\par}}
{
\begin {large}
1. Що буде виведено за виконання?
code
try:
    res = "8j"!="3.0j"
    if isinstance(res, bool):
        print(f'bool;{res}')
    elif isinstance(res, int):
        print(f'int;{res}')
    elif isinstance(res, float):
        print(f'float;{res:.2f}')
    else:
        print('???')
except:
    print('ERR')
code

2. Що буде виведено за виконання?
code
a = 4
b = 9
c = 1

def f(b):
    global c
    a = 3
    b *= 4
    c = 1
    return a + b + c

a = 2
b = 7
c = 3

try:
    print(f(b), a, b, c, sep=';')
except:
    print('ERR', a, b, c, sep=';')
code

3. Що буде виведено за виконання?
code
try:
    t = 5
    while t>=4:
        t += -3
    print(t, end=';')
except:
    print('ERR')
code

4. Що буде виведено за виконання?
code
try:
    for d in range(1, -13, 3):
        if d>=4:
            continue
        print(d, end=';')
    else:
        print('end',end=';')
    print(d)
except:
    print('ERR')
code

5. Що буде виведено за виконання?\par (Дійсні значення, якщо наявні, в цьому завданні будуть виводитись у форматі з фіксованою крапкою та, можливо, з доволі великою кількістю десяткових розрядів. У відповіді для дійсних значень у ЦЬОМУ завданні достатньо вказати один десятковий розряд після крапки, відкинувши решту розрядів без округлення. Наприклад, замість 13.66666666666667 достатньо вказати 13.6, замість 25.23 достатньо вказати 25.2)
code
try:
    print(7, end=';')
    print(int('b2'), end=';')
    print(6, end=';')
except KeyboardInterrupt:
    print(4, end=';')
except TypeError:
    print(3, end=';')
except:
    print('ERR', end=';')
else:
    print(5, end=';')
finally:
    print(7)
code

6. Що буде виведено за виконання?
code
def f(s):
    for el in s:
        el += 1
        return s
 
s = [1, 2, 3]
s = f(s)
print(s)
code

7. Що буде виведено за виконання?
code
class A:
    a = 1

    def f1(self): print(2, end=';')
    def _f2(self): print(3, end=';')

    def main(self):
        try: print(self.a, end=';')
        except: print('ERR', end=';')

        try: self.f1()
        except: print('ERR', end=';')

        try: self._f2()
        except: print('ERR', end=';')


class B(A):
    a = 4

    def f1(self): print(5, end=';')
    def _f2(self): print(6, end=';')

try: B().main()
except: print('ERR', end=';')
code

8. Що буде виведено за виконання?
code
def f(arg, n):
    arg[22] = 9
    n = 8
    arg = set(arg.items())

try:
    n = 1
    t = {22:6, 39:4, 42:9, 68:2}
    f(t, n)
    print(n, end=';')
    for x in t.values():
        print(x, sep=';', end=';')
except:
    print('ERR')
code

9. Що буде виведено за виконання?
code
class A:
    b = 1
    a = 2

    def __init__(self):
        a = 3

class B(A):
    b = 4

    def __init__(self):
        super().__init__()
        self.c = 7

obj = B()
obj.c = obj.b + 10
B.c = 13

try: print(obj.c, end= ';')
except: print('ERR', end=';')

try: print(A.c, end= ';')
except: print('ERR', end=';')

try: print(B.c, end= ';')
except: print('ERR', end=';')
code

10. 
Для графа на рисунку з вершини 6 запущено алгоритм пошуку в глибину, що будує кістякове дерево. Списки суміжності вершин переглядаються алгоритмом за зростанням міток вершин.\par
\begin{center}
	\adjustimage{max size={0.9\linewidth}{0.9\paperheight}}
	{../10/Traversals/127.png}
\end{center}
\par Які з наведених ребер входять до складу отриманого кістякового дерева?\par
1) (2, 3)\par
2) (3, 5)\par
3) (0, 5)\par
4) (1, 6)\par
5) (3, 4)\par
6) (0, 6)\par
{\vskip 15pt} У формі вибрати номери відповідних ребер.\par
%answer:1 2 4 5 6
\end {large}}
%end
{\smallskip \begin {small} Затверджено на засіданні кафедри теоретичної кібернетики, протокол \No 12 від   19.05.2020 р.\par\smallskip\smallskip
{\parbox {6.5 cm}{Зав. кафедри \dotfill Ю.В.~Крак}}\hspace {0.7cm}{\hbox to 6.5 cm{ Екзаменатор \dotfill Т.О.~Карнаух}} \end{small}}
%begin
{{\begin {small}\par
{ \center  Київський національний університет імені Тараса Шевченка \par}
{\parbox {5 cm}{Освітня програма \hfill}{\parbox {7 cm}{Прикладна математика, Системний аналіз \hfill}}\parbox {6 cm}{\hfill Семестр 2}}\par
{\parbox {5 cm} {Навчальний предмет \hfill}\parbox {7 cm}{Програмування \hfill}\parbox {6cm}{\hfill}\par}\vspace{-7pt} \end{small}}{\center Екзаменаційний білет \No \textbf 
{ 47 }\par}}
{
\begin {large}
1. Що буде виведено за виконання?
code
try:
    res = "8.0"!=2.0j
    if isinstance(res, bool):
        print(f'bool;{res}')
    elif isinstance(res, int):
        print(f'int;{res}')
    elif isinstance(res, float):
        print(f'float;{res:.2f}')
    else:
        print('???')
except:
    print('ERR')
code

2. Що буде виведено за виконання?
code
a = 4
b = 9
c = 0

def f(a):
    a = 1
    b = 4
    c = 5
    return a + b + c

a = 7
b = 2
c = 4

try:
    print(f(b), a, b, c, sep=';')
except:
    print('ERR', a, b, c, sep=';')
code

3. Що буде виведено за виконання?
code
try:
    i = 4
    while i<13:
        i += 2
    print(i, end=';')
except:
    print('ERR')
code

4. Що буде виведено за виконання?
code
try:
    for e in range(2, 8, -2):
        if e>7:
            continue
        print(e, end=';')
    else:
        print(e,end=';')
    print(e)
except:
    print('ERR')
code

5. Що буде виведено за виконання?\par (Дійсні значення, якщо наявні, в цьому завданні будуть виводитись у форматі з фіксованою крапкою та, можливо, з доволі великою кількістю десяткових розрядів. У відповіді для дійсних значень у ЦЬОМУ завданні достатньо вказати один десятковий розряд після крапки, відкинувши решту розрядів без округлення. Наприклад, замість 13.66666666666667 достатньо вказати 13.6, замість 25.23 достатньо вказати 25.2)
code
try:
    print(3, end=';')
    print(8/0.0, end=';')
    print(2, end=';')
except ZeroDivisionError:
    print(0, end=';')
except Exception:
    print(6, end=';')
except:
    print('ERR', end=';')
else:
    print(1, end=';')
finally:
    print(4)
code

6. Що буде виведено за виконання?
code
def f(n):
    if n>9: f(n//10)
    else:
        print(n%10, end='')
 
f(123456)
code

7. Що буде виведено за виконання?
code
class A:
    d = 1

    def __h1(self): print(2, end=';')
    def _h2(self): print(3, end=';')

    def main(self):
        try: print(self.d, end=';')
        except: print('ERR', end=';')

        try: self.__h1()
        except: print('ERR', end=';')

        try: self._h2()
        except: print('ERR', end=';')


class B(A):
    d = 4

    def __h1(self): print(5, end=';')
    def _h2(self): print(6, end=';')

try: B().main()
except: print('ERR', end=';')
code

8. Що буде виведено за виконання?
code
def f(arg, n):
    arg[31] = 6
    n = 5
    arg = list(arg.items())

try:
    n = 3
    d = {63:6, 31:4, 31:4}
    f(d, n)
    print(n, end=';')
    for x in d.values():
        print(x, sep=';', end=';')
except:
    print('ERR')
code

9. Що буде виведено за виконання?
code
class A:
    a = 1
    c = 2

    def __init__(self):
        c = 3

class B(A):
    c = 4

    def __init__(self):
        super().__init__()
        self.b = 7

obj = B()
obj.a = obj.a + 10
B.a = 13

try: print(obj.a, end= ';')
except: print('ERR', end=';')

try: print(A.a, end= ';')
except: print('ERR', end=';')

try: print(B.a, end= ';')
except: print('ERR', end=';')
code

10. 
Для графа на рисунку з вершини 5 запущено алгоритм пошуку в глибину, що будує кістякове дерево. Списки суміжності вершин переглядаються алгоритмом за зростанням міток вершин.\par
\begin{center}
	\adjustimage{max size={0.9\linewidth}{0.9\paperheight}}
	{../10/Traversals/196.png}
\end{center}
\par Які з наведених ребер входять до складу отриманого кістякового дерева?\par
1) (1, 5)\par
2) (0, 2)\par
3) (4, 5)\par
4) (3, 6)\par
5) (2, 5)\par
6) (0, 6)\par
{\vskip 15pt} У формі вибрати номери відповідних ребер.\par
%answer:2 4 6
\end {large}}
%end
{\smallskip \begin {small} Затверджено на засіданні кафедри теоретичної кібернетики, протокол \No 12 від   19.05.2020 р.\par\smallskip\smallskip
{\parbox {6.5 cm}{Зав. кафедри \dotfill Ю.В.~Крак}}\hspace {0.7cm}{\hbox to 6.5 cm{ Екзаменатор \dotfill Т.О.~Карнаух}} \end{small}}
%begin
{{\begin {small}\par
{ \center  Київський національний університет імені Тараса Шевченка \par}
{\parbox {5 cm}{Освітня програма \hfill}{\parbox {7 cm}{Прикладна математика, Системний аналіз \hfill}}\parbox {6 cm}{\hfill Семестр 2}}\par
{\parbox {5 cm} {Навчальний предмет \hfill}\parbox {7 cm}{Програмування \hfill}\parbox {6cm}{\hfill}\par}\vspace{-7pt} \end{small}}{\center Екзаменаційний білет \No \textbf 
{ 48 }\par}}
{
\begin {large}
1. Що буде виведено за виконання?
code
try:
    res = 5/3*9*5
    if isinstance(res, bool):
        print(f'bool;{res}')
    elif isinstance(res, int):
        print(f'int;{res}')
    elif isinstance(res, float):
        print(f'float;{res:.2f}')
    else:
        print('???')
except:
    print('ERR')
code

2. Що буде виведено за виконання?
code
a = 6
b = 8
c = 6

def g(a):
    a = 5
    b += 4
    c = 1
    return a + b + c

a = 7
b = 9
c = 5

try:
    print(g(a), a, b, c, sep=';')
except:
    print('ERR', a, b, c, sep=';')
code

3. Що буде виведено за виконання?
code
try:
    k = 8
    while k>1:
        k += -2
    print(k, end=';')
except:
    print('ERR')
code

4. Що буде виведено за виконання?
code
try:
    for f in range(11, 1, 2):
        if f<8:
            break
        print(f, end=';');f = 1
        if f<8:
            break
    else:
        print(f,end=';')
    print(f)
except:
    print('ERR')
code

5. Що буде виведено за виконання?\par (Дійсні значення, якщо наявні, в цьому завданні будуть виводитись у форматі з фіксованою крапкою та, можливо, з доволі великою кількістю десяткових розрядів. У відповіді для дійсних значень у ЦЬОМУ завданні достатньо вказати один десятковий розряд після крапки, відкинувши решту розрядів без округлення. Наприклад, замість 13.66666666666667 достатньо вказати 13.6, замість 25.23 достатньо вказати 25.2)
code
try:
    print(9, end=';')
    print(6j>=2j, end=';')
    print(5, end=';')
except ValueError:
    print(9, end=';')
except KeyboardInterrupt:
    print(1, end=';')
except:
    print('ERR', end=';')
else:
    print(8, end=';')
finally:
    print(3)
code

6. Що буде виведено за виконання?
code
def f(n):
    if n>9:
        f(n//100)
    print(n%10, end='')
 
f(123456)
code

7. Що буде виведено за виконання?
code
class A:
    b = 1

    def _g1(self): print(2, end=';')
    def g2(self): print(3, end=';')

    def main(self):
        try: print(self.b, end=';')
        except: print('ERR', end=';')

        try: self._g1()
        except: print('ERR', end=';')

        try: self.g2()
        except: print('ERR', end=';')


class B(A):
    b = 4

    def _g1(self): print(5, end=';')
    def g2(self): print(6, end=';')

try: B().main()
except: print('ERR', end=';')
code

8. Що буде виведено за виконання?
code
def f(arg, n):
    arg[83] = 7
    n = 9
    arg = tuple(arg.values())

try:
    n = 2
    t = {46:5, 68:9, 39:5, 39:3}
    f(t, n)
    print(n, end=';')
    for x in t.keys():
        print(x, sep=';', end=';')
except:
    print('ERR')
code

9. Що буде виведено за виконання?
code
class A:
    b = 1
    c = 2

    def __init__(self):
        b = 3

class B(A):
    b = 4

    def __init__(self):
        super().__init__()
        self.a = 7

obj = B()
obj.c = obj.b + 10
B.c = 13

try: print(obj.c, end= ';')
except: print('ERR', end=';')

try: print(A.c, end= ';')
except: print('ERR', end=';')

try: print(B.c, end= ';')
except: print('ERR', end=';')
code

10. 
Для графа на рисунку з вершини 5 запущено алгоритм пошуку в глибину, що будує кістякове дерево. Списки суміжності вершин переглядаються алгоритмом за зростанням міток вершин.\par
\begin{center}
	\adjustimage{max size={0.9\linewidth}{0.9\paperheight}}
	{../10/Traversals/140.png}
\end{center}
\par Які з наведених ребер входять до складу отриманого кістякового дерева?\par
1) (1, 6)\par
2) (0, 5)\par
3) (3, 5)\par
4) (3, 4)\par
5) (0, 4)\par
6) (4, 6)\par
{\vskip 15pt} У формі вибрати номери відповідних ребер.\par
%answer:1 2 4 5
\end {large}}
%end
{\smallskip \begin {small} Затверджено на засіданні кафедри теоретичної кібернетики, протокол \No 12 від   19.05.2020 р.\par\smallskip\smallskip
{\parbox {6.5 cm}{Зав. кафедри \dotfill Ю.В.~Крак}}\hspace {0.7cm}{\hbox to 6.5 cm{ Екзаменатор \dotfill Т.О.~Карнаух}} \end{small}}
%begin
{{\begin {small}\par
{ \center  Київський національний університет імені Тараса Шевченка \par}
{\parbox {5 cm}{Освітня програма \hfill}{\parbox {7 cm}{Прикладна математика, Системний аналіз \hfill}}\parbox {6 cm}{\hfill Семестр 2}}\par
{\parbox {5 cm} {Навчальний предмет \hfill}\parbox {7 cm}{Програмування \hfill}\parbox {6cm}{\hfill}\par}\vspace{-7pt} \end{small}}{\center Екзаменаційний білет \No \textbf 
{ 49 }\par}}
{
\begin {large}
1. Що буде виведено за виконання?
code
try:
    res = not 4<=True and 5>7 or 7.0<3.0
    if isinstance(res, bool):
        print(f'bool;{res}')
    elif isinstance(res, int):
        print(f'int;{res}')
    elif isinstance(res, float):
        print(f'float;{res:.2f}')
    else:
        print('???')
except:
    print('ERR')
code

2. Що буде виведено за виконання?
code
a = 5
b = 6
c = 3

def h(b):
    a = 2
    b = 1
    c = 3
    return a + b + c

a = 1
b = 8
c = 9

try:
    print(h(b), a, b, c, sep=';')
except:
    print('ERR', a, b, c, sep=';')
code

3. Що буде виведено за виконання?
code
try:
    n = 9
    while n>2:
        n += -2
    print(n, end=';')
except:
    print('ERR')
code

4. Що буде виведено за виконання?
code
try:
    for a in range(-9, -10, -3):
        if a>=7:
            continue
        print(a, end=';')
        if a<=5:
            break
    else:
        print(a,end=';')
    print(a)
except:
    print('ERR')
code

5. Що буде виведено за виконання?\par (Дійсні значення, якщо наявні, в цьому завданні будуть виводитись у форматі з фіксованою крапкою та, можливо, з доволі великою кількістю десяткових розрядів. У відповіді для дійсних значень у ЦЬОМУ завданні достатньо вказати один десятковий розряд після крапки, відкинувши решту розрядів без округлення. Наприклад, замість 13.66666666666667 достатньо вказати 13.6, замість 25.23 достатньо вказати 25.2)
code
try:
    print(4, end=';')
    print(0%2, end=';')
    print(7, end=';')
except TypeError:
    print(2, end=';')
except Exception:
    print(8, end=';')
except:
    print('ERR', end=';')
else:
    print(1, end=';')
finally:
    print(5)
code

6. Що буде виведено за виконання?
code
def f(n):
    print(n%10, end='')
    if n>9:
        f(n//100)
 
f(987654)
code

7. Що буде виведено за виконання?
code
class A:
    d = 1

    def _h1(self): print(2, end=';')
    def _h2(self): print(3, end=';')

    def main(self):
        try: print(self.d, end=';')
        except: print('ERR', end=';')

        try: self._h1()
        except: print('ERR', end=';')

        try: self._h2()
        except: print('ERR', end=';')


class B(A):
    d = 4

    def _h1(self): print(5, end=';')
    def _h2(self): print(6, end=';')

try: B().main()
except: print('ERR', end=';')
code

8. Що буде виведено за виконання?
code
def f(arg, n):
    arg[38] = 4
    n += 3

try:
    n = 6
    d = {45:5, 10:9, 38:1, 10:4}
    f(d, n)
    print(n, end=';')
    for x in d.items():
        print(*x, sep=';', end=';')
except:
    print('ERR')
code

9. Що буде виведено за виконання?
code
class A:
    c = 1
    b = 2

    def __init__(self):
        b = 3

class B(A):
    b = 4

    def __init__(self):
        super().__init__()
        self.a = 7

obj = B()
obj.b = obj.a + 10
B.b = 13

try: print(obj.b, end= ';')
except: print('ERR', end=';')

try: print(A.b, end= ';')
except: print('ERR', end=';')

try: print(B.b, end= ';')
except: print('ERR', end=';')
code

10. 
Для графа на рисунку з вершини 3 запущено алгоритм пошуку в глибину, що будує кістякове дерево. Списки суміжності вершин переглядаються алгоритмом за зростанням міток вершин.\par
\begin{center}
	\adjustimage{max size={0.9\linewidth}{0.9\paperheight}}
	{../10/Traversals/124.png}
\end{center}
\par Які з наведених ребер входять до складу отриманого кістякового дерева?\par
1) (1, 5)\par
2) (0, 6)\par
3) (1, 4)\par
4) (4, 5)\par
5) (2, 3)\par
6) (1, 2)\par
{\vskip 15pt} У формі вибрати номери відповідних ребер.\par
%answer:2 3 4 5 6
\end {large}}
%end
{\smallskip \begin {small} Затверджено на засіданні кафедри теоретичної кібернетики, протокол \No 12 від   19.05.2020 р.\par\smallskip\smallskip
{\parbox {6.5 cm}{Зав. кафедри \dotfill Ю.В.~Крак}}\hspace {0.7cm}{\hbox to 6.5 cm{ Екзаменатор \dotfill Т.О.~Карнаух}} \end{small}}
%begin
{{\begin {small}\par
{ \center  Київський національний університет імені Тараса Шевченка \par}
{\parbox {5 cm}{Освітня програма \hfill}{\parbox {7 cm}{Прикладна математика, Системний аналіз \hfill}}\parbox {6 cm}{\hfill Семестр 2}}\par
{\parbox {5 cm} {Навчальний предмет \hfill}\parbox {7 cm}{Програмування \hfill}\parbox {6cm}{\hfill}\par}\vspace{-7pt} \end{small}}{\center Екзаменаційний білет \No \textbf 
{ 50 }\par}}
{
\begin {large}
1. Що буде виведено за виконання?
code
try:
    res = 4/9//3*8
    if isinstance(res, bool):
        print(f'bool;{res}')
    elif isinstance(res, int):
        print(f'int;{res}')
    elif isinstance(res, float):
        print(f'float;{res:.2f}')
    else:
        print('???')
except:
    print('ERR')
code

2. Що буде виведено за виконання?
code
a = 3
b = 9
c = 4

def g(a):
    global c
    a = 4
    b = 3
    c = 2
    return a + b + c

a = 6
b = 1
c = 7

try:
    print(g(a), a, b, c, sep=';')
except:
    print('ERR', a, b, c, sep=';')
code

3. Що буде виведено за виконання?
code
try:
    m = 5
    while m<=8:
        m += 1
    print(m, end=';')
except:
    print('ERR')
code

4. Що буде виведено за виконання?
code
try:
    for c in range(15, 3, -1):
        if c>3:
            continue
        print(c, end=';')
    else:
        print(c,end=';')
    print(c)
except:
    print('ERR')
code

5. Що буде виведено за виконання?\par (Дійсні значення, якщо наявні, в цьому завданні будуть виводитись у форматі з фіксованою крапкою та, можливо, з доволі великою кількістю десяткових розрядів. У відповіді для дійсних значень у ЦЬОМУ завданні достатньо вказати один десятковий розряд після крапки, відкинувши решту розрядів без округлення. Наприклад, замість 13.66666666666667 достатньо вказати 13.6, замість 25.23 достатньо вказати 25.2)
code
try:
    print(6, end=';')
    print(int('9'), end=';')
    print(4, end=';')
except ZeroDivisionError:
    print(7, end=';')
except ValueError:
    print(0, end=';')
except:
    print('ERR', end=';')
else:
    print(3, end=';')
finally:
    print(2)
code

6. Що буде виведено за виконання?
code
def f(n):
    if n>9:
        f(n//100)
        print(n%10, end='')
 
f(123456)
code

7. Що буде виведено за виконання?
code
class A:
    b = 1

    def _f1(self): print(2, end=';')
    def __f2(self): print(3, end=';')

    def main(self):
        try: print(self.b, end=';')
        except: print('ERR', end=';')

        try: self._f1()
        except: print('ERR', end=';')

        try: self.__f2()
        except: print('ERR', end=';')


class B(A):
    b = 4

    def _f1(self): print(5, end=';')
    def __f2(self): print(6, end=';')

try: B().main()
except: print('ERR', end=';')
code

8. Що буде виведено за виконання?
code
def f(arg, n):
    arg[56] = 3
    n *= -2

try:
    n = 3
    t = {10:8, 28:3, 84:9, 84:0}
    f(t, n)
    print(n, end=';')
    for x, y in t.items():
        print(x, y, sep=';', end=';')
except:
    print('ERR')
code

9. Що буде виведено за виконання?
code
class A:
    c = 1
    a = 2

    def __init__(self):
        c = 3

class B(A):
    c = 4

    def __init__(self):
        super().__init__()
        self.b = 7

obj = B()
obj.c = obj.b + 10
B.c = 13

try: print(obj.c, end= ';')
except: print('ERR', end=';')

try: print(A.c, end= ';')
except: print('ERR', end=';')

try: print(B.c, end= ';')
except: print('ERR', end=';')
code

10. 
Для графа на рисунку з вершини 0 запущено алгоритм пошуку в глибину, що будує кістякове дерево. Списки суміжності вершин переглядаються алгоритмом за зростанням міток вершин.\par
\begin{center}
	\adjustimage{max size={0.9\linewidth}{0.9\paperheight}}
	{../10/Traversals/287.png}
\end{center}
\par Які з наведених ребер входять до складу отриманого кістякового дерева?\par
1) (0, 1)\par
2) (0, 2)\par
3) (0, 6)\par
4) (5, 6)\par
5) (4, 5)\par
6) (1, 5)\par
{\vskip 15pt} У формі вибрати номери відповідних ребер.\par
%answer:1 4 5
\end {large}}
%end
{\smallskip \begin {small} Затверджено на засіданні кафедри теоретичної кібернетики, протокол \No 12 від   19.05.2020 р.\par\smallskip\smallskip
{\parbox {6.5 cm}{Зав. кафедри \dotfill Ю.В.~Крак}}\hspace {0.7cm}{\hbox to 6.5 cm{ Екзаменатор \dotfill Т.О.~Карнаух}} \end{small}}
%begin
{{\begin {small}\par
{ \center  Київський національний університет імені Тараса Шевченка \par}
{\parbox {5 cm}{Освітня програма \hfill}{\parbox {7 cm}{Прикладна математика, Системний аналіз \hfill}}\parbox {6 cm}{\hfill Семестр 2}}\par
{\parbox {5 cm} {Навчальний предмет \hfill}\parbox {7 cm}{Програмування \hfill}\parbox {6cm}{\hfill}\par}\vspace{-7pt} \end{small}}{\center Екзаменаційний білет \No \textbf 
{ 51 }\par}}
{
\begin {large}
1. Що буде виведено за виконання?
code
try:
    res = 2<=True or not 6<5 and 8.0>7.0
    if isinstance(res, bool):
        print(f'bool;{res}')
    elif isinstance(res, int):
        print(f'int;{res}')
    elif isinstance(res, float):
        print(f'float;{res:.2f}')
    else:
        print('???')
except:
    print('ERR')
code

2. Що буде виведено за виконання?
code
a = 5
b = 0
c = 2

def f(b):
    a = 1
    b = 5
    c = 5
    return a + b + c

a = 8
b = 2
c = 8

try:
    print(f(b), a, b, c, sep=';')
except:
    print('ERR', a, b, c, sep=';')
code

3. Що буде виведено за виконання?
code
try:
    n = 3
    while n<5:
        n += 2
    print(n, end=';')
except:
    print('ERR')
code

4. Що буде виведено за виконання?
code
try:
    for f in range(-12, -5, -1):
        if f<=8:
            break
        print(f, end=';');f = -7
    else:
        print('end',end=';')
    print(f)
except:
    print('ERR')
code

5. Що буде виведено за виконання?\par (Дійсні значення, якщо наявні, в цьому завданні будуть виводитись у форматі з фіксованою крапкою та, можливо, з доволі великою кількістю десяткових розрядів. У відповіді для дійсних значень у ЦЬОМУ завданні достатньо вказати один десятковий розряд після крапки, відкинувши решту розрядів без округлення. Наприклад, замість 13.66666666666667 достатньо вказати 13.6, замість 25.23 достатньо вказати 25.2)
code
try:
    print(1, end=';')
    print(2//3, end=';')
    print(3, end=';')
except Exception:
    print(9, end=';')
except TypeError:
    print(8, end=';')
except:
    print('ERR', end=';')
else:
    print(0, end=';')
finally:
    print(4)
code

6. Що буде виведено за виконання?
code
def f(n):
    if n>9:
        return f(n//100)+1
    else:
        return 1

print(f(123456))
code

7. Що буде виведено за виконання?
code
class A:
    a = 1

    def _g1(self): print(2, end=';')
    def _g2(self): print(3, end=';')

    def main(self):
        try: print(self.a, end=';')
        except: print('ERR', end=';')

        try: self._g1()
        except: print('ERR', end=';')

        try: self._g2()
        except: print('ERR', end=';')


class B(A):
    a = 4

    def _g1(self): print(5, end=';')
    def _g2(self): print(6, end=';')

try: B().main()
except: print('ERR', end=';')
code

8. Що буде виведено за виконання?
code
def f(arg, n):
    n = 5
    tmp = arg
    tmp[-4:] = 'ab'
    return len(tmp)>3

try:
    g = ['0','1','2','3','4']
    n = 3
    f(g, n)
    print(n, end=';')
    for el in g:
        print(el, end=';')
except:
    print('ERR')
code

9. Що буде виведено за виконання?
code
class A:
    b = 1
    a = 2

    def __init__(self):
        b = 3

class B(A):
    a = 4

    def __init__(self):
        super().__init__()
        self.c = 7

obj = B()
obj.c = obj.a + 10
B.c = 13

try: print(obj.c, end= ';')
except: print('ERR', end=';')

try: print(A.c, end= ';')
except: print('ERR', end=';')

try: print(B.c, end= ';')
except: print('ERR', end=';')
code

10. 
Для графа на рисунку з вершини 2 запущено алгоритм пошуку в глибину, що будує кістякове дерево. Списки суміжності вершин переглядаються алгоритмом за зростанням міток вершин.\par
\begin{center}
	\adjustimage{max size={0.9\linewidth}{0.9\paperheight}}
	{../10/Traversals/102.png}
\end{center}
\par Які з наведених ребер входять до складу отриманого кістякового дерева?\par
1) (2, 3)\par
2) (2, 4)\par
3) (0, 5)\par
4) (4, 5)\par
5) (2, 5)\par
6) (0, 2)\par
{\vskip 15pt} У формі вибрати номери відповідних ребер.\par
%answer:4 6
\end {large}}
%end
{\smallskip \begin {small} Затверджено на засіданні кафедри теоретичної кібернетики, протокол \No 12 від   19.05.2020 р.\par\smallskip\smallskip
{\parbox {6.5 cm}{Зав. кафедри \dotfill Ю.В.~Крак}}\hspace {0.7cm}{\hbox to 6.5 cm{ Екзаменатор \dotfill Т.О.~Карнаух}} \end{small}}
%begin
{{\begin {small}\par
{ \center  Київський національний університет імені Тараса Шевченка \par}
{\parbox {5 cm}{Освітня програма \hfill}{\parbox {7 cm}{Прикладна математика, Системний аналіз \hfill}}\parbox {6 cm}{\hfill Семестр 2}}\par
{\parbox {5 cm} {Навчальний предмет \hfill}\parbox {7 cm}{Програмування \hfill}\parbox {6cm}{\hfill}\par}\vspace{-7pt} \end{small}}{\center Екзаменаційний білет \No \textbf 
{ 52 }\par}}
{
\begin {large}
1. Що буде виведено за виконання?
code
def f(n):
    print('f', end='')
    return n

try:
    print(f(-4) or f(1))
except:
    print('ERR')
code

2. Що буде виведено за виконання?
code
a = 3
b = 2
c = 1

def h(b):
    global c
    a = 3
    b += 5
    c = 2
    return a + b + c

a = 0
b = 4
c = 8

try:
    print(h(b), a, b, c, sep=';')
except:
    print('ERR', a, b, c, sep=';')
code

3. Що буде виведено за виконання?
code
try:
    m = 1
    while m<=10:
        m += 2
    print(m, end=';')
except:
    print('ERR')
code

4. Що буде виведено за виконання?
code
try:
    for e in range(12, -4, -2):
        if e<6:
            continue
        print(e, end=';')
    else:
        print(e,end=';')
    print(e)
except:
    print('ERR')
code

5. Що буде виведено за виконання?\par (Дійсні значення, якщо наявні, в цьому завданні будуть виводитись у форматі з фіксованою крапкою та, можливо, з доволі великою кількістю десяткових розрядів. У відповіді для дійсних значень у ЦЬОМУ завданні достатньо вказати один десятковий розряд після крапки, відкинувши решту розрядів без округлення. Наприклад, замість 13.66666666666667 достатньо вказати 13.6, замість 25.23 достатньо вказати 25.2)
code
try:
    print(6, end=';')
    print(int('5'), end=';')
    print(7, end=';')
except KeyboardInterrupt:
    print(9, end=';')
except ZeroDivisionError:
    print(1, end=';')
except:
    print('ERR', end=';')
else:
    print(4, end=';')
finally:
    print(8)
code

6. Що буде виведено за виконання?
code
def f(n):
    if n>2:
        f(n//2)
    print(n%2, end='')
 
f(16)
code

7. Що буде виведено за виконання?
code
class A:
    c = 1

    def __f1(self): print(2, end=';')
    def _f2(self): print(3, end=';')

    def main(self):
        try: print(self.c, end=';')
        except: print('ERR', end=';')

        try: self.__f1()
        except: print('ERR', end=';')

        try: self._f2()
        except: print('ERR', end=';')


class B(A):
    c = 4

    def __f1(self): print(5, end=';')
    def _f2(self): print(6, end=';')

try: B().main()
except: print('ERR', end=';')
code

8. Що буде виведено за виконання?
code
def f(arg, n):
    n = 8
    tmp = arg
    del tmp[:-6]
    return len(tmp)>3

try:
    d = ['0','1','2','3','4']
    n = 2
    f(d, n)
    print(n, end=';')
    for el in d:
        print(el, end=';')
except:
    print('ERR')
code

9. Що буде виведено за виконання?
code
class A:
    a = 1
    c = 2

    def __init__(self):
        c = 3

class B(A):
    a = 4

    def __init__(self):
        super().__init__()
        self.b = 7

obj = B()
obj.c = obj.a + 10
B.c = 13

try: print(obj.c, end= ';')
except: print('ERR', end=';')

try: print(A.c, end= ';')
except: print('ERR', end=';')

try: print(B.c, end= ';')
except: print('ERR', end=';')
code

10. 
Для графа на рисунку з вершини 2 запущено алгоритм пошуку в глибину, що будує кістякове дерево. Списки суміжності вершин переглядаються алгоритмом за зростанням міток вершин.\par
\begin{center}
	\adjustimage{max size={0.9\linewidth}{0.9\paperheight}}
	{../10/Traversals/268.png}
\end{center}
\par Які з наведених ребер входять до складу отриманого кістякового дерева?\par
1) (1, 4)\par
2) (0, 3)\par
3) (2, 3)\par
4) (1, 3)\par
5) (3, 4)\par
6) (4, 5)\par
{\vskip 15pt} У формі вибрати номери відповідних ребер.\par
%answer:1 2 3 4 6
\end {large}}
%end
{\smallskip \begin {small} Затверджено на засіданні кафедри теоретичної кібернетики, протокол \No 12 від   19.05.2020 р.\par\smallskip\smallskip
{\parbox {6.5 cm}{Зав. кафедри \dotfill Ю.В.~Крак}}\hspace {0.7cm}{\hbox to 6.5 cm{ Екзаменатор \dotfill Т.О.~Карнаух}} \end{small}}
%begin
{{\begin {small}\par
{ \center  Київський національний університет імені Тараса Шевченка \par}
{\parbox {5 cm}{Освітня програма \hfill}{\parbox {7 cm}{Прикладна математика, Системний аналіз \hfill}}\parbox {6 cm}{\hfill Семестр 2}}\par
{\parbox {5 cm} {Навчальний предмет \hfill}\parbox {7 cm}{Програмування \hfill}\parbox {6cm}{\hfill}\par}\vspace{-7pt} \end{small}}{\center Екзаменаційний білет \No \textbf 
{ 53 }\par}}
{
\begin {large}
1. Що буде виведено за виконання?
code
try:
    res = not 7==3 and 3.0!=False or 9.0>=8
    if isinstance(res, bool):
        print(f'bool;{res}')
    elif isinstance(res, int):
        print(f'int;{res}')
    elif isinstance(res, float):
        print(f'float;{res:.2f}')
    else:
        print('???')
except:
    print('ERR')
code

2. Що буде виведено за виконання?
code
a = 6
b = 7
c = 2

def g(a):
    a = 2
    b *= 4
    c = 3
    return a + b + c

a = 0
b = 3
c = 4

try:
    print(g(a), a, b, c, sep=';')
except:
    print('ERR', a, b, c, sep=';')
code

3. Що буде виведено за виконання?
code
try:
    n = 2
    while n<9:
        n += 1
    print(n, end=';')
except:
    print('ERR')
code

4. Що буде виведено за виконання?
code
try:
    for d in range(8, 4, 3):
        if d<5:
            continue
        print(d, end=';');d = -6
    else:
        print(d,end=';')
    print(d)
except:
    print('ERR')
code

5. Що буде виведено за виконання?\par (Дійсні значення, якщо наявні, в цьому завданні будуть виводитись у форматі з фіксованою крапкою та, можливо, з доволі великою кількістю десяткових розрядів. У відповіді для дійсних значень у ЦЬОМУ завданні достатньо вказати один десятковий розряд після крапки, відкинувши решту розрядів без округлення. Наприклад, замість 13.66666666666667 достатньо вказати 13.6, замість 25.23 достатньо вказати 25.2)
code
try:
    print(7, end=';')
    print(5j<=8j, end=';')
    print(0, end=';')
except ValueError:
    print(3, end=';')
except ZeroDivisionError:
    print(2, end=';')
except:
    print('ERR', end=';')
else:
    print(6, end=';')
finally:
    print(1)
code

6. Що буде виведено за виконання?
code
def f(n):
    if n>9:
        f(n//100)
    else:
        print(n%10,end='')

f(123456)
code

7. Що буде виведено за виконання?
code
class A:
    d = 1

    def _h1(self): print(2, end=';')
    def __h2(self): print(3, end=';')

    def main(self):
        try: print(self.d, end=';')
        except: print('ERR', end=';')

        try: self._h1()
        except: print('ERR', end=';')

        try: self.__h2()
        except: print('ERR', end=';')


class B(A):
    d = 4

    def _h1(self): print(5, end=';')
    def __h2(self): print(6, end=';')

try: B().main()
except: print('ERR', end=';')
code

8. Що буде виведено за виконання?
code
def f(arg, n):
    n = 5
    tmp = arg
    del tmp[-1:-8]
    return len(tmp)>3

try:
    a = [10,11,12,13,14]
    n = 1
    f(a, n)
    print(n, end=';')
    for el in a:
        print(el, end=';')
except:
    print('ERR')
code

9. Що буде виведено за виконання?
code
class A:
    a = 1
    b = 2

    def __init__(self):
        a = 3

class B(A):
    b = 4

    def __init__(self):
        super().__init__()
        self.c = 7

obj = B()
obj.b = obj.b + 10
B.b = 13

try: print(obj.b, end= ';')
except: print('ERR', end=';')

try: print(A.b, end= ';')
except: print('ERR', end=';')

try: print(B.b, end= ';')
except: print('ERR', end=';')
code

10. 
Для графа на рисунку з вершини 0 запущено алгоритм пошуку в глибину, що будує кістякове дерево. Списки суміжності вершин переглядаються алгоритмом за зростанням міток вершин.\par
\begin{center}
	\adjustimage{max size={0.9\linewidth}{0.9\paperheight}}
	{../10/Traversals/146.png}
\end{center}
\par Які з наведених ребер входять до складу отриманого кістякового дерева?\par
1) (1, 5)\par
2) (3, 5)\par
3) (0, 1)\par
4) (2, 6)\par
5) (5, 6)\par
6) (3, 4)\par
{\vskip 15pt} У формі вибрати номери відповідних ребер.\par
%answer:2 3 4 5 6
\end {large}}
%end
{\smallskip \begin {small} Затверджено на засіданні кафедри теоретичної кібернетики, протокол \No 12 від   19.05.2020 р.\par\smallskip\smallskip
{\parbox {6.5 cm}{Зав. кафедри \dotfill Ю.В.~Крак}}\hspace {0.7cm}{\hbox to 6.5 cm{ Екзаменатор \dotfill Т.О.~Карнаух}} \end{small}}
%begin
{{\begin {small}\par
{ \center  Київський національний університет імені Тараса Шевченка \par}
{\parbox {5 cm}{Освітня програма \hfill}{\parbox {7 cm}{Прикладна математика, Системний аналіз \hfill}}\parbox {6 cm}{\hfill Семестр 2}}\par
{\parbox {5 cm} {Навчальний предмет \hfill}\parbox {7 cm}{Програмування \hfill}\parbox {6cm}{\hfill}\par}\vspace{-7pt} \end{small}}{\center Екзаменаційний білет \No \textbf 
{ 54 }\par}}
{
\begin {large}
1. Що буде виведено за виконання?
code
def f(n):
    print('f', end='')
    return n>2

try:
    print(f(8) or f(-3))
except:
    print('ERR')
code

2. Що буде виведено за виконання?
code
a = 1
b = 5
c = 8

def f(b):
    global c
    a = 1
    b -= 5
    c = 2
    return a + b + c

a = 9
b = 6
c = 3

try:
    print(f(b), a, b, c, sep=';')
except:
    print('ERR', a, b, c, sep=';')
code

3. Що буде виведено за виконання?
code
try:
    j = 13
    while j>=8:
        j += -3
    print(j, end=';')
except:
    print('ERR')
code

4. Що буде виведено за виконання?
code
try:
    for b in range(3, -3, 2):
        if b<=4:
            break
        print(b, end=';')
    else:
        print(b,end=';')
    print(b)
except:
    print('ERR')
code

5. Що буде виведено за виконання?\par (Дійсні значення, якщо наявні, в цьому завданні будуть виводитись у форматі з фіксованою крапкою та, можливо, з доволі великою кількістю десяткових розрядів. У відповіді для дійсних значень у ЦЬОМУ завданні достатньо вказати один десятковий розряд після крапки, відкинувши решту розрядів без округлення. Наприклад, замість 13.66666666666667 достатньо вказати 13.6, замість 25.23 достатньо вказати 25.2)
code
try:
    print(6, end=';')
    print(int('a9'), end=';')
    print(7, end=';')
except KeyboardInterrupt:
    print(2, end=';')
except TypeError:
    print(4, end=';')
except:
    print('ERR', end=';')
else:
    print(0, end=';')
finally:
    print(8)
code

6. Що буде виведено за виконання?
code
def f(n):
    if n>9:
        f(n//10)
    print(n%10, end='')
 
f(543216)
code

7. Що буде виведено за виконання?
code
class A:
    c = 1

    def _g1(self): print(2, end=';')
    def g2(self): print(3, end=';')

    def main(self):
        try: print(self.c, end=';')
        except: print('ERR', end=';')

        try: self._g1()
        except: print('ERR', end=';')

        try: self.g2()
        except: print('ERR', end=';')


class B(A):
    c = 4

    def _g1(self): print(5, end=';')
    def g2(self): print(6, end=';')

try: B().main()
except: print('ERR', end=';')
code

8. Що буде виведено за виконання?
code
def f(arg, n):
    n = 7
    tmp = arg
    tmp[3:0] = 'bc'
    return len(tmp)>3

try:
    f = ['0','1','2','3','4']
    n = 8
    f(f, n)
    print(n, end=';')
    for el in f:
        print(el, end=';')
except:
    print('ERR')
code

9. Що буде виведено за виконання?
code
class A:
    b = 1
    c = 2

    def __init__(self):
        c = 3

class B(A):
    b = 4

    def __init__(self):
        super().__init__()
        self.a = 7

obj = B()
obj.c = obj.a + 10
B.c = 13

try: print(obj.c, end= ';')
except: print('ERR', end=';')

try: print(A.c, end= ';')
except: print('ERR', end=';')

try: print(B.c, end= ';')
except: print('ERR', end=';')
code

10. 
Для графа на рисунку з вершини 6 запущено алгоритм пошуку в глибину, що будує кістякове дерево. Списки суміжності вершин переглядаються алгоритмом за зростанням міток вершин.\par
\begin{center}
	\adjustimage{max size={0.9\linewidth}{0.9\paperheight}}
	{../10/Traversals/260.png}
\end{center}
\par Які з наведених ребер входять до складу отриманого кістякового дерева?\par
1) (3, 5)\par
2) (1, 3)\par
3) (0, 6)\par
4) (0, 5)\par
5) (0, 1)\par
6) (2, 5)\par
{\vskip 15pt} У формі вибрати номери відповідних ребер.\par
%answer:1 2 3 5 6
\end {large}}
%end
{\smallskip \begin {small} Затверджено на засіданні кафедри теоретичної кібернетики, протокол \No 12 від   19.05.2020 р.\par\smallskip\smallskip
{\parbox {6.5 cm}{Зав. кафедри \dotfill Ю.В.~Крак}}\hspace {0.7cm}{\hbox to 6.5 cm{ Екзаменатор \dotfill Т.О.~Карнаух}} \end{small}}
%begin
{{\begin {small}\par
{ \center  Київський національний університет імені Тараса Шевченка \par}
{\parbox {5 cm}{Освітня програма \hfill}{\parbox {7 cm}{Прикладна математика, Системний аналіз \hfill}}\parbox {6 cm}{\hfill Семестр 2}}\par
{\parbox {5 cm} {Навчальний предмет \hfill}\parbox {7 cm}{Програмування \hfill}\parbox {6cm}{\hfill}\par}\vspace{-7pt} \end{small}}{\center Екзаменаційний білет \No \textbf 
{ 55 }\par}}
{
\begin {large}
1. Що буде виведено за виконання?
code
try:
    res = 9==4 and not 5.0>6.0 and 2>=True
    if isinstance(res, bool):
        print(f'bool;{res}')
    elif isinstance(res, int):
        print(f'int;{res}')
    elif isinstance(res, float):
        print(f'float;{res:.2f}')
    else:
        print('???')
except:
    print('ERR')
code

2. Що буде виведено за виконання?
code
a = 0
b = 6
c = 9

def h(a):
    global c
    a = 3
    b = 2
    c = 1
    return a + b + c

a = 1
b = 4
c = 7

try:
    print(h(a), a, b, c, sep=';')
except:
    print('ERR', a, b, c, sep=';')
code

3. Що буде виведено за виконання?
code
try:
    m = 12
    while m>=3:
        m += -3
    print(m, end=';')
except:
    print('ERR')
code

4. Що буде виведено за виконання?
code
try:
    for b in range(14, -8, -3):
        if b>=6:
            continue
        print(b, end=';');b = -9
    else:
        print(b,end=';')
    print(b)
except:
    print('ERR')
code

5. Що буде виведено за виконання?\par (Дійсні значення, якщо наявні, в цьому завданні будуть виводитись у форматі з фіксованою крапкою та, можливо, з доволі великою кількістю десяткових розрядів. У відповіді для дійсних значень у ЦЬОМУ завданні достатньо вказати один десятковий розряд після крапки, відкинувши решту розрядів без округлення. Наприклад, замість 13.66666666666667 достатньо вказати 13.6, замість 25.23 достатньо вказати 25.2)
code
try:
    print(3, end=';')
    print(5/0.0, end=';')
    print(8, end=';')
except Exception:
    print(1, end=';')
except ValueError:
    print(0, end=';')
except:
    print('ERR', end=';')
else:
    print(4, end=';')
finally:
    print(5)
code

6. Що буде виведено за виконання?
code
def f(n):
    if n>9:
        f(n//100)
    else:
        print(n%10,end='')

f(987654)
code

7. Що буде виведено за виконання?
code
class A:
    a = 1

    def f1(self): print(2, end=';')
    def _f2(self): print(3, end=';')

    def main(self):
        try: print(self.a, end=';')
        except: print('ERR', end=';')

        try: self.f1()
        except: print('ERR', end=';')

        try: self._f2()
        except: print('ERR', end=';')


class B(A):
    a = 4

    def f1(self): print(5, end=';')
    def _f2(self): print(6, end=';')

try: B().main()
except: print('ERR', end=';')
code

8. Що буде виведено за виконання?
code
def f(arg, n):
    arg[20] = 8
    n -= -3

try:
    n = 5
    s = {72:2, 66:0, 55:3, 55:4}
    f(s, n)
    print(n, end=';')
    for x in s.keys():
        print(x, sep=';', end=';')
except:
    print('ERR')
code

9. Що буде виведено за виконання?
code
class A:
    c = 1
    a = 2

    def __init__(self):
        c = 3

class B(A):
    c = 4

    def __init__(self):
        super().__init__()
        self.b = 7

obj = B()
obj.c = obj.a + 10
B.c = 13

try: print(obj.c, end= ';')
except: print('ERR', end=';')

try: print(A.c, end= ';')
except: print('ERR', end=';')

try: print(B.c, end= ';')
except: print('ERR', end=';')
code

10. 
Для графа на рисунку з вершини 2 запущено алгоритм пошуку в глибину, що будує кістякове дерево. Списки суміжності вершин переглядаються алгоритмом за зростанням міток вершин.\par
\begin{center}
	\adjustimage{max size={0.9\linewidth}{0.9\paperheight}}
	{../10/Traversals/23.png}
\end{center}
\par Які з наведених ребер входять до складу отриманого кістякового дерева?\par
1) (1, 2)\par
2) (2, 3)\par
3) (3, 6)\par
4) (0, 1)\par
5) (2, 5)\par
6) (3, 5)\par
{\vskip 15pt} У формі вибрати номери відповідних ребер.\par
%answer:3 4 6
\end {large}}
%end
{\smallskip \begin {small} Затверджено на засіданні кафедри теоретичної кібернетики, протокол \No 12 від   19.05.2020 р.\par\smallskip\smallskip
{\parbox {6.5 cm}{Зав. кафедри \dotfill Ю.В.~Крак}}\hspace {0.7cm}{\hbox to 6.5 cm{ Екзаменатор \dotfill Т.О.~Карнаух}} \end{small}}
%begin
{{\begin {small}\par
{ \center  Київський національний університет імені Тараса Шевченка \par}
{\parbox {5 cm}{Освітня програма \hfill}{\parbox {7 cm}{Прикладна математика, Системний аналіз \hfill}}\parbox {6 cm}{\hfill Семестр 2}}\par
{\parbox {5 cm} {Навчальний предмет \hfill}\parbox {7 cm}{Програмування \hfill}\parbox {6cm}{\hfill}\par}\vspace{-7pt} \end{small}}{\center Екзаменаційний білет \No \textbf 
{ 56 }\par}}
{
\begin {large}
1. Що буде виведено за виконання?
code
try:
    res = not 8<=False or 4.0!=2.0 or 5<3
    if isinstance(res, bool):
        print(f'bool;{res}')
    elif isinstance(res, int):
        print(f'int;{res}')
    elif isinstance(res, float):
        print(f'float;{res:.2f}')
    else:
        print('???')
except:
    print('ERR')
code

2. Що буде виведено за виконання?
code
a = 5
b = 3
c = 0

def g(b):
    a = 4
    b = 3
    c = 5
    return a + b + c

a = 6
b = 3
c = 4

try:
    print(g(b), a, b, c, sep=';')
except:
    print('ERR', a, b, c, sep=';')
code

3. Що буде виведено за виконання?
code
try:
    i = 0
    while i<6:
        i += 2
    print(i, end=';')
except:
    print('ERR')
code

4. Що буде виведено за виконання?
code
try:
    for d in range(-1, -15, -3):
        if d>8:
            continue
        print(d, end=';')
    else:
        print(d,end=';')
    print(d)
except:
    print('ERR')
code

5. Що буде виведено за виконання?\par (Дійсні значення, якщо наявні, в цьому завданні будуть виводитись у форматі з фіксованою крапкою та, можливо, з доволі великою кількістю десяткових розрядів. У відповіді для дійсних значень у ЦЬОМУ завданні достатньо вказати один десятковий розряд після крапки, відкинувши решту розрядів без округлення. Наприклад, замість 13.66666666666667 достатньо вказати 13.6, замість 25.23 достатньо вказати 25.2)
code
try:
    print(3, end=';')
    print(int('d9'), end=';')
    print(2, end=';')
except ZeroDivisionError:
    print(7, end=';')
except KeyboardInterrupt:
    print(6, end=';')
except:
    print('ERR', end=';')
else:
    print(5, end=';')
finally:
    print(3)
code

6. Що буде виведено за виконання?
code
def f(n):
    print(n%10,end='')
    if n>9:
        f(n//100)
 
f(123456)
code

7. Що буде виведено за виконання?
code
class A:
    b = 1

    def __g1(self): print(2, end=';')
    def _g2(self): print(3, end=';')

    def main(self):
        try: print(self.b, end=';')
        except: print('ERR', end=';')

        try: self.__g1()
        except: print('ERR', end=';')

        try: self._g2()
        except: print('ERR', end=';')


class B(A):
    b = 4

    def __g1(self): print(5, end=';')
    def _g2(self): print(6, end=';')

try: B().main()
except: print('ERR', end=';')
code

8. Що буде виведено за виконання?
code
def f(arg, n):
    arg[42] = 8
    n += 2

try:
    n = 4
    d = {65:4, 14:9, 16:1, 16:8}
    f(d, n)
    print(n, end=';')
    for x in d.values():
        print(x, sep=';', end=';')
except:
    print('ERR')
code

9. Що буде виведено за виконання?
code
class A:
    c = 1
    b = 2

    def __init__(self):
        b = 3

class B(A):
    b = 4

    def __init__(self):
        super().__init__()
        self.a = 7

obj = B()
obj.b = obj.c + 10
B.b = 13

try: print(obj.b, end= ';')
except: print('ERR', end=';')

try: print(A.b, end= ';')
except: print('ERR', end=';')

try: print(B.b, end= ';')
except: print('ERR', end=';')
code

10. 
Для графа на рисунку з вершини 6 запущено алгоритм пошуку в глибину, що будує кістякове дерево. Списки суміжності вершин переглядаються алгоритмом за зростанням міток вершин.\par
\begin{center}
	\adjustimage{max size={0.9\linewidth}{0.9\paperheight}}
	{../10/Traversals/164.png}
\end{center}
\par Які з наведених ребер входять до складу отриманого кістякового дерева?\par
1) (4, 5)\par
2) (0, 5)\par
3) (0, 4)\par
4) (1, 6)\par
5) (0, 2)\par
6) (2, 4)\par
{\vskip 15pt} У формі вибрати номери відповідних ребер.\par
%answer:2 4 5 6
\end {large}}
%end
{\smallskip \begin {small} Затверджено на засіданні кафедри теоретичної кібернетики, протокол \No 12 від   19.05.2020 р.\par\smallskip\smallskip
{\parbox {6.5 cm}{Зав. кафедри \dotfill Ю.В.~Крак}}\hspace {0.7cm}{\hbox to 6.5 cm{ Екзаменатор \dotfill Т.О.~Карнаух}} \end{small}}
%begin
{{\begin {small}\par
{ \center  Київський національний університет імені Тараса Шевченка \par}
{\parbox {5 cm}{Освітня програма \hfill}{\parbox {7 cm}{Прикладна математика, Системний аналіз \hfill}}\parbox {6 cm}{\hfill Семестр 2}}\par
{\parbox {5 cm} {Навчальний предмет \hfill}\parbox {7 cm}{Програмування \hfill}\parbox {6cm}{\hfill}\par}\vspace{-7pt} \end{small}}{\center Екзаменаційний білет \No \textbf 
{ 57 }\par}}
{
\begin {large}
1. Що буде виведено за виконання?
code
def f(n):
    print('f', end='')
    return n

try:
    print(f(7) and f(-5))
except:
    print('ERR')
code

2. Що буде виведено за виконання?
code
a = 5
b = 4
c = 2

def f(a):
    a += 4
    b = 3
    c = 1
    return a + b + c

a = 8
b = 9
c = 1

try:
    print(f(a), a, b, c, sep=';')
except:
    print('ERR', a, b, c, sep=';')
code

3. Що буде виведено за виконання?
code
try:
    j = 6
    while j<=9:
        j += 1
    print(j, end=';')
except:
    print('ERR')
code

4. Що буде виведено за виконання?
code
try:
    for c in range(9, -9, -2):
        if c>3:
            break
        print(c, end=';')
        if c>3:
            break
    else:
        print(c,end=';')
    print(c)
except:
    print('ERR')
code

5. Що буде виведено за виконання?\par (Дійсні значення, якщо наявні, в цьому завданні будуть виводитись у форматі з фіксованою крапкою та, можливо, з доволі великою кількістю десяткових розрядів. У відповіді для дійсних значень у ЦЬОМУ завданні достатньо вказати один десятковий розряд після крапки, відкинувши решту розрядів без округлення. Наприклад, замість 13.66666666666667 достатньо вказати 13.6, замість 25.23 достатньо вказати 25.2)
code
try:
    print(1, end=';')
    print(4j==2j, end=';')
    print(2, end=';')
except Exception:
    print(9, end=';')
except TypeError:
    print(7, end=';')
except:
    print('ERR', end=';')
else:
    print(0, end=';')
finally:
    print(6)
code

6. Що буде виведено за виконання?
code
def f(n):
    if n>9:
        f(n//10)
    print(n%10, end='')
 
f(543216)
code

7. Що буде виведено за виконання?
code
class A:
    d = 1

    def _h1(self): print(2, end=';')
    def _h2(self): print(3, end=';')

    def main(self):
        try: print(self.d, end=';')
        except: print('ERR', end=';')

        try: self._h1()
        except: print('ERR', end=';')

        try: self._h2()
        except: print('ERR', end=';')


class B(A):
    d = 4

    def _h1(self): print(5, end=';')
    def _h2(self): print(6, end=';')

try: B().main()
except: print('ERR', end=';')
code

8. Що буде виведено за виконання?
code
def f(arg, n):
    arg[37] = 2
    n *= 3

try:
    n = 7
    t = {16:0, 26:2, 72:2, 16:9}
    f(t, n)
    print(n, end=';')
    for x in t.items():
        print(*x, sep=';', end=';')
except:
    print('ERR')
code

9. Що буде виведено за виконання?
code
class A:
    b = 1
    c = 2

    def __init__(self):
        b = 3

class B(A):
    b = 4

    def __init__(self):
        super().__init__()
        self.a = 7

obj = B()
obj.a = obj.b + 10
B.a = 13

try: print(obj.a, end= ';')
except: print('ERR', end=';')

try: print(A.a, end= ';')
except: print('ERR', end=';')

try: print(B.a, end= ';')
except: print('ERR', end=';')
code

10. 
Для графа на рисунку з вершини 5 запущено алгоритм пошуку в глибину, що будує кістякове дерево. Списки суміжності вершин переглядаються алгоритмом за зростанням міток вершин.\par
\begin{center}
	\adjustimage{max size={0.9\linewidth}{0.9\paperheight}}
	{../10/Traversals/235.png}
\end{center}
\par Які з наведених ребер входять до складу отриманого кістякового дерева?\par
1) (3, 4)\par
2) (1, 2)\par
3) (0, 1)\par
4) (0, 5)\par
5) (1, 4)\par
6) (1, 6)\par
{\vskip 15pt} У формі вибрати номери відповідних ребер.\par
%answer:1 2 3 4
\end {large}}
%end
{\smallskip \begin {small} Затверджено на засіданні кафедри теоретичної кібернетики, протокол \No 12 від   19.05.2020 р.\par\smallskip\smallskip
{\parbox {6.5 cm}{Зав. кафедри \dotfill Ю.В.~Крак}}\hspace {0.7cm}{\hbox to 6.5 cm{ Екзаменатор \dotfill Т.О.~Карнаух}} \end{small}}
%begin
{{\begin {small}\par
{ \center  Київський національний університет імені Тараса Шевченка \par}
{\parbox {5 cm}{Освітня програма \hfill}{\parbox {7 cm}{Прикладна математика, Системний аналіз \hfill}}\parbox {6 cm}{\hfill Семестр 2}}\par
{\parbox {5 cm} {Навчальний предмет \hfill}\parbox {7 cm}{Програмування \hfill}\parbox {6cm}{\hfill}\par}\vspace{-7pt} \end{small}}{\center Екзаменаційний білет \No \textbf 
{ 58 }\par}}
{
\begin {large}
1. Що буде виведено за виконання?
code
try:
    res = 6<True or not 2.0!=7 or 4==7.0
    if isinstance(res, bool):
        print(f'bool;{res}')
    elif isinstance(res, int):
        print(f'int;{res}')
    elif isinstance(res, float):
        print(f'float;{res:.2f}')
    else:
        print('???')
except:
    print('ERR')
code

2. Що буде виведено за виконання?
code
a = 7
b = 0
c = 9

def g(a):
    global c
    a -= 5
    b = 2
    c = 1
    return a + b + c

a = 2
b = 3
c = 1

try:
    print(g(a), a, b, c, sep=';')
except:
    print('ERR', a, b, c, sep=';')
code

3. Що буде виведено за виконання?
code
try:
    t = 8
    while t<=11:
        t += 2
    print(t, end=';')
except:
    print('ERR')
code

4. Що буде виведено за виконання?
code
try:
    for f in range(-5, 14, 2):
        if f<=4:
            continue
        print(f, end=';')
    else:
        print(f,end=';')
    print(f)
except:
    print('ERR')
code

5. Що буде виведено за виконання?\par (Дійсні значення, якщо наявні, в цьому завданні будуть виводитись у форматі з фіксованою крапкою та, можливо, з доволі великою кількістю десяткових розрядів. У відповіді для дійсних значень у ЦЬОМУ завданні достатньо вказати один десятковий розряд після крапки, відкинувши решту розрядів без округлення. Наприклад, замість 13.66666666666667 достатньо вказати 13.6, замість 25.23 достатньо вказати 25.2)
code
try:
    print(8, end=';')
    print(int('8'), end=';')
    print(4, end=';')
except ValueError:
    print(1, end=';')
except Exception:
    print(2, end=';')
except:
    print('ERR', end=';')
else:
    print(9, end=';')
finally:
    print(5)
code

6. Що буде виведено за виконання?
code
def f(n):
    if n>9:
        f(n//100)
    else:
        print(n%10,end='')

f(987654)
code

7. Що буде виведено за виконання?
code
class A:
    b = 1

    def h1(self): print(2, end=';')
    def _h2(self): print(3, end=';')

    def main(self):
        try: print(self.b, end=';')
        except: print('ERR', end=';')

        try: self.h1()
        except: print('ERR', end=';')

        try: self._h2()
        except: print('ERR', end=';')


class B(A):
    b = 4

    def h1(self): print(5, end=';')
    def _h2(self): print(6, end=';')

try: B().main()
except: print('ERR', end=';')
code

8. Що буде виведено за виконання?
code
def f(arg, n):
    arg[40] = 2
    n = 6
    arg = set(arg.keys())

try:
    n = 2
    s = {50:2, 18:5, 40:3, 40:0}
    f(s, n)
    print(n, end=';')
    for x in s.items():
        print(*x, sep=';', end=';')
except:
    print('ERR')
code

9. Що буде виведено за виконання?
code
class A:
    a = 1
    c = 2

    def __init__(self):
        c = 3

class B(A):
    c = 4

    def __init__(self):
        super().__init__()
        self.b = 7

obj = B()
obj.a = obj.b + 10
B.a = 13

try: print(obj.a, end= ';')
except: print('ERR', end=';')

try: print(A.a, end= ';')
except: print('ERR', end=';')

try: print(B.a, end= ';')
except: print('ERR', end=';')
code

10. 
Для графа на рисунку з вершини 3 запущено алгоритм пошуку в глибину, що будує кістякове дерево. Списки суміжності вершин переглядаються алгоритмом за зростанням міток вершин.\par
\begin{center}
	\adjustimage{max size={0.9\linewidth}{0.9\paperheight}}
	{../10/Traversals/145.png}
\end{center}
\par Які з наведених ребер входять до складу отриманого кістякового дерева?\par
1) (0, 3)\par
2) (0, 1)\par
3) (1, 2)\par
4) (2, 3)\par
5) (4, 6)\par
6) (1, 4)\par
{\vskip 15pt} У формі вибрати номери відповідних ребер.\par
%answer:1 2 3 5 6
\end {large}}
%end
{\smallskip \begin {small} Затверджено на засіданні кафедри теоретичної кібернетики, протокол \No 12 від   19.05.2020 р.\par\smallskip\smallskip
{\parbox {6.5 cm}{Зав. кафедри \dotfill Ю.В.~Крак}}\hspace {0.7cm}{\hbox to 6.5 cm{ Екзаменатор \dotfill Т.О.~Карнаух}} \end{small}}
%begin
{{\begin {small}\par
{ \center  Київський національний університет імені Тараса Шевченка \par}
{\parbox {5 cm}{Освітня програма \hfill}{\parbox {7 cm}{Прикладна математика, Системний аналіз \hfill}}\parbox {6 cm}{\hfill Семестр 2}}\par
{\parbox {5 cm} {Навчальний предмет \hfill}\parbox {7 cm}{Програмування \hfill}\parbox {6cm}{\hfill}\par}\vspace{-7pt} \end{small}}{\center Екзаменаційний білет \No \textbf 
{ 59 }\par}}
{
\begin {large}
1. Що буде виведено за виконання?
code
try:
    res = 8//4/5*7
    if isinstance(res, bool):
        print(f'bool;{res}')
    elif isinstance(res, int):
        print(f'int;{res}')
    elif isinstance(res, float):
        print(f'float;{res:.2f}')
    else:
        print('???')
except:
    print('ERR')
code

2. Що буде виведено за виконання?
code
a = 5
b = 6
c = 0

def h(b):
    a = 3
    b *= 4
    c = 5
    return a + b + c

a = 7
b = 4
c = 8

try:
    print(h(b), a, b, c, sep=';')
except:
    print('ERR', a, b, c, sep=';')
code

3. Що буде виведено за виконання?
code
try:
    i = 6
    while i>7:
        i += -2
    print(i, end=';')
except:
    print('ERR')
code

4. Що буде виведено за виконання?
code
try:
    for a in range(13, -14, 3):
        if a>=5:
            break
        print(a, end=';');a = 10
    else:
        print(a,end=';')
    print(a)
except:
    print('ERR')
code

5. Що буде виведено за виконання?\par (Дійсні значення, якщо наявні, в цьому завданні будуть виводитись у форматі з фіксованою крапкою та, можливо, з доволі великою кількістю десяткових розрядів. У відповіді для дійсних значень у ЦЬОМУ завданні достатньо вказати один десятковий розряд після крапки, відкинувши решту розрядів без округлення. Наприклад, замість 13.66666666666667 достатньо вказати 13.6, замість 25.23 достатньо вказати 25.2)
code
try:
    print(3, end=';')
    print(6%0, end=';')
    print(0, end=';')
except ValueError:
    print(7, end=';')
except KeyboardInterrupt:
    print(8, end=';')
except:
    print('ERR', end=';')
else:
    print(5, end=';')
finally:
    print(6)
code

6. Що буде виведено за виконання?
code
def f(s):
    for i in range(len(s)):
        s[i] += 1
        return
 
s = [1, 2, 3]
f(s)
print(s)
code

7. Що буде виведено за виконання?
code
class A:
    a = 1

    def _g1(self): print(2, end=';')
    def __g2(self): print(3, end=';')

    def main(self):
        try: print(self.a, end=';')
        except: print('ERR', end=';')

        try: self._g1()
        except: print('ERR', end=';')

        try: self.__g2()
        except: print('ERR', end=';')


class B(A):
    a = 4

    def _g1(self): print(5, end=';')
    def __g2(self): print(6, end=';')

try: B().main()
except: print('ERR', end=';')
code

8. Що буде виведено за виконання?
code
def f(arg, n):
    n = 2
    tmp = arg
    tmp[:1] = [13]
    return len(tmp)>3

try:
    d = [10,11,12,13,14]
    n = 4
    f(d, n)
    print(n, end=';')
    for el in d:
        print(el, end=';')
except:
    print('ERR')
code

9. Що буде виведено за виконання?
code
class A:
    b = 1
    a = 2

    def __init__(self):
        b = 3

class B(A):
    b = 4

    def __init__(self):
        super().__init__()
        self.c = 7

obj = B()
obj.c = obj.a + 10
B.c = 13

try: print(obj.c, end= ';')
except: print('ERR', end=';')

try: print(A.c, end= ';')
except: print('ERR', end=';')

try: print(B.c, end= ';')
except: print('ERR', end=';')
code

10. 
Для графа на рисунку з вершини 0 запущено алгоритм пошуку в глибину, що будує кістякове дерево. Списки суміжності вершин переглядаються алгоритмом за зростанням міток вершин.\par
\begin{center}
	\adjustimage{max size={0.9\linewidth}{0.9\paperheight}}
	{../10/Traversals/66.png}
\end{center}
\par Які з наведених ребер входять до складу отриманого кістякового дерева?\par
1) (0, 2)\par
2) (3, 4)\par
3) (3, 6)\par
4) (0, 3)\par
5) (4, 5)\par
6) (4, 6)\par
{\vskip 15pt} У формі вибрати номери відповідних ребер.\par
%answer:1 3 5
\end {large}}
%end
{\smallskip \begin {small} Затверджено на засіданні кафедри теоретичної кібернетики, протокол \No 12 від   19.05.2020 р.\par\smallskip\smallskip
{\parbox {6.5 cm}{Зав. кафедри \dotfill Ю.В.~Крак}}\hspace {0.7cm}{\hbox to 6.5 cm{ Екзаменатор \dotfill Т.О.~Карнаух}} \end{small}}
%begin
{{\begin {small}\par
{ \center  Київський національний університет імені Тараса Шевченка \par}
{\parbox {5 cm}{Освітня програма \hfill}{\parbox {7 cm}{Прикладна математика, Системний аналіз \hfill}}\parbox {6 cm}{\hfill Семестр 2}}\par
{\parbox {5 cm} {Навчальний предмет \hfill}\parbox {7 cm}{Програмування \hfill}\parbox {6cm}{\hfill}\par}\vspace{-7pt} \end{small}}{\center Екзаменаційний білет \No \textbf 
{ 60 }\par}}
{
\begin {large}
1. Що буде виведено за виконання?
code
try:
    res = 4**1.0*-1
    if isinstance(res, bool):
        print(f'bool;{res}')
    elif isinstance(res, int):
        print(f'int;{res}')
    elif isinstance(res, float):
        print(f'float;{res:.2f}')
    else:
        print('???')
except:
    print('ERR')
code

2. Що буде виведено за виконання?
code
a = 7
b = 8
c = 5

def f(a):
    global c
    a = 2
    b = 4
    c = 1
    return a + b + c

a = 1
b = 2
c = 9

try:
    print(f(a), a, b, c, sep=';')
except:
    print('ERR', a, b, c, sep=';')
code

3. Що буде виведено за виконання?
code
try:
    k = 7
    while k<5:
        k += 1
    print(k, end=';')
except:
    print('ERR')
code

4. Що буде виведено за виконання?
code
try:
    for e in range(6, 12, -1):
        if e<7:
            continue
        print(e, end=';')
    else:
        print('end',end=';')
    print(e)
except:
    print('ERR')
code

5. Що буде виведено за виконання?\par (Дійсні значення, якщо наявні, в цьому завданні будуть виводитись у форматі з фіксованою крапкою та, можливо, з доволі великою кількістю десяткових розрядів. У відповіді для дійсних значень у ЦЬОМУ завданні достатньо вказати один десятковий розряд після крапки, відкинувши решту розрядів без округлення. Наприклад, замість 13.66666666666667 достатньо вказати 13.6, замість 25.23 достатньо вказати 25.2)
code
try:
    print(3, end=';')
    print(1//1, end=';')
    print(4, end=';')
except ZeroDivisionError:
    print(9, end=';')
except TypeError:
    print(0, end=';')
except:
    print('ERR', end=';')
else:
    print(2, end=';')
finally:
    print(7)
code

6. Що буде виведено за виконання?
code
def f(n):
    if n>9:
        f(n//10)
    print(n%10, end='')
 
f(123456)
code

7. Що буде виведено за виконання?
code
class A:
    c = 1

    def _f1(self): print(2, end=';')
    def f2(self): print(3, end=';')

    def main(self):
        try: print(self.c, end=';')
        except: print('ERR', end=';')

        try: self._f1()
        except: print('ERR', end=';')

        try: self.f2()
        except: print('ERR', end=';')


class B(A):
    c = 4

    def _f1(self): print(5, end=';')
    def f2(self): print(6, end=';')

try: B().main()
except: print('ERR', end=';')
code

8. Що буде виведено за виконання?
code
def f(arg, n):
    arg[33] = 1
    n -= -3

try:
    n = 7
    s = {33:2, 50:4, 70:9, 89:5}
    f(s, n)
    print(n, end=';')
    for x in s :
        print(x, sep=';', end=';')
except:
    print('ERR')
code

9. Що буде виведено за виконання?
code
class A:
    a = 1
    b = 2

    def __init__(self):
        b = 3

class B(A):
    b = 4

    def __init__(self):
        super().__init__()
        self.c = 7

obj = B()
obj.b = obj.c + 10
B.b = 13

try: print(obj.b, end= ';')
except: print('ERR', end=';')

try: print(A.b, end= ';')
except: print('ERR', end=';')

try: print(B.b, end= ';')
except: print('ERR', end=';')
code

10. 
Для графа на рисунку з вершини 3 запущено алгоритм пошуку в глибину, що будує кістякове дерево. Списки суміжності вершин переглядаються алгоритмом за зростанням міток вершин.\par
\begin{center}
	\adjustimage{max size={0.9\linewidth}{0.9\paperheight}}
	{../10/Traversals/157.png}
\end{center}
\par Які з наведених ребер входять до складу отриманого кістякового дерева?\par
1) (0, 2)\par
2) (2, 4)\par
3) (1, 4)\par
4) (0, 5)\par
5) (3, 5)\par
6) (3, 6)\par
{\vskip 15pt} У формі вибрати номери відповідних ребер.\par
%answer:1 4
\end {large}}
%end
{\smallskip \begin {small} Затверджено на засіданні кафедри теоретичної кібернетики, протокол \No 12 від   19.05.2020 р.\par\smallskip\smallskip
{\parbox {6.5 cm}{Зав. кафедри \dotfill Ю.В.~Крак}}\hspace {0.7cm}{\hbox to 6.5 cm{ Екзаменатор \dotfill Т.О.~Карнаух}} \end{small}}
%begin
{{\begin {small}\par
{ \center  Київський національний університет імені Тараса Шевченка \par}
{\parbox {5 cm}{Освітня програма \hfill}{\parbox {7 cm}{Прикладна математика, Системний аналіз \hfill}}\parbox {6 cm}{\hfill Семестр 2}}\par
{\parbox {5 cm} {Навчальний предмет \hfill}\parbox {7 cm}{Програмування \hfill}\parbox {6cm}{\hfill}\par}\vspace{-7pt} \end{small}}{\center Екзаменаційний білет \No \textbf 
{ 61 }\par}}
{
\begin {large}
1. Що буде виведено за виконання?
code
def f(n):
    print('f', end='')
    return n

try:
    print(f(4) or f(9))
except:
    print('ERR')
code

2. Що буде виведено за виконання?
code
a = 7
b = 2
c = 4

def g(a):
    a -= 2
    b = 5
    c = 3
    return a + b + c

a = 8
b = 5
c = 1

try:
    print(g(b), a, b, c, sep=';')
except:
    print('ERR', a, b, c, sep=';')
code

3. Що буде виведено за виконання?
code
try:
    n = 4
    while n<=8:
        n += 1
    print(n, end=';')
except:
    print('ERR')
code

4. Що буде виведено за виконання?
code
try:
    for b in range(-11, 7, 2):
        if b>7:
            continue
        print(b, end=';')
    else:
        print(b,end=';')
    print(b)
except:
    print('ERR')
code

5. Що буде виведено за виконання?\par (Дійсні значення, якщо наявні, в цьому завданні будуть виводитись у форматі з фіксованою крапкою та, можливо, з доволі великою кількістю десяткових розрядів. У відповіді для дійсних значень у ЦЬОМУ завданні достатньо вказати один десятковий розряд після крапки, відкинувши решту розрядів без округлення. Наприклад, замість 13.66666666666667 достатньо вказати 13.6, замість 25.23 достатньо вказати 25.2)
code
try:
    print(4, end=';')
    print(int('8'), end=';')
    print(6, end=';')
except ValueError:
    print(2, end=';')
except ZeroDivisionError:
    print(5, end=';')
except:
    print('ERR', end=';')
else:
    print(0, end=';')
finally:
    print(9)
code

6. Що буде виведено за виконання?
code
def f(n):
    print(n%10,end='')
    if n>9:
        f(n//100)
 
f(123456)
code

7. Що буде виведено за виконання?
code
class A:
    a = 1

    def __f1(self): print(2, end=';')
    def _f2(self): print(3, end=';')

    def main(self):
        try: print(self.a, end=';')
        except: print('ERR', end=';')

        try: self.__f1()
        except: print('ERR', end=';')

        try: self._f2()
        except: print('ERR', end=';')


class B(A):
    a = 4

    def __f1(self): print(5, end=';')
    def _f2(self): print(6, end=';')

try: B().main()
except: print('ERR', end=';')
code

8. Що буде виведено за виконання?
code
def f(arg, n):
    arg[53] = 0
    n *= 3

try:
    n = 5
    s = {49:6, 53:9, 49:9}
    f(s, n)
    print(n, end=';')
    for x in s.keys():
        print(x, sep=';', end=';')
except:
    print('ERR')
code

9. Що буде виведено за виконання?
code
class A:
    b = 1
    a = 2

    def __init__(self):
        b = 3

class B(A):
    a = 4

    def __init__(self):
        super().__init__()
        self.c = 7

obj = B()
obj.c = obj.a + 10
B.c = 13

try: print(obj.c, end= ';')
except: print('ERR', end=';')

try: print(A.c, end= ';')
except: print('ERR', end=';')

try: print(B.c, end= ';')
except: print('ERR', end=';')
code

10. 
Для графа на рисунку з вершини 6 запущено алгоритм пошуку в глибину, що будує кістякове дерево. Списки суміжності вершин переглядаються алгоритмом за зростанням міток вершин.\par
\begin{center}
	\adjustimage{max size={0.9\linewidth}{0.9\paperheight}}
	{../10/Traversals/298.png}
\end{center}
\par Які з наведених ребер входять до складу отриманого кістякового дерева?\par
1) (1, 5)\par
2) (3, 5)\par
3) (4, 6)\par
4) (2, 3)\par
5) (0, 5)\par
6) (1, 3)\par
{\vskip 15pt} У формі вибрати номери відповідних ребер.\par
%answer:1 6
\end {large}}
%end
{\smallskip \begin {small} Затверджено на засіданні кафедри теоретичної кібернетики, протокол \No 12 від   19.05.2020 р.\par\smallskip\smallskip
{\parbox {6.5 cm}{Зав. кафедри \dotfill Ю.В.~Крак}}\hspace {0.7cm}{\hbox to 6.5 cm{ Екзаменатор \dotfill Т.О.~Карнаух}} \end{small}}
%begin
{{\begin {small}\par
{ \center  Київський національний університет імені Тараса Шевченка \par}
{\parbox {5 cm}{Освітня програма \hfill}{\parbox {7 cm}{Прикладна математика, Системний аналіз \hfill}}\parbox {6 cm}{\hfill Семестр 2}}\par
{\parbox {5 cm} {Навчальний предмет \hfill}\parbox {7 cm}{Програмування \hfill}\parbox {6cm}{\hfill}\par}\vspace{-7pt} \end{small}}{\center Екзаменаційний білет \No \textbf 
{ 62 }\par}}
{
\begin {large}
1. Що буде виведено за виконання?
code
try:
    res = "True"<=2j
    if isinstance(res, bool):
        print(f'bool;{res}')
    elif isinstance(res, int):
        print(f'int;{res}')
    elif isinstance(res, float):
        print(f'float;{res:.2f}')
    else:
        print('???')
except:
    print('ERR')
code

2. Що буде виведено за виконання?
code
a = 1
b = 6
c = 3

def f(a):
    global c
    a = 2
    b = 1
    c = 4
    return a + b + c

a = 8
b = 2
c = 0

try:
    print(f(a), a, b, c, sep=';')
except:
    print('ERR', a, b, c, sep=';')
code

3. Що буде виведено за виконання?
code
try:
    t = 10
    while t>4:
        t += -3
    print(t, end=';')
except:
    print('ERR')
code

4. Що буде виведено за виконання?
code
try:
    for f in range(-14, -2, -1):
        if f<3:
            break
        print(f, end=';')
    else:
        print('end',end=';')
    print(f)
except:
    print('ERR')
code

5. Що буде виведено за виконання?\par (Дійсні значення, якщо наявні, в цьому завданні будуть виводитись у форматі з фіксованою крапкою та, можливо, з доволі великою кількістю десяткових розрядів. У відповіді для дійсних значень у ЦЬОМУ завданні достатньо вказати один десятковий розряд після крапки, відкинувши решту розрядів без округлення. Наприклад, замість 13.66666666666667 достатньо вказати 13.6, замість 25.23 достатньо вказати 25.2)
code
try:
    print(3, end=';')
    print(7//1, end=';')
    print(1, end=';')
except TypeError:
    print(4, end=';')
except KeyboardInterrupt:
    print(1, end=';')
except:
    print('ERR', end=';')
else:
    print(8, end=';')
finally:
    print(3)
code

6. Що буде виведено за виконання?
code
def f(s1):
    s = s1[:]
    for i in range(len(s)):
        s[i] += 1
 
s = [1, 2, 3]
f(s)
print(s)
code

7. Що буде виведено за виконання?
code
class A:
    d = 1

    def _h1(self): print(2, end=';')
    def h2(self): print(3, end=';')

    def main(self):
        try: print(self.d, end=';')
        except: print('ERR', end=';')

        try: self._h1()
        except: print('ERR', end=';')

        try: self.h2()
        except: print('ERR', end=';')


class B(A):
    d = 4

    def _h1(self): print(5, end=';')
    def h2(self): print(6, end=';')

try: B().main()
except: print('ERR', end=';')
code

8. Що буде виведено за виконання?
code
def f(arg, n):
    n = 0
    tmp = arg
    tmp[-2:-7] = 'xy'
    return len(tmp)>3

try:
    b = ['0','1','2','3','4']
    n = 3
    f(b, n)
    print(n, end=';')
    for el in b:
        print(el, end=';')
except:
    print('ERR')
code

9. Що буде виведено за виконання?
code
class A:
    c = 1
    b = 2

    def __init__(self):
        b = 3

class B(A):
    c = 4

    def __init__(self):
        super().__init__()
        self.a = 7

obj = B()
obj.b = obj.b + 10
B.b = 13

try: print(obj.b, end= ';')
except: print('ERR', end=';')

try: print(A.b, end= ';')
except: print('ERR', end=';')

try: print(B.b, end= ';')
except: print('ERR', end=';')
code

10. 
Для графа на рисунку з вершини 2 запущено алгоритм пошуку в глибину, що будує кістякове дерево. Списки суміжності вершин переглядаються алгоритмом за зростанням міток вершин.\par
\begin{center}
	\adjustimage{max size={0.9\linewidth}{0.9\paperheight}}
	{../10/Traversals/189.png}
\end{center}
\par Які з наведених ребер входять до складу отриманого кістякового дерева?\par
1) (0, 2)\par
2) (0, 5)\par
3) (2, 3)\par
4) (3, 5)\par
5) (1, 5)\par
6) (1, 4)\par
{\vskip 15pt} У формі вибрати номери відповідних ребер.\par
%answer:1 4 5 6
\end {large}}
%end
{\smallskip \begin {small} Затверджено на засіданні кафедри теоретичної кібернетики, протокол \No 12 від   19.05.2020 р.\par\smallskip\smallskip
{\parbox {6.5 cm}{Зав. кафедри \dotfill Ю.В.~Крак}}\hspace {0.7cm}{\hbox to 6.5 cm{ Екзаменатор \dotfill Т.О.~Карнаух}} \end{small}}
%begin
{{\begin {small}\par
{ \center  Київський національний університет імені Тараса Шевченка \par}
{\parbox {5 cm}{Освітня програма \hfill}{\parbox {7 cm}{Прикладна математика, Системний аналіз \hfill}}\parbox {6 cm}{\hfill Семестр 2}}\par
{\parbox {5 cm} {Навчальний предмет \hfill}\parbox {7 cm}{Програмування \hfill}\parbox {6cm}{\hfill}\par}\vspace{-7pt} \end{small}}{\center Екзаменаційний білет \No \textbf 
{ 63 }\par}}
{
\begin {large}
1. Що буде виведено за виконання?
code
def f(n):
    print('f', end='')
    return n>=-3

try:
    print(f(-2) and f(-6))
except:
    print('ERR')
code

2. Що буде виведено за виконання?
code
a = 0
b = 9
c = 6

def h(b):
    global c
    a += 1
    b = 4
    c = 5
    return a + b + c

a = 3
b = 2
c = 0

try:
    print(h(a), a, b, c, sep=';')
except:
    print('ERR', a, b, c, sep=';')
code

3. Що буде виведено за виконання?
code
try:
    j = 0
    while j<7:
        j += 2
    print(j, end=';')
except:
    print('ERR')
code

4. Що буде виведено за виконання?
code
try:
    for c in range(-5, 14, -2):
        if c>=5:
            continue
        print(c, end=';');c = -4
    else:
        print(c,end=';')
    print(c)
except:
    print('ERR')
code

5. Що буде виведено за виконання?\par (Дійсні значення, якщо наявні, в цьому завданні будуть виводитись у форматі з фіксованою крапкою та, можливо, з доволі великою кількістю десяткових розрядів. У відповіді для дійсних значень у ЦЬОМУ завданні достатньо вказати один десятковий розряд після крапки, відкинувши решту розрядів без округлення. Наприклад, замість 13.66666666666667 достатньо вказати 13.6, замість 25.23 достатньо вказати 25.2)
code
try:
    print(0, end=';')
    print(9j<6j, end=';')
    print(7, end=';')
except Exception:
    print(5, end=';')
except TypeError:
    print(2, end=';')
except:
    print('ERR', end=';')
else:
    print(6, end=';')
finally:
    print(3)
code

6. Що буде виведено за виконання?
code
def f(n):
    if n>9:
        f(n//100)
    print(n%10,end='')
 
f(987654)
code

7. Що буде виведено за виконання?
code
class A:
    c = 1

    def g1(self): print(2, end=';')
    def _g2(self): print(3, end=';')

    def main(self):
        try: print(self.c, end=';')
        except: print('ERR', end=';')

        try: self.g1()
        except: print('ERR', end=';')

        try: self._g2()
        except: print('ERR', end=';')


class B(A):
    c = 4

    def g1(self): print(5, end=';')
    def _g2(self): print(6, end=';')

try: B().main()
except: print('ERR', end=';')
code

8. Що буде виведено за виконання?
code
def f(arg, n):
    arg[30] = 7
    n = 9
    arg = list(arg.items())

try:
    n = 0
    t = {64:4, 13:2, 83:3, 83:8}
    f(t, n)
    print(n, end=';')
    for x in t.values():
        print(x, sep=';', end=';')
except:
    print('ERR')
code

9. Що буде виведено за виконання?
code
class A:
    a = 1
    b = 2

    def __init__(self):
        b = 3

class B(A):
    b = 4

    def __init__(self):
        super().__init__()
        self.c = 7

obj = B()
obj.c = obj.a + 10
B.c = 13

try: print(obj.c, end= ';')
except: print('ERR', end=';')

try: print(A.c, end= ';')
except: print('ERR', end=';')

try: print(B.c, end= ';')
except: print('ERR', end=';')
code

10. 
Для графа на рисунку з вершини 5 запущено алгоритм пошуку в глибину, що будує кістякове дерево. Списки суміжності вершин переглядаються алгоритмом за зростанням міток вершин.\par
\begin{center}
	\adjustimage{max size={0.9\linewidth}{0.9\paperheight}}
	{../10/Traversals/230.png}
\end{center}
\par Які з наведених ребер входять до складу отриманого кістякового дерева?\par
1) (0, 2)\par
2) (1, 6)\par
3) (0, 4)\par
4) (2, 5)\par
5) (5, 6)\par
6) (0, 1)\par
{\vskip 15pt} У формі вибрати номери відповідних ребер.\par
%answer:1 2 3 4 6
\end {large}}
%end
{\smallskip \begin {small} Затверджено на засіданні кафедри теоретичної кібернетики, протокол \No 12 від   19.05.2020 р.\par\smallskip\smallskip
{\parbox {6.5 cm}{Зав. кафедри \dotfill Ю.В.~Крак}}\hspace {0.7cm}{\hbox to 6.5 cm{ Екзаменатор \dotfill Т.О.~Карнаух}} \end{small}}
%begin
{{\begin {small}\par
{ \center  Київський національний університет імені Тараса Шевченка \par}
{\parbox {5 cm}{Освітня програма \hfill}{\parbox {7 cm}{Прикладна математика, Системний аналіз \hfill}}\parbox {6 cm}{\hfill Семестр 2}}\par
{\parbox {5 cm} {Навчальний предмет \hfill}\parbox {7 cm}{Програмування \hfill}\parbox {6cm}{\hfill}\par}\vspace{-7pt} \end{small}}{\center Екзаменаційний білет \No \textbf 
{ 64 }\par}}
{
\begin {large}
1. Що буде виведено за виконання?
code
try:
    res = 4**-3-1
    if isinstance(res, bool):
        print(f'bool;{res}')
    elif isinstance(res, int):
        print(f'int;{res}')
    elif isinstance(res, float):
        print(f'float;{res:.2f}')
    else:
        print('???')
except:
    print('ERR')
code

2. Що буде виведено за виконання?
code
a = 7
b = 4
c = 9

def g(b):
    a = 5
    b = 3
    c = 1
    return a + b + c

a = 5
b = 9
c = 6

try:
    print(g(a), a, b, c, sep=';')
except:
    print('ERR', a, b, c, sep=';')
code

3. Що буде виведено за виконання?
code
try:
    k = 5
    while k<=13:
        k += 1
    print(k, end=';')
except:
    print('ERR')
code

4. Що буде виведено за виконання?
code
try:
    for e in range(6, -4, -3):
        if e<=8:
            continue
        print(e, end=';');e = 5
    else:
        print(e,end=';')
    print(e)
except:
    print('ERR')
code

5. Що буде виведено за виконання?\par (Дійсні значення, якщо наявні, в цьому завданні будуть виводитись у форматі з фіксованою крапкою та, можливо, з доволі великою кількістю десяткових розрядів. У відповіді для дійсних значень у ЦЬОМУ завданні достатньо вказати один десятковий розряд після крапки, відкинувши решту розрядів без округлення. Наприклад, замість 13.66666666666667 достатньо вказати 13.6, замість 25.23 достатньо вказати 25.2)
code
try:
    print(0, end=';')
    print(9/0.0, end=';')
    print(7, end=';')
except ZeroDivisionError:
    print(8, end=';')
except KeyboardInterrupt:
    print(2, end=';')
except:
    print('ERR', end=';')
else:
    print(4, end=';')
finally:
    print(1)
code

6. Що буде виведено за виконання?
code
def f(s):
    for i in range(len(s)):
        s[i] += 1
    return
 
s = [1, 2, 3]
f(s)
print(s)
code

7. Що буде виведено за виконання?
code
class A:
    b = 1

    def _h1(self): print(2, end=';')
    def _h2(self): print(3, end=';')

    def main(self):
        try: print(self.b, end=';')
        except: print('ERR', end=';')

        try: self._h1()
        except: print('ERR', end=';')

        try: self._h2()
        except: print('ERR', end=';')


class B(A):
    b = 4

    def _h1(self): print(5, end=';')
    def _h2(self): print(6, end=';')

try: B().main()
except: print('ERR', end=';')
code

8. Що буде виведено за виконання?
code
def f(arg, n):
    arg[47] = 4
    n = 8
    arg = tuple(arg.values())

try:
    n = 1
    d = {41:3, 72:6, 72:6}
    f(d, n)
    print(n, end=';')
    for x in d.items():
        print(*x, sep=';', end=';')
except:
    print('ERR')
code

9. Що буде виведено за виконання?
code
class A:
    a = 1
    c = 2

    def __init__(self):
        a = 3

class B(A):
    a = 4

    def __init__(self):
        super().__init__()
        self.b = 7

obj = B()
obj.a = obj.c + 10
B.a = 13

try: print(obj.a, end= ';')
except: print('ERR', end=';')

try: print(A.a, end= ';')
except: print('ERR', end=';')

try: print(B.a, end= ';')
except: print('ERR', end=';')
code

10. 
Для графа на рисунку з вершини 3 запущено алгоритм пошуку в глибину, що будує кістякове дерево. Списки суміжності вершин переглядаються алгоритмом за зростанням міток вершин.\par
\begin{center}
	\adjustimage{max size={0.9\linewidth}{0.9\paperheight}}
	{../10/Traversals/7.png}
\end{center}
\par Які з наведених ребер входять до складу отриманого кістякового дерева?\par
1) (1, 2)\par
2) (1, 5)\par
3) (5, 6)\par
4) (0, 1)\par
5) (0, 4)\par
6) (2, 6)\par
{\vskip 15pt} У формі вибрати номери відповідних ребер.\par
%answer:1 3 4 5
\end {large}}
%end
{\smallskip \begin {small} Затверджено на засіданні кафедри теоретичної кібернетики, протокол \No 12 від   19.05.2020 р.\par\smallskip\smallskip
{\parbox {6.5 cm}{Зав. кафедри \dotfill Ю.В.~Крак}}\hspace {0.7cm}{\hbox to 6.5 cm{ Екзаменатор \dotfill Т.О.~Карнаух}} \end{small}}
%begin
{{\begin {small}\par
{ \center  Київський національний університет імені Тараса Шевченка \par}
{\parbox {5 cm}{Освітня програма \hfill}{\parbox {7 cm}{Прикладна математика, Системний аналіз \hfill}}\parbox {6 cm}{\hfill Семестр 2}}\par
{\parbox {5 cm} {Навчальний предмет \hfill}\parbox {7 cm}{Програмування \hfill}\parbox {6cm}{\hfill}\par}\vspace{-7pt} \end{small}}{\center Екзаменаційний білет \No \textbf 
{ 65 }\par}}
{
\begin {large}
1. Що буде виведено за виконання?
code
try:
    res = not False<=9 and 6>8.0 and 3.0>=3
    if isinstance(res, bool):
        print(f'bool;{res}')
    elif isinstance(res, int):
        print(f'int;{res}')
    elif isinstance(res, float):
        print(f'float;{res:.2f}')
    else:
        print('???')
except:
    print('ERR')
code

2. Що буде виведено за виконання?
code
a = 7
b = 2
c = 0

def h(a):
    global c
    a = 4
    b = 5
    c = 2
    return a + b + c

a = 5
b = 8
c = 4

try:
    print(h(b), a, b, c, sep=';')
except:
    print('ERR', a, b, c, sep=';')
code

3. Що буде виведено за виконання?
code
try:
    i = 7
    while i>=0:
        i += -2
    print(i, end=';')
except:
    print('ERR')
code

4. Що буде виведено за виконання?
code
try:
    for a in range(-14, -2, 3):
        if a<=6:
            continue
        print(a, end=';')
    else:
        print(a,end=';')
    print(a)
except:
    print('ERR')
code

5. Що буде виведено за виконання?\par (Дійсні значення, якщо наявні, в цьому завданні будуть виводитись у форматі з фіксованою крапкою та, можливо, з доволі великою кількістю десяткових розрядів. У відповіді для дійсних значень у ЦЬОМУ завданні достатньо вказати один десятковий розряд після крапки, відкинувши решту розрядів без округлення. Наприклад, замість 13.66666666666667 достатньо вказати 13.6, замість 25.23 достатньо вказати 25.2)
code
try:
    print(6, end=';')
    print(int('d5'), end=';')
    print(2, end=';')
except ValueError:
    print(1, end=';')
except Exception:
    print(6, end=';')
except:
    print('ERR', end=';')
else:
    print(7, end=';')
finally:
    print(0)
code

6. Що буде виведено за виконання?
code
def f(s):
    for el in s:
        el += 1
        return s
 
s = [1, 2, 3]
s = f(s)
print(s)
code

7. Що буде виведено за виконання?
code
class A:
    a = 1

    def _g1(self): print(2, end=';')
    def __g2(self): print(3, end=';')

    def main(self):
        try: print(self.a, end=';')
        except: print('ERR', end=';')

        try: self._g1()
        except: print('ERR', end=';')

        try: self.__g2()
        except: print('ERR', end=';')


class B(A):
    a = 4

    def _g1(self): print(5, end=';')
    def __g2(self): print(6, end=';')

try: B().main()
except: print('ERR', end=';')
code

8. Що буде виведено за виконання?
code
def f(arg, n):
    n = 9
    tmp = arg
    tmp[-5:-9] = [1,0]
    return len(tmp)>3

try:
    g = [10,11,12,13,14]
    n = 6
    f(g, n)
    print(n, end=';')
    for el in g:
        print(el, end=';')
except:
    print('ERR')
code

9. Що буде виведено за виконання?
code
class A:
    b = 1
    c = 2

    def __init__(self):
        b = 3

class B(A):
    c = 4

    def __init__(self):
        super().__init__()
        self.a = 7

obj = B()
obj.b = obj.c + 10
B.b = 13

try: print(obj.b, end= ';')
except: print('ERR', end=';')

try: print(A.b, end= ';')
except: print('ERR', end=';')

try: print(B.b, end= ';')
except: print('ERR', end=';')
code

10. 
Для графа на рисунку з вершини 4 запущено алгоритм пошуку в глибину, що будує кістякове дерево. Списки суміжності вершин переглядаються алгоритмом за зростанням міток вершин.\par
\begin{center}
	\adjustimage{max size={0.9\linewidth}{0.9\paperheight}}
	{../10/Traversals/197.png}
\end{center}
\par Які з наведених ребер входять до складу отриманого кістякового дерева?\par
1) (2, 4)\par
2) (5, 6)\par
3) (4, 5)\par
4) (4, 6)\par
5) (2, 5)\par
6) (0, 4)\par
{\vskip 15pt} У формі вибрати номери відповідних ребер.\par
%answer:2 5 6
\end {large}}
%end
{\smallskip \begin {small} Затверджено на засіданні кафедри теоретичної кібернетики, протокол \No 12 від   19.05.2020 р.\par\smallskip\smallskip
{\parbox {6.5 cm}{Зав. кафедри \dotfill Ю.В.~Крак}}\hspace {0.7cm}{\hbox to 6.5 cm{ Екзаменатор \dotfill Т.О.~Карнаух}} \end{small}}
%begin
{{\begin {small}\par
{ \center  Київський національний університет імені Тараса Шевченка \par}
{\parbox {5 cm}{Освітня програма \hfill}{\parbox {7 cm}{Прикладна математика, Системний аналіз \hfill}}\parbox {6 cm}{\hfill Семестр 2}}\par
{\parbox {5 cm} {Навчальний предмет \hfill}\parbox {7 cm}{Програмування \hfill}\parbox {6cm}{\hfill}\par}\vspace{-7pt} \end{small}}{\center Екзаменаційний білет \No \textbf 
{ 66 }\par}}
{
\begin {large}
1. Що буде виведено за виконання?
code
try:
    res = 9**-1.0*2
    if isinstance(res, bool):
        print(f'bool;{res}')
    elif isinstance(res, int):
        print(f'int;{res}')
    elif isinstance(res, float):
        print(f'float;{res:.2f}')
    else:
        print('???')
except:
    print('ERR')
code

2. Що буде виведено за виконання?
code
a = 7
b = 9
c = 4

def f(a):
    global c
    a = 3
    b *= 4
    c = 1
    return a + b + c

a = 8
b = 3
c = 1

try:
    print(f(a), a, b, c, sep=';')
except:
    print('ERR', a, b, c, sep=';')
code

3. Що буде виведено за виконання?
code
try:
    m = 6
    while m<12:
        m += 2
    print(m, end=';')
except:
    print('ERR')
code

4. Що буде виведено за виконання?
code
try:
    for d in range(3, -8, -1):
        if d<4:
            break
        print(d, end=';')
    else:
        print(d,end=';')
    print(d)
except:
    print('ERR')
code

5. Що буде виведено за виконання?\par (Дійсні значення, якщо наявні, в цьому завданні будуть виводитись у форматі з фіксованою крапкою та, можливо, з доволі великою кількістю десяткових розрядів. У відповіді для дійсних значень у ЦЬОМУ завданні достатньо вказати один десятковий розряд після крапки, відкинувши решту розрядів без округлення. Наприклад, замість 13.66666666666667 достатньо вказати 13.6, замість 25.23 достатньо вказати 25.2)
code
try:
    print(3, end=';')
    print(9%2, end=';')
    print(5, end=';')
except Exception:
    print(4, end=';')
except KeyboardInterrupt:
    print(8, end=';')
except:
    print('ERR', end=';')
else:
    print(7, end=';')
finally:
    print(9)
code

6. Що буде виведено за виконання?
code
def f(s):
    s1 = s
    for i in range(len(s)):
        s1[i] += 1
        return s1

s = [1, 2, 3]
f(s)
print(s)
code

7. Що буде виведено за виконання?
code
class A:
    c = 1

    def __f1(self): print(2, end=';')
    def _f2(self): print(3, end=';')

    def main(self):
        try: print(self.c, end=';')
        except: print('ERR', end=';')

        try: self.__f1()
        except: print('ERR', end=';')

        try: self._f2()
        except: print('ERR', end=';')


class B(A):
    c = 4

    def __f1(self): print(5, end=';')
    def _f2(self): print(6, end=';')

try: B().main()
except: print('ERR', end=';')
code

8. Що буде виведено за виконання?
code
def f(arg, n):
    n = 4
    tmp = arg
    del tmp[1:-4]
    return len(tmp)>3

try:
    h = ['0','1','2','3','4']
    n = 8
    f(h, n)
    print(n, end=';')
    for el in h:
        print(el, end=';')
except:
    print('ERR')
code

9. Що буде виведено за виконання?
code
class A:
    c = 1
    a = 2

    def __init__(self):
        a = 3

class B(A):
    c = 4

    def __init__(self):
        super().__init__()
        self.b = 7

obj = B()
obj.a = obj.b + 10
B.a = 13

try: print(obj.a, end= ';')
except: print('ERR', end=';')

try: print(A.a, end= ';')
except: print('ERR', end=';')

try: print(B.a, end= ';')
except: print('ERR', end=';')
code

10. 
Для графа на рисунку з вершини 3 запущено алгоритм пошуку в глибину, що будує кістякове дерево. Списки суміжності вершин переглядаються алгоритмом за зростанням міток вершин.\par
\begin{center}
	\adjustimage{max size={0.9\linewidth}{0.9\paperheight}}
	{../10/Traversals/126.png}
\end{center}
\par Які з наведених ребер входять до складу отриманого кістякового дерева?\par
1) (3, 4)\par
2) (3, 6)\par
3) (3, 5)\par
4) (1, 3)\par
5) (1, 2)\par
6) (0, 5)\par
{\vskip 15pt} У формі вибрати номери відповідних ребер.\par
%answer:5 6
\end {large}}
%end
{\smallskip \begin {small} Затверджено на засіданні кафедри теоретичної кібернетики, протокол \No 12 від   19.05.2020 р.\par\smallskip\smallskip
{\parbox {6.5 cm}{Зав. кафедри \dotfill Ю.В.~Крак}}\hspace {0.7cm}{\hbox to 6.5 cm{ Екзаменатор \dotfill Т.О.~Карнаух}} \end{small}}
%begin
{{\begin {small}\par
{ \center  Київський національний університет імені Тараса Шевченка \par}
{\parbox {5 cm}{Освітня програма \hfill}{\parbox {7 cm}{Прикладна математика, Системний аналіз \hfill}}\parbox {6 cm}{\hfill Семестр 2}}\par
{\parbox {5 cm} {Навчальний предмет \hfill}\parbox {7 cm}{Програмування \hfill}\parbox {6cm}{\hfill}\par}\vspace{-7pt} \end{small}}{\center Екзаменаційний білет \No \textbf 
{ 67 }\par}}
{
\begin {large}
1. Що буде виведено за виконання?
code
try:
    res = 9**-0.5*2
    if isinstance(res, bool):
        print(f'bool;{res}')
    elif isinstance(res, int):
        print(f'int;{res}')
    elif isinstance(res, float):
        print(f'float;{res:.2f}')
    else:
        print('???')
except:
    print('ERR')
code

2. Що буде виведено за виконання?
code
a = 1
b = 3
c = 6

def f(b):
    a = 3
    b = 2
    c = 3
    return a + b + c

a = 3
b = 9
c = 8

try:
    print(f(b), a, b, c, sep=';')
except:
    print('ERR', a, b, c, sep=';')
code

3. Що буде виведено за виконання?
code
try:
    t = 1
    while t<=10:
        t += 2
    print(t, end=';')
except:
    print('ERR')
code

4. Що буде виведено за виконання?
code
try:
    for d in range(-9, -5, 2):
        if d>6:
            continue
        print(d, end=';')
    else:
        print(d,end=';')
    print(d)
except:
    print('ERR')
code

5. Що буде виведено за виконання?\par (Дійсні значення, якщо наявні, в цьому завданні будуть виводитись у форматі з фіксованою крапкою та, можливо, з доволі великою кількістю десяткових розрядів. У відповіді для дійсних значень у ЦЬОМУ завданні достатньо вказати один десятковий розряд після крапки, відкинувши решту розрядів без округлення. Наприклад, замість 13.66666666666667 достатньо вказати 13.6, замість 25.23 достатньо вказати 25.2)
code
try:
    print(4, end=';')
    print(int('c5'), end=';')
    print(6, end=';')
except ValueError:
    print(3, end=';')
except ZeroDivisionError:
    print(1, end=';')
except:
    print('ERR', end=';')
else:
    print(2, end=';')
finally:
    print(8)
code

6. Що буде виведено за виконання?
code
def f(n):
    print(n%10, end='')
    if n>9:
        f(n//10)
 
f(123456)
code

7. Що буде виведено за виконання?
code
class A:
    b = 1

    def _f1(self): print(2, end=';')
    def __f2(self): print(3, end=';')

    def main(self):
        try: print(self.b, end=';')
        except: print('ERR', end=';')

        try: self._f1()
        except: print('ERR', end=';')

        try: self.__f2()
        except: print('ERR', end=';')


class B(A):
    b = 4

    def _f1(self): print(5, end=';')
    def __f2(self): print(6, end=';')

try: B().main()
except: print('ERR', end=';')
code

8. Що буде виведено за виконання?
code
def f(arg, n):
    n = 2
    tmp = arg
    tmp[-8:7] = [12]
    return len(tmp)>3

try:
    a = [10,11,12,13,14]
    n = 1
    f(a, n)
    print(n, end=';')
    for el in a:
        print(el, end=';')
except:
    print('ERR')
code

9. Що буде виведено за виконання?
code
class A:
    a = 1
    c = 2

    def __init__(self):
        a = 3

class B(A):
    c = 4

    def __init__(self):
        super().__init__()
        self.b = 7

obj = B()
obj.c = obj.b + 10
B.c = 13

try: print(obj.c, end= ';')
except: print('ERR', end=';')

try: print(A.c, end= ';')
except: print('ERR', end=';')

try: print(B.c, end= ';')
except: print('ERR', end=';')
code

10. 
Для графа на рисунку з вершини 6 запущено алгоритм пошуку в глибину, що будує кістякове дерево. Списки суміжності вершин переглядаються алгоритмом за зростанням міток вершин.\par
\begin{center}
	\adjustimage{max size={0.9\linewidth}{0.9\paperheight}}
	{../10/Traversals/207.png}
\end{center}
\par Які з наведених ребер входять до складу отриманого кістякового дерева?\par
1) (0, 3)\par
2) (5, 6)\par
3) (1, 4)\par
4) (4, 5)\par
5) (0, 1)\par
6) (2, 4)\par
{\vskip 15pt} У формі вибрати номери відповідних ребер.\par
%answer:1 3 4 5 6
\end {large}}
%end
{\smallskip \begin {small} Затверджено на засіданні кафедри теоретичної кібернетики, протокол \No 12 від   19.05.2020 р.\par\smallskip\smallskip
{\parbox {6.5 cm}{Зав. кафедри \dotfill Ю.В.~Крак}}\hspace {0.7cm}{\hbox to 6.5 cm{ Екзаменатор \dotfill Т.О.~Карнаух}} \end{small}}
%begin
{{\begin {small}\par
{ \center  Київський національний університет імені Тараса Шевченка \par}
{\parbox {5 cm}{Освітня програма \hfill}{\parbox {7 cm}{Прикладна математика, Системний аналіз \hfill}}\parbox {6 cm}{\hfill Семестр 2}}\par
{\parbox {5 cm} {Навчальний предмет \hfill}\parbox {7 cm}{Програмування \hfill}\parbox {6cm}{\hfill}\par}\vspace{-7pt} \end{small}}{\center Екзаменаційний білет \No \textbf 
{ 68 }\par}}
{
\begin {large}
1. Що буде виведено за виконання?
code
try:
    res = "1">False
    if isinstance(res, bool):
        print(f'bool;{res}')
    elif isinstance(res, int):
        print(f'int;{res}')
    elif isinstance(res, float):
        print(f'float;{res:.2f}')
    else:
        print('???')
except:
    print('ERR')
code

2. Що буде виведено за виконання?
code
a = 7
b = 0
c = 1

def g(b):
    a = 5
    b = 1
    c = 4
    return a + b + c

a = 5
b = 4
c = 2

try:
    print(g(b), a, b, c, sep=';')
except:
    print('ERR', a, b, c, sep=';')
code

3. Що буде виведено за виконання?
code
try:
    n = 11
    while n>1:
        n += -3
    print(n, end=';')
except:
    print('ERR')
code

4. Що буде виведено за виконання?
code
try:
    for f in range(14, 9, 3):
        if f>=3:
            continue
        print(f, end=';')
    else:
        print('end',end=';')
    print(f)
except:
    print('ERR')
code

5. Що буде виведено за виконання?\par (Дійсні значення, якщо наявні, в цьому завданні будуть виводитись у форматі з фіксованою крапкою та, можливо, з доволі великою кількістю десяткових розрядів. У відповіді для дійсних значень у ЦЬОМУ завданні достатньо вказати один десятковий розряд після крапки, відкинувши решту розрядів без округлення. Наприклад, замість 13.66666666666667 достатньо вказати 13.6, замість 25.23 достатньо вказати 25.2)
code
try:
    print(0, end=';')
    print(int('9'), end=';')
    print(2, end=';')
except TypeError:
    print(8, end=';')
except ValueError:
    print(7, end=';')
except:
    print('ERR', end=';')
else:
    print(6, end=';')
finally:
    print(5)
code

6. Що буде виведено за виконання?
code
def f(n):
    print(n%10, end='')
    if n>9:
        f(n//10)
 
f(543216)
code

7. Що буде виведено за виконання?
code
class A:
    d = 1

    def _g1(self): print(2, end=';')
    def _g2(self): print(3, end=';')

    def main(self):
        try: print(self.d, end=';')
        except: print('ERR', end=';')

        try: self._g1()
        except: print('ERR', end=';')

        try: self._g2()
        except: print('ERR', end=';')


class B(A):
    d = 4

    def _g1(self): print(5, end=';')
    def _g2(self): print(6, end=';')

try: B().main()
except: print('ERR', end=';')
code

8. Що буде виведено за виконання?
code
def f(arg, n):
    arg[40] = 2
    n = 7
    arg = set(arg.keys())

try:
    n = 3
    s = {15:3, 30:1, 57:2, 57:5}
    f(s, n)
    print(n, end=';')
    for x in s.values():
        print(x, sep=';', end=';')
except:
    print('ERR')
code

9. Що буде виведено за виконання?
code
class A:
    c = 1
    b = 2

    def __init__(self):
        b = 3

class B(A):
    c = 4

    def __init__(self):
        super().__init__()
        self.a = 7

obj = B()
obj.a = obj.c + 10
B.a = 13

try: print(obj.a, end= ';')
except: print('ERR', end=';')

try: print(A.a, end= ';')
except: print('ERR', end=';')

try: print(B.a, end= ';')
except: print('ERR', end=';')
code

10. 
Для графа на рисунку з вершини 2 запущено алгоритм пошуку в глибину, що будує кістякове дерево. Списки суміжності вершин переглядаються алгоритмом за зростанням міток вершин.\par
\begin{center}
	\adjustimage{max size={0.9\linewidth}{0.9\paperheight}}
	{../10/Traversals/161.png}
\end{center}
\par Які з наведених ребер входять до складу отриманого кістякового дерева?\par
1) (2, 5)\par
2) (2, 4)\par
3) (4, 5)\par
4) (1, 3)\par
5) (0, 5)\par
6) (0, 1)\par
{\vskip 15pt} У формі вибрати номери відповідних ребер.\par
%answer:2 3 4 5 6
\end {large}}
%end
{\smallskip \begin {small} Затверджено на засіданні кафедри теоретичної кібернетики, протокол \No 12 від   19.05.2020 р.\par\smallskip\smallskip
{\parbox {6.5 cm}{Зав. кафедри \dotfill Ю.В.~Крак}}\hspace {0.7cm}{\hbox to 6.5 cm{ Екзаменатор \dotfill Т.О.~Карнаух}} \end{small}}
%begin
{{\begin {small}\par
{ \center  Київський національний університет імені Тараса Шевченка \par}
{\parbox {5 cm}{Освітня програма \hfill}{\parbox {7 cm}{Прикладна математика, Системний аналіз \hfill}}\parbox {6 cm}{\hfill Семестр 2}}\par
{\parbox {5 cm} {Навчальний предмет \hfill}\parbox {7 cm}{Програмування \hfill}\parbox {6cm}{\hfill}\par}\vspace{-7pt} \end{small}}{\center Екзаменаційний білет \No \textbf 
{ 69 }\par}}
{
\begin {large}
1. Що буде виведено за виконання?
code
try:
    res = "False"<6
    if isinstance(res, bool):
        print(f'bool;{res}')
    elif isinstance(res, int):
        print(f'int;{res}')
    elif isinstance(res, float):
        print(f'float;{res:.2f}')
    else:
        print('???')
except:
    print('ERR')
code

2. Що буде виведено за виконання?
code
a = 5
b = 6
c = 8

def f(b):
    a = 2
    b += 3
    c = 1
    return a + b + c

a = 5
b = 7
c = 2

try:
    print(f(b), a, b, c, sep=';')
except:
    print('ERR', a, b, c, sep=';')
code

3. Що буде виведено за виконання?
code
try:
    k = 5
    while k>=3:
        k += -2
    print(k, end=';')
except:
    print('ERR')
code

4. Що буде виведено за виконання?
code
try:
    for c in range(-11, 10, -3):
        if c>=4:
            break
        print(c, end=';');c = 6
    else:
        print(c,end=';')
    print(c)
except:
    print('ERR')
code

5. Що буде виведено за виконання?\par (Дійсні значення, якщо наявні, в цьому завданні будуть виводитись у форматі з фіксованою крапкою та, можливо, з доволі великою кількістю десяткових розрядів. У відповіді для дійсних значень у ЦЬОМУ завданні достатньо вказати один десятковий розряд після крапки, відкинувши решту розрядів без округлення. Наприклад, замість 13.66666666666667 достатньо вказати 13.6, замість 25.23 достатньо вказати 25.2)
code
try:
    print(3, end=';')
    print(1j>=5j, end=';')
    print(0, end=';')
except TypeError:
    print(4, end=';')
except KeyboardInterrupt:
    print(1, end=';')
except:
    print('ERR', end=';')
else:
    print(0, end=';')
finally:
    print(6)
code

6. Що буде виведено за виконання?
code
def f(n):
    if n>2:
        f(n//2)
    else:
        print(n%2, end='')

f(16)
code

7. Що буде виведено за виконання?
code
class A:
    c = 1

    def h1(self): print(2, end=';')
    def _h2(self): print(3, end=';')

    def main(self):
        try: print(self.c, end=';')
        except: print('ERR', end=';')

        try: self.h1()
        except: print('ERR', end=';')

        try: self._h2()
        except: print('ERR', end=';')


class B(A):
    c = 4

    def h1(self): print(5, end=';')
    def _h2(self): print(6, end=';')

try: B().main()
except: print('ERR', end=';')
code

8. Що буде виведено за виконання?
code
def f(arg, n):
    arg[83] = 0
    n = 5
    arg = set(arg.values())

try:
    n = 4
    s = {83:3, 36:7, 16:5, 36:8}
    f(s, n)
    print(n, end=';')
    for x in s.keys():
        print(x, sep=';', end=';')
except:
    print('ERR')
code

9. Що буде виведено за виконання?
code
class A:
    a = 1
    b = 2

    def __init__(self):
        b = 3

class B(A):
    b = 4

    def __init__(self):
        super().__init__()
        self.c = 7

obj = B()
obj.a = obj.b + 10
B.a = 13

try: print(obj.a, end= ';')
except: print('ERR', end=';')

try: print(A.a, end= ';')
except: print('ERR', end=';')

try: print(B.a, end= ';')
except: print('ERR', end=';')
code

10. 
Для графа на рисунку з вершини 6 запущено алгоритм пошуку в глибину, що будує кістякове дерево. Списки суміжності вершин переглядаються алгоритмом за зростанням міток вершин.\par
\begin{center}
	\adjustimage{max size={0.9\linewidth}{0.9\paperheight}}
	{../10/Traversals/292.png}
\end{center}
\par Які з наведених ребер входять до складу отриманого кістякового дерева?\par
1) (5, 6)\par
2) (4, 5)\par
3) (3, 4)\par
4) (1, 5)\par
5) (4, 6)\par
6) (2, 5)\par
{\vskip 15pt} У формі вибрати номери відповідних ребер.\par
%answer:2 3
\end {large}}
%end
{\smallskip \begin {small} Затверджено на засіданні кафедри теоретичної кібернетики, протокол \No 12 від   19.05.2020 р.\par\smallskip\smallskip
{\parbox {6.5 cm}{Зав. кафедри \dotfill Ю.В.~Крак}}\hspace {0.7cm}{\hbox to 6.5 cm{ Екзаменатор \dotfill Т.О.~Карнаух}} \end{small}}
%begin
{{\begin {small}\par
{ \center  Київський національний університет імені Тараса Шевченка \par}
{\parbox {5 cm}{Освітня програма \hfill}{\parbox {7 cm}{Прикладна математика, Системний аналіз \hfill}}\parbox {6 cm}{\hfill Семестр 2}}\par
{\parbox {5 cm} {Навчальний предмет \hfill}\parbox {7 cm}{Програмування \hfill}\parbox {6cm}{\hfill}\par}\vspace{-7pt} \end{small}}{\center Екзаменаційний білет \No \textbf 
{ 70 }\par}}
{
\begin {large}
1. Що буде виведено за виконання?
code
try:
    res = 3//5%4*6
    if isinstance(res, bool):
        print(f'bool;{res}')
    elif isinstance(res, int):
        print(f'int;{res}')
    elif isinstance(res, float):
        print(f'float;{res:.2f}')
    else:
        print('???')
except:
    print('ERR')
code

2. Що буде виведено за виконання?
code
a = 1
b = 4
c = 9

def h(b):
    a *= 5
    b = 4
    c = 2
    return a + b + c

a = 3
b = 0
c = 6

try:
    print(h(b), a, b, c, sep=';')
except:
    print('ERR', a, b, c, sep=';')
code

3. Що буде виведено за виконання?
code
try:
    i = 7
    while i<6:
        i += 1
    print(i, end=';')
except:
    print('ERR')
code

4. Що буде виведено за виконання?
code
try:
    for a in range(12, 8, -2):
        if a<=5:
            break
        print(a, end=';')
    else:
        print(a,end=';')
    print(a)
except:
    print('ERR')
code

5. Що буде виведено за виконання?\par (Дійсні значення, якщо наявні, в цьому завданні будуть виводитись у форматі з фіксованою крапкою та, можливо, з доволі великою кількістю десяткових розрядів. У відповіді для дійсних значень у ЦЬОМУ завданні достатньо вказати один десятковий розряд після крапки, відкинувши решту розрядів без округлення. Наприклад, замість 13.66666666666667 достатньо вказати 13.6, замість 25.23 достатньо вказати 25.2)
code
try:
    print(8, end=';')
    print(3/0, end=';')
    print(2, end=';')
except Exception:
    print(7, end=';')
except ZeroDivisionError:
    print(9, end=';')
except:
    print('ERR', end=';')
else:
    print(5, end=';')
finally:
    print(4)
code

6. Що буде виведено за виконання?
code
def f(n):
    if n>9:
        return f(n//100)+1
    else:
        return 1

print(f(123456))
code

7. Що буде виведено за виконання?
code
class A:
    b = 1

    def _h1(self): print(2, end=';')
    def h2(self): print(3, end=';')

    def main(self):
        try: print(self.b, end=';')
        except: print('ERR', end=';')

        try: self._h1()
        except: print('ERR', end=';')

        try: self.h2()
        except: print('ERR', end=';')


class B(A):
    b = 4

    def _h1(self): print(5, end=';')
    def h2(self): print(6, end=';')

try: B().main()
except: print('ERR', end=';')
code

8. Що буде виведено за виконання?
code
def f(arg, n):
    n = 5
    tmp = arg
    tmp[:4] = 'yz'
    return len(tmp)>3

try:
    c = [10,11,12,13,14]
    n = 7
    f(c, n)
    print(n, end=';')
    for el in c:
        print(el, end=';')
except:
    print('ERR')
code

9. Що буде виведено за виконання?
code
class A:
    c = 1
    a = 2

    def __init__(self):
        c = 3

class B(A):
    c = 4

    def __init__(self):
        super().__init__()
        self.b = 7

obj = B()
obj.c = obj.a + 10
B.c = 13

try: print(obj.c, end= ';')
except: print('ERR', end=';')

try: print(A.c, end= ';')
except: print('ERR', end=';')

try: print(B.c, end= ';')
except: print('ERR', end=';')
code

10. 
Для графа на рисунку з вершини 2 запущено алгоритм пошуку в глибину, що будує кістякове дерево. Списки суміжності вершин переглядаються алгоритмом за зростанням міток вершин.\par
\begin{center}
	\adjustimage{max size={0.9\linewidth}{0.9\paperheight}}
	{../10/Traversals/108.png}
\end{center}
\par Які з наведених ребер входять до складу отриманого кістякового дерева?\par
1) (3, 4)\par
2) (1, 6)\par
3) (0, 2)\par
4) (1, 5)\par
5) (3, 6)\par
6) (2, 3)\par
{\vskip 15pt} У формі вибрати номери відповідних ребер.\par
%answer:1 3 4 5
\end {large}}
%end
{\smallskip \begin {small} Затверджено на засіданні кафедри теоретичної кібернетики, протокол \No 12 від   19.05.2020 р.\par\smallskip\smallskip
{\parbox {6.5 cm}{Зав. кафедри \dotfill Ю.В.~Крак}}\hspace {0.7cm}{\hbox to 6.5 cm{ Екзаменатор \dotfill Т.О.~Карнаух}} \end{small}}
%begin
{{\begin {small}\par
{ \center  Київський національний університет імені Тараса Шевченка \par}
{\parbox {5 cm}{Освітня програма \hfill}{\parbox {7 cm}{Прикладна математика, Системний аналіз \hfill}}\parbox {6 cm}{\hfill Семестр 2}}\par
{\parbox {5 cm} {Навчальний предмет \hfill}\parbox {7 cm}{Програмування \hfill}\parbox {6cm}{\hfill}\par}\vspace{-7pt} \end{small}}{\center Екзаменаційний білет \No \textbf 
{ 71 }\par}}
{
\begin {large}
1. Що буде виведено за виконання?
code
try:
    res = 6/8*7*3
    if isinstance(res, bool):
        print(f'bool;{res}')
    elif isinstance(res, int):
        print(f'int;{res}')
    elif isinstance(res, float):
        print(f'float;{res:.2f}')
    else:
        print('???')
except:
    print('ERR')
code

2. Що буде виведено за виконання?
code
a = 1
b = 2
c = 4

def g(a):
    global c
    a -= 1
    b = 3
    c = 5
    return a + b + c

a = 8
b = 0
c = 6

try:
    print(g(a), a, b, c, sep=';')
except:
    print('ERR', a, b, c, sep=';')
code

3. Що буде виведено за виконання?
code
try:
    i = 3
    while i<=7:
        i += 1
    print(i, end=';')
except:
    print('ERR')
code

4. Що буде виведено за виконання?
code
try:
    for b in range(-3, 3, 2):
        if b<8:
            continue
        print(b, end=';');b = -3
    else:
        print(b,end=';')
    print(b)
except:
    print('ERR')
code

5. Що буде виведено за виконання?\par (Дійсні значення, якщо наявні, в цьому завданні будуть виводитись у форматі з фіксованою крапкою та, можливо, з доволі великою кількістю десяткових розрядів. У відповіді для дійсних значень у ЦЬОМУ завданні достатньо вказати один десятковий розряд після крапки, відкинувши решту розрядів без округлення. Наприклад, замість 13.66666666666667 достатньо вказати 13.6, замість 25.23 достатньо вказати 25.2)
code
try:
    print(1, end=';')
    print(3%3, end=';')
    print(9, end=';')
except ZeroDivisionError:
    print(2, end=';')
except TypeError:
    print(6, end=';')
except:
    print('ERR', end=';')
else:
    print(0, end=';')
finally:
    print(8)
code

6. Що буде виведено за виконання?
code
def f(n):
    print(n%10, end='')
    if n>9:
        f(n//100)
 
f(987654)
code

7. Що буде виведено за виконання?
code
class A:
    d = 1

    def f1(self): print(2, end=';')
    def _f2(self): print(3, end=';')

    def main(self):
        try: print(self.d, end=';')
        except: print('ERR', end=';')

        try: self.f1()
        except: print('ERR', end=';')

        try: self._f2()
        except: print('ERR', end=';')


class B(A):
    d = 4

    def f1(self): print(5, end=';')
    def _f2(self): print(6, end=';')

try: B().main()
except: print('ERR', end=';')
code

8. Що буде виведено за виконання?
code
def f(arg, n):
    arg[71] = 3
    n -= 2

try:
    n = 4
    d = {82:8, 78:0, 89:0, 89:8}
    f(d, n)
    print(n, end=';')
    for x, y in d.items():
        print(x, y, sep=';', end=';')
except:
    print('ERR')
code

9. Що буде виведено за виконання?
code
class A:
    b = 1
    c = 2

    def __init__(self):
        b = 3

class B(A):
    b = 4

    def __init__(self):
        super().__init__()
        self.a = 7

obj = B()
obj.b = obj.b + 10
B.b = 13

try: print(obj.b, end= ';')
except: print('ERR', end=';')

try: print(A.b, end= ';')
except: print('ERR', end=';')

try: print(B.b, end= ';')
except: print('ERR', end=';')
code

10. 
Для графа на рисунку з вершини 5 запущено алгоритм пошуку в глибину, що будує кістякове дерево. Списки суміжності вершин переглядаються алгоритмом за зростанням міток вершин.\par
\begin{center}
	\adjustimage{max size={0.9\linewidth}{0.9\paperheight}}
	{../10/Traversals/98.png}
\end{center}
\par Які з наведених ребер входять до складу отриманого кістякового дерева?\par
1) (3, 4)\par
2) (0, 3)\par
3) (1, 6)\par
4) (0, 1)\par
5) (4, 6)\par
6) (5, 6)\par
{\vskip 15pt} У формі вибрати номери відповідних ребер.\par
%answer:1 2 3 4 6
\end {large}}
%end
{\smallskip \begin {small} Затверджено на засіданні кафедри теоретичної кібернетики, протокол \No 12 від   19.05.2020 р.\par\smallskip\smallskip
{\parbox {6.5 cm}{Зав. кафедри \dotfill Ю.В.~Крак}}\hspace {0.7cm}{\hbox to 6.5 cm{ Екзаменатор \dotfill Т.О.~Карнаух}} \end{small}}
%begin
{{\begin {small}\par
{ \center  Київський національний університет імені Тараса Шевченка \par}
{\parbox {5 cm}{Освітня програма \hfill}{\parbox {7 cm}{Прикладна математика, Системний аналіз \hfill}}\parbox {6 cm}{\hfill Семестр 2}}\par
{\parbox {5 cm} {Навчальний предмет \hfill}\parbox {7 cm}{Програмування \hfill}\parbox {6cm}{\hfill}\par}\vspace{-7pt} \end{small}}{\center Екзаменаційний білет \No \textbf 
{ 72 }\par}}
{
\begin {large}
1. Що буде виведено за виконання?
code
def f(n):
    print('f', end='')
    return n

try:
    print(f(3) or f(6))
except:
    print('ERR')
code

2. Що буде виведено за виконання?
code
a = 9
b = 7
c = 5

def g(a):
    global c
    a -= 4
    b = 2
    c = 3
    return a + b + c

a = 3
b = 3
c = 4

try:
    print(g(b), a, b, c, sep=';')
except:
    print('ERR', a, b, c, sep=';')
code

3. Що буде виведено за виконання?
code
try:
    m = 8
    while m<10:
        m += 2
    print(m, end=';')
except:
    print('ERR')
code

4. Що буде виведено за виконання?
code
try:
    for e in range(-13, -1, -3):
        if e>7:
            continue
        print(e, end=';')
        if e<=4:
            break
    else:
        print(e,end=';')
    print(e)
except:
    print('ERR')
code

5. Що буде виведено за виконання?\par (Дійсні значення, якщо наявні, в цьому завданні будуть виводитись у форматі з фіксованою крапкою та, можливо, з доволі великою кількістю десяткових розрядів. У відповіді для дійсних значень у ЦЬОМУ завданні достатньо вказати один десятковий розряд після крапки, відкинувши решту розрядів без округлення. Наприклад, замість 13.66666666666667 достатньо вказати 13.6, замість 25.23 достатньо вказати 25.2)
code
try:
    print(7, end=';')
    print(int('4'), end=';')
    print(5, end=';')
except ValueError:
    print(9, end=';')
except KeyboardInterrupt:
    print(7, end=';')
except:
    print('ERR', end=';')
else:
    print(3, end=';')
finally:
    print(0)
code

6. Що буде виведено за виконання?
code
def f(n):
    if n>9:
        f(n//100)
    else:
        print(n%10,end='')

f(123456)
code

7. Що буде виведено за виконання?
code
class A:
    a = 1

    def _g1(self): print(2, end=';')
    def g2(self): print(3, end=';')

    def main(self):
        try: print(self.a, end=';')
        except: print('ERR', end=';')

        try: self._g1()
        except: print('ERR', end=';')

        try: self.g2()
        except: print('ERR', end=';')


class B(A):
    a = 4

    def _g1(self): print(5, end=';')
    def g2(self): print(6, end=';')

try: B().main()
except: print('ERR', end=';')
code

8. Що буде виведено за виконання?
code
def f(arg, n):
    n = 0
    tmp = arg
    del tmp[5:-3]
    return len(tmp)>3

try:
    e = ['0','1','2','3','4']
    n = 6
    f(e, n)
    print(n, end=';')
    for el in e:
        print(el, end=';')
except:
    print('ERR')
code

9. Що буде виведено за виконання?
code
class A:
    b = 1
    a = 2

    def __init__(self):
        a = 3

class B(A):
    a = 4

    def __init__(self):
        super().__init__()
        self.c = 7

obj = B()
obj.c = obj.a + 10
B.c = 13

try: print(obj.c, end= ';')
except: print('ERR', end=';')

try: print(A.c, end= ';')
except: print('ERR', end=';')

try: print(B.c, end= ';')
except: print('ERR', end=';')
code

10. 
Для графа на рисунку з вершини 5 запущено алгоритм пошуку в глибину, що будує кістякове дерево. Списки суміжності вершин переглядаються алгоритмом за зростанням міток вершин.\par
\begin{center}
	\adjustimage{max size={0.9\linewidth}{0.9\paperheight}}
	{../10/Traversals/25.png}
\end{center}
\par Які з наведених ребер входять до складу отриманого кістякового дерева?\par
1) (4, 5)\par
2) (1, 6)\par
3) (3, 6)\par
4) (0, 1)\par
5) (2, 6)\par
6) (1, 4)\par
{\vskip 15pt} У формі вибрати номери відповідних ребер.\par
%answer:1 3 4 5 6
\end {large}}
%end
{\smallskip \begin {small} Затверджено на засіданні кафедри теоретичної кібернетики, протокол \No 12 від   19.05.2020 р.\par\smallskip\smallskip
{\parbox {6.5 cm}{Зав. кафедри \dotfill Ю.В.~Крак}}\hspace {0.7cm}{\hbox to 6.5 cm{ Екзаменатор \dotfill Т.О.~Карнаух}} \end{small}}
%begin
{{\begin {small}\par
{ \center  Київський національний університет імені Тараса Шевченка \par}
{\parbox {5 cm}{Освітня програма \hfill}{\parbox {7 cm}{Прикладна математика, Системний аналіз \hfill}}\parbox {6 cm}{\hfill Семестр 2}}\par
{\parbox {5 cm} {Навчальний предмет \hfill}\parbox {7 cm}{Програмування \hfill}\parbox {6cm}{\hfill}\par}\vspace{-7pt} \end{small}}{\center Екзаменаційний білет \No \textbf 
{ 73 }\par}}
{
\begin {large}
1. Що буде виведено за виконання?
code
try:
    res = 6.0j=="8j"
    if isinstance(res, bool):
        print(f'bool;{res}')
    elif isinstance(res, int):
        print(f'int;{res}')
    elif isinstance(res, float):
        print(f'float;{res:.2f}')
    else:
        print('???')
except:
    print('ERR')
code

2. Що буде виведено за виконання?
code
a = 9
b = 7
c = 2

def h(b):
    a *= 1
    b = 4
    c = 5
    return a + b + c

a = 8
b = 1
c = 6

try:
    print(h(a), a, b, c, sep=';')
except:
    print('ERR', a, b, c, sep=';')
code

3. Що буде виведено за виконання?
code
try:
    j = 8
    while j>7:
        j += -3
    print(j, end=';')
except:
    print('ERR')
code

4. Що буде виведено за виконання?
code
try:
    for b in range(10, -10, 3):
        if b<6:
            continue
        print(b, end=';')
    else:
        print(b,end=';')
    print(b)
except:
    print('ERR')
code

5. Що буде виведено за виконання?\par (Дійсні значення, якщо наявні, в цьому завданні будуть виводитись у форматі з фіксованою крапкою та, можливо, з доволі великою кількістю десяткових розрядів. У відповіді для дійсних значень у ЦЬОМУ завданні достатньо вказати один десятковий розряд після крапки, відкинувши решту розрядів без округлення. Наприклад, замість 13.66666666666667 достатньо вказати 13.6, замість 25.23 достатньо вказати 25.2)
code
try:
    print(1, end=';')
    print(5//0.0, end=';')
    print(4, end=';')
except Exception:
    print(2, end=';')
except TypeError:
    print(6, end=';')
except:
    print('ERR', end=';')
else:
    print(8, end=';')
finally:
    print(8)
code

6. Що буде виведено за виконання?
code
def f(s):
    s1 = s
    for i in range(len(s)):
        s1[i] += 1
    return s1

s = [1, 2, 3]
f(s)
print(s)
code

7. Що буде виведено за виконання?
code
class A:
    c = 1

    def _h1(self): print(2, end=';')
    def _h2(self): print(3, end=';')

    def main(self):
        try: print(self.c, end=';')
        except: print('ERR', end=';')

        try: self._h1()
        except: print('ERR', end=';')

        try: self._h2()
        except: print('ERR', end=';')


class B(A):
    c = 4

    def _h1(self): print(5, end=';')
    def _h2(self): print(6, end=';')

try: B().main()
except: print('ERR', end=';')
code

8. Що буде виведено за виконання?
code
def f(arg, n):
    arg[63] = 9
    n = 7
    arg = tuple(arg.items())

try:
    n = 3
    d = {85:6, 28:0, 63:8, 63:6}
    f(d, n)
    print(n, end=';')
    for x in d.keys():
        print(x, sep=';', end=';')
except:
    print('ERR')
code

9. Що буде виведено за виконання?
code
class A:
    c = 1
    b = 2

    def __init__(self):
        b = 3

class B(A):
    c = 4

    def __init__(self):
        super().__init__()
        self.a = 7

obj = B()
obj.c = obj.b + 10
B.c = 13

try: print(obj.c, end= ';')
except: print('ERR', end=';')

try: print(A.c, end= ';')
except: print('ERR', end=';')

try: print(B.c, end= ';')
except: print('ERR', end=';')
code

10. 
Для графа на рисунку з вершини 1 запущено алгоритм пошуку в глибину, що будує кістякове дерево. Списки суміжності вершин переглядаються алгоритмом за зростанням міток вершин.\par
\begin{center}
	\adjustimage{max size={0.9\linewidth}{0.9\paperheight}}
	{../10/Traversals/52.png}
\end{center}
\par Які з наведених ребер входять до складу отриманого кістякового дерева?\par
1) (0, 4)\par
2) (0, 5)\par
3) (0, 2)\par
4) (1, 6)\par
5) (0, 6)\par
6) (0, 1)\par
{\vskip 15pt} У формі вибрати номери відповідних ребер.\par
%answer:1 3 5 6
\end {large}}
%end
{\smallskip \begin {small} Затверджено на засіданні кафедри теоретичної кібернетики, протокол \No 12 від   19.05.2020 р.\par\smallskip\smallskip
{\parbox {6.5 cm}{Зав. кафедри \dotfill Ю.В.~Крак}}\hspace {0.7cm}{\hbox to 6.5 cm{ Екзаменатор \dotfill Т.О.~Карнаух}} \end{small}}
%begin
{{\begin {small}\par
{ \center  Київський національний університет імені Тараса Шевченка \par}
{\parbox {5 cm}{Освітня програма \hfill}{\parbox {7 cm}{Прикладна математика, Системний аналіз \hfill}}\parbox {6 cm}{\hfill Семестр 2}}\par
{\parbox {5 cm} {Навчальний предмет \hfill}\parbox {7 cm}{Програмування \hfill}\parbox {6cm}{\hfill}\par}\vspace{-7pt} \end{small}}{\center Екзаменаційний білет \No \textbf 
{ 74 }\par}}
{
\begin {large}
1. Що буде виведено за виконання?
code
def f(n):
    print('f', end='')
    return n!=-4

try:
    print(f(0) and f(-4))
except:
    print('ERR')
code

2. Що буде виведено за виконання?
code
a = 0
b = 5
c = 6

def f(b):
    a += 2
    b = 5
    c = 1
    return a + b + c

a = 2
b = 7
c = 1

try:
    print(f(b), a, b, c, sep=';')
except:
    print('ERR', a, b, c, sep=';')
code

3. Що буде виведено за виконання?
code
try:
    j = 9
    while j<=8:
        j += 2
    print(j, end=';')
except:
    print('ERR')
code

4. Що буде виведено за виконання?
code
try:
    for f in range(-6, 5, -1):
        if f<=5:
            break
        print(f, end=';');f = -7
    else:
        print(f,end=';')
    print(f)
except:
    print('ERR')
code

5. Що буде виведено за виконання?\par (Дійсні значення, якщо наявні, в цьому завданні будуть виводитись у форматі з фіксованою крапкою та, можливо, з доволі великою кількістю десяткових розрядів. У відповіді для дійсних значень у ЦЬОМУ завданні достатньо вказати один десятковий розряд після крапки, відкинувши решту розрядів без округлення. Наприклад, замість 13.66666666666667 достатньо вказати 13.6, замість 25.23 достатньо вказати 25.2)
code
try:
    print(5, end=';')
    print(7j!=9j, end=';')
    print(4, end=';')
except ZeroDivisionError:
    print(3, end=';')
except ValueError:
    print(6, end=';')
except:
    print('ERR', end=';')
else:
    print(0, end=';')
finally:
    print(1)
code

6. Що буде виведено за виконання?
code
def f(n):
    if n<=9:
        print(n%10,end='')
    else:
        f(n//10)

f(543216)
code

7. Що буде виведено за виконання?
code
class A:
    a = 1

    def __g1(self): print(2, end=';')
    def _g2(self): print(3, end=';')

    def main(self):
        try: print(self.a, end=';')
        except: print('ERR', end=';')

        try: self.__g1()
        except: print('ERR', end=';')

        try: self._g2()
        except: print('ERR', end=';')


class B(A):
    a = 4

    def __g1(self): print(5, end=';')
    def _g2(self): print(6, end=';')

try: B().main()
except: print('ERR', end=';')
code

8. Що буде виведено за виконання?
code
def f(arg, n):
    arg[51] = 9
    n = 8
    arg = list(arg.keys())

try:
    n = 0
    t = {11:9, 77:7, 21:8, 21:1}
    f(t, n)
    print(n, end=';')
    for x in t.items():
        print(x, sep=';', end=';')
except:
    print('ERR')
code

9. Що буде виведено за виконання?
code
class A:
    a = 1
    b = 2

    def __init__(self):
        a = 3

class B(A):
    b = 4

    def __init__(self):
        super().__init__()
        self.c = 7

obj = B()
obj.a = obj.b + 10
B.a = 13

try: print(obj.a, end= ';')
except: print('ERR', end=';')

try: print(A.a, end= ';')
except: print('ERR', end=';')

try: print(B.a, end= ';')
except: print('ERR', end=';')
code

10. 
Для графа на рисунку з вершини 2 запущено алгоритм пошуку в глибину, що будує кістякове дерево. Списки суміжності вершин переглядаються алгоритмом за зростанням міток вершин.\par
\begin{center}
	\adjustimage{max size={0.9\linewidth}{0.9\paperheight}}
	{../10/Traversals/151.png}
\end{center}
\par Які з наведених ребер входять до складу отриманого кістякового дерева?\par
1) (1, 6)\par
2) (0, 6)\par
3) (2, 5)\par
4) (0, 5)\par
5) (1, 2)\par
6) (0, 1)\par
{\vskip 15pt} У формі вибрати номери відповідних ребер.\par
%answer:5 6
\end {large}}
%end
{\smallskip \begin {small} Затверджено на засіданні кафедри теоретичної кібернетики, протокол \No 12 від   19.05.2020 р.\par\smallskip\smallskip
{\parbox {6.5 cm}{Зав. кафедри \dotfill Ю.В.~Крак}}\hspace {0.7cm}{\hbox to 6.5 cm{ Екзаменатор \dotfill Т.О.~Карнаух}} \end{small}}
%begin
{{\begin {small}\par
{ \center  Київський національний університет імені Тараса Шевченка \par}
{\parbox {5 cm}{Освітня програма \hfill}{\parbox {7 cm}{Прикладна математика, Системний аналіз \hfill}}\parbox {6 cm}{\hfill Семестр 2}}\par
{\parbox {5 cm} {Навчальний предмет \hfill}\parbox {7 cm}{Програмування \hfill}\parbox {6cm}{\hfill}\par}\vspace{-7pt} \end{small}}{\center Екзаменаційний білет \No \textbf 
{ 75 }\par}}
{
\begin {large}
1. Що буде виведено за виконання?
code
try:
    res = not 7==False or 9!=5 and 6.0>=4.0
    if isinstance(res, bool):
        print(f'bool;{res}')
    elif isinstance(res, int):
        print(f'int;{res}')
    elif isinstance(res, float):
        print(f'float;{res:.2f}')
    else:
        print('???')
except:
    print('ERR')
code

2. Що буде виведено за виконання?
code
a = 7
b = 1
c = 5

def h(a):
    global c
    a = 2
    b = 4
    c = 1
    return a + b + c

a = 0
b = 6
c = 2

try:
    print(h(a), a, b, c, sep=';')
except:
    print('ERR', a, b, c, sep=';')
code

3. Що буде виведено за виконання?
code
try:
    n = 2
    while n<6:
        n += 1
    print(n, end=';')
except:
    print('ERR')
code

4. Що буде виведено за виконання?
code
try:
    for e in range(-2, -15, -2):
        if e>=3:
            continue
        print(e, end=';')
    else:
        print(e,end=';')
    print(e)
except:
    print('ERR')
code

5. Що буде виведено за виконання?\par (Дійсні значення, якщо наявні, в цьому завданні будуть виводитись у форматі з фіксованою крапкою та, можливо, з доволі великою кількістю десяткових розрядів. У відповіді для дійсних значень у ЦЬОМУ завданні достатньо вказати один десятковий розряд після крапки, відкинувши решту розрядів без округлення. Наприклад, замість 13.66666666666667 достатньо вказати 13.6, замість 25.23 достатньо вказати 25.2)
code
try:
    print(9, end=';')
    print(int('a2'), end=';')
    print(8, end=';')
except KeyboardInterrupt:
    print(3, end=';')
except Exception:
    print(2, end=';')
except:
    print('ERR', end=';')
else:
    print(6, end=';')
finally:
    print(4)
code

6. Що буде виведено за виконання?
code
def f(n):
    if n<=9:
        print(n%10, end='')
    else:
        f(n//10)

f(987651)
code

7. Що буде виведено за виконання?
code
class A:
    b = 1

    def _f1(self): print(2, end=';')
    def __f2(self): print(3, end=';')

    def main(self):
        try: print(self.b, end=';')
        except: print('ERR', end=';')

        try: self._f1()
        except: print('ERR', end=';')

        try: self.__f2()
        except: print('ERR', end=';')


class B(A):
    b = 4

    def _f1(self): print(5, end=';')
    def __f2(self): print(6, end=';')

try: B().main()
except: print('ERR', end=';')
code

8. Що буде виведено за виконання?
code
def f(arg, n):
    arg[75] = 0
    n += -2

try:
    n = 6
    t = {77:6, 75:8, 77:0}
    f(t, n)
    print(n, end=';')
    for x in t.keys():
        print(x, sep=';', end=';')
except:
    print('ERR')
code

9. Що буде виведено за виконання?
code
class A:
    a = 1
    c = 2

    def __init__(self):
        a = 3

class B(A):
    a = 4

    def __init__(self):
        super().__init__()
        self.b = 7

obj = B()
obj.a = obj.c + 10
B.a = 13

try: print(obj.a, end= ';')
except: print('ERR', end=';')

try: print(A.a, end= ';')
except: print('ERR', end=';')

try: print(B.a, end= ';')
except: print('ERR', end=';')
code

10. 
Для графа на рисунку з вершини 3 запущено алгоритм пошуку в глибину, що будує кістякове дерево. Списки суміжності вершин переглядаються алгоритмом за зростанням міток вершин.\par
\begin{center}
	\adjustimage{max size={0.9\linewidth}{0.9\paperheight}}
	{../10/Traversals/37.png}
\end{center}
\par Які з наведених ребер входять до складу отриманого кістякового дерева?\par
1) (0, 5)\par
2) (2, 4)\par
3) (0, 4)\par
4) (1, 2)\par
5) (0, 6)\par
6) (0, 3)\par
{\vskip 15pt} У формі вибрати номери відповідних ребер.\par
%answer:2 3 4 6
\end {large}}
%end
{\smallskip \begin {small} Затверджено на засіданні кафедри теоретичної кібернетики, протокол \No 12 від   19.05.2020 р.\par\smallskip\smallskip
{\parbox {6.5 cm}{Зав. кафедри \dotfill Ю.В.~Крак}}\hspace {0.7cm}{\hbox to 6.5 cm{ Екзаменатор \dotfill Т.О.~Карнаух}} \end{small}}
%begin
{{\begin {small}\par
{ \center  Київський національний університет імені Тараса Шевченка \par}
{\parbox {5 cm}{Освітня програма \hfill}{\parbox {7 cm}{Прикладна математика, Системний аналіз \hfill}}\parbox {6 cm}{\hfill Семестр 2}}\par
{\parbox {5 cm} {Навчальний предмет \hfill}\parbox {7 cm}{Програмування \hfill}\parbox {6cm}{\hfill}\par}\vspace{-7pt} \end{small}}{\center Екзаменаційний білет \No \textbf 
{ 76 }\par}}
{
\begin {large}
1. Що буде виведено за виконання?
code
try:
    res = 7//3//4*7
    if isinstance(res, bool):
        print(f'bool;{res}')
    elif isinstance(res, int):
        print(f'int;{res}')
    elif isinstance(res, float):
        print(f'float;{res:.2f}')
    else:
        print('???')
except:
    print('ERR')
code

2. Що буде виведено за виконання?
code
a = 9
b = 4
c = 3

def g(b):
    global c
    a = 5
    b = 3
    c = 5
    return a + b + c

a = 8
b = 9
c = 1

try:
    print(g(b), a, b, c, sep=';')
except:
    print('ERR', a, b, c, sep=';')
code

3. Що буде виведено за виконання?
code
try:
    m = 6
    while m>=5:
        m += -2
    print(m, end=';')
except:
    print('ERR')
code

4. Що буде виведено за виконання?
code
try:
    for a in range(-7, 15, 3):
        if a>4:
            break
        print(a, end=';');a = 0
    else:
        print(a,end=';')
    print(a)
except:
    print('ERR')
code

5. Що буде виведено за виконання?\par (Дійсні значення, якщо наявні, в цьому завданні будуть виводитись у форматі з фіксованою крапкою та, можливо, з доволі великою кількістю десяткових розрядів. У відповіді для дійсних значень у ЦЬОМУ завданні достатньо вказати один десятковий розряд після крапки, відкинувши решту розрядів без округлення. Наприклад, замість 13.66666666666667 достатньо вказати 13.6, замість 25.23 достатньо вказати 25.2)
code
try:
    print(0, end=';')
    print(1/1, end=';')
    print(7, end=';')
except ZeroDivisionError:
    print(9, end=';')
except KeyboardInterrupt:
    print(5, end=';')
except:
    print('ERR', end=';')
else:
    print(7, end=';')
finally:
    print(3)
code

6. Що буде виведено за виконання?
code
def f(s):
    for el in s:
        el += 1
    return s
 
s = [1, 2, 3]
s = f(s)
print(s)
code

7. Що буде виведено за виконання?
code
class A:
    d = 1

    def g1(self): print(2, end=';')
    def _g2(self): print(3, end=';')

    def main(self):
        try: print(self.d, end=';')
        except: print('ERR', end=';')

        try: self.g1()
        except: print('ERR', end=';')

        try: self._g2()
        except: print('ERR', end=';')


class B(A):
    d = 4

    def g1(self): print(5, end=';')
    def _g2(self): print(6, end=';')

try: B().main()
except: print('ERR', end=';')
code

8. Що буде виведено за виконання?
code
def f(arg, n):
    n = 9
    tmp = arg
    tmp[:-2] = 'ac'
    return len(tmp)>3

try:
    c = [10,11,12,13,14]
    n = 3
    f(c, n)
    print(n, end=';')
    for el in c:
        print(el, end=';')
except:
    print('ERR')
code

9. Що буде виведено за виконання?
code
class A:
    b = 1
    c = 2

    def __init__(self):
        c = 3

class B(A):
    c = 4

    def __init__(self):
        super().__init__()
        self.a = 7

obj = B()
obj.c = obj.a + 10
B.c = 13

try: print(obj.c, end= ';')
except: print('ERR', end=';')

try: print(A.c, end= ';')
except: print('ERR', end=';')

try: print(B.c, end= ';')
except: print('ERR', end=';')
code

10. 
Для графа на рисунку з вершини 4 запущено алгоритм пошуку в глибину, що будує кістякове дерево. Списки суміжності вершин переглядаються алгоритмом за зростанням міток вершин.\par
\begin{center}
	\adjustimage{max size={0.9\linewidth}{0.9\paperheight}}
	{../10/Traversals/84.png}
\end{center}
\par Які з наведених ребер входять до складу отриманого кістякового дерева?\par
1) (1, 5)\par
2) (5, 6)\par
3) (3, 6)\par
4) (0, 6)\par
5) (3, 5)\par
6) (2, 6)\par
{\vskip 15pt} У формі вибрати номери відповідних ребер.\par
%answer:2 4 6
\end {large}}
%end
{\smallskip \begin {small} Затверджено на засіданні кафедри теоретичної кібернетики, протокол \No 12 від   19.05.2020 р.\par\smallskip\smallskip
{\parbox {6.5 cm}{Зав. кафедри \dotfill Ю.В.~Крак}}\hspace {0.7cm}{\hbox to 6.5 cm{ Екзаменатор \dotfill Т.О.~Карнаух}} \end{small}}
%begin
{{\begin {small}\par
{ \center  Київський національний університет імені Тараса Шевченка \par}
{\parbox {5 cm}{Освітня програма \hfill}{\parbox {7 cm}{Прикладна математика, Системний аналіз \hfill}}\parbox {6 cm}{\hfill Семестр 2}}\par
{\parbox {5 cm} {Навчальний предмет \hfill}\parbox {7 cm}{Програмування \hfill}\parbox {6cm}{\hfill}\par}\vspace{-7pt} \end{small}}{\center Екзаменаційний білет \No \textbf 
{ 77 }\par}}
{
\begin {large}
1. Що буде виведено за виконання?
code
try:
    res = 2**1.0--1
    if isinstance(res, bool):
        print(f'bool;{res}')
    elif isinstance(res, int):
        print(f'int;{res}')
    elif isinstance(res, float):
        print(f'float;{res:.2f}')
    else:
        print('???')
except:
    print('ERR')
code

2. Що буде виведено за виконання?
code
a = 0
b = 8
c = 3

def f(a):
    a = 2
    b = 3
    c = 1
    return a + b + c

a = 5
b = 6
c = 2

try:
    print(f(a), a, b, c, sep=';')
except:
    print('ERR', a, b, c, sep=';')
code

3. Що буде виведено за виконання?
code
try:
    n = 15
    while n>=6:
        n += -3
    print(n, end=';')
except:
    print('ERR')
code

4. Що буде виведено за виконання?
code
try:
    for d in range(-12, -13, -3):
        if d>=7:
            continue
        print(d, end=';')
    else:
        print(d,end=';')
    print(d)
except:
    print('ERR')
code

5. Що буде виведено за виконання?\par (Дійсні значення, якщо наявні, в цьому завданні будуть виводитись у форматі з фіксованою крапкою та, можливо, з доволі великою кількістю десяткових розрядів. У відповіді для дійсних значень у ЦЬОМУ завданні достатньо вказати один десятковий розряд після крапки, відкинувши решту розрядів без округлення. Наприклад, замість 13.66666666666667 достатньо вказати 13.6, замість 25.23 достатньо вказати 25.2)
code
try:
    print(4, end=';')
    print(int('b6'), end=';')
    print(9, end=';')
except ValueError:
    print(1, end=';')
except Exception:
    print(5, end=';')
except:
    print('ERR', end=';')
else:
    print(8, end=';')
finally:
    print(2)
code

6. Що буде виведено за виконання?
code
def f(n):
    if n>9:
        f(n//100)
    print(n%10, end='')
 
f(123456)
code

7. Що буде виведено за виконання?
code
class A:
    a = 1

    def _f1(self): print(2, end=';')
    def __f2(self): print(3, end=';')

    def main(self):
        try: print(self.a, end=';')
        except: print('ERR', end=';')

        try: self._f1()
        except: print('ERR', end=';')

        try: self.__f2()
        except: print('ERR', end=';')


class B(A):
    a = 4

    def _f1(self): print(5, end=';')
    def __f2(self): print(6, end=';')

try: B().main()
except: print('ERR', end=';')
code

8. Що буде виведено за виконання?
code
def f(arg, n):
    arg[47] = 0
    n *= 3

try:
    n = 3
    d = {29:4, 47:1, 57:5, 47:4}
    f(d, n)
    print(n, end=';')
    for x in d.items():
        print(x, sep=';', end=';')
except:
    print('ERR')
code

9. Що буде виведено за виконання?
code
class A:
    c = 1
    a = 2

    def __init__(self):
        a = 3

class B(A):
    a = 4

    def __init__(self):
        super().__init__()
        self.b = 7

obj = B()
obj.b = obj.a + 10
B.b = 13

try: print(obj.b, end= ';')
except: print('ERR', end=';')

try: print(A.b, end= ';')
except: print('ERR', end=';')

try: print(B.b, end= ';')
except: print('ERR', end=';')
code

10. 
Для графа на рисунку з вершини 0 запущено алгоритм пошуку в глибину, що будує кістякове дерево. Списки суміжності вершин переглядаються алгоритмом за зростанням міток вершин.\par
\begin{center}
	\adjustimage{max size={0.9\linewidth}{0.9\paperheight}}
	{../10/Traversals/258.png}
\end{center}
\par Які з наведених ребер входять до складу отриманого кістякового дерева?\par
1) (2, 3)\par
2) (4, 5)\par
3) (0, 3)\par
4) (2, 6)\par
5) (0, 5)\par
6) (1, 5)\par
{\vskip 15pt} У формі вибрати номери відповідних ребер.\par
%answer:1 2 3 4 6
\end {large}}
%end
{\smallskip \begin {small} Затверджено на засіданні кафедри теоретичної кібернетики, протокол \No 12 від   19.05.2020 р.\par\smallskip\smallskip
{\parbox {6.5 cm}{Зав. кафедри \dotfill Ю.В.~Крак}}\hspace {0.7cm}{\hbox to 6.5 cm{ Екзаменатор \dotfill Т.О.~Карнаух}} \end{small}}
%begin
{{\begin {small}\par
{ \center  Київський національний університет імені Тараса Шевченка \par}
{\parbox {5 cm}{Освітня програма \hfill}{\parbox {7 cm}{Прикладна математика, Системний аналіз \hfill}}\parbox {6 cm}{\hfill Семестр 2}}\par
{\parbox {5 cm} {Навчальний предмет \hfill}\parbox {7 cm}{Програмування \hfill}\parbox {6cm}{\hfill}\par}\vspace{-7pt} \end{small}}{\center Екзаменаційний білет \No \textbf 
{ 78 }\par}}
{
\begin {large}
1. Що буде виведено за виконання?
code
try:
    res = not 5.0<2 and 9.0>8 or 4<=True
    if isinstance(res, bool):
        print(f'bool;{res}')
    elif isinstance(res, int):
        print(f'int;{res}')
    elif isinstance(res, float):
        print(f'float;{res:.2f}')
    else:
        print('???')
except:
    print('ERR')
code

2. Що буде виведено за виконання?
code
a = 4
b = 7
c = 3

def h(a):
    a = 4
    b = 1
    c = 2
    return a + b + c

a = 5
b = 2
c = 6

try:
    print(h(b), a, b, c, sep=';')
except:
    print('ERR', a, b, c, sep=';')
code

3. Що буде виведено за виконання?
code
try:
    k = 5
    while k<5:
        k += 2
    print(k, end=';')
except:
    print('ERR')
code

4. Що буде виведено за виконання?
code
try:
    for c in range(-8, 7, -2):
        if c<8:
            continue
        print(c, end=';')
        if c>6:
            break
    else:
        print('end',end=';')
    print(c)
except:
    print('ERR')
code

5. Що буде виведено за виконання?\par (Дійсні значення, якщо наявні, в цьому завданні будуть виводитись у форматі з фіксованою крапкою та, можливо, з доволі великою кількістю десяткових розрядів. У відповіді для дійсних значень у ЦЬОМУ завданні достатньо вказати один десятковий розряд після крапки, відкинувши решту розрядів без округлення. Наприклад, замість 13.66666666666667 достатньо вказати 13.6, замість 25.23 достатньо вказати 25.2)
code
try:
    print(0, end=';')
    print(3//0, end=';')
    print(7, end=';')
except TypeError:
    print(1, end=';')
except ValueError:
    print(4, end=';')
except:
    print('ERR', end=';')
else:
    print(9, end=';')
finally:
    print(2)
code

6. Що буде виведено за виконання?
code
def f(n):
    if n>9:
        return f(n//10)+1
    else:
        return 1

print(f(123456))
code

7. Що буде виведено за виконання?
code
class A:
    c = 1

    def _h1(self): print(2, end=';')
    def _h2(self): print(3, end=';')

    def main(self):
        try: print(self.c, end=';')
        except: print('ERR', end=';')

        try: self._h1()
        except: print('ERR', end=';')

        try: self._h2()
        except: print('ERR', end=';')


class B(A):
    c = 4

    def _h1(self): print(5, end=';')
    def _h2(self): print(6, end=';')

try: B().main()
except: print('ERR', end=';')
code

8. Що буде виведено за виконання?
code
def f(arg, n):
    n = 2
    tmp = arg
    tmp[-8:] = 'bd'
    return len(tmp)>3

try:
    b = ['0','1','2','3','4']
    n = 4
    f(b, n)
    print(n, end=';')
    for el in b:
        print(el, end=';')
except:
    print('ERR')
code

9. Що буде виведено за виконання?
code
class A:
    b = 1
    a = 2

    def __init__(self):
        b = 3

class B(A):
    b = 4

    def __init__(self):
        super().__init__()
        self.c = 7

obj = B()
obj.b = obj.c + 10
B.b = 13

try: print(obj.b, end= ';')
except: print('ERR', end=';')

try: print(A.b, end= ';')
except: print('ERR', end=';')

try: print(B.b, end= ';')
except: print('ERR', end=';')
code

10. 
Для графа на рисунку з вершини 6 запущено алгоритм пошуку в глибину, що будує кістякове дерево. Списки суміжності вершин переглядаються алгоритмом за зростанням міток вершин.\par
\begin{center}
	\adjustimage{max size={0.9\linewidth}{0.9\paperheight}}
	{../10/Traversals/87.png}
\end{center}
\par Які з наведених ребер входять до складу отриманого кістякового дерева?\par
1) (4, 6)\par
2) (0, 1)\par
3) (1, 2)\par
4) (2, 6)\par
5) (1, 5)\par
6) (3, 6)\par
{\vskip 15pt} У формі вибрати номери відповідних ребер.\par
%answer:2 3 5
\end {large}}
%end
{\smallskip \begin {small} Затверджено на засіданні кафедри теоретичної кібернетики, протокол \No 12 від   19.05.2020 р.\par\smallskip\smallskip
{\parbox {6.5 cm}{Зав. кафедри \dotfill Ю.В.~Крак}}\hspace {0.7cm}{\hbox to 6.5 cm{ Екзаменатор \dotfill Т.О.~Карнаух}} \end{small}}
%begin
{{\begin {small}\par
{ \center  Київський національний університет імені Тараса Шевченка \par}
{\parbox {5 cm}{Освітня програма \hfill}{\parbox {7 cm}{Прикладна математика, Системний аналіз \hfill}}\parbox {6 cm}{\hfill Семестр 2}}\par
{\parbox {5 cm} {Навчальний предмет \hfill}\parbox {7 cm}{Програмування \hfill}\parbox {6cm}{\hfill}\par}\vspace{-7pt} \end{small}}{\center Екзаменаційний білет \No \textbf 
{ 79 }\par}}
{
\begin {large}
1. Що буде виведено за виконання?
code
def f(n):
    print('f', end='')
    return n>0

try:
    print(f(-8) and f(-9))
except:
    print('ERR')
code

2. Що буде виведено за виконання?
code
a = 9
b = 1
c = 7

def f(b):
    global c
    a = 5
    b = 4
    c = 3
    return a + b + c

a = 4
b = 8
c = 0

try:
    print(f(b), a, b, c, sep=';')
except:
    print('ERR', a, b, c, sep=';')
code

3. Що буде виведено за виконання?
code
try:
    t = 8
    while t<=11:
        t += 1
    print(t, end=';')
except:
    print('ERR')
code

4. Що буде виведено за виконання?
code
try:
    for e in range(7, -7, -1):
        if e>6:
            continue
        print(e, end=';')
    else:
        print(e,end=';')
    print(e)
except:
    print('ERR')
code

5. Що буде виведено за виконання?\par (Дійсні значення, якщо наявні, в цьому завданні будуть виводитись у форматі з фіксованою крапкою та, можливо, з доволі великою кількістю десяткових розрядів. У відповіді для дійсних значень у ЦЬОМУ завданні достатньо вказати один десятковий розряд після крапки, відкинувши решту розрядів без округлення. Наприклад, замість 13.66666666666667 достатньо вказати 13.6, замість 25.23 достатньо вказати 25.2)
code
try:
    print(5, end=';')
    print(0j>4j, end=';')
    print(6, end=';')
except TypeError:
    print(8, end=';')
except Exception:
    print(2, end=';')
except:
    print('ERR', end=';')
else:
    print(0, end=';')
finally:
    print(4)
code

6. Що буде виведено за виконання?
code
def f(n):
    if n>9:
        f(n//100)
        print(n%10, end='')
 
f(123456)
code

7. Що буде виведено за виконання?
code
class A:
    d = 1

    def __g1(self): print(2, end=';')
    def _g2(self): print(3, end=';')

    def main(self):
        try: print(self.d, end=';')
        except: print('ERR', end=';')

        try: self.__g1()
        except: print('ERR', end=';')

        try: self._g2()
        except: print('ERR', end=';')


class B(A):
    d = 4

    def __g1(self): print(5, end=';')
    def _g2(self): print(6, end=';')

try: B().main()
except: print('ERR', end=';')
code

8. Що буде виведено за виконання?
code
def f(arg, n):
    arg[74] = 7
    n = 9
    arg = tuple(arg.items())

try:
    n = 1
    d = {74:3, 54:5, 10:0, 74:6}
    f(d, n)
    print(n, end=';')
    for x, y in d.items():
        print(x, y, sep=';', end=';')
except:
    print('ERR')
code

9. Що буде виведено за виконання?
code
class A:
    a = 1
    b = 2

    def __init__(self):
        a = 3

class B(A):
    b = 4

    def __init__(self):
        super().__init__()
        self.c = 7

obj = B()
obj.a = obj.c + 10
B.a = 13

try: print(obj.a, end= ';')
except: print('ERR', end=';')

try: print(A.a, end= ';')
except: print('ERR', end=';')

try: print(B.a, end= ';')
except: print('ERR', end=';')
code

10. 
Для графа на рисунку з вершини 0 запущено алгоритм пошуку в глибину, що будує кістякове дерево. Списки суміжності вершин переглядаються алгоритмом за зростанням міток вершин.\par
\begin{center}
	\adjustimage{max size={0.9\linewidth}{0.9\paperheight}}
	{../10/Traversals/204.png}
\end{center}
\par Які з наведених ребер входять до складу отриманого кістякового дерева?\par
1) (1, 3)\par
2) (5, 6)\par
3) (0, 2)\par
4) (3, 6)\par
5) (3, 4)\par
6) (2, 4)\par
{\vskip 15pt} У формі вибрати номери відповідних ребер.\par
%answer:1 2 3 5 6
\end {large}}
%end
{\smallskip \begin {small} Затверджено на засіданні кафедри теоретичної кібернетики, протокол \No 12 від   19.05.2020 р.\par\smallskip\smallskip
{\parbox {6.5 cm}{Зав. кафедри \dotfill Ю.В.~Крак}}\hspace {0.7cm}{\hbox to 6.5 cm{ Екзаменатор \dotfill Т.О.~Карнаух}} \end{small}}
%begin
{{\begin {small}\par
{ \center  Київський національний університет імені Тараса Шевченка \par}
{\parbox {5 cm}{Освітня програма \hfill}{\parbox {7 cm}{Прикладна математика, Системний аналіз \hfill}}\parbox {6 cm}{\hfill Семестр 2}}\par
{\parbox {5 cm} {Навчальний предмет \hfill}\parbox {7 cm}{Програмування \hfill}\parbox {6cm}{\hfill}\par}\vspace{-7pt} \end{small}}{\center Екзаменаційний білет \No \textbf 
{ 80 }\par}}
{
\begin {large}
1. Що буде виведено за виконання?
code
try:
    res = 4/5/5*5
    if isinstance(res, bool):
        print(f'bool;{res}')
    elif isinstance(res, int):
        print(f'int;{res}')
    elif isinstance(res, float):
        print(f'float;{res:.2f}')
    else:
        print('???')
except:
    print('ERR')
code

2. Що буде виведено за виконання?
code
a = 8
b = 5
c = 4

def g(a):
    global c
    a = 4
    b += 2
    c = 3
    return a + b + c

a = 3
b = 0
c = 9

try:
    print(g(a), a, b, c, sep=';')
except:
    print('ERR', a, b, c, sep=';')
code

3. Що буде виведено за виконання?
code
try:
    i = 11
    while i>8:
        i += -2
    print(i, end=';')
except:
    print('ERR')
code

4. Що буде виведено за виконання?
code
try:
    for d in range(5, -14, 2):
        if d<=4:
            break
        print(d, end=';')
    else:
        print('end',end=';')
    print(d)
except:
    print('ERR')
code

5. Що буде виведено за виконання?\par (Дійсні значення, якщо наявні, в цьому завданні будуть виводитись у форматі з фіксованою крапкою та, можливо, з доволі великою кількістю десяткових розрядів. У відповіді для дійсних значень у ЦЬОМУ завданні достатньо вказати один десятковий розряд після крапки, відкинувши решту розрядів без округлення. Наприклад, замість 13.66666666666667 достатньо вказати 13.6, замість 25.23 достатньо вказати 25.2)
code
try:
    print(6, end=';')
    print(int('8'), end=';')
    print(5, end=';')
except KeyboardInterrupt:
    print(1, end=';')
except ZeroDivisionError:
    print(7, end=';')
except:
    print('ERR', end=';')
else:
    print(9, end=';')
finally:
    print(3)
code

6. Що буде виведено за виконання?
code
def f(n):
    if n>2:
        f(n//2)
    print(n%2, end='')
 
f(16)
code

7. Що буде виведено за виконання?
code
class A:
    b = 1

    def _f1(self): print(2, end=';')
    def f2(self): print(3, end=';')

    def main(self):
        try: print(self.b, end=';')
        except: print('ERR', end=';')

        try: self._f1()
        except: print('ERR', end=';')

        try: self.f2()
        except: print('ERR', end=';')


class B(A):
    b = 4

    def _f1(self): print(5, end=';')
    def f2(self): print(6, end=';')

try: B().main()
except: print('ERR', end=';')
code

8. Що буде виведено за виконання?
code
def f(arg, n):
    n = 7
    tmp = arg
    del tmp[3:]
    return len(tmp)>3

try:
    e = [10,11,12,13,14]
    n = 0
    f(e, n)
    print(n, end=';')
    for el in e:
        print(el, end=';')
except:
    print('ERR')
code

9. Що буде виведено за виконання?
code
class A:
    c = 1
    b = 2

    def __init__(self):
        b = 3

class B(A):
    c = 4

    def __init__(self):
        super().__init__()
        self.a = 7

obj = B()
obj.b = obj.c + 10
B.b = 13

try: print(obj.b, end= ';')
except: print('ERR', end=';')

try: print(A.b, end= ';')
except: print('ERR', end=';')

try: print(B.b, end= ';')
except: print('ERR', end=';')
code

10. 
Для графа на рисунку з вершини 5 запущено алгоритм пошуку в глибину, що будує кістякове дерево. Списки суміжності вершин переглядаються алгоритмом за зростанням міток вершин.\par
\begin{center}
	\adjustimage{max size={0.9\linewidth}{0.9\paperheight}}
	{../10/Traversals/113.png}
\end{center}
\par Які з наведених ребер входять до складу отриманого кістякового дерева?\par
1) (0, 2)\par
2) (2, 4)\par
3) (4, 6)\par
4) (3, 4)\par
5) (0, 6)\par
6) (1, 6)\par
{\vskip 15pt} У формі вибрати номери відповідних ребер.\par
%answer:1 2 6
\end {large}}
%end
{\smallskip \begin {small} Затверджено на засіданні кафедри теоретичної кібернетики, протокол \No 12 від   19.05.2020 р.\par\smallskip\smallskip
{\parbox {6.5 cm}{Зав. кафедри \dotfill Ю.В.~Крак}}\hspace {0.7cm}{\hbox to 6.5 cm{ Екзаменатор \dotfill Т.О.~Карнаух}} \end{small}}
%begin
{{\begin {small}\par
{ \center  Київський національний університет імені Тараса Шевченка \par}
{\parbox {5 cm}{Освітня програма \hfill}{\parbox {7 cm}{Прикладна математика, Системний аналіз \hfill}}\parbox {6 cm}{\hfill Семестр 2}}\par
{\parbox {5 cm} {Навчальний предмет \hfill}\parbox {7 cm}{Програмування \hfill}\parbox {6cm}{\hfill}\par}\vspace{-7pt} \end{small}}{\center Екзаменаційний білет \No \textbf 
{ 81 }\par}}
{
\begin {large}
1. Що буде виведено за виконання?
code
def f(n):
    print('f', end='')
    return n

try:
    print(f(2) or f(5))
except:
    print('ERR')
code

2. Що буде виведено за виконання?
code
a = 0
b = 1
c = 9

def f(b):
    global c
    a = 3
    b -= 5
    c = 2
    return a + b + c

a = 5
b = 7
c = 4

try:
    print(f(a), a, b, c, sep=';')
except:
    print('ERR', a, b, c, sep=';')
code

3. Що буде виведено за виконання?
code
try:
    k = 9
    while k>=2:
        k += -2
    print(k, end=';')
except:
    print('ERR')
code

4. Що буде виведено за виконання?
code
try:
    for a in range(-4, 1, -2):
        if a<5:
            continue
        print(a, end=';');a = 2
    else:
        print(a,end=';')
    print(a)
except:
    print('ERR')
code

5. Що буде виведено за виконання?\par (Дійсні значення, якщо наявні, в цьому завданні будуть виводитись у форматі з фіксованою крапкою та, можливо, з доволі великою кількістю десяткових розрядів. У відповіді для дійсних значень у ЦЬОМУ завданні достатньо вказати один десятковий розряд після крапки, відкинувши решту розрядів без округлення. Наприклад, замість 13.66666666666667 достатньо вказати 13.6, замість 25.23 достатньо вказати 25.2)
code
try:
    print(6, end=';')
    print(2%3, end=';')
    print(7, end=';')
except ZeroDivisionError:
    print(9, end=';')
except KeyboardInterrupt:
    print(4, end=';')
except:
    print('ERR', end=';')
else:
    print(3, end=';')
finally:
    print(5)
code

6. Що буде виведено за виконання?
code
def f(n):
    print(n%2, end='')
    if n>2:
        f(n//2)
 
f(16)
code

7. Що буде виведено за виконання?
code
class A:
    a = 1

    def __h1(self): print(2, end=';')
    def _h2(self): print(3, end=';')

    def main(self):
        try: print(self.a, end=';')
        except: print('ERR', end=';')

        try: self.__h1()
        except: print('ERR', end=';')

        try: self._h2()
        except: print('ERR', end=';')


class B(A):
    a = 4

    def __h1(self): print(5, end=';')
    def _h2(self): print(6, end=';')

try: B().main()
except: print('ERR', end=';')
code

8. Що буде виведено за виконання?
code
def f(arg, n):
    n = 6
    tmp = arg
    tmp[:8] = [4,5]
    return len(tmp)>3

try:
    d = ['0','1','2','3','4']
    n = 8
    f(d, n)
    print(n, end=';')
    for el in d:
        print(el, end=';')
except:
    print('ERR')
code

9. Що буде виведено за виконання?
code
class A:
    b = 1
    c = 2

    def __init__(self):
        b = 3

class B(A):
    c = 4

    def __init__(self):
        super().__init__()
        self.a = 7

obj = B()
obj.a = obj.b + 10
B.a = 13

try: print(obj.a, end= ';')
except: print('ERR', end=';')

try: print(A.a, end= ';')
except: print('ERR', end=';')

try: print(B.a, end= ';')
except: print('ERR', end=';')
code

10. 
Для графа на рисунку з вершини 2 запущено алгоритм пошуку в глибину, що будує кістякове дерево. Списки суміжності вершин переглядаються алгоритмом за зростанням міток вершин.\par
\begin{center}
	\adjustimage{max size={0.9\linewidth}{0.9\paperheight}}
	{../10/Traversals/183.png}
\end{center}
\par Які з наведених ребер входять до складу отриманого кістякового дерева?\par
1) (0, 3)\par
2) (3, 6)\par
3) (3, 4)\par
4) (1, 4)\par
5) (2, 4)\par
6) (4, 6)\par
{\vskip 15pt} У формі вибрати номери відповідних ребер.\par
%answer:1 3 4
\end {large}}
%end
{\smallskip \begin {small} Затверджено на засіданні кафедри теоретичної кібернетики, протокол \No 12 від   19.05.2020 р.\par\smallskip\smallskip
{\parbox {6.5 cm}{Зав. кафедри \dotfill Ю.В.~Крак}}\hspace {0.7cm}{\hbox to 6.5 cm{ Екзаменатор \dotfill Т.О.~Карнаух}} \end{small}}
%begin
{{\begin {small}\par
{ \center  Київський національний університет імені Тараса Шевченка \par}
{\parbox {5 cm}{Освітня програма \hfill}{\parbox {7 cm}{Прикладна математика, Системний аналіз \hfill}}\parbox {6 cm}{\hfill Семестр 2}}\par
{\parbox {5 cm} {Навчальний предмет \hfill}\parbox {7 cm}{Програмування \hfill}\parbox {6cm}{\hfill}\par}\vspace{-7pt} \end{small}}{\center Екзаменаційний білет \No \textbf 
{ 82 }\par}}
{
\begin {large}
1. Що буде виведено за виконання?
code
def f(n):
    print('f', end='')
    return n

try:
    print(f(-1) or f(-7))
except:
    print('ERR')
code

2. Що буде виведено за виконання?
code
a = 8
b = 4
c = 5

def g(a):
    global c
    a = 3
    b = 1
    c = 2
    return a + b + c

a = 6
b = 2
c = 7

try:
    print(g(a), a, b, c, sep=';')
except:
    print('ERR', a, b, c, sep=';')
code

3. Що буде виведено за виконання?
code
try:
    t = 13
    while t>9:
        t += -3
    print(t, end=';')
except:
    print('ERR')
code

4. Що буде виведено за виконання?
code
try:
    for f in range(11, -12, 3):
        if f>=3:
            continue
        print(f, end=';');f = -6
    else:
        print(f,end=';')
    print(f)
except:
    print('ERR')
code

5. Що буде виведено за виконання?\par (Дійсні значення, якщо наявні, в цьому завданні будуть виводитись у форматі з фіксованою крапкою та, можливо, з доволі великою кількістю десяткових розрядів. У відповіді для дійсних значень у ЦЬОМУ завданні достатньо вказати один десятковий розряд після крапки, відкинувши решту розрядів без округлення. Наприклад, замість 13.66666666666667 достатньо вказати 13.6, замість 25.23 достатньо вказати 25.2)
code
try:
    print(8, end=';')
    print(0%0.0, end=';')
    print(1, end=';')
except Exception:
    print(9, end=';')
except TypeError:
    print(6, end=';')
except:
    print('ERR', end=';')
else:
    print(2, end=';')
finally:
    print(4)
code

6. Що буде виведено за виконання?
code
def f(n):
    if n>9: f(n//10)
    else:
        print(n%10, end='')
 
f(123456)
code

7. Що буде виведено за виконання?
code
class A:
    c = 1

    def _h1(self): print(2, end=';')
    def _h2(self): print(3, end=';')

    def main(self):
        try: print(self.c, end=';')
        except: print('ERR', end=';')

        try: self._h1()
        except: print('ERR', end=';')

        try: self._h2()
        except: print('ERR', end=';')


class B(A):
    c = 4

    def _h1(self): print(5, end=';')
    def _h2(self): print(6, end=';')

try: B().main()
except: print('ERR', end=';')
code

8. Що буде виведено за виконання?
code
def f(arg, n):
    n = 1
    tmp = arg
    tmp[7:-6] = [2,3]
    return len(tmp)>3

try:
    a = [10,11,12,13,14]
    n = 9
    f(a, n)
    print(n, end=';')
    for el in a:
        print(el, end=';')
except:
    print('ERR')
code

9. Що буде виведено за виконання?
code
class A:
    a = 1
    c = 2

    def __init__(self):
        c = 3

class B(A):
    a = 4

    def __init__(self):
        super().__init__()
        self.b = 7

obj = B()
obj.c = obj.b + 10
B.c = 13

try: print(obj.c, end= ';')
except: print('ERR', end=';')

try: print(A.c, end= ';')
except: print('ERR', end=';')

try: print(B.c, end= ';')
except: print('ERR', end=';')
code

10. 
Для графа на рисунку з вершини 3 запущено алгоритм пошуку в глибину, що будує кістякове дерево. Списки суміжності вершин переглядаються алгоритмом за зростанням міток вершин.\par
\begin{center}
	\adjustimage{max size={0.9\linewidth}{0.9\paperheight}}
	{../10/Traversals/82.png}
\end{center}
\par Які з наведених ребер входять до складу отриманого кістякового дерева?\par
1) (0, 1)\par
2) (0, 3)\par
3) (1, 5)\par
4) (2, 3)\par
5) (0, 4)\par
6) (2, 6)\par
{\vskip 15pt} У формі вибрати номери відповідних ребер.\par
%answer:1 2 3 5 6
\end {large}}
%end
{\smallskip \begin {small} Затверджено на засіданні кафедри теоретичної кібернетики, протокол \No 12 від   19.05.2020 р.\par\smallskip\smallskip
{\parbox {6.5 cm}{Зав. кафедри \dotfill Ю.В.~Крак}}\hspace {0.7cm}{\hbox to 6.5 cm{ Екзаменатор \dotfill Т.О.~Карнаух}} \end{small}}
%begin
{{\begin {small}\par
{ \center  Київський національний університет імені Тараса Шевченка \par}
{\parbox {5 cm}{Освітня програма \hfill}{\parbox {7 cm}{Прикладна математика, Системний аналіз \hfill}}\parbox {6 cm}{\hfill Семестр 2}}\par
{\parbox {5 cm} {Навчальний предмет \hfill}\parbox {7 cm}{Програмування \hfill}\parbox {6cm}{\hfill}\par}\vspace{-7pt} \end{small}}{\center Екзаменаційний білет \No \textbf 
{ 83 }\par}}
{
\begin {large}
1. Що буде виведено за виконання?
code
try:
    res = 4**0.5*False
    if isinstance(res, bool):
        print(f'bool;{res}')
    elif isinstance(res, int):
        print(f'int;{res}')
    elif isinstance(res, float):
        print(f'float;{res:.2f}')
    else:
        print('???')
except:
    print('ERR')
code

2. Що буде виведено за виконання?
code
a = 9
b = 3
c = 0

def h(b):
    a = 5
    b = 4
    c = 1
    return a + b + c

a = 1
b = 4
c = 9

try:
    print(h(a), a, b, c, sep=';')
except:
    print('ERR', a, b, c, sep=';')
code

3. Що буде виведено за виконання?
code
try:
    i = 1
    while i<=12:
        i += 1
    print(i, end=';')
except:
    print('ERR')
code

4. Що буде виведено за виконання?
code
try:
    for c in range(13, -11, -3):
        if c>7:
            continue
        print(c, end=';')
    else:
        print(c,end=';')
    print(c)
except:
    print('ERR')
code

5. Що буде виведено за виконання?\par (Дійсні значення, якщо наявні, в цьому завданні будуть виводитись у форматі з фіксованою крапкою та, можливо, з доволі великою кількістю десяткових розрядів. У відповіді для дійсних значень у ЦЬОМУ завданні достатньо вказати один десятковий розряд після крапки, відкинувши решту розрядів без округлення. Наприклад, замість 13.66666666666667 достатньо вказати 13.6, замість 25.23 достатньо вказати 25.2)
code
try:
    print(5, end=';')
    print(int('7'), end=';')
    print(8, end=';')
except ValueError:
    print(3, end=';')
except KeyboardInterrupt:
    print(0, end=';')
except:
    print('ERR', end=';')
else:
    print(1, end=';')
finally:
    print(8)
code

6. Що буде виведено за виконання?
code
def f(n):
    if n>9:
        f(n//10)
    else:
        print(n%10, end='')

f(543216)
code

7. Що буде виведено за виконання?
code
class A:
    b = 1

    def _g1(self): print(2, end=';')
    def g2(self): print(3, end=';')

    def main(self):
        try: print(self.b, end=';')
        except: print('ERR', end=';')

        try: self._g1()
        except: print('ERR', end=';')

        try: self.g2()
        except: print('ERR', end=';')


class B(A):
    b = 4

    def _g1(self): print(5, end=';')
    def g2(self): print(6, end=';')

try: B().main()
except: print('ERR', end=';')
code

8. Що буде виведено за виконання?
code
def f(arg, n):
    arg[14] = 3
    n = 5
    arg = set(arg.values())

try:
    n = 2
    t = {63:9, 81:8, 81:9}
    f(t, n)
    print(n, end=';')
    for x in t.values():
        print(x, sep=';', end=';')
except:
    print('ERR')
code

9. Що буде виведено за виконання?
code
class A:
    b = 1
    a = 2

    def __init__(self):
        a = 3

class B(A):
    b = 4

    def __init__(self):
        super().__init__()
        self.c = 7

obj = B()
obj.a = obj.a + 10
B.a = 13

try: print(obj.a, end= ';')
except: print('ERR', end=';')

try: print(A.a, end= ';')
except: print('ERR', end=';')

try: print(B.a, end= ';')
except: print('ERR', end=';')
code

10. 
Для графа на рисунку з вершини 1 запущено алгоритм пошуку в глибину, що будує кістякове дерево. Списки суміжності вершин переглядаються алгоритмом за зростанням міток вершин.\par
\begin{center}
	\adjustimage{max size={0.9\linewidth}{0.9\paperheight}}
	{../10/Traversals/10.png}
\end{center}
\par Які з наведених ребер входять до складу отриманого кістякового дерева?\par
1) (4, 6)\par
2) (4, 5)\par
3) (5, 6)\par
4) (2, 4)\par
5) (0, 3)\par
6) (0, 2)\par
{\vskip 15pt} У формі вибрати номери відповідних ребер.\par
%answer:2 3 4 5 6
\end {large}}
%end
{\smallskip \begin {small} Затверджено на засіданні кафедри теоретичної кібернетики, протокол \No 12 від   19.05.2020 р.\par\smallskip\smallskip
{\parbox {6.5 cm}{Зав. кафедри \dotfill Ю.В.~Крак}}\hspace {0.7cm}{\hbox to 6.5 cm{ Екзаменатор \dotfill Т.О.~Карнаух}} \end{small}}
%begin
{{\begin {small}\par
{ \center  Київський національний університет імені Тараса Шевченка \par}
{\parbox {5 cm}{Освітня програма \hfill}{\parbox {7 cm}{Прикладна математика, Системний аналіз \hfill}}\parbox {6 cm}{\hfill Семестр 2}}\par
{\parbox {5 cm} {Навчальний предмет \hfill}\parbox {7 cm}{Програмування \hfill}\parbox {6cm}{\hfill}\par}\vspace{-7pt} \end{small}}{\center Екзаменаційний білет \No \textbf 
{ 84 }\par}}
{
\begin {large}
1. Що буде виведено за виконання?
code
try:
    res = "9.0">=True
    if isinstance(res, bool):
        print(f'bool;{res}')
    elif isinstance(res, int):
        print(f'int;{res}')
    elif isinstance(res, float):
        print(f'float;{res:.2f}')
    else:
        print('???')
except:
    print('ERR')
code

2. Що буде виведено за виконання?
code
a = 2
b = 8
c = 6

def h(a):
    a *= 4
    b = 1
    c = 1
    return a + b + c

a = 3
b = 8
c = 9

try:
    print(h(b), a, b, c, sep=';')
except:
    print('ERR', a, b, c, sep=';')
code

3. Що буде виведено за виконання?
code
try:
    j = 7
    while j>2:
        j += -3
    print(j, end=';')
except:
    print('ERR')
code

4. Що буде виведено за виконання?
code
try:
    for b in range(2, 6, -1):
        if b<=8:
            break
        print(b, end=';')
    else:
        print(b,end=';')
    print(b)
except:
    print('ERR')
code

5. Що буде виведено за виконання?\par (Дійсні значення, якщо наявні, в цьому завданні будуть виводитись у форматі з фіксованою крапкою та, можливо, з доволі великою кількістю десяткових розрядів. У відповіді для дійсних значень у ЦЬОМУ завданні достатньо вказати один десятковий розряд після крапки, відкинувши решту розрядів без округлення. Наприклад, замість 13.66666666666667 достатньо вказати 13.6, замість 25.23 достатньо вказати 25.2)
code
try:
    print(1, end=';')
    print(6j==3j, end=';')
    print(7, end=';')
except ValueError:
    print(4, end=';')
except TypeError:
    print(3, end=';')
except:
    print('ERR', end=';')
else:
    print(5, end=';')
finally:
    print(2)
code

6. Що буде виведено за виконання?
code
def f(n):
    if n<=9:
        print(n%10, end='')
    else:
        f(n//10)

f(123456)
code

7. Що буде виведено за виконання?
code
class A:
    d = 1

    def f1(self): print(2, end=';')
    def _f2(self): print(3, end=';')

    def main(self):
        try: print(self.d, end=';')
        except: print('ERR', end=';')

        try: self.f1()
        except: print('ERR', end=';')

        try: self._f2()
        except: print('ERR', end=';')


class B(A):
    d = 4

    def f1(self): print(5, end=';')
    def _f2(self): print(6, end=';')

try: B().main()
except: print('ERR', end=';')
code

8. Що буде виведено за виконання?
code
def f(arg, n):
    arg[18] = 5
    n += 2

try:
    n = 5
    s = {57:6, 50:0, 18:7, 73:6}
    f(s, n)
    print(n, end=';')
    for x in s :
        print(x, sep=';', end=';')
except:
    print('ERR')
code

9. Що буде виведено за виконання?
code
class A:
    c = 1
    a = 2

    def __init__(self):
        c = 3

class B(A):
    a = 4

    def __init__(self):
        super().__init__()
        self.b = 7

obj = B()
obj.b = obj.c + 10
B.b = 13

try: print(obj.b, end= ';')
except: print('ERR', end=';')

try: print(A.b, end= ';')
except: print('ERR', end=';')

try: print(B.b, end= ';')
except: print('ERR', end=';')
code

10. 
Для графа на рисунку з вершини 0 запущено алгоритм пошуку в глибину, що будує кістякове дерево. Списки суміжності вершин переглядаються алгоритмом за зростанням міток вершин.\par
\begin{center}
	\adjustimage{max size={0.9\linewidth}{0.9\paperheight}}
	{../10/Traversals/265.png}
\end{center}
\par Які з наведених ребер входять до складу отриманого кістякового дерева?\par
1) (2, 6)\par
2) (1, 4)\par
3) (0, 6)\par
4) (1, 2)\par
5) (4, 6)\par
6) (4, 5)\par
{\vskip 15pt} У формі вибрати номери відповідних ребер.\par
%answer:4 5 6
\end {large}}
%end
{\smallskip \begin {small} Затверджено на засіданні кафедри теоретичної кібернетики, протокол \No 12 від   19.05.2020 р.\par\smallskip\smallskip
{\parbox {6.5 cm}{Зав. кафедри \dotfill Ю.В.~Крак}}\hspace {0.7cm}{\hbox to 6.5 cm{ Екзаменатор \dotfill Т.О.~Карнаух}} \end{small}}
%begin
{{\begin {small}\par
{ \center  Київський національний університет імені Тараса Шевченка \par}
{\parbox {5 cm}{Освітня програма \hfill}{\parbox {7 cm}{Прикладна математика, Системний аналіз \hfill}}\parbox {6 cm}{\hfill Семестр 2}}\par
{\parbox {5 cm} {Навчальний предмет \hfill}\parbox {7 cm}{Програмування \hfill}\parbox {6cm}{\hfill}\par}\vspace{-7pt} \end{small}}{\center Екзаменаційний білет \No \textbf 
{ 85 }\par}}
{
\begin {large}
1. Що буде виведено за виконання?
code
try:
    res = 3.0>=3 and not 8==4.0 or False<=2
    if isinstance(res, bool):
        print(f'bool;{res}')
    elif isinstance(res, int):
        print(f'int;{res}')
    elif isinstance(res, float):
        print(f'float;{res:.2f}')
    else:
        print('???')
except:
    print('ERR')
code

2. Що буде виведено за виконання?
code
a = 5
b = 6
c = 2

def f(b):
    global c
    a -= 4
    b = 5
    c = 2
    return a + b + c

a = 1
b = 7
c = 0

try:
    print(f(a), a, b, c, sep=';')
except:
    print('ERR', a, b, c, sep=';')
code

3. Що буде виведено за виконання?
code
try:
    m = 10
    while m>=1:
        m += -2
    print(m, end=';')
except:
    print('ERR')
code

4. Що буде виведено за виконання?
code
try:
    for d in range(15, 13, 2):
        if d<6:
            break
        print(d, end=';');d = 7
    else:
        print(d,end=';')
    print(d)
except:
    print('ERR')
code

5. Що буде виведено за виконання?\par (Дійсні значення, якщо наявні, в цьому завданні будуть виводитись у форматі з фіксованою крапкою та, можливо, з доволі великою кількістю десяткових розрядів. У відповіді для дійсних значень у ЦЬОМУ завданні достатньо вказати один десятковий розряд після крапки, відкинувши решту розрядів без округлення. Наприклад, замість 13.66666666666667 достатньо вказати 13.6, замість 25.23 достатньо вказати 25.2)
code
try:
    print(9, end=';')
    print(int('c0'), end=';')
    print(6, end=';')
except Exception:
    print(9, end=';')
except ZeroDivisionError:
    print(0, end=';')
except:
    print('ERR', end=';')
else:
    print(1, end=';')
finally:
    print(3)
code

6. Що буде виведено за виконання?
code
def f(s):
    for i in range(len(s)):
        s[i] += 1
    return
 
s = [1, 2, 3]
f(s)
print(s)
code

7. Що буде виведено за виконання?
code
class A:
    b = 1

    def _g1(self): print(2, end=';')
    def __g2(self): print(3, end=';')

    def main(self):
        try: print(self.b, end=';')
        except: print('ERR', end=';')

        try: self._g1()
        except: print('ERR', end=';')

        try: self.__g2()
        except: print('ERR', end=';')


class B(A):
    b = 4

    def _g1(self): print(5, end=';')
    def __g2(self): print(6, end=';')

try: B().main()
except: print('ERR', end=';')
code

8. Що буде виведено за виконання?
code
def f(arg, n):
    n = 3
    tmp = arg
    del tmp[-6:]
    return len(tmp)>3

try:
    g = ['0','1','2','3','4']
    n = 5
    f(g, n)
    print(n, end=';')
    for el in g:
        print(el, end=';')
except:
    print('ERR')
code

9. Що буде виведено за виконання?
code
class A:
    c = 1
    a = 2

    def __init__(self):
        a = 3

class B(A):
    a = 4

    def __init__(self):
        super().__init__()
        self.b = 7

obj = B()
obj.a = obj.b + 10
B.a = 13

try: print(obj.a, end= ';')
except: print('ERR', end=';')

try: print(A.a, end= ';')
except: print('ERR', end=';')

try: print(B.a, end= ';')
except: print('ERR', end=';')
code

10. 
Для графа на рисунку з вершини 4 запущено алгоритм пошуку в глибину, що будує кістякове дерево. Списки суміжності вершин переглядаються алгоритмом за зростанням міток вершин.\par
\begin{center}
	\adjustimage{max size={0.9\linewidth}{0.9\paperheight}}
	{../10/Traversals/284.png}
\end{center}
\par Які з наведених ребер входять до складу отриманого кістякового дерева?\par
1) (1, 2)\par
2) (2, 5)\par
3) (3, 4)\par
4) (4, 6)\par
5) (3, 5)\par
6) (2, 4)\par
{\vskip 15pt} У формі вибрати номери відповідних ребер.\par
%answer:5 6
\end {large}}
%end
{\smallskip \begin {small} Затверджено на засіданні кафедри теоретичної кібернетики, протокол \No 12 від   19.05.2020 р.\par\smallskip\smallskip
{\parbox {6.5 cm}{Зав. кафедри \dotfill Ю.В.~Крак}}\hspace {0.7cm}{\hbox to 6.5 cm{ Екзаменатор \dotfill Т.О.~Карнаух}} \end{small}}
%begin
{{\begin {small}\par
{ \center  Київський національний університет імені Тараса Шевченка \par}
{\parbox {5 cm}{Освітня програма \hfill}{\parbox {7 cm}{Прикладна математика, Системний аналіз \hfill}}\parbox {6 cm}{\hfill Семестр 2}}\par
{\parbox {5 cm} {Навчальний предмет \hfill}\parbox {7 cm}{Програмування \hfill}\parbox {6cm}{\hfill}\par}\vspace{-7pt} \end{small}}{\center Екзаменаційний білет \No \textbf 
{ 86 }\par}}
{
\begin {large}
1. Що буде виведено за виконання?
code
try:
    res = 3//4/8*4
    if isinstance(res, bool):
        print(f'bool;{res}')
    elif isinstance(res, int):
        print(f'int;{res}')
    elif isinstance(res, float):
        print(f'float;{res:.2f}')
    else:
        print('???')
except:
    print('ERR')
code

2. Що буде виведено за виконання?
code
a = 8
b = 2
c = 7

def g(a):
    global c
    a = 4
    b = 2
    c = 3
    return a + b + c

a = 0
b = 6
c = 5

try:
    print(g(b), a, b, c, sep=';')
except:
    print('ERR', a, b, c, sep=';')
code

3. Що буде виведено за виконання?
code
try:
    k = 14
    while k>7:
        k += -2
    print(k, end=';')
except:
    print('ERR')
code

4. Що буде виведено за виконання?
code
try:
    for f in range(4, 2, 2):
        if f>5:
            continue
        print(f, end=';')
    else:
        print(f,end=';')
    print(f)
except:
    print('ERR')
code

5. Що буде виведено за виконання?\par (Дійсні значення, якщо наявні, в цьому завданні будуть виводитись у форматі з фіксованою крапкою та, можливо, з доволі великою кількістю десяткових розрядів. У відповіді для дійсних значень у ЦЬОМУ завданні достатньо вказати один десятковий розряд після крапки, відкинувши решту розрядів без округлення. Наприклад, замість 13.66666666666667 достатньо вказати 13.6, замість 25.23 достатньо вказати 25.2)
code
try:
    print(4, end=';')
    print(5j<=0j, end=';')
    print(2, end=';')
except Exception:
    print(7, end=';')
except ValueError:
    print(8, end=';')
except:
    print('ERR', end=';')
else:
    print(5, end=';')
finally:
    print(9)
code

6. Що буде виведено за виконання?
code
def f(n):
    print(n%10, end='')
    if n>9:
        f(n//100)
 
f(987654)
code

7. Що буде виведено за виконання?
code
class A:
    d = 1

    def __f1(self): print(2, end=';')
    def _f2(self): print(3, end=';')

    def main(self):
        try: print(self.d, end=';')
        except: print('ERR', end=';')

        try: self.__f1()
        except: print('ERR', end=';')

        try: self._f2()
        except: print('ERR', end=';')


class B(A):
    d = 4

    def __f1(self): print(5, end=';')
    def _f2(self): print(6, end=';')

try: B().main()
except: print('ERR', end=';')
code

8. Що буде виведено за виконання?
code
def f(arg, n):
    arg[82] = 6
    n -= -3

try:
    n = 4
    t = {55:7, 75:3, 28:6, 75:4}
    f(t, n)
    print(n, end=';')
    for x in t.values():
        print(x, sep=';', end=';')
except:
    print('ERR')
code

9. Що буде виведено за виконання?
code
class A:
    c = 1
    b = 2

    def __init__(self):
        c = 3

class B(A):
    c = 4

    def __init__(self):
        super().__init__()
        self.a = 7

obj = B()
obj.c = obj.c + 10
B.c = 13

try: print(obj.c, end= ';')
except: print('ERR', end=';')

try: print(A.c, end= ';')
except: print('ERR', end=';')

try: print(B.c, end= ';')
except: print('ERR', end=';')
code

10. 
Для графа на рисунку з вершини 1 запущено алгоритм пошуку в глибину, що будує кістякове дерево. Списки суміжності вершин переглядаються алгоритмом за зростанням міток вершин.\par
\begin{center}
	\adjustimage{max size={0.9\linewidth}{0.9\paperheight}}
	{../10/Traversals/217.png}
\end{center}
\par Які з наведених ребер входять до складу отриманого кістякового дерева?\par
1) (1, 3)\par
2) (2, 5)\par
3) (0, 3)\par
4) (0, 5)\par
5) (1, 2)\par
6) (3, 6)\par
{\vskip 15pt} У формі вибрати номери відповідних ребер.\par
%answer:2 3 4 5 6
\end {large}}
%end
{\smallskip \begin {small} Затверджено на засіданні кафедри теоретичної кібернетики, протокол \No 12 від   19.05.2020 р.\par\smallskip\smallskip
{\parbox {6.5 cm}{Зав. кафедри \dotfill Ю.В.~Крак}}\hspace {0.7cm}{\hbox to 6.5 cm{ Екзаменатор \dotfill Т.О.~Карнаух}} \end{small}}
%begin
{{\begin {small}\par
{ \center  Київський національний університет імені Тараса Шевченка \par}
{\parbox {5 cm}{Освітня програма \hfill}{\parbox {7 cm}{Прикладна математика, Системний аналіз \hfill}}\parbox {6 cm}{\hfill Семестр 2}}\par
{\parbox {5 cm} {Навчальний предмет \hfill}\parbox {7 cm}{Програмування \hfill}\parbox {6cm}{\hfill}\par}\vspace{-7pt} \end{small}}{\center Екзаменаційний білет \No \textbf 
{ 87 }\par}}
{
\begin {large}
1. Що буде виведено за виконання?
code
try:
    res = "True"==4.0j
    if isinstance(res, bool):
        print(f'bool;{res}')
    elif isinstance(res, int):
        print(f'int;{res}')
    elif isinstance(res, float):
        print(f'float;{res:.2f}')
    else:
        print('???')
except:
    print('ERR')
code

2. Що буде виведено за виконання?
code
a = 4
b = 3
c = 5

def g(a):
    a *= 3
    b = 3
    c = 5
    return a + b + c

a = 0
b = 3
c = 6

try:
    print(g(a), a, b, c, sep=';')
except:
    print('ERR', a, b, c, sep=';')
code

3. Що буде виведено за виконання?
code
try:
    m = 3
    while m<13:
        m += 2
    print(m, end=';')
except:
    print('ERR')
code

4. Що буде виведено за виконання?
code
try:
    for b in range(8, 12, 3):
        if b<=3:
            continue
        print(b, end=';')
    else:
        print(b,end=';')
    print(b)
except:
    print('ERR')
code

5. Що буде виведено за виконання?\par (Дійсні значення, якщо наявні, в цьому завданні будуть виводитись у форматі з фіксованою крапкою та, можливо, з доволі великою кількістю десяткових розрядів. У відповіді для дійсних значень у ЦЬОМУ завданні достатньо вказати один десятковий розряд після крапки, відкинувши решту розрядів без округлення. Наприклад, замість 13.66666666666667 достатньо вказати 13.6, замість 25.23 достатньо вказати 25.2)
code
try:
    print(1, end=';')
    print(int('a6'), end=';')
    print(4, end=';')
except TypeError:
    print(2, end=';')
except ZeroDivisionError:
    print(8, end=';')
except:
    print('ERR', end=';')
else:
    print(0, end=';')
finally:
    print(3)
code

6. Що буде виведено за виконання?
code
def f(n):
    if n>9: f(n//10)
    else:
        print(n%10, end='')
 
f(123456)
code

7. Що буде виведено за виконання?
code
class A:
    a = 1

    def _h1(self): print(2, end=';')
    def _h2(self): print(3, end=';')

    def main(self):
        try: print(self.a, end=';')
        except: print('ERR', end=';')

        try: self._h1()
        except: print('ERR', end=';')

        try: self._h2()
        except: print('ERR', end=';')


class B(A):
    a = 4

    def _h1(self): print(5, end=';')
    def _h2(self): print(6, end=';')

try: B().main()
except: print('ERR', end=';')
code

8. Що буде виведено за виконання?
code
def f(arg, n):
    arg[77] = 9
    n += -2

try:
    n = 6
    s = {12:7, 77:5, 46:3, 38:7}
    f(s, n)
    print(n, end=';')
    for x in s :
        print(x, sep=';', end=';')
except:
    print('ERR')
code

9. Що буде виведено за виконання?
code
class A:
    b = 1
    a = 2

    def __init__(self):
        b = 3

class B(A):
    b = 4

    def __init__(self):
        super().__init__()
        self.c = 7

obj = B()
obj.b = obj.a + 10
B.b = 13

try: print(obj.b, end= ';')
except: print('ERR', end=';')

try: print(A.b, end= ';')
except: print('ERR', end=';')

try: print(B.b, end= ';')
except: print('ERR', end=';')
code

10. 
Для графа на рисунку з вершини 1 запущено алгоритм пошуку в глибину, що будує кістякове дерево. Списки суміжності вершин переглядаються алгоритмом за зростанням міток вершин.\par
\begin{center}
	\adjustimage{max size={0.9\linewidth}{0.9\paperheight}}
	{../10/Traversals/271.png}
\end{center}
\par Які з наведених ребер входять до складу отриманого кістякового дерева?\par
1) (0, 1)\par
2) (2, 3)\par
3) (1, 5)\par
4) (5, 6)\par
5) (1, 6)\par
6) (0, 4)\par
{\vskip 15pt} У формі вибрати номери відповідних ребер.\par
%answer:1 2 4 6
\end {large}}
%end
{\smallskip \begin {small} Затверджено на засіданні кафедри теоретичної кібернетики, протокол \No 12 від   19.05.2020 р.\par\smallskip\smallskip
{\parbox {6.5 cm}{Зав. кафедри \dotfill Ю.В.~Крак}}\hspace {0.7cm}{\hbox to 6.5 cm{ Екзаменатор \dotfill Т.О.~Карнаух}} \end{small}}
%begin
{{\begin {small}\par
{ \center  Київський національний університет імені Тараса Шевченка \par}
{\parbox {5 cm}{Освітня програма \hfill}{\parbox {7 cm}{Прикладна математика, Системний аналіз \hfill}}\parbox {6 cm}{\hfill Семестр 2}}\par
{\parbox {5 cm} {Навчальний предмет \hfill}\parbox {7 cm}{Програмування \hfill}\parbox {6cm}{\hfill}\par}\vspace{-7pt} \end{small}}{\center Екзаменаційний білет \No \textbf 
{ 88 }\par}}
{
\begin {large}
1. Що буде виведено за виконання?
code
try:
    res = "3.0"<="4j"
    if isinstance(res, bool):
        print(f'bool;{res}')
    elif isinstance(res, int):
        print(f'int;{res}')
    elif isinstance(res, float):
        print(f'float;{res:.2f}')
    else:
        print('???')
except:
    print('ERR')
code

2. Що буде виведено за виконання?
code
a = 8
b = 7
c = 1

def h(a):
    a = 1
    b += 4
    c = 2
    return a + b + c

a = 2
b = 4
c = 9

try:
    print(h(b), a, b, c, sep=';')
except:
    print('ERR', a, b, c, sep=';')
code

3. Що буде виведено за виконання?
code
try:
    k = 6
    while k<=9:
        k += 2
    print(k, end=';')
except:
    print('ERR')
code

4. Що буде виведено за виконання?
code
try:
    for a in range(-1, -9, -3):
        if a>=4:
            break
        print(a, end=';')
        if a<7:
            break
    else:
        print(a,end=';')
    print(a)
except:
    print('ERR')
code

5. Що буде виведено за виконання?\par (Дійсні значення, якщо наявні, в цьому завданні будуть виводитись у форматі з фіксованою крапкою та, можливо, з доволі великою кількістю десяткових розрядів. У відповіді для дійсних значень у ЦЬОМУ завданні достатньо вказати один десятковий розряд після крапки, відкинувши решту розрядів без округлення. Наприклад, замість 13.66666666666667 достатньо вказати 13.6, замість 25.23 достатньо вказати 25.2)
code
try:
    print(7, end=';')
    print(4/2, end=';')
    print(5, end=';')
except KeyboardInterrupt:
    print(3, end=';')
except ZeroDivisionError:
    print(1, end=';')
except:
    print('ERR', end=';')
else:
    print(6, end=';')
finally:
    print(8)
code

6. Що буде виведено за виконання?
code
def f(n):
    if n>2:
        f(n//2)
    else:
        print(n%2, end='')

f(16)
code

7. Що буде виведено за виконання?
code
class A:
    c = 1

    def _g1(self): print(2, end=';')
    def g2(self): print(3, end=';')

    def main(self):
        try: print(self.c, end=';')
        except: print('ERR', end=';')

        try: self._g1()
        except: print('ERR', end=';')

        try: self.g2()
        except: print('ERR', end=';')


class B(A):
    c = 4

    def _g1(self): print(5, end=';')
    def g2(self): print(6, end=';')

try: B().main()
except: print('ERR', end=';')
code

8. Що буде виведено за виконання?
code
def f(arg, n):
    arg[26] = 3
    n *= 3

try:
    n = 3
    d = {41:0, 33:5, 33:4}
    f(d, n)
    print(n, end=';')
    for x in d.values():
        print(x, sep=';', end=';')
except:
    print('ERR')
code

9. Що буде виведено за виконання?
code
class A:
    a = 1
    c = 2

    def __init__(self):
        c = 3

class B(A):
    c = 4

    def __init__(self):
        super().__init__()
        self.b = 7

obj = B()
obj.a = obj.b + 10
B.a = 13

try: print(obj.a, end= ';')
except: print('ERR', end=';')

try: print(A.a, end= ';')
except: print('ERR', end=';')

try: print(B.a, end= ';')
except: print('ERR', end=';')
code

10. 
Для графа на рисунку з вершини 1 запущено алгоритм пошуку в глибину, що будує кістякове дерево. Списки суміжності вершин переглядаються алгоритмом за зростанням міток вершин.\par
\begin{center}
	\adjustimage{max size={0.9\linewidth}{0.9\paperheight}}
	{../10/Traversals/233.png}
\end{center}
\par Які з наведених ребер входять до складу отриманого кістякового дерева?\par
1) (0, 5)\par
2) (1, 4)\par
3) (0, 3)\par
4) (2, 4)\par
5) (1, 6)\par
6) (1, 3)\par
{\vskip 15pt} У формі вибрати номери відповідних ребер.\par
%answer:1 3 4
\end {large}}
%end
{\smallskip \begin {small} Затверджено на засіданні кафедри теоретичної кібернетики, протокол \No 12 від   19.05.2020 р.\par\smallskip\smallskip
{\parbox {6.5 cm}{Зав. кафедри \dotfill Ю.В.~Крак}}\hspace {0.7cm}{\hbox to 6.5 cm{ Екзаменатор \dotfill Т.О.~Карнаух}} \end{small}}
%begin
{{\begin {small}\par
{ \center  Київський національний університет імені Тараса Шевченка \par}
{\parbox {5 cm}{Освітня програма \hfill}{\parbox {7 cm}{Прикладна математика, Системний аналіз \hfill}}\parbox {6 cm}{\hfill Семестр 2}}\par
{\parbox {5 cm} {Навчальний предмет \hfill}\parbox {7 cm}{Програмування \hfill}\parbox {6cm}{\hfill}\par}\vspace{-7pt} \end{small}}{\center Екзаменаційний білет \No \textbf 
{ 89 }\par}}
{
\begin {large}
1. Що буде виведено за виконання?
code
try:
    res = 6/9*9*3
    if isinstance(res, bool):
        print(f'bool;{res}')
    elif isinstance(res, int):
        print(f'int;{res}')
    elif isinstance(res, float):
        print(f'float;{res:.2f}')
    else:
        print('???')
except:
    print('ERR')
code

2. Що буде виведено за виконання?
code
a = 1
b = 3
c = 4

def h(b):
    a = 5
    b = 1
    c = 3
    return a + b + c

a = 9
b = 2
c = 1

try:
    print(h(a), a, b, c, sep=';')
except:
    print('ERR', a, b, c, sep=';')
code

3. Що буде виведено за виконання?
code
try:
    t = 12
    while t>=4:
        t += -3
    print(t, end=';')
except:
    print('ERR')
code

4. Що буде виведено за виконання?
code
try:
    for e in range(9, -3, -2):
        if e>7:
            continue
        print(e, end=';');e = -10
    else:
        print(e,end=';')
    print(e)
except:
    print('ERR')
code

5. Що буде виведено за виконання?\par (Дійсні значення, якщо наявні, в цьому завданні будуть виводитись у форматі з фіксованою крапкою та, можливо, з доволі великою кількістю десяткових розрядів. У відповіді для дійсних значень у ЦЬОМУ завданні достатньо вказати один десятковий розряд після крапки, відкинувши решту розрядів без округлення. Наприклад, замість 13.66666666666667 достатньо вказати 13.6, замість 25.23 достатньо вказати 25.2)
code
try:
    print(0, end=';')
    print(int('9'), end=';')
    print(7, end=';')
except Exception:
    print(2, end=';')
except TypeError:
    print(4, end=';')
except:
    print('ERR', end=';')
else:
    print(9, end=';')
finally:
    print(7)
code

6. Що буде виведено за виконання?
code
def f(n):
    print(n%2, end='')
    if n>2:
        f(n//2)
 
f(16)
code

7. Що буде виведено за виконання?
code
class A:
    c = 1

    def _h1(self): print(2, end=';')
    def __h2(self): print(3, end=';')

    def main(self):
        try: print(self.c, end=';')
        except: print('ERR', end=';')

        try: self._h1()
        except: print('ERR', end=';')

        try: self.__h2()
        except: print('ERR', end=';')


class B(A):
    c = 4

    def _h1(self): print(5, end=';')
    def __h2(self): print(6, end=';')

try: B().main()
except: print('ERR', end=';')
code

8. Що буде виведено за виконання?
code
def f(arg, n):
    arg[67] = 5
    n = 6
    arg = list(arg.keys())

try:
    n = 4
    s = {28:5, 14:9, 14:9}
    f(s, n)
    print(n, end=';')
    for x in s.items():
        print(x, sep=';', end=';')
except:
    print('ERR')
code

9. Що буде виведено за виконання?
code
class A:
    b = 1
    c = 2

    def __init__(self):
        c = 3

class B(A):
    c = 4

    def __init__(self):
        super().__init__()
        self.a = 7

obj = B()
obj.c = obj.a + 10
B.c = 13

try: print(obj.c, end= ';')
except: print('ERR', end=';')

try: print(A.c, end= ';')
except: print('ERR', end=';')

try: print(B.c, end= ';')
except: print('ERR', end=';')
code

10. 
Для графа на рисунку з вершини 5 запущено алгоритм пошуку в глибину, що будує кістякове дерево. Списки суміжності вершин переглядаються алгоритмом за зростанням міток вершин.\par
\begin{center}
	\adjustimage{max size={0.9\linewidth}{0.9\paperheight}}
	{../10/Traversals/277.png}
\end{center}
\par Які з наведених ребер входять до складу отриманого кістякового дерева?\par
1) (5, 6)\par
2) (3, 6)\par
3) (3, 5)\par
4) (2, 3)\par
5) (1, 3)\par
6) (4, 6)\par
{\vskip 15pt} У формі вибрати номери відповідних ребер.\par
%answer:3 5 6
\end {large}}
%end
{\smallskip \begin {small} Затверджено на засіданні кафедри теоретичної кібернетики, протокол \No 12 від   19.05.2020 р.\par\smallskip\smallskip
{\parbox {6.5 cm}{Зав. кафедри \dotfill Ю.В.~Крак}}\hspace {0.7cm}{\hbox to 6.5 cm{ Екзаменатор \dotfill Т.О.~Карнаух}} \end{small}}
%begin
{{\begin {small}\par
{ \center  Київський національний університет імені Тараса Шевченка \par}
{\parbox {5 cm}{Освітня програма \hfill}{\parbox {7 cm}{Прикладна математика, Системний аналіз \hfill}}\parbox {6 cm}{\hfill Семестр 2}}\par
{\parbox {5 cm} {Навчальний предмет \hfill}\parbox {7 cm}{Програмування \hfill}\parbox {6cm}{\hfill}\par}\vspace{-7pt} \end{small}}{\center Екзаменаційний білет \No \textbf 
{ 90 }\par}}
{
\begin {large}
1. Що буде виведено за виконання?
code
try:
    res = -2**2.0+-2
    if isinstance(res, bool):
        print(f'bool;{res}')
    elif isinstance(res, int):
        print(f'int;{res}')
    elif isinstance(res, float):
        print(f'float;{res:.2f}')
    else:
        print('???')
except:
    print('ERR')
code

2. Що буде виведено за виконання?
code
a = 3
b = 6
c = 5

def f(b):
    global c
    a = 4
    b = 2
    c = 5
    return a + b + c

a = 8
b = 0
c = 7

try:
    print(f(a), a, b, c, sep=';')
except:
    print('ERR', a, b, c, sep=';')
code

3. Що буде виведено за виконання?
code
try:
    j = 7
    while j<10:
        j += 1
    print(j, end=';')
except:
    print('ERR')
code

4. Що буде виведено за виконання?
code
try:
    for c in range(-10, 4, -1):
        if c>=8:
            continue
        print(c, end=';')
    else:
        print('end',end=';')
    print(c)
except:
    print('ERR')
code

5. Що буде виведено за виконання?\par (Дійсні значення, якщо наявні, в цьому завданні будуть виводитись у форматі з фіксованою крапкою та, можливо, з доволі великою кількістю десяткових розрядів. У відповіді для дійсних значень у ЦЬОМУ завданні достатньо вказати один десятковий розряд після крапки, відкинувши решту розрядів без округлення. Наприклад, замість 13.66666666666667 достатньо вказати 13.6, замість 25.23 достатньо вказати 25.2)
code
try:
    print(8, end=';')
    print(2//0, end=';')
    print(0, end=';')
except ValueError:
    print(3, end=';')
except KeyboardInterrupt:
    print(6, end=';')
except:
    print('ERR', end=';')
else:
    print(1, end=';')
finally:
    print(5)
code

6. Що буде виведено за виконання?
code
def f(n):
    if n>9:
        return f(n//100)+1
    else:
        return 1

print(f(123456))
code

7. Що буде виведено за виконання?
code
class A:
    d = 1

    def f1(self): print(2, end=';')
    def _f2(self): print(3, end=';')

    def main(self):
        try: print(self.d, end=';')
        except: print('ERR', end=';')

        try: self.f1()
        except: print('ERR', end=';')

        try: self._f2()
        except: print('ERR', end=';')


class B(A):
    d = 4

    def f1(self): print(5, end=';')
    def _f2(self): print(6, end=';')

try: B().main()
except: print('ERR', end=';')
code

8. Що буде виведено за виконання?
code
def f(arg, n):
    arg[49] = 9
    n -= -3

try:
    n = 7
    t = {16:1, 43:4, 14:2, 28:5}
    f(t, n)
    print(n, end=';')
    for x in t.keys():
        print(x, sep=';', end=';')
except:
    print('ERR')
code

9. Що буде виведено за виконання?
code
class A:
    a = 1
    b = 2

    def __init__(self):
        a = 3

class B(A):
    a = 4

    def __init__(self):
        super().__init__()
        self.c = 7

obj = B()
obj.c = obj.b + 10
B.c = 13

try: print(obj.c, end= ';')
except: print('ERR', end=';')

try: print(A.c, end= ';')
except: print('ERR', end=';')

try: print(B.c, end= ';')
except: print('ERR', end=';')
code

10. 
Для графа на рисунку з вершини 2 запущено алгоритм пошуку в глибину, що будує кістякове дерево. Списки суміжності вершин переглядаються алгоритмом за зростанням міток вершин.\par
\begin{center}
	\adjustimage{max size={0.9\linewidth}{0.9\paperheight}}
	{../10/Traversals/101.png}
\end{center}
\par Які з наведених ребер входять до складу отриманого кістякового дерева?\par
1) (2, 3)\par
2) (1, 3)\par
3) (2, 4)\par
4) (3, 6)\par
5) (3, 4)\par
6) (1, 4)\par
{\vskip 15pt} У формі вибрати номери відповідних ребер.\par
%answer:1 6
\end {large}}
%end
{\smallskip \begin {small} Затверджено на засіданні кафедри теоретичної кібернетики, протокол \No 12 від   19.05.2020 р.\par\smallskip\smallskip
{\parbox {6.5 cm}{Зав. кафедри \dotfill Ю.В.~Крак}}\hspace {0.7cm}{\hbox to 6.5 cm{ Екзаменатор \dotfill Т.О.~Карнаух}} \end{small}}
%begin
{{\begin {small}\par
{ \center  Київський національний університет імені Тараса Шевченка \par}
{\parbox {5 cm}{Освітня програма \hfill}{\parbox {7 cm}{Прикладна математика, Системний аналіз \hfill}}\parbox {6 cm}{\hfill Семестр 2}}\par
{\parbox {5 cm} {Навчальний предмет \hfill}\parbox {7 cm}{Програмування \hfill}\parbox {6cm}{\hfill}\par}\vspace{-7pt} \end{small}}{\center Екзаменаційний білет \No \textbf 
{ 91 }\par}}
{
\begin {large}
1. Що буде виведено за виконання?
code
try:
    res = 7!=7.0 or not 5.0<6 and True>9
    if isinstance(res, bool):
        print(f'bool;{res}')
    elif isinstance(res, int):
        print(f'int;{res}')
    elif isinstance(res, float):
        print(f'float;{res:.2f}')
    else:
        print('???')
except:
    print('ERR')
code

2. Що буде виведено за виконання?
code
a = 1
b = 4
c = 9

def g(b):
    global c
    a = 5
    b += 1
    c = 4
    return a + b + c

a = 5
b = 8
c = 7

try:
    print(g(b), a, b, c, sep=';')
except:
    print('ERR', a, b, c, sep=';')
code

3. Що буде виведено за виконання?
code
try:
    m = 14
    while m>=0:
        m += -2
    print(m, end=';')
except:
    print('ERR')
code

4. Що буде виведено за виконання?
code
try:
    for c in range(0, -6, -3):
        if c<4:
            continue
        print(c, end=';')
    else:
        print(c,end=';')
    print(c)
except:
    print('ERR')
code

5. Що буде виведено за виконання?\par (Дійсні значення, якщо наявні, в цьому завданні будуть виводитись у форматі з фіксованою крапкою та, можливо, з доволі великою кількістю десяткових розрядів. У відповіді для дійсних значень у ЦЬОМУ завданні достатньо вказати один десятковий розряд після крапки, відкинувши решту розрядів без округлення. Наприклад, замість 13.66666666666667 достатньо вказати 13.6, замість 25.23 достатньо вказати 25.2)
code
try:
    print(5, end=';')
    print(7//2, end=';')
    print(4, end=';')
except Exception:
    print(8, end=';')
except ValueError:
    print(2, end=';')
except:
    print('ERR', end=';')
else:
    print(1, end=';')
finally:
    print(6)
code

6. Що буде виведено за виконання?
code
def f(n):
    if n>9:
        f(n//10)
    else:
        print(n%10, end='')

f(543216)
code

7. Що буде виведено за виконання?
code
class A:
    a = 1

    def __h1(self): print(2, end=';')
    def _h2(self): print(3, end=';')

    def main(self):
        try: print(self.a, end=';')
        except: print('ERR', end=';')

        try: self.__h1()
        except: print('ERR', end=';')

        try: self._h2()
        except: print('ERR', end=';')


class B(A):
    a = 4

    def __h1(self): print(5, end=';')
    def _h2(self): print(6, end=';')

try: B().main()
except: print('ERR', end=';')
code

8. Що буде виведено за виконання?
code
def f(arg, n):
    arg[17] = 2
    n = 8
    arg = tuple(arg.items())

try:
    n = 0
    d = {58:3, 48:2, 23:7, 48:8}
    f(d, n)
    print(n, end=';')
    for x in d.values():
        print(x, sep=';', end=';')
except:
    print('ERR')
code

9. Що буде виведено за виконання?
code
class A:
    c = 1
    b = 2

    def __init__(self):
        b = 3

class B(A):
    c = 4

    def __init__(self):
        super().__init__()
        self.a = 7

obj = B()
obj.b = obj.c + 10
B.b = 13

try: print(obj.b, end= ';')
except: print('ERR', end=';')

try: print(A.b, end= ';')
except: print('ERR', end=';')

try: print(B.b, end= ';')
except: print('ERR', end=';')
code

10. 
Для графа на рисунку з вершини 1 запущено алгоритм пошуку в глибину, що будує кістякове дерево. Списки суміжності вершин переглядаються алгоритмом за зростанням міток вершин.\par
\begin{center}
	\adjustimage{max size={0.9\linewidth}{0.9\paperheight}}
	{../10/Traversals/246.png}
\end{center}
\par Які з наведених ребер входять до складу отриманого кістякового дерева?\par
1) (3, 4)\par
2) (2, 6)\par
3) (3, 5)\par
4) (1, 6)\par
5) (2, 5)\par
6) (1, 4)\par
{\vskip 15pt} У формі вибрати номери відповідних ребер.\par
%answer:2 5
\end {large}}
%end
{\smallskip \begin {small} Затверджено на засіданні кафедри теоретичної кібернетики, протокол \No 12 від   19.05.2020 р.\par\smallskip\smallskip
{\parbox {6.5 cm}{Зав. кафедри \dotfill Ю.В.~Крак}}\hspace {0.7cm}{\hbox to 6.5 cm{ Екзаменатор \dotfill Т.О.~Карнаух}} \end{small}}
%begin
{{\begin {small}\par
{ \center  Київський національний університет імені Тараса Шевченка \par}
{\parbox {5 cm}{Освітня програма \hfill}{\parbox {7 cm}{Прикладна математика, Системний аналіз \hfill}}\parbox {6 cm}{\hfill Семестр 2}}\par
{\parbox {5 cm} {Навчальний предмет \hfill}\parbox {7 cm}{Програмування \hfill}\parbox {6cm}{\hfill}\par}\vspace{-7pt} \end{small}}{\center Екзаменаційний білет \No \textbf 
{ 92 }\par}}
{
\begin {large}
1. Що буде виведено за виконання?
code
def f(n):
    print('f', end='')
    return n<-4

try:
    print(f(2) and f(6))
except:
    print('ERR')
code

2. Що буде виведено за виконання?
code
a = 6
b = 0
c = 2

def h(b):
    a = 3
    b -= 2
    c = 4
    return a + b + c

a = 3
b = 3
c = 0

try:
    print(h(a), a, b, c, sep=';')
except:
    print('ERR', a, b, c, sep=';')
code

3. Що буде виведено за виконання?
code
try:
    n = 15
    while n>3:
        n += -3
    print(n, end=';')
except:
    print('ERR')
code

4. Що буде виведено за виконання?
code
try:
    for f in range(1, 0, -2):
        if f<=8:
            continue
        print(f, end=';');f = -5
        if f>=8:
            break
    else:
        print(f,end=';')
    print(f)
except:
    print('ERR')
code

5. Що буде виведено за виконання?\par (Дійсні значення, якщо наявні, в цьому завданні будуть виводитись у форматі з фіксованою крапкою та, можливо, з доволі великою кількістю десяткових розрядів. У відповіді для дійсних значень у ЦЬОМУ завданні достатньо вказати один десятковий розряд після крапки, відкинувши решту розрядів без округлення. Наприклад, замість 13.66666666666667 достатньо вказати 13.6, замість 25.23 достатньо вказати 25.2)
code
try:
    print(3, end=';')
    print(0/0, end=';')
    print(9, end=';')
except TypeError:
    print(1, end=';')
except ZeroDivisionError:
    print(2, end=';')
except:
    print('ERR', end=';')
else:
    print(7, end=';')
finally:
    print(0)
code

6. Що буде виведено за виконання?
code
def f(n):
    print(n%10, end='')
    if n>9:
        f(n//10)
 
f(543216)
code

7. Що буде виведено за виконання?
code
class A:
    b = 1

    def g1(self): print(2, end=';')
    def _g2(self): print(3, end=';')

    def main(self):
        try: print(self.b, end=';')
        except: print('ERR', end=';')

        try: self.g1()
        except: print('ERR', end=';')

        try: self._g2()
        except: print('ERR', end=';')


class B(A):
    b = 4

    def g1(self): print(5, end=';')
    def _g2(self): print(6, end=';')

try: B().main()
except: print('ERR', end=';')
code

8. Що буде виведено за виконання?
code
def f(arg, n):
    arg[25] = 2
    n += -2

try:
    n = 5
    d = {33:7, 54:0, 33:7}
    f(d, n)
    print(n, end=';')
    for x in d :
        print(x, sep=';', end=';')
except:
    print('ERR')
code

9. Що буде виведено за виконання?
code
class A:
    b = 1
    a = 2

    def __init__(self):
        b = 3

class B(A):
    a = 4

    def __init__(self):
        super().__init__()
        self.c = 7

obj = B()
obj.a = obj.a + 10
B.a = 13

try: print(obj.a, end= ';')
except: print('ERR', end=';')

try: print(A.a, end= ';')
except: print('ERR', end=';')

try: print(B.a, end= ';')
except: print('ERR', end=';')
code

10. 
Для графа на рисунку з вершини 0 запущено алгоритм пошуку в глибину, що будує кістякове дерево. Списки суміжності вершин переглядаються алгоритмом за зростанням міток вершин.\par
\begin{center}
	\adjustimage{max size={0.9\linewidth}{0.9\paperheight}}
	{../10/Traversals/130.png}
\end{center}
\par Які з наведених ребер входять до складу отриманого кістякового дерева?\par
1) (1, 5)\par
2) (1, 2)\par
3) (1, 4)\par
4) (3, 4)\par
5) (0, 4)\par
6) (2, 5)\par
{\vskip 15pt} У формі вибрати номери відповідних ребер.\par
%answer:2 3 4 6
\end {large}}
%end
{\smallskip \begin {small} Затверджено на засіданні кафедри теоретичної кібернетики, протокол \No 12 від   19.05.2020 р.\par\smallskip\smallskip
{\parbox {6.5 cm}{Зав. кафедри \dotfill Ю.В.~Крак}}\hspace {0.7cm}{\hbox to 6.5 cm{ Екзаменатор \dotfill Т.О.~Карнаух}} \end{small}}
%begin
{{\begin {small}\par
{ \center  Київський національний університет імені Тараса Шевченка \par}
{\parbox {5 cm}{Освітня програма \hfill}{\parbox {7 cm}{Прикладна математика, Системний аналіз \hfill}}\parbox {6 cm}{\hfill Семестр 2}}\par
{\parbox {5 cm} {Навчальний предмет \hfill}\parbox {7 cm}{Програмування \hfill}\parbox {6cm}{\hfill}\par}\vspace{-7pt} \end{small}}{\center Екзаменаційний білет \No \textbf 
{ 93 }\par}}
{
\begin {large}
1. Що буде виведено за виконання?
code
try:
    res = not 8.0<4 or False>5 and 2.0!=6
    if isinstance(res, bool):
        print(f'bool;{res}')
    elif isinstance(res, int):
        print(f'int;{res}')
    elif isinstance(res, float):
        print(f'float;{res:.2f}')
    else:
        print('???')
except:
    print('ERR')
code

2. Що буде виведено за виконання?
code
a = 1
b = 7
c = 4

def f(a):
    global c
    a *= 3
    b = 2
    c = 1
    return a + b + c

a = 5
b = 8
c = 2

try:
    print(f(b), a, b, c, sep=';')
except:
    print('ERR', a, b, c, sep=';')
code

3. Що буде виведено за виконання?
code
try:
    i = 6
    while i>6:
        i += -3
    print(i, end=';')
except:
    print('ERR')
code

4. Що буде виведено за виконання?
code
try:
    for e in range(-15, 11, -1):
        if e>6:
            break
        print(e, end=';')
    else:
        print(e,end=';')
    print(e)
except:
    print('ERR')
code

5. Що буде виведено за виконання?\par (Дійсні значення, якщо наявні, в цьому завданні будуть виводитись у форматі з фіксованою крапкою та, можливо, з доволі великою кількістю десяткових розрядів. У відповіді для дійсних значень у ЦЬОМУ завданні достатньо вказати один десятковий розряд після крапки, відкинувши решту розрядів без округлення. Наприклад, замість 13.66666666666667 достатньо вказати 13.6, замість 25.23 достатньо вказати 25.2)
code
try:
    print(5, end=';')
    print(int('9'), end=';')
    print(4, end=';')
except KeyboardInterrupt:
    print(6, end=';')
except KeyboardInterrupt:
    print(3, end=';')
except:
    print('ERR', end=';')
else:
    print(8, end=';')
finally:
    print(3)
code

6. Що буде виведено за виконання?
code
def f(n):
    if n>9:
        f(n//100)
    else:
        print(n%10,end='')

f(987654)
code

7. Що буде виведено за виконання?
code
class A:
    b = 1

    def _f1(self): print(2, end=';')
    def _f2(self): print(3, end=';')

    def main(self):
        try: print(self.b, end=';')
        except: print('ERR', end=';')

        try: self._f1()
        except: print('ERR', end=';')

        try: self._f2()
        except: print('ERR', end=';')


class B(A):
    b = 4

    def _f1(self): print(5, end=';')
    def _f2(self): print(6, end=';')

try: B().main()
except: print('ERR', end=';')
code

8. Що буде виведено за виконання?
code
def f(arg, n):
    arg[12] = 0
    n = 9
    arg = list(arg.values())

try:
    n = 1
    t = {58:7, 43:4, 15:6, 15:5}
    f(t, n)
    print(n, end=';')
    for x in t.items():
        print(*x, sep=';', end=';')
except:
    print('ERR')
code

9. Що буде виведено за виконання?
code
class A:
    a = 1
    c = 2

    def __init__(self):
        a = 3

class B(A):
    a = 4

    def __init__(self):
        super().__init__()
        self.b = 7

obj = B()
obj.c = obj.b + 10
B.c = 13

try: print(obj.c, end= ';')
except: print('ERR', end=';')

try: print(A.c, end= ';')
except: print('ERR', end=';')

try: print(B.c, end= ';')
except: print('ERR', end=';')
code

10. 
Для графа на рисунку з вершини 4 запущено алгоритм пошуку в глибину, що будує кістякове дерево. Списки суміжності вершин переглядаються алгоритмом за зростанням міток вершин.\par
\begin{center}
	\adjustimage{max size={0.9\linewidth}{0.9\paperheight}}
	{../10/Traversals/238.png}
\end{center}
\par Які з наведених ребер входять до складу отриманого кістякового дерева?\par
1) (0, 6)\par
2) (0, 3)\par
3) (4, 5)\par
4) (0, 4)\par
5) (3, 4)\par
6) (2, 6)\par
{\vskip 15pt} У формі вибрати номери відповідних ребер.\par
%answer:2 4 6
\end {large}}
%end
{\smallskip \begin {small} Затверджено на засіданні кафедри теоретичної кібернетики, протокол \No 12 від   19.05.2020 р.\par\smallskip\smallskip
{\parbox {6.5 cm}{Зав. кафедри \dotfill Ю.В.~Крак}}\hspace {0.7cm}{\hbox to 6.5 cm{ Екзаменатор \dotfill Т.О.~Карнаух}} \end{small}}
%begin
{{\begin {small}\par
{ \center  Київський національний університет імені Тараса Шевченка \par}
{\parbox {5 cm}{Освітня програма \hfill}{\parbox {7 cm}{Прикладна математика, Системний аналіз \hfill}}\parbox {6 cm}{\hfill Семестр 2}}\par
{\parbox {5 cm} {Навчальний предмет \hfill}\parbox {7 cm}{Програмування \hfill}\parbox {6cm}{\hfill}\par}\vspace{-7pt} \end{small}}{\center Екзаменаційний білет \No \textbf 
{ 94 }\par}}
{
\begin {large}
1. Що буде виведено за виконання?
code
try:
    res = 3**True-1
    if isinstance(res, bool):
        print(f'bool;{res}')
    elif isinstance(res, int):
        print(f'int;{res}')
    elif isinstance(res, float):
        print(f'float;{res:.2f}')
    else:
        print('???')
except:
    print('ERR')
code

2. Що буде виведено за виконання?
code
a = 6
b = 9
c = 5

def f(b):
    global c
    a *= 5
    b = 3
    c = 2
    return a + b + c

a = 9
b = 4
c = 8

try:
    print(f(a), a, b, c, sep=';')
except:
    print('ERR', a, b, c, sep=';')
code

3. Що буде виведено за виконання?
code
try:
    t = 2
    while t<=8:
        t += 2
    print(t, end=';')
except:
    print('ERR')
code

4. Що буде виведено за виконання?
code
try:
    for b in range(-8, 7, 3):
        if b<5:
            break
        print(b, end=';');b = 9
    else:
        print(b,end=';')
    print(b)
except:
    print('ERR')
code

5. Що буде виведено за виконання?\par (Дійсні значення, якщо наявні, в цьому завданні будуть виводитись у форматі з фіксованою крапкою та, можливо, з доволі великою кількістю десяткових розрядів. У відповіді для дійсних значень у ЦЬОМУ завданні достатньо вказати один десятковий розряд після крапки, відкинувши решту розрядів без округлення. Наприклад, замість 13.66666666666667 достатньо вказати 13.6, замість 25.23 достатньо вказати 25.2)
code
try:
    print(8, end=';')
    print(int('d9'), end=';')
    print(6, end=';')
except ZeroDivisionError:
    print(4, end=';')
except ValueError:
    print(1, end=';')
except:
    print('ERR', end=';')
else:
    print(2, end=';')
finally:
    print(7)
code

6. Що буде виведено за виконання?
code
def f(n):
    if n<=9:
        print(n%10, end='')
    else:
        f(n//10)

f(123456)
code

7. Що буде виведено за виконання?
code
class A:
    c = 1

    def _g1(self): print(2, end=';')
    def __g2(self): print(3, end=';')

    def main(self):
        try: print(self.c, end=';')
        except: print('ERR', end=';')

        try: self._g1()
        except: print('ERR', end=';')

        try: self.__g2()
        except: print('ERR', end=';')


class B(A):
    c = 4

    def _g1(self): print(5, end=';')
    def __g2(self): print(6, end=';')

try: B().main()
except: print('ERR', end=';')
code

8. Що буде виведено за виконання?
code
def f(arg, n):
    n = 2
    tmp = arg
    tmp[:-4] = 'bc'
    return len(tmp)>3

try:
    f = ['0','1','2','3','4']
    n = 0
    f(f, n)
    print(n, end=';')
    for el in f:
        print(el, end=';')
except:
    print('ERR')
code

9. Що буде виведено за виконання?
code
class A:
    c = 1
    a = 2

    def __init__(self):
        a = 3

class B(A):
    a = 4

    def __init__(self):
        super().__init__()
        self.b = 7

obj = B()
obj.a = obj.b + 10
B.a = 13

try: print(obj.a, end= ';')
except: print('ERR', end=';')

try: print(A.a, end= ';')
except: print('ERR', end=';')

try: print(B.a, end= ';')
except: print('ERR', end=';')
code

10. 
Для графа на рисунку з вершини 4 запущено алгоритм пошуку в глибину, що будує кістякове дерево. Списки суміжності вершин переглядаються алгоритмом за зростанням міток вершин.\par
\begin{center}
	\adjustimage{max size={0.9\linewidth}{0.9\paperheight}}
	{../10/Traversals/210.png}
\end{center}
\par Які з наведених ребер входять до складу отриманого кістякового дерева?\par
1) (3, 4)\par
2) (3, 6)\par
3) (0, 1)\par
4) (0, 2)\par
5) (2, 5)\par
6) (0, 4)\par
{\vskip 15pt} У формі вибрати номери відповідних ребер.\par
%answer:2 3 5 6
\end {large}}
%end
{\smallskip \begin {small} Затверджено на засіданні кафедри теоретичної кібернетики, протокол \No 12 від   19.05.2020 р.\par\smallskip\smallskip
{\parbox {6.5 cm}{Зав. кафедри \dotfill Ю.В.~Крак}}\hspace {0.7cm}{\hbox to 6.5 cm{ Екзаменатор \dotfill Т.О.~Карнаух}} \end{small}}
%begin
{{\begin {small}\par
{ \center  Київський національний університет імені Тараса Шевченка \par}
{\parbox {5 cm}{Освітня програма \hfill}{\parbox {7 cm}{Прикладна математика, Системний аналіз \hfill}}\parbox {6 cm}{\hfill Семестр 2}}\par
{\parbox {5 cm} {Навчальний предмет \hfill}\parbox {7 cm}{Програмування \hfill}\parbox {6cm}{\hfill}\par}\vspace{-7pt} \end{small}}{\center Екзаменаційний білет \No \textbf 
{ 95 }\par}}
{
\begin {large}
1. Що буде виведено за виконання?
code
try:
    res = 1**1.0*True
    if isinstance(res, bool):
        print(f'bool;{res}')
    elif isinstance(res, int):
        print(f'int;{res}')
    elif isinstance(res, float):
        print(f'float;{res:.2f}')
    else:
        print('???')
except:
    print('ERR')
code

2. Що буде виведено за виконання?
code
a = 2
b = 1
c = 0

def h(a):
    a = 2
    b = 4
    c = 1
    return a + b + c

a = 6
b = 9
c = 8

try:
    print(h(b), a, b, c, sep=';')
except:
    print('ERR', a, b, c, sep=';')
code

3. Що буде виведено за виконання?
code
try:
    j = 11
    while j>=8:
        j += -2
    print(j, end=';')
except:
    print('ERR')
code

4. Що буде виведено за виконання?
code
try:
    for d in range(-1, 10, 2):
        if d>=7:
            continue
        print(d, end=';')
    else:
        print(d,end=';')
    print(d)
except:
    print('ERR')
code

5. Що буде виведено за виконання?\par (Дійсні значення, якщо наявні, в цьому завданні будуть виводитись у форматі з фіксованою крапкою та, можливо, з доволі великою кількістю десяткових розрядів. У відповіді для дійсних значень у ЦЬОМУ завданні достатньо вказати один десятковий розряд після крапки, відкинувши решту розрядів без округлення. Наприклад, замість 13.66666666666667 достатньо вказати 13.6, замість 25.23 достатньо вказати 25.2)
code
try:
    print(0, end=';')
    print(5j==7j, end=';')
    print(8, end=';')
except TypeError:
    print(3, end=';')
except Exception:
    print(6, end=';')
except:
    print('ERR', end=';')
else:
    print(4, end=';')
finally:
    print(0)
code

6. Що буде виведено за виконання?
code
def f(n):
    if n>9:
        f(n//10)
    print(n%10, end='')
 
f(123456)
code

7. Що буде виведено за виконання?
code
class A:
    d = 1

    def _f1(self): print(2, end=';')
    def f2(self): print(3, end=';')

    def main(self):
        try: print(self.d, end=';')
        except: print('ERR', end=';')

        try: self._f1()
        except: print('ERR', end=';')

        try: self.f2()
        except: print('ERR', end=';')


class B(A):
    d = 4

    def _f1(self): print(5, end=';')
    def f2(self): print(6, end=';')

try: B().main()
except: print('ERR', end=';')
code

8. Що буде виведено за виконання?
code
def f(arg, n):
    arg[16] = 4
    n = 7
    arg = set(arg.keys())

try:
    n = 2
    s = {16:3, 74:9, 74:1}
    f(s, n)
    print(n, end=';')
    for x in s.keys():
        print(x, sep=';', end=';')
except:
    print('ERR')
code

9. Що буде виведено за виконання?
code
class A:
    b = 1
    c = 2

    def __init__(self):
        c = 3

class B(A):
    c = 4

    def __init__(self):
        super().__init__()
        self.a = 7

obj = B()
obj.c = obj.b + 10
B.c = 13

try: print(obj.c, end= ';')
except: print('ERR', end=';')

try: print(A.c, end= ';')
except: print('ERR', end=';')

try: print(B.c, end= ';')
except: print('ERR', end=';')
code

10. 
Для графа на рисунку з вершини 2 запущено алгоритм пошуку в глибину, що будує кістякове дерево. Списки суміжності вершин переглядаються алгоритмом за зростанням міток вершин.\par
\begin{center}
	\adjustimage{max size={0.9\linewidth}{0.9\paperheight}}
	{../10/Traversals/191.png}
\end{center}
\par Які з наведених ребер входять до складу отриманого кістякового дерева?\par
1) (5, 6)\par
2) (2, 5)\par
3) (0, 3)\par
4) (3, 4)\par
5) (1, 3)\par
6) (4, 6)\par
{\vskip 15pt} У формі вибрати номери відповідних ребер.\par
%answer:1 5 6
\end {large}}
%end
{\smallskip \begin {small} Затверджено на засіданні кафедри теоретичної кібернетики, протокол \No 12 від   19.05.2020 р.\par\smallskip\smallskip
{\parbox {6.5 cm}{Зав. кафедри \dotfill Ю.В.~Крак}}\hspace {0.7cm}{\hbox to 6.5 cm{ Екзаменатор \dotfill Т.О.~Карнаух}} \end{small}}
%begin
{{\begin {small}\par
{ \center  Київський національний університет імені Тараса Шевченка \par}
{\parbox {5 cm}{Освітня програма \hfill}{\parbox {7 cm}{Прикладна математика, Системний аналіз \hfill}}\parbox {6 cm}{\hfill Семестр 2}}\par
{\parbox {5 cm} {Навчальний предмет \hfill}\parbox {7 cm}{Програмування \hfill}\parbox {6cm}{\hfill}\par}\vspace{-7pt} \end{small}}{\center Екзаменаційний білет \No \textbf 
{ 96 }\par}}
{
\begin {large}
1. Що буде виведено за виконання?
code
try:
    res = -1**True+-1
    if isinstance(res, bool):
        print(f'bool;{res}')
    elif isinstance(res, int):
        print(f'int;{res}')
    elif isinstance(res, float):
        print(f'float;{res:.2f}')
    else:
        print('???')
except:
    print('ERR')
code

2. Що буде виведено за виконання?
code
a = 7
b = 5
c = 4

def f(a):
    global c
    a = 3
    b = 5
    c = 5
    return a + b + c

a = 3
b = 8
c = 3

try:
    print(f(b), a, b, c, sep=';')
except:
    print('ERR', a, b, c, sep=';')
code

3. Що буде виведено за виконання?
code
try:
    n = 4
    while n<12:
        n += 1
    print(n, end=';')
except:
    print('ERR')
code

4. Що буде виведено за виконання?
code
try:
    for a in range(-5, -9, 3):
        if a<=3:
            continue
        print(a, end=';')
    else:
        print('end',end=';')
    print(a)
except:
    print('ERR')
code

5. Що буде виведено за виконання?\par (Дійсні значення, якщо наявні, в цьому завданні будуть виводитись у форматі з фіксованою крапкою та, можливо, з доволі великою кількістю десяткових розрядів. У відповіді для дійсних значень у ЦЬОМУ завданні достатньо вказати один десятковий розряд після крапки, відкинувши решту розрядів без округлення. Наприклад, замість 13.66666666666667 достатньо вказати 13.6, замість 25.23 достатньо вказати 25.2)
code
try:
    print(7, end=';')
    print(9%1, end=';')
    print(2, end=';')
except KeyboardInterrupt:
    print(5, end=';')
except ValueError:
    print(1, end=';')
except:
    print('ERR', end=';')
else:
    print(2, end=';')
finally:
    print(1)
code

6. Що буде виведено за виконання?
code
def f(n):
    if n<=9:
        print(n%10,end='')
    else:
        f(n//10)

f(543216)
code

7. Що буде виведено за виконання?
code
class A:
    a = 1

    def _h1(self): print(2, end=';')
    def __h2(self): print(3, end=';')

    def main(self):
        try: print(self.a, end=';')
        except: print('ERR', end=';')

        try: self._h1()
        except: print('ERR', end=';')

        try: self.__h2()
        except: print('ERR', end=';')


class B(A):
    a = 4

    def _h1(self): print(5, end=';')
    def __h2(self): print(6, end=';')

try: B().main()
except: print('ERR', end=';')
code

8. Що буде виведено за виконання?
code
def f(arg, n):
    arg[61] = 5
    n -= 2

try:
    n = 7
    s = {40:8, 44:8, 44:4}
    f(s, n)
    print(n, end=';')
    for x in s.values():
        print(x, sep=';', end=';')
except:
    print('ERR')
code

9. Що буде виведено за виконання?
code
class A:
    a = 1
    b = 2

    def __init__(self):
        a = 3

class B(A):
    a = 4

    def __init__(self):
        super().__init__()
        self.c = 7

obj = B()
obj.a = obj.c + 10
B.a = 13

try: print(obj.a, end= ';')
except: print('ERR', end=';')

try: print(A.a, end= ';')
except: print('ERR', end=';')

try: print(B.a, end= ';')
except: print('ERR', end=';')
code

10. 
Для графа на рисунку з вершини 3 запущено алгоритм пошуку в глибину, що будує кістякове дерево. Списки суміжності вершин переглядаються алгоритмом за зростанням міток вершин.\par
\begin{center}
	\adjustimage{max size={0.9\linewidth}{0.9\paperheight}}
	{../10/Traversals/91.png}
\end{center}
\par Які з наведених ребер входять до складу отриманого кістякового дерева?\par
1) (2, 5)\par
2) (0, 3)\par
3) (4, 5)\par
4) (2, 4)\par
5) (0, 1)\par
6) (5, 6)\par
{\vskip 15pt} У формі вибрати номери відповідних ребер.\par
%answer:2 3 4 5 6
\end {large}}
%end
{\smallskip \begin {small} Затверджено на засіданні кафедри теоретичної кібернетики, протокол \No 12 від   19.05.2020 р.\par\smallskip\smallskip
{\parbox {6.5 cm}{Зав. кафедри \dotfill Ю.В.~Крак}}\hspace {0.7cm}{\hbox to 6.5 cm{ Екзаменатор \dotfill Т.О.~Карнаух}} \end{small}}
%begin
{{\begin {small}\par
{ \center  Київський національний університет імені Тараса Шевченка \par}
{\parbox {5 cm}{Освітня програма \hfill}{\parbox {7 cm}{Прикладна математика, Системний аналіз \hfill}}\parbox {6 cm}{\hfill Семестр 2}}\par
{\parbox {5 cm} {Навчальний предмет \hfill}\parbox {7 cm}{Програмування \hfill}\parbox {6cm}{\hfill}\par}\vspace{-7pt} \end{small}}{\center Екзаменаційний білет \No \textbf 
{ 97 }\par}}
{
\begin {large}
1. Що буде виведено за виконання?
code
try:
    res = 7<=9.0 and not True>=5 or 4==6.0
    if isinstance(res, bool):
        print(f'bool;{res}')
    elif isinstance(res, int):
        print(f'int;{res}')
    elif isinstance(res, float):
        print(f'float;{res:.2f}')
    else:
        print('???')
except:
    print('ERR')
code

2. Що буде виведено за виконання?
code
a = 6
b = 9
c = 2

def g(b):
    a = 4
    b = 1
    c = 3
    return a + b + c

a = 1
b = 5
c = 0

try:
    print(g(a), a, b, c, sep=';')
except:
    print('ERR', a, b, c, sep=';')
code

3. Що буде виведено за виконання?
code
try:
    m = 9
    while m<=9:
        m += 1
    print(m, end=';')
except:
    print('ERR')
code

4. Що буде виведено за виконання?
code
try:
    for b in range(5, 12, -3):
        if b<5:
            continue
        print(b, end=';');b = 3
    else:
        print(b,end=';')
    print(b)
except:
    print('ERR')
code

5. Що буде виведено за виконання?\par (Дійсні значення, якщо наявні, в цьому завданні будуть виводитись у форматі з фіксованою крапкою та, можливо, з доволі великою кількістю десяткових розрядів. У відповіді для дійсних значень у ЦЬОМУ завданні достатньо вказати один десятковий розряд після крапки, відкинувши решту розрядів без округлення. Наприклад, замість 13.66666666666667 достатньо вказати 13.6, замість 25.23 достатньо вказати 25.2)
code
try:
    print(9, end=';')
    print(8%0.0, end=';')
    print(5, end=';')
except ZeroDivisionError:
    print(4, end=';')
except TypeError:
    print(7, end=';')
except:
    print('ERR', end=';')
else:
    print(6, end=';')
finally:
    print(3)
code

6. Що буде виведено за виконання?
code
def f(n):
    print(n%10, end='')
    if n>9:
        f(n//10)
 
f(123456)
code

7. Що буде виведено за виконання?
code
class A:
    d = 1

    def _f1(self): print(2, end=';')
    def f2(self): print(3, end=';')

    def main(self):
        try: print(self.d, end=';')
        except: print('ERR', end=';')

        try: self._f1()
        except: print('ERR', end=';')

        try: self.f2()
        except: print('ERR', end=';')


class B(A):
    d = 4

    def _f1(self): print(5, end=';')
    def f2(self): print(6, end=';')

try: B().main()
except: print('ERR', end=';')
code

8. Що буде виведено за виконання?
code
def f(arg, n):
    arg[61] = 2
    n *= -3

try:
    n = 4
    t = {63:1, 61:8, 63:0}
    f(t, n)
    print(n, end=';')
    for x in t.keys():
        print(x, sep=';', end=';')
except:
    print('ERR')
code

9. Що буде виведено за виконання?
code
class A:
    a = 1
    c = 2

    def __init__(self):
        a = 3

class B(A):
    a = 4

    def __init__(self):
        super().__init__()
        self.b = 7

obj = B()
obj.b = obj.a + 10
B.b = 13

try: print(obj.b, end= ';')
except: print('ERR', end=';')

try: print(A.b, end= ';')
except: print('ERR', end=';')

try: print(B.b, end= ';')
except: print('ERR', end=';')
code

10. 
Для графа на рисунку з вершини 0 запущено алгоритм пошуку в глибину, що будує кістякове дерево. Списки суміжності вершин переглядаються алгоритмом за зростанням міток вершин.\par
\begin{center}
	\adjustimage{max size={0.9\linewidth}{0.9\paperheight}}
	{../10/Traversals/165.png}
\end{center}
\par Які з наведених ребер входять до складу отриманого кістякового дерева?\par
1) (4, 6)\par
2) (0, 4)\par
3) (2, 4)\par
4) (2, 3)\par
5) (5, 6)\par
6) (1, 2)\par
{\vskip 15pt} У формі вибрати номери відповідних ребер.\par
%answer:2 3 4 5 6
\end {large}}
%end
{\smallskip \begin {small} Затверджено на засіданні кафедри теоретичної кібернетики, протокол \No 12 від   19.05.2020 р.\par\smallskip\smallskip
{\parbox {6.5 cm}{Зав. кафедри \dotfill Ю.В.~Крак}}\hspace {0.7cm}{\hbox to 6.5 cm{ Екзаменатор \dotfill Т.О.~Карнаух}} \end{small}}
%begin
{{\begin {small}\par
{ \center  Київський національний університет імені Тараса Шевченка \par}
{\parbox {5 cm}{Освітня програма \hfill}{\parbox {7 cm}{Прикладна математика, Системний аналіз \hfill}}\parbox {6 cm}{\hfill Семестр 2}}\par
{\parbox {5 cm} {Навчальний предмет \hfill}\parbox {7 cm}{Програмування \hfill}\parbox {6cm}{\hfill}\par}\vspace{-7pt} \end{small}}{\center Екзаменаційний білет \No \textbf 
{ 98 }\par}}
{
\begin {large}
1. Що буде виведено за виконання?
code
try:
    res = 8/8//6*6
    if isinstance(res, bool):
        print(f'bool;{res}')
    elif isinstance(res, int):
        print(f'int;{res}')
    elif isinstance(res, float):
        print(f'float;{res:.2f}')
    else:
        print('???')
except:
    print('ERR')
code

2. Що буде виведено за виконання?
code
a = 7
b = 4
c = 3

def f(a):
    a = 2
    b = 5
    c = 3
    return a + b + c

a = 7
b = 6
c = 2

try:
    print(f(b), a, b, c, sep=';')
except:
    print('ERR', a, b, c, sep=';')
code

3. Що буде виведено за виконання?
code
try:
    t = 7
    while t>9:
        t += -3
    print(t, end=';')
except:
    print('ERR')
code

4. Що буде виведено за виконання?
code
try:
    for d in range(6, 2, -2):
        if d>=4:
            break
        print(d, end=';')
    else:
        print(d,end=';')
    print(d)
except:
    print('ERR')
code

5. Що буде виведено за виконання?\par (Дійсні значення, якщо наявні, в цьому завданні будуть виводитись у форматі з фіксованою крапкою та, можливо, з доволі великою кількістю десяткових розрядів. У відповіді для дійсних значень у ЦЬОМУ завданні достатньо вказати один десятковий розряд після крапки, відкинувши решту розрядів без округлення. Наприклад, замість 13.66666666666667 достатньо вказати 13.6, замість 25.23 достатньо вказати 25.2)
code
try:
    print(0, end=';')
    print(int('7'), end=';')
    print(3, end=';')
except Exception:
    print(2, end=';')
except TypeError:
    print(6, end=';')
except:
    print('ERR', end=';')
else:
    print(9, end=';')
finally:
    print(5)
code

6. Що буде виведено за виконання?
code
def f(n):
    if n>9:
        f(n//100)
    print(n%10, end='')
 
f(123456)
code

7. Що буде виведено за виконання?
code
class A:
    c = 1

    def _g1(self): print(2, end=';')
    def _g2(self): print(3, end=';')

    def main(self):
        try: print(self.c, end=';')
        except: print('ERR', end=';')

        try: self._g1()
        except: print('ERR', end=';')

        try: self._g2()
        except: print('ERR', end=';')


class B(A):
    c = 4

    def _g1(self): print(5, end=';')
    def _g2(self): print(6, end=';')

try: B().main()
except: print('ERR', end=';')
code

8. Що буде виведено за виконання?
code
def f(arg, n):
    arg[35] = 5
    n = 6
    arg = tuple(arg.items())

try:
    n = 4
    d = {44:8, 84:3, 53:4}
    f(d, n)
    print(n, end=';')
    for x in d.keys():
        print(x, sep=';', end=';')
except:
    print('ERR')
code

9. Що буде виведено за виконання?
code
class A:
    a = 1
    b = 2

    def __init__(self):
        b = 3

class B(A):
    b = 4

    def __init__(self):
        super().__init__()
        self.c = 7

obj = B()
obj.c = obj.b + 10
B.c = 13

try: print(obj.c, end= ';')
except: print('ERR', end=';')

try: print(A.c, end= ';')
except: print('ERR', end=';')

try: print(B.c, end= ';')
except: print('ERR', end=';')
code

10. 
Для графа на рисунку з вершини 0 запущено алгоритм пошуку в глибину, що будує кістякове дерево. Списки суміжності вершин переглядаються алгоритмом за зростанням міток вершин.\par
\begin{center}
	\adjustimage{max size={0.9\linewidth}{0.9\paperheight}}
	{../10/Traversals/200.png}
\end{center}
\par Які з наведених ребер входять до складу отриманого кістякового дерева?\par
1) (5, 6)\par
2) (2, 4)\par
3) (1, 2)\par
4) (1, 6)\par
5) (2, 6)\par
6) (0, 4)\par
{\vskip 15pt} У формі вибрати номери відповідних ребер.\par
%answer:1 2 3
\end {large}}
%end
{\smallskip \begin {small} Затверджено на засіданні кафедри теоретичної кібернетики, протокол \No 12 від   19.05.2020 р.\par\smallskip\smallskip
{\parbox {6.5 cm}{Зав. кафедри \dotfill Ю.В.~Крак}}\hspace {0.7cm}{\hbox to 6.5 cm{ Екзаменатор \dotfill Т.О.~Карнаух}} \end{small}}
%begin
{{\begin {small}\par
{ \center  Київський національний університет імені Тараса Шевченка \par}
{\parbox {5 cm}{Освітня програма \hfill}{\parbox {7 cm}{Прикладна математика, Системний аналіз \hfill}}\parbox {6 cm}{\hfill Семестр 2}}\par
{\parbox {5 cm} {Навчальний предмет \hfill}\parbox {7 cm}{Програмування \hfill}\parbox {6cm}{\hfill}\par}\vspace{-7pt} \end{small}}{\center Екзаменаційний білет \No \textbf 
{ 99 }\par}}
{
\begin {large}
1. Що буде виведено за виконання?
code
try:
    res = -1**-4+2
    if isinstance(res, bool):
        print(f'bool;{res}')
    elif isinstance(res, int):
        print(f'int;{res}')
    elif isinstance(res, float):
        print(f'float;{res:.2f}')
    else:
        print('???')
except:
    print('ERR')
code

2. Що буде виведено за виконання?
code
a = 3
b = 2
c = 0

def g(a):
    a = 4
    b -= 1
    c = 5
    return a + b + c

a = 7
b = 1
c = 6

try:
    print(g(a), a, b, c, sep=';')
except:
    print('ERR', a, b, c, sep=';')
code

3. Що буде виведено за виконання?
code
try:
    t = 0
    while t<7:
        t += 2
    print(t, end=';')
except:
    print('ERR')
code

4. Що буде виведено за виконання?
code
try:
    for a in range(-7, 9, -1):
        if a>7:
            continue
        print(a, end=';')
    else:
        print(a,end=';')
    print(a)
except:
    print('ERR')
code

5. Що буде виведено за виконання?\par (Дійсні значення, якщо наявні, в цьому завданні будуть виводитись у форматі з фіксованою крапкою та, можливо, з доволі великою кількістю десяткових розрядів. У відповіді для дійсних значень у ЦЬОМУ завданні достатньо вказати один десятковий розряд після крапки, відкинувши решту розрядів без округлення. Наприклад, замість 13.66666666666667 достатньо вказати 13.6, замість 25.23 достатньо вказати 25.2)
code
try:
    print(0, end=';')
    print(1j<=1j, end=';')
    print(8, end=';')
except KeyboardInterrupt:
    print(4, end=';')
except ZeroDivisionError:
    print(4, end=';')
except:
    print('ERR', end=';')
else:
    print(1, end=';')
finally:
    print(3)
code

6. Що буде виведено за виконання?
code
def f(n):
    if n>9:
        f(n//100)
        print(n%10, end='')
 
f(123456)
code

7. Що буде виведено за виконання?
code
class A:
    b = 1

    def __h1(self): print(2, end=';')
    def _h2(self): print(3, end=';')

    def main(self):
        try: print(self.b, end=';')
        except: print('ERR', end=';')

        try: self.__h1()
        except: print('ERR', end=';')

        try: self._h2()
        except: print('ERR', end=';')


class B(A):
    b = 4

    def __h1(self): print(5, end=';')
    def _h2(self): print(6, end=';')

try: B().main()
except: print('ERR', end=';')
code

8. Що буде виведено за виконання?
code
def f(arg, n):
    arg[76] = 1
    n += -2

try:
    n = 3
    t = {13:5, 76:4, 76:4}
    f(t, n)
    print(n, end=';')
    for x, y in t.items():
        print(x, y, sep=';', end=';')
except:
    print('ERR')
code

9. Що буде виведено за виконання?
code
class A:
    c = 1
    a = 2

    def __init__(self):
        a = 3

class B(A):
    a = 4

    def __init__(self):
        super().__init__()
        self.b = 7

obj = B()
obj.a = obj.c + 10
B.a = 13

try: print(obj.a, end= ';')
except: print('ERR', end=';')

try: print(A.a, end= ';')
except: print('ERR', end=';')

try: print(B.a, end= ';')
except: print('ERR', end=';')
code

10. 
Для графа на рисунку з вершини 5 запущено алгоритм пошуку в глибину, що будує кістякове дерево. Списки суміжності вершин переглядаються алгоритмом за зростанням міток вершин.\par
\begin{center}
	\adjustimage{max size={0.9\linewidth}{0.9\paperheight}}
	{../10/Traversals/97.png}
\end{center}
\par Які з наведених ребер входять до складу отриманого кістякового дерева?\par
1) (0, 4)\par
2) (3, 4)\par
3) (2, 4)\par
4) (3, 6)\par
5) (0, 5)\par
6) (0, 1)\par
{\vskip 15pt} У формі вибрати номери відповідних ребер.\par
%answer:2 3 4 5 6
\end {large}}
%end
{\smallskip \begin {small} Затверджено на засіданні кафедри теоретичної кібернетики, протокол \No 12 від   19.05.2020 р.\par\smallskip\smallskip
{\parbox {6.5 cm}{Зав. кафедри \dotfill Ю.В.~Крак}}\hspace {0.7cm}{\hbox to 6.5 cm{ Екзаменатор \dotfill Т.О.~Карнаух}} \end{small}}
%begin
{{\begin {small}\par
{ \center  Київський національний університет імені Тараса Шевченка \par}
{\parbox {5 cm}{Освітня програма \hfill}{\parbox {7 cm}{Прикладна математика, Системний аналіз \hfill}}\parbox {6 cm}{\hfill Семестр 2}}\par
{\parbox {5 cm} {Навчальний предмет \hfill}\parbox {7 cm}{Програмування \hfill}\parbox {6cm}{\hfill}\par}\vspace{-7pt} \end{small}}{\center Екзаменаційний білет \No \textbf 
{ 100 }\par}}
{
\begin {large}
1. Що буде виведено за виконання?
code
try:
    res = 9//6%7*8
    if isinstance(res, bool):
        print(f'bool;{res}')
    elif isinstance(res, int):
        print(f'int;{res}')
    elif isinstance(res, float):
        print(f'float;{res:.2f}')
    else:
        print('???')
except:
    print('ERR')
code

2. Що буде виведено за виконання?
code
a = 1
b = 9
c = 8

def h(b):
    global c
    a = 2
    b = 4
    c = 1
    return a + b + c

a = 0
b = 5
c = 4

try:
    print(h(a), a, b, c, sep=';')
except:
    print('ERR', a, b, c, sep=';')
code

3. Що буде виведено за виконання?
code
try:
    n = 8
    while n<=13:
        n += 2
    print(n, end=';')
except:
    print('ERR')
code

4. Що буде виведено за виконання?
code
try:
    for e in range(-12, -7, 2):
        if e<=3:
            continue
        print(e, end=';');e = 4
    else:
        print(e,end=';')
    print(e)
except:
    print('ERR')
code

5. Що буде виведено за виконання?\par (Дійсні значення, якщо наявні, в цьому завданні будуть виводитись у форматі з фіксованою крапкою та, можливо, з доволі великою кількістю десяткових розрядів. У відповіді для дійсних значень у ЦЬОМУ завданні достатньо вказати один десятковий розряд після крапки, відкинувши решту розрядів без округлення. Наприклад, замість 13.66666666666667 достатньо вказати 13.6, замість 25.23 достатньо вказати 25.2)
code
try:
    print(7, end=';')
    print(int('b6'), end=';')
    print(5, end=';')
except Exception:
    print(0, end=';')
except ValueError:
    print(8, end=';')
except:
    print('ERR', end=';')
else:
    print(2, end=';')
finally:
    print(9)
code

6. Що буде виведено за виконання?
code
def f(s):
    s1 = s
    for i in range(len(s)):
        s1[i] += 1
        return s1

s = [1, 2, 3]
f(s)
print(s)
code

7. Що буде виведено за виконання?
code
class A:
    a = 1

    def h1(self): print(2, end=';')
    def _h2(self): print(3, end=';')

    def main(self):
        try: print(self.a, end=';')
        except: print('ERR', end=';')

        try: self.h1()
        except: print('ERR', end=';')

        try: self._h2()
        except: print('ERR', end=';')


class B(A):
    a = 4

    def h1(self): print(5, end=';')
    def _h2(self): print(6, end=';')

try: B().main()
except: print('ERR', end=';')
code

8. Що буде виведено за виконання?
code
def f(arg, n):
    arg[31] = 7
    n -= 2

try:
    n = 6
    d = {49:4, 85:9, 85:1}
    f(d, n)
    print(n, end=';')
    for x in d.items():
        print(x, sep=';', end=';')
except:
    print('ERR')
code

9. Що буде виведено за виконання?
code
class A:
    c = 1
    b = 2

    def __init__(self):
        c = 3

class B(A):
    c = 4

    def __init__(self):
        super().__init__()
        self.a = 7

obj = B()
obj.c = obj.a + 10
B.c = 13

try: print(obj.c, end= ';')
except: print('ERR', end=';')

try: print(A.c, end= ';')
except: print('ERR', end=';')

try: print(B.c, end= ';')
except: print('ERR', end=';')
code

10. 
Для графа на рисунку з вершини 5 запущено алгоритм пошуку в глибину, що будує кістякове дерево. Списки суміжності вершин переглядаються алгоритмом за зростанням міток вершин.\par
\begin{center}
	\adjustimage{max size={0.9\linewidth}{0.9\paperheight}}
	{../10/Traversals/175.png}
\end{center}
\par Які з наведених ребер входять до складу отриманого кістякового дерева?\par
1) (0, 5)\par
2) (5, 6)\par
3) (3, 6)\par
4) (0, 3)\par
5) (4, 5)\par
6) (1, 5)\par
{\vskip 15pt} У формі вибрати номери відповідних ребер.\par
%answer:1 3 4
\end {large}}
%end
{\smallskip \begin {small} Затверджено на засіданні кафедри теоретичної кібернетики, протокол \No 12 від   19.05.2020 р.\par\smallskip\smallskip
{\parbox {6.5 cm}{Зав. кафедри \dotfill Ю.В.~Крак}}\hspace {0.7cm}{\hbox to 6.5 cm{ Екзаменатор \dotfill Т.О.~Карнаух}} \end{small}}
%begin
{{\begin {small}\par
{ \center  Київський національний університет імені Тараса Шевченка \par}
{\parbox {5 cm}{Освітня програма \hfill}{\parbox {7 cm}{Прикладна математика, Системний аналіз \hfill}}\parbox {6 cm}{\hfill Семестр 2}}\par
{\parbox {5 cm} {Навчальний предмет \hfill}\parbox {7 cm}{Програмування \hfill}\parbox {6cm}{\hfill}\par}\vspace{-7pt} \end{small}}{\center Екзаменаційний білет \No \textbf 
{ 101 }\par}}
{
\begin {large}
1. Що буде виведено за виконання?
code
try:
    res = 5/7//3*9
    if isinstance(res, bool):
        print(f'bool;{res}')
    elif isinstance(res, int):
        print(f'int;{res}')
    elif isinstance(res, float):
        print(f'float;{res:.2f}')
    else:
        print('???')
except:
    print('ERR')
code

2. Що буде виведено за виконання?
code
a = 6
b = 1
c = 8

def g(a):
    a = 5
    b = 4
    c = 2
    return a + b + c

a = 0
b = 2
c = 3

try:
    print(g(b), a, b, c, sep=';')
except:
    print('ERR', a, b, c, sep=';')
code

3. Що буде виведено за виконання?
code
try:
    m = 9
    while m>=5:
        m += -2
    print(m, end=';')
except:
    print('ERR')
code

4. Що буде виведено за виконання?
code
try:
    for f in range(-13, -15, -2):
        if f<6:
            break
        print(f, end=';')
        if f>=8:
            break
    else:
        print(f,end=';')
    print(f)
except:
    print('ERR')
code

5. Що буде виведено за виконання?\par (Дійсні значення, якщо наявні, в цьому завданні будуть виводитись у форматі з фіксованою крапкою та, можливо, з доволі великою кількістю десяткових розрядів. У відповіді для дійсних значень у ЦЬОМУ завданні достатньо вказати один десятковий розряд після крапки, відкинувши решту розрядів без округлення. Наприклад, замість 13.66666666666667 достатньо вказати 13.6, замість 25.23 достатньо вказати 25.2)
code
try:
    print(1, end=';')
    print(5/0, end=';')
    print(6, end=';')
except ValueError:
    print(9, end=';')
except KeyboardInterrupt:
    print(7, end=';')
except:
    print('ERR', end=';')
else:
    print(3, end=';')
finally:
    print(2)
code

6. Що буде виведено за виконання?
code
def f(n):
    if n>9:
        f(n//100)
    else:
        print(n%10,end='')

f(123456)
code

7. Що буде виведено за виконання?
code
class A:
    c = 1

    def __f1(self): print(2, end=';')
    def _f2(self): print(3, end=';')

    def main(self):
        try: print(self.c, end=';')
        except: print('ERR', end=';')

        try: self.__f1()
        except: print('ERR', end=';')

        try: self._f2()
        except: print('ERR', end=';')


class B(A):
    c = 4

    def __f1(self): print(5, end=';')
    def _f2(self): print(6, end=';')

try: B().main()
except: print('ERR', end=';')
code

8. Що буде виведено за виконання?
code
def f(arg, n):
    arg[19] = 7
    n *= 3

try:
    n = 7
    s = {79:6, 35:8, 55:6}
    f(s, n)
    print(n, end=';')
    for x in s.items():
        print(*x, sep=';', end=';')
except:
    print('ERR')
code

9. Що буде виведено за виконання?
code
class A:
    b = 1
    a = 2

    def __init__(self):
        a = 3

class B(A):
    a = 4

    def __init__(self):
        super().__init__()
        self.c = 7

obj = B()
obj.b = obj.a + 10
B.b = 13

try: print(obj.b, end= ';')
except: print('ERR', end=';')

try: print(A.b, end= ';')
except: print('ERR', end=';')

try: print(B.b, end= ';')
except: print('ERR', end=';')
code

10. 
Для графа на рисунку з вершини 6 запущено алгоритм пошуку в глибину, що будує кістякове дерево. Списки суміжності вершин переглядаються алгоритмом за зростанням міток вершин.\par
\begin{center}
	\adjustimage{max size={0.9\linewidth}{0.9\paperheight}}
	{../10/Traversals/213.png}
\end{center}
\par Які з наведених ребер входять до складу отриманого кістякового дерева?\par
1) (0, 3)\par
2) (0, 1)\par
3) (1, 2)\par
4) (1, 6)\par
5) (3, 6)\par
6) (2, 5)\par
{\vskip 15pt} У формі вибрати номери відповідних ребер.\par
%answer:2 3
\end {large}}
%end
{\smallskip \begin {small} Затверджено на засіданні кафедри теоретичної кібернетики, протокол \No 12 від   19.05.2020 р.\par\smallskip\smallskip
{\parbox {6.5 cm}{Зав. кафедри \dotfill Ю.В.~Крак}}\hspace {0.7cm}{\hbox to 6.5 cm{ Екзаменатор \dotfill Т.О.~Карнаух}} \end{small}}
%begin
{{\begin {small}\par
{ \center  Київський національний університет імені Тараса Шевченка \par}
{\parbox {5 cm}{Освітня програма \hfill}{\parbox {7 cm}{Прикладна математика, Системний аналіз \hfill}}\parbox {6 cm}{\hfill Семестр 2}}\par
{\parbox {5 cm} {Навчальний предмет \hfill}\parbox {7 cm}{Програмування \hfill}\parbox {6cm}{\hfill}\par}\vspace{-7pt} \end{small}}{\center Екзаменаційний білет \No \textbf 
{ 102 }\par}}
{
\begin {large}
1. Що буде виведено за виконання?
code
try:
    res = 9//7%3*6
    if isinstance(res, bool):
        print(f'bool;{res}')
    elif isinstance(res, int):
        print(f'int;{res}')
    elif isinstance(res, float):
        print(f'float;{res:.2f}')
    else:
        print('???')
except:
    print('ERR')
code

2. Що буде виведено за виконання?
code
a = 2
b = 6
c = 3

def h(b):
    global c
    a += 3
    b = 1
    c = 2
    return a + b + c

a = 7
b = 1
c = 5

try:
    print(h(b), a, b, c, sep=';')
except:
    print('ERR', a, b, c, sep=';')
code

3. Що буде виведено за виконання?
code
try:
    n = 8
    while n>=3:
        n += -2
    print(n, end=';')
except:
    print('ERR')
code

4. Що буде виведено за виконання?
code
try:
    for c in range(9, -11, -1):
        if c<=8:
            continue
        print(c, end=';')
    else:
        print('end',end=';')
    print(c)
except:
    print('ERR')
code

5. Що буде виведено за виконання?\par (Дійсні значення, якщо наявні, в цьому завданні будуть виводитись у форматі з фіксованою крапкою та, можливо, з доволі великою кількістю десяткових розрядів. У відповіді для дійсних значень у ЦЬОМУ завданні достатньо вказати один десятковий розряд після крапки, відкинувши решту розрядів без округлення. Наприклад, замість 13.66666666666667 достатньо вказати 13.6, замість 25.23 достатньо вказати 25.2)
code
try:
    print(0, end=';')
    print(4j<3j, end=';')
    print(8, end=';')
except ZeroDivisionError:
    print(6, end=';')
except TypeError:
    print(7, end=';')
except:
    print('ERR', end=';')
else:
    print(3, end=';')
finally:
    print(2)
code

6. Що буде виведено за виконання?
code
def f(n):
    if n>9:
        return f(n//10)+1
    else:
        return 1

print(f(123456))
code

7. Що буде виведено за виконання?
code
class A:
    a = 1

    def g1(self): print(2, end=';')
    def _g2(self): print(3, end=';')

    def main(self):
        try: print(self.a, end=';')
        except: print('ERR', end=';')

        try: self.g1()
        except: print('ERR', end=';')

        try: self._g2()
        except: print('ERR', end=';')


class B(A):
    a = 4

    def g1(self): print(5, end=';')
    def _g2(self): print(6, end=';')

try: B().main()
except: print('ERR', end=';')
code

8. Що буде виведено за виконання?
code
def f(arg, n):
    n = 8
    tmp = arg
    del tmp[-9:]
    return len(tmp)>3

try:
    h = [10,11,12,13,14]
    n = 7
    f(h, n)
    print(n, end=';')
    for el in h:
        print(el, end=';')
except:
    print('ERR')
code

9. Що буде виведено за виконання?
code
class A:
    b = 1
    c = 2

    def __init__(self):
        b = 3

class B(A):
    b = 4

    def __init__(self):
        super().__init__()
        self.a = 7

obj = B()
obj.b = obj.c + 10
B.b = 13

try: print(obj.b, end= ';')
except: print('ERR', end=';')

try: print(A.b, end= ';')
except: print('ERR', end=';')

try: print(B.b, end= ';')
except: print('ERR', end=';')
code

10. 
Для графа на рисунку з вершини 3 запущено алгоритм пошуку в глибину, що будує кістякове дерево. Списки суміжності вершин переглядаються алгоритмом за зростанням міток вершин.\par
\begin{center}
	\adjustimage{max size={0.9\linewidth}{0.9\paperheight}}
	{../10/Traversals/251.png}
\end{center}
\par Які з наведених ребер входять до складу отриманого кістякового дерева?\par
1) (0, 4)\par
2) (2, 5)\par
3) (3, 4)\par
4) (2, 6)\par
5) (3, 6)\par
6) (1, 3)\par
{\vskip 15pt} У формі вибрати номери відповідних ребер.\par
%answer:1 2 4 6
\end {large}}
%end
{\smallskip \begin {small} Затверджено на засіданні кафедри теоретичної кібернетики, протокол \No 12 від   19.05.2020 р.\par\smallskip\smallskip
{\parbox {6.5 cm}{Зав. кафедри \dotfill Ю.В.~Крак}}\hspace {0.7cm}{\hbox to 6.5 cm{ Екзаменатор \dotfill Т.О.~Карнаух}} \end{small}}
%begin
{{\begin {small}\par
{ \center  Київський національний університет імені Тараса Шевченка \par}
{\parbox {5 cm}{Освітня програма \hfill}{\parbox {7 cm}{Прикладна математика, Системний аналіз \hfill}}\parbox {6 cm}{\hfill Семестр 2}}\par
{\parbox {5 cm} {Навчальний предмет \hfill}\parbox {7 cm}{Програмування \hfill}\parbox {6cm}{\hfill}\par}\vspace{-7pt} \end{small}}{\center Екзаменаційний білет \No \textbf 
{ 103 }\par}}
{
\begin {large}
1. Що буде виведено за виконання?
code
try:
    res = not 8>9 or 4.0>=7.0 or 2<=False
    if isinstance(res, bool):
        print(f'bool;{res}')
    elif isinstance(res, int):
        print(f'int;{res}')
    elif isinstance(res, float):
        print(f'float;{res:.2f}')
    else:
        print('???')
except:
    print('ERR')
code

2. Що буде виведено за виконання?
code
a = 9
b = 4
c = 7

def f(b):
    global c
    a = 3
    b = 1
    c = 3
    return a + b + c

a = 5
b = 4
c = 8

try:
    print(f(a), a, b, c, sep=';')
except:
    print('ERR', a, b, c, sep=';')
code

3. Що буде виведено за виконання?
code
try:
    i = 12
    while i>8:
        i += -3
    print(i, end=';')
except:
    print('ERR')
code

4. Що буде виведено за виконання?
code
try:
    for e in range(2, -4, -3):
        if e>=4:
            continue
        print(e, end=';');e = 10
    else:
        print(e,end=';')
    print(e)
except:
    print('ERR')
code

5. Що буде виведено за виконання?\par (Дійсні значення, якщо наявні, в цьому завданні будуть виводитись у форматі з фіксованою крапкою та, можливо, з доволі великою кількістю десяткових розрядів. У відповіді для дійсних значень у ЦЬОМУ завданні достатньо вказати один десятковий розряд після крапки, відкинувши решту розрядів без округлення. Наприклад, замість 13.66666666666667 достатньо вказати 13.6, замість 25.23 достатньо вказати 25.2)
code
try:
    print(8, end=';')
    print(5//3, end=';')
    print(9, end=';')
except Exception:
    print(4, end=';')
except ZeroDivisionError:
    print(1, end=';')
except:
    print('ERR', end=';')
else:
    print(0, end=';')
finally:
    print(2)
code

6. Що буде виведено за виконання?
code
def f(n):
    if n>9:
        f(n//10)
    print(n%10, end='')
 
f(543216)
code

7. Що буде виведено за виконання?
code
class A:
    b = 1

    def _f1(self): print(2, end=';')
    def _f2(self): print(3, end=';')

    def main(self):
        try: print(self.b, end=';')
        except: print('ERR', end=';')

        try: self._f1()
        except: print('ERR', end=';')

        try: self._f2()
        except: print('ERR', end=';')


class B(A):
    b = 4

    def _f1(self): print(5, end=';')
    def _f2(self): print(6, end=';')

try: B().main()
except: print('ERR', end=';')
code

8. Що буде виведено за виконання?
code
def f(arg, n):
    arg[89] = 7
    n *= -2

try:
    n = 6
    t = {51:7, 89:2, 47:4, 44:7}
    f(t, n)
    print(n, end=';')
    for x in t.items():
        print(x, sep=';', end=';')
except:
    print('ERR')
code

9. Що буде виведено за виконання?
code
class A:
    a = 1
    c = 2

    def __init__(self):
        c = 3

class B(A):
    c = 4

    def __init__(self):
        super().__init__()
        self.b = 7

obj = B()
obj.a = obj.c + 10
B.a = 13

try: print(obj.a, end= ';')
except: print('ERR', end=';')

try: print(A.a, end= ';')
except: print('ERR', end=';')

try: print(B.a, end= ';')
except: print('ERR', end=';')
code

10. 
Для графа на рисунку з вершини 6 запущено алгоритм пошуку в глибину, що будує кістякове дерево. Списки суміжності вершин переглядаються алгоритмом за зростанням міток вершин.\par
\begin{center}
	\adjustimage{max size={0.9\linewidth}{0.9\paperheight}}
	{../10/Traversals/86.png}
\end{center}
\par Які з наведених ребер входять до складу отриманого кістякового дерева?\par
1) (2, 3)\par
2) (0, 6)\par
3) (4, 5)\par
4) (0, 5)\par
5) (5, 6)\par
6) (2, 6)\par
{\vskip 15pt} У формі вибрати номери відповідних ребер.\par
%answer:1 2 4
\end {large}}
%end
{\smallskip \begin {small} Затверджено на засіданні кафедри теоретичної кібернетики, протокол \No 12 від   19.05.2020 р.\par\smallskip\smallskip
{\parbox {6.5 cm}{Зав. кафедри \dotfill Ю.В.~Крак}}\hspace {0.7cm}{\hbox to 6.5 cm{ Екзаменатор \dotfill Т.О.~Карнаух}} \end{small}}
%begin
{{\begin {small}\par
{ \center  Київський національний університет імені Тараса Шевченка \par}
{\parbox {5 cm}{Освітня програма \hfill}{\parbox {7 cm}{Прикладна математика, Системний аналіз \hfill}}\parbox {6 cm}{\hfill Семестр 2}}\par
{\parbox {5 cm} {Навчальний предмет \hfill}\parbox {7 cm}{Програмування \hfill}\parbox {6cm}{\hfill}\par}\vspace{-7pt} \end{small}}{\center Екзаменаційний білет \No \textbf 
{ 104 }\par}}
{
\begin {large}
1. Що буде виведено за виконання?
code
try:
    res = 1**0.5*False
    if isinstance(res, bool):
        print(f'bool;{res}')
    elif isinstance(res, int):
        print(f'int;{res}')
    elif isinstance(res, float):
        print(f'float;{res:.2f}')
    else:
        print('???')
except:
    print('ERR')
code

2. Що буде виведено за виконання?
code
a = 9
b = 0
c = 4

def g(a):
    a -= 5
    b = 4
    c = 4
    return a + b + c

a = 8
b = 4
c = 7

try:
    print(g(a), a, b, c, sep=';')
except:
    print('ERR', a, b, c, sep=';')
code

3. Що буде виведено за виконання?
code
try:
    j = 13
    while j>0:
        j += -2
    print(j, end=';')
except:
    print('ERR')
code

4. Що буде виведено за виконання?
code
try:
    for a in range(-2, -8, 2):
        if a>7:
            continue
        print(a, end=';')
    else:
        print(a,end=';')
    print(a)
except:
    print('ERR')
code

5. Що буде виведено за виконання?\par (Дійсні значення, якщо наявні, в цьому завданні будуть виводитись у форматі з фіксованою крапкою та, можливо, з доволі великою кількістю десяткових розрядів. У відповіді для дійсних значень у ЦЬОМУ завданні достатньо вказати один десятковий розряд після крапки, відкинувши решту розрядів без округлення. Наприклад, замість 13.66666666666667 достатньо вказати 13.6, замість 25.23 достатньо вказати 25.2)
code
try:
    print(8, end=';')
    print(int('7'), end=';')
    print(3, end=';')
except ValueError:
    print(4, end=';')
except TypeError:
    print(0, end=';')
except:
    print('ERR', end=';')
else:
    print(5, end=';')
finally:
    print(1)
code

6. Що буде виведено за виконання?
code
def f(n):
    if n<=9:
        print(n%10, end='')
    else:
        f(n//10)

f(987651)
code

7. Що буде виведено за виконання?
code
class A:
    d = 1

    def _h1(self): print(2, end=';')
    def __h2(self): print(3, end=';')

    def main(self):
        try: print(self.d, end=';')
        except: print('ERR', end=';')

        try: self._h1()
        except: print('ERR', end=';')

        try: self.__h2()
        except: print('ERR', end=';')


class B(A):
    d = 4

    def _h1(self): print(5, end=';')
    def __h2(self): print(6, end=';')

try: B().main()
except: print('ERR', end=';')
code

8. Що буде виведено за виконання?
code
def f(arg, n):
    arg[53] = 7
    n = 5
    arg = set(arg.values())

try:
    n = 3
    t = {52:5, 53:2, 52:0}
    f(t, n)
    print(n, end=';')
    for x, y in t.items():
        print(x, y, sep=';', end=';')
except:
    print('ERR')
code

9. Що буде виведено за виконання?
code
class A:
    b = 1
    a = 2

    def __init__(self):
        b = 3

class B(A):
    b = 4

    def __init__(self):
        super().__init__()
        self.c = 7

obj = B()
obj.b = obj.a + 10
B.b = 13

try: print(obj.b, end= ';')
except: print('ERR', end=';')

try: print(A.b, end= ';')
except: print('ERR', end=';')

try: print(B.b, end= ';')
except: print('ERR', end=';')
code

10. 
Для графа на рисунку з вершини 4 запущено алгоритм пошуку в глибину, що будує кістякове дерево. Списки суміжності вершин переглядаються алгоритмом за зростанням міток вершин.\par
\begin{center}
	\adjustimage{max size={0.9\linewidth}{0.9\paperheight}}
	{../10/Traversals/54.png}
\end{center}
\par Які з наведених ребер входять до складу отриманого кістякового дерева?\par
1) (0, 4)\par
2) (1, 2)\par
3) (3, 6)\par
4) (2, 4)\par
5) (3, 4)\par
6) (1, 6)\par
{\vskip 15pt} У формі вибрати номери відповідних ребер.\par
%answer:1 2 3 4 6
\end {large}}
%end
{\smallskip \begin {small} Затверджено на засіданні кафедри теоретичної кібернетики, протокол \No 12 від   19.05.2020 р.\par\smallskip\smallskip
{\parbox {6.5 cm}{Зав. кафедри \dotfill Ю.В.~Крак}}\hspace {0.7cm}{\hbox to 6.5 cm{ Екзаменатор \dotfill Т.О.~Карнаух}} \end{small}}
%begin
{{\begin {small}\par
{ \center  Київський національний університет імені Тараса Шевченка \par}
{\parbox {5 cm}{Освітня програма \hfill}{\parbox {7 cm}{Прикладна математика, Системний аналіз \hfill}}\parbox {6 cm}{\hfill Семестр 2}}\par
{\parbox {5 cm} {Навчальний предмет \hfill}\parbox {7 cm}{Програмування \hfill}\parbox {6cm}{\hfill}\par}\vspace{-7pt} \end{small}}{\center Екзаменаційний білет \No \textbf 
{ 105 }\par}}
{
\begin {large}
1. Що буде виведено за виконання?
code
try:
    res = False<3
    if isinstance(res, bool):
        print(f'bool;{res}')
    elif isinstance(res, int):
        print(f'int;{res}')
    elif isinstance(res, float):
        print(f'float;{res:.2f}')
    else:
        print('???')
except:
    print('ERR')
code

2. Що буде виведено за виконання?
code
a = 1
b = 3
c = 5

def h(b):
    global c
    a += 2
    b = 3
    c = 5
    return a + b + c

a = 6
b = 2
c = 9

try:
    print(h(b), a, b, c, sep=';')
except:
    print('ERR', a, b, c, sep=';')
code

3. Що буде виведено за виконання?
code
try:
    k = 5
    while k>=7:
        k += -3
    print(k, end=';')
except:
    print('ERR')
code

4. Що буде виведено за виконання?
code
try:
    for b in range(15, -14, 3):
        if b>=5:
            break
        print(b, end=';')
    else:
        print(b,end=';')
    print(b)
except:
    print('ERR')
code

5. Що буде виведено за виконання?\par (Дійсні значення, якщо наявні, в цьому завданні будуть виводитись у форматі з фіксованою крапкою та, можливо, з доволі великою кількістю десяткових розрядів. У відповіді для дійсних значень у ЦЬОМУ завданні достатньо вказати один десятковий розряд після крапки, відкинувши решту розрядів без округлення. Наприклад, замість 13.66666666666667 достатньо вказати 13.6, замість 25.23 достатньо вказати 25.2)
code
try:
    print(6, end=';')
    print(int('a9'), end=';')
    print(1, end=';')
except Exception:
    print(8, end=';')
except KeyboardInterrupt:
    print(6, end=';')
except:
    print('ERR', end=';')
else:
    print(5, end=';')
finally:
    print(2)
code

6. Що буде виведено за виконання?
code
def f(n):
    if n>2:
        f(n//2)
    print(n%2, end='')
 
f(16)
code

7. Що буде виведено за виконання?
code
class A:
    a = 1

    def _g1(self): print(2, end=';')
    def g2(self): print(3, end=';')

    def main(self):
        try: print(self.a, end=';')
        except: print('ERR', end=';')

        try: self._g1()
        except: print('ERR', end=';')

        try: self.g2()
        except: print('ERR', end=';')


class B(A):
    a = 4

    def _g1(self): print(5, end=';')
    def g2(self): print(6, end=';')

try: B().main()
except: print('ERR', end=';')
code

8. Що буде виведено за виконання?
code
def f(arg, n):
    n = 1
    tmp = arg
    tmp[-1:3] = (15)
    return len(tmp)>3

try:
    b = ['0','1','2','3','4']
    n = 3
    f(b, n)
    print(n, end=';')
    for el in b:
        print(el, end=';')
except:
    print('ERR')
code

9. Що буде виведено за виконання?
code
class A:
    b = 1
    c = 2

    def __init__(self):
        c = 3

class B(A):
    b = 4

    def __init__(self):
        super().__init__()
        self.a = 7

obj = B()
obj.c = obj.b + 10
B.c = 13

try: print(obj.c, end= ';')
except: print('ERR', end=';')

try: print(A.c, end= ';')
except: print('ERR', end=';')

try: print(B.c, end= ';')
except: print('ERR', end=';')
code

10. 
Для графа на рисунку з вершини 3 запущено алгоритм пошуку в глибину, що будує кістякове дерево. Списки суміжності вершин переглядаються алгоритмом за зростанням міток вершин.\par
\begin{center}
	\adjustimage{max size={0.9\linewidth}{0.9\paperheight}}
	{../10/Traversals/114.png}
\end{center}
\par Які з наведених ребер входять до складу отриманого кістякового дерева?\par
1) (2, 6)\par
2) (0, 2)\par
3) (3, 6)\par
4) (5, 6)\par
5) (0, 4)\par
6) (1, 4)\par
{\vskip 15pt} У формі вибрати номери відповідних ребер.\par
%answer:2 4 6
\end {large}}
%end
{\smallskip \begin {small} Затверджено на засіданні кафедри теоретичної кібернетики, протокол \No 12 від   19.05.2020 р.\par\smallskip\smallskip
{\parbox {6.5 cm}{Зав. кафедри \dotfill Ю.В.~Крак}}\hspace {0.7cm}{\hbox to 6.5 cm{ Екзаменатор \dotfill Т.О.~Карнаух}} \end{small}}
%begin
{{\begin {small}\par
{ \center  Київський національний університет імені Тараса Шевченка \par}
{\parbox {5 cm}{Освітня програма \hfill}{\parbox {7 cm}{Прикладна математика, Системний аналіз \hfill}}\parbox {6 cm}{\hfill Семестр 2}}\par
{\parbox {5 cm} {Навчальний предмет \hfill}\parbox {7 cm}{Програмування \hfill}\parbox {6cm}{\hfill}\par}\vspace{-7pt} \end{small}}{\center Екзаменаційний білет \No \textbf 
{ 106 }\par}}
{
\begin {large}
1. Що буде виведено за виконання?
code
try:
    res = 7//9*8*5
    if isinstance(res, bool):
        print(f'bool;{res}')
    elif isinstance(res, int):
        print(f'int;{res}')
    elif isinstance(res, float):
        print(f'float;{res:.2f}')
    else:
        print('???')
except:
    print('ERR')
code

2. Що буде виведено за виконання?
code
a = 8
b = 0
c = 3

def f(a):
    a = 1
    b *= 4
    c = 1
    return a + b + c

a = 9
b = 7
c = 6

try:
    print(f(b), a, b, c, sep=';')
except:
    print('ERR', a, b, c, sep=';')
code

3. Що буде виведено за виконання?
code
try:
    k = 5
    while k<11:
        k += 1
    print(k, end=';')
except:
    print('ERR')
code

4. Що буде виведено за виконання?
code
try:
    for d in range(11, 6, 2):
        if d<8:
            continue
        print(d, end=';');d = 1
        if d>4:
            break
    else:
        print(d,end=';')
    print(d)
except:
    print('ERR')
code

5. Що буде виведено за виконання?\par (Дійсні значення, якщо наявні, в цьому завданні будуть виводитись у форматі з фіксованою крапкою та, можливо, з доволі великою кількістю десяткових розрядів. У відповіді для дійсних значень у ЦЬОМУ завданні достатньо вказати один десятковий розряд після крапки, відкинувши решту розрядів без округлення. Наприклад, замість 13.66666666666667 достатньо вказати 13.6, замість 25.23 достатньо вказати 25.2)
code
try:
    print(0, end=';')
    print(int('c4'), end=';')
    print(3, end=';')
except Exception:
    print(9, end=';')
except TypeError:
    print(7, end=';')
except:
    print('ERR', end=';')
else:
    print(1, end=';')
finally:
    print(2)
code

6. Що буде виведено за виконання?
code
def f(s):
    s1 = s
    for i in range(len(s)):
        s1[i] += 1
    return s1

s = [1, 2, 3]
f(s)
print(s)
code

7. Що буде виведено за виконання?
code
class A:
    c = 1

    def _g1(self): print(2, end=';')
    def __g2(self): print(3, end=';')

    def main(self):
        try: print(self.c, end=';')
        except: print('ERR', end=';')

        try: self._g1()
        except: print('ERR', end=';')

        try: self.__g2()
        except: print('ERR', end=';')


class B(A):
    c = 4

    def _g1(self): print(5, end=';')
    def __g2(self): print(6, end=';')

try: B().main()
except: print('ERR', end=';')
code

8. Що буде виведено за виконання?
code
def f(arg, n):
    n = 4
    tmp = arg
    tmp[-2:-1] = 'xy'
    return len(tmp)>3

try:
    g = [10,11,12,13,14]
    n = 6
    f(g, n)
    print(n, end=';')
    for el in g:
        print(el, end=';')
except:
    print('ERR')
code

9. Що буде виведено за виконання?
code
class A:
    c = 1
    a = 2

    def __init__(self):
        c = 3

class B(A):
    a = 4

    def __init__(self):
        super().__init__()
        self.b = 7

obj = B()
obj.a = obj.b + 10
B.a = 13

try: print(obj.a, end= ';')
except: print('ERR', end=';')

try: print(A.a, end= ';')
except: print('ERR', end=';')

try: print(B.a, end= ';')
except: print('ERR', end=';')
code

10. 
Для графа на рисунку з вершини 1 запущено алгоритм пошуку в глибину, що будує кістякове дерево. Списки суміжності вершин переглядаються алгоритмом за зростанням міток вершин.\par
\begin{center}
	\adjustimage{max size={0.9\linewidth}{0.9\paperheight}}
	{../10/Traversals/218.png}
\end{center}
\par Які з наведених ребер входять до складу отриманого кістякового дерева?\par
1) (2, 5)\par
2) (0, 1)\par
3) (2, 4)\par
4) (0, 4)\par
5) (0, 6)\par
6) (3, 5)\par
{\vskip 15pt} У формі вибрати номери відповідних ребер.\par
%answer:1 2 3 4 6
\end {large}}
%end
{\smallskip \begin {small} Затверджено на засіданні кафедри теоретичної кібернетики, протокол \No 12 від   19.05.2020 р.\par\smallskip\smallskip
{\parbox {6.5 cm}{Зав. кафедри \dotfill Ю.В.~Крак}}\hspace {0.7cm}{\hbox to 6.5 cm{ Екзаменатор \dotfill Т.О.~Карнаух}} \end{small}}
%begin
{{\begin {small}\par
{ \center  Київський національний університет імені Тараса Шевченка \par}
{\parbox {5 cm}{Освітня програма \hfill}{\parbox {7 cm}{Прикладна математика, Системний аналіз \hfill}}\parbox {6 cm}{\hfill Семестр 2}}\par
{\parbox {5 cm} {Навчальний предмет \hfill}\parbox {7 cm}{Програмування \hfill}\parbox {6cm}{\hfill}\par}\vspace{-7pt} \end{small}}{\center Екзаменаційний білет \No \textbf 
{ 107 }\par}}
{
\begin {large}
1. Що буде виведено за виконання?
code
try:
    res = 9.0==8.0 and not 3<5 and 7!=True
    if isinstance(res, bool):
        print(f'bool;{res}')
    elif isinstance(res, int):
        print(f'int;{res}')
    elif isinstance(res, float):
        print(f'float;{res:.2f}')
    else:
        print('???')
except:
    print('ERR')
code

2. Що буде виведено за виконання?
code
a = 5
b = 8
c = 4

def g(b):
    a -= 3
    b = 2
    c = 5
    return a + b + c

a = 1
b = 0
c = 2

try:
    print(g(a), a, b, c, sep=';')
except:
    print('ERR', a, b, c, sep=';')
code

3. Що буде виведено за виконання?
code
try:
    t = 10
    while t>6:
        t += -3
    print(t, end=';')
except:
    print('ERR')
code

4. Що буде виведено за виконання?
code
try:
    for f in range(-9, -13, 3):
        if f>6:
            break
        print(f, end=';')
    else:
        print(f,end=';')
    print(f)
except:
    print('ERR')
code

5. Що буде виведено за виконання?\par (Дійсні значення, якщо наявні, в цьому завданні будуть виводитись у форматі з фіксованою крапкою та, можливо, з доволі великою кількістю десяткових розрядів. У відповіді для дійсних значень у ЦЬОМУ завданні достатньо вказати один десятковий розряд після крапки, відкинувши решту розрядів без округлення. Наприклад, замість 13.66666666666667 достатньо вказати 13.6, замість 25.23 достатньо вказати 25.2)
code
try:
    print(7, end=';')
    print(0//1, end=';')
    print(5, end=';')
except KeyboardInterrupt:
    print(4, end=';')
except ZeroDivisionError:
    print(8, end=';')
except:
    print('ERR', end=';')
else:
    print(6, end=';')
finally:
    print(9)
code

6. Що буде виведено за виконання?
code
def f(s):
    for i in range(len(s)):
        s[i] += 1
        return
 
s = [1, 2, 3]
f(s)
print(s)
code

7. Що буде виведено за виконання?
code
class A:
    d = 1

    def _f1(self): print(2, end=';')
    def _f2(self): print(3, end=';')

    def main(self):
        try: print(self.d, end=';')
        except: print('ERR', end=';')

        try: self._f1()
        except: print('ERR', end=';')

        try: self._f2()
        except: print('ERR', end=';')


class B(A):
    d = 4

    def _f1(self): print(5, end=';')
    def _f2(self): print(6, end=';')

try: B().main()
except: print('ERR', end=';')
code

8. Що буде виведено за виконання?
code
def f(arg, n):
    n = 9
    tmp = arg
    del tmp[8:0]
    return len(tmp)>3

try:
    e = ['0','1','2','3','4']
    n = 5
    f(e, n)
    print(n, end=';')
    for el in e:
        print(el, end=';')
except:
    print('ERR')
code

9. Що буде виведено за виконання?
code
class A:
    c = 1
    b = 2

    def __init__(self):
        b = 3

class B(A):
    c = 4

    def __init__(self):
        super().__init__()
        self.a = 7

obj = B()
obj.c = obj.a + 10
B.c = 13

try: print(obj.c, end= ';')
except: print('ERR', end=';')

try: print(A.c, end= ';')
except: print('ERR', end=';')

try: print(B.c, end= ';')
except: print('ERR', end=';')
code

10. 
Для графа на рисунку з вершини 4 запущено алгоритм пошуку в глибину, що будує кістякове дерево. Списки суміжності вершин переглядаються алгоритмом за зростанням міток вершин.\par
\begin{center}
	\adjustimage{max size={0.9\linewidth}{0.9\paperheight}}
	{../10/Traversals/28.png}
\end{center}
\par Які з наведених ребер входять до складу отриманого кістякового дерева?\par
1) (3, 5)\par
2) (1, 2)\par
3) (0, 1)\par
4) (2, 3)\par
5) (3, 4)\par
6) (2, 4)\par
{\vskip 15pt} У формі вибрати номери відповідних ребер.\par
%answer:1 2 3 6
\end {large}}
%end
{\smallskip \begin {small} Затверджено на засіданні кафедри теоретичної кібернетики, протокол \No 12 від   19.05.2020 р.\par\smallskip\smallskip
{\parbox {6.5 cm}{Зав. кафедри \dotfill Ю.В.~Крак}}\hspace {0.7cm}{\hbox to 6.5 cm{ Екзаменатор \dotfill Т.О.~Карнаух}} \end{small}}
%begin
{{\begin {small}\par
{ \center  Київський національний університет імені Тараса Шевченка \par}
{\parbox {5 cm}{Освітня програма \hfill}{\parbox {7 cm}{Прикладна математика, Системний аналіз \hfill}}\parbox {6 cm}{\hfill Семестр 2}}\par
{\parbox {5 cm} {Навчальний предмет \hfill}\parbox {7 cm}{Програмування \hfill}\parbox {6cm}{\hfill}\par}\vspace{-7pt} \end{small}}{\center Екзаменаційний білет \No \textbf 
{ 108 }\par}}
{
\begin {large}
1. Що буде виведено за виконання?
code
try:
    res = "5">"5.0j"
    if isinstance(res, bool):
        print(f'bool;{res}')
    elif isinstance(res, int):
        print(f'int;{res}')
    elif isinstance(res, float):
        print(f'float;{res:.2f}')
    else:
        print('???')
except:
    print('ERR')
code

2. Що буде виведено за виконання?
code
a = 9
b = 6
c = 2

def g(a):
    a = 2
    b = 1
    c = 5
    return a + b + c

a = 0
b = 5
c = 1

try:
    print(g(a), a, b, c, sep=';')
except:
    print('ERR', a, b, c, sep=';')
code

3. Що буде виведено за виконання?
code
try:
    i = 4
    while i<5:
        i += 1
    print(i, end=';')
except:
    print('ERR')
code

4. Що буде виведено за виконання?
code
try:
    for c in range(-10, 4, -1):
        if c<=3:
            continue
        print(c, end=';')
    else:
        print('end',end=';')
    print(c)
except:
    print('ERR')
code

5. Що буде виведено за виконання?\par (Дійсні значення, якщо наявні, в цьому завданні будуть виводитись у форматі з фіксованою крапкою та, можливо, з доволі великою кількістю десяткових розрядів. У відповіді для дійсних значень у ЦЬОМУ завданні достатньо вказати один десятковий розряд після крапки, відкинувши решту розрядів без округлення. Наприклад, замість 13.66666666666667 достатньо вказати 13.6, замість 25.23 достатньо вказати 25.2)
code
try:
    print(3, end=';')
    print(8j!=7j, end=';')
    print(4, end=';')
except ValueError:
    print(2, end=';')
except ZeroDivisionError:
    print(1, end=';')
except:
    print('ERR', end=';')
else:
    print(7, end=';')
finally:
    print(6)
code

6. Що буде виведено за виконання?
code
def f(n):
    if n>9:
        f(n//100)
    print(n%10,end='')
 
f(987654)
code

7. Що буде виведено за виконання?
code
class A:
    b = 1

    def _h1(self): print(2, end=';')
    def h2(self): print(3, end=';')

    def main(self):
        try: print(self.b, end=';')
        except: print('ERR', end=';')

        try: self._h1()
        except: print('ERR', end=';')

        try: self.h2()
        except: print('ERR', end=';')


class B(A):
    b = 4

    def _h1(self): print(5, end=';')
    def h2(self): print(6, end=';')

try: B().main()
except: print('ERR', end=';')
code

8. Що буде виведено за виконання?
code
def f(arg, n):
    arg[10] = 5
    n = 5
    arg = list(arg.keys())

try:
    n = 0
    s = {25:3, 10:2, 37:1, 37:9}
    f(s, n)
    print(n, end=';')
    for x in s.values():
        print(x, sep=';', end=';')
except:
    print('ERR')
code

9. Що буде виведено за виконання?
code
class A:
    a = 1
    b = 2

    def __init__(self):
        a = 3

class B(A):
    b = 4

    def __init__(self):
        super().__init__()
        self.c = 7

obj = B()
obj.b = obj.c + 10
B.b = 13

try: print(obj.b, end= ';')
except: print('ERR', end=';')

try: print(A.b, end= ';')
except: print('ERR', end=';')

try: print(B.b, end= ';')
except: print('ERR', end=';')
code

10. 
Для графа на рисунку з вершини 4 запущено алгоритм пошуку в глибину, що будує кістякове дерево. Списки суміжності вершин переглядаються алгоритмом за зростанням міток вершин.\par
\begin{center}
	\adjustimage{max size={0.9\linewidth}{0.9\paperheight}}
	{../10/Traversals/223.png}
\end{center}
\par Які з наведених ребер входять до складу отриманого кістякового дерева?\par
1) (1, 5)\par
2) (1, 2)\par
3) (2, 4)\par
4) (3, 6)\par
5) (0, 1)\par
6) (0, 3)\par
{\vskip 15pt} У формі вибрати номери відповідних ребер.\par
%answer:2 3 4 5 6
\end {large}}
%end
{\smallskip \begin {small} Затверджено на засіданні кафедри теоретичної кібернетики, протокол \No 12 від   19.05.2020 р.\par\smallskip\smallskip
{\parbox {6.5 cm}{Зав. кафедри \dotfill Ю.В.~Крак}}\hspace {0.7cm}{\hbox to 6.5 cm{ Екзаменатор \dotfill Т.О.~Карнаух}} \end{small}}
%begin
{{\begin {small}\par
{ \center  Київський національний університет імені Тараса Шевченка \par}
{\parbox {5 cm}{Освітня програма \hfill}{\parbox {7 cm}{Прикладна математика, Системний аналіз \hfill}}\parbox {6 cm}{\hfill Семестр 2}}\par
{\parbox {5 cm} {Навчальний предмет \hfill}\parbox {7 cm}{Програмування \hfill}\parbox {6cm}{\hfill}\par}\vspace{-7pt} \end{small}}{\center Екзаменаційний білет \No \textbf 
{ 109 }\par}}
{
\begin {large}
1. Що буде виведено за виконання?
code
try:
    res = 2**-1-True
    if isinstance(res, bool):
        print(f'bool;{res}')
    elif isinstance(res, int):
        print(f'int;{res}')
    elif isinstance(res, float):
        print(f'float;{res:.2f}')
    else:
        print('???')
except:
    print('ERR')
code

2. Що буде виведено за виконання?
code
a = 8
b = 4
c = 9

def f(a):
    global c
    a *= 5
    b = 4
    c = 1
    return a + b + c

a = 5
b = 2
c = 0

try:
    print(f(b), a, b, c, sep=';')
except:
    print('ERR', a, b, c, sep=';')
code

3. Що буде виведено за виконання?
code
try:
    j = 7
    while j>=2:
        j += -2
    print(j, end=';')
except:
    print('ERR')
code

4. Що буде виведено за виконання?
code
try:
    for c in range(-4, -3, -3):
        if c>8:
            continue
        print(c, end=';');c = -2
    else:
        print(c,end=';')
    print(c)
except:
    print('ERR')
code

5. Що буде виведено за виконання?\par (Дійсні значення, якщо наявні, в цьому завданні будуть виводитись у форматі з фіксованою крапкою та, можливо, з доволі великою кількістю десяткових розрядів. У відповіді для дійсних значень у ЦЬОМУ завданні достатньо вказати один десятковий розряд після крапки, відкинувши решту розрядів без округлення. Наприклад, замість 13.66666666666667 достатньо вказати 13.6, замість 25.23 достатньо вказати 25.2)
code
try:
    print(3, end=';')
    print(9/0.0, end=';')
    print(5, end=';')
except KeyboardInterrupt:
    print(0, end=';')
except TypeError:
    print(9, end=';')
except:
    print('ERR', end=';')
else:
    print(1, end=';')
finally:
    print(6)
code

6. Що буде виведено за виконання?
code
def f(s):
    for el in s:
        el += 1
        return s
 
s = [1, 2, 3]
s = f(s)
print(s)
code

7. Що буде виведено за виконання?
code
class A:
    d = 1

    def f1(self): print(2, end=';')
    def _f2(self): print(3, end=';')

    def main(self):
        try: print(self.d, end=';')
        except: print('ERR', end=';')

        try: self.f1()
        except: print('ERR', end=';')

        try: self._f2()
        except: print('ERR', end=';')


class B(A):
    d = 4

    def f1(self): print(5, end=';')
    def _f2(self): print(6, end=';')

try: B().main()
except: print('ERR', end=';')
code

8. Що буде виведено за виконання?
code
def f(arg, n):
    arg[18] = 6
    n = 8
    arg = list(arg.values())

try:
    n = 2
    t = {50:4, 40:4, 40:4}
    f(t, n)
    print(n, end=';')
    for x in t.keys():
        print(x, sep=';', end=';')
except:
    print('ERR')
code

9. Що буде виведено за виконання?
code
class A:
    b = 1
    c = 2

    def __init__(self):
        b = 3

class B(A):
    b = 4

    def __init__(self):
        super().__init__()
        self.a = 7

obj = B()
obj.c = obj.a + 10
B.c = 13

try: print(obj.c, end= ';')
except: print('ERR', end=';')

try: print(A.c, end= ';')
except: print('ERR', end=';')

try: print(B.c, end= ';')
except: print('ERR', end=';')
code

10. 
Для графа на рисунку з вершини 1 запущено алгоритм пошуку в глибину, що будує кістякове дерево. Списки суміжності вершин переглядаються алгоритмом за зростанням міток вершин.\par
\begin{center}
	\adjustimage{max size={0.9\linewidth}{0.9\paperheight}}
	{../10/Traversals/142.png}
\end{center}
\par Які з наведених ребер входять до складу отриманого кістякового дерева?\par
1) (3, 6)\par
2) (2, 5)\par
3) (4, 5)\par
4) (0, 6)\par
5) (0, 5)\par
6) (0, 3)\par
{\vskip 15pt} У формі вибрати номери відповідних ребер.\par
%answer:1 2 5 6
\end {large}}
%end
{\smallskip \begin {small} Затверджено на засіданні кафедри теоретичної кібернетики, протокол \No 12 від   19.05.2020 р.\par\smallskip\smallskip
{\parbox {6.5 cm}{Зав. кафедри \dotfill Ю.В.~Крак}}\hspace {0.7cm}{\hbox to 6.5 cm{ Екзаменатор \dotfill Т.О.~Карнаух}} \end{small}}
%begin
{{\begin {small}\par
{ \center  Київський національний університет імені Тараса Шевченка \par}
{\parbox {5 cm}{Освітня програма \hfill}{\parbox {7 cm}{Прикладна математика, Системний аналіз \hfill}}\parbox {6 cm}{\hfill Семестр 2}}\par
{\parbox {5 cm} {Навчальний предмет \hfill}\parbox {7 cm}{Програмування \hfill}\parbox {6cm}{\hfill}\par}\vspace{-7pt} \end{small}}{\center Екзаменаційний білет \No \textbf 
{ 110 }\par}}
{
\begin {large}
1. Що буде виведено за виконання?
code
try:
    res = not 6.0>2 or 8<False and 5.0<=6
    if isinstance(res, bool):
        print(f'bool;{res}')
    elif isinstance(res, int):
        print(f'int;{res}')
    elif isinstance(res, float):
        print(f'float;{res:.2f}')
    else:
        print('???')
except:
    print('ERR')
code

2. Що буде виведено за виконання?
code
a = 3
b = 7
c = 0

def h(b):
    global c
    a = 4
    b = 5
    c = 2
    return a + b + c

a = 9
b = 3
c = 5

try:
    print(h(b), a, b, c, sep=';')
except:
    print('ERR', a, b, c, sep=';')
code

3. Що буде виведено за виконання?
code
try:
    m = 5
    while m>9:
        m += -2
    print(m, end=';')
except:
    print('ERR')
code

4. Що буде виведено за виконання?
code
try:
    for f in range(13, 3, -2):
        if f>=3:
            break
        print(f, end=';')
    else:
        print(f,end=';')
    print(f)
except:
    print('ERR')
code

5. Що буде виведено за виконання?\par (Дійсні значення, якщо наявні, в цьому завданні будуть виводитись у форматі з фіксованою крапкою та, можливо, з доволі великою кількістю десяткових розрядів. У відповіді для дійсних значень у ЦЬОМУ завданні достатньо вказати один десятковий розряд після крапки, відкинувши решту розрядів без округлення. Наприклад, замість 13.66666666666667 достатньо вказати 13.6, замість 25.23 достатньо вказати 25.2)
code
try:
    print(0, end=';')
    print(int('4'), end=';')
    print(2, end=';')
except Exception:
    print(5, end=';')
except ValueError:
    print(7, end=';')
except:
    print('ERR', end=';')
else:
    print(8, end=';')
finally:
    print(3)
code

6. Що буде виведено за виконання?
code
def f(s):
    for el in s:
        el += 1
    return s
 
s = [1, 2, 3]
s = f(s)
print(s)
code

7. Що буде виведено за виконання?
code
class A:
    c = 1

    def __h1(self): print(2, end=';')
    def _h2(self): print(3, end=';')

    def main(self):
        try: print(self.c, end=';')
        except: print('ERR', end=';')

        try: self.__h1()
        except: print('ERR', end=';')

        try: self._h2()
        except: print('ERR', end=';')


class B(A):
    c = 4

    def __h1(self): print(5, end=';')
    def _h2(self): print(6, end=';')

try: B().main()
except: print('ERR', end=';')
code

8. Що буде виведено за виконання?
code
def f(arg, n):
    n = 2
    tmp = arg
    tmp[-3:1] = 'ab'
    return len(tmp)>3

try:
    c = [10,11,12,13,14]
    n = 4
    f(c, n)
    print(n, end=';')
    for el in c:
        print(el, end=';')
except:
    print('ERR')
code

9. Що буде виведено за виконання?
code
class A:
    a = 1
    b = 2

    def __init__(self):
        b = 3

class B(A):
    b = 4

    def __init__(self):
        super().__init__()
        self.c = 7

obj = B()
obj.b = obj.b + 10
B.b = 13

try: print(obj.b, end= ';')
except: print('ERR', end=';')

try: print(A.b, end= ';')
except: print('ERR', end=';')

try: print(B.b, end= ';')
except: print('ERR', end=';')
code

10. 
Для графа на рисунку з вершини 4 запущено алгоритм пошуку в глибину, що будує кістякове дерево. Списки суміжності вершин переглядаються алгоритмом за зростанням міток вершин.\par
\begin{center}
	\adjustimage{max size={0.9\linewidth}{0.9\paperheight}}
	{../10/Traversals/297.png}
\end{center}
\par Які з наведених ребер входять до складу отриманого кістякового дерева?\par
1) (2, 6)\par
2) (4, 5)\par
3) (0, 6)\par
4) (1, 3)\par
5) (1, 4)\par
6) (0, 5)\par
{\vskip 15pt} У формі вибрати номери відповідних ребер.\par
%answer:1 3 4 5 6
\end {large}}
%end
{\smallskip \begin {small} Затверджено на засіданні кафедри теоретичної кібернетики, протокол \No 12 від   19.05.2020 р.\par\smallskip\smallskip
{\parbox {6.5 cm}{Зав. кафедри \dotfill Ю.В.~Крак}}\hspace {0.7cm}{\hbox to 6.5 cm{ Екзаменатор \dotfill Т.О.~Карнаух}} \end{small}}
%begin
{{\begin {small}\par
{ \center  Київський національний університет імені Тараса Шевченка \par}
{\parbox {5 cm}{Освітня програма \hfill}{\parbox {7 cm}{Прикладна математика, Системний аналіз \hfill}}\parbox {6 cm}{\hfill Семестр 2}}\par
{\parbox {5 cm} {Навчальний предмет \hfill}\parbox {7 cm}{Програмування \hfill}\parbox {6cm}{\hfill}\par}\vspace{-7pt} \end{small}}{\center Екзаменаційний білет \No \textbf 
{ 111 }\par}}
{
\begin {large}
1. Що буде виведено за виконання?
code
try:
    res = 3.0==3 and not True!=2.0 or 9>=4
    if isinstance(res, bool):
        print(f'bool;{res}')
    elif isinstance(res, int):
        print(f'int;{res}')
    elif isinstance(res, float):
        print(f'float;{res:.2f}')
    else:
        print('???')
except:
    print('ERR')
code

2. Що буде виведено за виконання?
code
a = 3
b = 1
c = 7

def h(a):
    a += 2
    b = 3
    c = 5
    return a + b + c

a = 6
b = 2
c = 7

try:
    print(h(a), a, b, c, sep=';')
except:
    print('ERR', a, b, c, sep=';')
code

3. Що буде виведено за виконання?
code
try:
    j = 0
    while j<=6:
        j += 2
    print(j, end=';')
except:
    print('ERR')
code

4. Що буде виведено за виконання?
code
try:
    for e in range(4, -12, -2):
        if e<4:
            continue
        print(e, end=';')
    else:
        print(e,end=';')
    print(e)
except:
    print('ERR')
code

5. Що буде виведено за виконання?\par (Дійсні значення, якщо наявні, в цьому завданні будуть виводитись у форматі з фіксованою крапкою та, можливо, з доволі великою кількістю десяткових розрядів. У відповіді для дійсних значень у ЦЬОМУ завданні достатньо вказати один десятковий розряд після крапки, відкинувши решту розрядів без округлення. Наприклад, замість 13.66666666666667 достатньо вказати 13.6, замість 25.23 достатньо вказати 25.2)
code
try:
    print(7, end=';')
    print(0j>8j, end=';')
    print(1, end=';')
except KeyboardInterrupt:
    print(2, end=';')
except TypeError:
    print(3, end=';')
except:
    print('ERR', end=';')
else:
    print(6, end=';')
finally:
    print(4)
code

6. Що буде виведено за виконання?
code
def f(n):
    print(n%10,end='')
    if n>9:
        f(n//100)
 
f(123456)
code

7. Що буде виведено за виконання?
code
class A:
    a = 1

    def _g1(self): print(2, end=';')
    def __g2(self): print(3, end=';')

    def main(self):
        try: print(self.a, end=';')
        except: print('ERR', end=';')

        try: self._g1()
        except: print('ERR', end=';')

        try: self.__g2()
        except: print('ERR', end=';')


class B(A):
    a = 4

    def _g1(self): print(5, end=';')
    def __g2(self): print(6, end=';')

try: B().main()
except: print('ERR', end=';')
code

8. Що буде виведено за виконання?
code
def f(arg, n):
    arg[40] = 7
    n -= 2

try:
    n = 5
    d = {40:7, 36:2, 21:8, 40:3}
    f(d, n)
    print(n, end=';')
    for x in d.keys():
        print(x, sep=';', end=';')
except:
    print('ERR')
code

9. Що буде виведено за виконання?
code
class A:
    b = 1
    a = 2

    def __init__(self):
        a = 3

class B(A):
    b = 4

    def __init__(self):
        super().__init__()
        self.c = 7

obj = B()
obj.c = obj.a + 10
B.c = 13

try: print(obj.c, end= ';')
except: print('ERR', end=';')

try: print(A.c, end= ';')
except: print('ERR', end=';')

try: print(B.c, end= ';')
except: print('ERR', end=';')
code

10. 
Для графа на рисунку з вершини 0 запущено алгоритм пошуку в глибину, що будує кістякове дерево. Списки суміжності вершин переглядаються алгоритмом за зростанням міток вершин.\par
\begin{center}
	\adjustimage{max size={0.9\linewidth}{0.9\paperheight}}
	{../10/Traversals/74.png}
\end{center}
\par Які з наведених ребер входять до складу отриманого кістякового дерева?\par
1) (0, 1)\par
2) (1, 3)\par
3) (1, 2)\par
4) (3, 5)\par
5) (0, 6)\par
6) (1, 4)\par
{\vskip 15pt} У формі вибрати номери відповідних ребер.\par
%answer:1 3 4 5 6
\end {large}}
%end
{\smallskip \begin {small} Затверджено на засіданні кафедри теоретичної кібернетики, протокол \No 12 від   19.05.2020 р.\par\smallskip\smallskip
{\parbox {6.5 cm}{Зав. кафедри \dotfill Ю.В.~Крак}}\hspace {0.7cm}{\hbox to 6.5 cm{ Екзаменатор \dotfill Т.О.~Карнаух}} \end{small}}
%begin
{{\begin {small}\par
{ \center  Київський національний університет імені Тараса Шевченка \par}
{\parbox {5 cm}{Освітня програма \hfill}{\parbox {7 cm}{Прикладна математика, Системний аналіз \hfill}}\parbox {6 cm}{\hfill Семестр 2}}\par
{\parbox {5 cm} {Навчальний предмет \hfill}\parbox {7 cm}{Програмування \hfill}\parbox {6cm}{\hfill}\par}\vspace{-7pt} \end{small}}{\center Екзаменаційний білет \No \textbf 
{ 112 }\par}}
{
\begin {large}
1. Що буде виведено за виконання?
code
def f(n):
    print('f', end='')
    return n<=-3

try:
    print(f(3) or f(-6))
except:
    print('ERR')
code

2. Що буде виведено за виконання?
code
a = 2
b = 6
c = 8

def f(a):
    a = 1
    b = 3
    c = 4
    return a + b + c

a = 1
b = 7
c = 4

try:
    print(f(b), a, b, c, sep=';')
except:
    print('ERR', a, b, c, sep=';')
code

3. Що буде виведено за виконання?
code
try:
    t = 3
    while t<=13:
        t += 2
    print(t, end=';')
except:
    print('ERR')
code

4. Що буде виведено за виконання?
code
try:
    for b in range(1, 11, 2):
        if b<=7:
            continue
        print(b, end=';');b = 8
    else:
        print(b,end=';')
    print(b)
except:
    print('ERR')
code

5. Що буде виведено за виконання?\par (Дійсні значення, якщо наявні, в цьому завданні будуть виводитись у форматі з фіксованою крапкою та, можливо, з доволі великою кількістю десяткових розрядів. У відповіді для дійсних значень у ЦЬОМУ завданні достатньо вказати один десятковий розряд після крапки, відкинувши решту розрядів без округлення. Наприклад, замість 13.66666666666667 достатньо вказати 13.6, замість 25.23 достатньо вказати 25.2)
code
try:
    print(9, end=';')
    print(5%0, end=';')
    print(8, end=';')
except Exception:
    print(5, end=';')
except ValueError:
    print(9, end=';')
except:
    print('ERR', end=';')
else:
    print(0, end=';')
finally:
    print(4)
code

6. Що буде виведено за виконання?
code
def f(s1):
    s = s1[:]
    for i in range(len(s)):
        s[i] += 1
 
s = [1, 2, 3]
f(s)
print(s)
code

7. Що буде виведено за виконання?
code
class A:
    b = 1

    def _h1(self): print(2, end=';')
    def _h2(self): print(3, end=';')

    def main(self):
        try: print(self.b, end=';')
        except: print('ERR', end=';')

        try: self._h1()
        except: print('ERR', end=';')

        try: self._h2()
        except: print('ERR', end=';')


class B(A):
    b = 4

    def _h1(self): print(5, end=';')
    def _h2(self): print(6, end=';')

try: B().main()
except: print('ERR', end=';')
code

8. Що буде виведено за виконання?
code
def f(arg, n):
    n = 3
    tmp = arg
    del tmp[-5:-8]
    return len(tmp)>3

try:
    h = ['0','1','2','3','4']
    n = 0
    f(h, n)
    print(n, end=';')
    for el in h:
        print(el, end=';')
except:
    print('ERR')
code

9. Що буде виведено за виконання?
code
class A:
    a = 1
    c = 2

    def __init__(self):
        a = 3

class B(A):
    c = 4

    def __init__(self):
        super().__init__()
        self.b = 7

obj = B()
obj.a = obj.b + 10
B.a = 13

try: print(obj.a, end= ';')
except: print('ERR', end=';')

try: print(A.a, end= ';')
except: print('ERR', end=';')

try: print(B.a, end= ';')
except: print('ERR', end=';')
code

10. 
Для графа на рисунку з вершини 3 запущено алгоритм пошуку в глибину, що будує кістякове дерево. Списки суміжності вершин переглядаються алгоритмом за зростанням міток вершин.\par
\begin{center}
	\adjustimage{max size={0.9\linewidth}{0.9\paperheight}}
	{../10/Traversals/187.png}
\end{center}
\par Які з наведених ребер входять до складу отриманого кістякового дерева?\par
1) (1, 2)\par
2) (2, 3)\par
3) (3, 6)\par
4) (0, 6)\par
5) (0, 1)\par
6) (2, 6)\par
{\vskip 15pt} У формі вибрати номери відповідних ребер.\par
%answer:1 2 5
\end {large}}
%end
{\smallskip \begin {small} Затверджено на засіданні кафедри теоретичної кібернетики, протокол \No 12 від   19.05.2020 р.\par\smallskip\smallskip
{\parbox {6.5 cm}{Зав. кафедри \dotfill Ю.В.~Крак}}\hspace {0.7cm}{\hbox to 6.5 cm{ Екзаменатор \dotfill Т.О.~Карнаух}} \end{small}}
%begin
{{\begin {small}\par
{ \center  Київський національний університет імені Тараса Шевченка \par}
{\parbox {5 cm}{Освітня програма \hfill}{\parbox {7 cm}{Прикладна математика, Системний аналіз \hfill}}\parbox {6 cm}{\hfill Семестр 2}}\par
{\parbox {5 cm} {Навчальний предмет \hfill}\parbox {7 cm}{Програмування \hfill}\parbox {6cm}{\hfill}\par}\vspace{-7pt} \end{small}}{\center Екзаменаційний білет \No \textbf 
{ 113 }\par}}
{
\begin {large}
1. Що буде виведено за виконання?
code
try:
    res = 4/4/5*9
    if isinstance(res, bool):
        print(f'bool;{res}')
    elif isinstance(res, int):
        print(f'int;{res}')
    elif isinstance(res, float):
        print(f'float;{res:.2f}')
    else:
        print('???')
except:
    print('ERR')
code

2. Що буде виведено за виконання?
code
a = 8
b = 7
c = 1

def g(b):
    global c
    a = 4
    b = 5
    c = 3
    return a + b + c

a = 9
b = 2
c = 5

try:
    print(g(b), a, b, c, sep=';')
except:
    print('ERR', a, b, c, sep=';')
code

3. Що буде виведено за виконання?
code
try:
    j = 9
    while j<12:
        j += 1
    print(j, end=';')
except:
    print('ERR')
code

4. Що буде виведено за виконання?
code
try:
    for d in range(10, -1, -3):
        if d>=5:
            break
        print(d, end=';')
    else:
        print(d,end=';')
    print(d)
except:
    print('ERR')
code

5. Що буде виведено за виконання?\par (Дійсні значення, якщо наявні, в цьому завданні будуть виводитись у форматі з фіксованою крапкою та, можливо, з доволі великою кількістю десяткових розрядів. У відповіді для дійсних значень у ЦЬОМУ завданні достатньо вказати один десятковий розряд після крапки, відкинувши решту розрядів без округлення. Наприклад, замість 13.66666666666667 достатньо вказати 13.6, замість 25.23 достатньо вказати 25.2)
code
try:
    print(6, end=';')
    print(int('2'), end=';')
    print(8, end=';')
except ZeroDivisionError:
    print(1, end=';')
except KeyboardInterrupt:
    print(3, end=';')
except:
    print('ERR', end=';')
else:
    print(7, end=';')
finally:
    print(1)
code

6. Що буде виведено за виконання?
code
def f(n):
    print(n%10, end='')
    if n>9:
        f(n//10)
 
f(123456)
code

7. Що буде виведено за виконання?
code
class A:
    b = 1

    def _g1(self): print(2, end=';')
    def g2(self): print(3, end=';')

    def main(self):
        try: print(self.b, end=';')
        except: print('ERR', end=';')

        try: self._g1()
        except: print('ERR', end=';')

        try: self.g2()
        except: print('ERR', end=';')


class B(A):
    b = 4

    def _g1(self): print(5, end=';')
    def g2(self): print(6, end=';')

try: B().main()
except: print('ERR', end=';')
code

8. Що буде виведено за виконання?
code
def f(arg, n):
    n = 6
    tmp = arg
    tmp[:6] = [1,2]
    return len(tmp)>3

try:
    d = [10,11,12,13,14]
    n = 1
    f(d, n)
    print(n, end=';')
    for el in d:
        print(el, end=';')
except:
    print('ERR')
code

9. Що буде виведено за виконання?
code
class A:
    c = 1
    b = 2

    def __init__(self):
        c = 3

class B(A):
    b = 4

    def __init__(self):
        super().__init__()
        self.a = 7

obj = B()
obj.c = obj.c + 10
B.c = 13

try: print(obj.c, end= ';')
except: print('ERR', end=';')

try: print(A.c, end= ';')
except: print('ERR', end=';')

try: print(B.c, end= ';')
except: print('ERR', end=';')
code

10. 
Для графа на рисунку з вершини 4 запущено алгоритм пошуку в глибину, що будує кістякове дерево. Списки суміжності вершин переглядаються алгоритмом за зростанням міток вершин.\par
\begin{center}
	\adjustimage{max size={0.9\linewidth}{0.9\paperheight}}
	{../10/Traversals/171.png}
\end{center}
\par Які з наведених ребер входять до складу отриманого кістякового дерева?\par
1) (4, 6)\par
2) (0, 6)\par
3) (0, 5)\par
4) (2, 4)\par
5) (1, 3)\par
6) (3, 5)\par
{\vskip 15pt} У формі вибрати номери відповідних ребер.\par
%answer:2 3 4 5 6
\end {large}}
%end
{\smallskip \begin {small} Затверджено на засіданні кафедри теоретичної кібернетики, протокол \No 12 від   19.05.2020 р.\par\smallskip\smallskip
{\parbox {6.5 cm}{Зав. кафедри \dotfill Ю.В.~Крак}}\hspace {0.7cm}{\hbox to 6.5 cm{ Екзаменатор \dotfill Т.О.~Карнаух}} \end{small}}
%begin
{{\begin {small}\par
{ \center  Київський національний університет імені Тараса Шевченка \par}
{\parbox {5 cm}{Освітня програма \hfill}{\parbox {7 cm}{Прикладна математика, Системний аналіз \hfill}}\parbox {6 cm}{\hfill Семестр 2}}\par
{\parbox {5 cm} {Навчальний предмет \hfill}\parbox {7 cm}{Програмування \hfill}\parbox {6cm}{\hfill}\par}\vspace{-7pt} \end{small}}{\center Екзаменаційний білет \No \textbf 
{ 114 }\par}}
{
\begin {large}
1. Що буде виведено за виконання?
code
try:
    res = 6//5%4*4
    if isinstance(res, bool):
        print(f'bool;{res}')
    elif isinstance(res, int):
        print(f'int;{res}')
    elif isinstance(res, float):
        print(f'float;{res:.2f}')
    else:
        print('???')
except:
    print('ERR')
code

2. Що буде виведено за виконання?
code
a = 1
b = 4
c = 9

def h(b):
    global c
    a = 4
    b -= 1
    c = 2
    return a + b + c

a = 8
b = 5
c = 0

try:
    print(h(b), a, b, c, sep=';')
except:
    print('ERR', a, b, c, sep=';')
code

3. Що буде виведено за виконання?
code
try:
    k = 15
    while k>=5:
        k += -3
    print(k, end=';')
except:
    print('ERR')
code

4. Що буде виведено за виконання?
code
try:
    for a in range(14, -5, 3):
        if a>6:
            continue
        print(a, end=';')
    else:
        print('end',end=';')
    print(a)
except:
    print('ERR')
code

5. Що буде виведено за виконання?\par (Дійсні значення, якщо наявні, в цьому завданні будуть виводитись у форматі з фіксованою крапкою та, можливо, з доволі великою кількістю десяткових розрядів. У відповіді для дійсних значень у ЦЬОМУ завданні достатньо вказати один десятковий розряд після крапки, відкинувши решту розрядів без округлення. Наприклад, замість 13.66666666666667 достатньо вказати 13.6, замість 25.23 достатньо вказати 25.2)
code
try:
    print(5, end=';')
    print(int('d7'), end=';')
    print(8, end=';')
except TypeError:
    print(3, end=';')
except ValueError:
    print(4, end=';')
except:
    print('ERR', end=';')
else:
    print(9, end=';')
finally:
    print(0)
code

6. Що буде виведено за виконання?
code
def f(n):
    print(n%10, end='')
    if n>9:
        f(n//100)
 
f(987654)
code

7. Що буде виведено за виконання?
code
class A:
    c = 1

    def __f1(self): print(2, end=';')
    def _f2(self): print(3, end=';')

    def main(self):
        try: print(self.c, end=';')
        except: print('ERR', end=';')

        try: self.__f1()
        except: print('ERR', end=';')

        try: self._f2()
        except: print('ERR', end=';')


class B(A):
    c = 4

    def __f1(self): print(5, end=';')
    def _f2(self): print(6, end=';')

try: B().main()
except: print('ERR', end=';')
code

8. Що буде виведено за виконання?
code
def f(arg, n):
    arg[22] = 1
    n = 6
    arg = set(arg.items())

try:
    n = 1
    s = {75:6, 72:3, 75:3}
    f(s, n)
    print(n, end=';')
    for x in s.keys():
        print(x, sep=';', end=';')
except:
    print('ERR')
code

9. Що буде виведено за виконання?
code
class A:
    c = 1
    a = 2

    def __init__(self):
        a = 3

class B(A):
    c = 4

    def __init__(self):
        super().__init__()
        self.b = 7

obj = B()
obj.a = obj.b + 10
B.a = 13

try: print(obj.a, end= ';')
except: print('ERR', end=';')

try: print(A.a, end= ';')
except: print('ERR', end=';')

try: print(B.a, end= ';')
except: print('ERR', end=';')
code

10. 
Для графа на рисунку з вершини 5 запущено алгоритм пошуку в глибину, що будує кістякове дерево. Списки суміжності вершин переглядаються алгоритмом за зростанням міток вершин.\par
\begin{center}
	\adjustimage{max size={0.9\linewidth}{0.9\paperheight}}
	{../10/Traversals/155.png}
\end{center}
\par Які з наведених ребер входять до складу отриманого кістякового дерева?\par
1) (1, 3)\par
2) (2, 4)\par
3) (0, 5)\par
4) (0, 1)\par
5) (0, 2)\par
6) (0, 3)\par
{\vskip 15pt} У формі вибрати номери відповідних ребер.\par
%answer:1 2 3 4
\end {large}}
%end
{\smallskip \begin {small} Затверджено на засіданні кафедри теоретичної кібернетики, протокол \No 12 від   19.05.2020 р.\par\smallskip\smallskip
{\parbox {6.5 cm}{Зав. кафедри \dotfill Ю.В.~Крак}}\hspace {0.7cm}{\hbox to 6.5 cm{ Екзаменатор \dotfill Т.О.~Карнаух}} \end{small}}
%begin
{{\begin {small}\par
{ \center  Київський національний університет імені Тараса Шевченка \par}
{\parbox {5 cm}{Освітня програма \hfill}{\parbox {7 cm}{Прикладна математика, Системний аналіз \hfill}}\parbox {6 cm}{\hfill Семестр 2}}\par
{\parbox {5 cm} {Навчальний предмет \hfill}\parbox {7 cm}{Програмування \hfill}\parbox {6cm}{\hfill}\par}\vspace{-7pt} \end{small}}{\center Екзаменаційний білет \No \textbf 
{ 115 }\par}}
{
\begin {large}
1. Що буде виведено за виконання?
code
try:
    res = 7.0!="False"
    if isinstance(res, bool):
        print(f'bool;{res}')
    elif isinstance(res, int):
        print(f'int;{res}')
    elif isinstance(res, float):
        print(f'float;{res:.2f}')
    else:
        print('???')
except:
    print('ERR')
code

2. Що буде виведено за виконання?
code
a = 6
b = 3
c = 0

def h(a):
    global c
    a = 2
    b = 1
    c = 3
    return a + b + c

a = 4
b = 8
c = 9

try:
    print(h(a), a, b, c, sep=';')
except:
    print('ERR', a, b, c, sep=';')
code

3. Що буде виведено за виконання?
code
try:
    m = 7
    while m<9:
        m += 2
    print(m, end=';')
except:
    print('ERR')
code

4. Що буде виведено за виконання?
code
try:
    for c in range(-14, 13, -1):
        if c<=8:
            continue
        print(c, end=';')
    else:
        print('end',end=';')
    print(c)
except:
    print('ERR')
code

5. Що буде виведено за виконання?\par (Дійсні значення, якщо наявні, в цьому завданні будуть виводитись у форматі з фіксованою крапкою та, можливо, з доволі великою кількістю десяткових розрядів. У відповіді для дійсних значень у ЦЬОМУ завданні достатньо вказати один десятковий розряд після крапки, відкинувши решту розрядів без округлення. Наприклад, замість 13.66666666666667 достатньо вказати 13.6, замість 25.23 достатньо вказати 25.2)
code
try:
    print(2, end=';')
    print(6//2, end=';')
    print(9, end=';')
except ZeroDivisionError:
    print(0, end=';')
except Exception:
    print(3, end=';')
except:
    print('ERR', end=';')
else:
    print(1, end=';')
finally:
    print(2)
code

6. Що буде виведено за виконання?
code
def f(s):
    for i in range(len(s)):
        s[i] += 1
    return
 
s = [1, 2, 3]
f(s)
print(s)
code

7. Що буде виведено за виконання?
code
class A:
    d = 1

    def g1(self): print(2, end=';')
    def _g2(self): print(3, end=';')

    def main(self):
        try: print(self.d, end=';')
        except: print('ERR', end=';')

        try: self.g1()
        except: print('ERR', end=';')

        try: self._g2()
        except: print('ERR', end=';')


class B(A):
    d = 4

    def g1(self): print(5, end=';')
    def _g2(self): print(6, end=';')

try: B().main()
except: print('ERR', end=';')
code

8. Що буде виведено за виконання?
code
def f(arg, n):
    arg[35] = 0
    n = 7
    arg = tuple(arg.keys())

try:
    n = 4
    d = {57:7, 14:0, 36:7, 76:7}
    f(d, n)
    print(n, end=';')
    for x, y in d.items():
        print(x, y, sep=';', end=';')
except:
    print('ERR')
code

9. Що буде виведено за виконання?
code
class A:
    a = 1
    c = 2

    def __init__(self):
        c = 3

class B(A):
    c = 4

    def __init__(self):
        super().__init__()
        self.b = 7

obj = B()
obj.c = obj.a + 10
B.c = 13

try: print(obj.c, end= ';')
except: print('ERR', end=';')

try: print(A.c, end= ';')
except: print('ERR', end=';')

try: print(B.c, end= ';')
except: print('ERR', end=';')
code

10. 
Для графа на рисунку з вершини 3 запущено алгоритм пошуку в глибину, що будує кістякове дерево. Списки суміжності вершин переглядаються алгоритмом за зростанням міток вершин.\par
\begin{center}
	\adjustimage{max size={0.9\linewidth}{0.9\paperheight}}
	{../10/Traversals/123.png}
\end{center}
\par Які з наведених ребер входять до складу отриманого кістякового дерева?\par
1) (3, 4)\par
2) (1, 3)\par
3) (1, 5)\par
4) (2, 5)\par
5) (0, 6)\par
6) (1, 2)\par
{\vskip 15pt} У формі вибрати номери відповідних ребер.\par
%answer:1 2 4 5 6
\end {large}}
%end
{\smallskip \begin {small} Затверджено на засіданні кафедри теоретичної кібернетики, протокол \No 12 від   19.05.2020 р.\par\smallskip\smallskip
{\parbox {6.5 cm}{Зав. кафедри \dotfill Ю.В.~Крак}}\hspace {0.7cm}{\hbox to 6.5 cm{ Екзаменатор \dotfill Т.О.~Карнаух}} \end{small}}
%begin
{{\begin {small}\par
{ \center  Київський національний університет імені Тараса Шевченка \par}
{\parbox {5 cm}{Освітня програма \hfill}{\parbox {7 cm}{Прикладна математика, Системний аналіз \hfill}}\parbox {6 cm}{\hfill Семестр 2}}\par
{\parbox {5 cm} {Навчальний предмет \hfill}\parbox {7 cm}{Програмування \hfill}\parbox {6cm}{\hfill}\par}\vspace{-7pt} \end{small}}{\center Екзаменаційний білет \No \textbf 
{ 116 }\par}}
{
\begin {large}
1. Що буде виведено за виконання?
code
def f(n):
    print('f', end='')
    return n

try:
    print(f(-9) and f(5))
except:
    print('ERR')
code

2. Що буде виведено за виконання?
code
a = 6
b = 3
c = 2

def g(b):
    a += 3
    b = 4
    c = 1
    return a + b + c

a = 8
b = 5
c = 3

try:
    print(g(a), a, b, c, sep=';')
except:
    print('ERR', a, b, c, sep=';')
code

3. Що буде виведено за виконання?
code
try:
    i = 6
    while i<=6:
        i += 1
    print(i, end=';')
except:
    print('ERR')
code

4. Що буде виведено за виконання?
code
try:
    for a in range(12, 14, -2):
        if a<7:
            continue
        print(a, end=';')
    else:
        print(a,end=';')
    print(a)
except:
    print('ERR')
code

5. Що буде виведено за виконання?\par (Дійсні значення, якщо наявні, в цьому завданні будуть виводитись у форматі з фіксованою крапкою та, можливо, з доволі великою кількістю десяткових розрядів. У відповіді для дійсних значень у ЦЬОМУ завданні достатньо вказати один десятковий розряд після крапки, відкинувши решту розрядів без округлення. Наприклад, замість 13.66666666666667 достатньо вказати 13.6, замість 25.23 достатньо вказати 25.2)
code
try:
    print(4, end=';')
    print(int('b7'), end=';')
    print(8, end=';')
except ZeroDivisionError:
    print(5, end=';')
except TypeError:
    print(6, end=';')
except:
    print('ERR', end=';')
else:
    print(8, end=';')
finally:
    print(6)
code

6. Що буде виведено за виконання?
code
def f(n):
    if n>9:
        f(n//100)
        print(n%10, end='')
 
f(123456)
code

7. Що буде виведено за виконання?
code
class A:
    a = 1

    def _f1(self): print(2, end=';')
    def f2(self): print(3, end=';')

    def main(self):
        try: print(self.a, end=';')
        except: print('ERR', end=';')

        try: self._f1()
        except: print('ERR', end=';')

        try: self.f2()
        except: print('ERR', end=';')


class B(A):
    a = 4

    def _f1(self): print(5, end=';')
    def f2(self): print(6, end=';')

try: B().main()
except: print('ERR', end=';')
code

8. Що буде виведено за виконання?
code
def f(arg, n):
    arg[51] = 9
    n = 9
    arg = set(arg.keys())

try:
    n = 3
    t = {84:6, 30:2, 84:5}
    f(t, n)
    print(n, end=';')
    for x in t.items():
        print(*x, sep=';', end=';')
except:
    print('ERR')
code

9. Що буде виведено за виконання?
code
class A:
    c = 1
    a = 2

    def __init__(self):
        c = 3

class B(A):
    c = 4

    def __init__(self):
        super().__init__()
        self.b = 7

obj = B()
obj.b = obj.a + 10
B.b = 13

try: print(obj.b, end= ';')
except: print('ERR', end=';')

try: print(A.b, end= ';')
except: print('ERR', end=';')

try: print(B.b, end= ';')
except: print('ERR', end=';')
code

10. 
Для графа на рисунку з вершини 6 запущено алгоритм пошуку в глибину, що будує кістякове дерево. Списки суміжності вершин переглядаються алгоритмом за зростанням міток вершин.\par
\begin{center}
	\adjustimage{max size={0.9\linewidth}{0.9\paperheight}}
	{../10/Traversals/85.png}
\end{center}
\par Які з наведених ребер входять до складу отриманого кістякового дерева?\par
1) (0, 5)\par
2) (0, 6)\par
3) (1, 6)\par
4) (5, 6)\par
5) (1, 3)\par
6) (0, 2)\par
{\vskip 15pt} У формі вибрати номери відповідних ребер.\par
%answer:2 5 6
\end {large}}
%end
{\smallskip \begin {small} Затверджено на засіданні кафедри теоретичної кібернетики, протокол \No 12 від   19.05.2020 р.\par\smallskip\smallskip
{\parbox {6.5 cm}{Зав. кафедри \dotfill Ю.В.~Крак}}\hspace {0.7cm}{\hbox to 6.5 cm{ Екзаменатор \dotfill Т.О.~Карнаух}} \end{small}}
%begin
{{\begin {small}\par
{ \center  Київський національний університет імені Тараса Шевченка \par}
{\parbox {5 cm}{Освітня програма \hfill}{\parbox {7 cm}{Прикладна математика, Системний аналіз \hfill}}\parbox {6 cm}{\hfill Семестр 2}}\par
{\parbox {5 cm} {Навчальний предмет \hfill}\parbox {7 cm}{Програмування \hfill}\parbox {6cm}{\hfill}\par}\vspace{-7pt} \end{small}}{\center Екзаменаційний білет \No \textbf 
{ 117 }\par}}
{
\begin {large}
1. Що буде виведено за виконання?
code
try:
    res = 5.0==8 or not 2>4 or False<3.0
    if isinstance(res, bool):
        print(f'bool;{res}')
    elif isinstance(res, int):
        print(f'int;{res}')
    elif isinstance(res, float):
        print(f'float;{res:.2f}')
    else:
        print('???')
except:
    print('ERR')
code

2. Що буде виведено за виконання?
code
a = 0
b = 7
c = 2

def h(b):
    a = 4
    b = 1
    c = 2
    return a + b + c

a = 4
b = 3
c = 1

try:
    print(h(b), a, b, c, sep=';')
except:
    print('ERR', a, b, c, sep=';')
code

3. Що буде виведено за виконання?
code
try:
    i = 8
    while i>4:
        i += -3
    print(i, end=';')
except:
    print('ERR')
code

4. Що буде виведено за виконання?
code
try:
    for d in range(8, 5, 2):
        if d>=6:
            break
        print(d, end=';');d = -9
    else:
        print(d,end=';')
    print(d)
except:
    print('ERR')
code

5. Що буде виведено за виконання?\par (Дійсні значення, якщо наявні, в цьому завданні будуть виводитись у форматі з фіксованою крапкою та, можливо, з доволі великою кількістю десяткових розрядів. У відповіді для дійсних значень у ЦЬОМУ завданні достатньо вказати один десятковий розряд після крапки, відкинувши решту розрядів без округлення. Наприклад, замість 13.66666666666667 достатньо вказати 13.6, замість 25.23 достатньо вказати 25.2)
code
try:
    print(0, end=';')
    print(1/3, end=';')
    print(2, end=';')
except KeyboardInterrupt:
    print(5, end=';')
except Exception:
    print(9, end=';')
except:
    print('ERR', end=';')
else:
    print(3, end=';')
finally:
    print(4)
code

6. Що буде виведено за виконання?
code
def f(n):
    if n<=9:
        print(n%10,end='')
    else:
        f(n//10)

f(543216)
code

7. Що буде виведено за виконання?
code
class A:
    a = 1

    def _h1(self): print(2, end=';')
    def __h2(self): print(3, end=';')

    def main(self):
        try: print(self.a, end=';')
        except: print('ERR', end=';')

        try: self._h1()
        except: print('ERR', end=';')

        try: self.__h2()
        except: print('ERR', end=';')


class B(A):
    a = 4

    def _h1(self): print(5, end=';')
    def __h2(self): print(6, end=';')

try: B().main()
except: print('ERR', end=';')
code

8. Що буде виведено за виконання?
code
def f(arg, n):
    n = 7
    tmp = arg
    tmp[-7:] = ()
    return len(tmp)>3

try:
    f = [10,11,12,13,14]
    n = 9
    f(f, n)
    print(n, end=';')
    for el in f:
        print(el, end=';')
except:
    print('ERR')
code

9. Що буде виведено за виконання?
code
class A:
    c = 1
    b = 2

    def __init__(self):
        c = 3

class B(A):
    b = 4

    def __init__(self):
        super().__init__()
        self.a = 7

obj = B()
obj.b = obj.c + 10
B.b = 13

try: print(obj.b, end= ';')
except: print('ERR', end=';')

try: print(A.b, end= ';')
except: print('ERR', end=';')

try: print(B.b, end= ';')
except: print('ERR', end=';')
code

10. 
Для графа на рисунку з вершини 2 запущено алгоритм пошуку в глибину, що будує кістякове дерево. Списки суміжності вершин переглядаються алгоритмом за зростанням міток вершин.\par
\begin{center}
	\adjustimage{max size={0.9\linewidth}{0.9\paperheight}}
	{../10/Traversals/132.png}
\end{center}
\par Які з наведених ребер входять до складу отриманого кістякового дерева?\par
1) (0, 4)\par
2) (0, 5)\par
3) (0, 1)\par
4) (1, 2)\par
5) (3, 6)\par
6) (2, 4)\par
{\vskip 15pt} У формі вибрати номери відповідних ребер.\par
%answer:1 2 3 4 5
\end {large}}
%end
{\smallskip \begin {small} Затверджено на засіданні кафедри теоретичної кібернетики, протокол \No 12 від   19.05.2020 р.\par\smallskip\smallskip
{\parbox {6.5 cm}{Зав. кафедри \dotfill Ю.В.~Крак}}\hspace {0.7cm}{\hbox to 6.5 cm{ Екзаменатор \dotfill Т.О.~Карнаух}} \end{small}}
%begin
{{\begin {small}\par
{ \center  Київський національний університет імені Тараса Шевченка \par}
{\parbox {5 cm}{Освітня програма \hfill}{\parbox {7 cm}{Прикладна математика, Системний аналіз \hfill}}\parbox {6 cm}{\hfill Семестр 2}}\par
{\parbox {5 cm} {Навчальний предмет \hfill}\parbox {7 cm}{Програмування \hfill}\parbox {6cm}{\hfill}\par}\vspace{-7pt} \end{small}}{\center Екзаменаційний білет \No \textbf 
{ 118 }\par}}
{
\begin {large}
1. Що буде виведено за виконання?
code
try:
    res = 7j>=True
    if isinstance(res, bool):
        print(f'bool;{res}')
    elif isinstance(res, int):
        print(f'int;{res}')
    elif isinstance(res, float):
        print(f'float;{res:.2f}')
    else:
        print('???')
except:
    print('ERR')
code

2. Що буде виведено за виконання?
code
a = 6
b = 5
c = 6

def f(b):
    a = 5
    b = 2
    c = 3
    return a + b + c

a = 4
b = 1
c = 0

try:
    print(f(a), a, b, c, sep=';')
except:
    print('ERR', a, b, c, sep=';')
code

3. Що буде виведено за виконання?
code
try:
    n = 2
    while n<5:
        n += 2
    print(n, end=';')
except:
    print('ERR')
code

4. Що буде виведено за виконання?
code
try:
    for f in range(3, 0, 3):
        if f<4:
            break
        print(f, end=';')
        if f<=5:
            break
    else:
        print(f,end=';')
    print(f)
except:
    print('ERR')
code

5. Що буде виведено за виконання?\par (Дійсні значення, якщо наявні, в цьому завданні будуть виводитись у форматі з фіксованою крапкою та, можливо, з доволі великою кількістю десяткових розрядів. У відповіді для дійсних значень у ЦЬОМУ завданні достатньо вказати один десятковий розряд після крапки, відкинувши решту розрядів без округлення. Наприклад, замість 13.66666666666667 достатньо вказати 13.6, замість 25.23 достатньо вказати 25.2)
code
try:
    print(7, end=';')
    print(0%0.0, end=';')
    print(8, end=';')
except ValueError:
    print(9, end=';')
except ValueError:
    print(2, end=';')
except:
    print('ERR', end=';')
else:
    print(7, end=';')
finally:
    print(3)
code

6. Що буде виведено за виконання?
code
def f(n):
    if n>9:
        return f(n//10)+1
    else:
        return 1

print(f(123456))
code

7. Що буде виведено за виконання?
code
class A:
    d = 1

    def f1(self): print(2, end=';')
    def _f2(self): print(3, end=';')

    def main(self):
        try: print(self.d, end=';')
        except: print('ERR', end=';')

        try: self.f1()
        except: print('ERR', end=';')

        try: self._f2()
        except: print('ERR', end=';')


class B(A):
    d = 4

    def f1(self): print(5, end=';')
    def _f2(self): print(6, end=';')

try: B().main()
except: print('ERR', end=';')
code

8. Що буде виведено за виконання?
code
def f(arg, n):
    arg[42] = 3
    n = 5
    arg = tuple(arg.values())

try:
    n = 2
    d = {42:1, 11:9, 58:8, 58:7}
    f(d, n)
    print(n, end=';')
    for x in d.items():
        print(x, sep=';', end=';')
except:
    print('ERR')
code

9. Що буде виведено за виконання?
code
class A:
    b = 1
    a = 2

    def __init__(self):
        a = 3

class B(A):
    b = 4

    def __init__(self):
        super().__init__()
        self.c = 7

obj = B()
obj.a = obj.c + 10
B.a = 13

try: print(obj.a, end= ';')
except: print('ERR', end=';')

try: print(A.a, end= ';')
except: print('ERR', end=';')

try: print(B.a, end= ';')
except: print('ERR', end=';')
code

10. 
Для графа на рисунку з вершини 0 запущено алгоритм пошуку в глибину, що будує кістякове дерево. Списки суміжності вершин переглядаються алгоритмом за зростанням міток вершин.\par
\begin{center}
	\adjustimage{max size={0.9\linewidth}{0.9\paperheight}}
	{../10/Traversals/77.png}
\end{center}
\par Які з наведених ребер входять до складу отриманого кістякового дерева?\par
1) (0, 1)\par
2) (2, 6)\par
3) (2, 5)\par
4) (0, 3)\par
5) (1, 5)\par
6) (1, 6)\par
{\vskip 15pt} У формі вибрати номери відповідних ребер.\par
%answer:1 2 3
\end {large}}
%end
{\smallskip \begin {small} Затверджено на засіданні кафедри теоретичної кібернетики, протокол \No 12 від   19.05.2020 р.\par\smallskip\smallskip
{\parbox {6.5 cm}{Зав. кафедри \dotfill Ю.В.~Крак}}\hspace {0.7cm}{\hbox to 6.5 cm{ Екзаменатор \dotfill Т.О.~Карнаух}} \end{small}}
%begin
{{\begin {small}\par
{ \center  Київський національний університет імені Тараса Шевченка \par}
{\parbox {5 cm}{Освітня програма \hfill}{\parbox {7 cm}{Прикладна математика, Системний аналіз \hfill}}\parbox {6 cm}{\hfill Семестр 2}}\par
{\parbox {5 cm} {Навчальний предмет \hfill}\parbox {7 cm}{Програмування \hfill}\parbox {6cm}{\hfill}\par}\vspace{-7pt} \end{small}}{\center Екзаменаційний білет \No \textbf 
{ 119 }\par}}
{
\begin {large}
1. Що буде виведено за виконання?
code
try:
    res = 2j<1.0j
    if isinstance(res, bool):
        print(f'bool;{res}')
    elif isinstance(res, int):
        print(f'int;{res}')
    elif isinstance(res, float):
        print(f'float;{res:.2f}')
    else:
        print('???')
except:
    print('ERR')
code

2. Що буде виведено за виконання?
code
a = 3
b = 2
c = 5

def g(a):
    global c
    a = 1
    b = 4
    c = 5
    return a + b + c

a = 8
b = 9
c = 7

try:
    print(g(a), a, b, c, sep=';')
except:
    print('ERR', a, b, c, sep=';')
code

3. Що буде виведено за виконання?
code
try:
    k = 1
    while k<=7:
        k += 1
    print(k, end=';')
except:
    print('ERR')
code

4. Що буде виведено за виконання?
code
try:
    for b in range(0, 1, -1):
        if b>3:
            continue
        print(b, end=';')
    else:
        print(b,end=';')
    print(b)
except:
    print('ERR')
code

5. Що буде виведено за виконання?\par (Дійсні значення, якщо наявні, в цьому завданні будуть виводитись у форматі з фіксованою крапкою та, можливо, з доволі великою кількістю десяткових розрядів. У відповіді для дійсних значень у ЦЬОМУ завданні достатньо вказати один десятковий розряд після крапки, відкинувши решту розрядів без округлення. Наприклад, замість 13.66666666666667 достатньо вказати 13.6, замість 25.23 достатньо вказати 25.2)
code
try:
    print(6, end=';')
    print(int('4'), end=';')
    print(5, end=';')
except KeyboardInterrupt:
    print(1, end=';')
except Exception:
    print(9, end=';')
except:
    print('ERR', end=';')
else:
    print(3, end=';')
finally:
    print(6)
code

6. Що буде виведено за виконання?
code
def f(n):
    if n>2:
        f(n//2)
    else:
        print(n%2, end='')

f(16)
code

7. Що буде виведено за виконання?
code
class A:
    c = 1

    def _g1(self): print(2, end=';')
    def _g2(self): print(3, end=';')

    def main(self):
        try: print(self.c, end=';')
        except: print('ERR', end=';')

        try: self._g1()
        except: print('ERR', end=';')

        try: self._g2()
        except: print('ERR', end=';')


class B(A):
    c = 4

    def _g1(self): print(5, end=';')
    def _g2(self): print(6, end=';')

try: B().main()
except: print('ERR', end=';')
code

8. Що буде виведено за виконання?
code
def f(arg, n):
    n = 8
    tmp = arg
    tmp[4:2] = (13)
    return len(tmp)>3

try:
    a = ['0','1','2','3','4']
    n = 5
    f(a, n)
    print(n, end=';')
    for el in a:
        print(el, end=';')
except:
    print('ERR')
code

9. Що буде виведено за виконання?
code
class A:
    a = 1
    b = 2

    def __init__(self):
        b = 3

class B(A):
    a = 4

    def __init__(self):
        super().__init__()
        self.c = 7

obj = B()
obj.b = obj.a + 10
B.b = 13

try: print(obj.b, end= ';')
except: print('ERR', end=';')

try: print(A.b, end= ';')
except: print('ERR', end=';')

try: print(B.b, end= ';')
except: print('ERR', end=';')
code

10. 
Для графа на рисунку з вершини 0 запущено алгоритм пошуку в глибину, що будує кістякове дерево. Списки суміжності вершин переглядаються алгоритмом за зростанням міток вершин.\par
\begin{center}
	\adjustimage{max size={0.9\linewidth}{0.9\paperheight}}
	{../10/Traversals/245.png}
\end{center}
\par Які з наведених ребер входять до складу отриманого кістякового дерева?\par
1) (0, 6)\par
2) (0, 2)\par
3) (0, 4)\par
4) (1, 6)\par
5) (4, 5)\par
6) (0, 5)\par
{\vskip 15pt} У формі вибрати номери відповідних ребер.\par
%answer:2 4
\end {large}}
%end
{\smallskip \begin {small} Затверджено на засіданні кафедри теоретичної кібернетики, протокол \No 12 від   19.05.2020 р.\par\smallskip\smallskip
{\parbox {6.5 cm}{Зав. кафедри \dotfill Ю.В.~Крак}}\hspace {0.7cm}{\hbox to 6.5 cm{ Екзаменатор \dotfill Т.О.~Карнаух}} \end{small}}
%begin
{{\begin {small}\par
{ \center  Київський національний університет імені Тараса Шевченка \par}
{\parbox {5 cm}{Освітня програма \hfill}{\parbox {7 cm}{Прикладна математика, Системний аналіз \hfill}}\parbox {6 cm}{\hfill Семестр 2}}\par
{\parbox {5 cm} {Навчальний предмет \hfill}\parbox {7 cm}{Програмування \hfill}\parbox {6cm}{\hfill}\par}\vspace{-7pt} \end{small}}{\center Екзаменаційний білет \No \textbf 
{ 120 }\par}}
{
\begin {large}
1. Що буде виведено за виконання?
code
try:
    res = 3/6*6*8
    if isinstance(res, bool):
        print(f'bool;{res}')
    elif isinstance(res, int):
        print(f'int;{res}')
    elif isinstance(res, float):
        print(f'float;{res:.2f}')
    else:
        print('???')
except:
    print('ERR')
code

2. Що буде виведено за виконання?
code
a = 7
b = 6
c = 1

def f(a):
    global c
    a *= 3
    b = 5
    c = 2
    return a + b + c

a = 0
b = 4
c = 9

try:
    print(f(b), a, b, c, sep=';')
except:
    print('ERR', a, b, c, sep=';')
code

3. Що буде виведено за виконання?
code
try:
    n = 9
    while n>=1:
        n += -2
    print(n, end=';')
except:
    print('ERR')
code

4. Що буде виведено за виконання?
code
try:
    for e in range(-3, -10, -3):
        if e<=5:
            continue
        print(e, end=';');e = -8
    else:
        print(e,end=';')
    print(e)
except:
    print('ERR')
code

5. Що буде виведено за виконання?\par (Дійсні значення, якщо наявні, в цьому завданні будуть виводитись у форматі з фіксованою крапкою та, можливо, з доволі великою кількістю десяткових розрядів. У відповіді для дійсних значень у ЦЬОМУ завданні достатньо вказати один десятковий розряд після крапки, відкинувши решту розрядів без округлення. Наприклад, замість 13.66666666666667 достатньо вказати 13.6, замість 25.23 достатньо вказати 25.2)
code
try:
    print(5, end=';')
    print(8j>=4j, end=';')
    print(1, end=';')
except ZeroDivisionError:
    print(4, end=';')
except TypeError:
    print(2, end=';')
except:
    print('ERR', end=';')
else:
    print(7, end=';')
finally:
    print(0)
code

6. Що буде виведено за виконання?
code
def f(n):
    if n>9:
        f(n//10)
    print(n%10, end='')
 
f(123456)
code

7. Що буде виведено за виконання?
code
class A:
    b = 1

    def __h1(self): print(2, end=';')
    def _h2(self): print(3, end=';')

    def main(self):
        try: print(self.b, end=';')
        except: print('ERR', end=';')

        try: self.__h1()
        except: print('ERR', end=';')

        try: self._h2()
        except: print('ERR', end=';')


class B(A):
    b = 4

    def __h1(self): print(5, end=';')
    def _h2(self): print(6, end=';')

try: B().main()
except: print('ERR', end=';')
code

8. Що буде виведено за виконання?
code
def f(arg, n):
    arg[51] = 7
    n = 8
    arg = list(arg.items())

try:
    n = 3
    s = {51:1, 38:6, 38:1}
    f(s, n)
    print(n, end=';')
    for x in s.values():
        print(x, sep=';', end=';')
except:
    print('ERR')
code

9. Що буде виведено за виконання?
code
class A:
    b = 1
    c = 2

    def __init__(self):
        b = 3

class B(A):
    c = 4

    def __init__(self):
        super().__init__()
        self.a = 7

obj = B()
obj.c = obj.b + 10
B.c = 13

try: print(obj.c, end= ';')
except: print('ERR', end=';')

try: print(A.c, end= ';')
except: print('ERR', end=';')

try: print(B.c, end= ';')
except: print('ERR', end=';')
code

10. 
Для графа на рисунку з вершини 5 запущено алгоритм пошуку в глибину, що будує кістякове дерево. Списки суміжності вершин переглядаються алгоритмом за зростанням міток вершин.\par
\begin{center}
	\adjustimage{max size={0.9\linewidth}{0.9\paperheight}}
	{../10/Traversals/106.png}
\end{center}
\par Які з наведених ребер входять до складу отриманого кістякового дерева?\par
1) (1, 5)\par
2) (3, 5)\par
3) (4, 6)\par
4) (0, 4)\par
5) (1, 2)\par
6) (3, 6)\par
{\vskip 15pt} У формі вибрати номери відповідних ребер.\par
%answer:1 3 4 5 6
\end {large}}
%end
{\smallskip \begin {small} Затверджено на засіданні кафедри теоретичної кібернетики, протокол \No 12 від   19.05.2020 р.\par\smallskip\smallskip
{\parbox {6.5 cm}{Зав. кафедри \dotfill Ю.В.~Крак}}\hspace {0.7cm}{\hbox to 6.5 cm{ Екзаменатор \dotfill Т.О.~Карнаух}} \end{small}}
%begin
{{\begin {small}\par
{ \center  Київський національний університет імені Тараса Шевченка \par}
{\parbox {5 cm}{Освітня програма \hfill}{\parbox {7 cm}{Прикладна математика, Системний аналіз \hfill}}\parbox {6 cm}{\hfill Семестр 2}}\par
{\parbox {5 cm} {Навчальний предмет \hfill}\parbox {7 cm}{Програмування \hfill}\parbox {6cm}{\hfill}\par}\vspace{-7pt} \end{small}}{\center Екзаменаційний білет \No \textbf 
{ 121 }\par}}
{
\begin {large}
1. Що буде виведено за виконання?
code
try:
    res = 8/8/7*7
    if isinstance(res, bool):
        print(f'bool;{res}')
    elif isinstance(res, int):
        print(f'int;{res}')
    elif isinstance(res, float):
        print(f'float;{res:.2f}')
    else:
        print('???')
except:
    print('ERR')
code

2. Що буде виведено за виконання?
code
a = 8
b = 1
c = 0

def h(a):
    global c
    a = 4
    b = 2
    c = 1
    return a + b + c

a = 3
b = 7
c = 5

try:
    print(h(b), a, b, c, sep=';')
except:
    print('ERR', a, b, c, sep=';')
code

3. Що буде виведено за виконання?
code
try:
    j = 4
    while j<=10:
        j += 2
    print(j, end=';')
except:
    print('ERR')
code

4. Що буде виведено за виконання?
code
try:
    for e in range(-11, 15, -3):
        if e>4:
            break
        print(e, end=';')
    else:
        print(e,end=';')
    print(e)
except:
    print('ERR')
code

5. Що буде виведено за виконання?\par (Дійсні значення, якщо наявні, в цьому завданні будуть виводитись у форматі з фіксованою крапкою та, можливо, з доволі великою кількістю десяткових розрядів. У відповіді для дійсних значень у ЦЬОМУ завданні достатньо вказати один десятковий розряд після крапки, відкинувши решту розрядів без округлення. Наприклад, замість 13.66666666666667 достатньо вказати 13.6, замість 25.23 достатньо вказати 25.2)
code
try:
    print(4, end=';')
    print(2/2, end=';')
    print(7, end=';')
except TypeError:
    print(6, end=';')
except ValueError:
    print(8, end=';')
except:
    print('ERR', end=';')
else:
    print(5, end=';')
finally:
    print(0)
code

6. Що буде виведено за виконання?
code
def f(n):
    if n>9:
        f(n//100)
    print(n%10, end='')
 
f(123456)
code

7. Що буде виведено за виконання?
code
class A:
    d = 1

    def _h1(self): print(2, end=';')
    def __h2(self): print(3, end=';')

    def main(self):
        try: print(self.d, end=';')
        except: print('ERR', end=';')

        try: self._h1()
        except: print('ERR', end=';')

        try: self.__h2()
        except: print('ERR', end=';')


class B(A):
    d = 4

    def _h1(self): print(5, end=';')
    def __h2(self): print(6, end=';')

try: B().main()
except: print('ERR', end=';')
code

8. Що буде виведено за виконання?
code
def f(arg, n):
    n = 6
    tmp = arg
    tmp[:-7] = [1,0]
    return len(tmp)>3

try:
    g = [10,11,12,13,14]
    n = 1
    f(g, n)
    print(n, end=';')
    for el in g:
        print(el, end=';')
except:
    print('ERR')
code

9. Що буде виведено за виконання?
code
class A:
    b = 1
    a = 2

    def __init__(self):
        a = 3

class B(A):
    a = 4

    def __init__(self):
        super().__init__()
        self.c = 7

obj = B()
obj.c = obj.a + 10
B.c = 13

try: print(obj.c, end= ';')
except: print('ERR', end=';')

try: print(A.c, end= ';')
except: print('ERR', end=';')

try: print(B.c, end= ';')
except: print('ERR', end=';')
code

10. 
Для графа на рисунку з вершини 4 запущено алгоритм пошуку в глибину, що будує кістякове дерево. Списки суміжності вершин переглядаються алгоритмом за зростанням міток вершин.\par
\begin{center}
	\adjustimage{max size={0.9\linewidth}{0.9\paperheight}}
	{../10/Traversals/4.png}
\end{center}
\par Які з наведених ребер входять до складу отриманого кістякового дерева?\par
1) (0, 6)\par
2) (2, 3)\par
3) (1, 6)\par
4) (1, 3)\par
5) (0, 2)\par
6) (0, 3)\par
{\vskip 15pt} У формі вибрати номери відповідних ребер.\par
%answer:3 4 5 6
\end {large}}
%end
{\smallskip \begin {small} Затверджено на засіданні кафедри теоретичної кібернетики, протокол \No 12 від   19.05.2020 р.\par\smallskip\smallskip
{\parbox {6.5 cm}{Зав. кафедри \dotfill Ю.В.~Крак}}\hspace {0.7cm}{\hbox to 6.5 cm{ Екзаменатор \dotfill Т.О.~Карнаух}} \end{small}}
%begin
{{\begin {small}\par
{ \center  Київський національний університет імені Тараса Шевченка \par}
{\parbox {5 cm}{Освітня програма \hfill}{\parbox {7 cm}{Прикладна математика, Системний аналіз \hfill}}\parbox {6 cm}{\hfill Семестр 2}}\par
{\parbox {5 cm} {Навчальний предмет \hfill}\parbox {7 cm}{Програмування \hfill}\parbox {6cm}{\hfill}\par}\vspace{-7pt} \end{small}}{\center Екзаменаційний білет \No \textbf 
{ 122 }\par}}
{
\begin {large}
1. Що буде виведено за виконання?
code
try:
    res = "8.0"<="9"
    if isinstance(res, bool):
        print(f'bool;{res}')
    elif isinstance(res, int):
        print(f'int;{res}')
    elif isinstance(res, float):
        print(f'float;{res:.2f}')
    else:
        print('???')
except:
    print('ERR')
code

2. Що буде виведено за виконання?
code
a = 2
b = 9
c = 6

def g(b):
    a = 5
    b = 3
    c = 2
    return a + b + c

a = 4
b = 8
c = 9

try:
    print(g(a), a, b, c, sep=';')
except:
    print('ERR', a, b, c, sep=';')
code

3. Що буде виведено за виконання?
code
try:
    t = 0
    while t<11:
        t += 1
    print(t, end=';')
except:
    print('ERR')
code

4. Що буде виведено за виконання?
code
try:
    for f in range(-6, -2, 2):
        if f<=3:
            continue
        print(f, end=';')
    else:
        print(f,end=';')
    print(f)
except:
    print('ERR')
code

5. Що буде виведено за виконання?\par (Дійсні значення, якщо наявні, в цьому завданні будуть виводитись у форматі з фіксованою крапкою та, можливо, з доволі великою кількістю десяткових розрядів. У відповіді для дійсних значень у ЦЬОМУ завданні достатньо вказати один десятковий розряд після крапки, відкинувши решту розрядів без округлення. Наприклад, замість 13.66666666666667 достатньо вказати 13.6, замість 25.23 достатньо вказати 25.2)
code
try:
    print(3, end=';')
    print(int('d9'), end=';')
    print(1, end=';')
except Exception:
    print(1, end=';')
except KeyboardInterrupt:
    print(3, end=';')
except:
    print('ERR', end=';')
else:
    print(6, end=';')
finally:
    print(4)
code

6. Що буде виведено за виконання?
code
def f(n):
    if n>9:
        return f(n//100)+1
    else:
        return 1

print(f(123456))
code

7. Що буде виведено за виконання?
code
class A:
    b = 1

    def __g1(self): print(2, end=';')
    def _g2(self): print(3, end=';')

    def main(self):
        try: print(self.b, end=';')
        except: print('ERR', end=';')

        try: self.__g1()
        except: print('ERR', end=';')

        try: self._g2()
        except: print('ERR', end=';')


class B(A):
    b = 4

    def __g1(self): print(5, end=';')
    def _g2(self): print(6, end=';')

try: B().main()
except: print('ERR', end=';')
code

8. Що буде виведено за виконання?
code
def f(arg, n):
    arg[89] = 3
    n = 9
    arg = tuple(arg.keys())

try:
    n = 4
    d = {77:4, 20:4, 80:0, 80:9}
    f(d, n)
    print(n, end=';')
    for x, y in d.items():
        print(x, y, sep=';', end=';')
except:
    print('ERR')
code

9. Що буде виведено за виконання?
code
class A:
    c = 1
    a = 2

    def __init__(self):
        c = 3

class B(A):
    c = 4

    def __init__(self):
        super().__init__()
        self.b = 7

obj = B()
obj.b = obj.c + 10
B.b = 13

try: print(obj.b, end= ';')
except: print('ERR', end=';')

try: print(A.b, end= ';')
except: print('ERR', end=';')

try: print(B.b, end= ';')
except: print('ERR', end=';')
code

10. 
Для графа на рисунку з вершини 6 запущено алгоритм пошуку в глибину, що будує кістякове дерево. Списки суміжності вершин переглядаються алгоритмом за зростанням міток вершин.\par
\begin{center}
	\adjustimage{max size={0.9\linewidth}{0.9\paperheight}}
	{../10/Traversals/122.png}
\end{center}
\par Які з наведених ребер входять до складу отриманого кістякового дерева?\par
1) (3, 6)\par
2) (2, 3)\par
3) (2, 5)\par
4) (1, 5)\par
5) (3, 5)\par
6) (4, 5)\par
{\vskip 15pt} У формі вибрати номери відповідних ребер.\par
%answer:1 2 6
\end {large}}
%end
{\smallskip \begin {small} Затверджено на засіданні кафедри теоретичної кібернетики, протокол \No 12 від   19.05.2020 р.\par\smallskip\smallskip
{\parbox {6.5 cm}{Зав. кафедри \dotfill Ю.В.~Крак}}\hspace {0.7cm}{\hbox to 6.5 cm{ Екзаменатор \dotfill Т.О.~Карнаух}} \end{small}}
%begin
{{\begin {small}\par
{ \center  Київський національний університет імені Тараса Шевченка \par}
{\parbox {5 cm}{Освітня програма \hfill}{\parbox {7 cm}{Прикладна математика, Системний аналіз \hfill}}\parbox {6 cm}{\hfill Семестр 2}}\par
{\parbox {5 cm} {Навчальний предмет \hfill}\parbox {7 cm}{Програмування \hfill}\parbox {6cm}{\hfill}\par}\vspace{-7pt} \end{small}}{\center Екзаменаційний білет \No \textbf 
{ 123 }\par}}
{
\begin {large}
1. Що буде виведено за виконання?
code
try:
    res = 9**-0.5*-2
    if isinstance(res, bool):
        print(f'bool;{res}')
    elif isinstance(res, int):
        print(f'int;{res}')
    elif isinstance(res, float):
        print(f'float;{res:.2f}')
    else:
        print('???')
except:
    print('ERR')
code

2. Що буде виведено за виконання?
code
a = 3
b = 4
c = 0

def f(b):
    a = 4
    b = 1
    c = 3
    return a + b + c

a = 6
b = 1
c = 5

try:
    print(f(b), a, b, c, sep=';')
except:
    print('ERR', a, b, c, sep=';')
code

3. Що буде виведено за виконання?
code
try:
    n = 13
    while n>4:
        n += -2
    print(n, end=';')
except:
    print('ERR')
code

4. Що буде виведено за виконання?
code
try:
    for c in range(7, -6, -1):
        if c<8:
            continue
        print(c, end=';');c = -4
    else:
        print(c,end=';')
    print(c)
except:
    print('ERR')
code

5. Що буде виведено за виконання?\par (Дійсні значення, якщо наявні, в цьому завданні будуть виводитись у форматі з фіксованою крапкою та, можливо, з доволі великою кількістю десяткових розрядів. У відповіді для дійсних значень у ЦЬОМУ завданні достатньо вказати один десятковий розряд після крапки, відкинувши решту розрядів без округлення. Наприклад, замість 13.66666666666667 достатньо вказати 13.6, замість 25.23 достатньо вказати 25.2)
code
try:
    print(5, end=';')
    print(0j==9j, end=';')
    print(8, end=';')
except ZeroDivisionError:
    print(9, end=';')
except ValueError:
    print(7, end=';')
except:
    print('ERR', end=';')
else:
    print(2, end=';')
finally:
    print(5)
code

6. Що буде виведено за виконання?
code
def f(n):
    if n>9:
        f(n//10)
    else:
        print(n%10, end='')

f(543216)
code

7. Що буде виведено за виконання?
code
class A:
    a = 1

    def _f1(self): print(2, end=';')
    def _f2(self): print(3, end=';')

    def main(self):
        try: print(self.a, end=';')
        except: print('ERR', end=';')

        try: self._f1()
        except: print('ERR', end=';')

        try: self._f2()
        except: print('ERR', end=';')


class B(A):
    a = 4

    def _f1(self): print(5, end=';')
    def _f2(self): print(6, end=';')

try: B().main()
except: print('ERR', end=';')
code

8. Що буде виведено за виконання?
code
def f(arg, n):
    n = 3
    tmp = arg
    del tmp[6:7]
    return len(tmp)>3

try:
    b = ['0','1','2','3','4']
    n = 9
    f(b, n)
    print(n, end=';')
    for el in b:
        print(el, end=';')
except:
    print('ERR')
code

9. Що буде виведено за виконання?
code
class A:
    c = 1
    b = 2

    def __init__(self):
        c = 3

class B(A):
    c = 4

    def __init__(self):
        super().__init__()
        self.a = 7

obj = B()
obj.a = obj.b + 10
B.a = 13

try: print(obj.a, end= ';')
except: print('ERR', end=';')

try: print(A.a, end= ';')
except: print('ERR', end=';')

try: print(B.a, end= ';')
except: print('ERR', end=';')
code

10. 
Для графа на рисунку з вершини 6 запущено алгоритм пошуку в глибину, що будує кістякове дерево. Списки суміжності вершин переглядаються алгоритмом за зростанням міток вершин.\par
\begin{center}
	\adjustimage{max size={0.9\linewidth}{0.9\paperheight}}
	{../10/Traversals/99.png}
\end{center}
\par Які з наведених ребер входять до складу отриманого кістякового дерева?\par
1) (1, 4)\par
2) (4, 6)\par
3) (3, 4)\par
4) (2, 4)\par
5) (4, 5)\par
6) (2, 6)\par
{\vskip 15pt} У формі вибрати номери відповідних ребер.\par
%answer:3 5
\end {large}}
%end
{\smallskip \begin {small} Затверджено на засіданні кафедри теоретичної кібернетики, протокол \No 12 від   19.05.2020 р.\par\smallskip\smallskip
{\parbox {6.5 cm}{Зав. кафедри \dotfill Ю.В.~Крак}}\hspace {0.7cm}{\hbox to 6.5 cm{ Екзаменатор \dotfill Т.О.~Карнаух}} \end{small}}
%begin
{{\begin {small}\par
{ \center  Київський національний університет імені Тараса Шевченка \par}
{\parbox {5 cm}{Освітня програма \hfill}{\parbox {7 cm}{Прикладна математика, Системний аналіз \hfill}}\parbox {6 cm}{\hfill Семестр 2}}\par
{\parbox {5 cm} {Навчальний предмет \hfill}\parbox {7 cm}{Програмування \hfill}\parbox {6cm}{\hfill}\par}\vspace{-7pt} \end{small}}{\center Екзаменаційний білет \No \textbf 
{ 124 }\par}}
{
\begin {large}
1. Що буде виведено за виконання?
code
try:
    res = not 9.0<=6 and 9!=5 and True>=8.0
    if isinstance(res, bool):
        print(f'bool;{res}')
    elif isinstance(res, int):
        print(f'int;{res}')
    elif isinstance(res, float):
        print(f'float;{res:.2f}')
    else:
        print('???')
except:
    print('ERR')
code

2. Що буде виведено за виконання?
code
a = 2
b = 3
c = 7

def h(a):
    global c
    a *= 2
    b = 1
    c = 5
    return a + b + c

a = 0
b = 4
c = 1

try:
    print(h(a), a, b, c, sep=';')
except:
    print('ERR', a, b, c, sep=';')
code

3. Що буде виведено за виконання?
code
try:
    m = 12
    while m>=3:
        m += -3
    print(m, end=';')
except:
    print('ERR')
code

4. Що буде виведено за виконання?
code
try:
    for d in range(-15, 8, -2):
        if d>=7:
            continue
        print(d, end=';')
    else:
        print(d,end=';')
    print(d)
except:
    print('ERR')
code

5. Що буде виведено за виконання?\par (Дійсні значення, якщо наявні, в цьому завданні будуть виводитись у форматі з фіксованою крапкою та, можливо, з доволі великою кількістю десяткових розрядів. У відповіді для дійсних значень у ЦЬОМУ завданні достатньо вказати один десятковий розряд після крапки, відкинувши решту розрядів без округлення. Наприклад, замість 13.66666666666667 достатньо вказати 13.6, замість 25.23 достатньо вказати 25.2)
code
try:
    print(6, end=';')
    print(3//0.0, end=';')
    print(2, end=';')
except TypeError:
    print(1, end=';')
except Exception:
    print(8, end=';')
except:
    print('ERR', end=';')
else:
    print(9, end=';')
finally:
    print(7)
code

6. Що буде виведено за виконання?
code
def f(n):
    if n>9:
        f(n//100)
    print(n%10,end='')
 
f(987654)
code

7. Що буде виведено за виконання?
code
class A:
    c = 1

    def _f1(self): print(2, end=';')
    def f2(self): print(3, end=';')

    def main(self):
        try: print(self.c, end=';')
        except: print('ERR', end=';')

        try: self._f1()
        except: print('ERR', end=';')

        try: self.f2()
        except: print('ERR', end=';')


class B(A):
    c = 4

    def _f1(self): print(5, end=';')
    def f2(self): print(6, end=';')

try: B().main()
except: print('ERR', end=';')
code

8. Що буде виведено за виконання?
code
def f(arg, n):
    arg[83] = 1
    n += 3

try:
    n = 3
    s = {80:7, 83:1, 20:2, 80:4}
    f(s, n)
    print(n, end=';')
    for x in s.values():
        print(x, sep=';', end=';')
except:
    print('ERR')
code

9. Що буде виведено за виконання?
code
class A:
    b = 1
    c = 2

    def __init__(self):
        c = 3

class B(A):
    c = 4

    def __init__(self):
        super().__init__()
        self.a = 7

obj = B()
obj.c = obj.b + 10
B.c = 13

try: print(obj.c, end= ';')
except: print('ERR', end=';')

try: print(A.c, end= ';')
except: print('ERR', end=';')

try: print(B.c, end= ';')
except: print('ERR', end=';')
code

10. 
Для графа на рисунку з вершини 0 запущено алгоритм пошуку в глибину, що будує кістякове дерево. Списки суміжності вершин переглядаються алгоритмом за зростанням міток вершин.\par
\begin{center}
	\adjustimage{max size={0.9\linewidth}{0.9\paperheight}}
	{../10/Traversals/276.png}
\end{center}
\par Які з наведених ребер входять до складу отриманого кістякового дерева?\par
1) (1, 2)\par
2) (3, 5)\par
3) (0, 1)\par
4) (2, 4)\par
5) (3, 4)\par
6) (1, 4)\par
{\vskip 15pt} У формі вибрати номери відповідних ребер.\par
%answer:1 2 3 4 5
\end {large}}
%end
{\smallskip \begin {small} Затверджено на засіданні кафедри теоретичної кібернетики, протокол \No 12 від   19.05.2020 р.\par\smallskip\smallskip
{\parbox {6.5 cm}{Зав. кафедри \dotfill Ю.В.~Крак}}\hspace {0.7cm}{\hbox to 6.5 cm{ Екзаменатор \dotfill Т.О.~Карнаух}} \end{small}}
%begin
{{\begin {small}\par
{ \center  Київський національний університет імені Тараса Шевченка \par}
{\parbox {5 cm}{Освітня програма \hfill}{\parbox {7 cm}{Прикладна математика, Системний аналіз \hfill}}\parbox {6 cm}{\hfill Семестр 2}}\par
{\parbox {5 cm} {Навчальний предмет \hfill}\parbox {7 cm}{Програмування \hfill}\parbox {6cm}{\hfill}\par}\vspace{-7pt} \end{small}}{\center Екзаменаційний білет \No \textbf 
{ 125 }\par}}
{
\begin {large}
1. Що буде виведено за виконання?
code
try:
    res = not 2.0==7 or 4.0>False and 3<8
    if isinstance(res, bool):
        print(f'bool;{res}')
    elif isinstance(res, int):
        print(f'int;{res}')
    elif isinstance(res, float):
        print(f'float;{res:.2f}')
    else:
        print('???')
except:
    print('ERR')
code

2. Що буде виведено за виконання?
code
a = 8
b = 6
c = 5

def f(b):
    a = 3
    b -= 4
    c = 4
    return a + b + c

a = 9
b = 0
c = 5

try:
    print(f(a), a, b, c, sep=';')
except:
    print('ERR', a, b, c, sep=';')
code

3. Що буде виведено за виконання?
code
try:
    i = 11
    while i>2:
        i += -2
    print(i, end=';')
except:
    print('ERR')
code

4. Що буде виведено за виконання?
code
try:
    for a in range(-6, -12, 3):
        if a>=6:
            break
        print(a, end=';')
    else:
        print(a,end=';')
    print(a)
except:
    print('ERR')
code

5. Що буде виведено за виконання?\par (Дійсні значення, якщо наявні, в цьому завданні будуть виводитись у форматі з фіксованою крапкою та, можливо, з доволі великою кількістю десяткових розрядів. У відповіді для дійсних значень у ЦЬОМУ завданні достатньо вказати один десятковий розряд після крапки, відкинувши решту розрядів без округлення. Наприклад, замість 13.66666666666667 достатньо вказати 13.6, замість 25.23 достатньо вказати 25.2)
code
try:
    print(0, end=';')
    print(int('4'), end=';')
    print(1, end=';')
except ZeroDivisionError:
    print(8, end=';')
except KeyboardInterrupt:
    print(2, end=';')
except:
    print('ERR', end=';')
else:
    print(4, end=';')
finally:
    print(6)
code

6. Що буде виведено за виконання?
code
def f(n):
    if n>9:
        f(n//10)
    print(n%10, end='')
 
f(543216)
code

7. Що буде виведено за виконання?
code
class A:
    a = 1

    def g1(self): print(2, end=';')
    def _g2(self): print(3, end=';')

    def main(self):
        try: print(self.a, end=';')
        except: print('ERR', end=';')

        try: self.g1()
        except: print('ERR', end=';')

        try: self._g2()
        except: print('ERR', end=';')


class B(A):
    a = 4

    def g1(self): print(5, end=';')
    def _g2(self): print(6, end=';')

try: B().main()
except: print('ERR', end=';')
code

8. Що буде виведено за виконання?
code
def f(arg, n):
    n = 0
    tmp = arg
    tmp[-4:-9] = 'bd'
    return len(tmp)>3

try:
    a = ['0','1','2','3','4']
    n = 5
    f(a, n)
    print(n, end=';')
    for el in a:
        print(el, end=';')
except:
    print('ERR')
code

9. Що буде виведено за виконання?
code
class A:
    a = 1
    b = 2

    def __init__(self):
        a = 3

class B(A):
    a = 4

    def __init__(self):
        super().__init__()
        self.c = 7

obj = B()
obj.a = obj.a + 10
B.a = 13

try: print(obj.a, end= ';')
except: print('ERR', end=';')

try: print(A.a, end= ';')
except: print('ERR', end=';')

try: print(B.a, end= ';')
except: print('ERR', end=';')
code

10. 
Для графа на рисунку з вершини 1 запущено алгоритм пошуку в глибину, що будує кістякове дерево. Списки суміжності вершин переглядаються алгоритмом за зростанням міток вершин.\par
\begin{center}
	\adjustimage{max size={0.9\linewidth}{0.9\paperheight}}
	{../10/Traversals/212.png}
\end{center}
\par Які з наведених ребер входять до складу отриманого кістякового дерева?\par
1) (2, 3)\par
2) (0, 1)\par
3) (4, 6)\par
4) (3, 4)\par
5) (3, 6)\par
6) (4, 5)\par
{\vskip 15pt} У формі вибрати номери відповідних ребер.\par
%answer:1 2 4 6
\end {large}}
%end
{\smallskip \begin {small} Затверджено на засіданні кафедри теоретичної кібернетики, протокол \No 12 від   19.05.2020 р.\par\smallskip\smallskip
{\parbox {6.5 cm}{Зав. кафедри \dotfill Ю.В.~Крак}}\hspace {0.7cm}{\hbox to 6.5 cm{ Екзаменатор \dotfill Т.О.~Карнаух}} \end{small}}
%begin
{{\begin {small}\par
{ \center  Київський національний університет імені Тараса Шевченка \par}
{\parbox {5 cm}{Освітня програма \hfill}{\parbox {7 cm}{Прикладна математика, Системний аналіз \hfill}}\parbox {6 cm}{\hfill Семестр 2}}\par
{\parbox {5 cm} {Навчальний предмет \hfill}\parbox {7 cm}{Програмування \hfill}\parbox {6cm}{\hfill}\par}\vspace{-7pt} \end{small}}{\center Екзаменаційний білет \No \textbf 
{ 126 }\par}}
{
\begin {large}
1. Що буде виведено за виконання?
code
try:
    res = not 7.0>=3 and 7<=6.0 or True!=4
    if isinstance(res, bool):
        print(f'bool;{res}')
    elif isinstance(res, int):
        print(f'int;{res}')
    elif isinstance(res, float):
        print(f'float;{res:.2f}')
    else:
        print('???')
except:
    print('ERR')
code

2. Що буде виведено за виконання?
code
a = 6
b = 7
c = 1

def g(b):
    global c
    a = 1
    b += 2
    c = 3
    return a + b + c

a = 9
b = 3
c = 8

try:
    print(g(b), a, b, c, sep=';')
except:
    print('ERR', a, b, c, sep=';')
code

3. Що буде виведено за виконання?
code
try:
    m = 1
    while m<=8:
        m += 2
    print(m, end=';')
except:
    print('ERR')
code

4. Що буде виведено за виконання?
code
try:
    for b in range(-4, 15, 3):
        if b>5:
            continue
        print(b, end=';');b = -1
        if b<3:
            break
    else:
        print('end',end=';')
    print(b)
except:
    print('ERR')
code

5. Що буде виведено за виконання?\par (Дійсні значення, якщо наявні, в цьому завданні будуть виводитись у форматі з фіксованою крапкою та, можливо, з доволі великою кількістю десяткових розрядів. У відповіді для дійсних значень у ЦЬОМУ завданні достатньо вказати один десятковий розряд після крапки, відкинувши решту розрядів без округлення. Наприклад, замість 13.66666666666667 достатньо вказати 13.6, замість 25.23 достатньо вказати 25.2)
code
try:
    print(9, end=';')
    print(int('b0'), end=';')
    print(7, end=';')
except KeyboardInterrupt:
    print(5, end=';')
except TypeError:
    print(3, end=';')
except:
    print('ERR', end=';')
else:
    print(6, end=';')
finally:
    print(0)
code

6. Що буде виведено за виконання?
code
def f(n):
    print(n%2, end='')
    if n>2:
        f(n//2)
 
f(16)
code

7. Що буде виведено за виконання?
code
class A:
    d = 1

    def _h1(self): print(2, end=';')
    def __h2(self): print(3, end=';')

    def main(self):
        try: print(self.d, end=';')
        except: print('ERR', end=';')

        try: self._h1()
        except: print('ERR', end=';')

        try: self.__h2()
        except: print('ERR', end=';')


class B(A):
    d = 4

    def _h1(self): print(5, end=';')
    def __h2(self): print(6, end=';')

try: B().main()
except: print('ERR', end=';')
code

8. Що буде виведено за виконання?
code
def f(arg, n):
    arg[88] = 7
    n = 6
    arg = set(arg.values())

try:
    n = 0
    s = {88:7, 56:6, 65:0}
    f(s, n)
    print(n, end=';')
    for x in s.keys():
        print(x, sep=';', end=';')
except:
    print('ERR')
code

9. Що буде виведено за виконання?
code
class A:
    a = 1
    c = 2

    def __init__(self):
        c = 3

class B(A):
    c = 4

    def __init__(self):
        super().__init__()
        self.b = 7

obj = B()
obj.b = obj.c + 10
B.b = 13

try: print(obj.b, end= ';')
except: print('ERR', end=';')

try: print(A.b, end= ';')
except: print('ERR', end=';')

try: print(B.b, end= ';')
except: print('ERR', end=';')
code

10. 
Для графа на рисунку з вершини 2 запущено алгоритм пошуку в глибину, що будує кістякове дерево. Списки суміжності вершин переглядаються алгоритмом за зростанням міток вершин.\par
\begin{center}
	\adjustimage{max size={0.9\linewidth}{0.9\paperheight}}
	{../10/Traversals/216.png}
\end{center}
\par Які з наведених ребер входять до складу отриманого кістякового дерева?\par
1) (2, 5)\par
2) (0, 6)\par
3) (0, 1)\par
4) (1, 2)\par
5) (1, 4)\par
6) (4, 5)\par
{\vskip 15pt} У формі вибрати номери відповідних ребер.\par
%answer:2 3 4 5 6
\end {large}}
%end
{\smallskip \begin {small} Затверджено на засіданні кафедри теоретичної кібернетики, протокол \No 12 від   19.05.2020 р.\par\smallskip\smallskip
{\parbox {6.5 cm}{Зав. кафедри \dotfill Ю.В.~Крак}}\hspace {0.7cm}{\hbox to 6.5 cm{ Екзаменатор \dotfill Т.О.~Карнаух}} \end{small}}
%begin
{{\begin {small}\par
{ \center  Київський національний університет імені Тараса Шевченка \par}
{\parbox {5 cm}{Освітня програма \hfill}{\parbox {7 cm}{Прикладна математика, Системний аналіз \hfill}}\parbox {6 cm}{\hfill Семестр 2}}\par
{\parbox {5 cm} {Навчальний предмет \hfill}\parbox {7 cm}{Програмування \hfill}\parbox {6cm}{\hfill}\par}\vspace{-7pt} \end{small}}{\center Екзаменаційний білет \No \textbf 
{ 127 }\par}}
{
\begin {large}
1. Що буде виведено за виконання?
code
try:
    res = 2<6.0 or not 6>=False or 9.0>5
    if isinstance(res, bool):
        print(f'bool;{res}')
    elif isinstance(res, int):
        print(f'int;{res}')
    elif isinstance(res, float):
        print(f'float;{res:.2f}')
    else:
        print('???')
except:
    print('ERR')
code

2. Що буде виведено за виконання?
code
a = 2
b = 4
c = 0

def g(a):
    a += 5
    b = 5
    c = 3
    return a + b + c

a = 3
b = 5
c = 4

try:
    print(g(b), a, b, c, sep=';')
except:
    print('ERR', a, b, c, sep=';')
code

3. Що буде виведено за виконання?
code
try:
    k = 5
    while k<8:
        k += 1
    print(k, end=';')
except:
    print('ERR')
code

4. Що буде виведено за виконання?
code
try:
    for f in range(-13, -14, -1):
        if f<=6:
            break
        print(f, end=';')
    else:
        print(f,end=';')
    print(f)
except:
    print('ERR')
code

5. Що буде виведено за виконання?\par (Дійсні значення, якщо наявні, в цьому завданні будуть виводитись у форматі з фіксованою крапкою та, можливо, з доволі великою кількістю десяткових розрядів. У відповіді для дійсних значень у ЦЬОМУ завданні достатньо вказати один десятковий розряд після крапки, відкинувши решту розрядів без округлення. Наприклад, замість 13.66666666666667 достатньо вказати 13.6, замість 25.23 достатньо вказати 25.2)
code
try:
    print(4, end=';')
    print(7%0, end=';')
    print(8, end=';')
except ZeroDivisionError:
    print(1, end=';')
except ValueError:
    print(3, end=';')
except:
    print('ERR', end=';')
else:
    print(5, end=';')
finally:
    print(2)
code

6. Що буде виведено за виконання?
code
def f(s):
    s1 = s
    for i in range(len(s)):
        s1[i] += 1
        return s1

s = [1, 2, 3]
f(s)
print(s)
code

7. Що буде виведено за виконання?
code
class A:
    c = 1

    def h1(self): print(2, end=';')
    def _h2(self): print(3, end=';')

    def main(self):
        try: print(self.c, end=';')
        except: print('ERR', end=';')

        try: self.h1()
        except: print('ERR', end=';')

        try: self._h2()
        except: print('ERR', end=';')


class B(A):
    c = 4

    def h1(self): print(5, end=';')
    def _h2(self): print(6, end=';')

try: B().main()
except: print('ERR', end=';')
code

8. Що буде виведено за виконання?
code
def f(arg, n):
    n = 7
    tmp = arg
    del tmp[2:5]
    return len(tmp)>3

try:
    f = [10,11,12,13,14]
    n = 2
    f(f, n)
    print(n, end=';')
    for el in f:
        print(el, end=';')
except:
    print('ERR')
code

9. Що буде виведено за виконання?
code
class A:
    b = 1
    c = 2

    def __init__(self):
        b = 3

class B(A):
    b = 4

    def __init__(self):
        super().__init__()
        self.a = 7

obj = B()
obj.a = obj.b + 10
B.a = 13

try: print(obj.a, end= ';')
except: print('ERR', end=';')

try: print(A.a, end= ';')
except: print('ERR', end=';')

try: print(B.a, end= ';')
except: print('ERR', end=';')
code

10. 
Для графа на рисунку з вершини 5 запущено алгоритм пошуку в глибину, що будує кістякове дерево. Списки суміжності вершин переглядаються алгоритмом за зростанням міток вершин.\par
\begin{center}
	\adjustimage{max size={0.9\linewidth}{0.9\paperheight}}
	{../10/Traversals/78.png}
\end{center}
\par Які з наведених ребер входять до складу отриманого кістякового дерева?\par
1) (1, 5)\par
2) (1, 6)\par
3) (1, 2)\par
4) (0, 1)\par
5) (2, 5)\par
6) (5, 6)\par
{\vskip 15pt} У формі вибрати номери відповідних ребер.\par
%answer:1 4
\end {large}}
%end
{\smallskip \begin {small} Затверджено на засіданні кафедри теоретичної кібернетики, протокол \No 12 від   19.05.2020 р.\par\smallskip\smallskip
{\parbox {6.5 cm}{Зав. кафедри \dotfill Ю.В.~Крак}}\hspace {0.7cm}{\hbox to 6.5 cm{ Екзаменатор \dotfill Т.О.~Карнаух}} \end{small}}
%begin
{{\begin {small}\par
{ \center  Київський національний університет імені Тараса Шевченка \par}
{\parbox {5 cm}{Освітня програма \hfill}{\parbox {7 cm}{Прикладна математика, Системний аналіз \hfill}}\parbox {6 cm}{\hfill Семестр 2}}\par
{\parbox {5 cm} {Навчальний предмет \hfill}\parbox {7 cm}{Програмування \hfill}\parbox {6cm}{\hfill}\par}\vspace{-7pt} \end{small}}{\center Екзаменаційний білет \No \textbf 
{ 128 }\par}}
{
\begin {large}
1. Що буде виведено за виконання?
code
try:
    res = 4**-1.0*1
    if isinstance(res, bool):
        print(f'bool;{res}')
    elif isinstance(res, int):
        print(f'int;{res}')
    elif isinstance(res, float):
        print(f'float;{res:.2f}')
    else:
        print('???')
except:
    print('ERR')
code

2. Що буде виведено за виконання?
code
a = 2
b = 7
c = 5

def g(a):
    global c
    a = 5
    b = 5
    c = 2
    return a + b + c

a = 1
b = 6
c = 7

try:
    print(g(a), a, b, c, sep=';')
except:
    print('ERR', a, b, c, sep=';')
code

3. Що буде виведено за виконання?
code
try:
    i = 2
    while i<=13:
        i += 2
    print(i, end=';')
except:
    print('ERR')
code

4. Що буде виведено за виконання?
code
try:
    for e in range(12, -5, 2):
        if e<8:
            continue
        print(e, end=';')
    else:
        print(e,end=';')
    print(e)
except:
    print('ERR')
code

5. Що буде виведено за виконання?\par (Дійсні значення, якщо наявні, в цьому завданні будуть виводитись у форматі з фіксованою крапкою та, можливо, з доволі великою кількістю десяткових розрядів. У відповіді для дійсних значень у ЦЬОМУ завданні достатньо вказати один десятковий розряд після крапки, відкинувши решту розрядів без округлення. Наприклад, замість 13.66666666666667 достатньо вказати 13.6, замість 25.23 достатньо вказати 25.2)
code
try:
    print(9, end=';')
    print(6/3, end=';')
    print(5, end=';')
except Exception:
    print(3, end=';')
except TypeError:
    print(2, end=';')
except:
    print('ERR', end=';')
else:
    print(7, end=';')
finally:
    print(4)
code

6. Що буде виведено за виконання?
code
def f(n):
    if n>2:
        f(n//2)
    print(n%2, end='')
 
f(16)
code

7. Що буде виведено за виконання?
code
class A:
    b = 1

    def _f1(self): print(2, end=';')
    def _f2(self): print(3, end=';')

    def main(self):
        try: print(self.b, end=';')
        except: print('ERR', end=';')

        try: self._f1()
        except: print('ERR', end=';')

        try: self._f2()
        except: print('ERR', end=';')


class B(A):
    b = 4

    def _f1(self): print(5, end=';')
    def _f2(self): print(6, end=';')

try: B().main()
except: print('ERR', end=';')
code

8. Що буде виведено за виконання?
code
def f(arg, n):
    arg[40] = 3
    n *= -3

try:
    n = 4
    d = {40:4, 36:2, 45:3, 71:0}
    f(d, n)
    print(n, end=';')
    for x in d.values():
        print(x, sep=';', end=';')
except:
    print('ERR')
code

9. Що буде виведено за виконання?
code
class A:
    c = 1
    b = 2

    def __init__(self):
        b = 3

class B(A):
    b = 4

    def __init__(self):
        super().__init__()
        self.a = 7

obj = B()
obj.c = obj.b + 10
B.c = 13

try: print(obj.c, end= ';')
except: print('ERR', end=';')

try: print(A.c, end= ';')
except: print('ERR', end=';')

try: print(B.c, end= ';')
except: print('ERR', end=';')
code

10. 
Для графа на рисунку з вершини 5 запущено алгоритм пошуку в глибину, що будує кістякове дерево. Списки суміжності вершин переглядаються алгоритмом за зростанням міток вершин.\par
\begin{center}
	\adjustimage{max size={0.9\linewidth}{0.9\paperheight}}
	{../10/Traversals/180.png}
\end{center}
\par Які з наведених ребер входять до складу отриманого кістякового дерева?\par
1) (4, 6)\par
2) (2, 4)\par
3) (1, 5)\par
4) (0, 2)\par
5) (3, 5)\par
6) (2, 6)\par
{\vskip 15pt} У формі вибрати номери відповідних ребер.\par
%answer:1 2 3 4 5
\end {large}}
%end
{\smallskip \begin {small} Затверджено на засіданні кафедри теоретичної кібернетики, протокол \No 12 від   19.05.2020 р.\par\smallskip\smallskip
{\parbox {6.5 cm}{Зав. кафедри \dotfill Ю.В.~Крак}}\hspace {0.7cm}{\hbox to 6.5 cm{ Екзаменатор \dotfill Т.О.~Карнаух}} \end{small}}
%begin
{{\begin {small}\par
{ \center  Київський національний університет імені Тараса Шевченка \par}
{\parbox {5 cm}{Освітня програма \hfill}{\parbox {7 cm}{Прикладна математика, Системний аналіз \hfill}}\parbox {6 cm}{\hfill Семестр 2}}\par
{\parbox {5 cm} {Навчальний предмет \hfill}\parbox {7 cm}{Програмування \hfill}\parbox {6cm}{\hfill}\par}\vspace{-7pt} \end{small}}{\center Екзаменаційний білет \No \textbf 
{ 129 }\par}}
{
\begin {large}
1. Що буде виведено за виконання?
code
def f(n):
    print('f', end='')
    return n==4

try:
    print(f(4) and f(-1))
except:
    print('ERR')
code

2. Що буде виведено за виконання?
code
a = 9
b = 2
c = 0

def h(a):
    global c
    a = 3
    b = 1
    c = 4
    return a + b + c

a = 4
b = 3
c = 8

try:
    print(h(b), a, b, c, sep=';')
except:
    print('ERR', a, b, c, sep=';')
code

3. Що буде виведено за виконання?
code
try:
    n = 7
    while n<6:
        n += 1
    print(n, end=';')
except:
    print('ERR')
code

4. Що буде виведено за виконання?
code
try:
    for c in range(-8, -4, -3):
        if c>3:
            continue
        print(c, end=';');c = -7
        if c<=6:
            break
    else:
        print(c,end=';')
    print(c)
except:
    print('ERR')
code

5. Що буде виведено за виконання?\par (Дійсні значення, якщо наявні, в цьому завданні будуть виводитись у форматі з фіксованою крапкою та, можливо, з доволі великою кількістю десяткових розрядів. У відповіді для дійсних значень у ЦЬОМУ завданні достатньо вказати один десятковий розряд після крапки, відкинувши решту розрядів без округлення. Наприклад, замість 13.66666666666667 достатньо вказати 13.6, замість 25.23 достатньо вказати 25.2)
code
try:
    print(1, end=';')
    print(8j!=2j, end=';')
    print(0, end=';')
except KeyboardInterrupt:
    print(9, end=';')
except ZeroDivisionError:
    print(7, end=';')
except:
    print('ERR', end=';')
else:
    print(3, end=';')
finally:
    print(8)
code

6. Що буде виведено за виконання?
code
def f(n):
    print(n%10, end='')
    if n>9:
        f(n//10)
 
f(543216)
code

7. Що буде виведено за виконання?
code
class A:
    a = 1

    def __g1(self): print(2, end=';')
    def _g2(self): print(3, end=';')

    def main(self):
        try: print(self.a, end=';')
        except: print('ERR', end=';')

        try: self.__g1()
        except: print('ERR', end=';')

        try: self._g2()
        except: print('ERR', end=';')


class B(A):
    a = 4

    def __g1(self): print(5, end=';')
    def _g2(self): print(6, end=';')

try: B().main()
except: print('ERR', end=';')
code

8. Що буде виведено за виконання?
code
def f(arg, n):
    arg[51] = 3
    n -= 3

try:
    n = 3
    s = {88:0, 84:0, 46:7, 84:0}
    f(s, n)
    print(n, end=';')
    for x in s.items():
        print(*x, sep=';', end=';')
except:
    print('ERR')
code

9. Що буде виведено за виконання?
code
class A:
    b = 1
    a = 2

    def __init__(self):
        b = 3

class B(A):
    a = 4

    def __init__(self):
        super().__init__()
        self.c = 7

obj = B()
obj.c = obj.a + 10
B.c = 13

try: print(obj.c, end= ';')
except: print('ERR', end=';')

try: print(A.c, end= ';')
except: print('ERR', end=';')

try: print(B.c, end= ';')
except: print('ERR', end=';')
code

10. 
Для графа на рисунку з вершини 5 запущено алгоритм пошуку в глибину, що будує кістякове дерево. Списки суміжності вершин переглядаються алгоритмом за зростанням міток вершин.\par
\begin{center}
	\adjustimage{max size={0.9\linewidth}{0.9\paperheight}}
	{../10/Traversals/42.png}
\end{center}
\par Які з наведених ребер входять до складу отриманого кістякового дерева?\par
1) (0, 4)\par
2) (3, 4)\par
3) (3, 6)\par
4) (5, 6)\par
5) (0, 3)\par
6) (0, 1)\par
{\vskip 15pt} У формі вибрати номери відповідних ребер.\par
%answer:1 5 6
\end {large}}
%end
{\smallskip \begin {small} Затверджено на засіданні кафедри теоретичної кібернетики, протокол \No 12 від   19.05.2020 р.\par\smallskip\smallskip
{\parbox {6.5 cm}{Зав. кафедри \dotfill Ю.В.~Крак}}\hspace {0.7cm}{\hbox to 6.5 cm{ Екзаменатор \dotfill Т.О.~Карнаух}} \end{small}}
%begin
{{\begin {small}\par
{ \center  Київський національний університет імені Тараса Шевченка \par}
{\parbox {5 cm}{Освітня програма \hfill}{\parbox {7 cm}{Прикладна математика, Системний аналіз \hfill}}\parbox {6 cm}{\hfill Семестр 2}}\par
{\parbox {5 cm} {Навчальний предмет \hfill}\parbox {7 cm}{Програмування \hfill}\parbox {6cm}{\hfill}\par}\vspace{-7pt} \end{small}}{\center Екзаменаційний білет \No \textbf 
{ 130 }\par}}
{
\begin {large}
1. Що буде виведено за виконання?
code
try:
    res = 5//3//9*3
    if isinstance(res, bool):
        print(f'bool;{res}')
    elif isinstance(res, int):
        print(f'int;{res}')
    elif isinstance(res, float):
        print(f'float;{res:.2f}')
    else:
        print('???')
except:
    print('ERR')
code

2. Що буде виведено за виконання?
code
a = 1
b = 9
c = 6

def h(a):
    a = 4
    b -= 2
    c = 1
    return a + b + c

a = 2
b = 8
c = 7

try:
    print(h(a), a, b, c, sep=';')
except:
    print('ERR', a, b, c, sep=';')
code

3. Що буде виведено за виконання?
code
try:
    m = 6
    while m<=5:
        m += 2
    print(m, end=';')
except:
    print('ERR')
code

4. Що буде виведено за виконання?
code
try:
    for a in range(4, 5, -2):
        if a>=7:
            break
        print(a, end=';')
    else:
        print(a,end=';')
    print(a)
except:
    print('ERR')
code

5. Що буде виведено за виконання?\par (Дійсні значення, якщо наявні, в цьому завданні будуть виводитись у форматі з фіксованою крапкою та, можливо, з доволі великою кількістю десяткових розрядів. У відповіді для дійсних значень у ЦЬОМУ завданні достатньо вказати один десятковий розряд після крапки, відкинувши решту розрядів без округлення. Наприклад, замість 13.66666666666667 достатньо вказати 13.6, замість 25.23 достатньо вказати 25.2)
code
try:
    print(5, end=';')
    print(int('6'), end=';')
    print(9, end=';')
except Exception:
    print(1, end=';')
except ValueError:
    print(0, end=';')
except:
    print('ERR', end=';')
else:
    print(2, end=';')
finally:
    print(4)
code

6. Що буде виведено за виконання?
code
def f(s1):
    s = s1[:]
    for i in range(len(s)):
        s[i] += 1
 
s = [1, 2, 3]
f(s)
print(s)
code

7. Що буде виведено за виконання?
code
class A:
    c = 1

    def _g1(self): print(2, end=';')
    def g2(self): print(3, end=';')

    def main(self):
        try: print(self.c, end=';')
        except: print('ERR', end=';')

        try: self._g1()
        except: print('ERR', end=';')

        try: self.g2()
        except: print('ERR', end=';')


class B(A):
    c = 4

    def _g1(self): print(5, end=';')
    def g2(self): print(6, end=';')

try: B().main()
except: print('ERR', end=';')
code

8. Що буде виведено за виконання?
code
def f(arg, n):
    arg[18] = 7
    n = 7
    arg = list(arg.items())

try:
    n = 1
    t = {36:8, 48:1, 29:0}
    f(t, n)
    print(n, end=';')
    for x in t.keys():
        print(x, sep=';', end=';')
except:
    print('ERR')
code

9. Що буде виведено за виконання?
code
class A:
    c = 1
    a = 2

    def __init__(self):
        a = 3

class B(A):
    c = 4

    def __init__(self):
        super().__init__()
        self.b = 7

obj = B()
obj.a = obj.c + 10
B.a = 13

try: print(obj.a, end= ';')
except: print('ERR', end=';')

try: print(A.a, end= ';')
except: print('ERR', end=';')

try: print(B.a, end= ';')
except: print('ERR', end=';')
code

10. 
Для графа на рисунку з вершини 6 запущено алгоритм пошуку в глибину, що будує кістякове дерево. Списки суміжності вершин переглядаються алгоритмом за зростанням міток вершин.\par
\begin{center}
	\adjustimage{max size={0.9\linewidth}{0.9\paperheight}}
	{../10/Traversals/221.png}
\end{center}
\par Які з наведених ребер входять до складу отриманого кістякового дерева?\par
1) (4, 5)\par
2) (0, 5)\par
3) (1, 5)\par
4) (2, 4)\par
5) (2, 3)\par
6) (0, 1)\par
{\vskip 15pt} У формі вибрати номери відповідних ребер.\par
%answer:3 4 5 6
\end {large}}
%end
{\smallskip \begin {small} Затверджено на засіданні кафедри теоретичної кібернетики, протокол \No 12 від   19.05.2020 р.\par\smallskip\smallskip
{\parbox {6.5 cm}{Зав. кафедри \dotfill Ю.В.~Крак}}\hspace {0.7cm}{\hbox to 6.5 cm{ Екзаменатор \dotfill Т.О.~Карнаух}} \end{small}}
%begin
{{\begin {small}\par
{ \center  Київський національний університет імені Тараса Шевченка \par}
{\parbox {5 cm}{Освітня програма \hfill}{\parbox {7 cm}{Прикладна математика, Системний аналіз \hfill}}\parbox {6 cm}{\hfill Семестр 2}}\par
{\parbox {5 cm} {Навчальний предмет \hfill}\parbox {7 cm}{Програмування \hfill}\parbox {6cm}{\hfill}\par}\vspace{-7pt} \end{small}}{\center Екзаменаційний білет \No \textbf 
{ 131 }\par}}
{
\begin {large}
1. Що буде виведено за виконання?
code
try:
    res = 6//8/7*8
    if isinstance(res, bool):
        print(f'bool;{res}')
    elif isinstance(res, int):
        print(f'int;{res}')
    elif isinstance(res, float):
        print(f'float;{res:.2f}')
    else:
        print('???')
except:
    print('ERR')
code

2. Що буде виведено за виконання?
code
a = 8
b = 9
c = 0

def f(b):
    global c
    a *= 4
    b = 3
    c = 2
    return a + b + c

a = 4
b = 7
c = 1

try:
    print(f(b), a, b, c, sep=';')
except:
    print('ERR', a, b, c, sep=';')
code

3. Що буде виведено за виконання?
code
try:
    k = 10
    while k>=5:
        k += -3
    print(k, end=';')
except:
    print('ERR')
code

4. Що буде виведено за виконання?
code
try:
    for b in range(5, -15, 2):
        if b<=5:
            continue
        print(b, end=';')
    else:
        print('end',end=';')
    print(b)
except:
    print('ERR')
code

5. Що буде виведено за виконання?\par (Дійсні значення, якщо наявні, в цьому завданні будуть виводитись у форматі з фіксованою крапкою та, можливо, з доволі великою кількістю десяткових розрядів. У відповіді для дійсних значень у ЦЬОМУ завданні достатньо вказати один десятковий розряд після крапки, відкинувши решту розрядів без округлення. Наприклад, замість 13.66666666666667 достатньо вказати 13.6, замість 25.23 достатньо вказати 25.2)
code
try:
    print(2, end=';')
    print(int('a3'), end=';')
    print(6, end=';')
except Exception:
    print(7, end=';')
except ZeroDivisionError:
    print(4, end=';')
except:
    print('ERR', end=';')
else:
    print(5, end=';')
finally:
    print(9)
code

6. Що буде виведено за виконання?
code
def f(s):
    for el in s:
        el += 1
        return s
 
s = [1, 2, 3]
s = f(s)
print(s)
code

7. Що буде виведено за виконання?
code
class A:
    b = 1

    def __h1(self): print(2, end=';')
    def _h2(self): print(3, end=';')

    def main(self):
        try: print(self.b, end=';')
        except: print('ERR', end=';')

        try: self.__h1()
        except: print('ERR', end=';')

        try: self._h2()
        except: print('ERR', end=';')


class B(A):
    b = 4

    def __h1(self): print(5, end=';')
    def _h2(self): print(6, end=';')

try: B().main()
except: print('ERR', end=';')
code

8. Що буде виведено за виконання?
code
def f(arg, n):
    n = 8
    tmp = arg
    tmp[0:] = 'ab'
    return len(tmp)>3

try:
    c = ['0','1','2','3','4']
    n = 4
    f(c, n)
    print(n, end=';')
    for el in c:
        print(el, end=';')
except:
    print('ERR')
code

9. Що буде виведено за виконання?
code
class A:
    a = 1
    c = 2

    def __init__(self):
        a = 3

class B(A):
    a = 4

    def __init__(self):
        super().__init__()
        self.b = 7

obj = B()
obj.b = obj.a + 10
B.b = 13

try: print(obj.b, end= ';')
except: print('ERR', end=';')

try: print(A.b, end= ';')
except: print('ERR', end=';')

try: print(B.b, end= ';')
except: print('ERR', end=';')
code

10. 
Для графа на рисунку з вершини 1 запущено алгоритм пошуку в глибину, що будує кістякове дерево. Списки суміжності вершин переглядаються алгоритмом за зростанням міток вершин.\par
\begin{center}
	\adjustimage{max size={0.9\linewidth}{0.9\paperheight}}
	{../10/Traversals/285.png}
\end{center}
\par Які з наведених ребер входять до складу отриманого кістякового дерева?\par
1) (2, 5)\par
2) (3, 4)\par
3) (1, 6)\par
4) (1, 3)\par
5) (1, 5)\par
6) (1, 2)\par
{\vskip 15pt} У формі вибрати номери відповідних ребер.\par
%answer:1 2
\end {large}}
%end
{\smallskip \begin {small} Затверджено на засіданні кафедри теоретичної кібернетики, протокол \No 12 від   19.05.2020 р.\par\smallskip\smallskip
{\parbox {6.5 cm}{Зав. кафедри \dotfill Ю.В.~Крак}}\hspace {0.7cm}{\hbox to 6.5 cm{ Екзаменатор \dotfill Т.О.~Карнаух}} \end{small}}
%begin
{{\begin {small}\par
{ \center  Київський національний університет імені Тараса Шевченка \par}
{\parbox {5 cm}{Освітня програма \hfill}{\parbox {7 cm}{Прикладна математика, Системний аналіз \hfill}}\parbox {6 cm}{\hfill Семестр 2}}\par
{\parbox {5 cm} {Навчальний предмет \hfill}\parbox {7 cm}{Програмування \hfill}\parbox {6cm}{\hfill}\par}\vspace{-7pt} \end{small}}{\center Екзаменаційний білет \No \textbf 
{ 132 }\par}}
{
\begin {large}
1. Що буде виведено за виконання?
code
try:
    res = True!=9 and not 6==7.0 and 3<=4.0
    if isinstance(res, bool):
        print(f'bool;{res}')
    elif isinstance(res, int):
        print(f'int;{res}')
    elif isinstance(res, float):
        print(f'float;{res:.2f}')
    else:
        print('???')
except:
    print('ERR')
code

2. Що буде виведено за виконання?
code
a = 3
b = 5
c = 0

def f(b):
    a = 1
    b = 2
    c = 3
    return a + b + c

a = 6
b = 9
c = 8

try:
    print(f(a), a, b, c, sep=';')
except:
    print('ERR', a, b, c, sep=';')
code

3. Що буде виведено за виконання?
code
try:
    t = 14
    while t>8:
        t += -3
    print(t, end=';')
except:
    print('ERR')
code

4. Що буде виведено за виконання?
code
try:
    for d in range(15, 4, 3):
        if d<4:
            continue
        print(d, end=';');d = 5
    else:
        print(d,end=';')
    print(d)
except:
    print('ERR')
code

5. Що буде виведено за виконання?\par (Дійсні значення, якщо наявні, в цьому завданні будуть виводитись у форматі з фіксованою крапкою та, можливо, з доволі великою кількістю десяткових розрядів. У відповіді для дійсних значень у ЦЬОМУ завданні достатньо вказати один десятковий розряд після крапки, відкинувши решту розрядів без округлення. Наприклад, замість 13.66666666666667 достатньо вказати 13.6, замість 25.23 достатньо вказати 25.2)
code
try:
    print(1, end=';')
    print(8%1, end=';')
    print(0, end=';')
except ValueError:
    print(7, end=';')
except KeyboardInterrupt:
    print(4, end=';')
except:
    print('ERR', end=';')
else:
    print(8, end=';')
finally:
    print(0)
code

6. Що буде виведено за виконання?
code
def f(n):
    if n>9:
        f(n//100)
    else:
        print(n%10,end='')

f(987654)
code

7. Що буде виведено за виконання?
code
class A:
    d = 1

    def _f1(self): print(2, end=';')
    def f2(self): print(3, end=';')

    def main(self):
        try: print(self.d, end=';')
        except: print('ERR', end=';')

        try: self._f1()
        except: print('ERR', end=';')

        try: self.f2()
        except: print('ERR', end=';')


class B(A):
    d = 4

    def _f1(self): print(5, end=';')
    def f2(self): print(6, end=';')

try: B().main()
except: print('ERR', end=';')
code

8. Що буде виведено за виконання?
code
def f(arg, n):
    arg[23] = 8
    n += -2

try:
    n = 7
    t = {57:0, 59:1, 57:7}
    f(t, n)
    print(n, end=';')
    for x in t :
        print(x, sep=';', end=';')
except:
    print('ERR')
code

9. Що буде виведено за виконання?
code
class A:
    a = 1
    b = 2

    def __init__(self):
        b = 3

class B(A):
    b = 4

    def __init__(self):
        super().__init__()
        self.c = 7

obj = B()
obj.c = obj.b + 10
B.c = 13

try: print(obj.c, end= ';')
except: print('ERR', end=';')

try: print(A.c, end= ';')
except: print('ERR', end=';')

try: print(B.c, end= ';')
except: print('ERR', end=';')
code

10. 
Для графа на рисунку з вершини 0 запущено алгоритм пошуку в глибину, що будує кістякове дерево. Списки суміжності вершин переглядаються алгоритмом за зростанням міток вершин.\par
\begin{center}
	\adjustimage{max size={0.9\linewidth}{0.9\paperheight}}
	{../10/Traversals/134.png}
\end{center}
\par Які з наведених ребер входять до складу отриманого кістякового дерева?\par
1) (1, 3)\par
2) (2, 6)\par
3) (0, 6)\par
4) (2, 5)\par
5) (2, 4)\par
6) (0, 3)\par
{\vskip 15pt} У формі вибрати номери відповідних ребер.\par
%answer:1 2 4 5 6
\end {large}}
%end
{\smallskip \begin {small} Затверджено на засіданні кафедри теоретичної кібернетики, протокол \No 12 від   19.05.2020 р.\par\smallskip\smallskip
{\parbox {6.5 cm}{Зав. кафедри \dotfill Ю.В.~Крак}}\hspace {0.7cm}{\hbox to 6.5 cm{ Екзаменатор \dotfill Т.О.~Карнаух}} \end{small}}
%begin
{{\begin {small}\par
{ \center  Київський національний університет імені Тараса Шевченка \par}
{\parbox {5 cm}{Освітня програма \hfill}{\parbox {7 cm}{Прикладна математика, Системний аналіз \hfill}}\parbox {6 cm}{\hfill Семестр 2}}\par
{\parbox {5 cm} {Навчальний предмет \hfill}\parbox {7 cm}{Програмування \hfill}\parbox {6cm}{\hfill}\par}\vspace{-7pt} \end{small}}{\center Екзаменаційний білет \No \textbf 
{ 133 }\par}}
{
\begin {large}
1. Що буде виведено за виконання?
code
try:
    res = False>="True"
    if isinstance(res, bool):
        print(f'bool;{res}')
    elif isinstance(res, int):
        print(f'int;{res}')
    elif isinstance(res, float):
        print(f'float;{res:.2f}')
    else:
        print('???')
except:
    print('ERR')
code

2. Що буде виведено за виконання?
code
a = 4
b = 1
c = 2

def h(b):
    global c
    a = 5
    b = 4
    c = 1
    return a + b + c

a = 7
b = 4
c = 6

try:
    print(h(b), a, b, c, sep=';')
except:
    print('ERR', a, b, c, sep=';')
code

3. Що буде виведено за виконання?
code
try:
    j = 6
    while j>=1:
        j += -2
    print(j, end=';')
except:
    print('ERR')
code

4. Що буде виведено за виконання?
code
try:
    for e in range(10, 13, -1):
        if e<4:
            continue
        print(e, end=';')
    else:
        print(e,end=';')
    print(e)
except:
    print('ERR')
code

5. Що буде виведено за виконання?\par (Дійсні значення, якщо наявні, в цьому завданні будуть виводитись у форматі з фіксованою крапкою та, можливо, з доволі великою кількістю десяткових розрядів. У відповіді для дійсних значень у ЦЬОМУ завданні достатньо вказати один десятковий розряд після крапки, відкинувши решту розрядів без округлення. Наприклад, замість 13.66666666666667 достатньо вказати 13.6, замість 25.23 достатньо вказати 25.2)
code
try:
    print(9, end=';')
    print(int('2'), end=';')
    print(1, end=';')
except TypeError:
    print(6, end=';')
except ZeroDivisionError:
    print(5, end=';')
except:
    print('ERR', end=';')
else:
    print(3, end=';')
finally:
    print(5)
code

6. Що буде виведено за виконання?
code
def f(s):
    s1 = s
    for i in range(len(s)):
        s1[i] += 1
    return s1

s = [1, 2, 3]
f(s)
print(s)
code

7. Що буде виведено за виконання?
code
class A:
    c = 1

    def _h1(self): print(2, end=';')
    def __h2(self): print(3, end=';')

    def main(self):
        try: print(self.c, end=';')
        except: print('ERR', end=';')

        try: self._h1()
        except: print('ERR', end=';')

        try: self.__h2()
        except: print('ERR', end=';')


class B(A):
    c = 4

    def _h1(self): print(5, end=';')
    def __h2(self): print(6, end=';')

try: B().main()
except: print('ERR', end=';')
code

8. Що буде виведено за виконання?
code
def f(arg, n):
    arg[28] = 8
    n -= -3

try:
    n = 4
    t = {23:6, 82:8, 28:2, 87:6}
    f(t, n)
    print(n, end=';')
    for x in t.keys():
        print(x, sep=';', end=';')
except:
    print('ERR')
code

9. Що буде виведено за виконання?
code
class A:
    a = 1
    c = 2

    def __init__(self):
        a = 3

class B(A):
    a = 4

    def __init__(self):
        super().__init__()
        self.b = 7

obj = B()
obj.b = obj.a + 10
B.b = 13

try: print(obj.b, end= ';')
except: print('ERR', end=';')

try: print(A.b, end= ';')
except: print('ERR', end=';')

try: print(B.b, end= ';')
except: print('ERR', end=';')
code

10. 
Для графа на рисунку з вершини 1 запущено алгоритм пошуку в глибину, що будує кістякове дерево. Списки суміжності вершин переглядаються алгоритмом за зростанням міток вершин.\par
\begin{center}
	\adjustimage{max size={0.9\linewidth}{0.9\paperheight}}
	{../10/Traversals/18.png}
\end{center}
\par Які з наведених ребер входять до складу отриманого кістякового дерева?\par
1) (2, 5)\par
2) (0, 2)\par
3) (4, 6)\par
4) (4, 5)\par
5) (3, 4)\par
6) (1, 6)\par
{\vskip 15pt} У формі вибрати номери відповідних ребер.\par
%answer:1 2 3 4 5
\end {large}}
%end
{\smallskip \begin {small} Затверджено на засіданні кафедри теоретичної кібернетики, протокол \No 12 від   19.05.2020 р.\par\smallskip\smallskip
{\parbox {6.5 cm}{Зав. кафедри \dotfill Ю.В.~Крак}}\hspace {0.7cm}{\hbox to 6.5 cm{ Екзаменатор \dotfill Т.О.~Карнаух}} \end{small}}
%begin
{{\begin {small}\par
{ \center  Київський національний університет імені Тараса Шевченка \par}
{\parbox {5 cm}{Освітня програма \hfill}{\parbox {7 cm}{Прикладна математика, Системний аналіз \hfill}}\parbox {6 cm}{\hfill Семестр 2}}\par
{\parbox {5 cm} {Навчальний предмет \hfill}\parbox {7 cm}{Програмування \hfill}\parbox {6cm}{\hfill}\par}\vspace{-7pt} \end{small}}{\center Екзаменаційний білет \No \textbf 
{ 134 }\par}}
{
\begin {large}
1. Що буде виведено за виконання?
code
def f(n):
    print('f', end='')
    return n

try:
    print(f(1) or f(-5))
except:
    print('ERR')
code

2. Що буде виведено за виконання?
code
a = 8
b = 9
c = 5

def f(a):
    a = 4
    b = 5
    c = 2
    return a + b + c

a = 7
b = 1
c = 3

try:
    print(f(a), a, b, c, sep=';')
except:
    print('ERR', a, b, c, sep=';')
code

3. Що буде виведено за виконання?
code
try:
    k = 12
    while k>=7:
        k += -2
    print(k, end=';')
except:
    print('ERR')
code

4. Що буде виведено за виконання?
code
try:
    for d in range(13, -6, -3):
        if d<=8:
            break
        print(d, end=';')
    else:
        print(d,end=';')
    print(d)
except:
    print('ERR')
code

5. Що буде виведено за виконання?\par (Дійсні значення, якщо наявні, в цьому завданні будуть виводитись у форматі з фіксованою крапкою та, можливо, з доволі великою кількістю десяткових розрядів. У відповіді для дійсних значень у ЦЬОМУ завданні достатньо вказати один десятковий розряд після крапки, відкинувши решту розрядів без округлення. Наприклад, замість 13.66666666666667 достатньо вказати 13.6, замість 25.23 достатньо вказати 25.2)
code
try:
    print(6, end=';')
    print(4//0, end=';')
    print(2, end=';')
except KeyboardInterrupt:
    print(9, end=';')
except TypeError:
    print(8, end=';')
except:
    print('ERR', end=';')
else:
    print(7, end=';')
finally:
    print(1)
code

6. Що буде виведено за виконання?
code
def f(n):
    if n>9:
        f(n//100)
    else:
        print(n%10,end='')

f(123456)
code

7. Що буде виведено за виконання?
code
class A:
    b = 1

    def f1(self): print(2, end=';')
    def _f2(self): print(3, end=';')

    def main(self):
        try: print(self.b, end=';')
        except: print('ERR', end=';')

        try: self.f1()
        except: print('ERR', end=';')

        try: self._f2()
        except: print('ERR', end=';')


class B(A):
    b = 4

    def f1(self): print(5, end=';')
    def _f2(self): print(6, end=';')

try: B().main()
except: print('ERR', end=';')
code

8. Що буде виведено за виконання?
code
def f(arg, n):
    n = 1
    tmp = arg
    tmp[1:-5] = 'yz'
    return len(tmp)>3

try:
    h = [10,11,12,13,14]
    n = 8
    f(h, n)
    print(n, end=';')
    for el in h:
        print(el, end=';')
except:
    print('ERR')
code

9. Що буде виведено за виконання?
code
class A:
    c = 1
    a = 2

    def __init__(self):
        a = 3

class B(A):
    a = 4

    def __init__(self):
        super().__init__()
        self.b = 7

obj = B()
obj.c = obj.b + 10
B.c = 13

try: print(obj.c, end= ';')
except: print('ERR', end=';')

try: print(A.c, end= ';')
except: print('ERR', end=';')

try: print(B.c, end= ';')
except: print('ERR', end=';')
code

10. 
Для графа на рисунку з вершини 5 запущено алгоритм пошуку в глибину, що будує кістякове дерево. Списки суміжності вершин переглядаються алгоритмом за зростанням міток вершин.\par
\begin{center}
	\adjustimage{max size={0.9\linewidth}{0.9\paperheight}}
	{../10/Traversals/249.png}
\end{center}
\par Які з наведених ребер входять до складу отриманого кістякового дерева?\par
1) (2, 5)\par
2) (5, 6)\par
3) (2, 4)\par
4) (0, 1)\par
5) (4, 5)\par
6) (1, 6)\par
{\vskip 15pt} У формі вибрати номери відповідних ребер.\par
%answer:1 4 6
\end {large}}
%end
{\smallskip \begin {small} Затверджено на засіданні кафедри теоретичної кібернетики, протокол \No 12 від   19.05.2020 р.\par\smallskip\smallskip
{\parbox {6.5 cm}{Зав. кафедри \dotfill Ю.В.~Крак}}\hspace {0.7cm}{\hbox to 6.5 cm{ Екзаменатор \dotfill Т.О.~Карнаух}} \end{small}}
%begin
{{\begin {small}\par
{ \center  Київський національний університет імені Тараса Шевченка \par}
{\parbox {5 cm}{Освітня програма \hfill}{\parbox {7 cm}{Прикладна математика, Системний аналіз \hfill}}\parbox {6 cm}{\hfill Семестр 2}}\par
{\parbox {5 cm} {Навчальний предмет \hfill}\parbox {7 cm}{Програмування \hfill}\parbox {6cm}{\hfill}\par}\vspace{-7pt} \end{small}}{\center Екзаменаційний білет \No \textbf 
{ 135 }\par}}
{
\begin {large}
1. Що буде виведено за виконання?
code
try:
    res = 1**1.0*False
    if isinstance(res, bool):
        print(f'bool;{res}')
    elif isinstance(res, int):
        print(f'int;{res}')
    elif isinstance(res, float):
        print(f'float;{res:.2f}')
    else:
        print('???')
except:
    print('ERR')
code

2. Що буде виведено за виконання?
code
a = 3
b = 6
c = 5

def f(b):
    global c
    a += 1
    b = 5
    c = 1
    return a + b + c

a = 2
b = 0
c = 4

try:
    print(f(a), a, b, c, sep=';')
except:
    print('ERR', a, b, c, sep=';')
code

3. Що буде виведено за виконання?
code
try:
    j = 8
    while j<12:
        j += 1
    print(j, end=';')
except:
    print('ERR')
code

4. Що буде виведено за виконання?
code
try:
    for a in range(6, -2, -2):
        if a>=7:
            continue
        print(a, end=';');a = 7
    else:
        print(a,end=';')
    print(a)
except:
    print('ERR')
code

5. Що буде виведено за виконання?\par (Дійсні значення, якщо наявні, в цьому завданні будуть виводитись у форматі з фіксованою крапкою та, можливо, з доволі великою кількістю десяткових розрядів. У відповіді для дійсних значень у ЦЬОМУ завданні достатньо вказати один десятковий розряд після крапки, відкинувши решту розрядів без округлення. Наприклад, замість 13.66666666666667 достатньо вказати 13.6, замість 25.23 достатньо вказати 25.2)
code
try:
    print(0, end=';')
    print(3j>5j, end=';')
    print(3, end=';')
except Exception:
    print(5, end=';')
except ValueError:
    print(6, end=';')
except:
    print('ERR', end=';')
else:
    print(8, end=';')
finally:
    print(2)
code

6. Що буде виведено за виконання?
code
def f(s):
    for el in s:
        el += 1
    return s
 
s = [1, 2, 3]
s = f(s)
print(s)
code

7. Що буде виведено за виконання?
code
class A:
    a = 1

    def _g1(self): print(2, end=';')
    def _g2(self): print(3, end=';')

    def main(self):
        try: print(self.a, end=';')
        except: print('ERR', end=';')

        try: self._g1()
        except: print('ERR', end=';')

        try: self._g2()
        except: print('ERR', end=';')


class B(A):
    a = 4

    def _g1(self): print(5, end=';')
    def _g2(self): print(6, end=';')

try: B().main()
except: print('ERR', end=';')
code

8. Що буде виведено за виконання?
code
def f(arg, n):
    arg[22] = 7
    n *= 2

try:
    n = 6
    s = {30:1, 76:7, 75:7, 87:7}
    f(s, n)
    print(n, end=';')
    for x, y in s.items():
        print(x, y, sep=';', end=';')
except:
    print('ERR')
code

9. Що буде виведено за виконання?
code
class A:
    b = 1
    c = 2

    def __init__(self):
        b = 3

class B(A):
    b = 4

    def __init__(self):
        super().__init__()
        self.a = 7

obj = B()
obj.c = obj.a + 10
B.c = 13

try: print(obj.c, end= ';')
except: print('ERR', end=';')

try: print(A.c, end= ';')
except: print('ERR', end=';')

try: print(B.c, end= ';')
except: print('ERR', end=';')
code

10. 
Для графа на рисунку з вершини 1 запущено алгоритм пошуку в глибину, що будує кістякове дерево. Списки суміжності вершин переглядаються алгоритмом за зростанням міток вершин.\par
\begin{center}
	\adjustimage{max size={0.9\linewidth}{0.9\paperheight}}
	{../10/Traversals/193.png}
\end{center}
\par Які з наведених ребер входять до складу отриманого кістякового дерева?\par
1) (0, 5)\par
2) (0, 2)\par
3) (3, 5)\par
4) (0, 1)\par
5) (1, 2)\par
6) (2, 4)\par
{\vskip 15pt} У формі вибрати номери відповідних ребер.\par
%answer:2 3 4 6
\end {large}}
%end
{\smallskip \begin {small} Затверджено на засіданні кафедри теоретичної кібернетики, протокол \No 12 від   19.05.2020 р.\par\smallskip\smallskip
{\parbox {6.5 cm}{Зав. кафедри \dotfill Ю.В.~Крак}}\hspace {0.7cm}{\hbox to 6.5 cm{ Екзаменатор \dotfill Т.О.~Карнаух}} \end{small}}
%begin
{{\begin {small}\par
{ \center  Київський національний університет імені Тараса Шевченка \par}
{\parbox {5 cm}{Освітня програма \hfill}{\parbox {7 cm}{Прикладна математика, Системний аналіз \hfill}}\parbox {6 cm}{\hfill Семестр 2}}\par
{\parbox {5 cm} {Навчальний предмет \hfill}\parbox {7 cm}{Програмування \hfill}\parbox {6cm}{\hfill}\par}\vspace{-7pt} \end{small}}{\center Екзаменаційний білет \No \textbf 
{ 136 }\par}}
{
\begin {large}
1. Що буде виведено за виконання?
code
try:
    res = 4/6//5*4
    if isinstance(res, bool):
        print(f'bool;{res}')
    elif isinstance(res, int):
        print(f'int;{res}')
    elif isinstance(res, float):
        print(f'float;{res:.2f}')
    else:
        print('???')
except:
    print('ERR')
code

2. Що буде виведено за виконання?
code
a = 7
b = 5
c = 6

def h(a):
    a = 3
    b *= 5
    c = 4
    return a + b + c

a = 8
b = 2
c = 9

try:
    print(h(a), a, b, c, sep=';')
except:
    print('ERR', a, b, c, sep=';')
code

3. Що буде виведено за виконання?
code
try:
    i = 9
    while i<9:
        i += 1
    print(i, end=';')
except:
    print('ERR')
code

4. Що буде виведено за виконання?
code
try:
    for c in range(-12, 2, 2):
        if c>3:
            break
        print(c, end=';');c = 1
    else:
        print(c,end=';')
    print(c)
except:
    print('ERR')
code

5. Що буде виведено за виконання?\par (Дійсні значення, якщо наявні, в цьому завданні будуть виводитись у форматі з фіксованою крапкою та, можливо, з доволі великою кількістю десяткових розрядів. У відповіді для дійсних значень у ЦЬОМУ завданні достатньо вказати один десятковий розряд після крапки, відкинувши решту розрядів без округлення. Наприклад, замість 13.66666666666667 достатньо вказати 13.6, замість 25.23 достатньо вказати 25.2)
code
try:
    print(7, end=';')
    print(int('9'), end=';')
    print(4, end=';')
except TypeError:
    print(0, end=';')
except ZeroDivisionError:
    print(1, end=';')
except:
    print('ERR', end=';')
else:
    print(5, end=';')
finally:
    print(1)
code

6. Що буде виведено за виконання?
code
def f(n):
    if n<=9:
        print(n%10, end='')
    else:
        f(n//10)

f(123456)
code

7. Що буде виведено за виконання?
code
class A:
    d = 1

    def _h1(self): print(2, end=';')
    def h2(self): print(3, end=';')

    def main(self):
        try: print(self.d, end=';')
        except: print('ERR', end=';')

        try: self._h1()
        except: print('ERR', end=';')

        try: self.h2()
        except: print('ERR', end=';')


class B(A):
    d = 4

    def _h1(self): print(5, end=';')
    def h2(self): print(6, end=';')

try: B().main()
except: print('ERR', end=';')
code

8. Що буде виведено за виконання?
code
def f(arg, n):
    n = 9
    tmp = arg
    del tmp[:-3]
    return len(tmp)>3

try:
    e = [10,11,12,13,14]
    n = 3
    f(e, n)
    print(n, end=';')
    for el in e:
        print(el, end=';')
except:
    print('ERR')
code

9. Що буде виведено за виконання?
code
class A:
    c = 1
    b = 2

    def __init__(self):
        b = 3

class B(A):
    b = 4

    def __init__(self):
        super().__init__()
        self.a = 7

obj = B()
obj.a = obj.c + 10
B.a = 13

try: print(obj.a, end= ';')
except: print('ERR', end=';')

try: print(A.a, end= ';')
except: print('ERR', end=';')

try: print(B.a, end= ';')
except: print('ERR', end=';')
code

10. 
Для графа на рисунку з вершини 0 запущено алгоритм пошуку в глибину, що будує кістякове дерево. Списки суміжності вершин переглядаються алгоритмом за зростанням міток вершин.\par
\begin{center}
	\adjustimage{max size={0.9\linewidth}{0.9\paperheight}}
	{../10/Traversals/272.png}
\end{center}
\par Які з наведених ребер входять до складу отриманого кістякового дерева?\par
1) (3, 5)\par
2) (0, 2)\par
3) (3, 6)\par
4) (1, 3)\par
5) (4, 5)\par
6) (0, 1)\par
{\vskip 15pt} У формі вибрати номери відповідних ребер.\par
%answer:1 4 5 6
\end {large}}
%end
{\smallskip \begin {small} Затверджено на засіданні кафедри теоретичної кібернетики, протокол \No 12 від   19.05.2020 р.\par\smallskip\smallskip
{\parbox {6.5 cm}{Зав. кафедри \dotfill Ю.В.~Крак}}\hspace {0.7cm}{\hbox to 6.5 cm{ Екзаменатор \dotfill Т.О.~Карнаух}} \end{small}}
%begin
{{\begin {small}\par
{ \center  Київський національний університет імені Тараса Шевченка \par}
{\parbox {5 cm}{Освітня програма \hfill}{\parbox {7 cm}{Прикладна математика, Системний аналіз \hfill}}\parbox {6 cm}{\hfill Семестр 2}}\par
{\parbox {5 cm} {Навчальний предмет \hfill}\parbox {7 cm}{Програмування \hfill}\parbox {6cm}{\hfill}\par}\vspace{-7pt} \end{small}}{\center Екзаменаційний білет \No \textbf 
{ 137 }\par}}
{
\begin {large}
1. Що буде виведено за виконання?
code
try:
    res = -1**1.0+-1
    if isinstance(res, bool):
        print(f'bool;{res}')
    elif isinstance(res, int):
        print(f'int;{res}')
    elif isinstance(res, float):
        print(f'float;{res:.2f}')
    else:
        print('???')
except:
    print('ERR')
code

2. Що буде виведено за виконання?
code
a = 3
b = 1
c = 3

def g(b):
    global c
    a -= 2
    b = 4
    c = 3
    return a + b + c

a = 1
b = 8
c = 5

try:
    print(g(b), a, b, c, sep=';')
except:
    print('ERR', a, b, c, sep=';')
code

3. Що буде виведено за виконання?
code
try:
    k = 3
    while k<=7:
        k += 2
    print(k, end=';')
except:
    print('ERR')
code

4. Що буде виведено за виконання?
code
try:
    for b in range(3, 9, -1):
        if b<=5:
            continue
        print(b, end=';')
    else:
        print('end',end=';')
    print(b)
except:
    print('ERR')
code

5. Що буде виведено за виконання?\par (Дійсні значення, якщо наявні, в цьому завданні будуть виводитись у форматі з фіксованою крапкою та, можливо, з доволі великою кількістю десяткових розрядів. У відповіді для дійсних значень у ЦЬОМУ завданні достатньо вказати один десятковий розряд після крапки, відкинувши решту розрядів без округлення. Наприклад, замість 13.66666666666667 достатньо вказати 13.6, замість 25.23 достатньо вказати 25.2)
code
try:
    print(7, end=';')
    print(8/3, end=';')
    print(6, end=';')
except ValueError:
    print(4, end=';')
except KeyboardInterrupt:
    print(3, end=';')
except:
    print('ERR', end=';')
else:
    print(9, end=';')
finally:
    print(2)
code

6. Що буде виведено за виконання?
code
def f(s):
    for i in range(len(s)):
        s[i] += 1
        return
 
s = [1, 2, 3]
f(s)
print(s)
code

7. Що буде виведено за виконання?
code
class A:
    d = 1

    def _g1(self): print(2, end=';')
    def _g2(self): print(3, end=';')

    def main(self):
        try: print(self.d, end=';')
        except: print('ERR', end=';')

        try: self._g1()
        except: print('ERR', end=';')

        try: self._g2()
        except: print('ERR', end=';')


class B(A):
    d = 4

    def _g1(self): print(5, end=';')
    def _g2(self): print(6, end=';')

try: B().main()
except: print('ERR', end=';')
code

8. Що буде виведено за виконання?
code
def f(arg, n):
    n = 4
    tmp = arg
    tmp[:-7] = [2,3]
    return len(tmp)>3

try:
    d = ['0','1','2','3','4']
    n = 2
    f(d, n)
    print(n, end=';')
    for el in d:
        print(el, end=';')
except:
    print('ERR')
code

9. Що буде виведено за виконання?
code
class A:
    a = 1
    b = 2

    def __init__(self):
        b = 3

class B(A):
    b = 4

    def __init__(self):
        super().__init__()
        self.c = 7

obj = B()
obj.b = obj.c + 10
B.b = 13

try: print(obj.b, end= ';')
except: print('ERR', end=';')

try: print(A.b, end= ';')
except: print('ERR', end=';')

try: print(B.b, end= ';')
except: print('ERR', end=';')
code

10. 
Для графа на рисунку з вершини 1 запущено алгоритм пошуку в глибину, що будує кістякове дерево. Списки суміжності вершин переглядаються алгоритмом за зростанням міток вершин.\par
\begin{center}
	\adjustimage{max size={0.9\linewidth}{0.9\paperheight}}
	{../10/Traversals/136.png}
\end{center}
\par Які з наведених ребер входять до складу отриманого кістякового дерева?\par
1) (1, 4)\par
2) (0, 5)\par
3) (0, 6)\par
4) (2, 5)\par
5) (1, 3)\par
6) (0, 4)\par
{\vskip 15pt} У формі вибрати номери відповідних ребер.\par
%answer:2 4 6
\end {large}}
%end
{\smallskip \begin {small} Затверджено на засіданні кафедри теоретичної кібернетики, протокол \No 12 від   19.05.2020 р.\par\smallskip\smallskip
{\parbox {6.5 cm}{Зав. кафедри \dotfill Ю.В.~Крак}}\hspace {0.7cm}{\hbox to 6.5 cm{ Екзаменатор \dotfill Т.О.~Карнаух}} \end{small}}
%begin
{{\begin {small}\par
{ \center  Київський національний університет імені Тараса Шевченка \par}
{\parbox {5 cm}{Освітня програма \hfill}{\parbox {7 cm}{Прикладна математика, Системний аналіз \hfill}}\parbox {6 cm}{\hfill Семестр 2}}\par
{\parbox {5 cm} {Навчальний предмет \hfill}\parbox {7 cm}{Програмування \hfill}\parbox {6cm}{\hfill}\par}\vspace{-7pt} \end{small}}{\center Екзаменаційний білет \No \textbf 
{ 138 }\par}}
{
\begin {large}
1. Що буде виведено за виконання?
code
try:
    res = 8<2 and not True>5.0 or 5!=2.0
    if isinstance(res, bool):
        print(f'bool;{res}')
    elif isinstance(res, int):
        print(f'int;{res}')
    elif isinstance(res, float):
        print(f'float;{res:.2f}')
    else:
        print('???')
except:
    print('ERR')
code

2. Що буде виведено за виконання?
code
a = 2
b = 0
c = 4

def g(b):
    global c
    a = 3
    b = 5
    c = 2
    return a + b + c

a = 3
b = 6
c = 8

try:
    print(g(a), a, b, c, sep=';')
except:
    print('ERR', a, b, c, sep=';')
code

3. Що буде виведено за виконання?
code
try:
    j = 9
    while j>9:
        j += -3
    print(j, end=';')
except:
    print('ERR')
code

4. Що буде виведено за виконання?
code
try:
    for f in range(-5, -8, -2):
        if f>=6:
            continue
        print(f, end=';')
        if f<7:
            break
    else:
        print(f,end=';')
    print(f)
except:
    print('ERR')
code

5. Що буде виведено за виконання?\par (Дійсні значення, якщо наявні, в цьому завданні будуть виводитись у форматі з фіксованою крапкою та, можливо, з доволі великою кількістю десяткових розрядів. У відповіді для дійсних значень у ЦЬОМУ завданні достатньо вказати один десятковий розряд після крапки, відкинувши решту розрядів без округлення. Наприклад, замість 13.66666666666667 достатньо вказати 13.6, замість 25.23 достатньо вказати 25.2)
code
try:
    print(0, end=';')
    print(3%0.0, end=';')
    print(9, end=';')
except Exception:
    print(0, end=';')
except Exception:
    print(5, end=';')
except:
    print('ERR', end=';')
else:
    print(6, end=';')
finally:
    print(7)
code

6. Що буде виведено за виконання?
code
def f(n):
    print(n%10,end='')
    if n>9:
        f(n//100)
 
f(123456)
code

7. Що буде виведено за виконання?
code
class A:
    b = 1

    def f1(self): print(2, end=';')
    def _f2(self): print(3, end=';')

    def main(self):
        try: print(self.b, end=';')
        except: print('ERR', end=';')

        try: self.f1()
        except: print('ERR', end=';')

        try: self._f2()
        except: print('ERR', end=';')


class B(A):
    b = 4

    def f1(self): print(5, end=';')
    def _f2(self): print(6, end=';')

try: B().main()
except: print('ERR', end=';')
code

8. Що буде виведено за виконання?
code
def f(arg, n):
    arg[10] = 1
    n += 2

try:
    n = 5
    d = {10:5, 41:1, 10:2}
    f(d, n)
    print(n, end=';')
    for x in d :
        print(x, sep=';', end=';')
except:
    print('ERR')
code

9. Що буде виведено за виконання?
code
class A:
    b = 1
    a = 2

    def __init__(self):
        b = 3

class B(A):
    b = 4

    def __init__(self):
        super().__init__()
        self.c = 7

obj = B()
obj.a = obj.b + 10
B.a = 13

try: print(obj.a, end= ';')
except: print('ERR', end=';')

try: print(A.a, end= ';')
except: print('ERR', end=';')

try: print(B.a, end= ';')
except: print('ERR', end=';')
code

10. 
Для графа на рисунку з вершини 5 запущено алгоритм пошуку в глибину, що будує кістякове дерево. Списки суміжності вершин переглядаються алгоритмом за зростанням міток вершин.\par
\begin{center}
	\adjustimage{max size={0.9\linewidth}{0.9\paperheight}}
	{../10/Traversals/65.png}
\end{center}
\par Які з наведених ребер входять до складу отриманого кістякового дерева?\par
1) (3, 5)\par
2) (0, 1)\par
3) (2, 6)\par
4) (3, 4)\par
5) (0, 5)\par
6) (0, 2)\par
{\vskip 15pt} У формі вибрати номери відповідних ребер.\par
%answer:1 2 3 4 5
\end {large}}
%end
{\smallskip \begin {small} Затверджено на засіданні кафедри теоретичної кібернетики, протокол \No 12 від   19.05.2020 р.\par\smallskip\smallskip
{\parbox {6.5 cm}{Зав. кафедри \dotfill Ю.В.~Крак}}\hspace {0.7cm}{\hbox to 6.5 cm{ Екзаменатор \dotfill Т.О.~Карнаух}} \end{small}}
%begin
{{\begin {small}\par
{ \center  Київський національний університет імені Тараса Шевченка \par}
{\parbox {5 cm}{Освітня програма \hfill}{\parbox {7 cm}{Прикладна математика, Системний аналіз \hfill}}\parbox {6 cm}{\hfill Семестр 2}}\par
{\parbox {5 cm} {Навчальний предмет \hfill}\parbox {7 cm}{Програмування \hfill}\parbox {6cm}{\hfill}\par}\vspace{-7pt} \end{small}}{\center Екзаменаційний білет \No \textbf 
{ 139 }\par}}
{
\begin {large}
1. Що буде виведено за виконання?
code
try:
    res = 9**0.5*-2
    if isinstance(res, bool):
        print(f'bool;{res}')
    elif isinstance(res, int):
        print(f'int;{res}')
    elif isinstance(res, float):
        print(f'float;{res:.2f}')
    else:
        print('???')
except:
    print('ERR')
code

2. Що буде виведено за виконання?
code
a = 2
b = 5
c = 1

def h(a):
    a = 3
    b = 4
    c = 1
    return a + b + c

a = 7
b = 0
c = 9

try:
    print(h(b), a, b, c, sep=';')
except:
    print('ERR', a, b, c, sep=';')
code

3. Що буде виведено за виконання?
code
try:
    i = 5
    while i>=0:
        i += -3
    print(i, end=';')
except:
    print('ERR')
code

4. Що буде виведено за виконання?
code
try:
    for e in range(-1, -7, -3):
        if e<5:
            continue
        print(e, end=';');e = 9
    else:
        print('end',end=';')
    print(e)
except:
    print('ERR')
code

5. Що буде виведено за виконання?\par (Дійсні значення, якщо наявні, в цьому завданні будуть виводитись у форматі з фіксованою крапкою та, можливо, з доволі великою кількістю десяткових розрядів. У відповіді для дійсних значень у ЦЬОМУ завданні достатньо вказати один десятковий розряд після крапки, відкинувши решту розрядів без округлення. Наприклад, замість 13.66666666666667 достатньо вказати 13.6, замість 25.23 достатньо вказати 25.2)
code
try:
    print(1, end=';')
    print(2j<=6j, end=';')
    print(8, end=';')
except ZeroDivisionError:
    print(4, end=';')
except ValueError:
    print(5, end=';')
except:
    print('ERR', end=';')
else:
    print(0, end=';')
finally:
    print(6)
code

6. Що буде виведено за виконання?
code
def f(n):
    if n<=9:
        print(n%10, end='')
    else:
        f(n//10)

f(987651)
code

7. Що буде виведено за виконання?
code
class A:
    c = 1

    def _g1(self): print(2, end=';')
    def __g2(self): print(3, end=';')

    def main(self):
        try: print(self.c, end=';')
        except: print('ERR', end=';')

        try: self._g1()
        except: print('ERR', end=';')

        try: self.__g2()
        except: print('ERR', end=';')


class B(A):
    c = 4

    def _g1(self): print(5, end=';')
    def __g2(self): print(6, end=';')

try: B().main()
except: print('ERR', end=';')
code

8. Що буде виведено за виконання?
code
def f(arg, n):
    arg[52] = 9
    n *= -2

try:
    n = 7
    d = {11:0, 19:1, 11:9}
    f(d, n)
    print(n, end=';')
    for x in d.items():
        print(*x, sep=';', end=';')
except:
    print('ERR')
code

9. Що буде виведено за виконання?
code
class A:
    c = 1
    b = 2

    def __init__(self):
        c = 3

class B(A):
    b = 4

    def __init__(self):
        super().__init__()
        self.a = 7

obj = B()
obj.b = obj.c + 10
B.b = 13

try: print(obj.b, end= ';')
except: print('ERR', end=';')

try: print(A.b, end= ';')
except: print('ERR', end=';')

try: print(B.b, end= ';')
except: print('ERR', end=';')
code

10. 
Для графа на рисунку з вершини 1 запущено алгоритм пошуку в глибину, що будує кістякове дерево. Списки суміжності вершин переглядаються алгоритмом за зростанням міток вершин.\par
\begin{center}
	\adjustimage{max size={0.9\linewidth}{0.9\paperheight}}
	{../10/Traversals/279.png}
\end{center}
\par Які з наведених ребер входять до складу отриманого кістякового дерева?\par
1) (3, 4)\par
2) (1, 2)\par
3) (1, 6)\par
4) (0, 4)\par
5) (5, 6)\par
6) (2, 5)\par
{\vskip 15pt} У формі вибрати номери відповідних ребер.\par
%answer:1 2 4 5
\end {large}}
%end
{\smallskip \begin {small} Затверджено на засіданні кафедри теоретичної кібернетики, протокол \No 12 від   19.05.2020 р.\par\smallskip\smallskip
{\parbox {6.5 cm}{Зав. кафедри \dotfill Ю.В.~Крак}}\hspace {0.7cm}{\hbox to 6.5 cm{ Екзаменатор \dotfill Т.О.~Карнаух}} \end{small}}
%begin
{{\begin {small}\par
{ \center  Київський національний університет імені Тараса Шевченка \par}
{\parbox {5 cm}{Освітня програма \hfill}{\parbox {7 cm}{Прикладна математика, Системний аналіз \hfill}}\parbox {6 cm}{\hfill Семестр 2}}\par
{\parbox {5 cm} {Навчальний предмет \hfill}\parbox {7 cm}{Програмування \hfill}\parbox {6cm}{\hfill}\par}\vspace{-7pt} \end{small}}{\center Екзаменаційний білет \No \textbf 
{ 140 }\par}}
{
\begin {large}
1. Що буде виведено за виконання?
code
def f(n):
    print('f', end='')
    return n<-2

try:
    print(f(-2) and f(8))
except:
    print('ERR')
code

2. Що буде виведено за виконання?
code
a = 4
b = 2
c = 9

def f(a):
    a = 1
    b += 2
    c = 5
    return a + b + c

a = 6
b = 7
c = 0

try:
    print(f(a), a, b, c, sep=';')
except:
    print('ERR', a, b, c, sep=';')
code

3. Що буде виведено за виконання?
code
try:
    n = 15
    while n>6:
        n += -2
    print(n, end=';')
except:
    print('ERR')
code

4. Що буде виведено за виконання?
code
try:
    for b in range(-15, -9, 3):
        if b>7:
            break
        print(b, end=';')
    else:
        print(b,end=';')
    print(b)
except:
    print('ERR')
code

5. Що буде виведено за виконання?\par (Дійсні значення, якщо наявні, в цьому завданні будуть виводитись у форматі з фіксованою крапкою та, можливо, з доволі великою кількістю десяткових розрядів. У відповіді для дійсних значень у ЦЬОМУ завданні достатньо вказати один десятковий розряд після крапки, відкинувши решту розрядів без округлення. Наприклад, замість 13.66666666666667 достатньо вказати 13.6, замість 25.23 достатньо вказати 25.2)
code
try:
    print(3, end=';')
    print(int('c1'), end=';')
    print(7, end=';')
except KeyboardInterrupt:
    print(4, end=';')
except TypeError:
    print(2, end=';')
except:
    print('ERR', end=';')
else:
    print(9, end=';')
finally:
    print(8)
code

6. Що буде виведено за виконання?
code
def f(n):
    if n>9: f(n//10)
    else:
        print(n%10, end='')
 
f(123456)
code

7. Що буде виведено за виконання?
code
class A:
    a = 1

    def __f1(self): print(2, end=';')
    def _f2(self): print(3, end=';')

    def main(self):
        try: print(self.a, end=';')
        except: print('ERR', end=';')

        try: self.__f1()
        except: print('ERR', end=';')

        try: self._f2()
        except: print('ERR', end=';')


class B(A):
    a = 4

    def __f1(self): print(5, end=';')
    def _f2(self): print(6, end=';')

try: B().main()
except: print('ERR', end=';')
code

8. Що буде виведено за виконання?
code
def f(arg, n):
    arg[82] = 5
    n = 9
    arg = set(arg.values())

try:
    n = 4
    d = {10:1, 23:8, 74:2, 84:0}
    f(d, n)
    print(n, end=';')
    for x in d.values():
        print(x, sep=';', end=';')
except:
    print('ERR')
code

9. Що буде виведено за виконання?
code
class A:
    a = 1
    b = 2

    def __init__(self):
        b = 3

class B(A):
    a = 4

    def __init__(self):
        super().__init__()
        self.c = 7

obj = B()
obj.a = obj.b + 10
B.a = 13

try: print(obj.a, end= ';')
except: print('ERR', end=';')

try: print(A.a, end= ';')
except: print('ERR', end=';')

try: print(B.a, end= ';')
except: print('ERR', end=';')
code

10. 
Для графа на рисунку з вершини 4 запущено алгоритм пошуку в глибину, що будує кістякове дерево. Списки суміжності вершин переглядаються алгоритмом за зростанням міток вершин.\par
\begin{center}
	\adjustimage{max size={0.9\linewidth}{0.9\paperheight}}
	{../10/Traversals/219.png}
\end{center}
\par Які з наведених ребер входять до складу отриманого кістякового дерева?\par
1) (2, 4)\par
2) (1, 6)\par
3) (0, 6)\par
4) (2, 5)\par
5) (0, 3)\par
6) (3, 6)\par
{\vskip 15pt} У формі вибрати номери відповідних ребер.\par
%answer:2 4 5 6
\end {large}}
%end
{\smallskip \begin {small} Затверджено на засіданні кафедри теоретичної кібернетики, протокол \No 12 від   19.05.2020 р.\par\smallskip\smallskip
{\parbox {6.5 cm}{Зав. кафедри \dotfill Ю.В.~Крак}}\hspace {0.7cm}{\hbox to 6.5 cm{ Екзаменатор \dotfill Т.О.~Карнаух}} \end{small}}
%begin
{{\begin {small}\par
{ \center  Київський національний університет імені Тараса Шевченка \par}
{\parbox {5 cm}{Освітня програма \hfill}{\parbox {7 cm}{Прикладна математика, Системний аналіз \hfill}}\parbox {6 cm}{\hfill Семестр 2}}\par
{\parbox {5 cm} {Навчальний предмет \hfill}\parbox {7 cm}{Програмування \hfill}\parbox {6cm}{\hfill}\par}\vspace{-7pt} \end{small}}{\center Екзаменаційний білет \No \textbf 
{ 141 }\par}}
{
\begin {large}
1. Що буде виведено за виконання?
code
try:
    res = "True">1j
    if isinstance(res, bool):
        print(f'bool;{res}')
    elif isinstance(res, int):
        print(f'int;{res}')
    elif isinstance(res, float):
        print(f'float;{res:.2f}')
    else:
        print('???')
except:
    print('ERR')
code

2. Що буде виведено за виконання?
code
a = 3
b = 9
c = 2

def f(b):
    a = 3
    b = 4
    c = 1
    return a + b + c

a = 1
b = 4
c = 8

try:
    print(f(b), a, b, c, sep=';')
except:
    print('ERR', a, b, c, sep=';')
code

3. Що буде виведено за виконання?
code
try:
    m = 6
    while m>=9:
        m += -3
    print(m, end=';')
except:
    print('ERR')
code

4. Що буде виведено за виконання?
code
try:
    for c in range(1, -13, 3):
        if c<=8:
            continue
        print(c, end=';')
    else:
        print(c,end=';')
    print(c)
except:
    print('ERR')
code

5. Що буде виведено за виконання?\par (Дійсні значення, якщо наявні, в цьому завданні будуть виводитись у форматі з фіксованою крапкою та, можливо, з доволі великою кількістю десяткових розрядів. У відповіді для дійсних значень у ЦЬОМУ завданні достатньо вказати один десятковий розряд після крапки, відкинувши решту розрядів без округлення. Наприклад, замість 13.66666666666667 достатньо вказати 13.6, замість 25.23 достатньо вказати 25.2)
code
try:
    print(6, end=';')
    print(int('d7'), end=';')
    print(1, end=';')
except KeyboardInterrupt:
    print(3, end=';')
except ValueError:
    print(5, end=';')
except:
    print('ERR', end=';')
else:
    print(9, end=';')
finally:
    print(4)
code

6. Що буде виведено за виконання?
code
def f(n):
    print(n%10,end='')
    if n>9:
        f(n//100)
 
f(123456)
code

7. Що буде виведено за виконання?
code
class A:
    a = 1

    def _h1(self): print(2, end=';')
    def h2(self): print(3, end=';')

    def main(self):
        try: print(self.a, end=';')
        except: print('ERR', end=';')

        try: self._h1()
        except: print('ERR', end=';')

        try: self.h2()
        except: print('ERR', end=';')


class B(A):
    a = 4

    def _h1(self): print(5, end=';')
    def h2(self): print(6, end=';')

try: B().main()
except: print('ERR', end=';')
code

8. Що буде виведено за виконання?
code
def f(arg, n):
    arg[79] = 8
    n += 3

try:
    n = 6
    s = {79:8, 66:3, 79:7}
    f(s, n)
    print(n, end=';')
    for x in s.keys():
        print(x, sep=';', end=';')
except:
    print('ERR')
code

9. Що буде виведено за виконання?
code
class A:
    b = 1
    c = 2

    def __init__(self):
        b = 3

class B(A):
    b = 4

    def __init__(self):
        super().__init__()
        self.a = 7

obj = B()
obj.c = obj.a + 10
B.c = 13

try: print(obj.c, end= ';')
except: print('ERR', end=';')

try: print(A.c, end= ';')
except: print('ERR', end=';')

try: print(B.c, end= ';')
except: print('ERR', end=';')
code

10. 
Для графа на рисунку з вершини 1 запущено алгоритм пошуку в глибину, що будує кістякове дерево. Списки суміжності вершин переглядаються алгоритмом за зростанням міток вершин.\par
\begin{center}
	\adjustimage{max size={0.9\linewidth}{0.9\paperheight}}
	{../10/Traversals/20.png}
\end{center}
\par Які з наведених ребер входять до складу отриманого кістякового дерева?\par
1) (2, 6)\par
2) (4, 6)\par
3) (0, 5)\par
4) (5, 6)\par
5) (0, 4)\par
6) (1, 5)\par
{\vskip 15pt} У формі вибрати номери відповідних ребер.\par
%answer:1 2 3 5 6
\end {large}}
%end
{\smallskip \begin {small} Затверджено на засіданні кафедри теоретичної кібернетики, протокол \No 12 від   19.05.2020 р.\par\smallskip\smallskip
{\parbox {6.5 cm}{Зав. кафедри \dotfill Ю.В.~Крак}}\hspace {0.7cm}{\hbox to 6.5 cm{ Екзаменатор \dotfill Т.О.~Карнаух}} \end{small}}
%begin
{{\begin {small}\par
{ \center  Київський національний університет імені Тараса Шевченка \par}
{\parbox {5 cm}{Освітня програма \hfill}{\parbox {7 cm}{Прикладна математика, Системний аналіз \hfill}}\parbox {6 cm}{\hfill Семестр 2}}\par
{\parbox {5 cm} {Навчальний предмет \hfill}\parbox {7 cm}{Програмування \hfill}\parbox {6cm}{\hfill}\par}\vspace{-7pt} \end{small}}{\center Екзаменаційний білет \No \textbf 
{ 142 }\par}}
{
\begin {large}
1. Що буде виведено за виконання?
code
try:
    res = 3**4-1
    if isinstance(res, bool):
        print(f'bool;{res}')
    elif isinstance(res, int):
        print(f'int;{res}')
    elif isinstance(res, float):
        print(f'float;{res:.2f}')
    else:
        print('???')
except:
    print('ERR')
code

2. Що буде виведено за виконання?
code
a = 3
b = 8
c = 0

def h(b):
    global c
    a = 4
    b *= 2
    c = 1
    return a + b + c

a = 6
b = 5
c = 7

try:
    print(h(b), a, b, c, sep=';')
except:
    print('ERR', a, b, c, sep=';')
code

3. Що буде виведено за виконання?
code
try:
    t = 9
    while t<11:
        t += 2
    print(t, end=';')
except:
    print('ERR')
code

4. Що буде виведено за виконання?
code
try:
    for d in range(-3, 11, -1):
        if d<3:
            break
        print(d, end=';')
    else:
        print(d,end=';')
    print(d)
except:
    print('ERR')
code

5. Що буде виведено за виконання?\par (Дійсні значення, якщо наявні, в цьому завданні будуть виводитись у форматі з фіксованою крапкою та, можливо, з доволі великою кількістю десяткових розрядів. У відповіді для дійсних значень у ЦЬОМУ завданні достатньо вказати один десятковий розряд після крапки, відкинувши решту розрядів без округлення. Наприклад, замість 13.66666666666667 достатньо вказати 13.6, замість 25.23 достатньо вказати 25.2)
code
try:
    print(2, end=';')
    print(8j<0j, end=';')
    print(0, end=';')
except TypeError:
    print(7, end=';')
except Exception:
    print(9, end=';')
except:
    print('ERR', end=';')
else:
    print(2, end=';')
finally:
    print(8)
code

6. Що буде виведено за виконання?
code
def f(s):
    for i in range(len(s)):
        s[i] += 1
    return
 
s = [1, 2, 3]
f(s)
print(s)
code

7. Що буде виведено за виконання?
code
class A:
    b = 1

    def g1(self): print(2, end=';')
    def _g2(self): print(3, end=';')

    def main(self):
        try: print(self.b, end=';')
        except: print('ERR', end=';')

        try: self.g1()
        except: print('ERR', end=';')

        try: self._g2()
        except: print('ERR', end=';')


class B(A):
    b = 4

    def g1(self): print(5, end=';')
    def _g2(self): print(6, end=';')

try: B().main()
except: print('ERR', end=';')
code

8. Що буде виведено за виконання?
code
def f(arg, n):
    arg[47] = 6
    n -= -3

try:
    n = 3
    t = {62:7, 84:3, 83:6}
    f(t, n)
    print(n, end=';')
    for x, y in t.items():
        print(x, y, sep=';', end=';')
except:
    print('ERR')
code

9. Що буде виведено за виконання?
code
class A:
    b = 1
    a = 2

    def __init__(self):
        a = 3

class B(A):
    a = 4

    def __init__(self):
        super().__init__()
        self.c = 7

obj = B()
obj.b = obj.c + 10
B.b = 13

try: print(obj.b, end= ';')
except: print('ERR', end=';')

try: print(A.b, end= ';')
except: print('ERR', end=';')

try: print(B.b, end= ';')
except: print('ERR', end=';')
code

10. 
Для графа на рисунку з вершини 4 запущено алгоритм пошуку в глибину, що будує кістякове дерево. Списки суміжності вершин переглядаються алгоритмом за зростанням міток вершин.\par
\begin{center}
	\adjustimage{max size={0.9\linewidth}{0.9\paperheight}}
	{../10/Traversals/83.png}
\end{center}
\par Які з наведених ребер входять до складу отриманого кістякового дерева?\par
1) (1, 4)\par
2) (4, 5)\par
3) (5, 6)\par
4) (2, 4)\par
5) (1, 6)\par
6) (3, 4)\par
{\vskip 15pt} У формі вибрати номери відповідних ребер.\par
%answer:3 5
\end {large}}
%end
{\smallskip \begin {small} Затверджено на засіданні кафедри теоретичної кібернетики, протокол \No 12 від   19.05.2020 р.\par\smallskip\smallskip
{\parbox {6.5 cm}{Зав. кафедри \dotfill Ю.В.~Крак}}\hspace {0.7cm}{\hbox to 6.5 cm{ Екзаменатор \dotfill Т.О.~Карнаух}} \end{small}}
%begin
{{\begin {small}\par
{ \center  Київський національний університет імені Тараса Шевченка \par}
{\parbox {5 cm}{Освітня програма \hfill}{\parbox {7 cm}{Прикладна математика, Системний аналіз \hfill}}\parbox {6 cm}{\hfill Семестр 2}}\par
{\parbox {5 cm} {Навчальний предмет \hfill}\parbox {7 cm}{Програмування \hfill}\parbox {6cm}{\hfill}\par}\vspace{-7pt} \end{small}}{\center Екзаменаційний білет \No \textbf 
{ 143 }\par}}
{
\begin {large}
1. Що буде виведено за виконання?
code
try:
    res = "9"==3.0j
    if isinstance(res, bool):
        print(f'bool;{res}')
    elif isinstance(res, int):
        print(f'int;{res}')
    elif isinstance(res, float):
        print(f'float;{res:.2f}')
    else:
        print('???')
except:
    print('ERR')
code

2. Що буде виведено за виконання?
code
a = 4
b = 2
c = 9

def g(a):
    a = 3
    b -= 5
    c = 1
    return a + b + c

a = 1
b = 6
c = 9

try:
    print(g(b), a, b, c, sep=';')
except:
    print('ERR', a, b, c, sep=';')
code

3. Що буде виведено за виконання?
code
try:
    t = 13
    while t>1:
        t += -2
    print(t, end=';')
except:
    print('ERR')
code

4. Що буде виведено за виконання?
code
try:
    for a in range(8, 3, 2):
        if a>=4:
            continue
        print(a, end=';');a = 6
    else:
        print(a,end=';')
    print(a)
except:
    print('ERR')
code

5. Що буде виведено за виконання?\par (Дійсні значення, якщо наявні, в цьому завданні будуть виводитись у форматі з фіксованою крапкою та, можливо, з доволі великою кількістю десяткових розрядів. У відповіді для дійсних значень у ЦЬОМУ завданні достатньо вказати один десятковий розряд після крапки, відкинувши решту розрядів без округлення. Наприклад, замість 13.66666666666667 достатньо вказати 13.6, замість 25.23 достатньо вказати 25.2)
code
try:
    print(1, end=';')
    print(int('6'), end=';')
    print(3, end=';')
except ZeroDivisionError:
    print(5, end=';')
except ValueError:
    print(4, end=';')
except:
    print('ERR', end=';')
else:
    print(0, end=';')
finally:
    print(7)
code

6. Що буде виведено за виконання?
code
def f(s):
    for el in s:
        el += 1
    return s
 
s = [1, 2, 3]
s = f(s)
print(s)
code

7. Що буде виведено за виконання?
code
class A:
    d = 1

    def _f1(self): print(2, end=';')
    def _f2(self): print(3, end=';')

    def main(self):
        try: print(self.d, end=';')
        except: print('ERR', end=';')

        try: self._f1()
        except: print('ERR', end=';')

        try: self._f2()
        except: print('ERR', end=';')


class B(A):
    d = 4

    def _f1(self): print(5, end=';')
    def _f2(self): print(6, end=';')

try: B().main()
except: print('ERR', end=';')
code

8. Що буде виведено за виконання?
code
def f(arg, n):
    n = 6
    tmp = arg
    del tmp[7:]
    return len(tmp)>3

try:
    h = [10,11,12,13,14]
    n = 0
    f(h, n)
    print(n, end=';')
    for el in h:
        print(el, end=';')
except:
    print('ERR')
code

9. Що буде виведено за виконання?
code
class A:
    c = 1
    a = 2

    def __init__(self):
        c = 3

class B(A):
    a = 4

    def __init__(self):
        super().__init__()
        self.b = 7

obj = B()
obj.a = obj.c + 10
B.a = 13

try: print(obj.a, end= ';')
except: print('ERR', end=';')

try: print(A.a, end= ';')
except: print('ERR', end=';')

try: print(B.a, end= ';')
except: print('ERR', end=';')
code

10. 
Для графа на рисунку з вершини 0 запущено алгоритм пошуку в глибину, що будує кістякове дерево. Списки суміжності вершин переглядаються алгоритмом за зростанням міток вершин.\par
\begin{center}
	\adjustimage{max size={0.9\linewidth}{0.9\paperheight}}
	{../10/Traversals/228.png}
\end{center}
\par Які з наведених ребер входять до складу отриманого кістякового дерева?\par
1) (0, 1)\par
2) (0, 6)\par
3) (2, 5)\par
4) (2, 6)\par
5) (1, 2)\par
6) (4, 5)\par
{\vskip 15pt} У формі вибрати номери відповідних ребер.\par
%answer:1 3 5 6
\end {large}}
%end
{\smallskip \begin {small} Затверджено на засіданні кафедри теоретичної кібернетики, протокол \No 12 від   19.05.2020 р.\par\smallskip\smallskip
{\parbox {6.5 cm}{Зав. кафедри \dotfill Ю.В.~Крак}}\hspace {0.7cm}{\hbox to 6.5 cm{ Екзаменатор \dotfill Т.О.~Карнаух}} \end{small}}
%begin
{{\begin {small}\par
{ \center  Київський національний університет імені Тараса Шевченка \par}
{\parbox {5 cm}{Освітня програма \hfill}{\parbox {7 cm}{Прикладна математика, Системний аналіз \hfill}}\parbox {6 cm}{\hfill Семестр 2}}\par
{\parbox {5 cm} {Навчальний предмет \hfill}\parbox {7 cm}{Програмування \hfill}\parbox {6cm}{\hfill}\par}\vspace{-7pt} \end{small}}{\center Екзаменаційний білет \No \textbf 
{ 144 }\par}}
{
\begin {large}
1. Що буде виведено за виконання?
code
try:
    res = 4**-1.0*2
    if isinstance(res, bool):
        print(f'bool;{res}')
    elif isinstance(res, int):
        print(f'int;{res}')
    elif isinstance(res, float):
        print(f'float;{res:.2f}')
    else:
        print('???')
except:
    print('ERR')
code

2. Що буде виведено за виконання?
code
a = 5
b = 8
c = 3

def g(b):
    a = 5
    b += 3
    c = 4
    return a + b + c

a = 0
b = 7
c = 1

try:
    print(g(a), a, b, c, sep=';')
except:
    print('ERR', a, b, c, sep=';')
code

3. Що буде виведено за виконання?
code
try:
    n = 1
    while n<=10:
        n += 1
    print(n, end=';')
except:
    print('ERR')
code

4. Що буде виведено за виконання?
code
try:
    for f in range(11, -1, -2):
        if f>6:
            continue
        print(f, end=';')
    else:
        print(f,end=';')
    print(f)
except:
    print('ERR')
code

5. Що буде виведено за виконання?\par (Дійсні значення, якщо наявні, в цьому завданні будуть виводитись у форматі з фіксованою крапкою та, можливо, з доволі великою кількістю десяткових розрядів. У відповіді для дійсних значень у ЦЬОМУ завданні достатньо вказати один десятковий розряд після крапки, відкинувши решту розрядів без округлення. Наприклад, замість 13.66666666666667 достатньо вказати 13.6, замість 25.23 достатньо вказати 25.2)
code
try:
    print(8, end=';')
    print(2//0, end=';')
    print(3, end=';')
except TypeError:
    print(4, end=';')
except KeyboardInterrupt:
    print(0, end=';')
except:
    print('ERR', end=';')
else:
    print(9, end=';')
finally:
    print(6)
code

6. Що буде виведено за виконання?
code
def f(n):
    if n>9:
        return f(n//100)+1
    else:
        return 1

print(f(123456))
code

7. Що буде виведено за виконання?
code
class A:
    c = 1

    def _h1(self): print(2, end=';')
    def __h2(self): print(3, end=';')

    def main(self):
        try: print(self.c, end=';')
        except: print('ERR', end=';')

        try: self._h1()
        except: print('ERR', end=';')

        try: self.__h2()
        except: print('ERR', end=';')


class B(A):
    c = 4

    def _h1(self): print(5, end=';')
    def __h2(self): print(6, end=';')

try: B().main()
except: print('ERR', end=';')
code

8. Що буде виведено за виконання?
code
def f(arg, n):
    arg[53] = 6
    n += -2

try:
    n = 4
    t = {49:0, 68:2, 49:8}
    f(t, n)
    print(n, end=';')
    for x in t.values():
        print(x, sep=';', end=';')
except:
    print('ERR')
code

9. Що буде виведено за виконання?
code
class A:
    a = 1
    c = 2

    def __init__(self):
        c = 3

class B(A):
    a = 4

    def __init__(self):
        super().__init__()
        self.b = 7

obj = B()
obj.a = obj.b + 10
B.a = 13

try: print(obj.a, end= ';')
except: print('ERR', end=';')

try: print(A.a, end= ';')
except: print('ERR', end=';')

try: print(B.a, end= ';')
except: print('ERR', end=';')
code

10. 
Для графа на рисунку з вершини 2 запущено алгоритм пошуку в глибину, що будує кістякове дерево. Списки суміжності вершин переглядаються алгоритмом за зростанням міток вершин.\par
\begin{center}
	\adjustimage{max size={0.9\linewidth}{0.9\paperheight}}
	{../10/Traversals/259.png}
\end{center}
\par Які з наведених ребер входять до складу отриманого кістякового дерева?\par
1) (0, 5)\par
2) (4, 5)\par
3) (2, 5)\par
4) (4, 6)\par
5) (2, 6)\par
6) (0, 4)\par
{\vskip 15pt} У формі вибрати номери відповідних ребер.\par
%answer:1 6
\end {large}}
%end
{\smallskip \begin {small} Затверджено на засіданні кафедри теоретичної кібернетики, протокол \No 12 від   19.05.2020 р.\par\smallskip\smallskip
{\parbox {6.5 cm}{Зав. кафедри \dotfill Ю.В.~Крак}}\hspace {0.7cm}{\hbox to 6.5 cm{ Екзаменатор \dotfill Т.О.~Карнаух}} \end{small}}
%begin
{{\begin {small}\par
{ \center  Київський національний університет імені Тараса Шевченка \par}
{\parbox {5 cm}{Освітня програма \hfill}{\parbox {7 cm}{Прикладна математика, Системний аналіз \hfill}}\parbox {6 cm}{\hfill Семестр 2}}\par
{\parbox {5 cm} {Навчальний предмет \hfill}\parbox {7 cm}{Програмування \hfill}\parbox {6cm}{\hfill}\par}\vspace{-7pt} \end{small}}{\center Екзаменаційний білет \No \textbf 
{ 145 }\par}}
{
\begin {large}
1. Що буде виведено за виконання?
code
try:
    res = not 8.0<=7 or 9>=False and 3.0==4
    if isinstance(res, bool):
        print(f'bool;{res}')
    elif isinstance(res, int):
        print(f'int;{res}')
    elif isinstance(res, float):
        print(f'float;{res:.2f}')
    else:
        print('???')
except:
    print('ERR')
code

2. Що буде виведено за виконання?
code
a = 7
b = 0
c = 6

def g(a):
    global c
    a = 5
    b = 2
    c = 3
    return a + b + c

a = 5
b = 3
c = 2

try:
    print(g(a), a, b, c, sep=';')
except:
    print('ERR', a, b, c, sep=';')
code

3. Що буде виведено за виконання?
code
try:
    t = 6
    while t<10:
        t += 1
    print(t, end=';')
except:
    print('ERR')
code

4. Що буде виведено за виконання?
code
try:
    for a in range(-11, 1, -3):
        if a>7:
            continue
        print(a, end=';')
        if a>5:
            break
    else:
        print(a,end=';')
    print(a)
except:
    print('ERR')
code

5. Що буде виведено за виконання?\par (Дійсні значення, якщо наявні, в цьому завданні будуть виводитись у форматі з фіксованою крапкою та, можливо, з доволі великою кількістю десяткових розрядів. У відповіді для дійсних значень у ЦЬОМУ завданні достатньо вказати один десятковий розряд після крапки, відкинувши решту розрядів без округлення. Наприклад, замість 13.66666666666667 достатньо вказати 13.6, замість 25.23 достатньо вказати 25.2)
code
try:
    print(5, end=';')
    print(1%2, end=';')
    print(9, end=';')
except Exception:
    print(5, end=';')
except ZeroDivisionError:
    print(1, end=';')
except:
    print('ERR', end=';')
else:
    print(6, end=';')
finally:
    print(2)
code

6. Що буде виведено за виконання?
code
def f(s):
    s1 = s
    for i in range(len(s)):
        s1[i] += 1
    return s1

s = [1, 2, 3]
f(s)
print(s)
code

7. Що буде виведено за виконання?
code
class A:
    b = 1

    def __h1(self): print(2, end=';')
    def _h2(self): print(3, end=';')

    def main(self):
        try: print(self.b, end=';')
        except: print('ERR', end=';')

        try: self.__h1()
        except: print('ERR', end=';')

        try: self._h2()
        except: print('ERR', end=';')


class B(A):
    b = 4

    def __h1(self): print(5, end=';')
    def _h2(self): print(6, end=';')

try: B().main()
except: print('ERR', end=';')
code

8. Що буде виведено за виконання?
code
def f(arg, n):
    arg[22] = 6
    n *= 3

try:
    n = 5
    s = {22:0, 39:3, 22:6}
    f(s, n)
    print(n, end=';')
    for x in s :
        print(x, sep=';', end=';')
except:
    print('ERR')
code

9. Що буде виведено за виконання?
code
class A:
    a = 1
    b = 2

    def __init__(self):
        b = 3

class B(A):
    a = 4

    def __init__(self):
        super().__init__()
        self.c = 7

obj = B()
obj.a = obj.c + 10
B.a = 13

try: print(obj.a, end= ';')
except: print('ERR', end=';')

try: print(A.a, end= ';')
except: print('ERR', end=';')

try: print(B.a, end= ';')
except: print('ERR', end=';')
code

10. 
Для графа на рисунку з вершини 0 запущено алгоритм пошуку в глибину, що будує кістякове дерево. Списки суміжності вершин переглядаються алгоритмом за зростанням міток вершин.\par
\begin{center}
	\adjustimage{max size={0.9\linewidth}{0.9\paperheight}}
	{../10/Traversals/16.png}
\end{center}
\par Які з наведених ребер входять до складу отриманого кістякового дерева?\par
1) (0, 2)\par
2) (0, 4)\par
3) (3, 5)\par
4) (1, 2)\par
5) (1, 5)\par
6) (3, 6)\par
{\vskip 15pt} У формі вибрати номери відповідних ребер.\par
%answer:1 3 4 5 6
\end {large}}
%end
{\smallskip \begin {small} Затверджено на засіданні кафедри теоретичної кібернетики, протокол \No 12 від   19.05.2020 р.\par\smallskip\smallskip
{\parbox {6.5 cm}{Зав. кафедри \dotfill Ю.В.~Крак}}\hspace {0.7cm}{\hbox to 6.5 cm{ Екзаменатор \dotfill Т.О.~Карнаух}} \end{small}}
%begin
{{\begin {small}\par
{ \center  Київський національний університет імені Тараса Шевченка \par}
{\parbox {5 cm}{Освітня програма \hfill}{\parbox {7 cm}{Прикладна математика, Системний аналіз \hfill}}\parbox {6 cm}{\hfill Семестр 2}}\par
{\parbox {5 cm} {Навчальний предмет \hfill}\parbox {7 cm}{Програмування \hfill}\parbox {6cm}{\hfill}\par}\vspace{-7pt} \end{small}}{\center Екзаменаційний білет \No \textbf 
{ 146 }\par}}
{
\begin {large}
1. Що буде виведено за виконання?
code
def f(n):
    print('f', end='')
    return n

try:
    print(f(-4) or f(-3))
except:
    print('ERR')
code

2. Що буде виведено за виконання?
code
a = 7
b = 5
c = 6

def g(a):
    a = 2
    b = 4
    c = 1
    return a + b + c

a = 0
b = 8
c = 1

try:
    print(g(b), a, b, c, sep=';')
except:
    print('ERR', a, b, c, sep=';')
code

3. Що буде виведено за виконання?
code
try:
    j = 10
    while j>=0:
        j += -2
    print(j, end=';')
except:
    print('ERR')
code

4. Що буде виведено за виконання?
code
try:
    for d in range(14, 10, 3):
        if d<=8:
            break
        print(d, end=';')
    else:
        print(d,end=';')
    print(d)
except:
    print('ERR')
code

5. Що буде виведено за виконання?\par (Дійсні значення, якщо наявні, в цьому завданні будуть виводитись у форматі з фіксованою крапкою та, можливо, з доволі великою кількістю десяткових розрядів. У відповіді для дійсних значень у ЦЬОМУ завданні достатньо вказати один десятковий розряд після крапки, відкинувши решту розрядів без округлення. Наприклад, замість 13.66666666666667 достатньо вказати 13.6, замість 25.23 достатньо вказати 25.2)
code
try:
    print(3, end=';')
    print(int('b7'), end=';')
    print(8, end=';')
except Exception:
    print(0, end=';')
except ValueError:
    print(4, end=';')
except:
    print('ERR', end=';')
else:
    print(5, end=';')
finally:
    print(9)
code

6. Що буде виведено за виконання?
code
def f(n):
    if n>9:
        f(n//10)
    print(n%10, end='')
 
f(123456)
code

7. Що буде виведено за виконання?
code
class A:
    c = 1

    def _g1(self): print(2, end=';')
    def _g2(self): print(3, end=';')

    def main(self):
        try: print(self.c, end=';')
        except: print('ERR', end=';')

        try: self._g1()
        except: print('ERR', end=';')

        try: self._g2()
        except: print('ERR', end=';')


class B(A):
    c = 4

    def _g1(self): print(5, end=';')
    def _g2(self): print(6, end=';')

try: B().main()
except: print('ERR', end=';')
code

8. Що буде виведено за виконання?
code
def f(arg, n):
    n = 5
    tmp = arg
    tmp[6:] = [1,2]
    return len(tmp)>3

try:
    e = ['0','1','2','3','4']
    n = 7
    f(e, n)
    print(n, end=';')
    for el in e:
        print(el, end=';')
except:
    print('ERR')
code

9. Що буде виведено за виконання?
code
class A:
    a = 1
    c = 2

    def __init__(self):
        a = 3

class B(A):
    c = 4

    def __init__(self):
        super().__init__()
        self.b = 7

obj = B()
obj.b = obj.a + 10
B.b = 13

try: print(obj.b, end= ';')
except: print('ERR', end=';')

try: print(A.b, end= ';')
except: print('ERR', end=';')

try: print(B.b, end= ';')
except: print('ERR', end=';')
code

10. 
Для графа на рисунку з вершини 3 запущено алгоритм пошуку в глибину, що будує кістякове дерево. Списки суміжності вершин переглядаються алгоритмом за зростанням міток вершин.\par
\begin{center}
	\adjustimage{max size={0.9\linewidth}{0.9\paperheight}}
	{../10/Traversals/186.png}
\end{center}
\par Які з наведених ребер входять до складу отриманого кістякового дерева?\par
1) (0, 4)\par
2) (2, 3)\par
3) (3, 6)\par
4) (2, 6)\par
5) (4, 5)\par
6) (4, 6)\par
{\vskip 15pt} У формі вибрати номери відповідних ребер.\par
%answer:1 2 5 6
\end {large}}
%end
{\smallskip \begin {small} Затверджено на засіданні кафедри теоретичної кібернетики, протокол \No 12 від   19.05.2020 р.\par\smallskip\smallskip
{\parbox {6.5 cm}{Зав. кафедри \dotfill Ю.В.~Крак}}\hspace {0.7cm}{\hbox to 6.5 cm{ Екзаменатор \dotfill Т.О.~Карнаух}} \end{small}}
%begin
{{\begin {small}\par
{ \center  Київський національний університет імені Тараса Шевченка \par}
{\parbox {5 cm}{Освітня програма \hfill}{\parbox {7 cm}{Прикладна математика, Системний аналіз \hfill}}\parbox {6 cm}{\hfill Семестр 2}}\par
{\parbox {5 cm} {Навчальний предмет \hfill}\parbox {7 cm}{Програмування \hfill}\parbox {6cm}{\hfill}\par}\vspace{-7pt} \end{small}}{\center Екзаменаційний білет \No \textbf 
{ 147 }\par}}
{
\begin {large}
1. Що буде виведено за виконання?
code
try:
    res = not 9>2 or True!=6.0 or 7>=2.0
    if isinstance(res, bool):
        print(f'bool;{res}')
    elif isinstance(res, int):
        print(f'int;{res}')
    elif isinstance(res, float):
        print(f'float;{res:.2f}')
    else:
        print('???')
except:
    print('ERR')
code

2. Що буде виведено за виконання?
code
a = 4
b = 9
c = 3

def f(b):
    global c
    a = 5
    b = 5
    c = 1
    return a + b + c

a = 8
b = 9
c = 6

try:
    print(f(a), a, b, c, sep=';')
except:
    print('ERR', a, b, c, sep=';')
code

3. Що буде виведено за виконання?
code
try:
    m = 8
    while m>5:
        m += -3
    print(m, end=';')
except:
    print('ERR')
code

4. Що буде виведено за виконання?
code
try:
    for c in range(7, 14, -1):
        if c>=6:
            continue
        print(c, end=';');c = 8
    else:
        print(c,end=';')
    print(c)
except:
    print('ERR')
code

5. Що буде виведено за виконання?\par (Дійсні значення, якщо наявні, в цьому завданні будуть виводитись у форматі з фіксованою крапкою та, можливо, з доволі великою кількістю десяткових розрядів. У відповіді для дійсних значень у ЦЬОМУ завданні достатньо вказати один десятковий розряд після крапки, відкинувши решту розрядів без округлення. Наприклад, замість 13.66666666666667 достатньо вказати 13.6, замість 25.23 достатньо вказати 25.2)
code
try:
    print(3, end=';')
    print(int('4'), end=';')
    print(2, end=';')
except KeyboardInterrupt:
    print(8, end=';')
except ZeroDivisionError:
    print(7, end=';')
except:
    print('ERR', end=';')
else:
    print(6, end=';')
finally:
    print(0)
code

6. Що буде виведено за виконання?
code
def f(n):
    if n>9:
        f(n//100)
    print(n%10,end='')
 
f(987654)
code

7. Що буде виведено за виконання?
code
class A:
    d = 1

    def _f1(self): print(2, end=';')
    def __f2(self): print(3, end=';')

    def main(self):
        try: print(self.d, end=';')
        except: print('ERR', end=';')

        try: self._f1()
        except: print('ERR', end=';')

        try: self.__f2()
        except: print('ERR', end=';')


class B(A):
    d = 4

    def _f1(self): print(5, end=';')
    def __f2(self): print(6, end=';')

try: B().main()
except: print('ERR', end=';')
code

8. Що буде виведено за виконання?
code
def f(arg, n):
    n = 8
    tmp = arg
    tmp[-5:-4] = [12]
    return len(tmp)>3

try:
    g = ['0','1','2','3','4']
    n = 6
    f(g, n)
    print(n, end=';')
    for el in g:
        print(el, end=';')
except:
    print('ERR')
code

9. Що буде виведено за виконання?
code
class A:
    b = 1
    a = 2

    def __init__(self):
        a = 3

class B(A):
    b = 4

    def __init__(self):
        super().__init__()
        self.c = 7

obj = B()
obj.c = obj.b + 10
B.c = 13

try: print(obj.c, end= ';')
except: print('ERR', end=';')

try: print(A.c, end= ';')
except: print('ERR', end=';')

try: print(B.c, end= ';')
except: print('ERR', end=';')
code

10. 
Для графа на рисунку з вершини 0 запущено алгоритм пошуку в глибину, що будує кістякове дерево. Списки суміжності вершин переглядаються алгоритмом за зростанням міток вершин.\par
\begin{center}
	\adjustimage{max size={0.9\linewidth}{0.9\paperheight}}
	{../10/Traversals/6.png}
\end{center}
\par Які з наведених ребер входять до складу отриманого кістякового дерева?\par
1) (0, 4)\par
2) (2, 4)\par
3) (4, 5)\par
4) (5, 6)\par
5) (0, 5)\par
6) (3, 6)\par
{\vskip 15pt} У формі вибрати номери відповідних ребер.\par
%answer:2 4 6
\end {large}}
%end
{\smallskip \begin {small} Затверджено на засіданні кафедри теоретичної кібернетики, протокол \No 12 від   19.05.2020 р.\par\smallskip\smallskip
{\parbox {6.5 cm}{Зав. кафедри \dotfill Ю.В.~Крак}}\hspace {0.7cm}{\hbox to 6.5 cm{ Екзаменатор \dotfill Т.О.~Карнаух}} \end{small}}
%begin
{{\begin {small}\par
{ \center  Київський національний університет імені Тараса Шевченка \par}
{\parbox {5 cm}{Освітня програма \hfill}{\parbox {7 cm}{Прикладна математика, Системний аналіз \hfill}}\parbox {6 cm}{\hfill Семестр 2}}\par
{\parbox {5 cm} {Навчальний предмет \hfill}\parbox {7 cm}{Програмування \hfill}\parbox {6cm}{\hfill}\par}\vspace{-7pt} \end{small}}{\center Екзаменаційний білет \No \textbf 
{ 148 }\par}}
{
\begin {large}
1. Що буде виведено за виконання?
code
try:
    res = False!="1.0"
    if isinstance(res, bool):
        print(f'bool;{res}')
    elif isinstance(res, int):
        print(f'int;{res}')
    elif isinstance(res, float):
        print(f'float;{res:.2f}')
    else:
        print('???')
except:
    print('ERR')
code

2. Що буде виведено за виконання?
code
a = 2
b = 0
c = 1

def h(b):
    global c
    a = 3
    b = 2
    c = 4
    return a + b + c

a = 4
b = 5
c = 7

try:
    print(h(a), a, b, c, sep=';')
except:
    print('ERR', a, b, c, sep=';')
code

3. Що буде виведено за виконання?
code
try:
    k = 5
    while k<=9:
        k += 2
    print(k, end=';')
except:
    print('ERR')
code

4. Що буде виведено за виконання?
code
try:
    for e in range(-9, 0, -2):
        if e<3:
            continue
        print(e, end=';')
    else:
        print(e,end=';')
    print(e)
except:
    print('ERR')
code

5. Що буде виведено за виконання?\par (Дійсні значення, якщо наявні, в цьому завданні будуть виводитись у форматі з фіксованою крапкою та, можливо, з доволі великою кількістю десяткових розрядів. У відповіді для дійсних значень у ЦЬОМУ завданні достатньо вказати один десятковий розряд після крапки, відкинувши решту розрядів без округлення. Наприклад, замість 13.66666666666667 достатньо вказати 13.6, замість 25.23 достатньо вказати 25.2)
code
try:
    print(1, end=';')
    print(9//1, end=';')
    print(0, end=';')
except TypeError:
    print(3, end=';')
except KeyboardInterrupt:
    print(2, end=';')
except:
    print('ERR', end=';')
else:
    print(7, end=';')
finally:
    print(1)
code

6. Що буде виведено за виконання?
code
def f(n):
    print(n%10, end='')
    if n>9:
        f(n//100)
 
f(987654)
code

7. Що буде виведено за виконання?
code
class A:
    a = 1

    def __h1(self): print(2, end=';')
    def _h2(self): print(3, end=';')

    def main(self):
        try: print(self.a, end=';')
        except: print('ERR', end=';')

        try: self.__h1()
        except: print('ERR', end=';')

        try: self._h2()
        except: print('ERR', end=';')


class B(A):
    a = 4

    def __h1(self): print(5, end=';')
    def _h2(self): print(6, end=';')

try: B().main()
except: print('ERR', end=';')
code

8. Що буде виведено за виконання?
code
def f(arg, n):
    arg[30] = 9
    n -= 2

try:
    n = 7
    d = {90:2, 88:8, 90:4}
    f(d, n)
    print(n, end=';')
    for x in d.items():
        print(x, sep=';', end=';')
except:
    print('ERR')
code

9. Що буде виведено за виконання?
code
class A:
    c = 1
    a = 2

    def __init__(self):
        c = 3

class B(A):
    a = 4

    def __init__(self):
        super().__init__()
        self.b = 7

obj = B()
obj.b = obj.a + 10
B.b = 13

try: print(obj.b, end= ';')
except: print('ERR', end=';')

try: print(A.b, end= ';')
except: print('ERR', end=';')

try: print(B.b, end= ';')
except: print('ERR', end=';')
code

10. 
Для графа на рисунку з вершини 6 запущено алгоритм пошуку в глибину, що будує кістякове дерево. Списки суміжності вершин переглядаються алгоритмом за зростанням міток вершин.\par
\begin{center}
	\adjustimage{max size={0.9\linewidth}{0.9\paperheight}}
	{../10/Traversals/225.png}
\end{center}
\par Які з наведених ребер входять до складу отриманого кістякового дерева?\par
1) (1, 3)\par
2) (0, 2)\par
3) (1, 2)\par
4) (2, 4)\par
5) (1, 4)\par
6) (1, 5)\par
{\vskip 15pt} У формі вибрати номери відповідних ребер.\par
%answer:2 4
\end {large}}
%end
{\smallskip \begin {small} Затверджено на засіданні кафедри теоретичної кібернетики, протокол \No 12 від   19.05.2020 р.\par\smallskip\smallskip
{\parbox {6.5 cm}{Зав. кафедри \dotfill Ю.В.~Крак}}\hspace {0.7cm}{\hbox to 6.5 cm{ Екзаменатор \dotfill Т.О.~Карнаух}} \end{small}}
%begin
{{\begin {small}\par
{ \center  Київський національний університет імені Тараса Шевченка \par}
{\parbox {5 cm}{Освітня програма \hfill}{\parbox {7 cm}{Прикладна математика, Системний аналіз \hfill}}\parbox {6 cm}{\hfill Семестр 2}}\par
{\parbox {5 cm} {Навчальний предмет \hfill}\parbox {7 cm}{Програмування \hfill}\parbox {6cm}{\hfill}\par}\vspace{-7pt} \end{small}}{\center Екзаменаційний білет \No \textbf 
{ 149 }\par}}
{
\begin {large}
1. Що буде виведено за виконання?
code
def f(n):
    print('f', end='')
    return n>2

try:
    print(f(7) or f(9))
except:
    print('ERR')
code

2. Що буде виведено за виконання?
code
a = 2
b = 4
c = 3

def f(a):
    a = 1
    b = 2
    c = 3
    return a + b + c

a = 1
b = 5
c = 6

try:
    print(f(b), a, b, c, sep=';')
except:
    print('ERR', a, b, c, sep=';')
code

3. Що буде виведено за виконання?
code
try:
    i = 4
    while i<=12:
        i += 2
    print(i, end=';')
except:
    print('ERR')
code

4. Що буде виведено за виконання?
code
try:
    for b in range(-7, 12, -3):
        if b>=5:
            continue
        print(b, end=';');b = 0
    else:
        print('end',end=';')
    print(b)
except:
    print('ERR')
code

5. Що буде виведено за виконання?\par (Дійсні значення, якщо наявні, в цьому завданні будуть виводитись у форматі з фіксованою крапкою та, можливо, з доволі великою кількістю десяткових розрядів. У відповіді для дійсних значень у ЦЬОМУ завданні достатньо вказати один десятковий розряд після крапки, відкинувши решту розрядів без округлення. Наприклад, замість 13.66666666666667 достатньо вказати 13.6, замість 25.23 достатньо вказати 25.2)
code
try:
    print(5, end=';')
    print(6j>=1j, end=';')
    print(4, end=';')
except ValueError:
    print(8, end=';')
except Exception:
    print(5, end=';')
except:
    print('ERR', end=';')
else:
    print(3, end=';')
finally:
    print(6)
code

6. Що буде виведено за виконання?
code
def f(n):
    if n>9:
        f(n//10)
    print(n%10, end='')
 
f(543216)
code

7. Що буде виведено за виконання?
code
class A:
    a = 1

    def _g1(self): print(2, end=';')
    def g2(self): print(3, end=';')

    def main(self):
        try: print(self.a, end=';')
        except: print('ERR', end=';')

        try: self._g1()
        except: print('ERR', end=';')

        try: self.g2()
        except: print('ERR', end=';')


class B(A):
    a = 4

    def _g1(self): print(5, end=';')
    def g2(self): print(6, end=';')

try: B().main()
except: print('ERR', end=';')
code

8. Що буде виведено за виконання?
code
def f(arg, n):
    arg[63] = 9
    n = 5
    arg = tuple(arg.items())

try:
    n = 1
    t = {63:2, 89:1, 90:5, 89:3}
    f(t, n)
    print(n, end=';')
    for x in t.values():
        print(x, sep=';', end=';')
except:
    print('ERR')
code

9. Що буде виведено за виконання?
code
class A:
    b = 1
    c = 2

    def __init__(self):
        b = 3

class B(A):
    b = 4

    def __init__(self):
        super().__init__()
        self.a = 7

obj = B()
obj.c = obj.b + 10
B.c = 13

try: print(obj.c, end= ';')
except: print('ERR', end=';')

try: print(A.c, end= ';')
except: print('ERR', end=';')

try: print(B.c, end= ';')
except: print('ERR', end=';')
code

10. 
Для графа на рисунку з вершини 4 запущено алгоритм пошуку в глибину, що будує кістякове дерево. Списки суміжності вершин переглядаються алгоритмом за зростанням міток вершин.\par
\begin{center}
	\adjustimage{max size={0.9\linewidth}{0.9\paperheight}}
	{../10/Traversals/39.png}
\end{center}
\par Які з наведених ребер входять до складу отриманого кістякового дерева?\par
1) (2, 6)\par
2) (1, 5)\par
3) (3, 5)\par
4) (3, 6)\par
5) (1, 4)\par
6) (1, 2)\par
{\vskip 15pt} У формі вибрати номери відповідних ребер.\par
%answer:1 3 4 5 6
\end {large}}
%end
{\smallskip \begin {small} Затверджено на засіданні кафедри теоретичної кібернетики, протокол \No 12 від   19.05.2020 р.\par\smallskip\smallskip
{\parbox {6.5 cm}{Зав. кафедри \dotfill Ю.В.~Крак}}\hspace {0.7cm}{\hbox to 6.5 cm{ Екзаменатор \dotfill Т.О.~Карнаух}} \end{small}}
%begin
{{\begin {small}\par
{ \center  Київський національний університет імені Тараса Шевченка \par}
{\parbox {5 cm}{Освітня програма \hfill}{\parbox {7 cm}{Прикладна математика, Системний аналіз \hfill}}\parbox {6 cm}{\hfill Семестр 2}}\par
{\parbox {5 cm} {Навчальний предмет \hfill}\parbox {7 cm}{Програмування \hfill}\parbox {6cm}{\hfill}\par}\vspace{-7pt} \end{small}}{\center Екзаменаційний білет \No \textbf 
{ 150 }\par}}
{
\begin {large}
1. Що буде виведено за виконання?
code
try:
    res = 9.0<=5.0 and not False==6 and 3<4
    if isinstance(res, bool):
        print(f'bool;{res}')
    elif isinstance(res, int):
        print(f'int;{res}')
    elif isinstance(res, float):
        print(f'float;{res:.2f}')
    else:
        print('???')
except:
    print('ERR')
code

2. Що буде виведено за виконання?
code
a = 8
b = 0
c = 9

def h(b):
    global c
    a = 5
    b = 4
    c = 2
    return a + b + c

a = 7
b = 3
c = 4

try:
    print(h(a), a, b, c, sep=';')
except:
    print('ERR', a, b, c, sep=';')
code

3. Що буде виведено за виконання?
code
try:
    m = 3
    while m<7:
        m += 1
    print(m, end=';')
except:
    print('ERR')
code

4. Що буде виведено за виконання?
code
try:
    for f in range(-14, 7, 2):
        if f<=4:
            break
        print(f, end=';')
    else:
        print(f,end=';')
    print(f)
except:
    print('ERR')
code

5. Що буде виведено за виконання?\par (Дійсні значення, якщо наявні, в цьому завданні будуть виводитись у форматі з фіксованою крапкою та, можливо, з доволі великою кількістю десяткових розрядів. У відповіді для дійсних значень у ЦЬОМУ завданні достатньо вказати один десятковий розряд після крапки, відкинувши решту розрядів без округлення. Наприклад, замість 13.66666666666667 достатньо вказати 13.6, замість 25.23 достатньо вказати 25.2)
code
try:
    print(2, end=';')
    print(8/0.0, end=';')
    print(4, end=';')
except ZeroDivisionError:
    print(1, end=';')
except TypeError:
    print(9, end=';')
except:
    print('ERR', end=';')
else:
    print(0, end=';')
finally:
    print(7)
code

6. Що буде виведено за виконання?
code
def f(n):
    if n>2:
        f(n//2)
    print(n%2, end='')
 
f(16)
code

7. Що буде виведено за виконання?
code
class A:
    b = 1

    def f1(self): print(2, end=';')
    def _f2(self): print(3, end=';')

    def main(self):
        try: print(self.b, end=';')
        except: print('ERR', end=';')

        try: self.f1()
        except: print('ERR', end=';')

        try: self._f2()
        except: print('ERR', end=';')


class B(A):
    b = 4

    def f1(self): print(5, end=';')
    def _f2(self): print(6, end=';')

try: B().main()
except: print('ERR', end=';')
code

8. Що буде виведено за виконання?
code
def f(arg, n):
    n = 4
    tmp = arg
    del tmp[2:5]
    return len(tmp)>3

try:
    d = [10,11,12,13,14]
    n = 3
    f(d, n)
    print(n, end=';')
    for el in d:
        print(el, end=';')
except:
    print('ERR')
code

9. Що буде виведено за виконання?
code
class A:
    c = 1
    b = 2

    def __init__(self):
        b = 3

class B(A):
    b = 4

    def __init__(self):
        super().__init__()
        self.a = 7

obj = B()
obj.c = obj.a + 10
B.c = 13

try: print(obj.c, end= ';')
except: print('ERR', end=';')

try: print(A.c, end= ';')
except: print('ERR', end=';')

try: print(B.c, end= ';')
except: print('ERR', end=';')
code

10. 
Для графа на рисунку з вершини 6 запущено алгоритм пошуку в глибину, що будує кістякове дерево. Списки суміжності вершин переглядаються алгоритмом за зростанням міток вершин.\par
\begin{center}
	\adjustimage{max size={0.9\linewidth}{0.9\paperheight}}
	{../10/Traversals/243.png}
\end{center}
\par Які з наведених ребер входять до складу отриманого кістякового дерева?\par
1) (2, 4)\par
2) (2, 5)\par
3) (0, 2)\par
4) (0, 3)\par
5) (0, 1)\par
6) (1, 3)\par
{\vskip 15pt} У формі вибрати номери відповідних ребер.\par
%answer:1 3 5 6
\end {large}}
%end
{\smallskip \begin {small} Затверджено на засіданні кафедри теоретичної кібернетики, протокол \No 12 від   19.05.2020 р.\par\smallskip\smallskip
{\parbox {6.5 cm}{Зав. кафедри \dotfill Ю.В.~Крак}}\hspace {0.7cm}{\hbox to 6.5 cm{ Екзаменатор \dotfill Т.О.~Карнаух}} \end{small}}
%begin
{{\begin {small}\par
{ \center  Київський національний університет імені Тараса Шевченка \par}
{\parbox {5 cm}{Освітня програма \hfill}{\parbox {7 cm}{Прикладна математика, Системний аналіз \hfill}}\parbox {6 cm}{\hfill Семестр 2}}\par
{\parbox {5 cm} {Навчальний предмет \hfill}\parbox {7 cm}{Програмування \hfill}\parbox {6cm}{\hfill}\par}\vspace{-7pt} \end{small}}{\center Екзаменаційний білет \No \textbf 
{ 151 }\par}}
{
\begin {large}
1. Що буде виведено за виконання?
code
try:
    res = 1**-0.5*True
    if isinstance(res, bool):
        print(f'bool;{res}')
    elif isinstance(res, int):
        print(f'int;{res}')
    elif isinstance(res, float):
        print(f'float;{res:.2f}')
    else:
        print('???')
except:
    print('ERR')
code

2. Що буде виведено за виконання?
code
a = 5
b = 7
c = 6

def g(a):
    a = 3
    b = 4
    c = 1
    return a + b + c

a = 8
b = 1
c = 0

try:
    print(g(b), a, b, c, sep=';')
except:
    print('ERR', a, b, c, sep=';')
code

3. Що буде виведено за виконання?
code
try:
    k = 14
    while k>6:
        k += -2
    print(k, end=';')
except:
    print('ERR')
code

4. Що буде виведено за виконання?
code
try:
    for c in range(-2, 8, 3):
        if c>3:
            continue
        print(c, end=';')
    else:
        print(c,end=';')
    print(c)
except:
    print('ERR')
code

5. Що буде виведено за виконання?\par (Дійсні значення, якщо наявні, в цьому завданні будуть виводитись у форматі з фіксованою крапкою та, можливо, з доволі великою кількістю десяткових розрядів. У відповіді для дійсних значень у ЦЬОМУ завданні достатньо вказати один десятковий розряд після крапки, відкинувши решту розрядів без округлення. Наприклад, замість 13.66666666666667 достатньо вказати 13.6, замість 25.23 достатньо вказати 25.2)
code
try:
    print(6, end=';')
    print(int('a7'), end=';')
    print(4, end=';')
except ZeroDivisionError:
    print(2, end=';')
except KeyboardInterrupt:
    print(0, end=';')
except:
    print('ERR', end=';')
else:
    print(3, end=';')
finally:
    print(1)
code

6. Що буде виведено за виконання?
code
def f(n):
    if n<=9:
        print(n%10,end='')
    else:
        f(n//10)

f(543216)
code

7. Що буде виведено за виконання?
code
class A:
    c = 1

    def _h1(self): print(2, end=';')
    def __h2(self): print(3, end=';')

    def main(self):
        try: print(self.c, end=';')
        except: print('ERR', end=';')

        try: self._h1()
        except: print('ERR', end=';')

        try: self.__h2()
        except: print('ERR', end=';')


class B(A):
    c = 4

    def _h1(self): print(5, end=';')
    def __h2(self): print(6, end=';')

try: B().main()
except: print('ERR', end=';')
code

8. Що буде виведено за виконання?
code
def f(arg, n):
    n = 9
    tmp = arg
    tmp[1:7] = ()
    return len(tmp)>3

try:
    c = ['0','1','2','3','4']
    n = 1
    f(c, n)
    print(n, end=';')
    for el in c:
        print(el, end=';')
except:
    print('ERR')
code

9. Що буде виведено за виконання?
code
class A:
    c = 1
    a = 2

    def __init__(self):
        c = 3

class B(A):
    c = 4

    def __init__(self):
        super().__init__()
        self.b = 7

obj = B()
obj.a = obj.c + 10
B.a = 13

try: print(obj.a, end= ';')
except: print('ERR', end=';')

try: print(A.a, end= ';')
except: print('ERR', end=';')

try: print(B.a, end= ';')
except: print('ERR', end=';')
code

10. 
Для графа на рисунку з вершини 1 запущено алгоритм пошуку в глибину, що будує кістякове дерево. Списки суміжності вершин переглядаються алгоритмом за зростанням міток вершин.\par
\begin{center}
	\adjustimage{max size={0.9\linewidth}{0.9\paperheight}}
	{../10/Traversals/76.png}
\end{center}
\par Які з наведених ребер входять до складу отриманого кістякового дерева?\par
1) (3, 5)\par
2) (2, 4)\par
3) (3, 6)\par
4) (0, 3)\par
5) (2, 3)\par
6) (1, 4)\par
{\vskip 15pt} У формі вибрати номери відповідних ребер.\par
%answer:1 3 4 5 6
\end {large}}
%end
{\smallskip \begin {small} Затверджено на засіданні кафедри теоретичної кібернетики, протокол \No 12 від   19.05.2020 р.\par\smallskip\smallskip
{\parbox {6.5 cm}{Зав. кафедри \dotfill Ю.В.~Крак}}\hspace {0.7cm}{\hbox to 6.5 cm{ Екзаменатор \dotfill Т.О.~Карнаух}} \end{small}}
%begin
{{\begin {small}\par
{ \center  Київський національний університет імені Тараса Шевченка \par}
{\parbox {5 cm}{Освітня програма \hfill}{\parbox {7 cm}{Прикладна математика, Системний аналіз \hfill}}\parbox {6 cm}{\hfill Семестр 2}}\par
{\parbox {5 cm} {Навчальний предмет \hfill}\parbox {7 cm}{Програмування \hfill}\parbox {6cm}{\hfill}\par}\vspace{-7pt} \end{small}}{\center Екзаменаційний білет \No \textbf 
{ 152 }\par}}
{
\begin {large}
1. Що буде виведено за виконання?
code
try:
    res = 1**-2-1
    if isinstance(res, bool):
        print(f'bool;{res}')
    elif isinstance(res, int):
        print(f'int;{res}')
    elif isinstance(res, float):
        print(f'float;{res:.2f}')
    else:
        print('???')
except:
    print('ERR')
code

2. Що буде виведено за виконання?
code
a = 2
b = 4
c = 7

def f(a):
    global c
    a = 2
    b -= 4
    c = 2
    return a + b + c

a = 0
b = 9
c = 8

try:
    print(f(b), a, b, c, sep=';')
except:
    print('ERR', a, b, c, sep=';')
code

3. Що буде виведено за виконання?
code
try:
    j = 7
    while j<=6:
        j += 2
    print(j, end=';')
except:
    print('ERR')
code

4. Що буде виведено за виконання?
code
try:
    for a in range(9, -10, -1):
        if a<8:
            continue
        print(a, end=';')
    else:
        print('end',end=';')
    print(a)
except:
    print('ERR')
code

5. Що буде виведено за виконання?\par (Дійсні значення, якщо наявні, в цьому завданні будуть виводитись у форматі з фіксованою крапкою та, можливо, з доволі великою кількістю десяткових розрядів. У відповіді для дійсних значень у ЦЬОМУ завданні достатньо вказати один десятковий розряд після крапки, відкинувши решту розрядів без округлення. Наприклад, замість 13.66666666666667 достатньо вказати 13.6, замість 25.23 достатньо вказати 25.2)
code
try:
    print(9, end=';')
    print(5%2, end=';')
    print(8, end=';')
except ValueError:
    print(0, end=';')
except TypeError:
    print(1, end=';')
except:
    print('ERR', end=';')
else:
    print(7, end=';')
finally:
    print(6)
code

6. Що буде виведено за виконання?
code
def f(n):
    print(n%2, end='')
    if n>2:
        f(n//2)
 
f(16)
code

7. Що буде виведено за виконання?
code
class A:
    d = 1

    def f1(self): print(2, end=';')
    def _f2(self): print(3, end=';')

    def main(self):
        try: print(self.d, end=';')
        except: print('ERR', end=';')

        try: self.f1()
        except: print('ERR', end=';')

        try: self._f2()
        except: print('ERR', end=';')


class B(A):
    d = 4

    def f1(self): print(5, end=';')
    def _f2(self): print(6, end=';')

try: B().main()
except: print('ERR', end=';')
code

8. Що буде виведено за виконання?
code
def f(arg, n):
    arg[66] = 1
    n *= -3

try:
    n = 6
    s = {36:9, 66:5, 73:2, 30:2}
    f(s, n)
    print(n, end=';')
    for x in s.keys():
        print(x, sep=';', end=';')
except:
    print('ERR')
code

9. Що буде виведено за виконання?
code
class A:
    a = 1
    c = 2

    def __init__(self):
        c = 3

class B(A):
    c = 4

    def __init__(self):
        super().__init__()
        self.b = 7

obj = B()
obj.b = obj.b + 10
B.b = 13

try: print(obj.b, end= ';')
except: print('ERR', end=';')

try: print(A.b, end= ';')
except: print('ERR', end=';')

try: print(B.b, end= ';')
except: print('ERR', end=';')
code

10. 
Для графа на рисунку з вершини 3 запущено алгоритм пошуку в глибину, що будує кістякове дерево. Списки суміжності вершин переглядаються алгоритмом за зростанням міток вершин.\par
\begin{center}
	\adjustimage{max size={0.9\linewidth}{0.9\paperheight}}
	{../10/Traversals/110.png}
\end{center}
\par Які з наведених ребер входять до складу отриманого кістякового дерева?\par
1) (0, 3)\par
2) (1, 4)\par
3) (0, 4)\par
4) (1, 2)\par
5) (0, 2)\par
6) (2, 3)\par
{\vskip 15pt} У формі вибрати номери відповідних ребер.\par
%answer:1 2 4 5
\end {large}}
%end
{\smallskip \begin {small} Затверджено на засіданні кафедри теоретичної кібернетики, протокол \No 12 від   19.05.2020 р.\par\smallskip\smallskip
{\parbox {6.5 cm}{Зав. кафедри \dotfill Ю.В.~Крак}}\hspace {0.7cm}{\hbox to 6.5 cm{ Екзаменатор \dotfill Т.О.~Карнаух}} \end{small}}
%begin
{{\begin {small}\par
{ \center  Київський національний університет імені Тараса Шевченка \par}
{\parbox {5 cm}{Освітня програма \hfill}{\parbox {7 cm}{Прикладна математика, Системний аналіз \hfill}}\parbox {6 cm}{\hfill Семестр 2}}\par
{\parbox {5 cm} {Навчальний предмет \hfill}\parbox {7 cm}{Програмування \hfill}\parbox {6cm}{\hfill}\par}\vspace{-7pt} \end{small}}{\center Екзаменаційний білет \No \textbf 
{ 153 }\par}}
{
\begin {large}
1. Що буде виведено за виконання?
code
try:
    res = 3**-2+True
    if isinstance(res, bool):
        print(f'bool;{res}')
    elif isinstance(res, int):
        print(f'int;{res}')
    elif isinstance(res, float):
        print(f'float;{res:.2f}')
    else:
        print('???')
except:
    print('ERR')
code

2. Що буде виведено за виконання?
code
a = 2
b = 9
c = 8

def f(b):
    a = 5
    b = 5
    c = 2
    return a + b + c

a = 9
b = 5
c = 4

try:
    print(f(b), a, b, c, sep=';')
except:
    print('ERR', a, b, c, sep=';')
code

3. Що буде виведено за виконання?
code
try:
    t = 11
    while t>=2:
        t += -3
    print(t, end=';')
except:
    print('ERR')
code

4. Що буде виведено за виконання?
code
try:
    for f in range(0, -11, 2):
        if f<7:
            break
        print(f, end=';');f = -8
    else:
        print(f,end=';')
    print(f)
except:
    print('ERR')
code

5. Що буде виведено за виконання?\par (Дійсні значення, якщо наявні, в цьому завданні будуть виводитись у форматі з фіксованою крапкою та, можливо, з доволі великою кількістю десяткових розрядів. У відповіді для дійсних значень у ЦЬОМУ завданні достатньо вказати один десятковий розряд після крапки, відкинувши решту розрядів без округлення. Наприклад, замість 13.66666666666667 достатньо вказати 13.6, замість 25.23 достатньо вказати 25.2)
code
try:
    print(9, end=';')
    print(5j==5j, end=';')
    print(3, end=';')
except Exception:
    print(4, end=';')
except TypeError:
    print(8, end=';')
except:
    print('ERR', end=';')
else:
    print(2, end=';')
finally:
    print(8)
code

6. Що буде виведено за виконання?
code
def f(n):
    print(n%10, end='')
    if n>9:
        f(n//10)
 
f(543216)
code

7. Що буде виведено за виконання?
code
class A:
    b = 1

    def _g1(self): print(2, end=';')
    def _g2(self): print(3, end=';')

    def main(self):
        try: print(self.b, end=';')
        except: print('ERR', end=';')

        try: self._g1()
        except: print('ERR', end=';')

        try: self._g2()
        except: print('ERR', end=';')


class B(A):
    b = 4

    def _g1(self): print(5, end=';')
    def _g2(self): print(6, end=';')

try: B().main()
except: print('ERR', end=';')
code

8. Що буде виведено за виконання?
code
def f(arg, n):
    arg[19] = 3
    n = 7
    arg = list(arg.keys())

try:
    n = 3
    s = {67:9, 53:6, 19:6, 53:5}
    f(s, n)
    print(n, end=';')
    for x in s.items():
        print(*x, sep=';', end=';')
except:
    print('ERR')
code

9. Що буде виведено за виконання?
code
class A:
    c = 1
    b = 2

    def __init__(self):
        c = 3

class B(A):
    c = 4

    def __init__(self):
        super().__init__()
        self.a = 7

obj = B()
obj.c = obj.a + 10
B.c = 13

try: print(obj.c, end= ';')
except: print('ERR', end=';')

try: print(A.c, end= ';')
except: print('ERR', end=';')

try: print(B.c, end= ';')
except: print('ERR', end=';')
code

10. 
Для графа на рисунку з вершини 1 запущено алгоритм пошуку в глибину, що будує кістякове дерево. Списки суміжності вершин переглядаються алгоритмом за зростанням міток вершин.\par
\begin{center}
	\adjustimage{max size={0.9\linewidth}{0.9\paperheight}}
	{../10/Traversals/289.png}
\end{center}
\par Які з наведених ребер входять до складу отриманого кістякового дерева?\par
1) (3, 4)\par
2) (1, 3)\par
3) (2, 4)\par
4) (2, 6)\par
5) (0, 4)\par
6) (1, 5)\par
{\vskip 15pt} У формі вибрати номери відповідних ребер.\par
%answer:2 3 5
\end {large}}
%end
{\smallskip \begin {small} Затверджено на засіданні кафедри теоретичної кібернетики, протокол \No 12 від   19.05.2020 р.\par\smallskip\smallskip
{\parbox {6.5 cm}{Зав. кафедри \dotfill Ю.В.~Крак}}\hspace {0.7cm}{\hbox to 6.5 cm{ Екзаменатор \dotfill Т.О.~Карнаух}} \end{small}}
%begin
{{\begin {small}\par
{ \center  Київський національний університет імені Тараса Шевченка \par}
{\parbox {5 cm}{Освітня програма \hfill}{\parbox {7 cm}{Прикладна математика, Системний аналіз \hfill}}\parbox {6 cm}{\hfill Семестр 2}}\par
{\parbox {5 cm} {Навчальний предмет \hfill}\parbox {7 cm}{Програмування \hfill}\parbox {6cm}{\hfill}\par}\vspace{-7pt} \end{small}}{\center Екзаменаційний білет \No \textbf 
{ 154 }\par}}
{
\begin {large}
1. Що буде виведено за виконання?
code
def f(n):
    print('f', end='')
    return n

try:
    print(f(-8) and f(0))
except:
    print('ERR')
code

2. Що буде виведено за виконання?
code
a = 3
b = 0
c = 6

def g(a):
    global c
    a = 3
    b = 4
    c = 1
    return a + b + c

a = 7
b = 1
c = 2

try:
    print(g(a), a, b, c, sep=';')
except:
    print('ERR', a, b, c, sep=';')
code

3. Що буде виведено за виконання?
code
try:
    i = 7
    while i>=4:
        i += -3
    print(i, end=';')
except:
    print('ERR')
code

4. Що буде виведено за виконання?
code
try:
    for d in range(-10, 6, -2):
        if d>4:
            continue
        print(d, end=';')
    else:
        print(d,end=';')
    print(d)
except:
    print('ERR')
code

5. Що буде виведено за виконання?\par (Дійсні значення, якщо наявні, в цьому завданні будуть виводитись у форматі з фіксованою крапкою та, можливо, з доволі великою кількістю десяткових розрядів. У відповіді для дійсних значень у ЦЬОМУ завданні достатньо вказати один десятковий розряд після крапки, відкинувши решту розрядів без округлення. Наприклад, замість 13.66666666666667 достатньо вказати 13.6, замість 25.23 достатньо вказати 25.2)
code
try:
    print(4, end=';')
    print(int('0'), end=';')
    print(7, end=';')
except Exception:
    print(3, end=';')
except ZeroDivisionError:
    print(2, end=';')
except:
    print('ERR', end=';')
else:
    print(5, end=';')
finally:
    print(1)
code

6. Що буде виведено за виконання?
code
def f(s):
    for i in range(len(s)):
        s[i] += 1
        return
 
s = [1, 2, 3]
f(s)
print(s)
code

7. Що буде виведено за виконання?
code
class A:
    d = 1

    def _h1(self): print(2, end=';')
    def h2(self): print(3, end=';')

    def main(self):
        try: print(self.d, end=';')
        except: print('ERR', end=';')

        try: self._h1()
        except: print('ERR', end=';')

        try: self.h2()
        except: print('ERR', end=';')


class B(A):
    d = 4

    def _h1(self): print(5, end=';')
    def h2(self): print(6, end=';')

try: B().main()
except: print('ERR', end=';')
code

8. Що буде виведено за виконання?
code
def f(arg, n):
    n = 7
    tmp = arg
    tmp[4:8] = 'ac'
    return len(tmp)>3

try:
    f = [10,11,12,13,14]
    n = 5
    f(f, n)
    print(n, end=';')
    for el in f:
        print(el, end=';')
except:
    print('ERR')
code

9. Що буде виведено за виконання?
code
class A:
    b = 1
    a = 2

    def __init__(self):
        a = 3

class B(A):
    a = 4

    def __init__(self):
        super().__init__()
        self.c = 7

obj = B()
obj.c = obj.b + 10
B.c = 13

try: print(obj.c, end= ';')
except: print('ERR', end=';')

try: print(A.c, end= ';')
except: print('ERR', end=';')

try: print(B.c, end= ';')
except: print('ERR', end=';')
code

10. 
Для графа на рисунку з вершини 6 запущено алгоритм пошуку в глибину, що будує кістякове дерево. Списки суміжності вершин переглядаються алгоритмом за зростанням міток вершин.\par
\begin{center}
	\adjustimage{max size={0.9\linewidth}{0.9\paperheight}}
	{../10/Traversals/227.png}
\end{center}
\par Які з наведених ребер входять до складу отриманого кістякового дерева?\par
1) (2, 3)\par
2) (1, 2)\par
3) (1, 4)\par
4) (4, 6)\par
5) (4, 5)\par
6) (1, 5)\par
{\vskip 15pt} У формі вибрати номери відповідних ребер.\par
%answer:1 5
\end {large}}
%end
{\smallskip \begin {small} Затверджено на засіданні кафедри теоретичної кібернетики, протокол \No 12 від   19.05.2020 р.\par\smallskip\smallskip
{\parbox {6.5 cm}{Зав. кафедри \dotfill Ю.В.~Крак}}\hspace {0.7cm}{\hbox to 6.5 cm{ Екзаменатор \dotfill Т.О.~Карнаух}} \end{small}}
%begin
{{\begin {small}\par
{ \center  Київський національний університет імені Тараса Шевченка \par}
{\parbox {5 cm}{Освітня програма \hfill}{\parbox {7 cm}{Прикладна математика, Системний аналіз \hfill}}\parbox {6 cm}{\hfill Семестр 2}}\par
{\parbox {5 cm} {Навчальний предмет \hfill}\parbox {7 cm}{Програмування \hfill}\parbox {6cm}{\hfill}\par}\vspace{-7pt} \end{small}}{\center Екзаменаційний білет \No \textbf 
{ 155 }\par}}
{
\begin {large}
1. Що буде виведено за виконання?
code
try:
    res = 4**-1.0*False
    if isinstance(res, bool):
        print(f'bool;{res}')
    elif isinstance(res, int):
        print(f'int;{res}')
    elif isinstance(res, float):
        print(f'float;{res:.2f}')
    else:
        print('???')
except:
    print('ERR')
code

2. Що буде виведено за виконання?
code
a = 5
b = 4
c = 2

def h(a):
    a = 3
    b *= 5
    c = 1
    return a + b + c

a = 3
b = 6
c = 1

try:
    print(h(b), a, b, c, sep=';')
except:
    print('ERR', a, b, c, sep=';')
code

3. Що буде виведено за виконання?
code
try:
    n = 0
    while n<8:
        n += 1
    print(n, end=';')
except:
    print('ERR')
code

4. Що буде виведено за виконання?
code
try:
    for b in range(2, -3, -3):
        if b<=5:
            continue
        print(b, end=';')
    else:
        print(b,end=';')
    print(b)
except:
    print('ERR')
code

5. Що буде виведено за виконання?\par (Дійсні значення, якщо наявні, в цьому завданні будуть виводитись у форматі з фіксованою крапкою та, можливо, з доволі великою кількістю десяткових розрядів. У відповіді для дійсних значень у ЦЬОМУ завданні достатньо вказати один десятковий розряд після крапки, відкинувши решту розрядів без округлення. Наприклад, замість 13.66666666666667 достатньо вказати 13.6, замість 25.23 достатньо вказати 25.2)
code
try:
    print(6, end=';')
    print(9//0.0, end=';')
    print(9, end=';')
except ValueError:
    print(7, end=';')
except KeyboardInterrupt:
    print(5, end=';')
except:
    print('ERR', end=';')
else:
    print(1, end=';')
finally:
    print(0)
code

6. Що буде виведено за виконання?
code
def f(s):
    s1 = s
    for i in range(len(s)):
        s1[i] += 1
        return s1

s = [1, 2, 3]
f(s)
print(s)
code

7. Що буде виведено за виконання?
code
class A:
    a = 1

    def __g1(self): print(2, end=';')
    def _g2(self): print(3, end=';')

    def main(self):
        try: print(self.a, end=';')
        except: print('ERR', end=';')

        try: self.__g1()
        except: print('ERR', end=';')

        try: self._g2()
        except: print('ERR', end=';')


class B(A):
    a = 4

    def __g1(self): print(5, end=';')
    def _g2(self): print(6, end=';')

try: B().main()
except: print('ERR', end=';')
code

8. Що буде виведено за виконання?
code
def f(arg, n):
    arg[60] = 2
    n = 8
    arg = list(arg.keys())

try:
    n = 2
    t = {79:1, 75:6, 75:0}
    f(t, n)
    print(n, end=';')
    for x in t.items():
        print(x, sep=';', end=';')
except:
    print('ERR')
code

9. Що буде виведено за виконання?
code
class A:
    b = 1
    c = 2

    def __init__(self):
        c = 3

class B(A):
    c = 4

    def __init__(self):
        super().__init__()
        self.a = 7

obj = B()
obj.a = obj.b + 10
B.a = 13

try: print(obj.a, end= ';')
except: print('ERR', end=';')

try: print(A.a, end= ';')
except: print('ERR', end=';')

try: print(B.a, end= ';')
except: print('ERR', end=';')
code

10. 
Для графа на рисунку з вершини 3 запущено алгоритм пошуку в глибину, що будує кістякове дерево. Списки суміжності вершин переглядаються алгоритмом за зростанням міток вершин.\par
\begin{center}
	\adjustimage{max size={0.9\linewidth}{0.9\paperheight}}
	{../10/Traversals/261.png}
\end{center}
\par Які з наведених ребер входять до складу отриманого кістякового дерева?\par
1) (1, 4)\par
2) (0, 4)\par
3) (0, 5)\par
4) (2, 4)\par
5) (2, 5)\par
6) (2, 6)\par
{\vskip 15pt} У формі вибрати номери відповідних ребер.\par
%answer:1 6
\end {large}}
%end
{\smallskip \begin {small} Затверджено на засіданні кафедри теоретичної кібернетики, протокол \No 12 від   19.05.2020 р.\par\smallskip\smallskip
{\parbox {6.5 cm}{Зав. кафедри \dotfill Ю.В.~Крак}}\hspace {0.7cm}{\hbox to 6.5 cm{ Екзаменатор \dotfill Т.О.~Карнаух}} \end{small}}
%begin
{{\begin {small}\par
{ \center  Київський національний університет імені Тараса Шевченка \par}
{\parbox {5 cm}{Освітня програма \hfill}{\parbox {7 cm}{Прикладна математика, Системний аналіз \hfill}}\parbox {6 cm}{\hfill Семестр 2}}\par
{\parbox {5 cm} {Навчальний предмет \hfill}\parbox {7 cm}{Програмування \hfill}\parbox {6cm}{\hfill}\par}\vspace{-7pt} \end{small}}{\center Екзаменаційний білет \No \textbf 
{ 156 }\par}}
{
\begin {large}
1. Що буде виведено за виконання?
code
try:
    res = 5//5*8*7
    if isinstance(res, bool):
        print(f'bool;{res}')
    elif isinstance(res, int):
        print(f'int;{res}')
    elif isinstance(res, float):
        print(f'float;{res:.2f}')
    else:
        print('???')
except:
    print('ERR')
code

2. Що буде виведено за виконання?
code
a = 0
b = 1
c = 5

def f(b):
    global c
    a = 2
    b -= 1
    c = 3
    return a + b + c

a = 4
b = 6
c = 9

try:
    print(f(b), a, b, c, sep=';')
except:
    print('ERR', a, b, c, sep=';')
code

3. Що буде виведено за виконання?
code
try:
    i = 2
    while i<11:
        i += 2
    print(i, end=';')
except:
    print('ERR')
code

4. Що буде виведено за виконання?
code
try:
    for e in range(-9, -13, -2):
        if e>=6:
            break
        print(e, end=';');e = -5
    else:
        print(e,end=';')
    print(e)
except:
    print('ERR')
code

5. Що буде виведено за виконання?\par (Дійсні значення, якщо наявні, в цьому завданні будуть виводитись у форматі з фіксованою крапкою та, можливо, з доволі великою кількістю десяткових розрядів. У відповіді для дійсних значень у ЦЬОМУ завданні достатньо вказати один десятковий розряд після крапки, відкинувши решту розрядів без округлення. Наприклад, замість 13.66666666666667 достатньо вказати 13.6, замість 25.23 достатньо вказати 25.2)
code
try:
    print(3, end=';')
    print(int('c4'), end=';')
    print(2, end=';')
except ValueError:
    print(6, end=';')
except Exception:
    print(8, end=';')
except:
    print('ERR', end=';')
else:
    print(8, end=';')
finally:
    print(6)
code

6. Що буде виведено за виконання?
code
def f(n):
    if n>9:
        f(n//10)
    else:
        print(n%10, end='')

f(543216)
code

7. Що буде виведено за виконання?
code
class A:
    c = 1

    def __f1(self): print(2, end=';')
    def _f2(self): print(3, end=';')

    def main(self):
        try: print(self.c, end=';')
        except: print('ERR', end=';')

        try: self.__f1()
        except: print('ERR', end=';')

        try: self._f2()
        except: print('ERR', end=';')


class B(A):
    c = 4

    def __f1(self): print(5, end=';')
    def _f2(self): print(6, end=';')

try: B().main()
except: print('ERR', end=';')
code

8. Що буде виведено за виконання?
code
def f(arg, n):
    arg[36] = 6
    n = 6
    arg = tuple(arg.items())

try:
    n = 0
    d = {36:9, 69:2, 36:3}
    f(d, n)
    print(n, end=';')
    for x in d.items():
        print(x, sep=';', end=';')
except:
    print('ERR')
code

9. Що буде виведено за виконання?
code
class A:
    a = 1
    b = 2

    def __init__(self):
        a = 3

class B(A):
    a = 4

    def __init__(self):
        super().__init__()
        self.c = 7

obj = B()
obj.c = obj.a + 10
B.c = 13

try: print(obj.c, end= ';')
except: print('ERR', end=';')

try: print(A.c, end= ';')
except: print('ERR', end=';')

try: print(B.c, end= ';')
except: print('ERR', end=';')
code

10. 
Для графа на рисунку з вершини 0 запущено алгоритм пошуку в глибину, що будує кістякове дерево. Списки суміжності вершин переглядаються алгоритмом за зростанням міток вершин.\par
\begin{center}
	\adjustimage{max size={0.9\linewidth}{0.9\paperheight}}
	{../10/Traversals/205.png}
\end{center}
\par Які з наведених ребер входять до складу отриманого кістякового дерева?\par
1) (5, 6)\par
2) (0, 5)\par
3) (0, 4)\par
4) (1, 6)\par
5) (3, 6)\par
6) (0, 6)\par
{\vskip 15pt} У формі вибрати номери відповідних ребер.\par
%answer:1 3 4
\end {large}}
%end
{\smallskip \begin {small} Затверджено на засіданні кафедри теоретичної кібернетики, протокол \No 12 від   19.05.2020 р.\par\smallskip\smallskip
{\parbox {6.5 cm}{Зав. кафедри \dotfill Ю.В.~Крак}}\hspace {0.7cm}{\hbox to 6.5 cm{ Екзаменатор \dotfill Т.О.~Карнаух}} \end{small}}
%begin
{{\begin {small}\par
{ \center  Київський національний університет імені Тараса Шевченка \par}
{\parbox {5 cm}{Освітня програма \hfill}{\parbox {7 cm}{Прикладна математика, Системний аналіз \hfill}}\parbox {6 cm}{\hfill Семестр 2}}\par
{\parbox {5 cm} {Навчальний предмет \hfill}\parbox {7 cm}{Програмування \hfill}\parbox {6cm}{\hfill}\par}\vspace{-7pt} \end{small}}{\center Екзаменаційний білет \No \textbf 
{ 157 }\par}}
{
\begin {large}
1. Що буде виведено за виконання?
code
try:
    res = 6j>=4.0
    if isinstance(res, bool):
        print(f'bool;{res}')
    elif isinstance(res, int):
        print(f'int;{res}')
    elif isinstance(res, float):
        print(f'float;{res:.2f}')
    else:
        print('???')
except:
    print('ERR')
code

2. Що буде виведено за виконання?
code
a = 8
b = 3
c = 2

def g(b):
    global c
    a *= 4
    b = 5
    c = 1
    return a + b + c

a = 7
b = 7
c = 0

try:
    print(g(a), a, b, c, sep=';')
except:
    print('ERR', a, b, c, sep=';')
code

3. Що буде виведено за виконання?
code
try:
    n = 14
    while n>3:
        n += -2
    print(n, end=';')
except:
    print('ERR')
code

4. Що буде виведено за виконання?
code
try:
    for c in range(-1, -12, -1):
        if c>=5:
            continue
        print(c, end=';')
        if c>=7:
            break
    else:
        print(c,end=';')
    print(c)
except:
    print('ERR')
code

5. Що буде виведено за виконання?\par (Дійсні значення, якщо наявні, в цьому завданні будуть виводитись у форматі з фіксованою крапкою та, можливо, з доволі великою кількістю десяткових розрядів. У відповіді для дійсних значень у ЦЬОМУ завданні достатньо вказати один десятковий розряд після крапки, відкинувши решту розрядів без округлення. Наприклад, замість 13.66666666666667 достатньо вказати 13.6, замість 25.23 достатньо вказати 25.2)
code
try:
    print(3, end=';')
    print(1/3, end=';')
    print(4, end=';')
except ZeroDivisionError:
    print(2, end=';')
except KeyboardInterrupt:
    print(5, end=';')
except:
    print('ERR', end=';')
else:
    print(7, end=';')
finally:
    print(9)
code

6. Що буде виведено за виконання?
code
def f(n):
    if n>9:
        return f(n//10)+1
    else:
        return 1

print(f(123456))
code

7. Що буде виведено за виконання?
code
class A:
    d = 1

    def _g1(self): print(2, end=';')
    def g2(self): print(3, end=';')

    def main(self):
        try: print(self.d, end=';')
        except: print('ERR', end=';')

        try: self._g1()
        except: print('ERR', end=';')

        try: self.g2()
        except: print('ERR', end=';')


class B(A):
    d = 4

    def _g1(self): print(5, end=';')
    def g2(self): print(6, end=';')

try: B().main()
except: print('ERR', end=';')
code

8. Що буде виведено за виконання?
code
def f(arg, n):
    arg[51] = 4
    n = 7
    arg = set(arg.values())

try:
    n = 1
    s = {56:3, 54:2, 33:8, 56:0}
    f(s, n)
    print(n, end=';')
    for x in s.values():
        print(x, sep=';', end=';')
except:
    print('ERR')
code

9. Що буде виведено за виконання?
code
class A:
    b = 1
    a = 2

    def __init__(self):
        b = 3

class B(A):
    a = 4

    def __init__(self):
        super().__init__()
        self.c = 7

obj = B()
obj.b = obj.c + 10
B.b = 13

try: print(obj.b, end= ';')
except: print('ERR', end=';')

try: print(A.b, end= ';')
except: print('ERR', end=';')

try: print(B.b, end= ';')
except: print('ERR', end=';')
code

10. 
Для графа на рисунку з вершини 5 запущено алгоритм пошуку в глибину, що будує кістякове дерево. Списки суміжності вершин переглядаються алгоритмом за зростанням міток вершин.\par
\begin{center}
	\adjustimage{max size={0.9\linewidth}{0.9\paperheight}}
	{../10/Traversals/283.png}
\end{center}
\par Які з наведених ребер входять до складу отриманого кістякового дерева?\par
1) (3, 6)\par
2) (0, 3)\par
3) (4, 5)\par
4) (4, 6)\par
5) (3, 5)\par
6) (0, 4)\par
{\vskip 15pt} У формі вибрати номери відповідних ребер.\par
%answer:2 4 5 6
\end {large}}
%end
{\smallskip \begin {small} Затверджено на засіданні кафедри теоретичної кібернетики, протокол \No 12 від   19.05.2020 р.\par\smallskip\smallskip
{\parbox {6.5 cm}{Зав. кафедри \dotfill Ю.В.~Крак}}\hspace {0.7cm}{\hbox to 6.5 cm{ Екзаменатор \dotfill Т.О.~Карнаух}} \end{small}}
%begin
{{\begin {small}\par
{ \center  Київський національний університет імені Тараса Шевченка \par}
{\parbox {5 cm}{Освітня програма \hfill}{\parbox {7 cm}{Прикладна математика, Системний аналіз \hfill}}\parbox {6 cm}{\hfill Семестр 2}}\par
{\parbox {5 cm} {Навчальний предмет \hfill}\parbox {7 cm}{Програмування \hfill}\parbox {6cm}{\hfill}\par}\vspace{-7pt} \end{small}}{\center Екзаменаційний білет \No \textbf 
{ 158 }\par}}
{
\begin {large}
1. Що буде виведено за виконання?
code
try:
    res = 8/4%4*6
    if isinstance(res, bool):
        print(f'bool;{res}')
    elif isinstance(res, int):
        print(f'int;{res}')
    elif isinstance(res, float):
        print(f'float;{res:.2f}')
    else:
        print('???')
except:
    print('ERR')
code

2. Що буде виведено за виконання?
code
a = 2
b = 4
c = 9

def h(a):
    a = 3
    b += 4
    c = 2
    return a + b + c

a = 3
b = 6
c = 1

try:
    print(h(a), a, b, c, sep=';')
except:
    print('ERR', a, b, c, sep=';')
code

3. Що буде виведено за виконання?
code
try:
    k = 8
    while k<=5:
        k += 1
    print(k, end=';')
except:
    print('ERR')
code

4. Що буде виведено за виконання?
code
try:
    for e in range(-7, 4, 3):
        if e<=4:
            continue
        print(e, end=';')
    else:
        print(e,end=';')
    print(e)
except:
    print('ERR')
code

5. Що буде виведено за виконання?\par (Дійсні значення, якщо наявні, в цьому завданні будуть виводитись у форматі з фіксованою крапкою та, можливо, з доволі великою кількістю десяткових розрядів. У відповіді для дійсних значень у ЦЬОМУ завданні достатньо вказати один десятковий розряд після крапки, відкинувши решту розрядів без округлення. Наприклад, замість 13.66666666666667 достатньо вказати 13.6, замість 25.23 достатньо вказати 25.2)
code
try:
    print(0, end=';')
    print(2//0, end=';')
    print(9, end=';')
except TypeError:
    print(1, end=';')
except KeyboardInterrupt:
    print(5, end=';')
except:
    print('ERR', end=';')
else:
    print(8, end=';')
finally:
    print(0)
code

6. Що буде виведено за виконання?
code
def f(n):
    if n>2:
        f(n//2)
    else:
        print(n%2, end='')

f(16)
code

7. Що буде виведено за виконання?
code
class A:
    b = 1

    def _f1(self): print(2, end=';')
    def __f2(self): print(3, end=';')

    def main(self):
        try: print(self.b, end=';')
        except: print('ERR', end=';')

        try: self._f1()
        except: print('ERR', end=';')

        try: self.__f2()
        except: print('ERR', end=';')


class B(A):
    b = 4

    def _f1(self): print(5, end=';')
    def __f2(self): print(6, end=';')

try: B().main()
except: print('ERR', end=';')
code

8. Що буде виведено за виконання?
code
def f(arg, n):
    arg[50] = 3
    n -= 2

try:
    n = 4
    d = {30:6, 18:1, 87:8, 18:4}
    f(d, n)
    print(n, end=';')
    for x in d.values():
        print(x, sep=';', end=';')
except:
    print('ERR')
code

9. Що буде виведено за виконання?
code
class A:
    a = 1
    c = 2

    def __init__(self):
        c = 3

class B(A):
    a = 4

    def __init__(self):
        super().__init__()
        self.b = 7

obj = B()
obj.a = obj.b + 10
B.a = 13

try: print(obj.a, end= ';')
except: print('ERR', end=';')

try: print(A.a, end= ';')
except: print('ERR', end=';')

try: print(B.a, end= ';')
except: print('ERR', end=';')
code

10. 
Для графа на рисунку з вершини 1 запущено алгоритм пошуку в глибину, що будує кістякове дерево. Списки суміжності вершин переглядаються алгоритмом за зростанням міток вершин.\par
\begin{center}
	\adjustimage{max size={0.9\linewidth}{0.9\paperheight}}
	{../10/Traversals/234.png}
\end{center}
\par Які з наведених ребер входять до складу отриманого кістякового дерева?\par
1) (1, 6)\par
2) (1, 2)\par
3) (1, 4)\par
4) (1, 5)\par
5) (2, 4)\par
6) (4, 5)\par
{\vskip 15pt} У формі вибрати номери відповідних ребер.\par
%answer:2 5 6
\end {large}}
%end
{\smallskip \begin {small} Затверджено на засіданні кафедри теоретичної кібернетики, протокол \No 12 від   19.05.2020 р.\par\smallskip\smallskip
{\parbox {6.5 cm}{Зав. кафедри \dotfill Ю.В.~Крак}}\hspace {0.7cm}{\hbox to 6.5 cm{ Екзаменатор \dotfill Т.О.~Карнаух}} \end{small}}
%begin
{{\begin {small}\par
{ \center  Київський національний університет імені Тараса Шевченка \par}
{\parbox {5 cm}{Освітня програма \hfill}{\parbox {7 cm}{Прикладна математика, Системний аналіз \hfill}}\parbox {6 cm}{\hfill Семестр 2}}\par
{\parbox {5 cm} {Навчальний предмет \hfill}\parbox {7 cm}{Програмування \hfill}\parbox {6cm}{\hfill}\par}\vspace{-7pt} \end{small}}{\center Екзаменаційний білет \No \textbf 
{ 159 }\par}}
{
\begin {large}
1. Що буде виведено за виконання?
code
try:
    res = "True"<"3"
    if isinstance(res, bool):
        print(f'bool;{res}')
    elif isinstance(res, int):
        print(f'int;{res}')
    elif isinstance(res, float):
        print(f'float;{res:.2f}')
    else:
        print('???')
except:
    print('ERR')
code

2. Що буде виведено за виконання?
code
a = 8
b = 5
c = 6

def f(a):
    a *= 5
    b = 4
    c = 1
    return a + b + c

a = 7
b = 9
c = 0

try:
    print(f(b), a, b, c, sep=';')
except:
    print('ERR', a, b, c, sep=';')
code

3. Що буде виведено за виконання?
code
try:
    j = 9
    while j>=8:
        j += -3
    print(j, end=';')
except:
    print('ERR')
code

4. Що буде виведено за виконання?
code
try:
    for a in range(-13, 2, 2):
        if a>3:
            break
        print(a, end=';');a = -6
    else:
        print('end',end=';')
    print(a)
except:
    print('ERR')
code

5. Що буде виведено за виконання?\par (Дійсні значення, якщо наявні, в цьому завданні будуть виводитись у форматі з фіксованою крапкою та, можливо, з доволі великою кількістю десяткових розрядів. У відповіді для дійсних значень у ЦЬОМУ завданні достатньо вказати один десятковий розряд після крапки, відкинувши решту розрядів без округлення. Наприклад, замість 13.66666666666667 достатньо вказати 13.6, замість 25.23 достатньо вказати 25.2)
code
try:
    print(3, end=';')
    print(7j>9j, end=';')
    print(6, end=';')
except TypeError:
    print(4, end=';')
except ZeroDivisionError:
    print(9, end=';')
except:
    print('ERR', end=';')
else:
    print(3, end=';')
finally:
    print(1)
code

6. Що буде виведено за виконання?
code
def f(n):
    if n>9: f(n//10)
    else:
        print(n%10, end='')
 
f(123456)
code

7. Що буде виведено за виконання?
code
class A:
    c = 1

    def h1(self): print(2, end=';')
    def _h2(self): print(3, end=';')

    def main(self):
        try: print(self.c, end=';')
        except: print('ERR', end=';')

        try: self.h1()
        except: print('ERR', end=';')

        try: self._h2()
        except: print('ERR', end=';')


class B(A):
    c = 4

    def h1(self): print(5, end=';')
    def _h2(self): print(6, end=';')

try: B().main()
except: print('ERR', end=';')
code

8. Що буде виведено за виконання?
code
def f(arg, n):
    n = 0
    tmp = arg
    tmp[-1:] = 'bc'
    return len(tmp)>3

try:
    b = ['0','1','2','3','4']
    n = 2
    f(b, n)
    print(n, end=';')
    for el in b:
        print(el, end=';')
except:
    print('ERR')
code

9. Що буде виведено за виконання?
code
class A:
    a = 1
    b = 2

    def __init__(self):
        b = 3

class B(A):
    a = 4

    def __init__(self):
        super().__init__()
        self.c = 7

obj = B()
obj.a = obj.c + 10
B.a = 13

try: print(obj.a, end= ';')
except: print('ERR', end=';')

try: print(A.a, end= ';')
except: print('ERR', end=';')

try: print(B.a, end= ';')
except: print('ERR', end=';')
code

10. 
Для графа на рисунку з вершини 6 запущено алгоритм пошуку в глибину, що будує кістякове дерево. Списки суміжності вершин переглядаються алгоритмом за зростанням міток вершин.\par
\begin{center}
	\adjustimage{max size={0.9\linewidth}{0.9\paperheight}}
	{../10/Traversals/35.png}
\end{center}
\par Які з наведених ребер входять до складу отриманого кістякового дерева?\par
1) (2, 6)\par
2) (0, 5)\par
3) (3, 6)\par
4) (4, 5)\par
5) (0, 2)\par
6) (1, 5)\par
{\vskip 15pt} У формі вибрати номери відповідних ребер.\par
%answer:1 2 4 5 6
\end {large}}
%end
{\smallskip \begin {small} Затверджено на засіданні кафедри теоретичної кібернетики, протокол \No 12 від   19.05.2020 р.\par\smallskip\smallskip
{\parbox {6.5 cm}{Зав. кафедри \dotfill Ю.В.~Крак}}\hspace {0.7cm}{\hbox to 6.5 cm{ Екзаменатор \dotfill Т.О.~Карнаух}} \end{small}}
%begin
{{\begin {small}\par
{ \center  Київський національний університет імені Тараса Шевченка \par}
{\parbox {5 cm}{Освітня програма \hfill}{\parbox {7 cm}{Прикладна математика, Системний аналіз \hfill}}\parbox {6 cm}{\hfill Семестр 2}}\par
{\parbox {5 cm} {Навчальний предмет \hfill}\parbox {7 cm}{Програмування \hfill}\parbox {6cm}{\hfill}\par}\vspace{-7pt} \end{small}}{\center Екзаменаційний білет \No \textbf 
{ 160 }\par}}
{
\begin {large}
1. Що буде виведено за виконання?
code
def f(n):
    print('f', end='')
    return n!=3

try:
    print(f(-7) or f(-7))
except:
    print('ERR')
code

2. Що буде виведено за виконання?
code
a = 5
b = 4
c = 3

def g(b):
    global c
    a = 2
    b += 5
    c = 3
    return a + b + c

a = 8
b = 2
c = 1

try:
    print(g(a), a, b, c, sep=';')
except:
    print('ERR', a, b, c, sep=';')
code

3. Що буде виведено за виконання?
code
try:
    j = 9
    while j<=13:
        j += 2
    print(j, end=';')
except:
    print('ERR')
code

4. Що буде виведено за виконання?
code
try:
    for b in range(10, 9, -3):
        if b<8:
            break
        print(b, end=';')
    else:
        print(b,end=';')
    print(b)
except:
    print('ERR')
code

5. Що буде виведено за виконання?\par (Дійсні значення, якщо наявні, в цьому завданні будуть виводитись у форматі з фіксованою крапкою та, можливо, з доволі великою кількістю десяткових розрядів. У відповіді для дійсних значень у ЦЬОМУ завданні достатньо вказати один десятковий розряд після крапки, відкинувши решту розрядів без округлення. Наприклад, замість 13.66666666666667 достатньо вказати 13.6, замість 25.23 достатньо вказати 25.2)
code
try:
    print(2, end=';')
    print(int('8'), end=';')
    print(7, end=';')
except ValueError:
    print(0, end=';')
except Exception:
    print(6, end=';')
except:
    print('ERR', end=';')
else:
    print(5, end=';')
finally:
    print(4)
code

6. Що буде виведено за виконання?
code
def f(n):
    if n>9:
        f(n//100)
    else:
        print(n%10,end='')

f(123456)
code

7. Що буде виведено за виконання?
code
class A:
    a = 1

    def _g1(self): print(2, end=';')
    def _g2(self): print(3, end=';')

    def main(self):
        try: print(self.a, end=';')
        except: print('ERR', end=';')

        try: self._g1()
        except: print('ERR', end=';')

        try: self._g2()
        except: print('ERR', end=';')


class B(A):
    a = 4

    def _g1(self): print(5, end=';')
    def _g2(self): print(6, end=';')

try: B().main()
except: print('ERR', end=';')
code

8. Що буде виведено за виконання?
code
def f(arg, n):
    arg[84] = 9
    n += -2

try:
    n = 5
    t = {14:1, 17:8, 14:0}
    f(t, n)
    print(n, end=';')
    for x in t :
        print(x, sep=';', end=';')
except:
    print('ERR')
code

9. Що буде виведено за виконання?
code
class A:
    c = 1
    a = 2

    def __init__(self):
        c = 3

class B(A):
    a = 4

    def __init__(self):
        super().__init__()
        self.b = 7

obj = B()
obj.a = obj.c + 10
B.a = 13

try: print(obj.a, end= ';')
except: print('ERR', end=';')

try: print(A.a, end= ';')
except: print('ERR', end=';')

try: print(B.a, end= ';')
except: print('ERR', end=';')
code

10. 
Для графа на рисунку з вершини 6 запущено алгоритм пошуку в глибину, що будує кістякове дерево. Списки суміжності вершин переглядаються алгоритмом за зростанням міток вершин.\par
\begin{center}
	\adjustimage{max size={0.9\linewidth}{0.9\paperheight}}
	{../10/Traversals/194.png}
\end{center}
\par Які з наведених ребер входять до складу отриманого кістякового дерева?\par
1) (3, 5)\par
2) (0, 5)\par
3) (0, 6)\par
4) (1, 6)\par
5) (4, 5)\par
6) (4, 6)\par
{\vskip 15pt} У формі вибрати номери відповідних ребер.\par
%answer:1 3 5
\end {large}}
%end
{\smallskip \begin {small} Затверджено на засіданні кафедри теоретичної кібернетики, протокол \No 12 від   19.05.2020 р.\par\smallskip\smallskip
{\parbox {6.5 cm}{Зав. кафедри \dotfill Ю.В.~Крак}}\hspace {0.7cm}{\hbox to 6.5 cm{ Екзаменатор \dotfill Т.О.~Карнаух}} \end{small}}
%begin
{{\begin {small}\par
{ \center  Київський національний університет імені Тараса Шевченка \par}
{\parbox {5 cm}{Освітня програма \hfill}{\parbox {7 cm}{Прикладна математика, Системний аналіз \hfill}}\parbox {6 cm}{\hfill Семестр 2}}\par
{\parbox {5 cm} {Навчальний предмет \hfill}\parbox {7 cm}{Програмування \hfill}\parbox {6cm}{\hfill}\par}\vspace{-7pt} \end{small}}{\center Екзаменаційний білет \No \textbf 
{ 161 }\par}}
{
\begin {large}
1. Що буде виведено за виконання?
code
def f(n):
    print('f', end='')
    return n

try:
    print(f(-8) and f(3))
except:
    print('ERR')
code

2. Що буде виведено за виконання?
code
a = 8
b = 7
c = 2

def h(b):
    global c
    a = 1
    b -= 2
    c = 4
    return a + b + c

a = 3
b = 9
c = 4

try:
    print(h(b), a, b, c, sep=';')
except:
    print('ERR', a, b, c, sep=';')
code

3. Що буде виведено за виконання?
code
try:
    m = 0
    while m<12:
        m += 1
    print(m, end=';')
except:
    print('ERR')
code

4. Що буде виведено за виконання?
code
try:
    for f in range(11, -8, -3):
        if f<=6:
            continue
        print(f, end=';');f = 10
    else:
        print(f,end=';')
    print(f)
except:
    print('ERR')
code

5. Що буде виведено за виконання?\par (Дійсні значення, якщо наявні, в цьому завданні будуть виводитись у форматі з фіксованою крапкою та, можливо, з доволі великою кількістю десяткових розрядів. У відповіді для дійсних значень у ЦЬОМУ завданні достатньо вказати один десятковий розряд після крапки, відкинувши решту розрядів без округлення. Наприклад, замість 13.66666666666667 достатньо вказати 13.6, замість 25.23 достатньо вказати 25.2)
code
try:
    print(6, end=';')
    print(1%1, end=';')
    print(3, end=';')
except Exception:
    print(4, end=';')
except KeyboardInterrupt:
    print(9, end=';')
except:
    print('ERR', end=';')
else:
    print(0, end=';')
finally:
    print(5)
code

6. Що буде виведено за виконання?
code
def f(s1):
    s = s1[:]
    for i in range(len(s)):
        s[i] += 1
 
s = [1, 2, 3]
f(s)
print(s)
code

7. Що буде виведено за виконання?
code
class A:
    d = 1

    def _h1(self): print(2, end=';')
    def _h2(self): print(3, end=';')

    def main(self):
        try: print(self.d, end=';')
        except: print('ERR', end=';')

        try: self._h1()
        except: print('ERR', end=';')

        try: self._h2()
        except: print('ERR', end=';')


class B(A):
    d = 4

    def _h1(self): print(5, end=';')
    def _h2(self): print(6, end=';')

try: B().main()
except: print('ERR', end=';')
code

8. Що буде виведено за виконання?
code
def f(arg, n):
    arg[49] = 0
    n = 5
    arg = tuple(arg.keys())

try:
    n = 4
    t = {22:4, 10:8, 10:8}
    f(t, n)
    print(n, end=';')
    for x in t.items():
        print(*x, sep=';', end=';')
except:
    print('ERR')
code

9. Що буде виведено за виконання?
code
class A:
    c = 1
    b = 2

    def __init__(self):
        c = 3

class B(A):
    c = 4

    def __init__(self):
        super().__init__()
        self.a = 7

obj = B()
obj.b = obj.a + 10
B.b = 13

try: print(obj.b, end= ';')
except: print('ERR', end=';')

try: print(A.b, end= ';')
except: print('ERR', end=';')

try: print(B.b, end= ';')
except: print('ERR', end=';')
code

10. 
Для графа на рисунку з вершини 0 запущено алгоритм пошуку в глибину, що будує кістякове дерево. Списки суміжності вершин переглядаються алгоритмом за зростанням міток вершин.\par
\begin{center}
	\adjustimage{max size={0.9\linewidth}{0.9\paperheight}}
	{../10/Traversals/100.png}
\end{center}
\par Які з наведених ребер входять до складу отриманого кістякового дерева?\par
1) (4, 5)\par
2) (2, 5)\par
3) (0, 1)\par
4) (3, 4)\par
5) (1, 4)\par
6) (1, 2)\par
{\vskip 15pt} У формі вибрати номери відповідних ребер.\par
%answer:1 2 3 4 6
\end {large}}
%end
{\smallskip \begin {small} Затверджено на засіданні кафедри теоретичної кібернетики, протокол \No 12 від   19.05.2020 р.\par\smallskip\smallskip
{\parbox {6.5 cm}{Зав. кафедри \dotfill Ю.В.~Крак}}\hspace {0.7cm}{\hbox to 6.5 cm{ Екзаменатор \dotfill Т.О.~Карнаух}} \end{small}}
%begin
{{\begin {small}\par
{ \center  Київський національний університет імені Тараса Шевченка \par}
{\parbox {5 cm}{Освітня програма \hfill}{\parbox {7 cm}{Прикладна математика, Системний аналіз \hfill}}\parbox {6 cm}{\hfill Семестр 2}}\par
{\parbox {5 cm} {Навчальний предмет \hfill}\parbox {7 cm}{Програмування \hfill}\parbox {6cm}{\hfill}\par}\vspace{-7pt} \end{small}}{\center Екзаменаційний білет \No \textbf 
{ 162 }\par}}
{
\begin {large}
1. Що буде виведено за виконання?
code
try:
    res = -3**-2+-1
    if isinstance(res, bool):
        print(f'bool;{res}')
    elif isinstance(res, int):
        print(f'int;{res}')
    elif isinstance(res, float):
        print(f'float;{res:.2f}')
    else:
        print('???')
except:
    print('ERR')
code

2. Що буде виведено за виконання?
code
a = 2
b = 7
c = 5

def h(a):
    a = 3
    b = 1
    c = 5
    return a + b + c

a = 4
b = 6
c = 3

try:
    print(h(b), a, b, c, sep=';')
except:
    print('ERR', a, b, c, sep=';')
code

3. Що буде виведено за виконання?
code
try:
    t = 3
    while t<6:
        t += 2
    print(t, end=';')
except:
    print('ERR')
code

4. Що буде виведено за виконання?
code
try:
    for d in range(12, 8, 3):
        if d>=7:
            continue
        print(d, end=';')
    else:
        print(d,end=';')
    print(d)
except:
    print('ERR')
code

5. Що буде виведено за виконання?\par (Дійсні значення, якщо наявні, в цьому завданні будуть виводитись у форматі з фіксованою крапкою та, можливо, з доволі великою кількістю десяткових розрядів. У відповіді для дійсних значень у ЦЬОМУ завданні достатньо вказати один десятковий розряд після крапки, відкинувши решту розрядів без округлення. Наприклад, замість 13.66666666666667 достатньо вказати 13.6, замість 25.23 достатньо вказати 25.2)
code
try:
    print(2, end=';')
    print(7j!=1j, end=';')
    print(8, end=';')
except ValueError:
    print(2, end=';')
except ZeroDivisionError:
    print(5, end=';')
except:
    print('ERR', end=';')
else:
    print(1, end=';')
finally:
    print(3)
code

6. Що буде виведено за виконання?
code
def f(n):
    if n>9:
        f(n//100)
        print(n%10, end='')
 
f(123456)
code

7. Що буде виведено за виконання?
code
class A:
    a = 1

    def __f1(self): print(2, end=';')
    def _f2(self): print(3, end=';')

    def main(self):
        try: print(self.a, end=';')
        except: print('ERR', end=';')

        try: self.__f1()
        except: print('ERR', end=';')

        try: self._f2()
        except: print('ERR', end=';')


class B(A):
    a = 4

    def __f1(self): print(5, end=';')
    def _f2(self): print(6, end=';')

try: B().main()
except: print('ERR', end=';')
code

8. Що буде виведено за виконання?
code
def f(arg, n):
    arg[39] = 7
    n = 8
    arg = list(arg.items())

try:
    n = 0
    d = {49:9, 42:4, 68:7, 68:2}
    f(d, n)
    print(n, end=';')
    for x in d.keys():
        print(x, sep=';', end=';')
except:
    print('ERR')
code

9. Що буде виведено за виконання?
code
class A:
    b = 1
    c = 2

    def __init__(self):
        c = 3

class B(A):
    c = 4

    def __init__(self):
        super().__init__()
        self.a = 7

obj = B()
obj.b = obj.c + 10
B.b = 13

try: print(obj.b, end= ';')
except: print('ERR', end=';')

try: print(A.b, end= ';')
except: print('ERR', end=';')

try: print(B.b, end= ';')
except: print('ERR', end=';')
code

10. 
Для графа на рисунку з вершини 4 запущено алгоритм пошуку в глибину, що будує кістякове дерево. Списки суміжності вершин переглядаються алгоритмом за зростанням міток вершин.\par
\begin{center}
	\adjustimage{max size={0.9\linewidth}{0.9\paperheight}}
	{../10/Traversals/8.png}
\end{center}
\par Які з наведених ребер входять до складу отриманого кістякового дерева?\par
1) (2, 5)\par
2) (1, 3)\par
3) (0, 6)\par
4) (2, 6)\par
5) (4, 5)\par
6) (0, 3)\par
{\vskip 15pt} У формі вибрати номери відповідних ребер.\par
%answer:2 3 4 5
\end {large}}
%end
{\smallskip \begin {small} Затверджено на засіданні кафедри теоретичної кібернетики, протокол \No 12 від   19.05.2020 р.\par\smallskip\smallskip
{\parbox {6.5 cm}{Зав. кафедри \dotfill Ю.В.~Крак}}\hspace {0.7cm}{\hbox to 6.5 cm{ Екзаменатор \dotfill Т.О.~Карнаух}} \end{small}}
%begin
{{\begin {small}\par
{ \center  Київський національний університет імені Тараса Шевченка \par}
{\parbox {5 cm}{Освітня програма \hfill}{\parbox {7 cm}{Прикладна математика, Системний аналіз \hfill}}\parbox {6 cm}{\hfill Семестр 2}}\par
{\parbox {5 cm} {Навчальний предмет \hfill}\parbox {7 cm}{Програмування \hfill}\parbox {6cm}{\hfill}\par}\vspace{-7pt} \end{small}}{\center Екзаменаційний білет \No \textbf 
{ 163 }\par}}
{
\begin {large}
1. Що буде виведено за виконання?
code
try:
    res = "7.0j"<=False
    if isinstance(res, bool):
        print(f'bool;{res}')
    elif isinstance(res, int):
        print(f'int;{res}')
    elif isinstance(res, float):
        print(f'float;{res:.2f}')
    else:
        print('???')
except:
    print('ERR')
code

2. Що буде виведено за виконання?
code
a = 9
b = 8
c = 0

def g(b):
    global c
    a = 2
    b = 4
    c = 3
    return a + b + c

a = 1
b = 6
c = 8

try:
    print(g(a), a, b, c, sep=';')
except:
    print('ERR', a, b, c, sep=';')
code

3. Що буде виведено за виконання?
code
try:
    m = 13
    while m>7:
        m += -2
    print(m, end=';')
except:
    print('ERR')
code

4. Що буде виведено за виконання?
code
try:
    for d in range(-2, -11, 2):
        if d>8:
            continue
        print(d, end=';')
    else:
        print(d,end=';')
    print(d)
except:
    print('ERR')
code

5. Що буде виведено за виконання?\par (Дійсні значення, якщо наявні, в цьому завданні будуть виводитись у форматі з фіксованою крапкою та, можливо, з доволі великою кількістю десяткових розрядів. У відповіді для дійсних значень у ЦЬОМУ завданні достатньо вказати один десятковий розряд після крапки, відкинувши решту розрядів без округлення. Наприклад, замість 13.66666666666667 достатньо вказати 13.6, замість 25.23 достатньо вказати 25.2)
code
try:
    print(7, end=';')
    print(0/0, end=';')
    print(8, end=';')
except TypeError:
    print(4, end=';')
except ZeroDivisionError:
    print(9, end=';')
except:
    print('ERR', end=';')
else:
    print(6, end=';')
finally:
    print(9)
code

6. Що буде виведено за виконання?
code
def f(n):
    if n>9:
        f(n//100)
    else:
        print(n%10,end='')

f(987654)
code

7. Що буде виведено за виконання?
code
class A:
    c = 1

    def h1(self): print(2, end=';')
    def _h2(self): print(3, end=';')

    def main(self):
        try: print(self.c, end=';')
        except: print('ERR', end=';')

        try: self.h1()
        except: print('ERR', end=';')

        try: self._h2()
        except: print('ERR', end=';')


class B(A):
    c = 4

    def h1(self): print(5, end=';')
    def _h2(self): print(6, end=';')

try: B().main()
except: print('ERR', end=';')
code

8. Що буде виведено за виконання?
code
def f(arg, n):
    arg[19] = 8
    n = 6
    arg = set(arg.values())

try:
    n = 2
    s = {61:9, 56:5, 30:9}
    f(s, n)
    print(n, end=';')
    for x in s.keys():
        print(x, sep=';', end=';')
except:
    print('ERR')
code

9. Що буде виведено за виконання?
code
class A:
    a = 1
    b = 2

    def __init__(self):
        b = 3

class B(A):
    b = 4

    def __init__(self):
        super().__init__()
        self.c = 7

obj = B()
obj.a = obj.c + 10
B.a = 13

try: print(obj.a, end= ';')
except: print('ERR', end=';')

try: print(A.a, end= ';')
except: print('ERR', end=';')

try: print(B.a, end= ';')
except: print('ERR', end=';')
code

10. 
Для графа на рисунку з вершини 1 запущено алгоритм пошуку в глибину, що будує кістякове дерево. Списки суміжності вершин переглядаються алгоритмом за зростанням міток вершин.\par
\begin{center}
	\adjustimage{max size={0.9\linewidth}{0.9\paperheight}}
	{../10/Traversals/115.png}
\end{center}
\par Які з наведених ребер входять до складу отриманого кістякового дерева?\par
1) (3, 6)\par
2) (1, 2)\par
3) (4, 5)\par
4) (0, 5)\par
5) (5, 6)\par
6) (1, 4)\par
{\vskip 15pt} У формі вибрати номери відповідних ребер.\par
%answer:2 4 5
\end {large}}
%end
{\smallskip \begin {small} Затверджено на засіданні кафедри теоретичної кібернетики, протокол \No 12 від   19.05.2020 р.\par\smallskip\smallskip
{\parbox {6.5 cm}{Зав. кафедри \dotfill Ю.В.~Крак}}\hspace {0.7cm}{\hbox to 6.5 cm{ Екзаменатор \dotfill Т.О.~Карнаух}} \end{small}}
%begin
{{\begin {small}\par
{ \center  Київський національний університет імені Тараса Шевченка \par}
{\parbox {5 cm}{Освітня програма \hfill}{\parbox {7 cm}{Прикладна математика, Системний аналіз \hfill}}\parbox {6 cm}{\hfill Семестр 2}}\par
{\parbox {5 cm} {Навчальний предмет \hfill}\parbox {7 cm}{Програмування \hfill}\parbox {6cm}{\hfill}\par}\vspace{-7pt} \end{small}}{\center Екзаменаційний білет \No \textbf 
{ 164 }\par}}
{
\begin {large}
1. Що буде виведено за виконання?
code
def f(n):
    print('f', end='')
    return n

try:
    print(f(2) and f(-6))
except:
    print('ERR')
code

2. Що буде виведено за виконання?
code
a = 0
b = 5
c = 6

def f(a):
    a = 5
    b *= 3
    c = 4
    return a + b + c

a = 1
b = 8
c = 6

try:
    print(f(b), a, b, c, sep=';')
except:
    print('ERR', a, b, c, sep=';')
code

3. Що буде виведено за виконання?
code
try:
    t = 5
    while t>5:
        t += -2
    print(t, end=';')
except:
    print('ERR')
code

4. Що буде виведено за виконання?
code
try:
    for c in range(1, 14, -1):
        if c<6:
            continue
        print(c, end=';');c = 4
    else:
        print(c,end=';')
    print(c)
except:
    print('ERR')
code

5. Що буде виведено за виконання?\par (Дійсні значення, якщо наявні, в цьому завданні будуть виводитись у форматі з фіксованою крапкою та, можливо, з доволі великою кількістю десяткових розрядів. У відповіді для дійсних значень у ЦЬОМУ завданні достатньо вказати один десятковий розряд після крапки, відкинувши решту розрядів без округлення. Наприклад, замість 13.66666666666667 достатньо вказати 13.6, замість 25.23 достатньо вказати 25.2)
code
try:
    print(0, end=';')
    print(int('c4'), end=';')
    print(8, end=';')
except TypeError:
    print(1, end=';')
except Exception:
    print(2, end=';')
except:
    print('ERR', end=';')
else:
    print(3, end=';')
finally:
    print(7)
code

6. Що буде виведено за виконання?
code
def f(n):
    if n>9:
        f(n//100)
    print(n%10, end='')
 
f(123456)
code

7. Що буде виведено за виконання?
code
class A:
    b = 1

    def _f1(self): print(2, end=';')
    def f2(self): print(3, end=';')

    def main(self):
        try: print(self.b, end=';')
        except: print('ERR', end=';')

        try: self._f1()
        except: print('ERR', end=';')

        try: self.f2()
        except: print('ERR', end=';')


class B(A):
    b = 4

    def _f1(self): print(5, end=';')
    def f2(self): print(6, end=';')

try: B().main()
except: print('ERR', end=';')
code

8. Що буде виведено за виконання?
code
def f(arg, n):
    arg[49] = 3
    n *= 3

try:
    n = 3
    t = {49:3, 72:5, 40:3, 40:9}
    f(t, n)
    print(n, end=';')
    for x in t.keys():
        print(x, sep=';', end=';')
except:
    print('ERR')
code

9. Що буде виведено за виконання?
code
class A:
    b = 1
    a = 2

    def __init__(self):
        b = 3

class B(A):
    b = 4

    def __init__(self):
        super().__init__()
        self.c = 7

obj = B()
obj.b = obj.c + 10
B.b = 13

try: print(obj.b, end= ';')
except: print('ERR', end=';')

try: print(A.b, end= ';')
except: print('ERR', end=';')

try: print(B.b, end= ';')
except: print('ERR', end=';')
code

10. 
Для графа на рисунку з вершини 5 запущено алгоритм пошуку в глибину, що будує кістякове дерево. Списки суміжності вершин переглядаються алгоритмом за зростанням міток вершин.\par
\begin{center}
	\adjustimage{max size={0.9\linewidth}{0.9\paperheight}}
	{../10/Traversals/81.png}
\end{center}
\par Які з наведених ребер входять до складу отриманого кістякового дерева?\par
1) (0, 1)\par
2) (2, 5)\par
3) (1, 2)\par
4) (0, 6)\par
5) (5, 6)\par
6) (3, 5)\par
{\vskip 15pt} У формі вибрати номери відповідних ребер.\par
%answer:1 2
\end {large}}
%end
{\smallskip \begin {small} Затверджено на засіданні кафедри теоретичної кібернетики, протокол \No 12 від   19.05.2020 р.\par\smallskip\smallskip
{\parbox {6.5 cm}{Зав. кафедри \dotfill Ю.В.~Крак}}\hspace {0.7cm}{\hbox to 6.5 cm{ Екзаменатор \dotfill Т.О.~Карнаух}} \end{small}}
%begin
{{\begin {small}\par
{ \center  Київський національний університет імені Тараса Шевченка \par}
{\parbox {5 cm}{Освітня програма \hfill}{\parbox {7 cm}{Прикладна математика, Системний аналіз \hfill}}\parbox {6 cm}{\hfill Семестр 2}}\par
{\parbox {5 cm} {Навчальний предмет \hfill}\parbox {7 cm}{Програмування \hfill}\parbox {6cm}{\hfill}\par}\vspace{-7pt} \end{small}}{\center Екзаменаційний білет \No \textbf 
{ 165 }\par}}
{
\begin {large}
1. Що буде виведено за виконання?
code
try:
    res = 2.0j>True
    if isinstance(res, bool):
        print(f'bool;{res}')
    elif isinstance(res, int):
        print(f'int;{res}')
    elif isinstance(res, float):
        print(f'float;{res:.2f}')
    else:
        print('???')
except:
    print('ERR')
code

2. Що буде виведено за виконання?
code
a = 7
b = 5
c = 1

def h(b):
    global c
    a = 4
    b = 1
    c = 2
    return a + b + c

a = 4
b = 9
c = 3

try:
    print(h(a), a, b, c, sep=';')
except:
    print('ERR', a, b, c, sep=';')
code

3. Що буде виведено за виконання?
code
try:
    n = 8
    while n>=1:
        n += -3
    print(n, end=';')
except:
    print('ERR')
code

4. Що буде виведено за виконання?
code
try:
    for b in range(-12, 11, -2):
        if b>4:
            break
        print(b, end=';')
    else:
        print(b,end=';')
    print(b)
except:
    print('ERR')
code

5. Що буде виведено за виконання?\par (Дійсні значення, якщо наявні, в цьому завданні будуть виводитись у форматі з фіксованою крапкою та, можливо, з доволі великою кількістю десяткових розрядів. У відповіді для дійсних значень у ЦЬОМУ завданні достатньо вказати один десятковий розряд після крапки, відкинувши решту розрядів без округлення. Наприклад, замість 13.66666666666667 достатньо вказати 13.6, замість 25.23 достатньо вказати 25.2)
code
try:
    print(6, end=';')
    print(int('5'), end=';')
    print(1, end=';')
except ValueError:
    print(9, end=';')
except KeyboardInterrupt:
    print(6, end=';')
except:
    print('ERR', end=';')
else:
    print(5, end=';')
finally:
    print(7)
code

6. Що буде виведено за виконання?
code
def f(n):
    print(n%10, end='')
    if n>9:
        f(n//10)
 
f(123456)
code

7. Що буде виведено за виконання?
code
class A:
    d = 1

    def _g1(self): print(2, end=';')
    def __g2(self): print(3, end=';')

    def main(self):
        try: print(self.d, end=';')
        except: print('ERR', end=';')

        try: self._g1()
        except: print('ERR', end=';')

        try: self.__g2()
        except: print('ERR', end=';')


class B(A):
    d = 4

    def _g1(self): print(5, end=';')
    def __g2(self): print(6, end=';')

try: B().main()
except: print('ERR', end=';')
code

8. Що буде виведено за виконання?
code
def f(arg, n):
    n = 2
    tmp = arg
    del tmp[:-2]
    return len(tmp)>3

try:
    a = [10,11,12,13,14]
    n = 5
    f(a, n)
    print(n, end=';')
    for el in a:
        print(el, end=';')
except:
    print('ERR')
code

9. Що буде виведено за виконання?
code
class A:
    b = 1
    c = 2

    def __init__(self):
        c = 3

class B(A):
    c = 4

    def __init__(self):
        super().__init__()
        self.a = 7

obj = B()
obj.b = obj.a + 10
B.b = 13

try: print(obj.b, end= ';')
except: print('ERR', end=';')

try: print(A.b, end= ';')
except: print('ERR', end=';')

try: print(B.b, end= ';')
except: print('ERR', end=';')
code

10. 
Для графа на рисунку з вершини 3 запущено алгоритм пошуку в глибину, що будує кістякове дерево. Списки суміжності вершин переглядаються алгоритмом за зростанням міток вершин.\par
\begin{center}
	\adjustimage{max size={0.9\linewidth}{0.9\paperheight}}
	{../10/Traversals/120.png}
\end{center}
\par Які з наведених ребер входять до складу отриманого кістякового дерева?\par
1) (0, 1)\par
2) (0, 3)\par
3) (1, 6)\par
4) (2, 5)\par
5) (1, 3)\par
6) (4, 5)\par
{\vskip 15pt} У формі вибрати номери відповідних ребер.\par
%answer:1 2 6
\end {large}}
%end
{\smallskip \begin {small} Затверджено на засіданні кафедри теоретичної кібернетики, протокол \No 12 від   19.05.2020 р.\par\smallskip\smallskip
{\parbox {6.5 cm}{Зав. кафедри \dotfill Ю.В.~Крак}}\hspace {0.7cm}{\hbox to 6.5 cm{ Екзаменатор \dotfill Т.О.~Карнаух}} \end{small}}
%begin
{{\begin {small}\par
{ \center  Київський національний університет імені Тараса Шевченка \par}
{\parbox {5 cm}{Освітня програма \hfill}{\parbox {7 cm}{Прикладна математика, Системний аналіз \hfill}}\parbox {6 cm}{\hfill Семестр 2}}\par
{\parbox {5 cm} {Навчальний предмет \hfill}\parbox {7 cm}{Програмування \hfill}\parbox {6cm}{\hfill}\par}\vspace{-7pt} \end{small}}{\center Екзаменаційний білет \No \textbf 
{ 166 }\par}}
{
\begin {large}
1. Що буде виведено за виконання?
code
try:
    res = "6.0"!="False"
    if isinstance(res, bool):
        print(f'bool;{res}')
    elif isinstance(res, int):
        print(f'int;{res}')
    elif isinstance(res, float):
        print(f'float;{res:.2f}')
    else:
        print('???')
except:
    print('ERR')
code

2. Що буде виведено за виконання?
code
a = 0
b = 2
c = 5

def f(a):
    a = 5
    b = 1
    c = 2
    return a + b + c

a = 8
b = 6
c = 9

try:
    print(f(b), a, b, c, sep=';')
except:
    print('ERR', a, b, c, sep=';')
code

3. Що буде виведено за виконання?
code
try:
    n = 7
    while n<=10:
        n += 1
    print(n, end=';')
except:
    print('ERR')
code

4. Що буде виведено за виконання?
code
try:
    for f in range(0, -1, 2):
        if f>=7:
            continue
        print(f, end=';');f = 2
    else:
        print(f,end=';')
    print(f)
except:
    print('ERR')
code

5. Що буде виведено за виконання?\par (Дійсні значення, якщо наявні, в цьому завданні будуть виводитись у форматі з фіксованою крапкою та, можливо, з доволі великою кількістю десяткових розрядів. У відповіді для дійсних значень у ЦЬОМУ завданні достатньо вказати один десятковий розряд після крапки, відкинувши решту розрядів без округлення. Наприклад, замість 13.66666666666667 достатньо вказати 13.6, замість 25.23 достатньо вказати 25.2)
code
try:
    print(3, end=';')
    print(0//0.0, end=';')
    print(8, end=';')
except ValueError:
    print(2, end=';')
except KeyboardInterrupt:
    print(4, end=';')
except:
    print('ERR', end=';')
else:
    print(4, end=';')
finally:
    print(1)
code

6. Що буде виведено за виконання?
code
def f(n):
    if n<=9:
        print(n%10, end='')
    else:
        f(n//10)

f(987651)
code

7. Що буде виведено за виконання?
code
class A:
    a = 1

    def _h1(self): print(2, end=';')
    def __h2(self): print(3, end=';')

    def main(self):
        try: print(self.a, end=';')
        except: print('ERR', end=';')

        try: self._h1()
        except: print('ERR', end=';')

        try: self.__h2()
        except: print('ERR', end=';')


class B(A):
    a = 4

    def _h1(self): print(5, end=';')
    def __h2(self): print(6, end=';')

try: B().main()
except: print('ERR', end=';')
code

8. Що буде виведено за виконання?
code
def f(arg, n):
    arg[83] = 3
    n = 9
    arg = tuple(arg.items())

try:
    n = 3
    s = {39:3, 67:5, 39:5}
    f(s, n)
    print(n, end=';')
    for x in s.values():
        print(x, sep=';', end=';')
except:
    print('ERR')
code

9. Що буде виведено за виконання?
code
class A:
    c = 1
    b = 2

    def __init__(self):
        c = 3

class B(A):
    c = 4

    def __init__(self):
        super().__init__()
        self.a = 7

obj = B()
obj.a = obj.c + 10
B.a = 13

try: print(obj.a, end= ';')
except: print('ERR', end=';')

try: print(A.a, end= ';')
except: print('ERR', end=';')

try: print(B.a, end= ';')
except: print('ERR', end=';')
code

10. 
Для графа на рисунку з вершини 6 запущено алгоритм пошуку в глибину, що будує кістякове дерево. Списки суміжності вершин переглядаються алгоритмом за зростанням міток вершин.\par
\begin{center}
	\adjustimage{max size={0.9\linewidth}{0.9\paperheight}}
	{../10/Traversals/256.png}
\end{center}
\par Які з наведених ребер входять до складу отриманого кістякового дерева?\par
1) (0, 3)\par
2) (3, 4)\par
3) (1, 6)\par
4) (1, 3)\par
5) (4, 6)\par
6) (0, 6)\par
{\vskip 15pt} У формі вибрати номери відповідних ребер.\par
%answer:1 4 6
\end {large}}
%end
{\smallskip \begin {small} Затверджено на засіданні кафедри теоретичної кібернетики, протокол \No 12 від   19.05.2020 р.\par\smallskip\smallskip
{\parbox {6.5 cm}{Зав. кафедри \dotfill Ю.В.~Крак}}\hspace {0.7cm}{\hbox to 6.5 cm{ Екзаменатор \dotfill Т.О.~Карнаух}} \end{small}}
%begin
{{\begin {small}\par
{ \center  Київський національний університет імені Тараса Шевченка \par}
{\parbox {5 cm}{Освітня програма \hfill}{\parbox {7 cm}{Прикладна математика, Системний аналіз \hfill}}\parbox {6 cm}{\hfill Семестр 2}}\par
{\parbox {5 cm} {Навчальний предмет \hfill}\parbox {7 cm}{Програмування \hfill}\parbox {6cm}{\hfill}\par}\vspace{-7pt} \end{small}}{\center Екзаменаційний білет \No \textbf 
{ 167 }\par}}
{
\begin {large}
1. Що буде виведено за виконання?
code
try:
    res = not True!=4.0 or 5==8.0 and 8<=4
    if isinstance(res, bool):
        print(f'bool;{res}')
    elif isinstance(res, int):
        print(f'int;{res}')
    elif isinstance(res, float):
        print(f'float;{res:.2f}')
    else:
        print('???')
except:
    print('ERR')
code

2. Що буде виведено за виконання?
code
a = 7
b = 1
c = 3

def h(a):
    a += 3
    b = 1
    c = 5
    return a + b + c

a = 2
b = 5
c = 9

try:
    print(h(a), a, b, c, sep=';')
except:
    print('ERR', a, b, c, sep=';')
code

3. Що буде виведено за виконання?
code
try:
    k = 10
    while k>7:
        k += -3
    print(k, end=';')
except:
    print('ERR')
code

4. Що буде виведено за виконання?
code
try:
    for e in range(-10, -2, -3):
        if e<3:
            break
        print(e, end=';')
    else:
        print('end',end=';')
    print(e)
except:
    print('ERR')
code

5. Що буде виведено за виконання?\par (Дійсні значення, якщо наявні, в цьому завданні будуть виводитись у форматі з фіксованою крапкою та, можливо, з доволі великою кількістю десяткових розрядів. У відповіді для дійсних значень у ЦЬОМУ завданні достатньо вказати один десятковий розряд після крапки, відкинувши решту розрядів без округлення. Наприклад, замість 13.66666666666667 достатньо вказати 13.6, замість 25.23 достатньо вказати 25.2)
code
try:
    print(7, end=';')
    print(3j>=4j, end=';')
    print(9, end=';')
except TypeError:
    print(6, end=';')
except ZeroDivisionError:
    print(5, end=';')
except:
    print('ERR', end=';')
else:
    print(0, end=';')
finally:
    print(8)
code

6. Що буде виведено за виконання?
code
def f(n):
    if n<=9:
        print(n%10, end='')
    else:
        f(n//10)

f(123456)
code

7. Що буде виведено за виконання?
code
class A:
    b = 1

    def f1(self): print(2, end=';')
    def _f2(self): print(3, end=';')

    def main(self):
        try: print(self.b, end=';')
        except: print('ERR', end=';')

        try: self.f1()
        except: print('ERR', end=';')

        try: self._f2()
        except: print('ERR', end=';')


class B(A):
    b = 4

    def f1(self): print(5, end=';')
    def _f2(self): print(6, end=';')

try: B().main()
except: print('ERR', end=';')
code

8. Що буде виведено за виконання?
code
def f(arg, n):
    arg[85] = 8
    n -= -3

try:
    n = 4
    s = {28:7, 63:1, 63:5}
    f(s, n)
    print(n, end=';')
    for x in s.items():
        print(x, sep=';', end=';')
except:
    print('ERR')
code

9. Що буде виведено за виконання?
code
class A:
    a = 1
    c = 2

    def __init__(self):
        c = 3

class B(A):
    a = 4

    def __init__(self):
        super().__init__()
        self.b = 7

obj = B()
obj.b = obj.c + 10
B.b = 13

try: print(obj.b, end= ';')
except: print('ERR', end=';')

try: print(A.b, end= ';')
except: print('ERR', end=';')

try: print(B.b, end= ';')
except: print('ERR', end=';')
code

10. 
Для графа на рисунку з вершини 5 запущено алгоритм пошуку в глибину, що будує кістякове дерево. Списки суміжності вершин переглядаються алгоритмом за зростанням міток вершин.\par
\begin{center}
	\adjustimage{max size={0.9\linewidth}{0.9\paperheight}}
	{../10/Traversals/60.png}
\end{center}
\par Які з наведених ребер входять до складу отриманого кістякового дерева?\par
1) (1, 5)\par
2) (3, 6)\par
3) (2, 5)\par
4) (0, 2)\par
5) (1, 6)\par
6) (1, 2)\par
{\vskip 15pt} У формі вибрати номери відповідних ребер.\par
%answer:1 2 4 6
\end {large}}
%end
{\smallskip \begin {small} Затверджено на засіданні кафедри теоретичної кібернетики, протокол \No 12 від   19.05.2020 р.\par\smallskip\smallskip
{\parbox {6.5 cm}{Зав. кафедри \dotfill Ю.В.~Крак}}\hspace {0.7cm}{\hbox to 6.5 cm{ Екзаменатор \dotfill Т.О.~Карнаух}} \end{small}}
%begin
{{\begin {small}\par
{ \center  Київський національний університет імені Тараса Шевченка \par}
{\parbox {5 cm}{Освітня програма \hfill}{\parbox {7 cm}{Прикладна математика, Системний аналіз \hfill}}\parbox {6 cm}{\hfill Семестр 2}}\par
{\parbox {5 cm} {Навчальний предмет \hfill}\parbox {7 cm}{Програмування \hfill}\parbox {6cm}{\hfill}\par}\vspace{-7pt} \end{small}}{\center Екзаменаційний білет \No \textbf 
{ 168 }\par}}
{
\begin {large}
1. Що буде виведено за виконання?
code
try:
    res = "2"==5j
    if isinstance(res, bool):
        print(f'bool;{res}')
    elif isinstance(res, int):
        print(f'int;{res}')
    elif isinstance(res, float):
        print(f'float;{res:.2f}')
    else:
        print('???')
except:
    print('ERR')
code

2. Що буде виведено за виконання?
code
a = 0
b = 4
c = 5

def g(b):
    global c
    a -= 2
    b = 5
    c = 2
    return a + b + c

a = 2
b = 4
c = 0

try:
    print(g(a), a, b, c, sep=';')
except:
    print('ERR', a, b, c, sep=';')
code

3. Що буде виведено за виконання?
code
try:
    i = 7
    while i>=2:
        i += -2
    print(i, end=';')
except:
    print('ERR')
code

4. Що буде виведено за виконання?
code
try:
    for a in range(-3, -3, -2):
        if a<=5:
            continue
        print(a, end=';')
        if a>3:
            break
    else:
        print(a,end=';')
    print(a)
except:
    print('ERR')
code

5. Що буде виведено за виконання?\par (Дійсні значення, якщо наявні, в цьому завданні будуть виводитись у форматі з фіксованою крапкою та, можливо, з доволі великою кількістю десяткових розрядів. У відповіді для дійсних значень у ЦЬОМУ завданні достатньо вказати один десятковий розряд після крапки, відкинувши решту розрядів без округлення. Наприклад, замість 13.66666666666667 достатньо вказати 13.6, замість 25.23 достатньо вказати 25.2)
code
try:
    print(2, end=';')
    print(int('2'), end=';')
    print(8, end=';')
except Exception:
    print(5, end=';')
except ValueError:
    print(0, end=';')
except:
    print('ERR', end=';')
else:
    print(9, end=';')
finally:
    print(3)
code

6. Що буде виведено за виконання?
code
def f(s):
    for el in s:
        el += 1
        return s
 
s = [1, 2, 3]
s = f(s)
print(s)
code

7. Що буде виведено за виконання?
code
class A:
    c = 1

    def _g1(self): print(2, end=';')
    def g2(self): print(3, end=';')

    def main(self):
        try: print(self.c, end=';')
        except: print('ERR', end=';')

        try: self._g1()
        except: print('ERR', end=';')

        try: self.g2()
        except: print('ERR', end=';')


class B(A):
    c = 4

    def _g1(self): print(5, end=';')
    def g2(self): print(6, end=';')

try: B().main()
except: print('ERR', end=';')
code

8. Що буде виведено за виконання?
code
def f(arg, n):
    arg[56] = 3
    n += -2

try:
    n = 5
    d = {73:8, 32:3, 60:8, 52:3}
    f(d, n)
    print(n, end=';')
    for x in d.keys():
        print(x, sep=';', end=';')
except:
    print('ERR')
code

9. Що буде виведено за виконання?
code
class A:
    c = 1
    a = 2

    def __init__(self):
        c = 3

class B(A):
    a = 4

    def __init__(self):
        super().__init__()
        self.b = 7

obj = B()
obj.a = obj.b + 10
B.a = 13

try: print(obj.a, end= ';')
except: print('ERR', end=';')

try: print(A.a, end= ';')
except: print('ERR', end=';')

try: print(B.a, end= ';')
except: print('ERR', end=';')
code

10. 
Для графа на рисунку з вершини 4 запущено алгоритм пошуку в глибину, що будує кістякове дерево. Списки суміжності вершин переглядаються алгоритмом за зростанням міток вершин.\par
\begin{center}
	\adjustimage{max size={0.9\linewidth}{0.9\paperheight}}
	{../10/Traversals/109.png}
\end{center}
\par Які з наведених ребер входять до складу отриманого кістякового дерева?\par
1) (0, 3)\par
2) (1, 2)\par
3) (4, 6)\par
4) (0, 6)\par
5) (1, 4)\par
6) (2, 6)\par
{\vskip 15pt} У формі вибрати номери відповідних ребер.\par
%answer:1 2 4 5 6
\end {large}}
%end
{\smallskip \begin {small} Затверджено на засіданні кафедри теоретичної кібернетики, протокол \No 12 від   19.05.2020 р.\par\smallskip\smallskip
{\parbox {6.5 cm}{Зав. кафедри \dotfill Ю.В.~Крак}}\hspace {0.7cm}{\hbox to 6.5 cm{ Екзаменатор \dotfill Т.О.~Карнаух}} \end{small}}
%begin
{{\begin {small}\par
{ \center  Київський національний університет імені Тараса Шевченка \par}
{\parbox {5 cm}{Освітня програма \hfill}{\parbox {7 cm}{Прикладна математика, Системний аналіз \hfill}}\parbox {6 cm}{\hfill Семестр 2}}\par
{\parbox {5 cm} {Навчальний предмет \hfill}\parbox {7 cm}{Програмування \hfill}\parbox {6cm}{\hfill}\par}\vspace{-7pt} \end{small}}{\center Екзаменаційний білет \No \textbf 
{ 169 }\par}}
{
\begin {large}
1. Що буде виведено за виконання?
code
def f(n):
    print('f', end='')
    return n<=-1

try:
    print(f(8) or f(-1))
except:
    print('ERR')
code

2. Що буде виведено за виконання?
code
a = 6
b = 7
c = 9

def f(b):
    a = 1
    b *= 3
    c = 4
    return a + b + c

a = 1
b = 8
c = 3

try:
    print(f(a), a, b, c, sep=';')
except:
    print('ERR', a, b, c, sep=';')
code

3. Що буде виведено за виконання?
code
try:
    i = 11
    while i>=6:
        i += -3
    print(i, end=';')
except:
    print('ERR')
code

4. Що буде виведено за виконання?
code
try:
    for b in range(2, 6, -1):
        if b>=6:
            continue
        print(b, end=';')
    else:
        print(b,end=';')
    print(b)
except:
    print('ERR')
code

5. Що буде виведено за виконання?\par (Дійсні значення, якщо наявні, в цьому завданні будуть виводитись у форматі з фіксованою крапкою та, можливо, з доволі великою кількістю десяткових розрядів. У відповіді для дійсних значень у ЦЬОМУ завданні достатньо вказати один десятковий розряд після крапки, відкинувши решту розрядів без округлення. Наприклад, замість 13.66666666666667 достатньо вказати 13.6, замість 25.23 достатньо вказати 25.2)
code
try:
    print(4, end=';')
    print(int('a6'), end=';')
    print(1, end=';')
except ZeroDivisionError:
    print(7, end=';')
except KeyboardInterrupt:
    print(6, end=';')
except:
    print('ERR', end=';')
else:
    print(4, end=';')
finally:
    print(9)
code

6. Що буде виведено за виконання?
code
def f(n):
    print(n%10, end='')
    if n>9:
        f(n//10)
 
f(543216)
code

7. Що буде виведено за виконання?
code
class A:
    c = 1

    def __g1(self): print(2, end=';')
    def _g2(self): print(3, end=';')

    def main(self):
        try: print(self.c, end=';')
        except: print('ERR', end=';')

        try: self.__g1()
        except: print('ERR', end=';')

        try: self._g2()
        except: print('ERR', end=';')


class B(A):
    c = 4

    def __g1(self): print(5, end=';')
    def _g2(self): print(6, end=';')

try: B().main()
except: print('ERR', end=';')
code

8. Що буде виведено за виконання?
code
def f(arg, n):
    n = 4
    tmp = arg
    tmp[-3:3] = 'xy'
    return len(tmp)>3

try:
    e = ['0','1','2','3','4']
    n = 6
    f(e, n)
    print(n, end=';')
    for el in e:
        print(el, end=';')
except:
    print('ERR')
code

9. Що буде виведено за виконання?
code
class A:
    b = 1
    a = 2

    def __init__(self):
        b = 3

class B(A):
    a = 4

    def __init__(self):
        super().__init__()
        self.c = 7

obj = B()
obj.b = obj.a + 10
B.b = 13

try: print(obj.b, end= ';')
except: print('ERR', end=';')

try: print(A.b, end= ';')
except: print('ERR', end=';')

try: print(B.b, end= ';')
except: print('ERR', end=';')
code

10. 
Для графа на рисунку з вершини 3 запущено алгоритм пошуку в глибину, що будує кістякове дерево. Списки суміжності вершин переглядаються алгоритмом за зростанням міток вершин.\par
\begin{center}
	\adjustimage{max size={0.9\linewidth}{0.9\paperheight}}
	{../10/Traversals/50.png}
\end{center}
\par Які з наведених ребер входять до складу отриманого кістякового дерева?\par
1) (5, 6)\par
2) (0, 1)\par
3) (2, 3)\par
4) (2, 4)\par
5) (1, 3)\par
6) (4, 6)\par
{\vskip 15pt} У формі вибрати номери відповідних ребер.\par
%answer:1 2 4 5 6
\end {large}}
%end
{\smallskip \begin {small} Затверджено на засіданні кафедри теоретичної кібернетики, протокол \No 12 від   19.05.2020 р.\par\smallskip\smallskip
{\parbox {6.5 cm}{Зав. кафедри \dotfill Ю.В.~Крак}}\hspace {0.7cm}{\hbox to 6.5 cm{ Екзаменатор \dotfill Т.О.~Карнаух}} \end{small}}
%begin
{{\begin {small}\par
{ \center  Київський національний університет імені Тараса Шевченка \par}
{\parbox {5 cm}{Освітня програма \hfill}{\parbox {7 cm}{Прикладна математика, Системний аналіз \hfill}}\parbox {6 cm}{\hfill Семестр 2}}\par
{\parbox {5 cm} {Навчальний предмет \hfill}\parbox {7 cm}{Програмування \hfill}\parbox {6cm}{\hfill}\par}\vspace{-7pt} \end{small}}{\center Екзаменаційний білет \No \textbf 
{ 170 }\par}}
{
\begin {large}
1. Що буде виведено за виконання?
code
try:
    res = 9**-0.5*2
    if isinstance(res, bool):
        print(f'bool;{res}')
    elif isinstance(res, int):
        print(f'int;{res}')
    elif isinstance(res, float):
        print(f'float;{res:.2f}')
    else:
        print('???')
except:
    print('ERR')
code

2. Що буде виведено за виконання?
code
a = 4
b = 1
c = 5

def h(a):
    global c
    a -= 5
    b = 4
    c = 3
    return a + b + c

a = 7
b = 6
c = 0

try:
    print(h(b), a, b, c, sep=';')
except:
    print('ERR', a, b, c, sep=';')
code

3. Що буде виведено за виконання?
code
try:
    m = 2
    while m<9:
        m += 2
    print(m, end=';')
except:
    print('ERR')
code

4. Що буде виведено за виконання?
code
try:
    for c in range(15, -7, 3):
        if c<=8:
            break
        print(c, end=';');c = -9
    else:
        print(c,end=';')
    print(c)
except:
    print('ERR')
code

5. Що буде виведено за виконання?\par (Дійсні значення, якщо наявні, в цьому завданні будуть виводитись у форматі з фіксованою крапкою та, можливо, з доволі великою кількістю десяткових розрядів. У відповіді для дійсних значень у ЦЬОМУ завданні достатньо вказати один десятковий розряд після крапки, відкинувши решту розрядів без округлення. Наприклад, замість 13.66666666666667 достатньо вказати 13.6, замість 25.23 достатньо вказати 25.2)
code
try:
    print(7, end=';')
    print(0%3, end=';')
    print(1, end=';')
except TypeError:
    print(8, end=';')
except Exception:
    print(2, end=';')
except:
    print('ERR', end=';')
else:
    print(3, end=';')
finally:
    print(5)
code

6. Що буде виведено за виконання?
code
def f(n):
    if n>2:
        f(n//2)
    print(n%2, end='')
 
f(16)
code

7. Що буде виведено за виконання?
code
class A:
    a = 1

    def _f1(self): print(2, end=';')
    def _f2(self): print(3, end=';')

    def main(self):
        try: print(self.a, end=';')
        except: print('ERR', end=';')

        try: self._f1()
        except: print('ERR', end=';')

        try: self._f2()
        except: print('ERR', end=';')


class B(A):
    a = 4

    def _f1(self): print(5, end=';')
    def _f2(self): print(6, end=';')

try: B().main()
except: print('ERR', end=';')
code

8. Що буде виведено за виконання?
code
def f(arg, n):
    arg[62] = 0
    n += 3

try:
    n = 3
    d = {35:4, 86:8, 28:9}
    f(d, n)
    print(n, end=';')
    for x in d.values():
        print(x, sep=';', end=';')
except:
    print('ERR')
code

9. Що буде виведено за виконання?
code
class A:
    c = 1
    a = 2

    def __init__(self):
        a = 3

class B(A):
    c = 4

    def __init__(self):
        super().__init__()
        self.b = 7

obj = B()
obj.c = obj.a + 10
B.c = 13

try: print(obj.c, end= ';')
except: print('ERR', end=';')

try: print(A.c, end= ';')
except: print('ERR', end=';')

try: print(B.c, end= ';')
except: print('ERR', end=';')
code

10. 
Для графа на рисунку з вершини 5 запущено алгоритм пошуку в глибину, що будує кістякове дерево. Списки суміжності вершин переглядаються алгоритмом за зростанням міток вершин.\par
\begin{center}
	\adjustimage{max size={0.9\linewidth}{0.9\paperheight}}
	{../10/Traversals/239.png}
\end{center}
\par Які з наведених ребер входять до складу отриманого кістякового дерева?\par
1) (1, 4)\par
2) (3, 4)\par
3) (5, 6)\par
4) (0, 5)\par
5) (4, 6)\par
6) (4, 5)\par
{\vskip 15pt} У формі вибрати номери відповідних ребер.\par
%answer:1 2 4
\end {large}}
%end
{\smallskip \begin {small} Затверджено на засіданні кафедри теоретичної кібернетики, протокол \No 12 від   19.05.2020 р.\par\smallskip\smallskip
{\parbox {6.5 cm}{Зав. кафедри \dotfill Ю.В.~Крак}}\hspace {0.7cm}{\hbox to 6.5 cm{ Екзаменатор \dotfill Т.О.~Карнаух}} \end{small}}
%begin
{{\begin {small}\par
{ \center  Київський національний університет імені Тараса Шевченка \par}
{\parbox {5 cm}{Освітня програма \hfill}{\parbox {7 cm}{Прикладна математика, Системний аналіз \hfill}}\parbox {6 cm}{\hfill Семестр 2}}\par
{\parbox {5 cm} {Навчальний предмет \hfill}\parbox {7 cm}{Програмування \hfill}\parbox {6cm}{\hfill}\par}\vspace{-7pt} \end{small}}{\center Екзаменаційний білет \No \textbf 
{ 171 }\par}}
{
\begin {large}
1. Що буде виведено за виконання?
code
try:
    res = 5<7.0 and not 3.0>False or 9>=2
    if isinstance(res, bool):
        print(f'bool;{res}')
    elif isinstance(res, int):
        print(f'int;{res}')
    elif isinstance(res, float):
        print(f'float;{res:.2f}')
    else:
        print('???')
except:
    print('ERR')
code

2. Що буде виведено за виконання?
code
a = 2
b = 3
c = 8

def g(b):
    global c
    a = 1
    b += 2
    c = 4
    return a + b + c

a = 9
b = 2
c = 4

try:
    print(g(a), a, b, c, sep=';')
except:
    print('ERR', a, b, c, sep=';')
code

3. Що буде виведено за виконання?
code
try:
    k = 1
    while k<=11:
        k += 1
    print(k, end=';')
except:
    print('ERR')
code

4. Що буде виведено за виконання?
code
try:
    for d in range(5, 12, -2):
        if d>7:
            continue
        print(d, end=';')
    else:
        print(d,end=';')
    print(d)
except:
    print('ERR')
code

5. Що буде виведено за виконання?\par (Дійсні значення, якщо наявні, в цьому завданні будуть виводитись у форматі з фіксованою крапкою та, можливо, з доволі великою кількістю десяткових розрядів. У відповіді для дійсних значень у ЦЬОМУ завданні достатньо вказати один десятковий розряд після крапки, відкинувши решту розрядів без округлення. Наприклад, замість 13.66666666666667 достатньо вказати 13.6, замість 25.23 достатньо вказати 25.2)
code
try:
    print(8, end=';')
    print(int('d7'), end=';')
    print(3, end=';')
except TypeError:
    print(1, end=';')
except ZeroDivisionError:
    print(4, end=';')
except:
    print('ERR', end=';')
else:
    print(0, end=';')
finally:
    print(6)
code

6. Що буде виведено за виконання?
code
def f(n):
    if n>9:
        f(n//10)
    print(n%10, end='')
 
f(543216)
code

7. Що буде виведено за виконання?
code
class A:
    b = 1

    def __h1(self): print(2, end=';')
    def _h2(self): print(3, end=';')

    def main(self):
        try: print(self.b, end=';')
        except: print('ERR', end=';')

        try: self.__h1()
        except: print('ERR', end=';')

        try: self._h2()
        except: print('ERR', end=';')


class B(A):
    b = 4

    def __h1(self): print(5, end=';')
    def _h2(self): print(6, end=';')

try: B().main()
except: print('ERR', end=';')
code

8. Що буде виведено за виконання?
code
def f(arg, n):
    arg[11] = 6
    n = 5
    arg = list(arg.values())

try:
    n = 2
    d = {74:6, 11:5, 74:8}
    f(d, n)
    print(n, end=';')
    for x, y in d.items():
        print(x, y, sep=';', end=';')
except:
    print('ERR')
code

9. Що буде виведено за виконання?
code
class A:
    a = 1
    c = 2

    def __init__(self):
        c = 3

class B(A):
    a = 4

    def __init__(self):
        super().__init__()
        self.b = 7

obj = B()
obj.b = obj.c + 10
B.b = 13

try: print(obj.b, end= ';')
except: print('ERR', end=';')

try: print(A.b, end= ';')
except: print('ERR', end=';')

try: print(B.b, end= ';')
except: print('ERR', end=';')
code

10. 
Для графа на рисунку з вершини 2 запущено алгоритм пошуку в глибину, що будує кістякове дерево. Списки суміжності вершин переглядаються алгоритмом за зростанням міток вершин.\par
\begin{center}
	\adjustimage{max size={0.9\linewidth}{0.9\paperheight}}
	{../10/Traversals/47.png}
\end{center}
\par Які з наведених ребер входять до складу отриманого кістякового дерева?\par
1) (1, 4)\par
2) (0, 2)\par
3) (1, 3)\par
4) (0, 5)\par
5) (3, 4)\par
6) (0, 1)\par
{\vskip 15pt} У формі вибрати номери відповідних ребер.\par
%answer:2 3 5 6
\end {large}}
%end
{\smallskip \begin {small} Затверджено на засіданні кафедри теоретичної кібернетики, протокол \No 12 від   19.05.2020 р.\par\smallskip\smallskip
{\parbox {6.5 cm}{Зав. кафедри \dotfill Ю.В.~Крак}}\hspace {0.7cm}{\hbox to 6.5 cm{ Екзаменатор \dotfill Т.О.~Карнаух}} \end{small}}
%begin
{{\begin {small}\par
{ \center  Київський національний університет імені Тараса Шевченка \par}
{\parbox {5 cm}{Освітня програма \hfill}{\parbox {7 cm}{Прикладна математика, Системний аналіз \hfill}}\parbox {6 cm}{\hfill Семестр 2}}\par
{\parbox {5 cm} {Навчальний предмет \hfill}\parbox {7 cm}{Програмування \hfill}\parbox {6cm}{\hfill}\par}\vspace{-7pt} \end{small}}{\center Екзаменаційний білет \No \textbf 
{ 172 }\par}}
{
\begin {large}
1. Що буде виведено за виконання?
code
try:
    res = -3**4--2
    if isinstance(res, bool):
        print(f'bool;{res}')
    elif isinstance(res, int):
        print(f'int;{res}')
    elif isinstance(res, float):
        print(f'float;{res:.2f}')
    else:
        print('???')
except:
    print('ERR')
code

2. Що буде виведено за виконання?
code
a = 5
b = 7
c = 8

def g(a):
    a = 1
    b -= 5
    c = 3
    return a + b + c

a = 1
b = 0
c = 3

try:
    print(g(b), a, b, c, sep=';')
except:
    print('ERR', a, b, c, sep=';')
code

3. Що буде виведено за виконання?
code
try:
    m = 15
    while m>8:
        m += -2
    print(m, end=';')
except:
    print('ERR')
code

4. Що буде виведено за виконання?
code
try:
    for a in range(-6, -5, 2):
        if a<4:
            continue
        print(a, end=';')
        if a<=6:
            break
    else:
        print(a,end=';')
    print(a)
except:
    print('ERR')
code

5. Що буде виведено за виконання?\par (Дійсні значення, якщо наявні, в цьому завданні будуть виводитись у форматі з фіксованою крапкою та, можливо, з доволі великою кількістю десяткових розрядів. У відповіді для дійсних значень у ЦЬОМУ завданні достатньо вказати один десятковий розряд після крапки, відкинувши решту розрядів без округлення. Наприклад, замість 13.66666666666667 достатньо вказати 13.6, замість 25.23 достатньо вказати 25.2)
code
try:
    print(2, end=';')
    print(9/0, end=';')
    print(5, end=';')
except KeyboardInterrupt:
    print(6, end=';')
except ValueError:
    print(2, end=';')
except:
    print('ERR', end=';')
else:
    print(9, end=';')
finally:
    print(4)
code

6. Що буде виведено за виконання?
code
def f(n):
    if n>9:
        f(n//100)
        print(n%10, end='')
 
f(123456)
code

7. Що буде виведено за виконання?
code
class A:
    d = 1

    def f1(self): print(2, end=';')
    def _f2(self): print(3, end=';')

    def main(self):
        try: print(self.d, end=';')
        except: print('ERR', end=';')

        try: self.f1()
        except: print('ERR', end=';')

        try: self._f2()
        except: print('ERR', end=';')


class B(A):
    d = 4

    def f1(self): print(5, end=';')
    def _f2(self): print(6, end=';')

try: B().main()
except: print('ERR', end=';')
code

8. Що буде виведено за виконання?
code
def f(arg, n):
    arg[32] = 8
    n = 6
    arg = set(arg.keys())

try:
    n = 1
    t = {32:8, 26:2, 26:2}
    f(t, n)
    print(n, end=';')
    for x, y in t.items():
        print(x, y, sep=';', end=';')
except:
    print('ERR')
code

9. Що буде виведено за виконання?
code
class A:
    c = 1
    b = 2

    def __init__(self):
        c = 3

class B(A):
    b = 4

    def __init__(self):
        super().__init__()
        self.a = 7

obj = B()
obj.b = obj.a + 10
B.b = 13

try: print(obj.b, end= ';')
except: print('ERR', end=';')

try: print(A.b, end= ';')
except: print('ERR', end=';')

try: print(B.b, end= ';')
except: print('ERR', end=';')
code

10. 
Для графа на рисунку з вершини 6 запущено алгоритм пошуку в глибину, що будує кістякове дерево. Списки суміжності вершин переглядаються алгоритмом за зростанням міток вершин.\par
\begin{center}
	\adjustimage{max size={0.9\linewidth}{0.9\paperheight}}
	{../10/Traversals/24.png}
\end{center}
\par Які з наведених ребер входять до складу отриманого кістякового дерева?\par
1) (0, 1)\par
2) (3, 6)\par
3) (4, 5)\par
4) (2, 5)\par
5) (1, 4)\par
6) (5, 6)\par
{\vskip 15pt} У формі вибрати номери відповідних ребер.\par
%answer:1 2 5
\end {large}}
%end
{\smallskip \begin {small} Затверджено на засіданні кафедри теоретичної кібернетики, протокол \No 12 від   19.05.2020 р.\par\smallskip\smallskip
{\parbox {6.5 cm}{Зав. кафедри \dotfill Ю.В.~Крак}}\hspace {0.7cm}{\hbox to 6.5 cm{ Екзаменатор \dotfill Т.О.~Карнаух}} \end{small}}
%begin
{{\begin {small}\par
{ \center  Київський національний університет імені Тараса Шевченка \par}
{\parbox {5 cm}{Освітня програма \hfill}{\parbox {7 cm}{Прикладна математика, Системний аналіз \hfill}}\parbox {6 cm}{\hfill Семестр 2}}\par
{\parbox {5 cm} {Навчальний предмет \hfill}\parbox {7 cm}{Програмування \hfill}\parbox {6cm}{\hfill}\par}\vspace{-7pt} \end{small}}{\center Екзаменаційний білет \No \textbf 
{ 173 }\par}}
{
\begin {large}
1. Що буде виведено за виконання?
code
try:
    res = 3/9//9*3
    if isinstance(res, bool):
        print(f'bool;{res}')
    elif isinstance(res, int):
        print(f'int;{res}')
    elif isinstance(res, float):
        print(f'float;{res:.2f}')
    else:
        print('???')
except:
    print('ERR')
code

2. Що буде виведено за виконання?
code
a = 0
b = 2
c = 7

def f(a):
    a = 5
    b = 3
    c = 4
    return a + b + c

a = 3
b = 1
c = 4

try:
    print(f(b), a, b, c, sep=';')
except:
    print('ERR', a, b, c, sep=';')
code

3. Що буде виведено за виконання?
code
try:
    j = 6
    while j<=5:
        j += 1
    print(j, end=';')
except:
    print('ERR')
code

4. Що буде виведено за виконання?
code
try:
    for f in range(-14, -9, 3):
        if f>5:
            continue
        print(f, end=';')
    else:
        print('end',end=';')
    print(f)
except:
    print('ERR')
code

5. Що буде виведено за виконання?\par (Дійсні значення, якщо наявні, в цьому завданні будуть виводитись у форматі з фіксованою крапкою та, можливо, з доволі великою кількістю десяткових розрядів. У відповіді для дійсних значень у ЦЬОМУ завданні достатньо вказати один десятковий розряд після крапки, відкинувши решту розрядів без округлення. Наприклад, замість 13.66666666666667 достатньо вказати 13.6, замість 25.23 достатньо вказати 25.2)
code
try:
    print(5, end=';')
    print(int('3'), end=';')
    print(7, end=';')
except Exception:
    print(8, end=';')
except ZeroDivisionError:
    print(0, end=';')
except:
    print('ERR', end=';')
else:
    print(1, end=';')
finally:
    print(8)
code

6. Що буде виведено за виконання?
code
def f(n):
    print(n%10,end='')
    if n>9:
        f(n//100)
 
f(123456)
code

7. Що буде виведено за виконання?
code
class A:
    c = 1

    def _h1(self): print(2, end=';')
    def _h2(self): print(3, end=';')

    def main(self):
        try: print(self.c, end=';')
        except: print('ERR', end=';')

        try: self._h1()
        except: print('ERR', end=';')

        try: self._h2()
        except: print('ERR', end=';')


class B(A):
    c = 4

    def _h1(self): print(5, end=';')
    def _h2(self): print(6, end=';')

try: B().main()
except: print('ERR', end=';')
code

8. Що буде виведено за виконання?
code
def f(arg, n):
    arg[50] = 4
    n = 8
    arg = tuple(arg.values())

try:
    n = 0
    t = {50:4, 26:6, 26:3}
    f(t, n)
    print(n, end=';')
    for x in t.keys():
        print(x, sep=';', end=';')
except:
    print('ERR')
code

9. Що буде виведено за виконання?
code
class A:
    a = 1
    b = 2

    def __init__(self):
        a = 3

class B(A):
    b = 4

    def __init__(self):
        super().__init__()
        self.c = 7

obj = B()
obj.c = obj.c + 10
B.c = 13

try: print(obj.c, end= ';')
except: print('ERR', end=';')

try: print(A.c, end= ';')
except: print('ERR', end=';')

try: print(B.c, end= ';')
except: print('ERR', end=';')
code

10. 
Для графа на рисунку з вершини 2 запущено алгоритм пошуку в глибину, що будує кістякове дерево. Списки суміжності вершин переглядаються алгоритмом за зростанням міток вершин.\par
\begin{center}
	\adjustimage{max size={0.9\linewidth}{0.9\paperheight}}
	{../10/Traversals/257.png}
\end{center}
\par Які з наведених ребер входять до складу отриманого кістякового дерева?\par
1) (4, 5)\par
2) (5, 6)\par
3) (4, 6)\par
4) (1, 5)\par
5) (3, 4)\par
6) (0, 5)\par
{\vskip 15pt} У формі вибрати номери відповідних ребер.\par
%answer:3 5 6
\end {large}}
%end
{\smallskip \begin {small} Затверджено на засіданні кафедри теоретичної кібернетики, протокол \No 12 від   19.05.2020 р.\par\smallskip\smallskip
{\parbox {6.5 cm}{Зав. кафедри \dotfill Ю.В.~Крак}}\hspace {0.7cm}{\hbox to 6.5 cm{ Екзаменатор \dotfill Т.О.~Карнаух}} \end{small}}
%begin
{{\begin {small}\par
{ \center  Київський національний університет імені Тараса Шевченка \par}
{\parbox {5 cm}{Освітня програма \hfill}{\parbox {7 cm}{Прикладна математика, Системний аналіз \hfill}}\parbox {6 cm}{\hfill Семестр 2}}\par
{\parbox {5 cm} {Навчальний предмет \hfill}\parbox {7 cm}{Програмування \hfill}\parbox {6cm}{\hfill}\par}\vspace{-7pt} \end{small}}{\center Екзаменаційний білет \No \textbf 
{ 174 }\par}}
{
\begin {large}
1. Що буде виведено за виконання?
code
try:
    res = 9//7/6*9
    if isinstance(res, bool):
        print(f'bool;{res}')
    elif isinstance(res, int):
        print(f'int;{res}')
    elif isinstance(res, float):
        print(f'float;{res:.2f}')
    else:
        print('???')
except:
    print('ERR')
code

2. Що буде виведено за виконання?
code
a = 0
b = 8
c = 6

def g(b):
    global c
    a = 2
    b = 1
    c = 4
    return a + b + c

a = 4
b = 3
c = 2

try:
    print(g(b), a, b, c, sep=';')
except:
    print('ERR', a, b, c, sep=';')
code

3. Що буде виведено за виконання?
code
try:
    k = 6
    while k>9:
        k += -2
    print(k, end=';')
except:
    print('ERR')
code

4. Що буде виведено за виконання?
code
try:
    for e in range(9, 10, -1):
        if e<=3:
            break
        print(e, end=';');e = -1
    else:
        print(e,end=';')
    print(e)
except:
    print('ERR')
code

5. Що буде виведено за виконання?\par (Дійсні значення, якщо наявні, в цьому завданні будуть виводитись у форматі з фіксованою крапкою та, можливо, з доволі великою кількістю десяткових розрядів. У відповіді для дійсних значень у ЦЬОМУ завданні достатньо вказати один десятковий розряд після крапки, відкинувши решту розрядів без округлення. Наприклад, замість 13.66666666666667 достатньо вказати 13.6, замість 25.23 достатньо вказати 25.2)
code
try:
    print(3, end=';')
    print(1%2, end=';')
    print(0, end=';')
except TypeError:
    print(9, end=';')
except ValueError:
    print(5, end=';')
except:
    print('ERR', end=';')
else:
    print(7, end=';')
finally:
    print(4)
code

6. Що буде виведено за виконання?
code
def f(n):
    print(n%2, end='')
    if n>2:
        f(n//2)
 
f(16)
code

7. Що буде виведено за виконання?
code
class A:
    b = 1

    def _g1(self): print(2, end=';')
    def g2(self): print(3, end=';')

    def main(self):
        try: print(self.b, end=';')
        except: print('ERR', end=';')

        try: self._g1()
        except: print('ERR', end=';')

        try: self.g2()
        except: print('ERR', end=';')


class B(A):
    b = 4

    def _g1(self): print(5, end=';')
    def g2(self): print(6, end=';')

try: B().main()
except: print('ERR', end=';')
code

8. Що буде виведено за виконання?
code
def f(arg, n):
    arg[30] = 2
    n = 9
    arg = list(arg.items())

try:
    n = 3
    d = {86:5, 43:2, 86:5}
    f(d, n)
    print(n, end=';')
    for x in d.values():
        print(x, sep=';', end=';')
except:
    print('ERR')
code

9. Що буде виведено за виконання?
code
class A:
    b = 1
    c = 2

    def __init__(self):
        c = 3

class B(A):
    b = 4

    def __init__(self):
        super().__init__()
        self.a = 7

obj = B()
obj.b = obj.a + 10
B.b = 13

try: print(obj.b, end= ';')
except: print('ERR', end=';')

try: print(A.b, end= ';')
except: print('ERR', end=';')

try: print(B.b, end= ';')
except: print('ERR', end=';')
code

10. 
Для графа на рисунку з вершини 6 запущено алгоритм пошуку в глибину, що будує кістякове дерево. Списки суміжності вершин переглядаються алгоритмом за зростанням міток вершин.\par
\begin{center}
	\adjustimage{max size={0.9\linewidth}{0.9\paperheight}}
	{../10/Traversals/188.png}
\end{center}
\par Які з наведених ребер входять до складу отриманого кістякового дерева?\par
1) (1, 6)\par
2) (2, 4)\par
3) (1, 4)\par
4) (3, 6)\par
5) (1, 3)\par
6) (1, 2)\par
{\vskip 15pt} У формі вибрати номери відповідних ребер.\par
%answer:1 2 5
\end {large}}
%end
{\smallskip \begin {small} Затверджено на засіданні кафедри теоретичної кібернетики, протокол \No 12 від   19.05.2020 р.\par\smallskip\smallskip
{\parbox {6.5 cm}{Зав. кафедри \dotfill Ю.В.~Крак}}\hspace {0.7cm}{\hbox to 6.5 cm{ Екзаменатор \dotfill Т.О.~Карнаух}} \end{small}}
%begin
{{\begin {small}\par
{ \center  Київський національний університет імені Тараса Шевченка \par}
{\parbox {5 cm}{Освітня програма \hfill}{\parbox {7 cm}{Прикладна математика, Системний аналіз \hfill}}\parbox {6 cm}{\hfill Семестр 2}}\par
{\parbox {5 cm} {Навчальний предмет \hfill}\parbox {7 cm}{Програмування \hfill}\parbox {6cm}{\hfill}\par}\vspace{-7pt} \end{small}}{\center Екзаменаційний білет \No \textbf 
{ 175 }\par}}
{
\begin {large}
1. Що буде виведено за виконання?
code
def f(n):
    print('f', end='')
    return n>=0

try:
    print(f(4) and f(-2))
except:
    print('ERR')
code

2. Що буде виведено за виконання?
code
a = 6
b = 9
c = 3

def f(a):
    a = 2
    b += 2
    c = 5
    return a + b + c

a = 5
b = 7
c = 2

try:
    print(f(b), a, b, c, sep=';')
except:
    print('ERR', a, b, c, sep=';')
code

3. Що буде виведено за виконання?
code
try:
    n = 4
    while n<8:
        n += 2
    print(n, end=';')
except:
    print('ERR')
code

4. Що буде виведено за виконання?
code
try:
    for c in range(-5, -6, -3):
        if c<7:
            break
        print(c, end=';');c = -4
    else:
        print(c,end=';')
    print(c)
except:
    print('ERR')
code

5. Що буде виведено за виконання?\par (Дійсні значення, якщо наявні, в цьому завданні будуть виводитись у форматі з фіксованою крапкою та, можливо, з доволі великою кількістю десяткових розрядів. У відповіді для дійсних значень у ЦЬОМУ завданні достатньо вказати один десятковий розряд після крапки, відкинувши решту розрядів без округлення. Наприклад, замість 13.66666666666667 достатньо вказати 13.6, замість 25.23 достатньо вказати 25.2)
code
try:
    print(6, end=';')
    print(2j<2j, end=';')
    print(2, end=';')
except Exception:
    print(5, end=';')
except KeyboardInterrupt:
    print(1, end=';')
except:
    print('ERR', end=';')
else:
    print(9, end=';')
finally:
    print(7)
code

6. Що буде виведено за виконання?
code
def f(s):
    s1 = s
    for i in range(len(s)):
        s1[i] += 1
        return s1

s = [1, 2, 3]
f(s)
print(s)
code

7. Що буде виведено за виконання?
code
class A:
    a = 1

    def _h1(self): print(2, end=';')
    def __h2(self): print(3, end=';')

    def main(self):
        try: print(self.a, end=';')
        except: print('ERR', end=';')

        try: self._h1()
        except: print('ERR', end=';')

        try: self.__h2()
        except: print('ERR', end=';')


class B(A):
    a = 4

    def _h1(self): print(5, end=';')
    def __h2(self): print(6, end=';')

try: B().main()
except: print('ERR', end=';')
code

8. Що буде виведено за виконання?
code
def f(arg, n):
    arg[78] = 4
    n *= 2

try:
    n = 6
    t = {19:6, 64:3, 19:3}
    f(t, n)
    print(n, end=';')
    for x in t :
        print(x, sep=';', end=';')
except:
    print('ERR')
code

9. Що буде виведено за виконання?
code
class A:
    a = 1
    c = 2

    def __init__(self):
        c = 3

class B(A):
    c = 4

    def __init__(self):
        super().__init__()
        self.b = 7

obj = B()
obj.c = obj.b + 10
B.c = 13

try: print(obj.c, end= ';')
except: print('ERR', end=';')

try: print(A.c, end= ';')
except: print('ERR', end=';')

try: print(B.c, end= ';')
except: print('ERR', end=';')
code

10. 
Для графа на рисунку з вершини 5 запущено алгоритм пошуку в глибину, що будує кістякове дерево. Списки суміжності вершин переглядаються алгоритмом за зростанням міток вершин.\par
\begin{center}
	\adjustimage{max size={0.9\linewidth}{0.9\paperheight}}
	{../10/Traversals/247.png}
\end{center}
\par Які з наведених ребер входять до складу отриманого кістякового дерева?\par
1) (2, 4)\par
2) (3, 5)\par
3) (1, 4)\par
4) (0, 3)\par
5) (1, 6)\par
6) (2, 3)\par
{\vskip 15pt} У формі вибрати номери відповідних ребер.\par
%answer:1 4 6
\end {large}}
%end
{\smallskip \begin {small} Затверджено на засіданні кафедри теоретичної кібернетики, протокол \No 12 від   19.05.2020 р.\par\smallskip\smallskip
{\parbox {6.5 cm}{Зав. кафедри \dotfill Ю.В.~Крак}}\hspace {0.7cm}{\hbox to 6.5 cm{ Екзаменатор \dotfill Т.О.~Карнаух}} \end{small}}
%begin
{{\begin {small}\par
{ \center  Київський національний університет імені Тараса Шевченка \par}
{\parbox {5 cm}{Освітня програма \hfill}{\parbox {7 cm}{Прикладна математика, Системний аналіз \hfill}}\parbox {6 cm}{\hfill Семестр 2}}\par
{\parbox {5 cm} {Навчальний предмет \hfill}\parbox {7 cm}{Програмування \hfill}\parbox {6cm}{\hfill}\par}\vspace{-7pt} \end{small}}{\center Екзаменаційний білет \No \textbf 
{ 176 }\par}}
{
\begin {large}
1. Що буде виведено за виконання?
code
def f(n):
    print('f', end='')
    return n

try:
    print(f(-4) or f(6))
except:
    print('ERR')
code

2. Що буде виведено за виконання?
code
a = 5
b = 9
c = 7

def h(a):
    a = 3
    b = 5
    c = 1
    return a + b + c

a = 1
b = 1
c = 2

try:
    print(h(a), a, b, c, sep=';')
except:
    print('ERR', a, b, c, sep=';')
code

3. Що буде виведено за виконання?
code
try:
    i = 5
    while i<13:
        i += 2
    print(i, end=';')
except:
    print('ERR')
code

4. Що буде виведено за виконання?
code
try:
    for a in range(-15, 13, 2):
        if a>=4:
            continue
        print(a, end=';')
    else:
        print(a,end=';')
    print(a)
except:
    print('ERR')
code

5. Що буде виведено за виконання?\par (Дійсні значення, якщо наявні, в цьому завданні будуть виводитись у форматі з фіксованою крапкою та, можливо, з доволі великою кількістю десяткових розрядів. У відповіді для дійсних значень у ЦЬОМУ завданні достатньо вказати один десятковий розряд після крапки, відкинувши решту розрядів без округлення. Наприклад, замість 13.66666666666667 достатньо вказати 13.6, замість 25.23 достатньо вказати 25.2)
code
try:
    print(4, end=';')
    print(3//1, end=';')
    print(8, end=';')
except KeyboardInterrupt:
    print(6, end=';')
except ZeroDivisionError:
    print(0, end=';')
except:
    print('ERR', end=';')
else:
    print(6, end=';')
finally:
    print(5)
code

6. Що буде виведено за виконання?
code
def f(s):
    s1 = s
    for i in range(len(s)):
        s1[i] += 1
    return s1

s = [1, 2, 3]
f(s)
print(s)
code

7. Що буде виведено за виконання?
code
class A:
    d = 1

    def _f1(self): print(2, end=';')
    def f2(self): print(3, end=';')

    def main(self):
        try: print(self.d, end=';')
        except: print('ERR', end=';')

        try: self._f1()
        except: print('ERR', end=';')

        try: self.f2()
        except: print('ERR', end=';')


class B(A):
    d = 4

    def _f1(self): print(5, end=';')
    def f2(self): print(6, end=';')

try: B().main()
except: print('ERR', end=';')
code

8. Що буде виведено за виконання?
code
def f(arg, n):
    n = 3
    tmp = arg
    del tmp[-7:2]
    return len(tmp)>3

try:
    d = [10,11,12,13,14]
    n = 1
    f(d, n)
    print(n, end=';')
    for el in d:
        print(el, end=';')
except:
    print('ERR')
code

9. Що буде виведено за виконання?
code
class A:
    c = 1
    b = 2

    def __init__(self):
        c = 3

class B(A):
    c = 4

    def __init__(self):
        super().__init__()
        self.a = 7

obj = B()
obj.a = obj.a + 10
B.a = 13

try: print(obj.a, end= ';')
except: print('ERR', end=';')

try: print(A.a, end= ';')
except: print('ERR', end=';')

try: print(B.a, end= ';')
except: print('ERR', end=';')
code

10. 
Для графа на рисунку з вершини 0 запущено алгоритм пошуку в глибину, що будує кістякове дерево. Списки суміжності вершин переглядаються алгоритмом за зростанням міток вершин.\par
\begin{center}
	\adjustimage{max size={0.9\linewidth}{0.9\paperheight}}
	{../10/Traversals/184.png}
\end{center}
\par Які з наведених ребер входять до складу отриманого кістякового дерева?\par
1) (0, 5)\par
2) (4, 5)\par
3) (2, 5)\par
4) (4, 6)\par
5) (2, 6)\par
6) (0, 1)\par
{\vskip 15pt} У формі вибрати номери відповідних ребер.\par
%answer:1 3 4 6
\end {large}}
%end
{\smallskip \begin {small} Затверджено на засіданні кафедри теоретичної кібернетики, протокол \No 12 від   19.05.2020 р.\par\smallskip\smallskip
{\parbox {6.5 cm}{Зав. кафедри \dotfill Ю.В.~Крак}}\hspace {0.7cm}{\hbox to 6.5 cm{ Екзаменатор \dotfill Т.О.~Карнаух}} \end{small}}
%begin
{{\begin {small}\par
{ \center  Київський національний університет імені Тараса Шевченка \par}
{\parbox {5 cm}{Освітня програма \hfill}{\parbox {7 cm}{Прикладна математика, Системний аналіз \hfill}}\parbox {6 cm}{\hfill Семестр 2}}\par
{\parbox {5 cm} {Навчальний предмет \hfill}\parbox {7 cm}{Програмування \hfill}\parbox {6cm}{\hfill}\par}\vspace{-7pt} \end{small}}{\center Екзаменаційний білет \No \textbf 
{ 177 }\par}}
{
\begin {large}
1. Що буде виведено за виконання?
code
try:
    res = 8.0!=True
    if isinstance(res, bool):
        print(f'bool;{res}')
    elif isinstance(res, int):
        print(f'int;{res}')
    elif isinstance(res, float):
        print(f'float;{res:.2f}')
    else:
        print('???')
except:
    print('ERR')
code

2. Що буде виведено за виконання?
code
a = 6
b = 5
c = 0

def g(b):
    global c
    a = 2
    b = 4
    c = 5
    return a + b + c

a = 4
b = 3
c = 8

try:
    print(g(b), a, b, c, sep=';')
except:
    print('ERR', a, b, c, sep=';')
code

3. Що буде виведено за виконання?
code
try:
    j = 12
    while j>=3:
        j += -3
    print(j, end=';')
except:
    print('ERR')
code

4. Що буде виведено за виконання?
code
try:
    for f in range(8, 0, 3):
        if f>=6:
            continue
        print(f, end=';')
    else:
        print(f,end=';')
    print(f)
except:
    print('ERR')
code

5. Що буде виведено за виконання?\par (Дійсні значення, якщо наявні, в цьому завданні будуть виводитись у форматі з фіксованою крапкою та, можливо, з доволі великою кількістю десяткових розрядів. У відповіді для дійсних значень у ЦЬОМУ завданні достатньо вказати один десятковий розряд після крапки, відкинувши решту розрядів без округлення. Наприклад, замість 13.66666666666667 достатньо вказати 13.6, замість 25.23 достатньо вказати 25.2)
code
try:
    print(2, end=';')
    print(1/0.0, end=';')
    print(9, end=';')
except Exception:
    print(3, end=';')
except TypeError:
    print(8, end=';')
except:
    print('ERR', end=';')
else:
    print(4, end=';')
finally:
    print(7)
code

6. Що буде виведено за виконання?
code
def f(n):
    print(n%10, end='')
    if n>9:
        f(n//100)
 
f(987654)
code

7. Що буде виведено за виконання?
code
class A:
    d = 1

    def g1(self): print(2, end=';')
    def _g2(self): print(3, end=';')

    def main(self):
        try: print(self.d, end=';')
        except: print('ERR', end=';')

        try: self.g1()
        except: print('ERR', end=';')

        try: self._g2()
        except: print('ERR', end=';')


class B(A):
    d = 4

    def g1(self): print(5, end=';')
    def _g2(self): print(6, end=';')

try: B().main()
except: print('ERR', end=';')
code

8. Що буде виведено за виконання?
code
def f(arg, n):
    arg[19] = 9
    n -= -3

try:
    n = 7
    s = {45:8, 67:3, 15:8}
    f(s, n)
    print(n, end=';')
    for x in s.values():
        print(x, sep=';', end=';')
except:
    print('ERR')
code

9. Що буде виведено за виконання?
code
class A:
    b = 1
    a = 2

    def __init__(self):
        b = 3

class B(A):
    b = 4

    def __init__(self):
        super().__init__()
        self.c = 7

obj = B()
obj.c = obj.b + 10
B.c = 13

try: print(obj.c, end= ';')
except: print('ERR', end=';')

try: print(A.c, end= ';')
except: print('ERR', end=';')

try: print(B.c, end= ';')
except: print('ERR', end=';')
code

10. 
Для графа на рисунку з вершини 3 запущено алгоритм пошуку в глибину, що будує кістякове дерево. Списки суміжності вершин переглядаються алгоритмом за зростанням міток вершин.\par
\begin{center}
	\adjustimage{max size={0.9\linewidth}{0.9\paperheight}}
	{../10/Traversals/48.png}
\end{center}
\par Які з наведених ребер входять до складу отриманого кістякового дерева?\par
1) (2, 6)\par
2) (3, 4)\par
3) (0, 2)\par
4) (4, 6)\par
5) (0, 5)\par
6) (5, 6)\par
{\vskip 15pt} У формі вибрати номери відповідних ребер.\par
%answer:3 4 6
\end {large}}
%end
{\smallskip \begin {small} Затверджено на засіданні кафедри теоретичної кібернетики, протокол \No 12 від   19.05.2020 р.\par\smallskip\smallskip
{\parbox {6.5 cm}{Зав. кафедри \dotfill Ю.В.~Крак}}\hspace {0.7cm}{\hbox to 6.5 cm{ Екзаменатор \dotfill Т.О.~Карнаух}} \end{small}}
%begin
{{\begin {small}\par
{ \center  Київський національний університет імені Тараса Шевченка \par}
{\parbox {5 cm}{Освітня програма \hfill}{\parbox {7 cm}{Прикладна математика, Системний аналіз \hfill}}\parbox {6 cm}{\hfill Семестр 2}}\par
{\parbox {5 cm} {Навчальний предмет \hfill}\parbox {7 cm}{Програмування \hfill}\parbox {6cm}{\hfill}\par}\vspace{-7pt} \end{small}}{\center Екзаменаційний білет \No \textbf 
{ 178 }\par}}
{
\begin {large}
1. Що буде виведено за виконання?
code
try:
    res = "False"<"9.0j"
    if isinstance(res, bool):
        print(f'bool;{res}')
    elif isinstance(res, int):
        print(f'int;{res}')
    elif isinstance(res, float):
        print(f'float;{res:.2f}')
    else:
        print('???')
except:
    print('ERR')
code

2. Що буде виведено за виконання?
code
a = 7
b = 9
c = 7

def h(b):
    global c
    a = 3
    b = 4
    c = 2
    return a + b + c

a = 0
b = 9
c = 1

try:
    print(h(a), a, b, c, sep=';')
except:
    print('ERR', a, b, c, sep=';')
code

3. Що буде виведено за виконання?
code
try:
    n = 8
    while n>4:
        n += -3
    print(n, end=';')
except:
    print('ERR')
code

4. Що буде виведено за виконання?
code
try:
    for b in range(7, -15, -1):
        if b>8:
            continue
        print(b, end=';')
    else:
        print(b,end=';')
    print(b)
except:
    print('ERR')
code

5. Що буде виведено за виконання?\par (Дійсні значення, якщо наявні, в цьому завданні будуть виводитись у форматі з фіксованою крапкою та, можливо, з доволі великою кількістю десяткових розрядів. У відповіді для дійсних значень у ЦЬОМУ завданні достатньо вказати один десятковий розряд після крапки, відкинувши решту розрядів без округлення. Наприклад, замість 13.66666666666667 достатньо вказати 13.6, замість 25.23 достатньо вказати 25.2)
code
try:
    print(0, end=';')
    print(0j<=0j, end=';')
    print(4, end=';')
except ValueError:
    print(3, end=';')
except Exception:
    print(8, end=';')
except:
    print('ERR', end=';')
else:
    print(7, end=';')
finally:
    print(2)
code

6. Що буде виведено за виконання?
code
def f(n):
    if n>9:
        f(n//10)
    else:
        print(n%10, end='')

f(543216)
code

7. Що буде виведено за виконання?
code
class A:
    c = 1

    def __h1(self): print(2, end=';')
    def _h2(self): print(3, end=';')

    def main(self):
        try: print(self.c, end=';')
        except: print('ERR', end=';')

        try: self.__h1()
        except: print('ERR', end=';')

        try: self._h2()
        except: print('ERR', end=';')


class B(A):
    c = 4

    def __h1(self): print(5, end=';')
    def _h2(self): print(6, end=';')

try: B().main()
except: print('ERR', end=';')
code

8. Що буде виведено за виконання?
code
def f(arg, n):
    arg[68] = 6
    n = 7
    arg = set(arg.keys())

try:
    n = 4
    s = {29:6, 67:0, 48:6, 48:7}
    f(s, n)
    print(n, end=';')
    for x in s.values():
        print(x, sep=';', end=';')
except:
    print('ERR')
code

9. Що буде виведено за виконання?
code
class A:
    a = 1
    b = 2

    def __init__(self):
        b = 3

class B(A):
    b = 4

    def __init__(self):
        super().__init__()
        self.c = 7

obj = B()
obj.c = obj.b + 10
B.c = 13

try: print(obj.c, end= ';')
except: print('ERR', end=';')

try: print(A.c, end= ';')
except: print('ERR', end=';')

try: print(B.c, end= ';')
except: print('ERR', end=';')
code

10. 
Для графа на рисунку з вершини 3 запущено алгоритм пошуку в глибину, що будує кістякове дерево. Списки суміжності вершин переглядаються алгоритмом за зростанням міток вершин.\par
\begin{center}
	\adjustimage{max size={0.9\linewidth}{0.9\paperheight}}
	{../10/Traversals/93.png}
\end{center}
\par Які з наведених ребер входять до складу отриманого кістякового дерева?\par
1) (1, 6)\par
2) (4, 5)\par
3) (0, 2)\par
4) (0, 5)\par
5) (1, 5)\par
6) (1, 3)\par
{\vskip 15pt} У формі вибрати номери відповідних ребер.\par
%answer:2 3 4 5 6
\end {large}}
%end
{\smallskip \begin {small} Затверджено на засіданні кафедри теоретичної кібернетики, протокол \No 12 від   19.05.2020 р.\par\smallskip\smallskip
{\parbox {6.5 cm}{Зав. кафедри \dotfill Ю.В.~Крак}}\hspace {0.7cm}{\hbox to 6.5 cm{ Екзаменатор \dotfill Т.О.~Карнаух}} \end{small}}
%begin
{{\begin {small}\par
{ \center  Київський національний університет імені Тараса Шевченка \par}
{\parbox {5 cm}{Освітня програма \hfill}{\parbox {7 cm}{Прикладна математика, Системний аналіз \hfill}}\parbox {6 cm}{\hfill Семестр 2}}\par
{\parbox {5 cm} {Навчальний предмет \hfill}\parbox {7 cm}{Програмування \hfill}\parbox {6cm}{\hfill}\par}\vspace{-7pt} \end{small}}{\center Екзаменаційний білет \No \textbf 
{ 179 }\par}}
{
\begin {large}
1. Що буде виведено за виконання?
code
try:
    res = 7//3*3*5
    if isinstance(res, bool):
        print(f'bool;{res}')
    elif isinstance(res, int):
        print(f'int;{res}')
    elif isinstance(res, float):
        print(f'float;{res:.2f}')
    else:
        print('???')
except:
    print('ERR')
code

2. Що буде виведено за виконання?
code
a = 6
b = 4
c = 9

def h(b):
    global c
    a *= 1
    b = 3
    c = 4
    return a + b + c

a = 8
b = 0
c = 1

try:
    print(h(a), a, b, c, sep=';')
except:
    print('ERR', a, b, c, sep=';')
code

3. Що буде виведено за виконання?
code
try:
    t = 8
    while t<=7:
        t += 1
    print(t, end=';')
except:
    print('ERR')
code

4. Що буде виведено за виконання?
code
try:
    for d in range(-11, 1, -3):
        if d<=5:
            continue
        print(d, end=';')
        if d<4:
            break
    else:
        print('end',end=';')
    print(d)
except:
    print('ERR')
code

5. Що буде виведено за виконання?\par (Дійсні значення, якщо наявні, в цьому завданні будуть виводитись у форматі з фіксованою крапкою та, можливо, з доволі великою кількістю десяткових розрядів. У відповіді для дійсних значень у ЦЬОМУ завданні достатньо вказати один десятковий розряд після крапки, відкинувши решту розрядів без округлення. Наприклад, замість 13.66666666666667 достатньо вказати 13.6, замість 25.23 достатньо вказати 25.2)
code
try:
    print(9, end=';')
    print(int('b6'), end=';')
    print(5, end=';')
except ValueError:
    print(1, end=';')
except TypeError:
    print(7, end=';')
except:
    print('ERR', end=';')
else:
    print(0, end=';')
finally:
    print(2)
code

6. Що буде виведено за виконання?
code
def f(n):
    if n>9: f(n//10)
    else:
        print(n%10, end='')
 
f(123456)
code

7. Що буде виведено за виконання?
code
class A:
    a = 1

    def _g1(self): print(2, end=';')
    def __g2(self): print(3, end=';')

    def main(self):
        try: print(self.a, end=';')
        except: print('ERR', end=';')

        try: self._g1()
        except: print('ERR', end=';')

        try: self.__g2()
        except: print('ERR', end=';')


class B(A):
    a = 4

    def _g1(self): print(5, end=';')
    def __g2(self): print(6, end=';')

try: B().main()
except: print('ERR', end=';')
code

8. Що буде виведено за виконання?
code
def f(arg, n):
    n = 7
    tmp = arg
    tmp[8:] = [4,5]
    return len(tmp)>3

try:
    a = ['0','1','2','3','4']
    n = 0
    f(a, n)
    print(n, end=';')
    for el in a:
        print(el, end=';')
except:
    print('ERR')
code

9. Що буде виведено за виконання?
code
class A:
    b = 1
    c = 2

    def __init__(self):
        c = 3

class B(A):
    b = 4

    def __init__(self):
        super().__init__()
        self.a = 7

obj = B()
obj.a = obj.b + 10
B.a = 13

try: print(obj.a, end= ';')
except: print('ERR', end=';')

try: print(A.a, end= ';')
except: print('ERR', end=';')

try: print(B.a, end= ';')
except: print('ERR', end=';')
code

10. 
Для графа на рисунку з вершини 0 запущено алгоритм пошуку в глибину, що будує кістякове дерево. Списки суміжності вершин переглядаються алгоритмом за зростанням міток вершин.\par
\begin{center}
	\adjustimage{max size={0.9\linewidth}{0.9\paperheight}}
	{../10/Traversals/178.png}
\end{center}
\par Які з наведених ребер входять до складу отриманого кістякового дерева?\par
1) (0, 1)\par
2) (0, 6)\par
3) (3, 6)\par
4) (4, 6)\par
5) (0, 4)\par
6) (1, 6)\par
{\vskip 15pt} У формі вибрати номери відповідних ребер.\par
%answer:1 3 4
\end {large}}
%end
{\smallskip \begin {small} Затверджено на засіданні кафедри теоретичної кібернетики, протокол \No 12 від   19.05.2020 р.\par\smallskip\smallskip
{\parbox {6.5 cm}{Зав. кафедри \dotfill Ю.В.~Крак}}\hspace {0.7cm}{\hbox to 6.5 cm{ Екзаменатор \dotfill Т.О.~Карнаух}} \end{small}}
%begin
{{\begin {small}\par
{ \center  Київський національний університет імені Тараса Шевченка \par}
{\parbox {5 cm}{Освітня програма \hfill}{\parbox {7 cm}{Прикладна математика, Системний аналіз \hfill}}\parbox {6 cm}{\hfill Семестр 2}}\par
{\parbox {5 cm} {Навчальний предмет \hfill}\parbox {7 cm}{Програмування \hfill}\parbox {6cm}{\hfill}\par}\vspace{-7pt} \end{small}}{\center Екзаменаційний білет \No \textbf 
{ 180 }\par}}
{
\begin {large}
1. Що буде виведено за виконання?
code
try:
    res = 3<8.0 or not False>=5.0 and 7!=6
    if isinstance(res, bool):
        print(f'bool;{res}')
    elif isinstance(res, int):
        print(f'int;{res}')
    elif isinstance(res, float):
        print(f'float;{res:.2f}')
    else:
        print('???')
except:
    print('ERR')
code

2. Що буде виведено за виконання?
code
a = 8
b = 3
c = 6

def f(a):
    a = 3
    b = 1
    c = 5
    return a + b + c

a = 2
b = 5
c = 4

try:
    print(f(a), a, b, c, sep=';')
except:
    print('ERR', a, b, c, sep=';')
code

3. Що буде виведено за виконання?
code
try:
    t = 6
    while t>=0:
        t += -2
    print(t, end=';')
except:
    print('ERR')
code

4. Що буде виведено за виконання?
code
try:
    for e in range(14, -4, -2):
        if e<3:
            break
        print(e, end=';');e = -10
    else:
        print(e,end=';')
    print(e)
except:
    print('ERR')
code

5. Що буде виведено за виконання?\par (Дійсні значення, якщо наявні, в цьому завданні будуть виводитись у форматі з фіксованою крапкою та, можливо, з доволі великою кількістю десяткових розрядів. У відповіді для дійсних значень у ЦЬОМУ завданні достатньо вказати один десятковий розряд після крапки, відкинувши решту розрядів без округлення. Наприклад, замість 13.66666666666667 достатньо вказати 13.6, замість 25.23 достатньо вказати 25.2)
code
try:
    print(8, end=';')
    print(int('1'), end=';')
    print(3, end=';')
except KeyboardInterrupt:
    print(4, end=';')
except ZeroDivisionError:
    print(6, end=';')
except:
    print('ERR', end=';')
else:
    print(9, end=';')
finally:
    print(5)
code

6. Що буде виведено за виконання?
code
def f(s):
    for el in s:
        el += 1
        return s
 
s = [1, 2, 3]
s = f(s)
print(s)
code

7. Що буде виведено за виконання?
code
class A:
    b = 1

    def _f1(self): print(2, end=';')
    def _f2(self): print(3, end=';')

    def main(self):
        try: print(self.b, end=';')
        except: print('ERR', end=';')

        try: self._f1()
        except: print('ERR', end=';')

        try: self._f2()
        except: print('ERR', end=';')


class B(A):
    b = 4

    def _f1(self): print(5, end=';')
    def _f2(self): print(6, end=';')

try: B().main()
except: print('ERR', end=';')
code

8. Що буде виведено за виконання?
code
def f(arg, n):
    n = 9
    tmp = arg
    tmp[-6:-6] = [14]
    return len(tmp)>3

try:
    f = [10,11,12,13,14]
    n = 8
    f(f, n)
    print(n, end=';')
    for el in f:
        print(el, end=';')
except:
    print('ERR')
code

9. Що буде виведено за виконання?
code
class A:
    c = 1
    a = 2

    def __init__(self):
        c = 3

class B(A):
    a = 4

    def __init__(self):
        super().__init__()
        self.b = 7

obj = B()
obj.c = obj.a + 10
B.c = 13

try: print(obj.c, end= ';')
except: print('ERR', end=';')

try: print(A.c, end= ';')
except: print('ERR', end=';')

try: print(B.c, end= ';')
except: print('ERR', end=';')
code

10. 
Для графа на рисунку з вершини 5 запущено алгоритм пошуку в глибину, що будує кістякове дерево. Списки суміжності вершин переглядаються алгоритмом за зростанням міток вершин.\par
\begin{center}
	\adjustimage{max size={0.9\linewidth}{0.9\paperheight}}
	{../10/Traversals/181.png}
\end{center}
\par Які з наведених ребер входять до складу отриманого кістякового дерева?\par
1) (3, 5)\par
2) (1, 4)\par
3) (0, 5)\par
4) (2, 3)\par
5) (1, 6)\par
6) (2, 4)\par
{\vskip 15pt} У формі вибрати номери відповідних ребер.\par
%answer:2 3 4 5 6
\end {large}}
%end
{\smallskip \begin {small} Затверджено на засіданні кафедри теоретичної кібернетики, протокол \No 12 від   19.05.2020 р.\par\smallskip\smallskip
{\parbox {6.5 cm}{Зав. кафедри \dotfill Ю.В.~Крак}}\hspace {0.7cm}{\hbox to 6.5 cm{ Екзаменатор \dotfill Т.О.~Карнаух}} \end{small}}
%begin
{{\begin {small}\par
{ \center  Київський національний університет імені Тараса Шевченка \par}
{\parbox {5 cm}{Освітня програма \hfill}{\parbox {7 cm}{Прикладна математика, Системний аналіз \hfill}}\parbox {6 cm}{\hfill Семестр 2}}\par
{\parbox {5 cm} {Навчальний предмет \hfill}\parbox {7 cm}{Програмування \hfill}\parbox {6cm}{\hfill}\par}\vspace{-7pt} \end{small}}{\center Екзаменаційний білет \No \textbf 
{ 181 }\par}}
{
\begin {large}
1. Що буде виведено за виконання?
code
try:
    res = 4**-1-True
    if isinstance(res, bool):
        print(f'bool;{res}')
    elif isinstance(res, int):
        print(f'int;{res}')
    elif isinstance(res, float):
        print(f'float;{res:.2f}')
    else:
        print('???')
except:
    print('ERR')
code

2. Що буде виведено за виконання?
code
a = 8
b = 0
c = 1

def f(a):
    a = 5
    b = 2
    c = 3
    return a + b + c

a = 7
b = 6
c = 9

try:
    print(f(b), a, b, c, sep=';')
except:
    print('ERR', a, b, c, sep=';')
code

3. Що буде виведено за виконання?
code
try:
    j = 3
    while j<10:
        j += 1
    print(j, end=';')
except:
    print('ERR')
code

4. Що буде виведено за виконання?
code
try:
    for d in range(13, 3, -2):
        if d<5:
            continue
        print(d, end=';');d = -3
    else:
        print(d,end=';')
    print(d)
except:
    print('ERR')
code

5. Що буде виведено за виконання?\par (Дійсні значення, якщо наявні, в цьому завданні будуть виводитись у форматі з фіксованою крапкою та, можливо, з доволі великою кількістю десяткових розрядів. У відповіді для дійсних значень у ЦЬОМУ завданні достатньо вказати один десятковий розряд після крапки, відкинувши решту розрядів без округлення. Наприклад, замість 13.66666666666667 достатньо вказати 13.6, замість 25.23 достатньо вказати 25.2)
code
try:
    print(9, end=';')
    print(5%1, end=';')
    print(4, end=';')
except KeyboardInterrupt:
    print(2, end=';')
except ZeroDivisionError:
    print(3, end=';')
except:
    print('ERR', end=';')
else:
    print(1, end=';')
finally:
    print(0)
code

6. Що буде виведено за виконання?
code
def f(n):
    print(n%10, end='')
    if n>9:
        f(n//10)
 
f(123456)
code

7. Що буде виведено за виконання?
code
class A:
    c = 1

    def _f1(self): print(2, end=';')
    def f2(self): print(3, end=';')

    def main(self):
        try: print(self.c, end=';')
        except: print('ERR', end=';')

        try: self._f1()
        except: print('ERR', end=';')

        try: self.f2()
        except: print('ERR', end=';')


class B(A):
    c = 4

    def _f1(self): print(5, end=';')
    def f2(self): print(6, end=';')

try: B().main()
except: print('ERR', end=';')
code

8. Що буде виведено за виконання?
code
def f(arg, n):
    n = 8
    tmp = arg
    del tmp[-9:-9]
    return len(tmp)>3

try:
    h = ['0','1','2','3','4']
    n = 1
    f(h, n)
    print(n, end=';')
    for el in h:
        print(el, end=';')
except:
    print('ERR')
code

9. Що буде виведено за виконання?
code
class A:
    b = 1
    c = 2

    def __init__(self):
        c = 3

class B(A):
    c = 4

    def __init__(self):
        super().__init__()
        self.a = 7

obj = B()
obj.b = obj.c + 10
B.b = 13

try: print(obj.b, end= ';')
except: print('ERR', end=';')

try: print(A.b, end= ';')
except: print('ERR', end=';')

try: print(B.b, end= ';')
except: print('ERR', end=';')
code

10. 
Для графа на рисунку з вершини 2 запущено алгоритм пошуку в глибину, що будує кістякове дерево. Списки суміжності вершин переглядаються алгоритмом за зростанням міток вершин.\par
\begin{center}
	\adjustimage{max size={0.9\linewidth}{0.9\paperheight}}
	{../10/Traversals/275.png}
\end{center}
\par Які з наведених ребер входять до складу отриманого кістякового дерева?\par
1) (2, 4)\par
2) (2, 3)\par
3) (4, 6)\par
4) (2, 5)\par
5) (1, 2)\par
6) (0, 1)\par
{\vskip 15pt} У формі вибрати номери відповідних ребер.\par
%answer:3 6
\end {large}}
%end
{\smallskip \begin {small} Затверджено на засіданні кафедри теоретичної кібернетики, протокол \No 12 від   19.05.2020 р.\par\smallskip\smallskip
{\parbox {6.5 cm}{Зав. кафедри \dotfill Ю.В.~Крак}}\hspace {0.7cm}{\hbox to 6.5 cm{ Екзаменатор \dotfill Т.О.~Карнаух}} \end{small}}
%begin
{{\begin {small}\par
{ \center  Київський національний університет імені Тараса Шевченка \par}
{\parbox {5 cm}{Освітня програма \hfill}{\parbox {7 cm}{Прикладна математика, Системний аналіз \hfill}}\parbox {6 cm}{\hfill Семестр 2}}\par
{\parbox {5 cm} {Навчальний предмет \hfill}\parbox {7 cm}{Програмування \hfill}\parbox {6cm}{\hfill}\par}\vspace{-7pt} \end{small}}{\center Екзаменаційний білет \No \textbf 
{ 182 }\par}}
{
\begin {large}
1. Що буде виведено за виконання?
code
def f(n):
    print('f', end='')
    return n==1

try:
    print(f(9) or f(7))
except:
    print('ERR')
code

2. Що буде виведено за виконання?
code
a = 4
b = 2
c = 3

def g(b):
    global c
    a = 1
    b = 4
    c = 1
    return a + b + c

a = 5
b = 7
c = 5

try:
    print(g(a), a, b, c, sep=';')
except:
    print('ERR', a, b, c, sep=';')
code

3. Що буде виведено за виконання?
code
try:
    k = 5
    while k>4:
        k += -3
    print(k, end=';')
except:
    print('ERR')
code

4. Що буде виведено за виконання?
code
try:
    for c in range(6, 7, 2):
        if c<=4:
            break
        print(c, end=';')
    else:
        print(c,end=';')
    print(c)
except:
    print('ERR')
code

5. Що буде виведено за виконання?\par (Дійсні значення, якщо наявні, в цьому завданні будуть виводитись у форматі з фіксованою крапкою та, можливо, з доволі великою кількістю десяткових розрядів. У відповіді для дійсних значень у ЦЬОМУ завданні достатньо вказати один десятковий розряд після крапки, відкинувши решту розрядів без округлення. Наприклад, замість 13.66666666666667 достатньо вказати 13.6, замість 25.23 достатньо вказати 25.2)
code
try:
    print(8, end=';')
    print(int('6'), end=';')
    print(7, end=';')
except TypeError:
    print(1, end=';')
except Exception:
    print(0, end=';')
except:
    print('ERR', end=';')
else:
    print(6, end=';')
finally:
    print(4)
code

6. Що буде виведено за виконання?
code
def f(n):
    if n>9:
        f(n//100)
    else:
        print(n%10,end='')

f(987654)
code

7. Що буде виведено за виконання?
code
class A:
    d = 1

    def _g1(self): print(2, end=';')
    def __g2(self): print(3, end=';')

    def main(self):
        try: print(self.d, end=';')
        except: print('ERR', end=';')

        try: self._g1()
        except: print('ERR', end=';')

        try: self.__g2()
        except: print('ERR', end=';')


class B(A):
    d = 4

    def _g1(self): print(5, end=';')
    def __g2(self): print(6, end=';')

try: B().main()
except: print('ERR', end=';')
code

8. Що буде виведено за виконання?
code
def f(arg, n):
    arg[46] = 0
    n *= 2

try:
    n = 7
    s = {11:7, 55:7, 89:1, 11:5}
    f(s, n)
    print(n, end=';')
    for x in s :
        print(x, sep=';', end=';')
except:
    print('ERR')
code

9. Що буде виведено за виконання?
code
class A:
    a = 1
    c = 2

    def __init__(self):
        a = 3

class B(A):
    a = 4

    def __init__(self):
        super().__init__()
        self.b = 7

obj = B()
obj.a = obj.a + 10
B.a = 13

try: print(obj.a, end= ';')
except: print('ERR', end=';')

try: print(A.a, end= ';')
except: print('ERR', end=';')

try: print(B.a, end= ';')
except: print('ERR', end=';')
code

10. 
Для графа на рисунку з вершини 5 запущено алгоритм пошуку в глибину, що будує кістякове дерево. Списки суміжності вершин переглядаються алгоритмом за зростанням міток вершин.\par
\begin{center}
	\adjustimage{max size={0.9\linewidth}{0.9\paperheight}}
	{../10/Traversals/294.png}
\end{center}
\par Які з наведених ребер входять до складу отриманого кістякового дерева?\par
1) (0, 5)\par
2) (0, 4)\par
3) (0, 6)\par
4) (3, 5)\par
5) (2, 6)\par
6) (3, 4)\par
{\vskip 15pt} У формі вибрати номери відповідних ребер.\par
%answer:1 2 5 6
\end {large}}
%end
{\smallskip \begin {small} Затверджено на засіданні кафедри теоретичної кібернетики, протокол \No 12 від   19.05.2020 р.\par\smallskip\smallskip
{\parbox {6.5 cm}{Зав. кафедри \dotfill Ю.В.~Крак}}\hspace {0.7cm}{\hbox to 6.5 cm{ Екзаменатор \dotfill Т.О.~Карнаух}} \end{small}}
%begin
{{\begin {small}\par
{ \center  Київський національний університет імені Тараса Шевченка \par}
{\parbox {5 cm}{Освітня програма \hfill}{\parbox {7 cm}{Прикладна математика, Системний аналіз \hfill}}\parbox {6 cm}{\hfill Семестр 2}}\par
{\parbox {5 cm} {Навчальний предмет \hfill}\parbox {7 cm}{Програмування \hfill}\parbox {6cm}{\hfill}\par}\vspace{-7pt} \end{small}}{\center Екзаменаційний білет \No \textbf 
{ 183 }\par}}
{
\begin {large}
1. Що буде виведено за виконання?
code
try:
    res = 7j>"4"
    if isinstance(res, bool):
        print(f'bool;{res}')
    elif isinstance(res, int):
        print(f'int;{res}')
    elif isinstance(res, float):
        print(f'float;{res:.2f}')
    else:
        print('???')
except:
    print('ERR')
code

2. Що буде виведено за виконання?
code
a = 2
b = 1
c = 4

def h(b):
    global c
    a = 4
    b = 2
    c = 5
    return a + b + c

a = 3
b = 9
c = 6

try:
    print(h(a), a, b, c, sep=';')
except:
    print('ERR', a, b, c, sep=';')
code

3. Що буде виведено за виконання?
code
try:
    n = 7
    while n<=12:
        n += 2
    print(n, end=';')
except:
    print('ERR')
code

4. Що буде виведено за виконання?
code
try:
    for b in range(3, 15, -1):
        if b>=6:
            continue
        print(b, end=';')
    else:
        print(b,end=';')
    print(b)
except:
    print('ERR')
code

5. Що буде виведено за виконання?\par (Дійсні значення, якщо наявні, в цьому завданні будуть виводитись у форматі з фіксованою крапкою та, можливо, з доволі великою кількістю десяткових розрядів. У відповіді для дійсних значень у ЦЬОМУ завданні достатньо вказати один десятковий розряд після крапки, відкинувши решту розрядів без округлення. Наприклад, замість 13.66666666666667 достатньо вказати 13.6, замість 25.23 достатньо вказати 25.2)
code
try:
    print(7, end=';')
    print(8j>=8j, end=';')
    print(5, end=';')
except ValueError:
    print(3, end=';')
except Exception:
    print(9, end=';')
except:
    print('ERR', end=';')
else:
    print(2, end=';')
finally:
    print(5)
code

6. Що буде виведено за виконання?
code
def f(n):
    if n>9:
        f(n//100)
    else:
        print(n%10,end='')

f(123456)
code

7. Що буде виведено за виконання?
code
class A:
    a = 1

    def h1(self): print(2, end=';')
    def _h2(self): print(3, end=';')

    def main(self):
        try: print(self.a, end=';')
        except: print('ERR', end=';')

        try: self.h1()
        except: print('ERR', end=';')

        try: self._h2()
        except: print('ERR', end=';')


class B(A):
    a = 4

    def h1(self): print(5, end=';')
    def _h2(self): print(6, end=';')

try: B().main()
except: print('ERR', end=';')
code

8. Що буде виведено за виконання?
code
def f(arg, n):
    n = 7
    tmp = arg
    tmp[:0] = 'xy'
    return len(tmp)>3

try:
    c = [10,11,12,13,14]
    n = 4
    f(c, n)
    print(n, end=';')
    for el in c:
        print(el, end=';')
except:
    print('ERR')
code

9. Що буде виведено за виконання?
code
class A:
    b = 1
    a = 2

    def __init__(self):
        b = 3

class B(A):
    a = 4

    def __init__(self):
        super().__init__()
        self.c = 7

obj = B()
obj.c = obj.b + 10
B.c = 13

try: print(obj.c, end= ';')
except: print('ERR', end=';')

try: print(A.c, end= ';')
except: print('ERR', end=';')

try: print(B.c, end= ';')
except: print('ERR', end=';')
code

10. 
Для графа на рисунку з вершини 3 запущено алгоритм пошуку в глибину, що будує кістякове дерево. Списки суміжності вершин переглядаються алгоритмом за зростанням міток вершин.\par
\begin{center}
	\adjustimage{max size={0.9\linewidth}{0.9\paperheight}}
	{../10/Traversals/118.png}
\end{center}
\par Які з наведених ребер входять до складу отриманого кістякового дерева?\par
1) (1, 3)\par
2) (0, 4)\par
3) (3, 5)\par
4) (0, 3)\par
5) (1, 4)\par
6) (5, 6)\par
{\vskip 15pt} У формі вибрати номери відповідних ребер.\par
%answer:4 5 6
\end {large}}
%end
{\smallskip \begin {small} Затверджено на засіданні кафедри теоретичної кібернетики, протокол \No 12 від   19.05.2020 р.\par\smallskip\smallskip
{\parbox {6.5 cm}{Зав. кафедри \dotfill Ю.В.~Крак}}\hspace {0.7cm}{\hbox to 6.5 cm{ Екзаменатор \dotfill Т.О.~Карнаух}} \end{small}}
%begin
{{\begin {small}\par
{ \center  Київський національний університет імені Тараса Шевченка \par}
{\parbox {5 cm}{Освітня програма \hfill}{\parbox {7 cm}{Прикладна математика, Системний аналіз \hfill}}\parbox {6 cm}{\hfill Семестр 2}}\par
{\parbox {5 cm} {Навчальний предмет \hfill}\parbox {7 cm}{Програмування \hfill}\parbox {6cm}{\hfill}\par}\vspace{-7pt} \end{small}}{\center Екзаменаційний білет \No \textbf 
{ 184 }\par}}
{
\begin {large}
1. Що буде виведено за виконання?
code
def f(n):
    print('f', end='')
    return n

try:
    print(f(5) and f(0))
except:
    print('ERR')
code

2. Що буде виведено за виконання?
code
a = 8
b = 0
c = 8

def g(a):
    a = 3
    b = 5
    c = 3
    return a + b + c

a = 4
b = 9
c = 2

try:
    print(g(b), a, b, c, sep=';')
except:
    print('ERR', a, b, c, sep=';')
code

3. Що буде виведено за виконання?
code
try:
    i = 7
    while i>=1:
        i += -2
    print(i, end=';')
except:
    print('ERR')
code

4. Що буде виведено за виконання?
code
try:
    for a in range(-8, 5, 3):
        if a>3:
            continue
        print(a, end=';')
    else:
        print('end',end=';')
    print(a)
except:
    print('ERR')
code

5. Що буде виведено за виконання?\par (Дійсні значення, якщо наявні, в цьому завданні будуть виводитись у форматі з фіксованою крапкою та, можливо, з доволі великою кількістю десяткових розрядів. У відповіді для дійсних значень у ЦЬОМУ завданні достатньо вказати один десятковий розряд після крапки, відкинувши решту розрядів без округлення. Наприклад, замість 13.66666666666667 достатньо вказати 13.6, замість 25.23 достатньо вказати 25.2)
code
try:
    print(9, end=';')
    print(2//0.0, end=';')
    print(3, end=';')
except ZeroDivisionError:
    print(1, end=';')
except TypeError:
    print(4, end=';')
except:
    print('ERR', end=';')
else:
    print(8, end=';')
finally:
    print(6)
code

6. Що буде виведено за виконання?
code
def f(s1):
    s = s1[:]
    for i in range(len(s)):
        s[i] += 1
 
s = [1, 2, 3]
f(s)
print(s)
code

7. Що буде виведено за виконання?
code
class A:
    b = 1

    def _g1(self): print(2, end=';')
    def _g2(self): print(3, end=';')

    def main(self):
        try: print(self.b, end=';')
        except: print('ERR', end=';')

        try: self._g1()
        except: print('ERR', end=';')

        try: self._g2()
        except: print('ERR', end=';')


class B(A):
    b = 4

    def _g1(self): print(5, end=';')
    def _g2(self): print(6, end=';')

try: B().main()
except: print('ERR', end=';')
code

8. Що буде виведено за виконання?
code
def f(arg, n):
    arg[68] = 9
    n += -3

try:
    n = 5
    t = {68:9, 27:9, 11:2, 11:9}
    f(t, n)
    print(n, end=';')
    for x in t.items():
        print(*x, sep=';', end=';')
except:
    print('ERR')
code

9. Що буде виведено за виконання?
code
class A:
    c = 1
    b = 2

    def __init__(self):
        b = 3

class B(A):
    c = 4

    def __init__(self):
        super().__init__()
        self.a = 7

obj = B()
obj.b = obj.c + 10
B.b = 13

try: print(obj.b, end= ';')
except: print('ERR', end=';')

try: print(A.b, end= ';')
except: print('ERR', end=';')

try: print(B.b, end= ';')
except: print('ERR', end=';')
code

10. 
Для графа на рисунку з вершини 5 запущено алгоритм пошуку в глибину, що будує кістякове дерево. Списки суміжності вершин переглядаються алгоритмом за зростанням міток вершин.\par
\begin{center}
	\adjustimage{max size={0.9\linewidth}{0.9\paperheight}}
	{../10/Traversals/154.png}
\end{center}
\par Які з наведених ребер входять до складу отриманого кістякового дерева?\par
1) (1, 3)\par
2) (0, 3)\par
3) (4, 6)\par
4) (1, 5)\par
5) (3, 4)\par
6) (2, 5)\par
{\vskip 15pt} У формі вибрати номери відповідних ребер.\par
%answer:2 3 4 5
\end {large}}
%end
{\smallskip \begin {small} Затверджено на засіданні кафедри теоретичної кібернетики, протокол \No 12 від   19.05.2020 р.\par\smallskip\smallskip
{\parbox {6.5 cm}{Зав. кафедри \dotfill Ю.В.~Крак}}\hspace {0.7cm}{\hbox to 6.5 cm{ Екзаменатор \dotfill Т.О.~Карнаух}} \end{small}}
%begin
{{\begin {small}\par
{ \center  Київський національний університет імені Тараса Шевченка \par}
{\parbox {5 cm}{Освітня програма \hfill}{\parbox {7 cm}{Прикладна математика, Системний аналіз \hfill}}\parbox {6 cm}{\hfill Семестр 2}}\par
{\parbox {5 cm} {Навчальний предмет \hfill}\parbox {7 cm}{Програмування \hfill}\parbox {6cm}{\hfill}\par}\vspace{-7pt} \end{small}}{\center Екзаменаційний білет \No \textbf 
{ 185 }\par}}
{
\begin {large}
1. Що буде виведено за виконання?
code
def f(n):
    print('f', end='')
    return n<=-3

try:
    print(f(-9) and f(-5))
except:
    print('ERR')
code

2. Що буде виведено за виконання?
code
a = 9
b = 4
c = 6

def f(b):
    global c
    a = 4
    b -= 1
    c = 3
    return a + b + c

a = 8
b = 3
c = 2

try:
    print(f(b), a, b, c, sep=';')
except:
    print('ERR', a, b, c, sep=';')
code

3. Що буде виведено за виконання?
code
try:
    t = 14
    while t>3:
        t += -3
    print(t, end=';')
except:
    print('ERR')
code

4. Що буде виведено за виконання?
code
try:
    for e in range(4, -10, -3):
        if e<7:
            continue
        print(e, end=';');e = 3
    else:
        print(e,end=';')
    print(e)
except:
    print('ERR')
code

5. Що буде виведено за виконання?\par (Дійсні значення, якщо наявні, в цьому завданні будуть виводитись у форматі з фіксованою крапкою та, можливо, з доволі великою кількістю десяткових розрядів. У відповіді для дійсних значень у ЦЬОМУ завданні достатньо вказати один десятковий розряд після крапки, відкинувши решту розрядів без округлення. Наприклад, замість 13.66666666666667 достатньо вказати 13.6, замість 25.23 достатньо вказати 25.2)
code
try:
    print(0, end=';')
    print(int('a7'), end=';')
    print(0, end=';')
except KeyboardInterrupt:
    print(2, end=';')
except ValueError:
    print(6, end=';')
except:
    print('ERR', end=';')
else:
    print(5, end=';')
finally:
    print(9)
code

6. Що буде виведено за виконання?
code
def f(n):
    if n>2:
        f(n//2)
    else:
        print(n%2, end='')

f(16)
code

7. Що буде виведено за виконання?
code
class A:
    a = 1

    def __f1(self): print(2, end=';')
    def _f2(self): print(3, end=';')

    def main(self):
        try: print(self.a, end=';')
        except: print('ERR', end=';')

        try: self.__f1()
        except: print('ERR', end=';')

        try: self._f2()
        except: print('ERR', end=';')


class B(A):
    a = 4

    def __f1(self): print(5, end=';')
    def _f2(self): print(6, end=';')

try: B().main()
except: print('ERR', end=';')
code

8. Що буде виведено за виконання?
code
def f(arg, n):
    arg[77] = 2
    n = 5
    arg = list(arg.values())

try:
    n = 2
    d = {35:5, 60:2, 77:8, 60:3}
    f(d, n)
    print(n, end=';')
    for x in d.keys():
        print(x, sep=';', end=';')
except:
    print('ERR')
code

9. Що буде виведено за виконання?
code
class A:
    c = 1
    a = 2

    def __init__(self):
        a = 3

class B(A):
    a = 4

    def __init__(self):
        super().__init__()
        self.b = 7

obj = B()
obj.a = obj.a + 10
B.a = 13

try: print(obj.a, end= ';')
except: print('ERR', end=';')

try: print(A.a, end= ';')
except: print('ERR', end=';')

try: print(B.a, end= ';')
except: print('ERR', end=';')
code

10. 
Для графа на рисунку з вершини 4 запущено алгоритм пошуку в глибину, що будує кістякове дерево. Списки суміжності вершин переглядаються алгоритмом за зростанням міток вершин.\par
\begin{center}
	\adjustimage{max size={0.9\linewidth}{0.9\paperheight}}
	{../10/Traversals/149.png}
\end{center}
\par Які з наведених ребер входять до складу отриманого кістякового дерева?\par
1) (2, 6)\par
2) (3, 6)\par
3) (1, 3)\par
4) (2, 3)\par
5) (1, 4)\par
6) (0, 1)\par
{\vskip 15pt} У формі вибрати номери відповідних ребер.\par
%answer:1 3 4 5 6
\end {large}}
%end
{\smallskip \begin {small} Затверджено на засіданні кафедри теоретичної кібернетики, протокол \No 12 від   19.05.2020 р.\par\smallskip\smallskip
{\parbox {6.5 cm}{Зав. кафедри \dotfill Ю.В.~Крак}}\hspace {0.7cm}{\hbox to 6.5 cm{ Екзаменатор \dotfill Т.О.~Карнаух}} \end{small}}
%begin
{{\begin {small}\par
{ \center  Київський національний університет імені Тараса Шевченка \par}
{\parbox {5 cm}{Освітня програма \hfill}{\parbox {7 cm}{Прикладна математика, Системний аналіз \hfill}}\parbox {6 cm}{\hfill Семестр 2}}\par
{\parbox {5 cm} {Навчальний предмет \hfill}\parbox {7 cm}{Програмування \hfill}\parbox {6cm}{\hfill}\par}\vspace{-7pt} \end{small}}{\center Екзаменаційний білет \No \textbf 
{ 186 }\par}}
{
\begin {large}
1. Що буде виведено за виконання?
code
try:
    res = not 8>3.0 and 9.0==6 or 4<=True
    if isinstance(res, bool):
        print(f'bool;{res}')
    elif isinstance(res, int):
        print(f'int;{res}')
    elif isinstance(res, float):
        print(f'float;{res:.2f}')
    else:
        print('???')
except:
    print('ERR')
code

2. Що буде виведено за виконання?
code
a = 1
b = 5
c = 7

def g(a):
    a = 2
    b += 5
    c = 2
    return a + b + c

a = 0
b = 5
c = 3

try:
    print(g(a), a, b, c, sep=';')
except:
    print('ERR', a, b, c, sep=';')
code

3. Що буде виведено за виконання?
code
try:
    m = 15
    while m>=8:
        m += -2
    print(m, end=';')
except:
    print('ERR')
code

4. Що буде виведено за виконання?
code
try:
    for f in range(-4, -14, -1):
        if f>=8:
            break
        print(f, end=';')
        if f>=8:
            break
    else:
        print(f,end=';')
    print(f)
except:
    print('ERR')
code

5. Що буде виведено за виконання?\par (Дійсні значення, якщо наявні, в цьому завданні будуть виводитись у форматі з фіксованою крапкою та, можливо, з доволі великою кількістю десяткових розрядів. У відповіді для дійсних значень у ЦЬОМУ завданні достатньо вказати один десятковий розряд після крапки, відкинувши решту розрядів без округлення. Наприклад, замість 13.66666666666667 достатньо вказати 13.6, замість 25.23 достатньо вказати 25.2)
code
try:
    print(4, end=';')
    print(int('3'), end=';')
    print(8, end=';')
except ZeroDivisionError:
    print(1, end=';')
except ValueError:
    print(7, end=';')
except:
    print('ERR', end=';')
else:
    print(9, end=';')
finally:
    print(8)
code

6. Що буде виведено за виконання?
code
def f(s):
    for el in s:
        el += 1
    return s
 
s = [1, 2, 3]
s = f(s)
print(s)
code

7. Що буде виведено за виконання?
code
class A:
    d = 1

    def _h1(self): print(2, end=';')
    def _h2(self): print(3, end=';')

    def main(self):
        try: print(self.d, end=';')
        except: print('ERR', end=';')

        try: self._h1()
        except: print('ERR', end=';')

        try: self._h2()
        except: print('ERR', end=';')


class B(A):
    d = 4

    def _h1(self): print(5, end=';')
    def _h2(self): print(6, end=';')

try: B().main()
except: print('ERR', end=';')
code

8. Що буде виведено за виконання?
code
def f(arg, n):
    n = 3
    tmp = arg
    tmp[-8:-1] = [4,5]
    return len(tmp)>3

try:
    b = ['0','1','2','3','4']
    n = 0
    f(b, n)
    print(n, end=';')
    for el in b:
        print(el, end=';')
except:
    print('ERR')
code

9. Що буде виведено за виконання?
code
class A:
    a = 1
    b = 2

    def __init__(self):
        a = 3

class B(A):
    a = 4

    def __init__(self):
        super().__init__()
        self.c = 7

obj = B()
obj.b = obj.c + 10
B.b = 13

try: print(obj.b, end= ';')
except: print('ERR', end=';')

try: print(A.b, end= ';')
except: print('ERR', end=';')

try: print(B.b, end= ';')
except: print('ERR', end=';')
code

10. 
Для графа на рисунку з вершини 3 запущено алгоритм пошуку в глибину, що будує кістякове дерево. Списки суміжності вершин переглядаються алгоритмом за зростанням міток вершин.\par
\begin{center}
	\adjustimage{max size={0.9\linewidth}{0.9\paperheight}}
	{../10/Traversals/192.png}
\end{center}
\par Які з наведених ребер входять до складу отриманого кістякового дерева?\par
1) (3, 5)\par
2) (1, 4)\par
3) (1, 2)\par
4) (3, 4)\par
5) (1, 5)\par
6) (0, 1)\par
{\vskip 15pt} У формі вибрати номери відповідних ребер.\par
%answer:2 6
\end {large}}
%end
{\smallskip \begin {small} Затверджено на засіданні кафедри теоретичної кібернетики, протокол \No 12 від   19.05.2020 р.\par\smallskip\smallskip
{\parbox {6.5 cm}{Зав. кафедри \dotfill Ю.В.~Крак}}\hspace {0.7cm}{\hbox to 6.5 cm{ Екзаменатор \dotfill Т.О.~Карнаух}} \end{small}}
%begin
{{\begin {small}\par
{ \center  Київський національний університет імені Тараса Шевченка \par}
{\parbox {5 cm}{Освітня програма \hfill}{\parbox {7 cm}{Прикладна математика, Системний аналіз \hfill}}\parbox {6 cm}{\hfill Семестр 2}}\par
{\parbox {5 cm} {Навчальний предмет \hfill}\parbox {7 cm}{Програмування \hfill}\parbox {6cm}{\hfill}\par}\vspace{-7pt} \end{small}}{\center Екзаменаційний білет \No \textbf 
{ 187 }\par}}
{
\begin {large}
1. Що буде виведено за виконання?
code
try:
    res = 1**1.0*False
    if isinstance(res, bool):
        print(f'bool;{res}')
    elif isinstance(res, int):
        print(f'int;{res}')
    elif isinstance(res, float):
        print(f'float;{res:.2f}')
    else:
        print('???')
except:
    print('ERR')
code

2. Що буде виведено за виконання?
code
a = 8
b = 0
c = 2

def h(a):
    global c
    a *= 3
    b = 5
    c = 4
    return a + b + c

a = 7
b = 9
c = 6

try:
    print(h(a), a, b, c, sep=';')
except:
    print('ERR', a, b, c, sep=';')
code

3. Що буде виведено за виконання?
code
try:
    k = 5
    while k<=6:
        k += 1
    print(k, end=';')
except:
    print('ERR')
code

4. Що буде виведено за виконання?
code
try:
    for b in range(0, 9, 2):
        if b<=6:
            break
        print(b, end=';');b = -2
    else:
        print(b,end=';')
    print(b)
except:
    print('ERR')
code

5. Що буде виведено за виконання?\par (Дійсні значення, якщо наявні, в цьому завданні будуть виводитись у форматі з фіксованою крапкою та, можливо, з доволі великою кількістю десяткових розрядів. У відповіді для дійсних значень у ЦЬОМУ завданні достатньо вказати один десятковий розряд після крапки, відкинувши решту розрядів без округлення. Наприклад, замість 13.66666666666667 достатньо вказати 13.6, замість 25.23 достатньо вказати 25.2)
code
try:
    print(0, end=';')
    print(2/0, end=';')
    print(5, end=';')
except Exception:
    print(1, end=';')
except KeyboardInterrupt:
    print(4, end=';')
except:
    print('ERR', end=';')
else:
    print(6, end=';')
finally:
    print(3)
code

6. Що буде виведено за виконання?
code
def f(n):
    if n>9:
        f(n//10)
    print(n%10, end='')
 
f(123456)
code

7. Що буде виведено за виконання?
code
class A:
    c = 1

    def _g1(self): print(2, end=';')
    def __g2(self): print(3, end=';')

    def main(self):
        try: print(self.c, end=';')
        except: print('ERR', end=';')

        try: self._g1()
        except: print('ERR', end=';')

        try: self.__g2()
        except: print('ERR', end=';')


class B(A):
    c = 4

    def _g1(self): print(5, end=';')
    def __g2(self): print(6, end=';')

try: B().main()
except: print('ERR', end=';')
code

8. Що буде виведено за виконання?
code
def f(arg, n):
    arg[51] = 2
    n -= 3

try:
    n = 3
    d = {46:5, 54:7, 86:4, 57:0}
    f(d, n)
    print(n, end=';')
    for x, y in d.items():
        print(x, y, sep=';', end=';')
except:
    print('ERR')
code

9. Що буде виведено за виконання?
code
class A:
    c = 1
    a = 2

    def __init__(self):
        c = 3

class B(A):
    c = 4

    def __init__(self):
        super().__init__()
        self.b = 7

obj = B()
obj.b = obj.c + 10
B.b = 13

try: print(obj.b, end= ';')
except: print('ERR', end=';')

try: print(A.b, end= ';')
except: print('ERR', end=';')

try: print(B.b, end= ';')
except: print('ERR', end=';')
code

10. 
Для графа на рисунку з вершини 1 запущено алгоритм пошуку в глибину, що будує кістякове дерево. Списки суміжності вершин переглядаються алгоритмом за зростанням міток вершин.\par
\begin{center}
	\adjustimage{max size={0.9\linewidth}{0.9\paperheight}}
	{../10/Traversals/57.png}
\end{center}
\par Які з наведених ребер входять до складу отриманого кістякового дерева?\par
1) (1, 6)\par
2) (2, 5)\par
3) (4, 5)\par
4) (1, 3)\par
5) (1, 5)\par
6) (0, 4)\par
{\vskip 15pt} У формі вибрати номери відповідних ребер.\par
%answer:3 6
\end {large}}
%end
{\smallskip \begin {small} Затверджено на засіданні кафедри теоретичної кібернетики, протокол \No 12 від   19.05.2020 р.\par\smallskip\smallskip
{\parbox {6.5 cm}{Зав. кафедри \dotfill Ю.В.~Крак}}\hspace {0.7cm}{\hbox to 6.5 cm{ Екзаменатор \dotfill Т.О.~Карнаух}} \end{small}}
%begin
{{\begin {small}\par
{ \center  Київський національний університет імені Тараса Шевченка \par}
{\parbox {5 cm}{Освітня програма \hfill}{\parbox {7 cm}{Прикладна математика, Системний аналіз \hfill}}\parbox {6 cm}{\hfill Семестр 2}}\par
{\parbox {5 cm} {Навчальний предмет \hfill}\parbox {7 cm}{Програмування \hfill}\parbox {6cm}{\hfill}\par}\vspace{-7pt} \end{small}}{\center Екзаменаційний білет \No \textbf 
{ 188 }\par}}
{
\begin {large}
1. Що буде виведено за виконання?
code
try:
    res = 7/5%8*7
    if isinstance(res, bool):
        print(f'bool;{res}')
    elif isinstance(res, int):
        print(f'int;{res}')
    elif isinstance(res, float):
        print(f'float;{res:.2f}')
    else:
        print('???')
except:
    print('ERR')
code

2. Що буде виведено за виконання?
code
a = 4
b = 1
c = 7

def h(b):
    a *= 1
    b = 5
    c = 1
    return a + b + c

a = 3
b = 5
c = 1

try:
    print(h(b), a, b, c, sep=';')
except:
    print('ERR', a, b, c, sep=';')
code

3. Що буде виведено за виконання?
code
try:
    i = 2
    while i<13:
        i += 2
    print(i, end=';')
except:
    print('ERR')
code

4. Що буде виведено за виконання?
code
try:
    for d in range(-1, 2, -2):
        if d>7:
            continue
        print(d, end=';')
    else:
        print(d,end=';')
    print(d)
except:
    print('ERR')
code

5. Що буде виведено за виконання?\par (Дійсні значення, якщо наявні, в цьому завданні будуть виводитись у форматі з фіксованою крапкою та, можливо, з доволі великою кількістю десяткових розрядів. У відповіді для дійсних значень у ЦЬОМУ завданні достатньо вказати один десятковий розряд після крапки, відкинувши решту розрядів без округлення. Наприклад, замість 13.66666666666667 достатньо вказати 13.6, замість 25.23 достатньо вказати 25.2)
code
try:
    print(7, end=';')
    print(8/2, end=';')
    print(5, end=';')
except TypeError:
    print(3, end=';')
except KeyboardInterrupt:
    print(2, end=';')
except:
    print('ERR', end=';')
else:
    print(0, end=';')
finally:
    print(6)
code

6. Що буде виведено за виконання?
code
def f(s):
    for i in range(len(s)):
        s[i] += 1
    return
 
s = [1, 2, 3]
f(s)
print(s)
code

7. Що буде виведено за виконання?
code
class A:
    b = 1

    def f1(self): print(2, end=';')
    def _f2(self): print(3, end=';')

    def main(self):
        try: print(self.b, end=';')
        except: print('ERR', end=';')

        try: self.f1()
        except: print('ERR', end=';')

        try: self._f2()
        except: print('ERR', end=';')


class B(A):
    b = 4

    def f1(self): print(5, end=';')
    def _f2(self): print(6, end=';')

try: B().main()
except: print('ERR', end=';')
code

8. Що буде виведено за виконання?
code
def f(arg, n):
    n = 9
    tmp = arg
    tmp[-2:-5] = (15)
    return len(tmp)>3

try:
    g = [10,11,12,13,14]
    n = 6
    f(g, n)
    print(n, end=';')
    for el in g:
        print(el, end=';')
except:
    print('ERR')
code

9. Що буде виведено за виконання?
code
class A:
    b = 1
    c = 2

    def __init__(self):
        c = 3

class B(A):
    c = 4

    def __init__(self):
        super().__init__()
        self.a = 7

obj = B()
obj.a = obj.c + 10
B.a = 13

try: print(obj.a, end= ';')
except: print('ERR', end=';')

try: print(A.a, end= ';')
except: print('ERR', end=';')

try: print(B.a, end= ';')
except: print('ERR', end=';')
code

10. 
Для графа на рисунку з вершини 2 запущено алгоритм пошуку в глибину, що будує кістякове дерево. Списки суміжності вершин переглядаються алгоритмом за зростанням міток вершин.\par
\begin{center}
	\adjustimage{max size={0.9\linewidth}{0.9\paperheight}}
	{../10/Traversals/244.png}
\end{center}
\par Які з наведених ребер входять до складу отриманого кістякового дерева?\par
1) (5, 6)\par
2) (3, 6)\par
3) (4, 5)\par
4) (0, 2)\par
5) (0, 5)\par
6) (1, 2)\par
{\vskip 15pt} У формі вибрати номери відповідних ребер.\par
%answer:4 5
\end {large}}
%end
{\smallskip \begin {small} Затверджено на засіданні кафедри теоретичної кібернетики, протокол \No 12 від   19.05.2020 р.\par\smallskip\smallskip
{\parbox {6.5 cm}{Зав. кафедри \dotfill Ю.В.~Крак}}\hspace {0.7cm}{\hbox to 6.5 cm{ Екзаменатор \dotfill Т.О.~Карнаух}} \end{small}}
%begin
{{\begin {small}\par
{ \center  Київський національний університет імені Тараса Шевченка \par}
{\parbox {5 cm}{Освітня програма \hfill}{\parbox {7 cm}{Прикладна математика, Системний аналіз \hfill}}\parbox {6 cm}{\hfill Семестр 2}}\par
{\parbox {5 cm} {Навчальний предмет \hfill}\parbox {7 cm}{Програмування \hfill}\parbox {6cm}{\hfill}\par}\vspace{-7pt} \end{small}}{\center Екзаменаційний білет \No \textbf 
{ 189 }\par}}
{
\begin {large}
1. Що буде виведено за виконання?
code
try:
    res = 3/8%3*5
    if isinstance(res, bool):
        print(f'bool;{res}')
    elif isinstance(res, int):
        print(f'int;{res}')
    elif isinstance(res, float):
        print(f'float;{res:.2f}')
    else:
        print('???')
except:
    print('ERR')
code

2. Що буде виведено за виконання?
code
a = 7
b = 5
c = 1

def f(b):
    a = 1
    b = 2
    c = 4
    return a + b + c

a = 0
b = 3
c = 6

try:
    print(f(b), a, b, c, sep=';')
except:
    print('ERR', a, b, c, sep=';')
code

3. Що буде виведено за виконання?
code
try:
    n = 11
    while n>2:
        n += -2
    print(n, end=';')
except:
    print('ERR')
code

4. Що буде виведено за виконання?
code
try:
    for a in range(9, -14, -3):
        if a>4:
            continue
        print(a, end=';')
    else:
        print(a,end=';')
    print(a)
except:
    print('ERR')
code

5. Що буде виведено за виконання?\par (Дійсні значення, якщо наявні, в цьому завданні будуть виводитись у форматі з фіксованою крапкою та, можливо, з доволі великою кількістю десяткових розрядів. У відповіді для дійсних значень у ЦЬОМУ завданні достатньо вказати один десятковий розряд після крапки, відкинувши решту розрядів без округлення. Наприклад, замість 13.66666666666667 достатньо вказати 13.6, замість 25.23 достатньо вказати 25.2)
code
try:
    print(1, end=';')
    print(int('d9'), end=';')
    print(4, end=';')
except TypeError:
    print(7, end=';')
except Exception:
    print(2, end=';')
except:
    print('ERR', end=';')
else:
    print(6, end=';')
finally:
    print(5)
code

6. Що буде виведено за виконання?
code
def f(n):
    if n<=9:
        print(n%10,end='')
    else:
        f(n//10)

f(543216)
code

7. Що буде виведено за виконання?
code
class A:
    b = 1

    def __h1(self): print(2, end=';')
    def _h2(self): print(3, end=';')

    def main(self):
        try: print(self.b, end=';')
        except: print('ERR', end=';')

        try: self.__h1()
        except: print('ERR', end=';')

        try: self._h2()
        except: print('ERR', end=';')


class B(A):
    b = 4

    def __h1(self): print(5, end=';')
    def _h2(self): print(6, end=';')

try: B().main()
except: print('ERR', end=';')
code

8. Що буде виведено за виконання?
code
def f(arg, n):
    n = 2
    tmp = arg
    tmp[0:1] = 'ac'
    return len(tmp)>3

try:
    f = ['0','1','2','3','4']
    n = 5
    f(f, n)
    print(n, end=';')
    for el in f:
        print(el, end=';')
except:
    print('ERR')
code

9. Що буде виведено за виконання?
code
class A:
    c = 1
    b = 2

    def __init__(self):
        c = 3

class B(A):
    c = 4

    def __init__(self):
        super().__init__()
        self.a = 7

obj = B()
obj.b = obj.a + 10
B.b = 13

try: print(obj.b, end= ';')
except: print('ERR', end=';')

try: print(A.b, end= ';')
except: print('ERR', end=';')

try: print(B.b, end= ';')
except: print('ERR', end=';')
code

10. 
Для графа на рисунку з вершини 1 запущено алгоритм пошуку в глибину, що будує кістякове дерево. Списки суміжності вершин переглядаються алгоритмом за зростанням міток вершин.\par
\begin{center}
	\adjustimage{max size={0.9\linewidth}{0.9\paperheight}}
	{../10/Traversals/137.png}
\end{center}
\par Які з наведених ребер входять до складу отриманого кістякового дерева?\par
1) (1, 5)\par
2) (0, 3)\par
3) (4, 5)\par
4) (0, 6)\par
5) (1, 3)\par
6) (3, 5)\par
{\vskip 15pt} У формі вибрати номери відповідних ребер.\par
%answer:2 3 4 5
\end {large}}
%end
{\smallskip \begin {small} Затверджено на засіданні кафедри теоретичної кібернетики, протокол \No 12 від   19.05.2020 р.\par\smallskip\smallskip
{\parbox {6.5 cm}{Зав. кафедри \dotfill Ю.В.~Крак}}\hspace {0.7cm}{\hbox to 6.5 cm{ Екзаменатор \dotfill Т.О.~Карнаух}} \end{small}}
%begin
{{\begin {small}\par
{ \center  Київський національний університет імені Тараса Шевченка \par}
{\parbox {5 cm}{Освітня програма \hfill}{\parbox {7 cm}{Прикладна математика, Системний аналіз \hfill}}\parbox {6 cm}{\hfill Семестр 2}}\par
{\parbox {5 cm} {Навчальний предмет \hfill}\parbox {7 cm}{Програмування \hfill}\parbox {6cm}{\hfill}\par}\vspace{-7pt} \end{small}}{\center Екзаменаційний білет \No \textbf 
{ 190 }\par}}
{
\begin {large}
1. Що буде виведено за виконання?
code
def f(n):
    print('f', end='')
    return n

try:
    print(f(-3) or f(1))
except:
    print('ERR')
code

2. Що буде виведено за виконання?
code
a = 4
b = 7
c = 3

def h(a):
    global c
    a = 3
    b = 1
    c = 2
    return a + b + c

a = 2
b = 5
c = 0

try:
    print(h(a), a, b, c, sep=';')
except:
    print('ERR', a, b, c, sep=';')
code

3. Що буде виведено за виконання?
code
try:
    t = 1
    while t<7:
        t += 1
    print(t, end=';')
except:
    print('ERR')
code

4. Що буде виведено за виконання?
code
try:
    for e in range(-3, 15, 3):
        if e>=3:
            continue
        print(e, end=';');e = -3
    else:
        print('end',end=';')
    print(e)
except:
    print('ERR')
code

5. Що буде виведено за виконання?\par (Дійсні значення, якщо наявні, в цьому завданні будуть виводитись у форматі з фіксованою крапкою та, можливо, з доволі великою кількістю десяткових розрядів. У відповіді для дійсних значень у ЦЬОМУ завданні достатньо вказати один десятковий розряд після крапки, відкинувши решту розрядів без округлення. Наприклад, замість 13.66666666666667 достатньо вказати 13.6, замість 25.23 достатньо вказати 25.2)
code
try:
    print(9, end=';')
    print(1j<3j, end=';')
    print(0, end=';')
except ZeroDivisionError:
    print(3, end=';')
except ValueError:
    print(8, end=';')
except:
    print('ERR', end=';')
else:
    print(4, end=';')
finally:
    print(7)
code

6. Що буде виведено за виконання?
code
def f(s):
    for i in range(len(s)):
        s[i] += 1
        return
 
s = [1, 2, 3]
f(s)
print(s)
code

7. Що буде виведено за виконання?
code
class A:
    d = 1

    def _h1(self): print(2, end=';')
    def h2(self): print(3, end=';')

    def main(self):
        try: print(self.d, end=';')
        except: print('ERR', end=';')

        try: self._h1()
        except: print('ERR', end=';')

        try: self.h2()
        except: print('ERR', end=';')


class B(A):
    d = 4

    def _h1(self): print(5, end=';')
    def h2(self): print(6, end=';')

try: B().main()
except: print('ERR', end=';')
code

8. Що буде виведено за виконання?
code
def f(arg, n):
    arg[68] = 7
    n = 9
    arg = tuple(arg.keys())

try:
    n = 1
    s = {58:8, 18:5, 18:3}
    f(s, n)
    print(n, end=';')
    for x in s.items():
        print(*x, sep=';', end=';')
except:
    print('ERR')
code

9. Що буде виведено за виконання?
code
class A:
    a = 1
    c = 2

    def __init__(self):
        c = 3

class B(A):
    c = 4

    def __init__(self):
        super().__init__()
        self.b = 7

obj = B()
obj.b = obj.c + 10
B.b = 13

try: print(obj.b, end= ';')
except: print('ERR', end=';')

try: print(A.b, end= ';')
except: print('ERR', end=';')

try: print(B.b, end= ';')
except: print('ERR', end=';')
code

10. 
Для графа на рисунку з вершини 1 запущено алгоритм пошуку в глибину, що будує кістякове дерево. Списки суміжності вершин переглядаються алгоритмом за зростанням міток вершин.\par
\begin{center}
	\adjustimage{max size={0.9\linewidth}{0.9\paperheight}}
	{../10/Traversals/14.png}
\end{center}
\par Які з наведених ребер входять до складу отриманого кістякового дерева?\par
1) (0, 6)\par
2) (4, 5)\par
3) (3, 4)\par
4) (1, 2)\par
5) (2, 3)\par
6) (3, 5)\par
{\vskip 15pt} У формі вибрати номери відповідних ребер.\par
%answer:1 3 5 6
\end {large}}
%end
{\smallskip \begin {small} Затверджено на засіданні кафедри теоретичної кібернетики, протокол \No 12 від   19.05.2020 р.\par\smallskip\smallskip
{\parbox {6.5 cm}{Зав. кафедри \dotfill Ю.В.~Крак}}\hspace {0.7cm}{\hbox to 6.5 cm{ Екзаменатор \dotfill Т.О.~Карнаух}} \end{small}}
%begin
{{\begin {small}\par
{ \center  Київський національний університет імені Тараса Шевченка \par}
{\parbox {5 cm}{Освітня програма \hfill}{\parbox {7 cm}{Прикладна математика, Системний аналіз \hfill}}\parbox {6 cm}{\hfill Семестр 2}}\par
{\parbox {5 cm} {Навчальний предмет \hfill}\parbox {7 cm}{Програмування \hfill}\parbox {6cm}{\hfill}\par}\vspace{-7pt} \end{small}}{\center Екзаменаційний білет \No \textbf 
{ 191 }\par}}
{
\begin {large}
1. Що буде виведено за виконання?
code
try:
    res = 4**0.5*1
    if isinstance(res, bool):
        print(f'bool;{res}')
    elif isinstance(res, int):
        print(f'int;{res}')
    elif isinstance(res, float):
        print(f'float;{res:.2f}')
    else:
        print('???')
except:
    print('ERR')
code

2. Що буде виведено за виконання?
code
a = 4
b = 6
c = 9

def f(b):
    a += 3
    b = 2
    c = 4
    return a + b + c

a = 8
b = 0
c = 2

try:
    print(f(a), a, b, c, sep=';')
except:
    print('ERR', a, b, c, sep=';')
code

3. Що буде виведено за виконання?
code
try:
    j = 12
    while j>=0:
        j += -3
    print(j, end=';')
except:
    print('ERR')
code

4. Що буде виведено за виконання?
code
try:
    for c in range(1, 13, 3):
        if c<5:
            break
        print(c, end=';')
    else:
        print(c,end=';')
    print(c)
except:
    print('ERR')
code

5. Що буде виведено за виконання?\par (Дійсні значення, якщо наявні, в цьому завданні будуть виводитись у форматі з фіксованою крапкою та, можливо, з доволі великою кількістю десяткових розрядів. У відповіді для дійсних значень у ЦЬОМУ завданні достатньо вказати один десятковий розряд після крапки, відкинувши решту розрядів без округлення. Наприклад, замість 13.66666666666667 достатньо вказати 13.6, замість 25.23 достатньо вказати 25.2)
code
try:
    print(4, end=';')
    print(3%0.0, end=';')
    print(8, end=';')
except TypeError:
    print(0, end=';')
except ValueError:
    print(1, end=';')
except:
    print('ERR', end=';')
else:
    print(6, end=';')
finally:
    print(9)
code

6. Що буде виведено за виконання?
code
def f(n):
    if n>9:
        return f(n//10)+1
    else:
        return 1

print(f(123456))
code

7. Що буде виведено за виконання?
code
class A:
    a = 1

    def _g1(self): print(2, end=';')
    def _g2(self): print(3, end=';')

    def main(self):
        try: print(self.a, end=';')
        except: print('ERR', end=';')

        try: self._g1()
        except: print('ERR', end=';')

        try: self._g2()
        except: print('ERR', end=';')


class B(A):
    a = 4

    def _g1(self): print(5, end=';')
    def _g2(self): print(6, end=';')

try: B().main()
except: print('ERR', end=';')
code

8. Що буде виведено за виконання?
code
def f(arg, n):
    arg[76] = 0
    n = 8
    arg = set(arg.items())

try:
    n = 3
    t = {76:2, 72:0, 76:4}
    f(t, n)
    print(n, end=';')
    for x in t.items():
        print(x, sep=';', end=';')
except:
    print('ERR')
code

9. Що буде виведено за виконання?
code
class A:
    a = 1
    b = 2

    def __init__(self):
        a = 3

class B(A):
    a = 4

    def __init__(self):
        super().__init__()
        self.c = 7

obj = B()
obj.a = obj.a + 10
B.a = 13

try: print(obj.a, end= ';')
except: print('ERR', end=';')

try: print(A.a, end= ';')
except: print('ERR', end=';')

try: print(B.a, end= ';')
except: print('ERR', end=';')
code

10. 
Для графа на рисунку з вершини 1 запущено алгоритм пошуку в глибину, що будує кістякове дерево. Списки суміжності вершин переглядаються алгоритмом за зростанням міток вершин.\par
\begin{center}
	\adjustimage{max size={0.9\linewidth}{0.9\paperheight}}
	{../10/Traversals/27.png}
\end{center}
\par Які з наведених ребер входять до складу отриманого кістякового дерева?\par
1) (0, 3)\par
2) (0, 5)\par
3) (0, 2)\par
4) (0, 4)\par
5) (0, 1)\par
6) (2, 4)\par
{\vskip 15pt} У формі вибрати номери відповідних ребер.\par
%answer:3 5 6
\end {large}}
%end
{\smallskip \begin {small} Затверджено на засіданні кафедри теоретичної кібернетики, протокол \No 12 від   19.05.2020 р.\par\smallskip\smallskip
{\parbox {6.5 cm}{Зав. кафедри \dotfill Ю.В.~Крак}}\hspace {0.7cm}{\hbox to 6.5 cm{ Екзаменатор \dotfill Т.О.~Карнаух}} \end{small}}
%begin
{{\begin {small}\par
{ \center  Київський національний університет імені Тараса Шевченка \par}
{\parbox {5 cm}{Освітня програма \hfill}{\parbox {7 cm}{Прикладна математика, Системний аналіз \hfill}}\parbox {6 cm}{\hfill Семестр 2}}\par
{\parbox {5 cm} {Навчальний предмет \hfill}\parbox {7 cm}{Програмування \hfill}\parbox {6cm}{\hfill}\par}\vspace{-7pt} \end{small}}{\center Екзаменаційний білет \No \textbf 
{ 192 }\par}}
{
\begin {large}
1. Що буде виведено за виконання?
code
try:
    res = 9//9*9*8
    if isinstance(res, bool):
        print(f'bool;{res}')
    elif isinstance(res, int):
        print(f'int;{res}')
    elif isinstance(res, float):
        print(f'float;{res:.2f}')
    else:
        print('???')
except:
    print('ERR')
code

2. Що буде виведено за виконання?
code
a = 9
b = 6
c = 1

def f(b):
    a = 4
    b = 5
    c = 1
    return a + b + c

a = 8
b = 2
c = 0

try:
    print(f(a), a, b, c, sep=';')
except:
    print('ERR', a, b, c, sep=';')
code

3. Що буде виведено за виконання?
code
try:
    m = 0
    while m<=9:
        m += 2
    print(m, end=';')
except:
    print('ERR')
code

4. Що буде виведено за виконання?
code
try:
    for f in range(-2, -5, 2):
        if f<=8:
            continue
        print(f, end=';')
    else:
        print(f,end=';')
    print(f)
except:
    print('ERR')
code

5. Що буде виведено за виконання?\par (Дійсні значення, якщо наявні, в цьому завданні будуть виводитись у форматі з фіксованою крапкою та, можливо, з доволі великою кількістю десяткових розрядів. У відповіді для дійсних значень у ЦЬОМУ завданні достатньо вказати один десятковий розряд після крапки, відкинувши решту розрядів без округлення. Наприклад, замість 13.66666666666667 достатньо вказати 13.6, замість 25.23 достатньо вказати 25.2)
code
try:
    print(5, end=';')
    print(2//3, end=';')
    print(7, end=';')
except Exception:
    print(8, end=';')
except ZeroDivisionError:
    print(3, end=';')
except:
    print('ERR', end=';')
else:
    print(6, end=';')
finally:
    print(1)
code

6. Що буде виведено за виконання?
code
def f(n):
    if n>9:
        f(n//100)
    print(n%10, end='')
 
f(123456)
code

7. Що буде виведено за виконання?
code
class A:
    c = 1

    def __f1(self): print(2, end=';')
    def _f2(self): print(3, end=';')

    def main(self):
        try: print(self.c, end=';')
        except: print('ERR', end=';')

        try: self.__f1()
        except: print('ERR', end=';')

        try: self._f2()
        except: print('ERR', end=';')


class B(A):
    c = 4

    def __f1(self): print(5, end=';')
    def _f2(self): print(6, end=';')

try: B().main()
except: print('ERR', end=';')
code

8. Що буде виведено за виконання?
code
def f(arg, n):
    arg[35] = 6
    n = 7
    arg = tuple(arg.items())

try:
    n = 4
    s = {72:9, 85:7, 60:0, 60:7}
    f(s, n)
    print(n, end=';')
    for x in s.keys():
        print(x, sep=';', end=';')
except:
    print('ERR')
code

9. Що буде виведено за виконання?
code
class A:
    b = 1
    a = 2

    def __init__(self):
        a = 3

class B(A):
    a = 4

    def __init__(self):
        super().__init__()
        self.c = 7

obj = B()
obj.b = obj.c + 10
B.b = 13

try: print(obj.b, end= ';')
except: print('ERR', end=';')

try: print(A.b, end= ';')
except: print('ERR', end=';')

try: print(B.b, end= ';')
except: print('ERR', end=';')
code

10. 
Для графа на рисунку з вершини 5 запущено алгоритм пошуку в глибину, що будує кістякове дерево. Списки суміжності вершин переглядаються алгоритмом за зростанням міток вершин.\par
\begin{center}
	\adjustimage{max size={0.9\linewidth}{0.9\paperheight}}
	{../10/Traversals/224.png}
\end{center}
\par Які з наведених ребер входять до складу отриманого кістякового дерева?\par
1) (0, 3)\par
2) (0, 5)\par
3) (1, 5)\par
4) (3, 6)\par
5) (2, 6)\par
6) (5, 6)\par
{\vskip 15pt} У формі вибрати номери відповідних ребер.\par
%answer:2 5
\end {large}}
%end
{\smallskip \begin {small} Затверджено на засіданні кафедри теоретичної кібернетики, протокол \No 12 від   19.05.2020 р.\par\smallskip\smallskip
{\parbox {6.5 cm}{Зав. кафедри \dotfill Ю.В.~Крак}}\hspace {0.7cm}{\hbox to 6.5 cm{ Екзаменатор \dotfill Т.О.~Карнаух}} \end{small}}
%begin
{{\begin {small}\par
{ \center  Київський національний університет імені Тараса Шевченка \par}
{\parbox {5 cm}{Освітня програма \hfill}{\parbox {7 cm}{Прикладна математика, Системний аналіз \hfill}}\parbox {6 cm}{\hfill Семестр 2}}\par
{\parbox {5 cm} {Навчальний предмет \hfill}\parbox {7 cm}{Програмування \hfill}\parbox {6cm}{\hfill}\par}\vspace{-7pt} \end{small}}{\center Екзаменаційний білет \No \textbf 
{ 193 }\par}}
{
\begin {large}
1. Що буде виведено за виконання?
code
def f(n):
    print('f', end='')
    return n

try:
    print(f(-8) or f(-3))
except:
    print('ERR')
code

2. Що буде виведено за виконання?
code
a = 5
b = 8
c = 9

def g(a):
    global c
    a = 5
    b = 4
    c = 2
    return a + b + c

a = 6
b = 3
c = 4

try:
    print(g(b), a, b, c, sep=';')
except:
    print('ERR', a, b, c, sep=';')
code

3. Що буде виведено за виконання?
code
try:
    t = 13
    while t>7:
        t += -3
    print(t, end=';')
except:
    print('ERR')
code

4. Що буде виведено за виконання?
code
try:
    for e in range(-11, -1, -2):
        if e>4:
            continue
        print(e, end=';')
    else:
        print(e,end=';')
    print(e)
except:
    print('ERR')
code

5. Що буде виведено за виконання?\par (Дійсні значення, якщо наявні, в цьому завданні будуть виводитись у форматі з фіксованою крапкою та, можливо, з доволі великою кількістю десяткових розрядів. У відповіді для дійсних значень у ЦЬОМУ завданні достатньо вказати один десятковий розряд після крапки, відкинувши решту розрядів без округлення. Наприклад, замість 13.66666666666667 достатньо вказати 13.6, замість 25.23 достатньо вказати 25.2)
code
try:
    print(5, end=';')
    print(int('7'), end=';')
    print(4, end=';')
except KeyboardInterrupt:
    print(9, end=';')
except KeyboardInterrupt:
    print(0, end=';')
except:
    print('ERR', end=';')
else:
    print(2, end=';')
finally:
    print(7)
code

6. Що буде виведено за виконання?
code
def f(n):
    if n>9:
        return f(n//100)+1
    else:
        return 1

print(f(123456))
code

7. Що буде виведено за виконання?
code
class A:
    c = 1

    def _g1(self): print(2, end=';')
    def g2(self): print(3, end=';')

    def main(self):
        try: print(self.c, end=';')
        except: print('ERR', end=';')

        try: self._g1()
        except: print('ERR', end=';')

        try: self.g2()
        except: print('ERR', end=';')


class B(A):
    c = 4

    def _g1(self): print(5, end=';')
    def g2(self): print(6, end=';')

try: B().main()
except: print('ERR', end=';')
code

8. Що буде виведено за виконання?
code
def f(arg, n):
    n = 2
    tmp = arg
    del tmp[-4:4]
    return len(tmp)>3

try:
    e = [10,11,12,13,14]
    n = 0
    f(e, n)
    print(n, end=';')
    for el in e:
        print(el, end=';')
except:
    print('ERR')
code

9. Що буде виведено за виконання?
code
class A:
    b = 1
    c = 2

    def __init__(self):
        b = 3

class B(A):
    b = 4

    def __init__(self):
        super().__init__()
        self.a = 7

obj = B()
obj.b = obj.a + 10
B.b = 13

try: print(obj.b, end= ';')
except: print('ERR', end=';')

try: print(A.b, end= ';')
except: print('ERR', end=';')

try: print(B.b, end= ';')
except: print('ERR', end=';')
code

10. 
Для графа на рисунку з вершини 5 запущено алгоритм пошуку в глибину, що будує кістякове дерево. Списки суміжності вершин переглядаються алгоритмом за зростанням міток вершин.\par
\begin{center}
	\adjustimage{max size={0.9\linewidth}{0.9\paperheight}}
	{../10/Traversals/33.png}
\end{center}
\par Які з наведених ребер входять до складу отриманого кістякового дерева?\par
1) (4, 6)\par
2) (1, 5)\par
3) (1, 2)\par
4) (0, 3)\par
5) (0, 1)\par
6) (1, 4)\par
{\vskip 15pt} У формі вибрати номери відповідних ребер.\par
%answer:1 2 4 5
\end {large}}
%end
{\smallskip \begin {small} Затверджено на засіданні кафедри теоретичної кібернетики, протокол \No 12 від   19.05.2020 р.\par\smallskip\smallskip
{\parbox {6.5 cm}{Зав. кафедри \dotfill Ю.В.~Крак}}\hspace {0.7cm}{\hbox to 6.5 cm{ Екзаменатор \dotfill Т.О.~Карнаух}} \end{small}}
%begin
{{\begin {small}\par
{ \center  Київський національний університет імені Тараса Шевченка \par}
{\parbox {5 cm}{Освітня програма \hfill}{\parbox {7 cm}{Прикладна математика, Системний аналіз \hfill}}\parbox {6 cm}{\hfill Семестр 2}}\par
{\parbox {5 cm} {Навчальний предмет \hfill}\parbox {7 cm}{Програмування \hfill}\parbox {6cm}{\hfill}\par}\vspace{-7pt} \end{small}}{\center Екзаменаційний білет \No \textbf 
{ 194 }\par}}
{
\begin {large}
1. Що буде виведено за виконання?
code
try:
    res = 6/6/7*9
    if isinstance(res, bool):
        print(f'bool;{res}')
    elif isinstance(res, int):
        print(f'int;{res}')
    elif isinstance(res, float):
        print(f'float;{res:.2f}')
    else:
        print('???')
except:
    print('ERR')
code

2. Що буде виведено за виконання?
code
a = 1
b = 7
c = 6

def h(b):
    a = 3
    b = 4
    c = 2
    return a + b + c

a = 4
b = 7
c = 0

try:
    print(h(b), a, b, c, sep=';')
except:
    print('ERR', a, b, c, sep=';')
code

3. Що буде виведено за виконання?
code
try:
    n = 6
    while n<11:
        n += 2
    print(n, end=';')
except:
    print('ERR')
code

4. Що буде виведено за виконання?
code
try:
    for f in range(11, 3, -1):
        if f<6:
            continue
        print(f, end=';');f = 4
    else:
        print(f,end=';')
    print(f)
except:
    print('ERR')
code

5. Що буде виведено за виконання?\par (Дійсні значення, якщо наявні, в цьому завданні будуть виводитись у форматі з фіксованою крапкою та, можливо, з доволі великою кількістю десяткових розрядів. У відповіді для дійсних значень у ЦЬОМУ завданні достатньо вказати один десятковий розряд після крапки, відкинувши решту розрядів без округлення. Наприклад, замість 13.66666666666667 достатньо вказати 13.6, замість 25.23 достатньо вказати 25.2)
code
try:
    print(3, end=';')
    print(int('c5'), end=';')
    print(9, end=';')
except ValueError:
    print(8, end=';')
except ZeroDivisionError:
    print(0, end=';')
except:
    print('ERR', end=';')
else:
    print(2, end=';')
finally:
    print(4)
code

6. Що буде виведено за виконання?
code
def f(n):
    if n>9:
        f(n//100)
    print(n%10,end='')
 
f(987654)
code

7. Що буде виведено за виконання?
code
class A:
    b = 1

    def f1(self): print(2, end=';')
    def _f2(self): print(3, end=';')

    def main(self):
        try: print(self.b, end=';')
        except: print('ERR', end=';')

        try: self.f1()
        except: print('ERR', end=';')

        try: self._f2()
        except: print('ERR', end=';')


class B(A):
    b = 4

    def f1(self): print(5, end=';')
    def _f2(self): print(6, end=';')

try: B().main()
except: print('ERR', end=';')
code

8. Що буде виведено за виконання?
code
def f(arg, n):
    arg[72] = 2
    n *= -2

try:
    n = 4
    t = {11:3, 52:2, 72:9, 12:6}
    f(t, n)
    print(n, end=';')
    for x in t :
        print(x, sep=';', end=';')
except:
    print('ERR')
code

9. Що буде виведено за виконання?
code
class A:
    a = 1
    b = 2

    def __init__(self):
        b = 3

class B(A):
    b = 4

    def __init__(self):
        super().__init__()
        self.c = 7

obj = B()
obj.c = obj.c + 10
B.c = 13

try: print(obj.c, end= ';')
except: print('ERR', end=';')

try: print(A.c, end= ';')
except: print('ERR', end=';')

try: print(B.c, end= ';')
except: print('ERR', end=';')
code

10. 
Для графа на рисунку з вершини 5 запущено алгоритм пошуку в глибину, що будує кістякове дерево. Списки суміжності вершин переглядаються алгоритмом за зростанням міток вершин.\par
\begin{center}
	\adjustimage{max size={0.9\linewidth}{0.9\paperheight}}
	{../10/Traversals/159.png}
\end{center}
\par Які з наведених ребер входять до складу отриманого кістякового дерева?\par
1) (3, 6)\par
2) (2, 4)\par
3) (2, 3)\par
4) (3, 5)\par
5) (1, 5)\par
6) (1, 3)\par
{\vskip 15pt} У формі вибрати номери відповідних ребер.\par
%answer:2 3 5 6
\end {large}}
%end
{\smallskip \begin {small} Затверджено на засіданні кафедри теоретичної кібернетики, протокол \No 12 від   19.05.2020 р.\par\smallskip\smallskip
{\parbox {6.5 cm}{Зав. кафедри \dotfill Ю.В.~Крак}}\hspace {0.7cm}{\hbox to 6.5 cm{ Екзаменатор \dotfill Т.О.~Карнаух}} \end{small}}
%begin
{{\begin {small}\par
{ \center  Київський національний університет імені Тараса Шевченка \par}
{\parbox {5 cm}{Освітня програма \hfill}{\parbox {7 cm}{Прикладна математика, Системний аналіз \hfill}}\parbox {6 cm}{\hfill Семестр 2}}\par
{\parbox {5 cm} {Навчальний предмет \hfill}\parbox {7 cm}{Програмування \hfill}\parbox {6cm}{\hfill}\par}\vspace{-7pt} \end{small}}{\center Екзаменаційний білет \No \textbf 
{ 195 }\par}}
{
\begin {large}
1. Що буде виведено за виконання?
code
def f(n):
    print('f', end='')
    return n>=3

try:
    print(f(-9) and f(4))
except:
    print('ERR')
code

2. Що буде виведено за виконання?
code
a = 7
b = 9
c = 4

def g(a):
    global c
    a -= 3
    b = 5
    c = 4
    return a + b + c

a = 1
b = 8
c = 5

try:
    print(g(b), a, b, c, sep=';')
except:
    print('ERR', a, b, c, sep=';')
code

3. Що буде виведено за виконання?
code
try:
    k = 10
    while k>=5:
        k += -2
    print(k, end=';')
except:
    print('ERR')
code

4. Що буде виведено за виконання?
code
try:
    for c in range(4, -4, -3):
        if c>=7:
            break
        print(c, end=';')
    else:
        print(c,end=';')
    print(c)
except:
    print('ERR')
code

5. Що буде виведено за виконання?\par (Дійсні значення, якщо наявні, в цьому завданні будуть виводитись у форматі з фіксованою крапкою та, можливо, з доволі великою кількістю десяткових розрядів. У відповіді для дійсних значень у ЦЬОМУ завданні достатньо вказати один десятковий розряд після крапки, відкинувши решту розрядів без округлення. Наприклад, замість 13.66666666666667 достатньо вказати 13.6, замість 25.23 достатньо вказати 25.2)
code
try:
    print(6, end=';')
    print(1j==6j, end=';')
    print(2, end=';')
except Exception:
    print(0, end=';')
except TypeError:
    print(1, end=';')
except:
    print('ERR', end=';')
else:
    print(4, end=';')
finally:
    print(8)
code

6. Що буде виведено за виконання?
code
def f(n):
    if n<=9:
        print(n%10, end='')
    else:
        f(n//10)

f(123456)
code

7. Що буде виведено за виконання?
code
class A:
    d = 1

    def _h1(self): print(2, end=';')
    def __h2(self): print(3, end=';')

    def main(self):
        try: print(self.d, end=';')
        except: print('ERR', end=';')

        try: self._h1()
        except: print('ERR', end=';')

        try: self.__h2()
        except: print('ERR', end=';')


class B(A):
    d = 4

    def _h1(self): print(5, end=';')
    def __h2(self): print(6, end=';')

try: B().main()
except: print('ERR', end=';')
code

8. Що буде виведено за виконання?
code
def f(arg, n):
    arg[39] = 3
    n = 6
    arg = list(arg.keys())

try:
    n = 0
    t = {68:5, 39:8, 26:9, 86:3}
    f(t, n)
    print(n, end=';')
    for x in t.keys():
        print(x, sep=';', end=';')
except:
    print('ERR')
code

9. Що буде виведено за виконання?
code
class A:
    c = 1
    b = 2

    def __init__(self):
        b = 3

class B(A):
    c = 4

    def __init__(self):
        super().__init__()
        self.a = 7

obj = B()
obj.a = obj.b + 10
B.a = 13

try: print(obj.a, end= ';')
except: print('ERR', end=';')

try: print(A.a, end= ';')
except: print('ERR', end=';')

try: print(B.a, end= ';')
except: print('ERR', end=';')
code

10. 
Для графа на рисунку з вершини 2 запущено алгоритм пошуку в глибину, що будує кістякове дерево. Списки суміжності вершин переглядаються алгоритмом за зростанням міток вершин.\par
\begin{center}
	\adjustimage{max size={0.9\linewidth}{0.9\paperheight}}
	{../10/Traversals/190.png}
\end{center}
\par Які з наведених ребер входять до складу отриманого кістякового дерева?\par
1) (5, 6)\par
2) (0, 3)\par
3) (1, 2)\par
4) (4, 5)\par
5) (3, 4)\par
6) (1, 5)\par
{\vskip 15pt} У формі вибрати номери відповідних ребер.\par
%answer:1 2 3 5
\end {large}}
%end
{\smallskip \begin {small} Затверджено на засіданні кафедри теоретичної кібернетики, протокол \No 12 від   19.05.2020 р.\par\smallskip\smallskip
{\parbox {6.5 cm}{Зав. кафедри \dotfill Ю.В.~Крак}}\hspace {0.7cm}{\hbox to 6.5 cm{ Екзаменатор \dotfill Т.О.~Карнаух}} \end{small}}
%begin
{{\begin {small}\par
{ \center  Київський національний університет імені Тараса Шевченка \par}
{\parbox {5 cm}{Освітня програма \hfill}{\parbox {7 cm}{Прикладна математика, Системний аналіз \hfill}}\parbox {6 cm}{\hfill Семестр 2}}\par
{\parbox {5 cm} {Навчальний предмет \hfill}\parbox {7 cm}{Програмування \hfill}\parbox {6cm}{\hfill}\par}\vspace{-7pt} \end{small}}{\center Екзаменаційний білет \No \textbf 
{ 196 }\par}}
{
\begin {large}
1. Що буде виведено за виконання?
code
try:
    res = not 2.0<=7 and True==9 and 5!=7.0
    if isinstance(res, bool):
        print(f'bool;{res}')
    elif isinstance(res, int):
        print(f'int;{res}')
    elif isinstance(res, float):
        print(f'float;{res:.2f}')
    else:
        print('???')
except:
    print('ERR')
code

2. Що буде виведено за виконання?
code
a = 6
b = 2
c = 3

def f(a):
    global c
    a -= 1
    b = 2
    c = 3
    return a + b + c

a = 0
b = 2
c = 3

try:
    print(f(b), a, b, c, sep=';')
except:
    print('ERR', a, b, c, sep=';')
code

3. Що буде виведено за виконання?
code
try:
    i = 9
    while i>6:
        i += -2
    print(i, end=';')
except:
    print('ERR')
code

4. Що буде виведено за виконання?
code
try:
    for b in range(-13, -8, -2):
        if b<=8:
            break
        print(b, end=';')
    else:
        print(b,end=';')
    print(b)
except:
    print('ERR')
code

5. Що буде виведено за виконання?\par (Дійсні значення, якщо наявні, в цьому завданні будуть виводитись у форматі з фіксованою крапкою та, можливо, з доволі великою кількістю десяткових розрядів. У відповіді для дійсних значень у ЦЬОМУ завданні достатньо вказати один десятковий розряд після крапки, відкинувши решту розрядів без округлення. Наприклад, замість 13.66666666666667 достатньо вказати 13.6, замість 25.23 достатньо вказати 25.2)
code
try:
    print(7, end=';')
    print(3/1, end=';')
    print(5, end=';')
except Exception:
    print(9, end=';')
except TypeError:
    print(6, end=';')
except:
    print('ERR', end=';')
else:
    print(2, end=';')
finally:
    print(8)
code

6. Що буде виведено за виконання?
code
def f(n):
    if n<=9:
        print(n%10, end='')
    else:
        f(n//10)

f(987651)
code

7. Що буде виведено за виконання?
code
class A:
    a = 1

    def _h1(self): print(2, end=';')
    def _h2(self): print(3, end=';')

    def main(self):
        try: print(self.a, end=';')
        except: print('ERR', end=';')

        try: self._h1()
        except: print('ERR', end=';')

        try: self._h2()
        except: print('ERR', end=';')


class B(A):
    a = 4

    def _h1(self): print(5, end=';')
    def _h2(self): print(6, end=';')

try: B().main()
except: print('ERR', end=';')
code

8. Що буде виведено за виконання?
code
def f(arg, n):
    arg[43] = 4
    n += 3

try:
    n = 6
    s = {83:8, 71:7, 90:4, 71:0}
    f(s, n)
    print(n, end=';')
    for x, y in s.items():
        print(x, y, sep=';', end=';')
except:
    print('ERR')
code

9. Що буде виведено за виконання?
code
class A:
    b = 1
    a = 2

    def __init__(self):
        b = 3

class B(A):
    a = 4

    def __init__(self):
        super().__init__()
        self.c = 7

obj = B()
obj.b = obj.a + 10
B.b = 13

try: print(obj.b, end= ';')
except: print('ERR', end=';')

try: print(A.b, end= ';')
except: print('ERR', end=';')

try: print(B.b, end= ';')
except: print('ERR', end=';')
code

10. 
Для графа на рисунку з вершини 6 запущено алгоритм пошуку в глибину, що будує кістякове дерево. Списки суміжності вершин переглядаються алгоритмом за зростанням міток вершин.\par
\begin{center}
	\adjustimage{max size={0.9\linewidth}{0.9\paperheight}}
	{../10/Traversals/253.png}
\end{center}
\par Які з наведених ребер входять до складу отриманого кістякового дерева?\par
1) (1, 3)\par
2) (3, 5)\par
3) (3, 6)\par
4) (0, 3)\par
5) (1, 2)\par
6) (2, 6)\par
{\vskip 15pt} У формі вибрати номери відповідних ребер.\par
%answer:1 2 4 5 6
\end {large}}
%end
{\smallskip \begin {small} Затверджено на засіданні кафедри теоретичної кібернетики, протокол \No 12 від   19.05.2020 р.\par\smallskip\smallskip
{\parbox {6.5 cm}{Зав. кафедри \dotfill Ю.В.~Крак}}\hspace {0.7cm}{\hbox to 6.5 cm{ Екзаменатор \dotfill Т.О.~Карнаух}} \end{small}}
%begin
{{\begin {small}\par
{ \center  Київський національний університет імені Тараса Шевченка \par}
{\parbox {5 cm}{Освітня програма \hfill}{\parbox {7 cm}{Прикладна математика, Системний аналіз \hfill}}\parbox {6 cm}{\hfill Семестр 2}}\par
{\parbox {5 cm} {Навчальний предмет \hfill}\parbox {7 cm}{Програмування \hfill}\parbox {6cm}{\hfill}\par}\vspace{-7pt} \end{small}}{\center Екзаменаційний білет \No \textbf 
{ 197 }\par}}
{
\begin {large}
1. Що буде виведено за виконання?
code
try:
    res = 5//4//5*4
    if isinstance(res, bool):
        print(f'bool;{res}')
    elif isinstance(res, int):
        print(f'int;{res}')
    elif isinstance(res, float):
        print(f'float;{res:.2f}')
    else:
        print('???')
except:
    print('ERR')
code

2. Що буде виведено за виконання?
code
a = 4
b = 0
c = 1

def h(b):
    a = 5
    b *= 2
    c = 4
    return a + b + c

a = 7
b = 6
c = 9

try:
    print(h(b), a, b, c, sep=';')
except:
    print('ERR', a, b, c, sep=';')
code

3. Що буде виведено за виконання?
code
try:
    m = 5
    while m>=9:
        m += -3
    print(m, end=';')
except:
    print('ERR')
code

4. Що буде виведено за виконання?
code
try:
    for a in range(13, -12, 3):
        if a<=5:
            continue
        print(a, end=';')
        if a<6:
            break
    else:
        print(a,end=';')
    print(a)
except:
    print('ERR')
code

5. Що буде виведено за виконання?\par (Дійсні значення, якщо наявні, в цьому завданні будуть виводитись у форматі з фіксованою крапкою та, можливо, з доволі великою кількістю десяткових розрядів. У відповіді для дійсних значень у ЦЬОМУ завданні достатньо вказати один десятковий розряд після крапки, відкинувши решту розрядів без округлення. Наприклад, замість 13.66666666666667 достатньо вказати 13.6, замість 25.23 достатньо вказати 25.2)
code
try:
    print(0, end=';')
    print(int('5'), end=';')
    print(6, end=';')
except ValueError:
    print(7, end=';')
except KeyboardInterrupt:
    print(3, end=';')
except:
    print('ERR', end=';')
else:
    print(1, end=';')
finally:
    print(4)
code

6. Що буде виведено за виконання?
code
def f(s1):
    s = s1[:]
    for i in range(len(s)):
        s[i] += 1
 
s = [1, 2, 3]
f(s)
print(s)
code

7. Що буде виведено за виконання?
code
class A:
    c = 1

    def __g1(self): print(2, end=';')
    def _g2(self): print(3, end=';')

    def main(self):
        try: print(self.c, end=';')
        except: print('ERR', end=';')

        try: self.__g1()
        except: print('ERR', end=';')

        try: self._g2()
        except: print('ERR', end=';')


class B(A):
    c = 4

    def __g1(self): print(5, end=';')
    def _g2(self): print(6, end=';')

try: B().main()
except: print('ERR', end=';')
code

8. Що буде виведено за виконання?
code
def f(arg, n):
    arg[77] = 4
    n -= 2

try:
    n = 5
    d = {77:3, 21:4, 77:9}
    f(d, n)
    print(n, end=';')
    for x in d.values():
        print(x, sep=';', end=';')
except:
    print('ERR')
code

9. Що буде виведено за виконання?
code
class A:
    c = 1
    a = 2

    def __init__(self):
        c = 3

class B(A):
    c = 4

    def __init__(self):
        super().__init__()
        self.b = 7

obj = B()
obj.c = obj.b + 10
B.c = 13

try: print(obj.c, end= ';')
except: print('ERR', end=';')

try: print(A.c, end= ';')
except: print('ERR', end=';')

try: print(B.c, end= ';')
except: print('ERR', end=';')
code

10. 
Для графа на рисунку з вершини 5 запущено алгоритм пошуку в глибину, що будує кістякове дерево. Списки суміжності вершин переглядаються алгоритмом за зростанням міток вершин.\par
\begin{center}
	\adjustimage{max size={0.9\linewidth}{0.9\paperheight}}
	{../10/Traversals/248.png}
\end{center}
\par Які з наведених ребер входять до складу отриманого кістякового дерева?\par
1) (1, 4)\par
2) (3, 6)\par
3) (0, 4)\par
4) (2, 5)\par
5) (0, 1)\par
6) (3, 5)\par
{\vskip 15pt} У формі вибрати номери відповідних ребер.\par
%answer:2 5
\end {large}}
%end
{\smallskip \begin {small} Затверджено на засіданні кафедри теоретичної кібернетики, протокол \No 12 від   19.05.2020 р.\par\smallskip\smallskip
{\parbox {6.5 cm}{Зав. кафедри \dotfill Ю.В.~Крак}}\hspace {0.7cm}{\hbox to 6.5 cm{ Екзаменатор \dotfill Т.О.~Карнаух}} \end{small}}
%begin
{{\begin {small}\par
{ \center  Київський національний університет імені Тараса Шевченка \par}
{\parbox {5 cm}{Освітня програма \hfill}{\parbox {7 cm}{Прикладна математика, Системний аналіз \hfill}}\parbox {6 cm}{\hfill Семестр 2}}\par
{\parbox {5 cm} {Навчальний предмет \hfill}\parbox {7 cm}{Програмування \hfill}\parbox {6cm}{\hfill}\par}\vspace{-7pt} \end{small}}{\center Екзаменаційний білет \No \textbf 
{ 198 }\par}}
{
\begin {large}
1. Що буде виведено за виконання?
code
try:
    res = True>=5.0
    if isinstance(res, bool):
        print(f'bool;{res}')
    elif isinstance(res, int):
        print(f'int;{res}')
    elif isinstance(res, float):
        print(f'float;{res:.2f}')
    else:
        print('???')
except:
    print('ERR')
code

2. Що буде виведено за виконання?
code
a = 5
b = 8
c = 2

def g(b):
    global c
    a = 1
    b += 5
    c = 3
    return a + b + c

a = 7
b = 0
c = 6

try:
    print(g(a), a, b, c, sep=';')
except:
    print('ERR', a, b, c, sep=';')
code

3. Що буде виведено за виконання?
code
try:
    n = 11
    while n>=9:
        n += -3
    print(n, end=';')
except:
    print('ERR')
code

4. Що буде виведено за виконання?
code
try:
    for d in range(-14, -9, 2):
        if d<3:
            continue
        print(d, end=';');d = 0
    else:
        print('end',end=';')
    print(d)
except:
    print('ERR')
code

5. Що буде виведено за виконання?\par (Дійсні значення, якщо наявні, в цьому завданні будуть виводитись у форматі з фіксованою крапкою та, можливо, з доволі великою кількістю десяткових розрядів. У відповіді для дійсних значень у ЦЬОМУ завданні достатньо вказати один десятковий розряд після крапки, відкинувши решту розрядів без округлення. Наприклад, замість 13.66666666666667 достатньо вказати 13.6, замість 25.23 достатньо вказати 25.2)
code
try:
    print(9, end=';')
    print(7//0, end=';')
    print(4, end=';')
except ZeroDivisionError:
    print(3, end=';')
except ValueError:
    print(8, end=';')
except:
    print('ERR', end=';')
else:
    print(9, end=';')
finally:
    print(6)
code

6. Що буде виведено за виконання?
code
def f(s):
    s1 = s
    for i in range(len(s)):
        s1[i] += 1
        return s1

s = [1, 2, 3]
f(s)
print(s)
code

7. Що буде виведено за виконання?
code
class A:
    d = 1

    def _f1(self): print(2, end=';')
    def __f2(self): print(3, end=';')

    def main(self):
        try: print(self.d, end=';')
        except: print('ERR', end=';')

        try: self._f1()
        except: print('ERR', end=';')

        try: self.__f2()
        except: print('ERR', end=';')


class B(A):
    d = 4

    def _f1(self): print(5, end=';')
    def __f2(self): print(6, end=';')

try: B().main()
except: print('ERR', end=';')
code

8. Що буде виведено за виконання?
code
def f(arg, n):
    n = 7
    tmp = arg
    tmp[3:6] = 'yz'
    return len(tmp)>3

try:
    d = ['0','1','2','3','4']
    n = 3
    f(d, n)
    print(n, end=';')
    for el in d:
        print(el, end=';')
except:
    print('ERR')
code

9. Що буде виведено за виконання?
code
class A:
    a = 1
    c = 2

    def __init__(self):
        c = 3

class B(A):
    c = 4

    def __init__(self):
        super().__init__()
        self.b = 7

obj = B()
obj.a = obj.c + 10
B.a = 13

try: print(obj.a, end= ';')
except: print('ERR', end=';')

try: print(A.a, end= ';')
except: print('ERR', end=';')

try: print(B.a, end= ';')
except: print('ERR', end=';')
code

10. 
Для графа на рисунку з вершини 2 запущено алгоритм пошуку в глибину, що будує кістякове дерево. Списки суміжності вершин переглядаються алгоритмом за зростанням міток вершин.\par
\begin{center}
	\adjustimage{max size={0.9\linewidth}{0.9\paperheight}}
	{../10/Traversals/252.png}
\end{center}
\par Які з наведених ребер входять до складу отриманого кістякового дерева?\par
1) (2, 5)\par
2) (3, 4)\par
3) (1, 6)\par
4) (0, 4)\par
5) (0, 1)\par
6) (0, 5)\par
{\vskip 15pt} У формі вибрати номери відповідних ребер.\par
%answer:2 3 5
\end {large}}
%end
{\smallskip \begin {small} Затверджено на засіданні кафедри теоретичної кібернетики, протокол \No 12 від   19.05.2020 р.\par\smallskip\smallskip
{\parbox {6.5 cm}{Зав. кафедри \dotfill Ю.В.~Крак}}\hspace {0.7cm}{\hbox to 6.5 cm{ Екзаменатор \dotfill Т.О.~Карнаух}} \end{small}}
%begin
{{\begin {small}\par
{ \center  Київський національний університет імені Тараса Шевченка \par}
{\parbox {5 cm}{Освітня програма \hfill}{\parbox {7 cm}{Прикладна математика, Системний аналіз \hfill}}\parbox {6 cm}{\hfill Семестр 2}}\par
{\parbox {5 cm} {Навчальний предмет \hfill}\parbox {7 cm}{Програмування \hfill}\parbox {6cm}{\hfill}\par}\vspace{-7pt} \end{small}}{\center Екзаменаційний білет \No \textbf 
{ 199 }\par}}
{
\begin {large}
1. Що буде виведено за виконання?
code
try:
    res = 2**2.0+2
    if isinstance(res, bool):
        print(f'bool;{res}')
    elif isinstance(res, int):
        print(f'int;{res}')
    elif isinstance(res, float):
        print(f'float;{res:.2f}')
    else:
        print('???')
except:
    print('ERR')
code

2. Що буде виведено за виконання?
code
a = 2
b = 9
c = 1

def f(a):
    global c
    a = 3
    b = 1
    c = 5
    return a + b + c

a = 5
b = 3
c = 8

try:
    print(f(b), a, b, c, sep=';')
except:
    print('ERR', a, b, c, sep=';')
code

3. Що буде виведено за виконання?
code
try:
    j = 9
    while j>7:
        j += -2
    print(j, end=';')
except:
    print('ERR')
code

4. Що буде виведено за виконання?
code
try:
    for d in range(10, 11, -3):
        if d>=4:
            continue
        print(d, end=';')
    else:
        print(d,end=';')
    print(d)
except:
    print('ERR')
code

5. Що буде виведено за виконання?\par (Дійсні значення, якщо наявні, в цьому завданні будуть виводитись у форматі з фіксованою крапкою та, можливо, з доволі великою кількістю десяткових розрядів. У відповіді для дійсних значень у ЦЬОМУ завданні достатньо вказати один десятковий розряд після крапки, відкинувши решту розрядів без округлення. Наприклад, замість 13.66666666666667 достатньо вказати 13.6, замість 25.23 достатньо вказати 25.2)
code
try:
    print(5, end=';')
    print(int('b2'), end=';')
    print(1, end=';')
except Exception:
    print(0, end=';')
except KeyboardInterrupt:
    print(2, end=';')
except:
    print('ERR', end=';')
else:
    print(0, end=';')
finally:
    print(7)
code

6. Що буде виведено за виконання?
code
def f(n):
    if n>9:
        f(n//100)
    else:
        print(n%10,end='')

f(987654)
code

7. Що буде виведено за виконання?
code
class A:
    a = 1

    def f1(self): print(2, end=';')
    def _f2(self): print(3, end=';')

    def main(self):
        try: print(self.a, end=';')
        except: print('ERR', end=';')

        try: self.f1()
        except: print('ERR', end=';')

        try: self._f2()
        except: print('ERR', end=';')


class B(A):
    a = 4

    def f1(self): print(5, end=';')
    def _f2(self): print(6, end=';')

try: B().main()
except: print('ERR', end=';')
code

8. Що буде виведено за виконання?
code
def f(arg, n):
    arg[23] = 8
    n = 7
    arg = set(arg.values())

try:
    n = 1
    d = {25:1, 23:5, 57:2, 57:0}
    f(d, n)
    print(n, end=';')
    for x, y in d.items():
        print(x, y, sep=';', end=';')
except:
    print('ERR')
code

9. Що буде виведено за виконання?
code
class A:
    b = 1
    c = 2

    def __init__(self):
        b = 3

class B(A):
    c = 4

    def __init__(self):
        super().__init__()
        self.a = 7

obj = B()
obj.c = obj.b + 10
B.c = 13

try: print(obj.c, end= ';')
except: print('ERR', end=';')

try: print(A.c, end= ';')
except: print('ERR', end=';')

try: print(B.c, end= ';')
except: print('ERR', end=';')
code

10. 
Для графа на рисунку з вершини 5 запущено алгоритм пошуку в глибину, що будує кістякове дерево. Списки суміжності вершин переглядаються алгоритмом за зростанням міток вершин.\par
\begin{center}
	\adjustimage{max size={0.9\linewidth}{0.9\paperheight}}
	{../10/Traversals/176.png}
\end{center}
\par Які з наведених ребер входять до складу отриманого кістякового дерева?\par
1) (3, 5)\par
2) (0, 3)\par
3) (5, 6)\par
4) (4, 6)\par
5) (0, 1)\par
6) (1, 2)\par
{\vskip 15pt} У формі вибрати номери відповідних ребер.\par
%answer:1 2 4 5 6
\end {large}}
%end
{\smallskip \begin {small} Затверджено на засіданні кафедри теоретичної кібернетики, протокол \No 12 від   19.05.2020 р.\par\smallskip\smallskip
{\parbox {6.5 cm}{Зав. кафедри \dotfill Ю.В.~Крак}}\hspace {0.7cm}{\hbox to 6.5 cm{ Екзаменатор \dotfill Т.О.~Карнаух}} \end{small}}
%begin
{{\begin {small}\par
{ \center  Київський національний університет імені Тараса Шевченка \par}
{\parbox {5 cm}{Освітня програма \hfill}{\parbox {7 cm}{Прикладна математика, Системний аналіз \hfill}}\parbox {6 cm}{\hfill Семестр 2}}\par
{\parbox {5 cm} {Навчальний предмет \hfill}\parbox {7 cm}{Програмування \hfill}\parbox {6cm}{\hfill}\par}\vspace{-7pt} \end{small}}{\center Екзаменаційний білет \No \textbf 
{ 200 }\par}}
{
\begin {large}
1. Що буде виведено за виконання?
code
try:
    res = 8>False or not 2>=6.0 or 4.0<3
    if isinstance(res, bool):
        print(f'bool;{res}')
    elif isinstance(res, int):
        print(f'int;{res}')
    elif isinstance(res, float):
        print(f'float;{res:.2f}')
    else:
        print('???')
except:
    print('ERR')
code

2. Що буде виведено за виконання?
code
a = 9
b = 5
c = 8

def h(a):
    a += 2
    b = 1
    c = 4
    return a + b + c

a = 3
b = 1
c = 4

try:
    print(h(a), a, b, c, sep=';')
except:
    print('ERR', a, b, c, sep=';')
code

3. Що буде виведено за виконання?
code
try:
    n = 12
    while n>=2:
        n += -2
    print(n, end=';')
except:
    print('ERR')
code

4. Що буде виведено за виконання?
code
try:
    for c in range(-5, -15, -1):
        if c>3:
            continue
        print(c, end=';');c = 9
    else:
        print(c,end=';')
    print(c)
except:
    print('ERR')
code

5. Що буде виведено за виконання?\par (Дійсні значення, якщо наявні, в цьому завданні будуть виводитись у форматі з фіксованою крапкою та, можливо, з доволі великою кількістю десяткових розрядів. У відповіді для дійсних значень у ЦЬОМУ завданні достатньо вказати один десятковий розряд після крапки, відкинувши решту розрядів без округлення. Наприклад, замість 13.66666666666667 достатньо вказати 13.6, замість 25.23 достатньо вказати 25.2)
code
try:
    print(8, end=';')
    print(6j>7j, end=';')
    print(3, end=';')
except ZeroDivisionError:
    print(9, end=';')
except TypeError:
    print(4, end=';')
except:
    print('ERR', end=';')
else:
    print(1, end=';')
finally:
    print(5)
code

6. Що буде виведено за виконання?
code
def f(n):
    if n>9:
        f(n//10)
    print(n%10, end='')
 
f(543216)
code

7. Що буде виведено за виконання?
code
class A:
    b = 1

    def _h1(self): print(2, end=';')
    def h2(self): print(3, end=';')

    def main(self):
        try: print(self.b, end=';')
        except: print('ERR', end=';')

        try: self._h1()
        except: print('ERR', end=';')

        try: self.h2()
        except: print('ERR', end=';')


class B(A):
    b = 4

    def _h1(self): print(5, end=';')
    def h2(self): print(6, end=';')

try: B().main()
except: print('ERR', end=';')
code

8. Що буде виведено за виконання?
code
def f(arg, n):
    n = 9
    tmp = arg
    del tmp[5:]
    return len(tmp)>3

try:
    h = [10,11,12,13,14]
    n = 4
    f(h, n)
    print(n, end=';')
    for el in h:
        print(el, end=';')
except:
    print('ERR')
code

9. Що буде виведено за виконання?
code
class A:
    c = 1
    b = 2

    def __init__(self):
        b = 3

class B(A):
    c = 4

    def __init__(self):
        super().__init__()
        self.a = 7

obj = B()
obj.a = obj.b + 10
B.a = 13

try: print(obj.a, end= ';')
except: print('ERR', end=';')

try: print(A.a, end= ';')
except: print('ERR', end=';')

try: print(B.a, end= ';')
except: print('ERR', end=';')
code

10. 
Для графа на рисунку з вершини 5 запущено алгоритм пошуку в глибину, що будує кістякове дерево. Списки суміжності вершин переглядаються алгоритмом за зростанням міток вершин.\par
\begin{center}
	\adjustimage{max size={0.9\linewidth}{0.9\paperheight}}
	{../10/Traversals/143.png}
\end{center}
\par Які з наведених ребер входять до складу отриманого кістякового дерева?\par
1) (0, 3)\par
2) (1, 3)\par
3) (2, 3)\par
4) (5, 6)\par
5) (3, 5)\par
6) (1, 6)\par
{\vskip 15pt} У формі вибрати номери відповідних ребер.\par
%answer:1 5 6
\end {large}}
%end
{\smallskip \begin {small} Затверджено на засіданні кафедри теоретичної кібернетики, протокол \No 12 від   19.05.2020 р.\par\smallskip\smallskip
{\parbox {6.5 cm}{Зав. кафедри \dotfill Ю.В.~Крак}}\hspace {0.7cm}{\hbox to 6.5 cm{ Екзаменатор \dotfill Т.О.~Карнаух}} \end{small}}
%begin
{{\begin {small}\par
{ \center  Київський національний університет імені Тараса Шевченка \par}
{\parbox {5 cm}{Освітня програма \hfill}{\parbox {7 cm}{Прикладна математика, Системний аналіз \hfill}}\parbox {6 cm}{\hfill Семестр 2}}\par
{\parbox {5 cm} {Навчальний предмет \hfill}\parbox {7 cm}{Програмування \hfill}\parbox {6cm}{\hfill}\par}\vspace{-7pt} \end{small}}{\center Екзаменаційний білет \No \textbf 
{ 201 }\par}}
{
\begin {large}
1. Що буде виведено за виконання?
code
try:
    res = 2**-3--1
    if isinstance(res, bool):
        print(f'bool;{res}')
    elif isinstance(res, int):
        print(f'int;{res}')
    elif isinstance(res, float):
        print(f'float;{res:.2f}')
    else:
        print('???')
except:
    print('ERR')
code

2. Що буде виведено за виконання?
code
a = 2
b = 6
c = 5

def g(b):
    a = 5
    b = 2
    c = 3
    return a + b + c

a = 7
b = 9
c = 0

try:
    print(g(a), a, b, c, sep=';')
except:
    print('ERR', a, b, c, sep=';')
code

3. Що буде виведено за виконання?
code
try:
    i = 10
    while i>1:
        i += -3
    print(i, end=';')
except:
    print('ERR')
code

4. Що буде виведено за виконання?
code
try:
    for f in range(14, 8, 2):
        if f<7:
            break
        print(f, end=';')
    else:
        print(f,end=';')
    print(f)
except:
    print('ERR')
code

5. Що буде виведено за виконання?\par (Дійсні значення, якщо наявні, в цьому завданні будуть виводитись у форматі з фіксованою крапкою та, можливо, з доволі великою кількістю десяткових розрядів. У відповіді для дійсних значень у ЦЬОМУ завданні достатньо вказати один десятковий розряд після крапки, відкинувши решту розрядів без округлення. Наприклад, замість 13.66666666666667 достатньо вказати 13.6, замість 25.23 достатньо вказати 25.2)
code
try:
    print(0, end=';')
    print(7%0.0, end=';')
    print(3, end=';')
except TypeError:
    print(5, end=';')
except ZeroDivisionError:
    print(8, end=';')
except:
    print('ERR', end=';')
else:
    print(6, end=';')
finally:
    print(2)
code

6. Що буде виведено за виконання?
code
def f(n):
    print(n%2, end='')
    if n>2:
        f(n//2)
 
f(16)
code

7. Що буде виведено за виконання?
code
class A:
    b = 1

    def _g1(self): print(2, end=';')
    def _g2(self): print(3, end=';')

    def main(self):
        try: print(self.b, end=';')
        except: print('ERR', end=';')

        try: self._g1()
        except: print('ERR', end=';')

        try: self._g2()
        except: print('ERR', end=';')


class B(A):
    b = 4

    def _g1(self): print(5, end=';')
    def _g2(self): print(6, end=';')

try: B().main()
except: print('ERR', end=';')
code

8. Що буде виведено за виконання?
code
def f(arg, n):
    n = 1
    tmp = arg
    tmp[2:] = 'ab'
    return len(tmp)>3

try:
    a = ['0','1','2','3','4']
    n = 8
    f(a, n)
    print(n, end=';')
    for el in a:
        print(el, end=';')
except:
    print('ERR')
code

9. Що буде виведено за виконання?
code
class A:
    a = 1
    c = 2

    def __init__(self):
        c = 3

class B(A):
    a = 4

    def __init__(self):
        super().__init__()
        self.b = 7

obj = B()
obj.a = obj.c + 10
B.a = 13

try: print(obj.a, end= ';')
except: print('ERR', end=';')

try: print(A.a, end= ';')
except: print('ERR', end=';')

try: print(B.a, end= ';')
except: print('ERR', end=';')
code

10. 
Для графа на рисунку з вершини 1 запущено алгоритм пошуку в глибину, що будує кістякове дерево. Списки суміжності вершин переглядаються алгоритмом за зростанням міток вершин.\par
\begin{center}
	\adjustimage{max size={0.9\linewidth}{0.9\paperheight}}
	{../10/Traversals/198.png}
\end{center}
\par Які з наведених ребер входять до складу отриманого кістякового дерева?\par
1) (0, 6)\par
2) (1, 6)\par
3) (3, 5)\par
4) (2, 5)\par
5) (1, 3)\par
6) (1, 5)\par
{\vskip 15pt} У формі вибрати номери відповідних ребер.\par
%answer:1 3 5
\end {large}}
%end
{\smallskip \begin {small} Затверджено на засіданні кафедри теоретичної кібернетики, протокол \No 12 від   19.05.2020 р.\par\smallskip\smallskip
{\parbox {6.5 cm}{Зав. кафедри \dotfill Ю.В.~Крак}}\hspace {0.7cm}{\hbox to 6.5 cm{ Екзаменатор \dotfill Т.О.~Карнаух}} \end{small}}
%begin
{{\begin {small}\par
{ \center  Київський національний університет імені Тараса Шевченка \par}
{\parbox {5 cm}{Освітня програма \hfill}{\parbox {7 cm}{Прикладна математика, Системний аналіз \hfill}}\parbox {6 cm}{\hfill Семестр 2}}\par
{\parbox {5 cm} {Навчальний предмет \hfill}\parbox {7 cm}{Програмування \hfill}\parbox {6cm}{\hfill}\par}\vspace{-7pt} \end{small}}{\center Екзаменаційний білет \No \textbf 
{ 202 }\par}}
{
\begin {large}
1. Що буде виведено за виконання?
code
try:
    res = 9**1.0*-2
    if isinstance(res, bool):
        print(f'bool;{res}')
    elif isinstance(res, int):
        print(f'int;{res}')
    elif isinstance(res, float):
        print(f'float;{res:.2f}')
    else:
        print('???')
except:
    print('ERR')
code

2. Що буде виведено за виконання?
code
a = 8
b = 1
c = 4

def h(a):
    global c
    a = 4
    b = 1
    c = 1
    return a + b + c

a = 3
b = 3
c = 5

try:
    print(h(a), a, b, c, sep=';')
except:
    print('ERR', a, b, c, sep=';')
code

3. Що буде виведено за виконання?
code
try:
    i = 8
    while i<=5:
        i += 1
    print(i, end=';')
except:
    print('ERR')
code

4. Що буде виведено за виконання?
code
try:
    for b in range(15, -11, -1):
        if b>=8:
            break
        print(b, end=';')
    else:
        print(b,end=';')
    print(b)
except:
    print('ERR')
code

5. Що буде виведено за виконання?\par (Дійсні значення, якщо наявні, в цьому завданні будуть виводитись у форматі з фіксованою крапкою та, можливо, з доволі великою кількістю десяткових розрядів. У відповіді для дійсних значень у ЦЬОМУ завданні достатньо вказати один десятковий розряд після крапки, відкинувши решту розрядів без округлення. Наприклад, замість 13.66666666666667 достатньо вказати 13.6, замість 25.23 достатньо вказати 25.2)
code
try:
    print(9, end=';')
    print(4/3, end=';')
    print(1, end=';')
except ValueError:
    print(2, end=';')
except Exception:
    print(0, end=';')
except:
    print('ERR', end=';')
else:
    print(6, end=';')
finally:
    print(8)
code

6. Що буде виведено за виконання?
code
def f(s):
    for el in s:
        el += 1
        return s
 
s = [1, 2, 3]
s = f(s)
print(s)
code

7. Що буде виведено за виконання?
code
class A:
    d = 1

    def _h1(self): print(2, end=';')
    def __h2(self): print(3, end=';')

    def main(self):
        try: print(self.d, end=';')
        except: print('ERR', end=';')

        try: self._h1()
        except: print('ERR', end=';')

        try: self.__h2()
        except: print('ERR', end=';')


class B(A):
    d = 4

    def _h1(self): print(5, end=';')
    def __h2(self): print(6, end=';')

try: B().main()
except: print('ERR', end=';')
code

8. Що буде виведено за виконання?
code
def f(arg, n):
    arg[34] = 0
    n = 6
    arg = list(arg.keys())

try:
    n = 4
    t = {82:7, 53:1, 79:6, 53:6}
    f(t, n)
    print(n, end=';')
    for x in t.items():
        print(*x, sep=';', end=';')
except:
    print('ERR')
code

9. Що буде виведено за виконання?
code
class A:
    a = 1
    b = 2

    def __init__(self):
        a = 3

class B(A):
    b = 4

    def __init__(self):
        super().__init__()
        self.c = 7

obj = B()
obj.b = obj.c + 10
B.b = 13

try: print(obj.b, end= ';')
except: print('ERR', end=';')

try: print(A.b, end= ';')
except: print('ERR', end=';')

try: print(B.b, end= ';')
except: print('ERR', end=';')
code

10. 
Для графа на рисунку з вершини 2 запущено алгоритм пошуку в глибину, що будує кістякове дерево. Списки суміжності вершин переглядаються алгоритмом за зростанням міток вершин.\par
\begin{center}
	\adjustimage{max size={0.9\linewidth}{0.9\paperheight}}
	{../10/Traversals/116.png}
\end{center}
\par Які з наведених ребер входять до складу отриманого кістякового дерева?\par
1) (2, 4)\par
2) (2, 6)\par
3) (5, 6)\par
4) (3, 6)\par
5) (0, 2)\par
6) (0, 6)\par
{\vskip 15pt} У формі вибрати номери відповідних ребер.\par
%answer:3 4 5
\end {large}}
%end
{\smallskip \begin {small} Затверджено на засіданні кафедри теоретичної кібернетики, протокол \No 12 від   19.05.2020 р.\par\smallskip\smallskip
{\parbox {6.5 cm}{Зав. кафедри \dotfill Ю.В.~Крак}}\hspace {0.7cm}{\hbox to 6.5 cm{ Екзаменатор \dotfill Т.О.~Карнаух}} \end{small}}
%begin
{{\begin {small}\par
{ \center  Київський національний університет імені Тараса Шевченка \par}
{\parbox {5 cm}{Освітня програма \hfill}{\parbox {7 cm}{Прикладна математика, Системний аналіз \hfill}}\parbox {6 cm}{\hfill Семестр 2}}\par
{\parbox {5 cm} {Навчальний предмет \hfill}\parbox {7 cm}{Програмування \hfill}\parbox {6cm}{\hfill}\par}\vspace{-7pt} \end{small}}{\center Екзаменаційний білет \No \textbf 
{ 203 }\par}}
{
\begin {large}
1. Що буде виведено за виконання?
code
def f(n):
    print('f', end='')
    return n

try:
    print(f(3) or f(6))
except:
    print('ERR')
code

2. Що буде виведено за виконання?
code
a = 2
b = 7
c = 4

def h(a):
    a = 5
    b = 4
    c = 2
    return a + b + c

a = 8
b = 0
c = 1

try:
    print(h(a), a, b, c, sep=';')
except:
    print('ERR', a, b, c, sep=';')
code

3. Що буде виведено за виконання?
code
try:
    j = 4
    while j<=8:
        j += 1
    print(j, end=';')
except:
    print('ERR')
code

4. Що буде виведено за виконання?
code
try:
    for a in range(-15, 4, -2):
        if a<=5:
            continue
        print(a, end=';')
    else:
        print(a,end=';')
    print(a)
except:
    print('ERR')
code

5. Що буде виведено за виконання?\par (Дійсні значення, якщо наявні, в цьому завданні будуть виводитись у форматі з фіксованою крапкою та, можливо, з доволі великою кількістю десяткових розрядів. У відповіді для дійсних значень у ЦЬОМУ завданні достатньо вказати один десятковий розряд після крапки, відкинувши решту розрядів без округлення. Наприклад, замість 13.66666666666667 достатньо вказати 13.6, замість 25.23 достатньо вказати 25.2)
code
try:
    print(3, end=';')
    print(9j<=0j, end=';')
    print(7, end=';')
except KeyboardInterrupt:
    print(5, end=';')
except TypeError:
    print(1, end=';')
except:
    print('ERR', end=';')
else:
    print(4, end=';')
finally:
    print(6)
code

6. Що буде виведено за виконання?
code
def f(n):
    if n>9:
        f(n//10)
    print(n%10, end='')
 
f(123456)
code

7. Що буде виведено за виконання?
code
class A:
    c = 1

    def _g1(self): print(2, end=';')
    def g2(self): print(3, end=';')

    def main(self):
        try: print(self.c, end=';')
        except: print('ERR', end=';')

        try: self._g1()
        except: print('ERR', end=';')

        try: self.g2()
        except: print('ERR', end=';')


class B(A):
    c = 4

    def _g1(self): print(5, end=';')
    def g2(self): print(6, end=';')

try: B().main()
except: print('ERR', end=';')
code

8. Що буде виведено за виконання?
code
def f(arg, n):
    n = 5
    tmp = arg
    del tmp[:-8]
    return len(tmp)>3

try:
    b = [10,11,12,13,14]
    n = 6
    f(b, n)
    print(n, end=';')
    for el in b:
        print(el, end=';')
except:
    print('ERR')
code

9. Що буде виведено за виконання?
code
class A:
    c = 1
    a = 2

    def __init__(self):
        c = 3

class B(A):
    a = 4

    def __init__(self):
        super().__init__()
        self.b = 7

obj = B()
obj.a = obj.a + 10
B.a = 13

try: print(obj.a, end= ';')
except: print('ERR', end=';')

try: print(A.a, end= ';')
except: print('ERR', end=';')

try: print(B.a, end= ';')
except: print('ERR', end=';')
code

10. 
Для графа на рисунку з вершини 6 запущено алгоритм пошуку в глибину, що будує кістякове дерево. Списки суміжності вершин переглядаються алгоритмом за зростанням міток вершин.\par
\begin{center}
	\adjustimage{max size={0.9\linewidth}{0.9\paperheight}}
	{../10/Traversals/172.png}
\end{center}
\par Які з наведених ребер входять до складу отриманого кістякового дерева?\par
1) (0, 3)\par
2) (1, 6)\par
3) (1, 4)\par
4) (2, 3)\par
5) (0, 2)\par
6) (4, 6)\par
{\vskip 15pt} У формі вибрати номери відповідних ребер.\par
%answer:3 4 5
\end {large}}
%end
{\smallskip \begin {small} Затверджено на засіданні кафедри теоретичної кібернетики, протокол \No 12 від   19.05.2020 р.\par\smallskip\smallskip
{\parbox {6.5 cm}{Зав. кафедри \dotfill Ю.В.~Крак}}\hspace {0.7cm}{\hbox to 6.5 cm{ Екзаменатор \dotfill Т.О.~Карнаух}} \end{small}}
%begin
{{\begin {small}\par
{ \center  Київський національний університет імені Тараса Шевченка \par}
{\parbox {5 cm}{Освітня програма \hfill}{\parbox {7 cm}{Прикладна математика, Системний аналіз \hfill}}\parbox {6 cm}{\hfill Семестр 2}}\par
{\parbox {5 cm} {Навчальний предмет \hfill}\parbox {7 cm}{Програмування \hfill}\parbox {6cm}{\hfill}\par}\vspace{-7pt} \end{small}}{\center Екзаменаційний білет \No \textbf 
{ 204 }\par}}
{
\begin {large}
1. Що буде виведено за виконання?
code
try:
    res = 9**-1.0*True
    if isinstance(res, bool):
        print(f'bool;{res}')
    elif isinstance(res, int):
        print(f'int;{res}')
    elif isinstance(res, float):
        print(f'float;{res:.2f}')
    else:
        print('???')
except:
    print('ERR')
code

2. Що буде виведено за виконання?
code
a = 9
b = 6
c = 5

def f(b):
    global c
    a = 3
    b = 3
    c = 1
    return a + b + c

a = 3
b = 4
c = 2

try:
    print(f(b), a, b, c, sep=';')
except:
    print('ERR', a, b, c, sep=';')
code

3. Що буде виведено за виконання?
code
try:
    m = 13
    while m>0:
        m += -3
    print(m, end=';')
except:
    print('ERR')
code

4. Що буде виведено за виконання?
code
try:
    for e in range(-8, -6, -3):
        if e>6:
            continue
        print(e, end=';');e = -1
    else:
        print('end',end=';')
    print(e)
except:
    print('ERR')
code

5. Що буде виведено за виконання?\par (Дійсні значення, якщо наявні, в цьому завданні будуть виводитись у форматі з фіксованою крапкою та, можливо, з доволі великою кількістю десяткових розрядів. У відповіді для дійсних значень у ЦЬОМУ завданні достатньо вказати один десятковий розряд після крапки, відкинувши решту розрядів без округлення. Наприклад, замість 13.66666666666667 достатньо вказати 13.6, замість 25.23 достатньо вказати 25.2)
code
try:
    print(7, end=';')
    print(int('d3'), end=';')
    print(2, end=';')
except Exception:
    print(9, end=';')
except KeyboardInterrupt:
    print(8, end=';')
except:
    print('ERR', end=';')
else:
    print(4, end=';')
finally:
    print(5)
code

6. Що буде виведено за виконання?
code
def f(n):
    print(n%10, end='')
    if n>9:
        f(n//10)
 
f(123456)
code

7. Що буде виведено за виконання?
code
class A:
    a = 1

    def f1(self): print(2, end=';')
    def _f2(self): print(3, end=';')

    def main(self):
        try: print(self.a, end=';')
        except: print('ERR', end=';')

        try: self.f1()
        except: print('ERR', end=';')

        try: self._f2()
        except: print('ERR', end=';')


class B(A):
    a = 4

    def f1(self): print(5, end=';')
    def _f2(self): print(6, end=';')

try: B().main()
except: print('ERR', end=';')
code

8. Що буде виведено за виконання?
code
def f(arg, n):
    arg[74] = 6
    n = 9
    arg = set(arg.items())

try:
    n = 3
    s = {66:3, 18:4, 64:4, 58:4}
    f(s, n)
    print(n, end=';')
    for x in s.values():
        print(x, sep=';', end=';')
except:
    print('ERR')
code

9. Що буде виведено за виконання?
code
class A:
    b = 1
    a = 2

    def __init__(self):
        a = 3

class B(A):
    b = 4

    def __init__(self):
        super().__init__()
        self.c = 7

obj = B()
obj.b = obj.c + 10
B.b = 13

try: print(obj.b, end= ';')
except: print('ERR', end=';')

try: print(A.b, end= ';')
except: print('ERR', end=';')

try: print(B.b, end= ';')
except: print('ERR', end=';')
code

10. 
Для графа на рисунку з вершини 3 запущено алгоритм пошуку в глибину, що будує кістякове дерево. Списки суміжності вершин переглядаються алгоритмом за зростанням міток вершин.\par
\begin{center}
	\adjustimage{max size={0.9\linewidth}{0.9\paperheight}}
	{../10/Traversals/61.png}
\end{center}
\par Які з наведених ребер входять до складу отриманого кістякового дерева?\par
1) (2, 3)\par
2) (1, 4)\par
3) (2, 5)\par
4) (0, 5)\par
5) (3, 4)\par
6) (3, 6)\par
{\vskip 15pt} У формі вибрати номери відповідних ребер.\par
%answer:1 4
\end {large}}
%end
{\smallskip \begin {small} Затверджено на засіданні кафедри теоретичної кібернетики, протокол \No 12 від   19.05.2020 р.\par\smallskip\smallskip
{\parbox {6.5 cm}{Зав. кафедри \dotfill Ю.В.~Крак}}\hspace {0.7cm}{\hbox to 6.5 cm{ Екзаменатор \dotfill Т.О.~Карнаух}} \end{small}}
%begin
{{\begin {small}\par
{ \center  Київський національний університет імені Тараса Шевченка \par}
{\parbox {5 cm}{Освітня програма \hfill}{\parbox {7 cm}{Прикладна математика, Системний аналіз \hfill}}\parbox {6 cm}{\hfill Семестр 2}}\par
{\parbox {5 cm} {Навчальний предмет \hfill}\parbox {7 cm}{Програмування \hfill}\parbox {6cm}{\hfill}\par}\vspace{-7pt} \end{small}}{\center Екзаменаційний білет \No \textbf 
{ 205 }\par}}
{
\begin {large}
1. Що буде виведено за виконання?
code
try:
    res = "False"<="2.0j"
    if isinstance(res, bool):
        print(f'bool;{res}')
    elif isinstance(res, int):
        print(f'int;{res}')
    elif isinstance(res, float):
        print(f'float;{res:.2f}')
    else:
        print('???')
except:
    print('ERR')
code

2. Що буде виведено за виконання?
code
a = 7
b = 1
c = 3

def g(a):
    a -= 1
    b = 2
    c = 4
    return a + b + c

a = 0
b = 6
c = 9

try:
    print(g(a), a, b, c, sep=';')
except:
    print('ERR', a, b, c, sep=';')
code

3. Що буде виведено за виконання?
code
try:
    t = 9
    while t<7:
        t += 2
    print(t, end=';')
except:
    print('ERR')
code

4. Що буде виведено за виконання?
code
try:
    for c in range(12, -7, 3):
        if c>3:
            break
        print(c, end=';')
    else:
        print(c,end=';')
    print(c)
except:
    print('ERR')
code

5. Що буде виведено за виконання?\par (Дійсні значення, якщо наявні, в цьому завданні будуть виводитись у форматі з фіксованою крапкою та, можливо, з доволі великою кількістю десяткових розрядів. У відповіді для дійсних значень у ЦЬОМУ завданні достатньо вказати один десятковий розряд після крапки, відкинувши решту розрядів без округлення. Наприклад, замість 13.66666666666667 достатньо вказати 13.6, замість 25.23 достатньо вказати 25.2)
code
try:
    print(0, end=';')
    print(int('1'), end=';')
    print(0, end=';')
except ValueError:
    print(9, end=';')
except ZeroDivisionError:
    print(2, end=';')
except:
    print('ERR', end=';')
else:
    print(6, end=';')
finally:
    print(4)
code

6. Що буде виведено за виконання?
code
def f(n):
    if n>9:
        return f(n//10)+1
    else:
        return 1

print(f(123456))
code

7. Що буде виведено за виконання?
code
class A:
    c = 1

    def __f1(self): print(2, end=';')
    def _f2(self): print(3, end=';')

    def main(self):
        try: print(self.c, end=';')
        except: print('ERR', end=';')

        try: self.__f1()
        except: print('ERR', end=';')

        try: self._f2()
        except: print('ERR', end=';')


class B(A):
    c = 4

    def __f1(self): print(5, end=';')
    def _f2(self): print(6, end=';')

try: B().main()
except: print('ERR', end=';')
code

8. Що буде виведено за виконання?
code
def f(arg, n):
    n = 4
    tmp = arg
    tmp[:7] = 'bd'
    return len(tmp)>3

try:
    c = ['0','1','2','3','4']
    n = 0
    f(c, n)
    print(n, end=';')
    for el in c:
        print(el, end=';')
except:
    print('ERR')
code

9. Що буде виведено за виконання?
code
class A:
    c = 1
    a = 2

    def __init__(self):
        a = 3

class B(A):
    a = 4

    def __init__(self):
        super().__init__()
        self.b = 7

obj = B()
obj.c = obj.a + 10
B.c = 13

try: print(obj.c, end= ';')
except: print('ERR', end=';')

try: print(A.c, end= ';')
except: print('ERR', end=';')

try: print(B.c, end= ';')
except: print('ERR', end=';')
code

10. 
Для графа на рисунку з вершини 5 запущено алгоритм пошуку в глибину, що будує кістякове дерево. Списки суміжності вершин переглядаються алгоритмом за зростанням міток вершин.\par
\begin{center}
	\adjustimage{max size={0.9\linewidth}{0.9\paperheight}}
	{../10/Traversals/166.png}
\end{center}
\par Які з наведених ребер входять до складу отриманого кістякового дерева?\par
1) (4, 6)\par
2) (3, 6)\par
3) (1, 3)\par
4) (0, 5)\par
5) (0, 2)\par
6) (1, 5)\par
{\vskip 15pt} У формі вибрати номери відповідних ребер.\par
%answer:1 2 3 4 5
\end {large}}
%end
{\smallskip \begin {small} Затверджено на засіданні кафедри теоретичної кібернетики, протокол \No 12 від   19.05.2020 р.\par\smallskip\smallskip
{\parbox {6.5 cm}{Зав. кафедри \dotfill Ю.В.~Крак}}\hspace {0.7cm}{\hbox to 6.5 cm{ Екзаменатор \dotfill Т.О.~Карнаух}} \end{small}}
%begin
{{\begin {small}\par
{ \center  Київський національний університет імені Тараса Шевченка \par}
{\parbox {5 cm}{Освітня програма \hfill}{\parbox {7 cm}{Прикладна математика, Системний аналіз \hfill}}\parbox {6 cm}{\hfill Семестр 2}}\par
{\parbox {5 cm} {Навчальний предмет \hfill}\parbox {7 cm}{Програмування \hfill}\parbox {6cm}{\hfill}\par}\vspace{-7pt} \end{small}}{\center Екзаменаційний білет \No \textbf 
{ 206 }\par}}
{
\begin {large}
1. Що буде виведено за виконання?
code
try:
    res = "8"==5j
    if isinstance(res, bool):
        print(f'bool;{res}')
    elif isinstance(res, int):
        print(f'int;{res}')
    elif isinstance(res, float):
        print(f'float;{res:.2f}')
    else:
        print('???')
except:
    print('ERR')
code

2. Що буде виведено за виконання?
code
a = 8
b = 5
c = 4

def f(b):
    global c
    a = 3
    b *= 5
    c = 5
    return a + b + c

a = 2
b = 5
c = 7

try:
    print(f(b), a, b, c, sep=';')
except:
    print('ERR', a, b, c, sep=';')
code

3. Що буде виведено за виконання?
code
try:
    t = 15
    while t>=8:
        t += -2
    print(t, end=';')
except:
    print('ERR')
code

4. Що буде виведено за виконання?
code
try:
    for e in range(-9, 5, 2):
        if e>=5:
            continue
        print(e, end=';');e = -9
        if e<=3:
            break
    else:
        print(e,end=';')
    print(e)
except:
    print('ERR')
code

5. Що буде виведено за виконання?\par (Дійсні значення, якщо наявні, в цьому завданні будуть виводитись у форматі з фіксованою крапкою та, можливо, з доволі великою кількістю десяткових розрядів. У відповіді для дійсних значень у ЦЬОМУ завданні достатньо вказати один десятковий розряд після крапки, відкинувши решту розрядів без округлення. Наприклад, замість 13.66666666666667 достатньо вказати 13.6, замість 25.23 достатньо вказати 25.2)
code
try:
    print(8, end=';')
    print(int('1'), end=';')
    print(7, end=';')
except KeyboardInterrupt:
    print(5, end=';')
except TypeError:
    print(3, end=';')
except:
    print('ERR', end=';')
else:
    print(1, end=';')
finally:
    print(8)
code

6. Що буде виведено за виконання?
code
def f(n):
    if n<=9:
        print(n%10,end='')
    else:
        f(n//10)

f(543216)
code

7. Що буде виведено за виконання?
code
class A:
    d = 1

    def _h1(self): print(2, end=';')
    def __h2(self): print(3, end=';')

    def main(self):
        try: print(self.d, end=';')
        except: print('ERR', end=';')

        try: self._h1()
        except: print('ERR', end=';')

        try: self.__h2()
        except: print('ERR', end=';')


class B(A):
    d = 4

    def _h1(self): print(5, end=';')
    def __h2(self): print(6, end=';')

try: B().main()
except: print('ERR', end=';')
code

8. Що буде виведено за виконання?
code
def f(arg, n):
    n = 5
    tmp = arg
    tmp[5:-5] = [2,3]
    return len(tmp)>3

try:
    g = [10,11,12,13,14]
    n = 8
    f(g, n)
    print(n, end=';')
    for el in g:
        print(el, end=';')
except:
    print('ERR')
code

9. Що буде виведено за виконання?
code
class A:
    a = 1
    c = 2

    def __init__(self):
        a = 3

class B(A):
    a = 4

    def __init__(self):
        super().__init__()
        self.b = 7

obj = B()
obj.b = obj.c + 10
B.b = 13

try: print(obj.b, end= ';')
except: print('ERR', end=';')

try: print(A.b, end= ';')
except: print('ERR', end=';')

try: print(B.b, end= ';')
except: print('ERR', end=';')
code

10. 
Для графа на рисунку з вершини 0 запущено алгоритм пошуку в глибину, що будує кістякове дерево. Списки суміжності вершин переглядаються алгоритмом за зростанням міток вершин.\par
\begin{center}
	\adjustimage{max size={0.9\linewidth}{0.9\paperheight}}
	{../10/Traversals/270.png}
\end{center}
\par Які з наведених ребер входять до складу отриманого кістякового дерева?\par
1) (1, 4)\par
2) (3, 4)\par
3) (0, 2)\par
4) (3, 6)\par
5) (0, 1)\par
6) (5, 6)\par
{\vskip 15pt} У формі вибрати номери відповідних ребер.\par
%answer:1 4 5 6
\end {large}}
%end
{\smallskip \begin {small} Затверджено на засіданні кафедри теоретичної кібернетики, протокол \No 12 від   19.05.2020 р.\par\smallskip\smallskip
{\parbox {6.5 cm}{Зав. кафедри \dotfill Ю.В.~Крак}}\hspace {0.7cm}{\hbox to 6.5 cm{ Екзаменатор \dotfill Т.О.~Карнаух}} \end{small}}
%begin
{{\begin {small}\par
{ \center  Київський національний університет імені Тараса Шевченка \par}
{\parbox {5 cm}{Освітня програма \hfill}{\parbox {7 cm}{Прикладна математика, Системний аналіз \hfill}}\parbox {6 cm}{\hfill Семестр 2}}\par
{\parbox {5 cm} {Навчальний предмет \hfill}\parbox {7 cm}{Програмування \hfill}\parbox {6cm}{\hfill}\par}\vspace{-7pt} \end{small}}{\center Екзаменаційний білет \No \textbf 
{ 207 }\par}}
{
\begin {large}
1. Що буде виведено за виконання?
code
try:
    res = not True==3 and 9.0>=7 and 2<=6.0
    if isinstance(res, bool):
        print(f'bool;{res}')
    elif isinstance(res, int):
        print(f'int;{res}')
    elif isinstance(res, float):
        print(f'float;{res:.2f}')
    else:
        print('???')
except:
    print('ERR')
code

2. Що буде виведено за виконання?
code
a = 3
b = 9
c = 2

def g(b):
    a = 3
    b -= 2
    c = 4
    return a + b + c

a = 8
b = 4
c = 6

try:
    print(g(a), a, b, c, sep=';')
except:
    print('ERR', a, b, c, sep=';')
code

3. Що буде виведено за виконання?
code
try:
    k = 6
    while k<11:
        k += 1
    print(k, end=';')
except:
    print('ERR')
code

4. Що буде виведено за виконання?
code
try:
    for b in range(-6, 1, -2):
        if b<8:
            continue
        print(b, end=';')
    else:
        print(b,end=';')
    print(b)
except:
    print('ERR')
code

5. Що буде виведено за виконання?\par (Дійсні значення, якщо наявні, в цьому завданні будуть виводитись у форматі з фіксованою крапкою та, можливо, з доволі великою кількістю десяткових розрядів. У відповіді для дійсних значень у ЦЬОМУ завданні достатньо вказати один десятковий розряд після крапки, відкинувши решту розрядів без округлення. Наприклад, замість 13.66666666666667 достатньо вказати 13.6, замість 25.23 достатньо вказати 25.2)
code
try:
    print(7, end=';')
    print(5%2, end=';')
    print(2, end=';')
except ZeroDivisionError:
    print(9, end=';')
except Exception:
    print(3, end=';')
except:
    print('ERR', end=';')
else:
    print(6, end=';')
finally:
    print(4)
code

6. Що буде виведено за виконання?
code
def f(n):
    if n>2:
        f(n//2)
    else:
        print(n%2, end='')

f(16)
code

7. Що буде виведено за виконання?
code
class A:
    b = 1

    def _g1(self): print(2, end=';')
    def _g2(self): print(3, end=';')

    def main(self):
        try: print(self.b, end=';')
        except: print('ERR', end=';')

        try: self._g1()
        except: print('ERR', end=';')

        try: self._g2()
        except: print('ERR', end=';')


class B(A):
    b = 4

    def _g1(self): print(5, end=';')
    def _g2(self): print(6, end=';')

try: B().main()
except: print('ERR', end=';')
code

8. Що буде виведено за виконання?
code
def f(arg, n):
    arg[57] = 8
    n *= -2

try:
    n = 4
    d = {53:7, 57:6, 45:3, 45:9}
    f(d, n)
    print(n, end=';')
    for x in d.items():
        print(x, sep=';', end=';')
except:
    print('ERR')
code

9. Що буде виведено за виконання?
code
class A:
    a = 1
    b = 2

    def __init__(self):
        b = 3

class B(A):
    a = 4

    def __init__(self):
        super().__init__()
        self.c = 7

obj = B()
obj.b = obj.a + 10
B.b = 13

try: print(obj.b, end= ';')
except: print('ERR', end=';')

try: print(A.b, end= ';')
except: print('ERR', end=';')

try: print(B.b, end= ';')
except: print('ERR', end=';')
code

10. 
Для графа на рисунку з вершини 0 запущено алгоритм пошуку в глибину, що будує кістякове дерево. Списки суміжності вершин переглядаються алгоритмом за зростанням міток вершин.\par
\begin{center}
	\adjustimage{max size={0.9\linewidth}{0.9\paperheight}}
	{../10/Traversals/131.png}
\end{center}
\par Які з наведених ребер входять до складу отриманого кістякового дерева?\par
1) (2, 3)\par
2) (1, 3)\par
3) (0, 5)\par
4) (1, 6)\par
5) (1, 5)\par
6) (2, 6)\par
{\vskip 15pt} У формі вибрати номери відповідних ребер.\par
%answer:1 2 6
\end {large}}
%end
{\smallskip \begin {small} Затверджено на засіданні кафедри теоретичної кібернетики, протокол \No 12 від   19.05.2020 р.\par\smallskip\smallskip
{\parbox {6.5 cm}{Зав. кафедри \dotfill Ю.В.~Крак}}\hspace {0.7cm}{\hbox to 6.5 cm{ Екзаменатор \dotfill Т.О.~Карнаух}} \end{small}}
%begin
{{\begin {small}\par
{ \center  Київський національний університет імені Тараса Шевченка \par}
{\parbox {5 cm}{Освітня програма \hfill}{\parbox {7 cm}{Прикладна математика, Системний аналіз \hfill}}\parbox {6 cm}{\hfill Семестр 2}}\par
{\parbox {5 cm} {Навчальний предмет \hfill}\parbox {7 cm}{Програмування \hfill}\parbox {6cm}{\hfill}\par}\vspace{-7pt} \end{small}}{\center Екзаменаційний білет \No \textbf 
{ 208 }\par}}
{
\begin {large}
1. Що буде виведено за виконання?
code
try:
    res = -2**False+False
    if isinstance(res, bool):
        print(f'bool;{res}')
    elif isinstance(res, int):
        print(f'int;{res}')
    elif isinstance(res, float):
        print(f'float;{res:.2f}')
    else:
        print('???')
except:
    print('ERR')
code

2. Що буде виведено за виконання?
code
a = 1
b = 0
c = 7

def f(a):
    global c
    a += 1
    b = 1
    c = 4
    return a + b + c

a = 9
b = 6
c = 0

try:
    print(f(a), a, b, c, sep=';')
except:
    print('ERR', a, b, c, sep=';')
code

3. Що буде виведено за виконання?
code
try:
    m = 4
    while m<=12:
        m += 2
    print(m, end=';')
except:
    print('ERR')
code

4. Що буде виведено за виконання?
code
try:
    for d in range(7, 6, 3):
        if d<=6:
            continue
        print(d, end=';')
    else:
        print('end',end=';')
    print(d)
except:
    print('ERR')
code

5. Що буде виведено за виконання?\par (Дійсні значення, якщо наявні, в цьому завданні будуть виводитись у форматі з фіксованою крапкою та, можливо, з доволі великою кількістю десяткових розрядів. У відповіді для дійсних значень у ЦЬОМУ завданні достатньо вказати один десятковий розряд після крапки, відкинувши решту розрядів без округлення. Наприклад, замість 13.66666666666667 достатньо вказати 13.6, замість 25.23 достатньо вказати 25.2)
code
try:
    print(0, end=';')
    print(2//0, end=';')
    print(0, end=';')
except ValueError:
    print(1, end=';')
except Exception:
    print(4, end=';')
except:
    print('ERR', end=';')
else:
    print(6, end=';')
finally:
    print(8)
code

6. Що буде виведено за виконання?
code
def f(n):
    if n>9:
        f(n//100)
    print(n%10,end='')
 
f(987654)
code

7. Що буде виведено за виконання?
code
class A:
    a = 1

    def __f1(self): print(2, end=';')
    def _f2(self): print(3, end=';')

    def main(self):
        try: print(self.a, end=';')
        except: print('ERR', end=';')

        try: self.__f1()
        except: print('ERR', end=';')

        try: self._f2()
        except: print('ERR', end=';')


class B(A):
    a = 4

    def __f1(self): print(5, end=';')
    def _f2(self): print(6, end=';')

try: B().main()
except: print('ERR', end=';')
code

8. Що буде виведено за виконання?
code
def f(arg, n):
    arg[22] = 8
    n -= -3

try:
    n = 6
    t = {78:9, 22:3, 22:4}
    f(t, n)
    print(n, end=';')
    for x in t.keys():
        print(x, sep=';', end=';')
except:
    print('ERR')
code

9. Що буде виведено за виконання?
code
class A:
    b = 1
    c = 2

    def __init__(self):
        b = 3

class B(A):
    c = 4

    def __init__(self):
        super().__init__()
        self.a = 7

obj = B()
obj.a = obj.b + 10
B.a = 13

try: print(obj.a, end= ';')
except: print('ERR', end=';')

try: print(A.a, end= ';')
except: print('ERR', end=';')

try: print(B.a, end= ';')
except: print('ERR', end=';')
code

10. 
Для графа на рисунку з вершини 4 запущено алгоритм пошуку в глибину, що будує кістякове дерево. Списки суміжності вершин переглядаються алгоритмом за зростанням міток вершин.\par
\begin{center}
	\adjustimage{max size={0.9\linewidth}{0.9\paperheight}}
	{../10/Traversals/167.png}
\end{center}
\par Які з наведених ребер входять до складу отриманого кістякового дерева?\par
1) (1, 4)\par
2) (0, 3)\par
3) (1, 5)\par
4) (0, 2)\par
5) (1, 3)\par
6) (0, 1)\par
{\vskip 15pt} У формі вибрати номери відповідних ребер.\par
%answer:1 3 4 6
\end {large}}
%end
{\smallskip \begin {small} Затверджено на засіданні кафедри теоретичної кібернетики, протокол \No 12 від   19.05.2020 р.\par\smallskip\smallskip
{\parbox {6.5 cm}{Зав. кафедри \dotfill Ю.В.~Крак}}\hspace {0.7cm}{\hbox to 6.5 cm{ Екзаменатор \dotfill Т.О.~Карнаух}} \end{small}}
%begin
{{\begin {small}\par
{ \center  Київський національний університет імені Тараса Шевченка \par}
{\parbox {5 cm}{Освітня програма \hfill}{\parbox {7 cm}{Прикладна математика, Системний аналіз \hfill}}\parbox {6 cm}{\hfill Семестр 2}}\par
{\parbox {5 cm} {Навчальний предмет \hfill}\parbox {7 cm}{Програмування \hfill}\parbox {6cm}{\hfill}\par}\vspace{-7pt} \end{small}}{\center Екзаменаційний білет \No \textbf 
{ 209 }\par}}
{
\begin {large}
1. Що буде виведено за виконання?
code
try:
    res = 4//3%6*3
    if isinstance(res, bool):
        print(f'bool;{res}')
    elif isinstance(res, int):
        print(f'int;{res}')
    elif isinstance(res, float):
        print(f'float;{res:.2f}')
    else:
        print('???')
except:
    print('ERR')
code

2. Що буде виведено за виконання?
code
a = 5
b = 3
c = 4

def h(b):
    a = 2
    b *= 5
    c = 3
    return a + b + c

a = 1
b = 8
c = 2

try:
    print(h(b), a, b, c, sep=';')
except:
    print('ERR', a, b, c, sep=';')
code

3. Що буде виведено за виконання?
code
try:
    j = 7
    while j>6:
        j += -3
    print(j, end=';')
except:
    print('ERR')
code

4. Що буде виведено за виконання?
code
try:
    for a in range(-4, 12, -3):
        if a>4:
            continue
        print(a, end=';');a = 8
    else:
        print(a,end=';')
    print(a)
except:
    print('ERR')
code

5. Що буде виведено за виконання?\par (Дійсні значення, якщо наявні, в цьому завданні будуть виводитись у форматі з фіксованою крапкою та, можливо, з доволі великою кількістю десяткових розрядів. У відповіді для дійсних значень у ЦЬОМУ завданні достатньо вказати один десятковий розряд після крапки, відкинувши решту розрядів без округлення. Наприклад, замість 13.66666666666667 достатньо вказати 13.6, замість 25.23 достатньо вказати 25.2)
code
try:
    print(5, end=';')
    print(7j!=3j, end=';')
    print(3, end=';')
except TypeError:
    print(9, end=';')
except KeyboardInterrupt:
    print(3, end=';')
except:
    print('ERR', end=';')
else:
    print(2, end=';')
finally:
    print(4)
code

6. Що буде виведено за виконання?
code
def f(n):
    if n>9:
        f(n//100)
        print(n%10, end='')
 
f(123456)
code

7. Що буде виведено за виконання?
code
class A:
    d = 1

    def g1(self): print(2, end=';')
    def _g2(self): print(3, end=';')

    def main(self):
        try: print(self.d, end=';')
        except: print('ERR', end=';')

        try: self.g1()
        except: print('ERR', end=';')

        try: self._g2()
        except: print('ERR', end=';')


class B(A):
    d = 4

    def g1(self): print(5, end=';')
    def _g2(self): print(6, end=';')

try: B().main()
except: print('ERR', end=';')
code

8. Що буде виведено за виконання?
code
def f(arg, n):
    arg[34] = 4
    n += 2

try:
    n = 7
    s = {75:1, 34:4, 55:9, 34:5}
    f(s, n)
    print(n, end=';')
    for x in s.keys():
        print(x, sep=';', end=';')
except:
    print('ERR')
code

9. Що буде виведено за виконання?
code
class A:
    c = 1
    b = 2

    def __init__(self):
        b = 3

class B(A):
    b = 4

    def __init__(self):
        super().__init__()
        self.a = 7

obj = B()
obj.c = obj.a + 10
B.c = 13

try: print(obj.c, end= ';')
except: print('ERR', end=';')

try: print(A.c, end= ';')
except: print('ERR', end=';')

try: print(B.c, end= ';')
except: print('ERR', end=';')
code

10. 
Для графа на рисунку з вершини 2 запущено алгоритм пошуку в глибину, що будує кістякове дерево. Списки суміжності вершин переглядаються алгоритмом за зростанням міток вершин.\par
\begin{center}
	\adjustimage{max size={0.9\linewidth}{0.9\paperheight}}
	{../10/Traversals/104.png}
\end{center}
\par Які з наведених ребер входять до складу отриманого кістякового дерева?\par
1) (0, 4)\par
2) (3, 6)\par
3) (0, 5)\par
4) (1, 2)\par
5) (2, 6)\par
6) (0, 2)\par
{\vskip 15pt} У формі вибрати номери відповідних ребер.\par
%answer:1 2 3 6
\end {large}}
%end
{\smallskip \begin {small} Затверджено на засіданні кафедри теоретичної кібернетики, протокол \No 12 від   19.05.2020 р.\par\smallskip\smallskip
{\parbox {6.5 cm}{Зав. кафедри \dotfill Ю.В.~Крак}}\hspace {0.7cm}{\hbox to 6.5 cm{ Екзаменатор \dotfill Т.О.~Карнаух}} \end{small}}
%begin
{{\begin {small}\par
{ \center  Київський національний університет імені Тараса Шевченка \par}
{\parbox {5 cm}{Освітня програма \hfill}{\parbox {7 cm}{Прикладна математика, Системний аналіз \hfill}}\parbox {6 cm}{\hfill Семестр 2}}\par
{\parbox {5 cm} {Навчальний предмет \hfill}\parbox {7 cm}{Програмування \hfill}\parbox {6cm}{\hfill}\par}\vspace{-7pt} \end{small}}{\center Екзаменаційний білет \No \textbf 
{ 210 }\par}}
{
\begin {large}
1. Що буде виведено за виконання?
code
def f(n):
    print('f', end='')
    return n>4

try:
    print(f(0) and f(2))
except:
    print('ERR')
code

2. Що буде виведено за виконання?
code
a = 7
b = 9
c = 0

def g(b):
    a = 5
    b = 4
    c = 2
    return a + b + c

a = 6
b = 8
c = 1

try:
    print(g(b), a, b, c, sep=';')
except:
    print('ERR', a, b, c, sep=';')
code

3. Що буде виведено за виконання?
code
try:
    k = 6
    while k>=3:
        k += -2
    print(k, end=';')
except:
    print('ERR')
code

4. Що буде виведено за виконання?
code
try:
    for f in range(-7, 7, -1):
        if f>=7:
            break
        print(f, end=';')
    else:
        print(f,end=';')
    print(f)
except:
    print('ERR')
code

5. Що буде виведено за виконання?\par (Дійсні значення, якщо наявні, в цьому завданні будуть виводитись у форматі з фіксованою крапкою та, можливо, з доволі великою кількістю десяткових розрядів. У відповіді для дійсних значень у ЦЬОМУ завданні достатньо вказати один десятковий розряд після крапки, відкинувши решту розрядів без округлення. Наприклад, замість 13.66666666666667 достатньо вказати 13.6, замість 25.23 достатньо вказати 25.2)
code
try:
    print(0, end=';')
    print(int('a5'), end=';')
    print(7, end=';')
except ValueError:
    print(6, end=';')
except ZeroDivisionError:
    print(8, end=';')
except:
    print('ERR', end=';')
else:
    print(9, end=';')
finally:
    print(1)
code

6. Що буде виведено за виконання?
code
def f(n):
    if n<=9:
        print(n%10, end='')
    else:
        f(n//10)

f(987651)
code

7. Що буде виведено за виконання?
code
class A:
    c = 1

    def _h1(self): print(2, end=';')
    def h2(self): print(3, end=';')

    def main(self):
        try: print(self.c, end=';')
        except: print('ERR', end=';')

        try: self._h1()
        except: print('ERR', end=';')

        try: self.h2()
        except: print('ERR', end=';')


class B(A):
    c = 4

    def _h1(self): print(5, end=';')
    def h2(self): print(6, end=';')

try: B().main()
except: print('ERR', end=';')
code

8. Що буде виведено за виконання?
code
def f(arg, n):
    n = 9
    tmp = arg
    tmp[4:1] = 'bc'
    return len(tmp)>3

try:
    a = [10,11,12,13,14]
    n = 2
    f(a, n)
    print(n, end=';')
    for el in a:
        print(el, end=';')
except:
    print('ERR')
code

9. Що буде виведено за виконання?
code
class A:
    b = 1
    a = 2

    def __init__(self):
        b = 3

class B(A):
    b = 4

    def __init__(self):
        super().__init__()
        self.c = 7

obj = B()
obj.c = obj.b + 10
B.c = 13

try: print(obj.c, end= ';')
except: print('ERR', end=';')

try: print(A.c, end= ';')
except: print('ERR', end=';')

try: print(B.c, end= ';')
except: print('ERR', end=';')
code

10. 
Для графа на рисунку з вершини 1 запущено алгоритм пошуку в глибину, що будує кістякове дерево. Списки суміжності вершин переглядаються алгоритмом за зростанням міток вершин.\par
\begin{center}
	\adjustimage{max size={0.9\linewidth}{0.9\paperheight}}
	{../10/Traversals/11.png}
\end{center}
\par Які з наведених ребер входять до складу отриманого кістякового дерева?\par
1) (2, 4)\par
2) (4, 6)\par
3) (5, 6)\par
4) (0, 3)\par
5) (3, 6)\par
6) (0, 5)\par
{\vskip 15pt} У формі вибрати номери відповідних ребер.\par
%answer:3 4
\end {large}}
%end
{\smallskip \begin {small} Затверджено на засіданні кафедри теоретичної кібернетики, протокол \No 12 від   19.05.2020 р.\par\smallskip\smallskip
{\parbox {6.5 cm}{Зав. кафедри \dotfill Ю.В.~Крак}}\hspace {0.7cm}{\hbox to 6.5 cm{ Екзаменатор \dotfill Т.О.~Карнаух}} \end{small}}
%begin
{{\begin {small}\par
{ \center  Київський національний університет імені Тараса Шевченка \par}
{\parbox {5 cm}{Освітня програма \hfill}{\parbox {7 cm}{Прикладна математика, Системний аналіз \hfill}}\parbox {6 cm}{\hfill Семестр 2}}\par
{\parbox {5 cm} {Навчальний предмет \hfill}\parbox {7 cm}{Програмування \hfill}\parbox {6cm}{\hfill}\par}\vspace{-7pt} \end{small}}{\center Екзаменаційний білет \No \textbf 
{ 211 }\par}}
{
\begin {large}
1. Що буде виведено за виконання?
code
def f(n):
    print('f', end='')
    return n

try:
    print(f(-1) or f(-2))
except:
    print('ERR')
code

2. Що буде виведено за виконання?
code
a = 7
b = 0
c = 4

def f(a):
    global c
    a = 4
    b = 2
    c = 5
    return a + b + c

a = 9
b = 5
c = 8

try:
    print(f(a), a, b, c, sep=';')
except:
    print('ERR', a, b, c, sep=';')
code

3. Що буде виведено за виконання?
code
try:
    m = 8
    while m>4:
        m += -2
    print(m, end=';')
except:
    print('ERR')
code

4. Що буде виведено за виконання?
code
try:
    for a in range(5, -2, -3):
        if a<=6:
            break
        print(a, end=';')
    else:
        print(a,end=';')
    print(a)
except:
    print('ERR')
code

5. Що буде виведено за виконання?\par (Дійсні значення, якщо наявні, в цьому завданні будуть виводитись у форматі з фіксованою крапкою та, можливо, з доволі великою кількістю десяткових розрядів. У відповіді для дійсних значень у ЦЬОМУ завданні достатньо вказати один десятковий розряд після крапки, відкинувши решту розрядів без округлення. Наприклад, замість 13.66666666666667 достатньо вказати 13.6, замість 25.23 достатньо вказати 25.2)
code
try:
    print(4, end=';')
    print(1%3, end=';')
    print(2, end=';')
except KeyboardInterrupt:
    print(5, end=';')
except Exception:
    print(0, end=';')
except:
    print('ERR', end=';')
else:
    print(3, end=';')
finally:
    print(8)
code

6. Що буде виведено за виконання?
code
def f(n):
    if n<=9:
        print(n%10, end='')
    else:
        f(n//10)

f(123456)
code

7. Що буде виведено за виконання?
code
class A:
    a = 1

    def _g1(self): print(2, end=';')
    def __g2(self): print(3, end=';')

    def main(self):
        try: print(self.a, end=';')
        except: print('ERR', end=';')

        try: self._g1()
        except: print('ERR', end=';')

        try: self.__g2()
        except: print('ERR', end=';')


class B(A):
    a = 4

    def _g1(self): print(5, end=';')
    def __g2(self): print(6, end=';')

try: B().main()
except: print('ERR', end=';')
code

8. Що буде виведено за виконання?
code
def f(arg, n):
    n = 7
    tmp = arg
    del tmp[1:-7]
    return len(tmp)>3

try:
    g = ['0','1','2','3','4']
    n = 3
    f(g, n)
    print(n, end=';')
    for el in g:
        print(el, end=';')
except:
    print('ERR')
code

9. Що буде виведено за виконання?
code
class A:
    c = 1
    b = 2

    def __init__(self):
        b = 3

class B(A):
    c = 4

    def __init__(self):
        super().__init__()
        self.a = 7

obj = B()
obj.c = obj.a + 10
B.c = 13

try: print(obj.c, end= ';')
except: print('ERR', end=';')

try: print(A.c, end= ';')
except: print('ERR', end=';')

try: print(B.c, end= ';')
except: print('ERR', end=';')
code

10. 
Для графа на рисунку з вершини 4 запущено алгоритм пошуку в глибину, що будує кістякове дерево. Списки суміжності вершин переглядаються алгоритмом за зростанням міток вершин.\par
\begin{center}
	\adjustimage{max size={0.9\linewidth}{0.9\paperheight}}
	{../10/Traversals/226.png}
\end{center}
\par Які з наведених ребер входять до складу отриманого кістякового дерева?\par
1) (2, 3)\par
2) (4, 5)\par
3) (0, 6)\par
4) (1, 5)\par
5) (5, 6)\par
6) (0, 2)\par
{\vskip 15pt} У формі вибрати номери відповідних ребер.\par
%answer:4 5 6
\end {large}}
%end
{\smallskip \begin {small} Затверджено на засіданні кафедри теоретичної кібернетики, протокол \No 12 від   19.05.2020 р.\par\smallskip\smallskip
{\parbox {6.5 cm}{Зав. кафедри \dotfill Ю.В.~Крак}}\hspace {0.7cm}{\hbox to 6.5 cm{ Екзаменатор \dotfill Т.О.~Карнаух}} \end{small}}
%begin
{{\begin {small}\par
{ \center  Київський національний університет імені Тараса Шевченка \par}
{\parbox {5 cm}{Освітня програма \hfill}{\parbox {7 cm}{Прикладна математика, Системний аналіз \hfill}}\parbox {6 cm}{\hfill Семестр 2}}\par
{\parbox {5 cm} {Навчальний предмет \hfill}\parbox {7 cm}{Програмування \hfill}\parbox {6cm}{\hfill}\par}\vspace{-7pt} \end{small}}{\center Екзаменаційний білет \No \textbf 
{ 212 }\par}}
{
\begin {large}
1. Що буде виведено за виконання?
code
try:
    res = not 9>5 or 8.0!=5.0 or False<8
    if isinstance(res, bool):
        print(f'bool;{res}')
    elif isinstance(res, int):
        print(f'int;{res}')
    elif isinstance(res, float):
        print(f'float;{res:.2f}')
    else:
        print('???')
except:
    print('ERR')
code

2. Що буде виведено за виконання?
code
a = 3
b = 9
c = 8

def f(a):
    global c
    a += 5
    b = 1
    c = 2
    return a + b + c

a = 0
b = 6
c = 2

try:
    print(f(b), a, b, c, sep=';')
except:
    print('ERR', a, b, c, sep=';')
code

3. Що буде виведено за виконання?
code
try:
    k = 14
    while k>=5:
        k += -3
    print(k, end=';')
except:
    print('ERR')
code

4. Що буде виведено за виконання?
code
try:
    for f in range(2, -13, -2):
        if f<5:
            continue
        print(f, end=';')
    else:
        print(f,end=';')
    print(f)
except:
    print('ERR')
code

5. Що буде виведено за виконання?\par (Дійсні значення, якщо наявні, в цьому завданні будуть виводитись у форматі з фіксованою крапкою та, можливо, з доволі великою кількістю десяткових розрядів. У відповіді для дійсних значень у ЦЬОМУ завданні достатньо вказати один десятковий розряд після крапки, відкинувши решту розрядів без округлення. Наприклад, замість 13.66666666666667 достатньо вказати 13.6, замість 25.23 достатньо вказати 25.2)
code
try:
    print(7, end=';')
    print(int('c6'), end=';')
    print(9, end=';')
except TypeError:
    print(7, end=';')
except ValueError:
    print(5, end=';')
except:
    print('ERR', end=';')
else:
    print(1, end=';')
finally:
    print(3)
code

6. Що буде виведено за виконання?
code
def f(n):
    print(n%10, end='')
    if n>9:
        f(n//100)
 
f(987654)
code

7. Що буде виведено за виконання?
code
class A:
    b = 1

    def __f1(self): print(2, end=';')
    def _f2(self): print(3, end=';')

    def main(self):
        try: print(self.b, end=';')
        except: print('ERR', end=';')

        try: self.__f1()
        except: print('ERR', end=';')

        try: self._f2()
        except: print('ERR', end=';')


class B(A):
    b = 4

    def __f1(self): print(5, end=';')
    def _f2(self): print(6, end=';')

try: B().main()
except: print('ERR', end=';')
code

8. Що буде виведено за виконання?
code
def f(arg, n):
    n = 1
    tmp = arg
    tmp[-6:] = [1,0]
    return len(tmp)>3

try:
    f = [10,11,12,13,14]
    n = 6
    f(f, n)
    print(n, end=';')
    for el in f:
        print(el, end=';')
except:
    print('ERR')
code

9. Що буде виведено за виконання?
code
class A:
    b = 1
    c = 2

    def __init__(self):
        b = 3

class B(A):
    c = 4

    def __init__(self):
        super().__init__()
        self.a = 7

obj = B()
obj.b = obj.c + 10
B.b = 13

try: print(obj.b, end= ';')
except: print('ERR', end=';')

try: print(A.b, end= ';')
except: print('ERR', end=';')

try: print(B.b, end= ';')
except: print('ERR', end=';')
code

10. 
Для графа на рисунку з вершини 3 запущено алгоритм пошуку в глибину, що будує кістякове дерево. Списки суміжності вершин переглядаються алгоритмом за зростанням міток вершин.\par
\begin{center}
	\adjustimage{max size={0.9\linewidth}{0.9\paperheight}}
	{../10/Traversals/209.png}
\end{center}
\par Які з наведених ребер входять до складу отриманого кістякового дерева?\par
1) (4, 5)\par
2) (0, 4)\par
3) (1, 6)\par
4) (1, 3)\par
5) (1, 4)\par
6) (2, 3)\par
{\vskip 15pt} У формі вибрати номери відповідних ребер.\par
%answer:1 2 3 4 5
\end {large}}
%end
{\smallskip \begin {small} Затверджено на засіданні кафедри теоретичної кібернетики, протокол \No 12 від   19.05.2020 р.\par\smallskip\smallskip
{\parbox {6.5 cm}{Зав. кафедри \dotfill Ю.В.~Крак}}\hspace {0.7cm}{\hbox to 6.5 cm{ Екзаменатор \dotfill Т.О.~Карнаух}} \end{small}}
%begin
{{\begin {small}\par
{ \center  Київський національний університет імені Тараса Шевченка \par}
{\parbox {5 cm}{Освітня програма \hfill}{\parbox {7 cm}{Прикладна математика, Системний аналіз \hfill}}\parbox {6 cm}{\hfill Семестр 2}}\par
{\parbox {5 cm} {Навчальний предмет \hfill}\parbox {7 cm}{Програмування \hfill}\parbox {6cm}{\hfill}\par}\vspace{-7pt} \end{small}}{\center Екзаменаційний білет \No \textbf 
{ 213 }\par}}
{
\begin {large}
1. Що буде виведено за виконання?
code
try:
    res = 7.0!=4 or not 4.0<6 and 7>False
    if isinstance(res, bool):
        print(f'bool;{res}')
    elif isinstance(res, int):
        print(f'int;{res}')
    elif isinstance(res, float):
        print(f'float;{res:.2f}')
    else:
        print('???')
except:
    print('ERR')
code

2. Що буде виведено за виконання?
code
a = 4
b = 1
c = 5

def h(a):
    a *= 3
    b = 4
    c = 4
    return a + b + c

a = 7
b = 0
c = 2

try:
    print(h(a), a, b, c, sep=';')
except:
    print('ERR', a, b, c, sep=';')
code

3. Що буде виведено за виконання?
code
try:
    j = 2
    while j<8:
        j += 1
    print(j, end=';')
except:
    print('ERR')
code

4. Що буде виведено за виконання?
code
try:
    for e in range(3, -10, 3):
        if e<3:
            continue
        print(e, end=';');e = -4
        if e>8:
            break
    else:
        print('end',end=';')
    print(e)
except:
    print('ERR')
code

5. Що буде виведено за виконання?\par (Дійсні значення, якщо наявні, в цьому завданні будуть виводитись у форматі з фіксованою крапкою та, можливо, з доволі великою кількістю десяткових розрядів. У відповіді для дійсних значень у ЦЬОМУ завданні достатньо вказати один десятковий розряд після крапки, відкинувши решту розрядів без округлення. Наприклад, замість 13.66666666666667 достатньо вказати 13.6, замість 25.23 достатньо вказати 25.2)
code
try:
    print(6, end=';')
    print(int('9'), end=';')
    print(2, end=';')
except ZeroDivisionError:
    print(0, end=';')
except Exception:
    print(8, end=';')
except:
    print('ERR', end=';')
else:
    print(4, end=';')
finally:
    print(8)
code

6. Що буде виведено за виконання?
code
def f(n):
    if n>9:
        return f(n//100)+1
    else:
        return 1

print(f(123456))
code

7. Що буде виведено за виконання?
code
class A:
    b = 1

    def _h1(self): print(2, end=';')
    def h2(self): print(3, end=';')

    def main(self):
        try: print(self.b, end=';')
        except: print('ERR', end=';')

        try: self._h1()
        except: print('ERR', end=';')

        try: self.h2()
        except: print('ERR', end=';')


class B(A):
    b = 4

    def _h1(self): print(5, end=';')
    def h2(self): print(6, end=';')

try: B().main()
except: print('ERR', end=';')
code

8. Що буде виведено за виконання?
code
def f(arg, n):
    n = 2
    tmp = arg
    tmp[:3] = ()
    return len(tmp)>3

try:
    h = ['0','1','2','3','4']
    n = 1
    f(h, n)
    print(n, end=';')
    for el in h:
        print(el, end=';')
except:
    print('ERR')
code

9. Що буде виведено за виконання?
code
class A:
    a = 1
    b = 2

    def __init__(self):
        b = 3

class B(A):
    b = 4

    def __init__(self):
        super().__init__()
        self.c = 7

obj = B()
obj.b = obj.a + 10
B.b = 13

try: print(obj.b, end= ';')
except: print('ERR', end=';')

try: print(A.b, end= ';')
except: print('ERR', end=';')

try: print(B.b, end= ';')
except: print('ERR', end=';')
code

10. 
Для графа на рисунку з вершини 4 запущено алгоритм пошуку в глибину, що будує кістякове дерево. Списки суміжності вершин переглядаються алгоритмом за зростанням міток вершин.\par
\begin{center}
	\adjustimage{max size={0.9\linewidth}{0.9\paperheight}}
	{../10/Traversals/103.png}
\end{center}
\par Які з наведених ребер входять до складу отриманого кістякового дерева?\par
1) (2, 5)\par
2) (2, 6)\par
3) (1, 2)\par
4) (0, 4)\par
5) (0, 3)\par
6) (0, 6)\par
{\vskip 15pt} У формі вибрати номери відповідних ребер.\par
%answer:1 2 3 4 5
\end {large}}
%end
{\smallskip \begin {small} Затверджено на засіданні кафедри теоретичної кібернетики, протокол \No 12 від   19.05.2020 р.\par\smallskip\smallskip
{\parbox {6.5 cm}{Зав. кафедри \dotfill Ю.В.~Крак}}\hspace {0.7cm}{\hbox to 6.5 cm{ Екзаменатор \dotfill Т.О.~Карнаух}} \end{small}}
%begin
{{\begin {small}\par
{ \center  Київський національний університет імені Тараса Шевченка \par}
{\parbox {5 cm}{Освітня програма \hfill}{\parbox {7 cm}{Прикладна математика, Системний аналіз \hfill}}\parbox {6 cm}{\hfill Семестр 2}}\par
{\parbox {5 cm} {Навчальний предмет \hfill}\parbox {7 cm}{Програмування \hfill}\parbox {6cm}{\hfill}\par}\vspace{-7pt} \end{small}}{\center Екзаменаційний білет \No \textbf 
{ 214 }\par}}
{
\begin {large}
1. Що буде виведено за виконання?
code
try:
    res = 3.0==5 and not 2>=4 or 2.0<=True
    if isinstance(res, bool):
        print(f'bool;{res}')
    elif isinstance(res, int):
        print(f'int;{res}')
    elif isinstance(res, float):
        print(f'float;{res:.2f}')
    else:
        print('???')
except:
    print('ERR')
code

2. Що буде виведено за виконання?
code
a = 7
b = 6
c = 9

def g(b):
    global c
    a = 3
    b -= 2
    c = 5
    return a + b + c

a = 3
b = 1
c = 4

try:
    print(g(b), a, b, c, sep=';')
except:
    print('ERR', a, b, c, sep=';')
code

3. Що буде виведено за виконання?
code
try:
    t = 12
    while t>=4:
        t += -3
    print(t, end=';')
except:
    print('ERR')
code

4. Що буде виведено за виконання?
code
try:
    for d in range(6, -3, -1):
        if d<=8:
            break
        print(d, end=';');d = -10
    else:
        print(d,end=';')
    print(d)
except:
    print('ERR')
code

5. Що буде виведено за виконання?\par (Дійсні значення, якщо наявні, в цьому завданні будуть виводитись у форматі з фіксованою крапкою та, можливо, з доволі великою кількістю десяткових розрядів. У відповіді для дійсних значень у ЦЬОМУ завданні достатньо вказати один десятковий розряд після крапки, відкинувши решту розрядів без округлення. Наприклад, замість 13.66666666666667 достатньо вказати 13.6, замість 25.23 достатньо вказати 25.2)
code
try:
    print(2, end=';')
    print(7j>=8j, end=';')
    print(1, end=';')
except ValueError:
    print(9, end=';')
except TypeError:
    print(0, end=';')
except:
    print('ERR', end=';')
else:
    print(4, end=';')
finally:
    print(6)
code

6. Що буде виведено за виконання?
code
def f(n):
    if n>9: f(n//10)
    else:
        print(n%10, end='')
 
f(123456)
code

7. Що буде виведено за виконання?
code
class A:
    c = 1

    def f1(self): print(2, end=';')
    def _f2(self): print(3, end=';')

    def main(self):
        try: print(self.c, end=';')
        except: print('ERR', end=';')

        try: self.f1()
        except: print('ERR', end=';')

        try: self._f2()
        except: print('ERR', end=';')


class B(A):
    c = 4

    def f1(self): print(5, end=';')
    def _f2(self): print(6, end=';')

try: B().main()
except: print('ERR', end=';')
code

8. Що буде виведено за виконання?
code
def f(arg, n):
    arg[18] = 7
    n = 5
    arg = tuple(arg.values())

try:
    n = 2
    d = {70:7, 13:2, 34:7, 13:0}
    f(d, n)
    print(n, end=';')
    for x in d.values():
        print(x, sep=';', end=';')
except:
    print('ERR')
code

9. Що буде виведено за виконання?
code
class A:
    a = 1
    c = 2

    def __init__(self):
        a = 3

class B(A):
    a = 4

    def __init__(self):
        super().__init__()
        self.b = 7

obj = B()
obj.c = obj.a + 10
B.c = 13

try: print(obj.c, end= ';')
except: print('ERR', end=';')

try: print(A.c, end= ';')
except: print('ERR', end=';')

try: print(B.c, end= ';')
except: print('ERR', end=';')
code

10. 
Для графа на рисунку з вершини 1 запущено алгоритм пошуку в глибину, що будує кістякове дерево. Списки суміжності вершин переглядаються алгоритмом за зростанням міток вершин.\par
\begin{center}
	\adjustimage{max size={0.9\linewidth}{0.9\paperheight}}
	{../10/Traversals/12.png}
\end{center}
\par Які з наведених ребер входять до складу отриманого кістякового дерева?\par
1) (0, 4)\par
2) (1, 4)\par
3) (3, 6)\par
4) (0, 3)\par
5) (0, 6)\par
6) (0, 5)\par
{\vskip 15pt} У формі вибрати номери відповідних ребер.\par
%answer:1 3 4
\end {large}}
%end
{\smallskip \begin {small} Затверджено на засіданні кафедри теоретичної кібернетики, протокол \No 12 від   19.05.2020 р.\par\smallskip\smallskip
{\parbox {6.5 cm}{Зав. кафедри \dotfill Ю.В.~Крак}}\hspace {0.7cm}{\hbox to 6.5 cm{ Екзаменатор \dotfill Т.О.~Карнаух}} \end{small}}
%begin
{{\begin {small}\par
{ \center  Київський національний університет імені Тараса Шевченка \par}
{\parbox {5 cm}{Освітня програма \hfill}{\parbox {7 cm}{Прикладна математика, Системний аналіз \hfill}}\parbox {6 cm}{\hfill Семестр 2}}\par
{\parbox {5 cm} {Навчальний предмет \hfill}\parbox {7 cm}{Програмування \hfill}\parbox {6cm}{\hfill}\par}\vspace{-7pt} \end{small}}{\center Екзаменаційний білет \No \textbf 
{ 215 }\par}}
{
\begin {large}
1. Що буде виведено за виконання?
code
try:
    res = 8/7//4*6
    if isinstance(res, bool):
        print(f'bool;{res}')
    elif isinstance(res, int):
        print(f'int;{res}')
    elif isinstance(res, float):
        print(f'float;{res:.2f}')
    else:
        print('???')
except:
    print('ERR')
code

2. Що буде виведено за виконання?
code
a = 6
b = 2
c = 3

def h(a):
    global c
    a = 3
    b = 1
    c = 4
    return a + b + c

a = 1
b = 3
c = 1

try:
    print(h(a), a, b, c, sep=';')
except:
    print('ERR', a, b, c, sep=';')
code

3. Що буде виведено за виконання?
code
try:
    m = 9
    while m<=10:
        m += 2
    print(m, end=';')
except:
    print('ERR')
code

4. Що буде виведено за виконання?
code
try:
    for c in range(-12, 14, 2):
        if c>=7:
            continue
        print(c, end=';')
    else:
        print(c,end=';')
    print(c)
except:
    print('ERR')
code

5. Що буде виведено за виконання?\par (Дійсні значення, якщо наявні, в цьому завданні будуть виводитись у форматі з фіксованою крапкою та, можливо, з доволі великою кількістю десяткових розрядів. У відповіді для дійсних значень у ЦЬОМУ завданні достатньо вказати один десятковий розряд після крапки, відкинувши решту розрядів без округлення. Наприклад, замість 13.66666666666667 достатньо вказати 13.6, замість 25.23 достатньо вказати 25.2)
code
try:
    print(3, end=';')
    print(5//0.0, end=';')
    print(9, end=';')
except KeyboardInterrupt:
    print(0, end=';')
except ZeroDivisionError:
    print(1, end=';')
except:
    print('ERR', end=';')
else:
    print(8, end=';')
finally:
    print(3)
code

6. Що буде виведено за виконання?
code
def f(n):
    print(n%10,end='')
    if n>9:
        f(n//100)
 
f(123456)
code

7. Що буде виведено за виконання?
code
class A:
    d = 1

    def _h1(self): print(2, end=';')
    def _h2(self): print(3, end=';')

    def main(self):
        try: print(self.d, end=';')
        except: print('ERR', end=';')

        try: self._h1()
        except: print('ERR', end=';')

        try: self._h2()
        except: print('ERR', end=';')


class B(A):
    d = 4

    def _h1(self): print(5, end=';')
    def _h2(self): print(6, end=';')

try: B().main()
except: print('ERR', end=';')
code

8. Що буде виведено за виконання?
code
def f(arg, n):
    n = 9
    tmp = arg
    del tmp[-4:8]
    return len(tmp)>3

try:
    b = ['0','1','2','3','4']
    n = 5
    f(b, n)
    print(n, end=';')
    for el in b:
        print(el, end=';')
except:
    print('ERR')
code

9. Що буде виведено за виконання?
code
class A:
    b = 1
    a = 2

    def __init__(self):
        b = 3

class B(A):
    a = 4

    def __init__(self):
        super().__init__()
        self.c = 7

obj = B()
obj.b = obj.a + 10
B.b = 13

try: print(obj.b, end= ';')
except: print('ERR', end=';')

try: print(A.b, end= ';')
except: print('ERR', end=';')

try: print(B.b, end= ';')
except: print('ERR', end=';')
code

10. 
Для графа на рисунку з вершини 5 запущено алгоритм пошуку в глибину, що будує кістякове дерево. Списки суміжності вершин переглядаються алгоритмом за зростанням міток вершин.\par
\begin{center}
	\adjustimage{max size={0.9\linewidth}{0.9\paperheight}}
	{../10/Traversals/179.png}
\end{center}
\par Які з наведених ребер входять до складу отриманого кістякового дерева?\par
1) (2, 3)\par
2) (3, 5)\par
3) (2, 5)\par
4) (1, 6)\par
5) (0, 4)\par
6) (0, 1)\par
{\vskip 15pt} У формі вибрати номери відповідних ребер.\par
%answer:1 3 4 5 6
\end {large}}
%end
{\smallskip \begin {small} Затверджено на засіданні кафедри теоретичної кібернетики, протокол \No 12 від   19.05.2020 р.\par\smallskip\smallskip
{\parbox {6.5 cm}{Зав. кафедри \dotfill Ю.В.~Крак}}\hspace {0.7cm}{\hbox to 6.5 cm{ Екзаменатор \dotfill Т.О.~Карнаух}} \end{small}}
%begin
{{\begin {small}\par
{ \center  Київський національний університет імені Тараса Шевченка \par}
{\parbox {5 cm}{Освітня програма \hfill}{\parbox {7 cm}{Прикладна математика, Системний аналіз \hfill}}\parbox {6 cm}{\hfill Семестр 2}}\par
{\parbox {5 cm} {Навчальний предмет \hfill}\parbox {7 cm}{Програмування \hfill}\parbox {6cm}{\hfill}\par}\vspace{-7pt} \end{small}}{\center Екзаменаційний білет \No \textbf 
{ 216 }\par}}
{
\begin {large}
1. Що буде виведено за виконання?
code
try:
    res = 1!=1.0
    if isinstance(res, bool):
        print(f'bool;{res}')
    elif isinstance(res, int):
        print(f'int;{res}')
    elif isinstance(res, float):
        print(f'float;{res:.2f}')
    else:
        print('???')
except:
    print('ERR')
code

2. Що буде виведено за виконання?
code
a = 8
b = 5
c = 7

def f(b):
    global c
    a = 1
    b -= 3
    c = 1
    return a + b + c

a = 1
b = 6
c = 0

try:
    print(f(b), a, b, c, sep=';')
except:
    print('ERR', a, b, c, sep=';')
code

3. Що буде виведено за виконання?
code
try:
    n = 0
    while n<5:
        n += 2
    print(n, end=';')
except:
    print('ERR')
code

4. Що буде виведено за виконання?
code
try:
    for b in range(-10, 10, -1):
        if b>4:
            continue
        print(b, end=';')
    else:
        print(b,end=';')
    print(b)
except:
    print('ERR')
code

5. Що буде виведено за виконання?\par (Дійсні значення, якщо наявні, в цьому завданні будуть виводитись у форматі з фіксованою крапкою та, можливо, з доволі великою кількістю десяткових розрядів. У відповіді для дійсних значень у ЦЬОМУ завданні достатньо вказати один десятковий розряд після крапки, відкинувши решту розрядів без округлення. Наприклад, замість 13.66666666666667 достатньо вказати 13.6, замість 25.23 достатньо вказати 25.2)
code
try:
    print(2, end=';')
    print(6j==7j, end=';')
    print(5, end=';')
except Exception:
    print(7, end=';')
except ZeroDivisionError:
    print(4, end=';')
except:
    print('ERR', end=';')
else:
    print(2, end=';')
finally:
    print(4)
code

6. Що буде виведено за виконання?
code
def f(n):
    if n>9:
        f(n//10)
    else:
        print(n%10, end='')

f(543216)
code

7. Що буде виведено за виконання?
code
class A:
    a = 1

    def g1(self): print(2, end=';')
    def _g2(self): print(3, end=';')

    def main(self):
        try: print(self.a, end=';')
        except: print('ERR', end=';')

        try: self.g1()
        except: print('ERR', end=';')

        try: self._g2()
        except: print('ERR', end=';')


class B(A):
    a = 4

    def g1(self): print(5, end=';')
    def _g2(self): print(6, end=';')

try: B().main()
except: print('ERR', end=';')
code

8. Що буде виведено за виконання?
code
def f(arg, n):
    arg[71] = 7
    n += -3

try:
    n = 3
    d = {19:5, 87:7, 87:3}
    f(d, n)
    print(n, end=';')
    for x in d.items():
        print(*x, sep=';', end=';')
except:
    print('ERR')
code

9. Що буде виведено за виконання?
code
class A:
    c = 1
    a = 2

    def __init__(self):
        a = 3

class B(A):
    c = 4

    def __init__(self):
        super().__init__()
        self.b = 7

obj = B()
obj.b = obj.c + 10
B.b = 13

try: print(obj.b, end= ';')
except: print('ERR', end=';')

try: print(A.b, end= ';')
except: print('ERR', end=';')

try: print(B.b, end= ';')
except: print('ERR', end=';')
code

10. 
Для графа на рисунку з вершини 3 запущено алгоритм пошуку в глибину, що будує кістякове дерево. Списки суміжності вершин переглядаються алгоритмом за зростанням міток вершин.\par
\begin{center}
	\adjustimage{max size={0.9\linewidth}{0.9\paperheight}}
	{../10/Traversals/296.png}
\end{center}
\par Які з наведених ребер входять до складу отриманого кістякового дерева?\par
1) (1, 4)\par
2) (4, 6)\par
3) (2, 5)\par
4) (2, 4)\par
5) (0, 1)\par
6) (1, 3)\par
{\vskip 15pt} У формі вибрати номери відповідних ребер.\par
%answer:1 2 3 4 5
\end {large}}
%end
{\smallskip \begin {small} Затверджено на засіданні кафедри теоретичної кібернетики, протокол \No 12 від   19.05.2020 р.\par\smallskip\smallskip
{\parbox {6.5 cm}{Зав. кафедри \dotfill Ю.В.~Крак}}\hspace {0.7cm}{\hbox to 6.5 cm{ Екзаменатор \dotfill Т.О.~Карнаух}} \end{small}}
%begin
{{\begin {small}\par
{ \center  Київський національний університет імені Тараса Шевченка \par}
{\parbox {5 cm}{Освітня програма \hfill}{\parbox {7 cm}{Прикладна математика, Системний аналіз \hfill}}\parbox {6 cm}{\hfill Семестр 2}}\par
{\parbox {5 cm} {Навчальний предмет \hfill}\parbox {7 cm}{Програмування \hfill}\parbox {6cm}{\hfill}\par}\vspace{-7pt} \end{small}}{\center Екзаменаційний білет \No \textbf 
{ 217 }\par}}
{
\begin {large}
1. Що буде виведено за виконання?
code
def f(n):
    print('f', end='')
    return n!=-1

try:
    print(f(-4) and f(1))
except:
    print('ERR')
code

2. Що буде виведено за виконання?
code
a = 8
b = 5
c = 4

def g(b):
    a = 1
    b = 5
    c = 3
    return a + b + c

a = 6
b = 7
c = 0

try:
    print(g(b), a, b, c, sep=';')
except:
    print('ERR', a, b, c, sep=';')
code

3. Що буде виведено за виконання?
code
try:
    t = 8
    while t<=9:
        t += 1
    print(t, end=';')
except:
    print('ERR')
code

4. Що буде виведено за виконання?
code
try:
    for a in range(8, 0, 3):
        if a<=3:
            continue
        print(a, end=';')
    else:
        print(a,end=';')
    print(a)
except:
    print('ERR')
code

5. Що буде виведено за виконання?\par (Дійсні значення, якщо наявні, в цьому завданні будуть виводитись у форматі з фіксованою крапкою та, можливо, з доволі великою кількістю десяткових розрядів. У відповіді для дійсних значень у ЦЬОМУ завданні достатньо вказати один десятковий розряд після крапки, відкинувши решту розрядів без округлення. Наприклад, замість 13.66666666666667 достатньо вказати 13.6, замість 25.23 достатньо вказати 25.2)
code
try:
    print(0, end=';')
    print(7/0, end=';')
    print(9, end=';')
except ValueError:
    print(3, end=';')
except KeyboardInterrupt:
    print(5, end=';')
except:
    print('ERR', end=';')
else:
    print(1, end=';')
finally:
    print(8)
code

6. Що буде виведено за виконання?
code
def f(s):
    s1 = s
    for i in range(len(s)):
        s1[i] += 1
    return s1

s = [1, 2, 3]
f(s)
print(s)
code

7. Що буде виведено за виконання?
code
class A:
    a = 1

    def _g1(self): print(2, end=';')
    def g2(self): print(3, end=';')

    def main(self):
        try: print(self.a, end=';')
        except: print('ERR', end=';')

        try: self._g1()
        except: print('ERR', end=';')

        try: self.g2()
        except: print('ERR', end=';')


class B(A):
    a = 4

    def _g1(self): print(5, end=';')
    def g2(self): print(6, end=';')

try: B().main()
except: print('ERR', end=';')
code

8. Що буде виведено за виконання?
code
def f(arg, n):
    arg[77] = 2
    n *= -2

try:
    n = 7
    s = {77:2, 17:3, 12:4, 77:5}
    f(s, n)
    print(n, end=';')
    for x in s :
        print(x, sep=';', end=';')
except:
    print('ERR')
code

9. Що буде виведено за виконання?
code
class A:
    a = 1
    c = 2

    def __init__(self):
        a = 3

class B(A):
    a = 4

    def __init__(self):
        super().__init__()
        self.b = 7

obj = B()
obj.c = obj.b + 10
B.c = 13

try: print(obj.c, end= ';')
except: print('ERR', end=';')

try: print(A.c, end= ';')
except: print('ERR', end=';')

try: print(B.c, end= ';')
except: print('ERR', end=';')
code

10. 
Для графа на рисунку з вершини 2 запущено алгоритм пошуку в глибину, що будує кістякове дерево. Списки суміжності вершин переглядаються алгоритмом за зростанням міток вершин.\par
\begin{center}
	\adjustimage{max size={0.9\linewidth}{0.9\paperheight}}
	{../10/Traversals/141.png}
\end{center}
\par Які з наведених ребер входять до складу отриманого кістякового дерева?\par
1) (4, 6)\par
2) (0, 4)\par
3) (0, 5)\par
4) (0, 1)\par
5) (4, 5)\par
6) (2, 4)\par
{\vskip 15pt} У формі вибрати номери відповідних ребер.\par
%answer:1 4 5
\end {large}}
%end
{\smallskip \begin {small} Затверджено на засіданні кафедри теоретичної кібернетики, протокол \No 12 від   19.05.2020 р.\par\smallskip\smallskip
{\parbox {6.5 cm}{Зав. кафедри \dotfill Ю.В.~Крак}}\hspace {0.7cm}{\hbox to 6.5 cm{ Екзаменатор \dotfill Т.О.~Карнаух}} \end{small}}
%begin
{{\begin {small}\par
{ \center  Київський національний університет імені Тараса Шевченка \par}
{\parbox {5 cm}{Освітня програма \hfill}{\parbox {7 cm}{Прикладна математика, Системний аналіз \hfill}}\parbox {6 cm}{\hfill Семестр 2}}\par
{\parbox {5 cm} {Навчальний предмет \hfill}\parbox {7 cm}{Програмування \hfill}\parbox {6cm}{\hfill}\par}\vspace{-7pt} \end{small}}{\center Екзаменаційний білет \No \textbf 
{ 218 }\par}}
{
\begin {large}
1. Що буде виведено за виконання?
code
try:
    res = 3//4/9*9
    if isinstance(res, bool):
        print(f'bool;{res}')
    elif isinstance(res, int):
        print(f'int;{res}')
    elif isinstance(res, float):
        print(f'float;{res:.2f}')
    else:
        print('???')
except:
    print('ERR')
code

2. Що буде виведено за виконання?
code
a = 9
b = 2
c = 7

def f(a):
    a = 2
    b = 3
    c = 1
    return a + b + c

a = 6
b = 0
c = 8

try:
    print(f(b), a, b, c, sep=';')
except:
    print('ERR', a, b, c, sep=';')
code

3. Що буде виведено за виконання?
code
try:
    j = 6
    while j>3:
        j += -2
    print(j, end=';')
except:
    print('ERR')
code

4. Що буде виведено за виконання?
code
try:
    for b in range(-14, -4, 2):
        if b<7:
            break
        print(b, end=';');b = -6
    else:
        print('end',end=';')
    print(b)
except:
    print('ERR')
code

5. Що буде виведено за виконання?\par (Дійсні значення, якщо наявні, в цьому завданні будуть виводитись у форматі з фіксованою крапкою та, можливо, з доволі великою кількістю десяткових розрядів. У відповіді для дійсних значень у ЦЬОМУ завданні достатньо вказати один десятковий розряд після крапки, відкинувши решту розрядів без округлення. Наприклад, замість 13.66666666666667 достатньо вказати 13.6, замість 25.23 достатньо вказати 25.2)
code
try:
    print(6, end=';')
    print(int('9'), end=';')
    print(7, end=';')
except TypeError:
    print(6, end=';')
except Exception:
    print(0, end=';')
except:
    print('ERR', end=';')
else:
    print(1, end=';')
finally:
    print(8)
code

6. Що буде виведено за виконання?
code
def f(n):
    if n>9:
        f(n//100)
    print(n%10, end='')
 
f(123456)
code

7. Що буде виведено за виконання?
code
class A:
    d = 1

    def __h1(self): print(2, end=';')
    def _h2(self): print(3, end=';')

    def main(self):
        try: print(self.d, end=';')
        except: print('ERR', end=';')

        try: self.__h1()
        except: print('ERR', end=';')

        try: self._h2()
        except: print('ERR', end=';')


class B(A):
    d = 4

    def __h1(self): print(5, end=';')
    def _h2(self): print(6, end=';')

try: B().main()
except: print('ERR', end=';')
code

8. Що буде виведено за виконання?
code
def f(arg, n):
    arg[36] = 9
    n -= 3

try:
    n = 6
    t = {10:0, 66:7, 40:8, 62:7}
    f(t, n)
    print(n, end=';')
    for x in t.values():
        print(x, sep=';', end=';')
except:
    print('ERR')
code

9. Що буде виведено за виконання?
code
class A:
    b = 1
    c = 2

    def __init__(self):
        c = 3

class B(A):
    c = 4

    def __init__(self):
        super().__init__()
        self.a = 7

obj = B()
obj.a = obj.a + 10
B.a = 13

try: print(obj.a, end= ';')
except: print('ERR', end=';')

try: print(A.a, end= ';')
except: print('ERR', end=';')

try: print(B.a, end= ';')
except: print('ERR', end=';')
code

10. 
Для графа на рисунку з вершини 4 запущено алгоритм пошуку в глибину, що будує кістякове дерево. Списки суміжності вершин переглядаються алгоритмом за зростанням міток вершин.\par
\begin{center}
	\adjustimage{max size={0.9\linewidth}{0.9\paperheight}}
	{../10/Traversals/144.png}
\end{center}
\par Які з наведених ребер входять до складу отриманого кістякового дерева?\par
1) (3, 5)\par
2) (1, 6)\par
3) (1, 5)\par
4) (4, 6)\par
5) (0, 3)\par
6) (1, 3)\par
{\vskip 15pt} У формі вибрати номери відповідних ребер.\par
%answer:3 5 6
\end {large}}
%end
{\smallskip \begin {small} Затверджено на засіданні кафедри теоретичної кібернетики, протокол \No 12 від   19.05.2020 р.\par\smallskip\smallskip
{\parbox {6.5 cm}{Зав. кафедри \dotfill Ю.В.~Крак}}\hspace {0.7cm}{\hbox to 6.5 cm{ Екзаменатор \dotfill Т.О.~Карнаух}} \end{small}}
%begin
{{\begin {small}\par
{ \center  Київський національний університет імені Тараса Шевченка \par}
{\parbox {5 cm}{Освітня програма \hfill}{\parbox {7 cm}{Прикладна математика, Системний аналіз \hfill}}\parbox {6 cm}{\hfill Семестр 2}}\par
{\parbox {5 cm} {Навчальний предмет \hfill}\parbox {7 cm}{Програмування \hfill}\parbox {6cm}{\hfill}\par}\vspace{-7pt} \end{small}}{\center Екзаменаційний білет \No \textbf 
{ 219 }\par}}
{
\begin {large}
1. Що буде виведено за виконання?
code
try:
    res = False>"5.0j"
    if isinstance(res, bool):
        print(f'bool;{res}')
    elif isinstance(res, int):
        print(f'int;{res}')
    elif isinstance(res, float):
        print(f'float;{res:.2f}')
    else:
        print('???')
except:
    print('ERR')
code

2. Що буде виведено за виконання?
code
a = 8
b = 3
c = 4

def g(a):
    a += 2
    b = 4
    c = 5
    return a + b + c

a = 5
b = 9
c = 2

try:
    print(g(a), a, b, c, sep=';')
except:
    print('ERR', a, b, c, sep=';')
code

3. Що буде виведено за виконання?
code
try:
    k = 5
    while k<=13:
        k += 2
    print(k, end=';')
except:
    print('ERR')
code

4. Що буде виведено за виконання?
code
try:
    for c in range(-12, 9, -2):
        if c>8:
            continue
        print(c, end=';')
    else:
        print(c,end=';')
    print(c)
except:
    print('ERR')
code

5. Що буде виведено за виконання?\par (Дійсні значення, якщо наявні, в цьому завданні будуть виводитись у форматі з фіксованою крапкою та, можливо, з доволі великою кількістю десяткових розрядів. У відповіді для дійсних значень у ЦЬОМУ завданні достатньо вказати один десятковий розряд після крапки, відкинувши решту розрядів без округлення. Наприклад, замість 13.66666666666667 достатньо вказати 13.6, замість 25.23 достатньо вказати 25.2)
code
try:
    print(4, end=';')
    print(3//1, end=';')
    print(5, end=';')
except TypeError:
    print(2, end=';')
except ZeroDivisionError:
    print(0, end=';')
except:
    print('ERR', end=';')
else:
    print(7, end=';')
finally:
    print(6)
code

6. Що буде виведено за виконання?
code
def f(n):
    if n>9:
        f(n//100)
    else:
        print(n%10,end='')

f(123456)
code

7. Що буде виведено за виконання?
code
class A:
    c = 1

    def _f1(self): print(2, end=';')
    def __f2(self): print(3, end=';')

    def main(self):
        try: print(self.c, end=';')
        except: print('ERR', end=';')

        try: self._f1()
        except: print('ERR', end=';')

        try: self.__f2()
        except: print('ERR', end=';')


class B(A):
    c = 4

    def _f1(self): print(5, end=';')
    def __f2(self): print(6, end=';')

try: B().main()
except: print('ERR', end=';')
code

8. Що буде виведено за виконання?
code
def f(arg, n):
    arg[48] = 2
    n = 8
    arg = tuple(arg.values())

try:
    n = 0
    d = {47:8, 14:0, 83:6, 14:7}
    f(d, n)
    print(n, end=';')
    for x in d.items():
        print(x, sep=';', end=';')
except:
    print('ERR')
code

9. Що буде виведено за виконання?
code
class A:
    a = 1
    b = 2

    def __init__(self):
        a = 3

class B(A):
    a = 4

    def __init__(self):
        super().__init__()
        self.c = 7

obj = B()
obj.c = obj.b + 10
B.c = 13

try: print(obj.c, end= ';')
except: print('ERR', end=';')

try: print(A.c, end= ';')
except: print('ERR', end=';')

try: print(B.c, end= ';')
except: print('ERR', end=';')
code

10. 
Для графа на рисунку з вершини 2 запущено алгоритм пошуку в глибину, що будує кістякове дерево. Списки суміжності вершин переглядаються алгоритмом за зростанням міток вершин.\par
\begin{center}
	\adjustimage{max size={0.9\linewidth}{0.9\paperheight}}
	{../10/Traversals/5.png}
\end{center}
\par Які з наведених ребер входять до складу отриманого кістякового дерева?\par
1) (2, 6)\par
2) (3, 5)\par
3) (0, 4)\par
4) (2, 5)\par
5) (0, 3)\par
6) (0, 5)\par
{\vskip 15pt} У формі вибрати номери відповідних ребер.\par
%answer:2 3 5
\end {large}}
%end
{\smallskip \begin {small} Затверджено на засіданні кафедри теоретичної кібернетики, протокол \No 12 від   19.05.2020 р.\par\smallskip\smallskip
{\parbox {6.5 cm}{Зав. кафедри \dotfill Ю.В.~Крак}}\hspace {0.7cm}{\hbox to 6.5 cm{ Екзаменатор \dotfill Т.О.~Карнаух}} \end{small}}
%begin
{{\begin {small}\par
{ \center  Київський національний університет імені Тараса Шевченка \par}
{\parbox {5 cm}{Освітня програма \hfill}{\parbox {7 cm}{Прикладна математика, Системний аналіз \hfill}}\parbox {6 cm}{\hfill Семестр 2}}\par
{\parbox {5 cm} {Навчальний предмет \hfill}\parbox {7 cm}{Програмування \hfill}\parbox {6cm}{\hfill}\par}\vspace{-7pt} \end{small}}{\center Екзаменаційний білет \No \textbf 
{ 220 }\par}}
{
\begin {large}
1. Що буде виведено за виконання?
code
try:
    res = True<=3 or not 2.0!=8 and 9>=6.0
    if isinstance(res, bool):
        print(f'bool;{res}')
    elif isinstance(res, int):
        print(f'int;{res}')
    elif isinstance(res, float):
        print(f'float;{res:.2f}')
    else:
        print('???')
except:
    print('ERR')
code

2. Що буде виведено за виконання?
code
a = 2
b = 3
c = 5

def h(b):
    global c
    a = 5
    b = 2
    c = 4
    return a + b + c

a = 4
b = 1
c = 9

try:
    print(h(a), a, b, c, sep=';')
except:
    print('ERR', a, b, c, sep=';')
code

3. Що буде виведено за виконання?
code
try:
    n = 10
    while n>=9:
        n += -2
    print(n, end=';')
except:
    print('ERR')
code

4. Що буде виведено за виконання?
code
try:
    for d in range(13, 3, -3):
        if d>=6:
            break
        print(d, end=';')
        if d>=5:
            break
    else:
        print(d,end=';')
    print(d)
except:
    print('ERR')
code

5. Що буде виведено за виконання?\par (Дійсні значення, якщо наявні, в цьому завданні будуть виводитись у форматі з фіксованою крапкою та, можливо, з доволі великою кількістю десяткових розрядів. У відповіді для дійсних значень у ЦЬОМУ завданні достатньо вказати один десятковий розряд після крапки, відкинувши решту розрядів без округлення. Наприклад, замість 13.66666666666667 достатньо вказати 13.6, замість 25.23 достатньо вказати 25.2)
code
try:
    print(1, end=';')
    print(int('b2'), end=';')
    print(3, end=';')
except ValueError:
    print(8, end=';')
except KeyboardInterrupt:
    print(4, end=';')
except:
    print('ERR', end=';')
else:
    print(5, end=';')
finally:
    print(9)
code

6. Що буде виведено за виконання?
code
def f(s):
    for i in range(len(s)):
        s[i] += 1
        return
 
s = [1, 2, 3]
f(s)
print(s)
code

7. Що буде виведено за виконання?
code
class A:
    b = 1

    def _g1(self): print(2, end=';')
    def _g2(self): print(3, end=';')

    def main(self):
        try: print(self.b, end=';')
        except: print('ERR', end=';')

        try: self._g1()
        except: print('ERR', end=';')

        try: self._g2()
        except: print('ERR', end=';')


class B(A):
    b = 4

    def _g1(self): print(5, end=';')
    def _g2(self): print(6, end=';')

try: B().main()
except: print('ERR', end=';')
code

8. Що буде виведено за виконання?
code
def f(arg, n):
    arg[42] = 7
    n = 6
    arg = set(arg.keys())

try:
    n = 3
    s = {42:2, 73:5, 64:7, 59:7}
    f(s, n)
    print(n, end=';')
    for x in s.values():
        print(x, sep=';', end=';')
except:
    print('ERR')
code

9. Що буде виведено за виконання?
code
class A:
    c = 1
    b = 2

    def __init__(self):
        b = 3

class B(A):
    b = 4

    def __init__(self):
        super().__init__()
        self.a = 7

obj = B()
obj.c = obj.b + 10
B.c = 13

try: print(obj.c, end= ';')
except: print('ERR', end=';')

try: print(A.c, end= ';')
except: print('ERR', end=';')

try: print(B.c, end= ';')
except: print('ERR', end=';')
code

10. 
Для графа на рисунку з вершини 0 запущено алгоритм пошуку в глибину, що будує кістякове дерево. Списки суміжності вершин переглядаються алгоритмом за зростанням міток вершин.\par
\begin{center}
	\adjustimage{max size={0.9\linewidth}{0.9\paperheight}}
	{../10/Traversals/215.png}
\end{center}
\par Які з наведених ребер входять до складу отриманого кістякового дерева?\par
1) (0, 1)\par
2) (1, 6)\par
3) (2, 4)\par
4) (2, 3)\par
5) (0, 3)\par
6) (3, 5)\par
{\vskip 15pt} У формі вибрати номери відповідних ребер.\par
%answer:1 2 3 4 6
\end {large}}
%end
{\smallskip \begin {small} Затверджено на засіданні кафедри теоретичної кібернетики, протокол \No 12 від   19.05.2020 р.\par\smallskip\smallskip
{\parbox {6.5 cm}{Зав. кафедри \dotfill Ю.В.~Крак}}\hspace {0.7cm}{\hbox to 6.5 cm{ Екзаменатор \dotfill Т.О.~Карнаух}} \end{small}}
%begin
{{\begin {small}\par
{ \center  Київський національний університет імені Тараса Шевченка \par}
{\parbox {5 cm}{Освітня програма \hfill}{\parbox {7 cm}{Прикладна математика, Системний аналіз \hfill}}\parbox {6 cm}{\hfill Семестр 2}}\par
{\parbox {5 cm} {Навчальний предмет \hfill}\parbox {7 cm}{Програмування \hfill}\parbox {6cm}{\hfill}\par}\vspace{-7pt} \end{small}}{\center Екзаменаційний білет \No \textbf 
{ 221 }\par}}
{
\begin {large}
1. Що буде виведено за виконання?
code
try:
    res = 4/6*8*3
    if isinstance(res, bool):
        print(f'bool;{res}')
    elif isinstance(res, int):
        print(f'int;{res}')
    elif isinstance(res, float):
        print(f'float;{res:.2f}')
    else:
        print('???')
except:
    print('ERR')
code

2. Що буде виведено за виконання?
code
a = 7
b = 5
c = 1

def h(b):
    a *= 3
    b = 1
    c = 2
    return a + b + c

a = 9
b = 8
c = 3

try:
    print(h(b), a, b, c, sep=';')
except:
    print('ERR', a, b, c, sep=';')
code

3. Що буде виведено за виконання?
code
try:
    i = 15
    while i>2:
        i += -3
    print(i, end=';')
except:
    print('ERR')
code

4. Що буде виведено за виконання?
code
try:
    for f in range(-4, -2, 2):
        if f>4:
            continue
        print(f, end=';')
    else:
        print(f,end=';')
    print(f)
except:
    print('ERR')
code

5. Що буде виведено за виконання?\par (Дійсні значення, якщо наявні, в цьому завданні будуть виводитись у форматі з фіксованою крапкою та, можливо, з доволі великою кількістю десяткових розрядів. У відповіді для дійсних значень у ЦЬОМУ завданні достатньо вказати один десятковий розряд після крапки, відкинувши решту розрядів без округлення. Наприклад, замість 13.66666666666667 достатньо вказати 13.6, замість 25.23 достатньо вказати 25.2)
code
try:
    print(1, end=';')
    print(3/0.0, end=';')
    print(7, end=';')
except Exception:
    print(8, end=';')
except TypeError:
    print(9, end=';')
except:
    print('ERR', end=';')
else:
    print(5, end=';')
finally:
    print(4)
code

6. Що буде виведено за виконання?
code
def f(s):
    for el in s:
        el += 1
    return s
 
s = [1, 2, 3]
s = f(s)
print(s)
code

7. Що буде виведено за виконання?
code
class A:
    c = 1

    def _h1(self): print(2, end=';')
    def __h2(self): print(3, end=';')

    def main(self):
        try: print(self.c, end=';')
        except: print('ERR', end=';')

        try: self._h1()
        except: print('ERR', end=';')

        try: self.__h2()
        except: print('ERR', end=';')


class B(A):
    c = 4

    def _h1(self): print(5, end=';')
    def __h2(self): print(6, end=';')

try: B().main()
except: print('ERR', end=';')
code

8. Що буде виведено за виконання?
code
def f(arg, n):
    arg[47] = 6
    n = 9
    arg = list(arg.items())

try:
    n = 0
    t = {86:9, 47:0, 86:2}
    f(t, n)
    print(n, end=';')
    for x in t.keys():
        print(x, sep=';', end=';')
except:
    print('ERR')
code

9. Що буде виведено за виконання?
code
class A:
    c = 1
    a = 2

    def __init__(self):
        a = 3

class B(A):
    a = 4

    def __init__(self):
        super().__init__()
        self.b = 7

obj = B()
obj.a = obj.a + 10
B.a = 13

try: print(obj.a, end= ';')
except: print('ERR', end=';')

try: print(A.a, end= ';')
except: print('ERR', end=';')

try: print(B.a, end= ';')
except: print('ERR', end=';')
code

10. 
Для графа на рисунку з вершини 1 запущено алгоритм пошуку в глибину, що будує кістякове дерево. Списки суміжності вершин переглядаються алгоритмом за зростанням міток вершин.\par
\begin{center}
	\adjustimage{max size={0.9\linewidth}{0.9\paperheight}}
	{../10/Traversals/214.png}
\end{center}
\par Які з наведених ребер входять до складу отриманого кістякового дерева?\par
1) (1, 6)\par
2) (2, 4)\par
3) (0, 1)\par
4) (1, 3)\par
5) (0, 3)\par
6) (0, 6)\par
{\vskip 15pt} У формі вибрати номери відповідних ребер.\par
%answer:2 3 5
\end {large}}
%end
{\smallskip \begin {small} Затверджено на засіданні кафедри теоретичної кібернетики, протокол \No 12 від   19.05.2020 р.\par\smallskip\smallskip
{\parbox {6.5 cm}{Зав. кафедри \dotfill Ю.В.~Крак}}\hspace {0.7cm}{\hbox to 6.5 cm{ Екзаменатор \dotfill Т.О.~Карнаух}} \end{small}}
%begin
{{\begin {small}\par
{ \center  Київський національний університет імені Тараса Шевченка \par}
{\parbox {5 cm}{Освітня програма \hfill}{\parbox {7 cm}{Прикладна математика, Системний аналіз \hfill}}\parbox {6 cm}{\hfill Семестр 2}}\par
{\parbox {5 cm} {Навчальний предмет \hfill}\parbox {7 cm}{Програмування \hfill}\parbox {6cm}{\hfill}\par}\vspace{-7pt} \end{small}}{\center Екзаменаційний білет \No \textbf 
{ 222 }\par}}
{
\begin {large}
1. Що буде виведено за виконання?
code
try:
    res = "True">="2j"
    if isinstance(res, bool):
        print(f'bool;{res}')
    elif isinstance(res, int):
        print(f'int;{res}')
    elif isinstance(res, float):
        print(f'float;{res:.2f}')
    else:
        print('???')
except:
    print('ERR')
code

2. Що буде виведено за виконання?
code
a = 2
b = 1
c = 6

def g(b):
    a = 2
    b = 4
    c = 3
    return a + b + c

a = 3
b = 0
c = 7

try:
    print(g(b), a, b, c, sep=';')
except:
    print('ERR', a, b, c, sep=';')
code

3. Що буде виведено за виконання?
code
try:
    i = 3
    while i<6:
        i += 1
    print(i, end=';')
except:
    print('ERR')
code

4. Що буде виведено за виконання?
code
try:
    for e in range(-11, 7, -1):
        if e<5:
            continue
        print(e, end=';');e = 1
    else:
        print(e,end=';')
    print(e)
except:
    print('ERR')
code

5. Що буде виведено за виконання?\par (Дійсні значення, якщо наявні, в цьому завданні будуть виводитись у форматі з фіксованою крапкою та, можливо, з доволі великою кількістю десяткових розрядів. У відповіді для дійсних значень у ЦЬОМУ завданні достатньо вказати один десятковий розряд після крапки, відкинувши решту розрядів без округлення. Наприклад, замість 13.66666666666667 достатньо вказати 13.6, замість 25.23 достатньо вказати 25.2)
code
try:
    print(0, end=';')
    print(2j<=9j, end=';')
    print(6, end=';')
except ValueError:
    print(3, end=';')
except ZeroDivisionError:
    print(9, end=';')
except:
    print('ERR', end=';')
else:
    print(5, end=';')
finally:
    print(2)
code

6. Що буде виведено за виконання?
code
def f(s):
    for i in range(len(s)):
        s[i] += 1
    return
 
s = [1, 2, 3]
f(s)
print(s)
code

7. Що буде виведено за виконання?
code
class A:
    a = 1

    def _f1(self): print(2, end=';')
    def _f2(self): print(3, end=';')

    def main(self):
        try: print(self.a, end=';')
        except: print('ERR', end=';')

        try: self._f1()
        except: print('ERR', end=';')

        try: self._f2()
        except: print('ERR', end=';')


class B(A):
    a = 4

    def _f1(self): print(5, end=';')
    def _f2(self): print(6, end=';')

try: B().main()
except: print('ERR', end=';')
code

8. Що буде виведено за виконання?
code
def f(arg, n):
    n = 7
    tmp = arg
    tmp[-3:-6] = [1,2]
    return len(tmp)>3

try:
    d = [10,11,12,13,14]
    n = 8
    f(d, n)
    print(n, end=';')
    for el in d:
        print(el, end=';')
except:
    print('ERR')
code

9. Що буде виведено за виконання?
code
class A:
    b = 1
    a = 2

    def __init__(self):
        b = 3

class B(A):
    b = 4

    def __init__(self):
        super().__init__()
        self.c = 7

obj = B()
obj.b = obj.c + 10
B.b = 13

try: print(obj.b, end= ';')
except: print('ERR', end=';')

try: print(A.b, end= ';')
except: print('ERR', end=';')

try: print(B.b, end= ';')
except: print('ERR', end=';')
code

10. 
Для графа на рисунку з вершини 6 запущено алгоритм пошуку в глибину, що будує кістякове дерево. Списки суміжності вершин переглядаються алгоритмом за зростанням міток вершин.\par
\begin{center}
	\adjustimage{max size={0.9\linewidth}{0.9\paperheight}}
	{../10/Traversals/26.png}
\end{center}
\par Які з наведених ребер входять до складу отриманого кістякового дерева?\par
1) (3, 6)\par
2) (0, 6)\par
3) (4, 5)\par
4) (0, 1)\par
5) (5, 6)\par
6) (1, 6)\par
{\vskip 15pt} У формі вибрати номери відповідних ребер.\par
%answer:2 3 4
\end {large}}
%end
{\smallskip \begin {small} Затверджено на засіданні кафедри теоретичної кібернетики, протокол \No 12 від   19.05.2020 р.\par\smallskip\smallskip
{\parbox {6.5 cm}{Зав. кафедри \dotfill Ю.В.~Крак}}\hspace {0.7cm}{\hbox to 6.5 cm{ Екзаменатор \dotfill Т.О.~Карнаух}} \end{small}}
%begin
{{\begin {small}\par
{ \center  Київський національний університет імені Тараса Шевченка \par}
{\parbox {5 cm}{Освітня програма \hfill}{\parbox {7 cm}{Прикладна математика, Системний аналіз \hfill}}\parbox {6 cm}{\hfill Семестр 2}}\par
{\parbox {5 cm} {Навчальний предмет \hfill}\parbox {7 cm}{Програмування \hfill}\parbox {6cm}{\hfill}\par}\vspace{-7pt} \end{small}}{\center Екзаменаційний білет \No \textbf 
{ 223 }\par}}
{
\begin {large}
1. Що буде виведено за виконання?
code
def f(n):
    print('f', end='')
    return n

try:
    print(f(-6) and f(5))
except:
    print('ERR')
code

2. Що буде виведено за виконання?
code
a = 4
b = 5
c = 9

def g(a):
    global c
    a = 5
    b = 1
    c = 3
    return a + b + c

a = 8
b = 8
c = 6

try:
    print(g(a), a, b, c, sep=';')
except:
    print('ERR', a, b, c, sep=';')
code

3. Що буде виведено за виконання?
code
try:
    j = 1
    while j<13:
        j += 2
    print(j, end=';')
except:
    print('ERR')
code

4. Що буде виведено за виконання?
code
try:
    for f in range(9, 11, -3):
        if f<=5:
            break
        print(f, end=';')
    else:
        print(f,end=';')
    print(f)
except:
    print('ERR')
code

5. Що буде виведено за виконання?\par (Дійсні значення, якщо наявні, в цьому завданні будуть виводитись у форматі з фіксованою крапкою та, можливо, з доволі великою кількістю десяткових розрядів. У відповіді для дійсних значень у ЦЬОМУ завданні достатньо вказати один десятковий розряд після крапки, відкинувши решту розрядів без округлення. Наприклад, замість 13.66666666666667 достатньо вказати 13.6, замість 25.23 достатньо вказати 25.2)
code
try:
    print(4, end=';')
    print(int('0'), end=';')
    print(7, end=';')
except KeyboardInterrupt:
    print(1, end=';')
except TypeError:
    print(6, end=';')
except:
    print('ERR', end=';')
else:
    print(8, end=';')
finally:
    print(9)
code

6. Що буде виведено за виконання?
code
def f(n):
    print(n%10, end='')
    if n>9:
        f(n//10)
 
f(543216)
code

7. Що буде виведено за виконання?
code
class A:
    b = 1

    def f1(self): print(2, end=';')
    def _f2(self): print(3, end=';')

    def main(self):
        try: print(self.b, end=';')
        except: print('ERR', end=';')

        try: self.f1()
        except: print('ERR', end=';')

        try: self._f2()
        except: print('ERR', end=';')


class B(A):
    b = 4

    def f1(self): print(5, end=';')
    def _f2(self): print(6, end=';')

try: B().main()
except: print('ERR', end=';')
code

8. Що буде виведено за виконання?
code
def f(arg, n):
    arg[56] = 1
    n *= -2

try:
    n = 3
    s = {18:7, 25:5, 25:3}
    f(s, n)
    print(n, end=';')
    for x in s.keys():
        print(x, sep=';', end=';')
except:
    print('ERR')
code

9. Що буде виведено за виконання?
code
class A:
    b = 1
    c = 2

    def __init__(self):
        b = 3

class B(A):
    c = 4

    def __init__(self):
        super().__init__()
        self.a = 7

obj = B()
obj.b = obj.a + 10
B.b = 13

try: print(obj.b, end= ';')
except: print('ERR', end=';')

try: print(A.b, end= ';')
except: print('ERR', end=';')

try: print(B.b, end= ';')
except: print('ERR', end=';')
code

10. 
Для графа на рисунку з вершини 4 запущено алгоритм пошуку в глибину, що будує кістякове дерево. Списки суміжності вершин переглядаються алгоритмом за зростанням міток вершин.\par
\begin{center}
	\adjustimage{max size={0.9\linewidth}{0.9\paperheight}}
	{../10/Traversals/211.png}
\end{center}
\par Які з наведених ребер входять до складу отриманого кістякового дерева?\par
1) (4, 6)\par
2) (3, 6)\par
3) (1, 6)\par
4) (1, 2)\par
5) (0, 5)\par
6) (3, 4)\par
{\vskip 15pt} У формі вибрати номери відповідних ребер.\par
%answer:2 3 4 5
\end {large}}
%end
{\smallskip \begin {small} Затверджено на засіданні кафедри теоретичної кібернетики, протокол \No 12 від   19.05.2020 р.\par\smallskip\smallskip
{\parbox {6.5 cm}{Зав. кафедри \dotfill Ю.В.~Крак}}\hspace {0.7cm}{\hbox to 6.5 cm{ Екзаменатор \dotfill Т.О.~Карнаух}} \end{small}}
%begin
{{\begin {small}\par
{ \center  Київський національний університет імені Тараса Шевченка \par}
{\parbox {5 cm}{Освітня програма \hfill}{\parbox {7 cm}{Прикладна математика, Системний аналіз \hfill}}\parbox {6 cm}{\hfill Семестр 2}}\par
{\parbox {5 cm} {Навчальний предмет \hfill}\parbox {7 cm}{Програмування \hfill}\parbox {6cm}{\hfill}\par}\vspace{-7pt} \end{small}}{\center Екзаменаційний білет \No \textbf 
{ 224 }\par}}
{
\begin {large}
1. Що буде виведено за виконання?
code
try:
    res = "6.0j"<4.0
    if isinstance(res, bool):
        print(f'bool;{res}')
    elif isinstance(res, int):
        print(f'int;{res}')
    elif isinstance(res, float):
        print(f'float;{res:.2f}')
    else:
        print('???')
except:
    print('ERR')
code

2. Що буде виведено за виконання?
code
a = 3
b = 0
c = 7

def h(b):
    a = 1
    b = 4
    c = 5
    return a + b + c

a = 5
b = 9
c = 2

try:
    print(h(a), a, b, c, sep=';')
except:
    print('ERR', a, b, c, sep=';')
code

3. Що буде виведено за виконання?
code
try:
    i = 7
    while i<=12:
        i += 1
    print(i, end=';')
except:
    print('ERR')
code

4. Що буде виведено за виконання?
code
try:
    for d in range(11, 15, 3):
        if d>=8:
            continue
        print(d, end=';');d = 3
    else:
        print(d,end=';')
    print(d)
except:
    print('ERR')
code

5. Що буде виведено за виконання?\par (Дійсні значення, якщо наявні, в цьому завданні будуть виводитись у форматі з фіксованою крапкою та, можливо, з доволі великою кількістю десяткових розрядів. У відповіді для дійсних значень у ЦЬОМУ завданні достатньо вказати один десятковий розряд після крапки, відкинувши решту розрядів без округлення. Наприклад, замість 13.66666666666667 достатньо вказати 13.6, замість 25.23 достатньо вказати 25.2)
code
try:
    print(2, end=';')
    print(int('b1'), end=';')
    print(4, end=';')
except KeyboardInterrupt:
    print(3, end=';')
except ValueError:
    print(8, end=';')
except:
    print('ERR', end=';')
else:
    print(7, end=';')
finally:
    print(6)
code

6. Що буде виведено за виконання?
code
def f(n):
    if n>2:
        f(n//2)
    print(n%2, end='')
 
f(16)
code

7. Що буде виведено за виконання?
code
class A:
    d = 1

    def _g1(self): print(2, end=';')
    def g2(self): print(3, end=';')

    def main(self):
        try: print(self.d, end=';')
        except: print('ERR', end=';')

        try: self._g1()
        except: print('ERR', end=';')

        try: self.g2()
        except: print('ERR', end=';')


class B(A):
    d = 4

    def _g1(self): print(5, end=';')
    def g2(self): print(6, end=';')

try: B().main()
except: print('ERR', end=';')
code

8. Що буде виведено за виконання?
code
def f(arg, n):
    arg[33] = 3
    n -= 2

try:
    n = 5
    t = {82:5, 33:5, 33:7}
    f(t, n)
    print(n, end=';')
    for x in t.values():
        print(x, sep=';', end=';')
except:
    print('ERR')
code

9. Що буде виведено за виконання?
code
class A:
    a = 1
    b = 2

    def __init__(self):
        b = 3

class B(A):
    a = 4

    def __init__(self):
        super().__init__()
        self.c = 7

obj = B()
obj.c = obj.b + 10
B.c = 13

try: print(obj.c, end= ';')
except: print('ERR', end=';')

try: print(A.c, end= ';')
except: print('ERR', end=';')

try: print(B.c, end= ';')
except: print('ERR', end=';')
code

10. 
Для графа на рисунку з вершини 6 запущено алгоритм пошуку в глибину, що будує кістякове дерево. Списки суміжності вершин переглядаються алгоритмом за зростанням міток вершин.\par
\begin{center}
	\adjustimage{max size={0.9\linewidth}{0.9\paperheight}}
	{../10/Traversals/89.png}
\end{center}
\par Які з наведених ребер входять до складу отриманого кістякового дерева?\par
1) (0, 5)\par
2) (4, 6)\par
3) (3, 5)\par
4) (2, 4)\par
5) (0, 4)\par
6) (1, 2)\par
{\vskip 15pt} У формі вибрати номери відповідних ребер.\par
%answer:3 5
\end {large}}
%end
{\smallskip \begin {small} Затверджено на засіданні кафедри теоретичної кібернетики, протокол \No 12 від   19.05.2020 р.\par\smallskip\smallskip
{\parbox {6.5 cm}{Зав. кафедри \dotfill Ю.В.~Крак}}\hspace {0.7cm}{\hbox to 6.5 cm{ Екзаменатор \dotfill Т.О.~Карнаух}} \end{small}}
%begin
{{\begin {small}\par
{ \center  Київський національний університет імені Тараса Шевченка \par}
{\parbox {5 cm}{Освітня програма \hfill}{\parbox {7 cm}{Прикладна математика, Системний аналіз \hfill}}\parbox {6 cm}{\hfill Семестр 2}}\par
{\parbox {5 cm} {Навчальний предмет \hfill}\parbox {7 cm}{Програмування \hfill}\parbox {6cm}{\hfill}\par}\vspace{-7pt} \end{small}}{\center Екзаменаційний білет \No \textbf 
{ 225 }\par}}
{
\begin {large}
1. Що буде виведено за виконання?
code
try:
    res = 6//3%6*8
    if isinstance(res, bool):
        print(f'bool;{res}')
    elif isinstance(res, int):
        print(f'int;{res}')
    elif isinstance(res, float):
        print(f'float;{res:.2f}')
    else:
        print('???')
except:
    print('ERR')
code

2. Що буде виведено за виконання?
code
a = 0
b = 6
c = 4

def h(a):
    global c
    a = 4
    b *= 5
    c = 2
    return a + b + c

a = 2
b = 8
c = 1

try:
    print(h(a), a, b, c, sep=';')
except:
    print('ERR', a, b, c, sep=';')
code

3. Що буде виведено за виконання?
code
try:
    n = 13
    while n>=6:
        n += -3
    print(n, end=';')
except:
    print('ERR')
code

4. Що буде виведено за виконання?
code
try:
    for e in range(0, -15, -2):
        if e<3:
            continue
        print(e, end=';')
    else:
        print(e,end=';')
    print(e)
except:
    print('ERR')
code

5. Що буде виведено за виконання?\par (Дійсні значення, якщо наявні, в цьому завданні будуть виводитись у форматі з фіксованою крапкою та, можливо, з доволі великою кількістю десяткових розрядів. У відповіді для дійсних значень у ЦЬОМУ завданні достатньо вказати один десятковий розряд після крапки, відкинувши решту розрядів без округлення. Наприклад, замість 13.66666666666667 достатньо вказати 13.6, замість 25.23 достатньо вказати 25.2)
code
try:
    print(0, end=';')
    print(5%2, end=';')
    print(0, end=';')
except ZeroDivisionError:
    print(7, end=';')
except Exception:
    print(1, end=';')
except:
    print('ERR', end=';')
else:
    print(3, end=';')
finally:
    print(5)
code

6. Що буде виведено за виконання?
code
def f(s):
    for i in range(len(s)):
        s[i] += 1
    return
 
s = [1, 2, 3]
f(s)
print(s)
code

7. Що буде виведено за виконання?
code
class A:
    b = 1

    def __h1(self): print(2, end=';')
    def _h2(self): print(3, end=';')

    def main(self):
        try: print(self.b, end=';')
        except: print('ERR', end=';')

        try: self.__h1()
        except: print('ERR', end=';')

        try: self._h2()
        except: print('ERR', end=';')


class B(A):
    b = 4

    def __h1(self): print(5, end=';')
    def _h2(self): print(6, end=';')

try: B().main()
except: print('ERR', end=';')
code

8. Що буде виведено за виконання?
code
def f(arg, n):
    arg[77] = 1
    n += 3

try:
    n = 4
    d = {62:7, 12:8, 32:7}
    f(d, n)
    print(n, end=';')
    for x in d.items():
        print(x, sep=';', end=';')
except:
    print('ERR')
code

9. Що буде виведено за виконання?
code
class A:
    c = 1
    a = 2

    def __init__(self):
        c = 3

class B(A):
    a = 4

    def __init__(self):
        super().__init__()
        self.b = 7

obj = B()
obj.c = obj.a + 10
B.c = 13

try: print(obj.c, end= ';')
except: print('ERR', end=';')

try: print(A.c, end= ';')
except: print('ERR', end=';')

try: print(B.c, end= ';')
except: print('ERR', end=';')
code

10. 
Для графа на рисунку з вершини 0 запущено алгоритм пошуку в глибину, що будує кістякове дерево. Списки суміжності вершин переглядаються алгоритмом за зростанням міток вершин.\par
\begin{center}
	\adjustimage{max size={0.9\linewidth}{0.9\paperheight}}
	{../10/Traversals/150.png}
\end{center}
\par Які з наведених ребер входять до складу отриманого кістякового дерева?\par
1) (0, 5)\par
2) (1, 2)\par
3) (3, 5)\par
4) (5, 6)\par
5) (2, 6)\par
6) (0, 4)\par
{\vskip 15pt} У формі вибрати номери відповідних ребер.\par
%answer:2 4 5
\end {large}}
%end
{\smallskip \begin {small} Затверджено на засіданні кафедри теоретичної кібернетики, протокол \No 12 від   19.05.2020 р.\par\smallskip\smallskip
{\parbox {6.5 cm}{Зав. кафедри \dotfill Ю.В.~Крак}}\hspace {0.7cm}{\hbox to 6.5 cm{ Екзаменатор \dotfill Т.О.~Карнаух}} \end{small}}
%begin
{{\begin {small}\par
{ \center  Київський національний університет імені Тараса Шевченка \par}
{\parbox {5 cm}{Освітня програма \hfill}{\parbox {7 cm}{Прикладна математика, Системний аналіз \hfill}}\parbox {6 cm}{\hfill Семестр 2}}\par
{\parbox {5 cm} {Навчальний предмет \hfill}\parbox {7 cm}{Програмування \hfill}\parbox {6cm}{\hfill}\par}\vspace{-7pt} \end{small}}{\center Екзаменаційний білет \No \textbf 
{ 226 }\par}}
{
\begin {large}
1. Що буде виведено за виконання?
code
try:
    res = 4**0.5*2
    if isinstance(res, bool):
        print(f'bool;{res}')
    elif isinstance(res, int):
        print(f'int;{res}')
    elif isinstance(res, float):
        print(f'float;{res:.2f}')
    else:
        print('???')
except:
    print('ERR')
code

2. Що буде виведено за виконання?
code
a = 4
b = 6
c = 5

def g(b):
    global c
    a = 1
    b -= 3
    c = 4
    return a + b + c

a = 3
b = 2
c = 9

try:
    print(g(b), a, b, c, sep=';')
except:
    print('ERR', a, b, c, sep=';')
code

3. Що буде виведено за виконання?
code
try:
    t = 3
    while t<10:
        t += 1
    print(t, end=';')
except:
    print('ERR')
code

4. Що буде виведено за виконання?
code
try:
    for c in range(1, -8, 2):
        if c>=6:
            continue
        print(c, end=';')
    else:
        print('end',end=';')
    print(c)
except:
    print('ERR')
code

5. Що буде виведено за виконання?\par (Дійсні значення, якщо наявні, в цьому завданні будуть виводитись у форматі з фіксованою крапкою та, можливо, з доволі великою кількістю десяткових розрядів. У відповіді для дійсних значень у ЦЬОМУ завданні достатньо вказати один десятковий розряд після крапки, відкинувши решту розрядів без округлення. Наприклад, замість 13.66666666666667 достатньо вказати 13.6, замість 25.23 достатньо вказати 25.2)
code
try:
    print(2, end=';')
    print(int('4'), end=';')
    print(8, end=';')
except TypeError:
    print(9, end=';')
except KeyboardInterrupt:
    print(6, end=';')
except:
    print('ERR', end=';')
else:
    print(4, end=';')
finally:
    print(8)
code

6. Що буде виведено за виконання?
code
def f(n):
    if n>9:
        f(n//100)
    print(n%10, end='')
 
f(123456)
code

7. Що буде виведено за виконання?
code
class A:
    a = 1

    def _g1(self): print(2, end=';')
    def g2(self): print(3, end=';')

    def main(self):
        try: print(self.a, end=';')
        except: print('ERR', end=';')

        try: self._g1()
        except: print('ERR', end=';')

        try: self.g2()
        except: print('ERR', end=';')


class B(A):
    a = 4

    def _g1(self): print(5, end=';')
    def g2(self): print(6, end=';')

try: B().main()
except: print('ERR', end=';')
code

8. Що буде виведено за виконання?
code
def f(arg, n):
    arg[40] = 3
    n = 7
    arg = list(arg.items())

try:
    n = 2
    t = {56:7, 41:0, 72:5, 61:1}
    f(t, n)
    print(n, end=';')
    for x in t.keys():
        print(x, sep=';', end=';')
except:
    print('ERR')
code

9. Що буде виведено за виконання?
code
class A:
    a = 1
    c = 2

    def __init__(self):
        c = 3

class B(A):
    a = 4

    def __init__(self):
        super().__init__()
        self.b = 7

obj = B()
obj.a = obj.c + 10
B.a = 13

try: print(obj.a, end= ';')
except: print('ERR', end=';')

try: print(A.a, end= ';')
except: print('ERR', end=';')

try: print(B.a, end= ';')
except: print('ERR', end=';')
code

10. 
Для графа на рисунку з вершини 0 запущено алгоритм пошуку в глибину, що будує кістякове дерево. Списки суміжності вершин переглядаються алгоритмом за зростанням міток вершин.\par
\begin{center}
	\adjustimage{max size={0.9\linewidth}{0.9\paperheight}}
	{../10/Traversals/13.png}
\end{center}
\par Які з наведених ребер входять до складу отриманого кістякового дерева?\par
1) (2, 3)\par
2) (0, 5)\par
3) (3, 6)\par
4) (1, 3)\par
5) (2, 5)\par
6) (3, 4)\par
{\vskip 15pt} У формі вибрати номери відповідних ребер.\par
%answer:4 5
\end {large}}
%end
{\smallskip \begin {small} Затверджено на засіданні кафедри теоретичної кібернетики, протокол \No 12 від   19.05.2020 р.\par\smallskip\smallskip
{\parbox {6.5 cm}{Зав. кафедри \dotfill Ю.В.~Крак}}\hspace {0.7cm}{\hbox to 6.5 cm{ Екзаменатор \dotfill Т.О.~Карнаух}} \end{small}}
%begin
{{\begin {small}\par
{ \center  Київський національний університет імені Тараса Шевченка \par}
{\parbox {5 cm}{Освітня програма \hfill}{\parbox {7 cm}{Прикладна математика, Системний аналіз \hfill}}\parbox {6 cm}{\hfill Семестр 2}}\par
{\parbox {5 cm} {Навчальний предмет \hfill}\parbox {7 cm}{Програмування \hfill}\parbox {6cm}{\hfill}\par}\vspace{-7pt} \end{small}}{\center Екзаменаційний білет \No \textbf 
{ 227 }\par}}
{
\begin {large}
1. Що буде виведено за виконання?
code
try:
    res = "False"=="3j"
    if isinstance(res, bool):
        print(f'bool;{res}')
    elif isinstance(res, int):
        print(f'int;{res}')
    elif isinstance(res, float):
        print(f'float;{res:.2f}')
    else:
        print('???')
except:
    print('ERR')
code

2. Що буде виведено за виконання?
code
a = 4
b = 1
c = 2

def f(a):
    global c
    a = 2
    b = 1
    c = 5
    return a + b + c

a = 9
b = 1
c = 7

try:
    print(f(b), a, b, c, sep=';')
except:
    print('ERR', a, b, c, sep=';')
code

3. Що буде виведено за виконання?
code
try:
    m = 5
    while m>0:
        m += -2
    print(m, end=';')
except:
    print('ERR')
code

4. Що буде виведено за виконання?
code
try:
    for b in range(-9, 5, -3):
        if b<=7:
            continue
        print(b, end=';')
    else:
        print(b,end=';')
    print(b)
except:
    print('ERR')
code

5. Що буде виведено за виконання?\par (Дійсні значення, якщо наявні, в цьому завданні будуть виводитись у форматі з фіксованою крапкою та, можливо, з доволі великою кількістю десяткових розрядів. У відповіді для дійсних значень у ЦЬОМУ завданні достатньо вказати один десятковий розряд після крапки, відкинувши решту розрядів без округлення. Наприклад, замість 13.66666666666667 достатньо вказати 13.6, замість 25.23 достатньо вказати 25.2)
code
try:
    print(9, end=';')
    print(int('a3'), end=';')
    print(1, end=';')
except Exception:
    print(6, end=';')
except ValueError:
    print(5, end=';')
except:
    print('ERR', end=';')
else:
    print(2, end=';')
finally:
    print(7)
code

6. Що буде виведено за виконання?
code
def f(n):
    if n<=9:
        print(n%10, end='')
    else:
        f(n//10)

f(123456)
code

7. Що буде виведено за виконання?
code
class A:
    d = 1

    def __f1(self): print(2, end=';')
    def _f2(self): print(3, end=';')

    def main(self):
        try: print(self.d, end=';')
        except: print('ERR', end=';')

        try: self.__f1()
        except: print('ERR', end=';')

        try: self._f2()
        except: print('ERR', end=';')


class B(A):
    d = 4

    def __f1(self): print(5, end=';')
    def _f2(self): print(6, end=';')

try: B().main()
except: print('ERR', end=';')
code

8. Що буде виведено за виконання?
code
def f(arg, n):
    arg[21] = 2
    n = 8
    arg = set(arg.values())

try:
    n = 4
    d = {55:5, 60:1, 60:9}
    f(d, n)
    print(n, end=';')
    for x, y in d.items():
        print(x, y, sep=';', end=';')
except:
    print('ERR')
code

9. Що буде виведено за виконання?
code
class A:
    b = 1
    a = 2

    def __init__(self):
        a = 3

class B(A):
    a = 4

    def __init__(self):
        super().__init__()
        self.c = 7

obj = B()
obj.b = obj.a + 10
B.b = 13

try: print(obj.b, end= ';')
except: print('ERR', end=';')

try: print(A.b, end= ';')
except: print('ERR', end=';')

try: print(B.b, end= ';')
except: print('ERR', end=';')
code

10. 
Для графа на рисунку з вершини 4 запущено алгоритм пошуку в глибину, що будує кістякове дерево. Списки суміжності вершин переглядаються алгоритмом за зростанням міток вершин.\par
\begin{center}
	\adjustimage{max size={0.9\linewidth}{0.9\paperheight}}
	{../10/Traversals/222.png}
\end{center}
\par Які з наведених ребер входять до складу отриманого кістякового дерева?\par
1) (1, 4)\par
2) (3, 4)\par
3) (3, 5)\par
4) (0, 6)\par
5) (1, 5)\par
6) (1, 6)\par
{\vskip 15pt} У формі вибрати номери відповідних ребер.\par
%answer:1 3 4
\end {large}}
%end
{\smallskip \begin {small} Затверджено на засіданні кафедри теоретичної кібернетики, протокол \No 12 від   19.05.2020 р.\par\smallskip\smallskip
{\parbox {6.5 cm}{Зав. кафедри \dotfill Ю.В.~Крак}}\hspace {0.7cm}{\hbox to 6.5 cm{ Екзаменатор \dotfill Т.О.~Карнаух}} \end{small}}
%begin
{{\begin {small}\par
{ \center  Київський національний університет імені Тараса Шевченка \par}
{\parbox {5 cm}{Освітня програма \hfill}{\parbox {7 cm}{Прикладна математика, Системний аналіз \hfill}}\parbox {6 cm}{\hfill Семестр 2}}\par
{\parbox {5 cm} {Навчальний предмет \hfill}\parbox {7 cm}{Програмування \hfill}\parbox {6cm}{\hfill}\par}\vspace{-7pt} \end{small}}{\center Екзаменаційний білет \No \textbf 
{ 228 }\par}}
{
\begin {large}
1. Що буде виведено за виконання?
code
try:
    res = not 6<9 and False==7 or 9.0>4.0
    if isinstance(res, bool):
        print(f'bool;{res}')
    elif isinstance(res, int):
        print(f'int;{res}')
    elif isinstance(res, float):
        print(f'float;{res:.2f}')
    else:
        print('???')
except:
    print('ERR')
code

2. Що буде виведено за виконання?
code
a = 7
b = 0
c = 0

def f(a):
    a = 5
    b += 5
    c = 4
    return a + b + c

a = 5
b = 1
c = 3

try:
    print(f(a), a, b, c, sep=';')
except:
    print('ERR', a, b, c, sep=';')
code

3. Що буде виведено за виконання?
code
try:
    n = 1
    while n<=11:
        n += 2
    print(n, end=';')
except:
    print('ERR')
code

4. Що буде виведено за виконання?
code
try:
    for a in range(-6, 4, -1):
        if a>4:
            break
        print(a, end=';');a = -8
    else:
        print(a,end=';')
    print(a)
except:
    print('ERR')
code

5. Що буде виведено за виконання?\par (Дійсні значення, якщо наявні, в цьому завданні будуть виводитись у форматі з фіксованою крапкою та, можливо, з доволі великою кількістю десяткових розрядів. У відповіді для дійсних значень у ЦЬОМУ завданні достатньо вказати один десятковий розряд після крапки, відкинувши решту розрядів без округлення. Наприклад, замість 13.66666666666667 достатньо вказати 13.6, замість 25.23 достатньо вказати 25.2)
code
try:
    print(0, end=';')
    print(8%0, end=';')
    print(0, end=';')
except ZeroDivisionError:
    print(2, end=';')
except Exception:
    print(4, end=';')
except:
    print('ERR', end=';')
else:
    print(3, end=';')
finally:
    print(9)
code

6. Що буде виведено за виконання?
code
def f(n):
    print(n%10, end='')
    if n>9:
        f(n//10)
 
f(123456)
code

7. Що буде виведено за виконання?
code
class A:
    c = 1

    def _h1(self): print(2, end=';')
    def __h2(self): print(3, end=';')

    def main(self):
        try: print(self.c, end=';')
        except: print('ERR', end=';')

        try: self._h1()
        except: print('ERR', end=';')

        try: self.__h2()
        except: print('ERR', end=';')


class B(A):
    c = 4

    def _h1(self): print(5, end=';')
    def __h2(self): print(6, end=';')

try: B().main()
except: print('ERR', end=';')
code

8. Що буде виведено за виконання?
code
def f(arg, n):
    arg[62] = 3
    n += -3

try:
    n = 5
    t = {12:8, 62:1, 84:0}
    f(t, n)
    print(n, end=';')
    for x in t.keys():
        print(x, sep=';', end=';')
except:
    print('ERR')
code

9. Що буде виведено за виконання?
code
class A:
    c = 1
    b = 2

    def __init__(self):
        c = 3

class B(A):
    c = 4

    def __init__(self):
        super().__init__()
        self.a = 7

obj = B()
obj.c = obj.b + 10
B.c = 13

try: print(obj.c, end= ';')
except: print('ERR', end=';')

try: print(A.c, end= ';')
except: print('ERR', end=';')

try: print(B.c, end= ';')
except: print('ERR', end=';')
code

10. 
Для графа на рисунку з вершини 5 запущено алгоритм пошуку в глибину, що будує кістякове дерево. Списки суміжності вершин переглядаються алгоритмом за зростанням міток вершин.\par
\begin{center}
	\adjustimage{max size={0.9\linewidth}{0.9\paperheight}}
	{../10/Traversals/169.png}
\end{center}
\par Які з наведених ребер входять до складу отриманого кістякового дерева?\par
1) (0, 6)\par
2) (1, 6)\par
3) (5, 6)\par
4) (3, 6)\par
5) (0, 5)\par
6) (2, 5)\par
{\vskip 15pt} У формі вибрати номери відповідних ребер.\par
%answer:2 4 5
\end {large}}
%end
{\smallskip \begin {small} Затверджено на засіданні кафедри теоретичної кібернетики, протокол \No 12 від   19.05.2020 р.\par\smallskip\smallskip
{\parbox {6.5 cm}{Зав. кафедри \dotfill Ю.В.~Крак}}\hspace {0.7cm}{\hbox to 6.5 cm{ Екзаменатор \dotfill Т.О.~Карнаух}} \end{small}}
%begin
{{\begin {small}\par
{ \center  Київський національний університет імені Тараса Шевченка \par}
{\parbox {5 cm}{Освітня програма \hfill}{\parbox {7 cm}{Прикладна математика, Системний аналіз \hfill}}\parbox {6 cm}{\hfill Семестр 2}}\par
{\parbox {5 cm} {Навчальний предмет \hfill}\parbox {7 cm}{Програмування \hfill}\parbox {6cm}{\hfill}\par}\vspace{-7pt} \end{small}}{\center Екзаменаційний білет \No \textbf 
{ 229 }\par}}
{
\begin {large}
1. Що буде виведено за виконання?
code
try:
    res = not 2!=8.0 or True>=5.0 and 4<=3
    if isinstance(res, bool):
        print(f'bool;{res}')
    elif isinstance(res, int):
        print(f'int;{res}')
    elif isinstance(res, float):
        print(f'float;{res:.2f}')
    else:
        print('???')
except:
    print('ERR')
code

2. Що буде виведено за виконання?
code
a = 7
b = 9
c = 4

def g(b):
    a *= 3
    b = 2
    c = 1
    return a + b + c

a = 8
b = 6
c = 2

try:
    print(g(b), a, b, c, sep=';')
except:
    print('ERR', a, b, c, sep=';')
code

3. Що буде виведено за виконання?
code
try:
    k = 8
    while k>1:
        k += -2
    print(k, end=';')
except:
    print('ERR')
code

4. Що буде виведено за виконання?
code
try:
    for f in range(-2, -14, -2):
        if f<5:
            break
        print(f, end=';')
    else:
        print(f,end=';')
    print(f)
except:
    print('ERR')
code

5. Що буде виведено за виконання?\par (Дійсні значення, якщо наявні, в цьому завданні будуть виводитись у форматі з фіксованою крапкою та, можливо, з доволі великою кількістю десяткових розрядів. У відповіді для дійсних значень у ЦЬОМУ завданні достатньо вказати один десятковий розряд після крапки, відкинувши решту розрядів без округлення. Наприклад, замість 13.66666666666667 достатньо вказати 13.6, замість 25.23 достатньо вказати 25.2)
code
try:
    print(1, end=';')
    print(5j!=5j, end=';')
    print(6, end=';')
except TypeError:
    print(7, end=';')
except KeyboardInterrupt:
    print(6, end=';')
except:
    print('ERR', end=';')
else:
    print(7, end=';')
finally:
    print(5)
code

6. Що буде виведено за виконання?
code
def f(n):
    if n>9:
        f(n//10)
    print(n%10, end='')
 
f(123456)
code

7. Що буде виведено за виконання?
code
class A:
    b = 1

    def f1(self): print(2, end=';')
    def _f2(self): print(3, end=';')

    def main(self):
        try: print(self.b, end=';')
        except: print('ERR', end=';')

        try: self.f1()
        except: print('ERR', end=';')

        try: self._f2()
        except: print('ERR', end=';')


class B(A):
    b = 4

    def f1(self): print(5, end=';')
    def _f2(self): print(6, end=';')

try: B().main()
except: print('ERR', end=';')
code

8. Що буде виведено за виконання?
code
def f(arg, n):
    arg[43] = 0
    n = 5
    arg = tuple(arg.keys())

try:
    n = 1
    s = {83:6, 25:2, 37:1, 25:5}
    f(s, n)
    print(n, end=';')
    for x in s.items():
        print(x, sep=';', end=';')
except:
    print('ERR')
code

9. Що буде виведено за виконання?
code
class A:
    b = 1
    c = 2

    def __init__(self):
        b = 3

class B(A):
    b = 4

    def __init__(self):
        super().__init__()
        self.a = 7

obj = B()
obj.b = obj.c + 10
B.b = 13

try: print(obj.b, end= ';')
except: print('ERR', end=';')

try: print(A.b, end= ';')
except: print('ERR', end=';')

try: print(B.b, end= ';')
except: print('ERR', end=';')
code

10. 
Для графа на рисунку з вершини 5 запущено алгоритм пошуку в глибину, що будує кістякове дерево. Списки суміжності вершин переглядаються алгоритмом за зростанням міток вершин.\par
\begin{center}
	\adjustimage{max size={0.9\linewidth}{0.9\paperheight}}
	{../10/Traversals/64.png}
\end{center}
\par Які з наведених ребер входять до складу отриманого кістякового дерева?\par
1) (0, 4)\par
2) (2, 6)\par
3) (4, 6)\par
4) (5, 6)\par
5) (0, 6)\par
6) (1, 6)\par
{\vskip 15pt} У формі вибрати номери відповідних ребер.\par
%answer:5 6
\end {large}}
%end
{\smallskip \begin {small} Затверджено на засіданні кафедри теоретичної кібернетики, протокол \No 12 від   19.05.2020 р.\par\smallskip\smallskip
{\parbox {6.5 cm}{Зав. кафедри \dotfill Ю.В.~Крак}}\hspace {0.7cm}{\hbox to 6.5 cm{ Екзаменатор \dotfill Т.О.~Карнаух}} \end{small}}
%begin
{{\begin {small}\par
{ \center  Київський національний університет імені Тараса Шевченка \par}
{\parbox {5 cm}{Освітня програма \hfill}{\parbox {7 cm}{Прикладна математика, Системний аналіз \hfill}}\parbox {6 cm}{\hfill Семестр 2}}\par
{\parbox {5 cm} {Навчальний предмет \hfill}\parbox {7 cm}{Програмування \hfill}\parbox {6cm}{\hfill}\par}\vspace{-7pt} \end{small}}{\center Екзаменаційний білет \No \textbf 
{ 230 }\par}}
{
\begin {large}
1. Що буде виведено за виконання?
code
try:
    res = not False>7.0 and 3.0<5 or 8==6
    if isinstance(res, bool):
        print(f'bool;{res}')
    elif isinstance(res, int):
        print(f'int;{res}')
    elif isinstance(res, float):
        print(f'float;{res:.2f}')
    else:
        print('???')
except:
    print('ERR')
code

2. Що буде виведено за виконання?
code
a = 6
b = 3
c = 4

def f(a):
    a = 3
    b = 4
    c = 2
    return a + b + c

a = 8
b = 0
c = 5

try:
    print(f(b), a, b, c, sep=';')
except:
    print('ERR', a, b, c, sep=';')
code

3. Що буде виведено за виконання?
code
try:
    m = 5
    while m<=8:
        m += 2
    print(m, end=';')
except:
    print('ERR')
code

4. Що буде виведено за виконання?
code
try:
    for b in range(8, 10, 3):
        if b<=6:
            continue
        print(b, end=';')
        if b<4:
            break
    else:
        print(b,end=';')
    print(b)
except:
    print('ERR')
code

5. Що буде виведено за виконання?\par (Дійсні значення, якщо наявні, в цьому завданні будуть виводитись у форматі з фіксованою крапкою та, можливо, з доволі великою кількістю десяткових розрядів. У відповіді для дійсних значень у ЦЬОМУ завданні достатньо вказати один десятковий розряд після крапки, відкинувши решту розрядів без округлення. Наприклад, замість 13.66666666666667 достатньо вказати 13.6, замість 25.23 достатньо вказати 25.2)
code
try:
    print(0, end=';')
    print(8/1, end=';')
    print(3, end=';')
except ZeroDivisionError:
    print(2, end=';')
except ValueError:
    print(1, end=';')
except:
    print('ERR', end=';')
else:
    print(4, end=';')
finally:
    print(9)
code

6. Що буде виведено за виконання?
code
def f(s):
    for el in s:
        el += 1
        return s
 
s = [1, 2, 3]
s = f(s)
print(s)
code

7. Що буде виведено за виконання?
code
class A:
    d = 1

    def _h1(self): print(2, end=';')
    def _h2(self): print(3, end=';')

    def main(self):
        try: print(self.d, end=';')
        except: print('ERR', end=';')

        try: self._h1()
        except: print('ERR', end=';')

        try: self._h2()
        except: print('ERR', end=';')


class B(A):
    d = 4

    def _h1(self): print(5, end=';')
    def _h2(self): print(6, end=';')

try: B().main()
except: print('ERR', end=';')
code

8. Що буде виведено за виконання?
code
def f(arg, n):
    arg[67] = 0
    n *= -3

try:
    n = 6
    d = {67:2, 26:3, 67:4}
    f(d, n)
    print(n, end=';')
    for x in d.values():
        print(x, sep=';', end=';')
except:
    print('ERR')
code

9. Що буде виведено за виконання?
code
class A:
    c = 1
    b = 2

    def __init__(self):
        b = 3

class B(A):
    b = 4

    def __init__(self):
        super().__init__()
        self.a = 7

obj = B()
obj.a = obj.a + 10
B.a = 13

try: print(obj.a, end= ';')
except: print('ERR', end=';')

try: print(A.a, end= ';')
except: print('ERR', end=';')

try: print(B.a, end= ';')
except: print('ERR', end=';')
code

10. 
Для графа на рисунку з вершини 2 запущено алгоритм пошуку в глибину, що будує кістякове дерево. Списки суміжності вершин переглядаються алгоритмом за зростанням міток вершин.\par
\begin{center}
	\adjustimage{max size={0.9\linewidth}{0.9\paperheight}}
	{../10/Traversals/293.png}
\end{center}
\par Які з наведених ребер входять до складу отриманого кістякового дерева?\par
1) (1, 3)\par
2) (3, 5)\par
3) (4, 5)\par
4) (0, 1)\par
5) (0, 5)\par
6) (2, 6)\par
{\vskip 15pt} У формі вибрати номери відповідних ребер.\par
%answer:2 3 4 6
\end {large}}
%end
{\smallskip \begin {small} Затверджено на засіданні кафедри теоретичної кібернетики, протокол \No 12 від   19.05.2020 р.\par\smallskip\smallskip
{\parbox {6.5 cm}{Зав. кафедри \dotfill Ю.В.~Крак}}\hspace {0.7cm}{\hbox to 6.5 cm{ Екзаменатор \dotfill Т.О.~Карнаух}} \end{small}}
%begin
{{\begin {small}\par
{ \center  Київський національний університет імені Тараса Шевченка \par}
{\parbox {5 cm}{Освітня програма \hfill}{\parbox {7 cm}{Прикладна математика, Системний аналіз \hfill}}\parbox {6 cm}{\hfill Семестр 2}}\par
{\parbox {5 cm} {Навчальний предмет \hfill}\parbox {7 cm}{Програмування \hfill}\parbox {6cm}{\hfill}\par}\vspace{-7pt} \end{small}}{\center Екзаменаційний білет \No \textbf 
{ 231 }\par}}
{
\begin {large}
1. Що буде виведено за виконання?
code
try:
    res = 1**-0.5*-1
    if isinstance(res, bool):
        print(f'bool;{res}')
    elif isinstance(res, int):
        print(f'int;{res}')
    elif isinstance(res, float):
        print(f'float;{res:.2f}')
    else:
        print('???')
except:
    print('ERR')
code

2. Що буде виведено за виконання?
code
a = 8
b = 0
c = 4

def h(a):
    global c
    a = 1
    b += 4
    c = 5
    return a + b + c

a = 5
b = 2
c = 3

try:
    print(h(a), a, b, c, sep=';')
except:
    print('ERR', a, b, c, sep=';')
code

3. Що буде виведено за виконання?
code
try:
    k = 7
    while k<9:
        k += 1
    print(k, end=';')
except:
    print('ERR')
code

4. Що буде виведено за виконання?
code
try:
    for e in range(5, 13, -1):
        if e>3:
            continue
        print(e, end=';');e = -5
    else:
        print('end',end=';')
    print(e)
except:
    print('ERR')
code

5. Що буде виведено за виконання?\par (Дійсні значення, якщо наявні, в цьому завданні будуть виводитись у форматі з фіксованою крапкою та, можливо, з доволі великою кількістю десяткових розрядів. У відповіді для дійсних значень у ЦЬОМУ завданні достатньо вказати один десятковий розряд після крапки, відкинувши решту розрядів без округлення. Наприклад, замість 13.66666666666667 достатньо вказати 13.6, замість 25.23 достатньо вказати 25.2)
code
try:
    print(1, end=';')
    print(3//0, end=';')
    print(2, end=';')
except TypeError:
    print(7, end=';')
except ZeroDivisionError:
    print(0, end=';')
except:
    print('ERR', end=';')
else:
    print(9, end=';')
finally:
    print(6)
code

6. Що буде виведено за виконання?
code
def f(s):
    for i in range(len(s)):
        s[i] += 1
        return
 
s = [1, 2, 3]
f(s)
print(s)
code

7. Що буде виведено за виконання?
code
class A:
    a = 1

    def _g1(self): print(2, end=';')
    def __g2(self): print(3, end=';')

    def main(self):
        try: print(self.a, end=';')
        except: print('ERR', end=';')

        try: self._g1()
        except: print('ERR', end=';')

        try: self.__g2()
        except: print('ERR', end=';')


class B(A):
    a = 4

    def _g1(self): print(5, end=';')
    def __g2(self): print(6, end=';')

try: B().main()
except: print('ERR', end=';')
code

8. Що буде виведено за виконання?
code
def f(arg, n):
    arg[80] = 9
    n = 6
    arg = set(arg.items())

try:
    n = 1
    t = {80:1, 71:6, 12:3, 80:2}
    f(t, n)
    print(n, end=';')
    for x in t.values():
        print(x, sep=';', end=';')
except:
    print('ERR')
code

9. Що буде виведено за виконання?
code
class A:
    a = 1
    b = 2

    def __init__(self):
        a = 3

class B(A):
    b = 4

    def __init__(self):
        super().__init__()
        self.c = 7

obj = B()
obj.c = obj.b + 10
B.c = 13

try: print(obj.c, end= ';')
except: print('ERR', end=';')

try: print(A.c, end= ';')
except: print('ERR', end=';')

try: print(B.c, end= ';')
except: print('ERR', end=';')
code

10. 
Для графа на рисунку з вершини 1 запущено алгоритм пошуку в глибину, що будує кістякове дерево. Списки суміжності вершин переглядаються алгоритмом за зростанням міток вершин.\par
\begin{center}
	\adjustimage{max size={0.9\linewidth}{0.9\paperheight}}
	{../10/Traversals/282.png}
\end{center}
\par Які з наведених ребер входять до складу отриманого кістякового дерева?\par
1) (0, 3)\par
2) (0, 5)\par
3) (1, 2)\par
4) (0, 4)\par
5) (4, 5)\par
6) (3, 5)\par
{\vskip 15pt} У формі вибрати номери відповідних ребер.\par
%answer:1 3 5 6
\end {large}}
%end
{\smallskip \begin {small} Затверджено на засіданні кафедри теоретичної кібернетики, протокол \No 12 від   19.05.2020 р.\par\smallskip\smallskip
{\parbox {6.5 cm}{Зав. кафедри \dotfill Ю.В.~Крак}}\hspace {0.7cm}{\hbox to 6.5 cm{ Екзаменатор \dotfill Т.О.~Карнаух}} \end{small}}
%begin
{{\begin {small}\par
{ \center  Київський національний університет імені Тараса Шевченка \par}
{\parbox {5 cm}{Освітня програма \hfill}{\parbox {7 cm}{Прикладна математика, Системний аналіз \hfill}}\parbox {6 cm}{\hfill Семестр 2}}\par
{\parbox {5 cm} {Навчальний предмет \hfill}\parbox {7 cm}{Програмування \hfill}\parbox {6cm}{\hfill}\par}\vspace{-7pt} \end{small}}{\center Екзаменаційний білет \No \textbf 
{ 232 }\par}}
{
\begin {large}
1. Що буде виведено за виконання?
code
try:
    res = True<=6
    if isinstance(res, bool):
        print(f'bool;{res}')
    elif isinstance(res, int):
        print(f'int;{res}')
    elif isinstance(res, float):
        print(f'float;{res:.2f}')
    else:
        print('???')
except:
    print('ERR')
code

2. Що буде виведено за виконання?
code
a = 3
b = 7
c = 9

def g(b):
    global c
    a = 5
    b = 2
    c = 1
    return a + b + c

a = 0
b = 5
c = 1

try:
    print(g(b), a, b, c, sep=';')
except:
    print('ERR', a, b, c, sep=';')
code

3. Що буде виведено за виконання?
code
try:
    k = 0
    while k<6:
        k += 1
    print(k, end=';')
except:
    print('ERR')
code

4. Що буде виведено за виконання?
code
try:
    for d in range(-8, -11, 3):
        if d>=4:
            continue
        print(d, end=';')
    else:
        print(d,end=';')
    print(d)
except:
    print('ERR')
code

5. Що буде виведено за виконання?\par (Дійсні значення, якщо наявні, в цьому завданні будуть виводитись у форматі з фіксованою крапкою та, можливо, з доволі великою кількістю десяткових розрядів. У відповіді для дійсних значень у ЦЬОМУ завданні достатньо вказати один десятковий розряд після крапки, відкинувши решту розрядів без округлення. Наприклад, замість 13.66666666666667 достатньо вказати 13.6, замість 25.23 достатньо вказати 25.2)
code
try:
    print(8, end=';')
    print(int('d5'), end=';')
    print(4, end=';')
except Exception:
    print(3, end=';')
except ValueError:
    print(5, end=';')
except:
    print('ERR', end=';')
else:
    print(0, end=';')
finally:
    print(4)
code

6. Що буде виведено за виконання?
code
def f(s1):
    s = s1[:]
    for i in range(len(s)):
        s[i] += 1
 
s = [1, 2, 3]
f(s)
print(s)
code

7. Що буде виведено за виконання?
code
class A:
    c = 1

    def f1(self): print(2, end=';')
    def _f2(self): print(3, end=';')

    def main(self):
        try: print(self.c, end=';')
        except: print('ERR', end=';')

        try: self.f1()
        except: print('ERR', end=';')

        try: self._f2()
        except: print('ERR', end=';')


class B(A):
    c = 4

    def f1(self): print(5, end=';')
    def _f2(self): print(6, end=';')

try: B().main()
except: print('ERR', end=';')
code

8. Що буде виведено за виконання?
code
def f(arg, n):
    arg[17] = 1
    n = 7
    arg = tuple(arg.values())

try:
    n = 4
    s = {17:0, 58:8, 68:2, 68:8}
    f(s, n)
    print(n, end=';')
    for x in s.items():
        print(*x, sep=';', end=';')
except:
    print('ERR')
code

9. Що буде виведено за виконання?
code
class A:
    a = 1
    c = 2

    def __init__(self):
        c = 3

class B(A):
    a = 4

    def __init__(self):
        super().__init__()
        self.b = 7

obj = B()
obj.b = obj.a + 10
B.b = 13

try: print(obj.b, end= ';')
except: print('ERR', end=';')

try: print(A.b, end= ';')
except: print('ERR', end=';')

try: print(B.b, end= ';')
except: print('ERR', end=';')
code

10. 
Для графа на рисунку з вершини 3 запущено алгоритм пошуку в глибину, що будує кістякове дерево. Списки суміжності вершин переглядаються алгоритмом за зростанням міток вершин.\par
\begin{center}
	\adjustimage{max size={0.9\linewidth}{0.9\paperheight}}
	{../10/Traversals/163.png}
\end{center}
\par Які з наведених ребер входять до складу отриманого кістякового дерева?\par
1) (1, 5)\par
2) (1, 2)\par
3) (0, 3)\par
4) (1, 4)\par
5) (0, 5)\par
6) (0, 6)\par
{\vskip 15pt} У формі вибрати номери відповідних ребер.\par
%answer:1 2 3 4 6
\end {large}}
%end
{\smallskip \begin {small} Затверджено на засіданні кафедри теоретичної кібернетики, протокол \No 12 від   19.05.2020 р.\par\smallskip\smallskip
{\parbox {6.5 cm}{Зав. кафедри \dotfill Ю.В.~Крак}}\hspace {0.7cm}{\hbox to 6.5 cm{ Екзаменатор \dotfill Т.О.~Карнаух}} \end{small}}
%begin
{{\begin {small}\par
{ \center  Київський національний університет імені Тараса Шевченка \par}
{\parbox {5 cm}{Освітня програма \hfill}{\parbox {7 cm}{Прикладна математика, Системний аналіз \hfill}}\parbox {6 cm}{\hfill Семестр 2}}\par
{\parbox {5 cm} {Навчальний предмет \hfill}\parbox {7 cm}{Програмування \hfill}\parbox {6cm}{\hfill}\par}\vspace{-7pt} \end{small}}{\center Екзаменаційний білет \No \textbf 
{ 233 }\par}}
{
\begin {large}
1. Що буде виведено за виконання?
code
try:
    res = "False"=="9"
    if isinstance(res, bool):
        print(f'bool;{res}')
    elif isinstance(res, int):
        print(f'int;{res}')
    elif isinstance(res, float):
        print(f'float;{res:.2f}')
    else:
        print('???')
except:
    print('ERR')
code

2. Що буде виведено за виконання?
code
a = 6
b = 2
c = 8

def h(a):
    global c
    a = 3
    b = 4
    c = 5
    return a + b + c

a = 4
b = 6
c = 4

try:
    print(h(a), a, b, c, sep=';')
except:
    print('ERR', a, b, c, sep=';')
code

3. Що буде виведено за виконання?
code
try:
    m = 9
    while m<=5:
        m += 2
    print(m, end=';')
except:
    print('ERR')
code

4. Що буде виведено за виконання?
code
try:
    for a in range(14, 8, -2):
        if a>=8:
            continue
        print(a, end=';')
    else:
        print(a,end=';')
    print(a)
except:
    print('ERR')
code

5. Що буде виведено за виконання?\par (Дійсні значення, якщо наявні, в цьому завданні будуть виводитись у форматі з фіксованою крапкою та, можливо, з доволі великою кількістю десяткових розрядів. У відповіді для дійсних значень у ЦЬОМУ завданні достатньо вказати один десятковий розряд після крапки, відкинувши решту розрядів без округлення. Наприклад, замість 13.66666666666667 достатньо вказати 13.6, замість 25.23 достатньо вказати 25.2)
code
try:
    print(2, end=';')
    print(6//2, end=';')
    print(7, end=';')
except KeyboardInterrupt:
    print(9, end=';')
except ZeroDivisionError:
    print(8, end=';')
except:
    print('ERR', end=';')
else:
    print(1, end=';')
finally:
    print(6)
code

6. Що буде виведено за виконання?
code
def f(n):
    if n<=9:
        print(n%10, end='')
    else:
        f(n//10)

f(987651)
code

7. Що буде виведено за виконання?
code
class A:
    b = 1

    def _g1(self): print(2, end=';')
    def _g2(self): print(3, end=';')

    def main(self):
        try: print(self.b, end=';')
        except: print('ERR', end=';')

        try: self._g1()
        except: print('ERR', end=';')

        try: self._g2()
        except: print('ERR', end=';')


class B(A):
    b = 4

    def _g1(self): print(5, end=';')
    def _g2(self): print(6, end=';')

try: B().main()
except: print('ERR', end=';')
code

8. Що буде виведено за виконання?
code
def f(arg, n):
    n = 6
    tmp = arg
    tmp[-7:] = 'ab'
    return len(tmp)>3

try:
    e = [10,11,12,13,14]
    n = 4
    f(e, n)
    print(n, end=';')
    for el in e:
        print(el, end=';')
except:
    print('ERR')
code

9. Що буде виведено за виконання?
code
class A:
    c = 1
    a = 2

    def __init__(self):
        a = 3

class B(A):
    a = 4

    def __init__(self):
        super().__init__()
        self.b = 7

obj = B()
obj.c = obj.c + 10
B.c = 13

try: print(obj.c, end= ';')
except: print('ERR', end=';')

try: print(A.c, end= ';')
except: print('ERR', end=';')

try: print(B.c, end= ';')
except: print('ERR', end=';')
code

10. 
Для графа на рисунку з вершини 4 запущено алгоритм пошуку в глибину, що будує кістякове дерево. Списки суміжності вершин переглядаються алгоритмом за зростанням міток вершин.\par
\begin{center}
	\adjustimage{max size={0.9\linewidth}{0.9\paperheight}}
	{../10/Traversals/199.png}
\end{center}
\par Які з наведених ребер входять до складу отриманого кістякового дерева?\par
1) (5, 6)\par
2) (1, 3)\par
3) (0, 5)\par
4) (0, 2)\par
5) (3, 5)\par
6) (2, 4)\par
{\vskip 15pt} У формі вибрати номери відповідних ребер.\par
%answer:2 3 4 5 6
\end {large}}
%end
{\smallskip \begin {small} Затверджено на засіданні кафедри теоретичної кібернетики, протокол \No 12 від   19.05.2020 р.\par\smallskip\smallskip
{\parbox {6.5 cm}{Зав. кафедри \dotfill Ю.В.~Крак}}\hspace {0.7cm}{\hbox to 6.5 cm{ Екзаменатор \dotfill Т.О.~Карнаух}} \end{small}}
%begin
{{\begin {small}\par
{ \center  Київський національний університет імені Тараса Шевченка \par}
{\parbox {5 cm}{Освітня програма \hfill}{\parbox {7 cm}{Прикладна математика, Системний аналіз \hfill}}\parbox {6 cm}{\hfill Семестр 2}}\par
{\parbox {5 cm} {Навчальний предмет \hfill}\parbox {7 cm}{Програмування \hfill}\parbox {6cm}{\hfill}\par}\vspace{-7pt} \end{small}}{\center Екзаменаційний білет \No \textbf 
{ 234 }\par}}
{
\begin {large}
1. Що буде виведено за виконання?
code
try:
    res = 9/7*3*4
    if isinstance(res, bool):
        print(f'bool;{res}')
    elif isinstance(res, int):
        print(f'int;{res}')
    elif isinstance(res, float):
        print(f'float;{res:.2f}')
    else:
        print('???')
except:
    print('ERR')
code

2. Що буде виведено за виконання?
code
a = 6
b = 9
c = 1

def f(a):
    global c
    a = 3
    b -= 2
    c = 5
    return a + b + c

a = 7
b = 0
c = 8

try:
    print(f(b), a, b, c, sep=';')
except:
    print('ERR', a, b, c, sep=';')
code

3. Що буде виведено за виконання?
code
try:
    n = 4
    while n<7:
        n += 1
    print(n, end=';')
except:
    print('ERR')
code

4. Що буде виведено за виконання?
code
try:
    for c in range(-5, 14, -3):
        if c>7:
            break
        print(c, end=';');c = -2
    else:
        print(c,end=';')
    print(c)
except:
    print('ERR')
code

5. Що буде виведено за виконання?\par (Дійсні значення, якщо наявні, в цьому завданні будуть виводитись у форматі з фіксованою крапкою та, можливо, з доволі великою кількістю десяткових розрядів. У відповіді для дійсних значень у ЦЬОМУ завданні достатньо вказати один десятковий розряд після крапки, відкинувши решту розрядів без округлення. Наприклад, замість 13.66666666666667 достатньо вказати 13.6, замість 25.23 достатньо вказати 25.2)
code
try:
    print(2, end=';')
    print(int('0'), end=';')
    print(4, end=';')
except KeyboardInterrupt:
    print(7, end=';')
except Exception:
    print(3, end=';')
except:
    print('ERR', end=';')
else:
    print(9, end=';')
finally:
    print(5)
code

6. Що буде виведено за виконання?
code
def f(n):
    if n>9:
        f(n//100)
        print(n%10, end='')
 
f(123456)
code

7. Що буде виведено за виконання?
code
class A:
    a = 1

    def __h1(self): print(2, end=';')
    def _h2(self): print(3, end=';')

    def main(self):
        try: print(self.a, end=';')
        except: print('ERR', end=';')

        try: self.__h1()
        except: print('ERR', end=';')

        try: self._h2()
        except: print('ERR', end=';')


class B(A):
    a = 4

    def __h1(self): print(5, end=';')
    def _h2(self): print(6, end=';')

try: B().main()
except: print('ERR', end=';')
code

8. Що буде виведено за виконання?
code
def f(arg, n):
    arg[23] = 7
    n -= -2

try:
    n = 3
    s = {19:7, 82:5, 82:3}
    f(s, n)
    print(n, end=';')
    for x, y in s.items():
        print(x, y, sep=';', end=';')
except:
    print('ERR')
code

9. Що буде виведено за виконання?
code
class A:
    b = 1
    a = 2

    def __init__(self):
        b = 3

class B(A):
    b = 4

    def __init__(self):
        super().__init__()
        self.c = 7

obj = B()
obj.b = obj.a + 10
B.b = 13

try: print(obj.b, end= ';')
except: print('ERR', end=';')

try: print(A.b, end= ';')
except: print('ERR', end=';')

try: print(B.b, end= ';')
except: print('ERR', end=';')
code

10. 
Для графа на рисунку з вершини 4 запущено алгоритм пошуку в глибину, що будує кістякове дерево. Списки суміжності вершин переглядаються алгоритмом за зростанням міток вершин.\par
\begin{center}
	\adjustimage{max size={0.9\linewidth}{0.9\paperheight}}
	{../10/Traversals/59.png}
\end{center}
\par Які з наведених ребер входять до складу отриманого кістякового дерева?\par
1) (4, 5)\par
2) (2, 5)\par
3) (3, 5)\par
4) (1, 3)\par
5) (0, 5)\par
6) (1, 5)\par
{\vskip 15pt} У формі вибрати номери відповідних ребер.\par
%answer:2 3
\end {large}}
%end
{\smallskip \begin {small} Затверджено на засіданні кафедри теоретичної кібернетики, протокол \No 12 від   19.05.2020 р.\par\smallskip\smallskip
{\parbox {6.5 cm}{Зав. кафедри \dotfill Ю.В.~Крак}}\hspace {0.7cm}{\hbox to 6.5 cm{ Екзаменатор \dotfill Т.О.~Карнаух}} \end{small}}
%begin
{{\begin {small}\par
{ \center  Київський національний університет імені Тараса Шевченка \par}
{\parbox {5 cm}{Освітня програма \hfill}{\parbox {7 cm}{Прикладна математика, Системний аналіз \hfill}}\parbox {6 cm}{\hfill Семестр 2}}\par
{\parbox {5 cm} {Навчальний предмет \hfill}\parbox {7 cm}{Програмування \hfill}\parbox {6cm}{\hfill}\par}\vspace{-7pt} \end{small}}{\center Екзаменаційний білет \No \textbf 
{ 235 }\par}}
{
\begin {large}
1. Що буде виведено за виконання?
code
def f(n):
    print('f', end='')
    return n==-2

try:
    print(f(9) or f(-5))
except:
    print('ERR')
code

2. Що буде виведено за виконання?
code
a = 3
b = 0
c = 7

def f(b):
    a = 1
    b = 2
    c = 3
    return a + b + c

a = 2
b = 1
c = 9

try:
    print(f(b), a, b, c, sep=';')
except:
    print('ERR', a, b, c, sep=';')
code

3. Що буде виведено за виконання?
code
try:
    j = 6
    while j<=10:
        j += 2
    print(j, end=';')
except:
    print('ERR')
code

4. Що буде виведено за виконання?
code
try:
    for d in range(3, 2, 2):
        if d<=7:
            continue
        print(d, end=';')
    else:
        print(d,end=';')
    print(d)
except:
    print('ERR')
code

5. Що буде виведено за виконання?\par (Дійсні значення, якщо наявні, в цьому завданні будуть виводитись у форматі з фіксованою крапкою та, можливо, з доволі великою кількістю десяткових розрядів. У відповіді для дійсних значень у ЦЬОМУ завданні достатньо вказати один десятковий розряд після крапки, відкинувши решту розрядів без округлення. Наприклад, замість 13.66666666666667 достатньо вказати 13.6, замість 25.23 достатньо вказати 25.2)
code
try:
    print(8, end=';')
    print(1j<1j, end=';')
    print(2, end=';')
except ValueError:
    print(1, end=';')
except TypeError:
    print(6, end=';')
except:
    print('ERR', end=';')
else:
    print(9, end=';')
finally:
    print(7)
code

6. Що буде виведено за виконання?
code
def f(n):
    print(n%2, end='')
    if n>2:
        f(n//2)
 
f(16)
code

7. Що буде виведено за виконання?
code
class A:
    c = 1

    def _f1(self): print(2, end=';')
    def f2(self): print(3, end=';')

    def main(self):
        try: print(self.c, end=';')
        except: print('ERR', end=';')

        try: self._f1()
        except: print('ERR', end=';')

        try: self.f2()
        except: print('ERR', end=';')


class B(A):
    c = 4

    def _f1(self): print(5, end=';')
    def f2(self): print(6, end=';')

try: B().main()
except: print('ERR', end=';')
code

8. Що буде виведено за виконання?
code
def f(arg, n):
    arg[13] = 6
    n += 3

try:
    n = 7
    d = {13:1, 85:8, 18:5, 85:9}
    f(d, n)
    print(n, end=';')
    for x in d :
        print(x, sep=';', end=';')
except:
    print('ERR')
code

9. Що буде виведено за виконання?
code
class A:
    b = 1
    a = 2

    def __init__(self):
        b = 3

class B(A):
    b = 4

    def __init__(self):
        super().__init__()
        self.c = 7

obj = B()
obj.b = obj.a + 10
B.b = 13

try: print(obj.b, end= ';')
except: print('ERR', end=';')

try: print(A.b, end= ';')
except: print('ERR', end=';')

try: print(B.b, end= ';')
except: print('ERR', end=';')
code

10. 
Для графа на рисунку з вершини 3 запущено алгоритм пошуку в глибину, що будує кістякове дерево. Списки суміжності вершин переглядаються алгоритмом за зростанням міток вершин.\par
\begin{center}
	\adjustimage{max size={0.9\linewidth}{0.9\paperheight}}
	{../10/Traversals/30.png}
\end{center}
\par Які з наведених ребер входять до складу отриманого кістякового дерева?\par
1) (1, 4)\par
2) (2, 5)\par
3) (5, 6)\par
4) (1, 6)\par
5) (3, 6)\par
6) (0, 5)\par
{\vskip 15pt} У формі вибрати номери відповідних ребер.\par
%answer:1 2 4 6
\end {large}}
%end
{\smallskip \begin {small} Затверджено на засіданні кафедри теоретичної кібернетики, протокол \No 12 від   19.05.2020 р.\par\smallskip\smallskip
{\parbox {6.5 cm}{Зав. кафедри \dotfill Ю.В.~Крак}}\hspace {0.7cm}{\hbox to 6.5 cm{ Екзаменатор \dotfill Т.О.~Карнаух}} \end{small}}
%begin
{{\begin {small}\par
{ \center  Київський національний університет імені Тараса Шевченка \par}
{\parbox {5 cm}{Освітня програма \hfill}{\parbox {7 cm}{Прикладна математика, Системний аналіз \hfill}}\parbox {6 cm}{\hfill Семестр 2}}\par
{\parbox {5 cm} {Навчальний предмет \hfill}\parbox {7 cm}{Програмування \hfill}\parbox {6cm}{\hfill}\par}\vspace{-7pt} \end{small}}{\center Екзаменаційний білет \No \textbf 
{ 236 }\par}}
{
\begin {large}
1. Що буде виведено за виконання?
code
try:
    res = 4>=4.0 or not False>3.0 and 2==5
    if isinstance(res, bool):
        print(f'bool;{res}')
    elif isinstance(res, int):
        print(f'int;{res}')
    elif isinstance(res, float):
        print(f'float;{res:.2f}')
    else:
        print('???')
except:
    print('ERR')
code

2. Що буде виведено за виконання?
code
a = 1
b = 5
c = 9

def h(b):
    a -= 4
    b = 1
    c = 2
    return a + b + c

a = 6
b = 3
c = 4

try:
    print(h(a), a, b, c, sep=';')
except:
    print('ERR', a, b, c, sep=';')
code

3. Що буде виведено за виконання?
code
try:
    t = 8
    while t<=7:
        t += 1
    print(t, end=';')
except:
    print('ERR')
code

4. Що буде виведено за виконання?
code
try:
    for e in range(15, -10, -3):
        if e<5:
            continue
        print(e, end=';')
    else:
        print(e,end=';')
    print(e)
except:
    print('ERR')
code

5. Що буде виведено за виконання?\par (Дійсні значення, якщо наявні, в цьому завданні будуть виводитись у форматі з фіксованою крапкою та, можливо, з доволі великою кількістю десяткових розрядів. У відповіді для дійсних значень у ЦЬОМУ завданні достатньо вказати один десятковий розряд після крапки, відкинувши решту розрядів без округлення. Наприклад, замість 13.66666666666667 достатньо вказати 13.6, замість 25.23 достатньо вказати 25.2)
code
try:
    print(5, end=';')
    print(3/3, end=';')
    print(8, end=';')
except TypeError:
    print(4, end=';')
except ValueError:
    print(0, end=';')
except:
    print('ERR', end=';')
else:
    print(8, end=';')
finally:
    print(4)
code

6. Що буде виведено за виконання?
code
def f(n):
    if n>9:
        f(n//100)
    else:
        print(n%10,end='')

f(987654)
code

7. Що буде виведено за виконання?
code
class A:
    d = 1

    def _g1(self): print(2, end=';')
    def __g2(self): print(3, end=';')

    def main(self):
        try: print(self.d, end=';')
        except: print('ERR', end=';')

        try: self._g1()
        except: print('ERR', end=';')

        try: self.__g2()
        except: print('ERR', end=';')


class B(A):
    d = 4

    def _g1(self): print(5, end=';')
    def __g2(self): print(6, end=';')

try: B().main()
except: print('ERR', end=';')
code

8. Що буде виведено за виконання?
code
def f(arg, n):
    arg[67] = 2
    n = 8
    arg = list(arg.keys())

try:
    n = 2
    d = {89:8, 67:0, 89:5}
    f(d, n)
    print(n, end=';')
    for x, y in d.items():
        print(x, y, sep=';', end=';')
except:
    print('ERR')
code

9. Що буде виведено за виконання?
code
class A:
    a = 1
    b = 2

    def __init__(self):
        b = 3

class B(A):
    b = 4

    def __init__(self):
        super().__init__()
        self.c = 7

obj = B()
obj.c = obj.c + 10
B.c = 13

try: print(obj.c, end= ';')
except: print('ERR', end=';')

try: print(A.c, end= ';')
except: print('ERR', end=';')

try: print(B.c, end= ';')
except: print('ERR', end=';')
code

10. 
Для графа на рисунку з вершини 6 запущено алгоритм пошуку в глибину, що будує кістякове дерево. Списки суміжності вершин переглядаються алгоритмом за зростанням міток вершин.\par
\begin{center}
	\adjustimage{max size={0.9\linewidth}{0.9\paperheight}}
	{../10/Traversals/242.png}
\end{center}
\par Які з наведених ребер входять до складу отриманого кістякового дерева?\par
1) (4, 5)\par
2) (0, 2)\par
3) (1, 5)\par
4) (3, 4)\par
5) (1, 2)\par
6) (4, 6)\par
{\vskip 15pt} У формі вибрати номери відповідних ребер.\par
%answer:2 3 5 6
\end {large}}
%end
{\smallskip \begin {small} Затверджено на засіданні кафедри теоретичної кібернетики, протокол \No 12 від   19.05.2020 р.\par\smallskip\smallskip
{\parbox {6.5 cm}{Зав. кафедри \dotfill Ю.В.~Крак}}\hspace {0.7cm}{\hbox to 6.5 cm{ Екзаменатор \dotfill Т.О.~Карнаух}} \end{small}}
%begin
{{\begin {small}\par
{ \center  Київський національний університет імені Тараса Шевченка \par}
{\parbox {5 cm}{Освітня програма \hfill}{\parbox {7 cm}{Прикладна математика, Системний аналіз \hfill}}\parbox {6 cm}{\hfill Семестр 2}}\par
{\parbox {5 cm} {Навчальний предмет \hfill}\parbox {7 cm}{Програмування \hfill}\parbox {6cm}{\hfill}\par}\vspace{-7pt} \end{small}}{\center Екзаменаційний білет \No \textbf 
{ 237 }\par}}
{
\begin {large}
1. Що буде виведено за виконання?
code
try:
    res = True!=9 and not 9.0<6 or 5.0<=3
    if isinstance(res, bool):
        print(f'bool;{res}')
    elif isinstance(res, int):
        print(f'int;{res}')
    elif isinstance(res, float):
        print(f'float;{res:.2f}')
    else:
        print('???')
except:
    print('ERR')
code

2. Що буде виведено за виконання?
code
a = 8
b = 5
c = 7

def h(a):
    global c
    a = 4
    b = 3
    c = 2
    return a + b + c

a = 5
b = 3
c = 4

try:
    print(h(a), a, b, c, sep=';')
except:
    print('ERR', a, b, c, sep=';')
code

3. Що буде виведено за виконання?
code
try:
    i = 2
    while i<11:
        i += 2
    print(i, end=';')
except:
    print('ERR')
code

4. Що буде виведено за виконання?
code
try:
    for c in range(2, -5, 2):
        if c<=3:
            break
        print(c, end=';');c = 10
    else:
        print(c,end=';')
    print(c)
except:
    print('ERR')
code

5. Що буде виведено за виконання?\par (Дійсні значення, якщо наявні, в цьому завданні будуть виводитись у форматі з фіксованою крапкою та, можливо, з доволі великою кількістю десяткових розрядів. У відповіді для дійсних значень у ЦЬОМУ завданні достатньо вказати один десятковий розряд після крапки, відкинувши решту розрядів без округлення. Наприклад, замість 13.66666666666667 достатньо вказати 13.6, замість 25.23 достатньо вказати 25.2)
code
try:
    print(2, end=';')
    print(6j>6j, end=';')
    print(7, end=';')
except Exception:
    print(5, end=';')
except ZeroDivisionError:
    print(3, end=';')
except:
    print('ERR', end=';')
else:
    print(0, end=';')
finally:
    print(1)
code

6. Що буде виведено за виконання?
code
def f(n):
    if n>9:
        f(n//100)
    print(n%10,end='')
 
f(987654)
code

7. Що буде виведено за виконання?
code
class A:
    b = 1

    def _h1(self): print(2, end=';')
    def _h2(self): print(3, end=';')

    def main(self):
        try: print(self.b, end=';')
        except: print('ERR', end=';')

        try: self._h1()
        except: print('ERR', end=';')

        try: self._h2()
        except: print('ERR', end=';')


class B(A):
    b = 4

    def _h1(self): print(5, end=';')
    def _h2(self): print(6, end=';')

try: B().main()
except: print('ERR', end=';')
code

8. Що буде виведено за виконання?
code
def f(arg, n):
    arg[49] = 8
    n *= 2

try:
    n = 4
    s = {49:4, 81:8, 49:8}
    f(s, n)
    print(n, end=';')
    for x in s.items():
        print(*x, sep=';', end=';')
except:
    print('ERR')
code

9. Що буде виведено за виконання?
code
class A:
    c = 1
    b = 2

    def __init__(self):
        b = 3

class B(A):
    b = 4

    def __init__(self):
        super().__init__()
        self.a = 7

obj = B()
obj.a = obj.b + 10
B.a = 13

try: print(obj.a, end= ';')
except: print('ERR', end=';')

try: print(A.a, end= ';')
except: print('ERR', end=';')

try: print(B.a, end= ';')
except: print('ERR', end=';')
code

10. 
Для графа на рисунку з вершини 3 запущено алгоритм пошуку в глибину, що будує кістякове дерево. Списки суміжності вершин переглядаються алгоритмом за зростанням міток вершин.\par
\begin{center}
	\adjustimage{max size={0.9\linewidth}{0.9\paperheight}}
	{../10/Traversals/46.png}
\end{center}
\par Які з наведених ребер входять до складу отриманого кістякового дерева?\par
1) (1, 3)\par
2) (3, 5)\par
3) (0, 2)\par
4) (0, 1)\par
5) (4, 6)\par
6) (2, 6)\par
{\vskip 15pt} У формі вибрати номери відповідних ребер.\par
%answer:1 3 4 5 6
\end {large}}
%end
{\smallskip \begin {small} Затверджено на засіданні кафедри теоретичної кібернетики, протокол \No 12 від   19.05.2020 р.\par\smallskip\smallskip
{\parbox {6.5 cm}{Зав. кафедри \dotfill Ю.В.~Крак}}\hspace {0.7cm}{\hbox to 6.5 cm{ Екзаменатор \dotfill Т.О.~Карнаух}} \end{small}}
%begin
{{\begin {small}\par
{ \center  Київський національний університет імені Тараса Шевченка \par}
{\parbox {5 cm}{Освітня програма \hfill}{\parbox {7 cm}{Прикладна математика, Системний аналіз \hfill}}\parbox {6 cm}{\hfill Семестр 2}}\par
{\parbox {5 cm} {Навчальний предмет \hfill}\parbox {7 cm}{Програмування \hfill}\parbox {6cm}{\hfill}\par}\vspace{-7pt} \end{small}}{\center Екзаменаційний білет \No \textbf 
{ 238 }\par}}
{
\begin {large}
1. Що буде виведено за виконання?
code
try:
    res = 5//8/4*5
    if isinstance(res, bool):
        print(f'bool;{res}')
    elif isinstance(res, int):
        print(f'int;{res}')
    elif isinstance(res, float):
        print(f'float;{res:.2f}')
    else:
        print('???')
except:
    print('ERR')
code

2. Що буде виведено за виконання?
code
a = 0
b = 6
c = 1

def g(b):
    a = 4
    b = 5
    c = 1
    return a + b + c

a = 9
b = 8
c = 2

try:
    print(g(a), a, b, c, sep=';')
except:
    print('ERR', a, b, c, sep=';')
code

3. Що буде виведено за виконання?
code
try:
    j = 7
    while j>=7:
        j += -3
    print(j, end=';')
except:
    print('ERR')
code

4. Що буде виведено за виконання?
code
try:
    for b in range(10, 6, -1):
        if b>=6:
            continue
        print(b, end=';')
    else:
        print(b,end=';')
    print(b)
except:
    print('ERR')
code

5. Що буде виведено за виконання?\par (Дійсні значення, якщо наявні, в цьому завданні будуть виводитись у форматі з фіксованою крапкою та, можливо, з доволі великою кількістю десяткових розрядів. У відповіді для дійсних значень у ЦЬОМУ завданні достатньо вказати один десятковий розряд після крапки, відкинувши решту розрядів без округлення. Наприклад, замість 13.66666666666667 достатньо вказати 13.6, замість 25.23 достатньо вказати 25.2)
code
try:
    print(9, end=';')
    print(int('3'), end=';')
    print(7, end=';')
except KeyboardInterrupt:
    print(6, end=';')
except ValueError:
    print(5, end=';')
except:
    print('ERR', end=';')
else:
    print(9, end=';')
finally:
    print(0)
code

6. Що буде виведено за виконання?
code
def f(n):
    print(n%10, end='')
    if n>9:
        f(n//10)
 
f(543216)
code

7. Що буде виведено за виконання?
code
class A:
    a = 1

    def _g1(self): print(2, end=';')
    def g2(self): print(3, end=';')

    def main(self):
        try: print(self.a, end=';')
        except: print('ERR', end=';')

        try: self._g1()
        except: print('ERR', end=';')

        try: self.g2()
        except: print('ERR', end=';')


class B(A):
    a = 4

    def _g1(self): print(5, end=';')
    def g2(self): print(6, end=';')

try: B().main()
except: print('ERR', end=';')
code

8. Що буде виведено за виконання?
code
def f(arg, n):
    n = 0
    tmp = arg
    tmp[0:2] = [15]
    return len(tmp)>3

try:
    c = ['0','1','2','3','4']
    n = 3
    f(c, n)
    print(n, end=';')
    for el in c:
        print(el, end=';')
except:
    print('ERR')
code

9. Що буде виведено за виконання?
code
class A:
    b = 1
    c = 2

    def __init__(self):
        b = 3

class B(A):
    b = 4

    def __init__(self):
        super().__init__()
        self.a = 7

obj = B()
obj.b = obj.a + 10
B.b = 13

try: print(obj.b, end= ';')
except: print('ERR', end=';')

try: print(A.b, end= ';')
except: print('ERR', end=';')

try: print(B.b, end= ';')
except: print('ERR', end=';')
code

10. 
Для графа на рисунку з вершини 6 запущено алгоритм пошуку в глибину, що будує кістякове дерево. Списки суміжності вершин переглядаються алгоритмом за зростанням міток вершин.\par
\begin{center}
	\adjustimage{max size={0.9\linewidth}{0.9\paperheight}}
	{../10/Traversals/255.png}
\end{center}
\par Які з наведених ребер входять до складу отриманого кістякового дерева?\par
1) (2, 5)\par
2) (1, 3)\par
3) (5, 6)\par
4) (0, 5)\par
5) (1, 2)\par
6) (1, 5)\par
{\vskip 15pt} У формі вибрати номери відповідних ребер.\par
%answer:1 2
\end {large}}
%end
{\smallskip \begin {small} Затверджено на засіданні кафедри теоретичної кібернетики, протокол \No 12 від   19.05.2020 р.\par\smallskip\smallskip
{\parbox {6.5 cm}{Зав. кафедри \dotfill Ю.В.~Крак}}\hspace {0.7cm}{\hbox to 6.5 cm{ Екзаменатор \dotfill Т.О.~Карнаух}} \end{small}}
%begin
{{\begin {small}\par
{ \center  Київський національний університет імені Тараса Шевченка \par}
{\parbox {5 cm}{Освітня програма \hfill}{\parbox {7 cm}{Прикладна математика, Системний аналіз \hfill}}\parbox {6 cm}{\hfill Семестр 2}}\par
{\parbox {5 cm} {Навчальний предмет \hfill}\parbox {7 cm}{Програмування \hfill}\parbox {6cm}{\hfill}\par}\vspace{-7pt} \end{small}}{\center Екзаменаційний білет \No \textbf 
{ 239 }\par}}
{
\begin {large}
1. Що буде виведено за виконання?
code
try:
    res = 8.0<=8 and not 7==False and 8!=2.0
    if isinstance(res, bool):
        print(f'bool;{res}')
    elif isinstance(res, int):
        print(f'int;{res}')
    elif isinstance(res, float):
        print(f'float;{res:.2f}')
    else:
        print('???')
except:
    print('ERR')
code

2. Що буде виведено за виконання?
code
a = 2
b = 7
c = 1

def f(a):
    global c
    a = 3
    b *= 4
    c = 3
    return a + b + c

a = 8
b = 9
c = 3

try:
    print(f(a), a, b, c, sep=';')
except:
    print('ERR', a, b, c, sep=';')
code

3. Що буде виведено за виконання?
code
try:
    t = 9
    while t>=5:
        t += -3
    print(t, end=';')
except:
    print('ERR')
code

4. Що буде виведено за виконання?
code
try:
    for a in range(12, -7, 3):
        if a>8:
            continue
        print(a, end=';');a = -7
        if a>7:
            break
    else:
        print('end',end=';')
    print(a)
except:
    print('ERR')
code

5. Що буде виведено за виконання?\par (Дійсні значення, якщо наявні, в цьому завданні будуть виводитись у форматі з фіксованою крапкою та, можливо, з доволі великою кількістю десяткових розрядів. У відповіді для дійсних значень у ЦЬОМУ завданні достатньо вказати один десятковий розряд після крапки, відкинувши решту розрядів без округлення. Наприклад, замість 13.66666666666667 достатньо вказати 13.6, замість 25.23 достатньо вказати 25.2)
code
try:
    print(8, end=';')
    print(1%0.0, end=';')
    print(4, end=';')
except TypeError:
    print(2, end=';')
except KeyboardInterrupt:
    print(8, end=';')
except:
    print('ERR', end=';')
else:
    print(3, end=';')
finally:
    print(7)
code

6. Що буде виведено за виконання?
code
def f(n):
    if n>9:
        f(n//10)
    print(n%10, end='')
 
f(543216)
code

7. Що буде виведено за виконання?
code
class A:
    d = 1

    def __h1(self): print(2, end=';')
    def _h2(self): print(3, end=';')

    def main(self):
        try: print(self.d, end=';')
        except: print('ERR', end=';')

        try: self.__h1()
        except: print('ERR', end=';')

        try: self._h2()
        except: print('ERR', end=';')


class B(A):
    d = 4

    def __h1(self): print(5, end=';')
    def _h2(self): print(6, end=';')

try: B().main()
except: print('ERR', end=';')
code

8. Що буде виведено за виконання?
code
def f(arg, n):
    arg[53] = 9
    n -= -3

try:
    n = 5
    t = {86:2, 57:8, 71:2, 86:3}
    f(t, n)
    print(n, end=';')
    for x in t :
        print(x, sep=';', end=';')
except:
    print('ERR')
code

9. Що буде виведено за виконання?
code
class A:
    a = 1
    c = 2

    def __init__(self):
        c = 3

class B(A):
    c = 4

    def __init__(self):
        super().__init__()
        self.b = 7

obj = B()
obj.c = obj.c + 10
B.c = 13

try: print(obj.c, end= ';')
except: print('ERR', end=';')

try: print(A.c, end= ';')
except: print('ERR', end=';')

try: print(B.c, end= ';')
except: print('ERR', end=';')
code

10. 
Для графа на рисунку з вершини 0 запущено алгоритм пошуку в глибину, що будує кістякове дерево. Списки суміжності вершин переглядаються алгоритмом за зростанням міток вершин.\par
\begin{center}
	\adjustimage{max size={0.9\linewidth}{0.9\paperheight}}
	{../10/Traversals/288.png}
\end{center}
\par Які з наведених ребер входять до складу отриманого кістякового дерева?\par
1) (3, 4)\par
2) (5, 6)\par
3) (1, 2)\par
4) (0, 2)\par
5) (0, 3)\par
6) (1, 3)\par
{\vskip 15pt} У формі вибрати номери відповідних ребер.\par
%answer:1 2 3 4 6
\end {large}}
%end
{\smallskip \begin {small} Затверджено на засіданні кафедри теоретичної кібернетики, протокол \No 12 від   19.05.2020 р.\par\smallskip\smallskip
{\parbox {6.5 cm}{Зав. кафедри \dotfill Ю.В.~Крак}}\hspace {0.7cm}{\hbox to 6.5 cm{ Екзаменатор \dotfill Т.О.~Карнаух}} \end{small}}
%begin
{{\begin {small}\par
{ \center  Київський національний університет імені Тараса Шевченка \par}
{\parbox {5 cm}{Освітня програма \hfill}{\parbox {7 cm}{Прикладна математика, Системний аналіз \hfill}}\parbox {6 cm}{\hfill Семестр 2}}\par
{\parbox {5 cm} {Навчальний предмет \hfill}\parbox {7 cm}{Програмування \hfill}\parbox {6cm}{\hfill}\par}\vspace{-7pt} \end{small}}{\center Екзаменаційний білет \No \textbf 
{ 240 }\par}}
{
\begin {large}
1. Що буде виведено за виконання?
code
try:
    res = 8/9//5*7
    if isinstance(res, bool):
        print(f'bool;{res}')
    elif isinstance(res, int):
        print(f'int;{res}')
    elif isinstance(res, float):
        print(f'float;{res:.2f}')
    else:
        print('???')
except:
    print('ERR')
code

2. Що буде виведено за виконання?
code
a = 2
b = 6
c = 5

def g(b):
    a = 2
    b += 1
    c = 5
    return a + b + c

a = 4
b = 7
c = 0

try:
    print(g(b), a, b, c, sep=';')
except:
    print('ERR', a, b, c, sep=';')
code

3. Що буде виведено за виконання?
code
try:
    i = 14
    while i>8:
        i += -2
    print(i, end=';')
except:
    print('ERR')
code

4. Що буде виведено за виконання?
code
try:
    for f in range(4, -12, -2):
        if f<4:
            break
        print(f, end=';')
    else:
        print(f,end=';')
    print(f)
except:
    print('ERR')
code

5. Що буде виведено за виконання?\par (Дійсні значення, якщо наявні, в цьому завданні будуть виводитись у форматі з фіксованою крапкою та, можливо, з доволі великою кількістю десяткових розрядів. У відповіді для дійсних значень у ЦЬОМУ завданні достатньо вказати один десятковий розряд після крапки, відкинувши решту розрядів без округлення. Наприклад, замість 13.66666666666667 достатньо вказати 13.6, замість 25.23 достатньо вказати 25.2)
code
try:
    print(1, end=';')
    print(int('c0'), end=';')
    print(9, end=';')
except ZeroDivisionError:
    print(2, end=';')
except Exception:
    print(5, end=';')
except:
    print('ERR', end=';')
else:
    print(4, end=';')
finally:
    print(6)
code

6. Що буде виведено за виконання?
code
def f(n):
    if n>9: f(n//10)
    else:
        print(n%10, end='')
 
f(123456)
code

7. Що буде виведено за виконання?
code
class A:
    c = 1

    def f1(self): print(2, end=';')
    def _f2(self): print(3, end=';')

    def main(self):
        try: print(self.c, end=';')
        except: print('ERR', end=';')

        try: self.f1()
        except: print('ERR', end=';')

        try: self._f2()
        except: print('ERR', end=';')


class B(A):
    c = 4

    def f1(self): print(5, end=';')
    def _f2(self): print(6, end=';')

try: B().main()
except: print('ERR', end=';')
code

8. Що буде виведено за виконання?
code
def f(arg, n):
    n = 5
    tmp = arg
    del tmp[:-8]
    return len(tmp)>3

try:
    f = ['0','1','2','3','4']
    n = 6
    f(f, n)
    print(n, end=';')
    for el in f:
        print(el, end=';')
except:
    print('ERR')
code

9. Що буде виведено за виконання?
code
class A:
    c = 1
    a = 2

    def __init__(self):
        c = 3

class B(A):
    c = 4

    def __init__(self):
        super().__init__()
        self.b = 7

obj = B()
obj.b = obj.a + 10
B.b = 13

try: print(obj.b, end= ';')
except: print('ERR', end=';')

try: print(A.b, end= ';')
except: print('ERR', end=';')

try: print(B.b, end= ';')
except: print('ERR', end=';')
code

10. 
Для графа на рисунку з вершини 6 запущено алгоритм пошуку в глибину, що будує кістякове дерево. Списки суміжності вершин переглядаються алгоритмом за зростанням міток вершин.\par
\begin{center}
	\adjustimage{max size={0.9\linewidth}{0.9\paperheight}}
	{../10/Traversals/232.png}
\end{center}
\par Які з наведених ребер входять до складу отриманого кістякового дерева?\par
1) (3, 4)\par
2) (0, 5)\par
3) (0, 6)\par
4) (2, 4)\par
5) (2, 6)\par
6) (2, 5)\par
{\vskip 15pt} У формі вибрати номери відповідних ребер.\par
%answer:1 3 4 6
\end {large}}
%end
{\smallskip \begin {small} Затверджено на засіданні кафедри теоретичної кібернетики, протокол \No 12 від   19.05.2020 р.\par\smallskip\smallskip
{\parbox {6.5 cm}{Зав. кафедри \dotfill Ю.В.~Крак}}\hspace {0.7cm}{\hbox to 6.5 cm{ Екзаменатор \dotfill Т.О.~Карнаух}} \end{small}}
%begin
{{\begin {small}\par
{ \center  Київський національний університет імені Тараса Шевченка \par}
{\parbox {5 cm}{Освітня програма \hfill}{\parbox {7 cm}{Прикладна математика, Системний аналіз \hfill}}\parbox {6 cm}{\hfill Семестр 2}}\par
{\parbox {5 cm} {Навчальний предмет \hfill}\parbox {7 cm}{Програмування \hfill}\parbox {6cm}{\hfill}\par}\vspace{-7pt} \end{small}}{\center Екзаменаційний білет \No \textbf 
{ 241 }\par}}
{
\begin {large}
1. Що буде виведено за виконання?
code
try:
    res = 8j<"9.0"
    if isinstance(res, bool):
        print(f'bool;{res}')
    elif isinstance(res, int):
        print(f'int;{res}')
    elif isinstance(res, float):
        print(f'float;{res:.2f}')
    else:
        print('???')
except:
    print('ERR')
code

2. Що буде виведено за виконання?
code
a = 3
b = 8
c = 9

def g(a):
    global c
    a = 1
    b *= 5
    c = 3
    return a + b + c

a = 5
b = 7
c = 1

try:
    print(g(a), a, b, c, sep=';')
except:
    print('ERR', a, b, c, sep=';')
code

3. Що буде виведено за виконання?
code
try:
    t = 11
    while t>9:
        t += -3
    print(t, end=';')
except:
    print('ERR')
code

4. Що буде виведено за виконання?
code
try:
    for c in range(-10, 1, -2):
        if c>3:
            break
        print(c, end=';')
    else:
        print(c,end=';')
    print(c)
except:
    print('ERR')
code

5. Що буде виведено за виконання?\par (Дійсні значення, якщо наявні, в цьому завданні будуть виводитись у форматі з фіксованою крапкою та, можливо, з доволі великою кількістю десяткових розрядів. У відповіді для дійсних значень у ЦЬОМУ завданні достатньо вказати один десятковий розряд після крапки, відкинувши решту розрядів без округлення. Наприклад, замість 13.66666666666667 достатньо вказати 13.6, замість 25.23 достатньо вказати 25.2)
code
try:
    print(9, end=';')
    print(int('c0'), end=';')
    print(4, end=';')
except Exception:
    print(2, end=';')
except KeyboardInterrupt:
    print(8, end=';')
except:
    print('ERR', end=';')
else:
    print(6, end=';')
finally:
    print(5)
code

6. Що буде виведено за виконання?
code
def f(n):
    print(n%10, end='')
    if n>9:
        f(n//100)
 
f(987654)
code

7. Що буде виведено за виконання?
code
class A:
    d = 1

    def _f1(self): print(2, end=';')
    def f2(self): print(3, end=';')

    def main(self):
        try: print(self.d, end=';')
        except: print('ERR', end=';')

        try: self._f1()
        except: print('ERR', end=';')

        try: self.f2()
        except: print('ERR', end=';')


class B(A):
    d = 4

    def _f1(self): print(5, end=';')
    def f2(self): print(6, end=';')

try: B().main()
except: print('ERR', end=';')
code

8. Що буде виведено за виконання?
code
def f(arg, n):
    arg[24] = 0
    n *= 2

try:
    n = 6
    d = {16:7, 84:3, 44:3, 16:5}
    f(d, n)
    print(n, end=';')
    for x in d.keys():
        print(x, sep=';', end=';')
except:
    print('ERR')
code

9. Що буде виведено за виконання?
code
class A:
    a = 1
    b = 2

    def __init__(self):
        b = 3

class B(A):
    b = 4

    def __init__(self):
        super().__init__()
        self.c = 7

obj = B()
obj.c = obj.a + 10
B.c = 13

try: print(obj.c, end= ';')
except: print('ERR', end=';')

try: print(A.c, end= ';')
except: print('ERR', end=';')

try: print(B.c, end= ';')
except: print('ERR', end=';')
code

10. 
Для графа на рисунку з вершини 4 запущено алгоритм пошуку в глибину, що будує кістякове дерево. Списки суміжності вершин переглядаються алгоритмом за зростанням міток вершин.\par
\begin{center}
	\adjustimage{max size={0.9\linewidth}{0.9\paperheight}}
	{../10/Traversals/202.png}
\end{center}
\par Які з наведених ребер входять до складу отриманого кістякового дерева?\par
1) (1, 4)\par
2) (1, 3)\par
3) (4, 5)\par
4) (1, 2)\par
5) (0, 6)\par
6) (0, 5)\par
{\vskip 15pt} У формі вибрати номери відповідних ребер.\par
%answer:2 4 6
\end {large}}
%end
{\smallskip \begin {small} Затверджено на засіданні кафедри теоретичної кібернетики, протокол \No 12 від   19.05.2020 р.\par\smallskip\smallskip
{\parbox {6.5 cm}{Зав. кафедри \dotfill Ю.В.~Крак}}\hspace {0.7cm}{\hbox to 6.5 cm{ Екзаменатор \dotfill Т.О.~Карнаух}} \end{small}}
%begin
{{\begin {small}\par
{ \center  Київський національний університет імені Тараса Шевченка \par}
{\parbox {5 cm}{Освітня програма \hfill}{\parbox {7 cm}{Прикладна математика, Системний аналіз \hfill}}\parbox {6 cm}{\hfill Семестр 2}}\par
{\parbox {5 cm} {Навчальний предмет \hfill}\parbox {7 cm}{Програмування \hfill}\parbox {6cm}{\hfill}\par}\vspace{-7pt} \end{small}}{\center Екзаменаційний білет \No \textbf 
{ 242 }\par}}
{
\begin {large}
1. Що буде виведено за виконання?
code
try:
    res = True<=8.0j
    if isinstance(res, bool):
        print(f'bool;{res}')
    elif isinstance(res, int):
        print(f'int;{res}')
    elif isinstance(res, float):
        print(f'float;{res:.2f}')
    else:
        print('???')
except:
    print('ERR')
code

2. Що буде виведено за виконання?
code
a = 7
b = 8
c = 0

def g(a):
    a = 3
    b = 1
    c = 4
    return a + b + c

a = 4
b = 6
c = 2

try:
    print(g(b), a, b, c, sep=';')
except:
    print('ERR', a, b, c, sep=';')
code

3. Що буде виведено за виконання?
code
try:
    k = 6
    while k<12:
        k += 2
    print(k, end=';')
except:
    print('ERR')
code

4. Що буде виведено за виконання?
code
try:
    for f in range(-3, -9, 2):
        if f<=7:
            continue
        print(f, end=';')
        if f>=8:
            break
    else:
        print('end',end=';')
    print(f)
except:
    print('ERR')
code

5. Що буде виведено за виконання?\par (Дійсні значення, якщо наявні, в цьому завданні будуть виводитись у форматі з фіксованою крапкою та, можливо, з доволі великою кількістю десяткових розрядів. У відповіді для дійсних значень у ЦЬОМУ завданні достатньо вказати один десятковий розряд після крапки, відкинувши решту розрядів без округлення. Наприклад, замість 13.66666666666667 достатньо вказати 13.6, замість 25.23 достатньо вказати 25.2)
code
try:
    print(7, end=';')
    print(1j<4j, end=';')
    print(3, end=';')
except ValueError:
    print(5, end=';')
except TypeError:
    print(9, end=';')
except:
    print('ERR', end=';')
else:
    print(1, end=';')
finally:
    print(6)
code

6. Що буде виведено за виконання?
code
def f(n):
    if n>2:
        f(n//2)
    print(n%2, end='')
 
f(16)
code

7. Що буде виведено за виконання?
code
class A:
    a = 1

    def _h1(self): print(2, end=';')
    def __h2(self): print(3, end=';')

    def main(self):
        try: print(self.a, end=';')
        except: print('ERR', end=';')

        try: self._h1()
        except: print('ERR', end=';')

        try: self.__h2()
        except: print('ERR', end=';')


class B(A):
    a = 4

    def _h1(self): print(5, end=';')
    def __h2(self): print(6, end=';')

try: B().main()
except: print('ERR', end=';')
code

8. Що буде виведено за виконання?
code
def f(arg, n):
    arg[12] = 5
    n = 5
    arg = tuple(arg.keys())

try:
    n = 0
    d = {65:9, 49:2, 79:7, 22:6}
    f(d, n)
    print(n, end=';')
    for x in d.keys():
        print(x, sep=';', end=';')
except:
    print('ERR')
code

9. Що буде виведено за виконання?
code
class A:
    b = 1
    c = 2

    def __init__(self):
        b = 3

class B(A):
    b = 4

    def __init__(self):
        super().__init__()
        self.a = 7

obj = B()
obj.b = obj.b + 10
B.b = 13

try: print(obj.b, end= ';')
except: print('ERR', end=';')

try: print(A.b, end= ';')
except: print('ERR', end=';')

try: print(B.b, end= ';')
except: print('ERR', end=';')
code

10. 
Для графа на рисунку з вершини 2 запущено алгоритм пошуку в глибину, що будує кістякове дерево. Списки суміжності вершин переглядаються алгоритмом за зростанням міток вершин.\par
\begin{center}
	\adjustimage{max size={0.9\linewidth}{0.9\paperheight}}
	{../10/Traversals/266.png}
\end{center}
\par Які з наведених ребер входять до складу отриманого кістякового дерева?\par
1) (0, 6)\par
2) (0, 2)\par
3) (3, 6)\par
4) (5, 6)\par
5) (2, 6)\par
6) (2, 3)\par
{\vskip 15pt} У формі вибрати номери відповідних ребер.\par
%answer:2 3
\end {large}}
%end
{\smallskip \begin {small} Затверджено на засіданні кафедри теоретичної кібернетики, протокол \No 12 від   19.05.2020 р.\par\smallskip\smallskip
{\parbox {6.5 cm}{Зав. кафедри \dotfill Ю.В.~Крак}}\hspace {0.7cm}{\hbox to 6.5 cm{ Екзаменатор \dotfill Т.О.~Карнаух}} \end{small}}
%begin
{{\begin {small}\par
{ \center  Київський національний університет імені Тараса Шевченка \par}
{\parbox {5 cm}{Освітня програма \hfill}{\parbox {7 cm}{Прикладна математика, Системний аналіз \hfill}}\parbox {6 cm}{\hfill Семестр 2}}\par
{\parbox {5 cm} {Навчальний предмет \hfill}\parbox {7 cm}{Програмування \hfill}\parbox {6cm}{\hfill}\par}\vspace{-7pt} \end{small}}{\center Екзаменаційний білет \No \textbf 
{ 243 }\par}}
{
\begin {large}
1. Що буде виведено за виконання?
code
def f(n):
    print('f', end='')
    return n

try:
    print(f(-7) and f(8))
except:
    print('ERR')
code

2. Що буде виведено за виконання?
code
a = 0
b = 4
c = 6

def f(b):
    a -= 4
    b = 2
    c = 3
    return a + b + c

a = 2
b = 7
c = 8

try:
    print(f(b), a, b, c, sep=';')
except:
    print('ERR', a, b, c, sep=';')
code

3. Що буде виведено за виконання?
code
try:
    m = 11
    while m>=6:
        m += -2
    print(m, end=';')
except:
    print('ERR')
code

4. Що буде виведено за виконання?
code
try:
    for a in range(-1, -6, 3):
        if a>=6:
            continue
        print(a, end=';');a = 6
    else:
        print(a,end=';')
    print(a)
except:
    print('ERR')
code

5. Що буде виведено за виконання?\par (Дійсні значення, якщо наявні, в цьому завданні будуть виводитись у форматі з фіксованою крапкою та, можливо, з доволі великою кількістю десяткових розрядів. У відповіді для дійсних значень у ЦЬОМУ завданні достатньо вказати один десятковий розряд після крапки, відкинувши решту розрядів без округлення. Наприклад, замість 13.66666666666667 достатньо вказати 13.6, замість 25.23 достатньо вказати 25.2)
code
try:
    print(0, end=';')
    print(7//1, end=';')
    print(3, end=';')
except ZeroDivisionError:
    print(2, end=';')
except KeyboardInterrupt:
    print(4, end=';')
except:
    print('ERR', end=';')
else:
    print(8, end=';')
finally:
    print(7)
code

6. Що буде виведено за виконання?
code
def f(n):
    if n>9:
        return f(n//100)+1
    else:
        return 1

print(f(123456))
code

7. Що буде виведено за виконання?
code
class A:
    b = 1

    def g1(self): print(2, end=';')
    def _g2(self): print(3, end=';')

    def main(self):
        try: print(self.b, end=';')
        except: print('ERR', end=';')

        try: self.g1()
        except: print('ERR', end=';')

        try: self._g2()
        except: print('ERR', end=';')


class B(A):
    b = 4

    def g1(self): print(5, end=';')
    def _g2(self): print(6, end=';')

try: B().main()
except: print('ERR', end=';')
code

8. Що буде виведено за виконання?
code
def f(arg, n):
    n = 7
    tmp = arg
    tmp[3:] = 'bc'
    return len(tmp)>3

try:
    d = [10,11,12,13,14]
    n = 2
    f(d, n)
    print(n, end=';')
    for el in d:
        print(el, end=';')
except:
    print('ERR')
code

9. Що буде виведено за виконання?
code
class A:
    a = 1
    c = 2

    def __init__(self):
        a = 3

class B(A):
    c = 4

    def __init__(self):
        super().__init__()
        self.b = 7

obj = B()
obj.c = obj.a + 10
B.c = 13

try: print(obj.c, end= ';')
except: print('ERR', end=';')

try: print(A.c, end= ';')
except: print('ERR', end=';')

try: print(B.c, end= ';')
except: print('ERR', end=';')
code

10. 
Для графа на рисунку з вершини 6 запущено алгоритм пошуку в глибину, що будує кістякове дерево. Списки суміжності вершин переглядаються алгоритмом за зростанням міток вершин.\par
\begin{center}
	\adjustimage{max size={0.9\linewidth}{0.9\paperheight}}
	{../10/Traversals/71.png}
\end{center}
\par Які з наведених ребер входять до складу отриманого кістякового дерева?\par
1) (2, 3)\par
2) (0, 2)\par
3) (2, 4)\par
4) (4, 6)\par
5) (0, 4)\par
6) (0, 5)\par
{\vskip 15pt} У формі вибрати номери відповідних ребер.\par
%answer:1 2 3 6
\end {large}}
%end
{\smallskip \begin {small} Затверджено на засіданні кафедри теоретичної кібернетики, протокол \No 12 від   19.05.2020 р.\par\smallskip\smallskip
{\parbox {6.5 cm}{Зав. кафедри \dotfill Ю.В.~Крак}}\hspace {0.7cm}{\hbox to 6.5 cm{ Екзаменатор \dotfill Т.О.~Карнаух}} \end{small}}
%begin
{{\begin {small}\par
{ \center  Київський національний університет імені Тараса Шевченка \par}
{\parbox {5 cm}{Освітня програма \hfill}{\parbox {7 cm}{Прикладна математика, Системний аналіз \hfill}}\parbox {6 cm}{\hfill Семестр 2}}\par
{\parbox {5 cm} {Навчальний предмет \hfill}\parbox {7 cm}{Програмування \hfill}\parbox {6cm}{\hfill}\par}\vspace{-7pt} \end{small}}{\center Екзаменаційний білет \No \textbf 
{ 244 }\par}}
{
\begin {large}
1. Що буде виведено за виконання?
code
try:
    res = not True>=7.0 or 5>6.0 or 9<4
    if isinstance(res, bool):
        print(f'bool;{res}')
    elif isinstance(res, int):
        print(f'int;{res}')
    elif isinstance(res, float):
        print(f'float;{res:.2f}')
    else:
        print('???')
except:
    print('ERR')
code

2. Що буде виведено за виконання?
code
a = 6
b = 2
c = 3

def h(b):
    global c
    a += 4
    b = 2
    c = 5
    return a + b + c

a = 1
b = 0
c = 4

try:
    print(h(a), a, b, c, sep=';')
except:
    print('ERR', a, b, c, sep=';')
code

3. Що буде виведено за виконання?
code
try:
    j = 10
    while j>3:
        j += -3
    print(j, end=';')
except:
    print('ERR')
code

4. Що буде виведено за виконання?
code
try:
    for d in range(-7, -13, -3):
        if d<5:
            continue
        print(d, end=';')
    else:
        print(d,end=';')
    print(d)
except:
    print('ERR')
code

5. Що буде виведено за виконання?\par (Дійсні значення, якщо наявні, в цьому завданні будуть виводитись у форматі з фіксованою крапкою та, можливо, з доволі великою кількістю десяткових розрядів. У відповіді для дійсних значень у ЦЬОМУ завданні достатньо вказати один десятковий розряд після крапки, відкинувши решту розрядів без округлення. Наприклад, замість 13.66666666666667 достатньо вказати 13.6, замість 25.23 достатньо вказати 25.2)
code
try:
    print(5, end=';')
    print(int('9'), end=';')
    print(2, end=';')
except TypeError:
    print(1, end=';')
except ValueError:
    print(3, end=';')
except:
    print('ERR', end=';')
else:
    print(4, end=';')
finally:
    print(0)
code

6. Що буде виведено за виконання?
code
def f(n):
    if n>9:
        f(n//10)
    else:
        print(n%10, end='')

f(543216)
code

7. Що буде виведено за виконання?
code
class A:
    c = 1

    def __g1(self): print(2, end=';')
    def _g2(self): print(3, end=';')

    def main(self):
        try: print(self.c, end=';')
        except: print('ERR', end=';')

        try: self.__g1()
        except: print('ERR', end=';')

        try: self._g2()
        except: print('ERR', end=';')


class B(A):
    c = 4

    def __g1(self): print(5, end=';')
    def _g2(self): print(6, end=';')

try: B().main()
except: print('ERR', end=';')
code

8. Що буде виведено за виконання?
code
def f(arg, n):
    n = 3
    tmp = arg
    tmp[:-2] = 'ac'
    return len(tmp)>3

try:
    a = ['0','1','2','3','4']
    n = 1
    f(a, n)
    print(n, end=';')
    for el in a:
        print(el, end=';')
except:
    print('ERR')
code

9. Що буде виведено за виконання?
code
class A:
    c = 1
    a = 2

    def __init__(self):
        a = 3

class B(A):
    c = 4

    def __init__(self):
        super().__init__()
        self.b = 7

obj = B()
obj.b = obj.a + 10
B.b = 13

try: print(obj.b, end= ';')
except: print('ERR', end=';')

try: print(A.b, end= ';')
except: print('ERR', end=';')

try: print(B.b, end= ';')
except: print('ERR', end=';')
code

10. 
Для графа на рисунку з вершини 6 запущено алгоритм пошуку в глибину, що будує кістякове дерево. Списки суміжності вершин переглядаються алгоритмом за зростанням міток вершин.\par
\begin{center}
	\adjustimage{max size={0.9\linewidth}{0.9\paperheight}}
	{../10/Traversals/203.png}
\end{center}
\par Які з наведених ребер входять до складу отриманого кістякового дерева?\par
1) (0, 6)\par
2) (0, 2)\par
3) (1, 4)\par
4) (4, 5)\par
5) (2, 6)\par
6) (0, 3)\par
{\vskip 15pt} У формі вибрати номери відповідних ребер.\par
%answer:1 2 3 4 6
\end {large}}
%end
{\smallskip \begin {small} Затверджено на засіданні кафедри теоретичної кібернетики, протокол \No 12 від   19.05.2020 р.\par\smallskip\smallskip
{\parbox {6.5 cm}{Зав. кафедри \dotfill Ю.В.~Крак}}\hspace {0.7cm}{\hbox to 6.5 cm{ Екзаменатор \dotfill Т.О.~Карнаух}} \end{small}}
%begin
{{\begin {small}\par
{ \center  Київський національний університет імені Тараса Шевченка \par}
{\parbox {5 cm}{Освітня програма \hfill}{\parbox {7 cm}{Прикладна математика, Системний аналіз \hfill}}\parbox {6 cm}{\hfill Семестр 2}}\par
{\parbox {5 cm} {Навчальний предмет \hfill}\parbox {7 cm}{Програмування \hfill}\parbox {6cm}{\hfill}\par}\vspace{-7pt} \end{small}}{\center Екзаменаційний білет \No \textbf 
{ 245 }\par}}
{
\begin {large}
1. Що буде виведено за виконання?
code
try:
    res = -1**-4+1
    if isinstance(res, bool):
        print(f'bool;{res}')
    elif isinstance(res, int):
        print(f'int;{res}')
    elif isinstance(res, float):
        print(f'float;{res:.2f}')
    else:
        print('???')
except:
    print('ERR')
code

2. Що буде виведено за виконання?
code
a = 9
b = 5
c = 1

def f(a):
    a = 1
    b *= 4
    c = 2
    return a + b + c

a = 6
b = 2
c = 7

try:
    print(f(b), a, b, c, sep=';')
except:
    print('ERR', a, b, c, sep=';')
code

3. Що буде виведено за виконання?
code
try:
    n = 3
    while n<=9:
        n += 1
    print(n, end=';')
except:
    print('ERR')
code

4. Що буде виведено за виконання?
code
try:
    for b in range(7, -3, -1):
        if b<8:
            continue
        print(b, end=';')
    else:
        print(b,end=';')
    print(b)
except:
    print('ERR')
code

5. Що буде виведено за виконання?\par (Дійсні значення, якщо наявні, в цьому завданні будуть виводитись у форматі з фіксованою крапкою та, можливо, з доволі великою кількістю десяткових розрядів. У відповіді для дійсних значень у ЦЬОМУ завданні достатньо вказати один десятковий розряд після крапки, відкинувши решту розрядів без округлення. Наприклад, замість 13.66666666666667 достатньо вказати 13.6, замість 25.23 достатньо вказати 25.2)
code
try:
    print(8, end=';')
    print(6%0.0, end=';')
    print(7, end=';')
except ZeroDivisionError:
    print(6, end=';')
except Exception:
    print(5, end=';')
except:
    print('ERR', end=';')
else:
    print(1, end=';')
finally:
    print(2)
code

6. Що буде виведено за виконання?
code
def f(s):
    s1 = s
    for i in range(len(s)):
        s1[i] += 1
        return s1

s = [1, 2, 3]
f(s)
print(s)
code

7. Що буде виведено за виконання?
code
class A:
    a = 1

    def _f1(self): print(2, end=';')
    def _f2(self): print(3, end=';')

    def main(self):
        try: print(self.a, end=';')
        except: print('ERR', end=';')

        try: self._f1()
        except: print('ERR', end=';')

        try: self._f2()
        except: print('ERR', end=';')


class B(A):
    a = 4

    def _f1(self): print(5, end=';')
    def _f2(self): print(6, end=';')

try: B().main()
except: print('ERR', end=';')
code

8. Що буде виведено за виконання?
code
def f(arg, n):
    n = 0
    tmp = arg
    del tmp[:6]
    return len(tmp)>3

try:
    c = [10,11,12,13,14]
    n = 9
    f(c, n)
    print(n, end=';')
    for el in c:
        print(el, end=';')
except:
    print('ERR')
code

9. Що буде виведено за виконання?
code
class A:
    b = 1
    a = 2

    def __init__(self):
        b = 3

class B(A):
    b = 4

    def __init__(self):
        super().__init__()
        self.c = 7

obj = B()
obj.c = obj.c + 10
B.c = 13

try: print(obj.c, end= ';')
except: print('ERR', end=';')

try: print(A.c, end= ';')
except: print('ERR', end=';')

try: print(B.c, end= ';')
except: print('ERR', end=';')
code

10. 
Для графа на рисунку з вершини 4 запущено алгоритм пошуку в глибину, що будує кістякове дерево. Списки суміжності вершин переглядаються алгоритмом за зростанням міток вершин.\par
\begin{center}
	\adjustimage{max size={0.9\linewidth}{0.9\paperheight}}
	{../10/Traversals/162.png}
\end{center}
\par Які з наведених ребер входять до складу отриманого кістякового дерева?\par
1) (3, 6)\par
2) (1, 3)\par
3) (5, 6)\par
4) (1, 4)\par
5) (2, 4)\par
6) (0, 3)\par
{\vskip 15pt} У формі вибрати номери відповідних ребер.\par
%answer:2 3 6
\end {large}}
%end
{\smallskip \begin {small} Затверджено на засіданні кафедри теоретичної кібернетики, протокол \No 12 від   19.05.2020 р.\par\smallskip\smallskip
{\parbox {6.5 cm}{Зав. кафедри \dotfill Ю.В.~Крак}}\hspace {0.7cm}{\hbox to 6.5 cm{ Екзаменатор \dotfill Т.О.~Карнаух}} \end{small}}
%begin
{{\begin {small}\par
{ \center  Київський національний університет імені Тараса Шевченка \par}
{\parbox {5 cm}{Освітня програма \hfill}{\parbox {7 cm}{Прикладна математика, Системний аналіз \hfill}}\parbox {6 cm}{\hfill Семестр 2}}\par
{\parbox {5 cm} {Навчальний предмет \hfill}\parbox {7 cm}{Програмування \hfill}\parbox {6cm}{\hfill}\par}\vspace{-7pt} \end{small}}{\center Екзаменаційний білет \No \textbf 
{ 246 }\par}}
{
\begin {large}
1. Що буде виведено за виконання?
code
try:
    res = 4**2.0--2
    if isinstance(res, bool):
        print(f'bool;{res}')
    elif isinstance(res, int):
        print(f'int;{res}')
    elif isinstance(res, float):
        print(f'float;{res:.2f}')
    else:
        print('???')
except:
    print('ERR')
code

2. Що буде виведено за виконання?
code
a = 0
b = 8
c = 5

def g(a):
    global c
    a -= 1
    b = 3
    c = 5
    return a + b + c

a = 9
b = 3
c = 4

try:
    print(g(a), a, b, c, sep=';')
except:
    print('ERR', a, b, c, sep=';')
code

3. Що буде виведено за виконання?
code
try:
    k = 8
    while k>=7:
        k += -2
    print(k, end=';')
except:
    print('ERR')
code

4. Що буде виведено за виконання?
code
try:
    for e in range(-15, -1, -1):
        if e<=4:
            break
        print(e, end=';');e = 7
    else:
        print(e,end=';')
    print(e)
except:
    print('ERR')
code

5. Що буде виведено за виконання?\par (Дійсні значення, якщо наявні, в цьому завданні будуть виводитись у форматі з фіксованою крапкою та, можливо, з доволі великою кількістю десяткових розрядів. У відповіді для дійсних значень у ЦЬОМУ завданні достатньо вказати один десятковий розряд після крапки, відкинувши решту розрядів без округлення. Наприклад, замість 13.66666666666667 достатньо вказати 13.6, замість 25.23 достатньо вказати 25.2)
code
try:
    print(0, end=';')
    print(3/0, end=';')
    print(8, end=';')
except TypeError:
    print(9, end=';')
except ValueError:
    print(4, end=';')
except:
    print('ERR', end=';')
else:
    print(7, end=';')
finally:
    print(9)
code

6. Що буде виведено за виконання?
code
def f(s):
    for el in s:
        el += 1
    return s
 
s = [1, 2, 3]
s = f(s)
print(s)
code

7. Що буде виведено за виконання?
code
class A:
    b = 1

    def h1(self): print(2, end=';')
    def _h2(self): print(3, end=';')

    def main(self):
        try: print(self.b, end=';')
        except: print('ERR', end=';')

        try: self.h1()
        except: print('ERR', end=';')

        try: self._h2()
        except: print('ERR', end=';')


class B(A):
    b = 4

    def h1(self): print(5, end=';')
    def _h2(self): print(6, end=';')

try: B().main()
except: print('ERR', end=';')
code

8. Що буде виведено за виконання?
code
def f(arg, n):
    arg[42] = 7
    n = 9
    arg = list(arg.values())

try:
    n = 3
    s = {23:9, 42:8, 56:5, 42:4}
    f(s, n)
    print(n, end=';')
    for x in s.values():
        print(x, sep=';', end=';')
except:
    print('ERR')
code

9. Що буде виведено за виконання?
code
class A:
    c = 1
    b = 2

    def __init__(self):
        b = 3

class B(A):
    b = 4

    def __init__(self):
        super().__init__()
        self.a = 7

obj = B()
obj.b = obj.a + 10
B.b = 13

try: print(obj.b, end= ';')
except: print('ERR', end=';')

try: print(A.b, end= ';')
except: print('ERR', end=';')

try: print(B.b, end= ';')
except: print('ERR', end=';')
code

10. 
Для графа на рисунку з вершини 4 запущено алгоритм пошуку в глибину, що будує кістякове дерево. Списки суміжності вершин переглядаються алгоритмом за зростанням міток вершин.\par
\begin{center}
	\adjustimage{max size={0.9\linewidth}{0.9\paperheight}}
	{../10/Traversals/112.png}
\end{center}
\par Які з наведених ребер входять до складу отриманого кістякового дерева?\par
1) (0, 4)\par
2) (0, 5)\par
3) (2, 3)\par
4) (0, 1)\par
5) (2, 5)\par
6) (1, 3)\par
{\vskip 15pt} У формі вибрати номери відповідних ребер.\par
%answer:1 3 4 5 6
\end {large}}
%end
{\smallskip \begin {small} Затверджено на засіданні кафедри теоретичної кібернетики, протокол \No 12 від   19.05.2020 р.\par\smallskip\smallskip
{\parbox {6.5 cm}{Зав. кафедри \dotfill Ю.В.~Крак}}\hspace {0.7cm}{\hbox to 6.5 cm{ Екзаменатор \dotfill Т.О.~Карнаух}} \end{small}}
%begin
{{\begin {small}\par
{ \center  Київський національний університет імені Тараса Шевченка \par}
{\parbox {5 cm}{Освітня програма \hfill}{\parbox {7 cm}{Прикладна математика, Системний аналіз \hfill}}\parbox {6 cm}{\hfill Семестр 2}}\par
{\parbox {5 cm} {Навчальний предмет \hfill}\parbox {7 cm}{Програмування \hfill}\parbox {6cm}{\hfill}\par}\vspace{-7pt} \end{small}}{\center Екзаменаційний білет \No \textbf 
{ 247 }\par}}
{
\begin {large}
1. Що буде виведено за виконання?
code
def f(n):
    print('f', end='')
    return n<2

try:
    print(f(7) or f(6))
except:
    print('ERR')
code

2. Що буде виведено за виконання?
code
a = 5
b = 3
c = 1

def f(b):
    global c
    a = 5
    b = 2
    c = 1
    return a + b + c

a = 9
b = 8
c = 7

try:
    print(f(b), a, b, c, sep=';')
except:
    print('ERR', a, b, c, sep=';')
code

3. Що буде виведено за виконання?
code
try:
    t = 0
    while t<6:
        t += 1
    print(t, end=';')
except:
    print('ERR')
code

4. Що буде виведено за виконання?
code
try:
    for b in range(-13, 12, -3):
        if b>=3:
            break
        print(b, end=';');b = 5
    else:
        print('end',end=';')
    print(b)
except:
    print('ERR')
code

5. Що буде виведено за виконання?\par (Дійсні значення, якщо наявні, в цьому завданні будуть виводитись у форматі з фіксованою крапкою та, можливо, з доволі великою кількістю десяткових розрядів. У відповіді для дійсних значень у ЦЬОМУ завданні достатньо вказати один десятковий розряд після крапки, відкинувши решту розрядів без округлення. Наприклад, замість 13.66666666666667 достатньо вказати 13.6, замість 25.23 достатньо вказати 25.2)
code
try:
    print(5, end=';')
    print(8j>2j, end=';')
    print(6, end=';')
except ZeroDivisionError:
    print(2, end=';')
except KeyboardInterrupt:
    print(3, end=';')
except:
    print('ERR', end=';')
else:
    print(4, end=';')
finally:
    print(0)
code

6. Що буде виведено за виконання?
code
def f(n):
    if n>9:
        return f(n//10)+1
    else:
        return 1

print(f(123456))
code

7. Що буде виведено за виконання?
code
class A:
    c = 1

    def _h1(self): print(2, end=';')
    def __h2(self): print(3, end=';')

    def main(self):
        try: print(self.c, end=';')
        except: print('ERR', end=';')

        try: self._h1()
        except: print('ERR', end=';')

        try: self.__h2()
        except: print('ERR', end=';')


class B(A):
    c = 4

    def _h1(self): print(5, end=';')
    def __h2(self): print(6, end=';')

try: B().main()
except: print('ERR', end=';')
code

8. Що буде виведено за виконання?
code
def f(arg, n):
    n = 4
    tmp = arg
    del tmp[-5:-3]
    return len(tmp)>3

try:
    e = [10,11,12,13,14]
    n = 8
    f(e, n)
    print(n, end=';')
    for el in e:
        print(el, end=';')
except:
    print('ERR')
code

9. Що буде виведено за виконання?
code
class A:
    b = 1
    a = 2

    def __init__(self):
        b = 3

class B(A):
    b = 4

    def __init__(self):
        super().__init__()
        self.c = 7

obj = B()
obj.c = obj.b + 10
B.c = 13

try: print(obj.c, end= ';')
except: print('ERR', end=';')

try: print(A.c, end= ';')
except: print('ERR', end=';')

try: print(B.c, end= ';')
except: print('ERR', end=';')
code

10. 
Для графа на рисунку з вершини 2 запущено алгоритм пошуку в глибину, що будує кістякове дерево. Списки суміжності вершин переглядаються алгоритмом за зростанням міток вершин.\par
\begin{center}
	\adjustimage{max size={0.9\linewidth}{0.9\paperheight}}
	{../10/Traversals/45.png}
\end{center}
\par Які з наведених ребер входять до складу отриманого кістякового дерева?\par
1) (2, 5)\par
2) (0, 5)\par
3) (0, 6)\par
4) (3, 5)\par
5) (2, 3)\par
6) (0, 3)\par
{\vskip 15pt} У формі вибрати номери відповідних ребер.\par
%answer:3 4 6
\end {large}}
%end
{\smallskip \begin {small} Затверджено на засіданні кафедри теоретичної кібернетики, протокол \No 12 від   19.05.2020 р.\par\smallskip\smallskip
{\parbox {6.5 cm}{Зав. кафедри \dotfill Ю.В.~Крак}}\hspace {0.7cm}{\hbox to 6.5 cm{ Екзаменатор \dotfill Т.О.~Карнаух}} \end{small}}
%begin
{{\begin {small}\par
{ \center  Київський національний університет імені Тараса Шевченка \par}
{\parbox {5 cm}{Освітня програма \hfill}{\parbox {7 cm}{Прикладна математика, Системний аналіз \hfill}}\parbox {6 cm}{\hfill Семестр 2}}\par
{\parbox {5 cm} {Навчальний предмет \hfill}\parbox {7 cm}{Програмування \hfill}\parbox {6cm}{\hfill}\par}\vspace{-7pt} \end{small}}{\center Екзаменаційний білет \No \textbf 
{ 248 }\par}}
{
\begin {large}
1. Що буде виведено за виконання?
code
try:
    res = 3**1.0-2
    if isinstance(res, bool):
        print(f'bool;{res}')
    elif isinstance(res, int):
        print(f'int;{res}')
    elif isinstance(res, float):
        print(f'float;{res:.2f}')
    else:
        print('???')
except:
    print('ERR')
code

2. Що буде виведено за виконання?
code
a = 5
b = 6
c = 4

def h(b):
    a = 5
    b = 4
    c = 3
    return a + b + c

a = 9
b = 2
c = 3

try:
    print(h(b), a, b, c, sep=';')
except:
    print('ERR', a, b, c, sep=';')
code

3. Що буде виведено за виконання?
code
try:
    i = 7
    while i>=1:
        i += -2
    print(i, end=';')
except:
    print('ERR')
code

4. Що буде виведено за виконання?
code
try:
    for e in range(6, 0, 3):
        if e>7:
            continue
        print(e, end=';')
    else:
        print(e,end=';')
    print(e)
except:
    print('ERR')
code

5. Що буде виведено за виконання?\par (Дійсні значення, якщо наявні, в цьому завданні будуть виводитись у форматі з фіксованою крапкою та, можливо, з доволі великою кількістю десяткових розрядів. У відповіді для дійсних значень у ЦЬОМУ завданні достатньо вказати один десятковий розряд після крапки, відкинувши решту розрядів без округлення. Наприклад, замість 13.66666666666667 достатньо вказати 13.6, замість 25.23 достатньо вказати 25.2)
code
try:
    print(1, end=';')
    print(1%3, end=';')
    print(6, end=';')
except Exception:
    print(8, end=';')
except Exception:
    print(2, end=';')
except:
    print('ERR', end=';')
else:
    print(9, end=';')
finally:
    print(5)
code

6. Що буде виведено за виконання?
code
def f(n):
    if n>9:
        f(n//100)
    else:
        print(n%10,end='')

f(123456)
code

7. Що буде виведено за виконання?
code
class A:
    d = 1

    def _f1(self): print(2, end=';')
    def _f2(self): print(3, end=';')

    def main(self):
        try: print(self.d, end=';')
        except: print('ERR', end=';')

        try: self._f1()
        except: print('ERR', end=';')

        try: self._f2()
        except: print('ERR', end=';')


class B(A):
    d = 4

    def _f1(self): print(5, end=';')
    def _f2(self): print(6, end=';')

try: B().main()
except: print('ERR', end=';')
code

8. Що буде виведено за виконання?
code
def f(arg, n):
    arg[64] = 2
    n = 7
    arg = set(arg.items())

try:
    n = 0
    t = {46:5, 53:7, 84:0, 84:0}
    f(t, n)
    print(n, end=';')
    for x in t.items():
        print(x, sep=';', end=';')
except:
    print('ERR')
code

9. Що буде виведено за виконання?
code
class A:
    a = 1
    b = 2

    def __init__(self):
        b = 3

class B(A):
    b = 4

    def __init__(self):
        super().__init__()
        self.c = 7

obj = B()
obj.a = obj.b + 10
B.a = 13

try: print(obj.a, end= ';')
except: print('ERR', end=';')

try: print(A.a, end= ';')
except: print('ERR', end=';')

try: print(B.a, end= ';')
except: print('ERR', end=';')
code

10. 
Для графа на рисунку з вершини 5 запущено алгоритм пошуку в глибину, що будує кістякове дерево. Списки суміжності вершин переглядаються алгоритмом за зростанням міток вершин.\par
\begin{center}
	\adjustimage{max size={0.9\linewidth}{0.9\paperheight}}
	{../10/Traversals/90.png}
\end{center}
\par Які з наведених ребер входять до складу отриманого кістякового дерева?\par
1) (2, 5)\par
2) (0, 1)\par
3) (0, 5)\par
4) (3, 4)\par
5) (1, 4)\par
6) (3, 5)\par
{\vskip 15pt} У формі вибрати номери відповідних ребер.\par
%answer:2 3 5
\end {large}}
%end
{\smallskip \begin {small} Затверджено на засіданні кафедри теоретичної кібернетики, протокол \No 12 від   19.05.2020 р.\par\smallskip\smallskip
{\parbox {6.5 cm}{Зав. кафедри \dotfill Ю.В.~Крак}}\hspace {0.7cm}{\hbox to 6.5 cm{ Екзаменатор \dotfill Т.О.~Карнаух}} \end{small}}
%begin
{{\begin {small}\par
{ \center  Київський національний університет імені Тараса Шевченка \par}
{\parbox {5 cm}{Освітня програма \hfill}{\parbox {7 cm}{Прикладна математика, Системний аналіз \hfill}}\parbox {6 cm}{\hfill Семестр 2}}\par
{\parbox {5 cm} {Навчальний предмет \hfill}\parbox {7 cm}{Програмування \hfill}\parbox {6cm}{\hfill}\par}\vspace{-7pt} \end{small}}{\center Екзаменаційний білет \No \textbf 
{ 249 }\par}}
{
\begin {large}
1. Що буде виведено за виконання?
code
try:
    res = 3==7.0 and not 7<6.0 and 6>=True
    if isinstance(res, bool):
        print(f'bool;{res}')
    elif isinstance(res, int):
        print(f'int;{res}')
    elif isinstance(res, float):
        print(f'float;{res:.2f}')
    else:
        print('???')
except:
    print('ERR')
code

2. Що буде виведено за виконання?
code
a = 9
b = 7
c = 1

def h(b):
    a += 1
    b = 4
    c = 5
    return a + b + c

a = 6
b = 2
c = 5

try:
    print(h(b), a, b, c, sep=';')
except:
    print('ERR', a, b, c, sep=';')
code

3. Що буде виведено за виконання?
code
try:
    j = 8
    while j<=13:
        j += 2
    print(j, end=';')
except:
    print('ERR')
code

4. Що буде виведено за виконання?
code
try:
    for f in range(6, 2, -2):
        if f>8:
            continue
        print(f, end=';')
    else:
        print(f,end=';')
    print(f)
except:
    print('ERR')
code

5. Що буде виведено за виконання?\par (Дійсні значення, якщо наявні, в цьому завданні будуть виводитись у форматі з фіксованою крапкою та, можливо, з доволі великою кількістю десяткових розрядів. У відповіді для дійсних значень у ЦЬОМУ завданні достатньо вказати один десятковий розряд після крапки, відкинувши решту розрядів без округлення. Наприклад, замість 13.66666666666667 достатньо вказати 13.6, замість 25.23 достатньо вказати 25.2)
code
try:
    print(0, end=';')
    print(int('a4'), end=';')
    print(7, end=';')
except ValueError:
    print(3, end=';')
except ZeroDivisionError:
    print(0, end=';')
except:
    print('ERR', end=';')
else:
    print(3, end=';')
finally:
    print(4)
code

6. Що буде виведено за виконання?
code
def f(s):
    s1 = s
    for i in range(len(s)):
        s1[i] += 1
    return s1

s = [1, 2, 3]
f(s)
print(s)
code

7. Що буде виведено за виконання?
code
class A:
    a = 1

    def __g1(self): print(2, end=';')
    def _g2(self): print(3, end=';')

    def main(self):
        try: print(self.a, end=';')
        except: print('ERR', end=';')

        try: self.__g1()
        except: print('ERR', end=';')

        try: self._g2()
        except: print('ERR', end=';')


class B(A):
    a = 4

    def __g1(self): print(5, end=';')
    def _g2(self): print(6, end=';')

try: B().main()
except: print('ERR', end=';')
code

8. Що буде виведено за виконання?
code
def f(arg, n):
    arg[60] = 4
    n -= -2

try:
    n = 7
    s = {13:5, 85:4, 53:9, 53:1}
    f(s, n)
    print(n, end=';')
    for x in s :
        print(x, sep=';', end=';')
except:
    print('ERR')
code

9. Що буде виведено за виконання?
code
class A:
    c = 1
    a = 2

    def __init__(self):
        a = 3

class B(A):
    a = 4

    def __init__(self):
        super().__init__()
        self.b = 7

obj = B()
obj.c = obj.a + 10
B.c = 13

try: print(obj.c, end= ';')
except: print('ERR', end=';')

try: print(A.c, end= ';')
except: print('ERR', end=';')

try: print(B.c, end= ';')
except: print('ERR', end=';')
code

10. 
Для графа на рисунку з вершини 5 запущено алгоритм пошуку в глибину, що будує кістякове дерево. Списки суміжності вершин переглядаються алгоритмом за зростанням міток вершин.\par
\begin{center}
	\adjustimage{max size={0.9\linewidth}{0.9\paperheight}}
	{../10/Traversals/2.png}
\end{center}
\par Які з наведених ребер входять до складу отриманого кістякового дерева?\par
1) (3, 5)\par
2) (1, 2)\par
3) (0, 3)\par
4) (0, 2)\par
5) (1, 3)\par
6) (2, 6)\par
{\vskip 15pt} У формі вибрати номери відповідних ребер.\par
%answer:3 4 5
\end {large}}
%end
{\smallskip \begin {small} Затверджено на засіданні кафедри теоретичної кібернетики, протокол \No 12 від   19.05.2020 р.\par\smallskip\smallskip
{\parbox {6.5 cm}{Зав. кафедри \dotfill Ю.В.~Крак}}\hspace {0.7cm}{\hbox to 6.5 cm{ Екзаменатор \dotfill Т.О.~Карнаух}} \end{small}}
%begin
{{\begin {small}\par
{ \center  Київський національний університет імені Тараса Шевченка \par}
{\parbox {5 cm}{Освітня програма \hfill}{\parbox {7 cm}{Прикладна математика, Системний аналіз \hfill}}\parbox {6 cm}{\hfill Семестр 2}}\par
{\parbox {5 cm} {Навчальний предмет \hfill}\parbox {7 cm}{Програмування \hfill}\parbox {6cm}{\hfill}\par}\vspace{-7pt} \end{small}}{\center Екзаменаційний білет \No \textbf 
{ 250 }\par}}
{
\begin {large}
1. Що буде виведено за виконання?
code
try:
    res = 2!=3.0 or not 4<=4.0 or False>2
    if isinstance(res, bool):
        print(f'bool;{res}')
    elif isinstance(res, int):
        print(f'int;{res}')
    elif isinstance(res, float):
        print(f'float;{res:.2f}')
    else:
        print('???')
except:
    print('ERR')
code

2. Що буде виведено за виконання?
code
a = 3
b = 0
c = 4

def h(b):
    a = 2
    b += 3
    c = 5
    return a + b + c

a = 8
b = 0
c = 8

try:
    print(h(a), a, b, c, sep=';')
except:
    print('ERR', a, b, c, sep=';')
code

3. Що буде виведено за виконання?
code
try:
    n = 14
    while n>2:
        n += -3
    print(n, end=';')
except:
    print('ERR')
code

4. Що буде виведено за виконання?
code
try:
    for c in range(7, -14, 2):
        if c<6:
            continue
        print(c, end=';')
    else:
        print(c,end=';')
    print(c)
except:
    print('ERR')
code

5. Що буде виведено за виконання?\par (Дійсні значення, якщо наявні, в цьому завданні будуть виводитись у форматі з фіксованою крапкою та, можливо, з доволі великою кількістю десяткових розрядів. У відповіді для дійсних значень у ЦЬОМУ завданні достатньо вказати один десятковий розряд після крапки, відкинувши решту розрядів без округлення. Наприклад, замість 13.66666666666667 достатньо вказати 13.6, замість 25.23 достатньо вказати 25.2)
code
try:
    print(1, end=';')
    print(int('9'), end=';')
    print(8, end=';')
except KeyboardInterrupt:
    print(7, end=';')
except TypeError:
    print(2, end=';')
except:
    print('ERR', end=';')
else:
    print(6, end=';')
finally:
    print(5)
code

6. Що буде виведено за виконання?
code
def f(n):
    if n<=9:
        print(n%10,end='')
    else:
        f(n//10)

f(543216)
code

7. Що буде виведено за виконання?
code
class A:
    b = 1

    def _f1(self): print(2, end=';')
    def f2(self): print(3, end=';')

    def main(self):
        try: print(self.b, end=';')
        except: print('ERR', end=';')

        try: self._f1()
        except: print('ERR', end=';')

        try: self.f2()
        except: print('ERR', end=';')


class B(A):
    b = 4

    def _f1(self): print(5, end=';')
    def f2(self): print(6, end=';')

try: B().main()
except: print('ERR', end=';')
code

8. Що буде виведено за виконання?
code
def f(arg, n):
    n = 5
    tmp = arg
    tmp[-1:5] = [4,5]
    return len(tmp)>3

try:
    h = ['0','1','2','3','4']
    n = 1
    f(h, n)
    print(n, end=';')
    for el in h:
        print(el, end=';')
except:
    print('ERR')
code

9. Що буде виведено за виконання?
code
class A:
    a = 1
    c = 2

    def __init__(self):
        a = 3

class B(A):
    a = 4

    def __init__(self):
        super().__init__()
        self.b = 7

obj = B()
obj.c = obj.a + 10
B.c = 13

try: print(obj.c, end= ';')
except: print('ERR', end=';')

try: print(A.c, end= ';')
except: print('ERR', end=';')

try: print(B.c, end= ';')
except: print('ERR', end=';')
code

10. 
Для графа на рисунку з вершини 1 запущено алгоритм пошуку в глибину, що будує кістякове дерево. Списки суміжності вершин переглядаються алгоритмом за зростанням міток вершин.\par
\begin{center}
	\adjustimage{max size={0.9\linewidth}{0.9\paperheight}}
	{../10/Traversals/95.png}
\end{center}
\par Які з наведених ребер входять до складу отриманого кістякового дерева?\par
1) (2, 5)\par
2) (0, 6)\par
3) (2, 4)\par
4) (3, 5)\par
5) (4, 5)\par
6) (2, 6)\par
{\vskip 15pt} У формі вибрати номери відповідних ребер.\par
%answer:1 2 3
\end {large}}
%end
{\smallskip \begin {small} Затверджено на засіданні кафедри теоретичної кібернетики, протокол \No 12 від   19.05.2020 р.\par\smallskip\smallskip
{\parbox {6.5 cm}{Зав. кафедри \dotfill Ю.В.~Крак}}\hspace {0.7cm}{\hbox to 6.5 cm{ Екзаменатор \dotfill Т.О.~Карнаух}} \end{small}}
%begin
{{\begin {small}\par
{ \center  Київський національний університет імені Тараса Шевченка \par}
{\parbox {5 cm}{Освітня програма \hfill}{\parbox {7 cm}{Прикладна математика, Системний аналіз \hfill}}\parbox {6 cm}{\hfill Семестр 2}}\par
{\parbox {5 cm} {Навчальний предмет \hfill}\parbox {7 cm}{Програмування \hfill}\parbox {6cm}{\hfill}\par}\vspace{-7pt} \end{small}}{\center Екзаменаційний білет \No \textbf 
{ 251 }\par}}
{
\begin {large}
1. Що буде виведено за виконання?
code
try:
    res = -2**False+-2
    if isinstance(res, bool):
        print(f'bool;{res}')
    elif isinstance(res, int):
        print(f'int;{res}')
    elif isinstance(res, float):
        print(f'float;{res:.2f}')
    else:
        print('???')
except:
    print('ERR')
code

2. Що буде виведено за виконання?
code
a = 7
b = 9
c = 4

def g(a):
    global c
    a = 4
    b *= 2
    c = 1
    return a + b + c

a = 2
b = 5
c = 3

try:
    print(g(a), a, b, c, sep=';')
except:
    print('ERR', a, b, c, sep=';')
code

3. Що буде виведено за виконання?
code
try:
    t = 13
    while t>5:
        t += -3
    print(t, end=';')
except:
    print('ERR')
code

4. Що буде виведено за виконання?
code
try:
    for d in range(13, 10, -3):
        if d<=4:
            continue
        print(d, end=';')
    else:
        print(d,end=';')
    print(d)
except:
    print('ERR')
code

5. Що буде виведено за виконання?\par (Дійсні значення, якщо наявні, в цьому завданні будуть виводитись у форматі з фіксованою крапкою та, можливо, з доволі великою кількістю десяткових розрядів. У відповіді для дійсних значень у ЦЬОМУ завданні достатньо вказати один десятковий розряд після крапки, відкинувши решту розрядів без округлення. Наприклад, замість 13.66666666666667 достатньо вказати 13.6, замість 25.23 достатньо вказати 25.2)
code
try:
    print(4, end=';')
    print(6j==0j, end=';')
    print(2, end=';')
except ZeroDivisionError:
    print(3, end=';')
except ValueError:
    print(9, end=';')
except:
    print('ERR', end=';')
else:
    print(1, end=';')
finally:
    print(5)
code

6. Що буде виведено за виконання?
code
def f(n):
    if n>2:
        f(n//2)
    else:
        print(n%2, end='')

f(16)
code

7. Що буде виведено за виконання?
code
class A:
    c = 1

    def _h1(self): print(2, end=';')
    def h2(self): print(3, end=';')

    def main(self):
        try: print(self.c, end=';')
        except: print('ERR', end=';')

        try: self._h1()
        except: print('ERR', end=';')

        try: self.h2()
        except: print('ERR', end=';')


class B(A):
    c = 4

    def _h1(self): print(5, end=';')
    def h2(self): print(6, end=';')

try: B().main()
except: print('ERR', end=';')
code

8. Що буде виведено за виконання?
code
def f(arg, n):
    arg[75] = 7
    n += 3

try:
    n = 4
    t = {73:4, 67:4, 37:7, 85:9}
    f(t, n)
    print(n, end=';')
    for x in t.keys():
        print(x, sep=';', end=';')
except:
    print('ERR')
code

9. Що буде виведено за виконання?
code
class A:
    b = 1
    c = 2

    def __init__(self):
        b = 3

class B(A):
    b = 4

    def __init__(self):
        super().__init__()
        self.a = 7

obj = B()
obj.b = obj.a + 10
B.b = 13

try: print(obj.b, end= ';')
except: print('ERR', end=';')

try: print(A.b, end= ';')
except: print('ERR', end=';')

try: print(B.b, end= ';')
except: print('ERR', end=';')
code

10. 
Для графа на рисунку з вершини 0 запущено алгоритм пошуку в глибину, що будує кістякове дерево. Списки суміжності вершин переглядаються алгоритмом за зростанням міток вершин.\par
\begin{center}
	\adjustimage{max size={0.9\linewidth}{0.9\paperheight}}
	{../10/Traversals/168.png}
\end{center}
\par Які з наведених ребер входять до складу отриманого кістякового дерева?\par
1) (0, 3)\par
2) (2, 3)\par
3) (1, 2)\par
4) (0, 1)\par
5) (1, 5)\par
6) (0, 2)\par
{\vskip 15pt} У формі вибрати номери відповідних ребер.\par
%answer:2 3 4
\end {large}}
%end
{\smallskip \begin {small} Затверджено на засіданні кафедри теоретичної кібернетики, протокол \No 12 від   19.05.2020 р.\par\smallskip\smallskip
{\parbox {6.5 cm}{Зав. кафедри \dotfill Ю.В.~Крак}}\hspace {0.7cm}{\hbox to 6.5 cm{ Екзаменатор \dotfill Т.О.~Карнаух}} \end{small}}
%begin
{{\begin {small}\par
{ \center  Київський національний університет імені Тараса Шевченка \par}
{\parbox {5 cm}{Освітня програма \hfill}{\parbox {7 cm}{Прикладна математика, Системний аналіз \hfill}}\parbox {6 cm}{\hfill Семестр 2}}\par
{\parbox {5 cm} {Навчальний предмет \hfill}\parbox {7 cm}{Програмування \hfill}\parbox {6cm}{\hfill}\par}\vspace{-7pt} \end{small}}{\center Екзаменаційний білет \No \textbf 
{ 252 }\par}}
{
\begin {large}
1. Що буде виведено за виконання?
code
try:
    res = 5.0>=False or not 6<2.0 and 7==8
    if isinstance(res, bool):
        print(f'bool;{res}')
    elif isinstance(res, int):
        print(f'int;{res}')
    elif isinstance(res, float):
        print(f'float;{res:.2f}')
    else:
        print('???')
except:
    print('ERR')
code

2. Що буде виведено за виконання?
code
a = 1
b = 0
c = 2

def h(a):
    global c
    a = 2
    b = 2
    c = 5
    return a + b + c

a = 0
b = 1
c = 7

try:
    print(h(a), a, b, c, sep=';')
except:
    print('ERR', a, b, c, sep=';')
code

3. Що буде виведено за виконання?
code
try:
    m = 5
    while m<8:
        m += 2
    print(m, end=';')
except:
    print('ERR')
code

4. Що буде виведено за виконання?
code
try:
    for a in range(12, 11, -1):
        if a>=5:
            break
        print(a, end=';');a = 2
    else:
        print(a,end=';')
    print(a)
except:
    print('ERR')
code

5. Що буде виведено за виконання?\par (Дійсні значення, якщо наявні, в цьому завданні будуть виводитись у форматі з фіксованою крапкою та, можливо, з доволі великою кількістю десяткових розрядів. У відповіді для дійсних значень у ЦЬОМУ завданні достатньо вказати один десятковий розряд після крапки, відкинувши решту розрядів без округлення. Наприклад, замість 13.66666666666667 достатньо вказати 13.6, замість 25.23 достатньо вказати 25.2)
code
try:
    print(7, end=';')
    print(0/2, end=';')
    print(8, end=';')
except Exception:
    print(2, end=';')
except KeyboardInterrupt:
    print(9, end=';')
except:
    print('ERR', end=';')
else:
    print(0, end=';')
finally:
    print(4)
code

6. Що буде виведено за виконання?
code
def f(n):
    print(n%10,end='')
    if n>9:
        f(n//100)
 
f(123456)
code

7. Що буде виведено за виконання?
code
class A:
    d = 1

    def __g1(self): print(2, end=';')
    def _g2(self): print(3, end=';')

    def main(self):
        try: print(self.d, end=';')
        except: print('ERR', end=';')

        try: self.__g1()
        except: print('ERR', end=';')

        try: self._g2()
        except: print('ERR', end=';')


class B(A):
    d = 4

    def __g1(self): print(5, end=';')
    def _g2(self): print(6, end=';')

try: B().main()
except: print('ERR', end=';')
code

8. Що буде виведено за виконання?
code
def f(arg, n):
    arg[61] = 2
    n -= 2

try:
    n = 3
    d = {33:9, 34:0, 24:0}
    f(d, n)
    print(n, end=';')
    for x in d.values():
        print(x, sep=';', end=';')
except:
    print('ERR')
code

9. Що буде виведено за виконання?
code
class A:
    c = 1
    b = 2

    def __init__(self):
        b = 3

class B(A):
    b = 4

    def __init__(self):
        super().__init__()
        self.a = 7

obj = B()
obj.c = obj.b + 10
B.c = 13

try: print(obj.c, end= ';')
except: print('ERR', end=';')

try: print(A.c, end= ';')
except: print('ERR', end=';')

try: print(B.c, end= ';')
except: print('ERR', end=';')
code

10. 
Для графа на рисунку з вершини 0 запущено алгоритм пошуку в глибину, що будує кістякове дерево. Списки суміжності вершин переглядаються алгоритмом за зростанням міток вершин.\par
\begin{center}
	\adjustimage{max size={0.9\linewidth}{0.9\paperheight}}
	{../10/Traversals/295.png}
\end{center}
\par Які з наведених ребер входять до складу отриманого кістякового дерева?\par
1) (1, 2)\par
2) (3, 4)\par
3) (0, 5)\par
4) (0, 3)\par
5) (4, 5)\par
6) (3, 6)\par
{\vskip 15pt} У формі вибрати номери відповідних ребер.\par
%answer:1 2 5 6
\end {large}}
%end
{\smallskip \begin {small} Затверджено на засіданні кафедри теоретичної кібернетики, протокол \No 12 від   19.05.2020 р.\par\smallskip\smallskip
{\parbox {6.5 cm}{Зав. кафедри \dotfill Ю.В.~Крак}}\hspace {0.7cm}{\hbox to 6.5 cm{ Екзаменатор \dotfill Т.О.~Карнаух}} \end{small}}
%begin
{{\begin {small}\par
{ \center  Київський національний університет імені Тараса Шевченка \par}
{\parbox {5 cm}{Освітня програма \hfill}{\parbox {7 cm}{Прикладна математика, Системний аналіз \hfill}}\parbox {6 cm}{\hfill Семестр 2}}\par
{\parbox {5 cm} {Навчальний предмет \hfill}\parbox {7 cm}{Програмування \hfill}\parbox {6cm}{\hfill}\par}\vspace{-7pt} \end{small}}{\center Екзаменаційний білет \No \textbf 
{ 253 }\par}}
{
\begin {large}
1. Що буде виведено за виконання?
code
def f(n):
    print('f', end='')
    return n>=0

try:
    print(f(3) or f(7))
except:
    print('ERR')
code

2. Що буде виведено за виконання?
code
a = 5
b = 4
c = 8

def f(a):
    a = 4
    b = 3
    c = 1
    return a + b + c

a = 3
b = 6
c = 9

try:
    print(f(a), a, b, c, sep=';')
except:
    print('ERR', a, b, c, sep=';')
code

3. Що буде виведено за виконання?
code
try:
    m = 9
    while m>=0:
        m += -2
    print(m, end=';')
except:
    print('ERR')
code

4. Що буде виведено за виконання?
code
try:
    for a in range(1, 6, 3):
        if a>=6:
            continue
        print(a, end=';');a = -1
    else:
        print(a,end=';')
    print(a)
except:
    print('ERR')
code

5. Що буде виведено за виконання?\par (Дійсні значення, якщо наявні, в цьому завданні будуть виводитись у форматі з фіксованою крапкою та, можливо, з доволі великою кількістю десяткових розрядів. У відповіді для дійсних значень у ЦЬОМУ завданні достатньо вказати один десятковий розряд після крапки, відкинувши решту розрядів без округлення. Наприклад, замість 13.66666666666667 достатньо вказати 13.6, замість 25.23 достатньо вказати 25.2)
code
try:
    print(8, end=';')
    print(int('b7'), end=';')
    print(1, end=';')
except TypeError:
    print(5, end=';')
except ZeroDivisionError:
    print(6, end=';')
except:
    print('ERR', end=';')
else:
    print(3, end=';')
finally:
    print(5)
code

6. Що буде виведено за виконання?
code
def f(n):
    if n>9:
        f(n//100)
    else:
        print(n%10,end='')

f(123456)
code

7. Що буде виведено за виконання?
code
class A:
    a = 1

    def _h1(self): print(2, end=';')
    def __h2(self): print(3, end=';')

    def main(self):
        try: print(self.a, end=';')
        except: print('ERR', end=';')

        try: self._h1()
        except: print('ERR', end=';')

        try: self.__h2()
        except: print('ERR', end=';')


class B(A):
    a = 4

    def _h1(self): print(5, end=';')
    def __h2(self): print(6, end=';')

try: B().main()
except: print('ERR', end=';')
code

8. Що буде виведено за виконання?
code
def f(arg, n):
    arg[18] = 9
    n = 6
    arg = set(arg.values())

try:
    n = 4
    s = {18:8, 55:2, 18:2}
    f(s, n)
    print(n, end=';')
    for x in s.keys():
        print(x, sep=';', end=';')
except:
    print('ERR')
code

9. Що буде виведено за виконання?
code
class A:
    a = 1
    b = 2

    def __init__(self):
        b = 3

class B(A):
    a = 4

    def __init__(self):
        super().__init__()
        self.c = 7

obj = B()
obj.b = obj.a + 10
B.b = 13

try: print(obj.b, end= ';')
except: print('ERR', end=';')

try: print(A.b, end= ';')
except: print('ERR', end=';')

try: print(B.b, end= ';')
except: print('ERR', end=';')
code

10. 
Для графа на рисунку з вершини 0 запущено алгоритм пошуку в глибину, що будує кістякове дерево. Списки суміжності вершин переглядаються алгоритмом за зростанням міток вершин.\par
\begin{center}
	\adjustimage{max size={0.9\linewidth}{0.9\paperheight}}
	{../10/Traversals/1.png}
\end{center}
\par Які з наведених ребер входять до складу отриманого кістякового дерева?\par
1) (2, 5)\par
2) (2, 3)\par
3) (1, 5)\par
4) (0, 4)\par
5) (3, 5)\par
6) (0, 5)\par
{\vskip 15pt} У формі вибрати номери відповідних ребер.\par
%answer:2 3 5
\end {large}}
%end
{\smallskip \begin {small} Затверджено на засіданні кафедри теоретичної кібернетики, протокол \No 12 від   19.05.2020 р.\par\smallskip\smallskip
{\parbox {6.5 cm}{Зав. кафедри \dotfill Ю.В.~Крак}}\hspace {0.7cm}{\hbox to 6.5 cm{ Екзаменатор \dotfill Т.О.~Карнаух}} \end{small}}
%begin
{{\begin {small}\par
{ \center  Київський національний університет імені Тараса Шевченка \par}
{\parbox {5 cm}{Освітня програма \hfill}{\parbox {7 cm}{Прикладна математика, Системний аналіз \hfill}}\parbox {6 cm}{\hfill Семестр 2}}\par
{\parbox {5 cm} {Навчальний предмет \hfill}\parbox {7 cm}{Програмування \hfill}\parbox {6cm}{\hfill}\par}\vspace{-7pt} \end{small}}{\center Екзаменаційний білет \No \textbf 
{ 254 }\par}}
{
\begin {large}
1. Що буде виведено за виконання?
code
try:
    res = "7.0j">="True"
    if isinstance(res, bool):
        print(f'bool;{res}')
    elif isinstance(res, int):
        print(f'int;{res}')
    elif isinstance(res, float):
        print(f'float;{res:.2f}')
    else:
        print('???')
except:
    print('ERR')
code

2. Що буде виведено за виконання?
code
a = 6
b = 1
c = 6

def f(b):
    global c
    a = 3
    b -= 2
    c = 3
    return a + b + c

a = 3
b = 2
c = 0

try:
    print(f(a), a, b, c, sep=';')
except:
    print('ERR', a, b, c, sep=';')
code

3. Що буде виведено за виконання?
code
try:
    i = 4
    while i<=5:
        i += 1
    print(i, end=';')
except:
    print('ERR')
code

4. Що буде виведено за виконання?
code
try:
    for b in range(0, -15, -2):
        if b<=7:
            break
        print(b, end=';')
        if b<=4:
            break
    else:
        print(b,end=';')
    print(b)
except:
    print('ERR')
code

5. Що буде виведено за виконання?\par (Дійсні значення, якщо наявні, в цьому завданні будуть виводитись у форматі з фіксованою крапкою та, можливо, з доволі великою кількістю десяткових розрядів. У відповіді для дійсних значень у ЦЬОМУ завданні достатньо вказати один десятковий розряд після крапки, відкинувши решту розрядів без округлення. Наприклад, замість 13.66666666666667 достатньо вказати 13.6, замість 25.23 достатньо вказати 25.2)
code
try:
    print(3, end=';')
    print(int('4'), end=';')
    print(6, end=';')
except KeyboardInterrupt:
    print(9, end=';')
except ValueError:
    print(8, end=';')
except:
    print('ERR', end=';')
else:
    print(0, end=';')
finally:
    print(1)
code

6. Що буде виведено за виконання?
code
def f(n):
    if n>9: f(n//10)
    else:
        print(n%10, end='')
 
f(123456)
code

7. Що буде виведено за виконання?
code
class A:
    d = 1

    def f1(self): print(2, end=';')
    def _f2(self): print(3, end=';')

    def main(self):
        try: print(self.d, end=';')
        except: print('ERR', end=';')

        try: self.f1()
        except: print('ERR', end=';')

        try: self._f2()
        except: print('ERR', end=';')


class B(A):
    d = 4

    def f1(self): print(5, end=';')
    def _f2(self): print(6, end=';')

try: B().main()
except: print('ERR', end=';')
code

8. Що буде виведено за виконання?
code
def f(arg, n):
    arg[79] = 8
    n = 9
    arg = list(arg.keys())

try:
    n = 1
    d = {32:8, 79:9, 79:7}
    f(d, n)
    print(n, end=';')
    for x in d.values():
        print(x, sep=';', end=';')
except:
    print('ERR')
code

9. Що буде виведено за виконання?
code
class A:
    c = 1
    b = 2

    def __init__(self):
        c = 3

class B(A):
    b = 4

    def __init__(self):
        super().__init__()
        self.a = 7

obj = B()
obj.c = obj.b + 10
B.c = 13

try: print(obj.c, end= ';')
except: print('ERR', end=';')

try: print(A.c, end= ';')
except: print('ERR', end=';')

try: print(B.c, end= ';')
except: print('ERR', end=';')
code

10. 
Для графа на рисунку з вершини 6 запущено алгоритм пошуку в глибину, що будує кістякове дерево. Списки суміжності вершин переглядаються алгоритмом за зростанням міток вершин.\par
\begin{center}
	\adjustimage{max size={0.9\linewidth}{0.9\paperheight}}
	{../10/Traversals/281.png}
\end{center}
\par Які з наведених ребер входять до складу отриманого кістякового дерева?\par
1) (0, 6)\par
2) (4, 5)\par
3) (1, 6)\par
4) (3, 5)\par
5) (2, 4)\par
6) (0, 5)\par
{\vskip 15pt} У формі вибрати номери відповідних ребер.\par
%answer:1 3 4 5 6
\end {large}}
%end
{\smallskip \begin {small} Затверджено на засіданні кафедри теоретичної кібернетики, протокол \No 12 від   19.05.2020 р.\par\smallskip\smallskip
{\parbox {6.5 cm}{Зав. кафедри \dotfill Ю.В.~Крак}}\hspace {0.7cm}{\hbox to 6.5 cm{ Екзаменатор \dotfill Т.О.~Карнаух}} \end{small}}
%begin
{{\begin {small}\par
{ \center  Київський національний університет імені Тараса Шевченка \par}
{\parbox {5 cm}{Освітня програма \hfill}{\parbox {7 cm}{Прикладна математика, Системний аналіз \hfill}}\parbox {6 cm}{\hfill Семестр 2}}\par
{\parbox {5 cm} {Навчальний предмет \hfill}\parbox {7 cm}{Програмування \hfill}\parbox {6cm}{\hfill}\par}\vspace{-7pt} \end{small}}{\center Екзаменаційний білет \No \textbf 
{ 255 }\par}}
{
\begin {large}
1. Що буде виведено за виконання?
code
try:
    res = "7">3.0
    if isinstance(res, bool):
        print(f'bool;{res}')
    elif isinstance(res, int):
        print(f'int;{res}')
    elif isinstance(res, float):
        print(f'float;{res:.2f}')
    else:
        print('???')
except:
    print('ERR')
code

2. Що буде виведено за виконання?
code
a = 0
b = 9
c = 4

def g(b):
    global c
    a = 3
    b = 5
    c = 4
    return a + b + c

a = 7
b = 5
c = 8

try:
    print(g(a), a, b, c, sep=';')
except:
    print('ERR', a, b, c, sep=';')
code

3. Що буде виведено за виконання?
code
try:
    t = 7
    while t<6:
        t += 1
    print(t, end=';')
except:
    print('ERR')
code

4. Що буде виведено за виконання?
code
try:
    for d in range(-9, -11, 2):
        if d<3:
            continue
        print(d, end=';')
    else:
        print('end',end=';')
    print(d)
except:
    print('ERR')
code

5. Що буде виведено за виконання?\par (Дійсні значення, якщо наявні, в цьому завданні будуть виводитись у форматі з фіксованою крапкою та, можливо, з доволі великою кількістю десяткових розрядів. У відповіді для дійсних значень у ЦЬОМУ завданні достатньо вказати один десятковий розряд після крапки, відкинувши решту розрядів без округлення. Наприклад, замість 13.66666666666667 достатньо вказати 13.6, замість 25.23 достатньо вказати 25.2)
code
try:
    print(7, end=';')
    print(2//0, end=';')
    print(6, end=';')
except Exception:
    print(7, end=';')
except TypeError:
    print(3, end=';')
except:
    print('ERR', end=';')
else:
    print(0, end=';')
finally:
    print(9)
code

6. Що буде виведено за виконання?
code
def f(n):
    if n>9:
        f(n//100)
        print(n%10, end='')
 
f(123456)
code

7. Що буде виведено за виконання?
code
class A:
    b = 1

    def _g1(self): print(2, end=';')
    def _g2(self): print(3, end=';')

    def main(self):
        try: print(self.b, end=';')
        except: print('ERR', end=';')

        try: self._g1()
        except: print('ERR', end=';')

        try: self._g2()
        except: print('ERR', end=';')


class B(A):
    b = 4

    def _g1(self): print(5, end=';')
    def _g2(self): print(6, end=';')

try: B().main()
except: print('ERR', end=';')
code

8. Що буде виведено за виконання?
code
def f(arg, n):
    arg[39] = 2
    n = 8
    arg = tuple(arg.items())

try:
    n = 2
    t = {68:5, 39:8, 39:4}
    f(t, n)
    print(n, end=';')
    for x in t.items():
        print(*x, sep=';', end=';')
except:
    print('ERR')
code

9. Що буде виведено за виконання?
code
class A:
    b = 1
    c = 2

    def __init__(self):
        b = 3

class B(A):
    b = 4

    def __init__(self):
        super().__init__()
        self.a = 7

obj = B()
obj.c = obj.a + 10
B.c = 13

try: print(obj.c, end= ';')
except: print('ERR', end=';')

try: print(A.c, end= ';')
except: print('ERR', end=';')

try: print(B.c, end= ';')
except: print('ERR', end=';')
code

10. 
Для графа на рисунку з вершини 3 запущено алгоритм пошуку в глибину, що будує кістякове дерево. Списки суміжності вершин переглядаються алгоритмом за зростанням міток вершин.\par
\begin{center}
	\adjustimage{max size={0.9\linewidth}{0.9\paperheight}}
	{../10/Traversals/111.png}
\end{center}
\par Які з наведених ребер входять до складу отриманого кістякового дерева?\par
1) (2, 6)\par
2) (1, 6)\par
3) (2, 3)\par
4) (2, 5)\par
5) (2, 4)\par
6) (0, 5)\par
{\vskip 15pt} У формі вибрати номери відповідних ребер.\par
%answer:1 2 6
\end {large}}
%end
{\smallskip \begin {small} Затверджено на засіданні кафедри теоретичної кібернетики, протокол \No 12 від   19.05.2020 р.\par\smallskip\smallskip
{\parbox {6.5 cm}{Зав. кафедри \dotfill Ю.В.~Крак}}\hspace {0.7cm}{\hbox to 6.5 cm{ Екзаменатор \dotfill Т.О.~Карнаух}} \end{small}}
%begin
{{\begin {small}\par
{ \center  Київський національний університет імені Тараса Шевченка \par}
{\parbox {5 cm}{Освітня програма \hfill}{\parbox {7 cm}{Прикладна математика, Системний аналіз \hfill}}\parbox {6 cm}{\hfill Семестр 2}}\par
{\parbox {5 cm} {Навчальний предмет \hfill}\parbox {7 cm}{Програмування \hfill}\parbox {6cm}{\hfill}\par}\vspace{-7pt} \end{small}}{\center Екзаменаційний білет \No \textbf 
{ 256 }\par}}
{
\begin {large}
1. Що буде виведено за виконання?
code
try:
    res = 7/5//7*6
    if isinstance(res, bool):
        print(f'bool;{res}')
    elif isinstance(res, int):
        print(f'int;{res}')
    elif isinstance(res, float):
        print(f'float;{res:.2f}')
    else:
        print('???')
except:
    print('ERR')
code

2. Що буде виведено за виконання?
code
a = 2
b = 6
c = 3

def f(b):
    global c
    a = 2
    b = 1
    c = 5
    return a + b + c

a = 1
b = 6
c = 9

try:
    print(f(a), a, b, c, sep=';')
except:
    print('ERR', a, b, c, sep=';')
code

3. Що буде виведено за виконання?
code
try:
    k = 12
    while k>8:
        k += -2
    print(k, end=';')
except:
    print('ERR')
code

4. Що буде виведено за виконання?
code
try:
    for c in range(-10, -5, -3):
        if c>4:
            continue
        print(c, end=';');c = 2
    else:
        print(c,end=';')
    print(c)
except:
    print('ERR')
code

5. Що буде виведено за виконання?\par (Дійсні значення, якщо наявні, в цьому завданні будуть виводитись у форматі з фіксованою крапкою та, можливо, з доволі великою кількістю десяткових розрядів. У відповіді для дійсних значень у ЦЬОМУ завданні достатньо вказати один десятковий розряд після крапки, відкинувши решту розрядів без округлення. Наприклад, замість 13.66666666666667 достатньо вказати 13.6, замість 25.23 достатньо вказати 25.2)
code
try:
    print(2, end=';')
    print(int('d5'), end=';')
    print(4, end=';')
except Exception:
    print(8, end=';')
except ValueError:
    print(1, end=';')
except:
    print('ERR', end=';')
else:
    print(9, end=';')
finally:
    print(2)
code

6. Що буде виведено за виконання?
code
def f(n):
    if n>9:
        f(n//100)
    print(n%10, end='')
 
f(123456)
code

7. Що буде виведено за виконання?
code
class A:
    c = 1

    def f1(self): print(2, end=';')
    def _f2(self): print(3, end=';')

    def main(self):
        try: print(self.c, end=';')
        except: print('ERR', end=';')

        try: self.f1()
        except: print('ERR', end=';')

        try: self._f2()
        except: print('ERR', end=';')


class B(A):
    c = 4

    def f1(self): print(5, end=';')
    def _f2(self): print(6, end=';')

try: B().main()
except: print('ERR', end=';')
code

8. Що буде виведено за виконання?
code
def f(arg, n):
    arg[60] = 3
    n = 5
    arg = set(arg.items())

try:
    n = 3
    t = {84:9, 48:4, 48:6}
    f(t, n)
    print(n, end=';')
    for x in t.items():
        print(*x, sep=';', end=';')
except:
    print('ERR')
code

9. Що буде виведено за виконання?
code
class A:
    c = 1
    a = 2

    def __init__(self):
        a = 3

class B(A):
    a = 4

    def __init__(self):
        super().__init__()
        self.b = 7

obj = B()
obj.b = obj.c + 10
B.b = 13

try: print(obj.b, end= ';')
except: print('ERR', end=';')

try: print(A.b, end= ';')
except: print('ERR', end=';')

try: print(B.b, end= ';')
except: print('ERR', end=';')
code

10. 
Для графа на рисунку з вершини 3 запущено алгоритм пошуку в глибину, що будує кістякове дерево. Списки суміжності вершин переглядаються алгоритмом за зростанням міток вершин.\par
\begin{center}
	\adjustimage{max size={0.9\linewidth}{0.9\paperheight}}
	{../10/Traversals/53.png}
\end{center}
\par Які з наведених ребер входять до складу отриманого кістякового дерева?\par
1) (4, 5)\par
2) (2, 3)\par
3) (2, 5)\par
4) (2, 6)\par
5) (0, 4)\par
6) (0, 2)\par
{\vskip 15pt} У формі вибрати номери відповідних ребер.\par
%answer:1 2 4 5 6
\end {large}}
%end
{\smallskip \begin {small} Затверджено на засіданні кафедри теоретичної кібернетики, протокол \No 12 від   19.05.2020 р.\par\smallskip\smallskip
{\parbox {6.5 cm}{Зав. кафедри \dotfill Ю.В.~Крак}}\hspace {0.7cm}{\hbox to 6.5 cm{ Екзаменатор \dotfill Т.О.~Карнаух}} \end{small}}
%begin
{{\begin {small}\par
{ \center  Київський національний університет імені Тараса Шевченка \par}
{\parbox {5 cm}{Освітня програма \hfill}{\parbox {7 cm}{Прикладна математика, Системний аналіз \hfill}}\parbox {6 cm}{\hfill Семестр 2}}\par
{\parbox {5 cm} {Навчальний предмет \hfill}\parbox {7 cm}{Програмування \hfill}\parbox {6cm}{\hfill}\par}\vspace{-7pt} \end{small}}{\center Екзаменаційний білет \No \textbf 
{ 257 }\par}}
{
\begin {large}
1. Що буде виведено за виконання?
code
try:
    res = 1**1.0*False
    if isinstance(res, bool):
        print(f'bool;{res}')
    elif isinstance(res, int):
        print(f'int;{res}')
    elif isinstance(res, float):
        print(f'float;{res:.2f}')
    else:
        print('???')
except:
    print('ERR')
code

2. Що буде виведено за виконання?
code
a = 1
b = 2
c = 3

def h(a):
    a = 2
    b = 3
    c = 1
    return a + b + c

a = 5
b = 8
c = 0

try:
    print(h(a), a, b, c, sep=';')
except:
    print('ERR', a, b, c, sep=';')
code

3. Що буде виведено за виконання?
code
try:
    i = 15
    while i>=4:
        i += -3
    print(i, end=';')
except:
    print('ERR')
code

4. Що буде виведено за виконання?
code
try:
    for e in range(8, -2, -1):
        if e>5:
            break
        print(e, end=';')
    else:
        print(e,end=';')
    print(e)
except:
    print('ERR')
code

5. Що буде виведено за виконання?\par (Дійсні значення, якщо наявні, в цьому завданні будуть виводитись у форматі з фіксованою крапкою та, можливо, з доволі великою кількістю десяткових розрядів. У відповіді для дійсних значень у ЦЬОМУ завданні достатньо вказати один десятковий розряд після крапки, відкинувши решту розрядів без округлення. Наприклад, замість 13.66666666666667 достатньо вказати 13.6, замість 25.23 достатньо вказати 25.2)
code
try:
    print(3, end=';')
    print(7/2, end=';')
    print(6, end=';')
except ZeroDivisionError:
    print(8, end=';')
except KeyboardInterrupt:
    print(1, end=';')
except:
    print('ERR', end=';')
else:
    print(5, end=';')
finally:
    print(0)
code

6. Що буде виведено за виконання?
code
def f(n):
    if n>9:
        return f(n//10)+1
    else:
        return 1

print(f(123456))
code

7. Що буде виведено за виконання?
code
class A:
    b = 1

    def _g1(self): print(2, end=';')
    def __g2(self): print(3, end=';')

    def main(self):
        try: print(self.b, end=';')
        except: print('ERR', end=';')

        try: self._g1()
        except: print('ERR', end=';')

        try: self.__g2()
        except: print('ERR', end=';')


class B(A):
    b = 4

    def _g1(self): print(5, end=';')
    def __g2(self): print(6, end=';')

try: B().main()
except: print('ERR', end=';')
code

8. Що буде виведено за виконання?
code
def f(arg, n):
    n = 9
    tmp = arg
    tmp[8:-9] = [1,2]
    return len(tmp)>3

try:
    g = [10,11,12,13,14]
    n = 2
    f(g, n)
    print(n, end=';')
    for el in g:
        print(el, end=';')
except:
    print('ERR')
code

9. Що буде виведено за виконання?
code
class A:
    a = 1
    c = 2

    def __init__(self):
        c = 3

class B(A):
    a = 4

    def __init__(self):
        super().__init__()
        self.b = 7

obj = B()
obj.a = obj.a + 10
B.a = 13

try: print(obj.a, end= ';')
except: print('ERR', end=';')

try: print(A.a, end= ';')
except: print('ERR', end=';')

try: print(B.a, end= ';')
except: print('ERR', end=';')
code

10. 
Для графа на рисунку з вершини 4 запущено алгоритм пошуку в глибину, що будує кістякове дерево. Списки суміжності вершин переглядаються алгоритмом за зростанням міток вершин.\par
\begin{center}
	\adjustimage{max size={0.9\linewidth}{0.9\paperheight}}
	{../10/Traversals/121.png}
\end{center}
\par Які з наведених ребер входять до складу отриманого кістякового дерева?\par
1) (4, 5)\par
2) (0, 6)\par
3) (2, 5)\par
4) (2, 3)\par
5) (4, 6)\par
6) (2, 6)\par
{\vskip 15pt} У формі вибрати номери відповідних ребер.\par
%answer:2 4
\end {large}}
%end
{\smallskip \begin {small} Затверджено на засіданні кафедри теоретичної кібернетики, протокол \No 12 від   19.05.2020 р.\par\smallskip\smallskip
{\parbox {6.5 cm}{Зав. кафедри \dotfill Ю.В.~Крак}}\hspace {0.7cm}{\hbox to 6.5 cm{ Екзаменатор \dotfill Т.О.~Карнаух}} \end{small}}
%begin
{{\begin {small}\par
{ \center  Київський національний університет імені Тараса Шевченка \par}
{\parbox {5 cm}{Освітня програма \hfill}{\parbox {7 cm}{Прикладна математика, Системний аналіз \hfill}}\parbox {6 cm}{\hfill Семестр 2}}\par
{\parbox {5 cm} {Навчальний предмет \hfill}\parbox {7 cm}{Програмування \hfill}\parbox {6cm}{\hfill}\par}\vspace{-7pt} \end{small}}{\center Екзаменаційний білет \No \textbf 
{ 258 }\par}}
{
\begin {large}
1. Що буде виведено за виконання?
code
try:
    res = 7//8%5*8
    if isinstance(res, bool):
        print(f'bool;{res}')
    elif isinstance(res, int):
        print(f'int;{res}')
    elif isinstance(res, float):
        print(f'float;{res:.2f}')
    else:
        print('???')
except:
    print('ERR')
code

2. Що буде виведено за виконання?
code
a = 7
b = 4
c = 5

def g(b):
    global c
    a = 4
    b = 4
    c = 5
    return a + b + c

a = 4
b = 0
c = 8

try:
    print(g(a), a, b, c, sep=';')
except:
    print('ERR', a, b, c, sep=';')
code

3. Що буде виведено за виконання?
code
try:
    k = 2
    while k<=9:
        k += 2
    print(k, end=';')
except:
    print('ERR')
code

4. Що буде виведено за виконання?
code
try:
    for f in range(10, -12, 3):
        if f<=8:
            continue
        print(f, end=';')
    else:
        print(f,end=';')
    print(f)
except:
    print('ERR')
code

5. Що буде виведено за виконання?\par (Дійсні значення, якщо наявні, в цьому завданні будуть виводитись у форматі з фіксованою крапкою та, можливо, з доволі великою кількістю десяткових розрядів. У відповіді для дійсних значень у ЦЬОМУ завданні достатньо вказати один десятковий розряд після крапки, відкинувши решту розрядів без округлення. Наприклад, замість 13.66666666666667 достатньо вказати 13.6, замість 25.23 достатньо вказати 25.2)
code
try:
    print(4, end=';')
    print(2//0.0, end=';')
    print(8, end=';')
except TypeError:
    print(7, end=';')
except Exception:
    print(0, end=';')
except:
    print('ERR', end=';')
else:
    print(1, end=';')
finally:
    print(5)
code

6. Що буде виведено за виконання?
code
def f(n):
    if n>2:
        f(n//2)
    else:
        print(n%2, end='')

f(16)
code

7. Що буде виведено за виконання?
code
class A:
    d = 1

    def _h1(self): print(2, end=';')
    def _h2(self): print(3, end=';')

    def main(self):
        try: print(self.d, end=';')
        except: print('ERR', end=';')

        try: self._h1()
        except: print('ERR', end=';')

        try: self._h2()
        except: print('ERR', end=';')


class B(A):
    d = 4

    def _h1(self): print(5, end=';')
    def _h2(self): print(6, end=';')

try: B().main()
except: print('ERR', end=';')
code

8. Що буде виведено за виконання?
code
def f(arg, n):
    arg[42] = 7
    n += -3

try:
    n = 6
    s = {37:2, 45:5, 37:7}
    f(s, n)
    print(n, end=';')
    for x in s.values():
        print(x, sep=';', end=';')
except:
    print('ERR')
code

9. Що буде виведено за виконання?
code
class A:
    b = 1
    a = 2

    def __init__(self):
        b = 3

class B(A):
    a = 4

    def __init__(self):
        super().__init__()
        self.c = 7

obj = B()
obj.b = obj.c + 10
B.b = 13

try: print(obj.b, end= ';')
except: print('ERR', end=';')

try: print(A.b, end= ';')
except: print('ERR', end=';')

try: print(B.b, end= ';')
except: print('ERR', end=';')
code

10. 
Для графа на рисунку з вершини 4 запущено алгоритм пошуку в глибину, що будує кістякове дерево. Списки суміжності вершин переглядаються алгоритмом за зростанням міток вершин.\par
\begin{center}
	\adjustimage{max size={0.9\linewidth}{0.9\paperheight}}
	{../10/Traversals/92.png}
\end{center}
\par Які з наведених ребер входять до складу отриманого кістякового дерева?\par
1) (1, 6)\par
2) (0, 6)\par
3) (4, 5)\par
4) (1, 3)\par
5) (0, 3)\par
6) (0, 5)\par
{\vskip 15pt} У формі вибрати номери відповідних ребер.\par
%answer:1 5 6
\end {large}}
%end
{\smallskip \begin {small} Затверджено на засіданні кафедри теоретичної кібернетики, протокол \No 12 від   19.05.2020 р.\par\smallskip\smallskip
{\parbox {6.5 cm}{Зав. кафедри \dotfill Ю.В.~Крак}}\hspace {0.7cm}{\hbox to 6.5 cm{ Екзаменатор \dotfill Т.О.~Карнаух}} \end{small}}
%begin
{{\begin {small}\par
{ \center  Київський національний університет імені Тараса Шевченка \par}
{\parbox {5 cm}{Освітня програма \hfill}{\parbox {7 cm}{Прикладна математика, Системний аналіз \hfill}}\parbox {6 cm}{\hfill Семестр 2}}\par
{\parbox {5 cm} {Навчальний предмет \hfill}\parbox {7 cm}{Програмування \hfill}\parbox {6cm}{\hfill}\par}\vspace{-7pt} \end{small}}{\center Екзаменаційний білет \No \textbf 
{ 259 }\par}}
{
\begin {large}
1. Що буде виведено за виконання?
code
def f(n):
    print('f', end='')
    return n

try:
    print(f(-2) and f(0))
except:
    print('ERR')
code

2. Що буде виведено за виконання?
code
a = 5
b = 4
c = 7

def f(a):
    a = 1
    b -= 5
    c = 4
    return a + b + c

a = 9
b = 8
c = 1

try:
    print(f(b), a, b, c, sep=';')
except:
    print('ERR', a, b, c, sep=';')
code

3. Що буде виведено за виконання?
code
try:
    j = 9
    while j<7:
        j += 2
    print(j, end=';')
except:
    print('ERR')
code

4. Що буде виведено за виконання?
code
try:
    for a in range(-1, -3, 2):
        if a>=3:
            continue
        print(a, end=';')
    else:
        print(a,end=';')
    print(a)
except:
    print('ERR')
code

5. Що буде виведено за виконання?\par (Дійсні значення, якщо наявні, в цьому завданні будуть виводитись у форматі з фіксованою крапкою та, можливо, з доволі великою кількістю десяткових розрядів. У відповіді для дійсних значень у ЦЬОМУ завданні достатньо вказати один десятковий розряд після крапки, відкинувши решту розрядів без округлення. Наприклад, замість 13.66666666666667 достатньо вказати 13.6, замість 25.23 достатньо вказати 25.2)
code
try:
    print(9, end=';')
    print(int('4'), end=';')
    print(3, end=';')
except TypeError:
    print(6, end=';')
except ValueError:
    print(1, end=';')
except:
    print('ERR', end=';')
else:
    print(2, end=';')
finally:
    print(4)
code

6. Що буде виведено за виконання?
code
def f(s):
    for el in s:
        el += 1
        return s
 
s = [1, 2, 3]
s = f(s)
print(s)
code

7. Що буде виведено за виконання?
code
class A:
    c = 1

    def _f1(self): print(2, end=';')
    def f2(self): print(3, end=';')

    def main(self):
        try: print(self.c, end=';')
        except: print('ERR', end=';')

        try: self._f1()
        except: print('ERR', end=';')

        try: self.f2()
        except: print('ERR', end=';')


class B(A):
    c = 4

    def _f1(self): print(5, end=';')
    def f2(self): print(6, end=';')

try: B().main()
except: print('ERR', end=';')
code

8. Що буде виведено за виконання?
code
def f(arg, n):
    n = 0
    tmp = arg
    del tmp[-2:-1]
    return len(tmp)>3

try:
    b = ['0','1','2','3','4']
    n = 8
    f(b, n)
    print(n, end=';')
    for el in b:
        print(el, end=';')
except:
    print('ERR')
code

9. Що буде виведено за виконання?
code
class A:
    b = 1
    a = 2

    def __init__(self):
        b = 3

class B(A):
    a = 4

    def __init__(self):
        super().__init__()
        self.c = 7

obj = B()
obj.a = obj.c + 10
B.a = 13

try: print(obj.a, end= ';')
except: print('ERR', end=';')

try: print(A.a, end= ';')
except: print('ERR', end=';')

try: print(B.a, end= ';')
except: print('ERR', end=';')
code

10. 
Для графа на рисунку з вершини 5 запущено алгоритм пошуку в глибину, що будує кістякове дерево. Списки суміжності вершин переглядаються алгоритмом за зростанням міток вершин.\par
\begin{center}
	\adjustimage{max size={0.9\linewidth}{0.9\paperheight}}
	{../10/Traversals/170.png}
\end{center}
\par Які з наведених ребер входять до складу отриманого кістякового дерева?\par
1) (1, 5)\par
2) (0, 3)\par
3) (0, 4)\par
4) (4, 6)\par
5) (1, 2)\par
6) (1, 4)\par
{\vskip 15pt} У формі вибрати номери відповідних ребер.\par
%answer:1 2 3 4 5
\end {large}}
%end
{\smallskip \begin {small} Затверджено на засіданні кафедри теоретичної кібернетики, протокол \No 12 від   19.05.2020 р.\par\smallskip\smallskip
{\parbox {6.5 cm}{Зав. кафедри \dotfill Ю.В.~Крак}}\hspace {0.7cm}{\hbox to 6.5 cm{ Екзаменатор \dotfill Т.О.~Карнаух}} \end{small}}
%begin
{{\begin {small}\par
{ \center  Київський національний університет імені Тараса Шевченка \par}
{\parbox {5 cm}{Освітня програма \hfill}{\parbox {7 cm}{Прикладна математика, Системний аналіз \hfill}}\parbox {6 cm}{\hfill Семестр 2}}\par
{\parbox {5 cm} {Навчальний предмет \hfill}\parbox {7 cm}{Програмування \hfill}\parbox {6cm}{\hfill}\par}\vspace{-7pt} \end{small}}{\center Екзаменаційний білет \No \textbf 
{ 260 }\par}}
{
\begin {large}
1. Що буде виведено за виконання?
code
try:
    res = 5//4/3*3
    if isinstance(res, bool):
        print(f'bool;{res}')
    elif isinstance(res, int):
        print(f'int;{res}')
    elif isinstance(res, float):
        print(f'float;{res:.2f}')
    else:
        print('???')
except:
    print('ERR')
code

2. Що буде виведено за виконання?
code
a = 9
b = 6
c = 3

def f(a):
    a = 2
    b = 1
    c = 3
    return a + b + c

a = 1
b = 2
c = 7

try:
    print(f(b), a, b, c, sep=';')
except:
    print('ERR', a, b, c, sep=';')
code

3. Що буде виведено за виконання?
code
try:
    n = 1
    while n<=13:
        n += 1
    print(n, end=';')
except:
    print('ERR')
code

4. Що буде виведено за виконання?
code
try:
    for b in range(-6, -8, -2):
        if b<4:
            break
        print(b, end=';')
    else:
        print(b,end=';')
    print(b)
except:
    print('ERR')
code

5. Що буде виведено за виконання?\par (Дійсні значення, якщо наявні, в цьому завданні будуть виводитись у форматі з фіксованою крапкою та, можливо, з доволі великою кількістю десяткових розрядів. У відповіді для дійсних значень у ЦЬОМУ завданні достатньо вказати один десятковий розряд після крапки, відкинувши решту розрядів без округлення. Наприклад, замість 13.66666666666667 достатньо вказати 13.6, замість 25.23 достатньо вказати 25.2)
code
try:
    print(3, end=';')
    print(8j>=4j, end=';')
    print(9, end=';')
except KeyboardInterrupt:
    print(6, end=';')
except ZeroDivisionError:
    print(7, end=';')
except:
    print('ERR', end=';')
else:
    print(0, end=';')
finally:
    print(5)
code

6. Що буде виведено за виконання?
code
def f(n):
    if n>9:
        f(n//100)
    else:
        print(n%10,end='')

f(987654)
code

7. Що буде виведено за виконання?
code
class A:
    a = 1

    def __g1(self): print(2, end=';')
    def _g2(self): print(3, end=';')

    def main(self):
        try: print(self.a, end=';')
        except: print('ERR', end=';')

        try: self.__g1()
        except: print('ERR', end=';')

        try: self._g2()
        except: print('ERR', end=';')


class B(A):
    a = 4

    def __g1(self): print(5, end=';')
    def _g2(self): print(6, end=';')

try: B().main()
except: print('ERR', end=';')
code

8. Що буде виведено за виконання?
code
def f(arg, n):
    arg[24] = 4
    n = 7
    arg = tuple(arg.values())

try:
    n = 1
    d = {55:1, 24:8, 44:7, 55:5}
    f(d, n)
    print(n, end=';')
    for x in d.items():
        print(x, sep=';', end=';')
except:
    print('ERR')
code

9. Що буде виведено за виконання?
code
class A:
    c = 1
    b = 2

    def __init__(self):
        b = 3

class B(A):
    c = 4

    def __init__(self):
        super().__init__()
        self.a = 7

obj = B()
obj.b = obj.b + 10
B.b = 13

try: print(obj.b, end= ';')
except: print('ERR', end=';')

try: print(A.b, end= ';')
except: print('ERR', end=';')

try: print(B.b, end= ';')
except: print('ERR', end=';')
code

10. 
Для графа на рисунку з вершини 6 запущено алгоритм пошуку в глибину, що будує кістякове дерево. Списки суміжності вершин переглядаються алгоритмом за зростанням міток вершин.\par
\begin{center}
	\adjustimage{max size={0.9\linewidth}{0.9\paperheight}}
	{../10/Traversals/63.png}
\end{center}
\par Які з наведених ребер входять до складу отриманого кістякового дерева?\par
1) (4, 6)\par
2) (0, 2)\par
3) (0, 6)\par
4) (2, 6)\par
5) (1, 5)\par
6) (1, 4)\par
{\vskip 15pt} У формі вибрати номери відповідних ребер.\par
%answer:2 3 5 6
\end {large}}
%end
{\smallskip \begin {small} Затверджено на засіданні кафедри теоретичної кібернетики, протокол \No 12 від   19.05.2020 р.\par\smallskip\smallskip
{\parbox {6.5 cm}{Зав. кафедри \dotfill Ю.В.~Крак}}\hspace {0.7cm}{\hbox to 6.5 cm{ Екзаменатор \dotfill Т.О.~Карнаух}} \end{small}}
%begin
{{\begin {small}\par
{ \center  Київський національний університет імені Тараса Шевченка \par}
{\parbox {5 cm}{Освітня програма \hfill}{\parbox {7 cm}{Прикладна математика, Системний аналіз \hfill}}\parbox {6 cm}{\hfill Семестр 2}}\par
{\parbox {5 cm} {Навчальний предмет \hfill}\parbox {7 cm}{Програмування \hfill}\parbox {6cm}{\hfill}\par}\vspace{-7pt} \end{small}}{\center Екзаменаційний білет \No \textbf 
{ 261 }\par}}
{
\begin {large}
1. Що буде виведено за виконання?
code
try:
    res = not 8.0!=5 and 9>True or 3<=9.0
    if isinstance(res, bool):
        print(f'bool;{res}')
    elif isinstance(res, int):
        print(f'int;{res}')
    elif isinstance(res, float):
        print(f'float;{res:.2f}')
    else:
        print('???')
except:
    print('ERR')
code

2. Що буде виведено за виконання?
code
a = 0
b = 4
c = 3

def g(b):
    global c
    a = 3
    b *= 5
    c = 1
    return a + b + c

a = 5
b = 2
c = 1

try:
    print(g(b), a, b, c, sep=';')
except:
    print('ERR', a, b, c, sep=';')
code

3. Що буде виведено за виконання?
code
try:
    j = 5
    while j>=2:
        j += -2
    print(j, end=';')
except:
    print('ERR')
code

4. Що буде виведено за виконання?
code
try:
    for e in range(2, 8, -3):
        if e<8:
            continue
        print(e, end=';');e = 4
        if e<5:
            break
    else:
        print('end',end=';')
    print(e)
except:
    print('ERR')
code

5. Що буде виведено за виконання?\par (Дійсні значення, якщо наявні, в цьому завданні будуть виводитись у форматі з фіксованою крапкою та, можливо, з доволі великою кількістю десяткових розрядів. У відповіді для дійсних значень у ЦЬОМУ завданні достатньо вказати один десятковий розряд після крапки, відкинувши решту розрядів без округлення. Наприклад, замість 13.66666666666667 достатньо вказати 13.6, замість 25.23 достатньо вказати 25.2)
code
try:
    print(3, end=';')
    print(7j!=2j, end=';')
    print(4, end=';')
except ZeroDivisionError:
    print(8, end=';')
except KeyboardInterrupt:
    print(2, end=';')
except:
    print('ERR', end=';')
else:
    print(0, end=';')
finally:
    print(1)
code

6. Що буде виведено за виконання?
code
def f(s1):
    s = s1[:]
    for i in range(len(s)):
        s[i] += 1
 
s = [1, 2, 3]
f(s)
print(s)
code

7. Що буде виведено за виконання?
code
class A:
    d = 1

    def __h1(self): print(2, end=';')
    def _h2(self): print(3, end=';')

    def main(self):
        try: print(self.d, end=';')
        except: print('ERR', end=';')

        try: self.__h1()
        except: print('ERR', end=';')

        try: self._h2()
        except: print('ERR', end=';')


class B(A):
    d = 4

    def __h1(self): print(5, end=';')
    def _h2(self): print(6, end=';')

try: B().main()
except: print('ERR', end=';')
code

8. Що буде виведено за виконання?
code
def f(arg, n):
    arg[67] = 5
    n *= -2

try:
    n = 7
    t = {63:4, 67:3, 67:8}
    f(t, n)
    print(n, end=';')
    for x in t :
        print(x, sep=';', end=';')
except:
    print('ERR')
code

9. Що буде виведено за виконання?
code
class A:
    a = 1
    b = 2

    def __init__(self):
        a = 3

class B(A):
    b = 4

    def __init__(self):
        super().__init__()
        self.c = 7

obj = B()
obj.c = obj.a + 10
B.c = 13

try: print(obj.c, end= ';')
except: print('ERR', end=';')

try: print(A.c, end= ';')
except: print('ERR', end=';')

try: print(B.c, end= ';')
except: print('ERR', end=';')
code

10. 
Для графа на рисунку з вершини 1 запущено алгоритм пошуку в глибину, що будує кістякове дерево. Списки суміжності вершин переглядаються алгоритмом за зростанням міток вершин.\par
\begin{center}
	\adjustimage{max size={0.9\linewidth}{0.9\paperheight}}
	{../10/Traversals/254.png}
\end{center}
\par Які з наведених ребер входять до складу отриманого кістякового дерева?\par
1) (0, 2)\par
2) (0, 1)\par
3) (2, 4)\par
4) (0, 5)\par
5) (0, 3)\par
6) (1, 5)\par
{\vskip 15pt} У формі вибрати номери відповідних ребер.\par
%answer:1 2 3 4 5
\end {large}}
%end
{\smallskip \begin {small} Затверджено на засіданні кафедри теоретичної кібернетики, протокол \No 12 від   19.05.2020 р.\par\smallskip\smallskip
{\parbox {6.5 cm}{Зав. кафедри \dotfill Ю.В.~Крак}}\hspace {0.7cm}{\hbox to 6.5 cm{ Екзаменатор \dotfill Т.О.~Карнаух}} \end{small}}
%begin
{{\begin {small}\par
{ \center  Київський національний університет імені Тараса Шевченка \par}
{\parbox {5 cm}{Освітня програма \hfill}{\parbox {7 cm}{Прикладна математика, Системний аналіз \hfill}}\parbox {6 cm}{\hfill Семестр 2}}\par
{\parbox {5 cm} {Навчальний предмет \hfill}\parbox {7 cm}{Програмування \hfill}\parbox {6cm}{\hfill}\par}\vspace{-7pt} \end{small}}{\center Екзаменаційний білет \No \textbf 
{ 262 }\par}}
{
\begin {large}
1. Що буде виведено за виконання?
code
try:
    res = 9/9*8*9
    if isinstance(res, bool):
        print(f'bool;{res}')
    elif isinstance(res, int):
        print(f'int;{res}')
    elif isinstance(res, float):
        print(f'float;{res:.2f}')
    else:
        print('???')
except:
    print('ERR')
code

2. Що буде виведено за виконання?
code
a = 7
b = 8
c = 9

def h(a):
    a = 4
    b += 2
    c = 2
    return a + b + c

a = 6
b = 7
c = 6

try:
    print(h(b), a, b, c, sep=';')
except:
    print('ERR', a, b, c, sep=';')
code

3. Що буде виведено за виконання?
code
try:
    n = 6
    while n>1:
        n += -3
    print(n, end=';')
except:
    print('ERR')
code

4. Що буде виведено за виконання?
code
try:
    for f in range(-12, 5, 3):
        if f>7:
            continue
        print(f, end=';')
    else:
        print(f,end=';')
    print(f)
except:
    print('ERR')
code

5. Що буде виведено за виконання?\par (Дійсні значення, якщо наявні, в цьому завданні будуть виводитись у форматі з фіксованою крапкою та, можливо, з доволі великою кількістю десяткових розрядів. У відповіді для дійсних значень у ЦЬОМУ завданні достатньо вказати один десятковий розряд після крапки, відкинувши решту розрядів без округлення. Наприклад, замість 13.66666666666667 достатньо вказати 13.6, замість 25.23 достатньо вказати 25.2)
code
try:
    print(9, end=';')
    print(6%1, end=';')
    print(5, end=';')
except TypeError:
    print(7, end=';')
except Exception:
    print(3, end=';')
except:
    print('ERR', end=';')
else:
    print(4, end=';')
finally:
    print(6)
code

6. Що буде виведено за виконання?
code
def f(n):
    if n>9:
        return f(n//100)+1
    else:
        return 1

print(f(123456))
code

7. Що буде виведено за виконання?
code
class A:
    c = 1

    def _h1(self): print(2, end=';')
    def _h2(self): print(3, end=';')

    def main(self):
        try: print(self.c, end=';')
        except: print('ERR', end=';')

        try: self._h1()
        except: print('ERR', end=';')

        try: self._h2()
        except: print('ERR', end=';')


class B(A):
    c = 4

    def _h1(self): print(5, end=';')
    def _h2(self): print(6, end=';')

try: B().main()
except: print('ERR', end=';')
code

8. Що буде виведено за виконання?
code
def f(arg, n):
    arg[71] = 0
    n = 5
    arg = list(arg.keys())

try:
    n = 0
    s = {20:9, 81:5, 81:1}
    f(s, n)
    print(n, end=';')
    for x in s.values():
        print(x, sep=';', end=';')
except:
    print('ERR')
code

9. Що буде виведено за виконання?
code
class A:
    c = 1
    a = 2

    def __init__(self):
        a = 3

class B(A):
    c = 4

    def __init__(self):
        super().__init__()
        self.b = 7

obj = B()
obj.a = obj.b + 10
B.a = 13

try: print(obj.a, end= ';')
except: print('ERR', end=';')

try: print(A.a, end= ';')
except: print('ERR', end=';')

try: print(B.a, end= ';')
except: print('ERR', end=';')
code

10. 
Для графа на рисунку з вершини 3 запущено алгоритм пошуку в глибину, що будує кістякове дерево. Списки суміжності вершин переглядаються алгоритмом за зростанням міток вершин.\par
\begin{center}
	\adjustimage{max size={0.9\linewidth}{0.9\paperheight}}
	{../10/Traversals/41.png}
\end{center}
\par Які з наведених ребер входять до складу отриманого кістякового дерева?\par
1) (4, 6)\par
2) (2, 6)\par
3) (2, 4)\par
4) (1, 3)\par
5) (0, 6)\par
6) (2, 5)\par
{\vskip 15pt} У формі вибрати номери відповідних ребер.\par
%answer:1 4 5
\end {large}}
%end
{\smallskip \begin {small} Затверджено на засіданні кафедри теоретичної кібернетики, протокол \No 12 від   19.05.2020 р.\par\smallskip\smallskip
{\parbox {6.5 cm}{Зав. кафедри \dotfill Ю.В.~Крак}}\hspace {0.7cm}{\hbox to 6.5 cm{ Екзаменатор \dotfill Т.О.~Карнаух}} \end{small}}
%begin
{{\begin {small}\par
{ \center  Київський національний університет імені Тараса Шевченка \par}
{\parbox {5 cm}{Освітня програма \hfill}{\parbox {7 cm}{Прикладна математика, Системний аналіз \hfill}}\parbox {6 cm}{\hfill Семестр 2}}\par
{\parbox {5 cm} {Навчальний предмет \hfill}\parbox {7 cm}{Програмування \hfill}\parbox {6cm}{\hfill}\par}\vspace{-7pt} \end{small}}{\center Екзаменаційний білет \No \textbf 
{ 263 }\par}}
{
\begin {large}
1. Що буде виведено за виконання?
code
try:
    res = False!=4j
    if isinstance(res, bool):
        print(f'bool;{res}')
    elif isinstance(res, int):
        print(f'int;{res}')
    elif isinstance(res, float):
        print(f'float;{res:.2f}')
    else:
        print('???')
except:
    print('ERR')
code

2. Що буде виведено за виконання?
code
a = 3
b = 6
c = 5

def h(b):
    a = 2
    b = 1
    c = 4
    return a + b + c

a = 0
b = 2
c = 4

try:
    print(h(b), a, b, c, sep=';')
except:
    print('ERR', a, b, c, sep=';')
code

3. Що буде виведено за виконання?
code
try:
    n = 8
    while n>6:
        n += -2
    print(n, end=';')
except:
    print('ERR')
code

4. Що буде виведено за виконання?
code
try:
    for d in range(5, 12, 2):
        if d<=6:
            continue
        print(d, end=';');d = -3
    else:
        print(d,end=';')
    print(d)
except:
    print('ERR')
code

5. Що буде виведено за виконання?\par (Дійсні значення, якщо наявні, в цьому завданні будуть виводитись у форматі з фіксованою крапкою та, можливо, з доволі великою кількістю десяткових розрядів. У відповіді для дійсних значень у ЦЬОМУ завданні достатньо вказати один десятковий розряд після крапки, відкинувши решту розрядів без округлення. Наприклад, замість 13.66666666666667 достатньо вказати 13.6, замість 25.23 достатньо вказати 25.2)
code
try:
    print(8, end=';')
    print(int('a2'), end=';')
    print(1, end=';')
except ValueError:
    print(9, end=';')
except ValueError:
    print(5, end=';')
except:
    print('ERR', end=';')
else:
    print(0, end=';')
finally:
    print(0)
code

6. Що буде виведено за виконання?
code
def f(n):
    print(n%2, end='')
    if n>2:
        f(n//2)
 
f(16)
code

7. Що буде виведено за виконання?
code
class A:
    b = 1

    def _f1(self): print(2, end=';')
    def __f2(self): print(3, end=';')

    def main(self):
        try: print(self.b, end=';')
        except: print('ERR', end=';')

        try: self._f1()
        except: print('ERR', end=';')

        try: self.__f2()
        except: print('ERR', end=';')


class B(A):
    b = 4

    def _f1(self): print(5, end=';')
    def __f2(self): print(6, end=';')

try: B().main()
except: print('ERR', end=';')
code

8. Що буде виведено за виконання?
code
def f(arg, n):
    arg[24] = 1
    n = 9
    arg = set(arg.items())

try:
    n = 4
    t = {86:6, 84:2, 15:0, 86:5}
    f(t, n)
    print(n, end=';')
    for x in t.values():
        print(x, sep=';', end=';')
except:
    print('ERR')
code

9. Що буде виведено за виконання?
code
class A:
    a = 1
    c = 2

    def __init__(self):
        c = 3

class B(A):
    c = 4

    def __init__(self):
        super().__init__()
        self.b = 7

obj = B()
obj.c = obj.c + 10
B.c = 13

try: print(obj.c, end= ';')
except: print('ERR', end=';')

try: print(A.c, end= ';')
except: print('ERR', end=';')

try: print(B.c, end= ';')
except: print('ERR', end=';')
code

10. 
Для графа на рисунку з вершини 2 запущено алгоритм пошуку в глибину, що будує кістякове дерево. Списки суміжності вершин переглядаються алгоритмом за зростанням міток вершин.\par
\begin{center}
	\adjustimage{max size={0.9\linewidth}{0.9\paperheight}}
	{../10/Traversals/220.png}
\end{center}
\par Які з наведених ребер входять до складу отриманого кістякового дерева?\par
1) (1, 5)\par
2) (2, 6)\par
3) (1, 2)\par
4) (2, 3)\par
5) (0, 4)\par
6) (0, 6)\par
{\vskip 15pt} У формі вибрати номери відповідних ребер.\par
%answer:1 5
\end {large}}
%end
{\smallskip \begin {small} Затверджено на засіданні кафедри теоретичної кібернетики, протокол \No 12 від   19.05.2020 р.\par\smallskip\smallskip
{\parbox {6.5 cm}{Зав. кафедри \dotfill Ю.В.~Крак}}\hspace {0.7cm}{\hbox to 6.5 cm{ Екзаменатор \dotfill Т.О.~Карнаух}} \end{small}}
%begin
{{\begin {small}\par
{ \center  Київський національний університет імені Тараса Шевченка \par}
{\parbox {5 cm}{Освітня програма \hfill}{\parbox {7 cm}{Прикладна математика, Системний аналіз \hfill}}\parbox {6 cm}{\hfill Семестр 2}}\par
{\parbox {5 cm} {Навчальний предмет \hfill}\parbox {7 cm}{Програмування \hfill}\parbox {6cm}{\hfill}\par}\vspace{-7pt} \end{small}}{\center Екзаменаційний білет \No \textbf 
{ 264 }\par}}
{
\begin {large}
1. Що буде виведено за виконання?
code
def f(n):
    print('f', end='')
    return n

try:
    print(f(8) and f(-3))
except:
    print('ERR')
code

2. Що буде виведено за виконання?
code
a = 7
b = 1
c = 8

def g(a):
    global c
    a = 3
    b = 5
    c = 4
    return a + b + c

a = 9
b = 1
c = 7

try:
    print(g(b), a, b, c, sep=';')
except:
    print('ERR', a, b, c, sep=';')
code

3. Що буде виведено за виконання?
code
try:
    i = 1
    while i<=5:
        i += 1
    print(i, end=';')
except:
    print('ERR')
code

4. Що буде виведено за виконання?
code
try:
    for c in range(-5, 13, -2):
        if c>=5:
            break
        print(c, end=';')
    else:
        print(c,end=';')
    print(c)
except:
    print('ERR')
code

5. Що буде виведено за виконання?\par (Дійсні значення, якщо наявні, в цьому завданні будуть виводитись у форматі з фіксованою крапкою та, можливо, з доволі великою кількістю десяткових розрядів. У відповіді для дійсних значень у ЦЬОМУ завданні достатньо вказати один десятковий розряд після крапки, відкинувши решту розрядів без округлення. Наприклад, замість 13.66666666666667 достатньо вказати 13.6, замість 25.23 достатньо вказати 25.2)
code
try:
    print(2, end=';')
    print(9//0.0, end=';')
    print(1, end=';')
except KeyboardInterrupt:
    print(8, end=';')
except TypeError:
    print(7, end=';')
except:
    print('ERR', end=';')
else:
    print(3, end=';')
finally:
    print(4)
code

6. Що буде виведено за виконання?
code
def f(n):
    print(n%10,end='')
    if n>9:
        f(n//100)
 
f(123456)
code

7. Що буде виведено за виконання?
code
class A:
    a = 1

    def _g1(self): print(2, end=';')
    def g2(self): print(3, end=';')

    def main(self):
        try: print(self.a, end=';')
        except: print('ERR', end=';')

        try: self._g1()
        except: print('ERR', end=';')

        try: self.g2()
        except: print('ERR', end=';')


class B(A):
    a = 4

    def _g1(self): print(5, end=';')
    def g2(self): print(6, end=';')

try: B().main()
except: print('ERR', end=';')
code

8. Що буде виведено за виконання?
code
def f(arg, n):
    n = 7
    tmp = arg
    tmp[7:] = 'xy'
    return len(tmp)>3

try:
    g = ['0','1','2','3','4']
    n = 6
    f(g, n)
    print(n, end=';')
    for el in g:
        print(el, end=';')
except:
    print('ERR')
code

9. Що буде виведено за виконання?
code
class A:
    b = 1
    c = 2

    def __init__(self):
        b = 3

class B(A):
    b = 4

    def __init__(self):
        super().__init__()
        self.a = 7

obj = B()
obj.b = obj.a + 10
B.b = 13

try: print(obj.b, end= ';')
except: print('ERR', end=';')

try: print(A.b, end= ';')
except: print('ERR', end=';')

try: print(B.b, end= ';')
except: print('ERR', end=';')
code

10. 
Для графа на рисунку з вершини 6 запущено алгоритм пошуку в глибину, що будує кістякове дерево. Списки суміжності вершин переглядаються алгоритмом за зростанням міток вершин.\par
\begin{center}
	\adjustimage{max size={0.9\linewidth}{0.9\paperheight}}
	{../10/Traversals/138.png}
\end{center}
\par Які з наведених ребер входять до складу отриманого кістякового дерева?\par
1) (1, 3)\par
2) (1, 6)\par
3) (0, 3)\par
4) (0, 2)\par
5) (2, 4)\par
6) (1, 2)\par
{\vskip 15pt} У формі вибрати номери відповідних ребер.\par
%answer:2 3 4 6
\end {large}}
%end
{\smallskip \begin {small} Затверджено на засіданні кафедри теоретичної кібернетики, протокол \No 12 від   19.05.2020 р.\par\smallskip\smallskip
{\parbox {6.5 cm}{Зав. кафедри \dotfill Ю.В.~Крак}}\hspace {0.7cm}{\hbox to 6.5 cm{ Екзаменатор \dotfill Т.О.~Карнаух}} \end{small}}
%begin
{{\begin {small}\par
{ \center  Київський національний університет імені Тараса Шевченка \par}
{\parbox {5 cm}{Освітня програма \hfill}{\parbox {7 cm}{Прикладна математика, Системний аналіз \hfill}}\parbox {6 cm}{\hfill Семестр 2}}\par
{\parbox {5 cm} {Навчальний предмет \hfill}\parbox {7 cm}{Програмування \hfill}\parbox {6cm}{\hfill}\par}\vspace{-7pt} \end{small}}{\center Екзаменаційний білет \No \textbf 
{ 265 }\par}}
{
\begin {large}
1. Що буде виведено за виконання?
code
try:
    res = 3//6%9*6
    if isinstance(res, bool):
        print(f'bool;{res}')
    elif isinstance(res, int):
        print(f'int;{res}')
    elif isinstance(res, float):
        print(f'float;{res:.2f}')
    else:
        print('???')
except:
    print('ERR')
code

2. Що буде виведено за виконання?
code
a = 8
b = 6
c = 4

def h(a):
    global c
    a = 5
    b = 1
    c = 3
    return a + b + c

a = 2
b = 9
c = 0

try:
    print(h(b), a, b, c, sep=';')
except:
    print('ERR', a, b, c, sep=';')
code

3. Що буде виведено за виконання?
code
try:
    m = 8
    while m<12:
        m += 2
    print(m, end=';')
except:
    print('ERR')
code

4. Що буде виведено за виконання?
code
try:
    for f in range(-11, 0, -1):
        if f<8:
            break
        print(f, end=';')
    else:
        print(f,end=';')
    print(f)
except:
    print('ERR')
code

5. Що буде виведено за виконання?\par (Дійсні значення, якщо наявні, в цьому завданні будуть виводитись у форматі з фіксованою крапкою та, можливо, з доволі великою кількістю десяткових розрядів. У відповіді для дійсних значень у ЦЬОМУ завданні достатньо вказати один десятковий розряд після крапки, відкинувши решту розрядів без округлення. Наприклад, замість 13.66666666666667 достатньо вказати 13.6, замість 25.23 достатньо вказати 25.2)
code
try:
    print(6, end=';')
    print(int('5'), end=';')
    print(1, end=';')
except Exception:
    print(6, end=';')
except ZeroDivisionError:
    print(5, end=';')
except:
    print('ERR', end=';')
else:
    print(3, end=';')
finally:
    print(0)
code

6. Що буде виведено за виконання?
code
def f(n):
    print(n%10, end='')
    if n>9:
        f(n//10)
 
f(543216)
code

7. Що буде виведено за виконання?
code
class A:
    b = 1

    def f1(self): print(2, end=';')
    def _f2(self): print(3, end=';')

    def main(self):
        try: print(self.b, end=';')
        except: print('ERR', end=';')

        try: self.f1()
        except: print('ERR', end=';')

        try: self._f2()
        except: print('ERR', end=';')


class B(A):
    b = 4

    def f1(self): print(5, end=';')
    def _f2(self): print(6, end=';')

try: B().main()
except: print('ERR', end=';')
code

8. Що буде виведено за виконання?
code
def f(arg, n):
    arg[66] = 3
    n = 8
    arg = tuple(arg.values())

try:
    n = 2
    s = {82:5, 66:2, 82:3}
    f(s, n)
    print(n, end=';')
    for x in s.keys():
        print(x, sep=';', end=';')
except:
    print('ERR')
code

9. Що буде виведено за виконання?
code
class A:
    b = 1
    a = 2

    def __init__(self):
        b = 3

class B(A):
    a = 4

    def __init__(self):
        super().__init__()
        self.c = 7

obj = B()
obj.c = obj.b + 10
B.c = 13

try: print(obj.c, end= ';')
except: print('ERR', end=';')

try: print(A.c, end= ';')
except: print('ERR', end=';')

try: print(B.c, end= ';')
except: print('ERR', end=';')
code

10. 
Для графа на рисунку з вершини 4 запущено алгоритм пошуку в глибину, що будує кістякове дерево. Списки суміжності вершин переглядаються алгоритмом за зростанням міток вершин.\par
\begin{center}
	\adjustimage{max size={0.9\linewidth}{0.9\paperheight}}
	{../10/Traversals/177.png}
\end{center}
\par Які з наведених ребер входять до складу отриманого кістякового дерева?\par
1) (3, 5)\par
2) (0, 1)\par
3) (5, 6)\par
4) (1, 3)\par
5) (1, 5)\par
6) (0, 4)\par
{\vskip 15pt} У формі вибрати номери відповідних ребер.\par
%answer:1 2 3 4 6
\end {large}}
%end
{\smallskip \begin {small} Затверджено на засіданні кафедри теоретичної кібернетики, протокол \No 12 від   19.05.2020 р.\par\smallskip\smallskip
{\parbox {6.5 cm}{Зав. кафедри \dotfill Ю.В.~Крак}}\hspace {0.7cm}{\hbox to 6.5 cm{ Екзаменатор \dotfill Т.О.~Карнаух}} \end{small}}
%begin
{{\begin {small}\par
{ \center  Київський національний університет імені Тараса Шевченка \par}
{\parbox {5 cm}{Освітня програма \hfill}{\parbox {7 cm}{Прикладна математика, Системний аналіз \hfill}}\parbox {6 cm}{\hfill Семестр 2}}\par
{\parbox {5 cm} {Навчальний предмет \hfill}\parbox {7 cm}{Програмування \hfill}\parbox {6cm}{\hfill}\par}\vspace{-7pt} \end{small}}{\center Екзаменаційний білет \No \textbf 
{ 266 }\par}}
{
\begin {large}
1. Що буде виведено за виконання?
code
try:
    res = 8/5*4*4
    if isinstance(res, bool):
        print(f'bool;{res}')
    elif isinstance(res, int):
        print(f'int;{res}')
    elif isinstance(res, float):
        print(f'float;{res:.2f}')
    else:
        print('???')
except:
    print('ERR')
code

2. Що буде виведено за виконання?
code
a = 5
b = 3
c = 1

def g(b):
    a = 2
    b = 5
    c = 4
    return a + b + c

a = 8
b = 2
c = 7

try:
    print(g(a), a, b, c, sep=';')
except:
    print('ERR', a, b, c, sep=';')
code

3. Що буде виведено за виконання?
code
try:
    i = 14
    while i>=3:
        i += -3
    print(i, end=';')
except:
    print('ERR')
code

4. Що буде виведено за виконання?
code
try:
    for b in range(3, -4, 2):
        if b<=3:
            continue
        print(b, end=';');b = -9
    else:
        print(b,end=';')
    print(b)
except:
    print('ERR')
code

5. Що буде виведено за виконання?\par (Дійсні значення, якщо наявні, в цьому завданні будуть виводитись у форматі з фіксованою крапкою та, можливо, з доволі великою кількістю десяткових розрядів. У відповіді для дійсних значень у ЦЬОМУ завданні достатньо вказати один десятковий розряд після крапки, відкинувши решту розрядів без округлення. Наприклад, замість 13.66666666666667 достатньо вказати 13.6, замість 25.23 достатньо вказати 25.2)
code
try:
    print(7, end=';')
    print(2%3, end=';')
    print(9, end=';')
except ZeroDivisionError:
    print(4, end=';')
except Exception:
    print(8, end=';')
except:
    print('ERR', end=';')
else:
    print(7, end=';')
finally:
    print(8)
code

6. Що буде виведено за виконання?
code
def f(s):
    for i in range(len(s)):
        s[i] += 1
        return
 
s = [1, 2, 3]
f(s)
print(s)
code

7. Що буде виведено за виконання?
code
class A:
    c = 1

    def h1(self): print(2, end=';')
    def _h2(self): print(3, end=';')

    def main(self):
        try: print(self.c, end=';')
        except: print('ERR', end=';')

        try: self.h1()
        except: print('ERR', end=';')

        try: self._h2()
        except: print('ERR', end=';')


class B(A):
    c = 4

    def h1(self): print(5, end=';')
    def _h2(self): print(6, end=';')

try: B().main()
except: print('ERR', end=';')
code

8. Що буде виведено за виконання?
code
def f(arg, n):
    arg[13] = 6
    n = 6
    arg = list(arg.keys())

try:
    n = 3
    d = {38:0, 11:3, 70:5, 11:8}
    f(d, n)
    print(n, end=';')
    for x, y in d.items():
        print(x, y, sep=';', end=';')
except:
    print('ERR')
code

9. Що буде виведено за виконання?
code
class A:
    a = 1
    c = 2

    def __init__(self):
        c = 3

class B(A):
    a = 4

    def __init__(self):
        super().__init__()
        self.b = 7

obj = B()
obj.a = obj.c + 10
B.a = 13

try: print(obj.a, end= ';')
except: print('ERR', end=';')

try: print(A.a, end= ';')
except: print('ERR', end=';')

try: print(B.a, end= ';')
except: print('ERR', end=';')
code

10. 
Для графа на рисунку з вершини 2 запущено алгоритм пошуку в глибину, що будує кістякове дерево. Списки суміжності вершин переглядаються алгоритмом за зростанням міток вершин.\par
\begin{center}
	\adjustimage{max size={0.9\linewidth}{0.9\paperheight}}
	{../10/Traversals/195.png}
\end{center}
\par Які з наведених ребер входять до складу отриманого кістякового дерева?\par
1) (0, 2)\par
2) (1, 6)\par
3) (1, 5)\par
4) (1, 3)\par
5) (2, 5)\par
6) (2, 6)\par
{\vskip 15pt} У формі вибрати номери відповідних ребер.\par
%answer:1 3 4
\end {large}}
%end
{\smallskip \begin {small} Затверджено на засіданні кафедри теоретичної кібернетики, протокол \No 12 від   19.05.2020 р.\par\smallskip\smallskip
{\parbox {6.5 cm}{Зав. кафедри \dotfill Ю.В.~Крак}}\hspace {0.7cm}{\hbox to 6.5 cm{ Екзаменатор \dotfill Т.О.~Карнаух}} \end{small}}
%begin
{{\begin {small}\par
{ \center  Київський національний університет імені Тараса Шевченка \par}
{\parbox {5 cm}{Освітня програма \hfill}{\parbox {7 cm}{Прикладна математика, Системний аналіз \hfill}}\parbox {6 cm}{\hfill Семестр 2}}\par
{\parbox {5 cm} {Навчальний предмет \hfill}\parbox {7 cm}{Програмування \hfill}\parbox {6cm}{\hfill}\par}\vspace{-7pt} \end{small}}{\center Екзаменаційний білет \No \textbf 
{ 267 }\par}}
{
\begin {large}
1. Що буде виведено за виконання?
code
try:
    res = "5">"False"
    if isinstance(res, bool):
        print(f'bool;{res}')
    elif isinstance(res, int):
        print(f'int;{res}')
    elif isinstance(res, float):
        print(f'float;{res:.2f}')
    else:
        print('???')
except:
    print('ERR')
code

2. Що буде виведено за виконання?
code
a = 2
b = 9
c = 3

def f(b):
    global c
    a -= 4
    b = 3
    c = 1
    return a + b + c

a = 1
b = 5
c = 8

try:
    print(f(b), a, b, c, sep=';')
except:
    print('ERR', a, b, c, sep=';')
code

3. Що буде виведено за виконання?
code
try:
    k = 6
    while k<=10:
        k += 1
    print(k, end=';')
except:
    print('ERR')
code

4. Що буде виведено за виконання?
code
try:
    for d in range(-8, -13, -1):
        if d>=4:
            continue
        print(d, end=';')
    else:
        print(d,end=';')
    print(d)
except:
    print('ERR')
code

5. Що буде виведено за виконання?\par (Дійсні значення, якщо наявні, в цьому завданні будуть виводитись у форматі з фіксованою крапкою та, можливо, з доволі великою кількістю десяткових розрядів. У відповіді для дійсних значень у ЦЬОМУ завданні достатньо вказати один десятковий розряд після крапки, відкинувши решту розрядів без округлення. Наприклад, замість 13.66666666666667 достатньо вказати 13.6, замість 25.23 достатньо вказати 25.2)
code
try:
    print(0, end=';')
    print(3/0, end=';')
    print(4, end=';')
except KeyboardInterrupt:
    print(5, end=';')
except TypeError:
    print(9, end=';')
except:
    print('ERR', end=';')
else:
    print(1, end=';')
finally:
    print(6)
code

6. Що буде виведено за виконання?
code
def f(n):
    if n<=9:
        print(n%10,end='')
    else:
        f(n//10)

f(543216)
code

7. Що буде виведено за виконання?
code
class A:
    a = 1

    def _g1(self): print(2, end=';')
    def __g2(self): print(3, end=';')

    def main(self):
        try: print(self.a, end=';')
        except: print('ERR', end=';')

        try: self._g1()
        except: print('ERR', end=';')

        try: self.__g2()
        except: print('ERR', end=';')


class B(A):
    a = 4

    def _g1(self): print(5, end=';')
    def __g2(self): print(6, end=';')

try: B().main()
except: print('ERR', end=';')
code

8. Що буде виведено за виконання?
code
def f(arg, n):
    arg[16] = 0
    n *= 3

try:
    n = 4
    t = {82:9, 66:0, 66:9}
    f(t, n)
    print(n, end=';')
    for x, y in t.items():
        print(x, y, sep=';', end=';')
except:
    print('ERR')
code

9. Що буде виведено за виконання?
code
class A:
    c = 1
    b = 2

    def __init__(self):
        c = 3

class B(A):
    b = 4

    def __init__(self):
        super().__init__()
        self.a = 7

obj = B()
obj.b = obj.a + 10
B.b = 13

try: print(obj.b, end= ';')
except: print('ERR', end=';')

try: print(A.b, end= ';')
except: print('ERR', end=';')

try: print(B.b, end= ';')
except: print('ERR', end=';')
code

10. 
Для графа на рисунку з вершини 3 запущено алгоритм пошуку в глибину, що будує кістякове дерево. Списки суміжності вершин переглядаються алгоритмом за зростанням міток вершин.\par
\begin{center}
	\adjustimage{max size={0.9\linewidth}{0.9\paperheight}}
	{../10/Traversals/56.png}
\end{center}
\par Які з наведених ребер входять до складу отриманого кістякового дерева?\par
1) (2, 6)\par
2) (1, 2)\par
3) (0, 1)\par
4) (1, 4)\par
5) (0, 6)\par
6) (3, 4)\par
{\vskip 15pt} У формі вибрати номери відповідних ребер.\par
%answer:1 2 3 6
\end {large}}
%end
{\smallskip \begin {small} Затверджено на засіданні кафедри теоретичної кібернетики, протокол \No 12 від   19.05.2020 р.\par\smallskip\smallskip
{\parbox {6.5 cm}{Зав. кафедри \dotfill Ю.В.~Крак}}\hspace {0.7cm}{\hbox to 6.5 cm{ Екзаменатор \dotfill Т.О.~Карнаух}} \end{small}}
%begin
{{\begin {small}\par
{ \center  Київський національний університет імені Тараса Шевченка \par}
{\parbox {5 cm}{Освітня програма \hfill}{\parbox {7 cm}{Прикладна математика, Системний аналіз \hfill}}\parbox {6 cm}{\hfill Семестр 2}}\par
{\parbox {5 cm} {Навчальний предмет \hfill}\parbox {7 cm}{Програмування \hfill}\parbox {6cm}{\hfill}\par}\vspace{-7pt} \end{small}}{\center Екзаменаційний білет \No \textbf 
{ 268 }\par}}
{
\begin {large}
1. Що буде виведено за виконання?
code
try:
    res = -2**-2-1
    if isinstance(res, bool):
        print(f'bool;{res}')
    elif isinstance(res, int):
        print(f'int;{res}')
    elif isinstance(res, float):
        print(f'float;{res:.2f}')
    else:
        print('???')
except:
    print('ERR')
code

2. Що буде виведено за виконання?
code
a = 6
b = 5
c = 0

def f(a):
    global c
    a = 1
    b = 3
    c = 2
    return a + b + c

a = 3
b = 4
c = 9

try:
    print(f(a), a, b, c, sep=';')
except:
    print('ERR', a, b, c, sep=';')
code

3. Що буде виведено за виконання?
code
try:
    m = 13
    while m>=9:
        m += -3
    print(m, end=';')
except:
    print('ERR')
code

4. Що буде виведено за виконання?
code
try:
    for c in range(-2, 4, -3):
        if c>5:
            continue
        print(c, end=';')
        if c>=6:
            break
    else:
        print(c,end=';')
    print(c)
except:
    print('ERR')
code

5. Що буде виведено за виконання?\par (Дійсні значення, якщо наявні, в цьому завданні будуть виводитись у форматі з фіксованою крапкою та, можливо, з доволі великою кількістю десяткових розрядів. У відповіді для дійсних значень у ЦЬОМУ завданні достатньо вказати один десятковий розряд після крапки, відкинувши решту розрядів без округлення. Наприклад, замість 13.66666666666667 достатньо вказати 13.6, замість 25.23 достатньо вказати 25.2)
code
try:
    print(2, end=';')
    print(int('b2'), end=';')
    print(8, end=';')
except ValueError:
    print(5, end=';')
except ZeroDivisionError:
    print(4, end=';')
except:
    print('ERR', end=';')
else:
    print(6, end=';')
finally:
    print(3)
code

6. Що буде виведено за виконання?
code
def f(s):
    s1 = s
    for i in range(len(s)):
        s1[i] += 1
    return s1

s = [1, 2, 3]
f(s)
print(s)
code

7. Що буде виведено за виконання?
code
class A:
    d = 1

    def _g1(self): print(2, end=';')
    def g2(self): print(3, end=';')

    def main(self):
        try: print(self.d, end=';')
        except: print('ERR', end=';')

        try: self._g1()
        except: print('ERR', end=';')

        try: self.g2()
        except: print('ERR', end=';')


class B(A):
    d = 4

    def _g1(self): print(5, end=';')
    def g2(self): print(6, end=';')

try: B().main()
except: print('ERR', end=';')
code

8. Що буде виведено за виконання?
code
def f(arg, n):
    n = 3
    tmp = arg
    tmp[-9:4] = (14)
    return len(tmp)>3

try:
    h = [10,11,12,13,14]
    n = 4
    f(h, n)
    print(n, end=';')
    for el in h:
        print(el, end=';')
except:
    print('ERR')
code

9. Що буде виведено за виконання?
code
class A:
    b = 1
    c = 2

    def __init__(self):
        c = 3

class B(A):
    b = 4

    def __init__(self):
        super().__init__()
        self.a = 7

obj = B()
obj.b = obj.a + 10
B.b = 13

try: print(obj.b, end= ';')
except: print('ERR', end=';')

try: print(A.b, end= ';')
except: print('ERR', end=';')

try: print(B.b, end= ';')
except: print('ERR', end=';')
code

10. 
Для графа на рисунку з вершини 0 запущено алгоритм пошуку в глибину, що будує кістякове дерево. Списки суміжності вершин переглядаються алгоритмом за зростанням міток вершин.\par
\begin{center}
	\adjustimage{max size={0.9\linewidth}{0.9\paperheight}}
	{../10/Traversals/139.png}
\end{center}
\par Які з наведених ребер входять до складу отриманого кістякового дерева?\par
1) (4, 5)\par
2) (4, 6)\par
3) (1, 5)\par
4) (0, 4)\par
5) (1, 6)\par
6) (3, 4)\par
{\vskip 15pt} У формі вибрати номери відповідних ребер.\par
%answer:1 3 6
\end {large}}
%end
{\smallskip \begin {small} Затверджено на засіданні кафедри теоретичної кібернетики, протокол \No 12 від   19.05.2020 р.\par\smallskip\smallskip
{\parbox {6.5 cm}{Зав. кафедри \dotfill Ю.В.~Крак}}\hspace {0.7cm}{\hbox to 6.5 cm{ Екзаменатор \dotfill Т.О.~Карнаух}} \end{small}}
%begin
{{\begin {small}\par
{ \center  Київський національний університет імені Тараса Шевченка \par}
{\parbox {5 cm}{Освітня програма \hfill}{\parbox {7 cm}{Прикладна математика, Системний аналіз \hfill}}\parbox {6 cm}{\hfill Семестр 2}}\par
{\parbox {5 cm} {Навчальний предмет \hfill}\parbox {7 cm}{Програмування \hfill}\parbox {6cm}{\hfill}\par}\vspace{-7pt} \end{small}}{\center Екзаменаційний білет \No \textbf 
{ 269 }\par}}
{
\begin {large}
1. Що буде виведено за виконання?
code
def f(n):
    print('f', end='')
    return n==1

try:
    print(f(-8) or f(1))
except:
    print('ERR')
code

2. Що буде виведено за виконання?
code
a = 0
b = 2
c = 6

def g(b):
    a = 3
    b = 4
    c = 2
    return a + b + c

a = 5
b = 8
c = 9

try:
    print(g(a), a, b, c, sep=';')
except:
    print('ERR', a, b, c, sep=';')
code

3. Що буде виведено за виконання?
code
try:
    k = 11
    while k>0:
        k += -2
    print(k, end=';')
except:
    print('ERR')
code

4. Що буде виведено за виконання?
code
try:
    for e in range(-15, 14, -2):
        if e>=7:
            continue
        print(e, end=';')
    else:
        print(e,end=';')
    print(e)
except:
    print('ERR')
code

5. Що буде виведено за виконання?\par (Дійсні значення, якщо наявні, в цьому завданні будуть виводитись у форматі з фіксованою крапкою та, можливо, з доволі великою кількістю десяткових розрядів. У відповіді для дійсних значень у ЦЬОМУ завданні достатньо вказати один десятковий розряд після крапки, відкинувши решту розрядів без округлення. Наприклад, замість 13.66666666666667 достатньо вказати 13.6, замість 25.23 достатньо вказати 25.2)
code
try:
    print(7, end=';')
    print(9j<=1j, end=';')
    print(1, end=';')
except ValueError:
    print(0, end=';')
except TypeError:
    print(8, end=';')
except:
    print('ERR', end=';')
else:
    print(6, end=';')
finally:
    print(0)
code

6. Що буде виведено за виконання?
code
def f(s):
    for i in range(len(s)):
        s[i] += 1
    return
 
s = [1, 2, 3]
f(s)
print(s)
code

7. Що буде виведено за виконання?
code
class A:
    d = 1

    def _f1(self): print(2, end=';')
    def _f2(self): print(3, end=';')

    def main(self):
        try: print(self.d, end=';')
        except: print('ERR', end=';')

        try: self._f1()
        except: print('ERR', end=';')

        try: self._f2()
        except: print('ERR', end=';')


class B(A):
    d = 4

    def _f1(self): print(5, end=';')
    def _f2(self): print(6, end=';')

try: B().main()
except: print('ERR', end=';')
code

8. Що буде виведено за виконання?
code
def f(arg, n):
    arg[24] = 6
    n -= 2

try:
    n = 5
    s = {17:8, 24:7, 24:1}
    f(s, n)
    print(n, end=';')
    for x in s.items():
        print(*x, sep=';', end=';')
except:
    print('ERR')
code

9. Що буде виведено за виконання?
code
class A:
    c = 1
    a = 2

    def __init__(self):
        c = 3

class B(A):
    c = 4

    def __init__(self):
        super().__init__()
        self.b = 7

obj = B()
obj.c = obj.a + 10
B.c = 13

try: print(obj.c, end= ';')
except: print('ERR', end=';')

try: print(A.c, end= ';')
except: print('ERR', end=';')

try: print(B.c, end= ';')
except: print('ERR', end=';')
code

10. 
Для графа на рисунку з вершини 4 запущено алгоритм пошуку в глибину, що будує кістякове дерево. Списки суміжності вершин переглядаються алгоритмом за зростанням міток вершин.\par
\begin{center}
	\adjustimage{max size={0.9\linewidth}{0.9\paperheight}}
	{../10/Traversals/80.png}
\end{center}
\par Які з наведених ребер входять до складу отриманого кістякового дерева?\par
1) (2, 3)\par
2) (5, 6)\par
3) (0, 1)\par
4) (1, 5)\par
5) (2, 5)\par
6) (3, 4)\par
{\vskip 15pt} У формі вибрати номери відповідних ребер.\par
%answer:1 2 3 4 6
\end {large}}
%end
{\smallskip \begin {small} Затверджено на засіданні кафедри теоретичної кібернетики, протокол \No 12 від   19.05.2020 р.\par\smallskip\smallskip
{\parbox {6.5 cm}{Зав. кафедри \dotfill Ю.В.~Крак}}\hspace {0.7cm}{\hbox to 6.5 cm{ Екзаменатор \dotfill Т.О.~Карнаух}} \end{small}}
%begin
{{\begin {small}\par
{ \center  Київський національний університет імені Тараса Шевченка \par}
{\parbox {5 cm}{Освітня програма \hfill}{\parbox {7 cm}{Прикладна математика, Системний аналіз \hfill}}\parbox {6 cm}{\hfill Семестр 2}}\par
{\parbox {5 cm} {Навчальний предмет \hfill}\parbox {7 cm}{Програмування \hfill}\parbox {6cm}{\hfill}\par}\vspace{-7pt} \end{small}}{\center Екзаменаційний білет \No \textbf 
{ 270 }\par}}
{
\begin {large}
1. Що буде виведено за виконання?
code
try:
    res = 6//3//6*5
    if isinstance(res, bool):
        print(f'bool;{res}')
    elif isinstance(res, int):
        print(f'int;{res}')
    elif isinstance(res, float):
        print(f'float;{res:.2f}')
    else:
        print('???')
except:
    print('ERR')
code

2. Що буде виведено за виконання?
code
a = 1
b = 3
c = 4

def h(b):
    global c
    a = 5
    b = 1
    c = 2
    return a + b + c

a = 7
b = 4
c = 1

try:
    print(h(b), a, b, c, sep=';')
except:
    print('ERR', a, b, c, sep=';')
code

3. Що буде виведено за виконання?
code
try:
    j = 15
    while j>4:
        j += -3
    print(j, end=';')
except:
    print('ERR')
code

4. Що буде виведено за виконання?
code
try:
    for a in range(-4, -7, 3):
        if a<6:
            break
        print(a, end=';');a = 9
    else:
        print('end',end=';')
    print(a)
except:
    print('ERR')
code

5. Що буде виведено за виконання?\par (Дійсні значення, якщо наявні, в цьому завданні будуть виводитись у форматі з фіксованою крапкою та, можливо, з доволі великою кількістю десяткових розрядів. У відповіді для дійсних значень у ЦЬОМУ завданні достатньо вказати один десятковий розряд після крапки, відкинувши решту розрядів без округлення. Наприклад, замість 13.66666666666667 достатньо вказати 13.6, замість 25.23 достатньо вказати 25.2)
code
try:
    print(1, end=';')
    print(int('5'), end=';')
    print(9, end=';')
except KeyboardInterrupt:
    print(4, end=';')
except Exception:
    print(2, end=';')
except:
    print('ERR', end=';')
else:
    print(7, end=';')
finally:
    print(3)
code

6. Що буде виведено за виконання?
code
def f(s):
    s1 = s
    for i in range(len(s)):
        s1[i] += 1
        return s1

s = [1, 2, 3]
f(s)
print(s)
code

7. Що буде виведено за виконання?
code
class A:
    b = 1

    def __h1(self): print(2, end=';')
    def _h2(self): print(3, end=';')

    def main(self):
        try: print(self.b, end=';')
        except: print('ERR', end=';')

        try: self.__h1()
        except: print('ERR', end=';')

        try: self._h2()
        except: print('ERR', end=';')


class B(A):
    b = 4

    def __h1(self): print(5, end=';')
    def _h2(self): print(6, end=';')

try: B().main()
except: print('ERR', end=';')
code

8. Що буде виведено за виконання?
code
def f(arg, n):
    arg[24] = 2
    n += -2

try:
    n = 3
    d = {24:2, 27:3, 46:1}
    f(d, n)
    print(n, end=';')
    for x in d.items():
        print(x, sep=';', end=';')
except:
    print('ERR')
code

9. Що буде виведено за виконання?
code
class A:
    a = 1
    b = 2

    def __init__(self):
        b = 3

class B(A):
    b = 4

    def __init__(self):
        super().__init__()
        self.c = 7

obj = B()
obj.c = obj.b + 10
B.c = 13

try: print(obj.c, end= ';')
except: print('ERR', end=';')

try: print(A.c, end= ';')
except: print('ERR', end=';')

try: print(B.c, end= ';')
except: print('ERR', end=';')
code

10. 
Для графа на рисунку з вершини 3 запущено алгоритм пошуку в глибину, що будує кістякове дерево. Списки суміжності вершин переглядаються алгоритмом за зростанням міток вершин.\par
\begin{center}
	\adjustimage{max size={0.9\linewidth}{0.9\paperheight}}
	{../10/Traversals/286.png}
\end{center}
\par Які з наведених ребер входять до складу отриманого кістякового дерева?\par
1) (0, 5)\par
2) (3, 6)\par
3) (3, 5)\par
4) (4, 6)\par
5) (0, 4)\par
6) (2, 5)\par
{\vskip 15pt} У формі вибрати номери відповідних ребер.\par
%answer:1 4 5
\end {large}}
%end
{\smallskip \begin {small} Затверджено на засіданні кафедри теоретичної кібернетики, протокол \No 12 від   19.05.2020 р.\par\smallskip\smallskip
{\parbox {6.5 cm}{Зав. кафедри \dotfill Ю.В.~Крак}}\hspace {0.7cm}{\hbox to 6.5 cm{ Екзаменатор \dotfill Т.О.~Карнаух}} \end{small}}
%begin
{{\begin {small}\par
{ \center  Київський національний університет імені Тараса Шевченка \par}
{\parbox {5 cm}{Освітня програма \hfill}{\parbox {7 cm}{Прикладна математика, Системний аналіз \hfill}}\parbox {6 cm}{\hfill Семестр 2}}\par
{\parbox {5 cm} {Навчальний предмет \hfill}\parbox {7 cm}{Програмування \hfill}\parbox {6cm}{\hfill}\par}\vspace{-7pt} \end{small}}{\center Екзаменаційний білет \No \textbf 
{ 271 }\par}}
{
\begin {large}
1. Що буде виведено за виконання?
code
try:
    res = 4/7/7*7
    if isinstance(res, bool):
        print(f'bool;{res}')
    elif isinstance(res, int):
        print(f'int;{res}')
    elif isinstance(res, float):
        print(f'float;{res:.2f}')
    else:
        print('???')
except:
    print('ERR')
code

2. Що буде виведено за виконання?
code
a = 4
b = 0
c = 1

def g(a):
    a *= 5
    b = 2
    c = 5
    return a + b + c

a = 0
b = 4
c = 9

try:
    print(g(b), a, b, c, sep=';')
except:
    print('ERR', a, b, c, sep=';')
code

3. Що буде виведено за виконання?
code
try:
    n = 4
    while n<11:
        n += 2
    print(n, end=';')
except:
    print('ERR')
code

4. Що буде виведено за виконання?
code
try:
    for b in range(9, -6, 3):
        if b>3:
            continue
        print(b, end=';')
    else:
        print(b,end=';')
    print(b)
except:
    print('ERR')
code

5. Що буде виведено за виконання?\par (Дійсні значення, якщо наявні, в цьому завданні будуть виводитись у форматі з фіксованою крапкою та, можливо, з доволі великою кількістю десяткових розрядів. У відповіді для дійсних значень у ЦЬОМУ завданні достатньо вказати один десятковий розряд після крапки, відкинувши решту розрядів без округлення. Наприклад, замість 13.66666666666667 достатньо вказати 13.6, замість 25.23 достатньо вказати 25.2)
code
try:
    print(8, end=';')
    print(4j!=8j, end=';')
    print(6, end=';')
except TypeError:
    print(9, end=';')
except ZeroDivisionError:
    print(2, end=';')
except:
    print('ERR', end=';')
else:
    print(3, end=';')
finally:
    print(7)
code

6. Що буде виведено за виконання?
code
def f(n):
    if n<=9:
        print(n%10, end='')
    else:
        f(n//10)

f(123456)
code

7. Що буде виведено за виконання?
code
class A:
    c = 1

    def _f1(self): print(2, end=';')
    def _f2(self): print(3, end=';')

    def main(self):
        try: print(self.c, end=';')
        except: print('ERR', end=';')

        try: self._f1()
        except: print('ERR', end=';')

        try: self._f2()
        except: print('ERR', end=';')


class B(A):
    c = 4

    def _f1(self): print(5, end=';')
    def _f2(self): print(6, end=';')

try: B().main()
except: print('ERR', end=';')
code

8. Що буде виведено за виконання?
code
def f(arg, n):
    n = 9
    tmp = arg
    tmp[-8:0] = ()
    return len(tmp)>3

try:
    b = ['0','1','2','3','4']
    n = 3
    f(b, n)
    print(n, end=';')
    for el in b:
        print(el, end=';')
except:
    print('ERR')
code

9. Що буде виведено за виконання?
code
class A:
    b = 1
    c = 2

    def __init__(self):
        c = 3

class B(A):
    b = 4

    def __init__(self):
        super().__init__()
        self.a = 7

obj = B()
obj.a = obj.b + 10
B.a = 13

try: print(obj.a, end= ';')
except: print('ERR', end=';')

try: print(A.a, end= ';')
except: print('ERR', end=';')

try: print(B.a, end= ';')
except: print('ERR', end=';')
code

10. 
Для графа на рисунку з вершини 1 запущено алгоритм пошуку в глибину, що будує кістякове дерево. Списки суміжності вершин переглядаються алгоритмом за зростанням міток вершин.\par
\begin{center}
	\adjustimage{max size={0.9\linewidth}{0.9\paperheight}}
	{../10/Traversals/62.png}
\end{center}
\par Які з наведених ребер входять до складу отриманого кістякового дерева?\par
1) (4, 5)\par
2) (1, 5)\par
3) (0, 1)\par
4) (0, 5)\par
5) (1, 6)\par
6) (0, 4)\par
{\vskip 15pt} У формі вибрати номери відповідних ребер.\par
%answer:1 3 6
\end {large}}
%end
{\smallskip \begin {small} Затверджено на засіданні кафедри теоретичної кібернетики, протокол \No 12 від   19.05.2020 р.\par\smallskip\smallskip
{\parbox {6.5 cm}{Зав. кафедри \dotfill Ю.В.~Крак}}\hspace {0.7cm}{\hbox to 6.5 cm{ Екзаменатор \dotfill Т.О.~Карнаух}} \end{small}}
%begin
{{\begin {small}\par
{ \center  Київський національний університет імені Тараса Шевченка \par}
{\parbox {5 cm}{Освітня програма \hfill}{\parbox {7 cm}{Прикладна математика, Системний аналіз \hfill}}\parbox {6 cm}{\hfill Семестр 2}}\par
{\parbox {5 cm} {Навчальний предмет \hfill}\parbox {7 cm}{Програмування \hfill}\parbox {6cm}{\hfill}\par}\vspace{-7pt} \end{small}}{\center Екзаменаційний білет \No \textbf 
{ 272 }\par}}
{
\begin {large}
1. Що буде виведено за виконання?
code
try:
    res = True<=2.0 or not 4.0==4 or 2<5
    if isinstance(res, bool):
        print(f'bool;{res}')
    elif isinstance(res, int):
        print(f'int;{res}')
    elif isinstance(res, float):
        print(f'float;{res:.2f}')
    else:
        print('???')
except:
    print('ERR')
code

2. Що буде виведено за виконання?
code
a = 9
b = 6
c = 2

def f(a):
    a = 1
    b = 3
    c = 4
    return a + b + c

a = 5
b = 8
c = 3

try:
    print(f(b), a, b, c, sep=';')
except:
    print('ERR', a, b, c, sep=';')
code

3. Що буде виведено за виконання?
code
try:
    i = 2
    while i<8:
        i += 1
    print(i, end=';')
except:
    print('ERR')
code

4. Що буде виведено за виконання?
code
try:
    for e in range(15, 15, -2):
        if e<=5:
            break
        print(e, end=';');e = 10
    else:
        print('end',end=';')
    print(e)
except:
    print('ERR')
code

5. Що буде виведено за виконання?\par (Дійсні значення, якщо наявні, в цьому завданні будуть виводитись у форматі з фіксованою крапкою та, можливо, з доволі великою кількістю десяткових розрядів. У відповіді для дійсних значень у ЦЬОМУ завданні достатньо вказати один десятковий розряд після крапки, відкинувши решту розрядів без округлення. Наприклад, замість 13.66666666666667 достатньо вказати 13.6, замість 25.23 достатньо вказати 25.2)
code
try:
    print(0, end=';')
    print(1/0, end=';')
    print(5, end=';')
except KeyboardInterrupt:
    print(5, end=';')
except Exception:
    print(2, end=';')
except:
    print('ERR', end=';')
else:
    print(1, end=';')
finally:
    print(7)
code

6. Що буде виведено за виконання?
code
def f(n):
    if n>9:
        f(n//10)
    else:
        print(n%10, end='')

f(543216)
code

7. Що буде виведено за виконання?
code
class A:
    a = 1

    def g1(self): print(2, end=';')
    def _g2(self): print(3, end=';')

    def main(self):
        try: print(self.a, end=';')
        except: print('ERR', end=';')

        try: self.g1()
        except: print('ERR', end=';')

        try: self._g2()
        except: print('ERR', end=';')


class B(A):
    a = 4

    def g1(self): print(5, end=';')
    def _g2(self): print(6, end=';')

try: B().main()
except: print('ERR', end=';')
code

8. Що буде виведено за виконання?
code
def f(arg, n):
    arg[58] = 0
    n = 8
    arg = tuple(arg.keys())

try:
    n = 3
    s = {80:2, 58:1, 76:9, 76:5}
    f(s, n)
    print(n, end=';')
    for x in s.keys():
        print(x, sep=';', end=';')
except:
    print('ERR')
code

9. Що буде виведено за виконання?
code
class A:
    b = 1
    a = 2

    def __init__(self):
        b = 3

class B(A):
    a = 4

    def __init__(self):
        super().__init__()
        self.c = 7

obj = B()
obj.c = obj.c + 10
B.c = 13

try: print(obj.c, end= ';')
except: print('ERR', end=';')

try: print(A.c, end= ';')
except: print('ERR', end=';')

try: print(B.c, end= ';')
except: print('ERR', end=';')
code

10. 
Для графа на рисунку з вершини 0 запущено алгоритм пошуку в глибину, що будує кістякове дерево. Списки суміжності вершин переглядаються алгоритмом за зростанням міток вершин.\par
\begin{center}
	\adjustimage{max size={0.9\linewidth}{0.9\paperheight}}
	{../10/Traversals/236.png}
\end{center}
\par Які з наведених ребер входять до складу отриманого кістякового дерева?\par
1) (2, 6)\par
2) (0, 5)\par
3) (3, 6)\par
4) (1, 6)\par
5) (0, 1)\par
6) (4, 5)\par
{\vskip 15pt} У формі вибрати номери відповідних ребер.\par
%answer:1 3 4 5 6
\end {large}}
%end
{\smallskip \begin {small} Затверджено на засіданні кафедри теоретичної кібернетики, протокол \No 12 від   19.05.2020 р.\par\smallskip\smallskip
{\parbox {6.5 cm}{Зав. кафедри \dotfill Ю.В.~Крак}}\hspace {0.7cm}{\hbox to 6.5 cm{ Екзаменатор \dotfill Т.О.~Карнаух}} \end{small}}
%begin
{{\begin {small}\par
{ \center  Київський національний університет імені Тараса Шевченка \par}
{\parbox {5 cm}{Освітня програма \hfill}{\parbox {7 cm}{Прикладна математика, Системний аналіз \hfill}}\parbox {6 cm}{\hfill Семестр 2}}\par
{\parbox {5 cm} {Навчальний предмет \hfill}\parbox {7 cm}{Програмування \hfill}\parbox {6cm}{\hfill}\par}\vspace{-7pt} \end{small}}{\center Екзаменаційний білет \No \textbf 
{ 273 }\par}}
{
\begin {large}
1. Що буде виведено за виконання?
code
try:
    res = 4**-0.5*-1
    if isinstance(res, bool):
        print(f'bool;{res}')
    elif isinstance(res, int):
        print(f'int;{res}')
    elif isinstance(res, float):
        print(f'float;{res:.2f}')
    else:
        print('???')
except:
    print('ERR')
code

2. Що буде виведено за виконання?
code
a = 6
b = 3
c = 2

def h(a):
    global c
    a = 3
    b += 1
    c = 4
    return a + b + c

a = 7
b = 8
c = 5

try:
    print(h(a), a, b, c, sep=';')
except:
    print('ERR', a, b, c, sep=';')
code

3. Що буде виведено за виконання?
code
try:
    m = 0
    while m<=5:
        m += 2
    print(m, end=';')
except:
    print('ERR')
code

4. Що буде виведено за виконання?
code
try:
    for f in range(-13, 7, -1):
        if f<=6:
            continue
        print(f, end=';')
    else:
        print(f,end=';')
    print(f)
except:
    print('ERR')
code

5. Що буде виведено за виконання?\par (Дійсні значення, якщо наявні, в цьому завданні будуть виводитись у форматі з фіксованою крапкою та, можливо, з доволі великою кількістю десяткових розрядів. У відповіді для дійсних значень у ЦЬОМУ завданні достатньо вказати один десятковий розряд після крапки, відкинувши решту розрядів без округлення. Наприклад, замість 13.66666666666667 достатньо вказати 13.6, замість 25.23 достатньо вказати 25.2)
code
try:
    print(6, end=';')
    print(int('9'), end=';')
    print(8, end=';')
except ValueError:
    print(0, end=';')
except KeyboardInterrupt:
    print(3, end=';')
except:
    print('ERR', end=';')
else:
    print(4, end=';')
finally:
    print(4)
code

6. Що буде виведено за виконання?
code
def f(n):
    if n>2:
        f(n//2)
    print(n%2, end='')
 
f(16)
code

7. Що буде виведено за виконання?
code
class A:
    a = 1

    def __h1(self): print(2, end=';')
    def _h2(self): print(3, end=';')

    def main(self):
        try: print(self.a, end=';')
        except: print('ERR', end=';')

        try: self.__h1()
        except: print('ERR', end=';')

        try: self._h2()
        except: print('ERR', end=';')


class B(A):
    a = 4

    def __h1(self): print(5, end=';')
    def _h2(self): print(6, end=';')

try: B().main()
except: print('ERR', end=';')
code

8. Що буде виведено за виконання?
code
def f(arg, n):
    n = 2
    tmp = arg
    del tmp[6:]
    return len(tmp)>3

try:
    e = [10,11,12,13,14]
    n = 0
    f(e, n)
    print(n, end=';')
    for el in e:
        print(el, end=';')
except:
    print('ERR')
code

9. Що буде виведено за виконання?
code
class A:
    c = 1
    a = 2

    def __init__(self):
        a = 3

class B(A):
    a = 4

    def __init__(self):
        super().__init__()
        self.b = 7

obj = B()
obj.b = obj.a + 10
B.b = 13

try: print(obj.b, end= ';')
except: print('ERR', end=';')

try: print(A.b, end= ';')
except: print('ERR', end=';')

try: print(B.b, end= ';')
except: print('ERR', end=';')
code

10. 
Для графа на рисунку з вершини 6 запущено алгоритм пошуку в глибину, що будує кістякове дерево. Списки суміжності вершин переглядаються алгоритмом за зростанням міток вершин.\par
\begin{center}
	\adjustimage{max size={0.9\linewidth}{0.9\paperheight}}
	{../10/Traversals/264.png}
\end{center}
\par Які з наведених ребер входять до складу отриманого кістякового дерева?\par
1) (0, 1)\par
2) (4, 6)\par
3) (2, 4)\par
4) (0, 2)\par
5) (2, 3)\par
6) (1, 6)\par
{\vskip 15pt} У формі вибрати номери відповідних ребер.\par
%answer:1 3 4 5 6
\end {large}}
%end
{\smallskip \begin {small} Затверджено на засіданні кафедри теоретичної кібернетики, протокол \No 12 від   19.05.2020 р.\par\smallskip\smallskip
{\parbox {6.5 cm}{Зав. кафедри \dotfill Ю.В.~Крак}}\hspace {0.7cm}{\hbox to 6.5 cm{ Екзаменатор \dotfill Т.О.~Карнаух}} \end{small}}
%begin
{{\begin {small}\par
{ \center  Київський національний університет імені Тараса Шевченка \par}
{\parbox {5 cm}{Освітня програма \hfill}{\parbox {7 cm}{Прикладна математика, Системний аналіз \hfill}}\parbox {6 cm}{\hfill Семестр 2}}\par
{\parbox {5 cm} {Навчальний предмет \hfill}\parbox {7 cm}{Програмування \hfill}\parbox {6cm}{\hfill}\par}\vspace{-7pt} \end{small}}{\center Екзаменаційний білет \No \textbf 
{ 274 }\par}}
{
\begin {large}
1. Що буде виведено за виконання?
code
def f(n):
    print('f', end='')
    return n<-4

try:
    print(f(-9) and f(2))
except:
    print('ERR')
code

2. Що буде виведено за виконання?
code
a = 3
b = 5
c = 7

def f(b):
    a -= 3
    b = 2
    c = 1
    return a + b + c

a = 2
b = 1
c = 6

try:
    print(f(a), a, b, c, sep=';')
except:
    print('ERR', a, b, c, sep=';')
code

3. Що буде виведено за виконання?
code
try:
    t = 7
    while t<=7:
        t += 1
    print(t, end=';')
except:
    print('ERR')
code

4. Що буде виведено за виконання?
code
try:
    for a in range(4, 9, -3):
        if a<8:
            continue
        print(a, end=';')
    else:
        print(a,end=';')
    print(a)
except:
    print('ERR')
code

5. Що буде виведено за виконання?\par (Дійсні значення, якщо наявні, в цьому завданні будуть виводитись у форматі з фіксованою крапкою та, можливо, з доволі великою кількістю десяткових розрядів. У відповіді для дійсних значень у ЦЬОМУ завданні достатньо вказати один десятковий розряд після крапки, відкинувши решту розрядів без округлення. Наприклад, замість 13.66666666666667 достатньо вказати 13.6, замість 25.23 достатньо вказати 25.2)
code
try:
    print(3, end=';')
    print(int('d8'), end=';')
    print(2, end=';')
except TypeError:
    print(9, end=';')
except Exception:
    print(0, end=';')
except:
    print('ERR', end=';')
else:
    print(7, end=';')
finally:
    print(6)
code

6. Що буде виведено за виконання?
code
def f(n):
    if n>9:
        f(n//10)
    print(n%10, end='')
 
f(123456)
code

7. Що буде виведено за виконання?
code
class A:
    b = 1

    def _f1(self): print(2, end=';')
    def __f2(self): print(3, end=';')

    def main(self):
        try: print(self.b, end=';')
        except: print('ERR', end=';')

        try: self._f1()
        except: print('ERR', end=';')

        try: self.__f2()
        except: print('ERR', end=';')


class B(A):
    b = 4

    def _f1(self): print(5, end=';')
    def __f2(self): print(6, end=';')

try: B().main()
except: print('ERR', end=';')
code

8. Що буде виведено за виконання?
code
def f(arg, n):
    arg[33] = 8
    n = 6
    arg = list(arg.items())

try:
    n = 4
    t = {33:3, 34:8, 83:7, 34:9}
    f(t, n)
    print(n, end=';')
    for x in t.items():
        print(x, sep=';', end=';')
except:
    print('ERR')
code

9. Що буде виведено за виконання?
code
class A:
    a = 1
    b = 2

    def __init__(self):
        a = 3

class B(A):
    a = 4

    def __init__(self):
        super().__init__()
        self.c = 7

obj = B()
obj.a = obj.b + 10
B.a = 13

try: print(obj.a, end= ';')
except: print('ERR', end=';')

try: print(A.a, end= ';')
except: print('ERR', end=';')

try: print(B.a, end= ';')
except: print('ERR', end=';')
code

10. 
Для графа на рисунку з вершини 0 запущено алгоритм пошуку в глибину, що будує кістякове дерево. Списки суміжності вершин переглядаються алгоритмом за зростанням міток вершин.\par
\begin{center}
	\adjustimage{max size={0.9\linewidth}{0.9\paperheight}}
	{../10/Traversals/15.png}
\end{center}
\par Які з наведених ребер входять до складу отриманого кістякового дерева?\par
1) (1, 5)\par
2) (0, 5)\par
3) (1, 2)\par
4) (4, 5)\par
5) (3, 6)\par
6) (2, 6)\par
{\vskip 15pt} У формі вибрати номери відповідних ребер.\par
%answer:3 4 5
\end {large}}
%end
{\smallskip \begin {small} Затверджено на засіданні кафедри теоретичної кібернетики, протокол \No 12 від   19.05.2020 р.\par\smallskip\smallskip
{\parbox {6.5 cm}{Зав. кафедри \dotfill Ю.В.~Крак}}\hspace {0.7cm}{\hbox to 6.5 cm{ Екзаменатор \dotfill Т.О.~Карнаух}} \end{small}}
%begin
{{\begin {small}\par
{ \center  Київський національний університет імені Тараса Шевченка \par}
{\parbox {5 cm}{Освітня програма \hfill}{\parbox {7 cm}{Прикладна математика, Системний аналіз \hfill}}\parbox {6 cm}{\hfill Семестр 2}}\par
{\parbox {5 cm} {Навчальний предмет \hfill}\parbox {7 cm}{Програмування \hfill}\parbox {6cm}{\hfill}\par}\vspace{-7pt} \end{small}}{\center Екзаменаційний білет \No \textbf 
{ 275 }\par}}
{
\begin {large}
1. Що буде виведено за виконання?
code
try:
    res = 5//3%9*5
    if isinstance(res, bool):
        print(f'bool;{res}')
    elif isinstance(res, int):
        print(f'int;{res}')
    elif isinstance(res, float):
        print(f'float;{res:.2f}')
    else:
        print('???')
except:
    print('ERR')
code

2. Що буде виведено за виконання?
code
a = 4
b = 9
c = 0

def g(b):
    a = 5
    b *= 4
    c = 4
    return a + b + c

a = 8
b = 7
c = 3

try:
    print(g(a), a, b, c, sep=';')
except:
    print('ERR', a, b, c, sep=';')
code

3. Що буде виведено за виконання?
code
try:
    t = 12
    while t>=8:
        t += -2
    print(t, end=';')
except:
    print('ERR')
code

4. Що буде виведено за виконання?
code
try:
    for d in range(-7, -1, 2):
        if d>=7:
            continue
        print(d, end=';')
        if d>3:
            break
    else:
        print(d,end=';')
    print(d)
except:
    print('ERR')
code

5. Що буде виведено за виконання?\par (Дійсні значення, якщо наявні, в цьому завданні будуть виводитись у форматі з фіксованою крапкою та, можливо, з доволі великою кількістю десяткових розрядів. У відповіді для дійсних значень у ЦЬОМУ завданні достатньо вказати один десятковий розряд після крапки, відкинувши решту розрядів без округлення. Наприклад, замість 13.66666666666667 достатньо вказати 13.6, замість 25.23 достатньо вказати 25.2)
code
try:
    print(5, end=';')
    print(1%3, end=';')
    print(7, end=';')
except ZeroDivisionError:
    print(8, end=';')
except ValueError:
    print(2, end=';')
except:
    print('ERR', end=';')
else:
    print(1, end=';')
finally:
    print(6)
code

6. Що буде виведено за виконання?
code
def f(n):
    if n>9:
        f(n//10)
    print(n%10, end='')
 
f(543216)
code

7. Що буде виведено за виконання?
code
class A:
    d = 1

    def _g1(self): print(2, end=';')
    def g2(self): print(3, end=';')

    def main(self):
        try: print(self.d, end=';')
        except: print('ERR', end=';')

        try: self._g1()
        except: print('ERR', end=';')

        try: self.g2()
        except: print('ERR', end=';')


class B(A):
    d = 4

    def _g1(self): print(5, end=';')
    def g2(self): print(6, end=';')

try: B().main()
except: print('ERR', end=';')
code

8. Що буде виведено за виконання?
code
def f(arg, n):
    n = 6
    tmp = arg
    tmp[-7:-4] = [1,0]
    return len(tmp)>3

try:
    c = ['0','1','2','3','4']
    n = 4
    f(c, n)
    print(n, end=';')
    for el in c:
        print(el, end=';')
except:
    print('ERR')
code

9. Що буде виведено за виконання?
code
class A:
    c = 1
    b = 2

    def __init__(self):
        c = 3

class B(A):
    c = 4

    def __init__(self):
        super().__init__()
        self.a = 7

obj = B()
obj.c = obj.b + 10
B.c = 13

try: print(obj.c, end= ';')
except: print('ERR', end=';')

try: print(A.c, end= ';')
except: print('ERR', end=';')

try: print(B.c, end= ';')
except: print('ERR', end=';')
code

10. 
Для графа на рисунку з вершини 4 запущено алгоритм пошуку в глибину, що будує кістякове дерево. Списки суміжності вершин переглядаються алгоритмом за зростанням міток вершин.\par
\begin{center}
	\adjustimage{max size={0.9\linewidth}{0.9\paperheight}}
	{../10/Traversals/19.png}
\end{center}
\par Які з наведених ребер входять до складу отриманого кістякового дерева?\par
1) (2, 5)\par
2) (1, 2)\par
3) (2, 6)\par
4) (5, 6)\par
5) (2, 3)\par
6) (3, 4)\par
{\vskip 15pt} У формі вибрати номери відповідних ребер.\par
%answer:2 5
\end {large}}
%end
{\smallskip \begin {small} Затверджено на засіданні кафедри теоретичної кібернетики, протокол \No 12 від   19.05.2020 р.\par\smallskip\smallskip
{\parbox {6.5 cm}{Зав. кафедри \dotfill Ю.В.~Крак}}\hspace {0.7cm}{\hbox to 6.5 cm{ Екзаменатор \dotfill Т.О.~Карнаух}} \end{small}}
%begin
{{\begin {small}\par
{ \center  Київський національний університет імені Тараса Шевченка \par}
{\parbox {5 cm}{Освітня програма \hfill}{\parbox {7 cm}{Прикладна математика, Системний аналіз \hfill}}\parbox {6 cm}{\hfill Семестр 2}}\par
{\parbox {5 cm} {Навчальний предмет \hfill}\parbox {7 cm}{Програмування \hfill}\parbox {6cm}{\hfill}\par}\vspace{-7pt} \end{small}}{\center Екзаменаційний білет \No \textbf 
{ 276 }\par}}
{
\begin {large}
1. Що буде виведено за виконання?
code
try:
    res = 9/9//4*9
    if isinstance(res, bool):
        print(f'bool;{res}')
    elif isinstance(res, int):
        print(f'int;{res}')
    elif isinstance(res, float):
        print(f'float;{res:.2f}')
    else:
        print('???')
except:
    print('ERR')
code

2. Що буде виведено за виконання?
code
a = 7
b = 0
c = 6

def g(b):
    global c
    a = 5
    b = 2
    c = 3
    return a + b + c

a = 7
b = 3
c = 2

try:
    print(g(b), a, b, c, sep=';')
except:
    print('ERR', a, b, c, sep=';')
code

3. Що буде виведено за виконання?
code
try:
    j = 3
    while j<10:
        j += 2
    print(j, end=';')
except:
    print('ERR')
code

4. Що буде виведено за виконання?
code
try:
    for c in range(11, 3, -1):
        if c>4:
            break
        print(c, end=';');c = 3
    else:
        print(c,end=';')
    print(c)
except:
    print('ERR')
code

5. Що буде виведено за виконання?\par (Дійсні значення, якщо наявні, в цьому завданні будуть виводитись у форматі з фіксованою крапкою та, можливо, з доволі великою кількістю десяткових розрядів. У відповіді для дійсних значень у ЦЬОМУ завданні достатньо вказати один десятковий розряд після крапки, відкинувши решту розрядів без округлення. Наприклад, замість 13.66666666666667 достатньо вказати 13.6, замість 25.23 достатньо вказати 25.2)
code
try:
    print(3, end=';')
    print(5//0.0, end=';')
    print(0, end=';')
except TypeError:
    print(4, end=';')
except KeyboardInterrupt:
    print(9, end=';')
except:
    print('ERR', end=';')
else:
    print(8, end=';')
finally:
    print(0)
code

6. Що буде виведено за виконання?
code
def f(s):
    for el in s:
        el += 1
    return s
 
s = [1, 2, 3]
s = f(s)
print(s)
code

7. Що буде виведено за виконання?
code
class A:
    c = 1

    def _h1(self): print(2, end=';')
    def h2(self): print(3, end=';')

    def main(self):
        try: print(self.c, end=';')
        except: print('ERR', end=';')

        try: self._h1()
        except: print('ERR', end=';')

        try: self.h2()
        except: print('ERR', end=';')


class B(A):
    c = 4

    def _h1(self): print(5, end=';')
    def h2(self): print(6, end=';')

try: B().main()
except: print('ERR', end=';')
code

8. Що буде виведено за виконання?
code
def f(arg, n):
    n = 1
    tmp = arg
    del tmp[:-2]
    return len(tmp)>3

try:
    d = [10,11,12,13,14]
    n = 8
    f(d, n)
    print(n, end=';')
    for el in d:
        print(el, end=';')
except:
    print('ERR')
code

9. Що буде виведено за виконання?
code
class A:
    a = 1
    c = 2

    def __init__(self):
        c = 3

class B(A):
    c = 4

    def __init__(self):
        super().__init__()
        self.b = 7

obj = B()
obj.c = obj.a + 10
B.c = 13

try: print(obj.c, end= ';')
except: print('ERR', end=';')

try: print(A.c, end= ';')
except: print('ERR', end=';')

try: print(B.c, end= ';')
except: print('ERR', end=';')
code

10. 
Для графа на рисунку з вершини 0 запущено алгоритм пошуку в глибину, що будує кістякове дерево. Списки суміжності вершин переглядаються алгоритмом за зростанням міток вершин.\par
\begin{center}
	\adjustimage{max size={0.9\linewidth}{0.9\paperheight}}
	{../10/Traversals/241.png}
\end{center}
\par Які з наведених ребер входять до складу отриманого кістякового дерева?\par
1) (1, 2)\par
2) (1, 4)\par
3) (4, 6)\par
4) (2, 3)\par
5) (0, 1)\par
6) (2, 6)\par
{\vskip 15pt} У формі вибрати номери відповідних ребер.\par
%answer:1 3 4 5 6
\end {large}}
%end
{\smallskip \begin {small} Затверджено на засіданні кафедри теоретичної кібернетики, протокол \No 12 від   19.05.2020 р.\par\smallskip\smallskip
{\parbox {6.5 cm}{Зав. кафедри \dotfill Ю.В.~Крак}}\hspace {0.7cm}{\hbox to 6.5 cm{ Екзаменатор \dotfill Т.О.~Карнаух}} \end{small}}
%begin
{{\begin {small}\par
{ \center  Київський національний університет імені Тараса Шевченка \par}
{\parbox {5 cm}{Освітня програма \hfill}{\parbox {7 cm}{Прикладна математика, Системний аналіз \hfill}}\parbox {6 cm}{\hfill Семестр 2}}\par
{\parbox {5 cm} {Навчальний предмет \hfill}\parbox {7 cm}{Програмування \hfill}\parbox {6cm}{\hfill}\par}\vspace{-7pt} \end{small}}{\center Екзаменаційний білет \No \textbf 
{ 277 }\par}}
{
\begin {large}
1. Що буде виведено за виконання?
code
try:
    res = 9**-1.0*True
    if isinstance(res, bool):
        print(f'bool;{res}')
    elif isinstance(res, int):
        print(f'int;{res}')
    elif isinstance(res, float):
        print(f'float;{res:.2f}')
    else:
        print('???')
except:
    print('ERR')
code

2. Що буде виведено за виконання?
code
a = 1
b = 4
c = 6

def h(a):
    global c
    a += 3
    b = 5
    c = 2
    return a + b + c

a = 8
b = 9
c = 2

try:
    print(h(b), a, b, c, sep=';')
except:
    print('ERR', a, b, c, sep=';')
code

3. Що буде виведено за виконання?
code
try:
    i = 9
    while i>=5:
        i += -3
    print(i, end=';')
except:
    print('ERR')
code

4. Що буде виведено за виконання?
code
try:
    for f in range(-14, -10, -3):
        if f<5:
            continue
        print(f, end=';')
    else:
        print(f,end=';')
    print(f)
except:
    print('ERR')
code

5. Що буде виведено за виконання?\par (Дійсні значення, якщо наявні, в цьому завданні будуть виводитись у форматі з фіксованою крапкою та, можливо, з доволі великою кількістю десяткових розрядів. У відповіді для дійсних значень у ЦЬОМУ завданні достатньо вказати один десятковий розряд після крапки, відкинувши решту розрядів без округлення. Наприклад, замість 13.66666666666667 достатньо вказати 13.6, замість 25.23 достатньо вказати 25.2)
code
try:
    print(7, end=';')
    print(3//1, end=';')
    print(2, end=';')
except Exception:
    print(9, end=';')
except ZeroDivisionError:
    print(1, end=';')
except:
    print('ERR', end=';')
else:
    print(5, end=';')
finally:
    print(4)
code

6. Що буде виведено за виконання?
code
def f(n):
    if n<=9:
        print(n%10, end='')
    else:
        f(n//10)

f(987651)
code

7. Що буде виведено за виконання?
code
class A:
    b = 1

    def _g1(self): print(2, end=';')
    def _g2(self): print(3, end=';')

    def main(self):
        try: print(self.b, end=';')
        except: print('ERR', end=';')

        try: self._g1()
        except: print('ERR', end=';')

        try: self._g2()
        except: print('ERR', end=';')


class B(A):
    b = 4

    def _g1(self): print(5, end=';')
    def _g2(self): print(6, end=';')

try: B().main()
except: print('ERR', end=';')
code

8. Що буде виведено за виконання?
code
def f(arg, n):
    arg[68] = 9
    n -= 3

try:
    n = 5
    s = {28:2, 68:1, 70:4, 70:8}
    f(s, n)
    print(n, end=';')
    for x in s :
        print(x, sep=';', end=';')
except:
    print('ERR')
code

9. Що буде виведено за виконання?
code
class A:
    b = 1
    a = 2

    def __init__(self):
        a = 3

class B(A):
    a = 4

    def __init__(self):
        super().__init__()
        self.c = 7

obj = B()
obj.c = obj.b + 10
B.c = 13

try: print(obj.c, end= ';')
except: print('ERR', end=';')

try: print(A.c, end= ';')
except: print('ERR', end=';')

try: print(B.c, end= ';')
except: print('ERR', end=';')
code

10. 
Для графа на рисунку з вершини 5 запущено алгоритм пошуку в глибину, що будує кістякове дерево. Списки суміжності вершин переглядаються алгоритмом за зростанням міток вершин.\par
\begin{center}
	\adjustimage{max size={0.9\linewidth}{0.9\paperheight}}
	{../10/Traversals/206.png}
\end{center}
\par Які з наведених ребер входять до складу отриманого кістякового дерева?\par
1) (1, 6)\par
2) (3, 5)\par
3) (0, 1)\par
4) (3, 6)\par
5) (4, 6)\par
6) (0, 4)\par
{\vskip 15pt} У формі вибрати номери відповідних ребер.\par
%answer:1 3 5
\end {large}}
%end
{\smallskip \begin {small} Затверджено на засіданні кафедри теоретичної кібернетики, протокол \No 12 від   19.05.2020 р.\par\smallskip\smallskip
{\parbox {6.5 cm}{Зав. кафедри \dotfill Ю.В.~Крак}}\hspace {0.7cm}{\hbox to 6.5 cm{ Екзаменатор \dotfill Т.О.~Карнаух}} \end{small}}
%begin
{{\begin {small}\par
{ \center  Київський національний університет імені Тараса Шевченка \par}
{\parbox {5 cm}{Освітня програма \hfill}{\parbox {7 cm}{Прикладна математика, Системний аналіз \hfill}}\parbox {6 cm}{\hfill Семестр 2}}\par
{\parbox {5 cm} {Навчальний предмет \hfill}\parbox {7 cm}{Програмування \hfill}\parbox {6cm}{\hfill}\par}\vspace{-7pt} \end{small}}{\center Екзаменаційний білет \No \textbf 
{ 278 }\par}}
{
\begin {large}
1. Що буде виведено за виконання?
code
try:
    res = True==9.0
    if isinstance(res, bool):
        print(f'bool;{res}')
    elif isinstance(res, int):
        print(f'int;{res}')
    elif isinstance(res, float):
        print(f'float;{res:.2f}')
    else:
        print('???')
except:
    print('ERR')
code

2. Що буде виведено за виконання?
code
a = 0
b = 5
c = 5

def g(b):
    global c
    a += 1
    b = 1
    c = 3
    return a + b + c

a = 1
b = 7
c = 2

try:
    print(g(b), a, b, c, sep=';')
except:
    print('ERR', a, b, c, sep=';')
code

3. Що буде виведено за виконання?
code
try:
    t = 5
    while t<=13:
        t += 2
    print(t, end=';')
except:
    print('ERR')
code

4. Що буде виведено за виконання?
code
try:
    for d in range(14, 1, 3):
        if d>=7:
            break
        print(d, end=';')
    else:
        print('end',end=';')
    print(d)
except:
    print('ERR')
code

5. Що буде виведено за виконання?\par (Дійсні значення, якщо наявні, в цьому завданні будуть виводитись у форматі з фіксованою крапкою та, можливо, з доволі великою кількістю десяткових розрядів. У відповіді для дійсних значень у ЦЬОМУ завданні достатньо вказати один десятковий розряд після крапки, відкинувши решту розрядів без округлення. Наприклад, замість 13.66666666666667 достатньо вказати 13.6, замість 25.23 достатньо вказати 25.2)
code
try:
    print(6, end=';')
    print(4j>6j, end=';')
    print(1, end=';')
except ValueError:
    print(6, end=';')
except ZeroDivisionError:
    print(7, end=';')
except:
    print('ERR', end=';')
else:
    print(5, end=';')
finally:
    print(2)
code

6. Що буде виведено за виконання?
code
def f(n):
    print(n%10, end='')
    if n>9:
        f(n//100)
 
f(987654)
code

7. Що буде виведено за виконання?
code
class A:
    d = 1

    def _h1(self): print(2, end=';')
    def __h2(self): print(3, end=';')

    def main(self):
        try: print(self.d, end=';')
        except: print('ERR', end=';')

        try: self._h1()
        except: print('ERR', end=';')

        try: self.__h2()
        except: print('ERR', end=';')


class B(A):
    d = 4

    def _h1(self): print(5, end=';')
    def __h2(self): print(6, end=';')

try: B().main()
except: print('ERR', end=';')
code

8. Що буде виведено за виконання?
code
def f(arg, n):
    arg[60] = 1
    n += -3

try:
    n = 6
    t = {49:4, 71:0, 49:8}
    f(t, n)
    print(n, end=';')
    for x in t :
        print(x, sep=';', end=';')
except:
    print('ERR')
code

9. Що буде виведено за виконання?
code
class A:
    c = 1
    a = 2

    def __init__(self):
        c = 3

class B(A):
    c = 4

    def __init__(self):
        super().__init__()
        self.b = 7

obj = B()
obj.a = obj.b + 10
B.a = 13

try: print(obj.a, end= ';')
except: print('ERR', end=';')

try: print(A.a, end= ';')
except: print('ERR', end=';')

try: print(B.a, end= ';')
except: print('ERR', end=';')
code

10. 
Для графа на рисунку з вершини 6 запущено алгоритм пошуку в глибину, що будує кістякове дерево. Списки суміжності вершин переглядаються алгоритмом за зростанням міток вершин.\par
\begin{center}
	\adjustimage{max size={0.9\linewidth}{0.9\paperheight}}
	{../10/Traversals/34.png}
\end{center}
\par Які з наведених ребер входять до складу отриманого кістякового дерева?\par
1) (1, 2)\par
2) (5, 6)\par
3) (2, 3)\par
4) (0, 2)\par
5) (3, 5)\par
6) (3, 4)\par
{\vskip 15pt} У формі вибрати номери відповідних ребер.\par
%answer:1 3 4 5 6
\end {large}}
%end
{\smallskip \begin {small} Затверджено на засіданні кафедри теоретичної кібернетики, протокол \No 12 від   19.05.2020 р.\par\smallskip\smallskip
{\parbox {6.5 cm}{Зав. кафедри \dotfill Ю.В.~Крак}}\hspace {0.7cm}{\hbox to 6.5 cm{ Екзаменатор \dotfill Т.О.~Карнаух}} \end{small}}
%begin
{{\begin {small}\par
{ \center  Київський національний університет імені Тараса Шевченка \par}
{\parbox {5 cm}{Освітня програма \hfill}{\parbox {7 cm}{Прикладна математика, Системний аналіз \hfill}}\parbox {6 cm}{\hfill Семестр 2}}\par
{\parbox {5 cm} {Навчальний предмет \hfill}\parbox {7 cm}{Програмування \hfill}\parbox {6cm}{\hfill}\par}\vspace{-7pt} \end{small}}{\center Екзаменаційний білет \No \textbf 
{ 279 }\par}}
{
\begin {large}
1. Що буде виведено за виконання?
code
try:
    res = 2j>="3.0j"
    if isinstance(res, bool):
        print(f'bool;{res}')
    elif isinstance(res, int):
        print(f'int;{res}')
    elif isinstance(res, float):
        print(f'float;{res:.2f}')
    else:
        print('???')
except:
    print('ERR')
code

2. Що буде виведено за виконання?
code
a = 0
b = 6
c = 4

def f(a):
    a *= 4
    b = 5
    c = 2
    return a + b + c

a = 3
b = 8
c = 9

try:
    print(f(a), a, b, c, sep=';')
except:
    print('ERR', a, b, c, sep=';')
code

3. Що буде виведено за виконання?
code
try:
    t = 7
    while t>7:
        t += -2
    print(t, end=';')
except:
    print('ERR')
code

4. Що буде виведено за виконання?
code
try:
    for a in range(-3, -9, -2):
        if a<=8:
            continue
        print(a, end=';');a = -6
    else:
        print(a,end=';')
    print(a)
except:
    print('ERR')
code

5. Що буде виведено за виконання?\par (Дійсні значення, якщо наявні, в цьому завданні будуть виводитись у форматі з фіксованою крапкою та, можливо, з доволі великою кількістю десяткових розрядів. У відповіді для дійсних значень у ЦЬОМУ завданні достатньо вказати один десятковий розряд після крапки, відкинувши решту розрядів без округлення. Наприклад, замість 13.66666666666667 достатньо вказати 13.6, замість 25.23 достатньо вказати 25.2)
code
try:
    print(0, end=';')
    print(int('c8'), end=';')
    print(3, end=';')
except TypeError:
    print(9, end=';')
except ValueError:
    print(8, end=';')
except:
    print('ERR', end=';')
else:
    print(0, end=';')
finally:
    print(3)
code

6. Що буде виведено за виконання?
code
def f(n):
    if n>9:
        f(n//100)
    print(n%10,end='')
 
f(987654)
code

7. Що буде виведено за виконання?
code
class A:
    c = 1

    def f1(self): print(2, end=';')
    def _f2(self): print(3, end=';')

    def main(self):
        try: print(self.c, end=';')
        except: print('ERR', end=';')

        try: self.f1()
        except: print('ERR', end=';')

        try: self._f2()
        except: print('ERR', end=';')


class B(A):
    c = 4

    def f1(self): print(5, end=';')
    def _f2(self): print(6, end=';')

try: B().main()
except: print('ERR', end=';')
code

8. Що буде виведено за виконання?
code
def f(arg, n):
    arg[61] = 3
    n = 7
    arg = set(arg.values())

try:
    n = 2
    d = {60:3, 86:8, 86:1}
    f(d, n)
    print(n, end=';')
    for x in d.keys():
        print(x, sep=';', end=';')
except:
    print('ERR')
code

9. Що буде виведено за виконання?
code
class A:
    a = 1
    b = 2

    def __init__(self):
        a = 3

class B(A):
    b = 4

    def __init__(self):
        super().__init__()
        self.c = 7

obj = B()
obj.c = obj.a + 10
B.c = 13

try: print(obj.c, end= ';')
except: print('ERR', end=';')

try: print(A.c, end= ';')
except: print('ERR', end=';')

try: print(B.c, end= ';')
except: print('ERR', end=';')
code

10. 
Для графа на рисунку з вершини 0 запущено алгоритм пошуку в глибину, що будує кістякове дерево. Списки суміжності вершин переглядаються алгоритмом за зростанням міток вершин.\par
\begin{center}
	\adjustimage{max size={0.9\linewidth}{0.9\paperheight}}
	{../10/Traversals/240.png}
\end{center}
\par Які з наведених ребер входять до складу отриманого кістякового дерева?\par
1) (0, 6)\par
2) (2, 3)\par
3) (4, 5)\par
4) (2, 4)\par
5) (5, 6)\par
6) (1, 4)\par
{\vskip 15pt} У формі вибрати номери відповідних ребер.\par
%answer:2 3 4 5 6
\end {large}}
%end
{\smallskip \begin {small} Затверджено на засіданні кафедри теоретичної кібернетики, протокол \No 12 від   19.05.2020 р.\par\smallskip\smallskip
{\parbox {6.5 cm}{Зав. кафедри \dotfill Ю.В.~Крак}}\hspace {0.7cm}{\hbox to 6.5 cm{ Екзаменатор \dotfill Т.О.~Карнаух}} \end{small}}
%begin
{{\begin {small}\par
{ \center  Київський національний університет імені Тараса Шевченка \par}
{\parbox {5 cm}{Освітня програма \hfill}{\parbox {7 cm}{Прикладна математика, Системний аналіз \hfill}}\parbox {6 cm}{\hfill Семестр 2}}\par
{\parbox {5 cm} {Навчальний предмет \hfill}\parbox {7 cm}{Програмування \hfill}\parbox {6cm}{\hfill}\par}\vspace{-7pt} \end{small}}{\center Екзаменаційний білет \No \textbf 
{ 280 }\par}}
{
\begin {large}
1. Що буде виведено за виконання?
code
def f(n):
    print('f', end='')
    return n

try:
    print(f(4) or f(5))
except:
    print('ERR')
code

2. Що буде виведено за виконання?
code
a = 5
b = 8
c = 1

def f(a):
    a = 4
    b = 1
    c = 5
    return a + b + c

a = 0
b = 9
c = 4

try:
    print(f(b), a, b, c, sep=';')
except:
    print('ERR', a, b, c, sep=';')
code

3. Що буде виведено за виконання?
code
try:
    n = 9
    while n<11:
        n += 1
    print(n, end=';')
except:
    print('ERR')
code

4. Що буде виведено за виконання?
code
try:
    for c in range(6, 4, 2):
        if c>3:
            break
        print(c, end=';');c = -4
    else:
        print(c,end=';')
    print(c)
except:
    print('ERR')
code

5. Що буде виведено за виконання?\par (Дійсні значення, якщо наявні, в цьому завданні будуть виводитись у форматі з фіксованою крапкою та, можливо, з доволі великою кількістю десяткових розрядів. У відповіді для дійсних значень у ЦЬОМУ завданні достатньо вказати один десятковий розряд після крапки, відкинувши решту розрядів без округлення. Наприклад, замість 13.66666666666667 достатньо вказати 13.6, замість 25.23 достатньо вказати 25.2)
code
try:
    print(5, end=';')
    print(int('4'), end=';')
    print(1, end=';')
except KeyboardInterrupt:
    print(6, end=';')
except Exception:
    print(9, end=';')
except:
    print('ERR', end=';')
else:
    print(7, end=';')
finally:
    print(2)
code

6. Що буде виведено за виконання?
code
def f(n):
    print(n%10, end='')
    if n>9:
        f(n//10)
 
f(123456)
code

7. Що буде виведено за виконання?
code
class A:
    a = 1

    def __h1(self): print(2, end=';')
    def _h2(self): print(3, end=';')

    def main(self):
        try: print(self.a, end=';')
        except: print('ERR', end=';')

        try: self.__h1()
        except: print('ERR', end=';')

        try: self._h2()
        except: print('ERR', end=';')


class B(A):
    a = 4

    def __h1(self): print(5, end=';')
    def _h2(self): print(6, end=';')

try: B().main()
except: print('ERR', end=';')
code

8. Що буде виведено за виконання?
code
def f(arg, n):
    arg[56] = 5
    n = 9
    arg = set(arg.keys())

try:
    n = 1
    s = {20:6, 19:0, 79:2, 19:0}
    f(s, n)
    print(n, end=';')
    for x in s.keys():
        print(x, sep=';', end=';')
except:
    print('ERR')
code

9. Що буде виведено за виконання?
code
class A:
    c = 1
    b = 2

    def __init__(self):
        b = 3

class B(A):
    c = 4

    def __init__(self):
        super().__init__()
        self.a = 7

obj = B()
obj.b = obj.c + 10
B.b = 13

try: print(obj.b, end= ';')
except: print('ERR', end=';')

try: print(A.b, end= ';')
except: print('ERR', end=';')

try: print(B.b, end= ';')
except: print('ERR', end=';')
code

10. 
Для графа на рисунку з вершини 3 запущено алгоритм пошуку в глибину, що будує кістякове дерево. Списки суміжності вершин переглядаються алгоритмом за зростанням міток вершин.\par
\begin{center}
	\adjustimage{max size={0.9\linewidth}{0.9\paperheight}}
	{../10/Traversals/269.png}
\end{center}
\par Які з наведених ребер входять до складу отриманого кістякового дерева?\par
1) (0, 1)\par
2) (1, 5)\par
3) (2, 3)\par
4) (1, 6)\par
5) (1, 4)\par
6) (3, 6)\par
{\vskip 15pt} У формі вибрати номери відповідних ребер.\par
%answer:1 2 3 5
\end {large}}
%end
{\smallskip \begin {small} Затверджено на засіданні кафедри теоретичної кібернетики, протокол \No 12 від   19.05.2020 р.\par\smallskip\smallskip
{\parbox {6.5 cm}{Зав. кафедри \dotfill Ю.В.~Крак}}\hspace {0.7cm}{\hbox to 6.5 cm{ Екзаменатор \dotfill Т.О.~Карнаух}} \end{small}}
%begin
{{\begin {small}\par
{ \center  Київський національний університет імені Тараса Шевченка \par}
{\parbox {5 cm}{Освітня програма \hfill}{\parbox {7 cm}{Прикладна математика, Системний аналіз \hfill}}\parbox {6 cm}{\hfill Семестр 2}}\par
{\parbox {5 cm} {Навчальний предмет \hfill}\parbox {7 cm}{Програмування \hfill}\parbox {6cm}{\hfill}\par}\vspace{-7pt} \end{small}}{\center Екзаменаційний білет \No \textbf 
{ 281 }\par}}
{
\begin {large}
1. Що буде виведено за виконання?
code
def f(n):
    print('f', end='')
    return n!=-1

try:
    print(f(-6) and f(-5))
except:
    print('ERR')
code

2. Що буде виведено за виконання?
code
a = 4
b = 0
c = 5

def h(a):
    global c
    a = 4
    b = 5
    c = 1
    return a + b + c

a = 6
b = 2
c = 7

try:
    print(h(b), a, b, c, sep=';')
except:
    print('ERR', a, b, c, sep=';')
code

3. Що буде виведено за виконання?
code
try:
    i = 9
    while i<8:
        i += 2
    print(i, end=';')
except:
    print('ERR')
code

4. Що буде виведено за виконання?
code
try:
    for b in range(-7, -7, -1):
        if b<6:
            continue
        print(b, end=';')
    else:
        print(b,end=';')
    print(b)
except:
    print('ERR')
code

5. Що буде виведено за виконання?\par (Дійсні значення, якщо наявні, в цьому завданні будуть виводитись у форматі з фіксованою крапкою та, можливо, з доволі великою кількістю десяткових розрядів. У відповіді для дійсних значень у ЦЬОМУ завданні достатньо вказати один десятковий розряд після крапки, відкинувши решту розрядів без округлення. Наприклад, замість 13.66666666666667 достатньо вказати 13.6, замість 25.23 достатньо вказати 25.2)
code
try:
    print(1, end=';')
    print(int('8'), end=';')
    print(5, end=';')
except ZeroDivisionError:
    print(2, end=';')
except KeyboardInterrupt:
    print(4, end=';')
except:
    print('ERR', end=';')
else:
    print(9, end=';')
finally:
    print(0)
code

6. Що буде виведено за виконання?
code
def f(n):
    if n>9:
        f(n//100)
    print(n%10, end='')
 
f(123456)
code

7. Що буде виведено за виконання?
code
class A:
    b = 1

    def _g1(self): print(2, end=';')
    def __g2(self): print(3, end=';')

    def main(self):
        try: print(self.b, end=';')
        except: print('ERR', end=';')

        try: self._g1()
        except: print('ERR', end=';')

        try: self.__g2()
        except: print('ERR', end=';')


class B(A):
    b = 4

    def _g1(self): print(5, end=';')
    def __g2(self): print(6, end=';')

try: B().main()
except: print('ERR', end=';')
code

8. Що буде виведено за виконання?
code
def f(arg, n):
    arg[78] = 5
    n = 5
    arg = tuple(arg.values())

try:
    n = 0
    t = {78:1, 13:5, 78:1}
    f(t, n)
    print(n, end=';')
    for x in t.values():
        print(x, sep=';', end=';')
except:
    print('ERR')
code

9. Що буде виведено за виконання?
code
class A:
    b = 1
    c = 2

    def __init__(self):
        b = 3

class B(A):
    c = 4

    def __init__(self):
        super().__init__()
        self.a = 7

obj = B()
obj.a = obj.a + 10
B.a = 13

try: print(obj.a, end= ';')
except: print('ERR', end=';')

try: print(A.a, end= ';')
except: print('ERR', end=';')

try: print(B.a, end= ';')
except: print('ERR', end=';')
code

10. 
Для графа на рисунку з вершини 5 запущено алгоритм пошуку в глибину, що будує кістякове дерево. Списки суміжності вершин переглядаються алгоритмом за зростанням міток вершин.\par
\begin{center}
	\adjustimage{max size={0.9\linewidth}{0.9\paperheight}}
	{../10/Traversals/237.png}
\end{center}
\par Які з наведених ребер входять до складу отриманого кістякового дерева?\par
1) (3, 4)\par
2) (2, 3)\par
3) (4, 5)\par
4) (1, 3)\par
5) (4, 6)\par
6) (0, 5)\par
{\vskip 15pt} У формі вибрати номери відповідних ребер.\par
%answer:1 2 4 5 6
\end {large}}
%end
{\smallskip \begin {small} Затверджено на засіданні кафедри теоретичної кібернетики, протокол \No 12 від   19.05.2020 р.\par\smallskip\smallskip
{\parbox {6.5 cm}{Зав. кафедри \dotfill Ю.В.~Крак}}\hspace {0.7cm}{\hbox to 6.5 cm{ Екзаменатор \dotfill Т.О.~Карнаух}} \end{small}}
%begin
{{\begin {small}\par
{ \center  Київський національний університет імені Тараса Шевченка \par}
{\parbox {5 cm}{Освітня програма \hfill}{\parbox {7 cm}{Прикладна математика, Системний аналіз \hfill}}\parbox {6 cm}{\hfill Семестр 2}}\par
{\parbox {5 cm} {Навчальний предмет \hfill}\parbox {7 cm}{Програмування \hfill}\parbox {6cm}{\hfill}\par}\vspace{-7pt} \end{small}}{\center Екзаменаційний білет \No \textbf 
{ 282 }\par}}
{
\begin {large}
1. Що буде виведено за виконання?
code
try:
    res = 4**-2+2
    if isinstance(res, bool):
        print(f'bool;{res}')
    elif isinstance(res, int):
        print(f'int;{res}')
    elif isinstance(res, float):
        print(f'float;{res:.2f}')
    else:
        print('???')
except:
    print('ERR')
code

2. Що буде виведено за виконання?
code
a = 6
b = 1
c = 7

def h(b):
    global c
    a -= 1
    b = 5
    c = 3
    return a + b + c

a = 2
b = 9
c = 4

try:
    print(h(b), a, b, c, sep=';')
except:
    print('ERR', a, b, c, sep=';')
code

3. Що буде виведено за виконання?
code
try:
    m = 7
    while m<=6:
        m += 1
    print(m, end=';')
except:
    print('ERR')
code

4. Що буде виведено за виконання?
code
try:
    for e in range(5, -13, 2):
        if e<=4:
            continue
        print(e, end=';')
    else:
        print(e,end=';')
    print(e)
except:
    print('ERR')
code

5. Що буде виведено за виконання?\par (Дійсні значення, якщо наявні, в цьому завданні будуть виводитись у форматі з фіксованою крапкою та, можливо, з доволі великою кількістю десяткових розрядів. У відповіді для дійсних значень у ЦЬОМУ завданні достатньо вказати один десятковий розряд після крапки, відкинувши решту розрядів без округлення. Наприклад, замість 13.66666666666667 достатньо вказати 13.6, замість 25.23 достатньо вказати 25.2)
code
try:
    print(6, end=';')
    print(3j<=7j, end=';')
    print(7, end=';')
except TypeError:
    print(7, end=';')
except Exception:
    print(3, end=';')
except:
    print('ERR', end=';')
else:
    print(0, end=';')
finally:
    print(5)
code

6. Що буде виведено за виконання?
code
def f(n):
    if n>2:
        f(n//2)
    print(n%2, end='')
 
f(16)
code

7. Що буде виведено за виконання?
code
class A:
    d = 1

    def _f1(self): print(2, end=';')
    def _f2(self): print(3, end=';')

    def main(self):
        try: print(self.d, end=';')
        except: print('ERR', end=';')

        try: self._f1()
        except: print('ERR', end=';')

        try: self._f2()
        except: print('ERR', end=';')


class B(A):
    d = 4

    def _f1(self): print(5, end=';')
    def _f2(self): print(6, end=';')

try: B().main()
except: print('ERR', end=';')
code

8. Що буде виведено за виконання?
code
def f(arg, n):
    arg[27] = 2
    n = 9
    arg = list(arg.items())

try:
    n = 4
    d = {27:7, 21:0, 38:5, 21:7}
    f(d, n)
    print(n, end=';')
    for x in d.items():
        print(*x, sep=';', end=';')
except:
    print('ERR')
code

9. Що буде виведено за виконання?
code
class A:
    a = 1
    c = 2

    def __init__(self):
        c = 3

class B(A):
    a = 4

    def __init__(self):
        super().__init__()
        self.b = 7

obj = B()
obj.b = obj.c + 10
B.b = 13

try: print(obj.b, end= ';')
except: print('ERR', end=';')

try: print(A.b, end= ';')
except: print('ERR', end=';')

try: print(B.b, end= ';')
except: print('ERR', end=';')
code

10. 
Для графа на рисунку з вершини 2 запущено алгоритм пошуку в глибину, що будує кістякове дерево. Списки суміжності вершин переглядаються алгоритмом за зростанням міток вершин.\par
\begin{center}
	\adjustimage{max size={0.9\linewidth}{0.9\paperheight}}
	{../10/Traversals/117.png}
\end{center}
\par Які з наведених ребер входять до складу отриманого кістякового дерева?\par
1) (3, 5)\par
2) (1, 5)\par
3) (0, 1)\par
4) (0, 2)\par
5) (1, 6)\par
6) (1, 3)\par
{\vskip 15pt} У формі вибрати номери відповідних ребер.\par
%answer:1 3 4 5 6
\end {large}}
%end
{\smallskip \begin {small} Затверджено на засіданні кафедри теоретичної кібернетики, протокол \No 12 від   19.05.2020 р.\par\smallskip\smallskip
{\parbox {6.5 cm}{Зав. кафедри \dotfill Ю.В.~Крак}}\hspace {0.7cm}{\hbox to 6.5 cm{ Екзаменатор \dotfill Т.О.~Карнаух}} \end{small}}
%begin
{{\begin {small}\par
{ \center  Київський національний університет імені Тараса Шевченка \par}
{\parbox {5 cm}{Освітня програма \hfill}{\parbox {7 cm}{Прикладна математика, Системний аналіз \hfill}}\parbox {6 cm}{\hfill Семестр 2}}\par
{\parbox {5 cm} {Навчальний предмет \hfill}\parbox {7 cm}{Програмування \hfill}\parbox {6cm}{\hfill}\par}\vspace{-7pt} \end{small}}{\center Екзаменаційний білет \No \textbf 
{ 283 }\par}}
{
\begin {large}
1. Що буде виведено за виконання?
code
try:
    res = 9**0.5*True
    if isinstance(res, bool):
        print(f'bool;{res}')
    elif isinstance(res, int):
        print(f'int;{res}')
    elif isinstance(res, float):
        print(f'float;{res:.2f}')
    else:
        print('???')
except:
    print('ERR')
code

2. Що буде виведено за виконання?
code
a = 3
b = 5
c = 0

def h(a):
    a -= 2
    b = 4
    c = 3
    return a + b + c

a = 8
b = 6
c = 3

try:
    print(h(a), a, b, c, sep=';')
except:
    print('ERR', a, b, c, sep=';')
code

3. Що буде виведено за виконання?
code
try:
    n = 5
    while n>5:
        n += -3
    print(n, end=';')
except:
    print('ERR')
code

4. Що буде виведено за виконання?
code
try:
    for f in range(-14, 14, -2):
        if f>=7:
            continue
        print(f, end=';');f = -5
    else:
        print(f,end=';')
    print(f)
except:
    print('ERR')
code

5. Що буде виведено за виконання?\par (Дійсні значення, якщо наявні, в цьому завданні будуть виводитись у форматі з фіксованою крапкою та, можливо, з доволі великою кількістю десяткових розрядів. У відповіді для дійсних значень у ЦЬОМУ завданні достатньо вказати один десятковий розряд після крапки, відкинувши решту розрядів без округлення. Наприклад, замість 13.66666666666667 достатньо вказати 13.6, замість 25.23 достатньо вказати 25.2)
code
try:
    print(6, end=';')
    print(int('d2'), end=';')
    print(4, end=';')
except ValueError:
    print(8, end=';')
except ZeroDivisionError:
    print(1, end=';')
except:
    print('ERR', end=';')
else:
    print(9, end=';')
finally:
    print(7)
code

6. Що буде виведено за виконання?
code
def f(n):
    if n>9:
        f(n//100)
        print(n%10, end='')
 
f(123456)
code

7. Що буде виведено за виконання?
code
class A:
    a = 1

    def h1(self): print(2, end=';')
    def _h2(self): print(3, end=';')

    def main(self):
        try: print(self.a, end=';')
        except: print('ERR', end=';')

        try: self.h1()
        except: print('ERR', end=';')

        try: self._h2()
        except: print('ERR', end=';')


class B(A):
    a = 4

    def h1(self): print(5, end=';')
    def _h2(self): print(6, end=';')

try: B().main()
except: print('ERR', end=';')
code

8. Що буде виведено за виконання?
code
def f(arg, n):
    n = 5
    tmp = arg
    tmp[-5:] = 'yz'
    return len(tmp)>3

try:
    f = [10,11,12,13,14]
    n = 7
    f(f, n)
    print(n, end=';')
    for el in f:
        print(el, end=';')
except:
    print('ERR')
code

9. Що буде виведено за виконання?
code
class A:
    a = 1
    c = 2

    def __init__(self):
        a = 3

class B(A):
    a = 4

    def __init__(self):
        super().__init__()
        self.b = 7

obj = B()
obj.b = obj.c + 10
B.b = 13

try: print(obj.b, end= ';')
except: print('ERR', end=';')

try: print(A.b, end= ';')
except: print('ERR', end=';')

try: print(B.b, end= ';')
except: print('ERR', end=';')
code

10. 
Для графа на рисунку з вершини 2 запущено алгоритм пошуку в глибину, що будує кістякове дерево. Списки суміжності вершин переглядаються алгоритмом за зростанням міток вершин.\par
\begin{center}
	\adjustimage{max size={0.9\linewidth}{0.9\paperheight}}
	{../10/Traversals/125.png}
\end{center}
\par Які з наведених ребер входять до складу отриманого кістякового дерева?\par
1) (2, 4)\par
2) (3, 6)\par
3) (0, 5)\par
4) (4, 5)\par
5) (4, 6)\par
6) (1, 2)\par
{\vskip 15pt} У формі вибрати номери відповідних ребер.\par
%answer:2 3 4 5 6
\end {large}}
%end
{\smallskip \begin {small} Затверджено на засіданні кафедри теоретичної кібернетики, протокол \No 12 від   19.05.2020 р.\par\smallskip\smallskip
{\parbox {6.5 cm}{Зав. кафедри \dotfill Ю.В.~Крак}}\hspace {0.7cm}{\hbox to 6.5 cm{ Екзаменатор \dotfill Т.О.~Карнаух}} \end{small}}
%begin
{{\begin {small}\par
{ \center  Київський національний університет імені Тараса Шевченка \par}
{\parbox {5 cm}{Освітня програма \hfill}{\parbox {7 cm}{Прикладна математика, Системний аналіз \hfill}}\parbox {6 cm}{\hfill Семестр 2}}\par
{\parbox {5 cm} {Навчальний предмет \hfill}\parbox {7 cm}{Програмування \hfill}\parbox {6cm}{\hfill}\par}\vspace{-7pt} \end{small}}{\center Екзаменаційний білет \No \textbf 
{ 284 }\par}}
{
\begin {large}
1. Що буде виведено за виконання?
code
try:
    res = 6!=7.0 and not 9>9.0 and 7>=False
    if isinstance(res, bool):
        print(f'bool;{res}')
    elif isinstance(res, int):
        print(f'int;{res}')
    elif isinstance(res, float):
        print(f'float;{res:.2f}')
    else:
        print('???')
except:
    print('ERR')
code

2. Що буде виведено за виконання?
code
a = 9
b = 1
c = 3

def g(b):
    a = 3
    b = 2
    c = 4
    return a + b + c

a = 8
b = 4
c = 2

try:
    print(g(a), a, b, c, sep=';')
except:
    print('ERR', a, b, c, sep=';')
code

3. Що буде виведено за виконання?
code
try:
    j = 10
    while j>=6:
        j += -2
    print(j, end=';')
except:
    print('ERR')
code

4. Що буде виведено за виконання?
code
try:
    for c in range(-1, 6, -3):
        if c>6:
            continue
        print(c, end=';')
    else:
        print(c,end=';')
    print(c)
except:
    print('ERR')
code

5. Що буде виведено за виконання?\par (Дійсні значення, якщо наявні, в цьому завданні будуть виводитись у форматі з фіксованою крапкою та, можливо, з доволі великою кількістю десяткових розрядів. У відповіді для дійсних значень у ЦЬОМУ завданні достатньо вказати один десятковий розряд після крапки, відкинувши решту розрядів без округлення. Наприклад, замість 13.66666666666667 достатньо вказати 13.6, замість 25.23 достатньо вказати 25.2)
code
try:
    print(0, end=';')
    print(1/0, end=';')
    print(9, end=';')
except ValueError:
    print(8, end=';')
except KeyboardInterrupt:
    print(4, end=';')
except:
    print('ERR', end=';')
else:
    print(6, end=';')
finally:
    print(5)
code

6. Що буде виведено за виконання?
code
def f(n):
    print(n%10,end='')
    if n>9:
        f(n//100)
 
f(123456)
code

7. Що буде виведено за виконання?
code
class A:
    c = 1

    def _f1(self): print(2, end=';')
    def f2(self): print(3, end=';')

    def main(self):
        try: print(self.c, end=';')
        except: print('ERR', end=';')

        try: self._f1()
        except: print('ERR', end=';')

        try: self.f2()
        except: print('ERR', end=';')


class B(A):
    c = 4

    def _f1(self): print(5, end=';')
    def f2(self): print(6, end=';')

try: B().main()
except: print('ERR', end=';')
code

8. Що буде виведено за виконання?
code
def f(arg, n):
    arg[69] = 4
    n *= 3

try:
    n = 3
    d = {69:5, 61:9, 73:7}
    f(d, n)
    print(n, end=';')
    for x in d.items():
        print(x, sep=';', end=';')
except:
    print('ERR')
code

9. Що буде виведено за виконання?
code
class A:
    c = 1
    b = 2

    def __init__(self):
        b = 3

class B(A):
    b = 4

    def __init__(self):
        super().__init__()
        self.a = 7

obj = B()
obj.a = obj.a + 10
B.a = 13

try: print(obj.a, end= ';')
except: print('ERR', end=';')

try: print(A.a, end= ';')
except: print('ERR', end=';')

try: print(B.a, end= ';')
except: print('ERR', end=';')
code

10. 
Для графа на рисунку з вершини 6 запущено алгоритм пошуку в глибину, що будує кістякове дерево. Списки суміжності вершин переглядаються алгоритмом за зростанням міток вершин.\par
\begin{center}
	\adjustimage{max size={0.9\linewidth}{0.9\paperheight}}
	{../10/Traversals/68.png}
\end{center}
\par Які з наведених ребер входять до складу отриманого кістякового дерева?\par
1) (0, 5)\par
2) (1, 6)\par
3) (2, 5)\par
4) (0, 1)\par
5) (4, 5)\par
6) (1, 4)\par
{\vskip 15pt} У формі вибрати номери відповідних ребер.\par
%answer:1 2 3 4 5
\end {large}}
%end
{\smallskip \begin {small} Затверджено на засіданні кафедри теоретичної кібернетики, протокол \No 12 від   19.05.2020 р.\par\smallskip\smallskip
{\parbox {6.5 cm}{Зав. кафедри \dotfill Ю.В.~Крак}}\hspace {0.7cm}{\hbox to 6.5 cm{ Екзаменатор \dotfill Т.О.~Карнаух}} \end{small}}
%begin
{{\begin {small}\par
{ \center  Київський національний університет імені Тараса Шевченка \par}
{\parbox {5 cm}{Освітня програма \hfill}{\parbox {7 cm}{Прикладна математика, Системний аналіз \hfill}}\parbox {6 cm}{\hfill Семестр 2}}\par
{\parbox {5 cm} {Навчальний предмет \hfill}\parbox {7 cm}{Програмування \hfill}\parbox {6cm}{\hfill}\par}\vspace{-7pt} \end{small}}{\center Екзаменаційний білет \No \textbf 
{ 285 }\par}}
{
\begin {large}
1. Що буде виведено за виконання?
code
def f(n):
    print('f', end='')
    return n

try:
    print(f(-7) or f(9))
except:
    print('ERR')
code

2. Що буде виведено за виконання?
code
a = 2
b = 1
c = 4

def f(b):
    global c
    a = 5
    b += 4
    c = 1
    return a + b + c

a = 0
b = 5
c = 9

try:
    print(f(b), a, b, c, sep=';')
except:
    print('ERR', a, b, c, sep=';')
code

3. Що буде виведено за виконання?
code
try:
    k = 6
    while k>8:
        k += -3
    print(k, end=';')
except:
    print('ERR')
code

4. Що буде виведено за виконання?
code
try:
    for a in range(9, 0, 3):
        if a>=3:
            break
        print(a, end=';')
    else:
        print(a,end=';')
    print(a)
except:
    print('ERR')
code

5. Що буде виведено за виконання?\par (Дійсні значення, якщо наявні, в цьому завданні будуть виводитись у форматі з фіксованою крапкою та, можливо, з доволі великою кількістю десяткових розрядів. У відповіді для дійсних значень у ЦЬОМУ завданні достатньо вказати один десятковий розряд після крапки, відкинувши решту розрядів без округлення. Наприклад, замість 13.66666666666667 достатньо вказати 13.6, замість 25.23 достатньо вказати 25.2)
code
try:
    print(3, end=';')
    print(2%2, end=';')
    print(0, end=';')
except TypeError:
    print(6, end=';')
except Exception:
    print(2, end=';')
except:
    print('ERR', end=';')
else:
    print(3, end=';')
finally:
    print(4)
code

6. Що буде виведено за виконання?
code
def f(n):
    if n>9:
        f(n//100)
    print(n%10,end='')
 
f(987654)
code

7. Що буде виведено за виконання?
code
class A:
    b = 1

    def __g1(self): print(2, end=';')
    def _g2(self): print(3, end=';')

    def main(self):
        try: print(self.b, end=';')
        except: print('ERR', end=';')

        try: self.__g1()
        except: print('ERR', end=';')

        try: self._g2()
        except: print('ERR', end=';')


class B(A):
    b = 4

    def __g1(self): print(5, end=';')
    def _g2(self): print(6, end=';')

try: B().main()
except: print('ERR', end=';')
code

8. Що буде виведено за виконання?
code
def f(arg, n):
    arg[63] = 2
    n -= 2

try:
    n = 4
    s = {78:3, 79:5, 45:3, 78:0}
    f(s, n)
    print(n, end=';')
    for x in s.keys():
        print(x, sep=';', end=';')
except:
    print('ERR')
code

9. Що буде виведено за виконання?
code
class A:
    b = 1
    a = 2

    def __init__(self):
        b = 3

class B(A):
    a = 4

    def __init__(self):
        super().__init__()
        self.c = 7

obj = B()
obj.b = obj.c + 10
B.b = 13

try: print(obj.b, end= ';')
except: print('ERR', end=';')

try: print(A.b, end= ';')
except: print('ERR', end=';')

try: print(B.b, end= ';')
except: print('ERR', end=';')
code

10. 
Для графа на рисунку з вершини 1 запущено алгоритм пошуку в глибину, що будує кістякове дерево. Списки суміжності вершин переглядаються алгоритмом за зростанням міток вершин.\par
\begin{center}
	\adjustimage{max size={0.9\linewidth}{0.9\paperheight}}
	{../10/Traversals/43.png}
\end{center}
\par Які з наведених ребер входять до складу отриманого кістякового дерева?\par
1) (3, 4)\par
2) (1, 4)\par
3) (2, 5)\par
4) (0, 6)\par
5) (0, 2)\par
6) (0, 1)\par
{\vskip 15pt} У формі вибрати номери відповідних ребер.\par
%answer:1 2 3 5 6
\end {large}}
%end
{\smallskip \begin {small} Затверджено на засіданні кафедри теоретичної кібернетики, протокол \No 12 від   19.05.2020 р.\par\smallskip\smallskip
{\parbox {6.5 cm}{Зав. кафедри \dotfill Ю.В.~Крак}}\hspace {0.7cm}{\hbox to 6.5 cm{ Екзаменатор \dotfill Т.О.~Карнаух}} \end{small}}
%begin
{{\begin {small}\par
{ \center  Київський національний університет імені Тараса Шевченка \par}
{\parbox {5 cm}{Освітня програма \hfill}{\parbox {7 cm}{Прикладна математика, Системний аналіз \hfill}}\parbox {6 cm}{\hfill Семестр 2}}\par
{\parbox {5 cm} {Навчальний предмет \hfill}\parbox {7 cm}{Програмування \hfill}\parbox {6cm}{\hfill}\par}\vspace{-7pt} \end{small}}{\center Екзаменаційний білет \No \textbf 
{ 286 }\par}}
{
\begin {large}
1. Що буде виведено за виконання?
code
try:
    res = "True"<False
    if isinstance(res, bool):
        print(f'bool;{res}')
    elif isinstance(res, int):
        print(f'int;{res}')
    elif isinstance(res, float):
        print(f'float;{res:.2f}')
    else:
        print('???')
except:
    print('ERR')
code

2. Що буде виведено за виконання?
code
a = 1
b = 5
c = 8

def h(b):
    global c
    a = 5
    b = 2
    c = 3
    return a + b + c

a = 6
b = 3
c = 0

try:
    print(h(a), a, b, c, sep=';')
except:
    print('ERR', a, b, c, sep=';')
code

3. Що буде виведено за виконання?
code
try:
    m = 8
    while m>=0:
        m += -2
    print(m, end=';')
except:
    print('ERR')
code

4. Що буде виведено за виконання?
code
try:
    for e in range(2, 13, -2):
        if e<=8:
            break
        print(e, end=';')
        if e<=7:
            break
    else:
        print(e,end=';')
    print(e)
except:
    print('ERR')
code

5. Що буде виведено за виконання?\par (Дійсні значення, якщо наявні, в цьому завданні будуть виводитись у форматі з фіксованою крапкою та, можливо, з доволі великою кількістю десяткових розрядів. У відповіді для дійсних значень у ЦЬОМУ завданні достатньо вказати один десятковий розряд після крапки, відкинувши решту розрядів без округлення. Наприклад, замість 13.66666666666667 достатньо вказати 13.6, замість 25.23 достатньо вказати 25.2)
code
try:
    print(8, end=';')
    print(5//0.0, end=';')
    print(1, end=';')
except TypeError:
    print(9, end=';')
except ValueError:
    print(7, end=';')
except:
    print('ERR', end=';')
else:
    print(2, end=';')
finally:
    print(0)
code

6. Що буде виведено за виконання?
code
def f(n):
    print(n%10, end='')
    if n>9:
        f(n//10)
 
f(123456)
code

7. Що буде виведено за виконання?
code
class A:
    d = 1

    def h1(self): print(2, end=';')
    def _h2(self): print(3, end=';')

    def main(self):
        try: print(self.d, end=';')
        except: print('ERR', end=';')

        try: self.h1()
        except: print('ERR', end=';')

        try: self._h2()
        except: print('ERR', end=';')


class B(A):
    d = 4

    def h1(self): print(5, end=';')
    def _h2(self): print(6, end=';')

try: B().main()
except: print('ERR', end=';')
code

8. Що буде виведено за виконання?
code
def f(arg, n):
    arg[40] = 5
    n += -2

try:
    n = 7
    d = {52:0, 19:4, 18:7, 19:7}
    f(d, n)
    print(n, end=';')
    for x in d.keys():
        print(x, sep=';', end=';')
except:
    print('ERR')
code

9. Що буде виведено за виконання?
code
class A:
    c = 1
    a = 2

    def __init__(self):
        a = 3

class B(A):
    c = 4

    def __init__(self):
        super().__init__()
        self.b = 7

obj = B()
obj.c = obj.a + 10
B.c = 13

try: print(obj.c, end= ';')
except: print('ERR', end=';')

try: print(A.c, end= ';')
except: print('ERR', end=';')

try: print(B.c, end= ';')
except: print('ERR', end=';')
code

10. 
Для графа на рисунку з вершини 3 запущено алгоритм пошуку в глибину, що будує кістякове дерево. Списки суміжності вершин переглядаються алгоритмом за зростанням міток вершин.\par
\begin{center}
	\adjustimage{max size={0.9\linewidth}{0.9\paperheight}}
	{../10/Traversals/3.png}
\end{center}
\par Які з наведених ребер входять до складу отриманого кістякового дерева?\par
1) (4, 5)\par
2) (2, 3)\par
3) (0, 3)\par
4) (3, 6)\par
5) (1, 5)\par
6) (1, 6)\par
{\vskip 15pt} У формі вибрати номери відповідних ребер.\par
%answer:1 3 5 6
\end {large}}
%end
{\smallskip \begin {small} Затверджено на засіданні кафедри теоретичної кібернетики, протокол \No 12 від   19.05.2020 р.\par\smallskip\smallskip
{\parbox {6.5 cm}{Зав. кафедри \dotfill Ю.В.~Крак}}\hspace {0.7cm}{\hbox to 6.5 cm{ Екзаменатор \dotfill Т.О.~Карнаух}} \end{small}}
%begin
{{\begin {small}\par
{ \center  Київський національний університет імені Тараса Шевченка \par}
{\parbox {5 cm}{Освітня програма \hfill}{\parbox {7 cm}{Прикладна математика, Системний аналіз \hfill}}\parbox {6 cm}{\hfill Семестр 2}}\par
{\parbox {5 cm} {Навчальний предмет \hfill}\parbox {7 cm}{Програмування \hfill}\parbox {6cm}{\hfill}\par}\vspace{-7pt} \end{small}}{\center Екзаменаційний білет \No \textbf 
{ 287 }\par}}
{
\begin {large}
1. Що буде виведено за виконання?
code
try:
    res = 4**False--2
    if isinstance(res, bool):
        print(f'bool;{res}')
    elif isinstance(res, int):
        print(f'int;{res}')
    elif isinstance(res, float):
        print(f'float;{res:.2f}')
    else:
        print('???')
except:
    print('ERR')
code

2. Що буде виведено за виконання?
code
a = 9
b = 7
c = 3

def f(a):
    a = 1
    b = 2
    c = 5
    return a + b + c

a = 6
b = 5
c = 2

try:
    print(f(a), a, b, c, sep=';')
except:
    print('ERR', a, b, c, sep=';')
code

3. Що буде виведено за виконання?
code
try:
    j = 7
    while j>=4:
        j += -3
    print(j, end=';')
except:
    print('ERR')
code

4. Що буде виведено за виконання?
code
try:
    for b in range(-15, 3, -3):
        if b<4:
            continue
        print(b, end=';');b = -8
    else:
        print('end',end=';')
    print(b)
except:
    print('ERR')
code

5. Що буде виведено за виконання?\par (Дійсні значення, якщо наявні, в цьому завданні будуть виводитись у форматі з фіксованою крапкою та, можливо, з доволі великою кількістю десяткових розрядів. У відповіді для дійсних значень у ЦЬОМУ завданні достатньо вказати один десятковий розряд після крапки, відкинувши решту розрядів без округлення. Наприклад, замість 13.66666666666667 достатньо вказати 13.6, замість 25.23 достатньо вказати 25.2)
code
try:
    print(5, end=';')
    print(7/1, end=';')
    print(3, end=';')
except ZeroDivisionError:
    print(8, end=';')
except Exception:
    print(6, end=';')
except:
    print('ERR', end=';')
else:
    print(1, end=';')
finally:
    print(9)
code

6. Що буде виведено за виконання?
code
def f(s):
    s1 = s
    for i in range(len(s)):
        s1[i] += 1
    return s1

s = [1, 2, 3]
f(s)
print(s)
code

7. Що буде виведено за виконання?
code
class A:
    a = 1

    def __g1(self): print(2, end=';')
    def _g2(self): print(3, end=';')

    def main(self):
        try: print(self.a, end=';')
        except: print('ERR', end=';')

        try: self.__g1()
        except: print('ERR', end=';')

        try: self._g2()
        except: print('ERR', end=';')


class B(A):
    a = 4

    def __g1(self): print(5, end=';')
    def _g2(self): print(6, end=';')

try: B().main()
except: print('ERR', end=';')
code

8. Що буде виведено за виконання?
code
def f(arg, n):
    arg[71] = 5
    n = 7
    arg = list(arg.values())

try:
    n = 0
    s = {27:5, 88:0, 71:6, 71:4}
    f(s, n)
    print(n, end=';')
    for x, y in s.items():
        print(x, y, sep=';', end=';')
except:
    print('ERR')
code

9. Що буде виведено за виконання?
code
class A:
    a = 1
    b = 2

    def __init__(self):
        a = 3

class B(A):
    b = 4

    def __init__(self):
        super().__init__()
        self.c = 7

obj = B()
obj.b = obj.b + 10
B.b = 13

try: print(obj.b, end= ';')
except: print('ERR', end=';')

try: print(A.b, end= ';')
except: print('ERR', end=';')

try: print(B.b, end= ';')
except: print('ERR', end=';')
code

10. 
Для графа на рисунку з вершини 1 запущено алгоритм пошуку в глибину, що будує кістякове дерево. Списки суміжності вершин переглядаються алгоритмом за зростанням міток вершин.\par
\begin{center}
	\adjustimage{max size={0.9\linewidth}{0.9\paperheight}}
	{../10/Traversals/94.png}
\end{center}
\par Які з наведених ребер входять до складу отриманого кістякового дерева?\par
1) (1, 2)\par
2) (1, 6)\par
3) (5, 6)\par
4) (0, 5)\par
5) (0, 6)\par
6) (2, 5)\par
{\vskip 15pt} У формі вибрати номери відповідних ребер.\par
%answer:3 6
\end {large}}
%end
{\smallskip \begin {small} Затверджено на засіданні кафедри теоретичної кібернетики, протокол \No 12 від   19.05.2020 р.\par\smallskip\smallskip
{\parbox {6.5 cm}{Зав. кафедри \dotfill Ю.В.~Крак}}\hspace {0.7cm}{\hbox to 6.5 cm{ Екзаменатор \dotfill Т.О.~Карнаух}} \end{small}}
%begin
{{\begin {small}\par
{ \center  Київський національний університет імені Тараса Шевченка \par}
{\parbox {5 cm}{Освітня програма \hfill}{\parbox {7 cm}{Прикладна математика, Системний аналіз \hfill}}\parbox {6 cm}{\hfill Семестр 2}}\par
{\parbox {5 cm} {Навчальний предмет \hfill}\parbox {7 cm}{Програмування \hfill}\parbox {6cm}{\hfill}\par}\vspace{-7pt} \end{small}}{\center Екзаменаційний білет \No \textbf 
{ 288 }\par}}
{
\begin {large}
1. Що буде виведено за виконання?
code
try:
    res = "2.0j"!="7j"
    if isinstance(res, bool):
        print(f'bool;{res}')
    elif isinstance(res, int):
        print(f'int;{res}')
    elif isinstance(res, float):
        print(f'float;{res:.2f}')
    else:
        print('???')
except:
    print('ERR')
code

2. Що буде виведено за виконання?
code
a = 7
b = 8
c = 0

def g(a):
    a *= 2
    b = 3
    c = 1
    return a + b + c

a = 6
b = 4
c = 9

try:
    print(g(a), a, b, c, sep=';')
except:
    print('ERR', a, b, c, sep=';')
code

3. Що буде виведено за виконання?
code
try:
    m = 5
    while m>1:
        m += -2
    print(m, end=';')
except:
    print('ERR')
code

4. Що буде виведено за виконання?
code
try:
    for d in range(13, -9, 2):
        if d>5:
            continue
        print(d, end=';')
    else:
        print(d,end=';')
    print(d)
except:
    print('ERR')
code

5. Що буде виведено за виконання?\par (Дійсні значення, якщо наявні, в цьому завданні будуть виводитись у форматі з фіксованою крапкою та, можливо, з доволі великою кількістю десяткових розрядів. У відповіді для дійсних значень у ЦЬОМУ завданні достатньо вказати один десятковий розряд після крапки, відкинувши решту розрядів без округлення. Наприклад, замість 13.66666666666667 достатньо вказати 13.6, замість 25.23 достатньо вказати 25.2)
code
try:
    print(4, end=';')
    print(int('a9'), end=';')
    print(8, end=';')
except KeyboardInterrupt:
    print(5, end=';')
except Exception:
    print(7, end=';')
except:
    print('ERR', end=';')
else:
    print(0, end=';')
finally:
    print(1)
code

6. Що буде виведено за виконання?
code
def f(n):
    if n<=9:
        print(n%10,end='')
    else:
        f(n//10)

f(543216)
code

7. Що буде виведено за виконання?
code
class A:
    c = 1

    def _f1(self): print(2, end=';')
    def __f2(self): print(3, end=';')

    def main(self):
        try: print(self.c, end=';')
        except: print('ERR', end=';')

        try: self._f1()
        except: print('ERR', end=';')

        try: self.__f2()
        except: print('ERR', end=';')


class B(A):
    c = 4

    def _f1(self): print(5, end=';')
    def __f2(self): print(6, end=';')

try: B().main()
except: print('ERR', end=';')
code

8. Що буде виведено за виконання?
code
def f(arg, n):
    n = 8
    tmp = arg
    tmp[0:6] = [2,3]
    return len(tmp)>3

try:
    a = ['0','1','2','3','4']
    n = 3
    f(a, n)
    print(n, end=';')
    for el in a:
        print(el, end=';')
except:
    print('ERR')
code

9. Що буде виведено за виконання?
code
class A:
    b = 1
    c = 2

    def __init__(self):
        c = 3

class B(A):
    b = 4

    def __init__(self):
        super().__init__()
        self.a = 7

obj = B()
obj.c = obj.a + 10
B.c = 13

try: print(obj.c, end= ';')
except: print('ERR', end=';')

try: print(A.c, end= ';')
except: print('ERR', end=';')

try: print(B.c, end= ';')
except: print('ERR', end=';')
code

10. 
Для графа на рисунку з вершини 6 запущено алгоритм пошуку в глибину, що будує кістякове дерево. Списки суміжності вершин переглядаються алгоритмом за зростанням міток вершин.\par
\begin{center}
	\adjustimage{max size={0.9\linewidth}{0.9\paperheight}}
	{../10/Traversals/147.png}
\end{center}
\par Які з наведених ребер входять до складу отриманого кістякового дерева?\par
1) (0, 5)\par
2) (0, 4)\par
3) (5, 6)\par
4) (0, 6)\par
5) (2, 3)\par
6) (2, 4)\par
{\vskip 15pt} У формі вибрати номери відповідних ребер.\par
%answer:2 4 5 6
\end {large}}
%end
{\smallskip \begin {small} Затверджено на засіданні кафедри теоретичної кібернетики, протокол \No 12 від   19.05.2020 р.\par\smallskip\smallskip
{\parbox {6.5 cm}{Зав. кафедри \dotfill Ю.В.~Крак}}\hspace {0.7cm}{\hbox to 6.5 cm{ Екзаменатор \dotfill Т.О.~Карнаух}} \end{small}}
%begin
{{\begin {small}\par
{ \center  Київський національний університет імені Тараса Шевченка \par}
{\parbox {5 cm}{Освітня програма \hfill}{\parbox {7 cm}{Прикладна математика, Системний аналіз \hfill}}\parbox {6 cm}{\hfill Семестр 2}}\par
{\parbox {5 cm} {Навчальний предмет \hfill}\parbox {7 cm}{Програмування \hfill}\parbox {6cm}{\hfill}\par}\vspace{-7pt} \end{small}}{\center Екзаменаційний білет \No \textbf 
{ 289 }\par}}
{
\begin {large}
1. Що буде виведено за виконання?
code
try:
    res = 2**-1+1
    if isinstance(res, bool):
        print(f'bool;{res}')
    elif isinstance(res, int):
        print(f'int;{res}')
    elif isinstance(res, float):
        print(f'float;{res:.2f}')
    else:
        print('???')
except:
    print('ERR')
code

2. Що буде виведено за виконання?
code
a = 7
b = 0
c = 1

def g(b):
    global c
    a = 4
    b = 3
    c = 1
    return a + b + c

a = 8
b = 4
c = 9

try:
    print(g(a), a, b, c, sep=';')
except:
    print('ERR', a, b, c, sep=';')
code

3. Що буде виведено за виконання?
code
try:
    n = 12
    while n>9:
        n += -3
    print(n, end=';')
except:
    print('ERR')
code

4. Що буде виведено за виконання?
code
try:
    for f in range(0, -12, 3):
        if f<4:
            continue
        print(f, end=';');f = 5
    else:
        print('end',end=';')
    print(f)
except:
    print('ERR')
code

5. Що буде виведено за виконання?\par (Дійсні значення, якщо наявні, в цьому завданні будуть виводитись у форматі з фіксованою крапкою та, можливо, з доволі великою кількістю десяткових розрядів. У відповіді для дійсних значень у ЦЬОМУ завданні достатньо вказати один десятковий розряд після крапки, відкинувши решту розрядів без округлення. Наприклад, замість 13.66666666666667 достатньо вказати 13.6, замість 25.23 достатньо вказати 25.2)
code
try:
    print(2, end=';')
    print(3j<3j, end=';')
    print(4, end=';')
except KeyboardInterrupt:
    print(6, end=';')
except ValueError:
    print(7, end=';')
except:
    print('ERR', end=';')
else:
    print(4, end=';')
finally:
    print(1)
code

6. Що буде виведено за виконання?
code
def f(s):
    for el in s:
        el += 1
        return s
 
s = [1, 2, 3]
s = f(s)
print(s)
code

7. Що буде виведено за виконання?
code
class A:
    c = 1

    def _f1(self): print(2, end=';')
    def f2(self): print(3, end=';')

    def main(self):
        try: print(self.c, end=';')
        except: print('ERR', end=';')

        try: self._f1()
        except: print('ERR', end=';')

        try: self.f2()
        except: print('ERR', end=';')


class B(A):
    c = 4

    def _f1(self): print(5, end=';')
    def f2(self): print(6, end=';')

try: B().main()
except: print('ERR', end=';')
code

8. Що буде виведено за виконання?
code
def f(arg, n):
    arg[19] = 1
    n = 8
    arg = set(arg.items())

try:
    n = 3
    d = {40:2, 77:0, 78:4, 78:7}
    f(d, n)
    print(n, end=';')
    for x in d.values():
        print(x, sep=';', end=';')
except:
    print('ERR')
code

9. Що буде виведено за виконання?
code
class A:
    b = 1
    a = 2

    def __init__(self):
        b = 3

class B(A):
    a = 4

    def __init__(self):
        super().__init__()
        self.c = 7

obj = B()
obj.a = obj.b + 10
B.a = 13

try: print(obj.a, end= ';')
except: print('ERR', end=';')

try: print(A.a, end= ';')
except: print('ERR', end=';')

try: print(B.a, end= ';')
except: print('ERR', end=';')
code

10. 
Для графа на рисунку з вершини 2 запущено алгоритм пошуку в глибину, що будує кістякове дерево. Списки суміжності вершин переглядаються алгоритмом за зростанням міток вершин.\par
\begin{center}
	\adjustimage{max size={0.9\linewidth}{0.9\paperheight}}
	{../10/Traversals/152.png}
\end{center}
\par Які з наведених ребер входять до складу отриманого кістякового дерева?\par
1) (2, 4)\par
2) (0, 2)\par
3) (3, 4)\par
4) (0, 1)\par
5) (0, 3)\par
6) (0, 6)\par
{\vskip 15pt} У формі вибрати номери відповідних ребер.\par
%answer:2 3 4 5 6
\end {large}}
%end
{\smallskip \begin {small} Затверджено на засіданні кафедри теоретичної кібернетики, протокол \No 12 від   19.05.2020 р.\par\smallskip\smallskip
{\parbox {6.5 cm}{Зав. кафедри \dotfill Ю.В.~Крак}}\hspace {0.7cm}{\hbox to 6.5 cm{ Екзаменатор \dotfill Т.О.~Карнаух}} \end{small}}
%begin
{{\begin {small}\par
{ \center  Київський національний університет імені Тараса Шевченка \par}
{\parbox {5 cm}{Освітня програма \hfill}{\parbox {7 cm}{Прикладна математика, Системний аналіз \hfill}}\parbox {6 cm}{\hfill Семестр 2}}\par
{\parbox {5 cm} {Навчальний предмет \hfill}\parbox {7 cm}{Програмування \hfill}\parbox {6cm}{\hfill}\par}\vspace{-7pt} \end{small}}{\center Екзаменаційний білет \No \textbf 
{ 290 }\par}}
{
\begin {large}
1. Що буде виведено за виконання?
code
def f(n):
    print('f', end='')
    return n<=3

try:
    print(f(-4) and f(-1))
except:
    print('ERR')
code

2. Що буде виведено за виконання?
code
a = 9
b = 5
c = 4

def f(a):
    a = 1
    b = 3
    c = 2
    return a + b + c

a = 1
b = 3
c = 7

try:
    print(f(b), a, b, c, sep=';')
except:
    print('ERR', a, b, c, sep=';')
code

3. Що буде виведено за виконання?
code
try:
    k = 5
    while k<9:
        k += 1
    print(k, end=';')
except:
    print('ERR')
code

4. Що буде виведено за виконання?
code
try:
    for d in range(-9, 7, -1):
        if d<=8:
            break
        print(d, end=';')
    else:
        print(d,end=';')
    print(d)
except:
    print('ERR')
code

5. Що буде виведено за виконання?\par (Дійсні значення, якщо наявні, в цьому завданні будуть виводитись у форматі з фіксованою крапкою та, можливо, з доволі великою кількістю десяткових розрядів. У відповіді для дійсних значень у ЦЬОМУ завданні достатньо вказати один десятковий розряд після крапки, відкинувши решту розрядів без округлення. Наприклад, замість 13.66666666666667 достатньо вказати 13.6, замість 25.23 достатньо вказати 25.2)
code
try:
    print(3, end=';')
    print(int('8'), end=';')
    print(9, end=';')
except ZeroDivisionError:
    print(0, end=';')
except TypeError:
    print(5, end=';')
except:
    print('ERR', end=';')
else:
    print(6, end=';')
finally:
    print(2)
code

6. Що буде виведено за виконання?
code
def f(n):
    print(n%10, end='')
    if n>9:
        f(n//10)
 
f(543216)
code

7. Що буде виведено за виконання?
code
class A:
    d = 1

    def _h1(self): print(2, end=';')
    def _h2(self): print(3, end=';')

    def main(self):
        try: print(self.d, end=';')
        except: print('ERR', end=';')

        try: self._h1()
        except: print('ERR', end=';')

        try: self._h2()
        except: print('ERR', end=';')


class B(A):
    d = 4

    def _h1(self): print(5, end=';')
    def _h2(self): print(6, end=';')

try: B().main()
except: print('ERR', end=';')
code

8. Що буде виведено за виконання?
code
def f(arg, n):
    n = 1
    tmp = arg
    tmp[4:0] = 'bd'
    return len(tmp)>3

try:
    f = [10,11,12,13,14]
    n = 4
    f(f, n)
    print(n, end=';')
    for el in f:
        print(el, end=';')
except:
    print('ERR')
code

9. Що буде виведено за виконання?
code
class A:
    a = 1
    b = 2

    def __init__(self):
        b = 3

class B(A):
    a = 4

    def __init__(self):
        super().__init__()
        self.c = 7

obj = B()
obj.c = obj.a + 10
B.c = 13

try: print(obj.c, end= ';')
except: print('ERR', end=';')

try: print(A.c, end= ';')
except: print('ERR', end=';')

try: print(B.c, end= ';')
except: print('ERR', end=';')
code

10. 
Для графа на рисунку з вершини 2 запущено алгоритм пошуку в глибину, що будує кістякове дерево. Списки суміжності вершин переглядаються алгоритмом за зростанням міток вершин.\par
\begin{center}
	\adjustimage{max size={0.9\linewidth}{0.9\paperheight}}
	{../10/Traversals/263.png}
\end{center}
\par Які з наведених ребер входять до складу отриманого кістякового дерева?\par
1) (0, 3)\par
2) (1, 3)\par
3) (3, 4)\par
4) (2, 6)\par
5) (1, 4)\par
6) (1, 2)\par
{\vskip 15pt} У формі вибрати номери відповідних ребер.\par
%answer:1 2 3 4 6
\end {large}}
%end
{\smallskip \begin {small} Затверджено на засіданні кафедри теоретичної кібернетики, протокол \No 12 від   19.05.2020 р.\par\smallskip\smallskip
{\parbox {6.5 cm}{Зав. кафедри \dotfill Ю.В.~Крак}}\hspace {0.7cm}{\hbox to 6.5 cm{ Екзаменатор \dotfill Т.О.~Карнаух}} \end{small}}
%begin
{{\begin {small}\par
{ \center  Київський національний університет імені Тараса Шевченка \par}
{\parbox {5 cm}{Освітня програма \hfill}{\parbox {7 cm}{Прикладна математика, Системний аналіз \hfill}}\parbox {6 cm}{\hfill Семестр 2}}\par
{\parbox {5 cm} {Навчальний предмет \hfill}\parbox {7 cm}{Програмування \hfill}\parbox {6cm}{\hfill}\par}\vspace{-7pt} \end{small}}{\center Екзаменаційний білет \No \textbf 
{ 291 }\par}}
{
\begin {large}
1. Що буде виведено за виконання?
code
def f(n):
    print('f', end='')
    return n

try:
    print(f(4) or f(-4))
except:
    print('ERR')
code

2. Що буде виведено за виконання?
code
a = 2
b = 8
c = 0

def h(a):
    global c
    a = 5
    b = 4
    c = 3
    return a + b + c

a = 6
b = 1
c = 4

try:
    print(h(a), a, b, c, sep=';')
except:
    print('ERR', a, b, c, sep=';')
code

3. Що буде виведено за виконання?
code
try:
    j = 4
    while j<=12:
        j += 2
    print(j, end=';')
except:
    print('ERR')
code

4. Що буде виведено за виконання?
code
try:
    for b in range(7, -4, 3):
        if b>=7:
            continue
        print(b, end=';')
    else:
        print(b,end=';')
    print(b)
except:
    print('ERR')
code

5. Що буде виведено за виконання?\par (Дійсні значення, якщо наявні, в цьому завданні будуть виводитись у форматі з фіксованою крапкою та, можливо, з доволі великою кількістю десяткових розрядів. У відповіді для дійсних значень у ЦЬОМУ завданні достатньо вказати один десятковий розряд після крапки, відкинувши решту розрядів без округлення. Наприклад, замість 13.66666666666667 достатньо вказати 13.6, замість 25.23 достатньо вказати 25.2)
code
try:
    print(7, end=';')
    print(int('b1'), end=';')
    print(2, end=';')
except KeyboardInterrupt:
    print(9, end=';')
except Exception:
    print(0, end=';')
except:
    print('ERR', end=';')
else:
    print(5, end=';')
finally:
    print(3)
code

6. Що буде виведено за виконання?
code
def f(n):
    if n>9:
        f(n//10)
    print(n%10, end='')
 
f(543216)
code

7. Що буде виведено за виконання?
code
class A:
    b = 1

    def g1(self): print(2, end=';')
    def _g2(self): print(3, end=';')

    def main(self):
        try: print(self.b, end=';')
        except: print('ERR', end=';')

        try: self.g1()
        except: print('ERR', end=';')

        try: self._g2()
        except: print('ERR', end=';')


class B(A):
    b = 4

    def g1(self): print(5, end=';')
    def _g2(self): print(6, end=';')

try: B().main()
except: print('ERR', end=';')
code

8. Що буде виведено за виконання?
code
def f(arg, n):
    arg[31] = 1
    n *= -3

try:
    n = 5
    t = {11:5, 48:6, 20:2, 11:3}
    f(t, n)
    print(n, end=';')
    for x in t.items():
        print(*x, sep=';', end=';')
except:
    print('ERR')
code

9. Що буде виведено за виконання?
code
class A:
    a = 1
    c = 2

    def __init__(self):
        a = 3

class B(A):
    a = 4

    def __init__(self):
        super().__init__()
        self.b = 7

obj = B()
obj.b = obj.c + 10
B.b = 13

try: print(obj.b, end= ';')
except: print('ERR', end=';')

try: print(A.b, end= ';')
except: print('ERR', end=';')

try: print(B.b, end= ';')
except: print('ERR', end=';')
code

10. 
Для графа на рисунку з вершини 0 запущено алгоритм пошуку в глибину, що будує кістякове дерево. Списки суміжності вершин переглядаються алгоритмом за зростанням міток вершин.\par
\begin{center}
	\adjustimage{max size={0.9\linewidth}{0.9\paperheight}}
	{../10/Traversals/88.png}
\end{center}
\par Які з наведених ребер входять до складу отриманого кістякового дерева?\par
1) (0, 3)\par
2) (4, 6)\par
3) (0, 4)\par
4) (0, 2)\par
5) (0, 1)\par
6) (2, 6)\par
{\vskip 15pt} У формі вибрати номери відповідних ребер.\par
%answer:5 6
\end {large}}
%end
{\smallskip \begin {small} Затверджено на засіданні кафедри теоретичної кібернетики, протокол \No 12 від   19.05.2020 р.\par\smallskip\smallskip
{\parbox {6.5 cm}{Зав. кафедри \dotfill Ю.В.~Крак}}\hspace {0.7cm}{\hbox to 6.5 cm{ Екзаменатор \dotfill Т.О.~Карнаух}} \end{small}}
%begin
{{\begin {small}\par
{ \center  Київський національний університет імені Тараса Шевченка \par}
{\parbox {5 cm}{Освітня програма \hfill}{\parbox {7 cm}{Прикладна математика, Системний аналіз \hfill}}\parbox {6 cm}{\hfill Семестр 2}}\par
{\parbox {5 cm} {Навчальний предмет \hfill}\parbox {7 cm}{Програмування \hfill}\parbox {6cm}{\hfill}\par}\vspace{-7pt} \end{small}}{\center Екзаменаційний білет \No \textbf 
{ 292 }\par}}
{
\begin {large}
1. Що буде виведено за виконання?
code
try:
    res = 8.0>=6.0 and not 8==3 and False<8
    if isinstance(res, bool):
        print(f'bool;{res}')
    elif isinstance(res, int):
        print(f'int;{res}')
    elif isinstance(res, float):
        print(f'float;{res:.2f}')
    else:
        print('???')
except:
    print('ERR')
code

2. Що буде виведено за виконання?
code
a = 3
b = 6
c = 5

def g(b):
    a = 1
    b = 2
    c = 5
    return a + b + c

a = 8
b = 9
c = 0

try:
    print(g(b), a, b, c, sep=';')
except:
    print('ERR', a, b, c, sep=';')
code

3. Що буде виведено за виконання?
code
try:
    i = 2
    while i<=12:
        i += 2
    print(i, end=';')
except:
    print('ERR')
code

4. Що буде виведено за виконання?
code
try:
    for a in range(-3, -1, -3):
        if a>5:
            continue
        print(a, end=';');a = 7
    else:
        print(a,end=';')
    print(a)
except:
    print('ERR')
code

5. Що буде виведено за виконання?\par (Дійсні значення, якщо наявні, в цьому завданні будуть виводитись у форматі з фіксованою крапкою та, можливо, з доволі великою кількістю десяткових розрядів. У відповіді для дійсних значень у ЦЬОМУ завданні достатньо вказати один десятковий розряд після крапки, відкинувши решту розрядів без округлення. Наприклад, замість 13.66666666666667 достатньо вказати 13.6, замість 25.23 достатньо вказати 25.2)
code
try:
    print(4, end=';')
    print(int('6'), end=';')
    print(8, end=';')
except TypeError:
    print(0, end=';')
except ValueError:
    print(1, end=';')
except:
    print('ERR', end=';')
else:
    print(2, end=';')
finally:
    print(8)
code

6. Що буде виведено за виконання?
code
def f(s):
    s1 = s
    for i in range(len(s)):
        s1[i] += 1
        return s1

s = [1, 2, 3]
f(s)
print(s)
code

7. Що буде виведено за виконання?
code
class A:
    a = 1

    def _f1(self): print(2, end=';')
    def __f2(self): print(3, end=';')

    def main(self):
        try: print(self.a, end=';')
        except: print('ERR', end=';')

        try: self._f1()
        except: print('ERR', end=';')

        try: self.__f2()
        except: print('ERR', end=';')


class B(A):
    a = 4

    def _f1(self): print(5, end=';')
    def __f2(self): print(6, end=';')

try: B().main()
except: print('ERR', end=';')
code

8. Що буде виведено за виконання?
code
def f(arg, n):
    arg[56] = 0
    n *= -2

try:
    n = 7
    s = {39:4, 50:4, 50:4}
    f(s, n)
    print(n, end=';')
    for x in s.values():
        print(x, sep=';', end=';')
except:
    print('ERR')
code

9. Що буде виведено за виконання?
code
class A:
    b = 1
    c = 2

    def __init__(self):
        c = 3

class B(A):
    c = 4

    def __init__(self):
        super().__init__()
        self.a = 7

obj = B()
obj.b = obj.a + 10
B.b = 13

try: print(obj.b, end= ';')
except: print('ERR', end=';')

try: print(A.b, end= ';')
except: print('ERR', end=';')

try: print(B.b, end= ';')
except: print('ERR', end=';')
code

10. 
Для графа на рисунку з вершини 5 запущено алгоритм пошуку в глибину, що будує кістякове дерево. Списки суміжності вершин переглядаються алгоритмом за зростанням міток вершин.\par
\begin{center}
	\adjustimage{max size={0.9\linewidth}{0.9\paperheight}}
	{../10/Traversals/38.png}
\end{center}
\par Які з наведених ребер входять до складу отриманого кістякового дерева?\par
1) (1, 4)\par
2) (3, 6)\par
3) (0, 1)\par
4) (5, 6)\par
5) (1, 6)\par
6) (0, 2)\par
{\vskip 15pt} У формі вибрати номери відповідних ребер.\par
%answer:1 3 4
\end {large}}
%end
{\smallskip \begin {small} Затверджено на засіданні кафедри теоретичної кібернетики, протокол \No 12 від   19.05.2020 р.\par\smallskip\smallskip
{\parbox {6.5 cm}{Зав. кафедри \dotfill Ю.В.~Крак}}\hspace {0.7cm}{\hbox to 6.5 cm{ Екзаменатор \dotfill Т.О.~Карнаух}} \end{small}}
%begin
{{\begin {small}\par
{ \center  Київський національний університет імені Тараса Шевченка \par}
{\parbox {5 cm}{Освітня програма \hfill}{\parbox {7 cm}{Прикладна математика, Системний аналіз \hfill}}\parbox {6 cm}{\hfill Семестр 2}}\par
{\parbox {5 cm} {Навчальний предмет \hfill}\parbox {7 cm}{Програмування \hfill}\parbox {6cm}{\hfill}\par}\vspace{-7pt} \end{small}}{\center Екзаменаційний білет \No \textbf 
{ 293 }\par}}
{
\begin {large}
1. Що буде виведено за виконання?
code
try:
    res = 4<=5.0
    if isinstance(res, bool):
        print(f'bool;{res}')
    elif isinstance(res, int):
        print(f'int;{res}')
    elif isinstance(res, float):
        print(f'float;{res:.2f}')
    else:
        print('???')
except:
    print('ERR')
code

2. Що буде виведено за виконання?
code
a = 2
b = 7
c = 8

def h(a):
    global c
    a = 4
    b = 1
    c = 5
    return a + b + c

a = 6
b = 0
c = 5

try:
    print(h(a), a, b, c, sep=';')
except:
    print('ERR', a, b, c, sep=';')
code

3. Що буде виведено за виконання?
code
try:
    t = 3
    while t<6:
        t += 1
    print(t, end=';')
except:
    print('ERR')
code

4. Що буде виведено за виконання?
code
try:
    for c in range(8, 1, -1):
        if c>3:
            continue
        print(c, end=';')
    else:
        print(c,end=';')
    print(c)
except:
    print('ERR')
code

5. Що буде виведено за виконання?\par (Дійсні значення, якщо наявні, в цьому завданні будуть виводитись у форматі з фіксованою крапкою та, можливо, з доволі великою кількістю десяткових розрядів. У відповіді для дійсних значень у ЦЬОМУ завданні достатньо вказати один десятковий розряд після крапки, відкинувши решту розрядів без округлення. Наприклад, замість 13.66666666666667 достатньо вказати 13.6, замість 25.23 достатньо вказати 25.2)
code
try:
    print(9, end=';')
    print(7%0.0, end=';')
    print(6, end=';')
except ZeroDivisionError:
    print(4, end=';')
except ValueError:
    print(5, end=';')
except:
    print('ERR', end=';')
else:
    print(3, end=';')
finally:
    print(3)
code

6. Що буде виведено за виконання?
code
def f(n):
    if n<=9:
        print(n%10, end='')
    else:
        f(n//10)

f(987651)
code

7. Що буде виведено за виконання?
code
class A:
    d = 1

    def _g1(self): print(2, end=';')
    def g2(self): print(3, end=';')

    def main(self):
        try: print(self.d, end=';')
        except: print('ERR', end=';')

        try: self._g1()
        except: print('ERR', end=';')

        try: self.g2()
        except: print('ERR', end=';')


class B(A):
    d = 4

    def _g1(self): print(5, end=';')
    def g2(self): print(6, end=';')

try: B().main()
except: print('ERR', end=';')
code

8. Що буде виведено за виконання?
code
def f(arg, n):
    arg[69] = 8
    n += 3

try:
    n = 6
    t = {48:2, 64:4, 48:3}
    f(t, n)
    print(n, end=';')
    for x in t.values():
        print(x, sep=';', end=';')
except:
    print('ERR')
code

9. Що буде виведено за виконання?
code
class A:
    c = 1
    a = 2

    def __init__(self):
        c = 3

class B(A):
    a = 4

    def __init__(self):
        super().__init__()
        self.b = 7

obj = B()
obj.c = obj.c + 10
B.c = 13

try: print(obj.c, end= ';')
except: print('ERR', end=';')

try: print(A.c, end= ';')
except: print('ERR', end=';')

try: print(B.c, end= ';')
except: print('ERR', end=';')
code

10. 
Для графа на рисунку з вершини 1 запущено алгоритм пошуку в глибину, що будує кістякове дерево. Списки суміжності вершин переглядаються алгоритмом за зростанням міток вершин.\par
\begin{center}
	\adjustimage{max size={0.9\linewidth}{0.9\paperheight}}
	{../10/Traversals/107.png}
\end{center}
\par Які з наведених ребер входять до складу отриманого кістякового дерева?\par
1) (2, 4)\par
2) (0, 5)\par
3) (5, 6)\par
4) (0, 3)\par
5) (4, 5)\par
6) (4, 6)\par
{\vskip 15pt} У формі вибрати номери відповідних ребер.\par
%answer:1 3 4
\end {large}}
%end
{\smallskip \begin {small} Затверджено на засіданні кафедри теоретичної кібернетики, протокол \No 12 від   19.05.2020 р.\par\smallskip\smallskip
{\parbox {6.5 cm}{Зав. кафедри \dotfill Ю.В.~Крак}}\hspace {0.7cm}{\hbox to 6.5 cm{ Екзаменатор \dotfill Т.О.~Карнаух}} \end{small}}
%begin
{{\begin {small}\par
{ \center  Київський національний університет імені Тараса Шевченка \par}
{\parbox {5 cm}{Освітня програма \hfill}{\parbox {7 cm}{Прикладна математика, Системний аналіз \hfill}}\parbox {6 cm}{\hfill Семестр 2}}\par
{\parbox {5 cm} {Навчальний предмет \hfill}\parbox {7 cm}{Програмування \hfill}\parbox {6cm}{\hfill}\par}\vspace{-7pt} \end{small}}{\center Екзаменаційний білет \No \textbf 
{ 294 }\par}}
{
\begin {large}
1. Що буде виведено за виконання?
code
try:
    res = -1**4--1
    if isinstance(res, bool):
        print(f'bool;{res}')
    elif isinstance(res, int):
        print(f'int;{res}')
    elif isinstance(res, float):
        print(f'float;{res:.2f}')
    else:
        print('???')
except:
    print('ERR')
code

2. Що буде виведено за виконання?
code
a = 1
b = 2
c = 8

def g(b):
    a *= 4
    b = 2
    c = 5
    return a + b + c

a = 5
b = 7
c = 3

try:
    print(g(b), a, b, c, sep=';')
except:
    print('ERR', a, b, c, sep=';')
code

3. Що буде виведено за виконання?
code
try:
    i = 15
    while i>=2:
        i += -2
    print(i, end=';')
except:
    print('ERR')
code

4. Що буде виведено за виконання?
code
try:
    for e in range(-10, 9, -2):
        if e<6:
            break
        print(e, end=';')
        if e>=5:
            break
    else:
        print(e,end=';')
    print(e)
except:
    print('ERR')
code

5. Що буде виведено за виконання?\par (Дійсні значення, якщо наявні, в цьому завданні будуть виводитись у форматі з фіксованою крапкою та, можливо, з доволі великою кількістю десяткових розрядів. У відповіді для дійсних значень у ЦЬОМУ завданні достатньо вказати один десятковий розряд після крапки, відкинувши решту розрядів без округлення. Наприклад, замість 13.66666666666667 достатньо вказати 13.6, замість 25.23 достатньо вказати 25.2)
code
try:
    print(1, end=';')
    print(7j>=5j, end=';')
    print(4, end=';')
except ZeroDivisionError:
    print(6, end=';')
except Exception:
    print(8, end=';')
except:
    print('ERR', end=';')
else:
    print(5, end=';')
finally:
    print(2)
code

6. Що буде виведено за виконання?
code
def f(n):
    if n>9:
        return f(n//10)+1
    else:
        return 1

print(f(123456))
code

7. Що буде виведено за виконання?
code
class A:
    c = 1

    def _h1(self): print(2, end=';')
    def _h2(self): print(3, end=';')

    def main(self):
        try: print(self.c, end=';')
        except: print('ERR', end=';')

        try: self._h1()
        except: print('ERR', end=';')

        try: self._h2()
        except: print('ERR', end=';')


class B(A):
    c = 4

    def _h1(self): print(5, end=';')
    def _h2(self): print(6, end=';')

try: B().main()
except: print('ERR', end=';')
code

8. Що буде виведено за виконання?
code
def f(arg, n):
    arg[89] = 2
    n = 5
    arg = tuple(arg.keys())

try:
    n = 1
    t = {43:6, 85:4, 89:9, 43:5}
    f(t, n)
    print(n, end=';')
    for x in t.values():
        print(x, sep=';', end=';')
except:
    print('ERR')
code

9. Що буде виведено за виконання?
code
class A:
    c = 1
    b = 2

    def __init__(self):
        b = 3

class B(A):
    c = 4

    def __init__(self):
        super().__init__()
        self.a = 7

obj = B()
obj.a = obj.b + 10
B.a = 13

try: print(obj.a, end= ';')
except: print('ERR', end=';')

try: print(A.a, end= ';')
except: print('ERR', end=';')

try: print(B.a, end= ';')
except: print('ERR', end=';')
code

10. 
Для графа на рисунку з вершини 4 запущено алгоритм пошуку в глибину, що будує кістякове дерево. Списки суміжності вершин переглядаються алгоритмом за зростанням міток вершин.\par
\begin{center}
	\adjustimage{max size={0.9\linewidth}{0.9\paperheight}}
	{../10/Traversals/299.png}
\end{center}
\par Які з наведених ребер входять до складу отриманого кістякового дерева?\par
1) (2, 5)\par
2) (4, 6)\par
3) (1, 4)\par
4) (1, 2)\par
5) (2, 4)\par
6) (4, 5)\par
{\vskip 15pt} У формі вибрати номери відповідних ребер.\par
%answer:1 4
\end {large}}
%end
{\smallskip \begin {small} Затверджено на засіданні кафедри теоретичної кібернетики, протокол \No 12 від   19.05.2020 р.\par\smallskip\smallskip
{\parbox {6.5 cm}{Зав. кафедри \dotfill Ю.В.~Крак}}\hspace {0.7cm}{\hbox to 6.5 cm{ Екзаменатор \dotfill Т.О.~Карнаух}} \end{small}}
%begin
{{\begin {small}\par
{ \center  Київський національний університет імені Тараса Шевченка \par}
{\parbox {5 cm}{Освітня програма \hfill}{\parbox {7 cm}{Прикладна математика, Системний аналіз \hfill}}\parbox {6 cm}{\hfill Семестр 2}}\par
{\parbox {5 cm} {Навчальний предмет \hfill}\parbox {7 cm}{Програмування \hfill}\parbox {6cm}{\hfill}\par}\vspace{-7pt} \end{small}}{\center Екзаменаційний білет \No \textbf 
{ 295 }\par}}
{
\begin {large}
1. Що буде виведено за виконання?
code
def f(n):
    print('f', end='')
    return n>1

try:
    print(f(-6) or f(-8))
except:
    print('ERR')
code

2. Що буде виведено за виконання?
code
a = 6
b = 5
c = 7

def f(a):
    global c
    a -= 1
    b = 4
    c = 5
    return a + b + c

a = 2
b = 0
c = 4

try:
    print(f(a), a, b, c, sep=';')
except:
    print('ERR', a, b, c, sep=';')
code

3. Що буде виведено за виконання?
code
try:
    k = 13
    while k>=3:
        k += -2
    print(k, end=';')
except:
    print('ERR')
code

4. Що буде виведено за виконання?
code
try:
    for a in range(-2, 5, 2):
        if a>=8:
            break
        print(a, end=';')
    else:
        print('end',end=';')
    print(a)
except:
    print('ERR')
code

5. Що буде виведено за виконання?\par (Дійсні значення, якщо наявні, в цьому завданні будуть виводитись у форматі з фіксованою крапкою та, можливо, з доволі великою кількістю десяткових розрядів. У відповіді для дійсних значень у ЦЬОМУ завданні достатньо вказати один десятковий розряд після крапки, відкинувши решту розрядів без округлення. Наприклад, замість 13.66666666666667 достатньо вказати 13.6, замість 25.23 достатньо вказати 25.2)
code
try:
    print(9, end=';')
    print(0//3, end=';')
    print(4, end=';')
except TypeError:
    print(1, end=';')
except KeyboardInterrupt:
    print(6, end=';')
except:
    print('ERR', end=';')
else:
    print(7, end=';')
finally:
    print(5)
code

6. Що буде виведено за виконання?
code
def f(n):
    if n<=9:
        print(n%10, end='')
    else:
        f(n//10)

f(123456)
code

7. Що буде виведено за виконання?
code
class A:
    b = 1

    def __f1(self): print(2, end=';')
    def _f2(self): print(3, end=';')

    def main(self):
        try: print(self.b, end=';')
        except: print('ERR', end=';')

        try: self.__f1()
        except: print('ERR', end=';')

        try: self._f2()
        except: print('ERR', end=';')


class B(A):
    b = 4

    def __f1(self): print(5, end=';')
    def _f2(self): print(6, end=';')

try: B().main()
except: print('ERR', end=';')
code

8. Що буде виведено за виконання?
code
def f(arg, n):
    arg[39] = 0
    n = 6
    arg = tuple(arg.items())

try:
    n = 2
    s = {43:0, 53:8, 45:6, 43:7}
    f(s, n)
    print(n, end=';')
    for x in s.values():
        print(x, sep=';', end=';')
except:
    print('ERR')
code

9. Що буде виведено за виконання?
code
class A:
    c = 1
    a = 2

    def __init__(self):
        c = 3

class B(A):
    a = 4

    def __init__(self):
        super().__init__()
        self.b = 7

obj = B()
obj.a = obj.b + 10
B.a = 13

try: print(obj.a, end= ';')
except: print('ERR', end=';')

try: print(A.a, end= ';')
except: print('ERR', end=';')

try: print(B.a, end= ';')
except: print('ERR', end=';')
code

10. 
Для графа на рисунку з вершини 6 запущено алгоритм пошуку в глибину, що будує кістякове дерево. Списки суміжності вершин переглядаються алгоритмом за зростанням міток вершин.\par
\begin{center}
	\adjustimage{max size={0.9\linewidth}{0.9\paperheight}}
	{../10/Traversals/262.png}
\end{center}
\par Які з наведених ребер входять до складу отриманого кістякового дерева?\par
1) (4, 6)\par
2) (2, 5)\par
3) (0, 4)\par
4) (3, 5)\par
5) (1, 5)\par
6) (1, 6)\par
{\vskip 15pt} У формі вибрати номери відповідних ребер.\par
%answer:2 3 4 6
\end {large}}
%end
{\smallskip \begin {small} Затверджено на засіданні кафедри теоретичної кібернетики, протокол \No 12 від   19.05.2020 р.\par\smallskip\smallskip
{\parbox {6.5 cm}{Зав. кафедри \dotfill Ю.В.~Крак}}\hspace {0.7cm}{\hbox to 6.5 cm{ Екзаменатор \dotfill Т.О.~Карнаух}} \end{small}}
%begin
{{\begin {small}\par
{ \center  Київський національний університет імені Тараса Шевченка \par}
{\parbox {5 cm}{Освітня програма \hfill}{\parbox {7 cm}{Прикладна математика, Системний аналіз \hfill}}\parbox {6 cm}{\hfill Семестр 2}}\par
{\parbox {5 cm} {Навчальний предмет \hfill}\parbox {7 cm}{Програмування \hfill}\parbox {6cm}{\hfill}\par}\vspace{-7pt} \end{small}}{\center Екзаменаційний білет \No \textbf 
{ 296 }\par}}
{
\begin {large}
1. Що буде виведено за виконання?
code
try:
    res = 8//6*8*6
    if isinstance(res, bool):
        print(f'bool;{res}')
    elif isinstance(res, int):
        print(f'int;{res}')
    elif isinstance(res, float):
        print(f'float;{res:.2f}')
    else:
        print('???')
except:
    print('ERR')
code

2. Що буде виведено за виконання?
code
a = 9
b = 8
c = 3

def h(a):
    a += 3
    b = 2
    c = 1
    return a + b + c

a = 1
b = 0
c = 4

try:
    print(h(a), a, b, c, sep=';')
except:
    print('ERR', a, b, c, sep=';')
code

3. Що буде виведено за виконання?
code
try:
    j = 8
    while j<=5:
        j += 1
    print(j, end=';')
except:
    print('ERR')
code

4. Що буде виведено за виконання?
code
try:
    for c in range(12, 11, -2):
        if c<=6:
            continue
        print(c, end=';');c = 8
    else:
        print(c,end=';')
    print(c)
except:
    print('ERR')
code

5. Що буде виведено за виконання?\par (Дійсні значення, якщо наявні, в цьому завданні будуть виводитись у форматі з фіксованою крапкою та, можливо, з доволі великою кількістю десяткових розрядів. У відповіді для дійсних значень у ЦЬОМУ завданні достатньо вказати один десятковий розряд після крапки, відкинувши решту розрядів без округлення. Наприклад, замість 13.66666666666667 достатньо вказати 13.6, замість 25.23 достатньо вказати 25.2)
code
try:
    print(3, end=';')
    print(9/2, end=';')
    print(8, end=';')
except KeyboardInterrupt:
    print(0, end=';')
except ZeroDivisionError:
    print(2, end=';')
except:
    print('ERR', end=';')
else:
    print(2, end=';')
finally:
    print(7)
code

6. Що буде виведено за виконання?
code
def f(n):
    if n>9:
        return f(n//100)+1
    else:
        return 1

print(f(123456))
code

7. Що буде виведено за виконання?
code
class A:
    a = 1

    def _g1(self): print(2, end=';')
    def _g2(self): print(3, end=';')

    def main(self):
        try: print(self.a, end=';')
        except: print('ERR', end=';')

        try: self._g1()
        except: print('ERR', end=';')

        try: self._g2()
        except: print('ERR', end=';')


class B(A):
    a = 4

    def _g1(self): print(5, end=';')
    def _g2(self): print(6, end=';')

try: B().main()
except: print('ERR', end=';')
code

8. Що буде виведено за виконання?
code
def f(arg, n):
    arg[86] = 2
    n = 6
    arg = list(arg.keys())

try:
    n = 3
    d = {79:6, 10:3, 28:3, 79:8}
    f(d, n)
    print(n, end=';')
    for x, y in d.items():
        print(x, y, sep=';', end=';')
except:
    print('ERR')
code

9. Що буде виведено за виконання?
code
class A:
    c = 1
    b = 2

    def __init__(self):
        b = 3

class B(A):
    c = 4

    def __init__(self):
        super().__init__()
        self.a = 7

obj = B()
obj.c = obj.b + 10
B.c = 13

try: print(obj.c, end= ';')
except: print('ERR', end=';')

try: print(A.c, end= ';')
except: print('ERR', end=';')

try: print(B.c, end= ';')
except: print('ERR', end=';')
code

10. 
Для графа на рисунку з вершини 6 запущено алгоритм пошуку в глибину, що будує кістякове дерево. Списки суміжності вершин переглядаються алгоритмом за зростанням міток вершин.\par
\begin{center}
	\adjustimage{max size={0.9\linewidth}{0.9\paperheight}}
	{../10/Traversals/40.png}
\end{center}
\par Які з наведених ребер входять до складу отриманого кістякового дерева?\par
1) (0, 1)\par
2) (2, 3)\par
3) (2, 4)\par
4) (0, 4)\par
5) (1, 4)\par
6) (4, 5)\par
{\vskip 15pt} У формі вибрати номери відповідних ребер.\par
%answer:1 2 3 5 6
\end {large}}
%end
{\smallskip \begin {small} Затверджено на засіданні кафедри теоретичної кібернетики, протокол \No 12 від   19.05.2020 р.\par\smallskip\smallskip
{\parbox {6.5 cm}{Зав. кафедри \dotfill Ю.В.~Крак}}\hspace {0.7cm}{\hbox to 6.5 cm{ Екзаменатор \dotfill Т.О.~Карнаух}} \end{small}}
%begin
{{\begin {small}\par
{ \center  Київський національний університет імені Тараса Шевченка \par}
{\parbox {5 cm}{Освітня програма \hfill}{\parbox {7 cm}{Прикладна математика, Системний аналіз \hfill}}\parbox {6 cm}{\hfill Семестр 2}}\par
{\parbox {5 cm} {Навчальний предмет \hfill}\parbox {7 cm}{Програмування \hfill}\parbox {6cm}{\hfill}\par}\vspace{-7pt} \end{small}}{\center Екзаменаційний білет \No \textbf 
{ 297 }\par}}
{
\begin {large}
1. Що буде виведено за виконання?
code
try:
    res = 4/5/7*8
    if isinstance(res, bool):
        print(f'bool;{res}')
    elif isinstance(res, int):
        print(f'int;{res}')
    elif isinstance(res, float):
        print(f'float;{res:.2f}')
    else:
        print('???')
except:
    print('ERR')
code

2. Що буде виведено за виконання?
code
a = 7
b = 2
c = 1

def f(b):
    a = 4
    b = 2
    c = 3
    return a + b + c

a = 9
b = 4
c = 3

try:
    print(f(b), a, b, c, sep=';')
except:
    print('ERR', a, b, c, sep=';')
code

3. Що буде виведено за виконання?
code
try:
    m = 1
    while m<10:
        m += 2
    print(m, end=';')
except:
    print('ERR')
code

4. Що буде виведено за виконання?
code
try:
    for f in range(10, -5, -1):
        if f>=5:
            continue
        print(f, end=';')
    else:
        print(f,end=';')
    print(f)
except:
    print('ERR')
code

5. Що буде виведено за виконання?\par (Дійсні значення, якщо наявні, в цьому завданні будуть виводитись у форматі з фіксованою крапкою та, можливо, з доволі великою кількістю десяткових розрядів. У відповіді для дійсних значень у ЦЬОМУ завданні достатньо вказати один десятковий розряд після крапки, відкинувши решту розрядів без округлення. Наприклад, замість 13.66666666666667 достатньо вказати 13.6, замість 25.23 достатньо вказати 25.2)
code
try:
    print(1, end=';')
    print(4%0, end=';')
    print(8, end=';')
except Exception:
    print(0, end=';')
except TypeError:
    print(3, end=';')
except:
    print('ERR', end=';')
else:
    print(6, end=';')
finally:
    print(9)
code

6. Що буде виведено за виконання?
code
def f(n):
    print(n%2, end='')
    if n>2:
        f(n//2)
 
f(16)
code

7. Що буде виведено за виконання?
code
class A:
    b = 1

    def __h1(self): print(2, end=';')
    def _h2(self): print(3, end=';')

    def main(self):
        try: print(self.b, end=';')
        except: print('ERR', end=';')

        try: self.__h1()
        except: print('ERR', end=';')

        try: self._h2()
        except: print('ERR', end=';')


class B(A):
    b = 4

    def __h1(self): print(5, end=';')
    def _h2(self): print(6, end=';')

try: B().main()
except: print('ERR', end=';')
code

8. Що буде виведено за виконання?
code
def f(arg, n):
    arg[63] = 7
    n = 5
    arg = set(arg.values())

try:
    n = 4
    t = {57:7, 14:3, 77:6, 49:6}
    f(t, n)
    print(n, end=';')
    for x in t.keys():
        print(x, sep=';', end=';')
except:
    print('ERR')
code

9. Що буде виведено за виконання?
code
class A:
    b = 1
    a = 2

    def __init__(self):
        a = 3

class B(A):
    b = 4

    def __init__(self):
        super().__init__()
        self.c = 7

obj = B()
obj.a = obj.c + 10
B.a = 13

try: print(obj.a, end= ';')
except: print('ERR', end=';')

try: print(A.a, end= ';')
except: print('ERR', end=';')

try: print(B.a, end= ';')
except: print('ERR', end=';')
code

10. 
Для графа на рисунку з вершини 6 запущено алгоритм пошуку в глибину, що будує кістякове дерево. Списки суміжності вершин переглядаються алгоритмом за зростанням міток вершин.\par
\begin{center}
	\adjustimage{max size={0.9\linewidth}{0.9\paperheight}}
	{../10/Traversals/290.png}
\end{center}
\par Які з наведених ребер входять до складу отриманого кістякового дерева?\par
1) (2, 3)\par
2) (0, 3)\par
3) (1, 3)\par
4) (4, 6)\par
5) (2, 4)\par
6) (0, 6)\par
{\vskip 15pt} У формі вибрати номери відповідних ребер.\par
%answer:1 2 3 5 6
\end {large}}
%end
{\smallskip \begin {small} Затверджено на засіданні кафедри теоретичної кібернетики, протокол \No 12 від   19.05.2020 р.\par\smallskip\smallskip
{\parbox {6.5 cm}{Зав. кафедри \dotfill Ю.В.~Крак}}\hspace {0.7cm}{\hbox to 6.5 cm{ Екзаменатор \dotfill Т.О.~Карнаух}} \end{small}}
%begin
{{\begin {small}\par
{ \center  Київський національний університет імені Тараса Шевченка \par}
{\parbox {5 cm}{Освітня програма \hfill}{\parbox {7 cm}{Прикладна математика, Системний аналіз \hfill}}\parbox {6 cm}{\hfill Семестр 2}}\par
{\parbox {5 cm} {Навчальний предмет \hfill}\parbox {7 cm}{Програмування \hfill}\parbox {6cm}{\hfill}\par}\vspace{-7pt} \end{small}}{\center Екзаменаційний білет \No \textbf 
{ 298 }\par}}
{
\begin {large}
1. Що буде виведено за виконання?
code
def f(n):
    print('f', end='')
    return n

try:
    print(f(-7) and f(1))
except:
    print('ERR')
code

2. Що буде виведено за виконання?
code
a = 2
b = 6
c = 7

def h(b):
    global c
    a += 2
    b = 5
    c = 4
    return a + b + c

a = 8
b = 3
c = 9

try:
    print(h(b), a, b, c, sep=';')
except:
    print('ERR', a, b, c, sep=';')
code

3. Що буде виведено за виконання?
code
try:
    n = 0
    while n<13:
        n += 2
    print(n, end=';')
except:
    print('ERR')
code

4. Що буде виведено за виконання?
code
try:
    for b in range(3, 8, 2):
        if b>7:
            continue
        print(b, end=';')
    else:
        print(b,end=';')
    print(b)
except:
    print('ERR')
code

5. Що буде виведено за виконання?\par (Дійсні значення, якщо наявні, в цьому завданні будуть виводитись у форматі з фіксованою крапкою та, можливо, з доволі великою кількістю десяткових розрядів. У відповіді для дійсних значень у ЦЬОМУ завданні достатньо вказати один десятковий розряд після крапки, відкинувши решту розрядів без округлення. Наприклад, замість 13.66666666666667 достатньо вказати 13.6, замість 25.23 достатньо вказати 25.2)
code
try:
    print(5, end=';')
    print(8j==9j, end=';')
    print(7, end=';')
except ValueError:
    print(2, end=';')
except Exception:
    print(5, end=';')
except:
    print('ERR', end=';')
else:
    print(1, end=';')
finally:
    print(6)
code

6. Що буде виведено за виконання?
code
def f(s):
    for i in range(len(s)):
        s[i] += 1
        return
 
s = [1, 2, 3]
f(s)
print(s)
code

7. Що буде виведено за виконання?
code
class A:
    d = 1

    def _f1(self): print(2, end=';')
    def __f2(self): print(3, end=';')

    def main(self):
        try: print(self.d, end=';')
        except: print('ERR', end=';')

        try: self._f1()
        except: print('ERR', end=';')

        try: self.__f2()
        except: print('ERR', end=';')


class B(A):
    d = 4

    def _f1(self): print(5, end=';')
    def __f2(self): print(6, end=';')

try: B().main()
except: print('ERR', end=';')
code

8. Що буде виведено за виконання?
code
def f(arg, n):
    arg[72] = 3
    n = 7
    arg = list(arg.keys())

try:
    n = 2
    t = {36:4, 72:3, 72:7}
    f(t, n)
    print(n, end=';')
    for x in t.items():
        print(x, sep=';', end=';')
except:
    print('ERR')
code

9. Що буде виведено за виконання?
code
class A:
    a = 1
    b = 2

    def __init__(self):
        a = 3

class B(A):
    b = 4

    def __init__(self):
        super().__init__()
        self.c = 7

obj = B()
obj.a = obj.b + 10
B.a = 13

try: print(obj.a, end= ';')
except: print('ERR', end=';')

try: print(A.a, end= ';')
except: print('ERR', end=';')

try: print(B.a, end= ';')
except: print('ERR', end=';')
code

10. 
Для графа на рисунку з вершини 5 запущено алгоритм пошуку в глибину, що будує кістякове дерево. Списки суміжності вершин переглядаються алгоритмом за зростанням міток вершин.\par
\begin{center}
	\adjustimage{max size={0.9\linewidth}{0.9\paperheight}}
	{../10/Traversals/31.png}
\end{center}
\par Які з наведених ребер входять до складу отриманого кістякового дерева?\par
1) (0, 2)\par
2) (3, 6)\par
3) (0, 6)\par
4) (4, 6)\par
5) (2, 6)\par
6) (1, 6)\par
{\vskip 15pt} У формі вибрати номери відповідних ребер.\par
%answer:1 2 6
\end {large}}
%end
{\smallskip \begin {small} Затверджено на засіданні кафедри теоретичної кібернетики, протокол \No 12 від   19.05.2020 р.\par\smallskip\smallskip
{\parbox {6.5 cm}{Зав. кафедри \dotfill Ю.В.~Крак}}\hspace {0.7cm}{\hbox to 6.5 cm{ Екзаменатор \dotfill Т.О.~Карнаух}} \end{small}}
%begin
{{\begin {small}\par
{ \center  Київський національний університет імені Тараса Шевченка \par}
{\parbox {5 cm}{Освітня програма \hfill}{\parbox {7 cm}{Прикладна математика, Системний аналіз \hfill}}\parbox {6 cm}{\hfill Семестр 2}}\par
{\parbox {5 cm} {Навчальний предмет \hfill}\parbox {7 cm}{Програмування \hfill}\parbox {6cm}{\hfill}\par}\vspace{-7pt} \end{small}}{\center Екзаменаційний білет \No \textbf 
{ 299 }\par}}
{
\begin {large}
1. Що буде виведено за виконання?
code
try:
    res = 1**-1.0*False
    if isinstance(res, bool):
        print(f'bool;{res}')
    elif isinstance(res, int):
        print(f'int;{res}')
    elif isinstance(res, float):
        print(f'float;{res:.2f}')
    else:
        print('???')
except:
    print('ERR')
code

2. Що буде виведено за виконання?
code
a = 1
b = 5
c = 9

def g(a):
    global c
    a -= 3
    b = 1
    c = 5
    return a + b + c

a = 1
b = 7
c = 0

try:
    print(g(b), a, b, c, sep=';')
except:
    print('ERR', a, b, c, sep=';')
code

3. Що буде виведено за виконання?
code
try:
    k = 6
    while k<=8:
        k += 1
    print(k, end=';')
except:
    print('ERR')
code

4. Що буде виведено за виконання?
code
try:
    for d in range(-11, -3, -3):
        if d<=4:
            continue
        print(d, end=';');d = 6
    else:
        print(d,end=';')
    print(d)
except:
    print('ERR')
code

5. Що буде виведено за виконання?\par (Дійсні значення, якщо наявні, в цьому завданні будуть виводитись у форматі з фіксованою крапкою та, можливо, з доволі великою кількістю десяткових розрядів. У відповіді для дійсних значень у ЦЬОМУ завданні достатньо вказати один десятковий розряд після крапки, відкинувши решту розрядів без округлення. Наприклад, замість 13.66666666666667 достатньо вказати 13.6, замість 25.23 достатньо вказати 25.2)
code
try:
    print(0, end=';')
    print(int('4'), end=';')
    print(3, end=';')
except ValueError:
    print(9, end=';')
except ZeroDivisionError:
    print(8, end=';')
except:
    print('ERR', end=';')
else:
    print(7, end=';')
finally:
    print(9)
code

6. Що буде виведено за виконання?
code
def f(s1):
    s = s1[:]
    for i in range(len(s)):
        s[i] += 1
 
s = [1, 2, 3]
f(s)
print(s)
code

7. Що буде виведено за виконання?
code
class A:
    a = 1

    def _h1(self): print(2, end=';')
    def h2(self): print(3, end=';')

    def main(self):
        try: print(self.a, end=';')
        except: print('ERR', end=';')

        try: self._h1()
        except: print('ERR', end=';')

        try: self.h2()
        except: print('ERR', end=';')


class B(A):
    a = 4

    def _h1(self): print(5, end=';')
    def h2(self): print(6, end=';')

try: B().main()
except: print('ERR', end=';')
code

8. Що буде виведено за виконання?
code
def f(arg, n):
    n = 7
    tmp = arg
    tmp[5:4] = [2,3]
    return len(tmp)>3

try:
    h = ['0','1','2','3','4']
    n = 5
    f(h, n)
    print(n, end=';')
    for el in h:
        print(el, end=';')
except:
    print('ERR')
code

9. Що буде виведено за виконання?
code
class A:
    b = 1
    c = 2

    def __init__(self):
        b = 3

class B(A):
    b = 4

    def __init__(self):
        super().__init__()
        self.a = 7

obj = B()
obj.c = obj.c + 10
B.c = 13

try: print(obj.c, end= ';')
except: print('ERR', end=';')

try: print(A.c, end= ';')
except: print('ERR', end=';')

try: print(B.c, end= ';')
except: print('ERR', end=';')
code

10. 
Для графа на рисунку з вершини 0 запущено алгоритм пошуку в глибину, що будує кістякове дерево. Списки суміжності вершин переглядаються алгоритмом за зростанням міток вершин.\par
\begin{center}
	\adjustimage{max size={0.9\linewidth}{0.9\paperheight}}
	{../10/Traversals/156.png}
\end{center}
\par Які з наведених ребер входять до складу отриманого кістякового дерева?\par
1) (2, 4)\par
2) (0, 3)\par
3) (0, 4)\par
4) (0, 6)\par
5) (2, 5)\par
6) (1, 5)\par
{\vskip 15pt} У формі вибрати номери відповідних ребер.\par
%answer:1 4 5
\end {large}}
%end
{\smallskip \begin {small} Затверджено на засіданні кафедри теоретичної кібернетики, протокол \No 12 від   19.05.2020 р.\par\smallskip\smallskip
{\parbox {6.5 cm}{Зав. кафедри \dotfill Ю.В.~Крак}}\hspace {0.7cm}{\hbox to 6.5 cm{ Екзаменатор \dotfill Т.О.~Карнаух}} \end{small}}
%begin
{{\begin {small}\par
{ \center  Київський національний університет імені Тараса Шевченка \par}
{\parbox {5 cm}{Освітня програма \hfill}{\parbox {7 cm}{Прикладна математика, Системний аналіз \hfill}}\parbox {6 cm}{\hfill Семестр 2}}\par
{\parbox {5 cm} {Навчальний предмет \hfill}\parbox {7 cm}{Програмування \hfill}\parbox {6cm}{\hfill}\par}\vspace{-7pt} \end{small}}{\center Екзаменаційний білет \No \textbf 
{ 300 }\par}}
{
\begin {large}
1. Що буде виведено за виконання?
code
try:
    res = not 4!=True or 5.0>3 or 3.0<=6
    if isinstance(res, bool):
        print(f'bool;{res}')
    elif isinstance(res, int):
        print(f'int;{res}')
    elif isinstance(res, float):
        print(f'float;{res:.2f}')
    else:
        print('???')
except:
    print('ERR')
code

2. Що буде виведено за виконання?
code
a = 2
b = 4
c = 8

def f(b):
    a = 3
    b *= 2
    c = 4
    return a + b + c

a = 6
b = 5
c = 3

try:
    print(f(a), a, b, c, sep=';')
except:
    print('ERR', a, b, c, sep=';')
code

3. Що буде виведено за виконання?
code
try:
    t = 10
    while t>7:
        t += -3
    print(t, end=';')
except:
    print('ERR')
code

4. Що буде виведено за виконання?
code
try:
    for e in range(4, -8, 3):
        if e<3:
            break
        print(e, end=';')
    else:
        print(e,end=';')
    print(e)
except:
    print('ERR')
code

5. Що буде виведено за виконання?\par (Дійсні значення, якщо наявні, в цьому завданні будуть виводитись у форматі з фіксованою крапкою та, можливо, з доволі великою кількістю десяткових розрядів. У відповіді для дійсних значень у ЦЬОМУ завданні достатньо вказати один десятковий розряд після крапки, відкинувши решту розрядів без округлення. Наприклад, замість 13.66666666666667 достатньо вказати 13.6, замість 25.23 достатньо вказати 25.2)
code
try:
    print(6, end=';')
    print(int('c1'), end=';')
    print(5, end=';')
except TypeError:
    print(4, end=';')
except KeyboardInterrupt:
    print(3, end=';')
except:
    print('ERR', end=';')
else:
    print(2, end=';')
finally:
    print(0)
code

6. Що буде виведено за виконання?
code
def f(n):
    if n>9:
        f(n//10)
    print(n%10, end='')
 
f(123456)
code

7. Що буде виведено за виконання?
code
class A:
    c = 1

    def g1(self): print(2, end=';')
    def _g2(self): print(3, end=';')

    def main(self):
        try: print(self.c, end=';')
        except: print('ERR', end=';')

        try: self.g1()
        except: print('ERR', end=';')

        try: self._g2()
        except: print('ERR', end=';')


class B(A):
    c = 4

    def g1(self): print(5, end=';')
    def _g2(self): print(6, end=';')

try: B().main()
except: print('ERR', end=';')
code

8. Що буде виведено за виконання?
code
def f(arg, n):
    arg[55] = 3
    n = 8
    arg = set(arg.items())

try:
    n = 0
    d = {46:8, 52:2, 46:2}
    f(d, n)
    print(n, end=';')
    for x in d.keys():
        print(x, sep=';', end=';')
except:
    print('ERR')
code

9. Що буде виведено за виконання?
code
class A:
    a = 1
    c = 2

    def __init__(self):
        c = 3

class B(A):
    c = 4

    def __init__(self):
        super().__init__()
        self.b = 7

obj = B()
obj.a = obj.b + 10
B.a = 13

try: print(obj.a, end= ';')
except: print('ERR', end=';')

try: print(A.a, end= ';')
except: print('ERR', end=';')

try: print(B.a, end= ';')
except: print('ERR', end=';')
code

10. 
Для графа на рисунку з вершини 1 запущено алгоритм пошуку в глибину, що будує кістякове дерево. Списки суміжності вершин переглядаються алгоритмом за зростанням міток вершин.\par
\begin{center}
	\adjustimage{max size={0.9\linewidth}{0.9\paperheight}}
	{../10/Traversals/119.png}
\end{center}
\par Які з наведених ребер входять до складу отриманого кістякового дерева?\par
1) (2, 4)\par
2) (2, 6)\par
3) (1, 4)\par
4) (0, 2)\par
5) (0, 5)\par
6) (4, 6)\par
{\vskip 15pt} У формі вибрати номери відповідних ребер.\par
%answer:4 5 6
\end {large}}
%end
{\smallskip \begin {small} Затверджено на засіданні кафедри теоретичної кібернетики, протокол \No 12 від   19.05.2020 р.\par\smallskip\smallskip
{\parbox {6.5 cm}{Зав. кафедри \dotfill Ю.В.~Крак}}\hspace {0.7cm}{\hbox to 6.5 cm{ Екзаменатор \dotfill Т.О.~Карнаух}} \end{small}}
%begin
{{\begin {small}\par
{ \center  Київський національний університет імені Тараса Шевченка \par}
{\parbox {5 cm}{Освітня програма \hfill}{\parbox {7 cm}{Прикладна математика, Системний аналіз \hfill}}\parbox {6 cm}{\hfill Семестр 2}}\par
{\parbox {5 cm} {Навчальний предмет \hfill}\parbox {7 cm}{Програмування \hfill}\parbox {6cm}{\hfill}\par}\vspace{-7pt} \end{small}}{\center Екзаменаційний білет \No \textbf 
{ 301 }\par}}
{
\begin {large}
1. Що буде виведено за виконання?
code
try:
    res = 4**1.0*1
    if isinstance(res, bool):
        print(f'bool;{res}')
    elif isinstance(res, int):
        print(f'int;{res}')
    elif isinstance(res, float):
        print(f'float;{res:.2f}')
    else:
        print('???')
except:
    print('ERR')
code

2. Що буде виведено за виконання?
code
a = 6
b = 5
c = 9

def g(b):
    a += 4
    b = 5
    c = 1
    return a + b + c

a = 1
b = 3
c = 4

try:
    print(g(b), a, b, c, sep=';')
except:
    print('ERR', a, b, c, sep=';')
code

3. Що буде виведено за виконання?
code
try:
    n = 0
    while n<=9:
        n += 2
    print(n, end=';')
except:
    print('ERR')
code

4. Що буде виведено за виконання?
code
try:
    for f in range(11, -10, -3):
        if f>=4:
            continue
        print(f, end=';')
    else:
        print(f,end=';')
    print(f)
except:
    print('ERR')
code

5. Що буде виведено за виконання?\par (Дійсні значення, якщо наявні, в цьому завданні будуть виводитись у форматі з фіксованою крапкою та, можливо, з доволі великою кількістю десяткових розрядів. У відповіді для дійсних значень у ЦЬОМУ завданні достатньо вказати один десятковий розряд після крапки, відкинувши решту розрядів без округлення. Наприклад, замість 13.66666666666667 достатньо вказати 13.6, замість 25.23 достатньо вказати 25.2)
code
try:
    print(9, end=';')
    print(7j!=4j, end=';')
    print(5, end=';')
except TypeError:
    print(4, end=';')
except ValueError:
    print(2, end=';')
except:
    print('ERR', end=';')
else:
    print(3, end=';')
finally:
    print(0)
code

6. Що буде виведено за виконання?
code
def f(n):
    if n>9:
        f(n//10)
    else:
        print(n%10, end='')

f(543216)
code

7. Що буде виведено за виконання?
code
class A:
    b = 1

    def __g1(self): print(2, end=';')
    def _g2(self): print(3, end=';')

    def main(self):
        try: print(self.b, end=';')
        except: print('ERR', end=';')

        try: self.__g1()
        except: print('ERR', end=';')

        try: self._g2()
        except: print('ERR', end=';')


class B(A):
    b = 4

    def __g1(self): print(5, end=';')
    def _g2(self): print(6, end=';')

try: B().main()
except: print('ERR', end=';')
code

8. Що буде виведено за виконання?
code
def f(arg, n):
    n = 0
    tmp = arg
    del tmp[:1]
    return len(tmp)>3

try:
    d = [10,11,12,13,14]
    n = 2
    f(d, n)
    print(n, end=';')
    for el in d:
        print(el, end=';')
except:
    print('ERR')
code

9. Що буде виведено за виконання?
code
class A:
    b = 1
    c = 2

    def __init__(self):
        b = 3

class B(A):
    b = 4

    def __init__(self):
        super().__init__()
        self.a = 7

obj = B()
obj.c = obj.b + 10
B.c = 13

try: print(obj.c, end= ';')
except: print('ERR', end=';')

try: print(A.c, end= ';')
except: print('ERR', end=';')

try: print(B.c, end= ';')
except: print('ERR', end=';')
code

10. 
Для графа на рисунку з вершини 6 запущено алгоритм пошуку в глибину, що будує кістякове дерево. Списки суміжності вершин переглядаються алгоритмом за зростанням міток вершин.\par
\begin{center}
	\adjustimage{max size={0.9\linewidth}{0.9\paperheight}}
	{../10/Traversals/34.png}
\end{center}
\par Які з наведених ребер входять до складу отриманого кістякового дерева?\par
1) (1, 2)\par
2) (5, 6)\par
3) (2, 3)\par
4) (0, 2)\par
5) (3, 5)\par
6) (3, 4)\par
{\vskip 15pt} У формі вибрати номери відповідних ребер.\par
%answer:1 3 4 5 6
\end {large}}
%end
{\smallskip \begin {small} Затверджено на засіданні кафедри теоретичної кібернетики, протокол \No 12 від   19.05.2020 р.\par\smallskip\smallskip
{\parbox {6.5 cm}{Зав. кафедри \dotfill Ю.В.~Крак}}\hspace {0.7cm}{\hbox to 6.5 cm{ Екзаменатор \dotfill Т.О.~Карнаух}} \end{small}}
%begin
{{\begin {small}\par
{ \center  Київський національний університет імені Тараса Шевченка \par}
{\parbox {5 cm}{Освітня програма \hfill}{\parbox {7 cm}{Прикладна математика, Системний аналіз \hfill}}\parbox {6 cm}{\hfill Семестр 2}}\par
{\parbox {5 cm} {Навчальний предмет \hfill}\parbox {7 cm}{Програмування \hfill}\parbox {6cm}{\hfill}\par}\vspace{-7pt} \end{small}}{\center Екзаменаційний білет \No \textbf 
{ 302 }\par}}
{
\begin {large}
1. Що буде виведено за виконання?
code
try:
    res = "4.0j"!=1
    if isinstance(res, bool):
        print(f'bool;{res}')
    elif isinstance(res, int):
        print(f'int;{res}')
    elif isinstance(res, float):
        print(f'float;{res:.2f}')
    else:
        print('???')
except:
    print('ERR')
code

2. Що буде виведено за виконання?
code
a = 0
b = 2
c = 8

def h(a):
    global c
    a -= 2
    b = 3
    c = 5
    return a + b + c

a = 7
b = 3
c = 2

try:
    print(h(b), a, b, c, sep=';')
except:
    print('ERR', a, b, c, sep=';')
code

3. Що буде виведено за виконання?
code
try:
    j = 11
    while j>=5:
        j += -3
    print(j, end=';')
except:
    print('ERR')
code

4. Що буде виведено за виконання?
code
try:
    for c in range(-4, -2, 2):
        if c<=8:
            break
        print(c, end=';');c = -2
    else:
        print('end',end=';')
    print(c)
except:
    print('ERR')
code

5. Що буде виведено за виконання?\par (Дійсні значення, якщо наявні, в цьому завданні будуть виводитись у форматі з фіксованою крапкою та, можливо, з доволі великою кількістю десяткових розрядів. У відповіді для дійсних значень у ЦЬОМУ завданні достатньо вказати один десятковий розряд після крапки, відкинувши решту розрядів без округлення. Наприклад, замість 13.66666666666667 достатньо вказати 13.6, замість 25.23 достатньо вказати 25.2)
code
try:
    print(8, end=';')
    print(int('c1'), end=';')
    print(6, end=';')
except ZeroDivisionError:
    print(8, end=';')
except KeyboardInterrupt:
    print(6, end=';')
except:
    print('ERR', end=';')
else:
    print(1, end=';')
finally:
    print(2)
code

6. Що буде виведено за виконання?
code
def f(n):
    if n>2:
        f(n//2)
    else:
        print(n%2, end='')

f(16)
code

7. Що буде виведено за виконання?
code
class A:
    c = 1

    def _f1(self): print(2, end=';')
    def f2(self): print(3, end=';')

    def main(self):
        try: print(self.c, end=';')
        except: print('ERR', end=';')

        try: self._f1()
        except: print('ERR', end=';')

        try: self.f2()
        except: print('ERR', end=';')


class B(A):
    c = 4

    def _f1(self): print(5, end=';')
    def f2(self): print(6, end=';')

try: B().main()
except: print('ERR', end=';')
code

8. Що буде виведено за виконання?
code
def f(arg, n):
    n = 6
    tmp = arg
    tmp[2:-1] = [1,0]
    return len(tmp)>3

try:
    a = ['0','1','2','3','4']
    n = 9
    f(a, n)
    print(n, end=';')
    for el in a:
        print(el, end=';')
except:
    print('ERR')
code

9. Що буде виведено за виконання?
code
class A:
    a = 1
    b = 2

    def __init__(self):
        b = 3

class B(A):
    b = 4

    def __init__(self):
        super().__init__()
        self.c = 7

obj = B()
obj.a = obj.c + 10
B.a = 13

try: print(obj.a, end= ';')
except: print('ERR', end=';')

try: print(A.a, end= ';')
except: print('ERR', end=';')

try: print(B.a, end= ';')
except: print('ERR', end=';')
code

10. 
Для графа на рисунку з вершини 0 запущено алгоритм пошуку в глибину, що будує кістякове дерево. Списки суміжності вершин переглядаються алгоритмом за зростанням міток вершин.\par
\begin{center}
	\adjustimage{max size={0.9\linewidth}{0.9\paperheight}}
	{../10/Traversals/100.png}
\end{center}
\par Які з наведених ребер входять до складу отриманого кістякового дерева?\par
1) (4, 5)\par
2) (2, 5)\par
3) (0, 1)\par
4) (3, 4)\par
5) (1, 4)\par
6) (1, 2)\par
{\vskip 15pt} У формі вибрати номери відповідних ребер.\par
%answer:1 2 3 4 6
\end {large}}
%end
{\smallskip \begin {small} Затверджено на засіданні кафедри теоретичної кібернетики, протокол \No 12 від   19.05.2020 р.\par\smallskip\smallskip
{\parbox {6.5 cm}{Зав. кафедри \dotfill Ю.В.~Крак}}\hspace {0.7cm}{\hbox to 6.5 cm{ Екзаменатор \dotfill Т.О.~Карнаух}} \end{small}}
%begin
{{\begin {small}\par
{ \center  Київський національний університет імені Тараса Шевченка \par}
{\parbox {5 cm}{Освітня програма \hfill}{\parbox {7 cm}{Прикладна математика, Системний аналіз \hfill}}\parbox {6 cm}{\hfill Семестр 2}}\par
{\parbox {5 cm} {Навчальний предмет \hfill}\parbox {7 cm}{Програмування \hfill}\parbox {6cm}{\hfill}\par}\vspace{-7pt} \end{small}}{\center Екзаменаційний білет \No \textbf 
{ 303 }\par}}
{
\begin {large}
1. Що буде виведено за виконання?
code
try:
    res = 1**0.5*True
    if isinstance(res, bool):
        print(f'bool;{res}')
    elif isinstance(res, int):
        print(f'int;{res}')
    elif isinstance(res, float):
        print(f'float;{res:.2f}')
    else:
        print('???')
except:
    print('ERR')
code

2. Що буде виведено за виконання?
code
a = 5
b = 3
c = 9

def f(b):
    a = 3
    b = 5
    c = 1
    return a + b + c

a = 2
b = 8
c = 1

try:
    print(f(b), a, b, c, sep=';')
except:
    print('ERR', a, b, c, sep=';')
code

3. Що буде виведено за виконання?
code
try:
    t = 14
    while t>1:
        t += -2
    print(t, end=';')
except:
    print('ERR')
code

4. Що буде виведено за виконання?
code
try:
    for b in range(-5, -6, -1):
        if b>6:
            continue
        print(b, end=';')
        if b<4:
            break
    else:
        print(b,end=';')
    print(b)
except:
    print('ERR')
code

5. Що буде виведено за виконання?\par (Дійсні значення, якщо наявні, в цьому завданні будуть виводитись у форматі з фіксованою крапкою та, можливо, з доволі великою кількістю десяткових розрядів. У відповіді для дійсних значень у ЦЬОМУ завданні достатньо вказати один десятковий розряд після крапки, відкинувши решту розрядів без округлення. Наприклад, замість 13.66666666666667 достатньо вказати 13.6, замість 25.23 достатньо вказати 25.2)
code
try:
    print(4, end=';')
    print(5/2, end=';')
    print(7, end=';')
except Exception:
    print(3, end=';')
except Exception:
    print(0, end=';')
except:
    print('ERR', end=';')
else:
    print(9, end=';')
finally:
    print(1)
code

6. Що буде виведено за виконання?
code
def f(n):
    if n>9:
        f(n//100)
    else:
        print(n%10,end='')

f(123456)
code

7. Що буде виведено за виконання?
code
class A:
    d = 1

    def _h1(self): print(2, end=';')
    def __h2(self): print(3, end=';')

    def main(self):
        try: print(self.d, end=';')
        except: print('ERR', end=';')

        try: self._h1()
        except: print('ERR', end=';')

        try: self.__h2()
        except: print('ERR', end=';')


class B(A):
    d = 4

    def _h1(self): print(5, end=';')
    def __h2(self): print(6, end=';')

try: B().main()
except: print('ERR', end=';')
code

8. Що буде виведено за виконання?
code
def f(arg, n):
    arg[69] = 5
    n = 9
    arg = tuple(arg.values())

try:
    n = 1
    s = {15:5, 17:9, 15:5}
    f(s, n)
    print(n, end=';')
    for x in s.items():
        print(*x, sep=';', end=';')
except:
    print('ERR')
code

9. Що буде виведено за виконання?
code
class A:
    c = 1
    a = 2

    def __init__(self):
        c = 3

class B(A):
    a = 4

    def __init__(self):
        super().__init__()
        self.b = 7

obj = B()
obj.b = obj.a + 10
B.b = 13

try: print(obj.b, end= ';')
except: print('ERR', end=';')

try: print(A.b, end= ';')
except: print('ERR', end=';')

try: print(B.b, end= ';')
except: print('ERR', end=';')
code

10. 
Для графа на рисунку з вершини 5 запущено алгоритм пошуку в глибину, що будує кістякове дерево. Списки суміжності вершин переглядаються алгоритмом за зростанням міток вершин.\par
\begin{center}
	\adjustimage{max size={0.9\linewidth}{0.9\paperheight}}
	{../10/Traversals/49.png}
\end{center}
\par Які з наведених ребер входять до складу отриманого кістякового дерева?\par
1) (4, 5)\par
2) (4, 6)\par
3) (2, 6)\par
4) (0, 1)\par
5) (1, 5)\par
6) (3, 6)\par
{\vskip 15pt} У формі вибрати номери відповідних ребер.\par
%answer:2 3 4 5
\end {large}}
%end
{\smallskip \begin {small} Затверджено на засіданні кафедри теоретичної кібернетики, протокол \No 12 від   19.05.2020 р.\par\smallskip\smallskip
{\parbox {6.5 cm}{Зав. кафедри \dotfill Ю.В.~Крак}}\hspace {0.7cm}{\hbox to 6.5 cm{ Екзаменатор \dotfill Т.О.~Карнаух}} \end{small}}
%begin
{{\begin {small}\par
{ \center  Київський національний університет імені Тараса Шевченка \par}
{\parbox {5 cm}{Освітня програма \hfill}{\parbox {7 cm}{Прикладна математика, Системний аналіз \hfill}}\parbox {6 cm}{\hfill Семестр 2}}\par
{\parbox {5 cm} {Навчальний предмет \hfill}\parbox {7 cm}{Програмування \hfill}\parbox {6cm}{\hfill}\par}\vspace{-7pt} \end{small}}{\center Екзаменаційний білет \No \textbf 
{ 304 }\par}}
{
\begin {large}
1. Що буде виведено за виконання?
code
try:
    res = "False">=6.0
    if isinstance(res, bool):
        print(f'bool;{res}')
    elif isinstance(res, int):
        print(f'int;{res}')
    elif isinstance(res, float):
        print(f'float;{res:.2f}')
    else:
        print('???')
except:
    print('ERR')
code

2. Що буде виведено за виконання?
code
a = 0
b = 7
c = 4

def h(a):
    global c
    a = 4
    b = 2
    c = 4
    return a + b + c

a = 6
b = 0
c = 9

try:
    print(h(b), a, b, c, sep=';')
except:
    print('ERR', a, b, c, sep=';')
code

3. Що буде виведено за виконання?
code
try:
    j = 4
    while j<11:
        j += 1
    print(j, end=';')
except:
    print('ERR')
code

4. Що буде виведено за виконання?
code
try:
    for e in range(-6, -11, 3):
        if e<5:
            continue
        print(e, end=';')
    else:
        print(e,end=';')
    print(e)
except:
    print('ERR')
code

5. Що буде виведено за виконання?\par (Дійсні значення, якщо наявні, в цьому завданні будуть виводитись у форматі з фіксованою крапкою та, можливо, з доволі великою кількістю десяткових розрядів. У відповіді для дійсних значень у ЦЬОМУ завданні достатньо вказати один десятковий розряд після крапки, відкинувши решту розрядів без округлення. Наприклад, замість 13.66666666666667 достатньо вказати 13.6, замість 25.23 достатньо вказати 25.2)
code
try:
    print(2, end=';')
    print(int('3'), end=';')
    print(0, end=';')
except KeyboardInterrupt:
    print(7, end=';')
except TypeError:
    print(8, end=';')
except:
    print('ERR', end=';')
else:
    print(9, end=';')
finally:
    print(5)
code

6. Що буде виведено за виконання?
code
def f(s):
    for i in range(len(s)):
        s[i] += 1
    return
 
s = [1, 2, 3]
f(s)
print(s)
code

7. Що буде виведено за виконання?
code
class A:
    a = 1

    def g1(self): print(2, end=';')
    def _g2(self): print(3, end=';')

    def main(self):
        try: print(self.a, end=';')
        except: print('ERR', end=';')

        try: self.g1()
        except: print('ERR', end=';')

        try: self._g2()
        except: print('ERR', end=';')


class B(A):
    a = 4

    def g1(self): print(5, end=';')
    def _g2(self): print(6, end=';')

try: B().main()
except: print('ERR', end=';')
code

8. Що буде виведено за виконання?
code
def f(arg, n):
    arg[71] = 7
    n -= -3

try:
    n = 3
    d = {75:0, 41:3, 47:2, 47:5}
    f(d, n)
    print(n, end=';')
    for x, y in d.items():
        print(x, y, sep=';', end=';')
except:
    print('ERR')
code

9. Що буде виведено за виконання?
code
class A:
    a = 1
    c = 2

    def __init__(self):
        c = 3

class B(A):
    a = 4

    def __init__(self):
        super().__init__()
        self.b = 7

obj = B()
obj.b = obj.c + 10
B.b = 13

try: print(obj.b, end= ';')
except: print('ERR', end=';')

try: print(A.b, end= ';')
except: print('ERR', end=';')

try: print(B.b, end= ';')
except: print('ERR', end=';')
code

10. 
Для графа на рисунку з вершини 5 запущено алгоритм пошуку в глибину, що будує кістякове дерево. Списки суміжності вершин переглядаються алгоритмом за зростанням міток вершин.\par
\begin{center}
	\adjustimage{max size={0.9\linewidth}{0.9\paperheight}}
	{../10/Traversals/237.png}
\end{center}
\par Які з наведених ребер входять до складу отриманого кістякового дерева?\par
1) (3, 4)\par
2) (2, 3)\par
3) (4, 5)\par
4) (1, 3)\par
5) (4, 6)\par
6) (0, 5)\par
{\vskip 15pt} У формі вибрати номери відповідних ребер.\par
%answer:1 2 4 5 6
\end {large}}
%end
{\smallskip \begin {small} Затверджено на засіданні кафедри теоретичної кібернетики, протокол \No 12 від   19.05.2020 р.\par\smallskip\smallskip
{\parbox {6.5 cm}{Зав. кафедри \dotfill Ю.В.~Крак}}\hspace {0.7cm}{\hbox to 6.5 cm{ Екзаменатор \dotfill Т.О.~Карнаух}} \end{small}}
%begin
{{\begin {small}\par
{ \center  Київський національний університет імені Тараса Шевченка \par}
{\parbox {5 cm}{Освітня програма \hfill}{\parbox {7 cm}{Прикладна математика, Системний аналіз \hfill}}\parbox {6 cm}{\hfill Семестр 2}}\par
{\parbox {5 cm} {Навчальний предмет \hfill}\parbox {7 cm}{Програмування \hfill}\parbox {6cm}{\hfill}\par}\vspace{-7pt} \end{small}}{\center Екзаменаційний білет \No \textbf 
{ 305 }\par}}
{
\begin {large}
1. Що буде виведено за виконання?
code
try:
    res = 2.0!=6.0 and not 5>=9 or 2<=False
    if isinstance(res, bool):
        print(f'bool;{res}')
    elif isinstance(res, int):
        print(f'int;{res}')
    elif isinstance(res, float):
        print(f'float;{res:.2f}')
    else:
        print('???')
except:
    print('ERR')
code

2. Що буде виведено за виконання?
code
a = 0
b = 9
c = 7

def f(b):
    a *= 4
    b = 2
    c = 1
    return a + b + c

a = 8
b = 4
c = 6

try:
    print(f(b), a, b, c, sep=';')
except:
    print('ERR', a, b, c, sep=';')
code

3. Що буде виведено за виконання?
code
try:
    k = 1
    while k<=7:
        k += 2
    print(k, end=';')
except:
    print('ERR')
code

4. Що буде виведено за виконання?
code
try:
    for a in range(15, 12, -2):
        if a<3:
            continue
        print(a, end=';');a = -10
    else:
        print(a,end=';')
    print(a)
except:
    print('ERR')
code

5. Що буде виведено за виконання?\par (Дійсні значення, якщо наявні, в цьому завданні будуть виводитись у форматі з фіксованою крапкою та, можливо, з доволі великою кількістю десяткових розрядів. У відповіді для дійсних значень у ЦЬОМУ завданні достатньо вказати один десятковий розряд після крапки, відкинувши решту розрядів без округлення. Наприклад, замість 13.66666666666667 достатньо вказати 13.6, замість 25.23 достатньо вказати 25.2)
code
try:
    print(4, end=';')
    print(6%0.0, end=';')
    print(1, end=';')
except ValueError:
    print(4, end=';')
except ZeroDivisionError:
    print(0, end=';')
except:
    print('ERR', end=';')
else:
    print(5, end=';')
finally:
    print(6)
code

6. Що буде виведено за виконання?
code
def f(n):
    if n>9:
        f(n//100)
    else:
        print(n%10,end='')

f(987654)
code

7. Що буде виведено за виконання?
code
class A:
    c = 1

    def _h1(self): print(2, end=';')
    def _h2(self): print(3, end=';')

    def main(self):
        try: print(self.c, end=';')
        except: print('ERR', end=';')

        try: self._h1()
        except: print('ERR', end=';')

        try: self._h2()
        except: print('ERR', end=';')


class B(A):
    c = 4

    def _h1(self): print(5, end=';')
    def _h2(self): print(6, end=';')

try: B().main()
except: print('ERR', end=';')
code

8. Що буде виведено за виконання?
code
def f(arg, n):
    n = 5
    tmp = arg
    tmp[3:-3] = [4,5]
    return len(tmp)>3

try:
    c = ['0','1','2','3','4']
    n = 7
    f(c, n)
    print(n, end=';')
    for el in c:
        print(el, end=';')
except:
    print('ERR')
code

9. Що буде виведено за виконання?
code
class A:
    b = 1
    a = 2

    def __init__(self):
        b = 3

class B(A):
    a = 4

    def __init__(self):
        super().__init__()
        self.c = 7

obj = B()
obj.a = obj.b + 10
B.a = 13

try: print(obj.a, end= ';')
except: print('ERR', end=';')

try: print(A.a, end= ';')
except: print('ERR', end=';')

try: print(B.a, end= ';')
except: print('ERR', end=';')
code

10. 
Для графа на рисунку з вершини 3 запущено алгоритм пошуку в глибину, що будує кістякове дерево. Списки суміжності вершин переглядаються алгоритмом за зростанням міток вершин.\par
\begin{center}
	\adjustimage{max size={0.9\linewidth}{0.9\paperheight}}
	{../10/Traversals/124.png}
\end{center}
\par Які з наведених ребер входять до складу отриманого кістякового дерева?\par
1) (1, 5)\par
2) (0, 6)\par
3) (1, 4)\par
4) (4, 5)\par
5) (2, 3)\par
6) (1, 2)\par
{\vskip 15pt} У формі вибрати номери відповідних ребер.\par
%answer:2 3 4 5 6
\end {large}}
%end
{\smallskip \begin {small} Затверджено на засіданні кафедри теоретичної кібернетики, протокол \No 12 від   19.05.2020 р.\par\smallskip\smallskip
{\parbox {6.5 cm}{Зав. кафедри \dotfill Ю.В.~Крак}}\hspace {0.7cm}{\hbox to 6.5 cm{ Екзаменатор \dotfill Т.О.~Карнаух}} \end{small}}
%begin
{{\begin {small}\par
{ \center  Київський національний університет імені Тараса Шевченка \par}
{\parbox {5 cm}{Освітня програма \hfill}{\parbox {7 cm}{Прикладна математика, Системний аналіз \hfill}}\parbox {6 cm}{\hfill Семестр 2}}\par
{\parbox {5 cm} {Навчальний предмет \hfill}\parbox {7 cm}{Програмування \hfill}\parbox {6cm}{\hfill}\par}\vspace{-7pt} \end{small}}{\center Екзаменаційний білет \No \textbf 
{ 306 }\par}}
{
\begin {large}
1. Що буде виведено за виконання?
code
try:
    res = 7<5.0 or not 9>7 and True==7.0
    if isinstance(res, bool):
        print(f'bool;{res}')
    elif isinstance(res, int):
        print(f'int;{res}')
    elif isinstance(res, float):
        print(f'float;{res:.2f}')
    else:
        print('???')
except:
    print('ERR')
code

2. Що буде виведено за виконання?
code
a = 2
b = 4
c = 5

def g(a):
    global c
    a = 2
    b = 1
    c = 3
    return a + b + c

a = 7
b = 1
c = 8

try:
    print(g(b), a, b, c, sep=';')
except:
    print('ERR', a, b, c, sep=';')
code

3. Що буде виведено за виконання?
code
try:
    m = 8
    while m<13:
        m += 1
    print(m, end=';')
except:
    print('ERR')
code

4. Що буде виведено за виконання?
code
try:
    for d in range(-8, 2, -3):
        if d<=7:
            break
        print(d, end=';')
    else:
        print(d,end=';')
    print(d)
except:
    print('ERR')
code

5. Що буде виведено за виконання?\par (Дійсні значення, якщо наявні, в цьому завданні будуть виводитись у форматі з фіксованою крапкою та, можливо, з доволі великою кількістю десяткових розрядів. У відповіді для дійсних значень у ЦЬОМУ завданні достатньо вказати один десятковий розряд після крапки, відкинувши решту розрядів без округлення. Наприклад, замість 13.66666666666667 достатньо вказати 13.6, замість 25.23 достатньо вказати 25.2)
code
try:
    print(9, end=';')
    print(8//3, end=';')
    print(7, end=';')
except Exception:
    print(3, end=';')
except TypeError:
    print(2, end=';')
except:
    print('ERR', end=';')
else:
    print(8, end=';')
finally:
    print(7)
code

6. Що буде виведено за виконання?
code
def f(s):
    for el in s:
        el += 1
    return s
 
s = [1, 2, 3]
s = f(s)
print(s)
code

7. Що буде виведено за виконання?
code
class A:
    d = 1

    def _f1(self): print(2, end=';')
    def f2(self): print(3, end=';')

    def main(self):
        try: print(self.d, end=';')
        except: print('ERR', end=';')

        try: self._f1()
        except: print('ERR', end=';')

        try: self.f2()
        except: print('ERR', end=';')


class B(A):
    d = 4

    def _f1(self): print(5, end=';')
    def f2(self): print(6, end=';')

try: B().main()
except: print('ERR', end=';')
code

8. Що буде виведено за виконання?
code
def f(arg, n):
    arg[46] = 8
    n = 9
    arg = set(arg.keys())

try:
    n = 4
    s = {46:1, 67:0, 16:6, 46:9}
    f(s, n)
    print(n, end=';')
    for x in s.items():
        print(x, sep=';', end=';')
except:
    print('ERR')
code

9. Що буде виведено за виконання?
code
class A:
    c = 1
    b = 2

    def __init__(self):
        b = 3

class B(A):
    c = 4

    def __init__(self):
        super().__init__()
        self.a = 7

obj = B()
obj.a = obj.c + 10
B.a = 13

try: print(obj.a, end= ';')
except: print('ERR', end=';')

try: print(A.a, end= ';')
except: print('ERR', end=';')

try: print(B.a, end= ';')
except: print('ERR', end=';')
code

10. 
Для графа на рисунку з вершини 3 запущено алгоритм пошуку в глибину, що будує кістякове дерево. Списки суміжності вершин переглядаються алгоритмом за зростанням міток вершин.\par
\begin{center}
	\adjustimage{max size={0.9\linewidth}{0.9\paperheight}}
	{../10/Traversals/46.png}
\end{center}
\par Які з наведених ребер входять до складу отриманого кістякового дерева?\par
1) (1, 3)\par
2) (3, 5)\par
3) (0, 2)\par
4) (0, 1)\par
5) (4, 6)\par
6) (2, 6)\par
{\vskip 15pt} У формі вибрати номери відповідних ребер.\par
%answer:1 3 4 5 6
\end {large}}
%end
{\smallskip \begin {small} Затверджено на засіданні кафедри теоретичної кібернетики, протокол \No 12 від   19.05.2020 р.\par\smallskip\smallskip
{\parbox {6.5 cm}{Зав. кафедри \dotfill Ю.В.~Крак}}\hspace {0.7cm}{\hbox to 6.5 cm{ Екзаменатор \dotfill Т.О.~Карнаух}} \end{small}}
%begin
{{\begin {small}\par
{ \center  Київський національний університет імені Тараса Шевченка \par}
{\parbox {5 cm}{Освітня програма \hfill}{\parbox {7 cm}{Прикладна математика, Системний аналіз \hfill}}\parbox {6 cm}{\hfill Семестр 2}}\par
{\parbox {5 cm} {Навчальний предмет \hfill}\parbox {7 cm}{Програмування \hfill}\parbox {6cm}{\hfill}\par}\vspace{-7pt} \end{small}}{\center Екзаменаційний білет \No \textbf 
{ 307 }\par}}
{
\begin {large}
1. Що буде виведено за виконання?
code
try:
    res = 9**-0.5*2
    if isinstance(res, bool):
        print(f'bool;{res}')
    elif isinstance(res, int):
        print(f'int;{res}')
    elif isinstance(res, float):
        print(f'float;{res:.2f}')
    else:
        print('???')
except:
    print('ERR')
code

2. Що буде виведено за виконання?
code
a = 5
b = 1
c = 0

def g(a):
    global c
    a += 3
    b = 1
    c = 2
    return a + b + c

a = 6
b = 7
c = 3

try:
    print(g(a), a, b, c, sep=';')
except:
    print('ERR', a, b, c, sep=';')
code

3. Що буде виведено за виконання?
code
try:
    m = 6
    while m>7:
        m += -2
    print(m, end=';')
except:
    print('ERR')
code

4. Що буде виведено за виконання?
code
try:
    for b in range(14, 15, 2):
        if b>=5:
            continue
        print(b, end=';')
    else:
        print('end',end=';')
    print(b)
except:
    print('ERR')
code

5. Що буде виведено за виконання?\par (Дійсні значення, якщо наявні, в цьому завданні будуть виводитись у форматі з фіксованою крапкою та, можливо, з доволі великою кількістю десяткових розрядів. У відповіді для дійсних значень у ЦЬОМУ завданні достатньо вказати один десятковий розряд після крапки, відкинувши решту розрядів без округлення. Наприклад, замість 13.66666666666667 достатньо вказати 13.6, замість 25.23 достатньо вказати 25.2)
code
try:
    print(2, end=';')
    print(int('9'), end=';')
    print(3, end=';')
except ZeroDivisionError:
    print(1, end=';')
except KeyboardInterrupt:
    print(5, end=';')
except:
    print('ERR', end=';')
else:
    print(6, end=';')
finally:
    print(4)
code

6. Що буде виведено за виконання?
code
def f(n):
    if n>9: f(n//10)
    else:
        print(n%10, end='')
 
f(123456)
code

7. Що буде виведено за виконання?
code
class A:
    b = 1

    def _g1(self): print(2, end=';')
    def __g2(self): print(3, end=';')

    def main(self):
        try: print(self.b, end=';')
        except: print('ERR', end=';')

        try: self._g1()
        except: print('ERR', end=';')

        try: self.__g2()
        except: print('ERR', end=';')


class B(A):
    b = 4

    def _g1(self): print(5, end=';')
    def __g2(self): print(6, end=';')

try: B().main()
except: print('ERR', end=';')
code

8. Що буде виведено за виконання?
code
def f(arg, n):
    n = 2
    tmp = arg
    del tmp[:7]
    return len(tmp)>3

try:
    e = [10,11,12,13,14]
    n = 8
    f(e, n)
    print(n, end=';')
    for el in e:
        print(el, end=';')
except:
    print('ERR')
code

9. Що буде виведено за виконання?
code
class A:
    b = 1
    c = 2

    def __init__(self):
        b = 3

class B(A):
    c = 4

    def __init__(self):
        super().__init__()
        self.a = 7

obj = B()
obj.c = obj.a + 10
B.c = 13

try: print(obj.c, end= ';')
except: print('ERR', end=';')

try: print(A.c, end= ';')
except: print('ERR', end=';')

try: print(B.c, end= ';')
except: print('ERR', end=';')
code

10. 
Для графа на рисунку з вершини 4 запущено алгоритм пошуку в глибину, що будує кістякове дерево. Списки суміжності вершин переглядаються алгоритмом за зростанням міток вершин.\par
\begin{center}
	\adjustimage{max size={0.9\linewidth}{0.9\paperheight}}
	{../10/Traversals/112.png}
\end{center}
\par Які з наведених ребер входять до складу отриманого кістякового дерева?\par
1) (0, 4)\par
2) (0, 5)\par
3) (2, 3)\par
4) (0, 1)\par
5) (2, 5)\par
6) (1, 3)\par
{\vskip 15pt} У формі вибрати номери відповідних ребер.\par
%answer:1 3 4 5 6
\end {large}}
%end
{\smallskip \begin {small} Затверджено на засіданні кафедри теоретичної кібернетики, протокол \No 12 від   19.05.2020 р.\par\smallskip\smallskip
{\parbox {6.5 cm}{Зав. кафедри \dotfill Ю.В.~Крак}}\hspace {0.7cm}{\hbox to 6.5 cm{ Екзаменатор \dotfill Т.О.~Карнаух}} \end{small}}
%begin
{{\begin {small}\par
{ \center  Київський національний університет імені Тараса Шевченка \par}
{\parbox {5 cm}{Освітня програма \hfill}{\parbox {7 cm}{Прикладна математика, Системний аналіз \hfill}}\parbox {6 cm}{\hfill Семестр 2}}\par
{\parbox {5 cm} {Навчальний предмет \hfill}\parbox {7 cm}{Програмування \hfill}\parbox {6cm}{\hfill}\par}\vspace{-7pt} \end{small}}{\center Екзаменаційний білет \No \textbf 
{ 308 }\par}}
{
\begin {large}
1. Що буде виведено за виконання?
code
def f(n):
    print('f', end='')
    return n!=4

try:
    print(f(-2) or f(2))
except:
    print('ERR')
code

2. Що буде виведено за виконання?
code
a = 6
b = 3
c = 1

def h(b):
    a = 5
    b = 3
    c = 5
    return a + b + c

a = 5
b = 2
c = 9

try:
    print(h(a), a, b, c, sep=';')
except:
    print('ERR', a, b, c, sep=';')
code

3. Що буде виведено за виконання?
code
try:
    t = 2
    while t<7:
        t += 2
    print(t, end=';')
except:
    print('ERR')
code

4. Що буде виведено за виконання?
code
try:
    for e in range(1, -15, 3):
        if e>8:
            break
        print(e, end=';');e = -7
    else:
        print(e,end=';')
    print(e)
except:
    print('ERR')
code

5. Що буде виведено за виконання?\par (Дійсні значення, якщо наявні, в цьому завданні будуть виводитись у форматі з фіксованою крапкою та, можливо, з доволі великою кількістю десяткових розрядів. У відповіді для дійсних значень у ЦЬОМУ завданні достатньо вказати один десятковий розряд після крапки, відкинувши решту розрядів без округлення. Наприклад, замість 13.66666666666667 достатньо вказати 13.6, замість 25.23 достатньо вказати 25.2)
code
try:
    print(0, end=';')
    print(int('d9'), end=';')
    print(3, end=';')
except ValueError:
    print(1, end=';')
except ValueError:
    print(5, end=';')
except:
    print('ERR', end=';')
else:
    print(4, end=';')
finally:
    print(6)
code

6. Що буде виведено за виконання?
code
def f(n):
    print(n%10, end='')
    if n>9:
        f(n//100)
 
f(987654)
code

7. Що буде виведено за виконання?
code
class A:
    a = 1

    def __h1(self): print(2, end=';')
    def _h2(self): print(3, end=';')

    def main(self):
        try: print(self.a, end=';')
        except: print('ERR', end=';')

        try: self.__h1()
        except: print('ERR', end=';')

        try: self._h2()
        except: print('ERR', end=';')


class B(A):
    a = 4

    def __h1(self): print(5, end=';')
    def _h2(self): print(6, end=';')

try: B().main()
except: print('ERR', end=';')
code

8. Що буде виведено за виконання?
code
def f(arg, n):
    arg[64] = 7
    n = 6
    arg = list(arg.values())

try:
    n = 2
    d = {39:8, 82:5, 64:9, 39:2}
    f(d, n)
    print(n, end=';')
    for x in d.values():
        print(x, sep=';', end=';')
except:
    print('ERR')
code

9. Що буде виведено за виконання?
code
class A:
    a = 1
    b = 2

    def __init__(self):
        b = 3

class B(A):
    a = 4

    def __init__(self):
        super().__init__()
        self.c = 7

obj = B()
obj.b = obj.b + 10
B.b = 13

try: print(obj.b, end= ';')
except: print('ERR', end=';')

try: print(A.b, end= ';')
except: print('ERR', end=';')

try: print(B.b, end= ';')
except: print('ERR', end=';')
code

10. 
Для графа на рисунку з вершини 3 запущено алгоритм пошуку в глибину, що будує кістякове дерево. Списки суміжності вершин переглядаються алгоритмом за зростанням міток вершин.\par
\begin{center}
	\adjustimage{max size={0.9\linewidth}{0.9\paperheight}}
	{../10/Traversals/91.png}
\end{center}
\par Які з наведених ребер входять до складу отриманого кістякового дерева?\par
1) (2, 5)\par
2) (0, 3)\par
3) (4, 5)\par
4) (2, 4)\par
5) (0, 1)\par
6) (5, 6)\par
{\vskip 15pt} У формі вибрати номери відповідних ребер.\par
%answer:2 3 4 5 6
\end {large}}
%end
{\smallskip \begin {small} Затверджено на засіданні кафедри теоретичної кібернетики, протокол \No 12 від   19.05.2020 р.\par\smallskip\smallskip
{\parbox {6.5 cm}{Зав. кафедри \dotfill Ю.В.~Крак}}\hspace {0.7cm}{\hbox to 6.5 cm{ Екзаменатор \dotfill Т.О.~Карнаух}} \end{small}}
%begin
{{\begin {small}\par
{ \center  Київський національний університет імені Тараса Шевченка \par}
{\parbox {5 cm}{Освітня програма \hfill}{\parbox {7 cm}{Прикладна математика, Системний аналіз \hfill}}\parbox {6 cm}{\hfill Семестр 2}}\par
{\parbox {5 cm} {Навчальний предмет \hfill}\parbox {7 cm}{Програмування \hfill}\parbox {6cm}{\hfill}\par}\vspace{-7pt} \end{small}}{\center Екзаменаційний білет \No \textbf 
{ 309 }\par}}
{
\begin {large}
1. Що буде виведено за виконання?
code
try:
    res = 6/8//6*3
    if isinstance(res, bool):
        print(f'bool;{res}')
    elif isinstance(res, int):
        print(f'int;{res}')
    elif isinstance(res, float):
        print(f'float;{res:.2f}')
    else:
        print('???')
except:
    print('ERR')
code

2. Що буде виведено за виконання?
code
a = 7
b = 4
c = 0

def f(b):
    global c
    a = 2
    b = 4
    c = 1
    return a + b + c

a = 3
b = 6
c = 8

try:
    print(f(a), a, b, c, sep=';')
except:
    print('ERR', a, b, c, sep=';')
code

3. Що буде виведено за виконання?
code
try:
    i = 9
    while i<=6:
        i += 1
    print(i, end=';')
except:
    print('ERR')
code

4. Що буде виведено за виконання?
code
try:
    for f in range(-13, 10, -1):
        if f>4:
            continue
        print(f, end=';')
    else:
        print(f,end=';')
    print(f)
except:
    print('ERR')
code

5. Що буде виведено за виконання?\par (Дійсні значення, якщо наявні, в цьому завданні будуть виводитись у форматі з фіксованою крапкою та, можливо, з доволі великою кількістю десяткових розрядів. У відповіді для дійсних значень у ЦЬОМУ завданні достатньо вказати один десятковий розряд після крапки, відкинувши решту розрядів без округлення. Наприклад, замість 13.66666666666667 достатньо вказати 13.6, замість 25.23 достатньо вказати 25.2)
code
try:
    print(8, end=';')
    print(2j<=0j, end=';')
    print(0, end=';')
except ZeroDivisionError:
    print(7, end=';')
except Exception:
    print(8, end=';')
except:
    print('ERR', end=';')
else:
    print(5, end=';')
finally:
    print(3)
code

6. Що буде виведено за виконання?
code
def f(n):
    if n>9: f(n//10)
    else:
        print(n%10, end='')
 
f(123456)
code

7. Що буде виведено за виконання?
code
class A:
    c = 1

    def _f1(self): print(2, end=';')
    def _f2(self): print(3, end=';')

    def main(self):
        try: print(self.c, end=';')
        except: print('ERR', end=';')

        try: self._f1()
        except: print('ERR', end=';')

        try: self._f2()
        except: print('ERR', end=';')


class B(A):
    c = 4

    def _f1(self): print(5, end=';')
    def _f2(self): print(6, end=';')

try: B().main()
except: print('ERR', end=';')
code

8. Що буде виведено за виконання?
code
def f(arg, n):
    n = 3
    tmp = arg
    tmp[-4:] = 'yz'
    return len(tmp)>3

try:
    b = ['0','1','2','3','4']
    n = 9
    f(b, n)
    print(n, end=';')
    for el in b:
        print(el, end=';')
except:
    print('ERR')
code

9. Що буде виведено за виконання?
code
class A:
    c = 1
    b = 2

    def __init__(self):
        b = 3

class B(A):
    c = 4

    def __init__(self):
        super().__init__()
        self.a = 7

obj = B()
obj.a = obj.c + 10
B.a = 13

try: print(obj.a, end= ';')
except: print('ERR', end=';')

try: print(A.a, end= ';')
except: print('ERR', end=';')

try: print(B.a, end= ';')
except: print('ERR', end=';')
code

10. 
Для графа на рисунку з вершини 6 запущено алгоритм пошуку в глибину, що будує кістякове дерево. Списки суміжності вершин переглядаються алгоритмом за зростанням міток вершин.\par
\begin{center}
	\adjustimage{max size={0.9\linewidth}{0.9\paperheight}}
	{../10/Traversals/225.png}
\end{center}
\par Які з наведених ребер входять до складу отриманого кістякового дерева?\par
1) (1, 3)\par
2) (0, 2)\par
3) (1, 2)\par
4) (2, 4)\par
5) (1, 4)\par
6) (1, 5)\par
{\vskip 15pt} У формі вибрати номери відповідних ребер.\par
%answer:2 4
\end {large}}
%end
{\smallskip \begin {small} Затверджено на засіданні кафедри теоретичної кібернетики, протокол \No 12 від   19.05.2020 р.\par\smallskip\smallskip
{\parbox {6.5 cm}{Зав. кафедри \dotfill Ю.В.~Крак}}\hspace {0.7cm}{\hbox to 6.5 cm{ Екзаменатор \dotfill Т.О.~Карнаух}} \end{small}}
%begin
{{\begin {small}\par
{ \center  Київський національний університет імені Тараса Шевченка \par}
{\parbox {5 cm}{Освітня програма \hfill}{\parbox {7 cm}{Прикладна математика, Системний аналіз \hfill}}\parbox {6 cm}{\hfill Семестр 2}}\par
{\parbox {5 cm} {Навчальний предмет \hfill}\parbox {7 cm}{Програмування \hfill}\parbox {6cm}{\hfill}\par}\vspace{-7pt} \end{small}}{\center Екзаменаційний білет \No \textbf 
{ 310 }\par}}
{
\begin {large}
1. Що буде виведено за виконання?
code
try:
    res = True=="9j"
    if isinstance(res, bool):
        print(f'bool;{res}')
    elif isinstance(res, int):
        print(f'int;{res}')
    elif isinstance(res, float):
        print(f'float;{res:.2f}')
    else:
        print('???')
except:
    print('ERR')
code

2. Що буде виведено за виконання?
code
a = 8
b = 2
c = 5

def f(a):
    global c
    a *= 3
    b = 4
    c = 5
    return a + b + c

a = 4
b = 1
c = 9

try:
    print(f(a), a, b, c, sep=';')
except:
    print('ERR', a, b, c, sep=';')
code

3. Що буде виведено за виконання?
code
try:
    n = 9
    while n>=8:
        n += -3
    print(n, end=';')
except:
    print('ERR')
code

4. Що буде виведено за виконання?
code
try:
    for d in range(-12, -14, -2):
        if d>=6:
            break
        print(d, end=';');d = 1
        if d>3:
            break
    else:
        print(d,end=';')
    print(d)
except:
    print('ERR')
code

5. Що буде виведено за виконання?\par (Дійсні значення, якщо наявні, в цьому завданні будуть виводитись у форматі з фіксованою крапкою та, можливо, з доволі великою кількістю десяткових розрядів. У відповіді для дійсних значень у ЦЬОМУ завданні достатньо вказати один десятковий розряд після крапки, відкинувши решту розрядів без округлення. Наприклад, замість 13.66666666666667 достатньо вказати 13.6, замість 25.23 достатньо вказати 25.2)
code
try:
    print(9, end=';')
    print(2/0, end=';')
    print(4, end=';')
except TypeError:
    print(1, end=';')
except KeyboardInterrupt:
    print(7, end=';')
except:
    print('ERR', end=';')
else:
    print(0, end=';')
finally:
    print(6)
code

6. Що буде виведено за виконання?
code
def f(s):
    for el in s:
        el += 1
        return s
 
s = [1, 2, 3]
s = f(s)
print(s)
code

7. Що буде виведено за виконання?
code
class A:
    b = 1

    def f1(self): print(2, end=';')
    def _f2(self): print(3, end=';')

    def main(self):
        try: print(self.b, end=';')
        except: print('ERR', end=';')

        try: self.f1()
        except: print('ERR', end=';')

        try: self._f2()
        except: print('ERR', end=';')


class B(A):
    b = 4

    def f1(self): print(5, end=';')
    def _f2(self): print(6, end=';')

try: B().main()
except: print('ERR', end=';')
code

8. Що буде виведено за виконання?
code
def f(arg, n):
    arg[60] = 8
    n -= 2

try:
    n = 4
    d = {73:0, 27:8, 73:8}
    f(d, n)
    print(n, end=';')
    for x, y in d.items():
        print(x, y, sep=';', end=';')
except:
    print('ERR')
code

9. Що буде виведено за виконання?
code
class A:
    b = 1
    a = 2

    def __init__(self):
        b = 3

class B(A):
    a = 4

    def __init__(self):
        super().__init__()
        self.c = 7

obj = B()
obj.b = obj.c + 10
B.b = 13

try: print(obj.b, end= ';')
except: print('ERR', end=';')

try: print(A.b, end= ';')
except: print('ERR', end=';')

try: print(B.b, end= ';')
except: print('ERR', end=';')
code

10. 
Для графа на рисунку з вершини 3 запущено алгоритм пошуку в глибину, що будує кістякове дерево. Списки суміжності вершин переглядаються алгоритмом за зростанням міток вершин.\par
\begin{center}
	\adjustimage{max size={0.9\linewidth}{0.9\paperheight}}
	{../10/Traversals/261.png}
\end{center}
\par Які з наведених ребер входять до складу отриманого кістякового дерева?\par
1) (1, 4)\par
2) (0, 4)\par
3) (0, 5)\par
4) (2, 4)\par
5) (2, 5)\par
6) (2, 6)\par
{\vskip 15pt} У формі вибрати номери відповідних ребер.\par
%answer:1 6
\end {large}}
%end
{\smallskip \begin {small} Затверджено на засіданні кафедри теоретичної кібернетики, протокол \No 12 від   19.05.2020 р.\par\smallskip\smallskip
{\parbox {6.5 cm}{Зав. кафедри \dotfill Ю.В.~Крак}}\hspace {0.7cm}{\hbox to 6.5 cm{ Екзаменатор \dotfill Т.О.~Карнаух}} \end{small}}
%begin
{{\begin {small}\par
{ \center  Київський національний університет імені Тараса Шевченка \par}
{\parbox {5 cm}{Освітня програма \hfill}{\parbox {7 cm}{Прикладна математика, Системний аналіз \hfill}}\parbox {6 cm}{\hfill Семестр 2}}\par
{\parbox {5 cm} {Навчальний предмет \hfill}\parbox {7 cm}{Програмування \hfill}\parbox {6cm}{\hfill}\par}\vspace{-7pt} \end{small}}{\center Екзаменаційний білет \No \textbf 
{ 311 }\par}}
{
\begin {large}
1. Що буде виведено за виконання?
code
try:
    res = 1**-4+-2
    if isinstance(res, bool):
        print(f'bool;{res}')
    elif isinstance(res, int):
        print(f'int;{res}')
    elif isinstance(res, float):
        print(f'float;{res:.2f}')
    else:
        print('???')
except:
    print('ERR')
code

2. Що буде виведено за виконання?
code
a = 2
b = 8
c = 7

def g(a):
    a = 5
    b = 1
    c = 4
    return a + b + c

a = 5
b = 1
c = 4

try:
    print(g(a), a, b, c, sep=';')
except:
    print('ERR', a, b, c, sep=';')
code

3. Що буде виведено за виконання?
code
try:
    i = 10
    while i>=0:
        i += -3
    print(i, end=';')
except:
    print('ERR')
code

4. Що буде виведено за виконання?
code
try:
    for a in range(4, 12, 2):
        if a<=3:
            continue
        print(a, end=';')
    else:
        print(a,end=';')
    print(a)
except:
    print('ERR')
code

5. Що буде виведено за виконання?\par (Дійсні значення, якщо наявні, в цьому завданні будуть виводитись у форматі з фіксованою крапкою та, можливо, з доволі великою кількістю десяткових розрядів. У відповіді для дійсних значень у ЦЬОМУ завданні достатньо вказати один десятковий розряд після крапки, відкинувши решту розрядів без округлення. Наприклад, замість 13.66666666666667 достатньо вказати 13.6, замість 25.23 достатньо вказати 25.2)
code
try:
    print(5, end=';')
    print(int('a6'), end=';')
    print(4, end=';')
except Exception:
    print(1, end=';')
except KeyboardInterrupt:
    print(8, end=';')
except:
    print('ERR', end=';')
else:
    print(3, end=';')
finally:
    print(9)
code

6. Що буде виведено за виконання?
code
def f(n):
    if n>9:
        f(n//10)
    print(n%10, end='')
 
f(543216)
code

7. Що буде виведено за виконання?
code
class A:
    a = 1

    def _h1(self): print(2, end=';')
    def _h2(self): print(3, end=';')

    def main(self):
        try: print(self.a, end=';')
        except: print('ERR', end=';')

        try: self._h1()
        except: print('ERR', end=';')

        try: self._h2()
        except: print('ERR', end=';')


class B(A):
    a = 4

    def _h1(self): print(5, end=';')
    def _h2(self): print(6, end=';')

try: B().main()
except: print('ERR', end=';')
code

8. Що буде виведено за виконання?
code
def f(arg, n):
    n = 1
    tmp = arg
    del tmp[6:-6]
    return len(tmp)>3

try:
    g = [10,11,12,13,14]
    n = 4
    f(g, n)
    print(n, end=';')
    for el in g:
        print(el, end=';')
except:
    print('ERR')
code

9. Що буде виведено за виконання?
code
class A:
    a = 1
    c = 2

    def __init__(self):
        c = 3

class B(A):
    a = 4

    def __init__(self):
        super().__init__()
        self.b = 7

obj = B()
obj.a = obj.a + 10
B.a = 13

try: print(obj.a, end= ';')
except: print('ERR', end=';')

try: print(A.a, end= ';')
except: print('ERR', end=';')

try: print(B.a, end= ';')
except: print('ERR', end=';')
code

10. 
Для графа на рисунку з вершини 0 запущено алгоритм пошуку в глибину, що будує кістякове дерево. Списки суміжності вершин переглядаються алгоритмом за зростанням міток вершин.\par
\begin{center}
	\adjustimage{max size={0.9\linewidth}{0.9\paperheight}}
	{../10/Traversals/276.png}
\end{center}
\par Які з наведених ребер входять до складу отриманого кістякового дерева?\par
1) (1, 2)\par
2) (3, 5)\par
3) (0, 1)\par
4) (2, 4)\par
5) (3, 4)\par
6) (1, 4)\par
{\vskip 15pt} У формі вибрати номери відповідних ребер.\par
%answer:1 2 3 4 5
\end {large}}
%end
{\smallskip \begin {small} Затверджено на засіданні кафедри теоретичної кібернетики, протокол \No 12 від   19.05.2020 р.\par\smallskip\smallskip
{\parbox {6.5 cm}{Зав. кафедри \dotfill Ю.В.~Крак}}\hspace {0.7cm}{\hbox to 6.5 cm{ Екзаменатор \dotfill Т.О.~Карнаух}} \end{small}}
%begin
{{\begin {small}\par
{ \center  Київський національний університет імені Тараса Шевченка \par}
{\parbox {5 cm}{Освітня програма \hfill}{\parbox {7 cm}{Прикладна математика, Системний аналіз \hfill}}\parbox {6 cm}{\hfill Семестр 2}}\par
{\parbox {5 cm} {Навчальний предмет \hfill}\parbox {7 cm}{Програмування \hfill}\parbox {6cm}{\hfill}\par}\vspace{-7pt} \end{small}}{\center Екзаменаційний білет \No \textbf 
{ 312 }\par}}
{
\begin {large}
1. Що буде виведено за виконання?
code
try:
    res = 7//4%5*4
    if isinstance(res, bool):
        print(f'bool;{res}')
    elif isinstance(res, int):
        print(f'int;{res}')
    elif isinstance(res, float):
        print(f'float;{res:.2f}')
    else:
        print('???')
except:
    print('ERR')
code

2. Що буде виведено за виконання?
code
a = 6
b = 2
c = 9

def h(b):
    a -= 5
    b = 4
    c = 2
    return a + b + c

a = 8
b = 5
c = 0

try:
    print(h(a), a, b, c, sep=';')
except:
    print('ERR', a, b, c, sep=';')
code

3. Що буде виведено за виконання?
code
try:
    k = 9
    while k>4:
        k += -2
    print(k, end=';')
except:
    print('ERR')
code

4. Що буде виведено за виконання?
code
try:
    for c in range(-13, 1, -1):
        if c<7:
            continue
        print(c, end=';')
    else:
        print(c,end=';')
    print(c)
except:
    print('ERR')
code

5. Що буде виведено за виконання?\par (Дійсні значення, якщо наявні, в цьому завданні будуть виводитись у форматі з фіксованою крапкою та, можливо, з доволі великою кількістю десяткових розрядів. У відповіді для дійсних значень у ЦЬОМУ завданні достатньо вказати один десятковий розряд після крапки, відкинувши решту розрядів без округлення. Наприклад, замість 13.66666666666667 достатньо вказати 13.6, замість 25.23 достатньо вказати 25.2)
code
try:
    print(7, end=';')
    print(int('2'), end=';')
    print(0, end=';')
except TypeError:
    print(4, end=';')
except ValueError:
    print(7, end=';')
except:
    print('ERR', end=';')
else:
    print(3, end=';')
finally:
    print(1)
code

6. Що буде виведено за виконання?
code
def f(n):
    print(n%10,end='')
    if n>9:
        f(n//100)
 
f(123456)
code

7. Що буде виведено за виконання?
code
class A:
    d = 1

    def __g1(self): print(2, end=';')
    def _g2(self): print(3, end=';')

    def main(self):
        try: print(self.d, end=';')
        except: print('ERR', end=';')

        try: self.__g1()
        except: print('ERR', end=';')

        try: self._g2()
        except: print('ERR', end=';')


class B(A):
    d = 4

    def __g1(self): print(5, end=';')
    def _g2(self): print(6, end=';')

try: B().main()
except: print('ERR', end=';')
code

8. Що буде виведено за виконання?
code
def f(arg, n):
    n = 6
    tmp = arg
    tmp[-3:] = 'bd'
    return len(tmp)>3

try:
    e = [10,11,12,13,14]
    n = 0
    f(e, n)
    print(n, end=';')
    for el in e:
        print(el, end=';')
except:
    print('ERR')
code

9. Що буде виведено за виконання?
code
class A:
    c = 1
    a = 2

    def __init__(self):
        c = 3

class B(A):
    a = 4

    def __init__(self):
        super().__init__()
        self.b = 7

obj = B()
obj.b = obj.c + 10
B.b = 13

try: print(obj.b, end= ';')
except: print('ERR', end=';')

try: print(A.b, end= ';')
except: print('ERR', end=';')

try: print(B.b, end= ';')
except: print('ERR', end=';')
code

10. 
Для графа на рисунку з вершини 6 запущено алгоритм пошуку в глибину, що будує кістякове дерево. Списки суміжності вершин переглядаються алгоритмом за зростанням міток вершин.\par
\begin{center}
	\adjustimage{max size={0.9\linewidth}{0.9\paperheight}}
	{../10/Traversals/255.png}
\end{center}
\par Які з наведених ребер входять до складу отриманого кістякового дерева?\par
1) (2, 5)\par
2) (1, 3)\par
3) (5, 6)\par
4) (0, 5)\par
5) (1, 2)\par
6) (1, 5)\par
{\vskip 15pt} У формі вибрати номери відповідних ребер.\par
%answer:1 2
\end {large}}
%end
{\smallskip \begin {small} Затверджено на засіданні кафедри теоретичної кібернетики, протокол \No 12 від   19.05.2020 р.\par\smallskip\smallskip
{\parbox {6.5 cm}{Зав. кафедри \dotfill Ю.В.~Крак}}\hspace {0.7cm}{\hbox to 6.5 cm{ Екзаменатор \dotfill Т.О.~Карнаух}} \end{small}}
%begin
{{\begin {small}\par
{ \center  Київський національний університет імені Тараса Шевченка \par}
{\parbox {5 cm}{Освітня програма \hfill}{\parbox {7 cm}{Прикладна математика, Системний аналіз \hfill}}\parbox {6 cm}{\hfill Семестр 2}}\par
{\parbox {5 cm} {Навчальний предмет \hfill}\parbox {7 cm}{Програмування \hfill}\parbox {6cm}{\hfill}\par}\vspace{-7pt} \end{small}}{\center Екзаменаційний білет \No \textbf 
{ 313 }\par}}
{
\begin {large}
1. Що буде виведено за виконання?
code
try:
    res = not 8>=6 and 8.0>True and 3.0<=4
    if isinstance(res, bool):
        print(f'bool;{res}')
    elif isinstance(res, int):
        print(f'int;{res}')
    elif isinstance(res, float):
        print(f'float;{res:.2f}')
    else:
        print('???')
except:
    print('ERR')
code

2. Що буде виведено за виконання?
code
a = 9
b = 6
c = 0

def f(a):
    global c
    a = 2
    b = 3
    c = 5
    return a + b + c

a = 3
b = 3
c = 8

try:
    print(f(a), a, b, c, sep=';')
except:
    print('ERR', a, b, c, sep=';')
code

3. Що буде виведено за виконання?
code
try:
    m = 5
    while m<=11:
        m += 2
    print(m, end=';')
except:
    print('ERR')
code

4. Що буде виведено за виконання?
code
try:
    for f in range(8, 6, -3):
        if f>4:
            continue
        print(f, end=';')
    else:
        print(f,end=';')
    print(f)
except:
    print('ERR')
code

5. Що буде виведено за виконання?\par (Дійсні значення, якщо наявні, в цьому завданні будуть виводитись у форматі з фіксованою крапкою та, можливо, з доволі великою кількістю десяткових розрядів. У відповіді для дійсних значень у ЦЬОМУ завданні достатньо вказати один десятковий розряд після крапки, відкинувши решту розрядів без округлення. Наприклад, замість 13.66666666666667 достатньо вказати 13.6, замість 25.23 достатньо вказати 25.2)
code
try:
    print(8, end=';')
    print(5j==3j, end=';')
    print(9, end=';')
except ZeroDivisionError:
    print(0, end=';')
except Exception:
    print(2, end=';')
except:
    print('ERR', end=';')
else:
    print(6, end=';')
finally:
    print(6)
code

6. Що буде виведено за виконання?
code
def f(n):
    if n<=9:
        print(n%10, end='')
    else:
        f(n//10)

f(123456)
code

7. Що буде виведено за виконання?
code
class A:
    b = 1

    def _g1(self): print(2, end=';')
    def g2(self): print(3, end=';')

    def main(self):
        try: print(self.b, end=';')
        except: print('ERR', end=';')

        try: self._g1()
        except: print('ERR', end=';')

        try: self.g2()
        except: print('ERR', end=';')


class B(A):
    b = 4

    def _g1(self): print(5, end=';')
    def g2(self): print(6, end=';')

try: B().main()
except: print('ERR', end=';')
code

8. Що буде виведено за виконання?
code
def f(arg, n):
    n = 1
    tmp = arg
    del tmp[8:-8]
    return len(tmp)>3

try:
    c = ['0','1','2','3','4']
    n = 7
    f(c, n)
    print(n, end=';')
    for el in c:
        print(el, end=';')
except:
    print('ERR')
code

9. Що буде виведено за виконання?
code
class A:
    c = 1
    a = 2

    def __init__(self):
        c = 3

class B(A):
    a = 4

    def __init__(self):
        super().__init__()
        self.b = 7

obj = B()
obj.a = obj.b + 10
B.a = 13

try: print(obj.a, end= ';')
except: print('ERR', end=';')

try: print(A.a, end= ';')
except: print('ERR', end=';')

try: print(B.a, end= ';')
except: print('ERR', end=';')
code

10. 
Для графа на рисунку з вершини 6 запущено алгоритм пошуку в глибину, що будує кістякове дерево. Списки суміжності вершин переглядаються алгоритмом за зростанням міток вершин.\par
\begin{center}
	\adjustimage{max size={0.9\linewidth}{0.9\paperheight}}
	{../10/Traversals/242.png}
\end{center}
\par Які з наведених ребер входять до складу отриманого кістякового дерева?\par
1) (4, 5)\par
2) (0, 2)\par
3) (1, 5)\par
4) (3, 4)\par
5) (1, 2)\par
6) (4, 6)\par
{\vskip 15pt} У формі вибрати номери відповідних ребер.\par
%answer:2 3 5 6
\end {large}}
%end
{\smallskip \begin {small} Затверджено на засіданні кафедри теоретичної кібернетики, протокол \No 12 від   19.05.2020 р.\par\smallskip\smallskip
{\parbox {6.5 cm}{Зав. кафедри \dotfill Ю.В.~Крак}}\hspace {0.7cm}{\hbox to 6.5 cm{ Екзаменатор \dotfill Т.О.~Карнаух}} \end{small}}
%begin
{{\begin {small}\par
{ \center  Київський національний університет імені Тараса Шевченка \par}
{\parbox {5 cm}{Освітня програма \hfill}{\parbox {7 cm}{Прикладна математика, Системний аналіз \hfill}}\parbox {6 cm}{\hfill Семестр 2}}\par
{\parbox {5 cm} {Навчальний предмет \hfill}\parbox {7 cm}{Програмування \hfill}\parbox {6cm}{\hfill}\par}\vspace{-7pt} \end{small}}{\center Екзаменаційний білет \No \textbf 
{ 314 }\par}}
{
\begin {large}
1. Що буде виведено за виконання?
code
try:
    res = "True"<=3j
    if isinstance(res, bool):
        print(f'bool;{res}')
    elif isinstance(res, int):
        print(f'int;{res}')
    elif isinstance(res, float):
        print(f'float;{res:.2f}')
    else:
        print('???')
except:
    print('ERR')
code

2. Що буде виведено за виконання?
code
a = 7
b = 4
c = 1

def h(a):
    a = 3
    b += 1
    c = 2
    return a + b + c

a = 3
b = 2
c = 6

try:
    print(h(a), a, b, c, sep=';')
except:
    print('ERR', a, b, c, sep=';')
code

3. Що буде виведено за виконання?
code
try:
    i = 15
    while i>=9:
        i += -2
    print(i, end=';')
except:
    print('ERR')
code

4. Що буде виведено за виконання?
code
try:
    for a in range(3, -8, 3):
        if a<=8:
            break
        print(a, end=';');a = 0
    else:
        print(a,end=';')
    print(a)
except:
    print('ERR')
code

5. Що буде виведено за виконання?\par (Дійсні значення, якщо наявні, в цьому завданні будуть виводитись у форматі з фіксованою крапкою та, можливо, з доволі великою кількістю десяткових розрядів. У відповіді для дійсних значень у ЦЬОМУ завданні достатньо вказати один десятковий розряд після крапки, відкинувши решту розрядів без округлення. Наприклад, замість 13.66666666666667 достатньо вказати 13.6, замість 25.23 достатньо вказати 25.2)
code
try:
    print(8, end=';')
    print(1%1, end=';')
    print(7, end=';')
except TypeError:
    print(4, end=';')
except KeyboardInterrupt:
    print(3, end=';')
except:
    print('ERR', end=';')
else:
    print(9, end=';')
finally:
    print(0)
code

6. Що буде виведено за виконання?
code
def f(s):
    s1 = s
    for i in range(len(s)):
        s1[i] += 1
        return s1

s = [1, 2, 3]
f(s)
print(s)
code

7. Що буде виведено за виконання?
code
class A:
    d = 1

    def h1(self): print(2, end=';')
    def _h2(self): print(3, end=';')

    def main(self):
        try: print(self.d, end=';')
        except: print('ERR', end=';')

        try: self.h1()
        except: print('ERR', end=';')

        try: self._h2()
        except: print('ERR', end=';')


class B(A):
    d = 4

    def h1(self): print(5, end=';')
    def _h2(self): print(6, end=';')

try: B().main()
except: print('ERR', end=';')
code

8. Що буде виведено за виконання?
code
def f(arg, n):
    arg[32] = 3
    n *= 2

try:
    n = 4
    t = {21:0, 32:7, 32:6}
    f(t, n)
    print(n, end=';')
    for x in t :
        print(x, sep=';', end=';')
except:
    print('ERR')
code

9. Що буде виведено за виконання?
code
class A:
    b = 1
    a = 2

    def __init__(self):
        a = 3

class B(A):
    b = 4

    def __init__(self):
        super().__init__()
        self.c = 7

obj = B()
obj.c = obj.c + 10
B.c = 13

try: print(obj.c, end= ';')
except: print('ERR', end=';')

try: print(A.c, end= ';')
except: print('ERR', end=';')

try: print(B.c, end= ';')
except: print('ERR', end=';')
code

10. 
Для графа на рисунку з вершини 4 запущено алгоритм пошуку в глибину, що будує кістякове дерево. Списки суміжності вершин переглядаються алгоритмом за зростанням міток вершин.\par
\begin{center}
	\adjustimage{max size={0.9\linewidth}{0.9\paperheight}}
	{../10/Traversals/144.png}
\end{center}
\par Які з наведених ребер входять до складу отриманого кістякового дерева?\par
1) (3, 5)\par
2) (1, 6)\par
3) (1, 5)\par
4) (4, 6)\par
5) (0, 3)\par
6) (1, 3)\par
{\vskip 15pt} У формі вибрати номери відповідних ребер.\par
%answer:3 5 6
\end {large}}
%end
{\smallskip \begin {small} Затверджено на засіданні кафедри теоретичної кібернетики, протокол \No 12 від   19.05.2020 р.\par\smallskip\smallskip
{\parbox {6.5 cm}{Зав. кафедри \dotfill Ю.В.~Крак}}\hspace {0.7cm}{\hbox to 6.5 cm{ Екзаменатор \dotfill Т.О.~Карнаух}} \end{small}}
%begin
{{\begin {small}\par
{ \center  Київський національний університет імені Тараса Шевченка \par}
{\parbox {5 cm}{Освітня програма \hfill}{\parbox {7 cm}{Прикладна математика, Системний аналіз \hfill}}\parbox {6 cm}{\hfill Семестр 2}}\par
{\parbox {5 cm} {Навчальний предмет \hfill}\parbox {7 cm}{Програмування \hfill}\parbox {6cm}{\hfill}\par}\vspace{-7pt} \end{small}}{\center Екзаменаційний білет \No \textbf 
{ 315 }\par}}
{
\begin {large}
1. Що буде виведено за виконання?
code
try:
    res = not False!=2 or 5<9.0 or 3==4.0
    if isinstance(res, bool):
        print(f'bool;{res}')
    elif isinstance(res, int):
        print(f'int;{res}')
    elif isinstance(res, float):
        print(f'float;{res:.2f}')
    else:
        print('???')
except:
    print('ERR')
code

2. Що буде виведено за виконання?
code
a = 4
b = 7
c = 0

def g(b):
    a = 2
    b = 3
    c = 1
    return a + b + c

a = 5
b = 9
c = 6

try:
    print(g(b), a, b, c, sep=';')
except:
    print('ERR', a, b, c, sep=';')
code

3. Що буде виведено за виконання?
code
try:
    k = 13
    while k>3:
        k += -3
    print(k, end=';')
except:
    print('ERR')
code

4. Що буде виведено за виконання?
code
try:
    for b in range(6, 13, -2):
        if b>=3:
            continue
        print(b, end=';')
    else:
        print(b,end=';')
    print(b)
except:
    print('ERR')
code

5. Що буде виведено за виконання?\par (Дійсні значення, якщо наявні, в цьому завданні будуть виводитись у форматі з фіксованою крапкою та, можливо, з доволі великою кількістю десяткових розрядів. У відповіді для дійсних значень у ЦЬОМУ завданні достатньо вказати один десятковий розряд після крапки, відкинувши решту розрядів без округлення. Наприклад, замість 13.66666666666667 достатньо вказати 13.6, замість 25.23 достатньо вказати 25.2)
code
try:
    print(2, end=';')
    print(5//0.0, end=';')
    print(1, end=';')
except ValueError:
    print(4, end=';')
except ZeroDivisionError:
    print(0, end=';')
except:
    print('ERR', end=';')
else:
    print(8, end=';')
finally:
    print(3)
code

6. Що буде виведено за виконання?
code
def f(n):
    if n<=9:
        print(n%10, end='')
    else:
        f(n//10)

f(987651)
code

7. Що буде виведено за виконання?
code
class A:
    c = 1

    def _f1(self): print(2, end=';')
    def __f2(self): print(3, end=';')

    def main(self):
        try: print(self.c, end=';')
        except: print('ERR', end=';')

        try: self._f1()
        except: print('ERR', end=';')

        try: self.__f2()
        except: print('ERR', end=';')


class B(A):
    c = 4

    def _f1(self): print(5, end=';')
    def __f2(self): print(6, end=';')

try: B().main()
except: print('ERR', end=';')
code

8. Що буде виведено за виконання?
code
def f(arg, n):
    arg[18] = 4
    n += 3

try:
    n = 7
    s = {82:7, 74:5, 74:4}
    f(s, n)
    print(n, end=';')
    for x in s.keys():
        print(x, sep=';', end=';')
except:
    print('ERR')
code

9. Що буде виведено за виконання?
code
class A:
    a = 1
    c = 2

    def __init__(self):
        c = 3

class B(A):
    a = 4

    def __init__(self):
        super().__init__()
        self.b = 7

obj = B()
obj.a = obj.b + 10
B.a = 13

try: print(obj.a, end= ';')
except: print('ERR', end=';')

try: print(A.a, end= ';')
except: print('ERR', end=';')

try: print(B.a, end= ';')
except: print('ERR', end=';')
code

10. 
Для графа на рисунку з вершини 2 запущено алгоритм пошуку в глибину, що будує кістякове дерево. Списки суміжності вершин переглядаються алгоритмом за зростанням міток вершин.\par
\begin{center}
	\adjustimage{max size={0.9\linewidth}{0.9\paperheight}}
	{../10/Traversals/73.png}
\end{center}
\par Які з наведених ребер входять до складу отриманого кістякового дерева?\par
1) (0, 6)\par
2) (3, 6)\par
3) (4, 6)\par
4) (2, 4)\par
5) (1, 2)\par
6) (0, 1)\par
{\vskip 15pt} У формі вибрати номери відповідних ребер.\par
%answer:2 3 5 6
\end {large}}
%end
{\smallskip \begin {small} Затверджено на засіданні кафедри теоретичної кібернетики, протокол \No 12 від   19.05.2020 р.\par\smallskip\smallskip
{\parbox {6.5 cm}{Зав. кафедри \dotfill Ю.В.~Крак}}\hspace {0.7cm}{\hbox to 6.5 cm{ Екзаменатор \dotfill Т.О.~Карнаух}} \end{small}}
%begin
{{\begin {small}\par
{ \center  Київський національний університет імені Тараса Шевченка \par}
{\parbox {5 cm}{Освітня програма \hfill}{\parbox {7 cm}{Прикладна математика, Системний аналіз \hfill}}\parbox {6 cm}{\hfill Семестр 2}}\par
{\parbox {5 cm} {Навчальний предмет \hfill}\parbox {7 cm}{Програмування \hfill}\parbox {6cm}{\hfill}\par}\vspace{-7pt} \end{small}}{\center Екзаменаційний білет \No \textbf 
{ 316 }\par}}
{
\begin {large}
1. Що буде виведено за виконання?
code
try:
    res = 2**True+False
    if isinstance(res, bool):
        print(f'bool;{res}')
    elif isinstance(res, int):
        print(f'int;{res}')
    elif isinstance(res, float):
        print(f'float;{res:.2f}')
    else:
        print('???')
except:
    print('ERR')
code

2. Що буде виведено за виконання?
code
a = 9
b = 4
c = 3

def f(b):
    global c
    a = 4
    b -= 1
    c = 5
    return a + b + c

a = 5
b = 7
c = 8

try:
    print(f(b), a, b, c, sep=';')
except:
    print('ERR', a, b, c, sep=';')
code

3. Що буде виведено за виконання?
code
try:
    i = 3
    while i<10:
        i += 1
    print(i, end=';')
except:
    print('ERR')
code

4. Що буде виведено за виконання?
code
try:
    for d in range(-15, -1, 2):
        if d<7:
            continue
        print(d, end=';');d = 3
    else:
        print(d,end=';')
    print(d)
except:
    print('ERR')
code

5. Що буде виведено за виконання?\par (Дійсні значення, якщо наявні, в цьому завданні будуть виводитись у форматі з фіксованою крапкою та, можливо, з доволі великою кількістю десяткових розрядів. У відповіді для дійсних значень у ЦЬОМУ завданні достатньо вказати один десятковий розряд після крапки, відкинувши решту розрядів без округлення. Наприклад, замість 13.66666666666667 достатньо вказати 13.6, замість 25.23 достатньо вказати 25.2)
code
try:
    print(5, end=';')
    print(2j<1j, end=';')
    print(7, end=';')
except ValueError:
    print(6, end=';')
except TypeError:
    print(9, end=';')
except:
    print('ERR', end=';')
else:
    print(5, end=';')
finally:
    print(0)
code

6. Що буде виведено за виконання?
code
def f(s):
    for i in range(len(s)):
        s[i] += 1
        return
 
s = [1, 2, 3]
f(s)
print(s)
code

7. Що буде виведено за виконання?
code
class A:
    a = 1

    def __f1(self): print(2, end=';')
    def _f2(self): print(3, end=';')

    def main(self):
        try: print(self.a, end=';')
        except: print('ERR', end=';')

        try: self.__f1()
        except: print('ERR', end=';')

        try: self._f2()
        except: print('ERR', end=';')


class B(A):
    a = 4

    def __f1(self): print(5, end=';')
    def _f2(self): print(6, end=';')

try: B().main()
except: print('ERR', end=';')
code

8. Що буде виведено за виконання?
code
def f(arg, n):
    arg[84] = 4
    n += -3

try:
    n = 5
    s = {25:5, 10:0, 10:0}
    f(s, n)
    print(n, end=';')
    for x in s.items():
        print(*x, sep=';', end=';')
except:
    print('ERR')
code

9. Що буде виведено за виконання?
code
class A:
    c = 1
    b = 2

    def __init__(self):
        c = 3

class B(A):
    b = 4

    def __init__(self):
        super().__init__()
        self.a = 7

obj = B()
obj.a = obj.c + 10
B.a = 13

try: print(obj.a, end= ';')
except: print('ERR', end=';')

try: print(A.a, end= ';')
except: print('ERR', end=';')

try: print(B.a, end= ';')
except: print('ERR', end=';')
code

10. 
Для графа на рисунку з вершини 2 запущено алгоритм пошуку в глибину, що будує кістякове дерево. Списки суміжності вершин переглядаються алгоритмом за зростанням міток вершин.\par
\begin{center}
	\adjustimage{max size={0.9\linewidth}{0.9\paperheight}}
	{../10/Traversals/259.png}
\end{center}
\par Які з наведених ребер входять до складу отриманого кістякового дерева?\par
1) (0, 5)\par
2) (4, 5)\par
3) (2, 5)\par
4) (4, 6)\par
5) (2, 6)\par
6) (0, 4)\par
{\vskip 15pt} У формі вибрати номери відповідних ребер.\par
%answer:1 6
\end {large}}
%end
{\smallskip \begin {small} Затверджено на засіданні кафедри теоретичної кібернетики, протокол \No 12 від   19.05.2020 р.\par\smallskip\smallskip
{\parbox {6.5 cm}{Зав. кафедри \dotfill Ю.В.~Крак}}\hspace {0.7cm}{\hbox to 6.5 cm{ Екзаменатор \dotfill Т.О.~Карнаух}} \end{small}}
%begin
{{\begin {small}\par
{ \center  Київський національний університет імені Тараса Шевченка \par}
{\parbox {5 cm}{Освітня програма \hfill}{\parbox {7 cm}{Прикладна математика, Системний аналіз \hfill}}\parbox {6 cm}{\hfill Семестр 2}}\par
{\parbox {5 cm} {Навчальний предмет \hfill}\parbox {7 cm}{Програмування \hfill}\parbox {6cm}{\hfill}\par}\vspace{-7pt} \end{small}}{\center Екзаменаційний білет \No \textbf 
{ 317 }\par}}
{
\begin {large}
1. Що буде виведено за виконання?
code
try:
    res = False<"8"
    if isinstance(res, bool):
        print(f'bool;{res}')
    elif isinstance(res, int):
        print(f'int;{res}')
    elif isinstance(res, float):
        print(f'float;{res:.2f}')
    else:
        print('???')
except:
    print('ERR')
code

2. Що буде виведено за виконання?
code
a = 1
b = 0
c = 6

def g(a):
    global c
    a = 3
    b *= 4
    c = 1
    return a + b + c

a = 4
b = 3
c = 0

try:
    print(g(a), a, b, c, sep=';')
except:
    print('ERR', a, b, c, sep=';')
code

3. Що буде виведено за виконання?
code
try:
    t = 7
    while t<=12:
        t += 1
    print(t, end=';')
except:
    print('ERR')
code

4. Що буде виведено за виконання?
code
try:
    for c in range(-14, 3, -2):
        if c<=5:
            continue
        print(c, end=';')
    else:
        print('end',end=';')
    print(c)
except:
    print('ERR')
code

5. Що буде виведено за виконання?\par (Дійсні значення, якщо наявні, в цьому завданні будуть виводитись у форматі з фіксованою крапкою та, можливо, з доволі великою кількістю десяткових розрядів. У відповіді для дійсних значень у ЦЬОМУ завданні достатньо вказати один десятковий розряд після крапки, відкинувши решту розрядів без округлення. Наприклад, замість 13.66666666666667 достатньо вказати 13.6, замість 25.23 достатньо вказати 25.2)
code
try:
    print(3, end=';')
    print(8//3, end=';')
    print(1, end=';')
except ZeroDivisionError:
    print(9, end=';')
except Exception:
    print(4, end=';')
except:
    print('ERR', end=';')
else:
    print(2, end=';')
finally:
    print(7)
code

6. Що буде виведено за виконання?
code
def f(n):
    if n>9:
        return f(n//100)+1
    else:
        return 1

print(f(123456))
code

7. Що буде виведено за виконання?
code
class A:
    d = 1

    def _g1(self): print(2, end=';')
    def __g2(self): print(3, end=';')

    def main(self):
        try: print(self.d, end=';')
        except: print('ERR', end=';')

        try: self._g1()
        except: print('ERR', end=';')

        try: self.__g2()
        except: print('ERR', end=';')


class B(A):
    d = 4

    def _g1(self): print(5, end=';')
    def __g2(self): print(6, end=';')

try: B().main()
except: print('ERR', end=';')
code

8. Що буде виведено за виконання?
code
def f(arg, n):
    arg[34] = 7
    n = 7
    arg = tuple(arg.items())

try:
    n = 1
    t = {17:1, 46:2, 17:4}
    f(t, n)
    print(n, end=';')
    for x, y in t.items():
        print(x, y, sep=';', end=';')
except:
    print('ERR')
code

9. Що буде виведено за виконання?
code
class A:
    b = 1
    c = 2

    def __init__(self):
        c = 3

class B(A):
    c = 4

    def __init__(self):
        super().__init__()
        self.a = 7

obj = B()
obj.b = obj.b + 10
B.b = 13

try: print(obj.b, end= ';')
except: print('ERR', end=';')

try: print(A.b, end= ';')
except: print('ERR', end=';')

try: print(B.b, end= ';')
except: print('ERR', end=';')
code

10. 
Для графа на рисунку з вершини 5 запущено алгоритм пошуку в глибину, що будує кістякове дерево. Списки суміжності вершин переглядаються алгоритмом за зростанням міток вершин.\par
\begin{center}
	\adjustimage{max size={0.9\linewidth}{0.9\paperheight}}
	{../10/Traversals/180.png}
\end{center}
\par Які з наведених ребер входять до складу отриманого кістякового дерева?\par
1) (4, 6)\par
2) (2, 4)\par
3) (1, 5)\par
4) (0, 2)\par
5) (3, 5)\par
6) (2, 6)\par
{\vskip 15pt} У формі вибрати номери відповідних ребер.\par
%answer:1 2 3 4 5
\end {large}}
%end
{\smallskip \begin {small} Затверджено на засіданні кафедри теоретичної кібернетики, протокол \No 12 від   19.05.2020 р.\par\smallskip\smallskip
{\parbox {6.5 cm}{Зав. кафедри \dotfill Ю.В.~Крак}}\hspace {0.7cm}{\hbox to 6.5 cm{ Екзаменатор \dotfill Т.О.~Карнаух}} \end{small}}
%begin
{{\begin {small}\par
{ \center  Київський національний університет імені Тараса Шевченка \par}
{\parbox {5 cm}{Освітня програма \hfill}{\parbox {7 cm}{Прикладна математика, Системний аналіз \hfill}}\parbox {6 cm}{\hfill Семестр 2}}\par
{\parbox {5 cm} {Навчальний предмет \hfill}\parbox {7 cm}{Програмування \hfill}\parbox {6cm}{\hfill}\par}\vspace{-7pt} \end{small}}{\center Екзаменаційний білет \No \textbf 
{ 318 }\par}}
{
\begin {large}
1. Що буде виведено за виконання?
code
try:
    res = not 9!=8 and 4<3.0 and 4.0==False
    if isinstance(res, bool):
        print(f'bool;{res}')
    elif isinstance(res, int):
        print(f'int;{res}')
    elif isinstance(res, float):
        print(f'float;{res:.2f}')
    else:
        print('???')
except:
    print('ERR')
code

2. Що буде виведено за виконання?
code
a = 1
b = 2
c = 9

def h(b):
    global c
    a = 4
    b = 1
    c = 2
    return a + b + c

a = 7
b = 6
c = 1

try:
    print(h(b), a, b, c, sep=';')
except:
    print('ERR', a, b, c, sep=';')
code

3. Що буде виведено за виконання?
code
try:
    m = 12
    while m>2:
        m += -3
    print(m, end=';')
except:
    print('ERR')
code

4. Що буде виведено за виконання?
code
try:
    for e in range(7, 11, -3):
        if e>=6:
            break
        print(e, end=';')
        if e<=6:
            break
    else:
        print(e,end=';')
    print(e)
except:
    print('ERR')
code

5. Що буде виведено за виконання?\par (Дійсні значення, якщо наявні, в цьому завданні будуть виводитись у форматі з фіксованою крапкою та, можливо, з доволі великою кількістю десяткових розрядів. У відповіді для дійсних значень у ЦЬОМУ завданні достатньо вказати один десятковий розряд після крапки, відкинувши решту розрядів без округлення. Наприклад, замість 13.66666666666667 достатньо вказати 13.6, замість 25.23 достатньо вказати 25.2)
code
try:
    print(6, end=';')
    print(1%0, end=';')
    print(2, end=';')
except KeyboardInterrupt:
    print(6, end=';')
except TypeError:
    print(9, end=';')
except:
    print('ERR', end=';')
else:
    print(4, end=';')
finally:
    print(3)
code

6. Що буде виведено за виконання?
code
def f(n):
    print(n%10, end='')
    if n>9:
        f(n//100)
 
f(987654)
code

7. Що буде виведено за виконання?
code
class A:
    b = 1

    def h1(self): print(2, end=';')
    def _h2(self): print(3, end=';')

    def main(self):
        try: print(self.b, end=';')
        except: print('ERR', end=';')

        try: self.h1()
        except: print('ERR', end=';')

        try: self._h2()
        except: print('ERR', end=';')


class B(A):
    b = 4

    def h1(self): print(5, end=';')
    def _h2(self): print(6, end=';')

try: B().main()
except: print('ERR', end=';')
code

8. Що буде виведено за виконання?
code
def f(arg, n):
    n = 2
    tmp = arg
    tmp[:3] = 'xy'
    return len(tmp)>3

try:
    d = ['0','1','2','3','4']
    n = 5
    f(d, n)
    print(n, end=';')
    for el in d:
        print(el, end=';')
except:
    print('ERR')
code

9. Що буде виведено за виконання?
code
class A:
    a = 1
    b = 2

    def __init__(self):
        a = 3

class B(A):
    a = 4

    def __init__(self):
        super().__init__()
        self.c = 7

obj = B()
obj.a = obj.c + 10
B.a = 13

try: print(obj.a, end= ';')
except: print('ERR', end=';')

try: print(A.a, end= ';')
except: print('ERR', end=';')

try: print(B.a, end= ';')
except: print('ERR', end=';')
code

10. 
Для графа на рисунку з вершини 0 запущено алгоритм пошуку в глибину, що будує кістякове дерево. Списки суміжності вершин переглядаються алгоритмом за зростанням міток вершин.\par
\begin{center}
	\adjustimage{max size={0.9\linewidth}{0.9\paperheight}}
	{../10/Traversals/168.png}
\end{center}
\par Які з наведених ребер входять до складу отриманого кістякового дерева?\par
1) (0, 3)\par
2) (2, 3)\par
3) (1, 2)\par
4) (0, 1)\par
5) (1, 5)\par
6) (0, 2)\par
{\vskip 15pt} У формі вибрати номери відповідних ребер.\par
%answer:2 3 4
\end {large}}
%end
{\smallskip \begin {small} Затверджено на засіданні кафедри теоретичної кібернетики, протокол \No 12 від   19.05.2020 р.\par\smallskip\smallskip
{\parbox {6.5 cm}{Зав. кафедри \dotfill Ю.В.~Крак}}\hspace {0.7cm}{\hbox to 6.5 cm{ Екзаменатор \dotfill Т.О.~Карнаух}} \end{small}}
%begin
{{\begin {small}\par
{ \center  Київський національний університет імені Тараса Шевченка \par}
{\parbox {5 cm}{Освітня програма \hfill}{\parbox {7 cm}{Прикладна математика, Системний аналіз \hfill}}\parbox {6 cm}{\hfill Семестр 2}}\par
{\parbox {5 cm} {Навчальний предмет \hfill}\parbox {7 cm}{Програмування \hfill}\parbox {6cm}{\hfill}\par}\vspace{-7pt} \end{small}}{\center Екзаменаційний білет \No \textbf 
{ 319 }\par}}
{
\begin {large}
1. Що буде виведено за виконання?
code
def f(n):
    print('f', end='')
    return n

try:
    print(f(8) and f(9))
except:
    print('ERR')
code

2. Що буде виведено за виконання?
code
a = 2
b = 8
c = 0

def h(a):
    a = 3
    b = 5
    c = 4
    return a + b + c

a = 5
b = 4
c = 3

try:
    print(h(a), a, b, c, sep=';')
except:
    print('ERR', a, b, c, sep=';')
code

3. Що буде виведено за виконання?
code
try:
    n = 8
    while n>=6:
        n += -2
    print(n, end=';')
except:
    print('ERR')
code

4. Що буде виведено за виконання?
code
try:
    for e in range(-12, -11, 3):
        if e<6:
            continue
        print(e, end=';');e = 8
    else:
        print('end',end=';')
    print(e)
except:
    print('ERR')
code

5. Що буде виведено за виконання?\par (Дійсні значення, якщо наявні, в цьому завданні будуть виводитись у форматі з фіксованою крапкою та, можливо, з доволі великою кількістю десяткових розрядів. У відповіді для дійсних значень у ЦЬОМУ завданні достатньо вказати один десятковий розряд після крапки, відкинувши решту розрядів без округлення. Наприклад, замість 13.66666666666667 достатньо вказати 13.6, замість 25.23 достатньо вказати 25.2)
code
try:
    print(0, end=';')
    print(int('b5'), end=';')
    print(8, end=';')
except KeyboardInterrupt:
    print(7, end=';')
except Exception:
    print(6, end=';')
except:
    print('ERR', end=';')
else:
    print(8, end=';')
finally:
    print(9)
code

6. Що буде виведено за виконання?
code
def f(n):
    if n>9:
        f(n//100)
    else:
        print(n%10,end='')

f(987654)
code

7. Що буде виведено за виконання?
code
class A:
    a = 1

    def _g1(self): print(2, end=';')
    def _g2(self): print(3, end=';')

    def main(self):
        try: print(self.a, end=';')
        except: print('ERR', end=';')

        try: self._g1()
        except: print('ERR', end=';')

        try: self._g2()
        except: print('ERR', end=';')


class B(A):
    a = 4

    def _g1(self): print(5, end=';')
    def _g2(self): print(6, end=';')

try: B().main()
except: print('ERR', end=';')
code

8. Що буде виведено за виконання?
code
def f(arg, n):
    arg[46] = 6
    n = 8
    arg = set(arg.keys())

try:
    n = 3
    s = {65:3, 48:0, 76:8, 68:0}
    f(s, n)
    print(n, end=';')
    for x in s.values():
        print(x, sep=';', end=';')
except:
    print('ERR')
code

9. Що буде виведено за виконання?
code
class A:
    b = 1
    c = 2

    def __init__(self):
        c = 3

class B(A):
    b = 4

    def __init__(self):
        super().__init__()
        self.a = 7

obj = B()
obj.c = obj.a + 10
B.c = 13

try: print(obj.c, end= ';')
except: print('ERR', end=';')

try: print(A.c, end= ';')
except: print('ERR', end=';')

try: print(B.c, end= ';')
except: print('ERR', end=';')
code

10. 
Для графа на рисунку з вершини 1 запущено алгоритм пошуку в глибину, що будує кістякове дерево. Списки суміжності вершин переглядаються алгоритмом за зростанням міток вершин.\par
\begin{center}
	\adjustimage{max size={0.9\linewidth}{0.9\paperheight}}
	{../10/Traversals/76.png}
\end{center}
\par Які з наведених ребер входять до складу отриманого кістякового дерева?\par
1) (3, 5)\par
2) (2, 4)\par
3) (3, 6)\par
4) (0, 3)\par
5) (2, 3)\par
6) (1, 4)\par
{\vskip 15pt} У формі вибрати номери відповідних ребер.\par
%answer:1 3 4 5 6
\end {large}}
%end
{\smallskip \begin {small} Затверджено на засіданні кафедри теоретичної кібернетики, протокол \No 12 від   19.05.2020 р.\par\smallskip\smallskip
{\parbox {6.5 cm}{Зав. кафедри \dotfill Ю.В.~Крак}}\hspace {0.7cm}{\hbox to 6.5 cm{ Екзаменатор \dotfill Т.О.~Карнаух}} \end{small}}
%begin
{{\begin {small}\par
{ \center  Київський національний університет імені Тараса Шевченка \par}
{\parbox {5 cm}{Освітня програма \hfill}{\parbox {7 cm}{Прикладна математика, Системний аналіз \hfill}}\parbox {6 cm}{\hfill Семестр 2}}\par
{\parbox {5 cm} {Навчальний предмет \hfill}\parbox {7 cm}{Програмування \hfill}\parbox {6cm}{\hfill}\par}\vspace{-7pt} \end{small}}{\center Екзаменаційний білет \No \textbf 
{ 320 }\par}}
{
\begin {large}
1. Що буде виведено за виконання?
code
try:
    res = -3**True--1
    if isinstance(res, bool):
        print(f'bool;{res}')
    elif isinstance(res, int):
        print(f'int;{res}')
    elif isinstance(res, float):
        print(f'float;{res:.2f}')
    else:
        print('???')
except:
    print('ERR')
code

2. Що буде виведено за виконання?
code
a = 1
b = 9
c = 8

def f(b):
    a = 2
    b -= 5
    c = 3
    return a + b + c

a = 7
b = 5
c = 2

try:
    print(f(b), a, b, c, sep=';')
except:
    print('ERR', a, b, c, sep=';')
code

3. Що буде виведено за виконання?
code
try:
    j = 6
    while j>=6:
        j += -3
    print(j, end=';')
except:
    print('ERR')
code

4. Що буде виведено за виконання?
code
try:
    for a in range(11, -12, -1):
        if a>3:
            break
        print(a, end=';')
    else:
        print(a,end=';')
    print(a)
except:
    print('ERR')
code

5. Що буде виведено за виконання?\par (Дійсні значення, якщо наявні, в цьому завданні будуть виводитись у форматі з фіксованою крапкою та, можливо, з доволі великою кількістю десяткових розрядів. У відповіді для дійсних значень у ЦЬОМУ завданні достатньо вказати один десятковий розряд після крапки, відкинувши решту розрядів без округлення. Наприклад, замість 13.66666666666667 достатньо вказати 13.6, замість 25.23 достатньо вказати 25.2)
code
try:
    print(4, end=';')
    print(int('5'), end=';')
    print(2, end=';')
except ZeroDivisionError:
    print(7, end=';')
except ValueError:
    print(3, end=';')
except:
    print('ERR', end=';')
else:
    print(1, end=';')
finally:
    print(0)
code

6. Що буде виведено за виконання?
code
def f(n):
    if n>9:
        f(n//100)
        print(n%10, end='')
 
f(123456)
code

7. Що буде виведено за виконання?
code
class A:
    c = 1

    def _h1(self): print(2, end=';')
    def h2(self): print(3, end=';')

    def main(self):
        try: print(self.c, end=';')
        except: print('ERR', end=';')

        try: self._h1()
        except: print('ERR', end=';')

        try: self.h2()
        except: print('ERR', end=';')


class B(A):
    c = 4

    def _h1(self): print(5, end=';')
    def h2(self): print(6, end=';')

try: B().main()
except: print('ERR', end=';')
code

8. Що буде виведено за виконання?
code
def f(arg, n):
    n = 0
    tmp = arg
    tmp[-2:-5] = ()
    return len(tmp)>3

try:
    h = [10,11,12,13,14]
    n = 6
    f(h, n)
    print(n, end=';')
    for el in h:
        print(el, end=';')
except:
    print('ERR')
code

9. Що буде виведено за виконання?
code
class A:
    a = 1
    c = 2

    def __init__(self):
        a = 3

class B(A):
    c = 4

    def __init__(self):
        super().__init__()
        self.b = 7

obj = B()
obj.b = obj.a + 10
B.b = 13

try: print(obj.b, end= ';')
except: print('ERR', end=';')

try: print(A.b, end= ';')
except: print('ERR', end=';')

try: print(B.b, end= ';')
except: print('ERR', end=';')
code

10. 
Для графа на рисунку з вершини 1 запущено алгоритм пошуку в глибину, що будує кістякове дерево. Списки суміжності вершин переглядаються алгоритмом за зростанням міток вершин.\par
\begin{center}
	\adjustimage{max size={0.9\linewidth}{0.9\paperheight}}
	{../10/Traversals/52.png}
\end{center}
\par Які з наведених ребер входять до складу отриманого кістякового дерева?\par
1) (0, 4)\par
2) (0, 5)\par
3) (0, 2)\par
4) (1, 6)\par
5) (0, 6)\par
6) (0, 1)\par
{\vskip 15pt} У формі вибрати номери відповідних ребер.\par
%answer:1 3 5 6
\end {large}}
%end
{\smallskip \begin {small} Затверджено на засіданні кафедри теоретичної кібернетики, протокол \No 12 від   19.05.2020 р.\par\smallskip\smallskip
{\parbox {6.5 cm}{Зав. кафедри \dotfill Ю.В.~Крак}}\hspace {0.7cm}{\hbox to 6.5 cm{ Екзаменатор \dotfill Т.О.~Карнаух}} \end{small}}
%begin
{{\begin {small}\par
{ \center  Київський національний університет імені Тараса Шевченка \par}
{\parbox {5 cm}{Освітня програма \hfill}{\parbox {7 cm}{Прикладна математика, Системний аналіз \hfill}}\parbox {6 cm}{\hfill Семестр 2}}\par
{\parbox {5 cm} {Навчальний предмет \hfill}\parbox {7 cm}{Програмування \hfill}\parbox {6cm}{\hfill}\par}\vspace{-7pt} \end{small}}{\center Екзаменаційний білет \No \textbf 
{ 321 }\par}}
{
\begin {large}
1. Що буде виведено за виконання?
code
def f(n):
    print('f', end='')
    return n

try:
    print(f(3) or f(-5))
except:
    print('ERR')
code

2. Що буде виведено за виконання?
code
a = 5
b = 1
c = 7

def g(a):
    a = 4
    b = 5
    c = 3
    return a + b + c

a = 4
b = 3
c = 8

try:
    print(g(b), a, b, c, sep=';')
except:
    print('ERR', a, b, c, sep=';')
code

3. Що буде виведено за виконання?
code
try:
    t = 14
    while t>5:
        t += -2
    print(t, end=';')
except:
    print('ERR')
code

4. Що буде виведено за виконання?
code
try:
    for f in range(-3, -4, -1):
        if f>5:
            continue
        print(f, end=';')
    else:
        print(f,end=';')
    print(f)
except:
    print('ERR')
code

5. Що буде виведено за виконання?\par (Дійсні значення, якщо наявні, в цьому завданні будуть виводитись у форматі з фіксованою крапкою та, можливо, з доволі великою кількістю десяткових розрядів. У відповіді для дійсних значень у ЦЬОМУ завданні достатньо вказати один десятковий розряд після крапки, відкинувши решту розрядів без округлення. Наприклад, замість 13.66666666666667 достатньо вказати 13.6, замість 25.23 достатньо вказати 25.2)
code
try:
    print(9, end=';')
    print(0/1, end=';')
    print(1, end=';')
except TypeError:
    print(2, end=';')
except Exception:
    print(8, end=';')
except:
    print('ERR', end=';')
else:
    print(6, end=';')
finally:
    print(7)
code

6. Що буде виведено за виконання?
code
def f(s):
    for i in range(len(s)):
        s[i] += 1
    return
 
s = [1, 2, 3]
f(s)
print(s)
code

7. Що буде виведено за виконання?
code
class A:
    d = 1

    def _f1(self): print(2, end=';')
    def f2(self): print(3, end=';')

    def main(self):
        try: print(self.d, end=';')
        except: print('ERR', end=';')

        try: self._f1()
        except: print('ERR', end=';')

        try: self.f2()
        except: print('ERR', end=';')


class B(A):
    d = 4

    def _f1(self): print(5, end=';')
    def f2(self): print(6, end=';')

try: B().main()
except: print('ERR', end=';')
code

8. Що буде виведено за виконання?
code
def f(arg, n):
    arg[72] = 1
    n -= -2

try:
    n = 6
    d = {14:3, 61:8, 28:9, 79:6}
    f(d, n)
    print(n, end=';')
    for x in d.keys():
        print(x, sep=';', end=';')
except:
    print('ERR')
code

9. Що буде виведено за виконання?
code
class A:
    c = 1
    a = 2

    def __init__(self):
        c = 3

class B(A):
    a = 4

    def __init__(self):
        super().__init__()
        self.b = 7

obj = B()
obj.b = obj.c + 10
B.b = 13

try: print(obj.b, end= ';')
except: print('ERR', end=';')

try: print(A.b, end= ';')
except: print('ERR', end=';')

try: print(B.b, end= ';')
except: print('ERR', end=';')
code

10. 
Для графа на рисунку з вершини 2 запущено алгоритм пошуку в глибину, що будує кістякове дерево. Списки суміжності вершин переглядаються алгоритмом за зростанням міток вершин.\par
\begin{center}
	\adjustimage{max size={0.9\linewidth}{0.9\paperheight}}
	{../10/Traversals/47.png}
\end{center}
\par Які з наведених ребер входять до складу отриманого кістякового дерева?\par
1) (1, 4)\par
2) (0, 2)\par
3) (1, 3)\par
4) (0, 5)\par
5) (3, 4)\par
6) (0, 1)\par
{\vskip 15pt} У формі вибрати номери відповідних ребер.\par
%answer:2 3 5 6
\end {large}}
%end
{\smallskip \begin {small} Затверджено на засіданні кафедри теоретичної кібернетики, протокол \No 12 від   19.05.2020 р.\par\smallskip\smallskip
{\parbox {6.5 cm}{Зав. кафедри \dotfill Ю.В.~Крак}}\hspace {0.7cm}{\hbox to 6.5 cm{ Екзаменатор \dotfill Т.О.~Карнаух}} \end{small}}
%begin
{{\begin {small}\par
{ \center  Київський національний університет імені Тараса Шевченка \par}
{\parbox {5 cm}{Освітня програма \hfill}{\parbox {7 cm}{Прикладна математика, Системний аналіз \hfill}}\parbox {6 cm}{\hfill Семестр 2}}\par
{\parbox {5 cm} {Навчальний предмет \hfill}\parbox {7 cm}{Програмування \hfill}\parbox {6cm}{\hfill}\par}\vspace{-7pt} \end{small}}{\center Екзаменаційний білет \No \textbf 
{ 322 }\par}}
{
\begin {large}
1. Що буде виведено за виконання?
code
try:
    res = "8.0">7.0j
    if isinstance(res, bool):
        print(f'bool;{res}')
    elif isinstance(res, int):
        print(f'int;{res}')
    elif isinstance(res, float):
        print(f'float;{res:.2f}')
    else:
        print('???')
except:
    print('ERR')
code

2. Що буде виведено за виконання?
code
a = 4
b = 0
c = 7

def h(a):
    global c
    a *= 2
    b = 1
    c = 5
    return a + b + c

a = 6
b = 5
c = 9

try:
    print(h(a), a, b, c, sep=';')
except:
    print('ERR', a, b, c, sep=';')
code

3. Що буде виведено за виконання?
code
try:
    n = 6
    while n<9:
        n += 2
    print(n, end=';')
except:
    print('ERR')
code

4. Що буде виведено за виконання?
code
try:
    for d in range(2, -13, -2):
        if d<=4:
            continue
        print(d, end=';')
    else:
        print(d,end=';')
    print(d)
except:
    print('ERR')
code

5. Що буде виведено за виконання?\par (Дійсні значення, якщо наявні, в цьому завданні будуть виводитись у форматі з фіксованою крапкою та, можливо, з доволі великою кількістю десяткових розрядів. У відповіді для дійсних значень у ЦЬОМУ завданні достатньо вказати один десятковий розряд після крапки, відкинувши решту розрядів без округлення. Наприклад, замість 13.66666666666667 достатньо вказати 13.6, замість 25.23 достатньо вказати 25.2)
code
try:
    print(3, end=';')
    print(4//0.0, end=';')
    print(5, end=';')
except KeyboardInterrupt:
    print(4, end=';')
except ValueError:
    print(3, end=';')
except:
    print('ERR', end=';')
else:
    print(2, end=';')
finally:
    print(7)
code

6. Що буде виведено за виконання?
code
def f(n):
    print(n%10, end='')
    if n>9:
        f(n//10)
 
f(543216)
code

7. Що буде виведено за виконання?
code
class A:
    c = 1

    def __g1(self): print(2, end=';')
    def _g2(self): print(3, end=';')

    def main(self):
        try: print(self.c, end=';')
        except: print('ERR', end=';')

        try: self.__g1()
        except: print('ERR', end=';')

        try: self._g2()
        except: print('ERR', end=';')


class B(A):
    c = 4

    def __g1(self): print(5, end=';')
    def _g2(self): print(6, end=';')

try: B().main()
except: print('ERR', end=';')
code

8. Що буде виведено за виконання?
code
def f(arg, n):
    arg[45] = 4
    n *= 2

try:
    n = 3
    t = {29:4, 45:1, 45:3}
    f(t, n)
    print(n, end=';')
    for x in t.values():
        print(x, sep=';', end=';')
except:
    print('ERR')
code

9. Що буде виведено за виконання?
code
class A:
    b = 1
    a = 2

    def __init__(self):
        a = 3

class B(A):
    b = 4

    def __init__(self):
        super().__init__()
        self.c = 7

obj = B()
obj.c = obj.b + 10
B.c = 13

try: print(obj.c, end= ';')
except: print('ERR', end=';')

try: print(A.c, end= ';')
except: print('ERR', end=';')

try: print(B.c, end= ';')
except: print('ERR', end=';')
code

10. 
Для графа на рисунку з вершини 0 запущено алгоритм пошуку в глибину, що будує кістякове дерево. Списки суміжності вершин переглядаються алгоритмом за зростанням міток вершин.\par
\begin{center}
	\adjustimage{max size={0.9\linewidth}{0.9\paperheight}}
	{../10/Traversals/204.png}
\end{center}
\par Які з наведених ребер входять до складу отриманого кістякового дерева?\par
1) (1, 3)\par
2) (5, 6)\par
3) (0, 2)\par
4) (3, 6)\par
5) (3, 4)\par
6) (2, 4)\par
{\vskip 15pt} У формі вибрати номери відповідних ребер.\par
%answer:1 2 3 5 6
\end {large}}
%end
{\smallskip \begin {small} Затверджено на засіданні кафедри теоретичної кібернетики, протокол \No 12 від   19.05.2020 р.\par\smallskip\smallskip
{\parbox {6.5 cm}{Зав. кафедри \dotfill Ю.В.~Крак}}\hspace {0.7cm}{\hbox to 6.5 cm{ Екзаменатор \dotfill Т.О.~Карнаух}} \end{small}}
%begin
{{\begin {small}\par
{ \center  Київський національний університет імені Тараса Шевченка \par}
{\parbox {5 cm}{Освітня програма \hfill}{\parbox {7 cm}{Прикладна математика, Системний аналіз \hfill}}\parbox {6 cm}{\hfill Семестр 2}}\par
{\parbox {5 cm} {Навчальний предмет \hfill}\parbox {7 cm}{Програмування \hfill}\parbox {6cm}{\hfill}\par}\vspace{-7pt} \end{small}}{\center Екзаменаційний білет \No \textbf 
{ 323 }\par}}
{
\begin {large}
1. Що буде виведено за виконання?
code
def f(n):
    print('f', end='')
    return n>-2

try:
    print(f(-1) and f(7))
except:
    print('ERR')
code

2. Що буде виведено за виконання?
code
a = 6
b = 9
c = 2

def f(b):
    global c
    a = 1
    b = 2
    c = 1
    return a + b + c

a = 0
b = 5
c = 3

try:
    print(f(a), a, b, c, sep=';')
except:
    print('ERR', a, b, c, sep=';')
code

3. Що буде виведено за виконання?
code
try:
    j = 3
    while j<8:
        j += 2
    print(j, end=';')
except:
    print('ERR')
code

4. Що буде виведено за виконання?
code
try:
    for b in range(-8, 8, 3):
        if b>=7:
            continue
        print(b, end=';')
    else:
        print(b,end=';')
    print(b)
except:
    print('ERR')
code

5. Що буде виведено за виконання?\par (Дійсні значення, якщо наявні, в цьому завданні будуть виводитись у форматі з фіксованою крапкою та, можливо, з доволі великою кількістю десяткових розрядів. У відповіді для дійсних значень у ЦЬОМУ завданні достатньо вказати один десятковий розряд після крапки, відкинувши решту розрядів без округлення. Наприклад, замість 13.66666666666667 достатньо вказати 13.6, замість 25.23 достатньо вказати 25.2)
code
try:
    print(0, end=';')
    print(int('b8'), end=';')
    print(1, end=';')
except ZeroDivisionError:
    print(5, end=';')
except ValueError:
    print(6, end=';')
except:
    print('ERR', end=';')
else:
    print(9, end=';')
finally:
    print(7)
code

6. Що буде виведено за виконання?
code
def f(n):
    if n>2:
        f(n//2)
    print(n%2, end='')
 
f(16)
code

7. Що буде виведено за виконання?
code
class A:
    a = 1

    def _h1(self): print(2, end=';')
    def __h2(self): print(3, end=';')

    def main(self):
        try: print(self.a, end=';')
        except: print('ERR', end=';')

        try: self._h1()
        except: print('ERR', end=';')

        try: self.__h2()
        except: print('ERR', end=';')


class B(A):
    a = 4

    def _h1(self): print(5, end=';')
    def __h2(self): print(6, end=';')

try: B().main()
except: print('ERR', end=';')
code

8. Що буде виведено за виконання?
code
def f(arg, n):
    n = 4
    tmp = arg
    del tmp[:-9]
    return len(tmp)>3

try:
    a = ['0','1','2','3','4']
    n = 9
    f(a, n)
    print(n, end=';')
    for el in a:
        print(el, end=';')
except:
    print('ERR')
code

9. Що буде виведено за виконання?
code
class A:
    c = 1
    b = 2

    def __init__(self):
        b = 3

class B(A):
    c = 4

    def __init__(self):
        super().__init__()
        self.a = 7

obj = B()
obj.a = obj.b + 10
B.a = 13

try: print(obj.a, end= ';')
except: print('ERR', end=';')

try: print(A.a, end= ';')
except: print('ERR', end=';')

try: print(B.a, end= ';')
except: print('ERR', end=';')
code

10. 
Для графа на рисунку з вершини 5 запущено алгоритм пошуку в глибину, що будує кістякове дерево. Списки суміжності вершин переглядаються алгоритмом за зростанням міток вершин.\par
\begin{center}
	\adjustimage{max size={0.9\linewidth}{0.9\paperheight}}
	{../10/Traversals/33.png}
\end{center}
\par Які з наведених ребер входять до складу отриманого кістякового дерева?\par
1) (4, 6)\par
2) (1, 5)\par
3) (1, 2)\par
4) (0, 3)\par
5) (0, 1)\par
6) (1, 4)\par
{\vskip 15pt} У формі вибрати номери відповідних ребер.\par
%answer:1 2 4 5
\end {large}}
%end
{\smallskip \begin {small} Затверджено на засіданні кафедри теоретичної кібернетики, протокол \No 12 від   19.05.2020 р.\par\smallskip\smallskip
{\parbox {6.5 cm}{Зав. кафедри \dotfill Ю.В.~Крак}}\hspace {0.7cm}{\hbox to 6.5 cm{ Екзаменатор \dotfill Т.О.~Карнаух}} \end{small}}
%begin
{{\begin {small}\par
{ \center  Київський національний університет імені Тараса Шевченка \par}
{\parbox {5 cm}{Освітня програма \hfill}{\parbox {7 cm}{Прикладна математика, Системний аналіз \hfill}}\parbox {6 cm}{\hfill Семестр 2}}\par
{\parbox {5 cm} {Навчальний предмет \hfill}\parbox {7 cm}{Програмування \hfill}\parbox {6cm}{\hfill}\par}\vspace{-7pt} \end{small}}{\center Екзаменаційний білет \No \textbf 
{ 324 }\par}}
{
\begin {large}
1. Що буде виведено за виконання?
code
def f(n):
    print('f', end='')
    return n<0

try:
    print(f(-3) and f(-9))
except:
    print('ERR')
code

2. Що буде виведено за виконання?
code
a = 6
b = 1
c = 2

def h(a):
    global c
    a = 4
    b = 3
    c = 5
    return a + b + c

a = 7
b = 8
c = 9

try:
    print(h(b), a, b, c, sep=';')
except:
    print('ERR', a, b, c, sep=';')
code

3. Що буде виведено за виконання?
code
try:
    k = 9
    while k<=5:
        k += 1
    print(k, end=';')
except:
    print('ERR')
code

4. Що буде виведено за виконання?
code
try:
    for c in range(9, -14, -3):
        if c<8:
            break
        print(c, end=';');c = -3
    else:
        print(c,end=';')
    print(c)
except:
    print('ERR')
code

5. Що буде виведено за виконання?\par (Дійсні значення, якщо наявні, в цьому завданні будуть виводитись у форматі з фіксованою крапкою та, можливо, з доволі великою кількістю десяткових розрядів. У відповіді для дійсних значень у ЦЬОМУ завданні достатньо вказати один десятковий розряд після крапки, відкинувши решту розрядів без округлення. Наприклад, замість 13.66666666666667 достатньо вказати 13.6, замість 25.23 достатньо вказати 25.2)
code
try:
    print(4, end=';')
    print(int('8'), end=';')
    print(5, end=';')
except ZeroDivisionError:
    print(3, end=';')
except KeyboardInterrupt:
    print(2, end=';')
except:
    print('ERR', end=';')
else:
    print(9, end=';')
finally:
    print(1)
code

6. Що буде виведено за виконання?
code
def f(n):
    print(n%10, end='')
    if n>9:
        f(n//10)
 
f(123456)
code

7. Що буде виведено за виконання?
code
class A:
    b = 1

    def _f1(self): print(2, end=';')
    def _f2(self): print(3, end=';')

    def main(self):
        try: print(self.b, end=';')
        except: print('ERR', end=';')

        try: self._f1()
        except: print('ERR', end=';')

        try: self._f2()
        except: print('ERR', end=';')


class B(A):
    b = 4

    def _f1(self): print(5, end=';')
    def _f2(self): print(6, end=';')

try: B().main()
except: print('ERR', end=';')
code

8. Що буде виведено за виконання?
code
def f(arg, n):
    arg[58] = 4
    n -= -3

try:
    n = 7
    d = {75:8, 71:1, 56:0, 47:0}
    f(d, n)
    print(n, end=';')
    for x in d.values():
        print(x, sep=';', end=';')
except:
    print('ERR')
code

9. Що буде виведено за виконання?
code
class A:
    a = 1
    b = 2

    def __init__(self):
        a = 3

class B(A):
    b = 4

    def __init__(self):
        super().__init__()
        self.c = 7

obj = B()
obj.a = obj.c + 10
B.a = 13

try: print(obj.a, end= ';')
except: print('ERR', end=';')

try: print(A.a, end= ';')
except: print('ERR', end=';')

try: print(B.a, end= ';')
except: print('ERR', end=';')
code

10. 
Для графа на рисунку з вершини 6 запущено алгоритм пошуку в глибину, що будує кістякове дерево. Списки суміжності вершин переглядаються алгоритмом за зростанням міток вершин.\par
\begin{center}
	\adjustimage{max size={0.9\linewidth}{0.9\paperheight}}
	{../10/Traversals/172.png}
\end{center}
\par Які з наведених ребер входять до складу отриманого кістякового дерева?\par
1) (0, 3)\par
2) (1, 6)\par
3) (1, 4)\par
4) (2, 3)\par
5) (0, 2)\par
6) (4, 6)\par
{\vskip 15pt} У формі вибрати номери відповідних ребер.\par
%answer:3 4 5
\end {large}}
%end
{\smallskip \begin {small} Затверджено на засіданні кафедри теоретичної кібернетики, протокол \No 12 від   19.05.2020 р.\par\smallskip\smallskip
{\parbox {6.5 cm}{Зав. кафедри \dotfill Ю.В.~Крак}}\hspace {0.7cm}{\hbox to 6.5 cm{ Екзаменатор \dotfill Т.О.~Карнаух}} \end{small}}
%begin
{{\begin {small}\par
{ \center  Київський національний університет імені Тараса Шевченка \par}
{\parbox {5 cm}{Освітня програма \hfill}{\parbox {7 cm}{Прикладна математика, Системний аналіз \hfill}}\parbox {6 cm}{\hfill Семестр 2}}\par
{\parbox {5 cm} {Навчальний предмет \hfill}\parbox {7 cm}{Програмування \hfill}\parbox {6cm}{\hfill}\par}\vspace{-7pt} \end{small}}{\center Екзаменаційний білет \No \textbf 
{ 325 }\par}}
{
\begin {large}
1. Що буде виведено за виконання?
code
try:
    res = 3**-3-True
    if isinstance(res, bool):
        print(f'bool;{res}')
    elif isinstance(res, int):
        print(f'int;{res}')
    elif isinstance(res, float):
        print(f'float;{res:.2f}')
    else:
        print('???')
except:
    print('ERR')
code

2. Що буде виведено за виконання?
code
a = 4
b = 0
c = 8

def g(b):
    a = 2
    b = 4
    c = 1
    return a + b + c

a = 6
b = 0
c = 9

try:
    print(g(a), a, b, c, sep=';')
except:
    print('ERR', a, b, c, sep=';')
code

3. Що буде виведено за виконання?
code
try:
    j = 5
    while j<6:
        j += 2
    print(j, end=';')
except:
    print('ERR')
code

4. Що буде виведено за виконання?
code
try:
    for c in range(1, -2, 2):
        if c>4:
            continue
        print(c, end=';')
    else:
        print(c,end=';')
    print(c)
except:
    print('ERR')
code

5. Що буде виведено за виконання?\par (Дійсні значення, якщо наявні, в цьому завданні будуть виводитись у форматі з фіксованою крапкою та, можливо, з доволі великою кількістю десяткових розрядів. У відповіді для дійсних значень у ЦЬОМУ завданні достатньо вказати один десятковий розряд після крапки, відкинувши решту розрядів без округлення. Наприклад, замість 13.66666666666667 достатньо вказати 13.6, замість 25.23 достатньо вказати 25.2)
code
try:
    print(6, end=';')
    print(0j>=9j, end=';')
    print(2, end=';')
except TypeError:
    print(9, end=';')
except Exception:
    print(6, end=';')
except:
    print('ERR', end=';')
else:
    print(8, end=';')
finally:
    print(5)
code

6. Що буде виведено за виконання?
code
def f(n):
    if n>9:
        f(n//100)
    print(n%10, end='')
 
f(123456)
code

7. Що буде виведено за виконання?
code
class A:
    b = 1

    def g1(self): print(2, end=';')
    def _g2(self): print(3, end=';')

    def main(self):
        try: print(self.b, end=';')
        except: print('ERR', end=';')

        try: self.g1()
        except: print('ERR', end=';')

        try: self._g2()
        except: print('ERR', end=';')


class B(A):
    b = 4

    def g1(self): print(5, end=';')
    def _g2(self): print(6, end=';')

try: B().main()
except: print('ERR', end=';')
code

8. Що буде виведено за виконання?
code
def f(arg, n):
    arg[74] = 1
    n *= 3

try:
    n = 3
    t = {15:8, 50:6, 58:5, 50:0}
    f(t, n)
    print(n, end=';')
    for x in t.items():
        print(x, sep=';', end=';')
except:
    print('ERR')
code

9. Що буде виведено за виконання?
code
class A:
    c = 1
    a = 2

    def __init__(self):
        a = 3

class B(A):
    c = 4

    def __init__(self):
        super().__init__()
        self.b = 7

obj = B()
obj.c = obj.a + 10
B.c = 13

try: print(obj.c, end= ';')
except: print('ERR', end=';')

try: print(A.c, end= ';')
except: print('ERR', end=';')

try: print(B.c, end= ';')
except: print('ERR', end=';')
code

10. 
Для графа на рисунку з вершини 4 запущено алгоритм пошуку в глибину, що будує кістякове дерево. Списки суміжності вершин переглядаються алгоритмом за зростанням міток вершин.\par
\begin{center}
	\adjustimage{max size={0.9\linewidth}{0.9\paperheight}}
	{../10/Traversals/226.png}
\end{center}
\par Які з наведених ребер входять до складу отриманого кістякового дерева?\par
1) (2, 3)\par
2) (4, 5)\par
3) (0, 6)\par
4) (1, 5)\par
5) (5, 6)\par
6) (0, 2)\par
{\vskip 15pt} У формі вибрати номери відповідних ребер.\par
%answer:4 5 6
\end {large}}
%end
{\smallskip \begin {small} Затверджено на засіданні кафедри теоретичної кібернетики, протокол \No 12 від   19.05.2020 р.\par\smallskip\smallskip
{\parbox {6.5 cm}{Зав. кафедри \dotfill Ю.В.~Крак}}\hspace {0.7cm}{\hbox to 6.5 cm{ Екзаменатор \dotfill Т.О.~Карнаух}} \end{small}}
%begin
{{\begin {small}\par
{ \center  Київський національний університет імені Тараса Шевченка \par}
{\parbox {5 cm}{Освітня програма \hfill}{\parbox {7 cm}{Прикладна математика, Системний аналіз \hfill}}\parbox {6 cm}{\hfill Семестр 2}}\par
{\parbox {5 cm} {Навчальний предмет \hfill}\parbox {7 cm}{Програмування \hfill}\parbox {6cm}{\hfill}\par}\vspace{-7pt} \end{small}}{\center Екзаменаційний білет \No \textbf 
{ 326 }\par}}
{
\begin {large}
1. Що буде виведено за виконання?
code
def f(n):
    print('f', end='')
    return n

try:
    print(f(0) or f(5))
except:
    print('ERR')
code

2. Що буде виведено за виконання?
code
a = 8
b = 2
c = 3

def g(b):
    a = 3
    b += 4
    c = 1
    return a + b + c

a = 1
b = 4
c = 3

try:
    print(g(b), a, b, c, sep=';')
except:
    print('ERR', a, b, c, sep=';')
code

3. Що буде виведено за виконання?
code
try:
    t = 7
    while t>1:
        t += -3
    print(t, end=';')
except:
    print('ERR')
code

4. Що буде виведено за виконання?
code
try:
    for b in range(-4, -9, -1):
        if b>=6:
            continue
        print(b, end=';');b = 4
        if b<=7:
            break
    else:
        print(b,end=';')
    print(b)
except:
    print('ERR')
code

5. Що буде виведено за виконання?\par (Дійсні значення, якщо наявні, в цьому завданні будуть виводитись у форматі з фіксованою крапкою та, можливо, з доволі великою кількістю десяткових розрядів. У відповіді для дійсних значень у ЦЬОМУ завданні достатньо вказати один десятковий розряд після крапки, відкинувши решту розрядів без округлення. Наприклад, замість 13.66666666666667 достатньо вказати 13.6, замість 25.23 достатньо вказати 25.2)
code
try:
    print(3, end=';')
    print(4j>7j, end=';')
    print(7, end=';')
except ValueError:
    print(0, end=';')
except ZeroDivisionError:
    print(1, end=';')
except:
    print('ERR', end=';')
else:
    print(0, end=';')
finally:
    print(4)
code

6. Що буде виведено за виконання?
code
def f(n):
    if n>9:
        return f(n//10)+1
    else:
        return 1

print(f(123456))
code

7. Що буде виведено за виконання?
code
class A:
    c = 1

    def _h1(self): print(2, end=';')
    def h2(self): print(3, end=';')

    def main(self):
        try: print(self.c, end=';')
        except: print('ERR', end=';')

        try: self._h1()
        except: print('ERR', end=';')

        try: self.h2()
        except: print('ERR', end=';')


class B(A):
    c = 4

    def _h1(self): print(5, end=';')
    def h2(self): print(6, end=';')

try: B().main()
except: print('ERR', end=';')
code

8. Що буде виведено за виконання?
code
def f(arg, n):
    arg[34] = 0
    n = 5
    arg = tuple(arg.items())

try:
    n = 0
    d = {56:2, 23:2, 34:6, 56:2}
    f(d, n)
    print(n, end=';')
    for x in d.keys():
        print(x, sep=';', end=';')
except:
    print('ERR')
code

9. Що буде виведено за виконання?
code
class A:
    a = 1
    b = 2

    def __init__(self):
        a = 3

class B(A):
    b = 4

    def __init__(self):
        super().__init__()
        self.c = 7

obj = B()
obj.b = obj.b + 10
B.b = 13

try: print(obj.b, end= ';')
except: print('ERR', end=';')

try: print(A.b, end= ';')
except: print('ERR', end=';')

try: print(B.b, end= ';')
except: print('ERR', end=';')
code

10. 
Для графа на рисунку з вершини 3 запущено алгоритм пошуку в глибину, що будує кістякове дерево. Списки суміжності вершин переглядаються алгоритмом за зростанням міток вершин.\par
\begin{center}
	\adjustimage{max size={0.9\linewidth}{0.9\paperheight}}
	{../10/Traversals/157.png}
\end{center}
\par Які з наведених ребер входять до складу отриманого кістякового дерева?\par
1) (0, 2)\par
2) (2, 4)\par
3) (1, 4)\par
4) (0, 5)\par
5) (3, 5)\par
6) (3, 6)\par
{\vskip 15pt} У формі вибрати номери відповідних ребер.\par
%answer:1 4
\end {large}}
%end
{\smallskip \begin {small} Затверджено на засіданні кафедри теоретичної кібернетики, протокол \No 12 від   19.05.2020 р.\par\smallskip\smallskip
{\parbox {6.5 cm}{Зав. кафедри \dotfill Ю.В.~Крак}}\hspace {0.7cm}{\hbox to 6.5 cm{ Екзаменатор \dotfill Т.О.~Карнаух}} \end{small}}
%begin
{{\begin {small}\par
{ \center  Київський національний університет імені Тараса Шевченка \par}
{\parbox {5 cm}{Освітня програма \hfill}{\parbox {7 cm}{Прикладна математика, Системний аналіз \hfill}}\parbox {6 cm}{\hfill Семестр 2}}\par
{\parbox {5 cm} {Навчальний предмет \hfill}\parbox {7 cm}{Програмування \hfill}\parbox {6cm}{\hfill}\par}\vspace{-7pt} \end{small}}{\center Екзаменаційний білет \No \textbf 
{ 327 }\par}}
{
\begin {large}
1. Що буде виведено за виконання?
code
try:
    res = -1**-2+False
    if isinstance(res, bool):
        print(f'bool;{res}')
    elif isinstance(res, int):
        print(f'int;{res}')
    elif isinstance(res, float):
        print(f'float;{res:.2f}')
    else:
        print('???')
except:
    print('ERR')
code

2. Що буде виведено за виконання?
code
a = 2
b = 5
c = 4

def f(b):
    a = 2
    b = 5
    c = 3
    return a + b + c

a = 7
b = 1
c = 3

try:
    print(f(a), a, b, c, sep=';')
except:
    print('ERR', a, b, c, sep=';')
code

3. Що буде виведено за виконання?
code
try:
    k = 8
    while k<=12:
        k += 1
    print(k, end=';')
except:
    print('ERR')
code

4. Що буде виведено за виконання?
code
try:
    for a in range(-11, 10, -3):
        if a<8:
            break
        print(a, end=';')
    else:
        print('end',end=';')
    print(a)
except:
    print('ERR')
code

5. Що буде виведено за виконання?\par (Дійсні значення, якщо наявні, в цьому завданні будуть виводитись у форматі з фіксованою крапкою та, можливо, з доволі великою кількістю десяткових розрядів. У відповіді для дійсних значень у ЦЬОМУ завданні достатньо вказати один десятковий розряд після крапки, відкинувши решту розрядів без округлення. Наприклад, замість 13.66666666666667 достатньо вказати 13.6, замість 25.23 достатньо вказати 25.2)
code
try:
    print(7, end=';')
    print(2/0, end=';')
    print(9, end=';')
except TypeError:
    print(1, end=';')
except Exception:
    print(6, end=';')
except:
    print('ERR', end=';')
else:
    print(3, end=';')
finally:
    print(5)
code

6. Що буде виведено за виконання?
code
def f(s):
    for el in s:
        el += 1
    return s
 
s = [1, 2, 3]
s = f(s)
print(s)
code

7. Що буде виведено за виконання?
code
class A:
    d = 1

    def f1(self): print(2, end=';')
    def _f2(self): print(3, end=';')

    def main(self):
        try: print(self.d, end=';')
        except: print('ERR', end=';')

        try: self.f1()
        except: print('ERR', end=';')

        try: self._f2()
        except: print('ERR', end=';')


class B(A):
    d = 4

    def f1(self): print(5, end=';')
    def _f2(self): print(6, end=';')

try: B().main()
except: print('ERR', end=';')
code

8. Що буде виведено за виконання?
code
def f(arg, n):
    arg[77] = 2
    n += -2

try:
    n = 5
    s = {15:3, 63:8, 12:0, 63:3}
    f(s, n)
    print(n, end=';')
    for x in s :
        print(x, sep=';', end=';')
except:
    print('ERR')
code

9. Що буде виведено за виконання?
code
class A:
    b = 1
    a = 2

    def __init__(self):
        a = 3

class B(A):
    b = 4

    def __init__(self):
        super().__init__()
        self.c = 7

obj = B()
obj.c = obj.a + 10
B.c = 13

try: print(obj.c, end= ';')
except: print('ERR', end=';')

try: print(A.c, end= ';')
except: print('ERR', end=';')

try: print(B.c, end= ';')
except: print('ERR', end=';')
code

10. 
Для графа на рисунку з вершини 4 запущено алгоритм пошуку в глибину, що будує кістякове дерево. Списки суміжності вершин переглядаються алгоритмом за зростанням міток вершин.\par
\begin{center}
	\adjustimage{max size={0.9\linewidth}{0.9\paperheight}}
	{../10/Traversals/121.png}
\end{center}
\par Які з наведених ребер входять до складу отриманого кістякового дерева?\par
1) (4, 5)\par
2) (0, 6)\par
3) (2, 5)\par
4) (2, 3)\par
5) (4, 6)\par
6) (2, 6)\par
{\vskip 15pt} У формі вибрати номери відповідних ребер.\par
%answer:2 4
\end {large}}
%end
{\smallskip \begin {small} Затверджено на засіданні кафедри теоретичної кібернетики, протокол \No 12 від   19.05.2020 р.\par\smallskip\smallskip
{\parbox {6.5 cm}{Зав. кафедри \dotfill Ю.В.~Крак}}\hspace {0.7cm}{\hbox to 6.5 cm{ Екзаменатор \dotfill Т.О.~Карнаух}} \end{small}}
%begin
{{\begin {small}\par
{ \center  Київський національний університет імені Тараса Шевченка \par}
{\parbox {5 cm}{Освітня програма \hfill}{\parbox {7 cm}{Прикладна математика, Системний аналіз \hfill}}\parbox {6 cm}{\hfill Семестр 2}}\par
{\parbox {5 cm} {Навчальний предмет \hfill}\parbox {7 cm}{Програмування \hfill}\parbox {6cm}{\hfill}\par}\vspace{-7pt} \end{small}}{\center Екзаменаційний білет \No \textbf 
{ 328 }\par}}
{
\begin {large}
1. Що буде виведено за виконання?
code
try:
    res = 3/7*3*7
    if isinstance(res, bool):
        print(f'bool;{res}')
    elif isinstance(res, int):
        print(f'int;{res}')
    elif isinstance(res, float):
        print(f'float;{res:.2f}')
    else:
        print('???')
except:
    print('ERR')
code

2. Що буде виведено за виконання?
code
a = 2
b = 9
c = 1

def g(b):
    global c
    a *= 5
    b = 2
    c = 4
    return a + b + c

a = 8
b = 6
c = 7

try:
    print(g(a), a, b, c, sep=';')
except:
    print('ERR', a, b, c, sep=';')
code

3. Що буде виведено за виконання?
code
try:
    m = 5
    while m>=4:
        m += -2
    print(m, end=';')
except:
    print('ERR')
code

4. Що буде виведено за виконання?
code
try:
    for d in range(-5, 15, -2):
        if d<=5:
            break
        print(d, end=';');d = 7
    else:
        print(d,end=';')
    print(d)
except:
    print('ERR')
code

5. Що буде виведено за виконання?\par (Дійсні значення, якщо наявні, в цьому завданні будуть виводитись у форматі з фіксованою крапкою та, можливо, з доволі великою кількістю десяткових розрядів. У відповіді для дійсних значень у ЦЬОМУ завданні достатньо вказати один десятковий розряд після крапки, відкинувши решту розрядів без округлення. Наприклад, замість 13.66666666666667 достатньо вказати 13.6, замість 25.23 достатньо вказати 25.2)
code
try:
    print(8, end=';')
    print(int('a8'), end=';')
    print(7, end=';')
except KeyboardInterrupt:
    print(0, end=';')
except ValueError:
    print(2, end=';')
except:
    print('ERR', end=';')
else:
    print(3, end=';')
finally:
    print(9)
code

6. Що буде виведено за виконання?
code
def f(n):
    if n<=9:
        print(n%10,end='')
    else:
        f(n//10)

f(543216)
code

7. Що буде виведено за виконання?
code
class A:
    a = 1

    def _h1(self): print(2, end=';')
    def __h2(self): print(3, end=';')

    def main(self):
        try: print(self.a, end=';')
        except: print('ERR', end=';')

        try: self._h1()
        except: print('ERR', end=';')

        try: self.__h2()
        except: print('ERR', end=';')


class B(A):
    a = 4

    def _h1(self): print(5, end=';')
    def __h2(self): print(6, end=';')

try: B().main()
except: print('ERR', end=';')
code

8. Що буде виведено за виконання?
code
def f(arg, n):
    arg[56] = 2
    n -= 2

try:
    n = 4
    t = {52:6, 54:0, 77:8}
    f(t, n)
    print(n, end=';')
    for x in t.keys():
        print(x, sep=';', end=';')
except:
    print('ERR')
code

9. Що буде виведено за виконання?
code
class A:
    a = 1
    c = 2

    def __init__(self):
        a = 3

class B(A):
    c = 4

    def __init__(self):
        super().__init__()
        self.b = 7

obj = B()
obj.b = obj.c + 10
B.b = 13

try: print(obj.b, end= ';')
except: print('ERR', end=';')

try: print(A.b, end= ';')
except: print('ERR', end=';')

try: print(B.b, end= ';')
except: print('ERR', end=';')
code

10. 
Для графа на рисунку з вершини 5 запущено алгоритм пошуку в глибину, що будує кістякове дерево. Списки суміжності вершин переглядаються алгоритмом за зростанням міток вершин.\par
\begin{center}
	\adjustimage{max size={0.9\linewidth}{0.9\paperheight}}
	{../10/Traversals/239.png}
\end{center}
\par Які з наведених ребер входять до складу отриманого кістякового дерева?\par
1) (1, 4)\par
2) (3, 4)\par
3) (5, 6)\par
4) (0, 5)\par
5) (4, 6)\par
6) (4, 5)\par
{\vskip 15pt} У формі вибрати номери відповідних ребер.\par
%answer:1 2 4
\end {large}}
%end
{\smallskip \begin {small} Затверджено на засіданні кафедри теоретичної кібернетики, протокол \No 12 від   19.05.2020 р.\par\smallskip\smallskip
{\parbox {6.5 cm}{Зав. кафедри \dotfill Ю.В.~Крак}}\hspace {0.7cm}{\hbox to 6.5 cm{ Екзаменатор \dotfill Т.О.~Карнаух}} \end{small}}
%begin
{{\begin {small}\par
{ \center  Київський національний університет імені Тараса Шевченка \par}
{\parbox {5 cm}{Освітня програма \hfill}{\parbox {7 cm}{Прикладна математика, Системний аналіз \hfill}}\parbox {6 cm}{\hfill Семестр 2}}\par
{\parbox {5 cm} {Навчальний предмет \hfill}\parbox {7 cm}{Програмування \hfill}\parbox {6cm}{\hfill}\par}\vspace{-7pt} \end{small}}{\center Екзаменаційний білет \No \textbf 
{ 329 }\par}}
{
\begin {large}
1. Що буде виведено за виконання?
code
try:
    res = 3//5/4*9
    if isinstance(res, bool):
        print(f'bool;{res}')
    elif isinstance(res, int):
        print(f'int;{res}')
    elif isinstance(res, float):
        print(f'float;{res:.2f}')
    else:
        print('???')
except:
    print('ERR')
code

2. Що буде виведено за виконання?
code
a = 5
b = 0
c = 0

def h(a):
    a = 3
    b -= 1
    c = 2
    return a + b + c

a = 1
b = 6
c = 4

try:
    print(h(a), a, b, c, sep=';')
except:
    print('ERR', a, b, c, sep=';')
code

3. Що буде виведено за виконання?
code
try:
    t = 6
    while t<=9:
        t += 2
    print(t, end=';')
except:
    print('ERR')
code

4. Що буде виведено за виконання?
code
try:
    for f in range(-9, -7, 2):
        if f>7:
            continue
        print(f, end=';')
    else:
        print(f,end=';')
    print(f)
except:
    print('ERR')
code

5. Що буде виведено за виконання?\par (Дійсні значення, якщо наявні, в цьому завданні будуть виводитись у форматі з фіксованою крапкою та, можливо, з доволі великою кількістю десяткових розрядів. У відповіді для дійсних значень у ЦЬОМУ завданні достатньо вказати один десятковий розряд після крапки, відкинувши решту розрядів без округлення. Наприклад, замість 13.66666666666667 достатньо вказати 13.6, замість 25.23 достатньо вказати 25.2)
code
try:
    print(4, end=';')
    print(int('6'), end=';')
    print(1, end=';')
except KeyboardInterrupt:
    print(5, end=';')
except TypeError:
    print(2, end=';')
except:
    print('ERR', end=';')
else:
    print(0, end=';')
finally:
    print(9)
code

6. Що буде виведено за виконання?
code
def f(s):
    s1 = s
    for i in range(len(s)):
        s1[i] += 1
    return s1

s = [1, 2, 3]
f(s)
print(s)
code

7. Що буде виведено за виконання?
code
class A:
    a = 1

    def _g1(self): print(2, end=';')
    def _g2(self): print(3, end=';')

    def main(self):
        try: print(self.a, end=';')
        except: print('ERR', end=';')

        try: self._g1()
        except: print('ERR', end=';')

        try: self._g2()
        except: print('ERR', end=';')


class B(A):
    a = 4

    def _g1(self): print(5, end=';')
    def _g2(self): print(6, end=';')

try: B().main()
except: print('ERR', end=';')
code

8. Що буде виведено за виконання?
code
def f(arg, n):
    arg[73] = 4
    n *= -3

try:
    n = 6
    d = {24:4, 52:2, 47:6, 24:6}
    f(d, n)
    print(n, end=';')
    for x, y in d.items():
        print(x, y, sep=';', end=';')
except:
    print('ERR')
code

9. Що буде виведено за виконання?
code
class A:
    c = 1
    b = 2

    def __init__(self):
        b = 3

class B(A):
    b = 4

    def __init__(self):
        super().__init__()
        self.a = 7

obj = B()
obj.a = obj.b + 10
B.a = 13

try: print(obj.a, end= ';')
except: print('ERR', end=';')

try: print(A.a, end= ';')
except: print('ERR', end=';')

try: print(B.a, end= ';')
except: print('ERR', end=';')
code

10. 
Для графа на рисунку з вершини 1 запущено алгоритм пошуку в глибину, що будує кістякове дерево. Списки суміжності вершин переглядаються алгоритмом за зростанням міток вершин.\par
\begin{center}
	\adjustimage{max size={0.9\linewidth}{0.9\paperheight}}
	{../10/Traversals/218.png}
\end{center}
\par Які з наведених ребер входять до складу отриманого кістякового дерева?\par
1) (2, 5)\par
2) (0, 1)\par
3) (2, 4)\par
4) (0, 4)\par
5) (0, 6)\par
6) (3, 5)\par
{\vskip 15pt} У формі вибрати номери відповідних ребер.\par
%answer:1 2 3 4 6
\end {large}}
%end
{\smallskip \begin {small} Затверджено на засіданні кафедри теоретичної кібернетики, протокол \No 12 від   19.05.2020 р.\par\smallskip\smallskip
{\parbox {6.5 cm}{Зав. кафедри \dotfill Ю.В.~Крак}}\hspace {0.7cm}{\hbox to 6.5 cm{ Екзаменатор \dotfill Т.О.~Карнаух}} \end{small}}
%begin
{{\begin {small}\par
{ \center  Київський національний університет імені Тараса Шевченка \par}
{\parbox {5 cm}{Освітня програма \hfill}{\parbox {7 cm}{Прикладна математика, Системний аналіз \hfill}}\parbox {6 cm}{\hfill Семестр 2}}\par
{\parbox {5 cm} {Навчальний предмет \hfill}\parbox {7 cm}{Програмування \hfill}\parbox {6cm}{\hfill}\par}\vspace{-7pt} \end{small}}{\center Екзаменаційний білет \No \textbf 
{ 330 }\par}}
{
\begin {large}
1. Що буде виведено за виконання?
code
try:
    res = 4**-1.0*2
    if isinstance(res, bool):
        print(f'bool;{res}')
    elif isinstance(res, int):
        print(f'int;{res}')
    elif isinstance(res, float):
        print(f'float;{res:.2f}')
    else:
        print('???')
except:
    print('ERR')
code

2. Що буде виведено за виконання?
code
a = 5
b = 4
c = 0

def g(a):
    global c
    a = 1
    b = 5
    c = 3
    return a + b + c

a = 9
b = 7
c = 1

try:
    print(g(b), a, b, c, sep=';')
except:
    print('ERR', a, b, c, sep=';')
code

3. Що буде виведено за виконання?
code
try:
    n = 11
    while n>0:
        n += -3
    print(n, end=';')
except:
    print('ERR')
code

4. Що буде виведено за виконання?
code
try:
    for e in range(5, 5, 3):
        if e<=3:
            continue
        print(e, end=';')
        if e<8:
            break
    else:
        print(e,end=';')
    print(e)
except:
    print('ERR')
code

5. Що буде виведено за виконання?\par (Дійсні значення, якщо наявні, в цьому завданні будуть виводитись у форматі з фіксованою крапкою та, можливо, з доволі великою кількістю десяткових розрядів. У відповіді для дійсних значень у ЦЬОМУ завданні достатньо вказати один десятковий розряд після крапки, відкинувши решту розрядів без округлення. Наприклад, замість 13.66666666666667 достатньо вказати 13.6, замість 25.23 достатньо вказати 25.2)
code
try:
    print(8, end=';')
    print(1%2, end=';')
    print(6, end=';')
except ZeroDivisionError:
    print(7, end=';')
except Exception:
    print(3, end=';')
except:
    print('ERR', end=';')
else:
    print(5, end=';')
finally:
    print(4)
code

6. Що буде виведено за виконання?
code
def f(n):
    if n>9:
        f(n//100)
    print(n%10,end='')
 
f(987654)
code

7. Що буде виведено за виконання?
code
class A:
    c = 1

    def __f1(self): print(2, end=';')
    def _f2(self): print(3, end=';')

    def main(self):
        try: print(self.c, end=';')
        except: print('ERR', end=';')

        try: self.__f1()
        except: print('ERR', end=';')

        try: self._f2()
        except: print('ERR', end=';')


class B(A):
    c = 4

    def __f1(self): print(5, end=';')
    def _f2(self): print(6, end=';')

try: B().main()
except: print('ERR', end=';')
code

8. Що буде виведено за виконання?
code
def f(arg, n):
    n = 3
    tmp = arg
    tmp[-6:-4] = [1,2]
    return len(tmp)>3

try:
    b = [10,11,12,13,14]
    n = 8
    f(b, n)
    print(n, end=';')
    for el in b:
        print(el, end=';')
except:
    print('ERR')
code

9. Що буде виведено за виконання?
code
class A:
    b = 1
    c = 2

    def __init__(self):
        b = 3

class B(A):
    b = 4

    def __init__(self):
        super().__init__()
        self.a = 7

obj = B()
obj.c = obj.a + 10
B.c = 13

try: print(obj.c, end= ';')
except: print('ERR', end=';')

try: print(A.c, end= ';')
except: print('ERR', end=';')

try: print(B.c, end= ';')
except: print('ERR', end=';')
code

10. 
Для графа на рисунку з вершини 5 запущено алгоритм пошуку в глибину, що будує кістякове дерево. Списки суміжності вершин переглядаються алгоритмом за зростанням міток вершин.\par
\begin{center}
	\adjustimage{max size={0.9\linewidth}{0.9\paperheight}}
	{../10/Traversals/113.png}
\end{center}
\par Які з наведених ребер входять до складу отриманого кістякового дерева?\par
1) (0, 2)\par
2) (2, 4)\par
3) (4, 6)\par
4) (3, 4)\par
5) (0, 6)\par
6) (1, 6)\par
{\vskip 15pt} У формі вибрати номери відповідних ребер.\par
%answer:1 2 6
\end {large}}
%end
{\smallskip \begin {small} Затверджено на засіданні кафедри теоретичної кібернетики, протокол \No 12 від   19.05.2020 р.\par\smallskip\smallskip
{\parbox {6.5 cm}{Зав. кафедри \dotfill Ю.В.~Крак}}\hspace {0.7cm}{\hbox to 6.5 cm{ Екзаменатор \dotfill Т.О.~Карнаух}} \end{small}}
%begin
{{\begin {small}\par
{ \center  Київський національний університет імені Тараса Шевченка \par}
{\parbox {5 cm}{Освітня програма \hfill}{\parbox {7 cm}{Прикладна математика, Системний аналіз \hfill}}\parbox {6 cm}{\hfill Семестр 2}}\par
{\parbox {5 cm} {Навчальний предмет \hfill}\parbox {7 cm}{Програмування \hfill}\parbox {6cm}{\hfill}\par}\vspace{-7pt} \end{small}}{\center Екзаменаційний білет \No \textbf 
{ 331 }\par}}
{
\begin {large}
1. Що буде виведено за виконання?
code
try:
    res = 1**0.5*-2
    if isinstance(res, bool):
        print(f'bool;{res}')
    elif isinstance(res, int):
        print(f'int;{res}')
    elif isinstance(res, float):
        print(f'float;{res:.2f}')
    else:
        print('???')
except:
    print('ERR')
code

2. Що буде виведено за виконання?
code
a = 5
b = 7
c = 8

def f(a):
    a += 3
    b = 5
    c = 4
    return a + b + c

a = 3
b = 2
c = 9

try:
    print(f(a), a, b, c, sep=';')
except:
    print('ERR', a, b, c, sep=';')
code

3. Що буде виведено за виконання?
code
try:
    k = 5
    while k>=2:
        k += -2
    print(k, end=';')
except:
    print('ERR')
code

4. Що буде виведено за виконання?
code
try:
    for d in range(13, 2, 2):
        if d>=3:
            continue
        print(d, end=';');d = -8
    else:
        print('end',end=';')
    print(d)
except:
    print('ERR')
code

5. Що буде виведено за виконання?\par (Дійсні значення, якщо наявні, в цьому завданні будуть виводитись у форматі з фіксованою крапкою та, можливо, з доволі великою кількістю десяткових розрядів. У відповіді для дійсних значень у ЦЬОМУ завданні достатньо вказати один десятковий розряд після крапки, відкинувши решту розрядів без округлення. Наприклад, замість 13.66666666666667 достатньо вказати 13.6, замість 25.23 достатньо вказати 25.2)
code
try:
    print(4, end=';')
    print(2%0.0, end=';')
    print(8, end=';')
except Exception:
    print(3, end=';')
except KeyboardInterrupt:
    print(0, end=';')
except:
    print('ERR', end=';')
else:
    print(6, end=';')
finally:
    print(7)
code

6. Що буде виведено за виконання?
code
def f(n):
    if n>9:
        f(n//10)
    print(n%10, end='')
 
f(123456)
code

7. Що буде виведено за виконання?
code
class A:
    d = 1

    def _h1(self): print(2, end=';')
    def _h2(self): print(3, end=';')

    def main(self):
        try: print(self.d, end=';')
        except: print('ERR', end=';')

        try: self._h1()
        except: print('ERR', end=';')

        try: self._h2()
        except: print('ERR', end=';')


class B(A):
    d = 4

    def _h1(self): print(5, end=';')
    def _h2(self): print(6, end=';')

try: B().main()
except: print('ERR', end=';')
code

8. Що буде виведено за виконання?
code
def f(arg, n):
    arg[70] = 3
    n += -2

try:
    n = 6
    s = {70:3, 66:9, 70:9}
    f(s, n)
    print(n, end=';')
    for x in s.items():
        print(x, sep=';', end=';')
except:
    print('ERR')
code

9. Що буде виведено за виконання?
code
class A:
    b = 1
    a = 2

    def __init__(self):
        a = 3

class B(A):
    a = 4

    def __init__(self):
        super().__init__()
        self.c = 7

obj = B()
obj.b = obj.c + 10
B.b = 13

try: print(obj.b, end= ';')
except: print('ERR', end=';')

try: print(A.b, end= ';')
except: print('ERR', end=';')

try: print(B.b, end= ';')
except: print('ERR', end=';')
code

10. 
Для графа на рисунку з вершини 1 запущено алгоритм пошуку в глибину, що будує кістякове дерево. Списки суміжності вершин переглядаються алгоритмом за зростанням міток вершин.\par
\begin{center}
	\adjustimage{max size={0.9\linewidth}{0.9\paperheight}}
	{../10/Traversals/18.png}
\end{center}
\par Які з наведених ребер входять до складу отриманого кістякового дерева?\par
1) (2, 5)\par
2) (0, 2)\par
3) (4, 6)\par
4) (4, 5)\par
5) (3, 4)\par
6) (1, 6)\par
{\vskip 15pt} У формі вибрати номери відповідних ребер.\par
%answer:1 2 3 4 5
\end {large}}
%end
{\smallskip \begin {small} Затверджено на засіданні кафедри теоретичної кібернетики, протокол \No 12 від   19.05.2020 р.\par\smallskip\smallskip
{\parbox {6.5 cm}{Зав. кафедри \dotfill Ю.В.~Крак}}\hspace {0.7cm}{\hbox to 6.5 cm{ Екзаменатор \dotfill Т.О.~Карнаух}} \end{small}}
%begin
{{\begin {small}\par
{ \center  Київський національний університет імені Тараса Шевченка \par}
{\parbox {5 cm}{Освітня програма \hfill}{\parbox {7 cm}{Прикладна математика, Системний аналіз \hfill}}\parbox {6 cm}{\hfill Семестр 2}}\par
{\parbox {5 cm} {Навчальний предмет \hfill}\parbox {7 cm}{Програмування \hfill}\parbox {6cm}{\hfill}\par}\vspace{-7pt} \end{small}}{\center Екзаменаційний білет \No \textbf 
{ 332 }\par}}
{
\begin {large}
1. Що буде виведено за виконання?
code
try:
    res = 8/3/6*7
    if isinstance(res, bool):
        print(f'bool;{res}')
    elif isinstance(res, int):
        print(f'int;{res}')
    elif isinstance(res, float):
        print(f'float;{res:.2f}')
    else:
        print('???')
except:
    print('ERR')
code

2. Що буде виведено за виконання?
code
a = 3
b = 6
c = 8

def h(a):
    a = 2
    b = 4
    c = 5
    return a + b + c

a = 2
b = 0
c = 6

try:
    print(h(b), a, b, c, sep=';')
except:
    print('ERR', a, b, c, sep=';')
code

3. Що буде виведено за виконання?
code
try:
    n = 0
    while n<10:
        n += 1
    print(n, end=';')
except:
    print('ERR')
code

4. Що буде виведено за виконання?
code
try:
    for e in range(10, -15, -3):
        if e<5:
            break
        print(e, end=';')
    else:
        print(e,end=';')
    print(e)
except:
    print('ERR')
code

5. Що буде виведено за виконання?\par (Дійсні значення, якщо наявні, в цьому завданні будуть виводитись у форматі з фіксованою крапкою та, можливо, з доволі великою кількістю десяткових розрядів. У відповіді для дійсних значень у ЦЬОМУ завданні достатньо вказати один десятковий розряд після крапки, відкинувши решту розрядів без округлення. Наприклад, замість 13.66666666666667 достатньо вказати 13.6, замість 25.23 достатньо вказати 25.2)
code
try:
    print(9, end=';')
    print(1//1, end=';')
    print(5, end=';')
except TypeError:
    print(3, end=';')
except ZeroDivisionError:
    print(2, end=';')
except:
    print('ERR', end=';')
else:
    print(6, end=';')
finally:
    print(7)
code

6. Що буде виведено за виконання?
code
def f(n):
    if n>2:
        f(n//2)
    else:
        print(n%2, end='')

f(16)
code

7. Що буде виведено за виконання?
code
class A:
    b = 1

    def __g1(self): print(2, end=';')
    def _g2(self): print(3, end=';')

    def main(self):
        try: print(self.b, end=';')
        except: print('ERR', end=';')

        try: self.__g1()
        except: print('ERR', end=';')

        try: self._g2()
        except: print('ERR', end=';')


class B(A):
    b = 4

    def __g1(self): print(5, end=';')
    def _g2(self): print(6, end=';')

try: B().main()
except: print('ERR', end=';')
code

8. Що буде виведено за виконання?
code
def f(arg, n):
    arg[69] = 5
    n *= 3

try:
    n = 7
    d = {69:0, 39:6, 29:5}
    f(d, n)
    print(n, end=';')
    for x in d :
        print(x, sep=';', end=';')
except:
    print('ERR')
code

9. Що буде виведено за виконання?
code
class A:
    c = 1
    b = 2

    def __init__(self):
        c = 3

class B(A):
    c = 4

    def __init__(self):
        super().__init__()
        self.a = 7

obj = B()
obj.a = obj.b + 10
B.a = 13

try: print(obj.a, end= ';')
except: print('ERR', end=';')

try: print(A.a, end= ';')
except: print('ERR', end=';')

try: print(B.a, end= ';')
except: print('ERR', end=';')
code

10. 
Для графа на рисунку з вершини 1 запущено алгоритм пошуку в глибину, що будує кістякове дерево. Списки суміжності вершин переглядаються алгоритмом за зростанням міток вершин.\par
\begin{center}
	\adjustimage{max size={0.9\linewidth}{0.9\paperheight}}
	{../10/Traversals/43.png}
\end{center}
\par Які з наведених ребер входять до складу отриманого кістякового дерева?\par
1) (3, 4)\par
2) (1, 4)\par
3) (2, 5)\par
4) (0, 6)\par
5) (0, 2)\par
6) (0, 1)\par
{\vskip 15pt} У формі вибрати номери відповідних ребер.\par
%answer:1 2 3 5 6
\end {large}}
%end
{\smallskip \begin {small} Затверджено на засіданні кафедри теоретичної кібернетики, протокол \No 12 від   19.05.2020 р.\par\smallskip\smallskip
{\parbox {6.5 cm}{Зав. кафедри \dotfill Ю.В.~Крак}}\hspace {0.7cm}{\hbox to 6.5 cm{ Екзаменатор \dotfill Т.О.~Карнаух}} \end{small}}
%begin
{{\begin {small}\par
{ \center  Київський національний університет імені Тараса Шевченка \par}
{\parbox {5 cm}{Освітня програма \hfill}{\parbox {7 cm}{Прикладна математика, Системний аналіз \hfill}}\parbox {6 cm}{\hfill Семестр 2}}\par
{\parbox {5 cm} {Навчальний предмет \hfill}\parbox {7 cm}{Програмування \hfill}\parbox {6cm}{\hfill}\par}\vspace{-7pt} \end{small}}{\center Екзаменаційний білет \No \textbf 
{ 333 }\par}}
{
\begin {large}
1. Що буде виведено за виконання?
code
def f(n):
    print('f', end='')
    return n

try:
    print(f(6) and f(-3))
except:
    print('ERR')
code

2. Що буде виведено за виконання?
code
a = 8
b = 4
c = 5

def f(b):
    global c
    a = 4
    b = 1
    c = 3
    return a + b + c

a = 2
b = 9
c = 3

try:
    print(f(a), a, b, c, sep=';')
except:
    print('ERR', a, b, c, sep=';')
code

3. Що буде виведено за виконання?
code
try:
    m = 2
    while m<=13:
        m += 2
    print(m, end=';')
except:
    print('ERR')
code

4. Що буде виведено за виконання?
code
try:
    for b in range(-2, -10, 3):
        if b>8:
            continue
        print(b, end=';')
    else:
        print(b,end=';')
    print(b)
except:
    print('ERR')
code

5. Що буде виведено за виконання?\par (Дійсні значення, якщо наявні, в цьому завданні будуть виводитись у форматі з фіксованою крапкою та, можливо, з доволі великою кількістю десяткових розрядів. У відповіді для дійсних значень у ЦЬОМУ завданні достатньо вказати один десятковий розряд після крапки, відкинувши решту розрядів без округлення. Наприклад, замість 13.66666666666667 достатньо вказати 13.6, замість 25.23 достатньо вказати 25.2)
code
try:
    print(1, end=';')
    print(int('c8'), end=';')
    print(4, end=';')
except ValueError:
    print(9, end=';')
except ValueError:
    print(5, end=';')
except:
    print('ERR', end=';')
else:
    print(0, end=';')
finally:
    print(5)
code

6. Що буде виведено за виконання?
code
def f(n):
    print(n%2, end='')
    if n>2:
        f(n//2)
 
f(16)
code

7. Що буде виведено за виконання?
code
class A:
    b = 1

    def f1(self): print(2, end=';')
    def _f2(self): print(3, end=';')

    def main(self):
        try: print(self.b, end=';')
        except: print('ERR', end=';')

        try: self.f1()
        except: print('ERR', end=';')

        try: self._f2()
        except: print('ERR', end=';')


class B(A):
    b = 4

    def f1(self): print(5, end=';')
    def _f2(self): print(6, end=';')

try: B().main()
except: print('ERR', end=';')
code

8. Що буде виведено за виконання?
code
def f(arg, n):
    n = 8
    tmp = arg
    tmp[-9:] = 'bc'
    return len(tmp)>3

try:
    g = ['0','1','2','3','4']
    n = 9
    f(g, n)
    print(n, end=';')
    for el in g:
        print(el, end=';')
except:
    print('ERR')
code

9. Що буде виведено за виконання?
code
class A:
    a = 1
    b = 2

    def __init__(self):
        b = 3

class B(A):
    b = 4

    def __init__(self):
        super().__init__()
        self.c = 7

obj = B()
obj.c = obj.a + 10
B.c = 13

try: print(obj.c, end= ';')
except: print('ERR', end=';')

try: print(A.c, end= ';')
except: print('ERR', end=';')

try: print(B.c, end= ';')
except: print('ERR', end=';')
code

10. 
Для графа на рисунку з вершини 2 запущено алгоритм пошуку в глибину, що будує кістякове дерево. Списки суміжності вершин переглядаються алгоритмом за зростанням міток вершин.\par
\begin{center}
	\adjustimage{max size={0.9\linewidth}{0.9\paperheight}}
	{../10/Traversals/132.png}
\end{center}
\par Які з наведених ребер входять до складу отриманого кістякового дерева?\par
1) (0, 4)\par
2) (0, 5)\par
3) (0, 1)\par
4) (1, 2)\par
5) (3, 6)\par
6) (2, 4)\par
{\vskip 15pt} У формі вибрати номери відповідних ребер.\par
%answer:1 2 3 4 5
\end {large}}
%end
{\smallskip \begin {small} Затверджено на засіданні кафедри теоретичної кібернетики, протокол \No 12 від   19.05.2020 р.\par\smallskip\smallskip
{\parbox {6.5 cm}{Зав. кафедри \dotfill Ю.В.~Крак}}\hspace {0.7cm}{\hbox to 6.5 cm{ Екзаменатор \dotfill Т.О.~Карнаух}} \end{small}}
%begin
{{\begin {small}\par
{ \center  Київський національний університет імені Тараса Шевченка \par}
{\parbox {5 cm}{Освітня програма \hfill}{\parbox {7 cm}{Прикладна математика, Системний аналіз \hfill}}\parbox {6 cm}{\hfill Семестр 2}}\par
{\parbox {5 cm} {Навчальний предмет \hfill}\parbox {7 cm}{Програмування \hfill}\parbox {6cm}{\hfill}\par}\vspace{-7pt} \end{small}}{\center Екзаменаційний білет \No \textbf 
{ 334 }\par}}
{
\begin {large}
1. Що буде виведено за виконання?
code
try:
    res = 1.0j!="8.0"
    if isinstance(res, bool):
        print(f'bool;{res}')
    elif isinstance(res, int):
        print(f'int;{res}')
    elif isinstance(res, float):
        print(f'float;{res:.2f}')
    else:
        print('???')
except:
    print('ERR')
code

2. Що буде виведено за виконання?
code
a = 7
b = 4
c = 2

def g(b):
    global c
    a = 2
    b -= 5
    c = 1
    return a + b + c

a = 9
b = 8
c = 6

try:
    print(g(b), a, b, c, sep=';')
except:
    print('ERR', a, b, c, sep=';')
code

3. Що буде виведено за виконання?
code
try:
    j = 7
    while j>=7:
        j += -2
    print(j, end=';')
except:
    print('ERR')
code

4. Що буде виведено за виконання?
code
try:
    for c in range(-7, 4, -1):
        if c>=4:
            continue
        print(c, end=';');c = 0
    else:
        print(c,end=';')
    print(c)
except:
    print('ERR')
code

5. Що буде виведено за виконання?\par (Дійсні значення, якщо наявні, в цьому завданні будуть виводитись у форматі з фіксованою крапкою та, можливо, з доволі великою кількістю десяткових розрядів. У відповіді для дійсних значень у ЦЬОМУ завданні достатньо вказати один десятковий розряд після крапки, відкинувши решту розрядів без округлення. Наприклад, замість 13.66666666666667 достатньо вказати 13.6, замість 25.23 достатньо вказати 25.2)
code
try:
    print(6, end=';')
    print(int('4'), end=';')
    print(1, end=';')
except KeyboardInterrupt:
    print(3, end=';')
except ZeroDivisionError:
    print(0, end=';')
except:
    print('ERR', end=';')
else:
    print(2, end=';')
finally:
    print(9)
code

6. Що буде виведено за виконання?
code
def f(n):
    if n>9:
        f(n//100)
    else:
        print(n%10,end='')

f(123456)
code

7. Що буде виведено за виконання?
code
class A:
    d = 1

    def _f1(self): print(2, end=';')
    def __f2(self): print(3, end=';')

    def main(self):
        try: print(self.d, end=';')
        except: print('ERR', end=';')

        try: self._f1()
        except: print('ERR', end=';')

        try: self.__f2()
        except: print('ERR', end=';')


class B(A):
    d = 4

    def _f1(self): print(5, end=';')
    def __f2(self): print(6, end=';')

try: B().main()
except: print('ERR', end=';')
code

8. Що буде виведено за виконання?
code
def f(arg, n):
    arg[40] = 0
    n = 9
    arg = list(arg.values())

try:
    n = 4
    t = {21:4, 40:7, 40:6}
    f(t, n)
    print(n, end=';')
    for x in t.items():
        print(*x, sep=';', end=';')
except:
    print('ERR')
code

9. Що буде виведено за виконання?
code
class A:
    a = 1
    c = 2

    def __init__(self):
        a = 3

class B(A):
    a = 4

    def __init__(self):
        super().__init__()
        self.b = 7

obj = B()
obj.b = obj.c + 10
B.b = 13

try: print(obj.b, end= ';')
except: print('ERR', end=';')

try: print(A.b, end= ';')
except: print('ERR', end=';')

try: print(B.b, end= ';')
except: print('ERR', end=';')
code

10. 
Для графа на рисунку з вершини 3 запущено алгоритм пошуку в глибину, що будує кістякове дерево. Списки суміжності вершин переглядаються алгоритмом за зростанням міток вершин.\par
\begin{center}
	\adjustimage{max size={0.9\linewidth}{0.9\paperheight}}
	{../10/Traversals/3.png}
\end{center}
\par Які з наведених ребер входять до складу отриманого кістякового дерева?\par
1) (4, 5)\par
2) (2, 3)\par
3) (0, 3)\par
4) (3, 6)\par
5) (1, 5)\par
6) (1, 6)\par
{\vskip 15pt} У формі вибрати номери відповідних ребер.\par
%answer:1 3 5 6
\end {large}}
%end
{\smallskip \begin {small} Затверджено на засіданні кафедри теоретичної кібернетики, протокол \No 12 від   19.05.2020 р.\par\smallskip\smallskip
{\parbox {6.5 cm}{Зав. кафедри \dotfill Ю.В.~Крак}}\hspace {0.7cm}{\hbox to 6.5 cm{ Екзаменатор \dotfill Т.О.~Карнаух}} \end{small}}
%begin
{{\begin {small}\par
{ \center  Київський національний університет імені Тараса Шевченка \par}
{\parbox {5 cm}{Освітня програма \hfill}{\parbox {7 cm}{Прикладна математика, Системний аналіз \hfill}}\parbox {6 cm}{\hfill Семестр 2}}\par
{\parbox {5 cm} {Навчальний предмет \hfill}\parbox {7 cm}{Програмування \hfill}\parbox {6cm}{\hfill}\par}\vspace{-7pt} \end{small}}{\center Екзаменаційний білет \No \textbf 
{ 335 }\par}}
{
\begin {large}
1. Що буде виведено за виконання?
code
def f(n):
    print('f', end='')
    return n<=2

try:
    print(f(-2) or f(-8))
except:
    print('ERR')
code

2. Що буде виведено за виконання?
code
a = 7
b = 1
c = 5

def f(a):
    a = 2
    b = 5
    c = 3
    return a + b + c

a = 8
b = 4
c = 0

try:
    print(f(a), a, b, c, sep=';')
except:
    print('ERR', a, b, c, sep=';')
code

3. Що буде виведено за виконання?
code
try:
    i = 11
    while i>9:
        i += -3
    print(i, end=';')
except:
    print('ERR')
code

4. Що буде виведено за виконання?
code
try:
    for f in range(-6, 7, -2):
        if f<=6:
            continue
        print(f, end=';')
    else:
        print(f,end=';')
    print(f)
except:
    print('ERR')
code

5. Що буде виведено за виконання?\par (Дійсні значення, якщо наявні, в цьому завданні будуть виводитись у форматі з фіксованою крапкою та, можливо, з доволі великою кількістю десяткових розрядів. У відповіді для дійсних значень у ЦЬОМУ завданні достатньо вказати один десятковий розряд після крапки, відкинувши решту розрядів без округлення. Наприклад, замість 13.66666666666667 достатньо вказати 13.6, замість 25.23 достатньо вказати 25.2)
code
try:
    print(7, end=';')
    print(8j<2j, end=';')
    print(5, end=';')
except Exception:
    print(0, end=';')
except TypeError:
    print(1, end=';')
except:
    print('ERR', end=';')
else:
    print(7, end=';')
finally:
    print(6)
code

6. Що буде виведено за виконання?
code
def f(s1):
    s = s1[:]
    for i in range(len(s)):
        s[i] += 1
 
s = [1, 2, 3]
f(s)
print(s)
code

7. Що буде виведено за виконання?
code
class A:
    c = 1

    def _g1(self): print(2, end=';')
    def g2(self): print(3, end=';')

    def main(self):
        try: print(self.c, end=';')
        except: print('ERR', end=';')

        try: self._g1()
        except: print('ERR', end=';')

        try: self.g2()
        except: print('ERR', end=';')


class B(A):
    c = 4

    def _g1(self): print(5, end=';')
    def g2(self): print(6, end=';')

try: B().main()
except: print('ERR', end=';')
code

8. Що буде виведено за виконання?
code
def f(arg, n):
    arg[74] = 1
    n -= -2

try:
    n = 5
    t = {39:2, 74:1, 57:7, 57:1}
    f(t, n)
    print(n, end=';')
    for x in t.items():
        print(*x, sep=';', end=';')
except:
    print('ERR')
code

9. Що буде виведено за виконання?
code
class A:
    b = 1
    c = 2

    def __init__(self):
        c = 3

class B(A):
    b = 4

    def __init__(self):
        super().__init__()
        self.a = 7

obj = B()
obj.a = obj.c + 10
B.a = 13

try: print(obj.a, end= ';')
except: print('ERR', end=';')

try: print(A.a, end= ';')
except: print('ERR', end=';')

try: print(B.a, end= ';')
except: print('ERR', end=';')
code

10. 
Для графа на рисунку з вершини 3 запущено алгоритм пошуку в глибину, що будує кістякове дерево. Списки суміжності вершин переглядаються алгоритмом за зростанням міток вершин.\par
\begin{center}
	\adjustimage{max size={0.9\linewidth}{0.9\paperheight}}
	{../10/Traversals/61.png}
\end{center}
\par Які з наведених ребер входять до складу отриманого кістякового дерева?\par
1) (2, 3)\par
2) (1, 4)\par
3) (2, 5)\par
4) (0, 5)\par
5) (3, 4)\par
6) (3, 6)\par
{\vskip 15pt} У формі вибрати номери відповідних ребер.\par
%answer:1 4
\end {large}}
%end
{\smallskip \begin {small} Затверджено на засіданні кафедри теоретичної кібернетики, протокол \No 12 від   19.05.2020 р.\par\smallskip\smallskip
{\parbox {6.5 cm}{Зав. кафедри \dotfill Ю.В.~Крак}}\hspace {0.7cm}{\hbox to 6.5 cm{ Екзаменатор \dotfill Т.О.~Карнаух}} \end{small}}
%begin
{{\begin {small}\par
{ \center  Київський національний університет імені Тараса Шевченка \par}
{\parbox {5 cm}{Освітня програма \hfill}{\parbox {7 cm}{Прикладна математика, Системний аналіз \hfill}}\parbox {6 cm}{\hfill Семестр 2}}\par
{\parbox {5 cm} {Навчальний предмет \hfill}\parbox {7 cm}{Програмування \hfill}\parbox {6cm}{\hfill}\par}\vspace{-7pt} \end{small}}{\center Екзаменаційний білет \No \textbf 
{ 336 }\par}}
{
\begin {large}
1. Що буде виведено за виконання?
code
def f(n):
    print('f', end='')
    return n==-3

try:
    print(f(5) and f(-7))
except:
    print('ERR')
code

2. Що буде виведено за виконання?
code
a = 3
b = 1
c = 0

def h(a):
    global c
    a = 4
    b *= 3
    c = 3
    return a + b + c

a = 5
b = 7
c = 1

try:
    print(h(b), a, b, c, sep=';')
except:
    print('ERR', a, b, c, sep=';')
code

3. Що буде виведено за виконання?
code
try:
    i = 1
    while i<5:
        i += 1
    print(i, end=';')
except:
    print('ERR')
code

4. Що буде виведено за виконання?
code
try:
    for a in range(-1, -3, -1):
        if a<7:
            break
        print(a, end=';')
    else:
        print(a,end=';')
    print(a)
except:
    print('ERR')
code

5. Що буде виведено за виконання?\par (Дійсні значення, якщо наявні, в цьому завданні будуть виводитись у форматі з фіксованою крапкою та, можливо, з доволі великою кількістю десяткових розрядів. У відповіді для дійсних значень у ЦЬОМУ завданні достатньо вказати один десятковий розряд після крапки, відкинувши решту розрядів без округлення. Наприклад, замість 13.66666666666667 достатньо вказати 13.6, замість 25.23 достатньо вказати 25.2)
code
try:
    print(9, end=';')
    print(8j==6j, end=';')
    print(2, end=';')
except ValueError:
    print(3, end=';')
except ZeroDivisionError:
    print(4, end=';')
except:
    print('ERR', end=';')
else:
    print(2, end=';')
finally:
    print(4)
code

6. Що буде виведено за виконання?
code
def f(n):
    if n>9:
        f(n//10)
    else:
        print(n%10, end='')

f(543216)
code

7. Що буде виведено за виконання?
code
class A:
    a = 1

    def h1(self): print(2, end=';')
    def _h2(self): print(3, end=';')

    def main(self):
        try: print(self.a, end=';')
        except: print('ERR', end=';')

        try: self.h1()
        except: print('ERR', end=';')

        try: self._h2()
        except: print('ERR', end=';')


class B(A):
    a = 4

    def h1(self): print(5, end=';')
    def _h2(self): print(6, end=';')

try: B().main()
except: print('ERR', end=';')
code

8. Що буде виведено за виконання?
code
def f(arg, n):
    arg[86] = 4
    n = 6
    arg = tuple(arg.keys())

try:
    n = 0
    s = {43:4, 52:7, 52:5}
    f(s, n)
    print(n, end=';')
    for x in s.keys():
        print(x, sep=';', end=';')
except:
    print('ERR')
code

9. Що буде виведено за виконання?
code
class A:
    c = 1
    a = 2

    def __init__(self):
        c = 3

class B(A):
    a = 4

    def __init__(self):
        super().__init__()
        self.b = 7

obj = B()
obj.b = obj.a + 10
B.b = 13

try: print(obj.b, end= ';')
except: print('ERR', end=';')

try: print(A.b, end= ';')
except: print('ERR', end=';')

try: print(B.b, end= ';')
except: print('ERR', end=';')
code

10. 
Для графа на рисунку з вершини 4 запущено алгоритм пошуку в глибину, що будує кістякове дерево. Списки суміжності вершин переглядаються алгоритмом за зростанням міток вершин.\par
\begin{center}
	\adjustimage{max size={0.9\linewidth}{0.9\paperheight}}
	{../10/Traversals/297.png}
\end{center}
\par Які з наведених ребер входять до складу отриманого кістякового дерева?\par
1) (2, 6)\par
2) (4, 5)\par
3) (0, 6)\par
4) (1, 3)\par
5) (1, 4)\par
6) (0, 5)\par
{\vskip 15pt} У формі вибрати номери відповідних ребер.\par
%answer:1 3 4 5 6
\end {large}}
%end
{\smallskip \begin {small} Затверджено на засіданні кафедри теоретичної кібернетики, протокол \No 12 від   19.05.2020 р.\par\smallskip\smallskip
{\parbox {6.5 cm}{Зав. кафедри \dotfill Ю.В.~Крак}}\hspace {0.7cm}{\hbox to 6.5 cm{ Екзаменатор \dotfill Т.О.~Карнаух}} \end{small}}
%begin
{{\begin {small}\par
{ \center  Київський національний університет імені Тараса Шевченка \par}
{\parbox {5 cm}{Освітня програма \hfill}{\parbox {7 cm}{Прикладна математика, Системний аналіз \hfill}}\parbox {6 cm}{\hfill Семестр 2}}\par
{\parbox {5 cm} {Навчальний предмет \hfill}\parbox {7 cm}{Програмування \hfill}\parbox {6cm}{\hfill}\par}\vspace{-7pt} \end{small}}{\center Екзаменаційний білет \No \textbf 
{ 337 }\par}}
{
\begin {large}
1. Що буде виведено за виконання?
code
try:
    res = 6//4*5*6
    if isinstance(res, bool):
        print(f'bool;{res}')
    elif isinstance(res, int):
        print(f'int;{res}')
    elif isinstance(res, float):
        print(f'float;{res:.2f}')
    else:
        print('???')
except:
    print('ERR')
code

2. Що буде виведено за виконання?
code
a = 6
b = 9
c = 4

def f(b):
    a = 5
    b += 1
    c = 4
    return a + b + c

a = 3
b = 0
c = 2

try:
    print(f(b), a, b, c, sep=';')
except:
    print('ERR', a, b, c, sep=';')
code

3. Що буде виведено за виконання?
code
try:
    m = 7
    while m<=8:
        m += 1
    print(m, end=';')
except:
    print('ERR')
code

4. Що буде виведено за виконання?
code
try:
    for b in range(14, 9, -2):
        if b<4:
            continue
        print(b, end=';')
    else:
        print('end',end=';')
    print(b)
except:
    print('ERR')
code

5. Що буде виведено за виконання?\par (Дійсні значення, якщо наявні, в цьому завданні будуть виводитись у форматі з фіксованою крапкою та, можливо, з доволі великою кількістю десяткових розрядів. У відповіді для дійсних значень у ЦЬОМУ завданні достатньо вказати один десятковий розряд після крапки, відкинувши решту розрядів без округлення. Наприклад, замість 13.66666666666667 достатньо вказати 13.6, замість 25.23 достатньо вказати 25.2)
code
try:
    print(8, end=';')
    print(1/0, end=';')
    print(9, end=';')
except KeyboardInterrupt:
    print(6, end=';')
except Exception:
    print(0, end=';')
except:
    print('ERR', end=';')
else:
    print(7, end=';')
finally:
    print(5)
code

6. Що буде виведено за виконання?
code
def f(n):
    if n<=9:
        print(n%10, end='')
    else:
        f(n//10)

f(123456)
code

7. Що буде виведено за виконання?
code
class A:
    b = 1

    def _h1(self): print(2, end=';')
    def _h2(self): print(3, end=';')

    def main(self):
        try: print(self.b, end=';')
        except: print('ERR', end=';')

        try: self._h1()
        except: print('ERR', end=';')

        try: self._h2()
        except: print('ERR', end=';')


class B(A):
    b = 4

    def _h1(self): print(5, end=';')
    def _h2(self): print(6, end=';')

try: B().main()
except: print('ERR', end=';')
code

8. Що буде виведено за виконання?
code
def f(arg, n):
    n = 4
    tmp = arg
    tmp[-1:] = 'ab'
    return len(tmp)>3

try:
    f = [10,11,12,13,14]
    n = 6
    f(f, n)
    print(n, end=';')
    for el in f:
        print(el, end=';')
except:
    print('ERR')
code

9. Що буде виведено за виконання?
code
class A:
    c = 1
    a = 2

    def __init__(self):
        a = 3

class B(A):
    c = 4

    def __init__(self):
        super().__init__()
        self.b = 7

obj = B()
obj.c = obj.a + 10
B.c = 13

try: print(obj.c, end= ';')
except: print('ERR', end=';')

try: print(A.c, end= ';')
except: print('ERR', end=';')

try: print(B.c, end= ';')
except: print('ERR', end=';')
code

10. 
Для графа на рисунку з вершини 2 запущено алгоритм пошуку в глибину, що будує кістякове дерево. Списки суміжності вершин переглядаються алгоритмом за зростанням міток вершин.\par
\begin{center}
	\adjustimage{max size={0.9\linewidth}{0.9\paperheight}}
	{../10/Traversals/293.png}
\end{center}
\par Які з наведених ребер входять до складу отриманого кістякового дерева?\par
1) (1, 3)\par
2) (3, 5)\par
3) (4, 5)\par
4) (0, 1)\par
5) (0, 5)\par
6) (2, 6)\par
{\vskip 15pt} У формі вибрати номери відповідних ребер.\par
%answer:2 3 4 6
\end {large}}
%end
{\smallskip \begin {small} Затверджено на засіданні кафедри теоретичної кібернетики, протокол \No 12 від   19.05.2020 р.\par\smallskip\smallskip
{\parbox {6.5 cm}{Зав. кафедри \dotfill Ю.В.~Крак}}\hspace {0.7cm}{\hbox to 6.5 cm{ Екзаменатор \dotfill Т.О.~Карнаух}} \end{small}}
%begin
{{\begin {small}\par
{ \center  Київський національний університет імені Тараса Шевченка \par}
{\parbox {5 cm}{Освітня програма \hfill}{\parbox {7 cm}{Прикладна математика, Системний аналіз \hfill}}\parbox {6 cm}{\hfill Семестр 2}}\par
{\parbox {5 cm} {Навчальний предмет \hfill}\parbox {7 cm}{Програмування \hfill}\parbox {6cm}{\hfill}\par}\vspace{-7pt} \end{small}}{\center Екзаменаційний білет \No \textbf 
{ 338 }\par}}
{
\begin {large}
1. Що буде виведено за виконання?
code
try:
    res = 4//8%7*3
    if isinstance(res, bool):
        print(f'bool;{res}')
    elif isinstance(res, int):
        print(f'int;{res}')
    elif isinstance(res, float):
        print(f'float;{res:.2f}')
    else:
        print('???')
except:
    print('ERR')
code

2. Що буде виведено за виконання?
code
a = 5
b = 8
c = 8

def f(b):
    global c
    a *= 2
    b = 2
    c = 4
    return a + b + c

a = 3
b = 2
c = 6

try:
    print(f(b), a, b, c, sep=';')
except:
    print('ERR', a, b, c, sep=';')
code

3. Що буде виведено за виконання?
code
try:
    t = 4
    while t<7:
        t += 2
    print(t, end=';')
except:
    print('ERR')
code

4. Що буде виведено за виконання?
code
try:
    for c in range(12, 0, -3):
        if c>5:
            continue
        print(c, end=';')
    else:
        print(c,end=';')
    print(c)
except:
    print('ERR')
code

5. Що буде виведено за виконання?\par (Дійсні значення, якщо наявні, в цьому завданні будуть виводитись у форматі з фіксованою крапкою та, можливо, з доволі великою кількістю десяткових розрядів. У відповіді для дійсних значень у ЦЬОМУ завданні достатньо вказати один десятковий розряд після крапки, відкинувши решту розрядів без округлення. Наприклад, замість 13.66666666666667 достатньо вказати 13.6, замість 25.23 достатньо вказати 25.2)
code
try:
    print(3, end=';')
    print(int('6'), end=';')
    print(7, end=';')
except TypeError:
    print(5, end=';')
except TypeError:
    print(1, end=';')
except:
    print('ERR', end=';')
else:
    print(0, end=';')
finally:
    print(2)
code

6. Що буде виведено за виконання?
code
def f(s1):
    s = s1[:]
    for i in range(len(s)):
        s[i] += 1
 
s = [1, 2, 3]
f(s)
print(s)
code

7. Що буде виведено за виконання?
code
class A:
    d = 1

    def __f1(self): print(2, end=';')
    def _f2(self): print(3, end=';')

    def main(self):
        try: print(self.d, end=';')
        except: print('ERR', end=';')

        try: self.__f1()
        except: print('ERR', end=';')

        try: self._f2()
        except: print('ERR', end=';')


class B(A):
    d = 4

    def __f1(self): print(5, end=';')
    def _f2(self): print(6, end=';')

try: B().main()
except: print('ERR', end=';')
code

8. Що буде виведено за виконання?
code
def f(arg, n):
    arg[77] = 8
    n = 5
    arg = set(arg.items())

try:
    n = 1
    d = {50:8, 70:8, 70:7}
    f(d, n)
    print(n, end=';')
    for x in d.values():
        print(x, sep=';', end=';')
except:
    print('ERR')
code

9. Що буде виведено за виконання?
code
class A:
    a = 1
    b = 2

    def __init__(self):
        a = 3

class B(A):
    b = 4

    def __init__(self):
        super().__init__()
        self.c = 7

obj = B()
obj.b = obj.a + 10
B.b = 13

try: print(obj.b, end= ';')
except: print('ERR', end=';')

try: print(A.b, end= ';')
except: print('ERR', end=';')

try: print(B.b, end= ';')
except: print('ERR', end=';')
code

10. 
Для графа на рисунку з вершини 6 запущено алгоритм пошуку в глибину, що будує кістякове дерево. Списки суміжності вершин переглядаються алгоритмом за зростанням міток вершин.\par
\begin{center}
	\adjustimage{max size={0.9\linewidth}{0.9\paperheight}}
	{../10/Traversals/253.png}
\end{center}
\par Які з наведених ребер входять до складу отриманого кістякового дерева?\par
1) (1, 3)\par
2) (3, 5)\par
3) (3, 6)\par
4) (0, 3)\par
5) (1, 2)\par
6) (2, 6)\par
{\vskip 15pt} У формі вибрати номери відповідних ребер.\par
%answer:1 2 4 5 6
\end {large}}
%end
{\smallskip \begin {small} Затверджено на засіданні кафедри теоретичної кібернетики, протокол \No 12 від   19.05.2020 р.\par\smallskip\smallskip
{\parbox {6.5 cm}{Зав. кафедри \dotfill Ю.В.~Крак}}\hspace {0.7cm}{\hbox to 6.5 cm{ Екзаменатор \dotfill Т.О.~Карнаух}} \end{small}}
%begin
{{\begin {small}\par
{ \center  Київський національний університет імені Тараса Шевченка \par}
{\parbox {5 cm}{Освітня програма \hfill}{\parbox {7 cm}{Прикладна математика, Системний аналіз \hfill}}\parbox {6 cm}{\hfill Семестр 2}}\par
{\parbox {5 cm} {Навчальний предмет \hfill}\parbox {7 cm}{Програмування \hfill}\parbox {6cm}{\hfill}\par}\vspace{-7pt} \end{small}}{\center Екзаменаційний білет \No \textbf 
{ 339 }\par}}
{
\begin {large}
1. Що буде виведено за виконання?
code
def f(n):
    print('f', end='')
    return n

try:
    print(f(6) or f(3))
except:
    print('ERR')
code

2. Що буде виведено за виконання?
code
a = 0
b = 5
c = 4

def g(a):
    a = 1
    b += 3
    c = 5
    return a + b + c

a = 1
b = 9
c = 7

try:
    print(g(b), a, b, c, sep=';')
except:
    print('ERR', a, b, c, sep=';')
code

3. Що буде виведено за виконання?
code
try:
    n = 8
    while n>8:
        n += -2
    print(n, end=';')
except:
    print('ERR')
code

4. Що буде виведено за виконання?
code
try:
    for f in range(0, 14, 2):
        if f>=8:
            break
        print(f, end=';');f = -6
        if f>=5:
            break
    else:
        print(f,end=';')
    print(f)
except:
    print('ERR')
code

5. Що буде виведено за виконання?\par (Дійсні значення, якщо наявні, в цьому завданні будуть виводитись у форматі з фіксованою крапкою та, можливо, з доволі великою кількістю десяткових розрядів. У відповіді для дійсних значень у ЦЬОМУ завданні достатньо вказати один десятковий розряд після крапки, відкинувши решту розрядів без округлення. Наприклад, замість 13.66666666666667 достатньо вказати 13.6, замість 25.23 достатньо вказати 25.2)
code
try:
    print(4, end=';')
    print(int('d9'), end=';')
    print(8, end=';')
except KeyboardInterrupt:
    print(3, end=';')
except ZeroDivisionError:
    print(0, end=';')
except:
    print('ERR', end=';')
else:
    print(1, end=';')
finally:
    print(3)
code

6. Що буде виведено за виконання?
code
def f(n):
    if n>9:
        f(n//10)
    print(n%10, end='')
 
f(543216)
code

7. Що буде виведено за виконання?
code
class A:
    a = 1

    def _g1(self): print(2, end=';')
    def g2(self): print(3, end=';')

    def main(self):
        try: print(self.a, end=';')
        except: print('ERR', end=';')

        try: self._g1()
        except: print('ERR', end=';')

        try: self.g2()
        except: print('ERR', end=';')


class B(A):
    a = 4

    def _g1(self): print(5, end=';')
    def g2(self): print(6, end=';')

try: B().main()
except: print('ERR', end=';')
code

8. Що буде виведено за виконання?
code
def f(arg, n):
    n = 7
    tmp = arg
    tmp[7:5] = [13]
    return len(tmp)>3

try:
    a = [10,11,12,13,14]
    n = 2
    f(a, n)
    print(n, end=';')
    for el in a:
        print(el, end=';')
except:
    print('ERR')
code

9. Що буде виведено за виконання?
code
class A:
    b = 1
    a = 2

    def __init__(self):
        b = 3

class B(A):
    a = 4

    def __init__(self):
        super().__init__()
        self.c = 7

obj = B()
obj.c = obj.b + 10
B.c = 13

try: print(obj.c, end= ';')
except: print('ERR', end=';')

try: print(A.c, end= ';')
except: print('ERR', end=';')

try: print(B.c, end= ';')
except: print('ERR', end=';')
code

10. 
Для графа на рисунку з вершини 2 запущено алгоритм пошуку в глибину, що будує кістякове дерево. Списки суміжності вершин переглядаються алгоритмом за зростанням міток вершин.\par
\begin{center}
	\adjustimage{max size={0.9\linewidth}{0.9\paperheight}}
	{../10/Traversals/116.png}
\end{center}
\par Які з наведених ребер входять до складу отриманого кістякового дерева?\par
1) (2, 4)\par
2) (2, 6)\par
3) (5, 6)\par
4) (3, 6)\par
5) (0, 2)\par
6) (0, 6)\par
{\vskip 15pt} У формі вибрати номери відповідних ребер.\par
%answer:3 4 5
\end {large}}
%end
{\smallskip \begin {small} Затверджено на засіданні кафедри теоретичної кібернетики, протокол \No 12 від   19.05.2020 р.\par\smallskip\smallskip
{\parbox {6.5 cm}{Зав. кафедри \dotfill Ю.В.~Крак}}\hspace {0.7cm}{\hbox to 6.5 cm{ Екзаменатор \dotfill Т.О.~Карнаух}} \end{small}}
%begin
{{\begin {small}\par
{ \center  Київський національний університет імені Тараса Шевченка \par}
{\parbox {5 cm}{Освітня програма \hfill}{\parbox {7 cm}{Прикладна математика, Системний аналіз \hfill}}\parbox {6 cm}{\hfill Семестр 2}}\par
{\parbox {5 cm} {Навчальний предмет \hfill}\parbox {7 cm}{Програмування \hfill}\parbox {6cm}{\hfill}\par}\vspace{-7pt} \end{small}}{\center Екзаменаційний білет \No \textbf 
{ 340 }\par}}
{
\begin {large}
1. Що буде виведено за виконання?
code
try:
    res = 7/7//9*8
    if isinstance(res, bool):
        print(f'bool;{res}')
    elif isinstance(res, int):
        print(f'int;{res}')
    elif isinstance(res, float):
        print(f'float;{res:.2f}')
    else:
        print('???')
except:
    print('ERR')
code

2. Що буде виведено за виконання?
code
a = 6
b = 1
c = 7

def h(b):
    global c
    a = 1
    b = 4
    c = 2
    return a + b + c

a = 9
b = 2
c = 3

try:
    print(h(a), a, b, c, sep=';')
except:
    print('ERR', a, b, c, sep=';')
code

3. Що буде виведено за виконання?
code
try:
    i = 6
    while i>=3:
        i += -3
    print(i, end=';')
except:
    print('ERR')
code

4. Що буде виведено за виконання?
code
try:
    for e in range(-10, -5, 3):
        if e<=7:
            continue
        print(e, end=';')
    else:
        print(e,end=';')
    print(e)
except:
    print('ERR')
code

5. Що буде виведено за виконання?\par (Дійсні значення, якщо наявні, в цьому завданні будуть виводитись у форматі з фіксованою крапкою та, можливо, з доволі великою кількістю десяткових розрядів. У відповіді для дійсних значень у ЦЬОМУ завданні достатньо вказати один десятковий розряд після крапки, відкинувши решту розрядів без округлення. Наприклад, замість 13.66666666666667 достатньо вказати 13.6, замість 25.23 достатньо вказати 25.2)
code
try:
    print(6, end=';')
    print(9%2, end=';')
    print(8, end=';')
except ValueError:
    print(7, end=';')
except Exception:
    print(2, end=';')
except:
    print('ERR', end=';')
else:
    print(5, end=';')
finally:
    print(4)
code

6. Що буде виведено за виконання?
code
def f(s):
    for i in range(len(s)):
        s[i] += 1
        return
 
s = [1, 2, 3]
f(s)
print(s)
code

7. Що буде виведено за виконання?
code
class A:
    c = 1

    def _f1(self): print(2, end=';')
    def __f2(self): print(3, end=';')

    def main(self):
        try: print(self.c, end=';')
        except: print('ERR', end=';')

        try: self._f1()
        except: print('ERR', end=';')

        try: self.__f2()
        except: print('ERR', end=';')


class B(A):
    c = 4

    def _f1(self): print(5, end=';')
    def __f2(self): print(6, end=';')

try: B().main()
except: print('ERR', end=';')
code

8. Що буде виведено за виконання?
code
def f(arg, n):
    n = 1
    tmp = arg
    del tmp[:2]
    return len(tmp)>3

try:
    b = ['0','1','2','3','4']
    n = 3
    f(b, n)
    print(n, end=';')
    for el in b:
        print(el, end=';')
except:
    print('ERR')
code

9. Що буде виведено за виконання?
code
class A:
    a = 1
    c = 2

    def __init__(self):
        c = 3

class B(A):
    a = 4

    def __init__(self):
        super().__init__()
        self.b = 7

obj = B()
obj.c = obj.a + 10
B.c = 13

try: print(obj.c, end= ';')
except: print('ERR', end=';')

try: print(A.c, end= ';')
except: print('ERR', end=';')

try: print(B.c, end= ';')
except: print('ERR', end=';')
code

10. 
Для графа на рисунку з вершини 5 запущено алгоритм пошуку в глибину, що будує кістякове дерево. Списки суміжності вершин переглядаються алгоритмом за зростанням міток вершин.\par
\begin{center}
	\adjustimage{max size={0.9\linewidth}{0.9\paperheight}}
	{../10/Traversals/21.png}
\end{center}
\par Які з наведених ребер входять до складу отриманого кістякового дерева?\par
1) (1, 5)\par
2) (4, 6)\par
3) (3, 4)\par
4) (1, 3)\par
5) (5, 6)\par
6) (0, 3)\par
{\vskip 15pt} У формі вибрати номери відповідних ребер.\par
%answer:2 4 6
\end {large}}
%end
{\smallskip \begin {small} Затверджено на засіданні кафедри теоретичної кібернетики, протокол \No 12 від   19.05.2020 р.\par\smallskip\smallskip
{\parbox {6.5 cm}{Зав. кафедри \dotfill Ю.В.~Крак}}\hspace {0.7cm}{\hbox to 6.5 cm{ Екзаменатор \dotfill Т.О.~Карнаух}} \end{small}}
%begin
{{\begin {small}\par
{ \center  Київський національний університет імені Тараса Шевченка \par}
{\parbox {5 cm}{Освітня програма \hfill}{\parbox {7 cm}{Прикладна математика, Системний аналіз \hfill}}\parbox {6 cm}{\hfill Семестр 2}}\par
{\parbox {5 cm} {Навчальний предмет \hfill}\parbox {7 cm}{Програмування \hfill}\parbox {6cm}{\hfill}\par}\vspace{-7pt} \end{small}}{\center Екзаменаційний білет \No \textbf 
{ 341 }\par}}
{
\begin {large}
1. Що буде виведено за виконання?
code
try:
    res = 2>7 or not True<=8.0 or 6>=2.0
    if isinstance(res, bool):
        print(f'bool;{res}')
    elif isinstance(res, int):
        print(f'int;{res}')
    elif isinstance(res, float):
        print(f'float;{res:.2f}')
    else:
        print('???')
except:
    print('ERR')
code

2. Що буде виведено за виконання?
code
a = 0
b = 1
c = 5

def h(b):
    global c
    a = 5
    b -= 4
    c = 2
    return a + b + c

a = 9
b = 6
c = 3

try:
    print(h(a), a, b, c, sep=';')
except:
    print('ERR', a, b, c, sep=';')
code

3. Що буде виведено за виконання?
code
try:
    m = 15
    while m>=8:
        m += -2
    print(m, end=';')
except:
    print('ERR')
code

4. Що буде виведено за виконання?
code
try:
    for d in range(15, -6, 3):
        if d<=3:
            continue
        print(d, end=';')
    else:
        print(d,end=';')
    print(d)
except:
    print('ERR')
code

5. Що буде виведено за виконання?\par (Дійсні значення, якщо наявні, в цьому завданні будуть виводитись у форматі з фіксованою крапкою та, можливо, з доволі великою кількістю десяткових розрядів. У відповіді для дійсних значень у ЦЬОМУ завданні достатньо вказати один десятковий розряд після крапки, відкинувши решту розрядів без округлення. Наприклад, замість 13.66666666666667 достатньо вказати 13.6, замість 25.23 достатньо вказати 25.2)
code
try:
    print(5, end=';')
    print(2j>=5j, end=';')
    print(3, end=';')
except Exception:
    print(0, end=';')
except ZeroDivisionError:
    print(1, end=';')
except:
    print('ERR', end=';')
else:
    print(9, end=';')
finally:
    print(4)
code

6. Що буде виведено за виконання?
code
def f(n):
    if n<=9:
        print(n%10, end='')
    else:
        f(n//10)

f(987651)
code

7. Що буде виведено за виконання?
code
class A:
    c = 1

    def __g1(self): print(2, end=';')
    def _g2(self): print(3, end=';')

    def main(self):
        try: print(self.c, end=';')
        except: print('ERR', end=';')

        try: self.__g1()
        except: print('ERR', end=';')

        try: self._g2()
        except: print('ERR', end=';')


class B(A):
    c = 4

    def __g1(self): print(5, end=';')
    def _g2(self): print(6, end=';')

try: B().main()
except: print('ERR', end=';')
code

8. Що буде виведено за виконання?
code
def f(arg, n):
    arg[42] = 5
    n += 2

try:
    n = 3
    s = {81:9, 42:7, 73:9}
    f(s, n)
    print(n, end=';')
    for x in s.keys():
        print(x, sep=';', end=';')
except:
    print('ERR')
code

9. Що буде виведено за виконання?
code
class A:
    c = 1
    b = 2

    def __init__(self):
        b = 3

class B(A):
    c = 4

    def __init__(self):
        super().__init__()
        self.a = 7

obj = B()
obj.b = obj.a + 10
B.b = 13

try: print(obj.b, end= ';')
except: print('ERR', end=';')

try: print(A.b, end= ';')
except: print('ERR', end=';')

try: print(B.b, end= ';')
except: print('ERR', end=';')
code

10. 
Для графа на рисунку з вершини 4 запущено алгоритм пошуку в глибину, що будує кістякове дерево. Списки суміжності вершин переглядаються алгоритмом за зростанням міток вершин.\par
\begin{center}
	\adjustimage{max size={0.9\linewidth}{0.9\paperheight}}
	{../10/Traversals/28.png}
\end{center}
\par Які з наведених ребер входять до складу отриманого кістякового дерева?\par
1) (3, 5)\par
2) (1, 2)\par
3) (0, 1)\par
4) (2, 3)\par
5) (3, 4)\par
6) (2, 4)\par
{\vskip 15pt} У формі вибрати номери відповідних ребер.\par
%answer:1 2 3 6
\end {large}}
%end
{\smallskip \begin {small} Затверджено на засіданні кафедри теоретичної кібернетики, протокол \No 12 від   19.05.2020 р.\par\smallskip\smallskip
{\parbox {6.5 cm}{Зав. кафедри \dotfill Ю.В.~Крак}}\hspace {0.7cm}{\hbox to 6.5 cm{ Екзаменатор \dotfill Т.О.~Карнаух}} \end{small}}
%begin
{{\begin {small}\par
{ \center  Київський національний університет імені Тараса Шевченка \par}
{\parbox {5 cm}{Освітня програма \hfill}{\parbox {7 cm}{Прикладна математика, Системний аналіз \hfill}}\parbox {6 cm}{\hfill Семестр 2}}\par
{\parbox {5 cm} {Навчальний предмет \hfill}\parbox {7 cm}{Програмування \hfill}\parbox {6cm}{\hfill}\par}\vspace{-7pt} \end{small}}{\center Екзаменаційний білет \No \textbf 
{ 342 }\par}}
{
\begin {large}
1. Що буде виведено за виконання?
code
try:
    res = 9/6%8*5
    if isinstance(res, bool):
        print(f'bool;{res}')
    elif isinstance(res, int):
        print(f'int;{res}')
    elif isinstance(res, float):
        print(f'float;{res:.2f}')
    else:
        print('???')
except:
    print('ERR')
code

2. Що буде виведено за виконання?
code
a = 2
b = 7
c = 4

def f(a):
    a *= 3
    b = 1
    c = 4
    return a + b + c

a = 8
b = 6
c = 1

try:
    print(f(a), a, b, c, sep=';')
except:
    print('ERR', a, b, c, sep=';')
code

3. Що буде виведено за виконання?
code
try:
    t = 9
    while t>9:
        t += -3
    print(t, end=';')
except:
    print('ERR')
code

4. Що буде виведено за виконання?
code
try:
    for a in range(7, 6, 2):
        if a>6:
            break
        print(a, end=';');a = -7
    else:
        print(a,end=';')
    print(a)
except:
    print('ERR')
code

5. Що буде виведено за виконання?\par (Дійсні значення, якщо наявні, в цьому завданні будуть виводитись у форматі з фіксованою крапкою та, можливо, з доволі великою кількістю десяткових розрядів. У відповіді для дійсних значень у ЦЬОМУ завданні достатньо вказати один десятковий розряд після крапки, відкинувши решту розрядів без округлення. Наприклад, замість 13.66666666666667 достатньо вказати 13.6, замість 25.23 достатньо вказати 25.2)
code
try:
    print(8, end=';')
    print(int('6'), end=';')
    print(7, end=';')
except ValueError:
    print(8, end=';')
except TypeError:
    print(9, end=';')
except:
    print('ERR', end=';')
else:
    print(6, end=';')
finally:
    print(1)
code

6. Що буде виведено за виконання?
code
def f(n):
    if n>9:
        f(n//100)
    print(n%10, end='')
 
f(123456)
code

7. Що буде виведено за виконання?
code
class A:
    b = 1

    def h1(self): print(2, end=';')
    def _h2(self): print(3, end=';')

    def main(self):
        try: print(self.b, end=';')
        except: print('ERR', end=';')

        try: self.h1()
        except: print('ERR', end=';')

        try: self._h2()
        except: print('ERR', end=';')


class B(A):
    b = 4

    def h1(self): print(5, end=';')
    def _h2(self): print(6, end=';')

try: B().main()
except: print('ERR', end=';')
code

8. Що буде виведено за виконання?
code
def f(arg, n):
    n = 0
    tmp = arg
    del tmp[-8:8]
    return len(tmp)>3

try:
    d = ['0','1','2','3','4']
    n = 5
    f(d, n)
    print(n, end=';')
    for el in d:
        print(el, end=';')
except:
    print('ERR')
code

9. Що буде виведено за виконання?
code
class A:
    b = 1
    c = 2

    def __init__(self):
        b = 3

class B(A):
    c = 4

    def __init__(self):
        super().__init__()
        self.a = 7

obj = B()
obj.c = obj.b + 10
B.c = 13

try: print(obj.c, end= ';')
except: print('ERR', end=';')

try: print(A.c, end= ';')
except: print('ERR', end=';')

try: print(B.c, end= ';')
except: print('ERR', end=';')
code

10. 
Для графа на рисунку з вершини 1 запущено алгоритм пошуку в глибину, що будує кістякове дерево. Списки суміжності вершин переглядаються алгоритмом за зростанням міток вершин.\par
\begin{center}
	\adjustimage{max size={0.9\linewidth}{0.9\paperheight}}
	{../10/Traversals/234.png}
\end{center}
\par Які з наведених ребер входять до складу отриманого кістякового дерева?\par
1) (1, 6)\par
2) (1, 2)\par
3) (1, 4)\par
4) (1, 5)\par
5) (2, 4)\par
6) (4, 5)\par
{\vskip 15pt} У формі вибрати номери відповідних ребер.\par
%answer:2 5 6
\end {large}}
%end
{\smallskip \begin {small} Затверджено на засіданні кафедри теоретичної кібернетики, протокол \No 12 від   19.05.2020 р.\par\smallskip\smallskip
{\parbox {6.5 cm}{Зав. кафедри \dotfill Ю.В.~Крак}}\hspace {0.7cm}{\hbox to 6.5 cm{ Екзаменатор \dotfill Т.О.~Карнаух}} \end{small}}
%begin
{{\begin {small}\par
{ \center  Київський національний університет імені Тараса Шевченка \par}
{\parbox {5 cm}{Освітня програма \hfill}{\parbox {7 cm}{Прикладна математика, Системний аналіз \hfill}}\parbox {6 cm}{\hfill Семестр 2}}\par
{\parbox {5 cm} {Навчальний предмет \hfill}\parbox {7 cm}{Програмування \hfill}\parbox {6cm}{\hfill}\par}\vspace{-7pt} \end{small}}{\center Екзаменаційний білет \No \textbf 
{ 343 }\par}}
{
\begin {large}
1. Що буде виведено за виконання?
code
try:
    res = True<="False"
    if isinstance(res, bool):
        print(f'bool;{res}')
    elif isinstance(res, int):
        print(f'int;{res}')
    elif isinstance(res, float):
        print(f'float;{res:.2f}')
    else:
        print('???')
except:
    print('ERR')
code

2. Що буде виведено за виконання?
code
a = 2
b = 1
c = 4

def g(b):
    global c
    a = 4
    b = 3
    c = 1
    return a + b + c

a = 9
b = 3
c = 0

try:
    print(g(b), a, b, c, sep=';')
except:
    print('ERR', a, b, c, sep=';')
code

3. Що буде виведено за виконання?
code
try:
    j = 12
    while j>4:
        j += -3
    print(j, end=';')
except:
    print('ERR')
code

4. Що буде виведено за виконання?
code
try:
    for d in range(2, 12, -3):
        if d>=4:
            continue
        print(d, end=';')
    else:
        print('end',end=';')
    print(d)
except:
    print('ERR')
code

5. Що буде виведено за виконання?\par (Дійсні значення, якщо наявні, в цьому завданні будуть виводитись у форматі з фіксованою крапкою та, можливо, з доволі великою кількістю десяткових розрядів. У відповіді для дійсних значень у ЦЬОМУ завданні достатньо вказати один десятковий розряд після крапки, відкинувши решту розрядів без округлення. Наприклад, замість 13.66666666666667 достатньо вказати 13.6, замість 25.23 достатньо вказати 25.2)
code
try:
    print(4, end=';')
    print(2/0, end=';')
    print(5, end=';')
except KeyboardInterrupt:
    print(0, end=';')
except KeyboardInterrupt:
    print(3, end=';')
except:
    print('ERR', end=';')
else:
    print(7, end=';')
finally:
    print(4)
code

6. Що буде виведено за виконання?
code
def f(n):
    if n<=9:
        print(n%10,end='')
    else:
        f(n//10)

f(543216)
code

7. Що буде виведено за виконання?
code
class A:
    a = 1

    def _f1(self): print(2, end=';')
    def _f2(self): print(3, end=';')

    def main(self):
        try: print(self.a, end=';')
        except: print('ERR', end=';')

        try: self._f1()
        except: print('ERR', end=';')

        try: self._f2()
        except: print('ERR', end=';')


class B(A):
    a = 4

    def _f1(self): print(5, end=';')
    def _f2(self): print(6, end=';')

try: B().main()
except: print('ERR', end=';')
code

8. Що буде виведено за виконання?
code
def f(arg, n):
    arg[50] = 3
    n = 8
    arg = list(arg.values())

try:
    n = 2
    t = {50:5, 20:6, 20:4}
    f(t, n)
    print(n, end=';')
    for x, y in t.items():
        print(x, y, sep=';', end=';')
except:
    print('ERR')
code

9. Що буде виведено за виконання?
code
class A:
    c = 1
    b = 2

    def __init__(self):
        c = 3

class B(A):
    c = 4

    def __init__(self):
        super().__init__()
        self.a = 7

obj = B()
obj.b = obj.a + 10
B.b = 13

try: print(obj.b, end= ';')
except: print('ERR', end=';')

try: print(A.b, end= ';')
except: print('ERR', end=';')

try: print(B.b, end= ';')
except: print('ERR', end=';')
code

10. 
Для графа на рисунку з вершини 3 запущено алгоритм пошуку в глибину, що будує кістякове дерево. Списки суміжності вершин переглядаються алгоритмом за зростанням міток вершин.\par
\begin{center}
	\adjustimage{max size={0.9\linewidth}{0.9\paperheight}}
	{../10/Traversals/231.png}
\end{center}
\par Які з наведених ребер входять до складу отриманого кістякового дерева?\par
1) (1, 4)\par
2) (3, 4)\par
3) (4, 5)\par
4) (1, 5)\par
5) (0, 6)\par
6) (4, 6)\par
{\vskip 15pt} У формі вибрати номери відповідних ребер.\par
%answer:2 4
\end {large}}
%end
{\smallskip \begin {small} Затверджено на засіданні кафедри теоретичної кібернетики, протокол \No 12 від   19.05.2020 р.\par\smallskip\smallskip
{\parbox {6.5 cm}{Зав. кафедри \dotfill Ю.В.~Крак}}\hspace {0.7cm}{\hbox to 6.5 cm{ Екзаменатор \dotfill Т.О.~Карнаух}} \end{small}}
%begin
{{\begin {small}\par
{ \center  Київський національний університет імені Тараса Шевченка \par}
{\parbox {5 cm}{Освітня програма \hfill}{\parbox {7 cm}{Прикладна математика, Системний аналіз \hfill}}\parbox {6 cm}{\hfill Семестр 2}}\par
{\parbox {5 cm} {Навчальний предмет \hfill}\parbox {7 cm}{Програмування \hfill}\parbox {6cm}{\hfill}\par}\vspace{-7pt} \end{small}}{\center Екзаменаційний білет \No \textbf 
{ 344 }\par}}
{
\begin {large}
1. Що буде виведено за виконання?
code
try:
    res = 4**1.0*-1
    if isinstance(res, bool):
        print(f'bool;{res}')
    elif isinstance(res, int):
        print(f'int;{res}')
    elif isinstance(res, float):
        print(f'float;{res:.2f}')
    else:
        print('???')
except:
    print('ERR')
code

2. Що буде виведено за виконання?
code
a = 5
b = 4
c = 9

def g(a):
    global c
    a += 3
    b = 1
    c = 5
    return a + b + c

a = 0
b = 7
c = 3

try:
    print(g(b), a, b, c, sep=';')
except:
    print('ERR', a, b, c, sep=';')
code

3. Що буде виведено за виконання?
code
try:
    k = 14
    while k>=2:
        k += -2
    print(k, end=';')
except:
    print('ERR')
code

4. Що буде виведено за виконання?
code
try:
    for b in range(-11, -6, -2):
        if b<6:
            break
        print(b, end=';');b = -10
    else:
        print(b,end=';')
    print(b)
except:
    print('ERR')
code

5. Що буде виведено за виконання?\par (Дійсні значення, якщо наявні, в цьому завданні будуть виводитись у форматі з фіксованою крапкою та, можливо, з доволі великою кількістю десяткових розрядів. У відповіді для дійсних значень у ЦЬОМУ завданні достатньо вказати один десятковий розряд після крапки, відкинувши решту розрядів без округлення. Наприклад, замість 13.66666666666667 достатньо вказати 13.6, замість 25.23 достатньо вказати 25.2)
code
try:
    print(8, end=';')
    print(5//3, end=';')
    print(3, end=';')
except TypeError:
    print(0, end=';')
except Exception:
    print(9, end=';')
except:
    print('ERR', end=';')
else:
    print(2, end=';')
finally:
    print(7)
code

6. Що буде виведено за виконання?
code
def f(n):
    if n>9: f(n//10)
    else:
        print(n%10, end='')
 
f(123456)
code

7. Що буде виведено за виконання?
code
class A:
    d = 1

    def _h1(self): print(2, end=';')
    def __h2(self): print(3, end=';')

    def main(self):
        try: print(self.d, end=';')
        except: print('ERR', end=';')

        try: self._h1()
        except: print('ERR', end=';')

        try: self.__h2()
        except: print('ERR', end=';')


class B(A):
    d = 4

    def _h1(self): print(5, end=';')
    def __h2(self): print(6, end=';')

try: B().main()
except: print('ERR', end=';')
code

8. Що буде виведено за виконання?
code
def f(arg, n):
    n = 3
    tmp = arg
    tmp[1:-7] = (12)
    return len(tmp)>3

try:
    e = [10,11,12,13,14]
    n = 1
    f(e, n)
    print(n, end=';')
    for el in e:
        print(el, end=';')
except:
    print('ERR')
code

9. Що буде виведено за виконання?
code
class A:
    b = 1
    c = 2

    def __init__(self):
        c = 3

class B(A):
    c = 4

    def __init__(self):
        super().__init__()
        self.a = 7

obj = B()
obj.c = obj.a + 10
B.c = 13

try: print(obj.c, end= ';')
except: print('ERR', end=';')

try: print(A.c, end= ';')
except: print('ERR', end=';')

try: print(B.c, end= ';')
except: print('ERR', end=';')
code

10. 
Для графа на рисунку з вершини 2 запущено алгоритм пошуку в глибину, що будує кістякове дерево. Списки суміжності вершин переглядаються алгоритмом за зростанням міток вершин.\par
\begin{center}
	\adjustimage{max size={0.9\linewidth}{0.9\paperheight}}
	{../10/Traversals/190.png}
\end{center}
\par Які з наведених ребер входять до складу отриманого кістякового дерева?\par
1) (5, 6)\par
2) (0, 3)\par
3) (1, 2)\par
4) (4, 5)\par
5) (3, 4)\par
6) (1, 5)\par
{\vskip 15pt} У формі вибрати номери відповідних ребер.\par
%answer:1 2 3 5
\end {large}}
%end
{\smallskip \begin {small} Затверджено на засіданні кафедри теоретичної кібернетики, протокол \No 12 від   19.05.2020 р.\par\smallskip\smallskip
{\parbox {6.5 cm}{Зав. кафедри \dotfill Ю.В.~Крак}}\hspace {0.7cm}{\hbox to 6.5 cm{ Екзаменатор \dotfill Т.О.~Карнаух}} \end{small}}
%begin
{{\begin {small}\par
{ \center  Київський національний університет імені Тараса Шевченка \par}
{\parbox {5 cm}{Освітня програма \hfill}{\parbox {7 cm}{Прикладна математика, Системний аналіз \hfill}}\parbox {6 cm}{\hfill Семестр 2}}\par
{\parbox {5 cm} {Навчальний предмет \hfill}\parbox {7 cm}{Програмування \hfill}\parbox {6cm}{\hfill}\par}\vspace{-7pt} \end{small}}{\center Екзаменаційний білет \No \textbf 
{ 345 }\par}}
{
\begin {large}
1. Що буде виведено за виконання?
code
def f(n):
    print('f', end='')
    return n

try:
    print(f(7) or f(8))
except:
    print('ERR')
code

2. Що буде виведено за виконання?
code
a = 8
b = 2
c = 7

def h(b):
    a -= 2
    b = 1
    c = 4
    return a + b + c

a = 9
b = 2
c = 4

try:
    print(h(a), a, b, c, sep=';')
except:
    print('ERR', a, b, c, sep=';')
code

3. Що буде виведено за виконання?
code
try:
    k = 4
    while k<11:
        k += 2
    print(k, end=';')
except:
    print('ERR')
code

4. Що буде виведено за виконання?
code
try:
    for e in range(12, 13, -1):
        if e<7:
            continue
        print(e, end=';')
    else:
        print(e,end=';')
    print(e)
except:
    print('ERR')
code

5. Що буде виведено за виконання?\par (Дійсні значення, якщо наявні, в цьому завданні будуть виводитись у форматі з фіксованою крапкою та, можливо, з доволі великою кількістю десяткових розрядів. У відповіді для дійсних значень у ЦЬОМУ завданні достатньо вказати один десятковий розряд після крапки, відкинувши решту розрядів без округлення. Наприклад, замість 13.66666666666667 достатньо вказати 13.6, замість 25.23 достатньо вказати 25.2)
code
try:
    print(1, end=';')
    print(int('c6'), end=';')
    print(4, end=';')
except ZeroDivisionError:
    print(2, end=';')
except ValueError:
    print(8, end=';')
except:
    print('ERR', end=';')
else:
    print(3, end=';')
finally:
    print(9)
code

6. Що буде виведено за виконання?
code
def f(s):
    s1 = s
    for i in range(len(s)):
        s1[i] += 1
    return s1

s = [1, 2, 3]
f(s)
print(s)
code

7. Що буде виведено за виконання?
code
class A:
    b = 1

    def _g1(self): print(2, end=';')
    def g2(self): print(3, end=';')

    def main(self):
        try: print(self.b, end=';')
        except: print('ERR', end=';')

        try: self._g1()
        except: print('ERR', end=';')

        try: self.g2()
        except: print('ERR', end=';')


class B(A):
    b = 4

    def _g1(self): print(5, end=';')
    def g2(self): print(6, end=';')

try: B().main()
except: print('ERR', end=';')
code

8. Що буде виведено за виконання?
code
def f(arg, n):
    arg[37] = 9
    n += 3

try:
    n = 4
    t = {69:1, 60:5, 36:8, 23:1}
    f(t, n)
    print(n, end=';')
    for x in t.values():
        print(x, sep=';', end=';')
except:
    print('ERR')
code

9. Що буде виведено за виконання?
code
class A:
    a = 1
    c = 2

    def __init__(self):
        c = 3

class B(A):
    c = 4

    def __init__(self):
        super().__init__()
        self.b = 7

obj = B()
obj.c = obj.b + 10
B.c = 13

try: print(obj.c, end= ';')
except: print('ERR', end=';')

try: print(A.c, end= ';')
except: print('ERR', end=';')

try: print(B.c, end= ';')
except: print('ERR', end=';')
code

10. 
Для графа на рисунку з вершини 0 запущено алгоритм пошуку в глибину, що будує кістякове дерево. Списки суміжності вершин переглядаються алгоритмом за зростанням міток вершин.\par
\begin{center}
	\adjustimage{max size={0.9\linewidth}{0.9\paperheight}}
	{../10/Traversals/146.png}
\end{center}
\par Які з наведених ребер входять до складу отриманого кістякового дерева?\par
1) (1, 5)\par
2) (3, 5)\par
3) (0, 1)\par
4) (2, 6)\par
5) (5, 6)\par
6) (3, 4)\par
{\vskip 15pt} У формі вибрати номери відповідних ребер.\par
%answer:2 3 4 5 6
\end {large}}
%end
{\smallskip \begin {small} Затверджено на засіданні кафедри теоретичної кібернетики, протокол \No 12 від   19.05.2020 р.\par\smallskip\smallskip
{\parbox {6.5 cm}{Зав. кафедри \dotfill Ю.В.~Крак}}\hspace {0.7cm}{\hbox to 6.5 cm{ Екзаменатор \dotfill Т.О.~Карнаух}} \end{small}}
%begin
{{\begin {small}\par
{ \center  Київський національний університет імені Тараса Шевченка \par}
{\parbox {5 cm}{Освітня програма \hfill}{\parbox {7 cm}{Прикладна математика, Системний аналіз \hfill}}\parbox {6 cm}{\hfill Семестр 2}}\par
{\parbox {5 cm} {Навчальний предмет \hfill}\parbox {7 cm}{Програмування \hfill}\parbox {6cm}{\hfill}\par}\vspace{-7pt} \end{small}}{\center Екзаменаційний білет \No \textbf 
{ 346 }\par}}
{
\begin {large}
1. Що буде виведено за виконання?
code
try:
    res = 5//9*3*4
    if isinstance(res, bool):
        print(f'bool;{res}')
    elif isinstance(res, int):
        print(f'int;{res}')
    elif isinstance(res, float):
        print(f'float;{res:.2f}')
    else:
        print('???')
except:
    print('ERR')
code

2. Що буде виведено за виконання?
code
a = 7
b = 8
c = 5

def h(a):
    a = 5
    b = 2
    c = 4
    return a + b + c

a = 6
b = 9
c = 5

try:
    print(h(b), a, b, c, sep=';')
except:
    print('ERR', a, b, c, sep=';')
code

3. Що буде виведено за виконання?
code
try:
    i = 10
    while i>1:
        i += -3
    print(i, end=';')
except:
    print('ERR')
code

4. Що буде виведено за виконання?
code
try:
    for c in range(14, -8, -1):
        if c>3:
            break
        print(c, end=';')
    else:
        print(c,end=';')
    print(c)
except:
    print('ERR')
code

5. Що буде виведено за виконання?\par (Дійсні значення, якщо наявні, в цьому завданні будуть виводитись у форматі з фіксованою крапкою та, можливо, з доволі великою кількістю десяткових розрядів. У відповіді для дійсних значень у ЦЬОМУ завданні достатньо вказати один десятковий розряд після крапки, відкинувши решту розрядів без округлення. Наприклад, замість 13.66666666666667 достатньо вказати 13.6, замість 25.23 достатньо вказати 25.2)
code
try:
    print(5, end=';')
    print(int('d0'), end=';')
    print(6, end=';')
except TypeError:
    print(7, end=';')
except KeyboardInterrupt:
    print(1, end=';')
except:
    print('ERR', end=';')
else:
    print(9, end=';')
finally:
    print(4)
code

6. Що буде виведено за виконання?
code
def f(n):
    if n>9:
        f(n//100)
    else:
        print(n%10,end='')

f(123456)
code

7. Що буде виведено за виконання?
code
class A:
    d = 1

    def _f1(self): print(2, end=';')
    def _f2(self): print(3, end=';')

    def main(self):
        try: print(self.d, end=';')
        except: print('ERR', end=';')

        try: self._f1()
        except: print('ERR', end=';')

        try: self._f2()
        except: print('ERR', end=';')


class B(A):
    d = 4

    def _f1(self): print(5, end=';')
    def _f2(self): print(6, end=';')

try: B().main()
except: print('ERR', end=';')
code

8. Що буде виведено за виконання?
code
def f(arg, n):
    arg[51] = 9
    n *= -3

try:
    n = 4
    d = {70:0, 90:3, 43:1, 43:4}
    f(d, n)
    print(n, end=';')
    for x in d :
        print(x, sep=';', end=';')
except:
    print('ERR')
code

9. Що буде виведено за виконання?
code
class A:
    c = 1
    a = 2

    def __init__(self):
        c = 3

class B(A):
    c = 4

    def __init__(self):
        super().__init__()
        self.b = 7

obj = B()
obj.a = obj.c + 10
B.a = 13

try: print(obj.a, end= ';')
except: print('ERR', end=';')

try: print(A.a, end= ';')
except: print('ERR', end=';')

try: print(B.a, end= ';')
except: print('ERR', end=';')
code

10. 
Для графа на рисунку з вершини 5 запущено алгоритм пошуку в глибину, що будує кістякове дерево. Списки суміжності вершин переглядаються алгоритмом за зростанням міток вершин.\par
\begin{center}
	\adjustimage{max size={0.9\linewidth}{0.9\paperheight}}
	{../10/Traversals/248.png}
\end{center}
\par Які з наведених ребер входять до складу отриманого кістякового дерева?\par
1) (1, 4)\par
2) (3, 6)\par
3) (0, 4)\par
4) (2, 5)\par
5) (0, 1)\par
6) (3, 5)\par
{\vskip 15pt} У формі вибрати номери відповідних ребер.\par
%answer:2 5
\end {large}}
%end
{\smallskip \begin {small} Затверджено на засіданні кафедри теоретичної кібернетики, протокол \No 12 від   19.05.2020 р.\par\smallskip\smallskip
{\parbox {6.5 cm}{Зав. кафедри \dotfill Ю.В.~Крак}}\hspace {0.7cm}{\hbox to 6.5 cm{ Екзаменатор \dotfill Т.О.~Карнаух}} \end{small}}
%begin
{{\begin {small}\par
{ \center  Київський національний університет імені Тараса Шевченка \par}
{\parbox {5 cm}{Освітня програма \hfill}{\parbox {7 cm}{Прикладна математика, Системний аналіз \hfill}}\parbox {6 cm}{\hfill Семестр 2}}\par
{\parbox {5 cm} {Навчальний предмет \hfill}\parbox {7 cm}{Програмування \hfill}\parbox {6cm}{\hfill}\par}\vspace{-7pt} \end{small}}{\center Екзаменаційний білет \No \textbf 
{ 347 }\par}}
{
\begin {large}
1. Що буде виведено за виконання?
code
try:
    res = not 3>=False and 9.0<6.0 and 5!=8
    if isinstance(res, bool):
        print(f'bool;{res}')
    elif isinstance(res, int):
        print(f'int;{res}')
    elif isinstance(res, float):
        print(f'float;{res:.2f}')
    else:
        print('???')
except:
    print('ERR')
code

2. Що буде виведено за виконання?
code
a = 2
b = 0
c = 6

def f(a):
    global c
    a = 3
    b = 5
    c = 2
    return a + b + c

a = 7
b = 3
c = 1

try:
    print(f(a), a, b, c, sep=';')
except:
    print('ERR', a, b, c, sep=';')
code

3. Що буде виведено за виконання?
code
try:
    i = 1
    while i<=5:
        i += 1
    print(i, end=';')
except:
    print('ERR')
code

4. Що буде виведено за виконання?
code
try:
    for f in range(-6, -10, -2):
        if f<=5:
            continue
        print(f, end=';')
    else:
        print(f,end=';')
    print(f)
except:
    print('ERR')
code

5. Що буде виведено за виконання?\par (Дійсні значення, якщо наявні, в цьому завданні будуть виводитись у форматі з фіксованою крапкою та, можливо, з доволі великою кількістю десяткових розрядів. У відповіді для дійсних значень у ЦЬОМУ завданні достатньо вказати один десятковий розряд після крапки, відкинувши решту розрядів без округлення. Наприклад, замість 13.66666666666667 достатньо вказати 13.6, замість 25.23 достатньо вказати 25.2)
code
try:
    print(6, end=';')
    print(3/0.0, end=';')
    print(7, end=';')
except ZeroDivisionError:
    print(5, end=';')
except Exception:
    print(1, end=';')
except:
    print('ERR', end=';')
else:
    print(2, end=';')
finally:
    print(0)
code

6. Що буде виведено за виконання?
code
def f(n):
    if n>2:
        f(n//2)
    print(n%2, end='')
 
f(16)
code

7. Що буде виведено за виконання?
code
class A:
    c = 1

    def __g1(self): print(2, end=';')
    def _g2(self): print(3, end=';')

    def main(self):
        try: print(self.c, end=';')
        except: print('ERR', end=';')

        try: self.__g1()
        except: print('ERR', end=';')

        try: self._g2()
        except: print('ERR', end=';')


class B(A):
    c = 4

    def __g1(self): print(5, end=';')
    def _g2(self): print(6, end=';')

try: B().main()
except: print('ERR', end=';')
code

8. Що буде виведено за виконання?
code
def f(arg, n):
    arg[71] = 5
    n -= -2

try:
    n = 5
    s = {87:3, 90:1, 90:8}
    f(s, n)
    print(n, end=';')
    for x in s.values():
        print(x, sep=';', end=';')
except:
    print('ERR')
code

9. Що буде виведено за виконання?
code
class A:
    b = 1
    a = 2

    def __init__(self):
        a = 3

class B(A):
    b = 4

    def __init__(self):
        super().__init__()
        self.c = 7

obj = B()
obj.b = obj.c + 10
B.b = 13

try: print(obj.b, end= ';')
except: print('ERR', end=';')

try: print(A.b, end= ';')
except: print('ERR', end=';')

try: print(B.b, end= ';')
except: print('ERR', end=';')
code

10. 
Для графа на рисунку з вершини 2 запущено алгоритм пошуку в глибину, що будує кістякове дерево. Списки суміжності вершин переглядаються алгоритмом за зростанням міток вершин.\par
\begin{center}
	\adjustimage{max size={0.9\linewidth}{0.9\paperheight}}
	{../10/Traversals/244.png}
\end{center}
\par Які з наведених ребер входять до складу отриманого кістякового дерева?\par
1) (5, 6)\par
2) (3, 6)\par
3) (4, 5)\par
4) (0, 2)\par
5) (0, 5)\par
6) (1, 2)\par
{\vskip 15pt} У формі вибрати номери відповідних ребер.\par
%answer:4 5
\end {large}}
%end
{\smallskip \begin {small} Затверджено на засіданні кафедри теоретичної кібернетики, протокол \No 12 від   19.05.2020 р.\par\smallskip\smallskip
{\parbox {6.5 cm}{Зав. кафедри \dotfill Ю.В.~Крак}}\hspace {0.7cm}{\hbox to 6.5 cm{ Екзаменатор \dotfill Т.О.~Карнаух}} \end{small}}
%begin
{{\begin {small}\par
{ \center  Київський національний університет імені Тараса Шевченка \par}
{\parbox {5 cm}{Освітня програма \hfill}{\parbox {7 cm}{Прикладна математика, Системний аналіз \hfill}}\parbox {6 cm}{\hfill Семестр 2}}\par
{\parbox {5 cm} {Навчальний предмет \hfill}\parbox {7 cm}{Програмування \hfill}\parbox {6cm}{\hfill}\par}\vspace{-7pt} \end{small}}{\center Екзаменаційний білет \No \textbf 
{ 348 }\par}}
{
\begin {large}
1. Що буде виведено за виконання?
code
try:
    res = 9**-0.5*1
    if isinstance(res, bool):
        print(f'bool;{res}')
    elif isinstance(res, int):
        print(f'int;{res}')
    elif isinstance(res, float):
        print(f'float;{res:.2f}')
    else:
        print('???')
except:
    print('ERR')
code

2. Що буде виведено за виконання?
code
a = 8
b = 1
c = 5

def h(a):
    global c
    a -= 5
    b = 3
    c = 2
    return a + b + c

a = 0
b = 6
c = 3

try:
    print(h(a), a, b, c, sep=';')
except:
    print('ERR', a, b, c, sep=';')
code

3. Що буде виведено за виконання?
code
try:
    n = 3
    while n<6:
        n += 1
    print(n, end=';')
except:
    print('ERR')
code

4. Що буде виведено за виконання?
code
try:
    for a in range(-4, 1, 2):
        if a>=8:
            continue
        print(a, end=';');a = 6
        if a>6:
            break
    else:
        print(a,end=';')
    print(a)
except:
    print('ERR')
code

5. Що буде виведено за виконання?\par (Дійсні значення, якщо наявні, в цьому завданні будуть виводитись у форматі з фіксованою крапкою та, можливо, з доволі великою кількістю десяткових розрядів. У відповіді для дійсних значень у ЦЬОМУ завданні достатньо вказати один десятковий розряд після крапки, відкинувши решту розрядів без округлення. Наприклад, замість 13.66666666666667 достатньо вказати 13.6, замість 25.23 достатньо вказати 25.2)
code
try:
    print(8, end=';')
    print(int('7'), end=';')
    print(6, end=';')
except ValueError:
    print(1, end=';')
except ValueError:
    print(3, end=';')
except:
    print('ERR', end=';')
else:
    print(9, end=';')
finally:
    print(2)
code

6. Що буде виведено за виконання?
code
def f(s):
    for el in s:
        el += 1
        return s
 
s = [1, 2, 3]
s = f(s)
print(s)
code

7. Що буде виведено за виконання?
code
class A:
    a = 1

    def _h1(self): print(2, end=';')
    def __h2(self): print(3, end=';')

    def main(self):
        try: print(self.a, end=';')
        except: print('ERR', end=';')

        try: self._h1()
        except: print('ERR', end=';')

        try: self.__h2()
        except: print('ERR', end=';')


class B(A):
    a = 4

    def _h1(self): print(5, end=';')
    def __h2(self): print(6, end=';')

try: B().main()
except: print('ERR', end=';')
code

8. Що буде виведено за виконання?
code
def f(arg, n):
    arg[33] = 3
    n = 7
    arg = tuple(arg.keys())

try:
    n = 3
    s = {12:4, 25:2, 12:4}
    f(s, n)
    print(n, end=';')
    for x in s.keys():
        print(x, sep=';', end=';')
except:
    print('ERR')
code

9. Що буде виведено за виконання?
code
class A:
    a = 1
    b = 2

    def __init__(self):
        a = 3

class B(A):
    b = 4

    def __init__(self):
        super().__init__()
        self.c = 7

obj = B()
obj.b = obj.a + 10
B.b = 13

try: print(obj.b, end= ';')
except: print('ERR', end=';')

try: print(A.b, end= ';')
except: print('ERR', end=';')

try: print(B.b, end= ';')
except: print('ERR', end=';')
code

10. 
Для графа на рисунку з вершини 0 запущено алгоритм пошуку в глибину, що будує кістякове дерево. Списки суміжності вершин переглядаються алгоритмом за зростанням міток вершин.\par
\begin{center}
	\adjustimage{max size={0.9\linewidth}{0.9\paperheight}}
	{../10/Traversals/130.png}
\end{center}
\par Які з наведених ребер входять до складу отриманого кістякового дерева?\par
1) (1, 5)\par
2) (1, 2)\par
3) (1, 4)\par
4) (3, 4)\par
5) (0, 4)\par
6) (2, 5)\par
{\vskip 15pt} У формі вибрати номери відповідних ребер.\par
%answer:2 3 4 6
\end {large}}
%end
{\smallskip \begin {small} Затверджено на засіданні кафедри теоретичної кібернетики, протокол \No 12 від   19.05.2020 р.\par\smallskip\smallskip
{\parbox {6.5 cm}{Зав. кафедри \dotfill Ю.В.~Крак}}\hspace {0.7cm}{\hbox to 6.5 cm{ Екзаменатор \dotfill Т.О.~Карнаух}} \end{small}}
%begin
{{\begin {small}\par
{ \center  Київський національний університет імені Тараса Шевченка \par}
{\parbox {5 cm}{Освітня програма \hfill}{\parbox {7 cm}{Прикладна математика, Системний аналіз \hfill}}\parbox {6 cm}{\hfill Семестр 2}}\par
{\parbox {5 cm} {Навчальний предмет \hfill}\parbox {7 cm}{Програмування \hfill}\parbox {6cm}{\hfill}\par}\vspace{-7pt} \end{small}}{\center Екзаменаційний білет \No \textbf 
{ 349 }\par}}
{
\begin {large}
1. Що буде виведено за виконання?
code
try:
    res = not 5<=7.0 or 3==True or 9>5.0
    if isinstance(res, bool):
        print(f'bool;{res}')
    elif isinstance(res, int):
        print(f'int;{res}')
    elif isinstance(res, float):
        print(f'float;{res:.2f}')
    else:
        print('???')
except:
    print('ERR')
code

2. Що буде виведено за виконання?
code
a = 8
b = 4
c = 8

def g(b):
    a = 1
    b = 4
    c = 3
    return a + b + c

a = 0
b = 7
c = 3

try:
    print(g(b), a, b, c, sep=';')
except:
    print('ERR', a, b, c, sep=';')
code

3. Що буде виведено за виконання?
code
try:
    t = 13
    while t>=7:
        t += -2
    print(t, end=';')
except:
    print('ERR')
code

4. Що буде виведено за виконання?
code
try:
    for d in range(11, -15, -3):
        if d<4:
            continue
        print(d, end=';');d = 5
    else:
        print(d,end=';')
    print(d)
except:
    print('ERR')
code

5. Що буде виведено за виконання?\par (Дійсні значення, якщо наявні, в цьому завданні будуть виводитись у форматі з фіксованою крапкою та, можливо, з доволі великою кількістю десяткових розрядів. У відповіді для дійсних значень у ЦЬОМУ завданні достатньо вказати один десятковий розряд після крапки, відкинувши решту розрядів без округлення. Наприклад, замість 13.66666666666667 достатньо вказати 13.6, замість 25.23 достатньо вказати 25.2)
code
try:
    print(8, end=';')
    print(5j>8j, end=';')
    print(0, end=';')
except ZeroDivisionError:
    print(4, end=';')
except KeyboardInterrupt:
    print(0, end=';')
except:
    print('ERR', end=';')
else:
    print(3, end=';')
finally:
    print(7)
code

6. Що буде виведено за виконання?
code
def f(n):
    if n>9:
        f(n//100)
    print(n%10,end='')
 
f(987654)
code

7. Що буде виведено за виконання?
code
class A:
    b = 1

    def _h1(self): print(2, end=';')
    def h2(self): print(3, end=';')

    def main(self):
        try: print(self.b, end=';')
        except: print('ERR', end=';')

        try: self._h1()
        except: print('ERR', end=';')

        try: self.h2()
        except: print('ERR', end=';')


class B(A):
    b = 4

    def _h1(self): print(5, end=';')
    def h2(self): print(6, end=';')

try: B().main()
except: print('ERR', end=';')
code

8. Що буде виведено за виконання?
code
def f(arg, n):
    n = 8
    tmp = arg
    tmp[-1:2] = 'ac'
    return len(tmp)>3

try:
    h = ['0','1','2','3','4']
    n = 2
    f(h, n)
    print(n, end=';')
    for el in h:
        print(el, end=';')
except:
    print('ERR')
code

9. Що буде виведено за виконання?
code
class A:
    b = 1
    a = 2

    def __init__(self):
        a = 3

class B(A):
    a = 4

    def __init__(self):
        super().__init__()
        self.c = 7

obj = B()
obj.b = obj.c + 10
B.b = 13

try: print(obj.b, end= ';')
except: print('ERR', end=';')

try: print(A.b, end= ';')
except: print('ERR', end=';')

try: print(B.b, end= ';')
except: print('ERR', end=';')
code

10. 
Для графа на рисунку з вершини 2 запущено алгоритм пошуку в глибину, що будує кістякове дерево. Списки суміжності вершин переглядаються алгоритмом за зростанням міток вершин.\par
\begin{center}
	\adjustimage{max size={0.9\linewidth}{0.9\paperheight}}
	{../10/Traversals/96.png}
\end{center}
\par Які з наведених ребер входять до складу отриманого кістякового дерева?\par
1) (1, 6)\par
2) (0, 5)\par
3) (1, 2)\par
4) (0, 1)\par
5) (2, 3)\par
6) (3, 5)\par
{\vskip 15pt} У формі вибрати номери відповідних ребер.\par
%answer:1 2 3 4 6
\end {large}}
%end
{\smallskip \begin {small} Затверджено на засіданні кафедри теоретичної кібернетики, протокол \No 12 від   19.05.2020 р.\par\smallskip\smallskip
{\parbox {6.5 cm}{Зав. кафедри \dotfill Ю.В.~Крак}}\hspace {0.7cm}{\hbox to 6.5 cm{ Екзаменатор \dotfill Т.О.~Карнаух}} \end{small}}
%begin
{{\begin {small}\par
{ \center  Київський національний університет імені Тараса Шевченка \par}
{\parbox {5 cm}{Освітня програма \hfill}{\parbox {7 cm}{Прикладна математика, Системний аналіз \hfill}}\parbox {6 cm}{\hfill Семестр 2}}\par
{\parbox {5 cm} {Навчальний предмет \hfill}\parbox {7 cm}{Програмування \hfill}\parbox {6cm}{\hfill}\par}\vspace{-7pt} \end{small}}{\center Екзаменаційний білет \No \textbf 
{ 350 }\par}}
{
\begin {large}
1. Що буде виведено за виконання?
code
try:
    res = 6j>="7"
    if isinstance(res, bool):
        print(f'bool;{res}')
    elif isinstance(res, int):
        print(f'int;{res}')
    elif isinstance(res, float):
        print(f'float;{res:.2f}')
    else:
        print('???')
except:
    print('ERR')
code

2. Що буде виведено за виконання?
code
a = 5
b = 1
c = 2

def g(a):
    a = 1
    b = 5
    c = 2
    return a + b + c

a = 9
b = 6
c = 4

try:
    print(g(b), a, b, c, sep=';')
except:
    print('ERR', a, b, c, sep=';')
code

3. Що буде виведено за виконання?
code
try:
    j = 5
    while j<=7:
        j += 2
    print(j, end=';')
except:
    print('ERR')
code

4. Що буде виведено за виконання?
code
try:
    for f in range(-15, 9, 3):
        if f<=8:
            continue
        print(f, end=';')
    else:
        print(f,end=';')
    print(f)
except:
    print('ERR')
code

5. Що буде виведено за виконання?\par (Дійсні значення, якщо наявні, в цьому завданні будуть виводитись у форматі з фіксованою крапкою та, можливо, з доволі великою кількістю десяткових розрядів. У відповіді для дійсних значень у ЦЬОМУ завданні достатньо вказати один десятковий розряд після крапки, відкинувши решту розрядів без округлення. Наприклад, замість 13.66666666666667 достатньо вказати 13.6, замість 25.23 достатньо вказати 25.2)
code
try:
    print(8, end=';')
    print(4%1, end=';')
    print(5, end=';')
except TypeError:
    print(2, end=';')
except Exception:
    print(6, end=';')
except:
    print('ERR', end=';')
else:
    print(9, end=';')
finally:
    print(1)
code

6. Що буде виведено за виконання?
code
def f(n):
    print(n%10, end='')
    if n>9:
        f(n//10)
 
f(543216)
code

7. Що буде виведено за виконання?
code
class A:
    c = 1

    def g1(self): print(2, end=';')
    def _g2(self): print(3, end=';')

    def main(self):
        try: print(self.c, end=';')
        except: print('ERR', end=';')

        try: self.g1()
        except: print('ERR', end=';')

        try: self._g2()
        except: print('ERR', end=';')


class B(A):
    c = 4

    def g1(self): print(5, end=';')
    def _g2(self): print(6, end=';')

try: B().main()
except: print('ERR', end=';')
code

8. Що буде виведено за виконання?
code
def f(arg, n):
    arg[18] = 4
    n = 7
    arg = list(arg.items())

try:
    n = 4
    t = {86:4, 88:2, 68:2, 71:4}
    f(t, n)
    print(n, end=';')
    for x in t.values():
        print(x, sep=';', end=';')
except:
    print('ERR')
code

9. Що буде виведено за виконання?
code
class A:
    c = 1
    b = 2

    def __init__(self):
        c = 3

class B(A):
    c = 4

    def __init__(self):
        super().__init__()
        self.a = 7

obj = B()
obj.a = obj.c + 10
B.a = 13

try: print(obj.a, end= ';')
except: print('ERR', end=';')

try: print(A.a, end= ';')
except: print('ERR', end=';')

try: print(B.a, end= ';')
except: print('ERR', end=';')
code

10. 
Для графа на рисунку з вершини 0 запущено алгоритм пошуку в глибину, що будує кістякове дерево. Списки суміжності вершин переглядаються алгоритмом за зростанням міток вершин.\par
\begin{center}
	\adjustimage{max size={0.9\linewidth}{0.9\paperheight}}
	{../10/Traversals/139.png}
\end{center}
\par Які з наведених ребер входять до складу отриманого кістякового дерева?\par
1) (4, 5)\par
2) (4, 6)\par
3) (1, 5)\par
4) (0, 4)\par
5) (1, 6)\par
6) (3, 4)\par
{\vskip 15pt} У формі вибрати номери відповідних ребер.\par
%answer:1 3 6
\end {large}}
%end
{\smallskip \begin {small} Затверджено на засіданні кафедри теоретичної кібернетики, протокол \No 12 від   19.05.2020 р.\par\smallskip\smallskip
{\parbox {6.5 cm}{Зав. кафедри \dotfill Ю.В.~Крак}}\hspace {0.7cm}{\hbox to 6.5 cm{ Екзаменатор \dotfill Т.О.~Карнаух}} \end{small}}
%begin
{{\begin {small}\par
{ \center  Київський національний університет імені Тараса Шевченка \par}
{\parbox {5 cm}{Освітня програма \hfill}{\parbox {7 cm}{Прикладна математика, Системний аналіз \hfill}}\parbox {6 cm}{\hfill Семестр 2}}\par
{\parbox {5 cm} {Навчальний предмет \hfill}\parbox {7 cm}{Програмування \hfill}\parbox {6cm}{\hfill}\par}\vspace{-7pt} \end{small}}{\center Екзаменаційний білет \No \textbf 
{ 351 }\par}}
{
\begin {large}
1. Що буде виведено за виконання?
code
def f(n):
    print('f', end='')
    return n>=-4

try:
    print(f(4) and f(0))
except:
    print('ERR')
code

2. Що буде виведено за виконання?
code
a = 4
b = 9
c = 8

def g(b):
    a = 3
    b *= 5
    c = 1
    return a + b + c

a = 7
b = 0
c = 2

try:
    print(g(a), a, b, c, sep=';')
except:
    print('ERR', a, b, c, sep=';')
code

3. Що буде виведено за виконання?
code
try:
    m = 7
    while m<11:
        m += 1
    print(m, end=';')
except:
    print('ERR')
code

4. Що буде виведено за виконання?
code
try:
    for a in range(9, -2, -2):
        if a>3:
            break
        print(a, end=';')
    else:
        print(a,end=';')
    print(a)
except:
    print('ERR')
code

5. Що буде виведено за виконання?\par (Дійсні значення, якщо наявні, в цьому завданні будуть виводитись у форматі з фіксованою крапкою та, можливо, з доволі великою кількістю десяткових розрядів. У відповіді для дійсних значень у ЦЬОМУ завданні достатньо вказати один десятковий розряд після крапки, відкинувши решту розрядів без округлення. Наприклад, замість 13.66666666666667 достатньо вказати 13.6, замість 25.23 достатньо вказати 25.2)
code
try:
    print(2, end=';')
    print(4//0, end=';')
    print(1, end=';')
except TypeError:
    print(9, end=';')
except Exception:
    print(6, end=';')
except:
    print('ERR', end=';')
else:
    print(0, end=';')
finally:
    print(3)
code

6. Що буде виведено за виконання?
code
def f(n):
    if n>9:
        return f(n//10)+1
    else:
        return 1

print(f(123456))
code

7. Що буде виведено за виконання?
code
class A:
    a = 1

    def _f1(self): print(2, end=';')
    def f2(self): print(3, end=';')

    def main(self):
        try: print(self.a, end=';')
        except: print('ERR', end=';')

        try: self._f1()
        except: print('ERR', end=';')

        try: self.f2()
        except: print('ERR', end=';')


class B(A):
    a = 4

    def _f1(self): print(5, end=';')
    def f2(self): print(6, end=';')

try: B().main()
except: print('ERR', end=';')
code

8. Що буде виведено за виконання?
code
def f(arg, n):
    arg[39] = 2
    n = 5
    arg = set(arg.values())

try:
    n = 3
    d = {69:9, 43:7, 29:1, 69:0}
    f(d, n)
    print(n, end=';')
    for x in d.items():
        print(x, sep=';', end=';')
except:
    print('ERR')
code

9. Що буде виведено за виконання?
code
class A:
    a = 1
    b = 2

    def __init__(self):
        b = 3

class B(A):
    b = 4

    def __init__(self):
        super().__init__()
        self.c = 7

obj = B()
obj.a = obj.b + 10
B.a = 13

try: print(obj.a, end= ';')
except: print('ERR', end=';')

try: print(A.a, end= ';')
except: print('ERR', end=';')

try: print(B.a, end= ';')
except: print('ERR', end=';')
code

10. 
Для графа на рисунку з вершини 6 запущено алгоритм пошуку в глибину, що будує кістякове дерево. Списки суміжності вершин переглядаються алгоритмом за зростанням міток вершин.\par
\begin{center}
	\adjustimage{max size={0.9\linewidth}{0.9\paperheight}}
	{../10/Traversals/164.png}
\end{center}
\par Які з наведених ребер входять до складу отриманого кістякового дерева?\par
1) (4, 5)\par
2) (0, 5)\par
3) (0, 4)\par
4) (1, 6)\par
5) (0, 2)\par
6) (2, 4)\par
{\vskip 15pt} У формі вибрати номери відповідних ребер.\par
%answer:2 4 5 6
\end {large}}
%end
{\smallskip \begin {small} Затверджено на засіданні кафедри теоретичної кібернетики, протокол \No 12 від   19.05.2020 р.\par\smallskip\smallskip
{\parbox {6.5 cm}{Зав. кафедри \dotfill Ю.В.~Крак}}\hspace {0.7cm}{\hbox to 6.5 cm{ Екзаменатор \dotfill Т.О.~Карнаух}} \end{small}}
%begin
{{\begin {small}\par
{ \center  Київський національний університет імені Тараса Шевченка \par}
{\parbox {5 cm}{Освітня програма \hfill}{\parbox {7 cm}{Прикладна математика, Системний аналіз \hfill}}\parbox {6 cm}{\hfill Семестр 2}}\par
{\parbox {5 cm} {Навчальний предмет \hfill}\parbox {7 cm}{Програмування \hfill}\parbox {6cm}{\hfill}\par}\vspace{-7pt} \end{small}}{\center Екзаменаційний білет \No \textbf 
{ 352 }\par}}
{
\begin {large}
1. Що буде виведено за виконання?
code
def f(n):
    print('f', end='')
    return n!=-2

try:
    print(f(1) and f(-4))
except:
    print('ERR')
code

2. Що буде виведено за виконання?
code
a = 6
b = 3
c = 0

def h(b):
    global c
    a = 5
    b = 4
    c = 3
    return a + b + c

a = 4
b = 1
c = 9

try:
    print(h(b), a, b, c, sep=';')
except:
    print('ERR', a, b, c, sep=';')
code

3. Що буде виведено за виконання?
code
try:
    j = 9
    while j<=9:
        j += 2
    print(j, end=';')
except:
    print('ERR')
code

4. Що буде виведено за виконання?
code
try:
    for b in range(-5, 10, 3):
        if b>=6:
            break
        print(b, end=';')
    else:
        print(b,end=';')
    print(b)
except:
    print('ERR')
code

5. Що буде виведено за виконання?\par (Дійсні значення, якщо наявні, в цьому завданні будуть виводитись у форматі з фіксованою крапкою та, можливо, з доволі великою кількістю десяткових розрядів. У відповіді для дійсних значень у ЦЬОМУ завданні достатньо вказати один десятковий розряд після крапки, відкинувши решту розрядів без округлення. Наприклад, замість 13.66666666666667 достатньо вказати 13.6, замість 25.23 достатньо вказати 25.2)
code
try:
    print(8, end=';')
    print(5j<=8j, end=';')
    print(7, end=';')
except ZeroDivisionError:
    print(3, end=';')
except ValueError:
    print(4, end=';')
except:
    print('ERR', end=';')
else:
    print(5, end=';')
finally:
    print(0)
code

6. Що буде виведено за виконання?
code
def f(n):
    if n>9:
        f(n//10)
    else:
        print(n%10, end='')

f(543216)
code

7. Що буде виведено за виконання?
code
class A:
    d = 1

    def _g1(self): print(2, end=';')
    def __g2(self): print(3, end=';')

    def main(self):
        try: print(self.d, end=';')
        except: print('ERR', end=';')

        try: self._g1()
        except: print('ERR', end=';')

        try: self.__g2()
        except: print('ERR', end=';')


class B(A):
    d = 4

    def _g1(self): print(5, end=';')
    def __g2(self): print(6, end=';')

try: B().main()
except: print('ERR', end=';')
code

8. Що буде виведено за виконання?
code
def f(arg, n):
    n = 6
    tmp = arg
    tmp[:-7] = [1,2]
    return len(tmp)>3

try:
    c = [10,11,12,13,14]
    n = 5
    f(c, n)
    print(n, end=';')
    for el in c:
        print(el, end=';')
except:
    print('ERR')
code

9. Що буде виведено за виконання?
code
class A:
    c = 1
    a = 2

    def __init__(self):
        c = 3

class B(A):
    c = 4

    def __init__(self):
        super().__init__()
        self.b = 7

obj = B()
obj.b = obj.a + 10
B.b = 13

try: print(obj.b, end= ';')
except: print('ERR', end=';')

try: print(A.b, end= ';')
except: print('ERR', end=';')

try: print(B.b, end= ';')
except: print('ERR', end=';')
code

10. 
Для графа на рисунку з вершини 4 запущено алгоритм пошуку в глибину, що будує кістякове дерево. Списки суміжності вершин переглядаються алгоритмом за зростанням міток вершин.\par
\begin{center}
	\adjustimage{max size={0.9\linewidth}{0.9\paperheight}}
	{../10/Traversals/202.png}
\end{center}
\par Які з наведених ребер входять до складу отриманого кістякового дерева?\par
1) (1, 4)\par
2) (1, 3)\par
3) (4, 5)\par
4) (1, 2)\par
5) (0, 6)\par
6) (0, 5)\par
{\vskip 15pt} У формі вибрати номери відповідних ребер.\par
%answer:2 4 6
\end {large}}
%end
{\smallskip \begin {small} Затверджено на засіданні кафедри теоретичної кібернетики, протокол \No 12 від   19.05.2020 р.\par\smallskip\smallskip
{\parbox {6.5 cm}{Зав. кафедри \dotfill Ю.В.~Крак}}\hspace {0.7cm}{\hbox to 6.5 cm{ Екзаменатор \dotfill Т.О.~Карнаух}} \end{small}}
%begin
{{\begin {small}\par
{ \center  Київський національний університет імені Тараса Шевченка \par}
{\parbox {5 cm}{Освітня програма \hfill}{\parbox {7 cm}{Прикладна математика, Системний аналіз \hfill}}\parbox {6 cm}{\hfill Семестр 2}}\par
{\parbox {5 cm} {Навчальний предмет \hfill}\parbox {7 cm}{Програмування \hfill}\parbox {6cm}{\hfill}\par}\vspace{-7pt} \end{small}}{\center Екзаменаційний білет \No \textbf 
{ 353 }\par}}
{
\begin {large}
1. Що буде виведено за виконання?
code
try:
    res = 6.0!=2 and not 4>2.0 or False>=7
    if isinstance(res, bool):
        print(f'bool;{res}')
    elif isinstance(res, int):
        print(f'int;{res}')
    elif isinstance(res, float):
        print(f'float;{res:.2f}')
    else:
        print('???')
except:
    print('ERR')
code

2. Що буде виведено за виконання?
code
a = 2
b = 7
c = 5

def f(b):
    a = 2
    b = 1
    c = 3
    return a + b + c

a = 8
b = 8
c = 9

try:
    print(f(a), a, b, c, sep=';')
except:
    print('ERR', a, b, c, sep=';')
code

3. Що буде виведено за виконання?
code
try:
    j = 12
    while j>5:
        j += -3
    print(j, end=';')
except:
    print('ERR')
code

4. Що буде виведено за виконання?
code
try:
    for e in range(4, -5, -1):
        if e<=7:
            continue
        print(e, end=';');e = 10
    else:
        print(e,end=';')
    print(e)
except:
    print('ERR')
code

5. Що буде виведено за виконання?\par (Дійсні значення, якщо наявні, в цьому завданні будуть виводитись у форматі з фіксованою крапкою та, можливо, з доволі великою кількістю десяткових розрядів. У відповіді для дійсних значень у ЦЬОМУ завданні достатньо вказати один десятковий розряд після крапки, відкинувши решту розрядів без округлення. Наприклад, замість 13.66666666666667 достатньо вказати 13.6, замість 25.23 достатньо вказати 25.2)
code
try:
    print(9, end=';')
    print(int('2'), end=';')
    print(8, end=';')
except KeyboardInterrupt:
    print(6, end=';')
except Exception:
    print(1, end=';')
except:
    print('ERR', end=';')
else:
    print(7, end=';')
finally:
    print(4)
code

6. Що буде виведено за виконання?
code
def f(n):
    print(n%2, end='')
    if n>2:
        f(n//2)
 
f(16)
code

7. Що буде виведено за виконання?
code
class A:
    c = 1

    def _f1(self): print(2, end=';')
    def _f2(self): print(3, end=';')

    def main(self):
        try: print(self.c, end=';')
        except: print('ERR', end=';')

        try: self._f1()
        except: print('ERR', end=';')

        try: self._f2()
        except: print('ERR', end=';')


class B(A):
    c = 4

    def _f1(self): print(5, end=';')
    def _f2(self): print(6, end=';')

try: B().main()
except: print('ERR', end=';')
code

8. Що буде виведено за виконання?
code
def f(arg, n):
    n = 7
    tmp = arg
    tmp[-2:0] = [13]
    return len(tmp)>3

try:
    g = [10,11,12,13,14]
    n = 9
    f(g, n)
    print(n, end=';')
    for el in g:
        print(el, end=';')
except:
    print('ERR')
code

9. Що буде виведено за виконання?
code
class A:
    a = 1
    c = 2

    def __init__(self):
        c = 3

class B(A):
    c = 4

    def __init__(self):
        super().__init__()
        self.b = 7

obj = B()
obj.c = obj.a + 10
B.c = 13

try: print(obj.c, end= ';')
except: print('ERR', end=';')

try: print(A.c, end= ';')
except: print('ERR', end=';')

try: print(B.c, end= ';')
except: print('ERR', end=';')
code

10. 
Для графа на рисунку з вершини 6 запущено алгоритм пошуку в глибину, що будує кістякове дерево. Списки суміжності вершин переглядаються алгоритмом за зростанням міток вершин.\par
\begin{center}
	\adjustimage{max size={0.9\linewidth}{0.9\paperheight}}
	{../10/Traversals/203.png}
\end{center}
\par Які з наведених ребер входять до складу отриманого кістякового дерева?\par
1) (0, 6)\par
2) (0, 2)\par
3) (1, 4)\par
4) (4, 5)\par
5) (2, 6)\par
6) (0, 3)\par
{\vskip 15pt} У формі вибрати номери відповідних ребер.\par
%answer:1 2 3 4 6
\end {large}}
%end
{\smallskip \begin {small} Затверджено на засіданні кафедри теоретичної кібернетики, протокол \No 12 від   19.05.2020 р.\par\smallskip\smallskip
{\parbox {6.5 cm}{Зав. кафедри \dotfill Ю.В.~Крак}}\hspace {0.7cm}{\hbox to 6.5 cm{ Екзаменатор \dotfill Т.О.~Карнаух}} \end{small}}
%begin
{{\begin {small}\par
{ \center  Київський національний університет імені Тараса Шевченка \par}
{\parbox {5 cm}{Освітня програма \hfill}{\parbox {7 cm}{Прикладна математика, Системний аналіз \hfill}}\parbox {6 cm}{\hfill Семестр 2}}\par
{\parbox {5 cm} {Навчальний предмет \hfill}\parbox {7 cm}{Програмування \hfill}\parbox {6cm}{\hfill}\par}\vspace{-7pt} \end{small}}{\center Екзаменаційний білет \No \textbf 
{ 354 }\par}}
{
\begin {large}
1. Що буде виведено за виконання?
code
try:
    res = 1**-4-False
    if isinstance(res, bool):
        print(f'bool;{res}')
    elif isinstance(res, int):
        print(f'int;{res}')
    elif isinstance(res, float):
        print(f'float;{res:.2f}')
    else:
        print('???')
except:
    print('ERR')
code

2. Що буде виведено за виконання?
code
a = 5
b = 1
c = 3

def f(a):
    global c
    a = 2
    b += 4
    c = 1
    return a + b + c

a = 6
b = 4
c = 1

try:
    print(f(b), a, b, c, sep=';')
except:
    print('ERR', a, b, c, sep=';')
code

3. Що буде виведено за виконання?
code
try:
    t = 2
    while t<=13:
        t += 1
    print(t, end=';')
except:
    print('ERR')
code

4. Що буде виведено за виконання?
code
try:
    for c in range(5, 14, 2):
        if c>=5:
            continue
        print(c, end=';')
    else:
        print('end',end=';')
    print(c)
except:
    print('ERR')
code

5. Що буде виведено за виконання?\par (Дійсні значення, якщо наявні, в цьому завданні будуть виводитись у форматі з фіксованою крапкою та, можливо, з доволі великою кількістю десяткових розрядів. У відповіді для дійсних значень у ЦЬОМУ завданні достатньо вказати один десятковий розряд після крапки, відкинувши решту розрядів без округлення. Наприклад, замість 13.66666666666667 достатньо вказати 13.6, замість 25.23 достатньо вказати 25.2)
code
try:
    print(7, end=';')
    print(9%2, end=';')
    print(5, end=';')
except TypeError:
    print(8, end=';')
except KeyboardInterrupt:
    print(0, end=';')
except:
    print('ERR', end=';')
else:
    print(3, end=';')
finally:
    print(2)
code

6. Що буде виведено за виконання?
code
def f(n):
    print(n%10,end='')
    if n>9:
        f(n//100)
 
f(123456)
code

7. Що буде виведено за виконання?
code
class A:
    a = 1

    def __h1(self): print(2, end=';')
    def _h2(self): print(3, end=';')

    def main(self):
        try: print(self.a, end=';')
        except: print('ERR', end=';')

        try: self.__h1()
        except: print('ERR', end=';')

        try: self._h2()
        except: print('ERR', end=';')


class B(A):
    a = 4

    def __h1(self): print(5, end=';')
    def _h2(self): print(6, end=';')

try: B().main()
except: print('ERR', end=';')
code

8. Що буде виведено за виконання?
code
def f(arg, n):
    n = 0
    tmp = arg
    del tmp[7:-3]
    return len(tmp)>3

try:
    f = ['0','1','2','3','4']
    n = 4
    f(f, n)
    print(n, end=';')
    for el in f:
        print(el, end=';')
except:
    print('ERR')
code

9. Що буде виведено за виконання?
code
class A:
    b = 1
    c = 2

    def __init__(self):
        b = 3

class B(A):
    b = 4

    def __init__(self):
        super().__init__()
        self.a = 7

obj = B()
obj.b = obj.c + 10
B.b = 13

try: print(obj.b, end= ';')
except: print('ERR', end=';')

try: print(A.b, end= ';')
except: print('ERR', end=';')

try: print(B.b, end= ';')
except: print('ERR', end=';')
code

10. 
Для графа на рисунку з вершини 4 запущено алгоритм пошуку в глибину, що будує кістякове дерево. Списки суміжності вершин переглядаються алгоритмом за зростанням міток вершин.\par
\begin{center}
	\adjustimage{max size={0.9\linewidth}{0.9\paperheight}}
	{../10/Traversals/84.png}
\end{center}
\par Які з наведених ребер входять до складу отриманого кістякового дерева?\par
1) (1, 5)\par
2) (5, 6)\par
3) (3, 6)\par
4) (0, 6)\par
5) (3, 5)\par
6) (2, 6)\par
{\vskip 15pt} У формі вибрати номери відповідних ребер.\par
%answer:2 4 6
\end {large}}
%end
{\smallskip \begin {small} Затверджено на засіданні кафедри теоретичної кібернетики, протокол \No 12 від   19.05.2020 р.\par\smallskip\smallskip
{\parbox {6.5 cm}{Зав. кафедри \dotfill Ю.В.~Крак}}\hspace {0.7cm}{\hbox to 6.5 cm{ Екзаменатор \dotfill Т.О.~Карнаух}} \end{small}}
%begin
{{\begin {small}\par
{ \center  Київський національний університет імені Тараса Шевченка \par}
{\parbox {5 cm}{Освітня програма \hfill}{\parbox {7 cm}{Прикладна математика, Системний аналіз \hfill}}\parbox {6 cm}{\hfill Семестр 2}}\par
{\parbox {5 cm} {Навчальний предмет \hfill}\parbox {7 cm}{Програмування \hfill}\parbox {6cm}{\hfill}\par}\vspace{-7pt} \end{small}}{\center Екзаменаційний білет \No \textbf 
{ 355 }\par}}
{
\begin {large}
1. Що буде виведено за виконання?
code
try:
    res = 9**-1.0*1
    if isinstance(res, bool):
        print(f'bool;{res}')
    elif isinstance(res, int):
        print(f'int;{res}')
    elif isinstance(res, float):
        print(f'float;{res:.2f}')
    else:
        print('???')
except:
    print('ERR')
code

2. Що буде виведено за виконання?
code
a = 5
b = 7
c = 3

def f(a):
    global c
    a = 4
    b = 2
    c = 1
    return a + b + c

a = 2
b = 6
c = 4

try:
    print(f(a), a, b, c, sep=';')
except:
    print('ERR', a, b, c, sep=';')
code

3. Що буде виведено за виконання?
code
try:
    m = 5
    while m>=6:
        m += -2
    print(m, end=';')
except:
    print('ERR')
code

4. Що буде виведено за виконання?
code
try:
    for f in range(15, 2, -3):
        if f<6:
            continue
        print(f, end=';');f = -9
        if f>4:
            break
    else:
        print(f,end=';')
    print(f)
except:
    print('ERR')
code

5. Що буде виведено за виконання?\par (Дійсні значення, якщо наявні, в цьому завданні будуть виводитись у форматі з фіксованою крапкою та, можливо, з доволі великою кількістю десяткових розрядів. У відповіді для дійсних значень у ЦЬОМУ завданні достатньо вказати один десятковий розряд після крапки, відкинувши решту розрядів без округлення. Наприклад, замість 13.66666666666667 достатньо вказати 13.6, замість 25.23 достатньо вказати 25.2)
code
try:
    print(1, end=';')
    print(int('a6'), end=';')
    print(8, end=';')
except ZeroDivisionError:
    print(7, end=';')
except ValueError:
    print(3, end=';')
except:
    print('ERR', end=';')
else:
    print(6, end=';')
finally:
    print(4)
code

6. Що буде виведено за виконання?
code
def f(n):
    if n>9:
        f(n//100)
    else:
        print(n%10,end='')

f(987654)
code

7. Що буде виведено за виконання?
code
class A:
    b = 1

    def f1(self): print(2, end=';')
    def _f2(self): print(3, end=';')

    def main(self):
        try: print(self.b, end=';')
        except: print('ERR', end=';')

        try: self.f1()
        except: print('ERR', end=';')

        try: self._f2()
        except: print('ERR', end=';')


class B(A):
    b = 4

    def f1(self): print(5, end=';')
    def _f2(self): print(6, end=';')

try: B().main()
except: print('ERR', end=';')
code

8. Що буде виведено за виконання?
code
def f(arg, n):
    arg[20] = 6
    n = 6
    arg = list(arg.values())

try:
    n = 1
    d = {20:5, 34:8, 57:2, 20:8}
    f(d, n)
    print(n, end=';')
    for x in d.keys():
        print(x, sep=';', end=';')
except:
    print('ERR')
code

9. Що буде виведено за виконання?
code
class A:
    a = 1
    b = 2

    def __init__(self):
        a = 3

class B(A):
    a = 4

    def __init__(self):
        super().__init__()
        self.c = 7

obj = B()
obj.a = obj.b + 10
B.a = 13

try: print(obj.a, end= ';')
except: print('ERR', end=';')

try: print(A.a, end= ';')
except: print('ERR', end=';')

try: print(B.a, end= ';')
except: print('ERR', end=';')
code

10. 
Для графа на рисунку з вершини 0 запущено алгоритм пошуку в глибину, що будує кістякове дерево. Списки суміжності вершин переглядаються алгоритмом за зростанням міток вершин.\par
\begin{center}
	\adjustimage{max size={0.9\linewidth}{0.9\paperheight}}
	{../10/Traversals/6.png}
\end{center}
\par Які з наведених ребер входять до складу отриманого кістякового дерева?\par
1) (0, 4)\par
2) (2, 4)\par
3) (4, 5)\par
4) (5, 6)\par
5) (0, 5)\par
6) (3, 6)\par
{\vskip 15pt} У формі вибрати номери відповідних ребер.\par
%answer:2 4 6
\end {large}}
%end
{\smallskip \begin {small} Затверджено на засіданні кафедри теоретичної кібернетики, протокол \No 12 від   19.05.2020 р.\par\smallskip\smallskip
{\parbox {6.5 cm}{Зав. кафедри \dotfill Ю.В.~Крак}}\hspace {0.7cm}{\hbox to 6.5 cm{ Екзаменатор \dotfill Т.О.~Карнаух}} \end{small}}
%begin
{{\begin {small}\par
{ \center  Київський національний університет імені Тараса Шевченка \par}
{\parbox {5 cm}{Освітня програма \hfill}{\parbox {7 cm}{Прикладна математика, Системний аналіз \hfill}}\parbox {6 cm}{\hfill Семестр 2}}\par
{\parbox {5 cm} {Навчальний предмет \hfill}\parbox {7 cm}{Програмування \hfill}\parbox {6cm}{\hfill}\par}\vspace{-7pt} \end{small}}{\center Екзаменаційний білет \No \textbf 
{ 356 }\par}}
{
\begin {large}
1. Що буде виведено за виконання?
code
try:
    res = 1**1.0*True
    if isinstance(res, bool):
        print(f'bool;{res}')
    elif isinstance(res, int):
        print(f'int;{res}')
    elif isinstance(res, float):
        print(f'float;{res:.2f}')
    else:
        print('???')
except:
    print('ERR')
code

2. Що буде виведено за виконання?
code
a = 0
b = 1
c = 0

def g(b):
    a = 5
    b = 4
    c = 2
    return a + b + c

a = 6
b = 3
c = 9

try:
    print(g(b), a, b, c, sep=';')
except:
    print('ERR', a, b, c, sep=';')
code

3. Що буде виведено за виконання?
code
try:
    n = 6
    while n<10:
        n += 2
    print(n, end=';')
except:
    print('ERR')
code

4. Що буде виведено за виконання?
code
try:
    for b in range(-9, 7, -2):
        if b>7:
            continue
        print(b, end=';')
    else:
        print('end',end=';')
    print(b)
except:
    print('ERR')
code

5. Що буде виведено за виконання?\par (Дійсні значення, якщо наявні, в цьому завданні будуть виводитись у форматі з фіксованою крапкою та, можливо, з доволі великою кількістю десяткових розрядів. У відповіді для дійсних значень у ЦЬОМУ завданні достатньо вказати один десятковий розряд після крапки, відкинувши решту розрядів без округлення. Наприклад, замість 13.66666666666667 достатньо вказати 13.6, замість 25.23 достатньо вказати 25.2)
code
try:
    print(9, end=';')
    print(5j!=3j, end=';')
    print(1, end=';')
except KeyboardInterrupt:
    print(2, end=';')
except ValueError:
    print(0, end=';')
except:
    print('ERR', end=';')
else:
    print(9, end=';')
finally:
    print(1)
code

6. Що буде виведено за виконання?
code
def f(n):
    if n>9:
        f(n//10)
    print(n%10, end='')
 
f(123456)
code

7. Що буде виведено за виконання?
code
class A:
    d = 1

    def _g1(self): print(2, end=';')
    def __g2(self): print(3, end=';')

    def main(self):
        try: print(self.d, end=';')
        except: print('ERR', end=';')

        try: self._g1()
        except: print('ERR', end=';')

        try: self.__g2()
        except: print('ERR', end=';')


class B(A):
    d = 4

    def _g1(self): print(5, end=';')
    def __g2(self): print(6, end=';')

try: B().main()
except: print('ERR', end=';')
code

8. Що буде виведено за виконання?
code
def f(arg, n):
    arg[87] = 6
    n = 8
    arg = set(arg.keys())

try:
    n = 0
    t = {42:7, 87:7, 89:0}
    f(t, n)
    print(n, end=';')
    for x in t.items():
        print(*x, sep=';', end=';')
except:
    print('ERR')
code

9. Що буде виведено за виконання?
code
class A:
    b = 1
    a = 2

    def __init__(self):
        a = 3

class B(A):
    a = 4

    def __init__(self):
        super().__init__()
        self.c = 7

obj = B()
obj.c = obj.a + 10
B.c = 13

try: print(obj.c, end= ';')
except: print('ERR', end=';')

try: print(A.c, end= ';')
except: print('ERR', end=';')

try: print(B.c, end= ';')
except: print('ERR', end=';')
code

10. 
Для графа на рисунку з вершини 6 запущено алгоритм пошуку в глибину, що будує кістякове дерево. Списки суміжності вершин переглядаються алгоритмом за зростанням міток вершин.\par
\begin{center}
	\adjustimage{max size={0.9\linewidth}{0.9\paperheight}}
	{../10/Traversals/188.png}
\end{center}
\par Які з наведених ребер входять до складу отриманого кістякового дерева?\par
1) (1, 6)\par
2) (2, 4)\par
3) (1, 4)\par
4) (3, 6)\par
5) (1, 3)\par
6) (1, 2)\par
{\vskip 15pt} У формі вибрати номери відповідних ребер.\par
%answer:1 2 5
\end {large}}
%end
{\smallskip \begin {small} Затверджено на засіданні кафедри теоретичної кібернетики, протокол \No 12 від   19.05.2020 р.\par\smallskip\smallskip
{\parbox {6.5 cm}{Зав. кафедри \dotfill Ю.В.~Крак}}\hspace {0.7cm}{\hbox to 6.5 cm{ Екзаменатор \dotfill Т.О.~Карнаух}} \end{small}}
%begin
{{\begin {small}\par
{ \center  Київський національний університет імені Тараса Шевченка \par}
{\parbox {5 cm}{Освітня програма \hfill}{\parbox {7 cm}{Прикладна математика, Системний аналіз \hfill}}\parbox {6 cm}{\hfill Семестр 2}}\par
{\parbox {5 cm} {Навчальний предмет \hfill}\parbox {7 cm}{Програмування \hfill}\parbox {6cm}{\hfill}\par}\vspace{-7pt} \end{small}}{\center Екзаменаційний білет \No \textbf 
{ 357 }\par}}
{
\begin {large}
1. Що буде виведено за виконання?
code
def f(n):
    print('f', end='')
    return n

try:
    print(f(9) or f(2))
except:
    print('ERR')
code

2. Що буде виведено за виконання?
code
a = 0
b = 6
c = 3

def h(b):
    a += 3
    b = 5
    c = 2
    return a + b + c

a = 8
b = 9
c = 2

try:
    print(h(b), a, b, c, sep=';')
except:
    print('ERR', a, b, c, sep=';')
code

3. Що буде виведено за виконання?
code
try:
    n = 6
    while n>3:
        n += -3
    print(n, end=';')
except:
    print('ERR')
code

4. Що буде виведено за виконання?
code
try:
    for c in range(3, 3, -3):
        if c>3:
            break
        print(c, end=';')
    else:
        print(c,end=';')
    print(c)
except:
    print('ERR')
code

5. Що буде виведено за виконання?\par (Дійсні значення, якщо наявні, в цьому завданні будуть виводитись у форматі з фіксованою крапкою та, можливо, з доволі великою кількістю десяткових розрядів. У відповіді для дійсних значень у ЦЬОМУ завданні достатньо вказати один десятковий розряд після крапки, відкинувши решту розрядів без округлення. Наприклад, замість 13.66666666666667 достатньо вказати 13.6, замість 25.23 достатньо вказати 25.2)
code
try:
    print(4, end=';')
    print(int('6'), end=';')
    print(2, end=';')
except TypeError:
    print(8, end=';')
except ZeroDivisionError:
    print(0, end=';')
except:
    print('ERR', end=';')
else:
    print(5, end=';')
finally:
    print(7)
code

6. Що буде виведено за виконання?
code
def f(n):
    print(n%10, end='')
    if n>9:
        f(n//10)
 
f(123456)
code

7. Що буде виведено за виконання?
code
class A:
    b = 1

    def __h1(self): print(2, end=';')
    def _h2(self): print(3, end=';')

    def main(self):
        try: print(self.b, end=';')
        except: print('ERR', end=';')

        try: self.__h1()
        except: print('ERR', end=';')

        try: self._h2()
        except: print('ERR', end=';')


class B(A):
    b = 4

    def __h1(self): print(5, end=';')
    def _h2(self): print(6, end=';')

try: B().main()
except: print('ERR', end=';')
code

8. Що буде виведено за виконання?
code
def f(arg, n):
    arg[51] = 9
    n += 2

try:
    n = 3
    d = {11:7, 60:2, 11:1}
    f(d, n)
    print(n, end=';')
    for x in d.keys():
        print(x, sep=';', end=';')
except:
    print('ERR')
code

9. Що буде виведено за виконання?
code
class A:
    c = 1
    b = 2

    def __init__(self):
        c = 3

class B(A):
    b = 4

    def __init__(self):
        super().__init__()
        self.a = 7

obj = B()
obj.c = obj.b + 10
B.c = 13

try: print(obj.c, end= ';')
except: print('ERR', end=';')

try: print(A.c, end= ';')
except: print('ERR', end=';')

try: print(B.c, end= ';')
except: print('ERR', end=';')
code

10. 
Для графа на рисунку з вершини 1 запущено алгоритм пошуку в глибину, що будує кістякове дерево. Списки суміжності вершин переглядаються алгоритмом за зростанням міток вершин.\par
\begin{center}
	\adjustimage{max size={0.9\linewidth}{0.9\paperheight}}
	{../10/Traversals/94.png}
\end{center}
\par Які з наведених ребер входять до складу отриманого кістякового дерева?\par
1) (1, 2)\par
2) (1, 6)\par
3) (5, 6)\par
4) (0, 5)\par
5) (0, 6)\par
6) (2, 5)\par
{\vskip 15pt} У формі вибрати номери відповідних ребер.\par
%answer:3 6
\end {large}}
%end
{\smallskip \begin {small} Затверджено на засіданні кафедри теоретичної кібернетики, протокол \No 12 від   19.05.2020 р.\par\smallskip\smallskip
{\parbox {6.5 cm}{Зав. кафедри \dotfill Ю.В.~Крак}}\hspace {0.7cm}{\hbox to 6.5 cm{ Екзаменатор \dotfill Т.О.~Карнаух}} \end{small}}
%begin
{{\begin {small}\par
{ \center  Київський національний університет імені Тараса Шевченка \par}
{\parbox {5 cm}{Освітня програма \hfill}{\parbox {7 cm}{Прикладна математика, Системний аналіз \hfill}}\parbox {6 cm}{\hfill Семестр 2}}\par
{\parbox {5 cm} {Навчальний предмет \hfill}\parbox {7 cm}{Програмування \hfill}\parbox {6cm}{\hfill}\par}\vspace{-7pt} \end{small}}{\center Екзаменаційний білет \No \textbf 
{ 358 }\par}}
{
\begin {large}
1. Що буде виведено за виконання?
code
try:
    res = 4**-0.5*-1
    if isinstance(res, bool):
        print(f'bool;{res}')
    elif isinstance(res, int):
        print(f'int;{res}')
    elif isinstance(res, float):
        print(f'float;{res:.2f}')
    else:
        print('???')
except:
    print('ERR')
code

2. Що буде виведено за виконання?
code
a = 8
b = 2
c = 7

def h(a):
    global c
    a = 5
    b = 3
    c = 1
    return a + b + c

a = 5
b = 1
c = 4

try:
    print(h(a), a, b, c, sep=';')
except:
    print('ERR', a, b, c, sep=';')
code

3. Що буде виведено за виконання?
code
try:
    k = 0
    while k<12:
        k += 2
    print(k, end=';')
except:
    print('ERR')
code

4. Що буде виведено за виконання?
code
try:
    for d in range(13, -11, 3):
        if d<8:
            break
        print(d, end=';')
    else:
        print(d,end=';')
    print(d)
except:
    print('ERR')
code

5. Що буде виведено за виконання?\par (Дійсні значення, якщо наявні, в цьому завданні будуть виводитись у форматі з фіксованою крапкою та, можливо, з доволі великою кількістю десяткових розрядів. У відповіді для дійсних значень у ЦЬОМУ завданні достатньо вказати один десятковий розряд після крапки, відкинувши решту розрядів без округлення. Наприклад, замість 13.66666666666667 достатньо вказати 13.6, замість 25.23 достатньо вказати 25.2)
code
try:
    print(3, end=';')
    print(6//0.0, end=';')
    print(3, end=';')
except Exception:
    print(7, end=';')
except ZeroDivisionError:
    print(1, end=';')
except:
    print('ERR', end=';')
else:
    print(9, end=';')
finally:
    print(2)
code

6. Що буде виведено за виконання?
code
def f(n):
    if n>9:
        return f(n//100)+1
    else:
        return 1

print(f(123456))
code

7. Що буде виведено за виконання?
code
class A:
    d = 1

    def _g1(self): print(2, end=';')
    def g2(self): print(3, end=';')

    def main(self):
        try: print(self.d, end=';')
        except: print('ERR', end=';')

        try: self._g1()
        except: print('ERR', end=';')

        try: self.g2()
        except: print('ERR', end=';')


class B(A):
    d = 4

    def _g1(self): print(5, end=';')
    def g2(self): print(6, end=';')

try: B().main()
except: print('ERR', end=';')
code

8. Що буде виведено за виконання?
code
def f(arg, n):
    arg[27] = 6
    n *= 3

try:
    n = 6
    t = {34:4, 26:0, 88:1, 26:2}
    f(t, n)
    print(n, end=';')
    for x in t.keys():
        print(x, sep=';', end=';')
except:
    print('ERR')
code

9. Що буде виведено за виконання?
code
class A:
    b = 1
    c = 2

    def __init__(self):
        c = 3

class B(A):
    b = 4

    def __init__(self):
        super().__init__()
        self.a = 7

obj = B()
obj.a = obj.c + 10
B.a = 13

try: print(obj.a, end= ';')
except: print('ERR', end=';')

try: print(A.a, end= ';')
except: print('ERR', end=';')

try: print(B.a, end= ';')
except: print('ERR', end=';')
code

10. 
Для графа на рисунку з вершини 4 запущено алгоритм пошуку в глибину, що будує кістякове дерево. Списки суміжності вершин переглядаються алгоритмом за зростанням міток вершин.\par
\begin{center}
	\adjustimage{max size={0.9\linewidth}{0.9\paperheight}}
	{../10/Traversals/201.png}
\end{center}
\par Які з наведених ребер входять до складу отриманого кістякового дерева?\par
1) (1, 4)\par
2) (1, 5)\par
3) (3, 5)\par
4) (0, 6)\par
5) (0, 2)\par
6) (3, 4)\par
{\vskip 15pt} У формі вибрати номери відповідних ребер.\par
%answer:1 2 3 4 5
\end {large}}
%end
{\smallskip \begin {small} Затверджено на засіданні кафедри теоретичної кібернетики, протокол \No 12 від   19.05.2020 р.\par\smallskip\smallskip
{\parbox {6.5 cm}{Зав. кафедри \dotfill Ю.В.~Крак}}\hspace {0.7cm}{\hbox to 6.5 cm{ Екзаменатор \dotfill Т.О.~Карнаух}} \end{small}}
%begin
{{\begin {small}\par
{ \center  Київський національний університет імені Тараса Шевченка \par}
{\parbox {5 cm}{Освітня програма \hfill}{\parbox {7 cm}{Прикладна математика, Системний аналіз \hfill}}\parbox {6 cm}{\hfill Семестр 2}}\par
{\parbox {5 cm} {Навчальний предмет \hfill}\parbox {7 cm}{Програмування \hfill}\parbox {6cm}{\hfill}\par}\vspace{-7pt} \end{small}}{\center Екзаменаційний білет \No \textbf 
{ 359 }\par}}
{
\begin {large}
1. Що буде виведено за виконання?
code
def f(n):
    print('f', end='')
    return n

try:
    print(f(-6) or f(-9))
except:
    print('ERR')
code

2. Що буде виведено за виконання?
code
a = 7
b = 5
c = 9

def f(b):
    a = 4
    b *= 2
    c = 4
    return a + b + c

a = 8
b = 1
c = 6

try:
    print(f(b), a, b, c, sep=';')
except:
    print('ERR', a, b, c, sep=';')
code

3. Що буде виведено за виконання?
code
try:
    i = 8
    while i<=8:
        i += 1
    print(i, end=';')
except:
    print('ERR')
code

4. Що буде виведено за виконання?
code
try:
    for e in range(-7, -7, 2):
        if e<=4:
            continue
        print(e, end=';');e = 1
    else:
        print(e,end=';')
    print(e)
except:
    print('ERR')
code

5. Що буде виведено за виконання?\par (Дійсні значення, якщо наявні, в цьому завданні будуть виводитись у форматі з фіксованою крапкою та, можливо, з доволі великою кількістю десяткових розрядів. У відповіді для дійсних значень у ЦЬОМУ завданні достатньо вказати один десятковий розряд після крапки, відкинувши решту розрядів без округлення. Наприклад, замість 13.66666666666667 достатньо вказати 13.6, замість 25.23 достатньо вказати 25.2)
code
try:
    print(8, end=';')
    print(4/3, end=';')
    print(5, end=';')
except Exception:
    print(0, end=';')
except ValueError:
    print(4, end=';')
except:
    print('ERR', end=';')
else:
    print(7, end=';')
finally:
    print(9)
code

6. Що буде виведено за виконання?
code
def f(n):
    print(n%10, end='')
    if n>9:
        f(n//100)
 
f(987654)
code

7. Що буде виведено за виконання?
code
class A:
    a = 1

    def _h1(self): print(2, end=';')
    def _h2(self): print(3, end=';')

    def main(self):
        try: print(self.a, end=';')
        except: print('ERR', end=';')

        try: self._h1()
        except: print('ERR', end=';')

        try: self._h2()
        except: print('ERR', end=';')


class B(A):
    a = 4

    def _h1(self): print(5, end=';')
    def _h2(self): print(6, end=';')

try: B().main()
except: print('ERR', end=';')
code

8. Що буде виведено за виконання?
code
def f(arg, n):
    arg[72] = 1
    n = 9
    arg = tuple(arg.items())

try:
    n = 2
    s = {72:9, 34:2, 72:2}
    f(s, n)
    print(n, end=';')
    for x in s.values():
        print(x, sep=';', end=';')
except:
    print('ERR')
code

9. Що буде виведено за виконання?
code
class A:
    c = 1
    a = 2

    def __init__(self):
        a = 3

class B(A):
    a = 4

    def __init__(self):
        super().__init__()
        self.b = 7

obj = B()
obj.b = obj.c + 10
B.b = 13

try: print(obj.b, end= ';')
except: print('ERR', end=';')

try: print(A.b, end= ';')
except: print('ERR', end=';')

try: print(B.b, end= ';')
except: print('ERR', end=';')
code

10. 
Для графа на рисунку з вершини 5 запущено алгоритм пошуку в глибину, що будує кістякове дерево. Списки суміжності вершин переглядаються алгоритмом за зростанням міток вершин.\par
\begin{center}
	\adjustimage{max size={0.9\linewidth}{0.9\paperheight}}
	{../10/Traversals/174.png}
\end{center}
\par Які з наведених ребер входять до складу отриманого кістякового дерева?\par
1) (5, 6)\par
2) (1, 3)\par
3) (0, 4)\par
4) (3, 5)\par
5) (2, 5)\par
6) (1, 6)\par
{\vskip 15pt} У формі вибрати номери відповідних ребер.\par
%answer:2 5 6
\end {large}}
%end
{\smallskip \begin {small} Затверджено на засіданні кафедри теоретичної кібернетики, протокол \No 12 від   19.05.2020 р.\par\smallskip\smallskip
{\parbox {6.5 cm}{Зав. кафедри \dotfill Ю.В.~Крак}}\hspace {0.7cm}{\hbox to 6.5 cm{ Екзаменатор \dotfill Т.О.~Карнаух}} \end{small}}
%begin
{{\begin {small}\par
{ \center  Київський національний університет імені Тараса Шевченка \par}
{\parbox {5 cm}{Освітня програма \hfill}{\parbox {7 cm}{Прикладна математика, Системний аналіз \hfill}}\parbox {6 cm}{\hfill Семестр 2}}\par
{\parbox {5 cm} {Навчальний предмет \hfill}\parbox {7 cm}{Програмування \hfill}\parbox {6cm}{\hfill}\par}\vspace{-7pt} \end{small}}{\center Екзаменаційний білет \No \textbf 
{ 360 }\par}}
{
\begin {large}
1. Що буде виведено за виконання?
code
def f(n):
    print('f', end='')
    return n>=-4

try:
    print(f(-5) and f(-1))
except:
    print('ERR')
code

2. Що буде виведено за виконання?
code
a = 4
b = 7
c = 0

def g(a):
    global c
    a = 3
    b -= 5
    c = 1
    return a + b + c

a = 2
b = 5
c = 3

try:
    print(g(a), a, b, c, sep=';')
except:
    print('ERR', a, b, c, sep=';')
code

3. Що буде виведено за виконання?
code
try:
    k = 7
    while k>=0:
        k += -2
    print(k, end=';')
except:
    print('ERR')
code

4. Що буде виведено за виконання?
code
try:
    for a in range(-1, -13, -1):
        if a>=5:
            continue
        print(a, end=';')
    else:
        print(a,end=';')
    print(a)
except:
    print('ERR')
code

5. Що буде виведено за виконання?\par (Дійсні значення, якщо наявні, в цьому завданні будуть виводитись у форматі з фіксованою крапкою та, можливо, з доволі великою кількістю десяткових розрядів. У відповіді для дійсних значень у ЦЬОМУ завданні достатньо вказати один десятковий розряд після крапки, відкинувши решту розрядів без округлення. Наприклад, замість 13.66666666666667 достатньо вказати 13.6, замість 25.23 достатньо вказати 25.2)
code
try:
    print(6, end=';')
    print(int('b3'), end=';')
    print(8, end=';')
except KeyboardInterrupt:
    print(1, end=';')
except TypeError:
    print(5, end=';')
except:
    print('ERR', end=';')
else:
    print(2, end=';')
finally:
    print(0)
code

6. Що буде виведено за виконання?
code
def f(s):
    for i in range(len(s)):
        s[i] += 1
    return
 
s = [1, 2, 3]
f(s)
print(s)
code

7. Що буде виведено за виконання?
code
class A:
    c = 1

    def f1(self): print(2, end=';')
    def _f2(self): print(3, end=';')

    def main(self):
        try: print(self.c, end=';')
        except: print('ERR', end=';')

        try: self.f1()
        except: print('ERR', end=';')

        try: self._f2()
        except: print('ERR', end=';')


class B(A):
    c = 4

    def f1(self): print(5, end=';')
    def _f2(self): print(6, end=';')

try: B().main()
except: print('ERR', end=';')
code

8. Що буде виведено за виконання?
code
def f(arg, n):
    arg[22] = 1
    n -= -3

try:
    n = 7
    s = {12:8, 22:9, 76:7, 76:7}
    f(s, n)
    print(n, end=';')
    for x, y in s.items():
        print(x, y, sep=';', end=';')
except:
    print('ERR')
code

9. Що буде виведено за виконання?
code
class A:
    a = 1
    c = 2

    def __init__(self):
        a = 3

class B(A):
    a = 4

    def __init__(self):
        super().__init__()
        self.b = 7

obj = B()
obj.a = obj.b + 10
B.a = 13

try: print(obj.a, end= ';')
except: print('ERR', end=';')

try: print(A.a, end= ';')
except: print('ERR', end=';')

try: print(B.a, end= ';')
except: print('ERR', end=';')
code

10. 
Для графа на рисунку з вершини 2 запущено алгоритм пошуку в глибину, що будує кістякове дерево. Списки суміжності вершин переглядаються алгоритмом за зростанням міток вершин.\par
\begin{center}
	\adjustimage{max size={0.9\linewidth}{0.9\paperheight}}
	{../10/Traversals/257.png}
\end{center}
\par Які з наведених ребер входять до складу отриманого кістякового дерева?\par
1) (4, 5)\par
2) (5, 6)\par
3) (4, 6)\par
4) (1, 5)\par
5) (3, 4)\par
6) (0, 5)\par
{\vskip 15pt} У формі вибрати номери відповідних ребер.\par
%answer:3 5 6
\end {large}}
%end
{\smallskip \begin {small} Затверджено на засіданні кафедри теоретичної кібернетики, протокол \No 12 від   19.05.2020 р.\par\smallskip\smallskip
{\parbox {6.5 cm}{Зав. кафедри \dotfill Ю.В.~Крак}}\hspace {0.7cm}{\hbox to 6.5 cm{ Екзаменатор \dotfill Т.О.~Карнаух}} \end{small}}
%begin
{{\begin {small}\par
{ \center  Київський національний університет імені Тараса Шевченка \par}
{\parbox {5 cm}{Освітня програма \hfill}{\parbox {7 cm}{Прикладна математика, Системний аналіз \hfill}}\parbox {6 cm}{\hfill Семестр 2}}\par
{\parbox {5 cm} {Навчальний предмет \hfill}\parbox {7 cm}{Програмування \hfill}\parbox {6cm}{\hfill}\par}\vspace{-7pt} \end{small}}{\center Екзаменаційний білет \No \textbf 
{ 361 }\par}}
{
\begin {large}
1. Що буде виведено за виконання?
code
try:
    res = "True"=="3.0"
    if isinstance(res, bool):
        print(f'bool;{res}')
    elif isinstance(res, int):
        print(f'int;{res}')
    elif isinstance(res, float):
        print(f'float;{res:.2f}')
    else:
        print('???')
except:
    print('ERR')
code

2. Що буде виведено за виконання?
code
a = 0
b = 2
c = 8

def h(b):
    global c
    a = 5
    b = 1
    c = 2
    return a + b + c

a = 7
b = 4
c = 3

try:
    print(h(b), a, b, c, sep=';')
except:
    print('ERR', a, b, c, sep=';')
code

3. Що буде виведено за виконання?
code
try:
    m = 13
    while m>3:
        m += -3
    print(m, end=';')
except:
    print('ERR')
code

4. Що буде виведено за виконання?
code
try:
    for a in range(8, 15, 2):
        if a<6:
            break
        print(a, end=';');a = -1
    else:
        print(a,end=';')
    print(a)
except:
    print('ERR')
code

5. Що буде виведено за виконання?\par (Дійсні значення, якщо наявні, в цьому завданні будуть виводитись у форматі з фіксованою крапкою та, можливо, з доволі великою кількістю десяткових розрядів. У відповіді для дійсних значень у ЦЬОМУ завданні достатньо вказати один десятковий розряд після крапки, відкинувши решту розрядів без округлення. Наприклад, замість 13.66666666666667 достатньо вказати 13.6, замість 25.23 достатньо вказати 25.2)
code
try:
    print(9, end=';')
    print(5j>5j, end=';')
    print(6, end=';')
except ValueError:
    print(0, end=';')
except KeyboardInterrupt:
    print(3, end=';')
except:
    print('ERR', end=';')
else:
    print(8, end=';')
finally:
    print(2)
code

6. Що буде виведено за виконання?
code
def f(n):
    if n>9:
        f(n//100)
        print(n%10, end='')
 
f(123456)
code

7. Що буде виведено за виконання?
code
class A:
    b = 1

    def _h1(self): print(2, end=';')
    def _h2(self): print(3, end=';')

    def main(self):
        try: print(self.b, end=';')
        except: print('ERR', end=';')

        try: self._h1()
        except: print('ERR', end=';')

        try: self._h2()
        except: print('ERR', end=';')


class B(A):
    b = 4

    def _h1(self): print(5, end=';')
    def _h2(self): print(6, end=';')

try: B().main()
except: print('ERR', end=';')
code

8. Що буде виведено за виконання?
code
def f(arg, n):
    arg[30] = 1
    n = 9
    arg = tuple(arg.values())

try:
    n = 2
    s = {30:6, 38:9, 30:1}
    f(s, n)
    print(n, end=';')
    for x in s.keys():
        print(x, sep=';', end=';')
except:
    print('ERR')
code

9. Що буде виведено за виконання?
code
class A:
    c = 1
    a = 2

    def __init__(self):
        c = 3

class B(A):
    a = 4

    def __init__(self):
        super().__init__()
        self.b = 7

obj = B()
obj.b = obj.c + 10
B.b = 13

try: print(obj.b, end= ';')
except: print('ERR', end=';')

try: print(A.b, end= ';')
except: print('ERR', end=';')

try: print(B.b, end= ';')
except: print('ERR', end=';')
code

10. 
Для графа на рисунку з вершини 1 запущено алгоритм пошуку в глибину, що будує кістякове дерево. Списки суміжності вершин переглядаються алгоритмом за зростанням міток вершин.\par
\begin{center}
	\adjustimage{max size={0.9\linewidth}{0.9\paperheight}}
	{../10/Traversals/279.png}
\end{center}
\par Які з наведених ребер входять до складу отриманого кістякового дерева?\par
1) (3, 4)\par
2) (1, 2)\par
3) (1, 6)\par
4) (0, 4)\par
5) (5, 6)\par
6) (2, 5)\par
{\vskip 15pt} У формі вибрати номери відповідних ребер.\par
%answer:1 2 4 5
\end {large}}
%end
{\smallskip \begin {small} Затверджено на засіданні кафедри теоретичної кібернетики, протокол \No 12 від   19.05.2020 р.\par\smallskip\smallskip
{\parbox {6.5 cm}{Зав. кафедри \dotfill Ю.В.~Крак}}\hspace {0.7cm}{\hbox to 6.5 cm{ Екзаменатор \dotfill Т.О.~Карнаух}} \end{small}}
%begin
{{\begin {small}\par
{ \center  Київський національний університет імені Тараса Шевченка \par}
{\parbox {5 cm}{Освітня програма \hfill}{\parbox {7 cm}{Прикладна математика, Системний аналіз \hfill}}\parbox {6 cm}{\hfill Семестр 2}}\par
{\parbox {5 cm} {Навчальний предмет \hfill}\parbox {7 cm}{Програмування \hfill}\parbox {6cm}{\hfill}\par}\vspace{-7pt} \end{small}}{\center Екзаменаційний білет \No \textbf 
{ 362 }\par}}
{
\begin {large}
1. Що буде виведено за виконання?
code
try:
    res = 8/8/4*6
    if isinstance(res, bool):
        print(f'bool;{res}')
    elif isinstance(res, int):
        print(f'int;{res}')
    elif isinstance(res, float):
        print(f'float;{res:.2f}')
    else:
        print('???')
except:
    print('ERR')
code

2. Що буде виведено за виконання?
code
a = 6
b = 5
c = 9

def f(a):
    a = 3
    b = 4
    c = 5
    return a + b + c

a = 1
b = 2
c = 6

try:
    print(f(a), a, b, c, sep=';')
except:
    print('ERR', a, b, c, sep=';')
code

3. Що буде виведено за виконання?
code
try:
    m = 1
    while m<=10:
        m += 2
    print(m, end=';')
except:
    print('ERR')
code

4. Що буде виведено за виконання?
code
try:
    for c in range(-12, -12, -2):
        if c>=7:
            continue
        print(c, end=';')
    else:
        print(c,end=';')
    print(c)
except:
    print('ERR')
code

5. Що буде виведено за виконання?\par (Дійсні значення, якщо наявні, в цьому завданні будуть виводитись у форматі з фіксованою крапкою та, можливо, з доволі великою кількістю десяткових розрядів. У відповіді для дійсних значень у ЦЬОМУ завданні достатньо вказати один десятковий розряд після крапки, відкинувши решту розрядів без округлення. Наприклад, замість 13.66666666666667 достатньо вказати 13.6, замість 25.23 достатньо вказати 25.2)
code
try:
    print(7, end=';')
    print(1//0.0, end=';')
    print(4, end=';')
except TypeError:
    print(7, end=';')
except Exception:
    print(3, end=';')
except:
    print('ERR', end=';')
else:
    print(8, end=';')
finally:
    print(1)
code

6. Що буде виведено за виконання?
code
def f(s):
    s1 = s
    for i in range(len(s)):
        s1[i] += 1
        return s1

s = [1, 2, 3]
f(s)
print(s)
code

7. Що буде виведено за виконання?
code
class A:
    c = 1

    def _g1(self): print(2, end=';')
    def g2(self): print(3, end=';')

    def main(self):
        try: print(self.c, end=';')
        except: print('ERR', end=';')

        try: self._g1()
        except: print('ERR', end=';')

        try: self.g2()
        except: print('ERR', end=';')


class B(A):
    c = 4

    def _g1(self): print(5, end=';')
    def g2(self): print(6, end=';')

try: B().main()
except: print('ERR', end=';')
code

8. Що буде виведено за виконання?
code
def f(arg, n):
    arg[58] = 9
    n += 2

try:
    n = 6
    t = {37:7, 25:9, 25:0}
    f(t, n)
    print(n, end=';')
    for x in t.items():
        print(x, sep=';', end=';')
except:
    print('ERR')
code

9. Що буде виведено за виконання?
code
class A:
    a = 1
    b = 2

    def __init__(self):
        b = 3

class B(A):
    a = 4

    def __init__(self):
        super().__init__()
        self.c = 7

obj = B()
obj.a = obj.b + 10
B.a = 13

try: print(obj.a, end= ';')
except: print('ERR', end=';')

try: print(A.a, end= ';')
except: print('ERR', end=';')

try: print(B.a, end= ';')
except: print('ERR', end=';')
code

10. 
Для графа на рисунку з вершини 6 запущено алгоритм пошуку в глибину, що будує кістякове дерево. Списки суміжності вершин переглядаються алгоритмом за зростанням міток вершин.\par
\begin{center}
	\adjustimage{max size={0.9\linewidth}{0.9\paperheight}}
	{../10/Traversals/129.png}
\end{center}
\par Які з наведених ребер входять до складу отриманого кістякового дерева?\par
1) (2, 3)\par
2) (1, 2)\par
3) (2, 6)\par
4) (3, 6)\par
5) (4, 5)\par
6) (5, 6)\par
{\vskip 15pt} У формі вибрати номери відповідних ребер.\par
%answer:1 3 5
\end {large}}
%end
{\smallskip \begin {small} Затверджено на засіданні кафедри теоретичної кібернетики, протокол \No 12 від   19.05.2020 р.\par\smallskip\smallskip
{\parbox {6.5 cm}{Зав. кафедри \dotfill Ю.В.~Крак}}\hspace {0.7cm}{\hbox to 6.5 cm{ Екзаменатор \dotfill Т.О.~Карнаух}} \end{small}}
%begin
{{\begin {small}\par
{ \center  Київський національний університет імені Тараса Шевченка \par}
{\parbox {5 cm}{Освітня програма \hfill}{\parbox {7 cm}{Прикладна математика, Системний аналіз \hfill}}\parbox {6 cm}{\hfill Семестр 2}}\par
{\parbox {5 cm} {Навчальний предмет \hfill}\parbox {7 cm}{Програмування \hfill}\parbox {6cm}{\hfill}\par}\vspace{-7pt} \end{small}}{\center Екзаменаційний білет \No \textbf 
{ 363 }\par}}
{
\begin {large}
1. Що буде виведено за виконання?
code
try:
    res = 6//9//8*3
    if isinstance(res, bool):
        print(f'bool;{res}')
    elif isinstance(res, int):
        print(f'int;{res}')
    elif isinstance(res, float):
        print(f'float;{res:.2f}')
    else:
        print('???')
except:
    print('ERR')
code

2. Що буде виведено за виконання?
code
a = 7
b = 1
c = 0

def g(b):
    a = 5
    b *= 2
    c = 1
    return a + b + c

a = 9
b = 2
c = 6

try:
    print(g(b), a, b, c, sep=';')
except:
    print('ERR', a, b, c, sep=';')
code

3. Що буде виведено за виконання?
code
try:
    t = 9
    while t>=6:
        t += -2
    print(t, end=';')
except:
    print('ERR')
code

4. Що буде виведено за виконання?
code
try:
    for e in range(10, 8, -1):
        if e>4:
            continue
        print(e, end=';')
    else:
        print(e,end=';')
    print(e)
except:
    print('ERR')
code

5. Що буде виведено за виконання?\par (Дійсні значення, якщо наявні, в цьому завданні будуть виводитись у форматі з фіксованою крапкою та, можливо, з доволі великою кількістю десяткових розрядів. У відповіді для дійсних значень у ЦЬОМУ завданні достатньо вказати один десятковий розряд після крапки, відкинувши решту розрядів без округлення. Наприклад, замість 13.66666666666667 достатньо вказати 13.6, замість 25.23 достатньо вказати 25.2)
code
try:
    print(2, end=';')
    print(int('b5'), end=';')
    print(6, end=';')
except ZeroDivisionError:
    print(0, end=';')
except KeyboardInterrupt:
    print(4, end=';')
except:
    print('ERR', end=';')
else:
    print(9, end=';')
finally:
    print(5)
code

6. Що буде виведено за виконання?
code
def f(n):
    if n>2:
        f(n//2)
    else:
        print(n%2, end='')

f(16)
code

7. Що буде виведено за виконання?
code
class A:
    d = 1

    def __f1(self): print(2, end=';')
    def _f2(self): print(3, end=';')

    def main(self):
        try: print(self.d, end=';')
        except: print('ERR', end=';')

        try: self.__f1()
        except: print('ERR', end=';')

        try: self._f2()
        except: print('ERR', end=';')


class B(A):
    d = 4

    def __f1(self): print(5, end=';')
    def _f2(self): print(6, end=';')

try: B().main()
except: print('ERR', end=';')
code

8. Що буде виведено за виконання?
code
def f(arg, n):
    arg[12] = 9
    n = 5
    arg = set(arg.keys())

try:
    n = 0
    d = {43:8, 88:9, 76:7}
    f(d, n)
    print(n, end=';')
    for x in d.keys():
        print(x, sep=';', end=';')
except:
    print('ERR')
code

9. Що буде виведено за виконання?
code
class A:
    b = 1
    c = 2

    def __init__(self):
        c = 3

class B(A):
    c = 4

    def __init__(self):
        super().__init__()
        self.a = 7

obj = B()
obj.c = obj.a + 10
B.c = 13

try: print(obj.c, end= ';')
except: print('ERR', end=';')

try: print(A.c, end= ';')
except: print('ERR', end=';')

try: print(B.c, end= ';')
except: print('ERR', end=';')
code

10. 
Для графа на рисунку з вершини 2 запущено алгоритм пошуку в глибину, що будує кістякове дерево. Списки суміжності вершин переглядаються алгоритмом за зростанням міток вершин.\par
\begin{center}
	\adjustimage{max size={0.9\linewidth}{0.9\paperheight}}
	{../10/Traversals/23.png}
\end{center}
\par Які з наведених ребер входять до складу отриманого кістякового дерева?\par
1) (1, 2)\par
2) (2, 3)\par
3) (3, 6)\par
4) (0, 1)\par
5) (2, 5)\par
6) (3, 5)\par
{\vskip 15pt} У формі вибрати номери відповідних ребер.\par
%answer:3 4 6
\end {large}}
%end
{\smallskip \begin {small} Затверджено на засіданні кафедри теоретичної кібернетики, протокол \No 12 від   19.05.2020 р.\par\smallskip\smallskip
{\parbox {6.5 cm}{Зав. кафедри \dotfill Ю.В.~Крак}}\hspace {0.7cm}{\hbox to 6.5 cm{ Екзаменатор \dotfill Т.О.~Карнаух}} \end{small}}
%begin
{{\begin {small}\par
{ \center  Київський національний університет імені Тараса Шевченка \par}
{\parbox {5 cm}{Освітня програма \hfill}{\parbox {7 cm}{Прикладна математика, Системний аналіз \hfill}}\parbox {6 cm}{\hfill Семестр 2}}\par
{\parbox {5 cm} {Навчальний предмет \hfill}\parbox {7 cm}{Програмування \hfill}\parbox {6cm}{\hfill}\par}\vspace{-7pt} \end{small}}{\center Екзаменаційний білет \No \textbf 
{ 364 }\par}}
{
\begin {large}
1. Що буде виведено за виконання?
code
try:
    res = 7.0j<4
    if isinstance(res, bool):
        print(f'bool;{res}')
    elif isinstance(res, int):
        print(f'int;{res}')
    elif isinstance(res, float):
        print(f'float;{res:.2f}')
    else:
        print('???')
except:
    print('ERR')
code

2. Що буде виведено за виконання?
code
a = 7
b = 3
c = 8

def g(a):
    global c
    a = 1
    b = 4
    c = 2
    return a + b + c

a = 5
b = 1
c = 4

try:
    print(g(a), a, b, c, sep=';')
except:
    print('ERR', a, b, c, sep=';')
code

3. Що буде виведено за виконання?
code
try:
    t = 9
    while t<8:
        t += 1
    print(t, end=';')
except:
    print('ERR')
code

4. Що буде виведено за виконання?
code
try:
    for f in range(-2, -9, 3):
        if f<=8:
            break
        print(f, end=';');f = 9
        if f>=3:
            break
    else:
        print('end',end=';')
    print(f)
except:
    print('ERR')
code

5. Що буде виведено за виконання?\par (Дійсні значення, якщо наявні, в цьому завданні будуть виводитись у форматі з фіксованою крапкою та, можливо, з доволі великою кількістю десяткових розрядів. У відповіді для дійсних значень у ЦЬОМУ завданні достатньо вказати один десятковий розряд після крапки, відкинувши решту розрядів без округлення. Наприклад, замість 13.66666666666667 достатньо вказати 13.6, замість 25.23 достатньо вказати 25.2)
code
try:
    print(8, end=';')
    print(int('0'), end=';')
    print(6, end=';')
except ValueError:
    print(3, end=';')
except ZeroDivisionError:
    print(2, end=';')
except:
    print('ERR', end=';')
else:
    print(7, end=';')
finally:
    print(1)
code

6. Що буде виведено за виконання?
code
def f(s):
    for el in s:
        el += 1
    return s
 
s = [1, 2, 3]
s = f(s)
print(s)
code

7. Що буде виведено за виконання?
code
class A:
    a = 1

    def _h1(self): print(2, end=';')
    def __h2(self): print(3, end=';')

    def main(self):
        try: print(self.a, end=';')
        except: print('ERR', end=';')

        try: self._h1()
        except: print('ERR', end=';')

        try: self.__h2()
        except: print('ERR', end=';')


class B(A):
    a = 4

    def _h1(self): print(5, end=';')
    def __h2(self): print(6, end=';')

try: B().main()
except: print('ERR', end=';')
code

8. Що буде виведено за виконання?
code
def f(arg, n):
    arg[28] = 8
    n -= -2

try:
    n = 5
    d = {11:1, 28:2, 37:9, 51:6}
    f(d, n)
    print(n, end=';')
    for x in d.values():
        print(x, sep=';', end=';')
except:
    print('ERR')
code

9. Що буде виведено за виконання?
code
class A:
    a = 1
    c = 2

    def __init__(self):
        a = 3

class B(A):
    a = 4

    def __init__(self):
        super().__init__()
        self.b = 7

obj = B()
obj.b = obj.c + 10
B.b = 13

try: print(obj.b, end= ';')
except: print('ERR', end=';')

try: print(A.b, end= ';')
except: print('ERR', end=';')

try: print(B.b, end= ';')
except: print('ERR', end=';')
code

10. 
Для графа на рисунку з вершини 5 запущено алгоритм пошуку в глибину, що будує кістякове дерево. Списки суміжності вершин переглядаються алгоритмом за зростанням міток вершин.\par
\begin{center}
	\adjustimage{max size={0.9\linewidth}{0.9\paperheight}}
	{../10/Traversals/206.png}
\end{center}
\par Які з наведених ребер входять до складу отриманого кістякового дерева?\par
1) (1, 6)\par
2) (3, 5)\par
3) (0, 1)\par
4) (3, 6)\par
5) (4, 6)\par
6) (0, 4)\par
{\vskip 15pt} У формі вибрати номери відповідних ребер.\par
%answer:1 3 5
\end {large}}
%end
{\smallskip \begin {small} Затверджено на засіданні кафедри теоретичної кібернетики, протокол \No 12 від   19.05.2020 р.\par\smallskip\smallskip
{\parbox {6.5 cm}{Зав. кафедри \dotfill Ю.В.~Крак}}\hspace {0.7cm}{\hbox to 6.5 cm{ Екзаменатор \dotfill Т.О.~Карнаух}} \end{small}}
%begin
{{\begin {small}\par
{ \center  Київський національний університет імені Тараса Шевченка \par}
{\parbox {5 cm}{Освітня програма \hfill}{\parbox {7 cm}{Прикладна математика, Системний аналіз \hfill}}\parbox {6 cm}{\hfill Семестр 2}}\par
{\parbox {5 cm} {Навчальний предмет \hfill}\parbox {7 cm}{Програмування \hfill}\parbox {6cm}{\hfill}\par}\vspace{-7pt} \end{small}}{\center Екзаменаційний білет \No \textbf 
{ 365 }\par}}
{
\begin {large}
1. Що буде виведено за виконання?
code
def f(n):
    print('f', end='')
    return n<=4

try:
    print(f(-9) or f(3))
except:
    print('ERR')
code

2. Що буде виведено за виконання?
code
a = 9
b = 0
c = 9

def f(b):
    a = 3
    b = 1
    c = 2
    return a + b + c

a = 1
b = 7
c = 8

try:
    print(f(a), a, b, c, sep=';')
except:
    print('ERR', a, b, c, sep=';')
code

3. Що буде виведено за виконання?
code
try:
    k = 14
    while k>2:
        k += -3
    print(k, end=';')
except:
    print('ERR')
code

4. Що буде виведено за виконання?
code
try:
    for d in range(-13, 5, -3):
        if d>=3:
            continue
        print(d, end=';')
    else:
        print(d,end=';')
    print(d)
except:
    print('ERR')
code

5. Що буде виведено за виконання?\par (Дійсні значення, якщо наявні, в цьому завданні будуть виводитись у форматі з фіксованою крапкою та, можливо, з доволі великою кількістю десяткових розрядів. У відповіді для дійсних значень у ЦЬОМУ завданні достатньо вказати один десятковий розряд після крапки, відкинувши решту розрядів без округлення. Наприклад, замість 13.66666666666667 достатньо вказати 13.6, замість 25.23 достатньо вказати 25.2)
code
try:
    print(4, end=';')
    print(9/1, end=';')
    print(6, end=';')
except Exception:
    print(4, end=';')
except TypeError:
    print(1, end=';')
except:
    print('ERR', end=';')
else:
    print(7, end=';')
finally:
    print(2)
code

6. Що буде виведено за виконання?
code
def f(s):
    for el in s:
        el += 1
    return s
 
s = [1, 2, 3]
s = f(s)
print(s)
code

7. Що буде виведено за виконання?
code
class A:
    a = 1

    def f1(self): print(2, end=';')
    def _f2(self): print(3, end=';')

    def main(self):
        try: print(self.a, end=';')
        except: print('ERR', end=';')

        try: self.f1()
        except: print('ERR', end=';')

        try: self._f2()
        except: print('ERR', end=';')


class B(A):
    a = 4

    def f1(self): print(5, end=';')
    def _f2(self): print(6, end=';')

try: B().main()
except: print('ERR', end=';')
code

8. Що буде виведено за виконання?
code
def f(arg, n):
    arg[73] = 5
    n = 8
    arg = list(arg.items())

try:
    n = 3
    t = {73:5, 57:2, 73:0}
    f(t, n)
    print(n, end=';')
    for x in t.items():
        print(x, sep=';', end=';')
except:
    print('ERR')
code

9. Що буде виведено за виконання?
code
class A:
    b = 1
    a = 2

    def __init__(self):
        a = 3

class B(A):
    b = 4

    def __init__(self):
        super().__init__()
        self.c = 7

obj = B()
obj.a = obj.a + 10
B.a = 13

try: print(obj.a, end= ';')
except: print('ERR', end=';')

try: print(A.a, end= ';')
except: print('ERR', end=';')

try: print(B.a, end= ';')
except: print('ERR', end=';')
code

10. 
Для графа на рисунку з вершини 2 запущено алгоритм пошуку в глибину, що будує кістякове дерево. Списки суміжності вершин переглядаються алгоритмом за зростанням міток вершин.\par
\begin{center}
	\adjustimage{max size={0.9\linewidth}{0.9\paperheight}}
	{../10/Traversals/102.png}
\end{center}
\par Які з наведених ребер входять до складу отриманого кістякового дерева?\par
1) (2, 3)\par
2) (2, 4)\par
3) (0, 5)\par
4) (4, 5)\par
5) (2, 5)\par
6) (0, 2)\par
{\vskip 15pt} У формі вибрати номери відповідних ребер.\par
%answer:4 6
\end {large}}
%end
{\smallskip \begin {small} Затверджено на засіданні кафедри теоретичної кібернетики, протокол \No 12 від   19.05.2020 р.\par\smallskip\smallskip
{\parbox {6.5 cm}{Зав. кафедри \dotfill Ю.В.~Крак}}\hspace {0.7cm}{\hbox to 6.5 cm{ Екзаменатор \dotfill Т.О.~Карнаух}} \end{small}}
%begin
{{\begin {small}\par
{ \center  Київський національний університет імені Тараса Шевченка \par}
{\parbox {5 cm}{Освітня програма \hfill}{\parbox {7 cm}{Прикладна математика, Системний аналіз \hfill}}\parbox {6 cm}{\hfill Семестр 2}}\par
{\parbox {5 cm} {Навчальний предмет \hfill}\parbox {7 cm}{Програмування \hfill}\parbox {6cm}{\hfill}\par}\vspace{-7pt} \end{small}}{\center Екзаменаційний білет \No \textbf 
{ 366 }\par}}
{
\begin {large}
1. Що буде виведено за виконання?
code
try:
    res = 3//7*9*8
    if isinstance(res, bool):
        print(f'bool;{res}')
    elif isinstance(res, int):
        print(f'int;{res}')
    elif isinstance(res, float):
        print(f'float;{res:.2f}')
    else:
        print('???')
except:
    print('ERR')
code

2. Що буде виведено за виконання?
code
a = 5
b = 3
c = 2

def h(b):
    global c
    a = 4
    b = 3
    c = 5
    return a + b + c

a = 6
b = 0
c = 4

try:
    print(h(b), a, b, c, sep=';')
except:
    print('ERR', a, b, c, sep=';')
code

3. Що буде виведено за виконання?
code
try:
    k = 0
    while k<=11:
        k += 1
    print(k, end=';')
except:
    print('ERR')
code

4. Що буде виведено за виконання?
code
try:
    for b in range(-10, 4, 3):
        if b>5:
            continue
        print(b, end=';')
    else:
        print(b,end=';')
    print(b)
except:
    print('ERR')
code

5. Що буде виведено за виконання?\par (Дійсні значення, якщо наявні, в цьому завданні будуть виводитись у форматі з фіксованою крапкою та, можливо, з доволі великою кількістю десяткових розрядів. У відповіді для дійсних значень у ЦЬОМУ завданні достатньо вказати один десятковий розряд після крапки, відкинувши решту розрядів без округлення. Наприклад, замість 13.66666666666667 достатньо вказати 13.6, замість 25.23 достатньо вказати 25.2)
code
try:
    print(0, end=';')
    print(int('a9'), end=';')
    print(8, end=';')
except ValueError:
    print(3, end=';')
except KeyboardInterrupt:
    print(5, end=';')
except:
    print('ERR', end=';')
else:
    print(4, end=';')
finally:
    print(1)
code

6. Що буде виведено за виконання?
code
def f(n):
    if n>9:
        f(n//100)
    print(n%10, end='')
 
f(123456)
code

7. Що буде виведено за виконання?
code
class A:
    c = 1

    def _g1(self): print(2, end=';')
    def g2(self): print(3, end=';')

    def main(self):
        try: print(self.c, end=';')
        except: print('ERR', end=';')

        try: self._g1()
        except: print('ERR', end=';')

        try: self.g2()
        except: print('ERR', end=';')


class B(A):
    c = 4

    def _g1(self): print(5, end=';')
    def g2(self): print(6, end=';')

try: B().main()
except: print('ERR', end=';')
code

8. Що буде виведено за виконання?
code
def f(arg, n):
    arg[60] = 8
    n = 7
    arg = tuple(arg.items())

try:
    n = 1
    s = {18:6, 83:8, 18:3}
    f(s, n)
    print(n, end=';')
    for x in s.values():
        print(x, sep=';', end=';')
except:
    print('ERR')
code

9. Що буде виведено за виконання?
code
class A:
    c = 1
    b = 2

    def __init__(self):
        c = 3

class B(A):
    b = 4

    def __init__(self):
        super().__init__()
        self.a = 7

obj = B()
obj.c = obj.b + 10
B.c = 13

try: print(obj.c, end= ';')
except: print('ERR', end=';')

try: print(A.c, end= ';')
except: print('ERR', end=';')

try: print(B.c, end= ';')
except: print('ERR', end=';')
code

10. 
Для графа на рисунку з вершини 2 запущено алгоритм пошуку в глибину, що будує кістякове дерево. Списки суміжності вершин переглядаються алгоритмом за зростанням міток вершин.\par
\begin{center}
	\adjustimage{max size={0.9\linewidth}{0.9\paperheight}}
	{../10/Traversals/125.png}
\end{center}
\par Які з наведених ребер входять до складу отриманого кістякового дерева?\par
1) (2, 4)\par
2) (3, 6)\par
3) (0, 5)\par
4) (4, 5)\par
5) (4, 6)\par
6) (1, 2)\par
{\vskip 15pt} У формі вибрати номери відповідних ребер.\par
%answer:2 3 4 5 6
\end {large}}
%end
{\smallskip \begin {small} Затверджено на засіданні кафедри теоретичної кібернетики, протокол \No 12 від   19.05.2020 р.\par\smallskip\smallskip
{\parbox {6.5 cm}{Зав. кафедри \dotfill Ю.В.~Крак}}\hspace {0.7cm}{\hbox to 6.5 cm{ Екзаменатор \dotfill Т.О.~Карнаух}} \end{small}}
%begin
{{\begin {small}\par
{ \center  Київський національний університет імені Тараса Шевченка \par}
{\parbox {5 cm}{Освітня програма \hfill}{\parbox {7 cm}{Прикладна математика, Системний аналіз \hfill}}\parbox {6 cm}{\hfill Семестр 2}}\par
{\parbox {5 cm} {Навчальний предмет \hfill}\parbox {7 cm}{Програмування \hfill}\parbox {6cm}{\hfill}\par}\vspace{-7pt} \end{small}}{\center Екзаменаційний білет \No \textbf 
{ 367 }\par}}
{
\begin {large}
1. Що буде виведено за виконання?
code
try:
    res = True==6 or not 3<=5 and 9.0<5.0
    if isinstance(res, bool):
        print(f'bool;{res}')
    elif isinstance(res, int):
        print(f'int;{res}')
    elif isinstance(res, float):
        print(f'float;{res:.2f}')
    else:
        print('???')
except:
    print('ERR')
code

2. Що буде виведено за виконання?
code
a = 8
b = 5
c = 3

def h(a):
    global c
    a = 3
    b -= 4
    c = 4
    return a + b + c

a = 4
b = 6
c = 2

try:
    print(h(b), a, b, c, sep=';')
except:
    print('ERR', a, b, c, sep=';')
code

3. Що буде виведено за виконання?
code
try:
    n = 11
    while n>=8:
        n += -2
    print(n, end=';')
except:
    print('ERR')
code

4. Що буде виведено за виконання?
code
try:
    for a in range(6, -1, -3):
        if a<6:
            continue
        print(a, end=';')
    else:
        print(a,end=';')
    print(a)
except:
    print('ERR')
code

5. Що буде виведено за виконання?\par (Дійсні значення, якщо наявні, в цьому завданні будуть виводитись у форматі з фіксованою крапкою та, можливо, з доволі великою кількістю десяткових розрядів. У відповіді для дійсних значень у ЦЬОМУ завданні достатньо вказати один десятковий розряд після крапки, відкинувши решту розрядів без округлення. Наприклад, замість 13.66666666666667 достатньо вказати 13.6, замість 25.23 достатньо вказати 25.2)
code
try:
    print(9, end=';')
    print(8%3, end=';')
    print(2, end=';')
except TypeError:
    print(7, end=';')
except Exception:
    print(5, end=';')
except:
    print('ERR', end=';')
else:
    print(6, end=';')
finally:
    print(3)
code

6. Що буде виведено за виконання?
code
def f(n):
    if n>9:
        f(n//100)
    print(n%10,end='')
 
f(987654)
code

7. Що буде виведено за виконання?
code
class A:
    d = 1

    def _f1(self): print(2, end=';')
    def __f2(self): print(3, end=';')

    def main(self):
        try: print(self.d, end=';')
        except: print('ERR', end=';')

        try: self._f1()
        except: print('ERR', end=';')

        try: self.__f2()
        except: print('ERR', end=';')


class B(A):
    d = 4

    def _f1(self): print(5, end=';')
    def __f2(self): print(6, end=';')

try: B().main()
except: print('ERR', end=';')
code

8. Що буде виведено за виконання?
code
def f(arg, n):
    arg[88] = 7
    n *= -3

try:
    n = 3
    s = {45:7, 86:7, 89:7, 89:7}
    f(s, n)
    print(n, end=';')
    for x in s :
        print(x, sep=';', end=';')
except:
    print('ERR')
code

9. Що буде виведено за виконання?
code
class A:
    c = 1
    b = 2

    def __init__(self):
        b = 3

class B(A):
    c = 4

    def __init__(self):
        super().__init__()
        self.a = 7

obj = B()
obj.b = obj.a + 10
B.b = 13

try: print(obj.b, end= ';')
except: print('ERR', end=';')

try: print(A.b, end= ';')
except: print('ERR', end=';')

try: print(B.b, end= ';')
except: print('ERR', end=';')
code

10. 
Для графа на рисунку з вершини 5 запущено алгоритм пошуку в глибину, що будує кістякове дерево. Списки суміжності вершин переглядаються алгоритмом за зростанням міток вершин.\par
\begin{center}
	\adjustimage{max size={0.9\linewidth}{0.9\paperheight}}
	{../10/Traversals/181.png}
\end{center}
\par Які з наведених ребер входять до складу отриманого кістякового дерева?\par
1) (3, 5)\par
2) (1, 4)\par
3) (0, 5)\par
4) (2, 3)\par
5) (1, 6)\par
6) (2, 4)\par
{\vskip 15pt} У формі вибрати номери відповідних ребер.\par
%answer:2 3 4 5 6
\end {large}}
%end
{\smallskip \begin {small} Затверджено на засіданні кафедри теоретичної кібернетики, протокол \No 12 від   19.05.2020 р.\par\smallskip\smallskip
{\parbox {6.5 cm}{Зав. кафедри \dotfill Ю.В.~Крак}}\hspace {0.7cm}{\hbox to 6.5 cm{ Екзаменатор \dotfill Т.О.~Карнаух}} \end{small}}
%begin
{{\begin {small}\par
{ \center  Київський національний університет імені Тараса Шевченка \par}
{\parbox {5 cm}{Освітня програма \hfill}{\parbox {7 cm}{Прикладна математика, Системний аналіз \hfill}}\parbox {6 cm}{\hfill Семестр 2}}\par
{\parbox {5 cm} {Навчальний предмет \hfill}\parbox {7 cm}{Програмування \hfill}\parbox {6cm}{\hfill}\par}\vspace{-7pt} \end{small}}{\center Екзаменаційний білет \No \textbf 
{ 368 }\par}}
{
\begin {large}
1. Що буде виведено за виконання?
code
try:
    res = 8!=4 and not 9<7.0 and False==4.0
    if isinstance(res, bool):
        print(f'bool;{res}')
    elif isinstance(res, int):
        print(f'int;{res}')
    elif isinstance(res, float):
        print(f'float;{res:.2f}')
    else:
        print('???')
except:
    print('ERR')
code

2. Що буде виведено за виконання?
code
a = 7
b = 1
c = 5

def g(a):
    a = 2
    b = 1
    c = 3
    return a + b + c

a = 8
b = 6
c = 3

try:
    print(g(b), a, b, c, sep=';')
except:
    print('ERR', a, b, c, sep=';')
code

3. Що буде виведено за виконання?
code
try:
    j = 15
    while j>=5:
        j += -2
    print(j, end=';')
except:
    print('ERR')
code

4. Що буде виведено за виконання?
code
try:
    for b in range(0, -3, -1):
        if b<=5:
            break
        print(b, end=';');b = -5
    else:
        print(b,end=';')
    print(b)
except:
    print('ERR')
code

5. Що буде виведено за виконання?\par (Дійсні значення, якщо наявні, в цьому завданні будуть виводитись у форматі з фіксованою крапкою та, можливо, з доволі великою кількістю десяткових розрядів. У відповіді для дійсних значень у ЦЬОМУ завданні достатньо вказати один десятковий розряд після крапки, відкинувши решту розрядів без округлення. Наприклад, замість 13.66666666666667 достатньо вказати 13.6, замість 25.23 достатньо вказати 25.2)
code
try:
    print(0, end=';')
    print(1/0, end=';')
    print(5, end=';')
except ZeroDivisionError:
    print(6, end=';')
except Exception:
    print(0, end=';')
except:
    print('ERR', end=';')
else:
    print(8, end=';')
finally:
    print(4)
code

6. Що буде виведено за виконання?
code
def f(n):
    if n>9:
        return f(n//100)+1
    else:
        return 1

print(f(123456))
code

7. Що буде виведено за виконання?
code
class A:
    b = 1

    def __h1(self): print(2, end=';')
    def _h2(self): print(3, end=';')

    def main(self):
        try: print(self.b, end=';')
        except: print('ERR', end=';')

        try: self.__h1()
        except: print('ERR', end=';')

        try: self._h2()
        except: print('ERR', end=';')


class B(A):
    b = 4

    def __h1(self): print(5, end=';')
    def _h2(self): print(6, end=';')

try: B().main()
except: print('ERR', end=';')
code

8. Що буде виведено за виконання?
code
def f(arg, n):
    arg[15] = 3
    n -= 3

try:
    n = 4
    s = {65:8, 22:7, 65:9}
    f(s, n)
    print(n, end=';')
    for x in s :
        print(x, sep=';', end=';')
except:
    print('ERR')
code

9. Що буде виведено за виконання?
code
class A:
    b = 1
    c = 2

    def __init__(self):
        b = 3

class B(A):
    c = 4

    def __init__(self):
        super().__init__()
        self.a = 7

obj = B()
obj.c = obj.b + 10
B.c = 13

try: print(obj.c, end= ';')
except: print('ERR', end=';')

try: print(A.c, end= ';')
except: print('ERR', end=';')

try: print(B.c, end= ';')
except: print('ERR', end=';')
code

10. 
Для графа на рисунку з вершини 4 запущено алгоритм пошуку в глибину, що будує кістякове дерево. Списки суміжності вершин переглядаються алгоритмом за зростанням міток вершин.\par
\begin{center}
	\adjustimage{max size={0.9\linewidth}{0.9\paperheight}}
	{../10/Traversals/222.png}
\end{center}
\par Які з наведених ребер входять до складу отриманого кістякового дерева?\par
1) (1, 4)\par
2) (3, 4)\par
3) (3, 5)\par
4) (0, 6)\par
5) (1, 5)\par
6) (1, 6)\par
{\vskip 15pt} У формі вибрати номери відповідних ребер.\par
%answer:1 3 4
\end {large}}
%end
{\smallskip \begin {small} Затверджено на засіданні кафедри теоретичної кібернетики, протокол \No 12 від   19.05.2020 р.\par\smallskip\smallskip
{\parbox {6.5 cm}{Зав. кафедри \dotfill Ю.В.~Крак}}\hspace {0.7cm}{\hbox to 6.5 cm{ Екзаменатор \dotfill Т.О.~Карнаух}} \end{small}}
%begin
{{\begin {small}\par
{ \center  Київський національний університет імені Тараса Шевченка \par}
{\parbox {5 cm}{Освітня програма \hfill}{\parbox {7 cm}{Прикладна математика, Системний аналіз \hfill}}\parbox {6 cm}{\hfill Семестр 2}}\par
{\parbox {5 cm} {Навчальний предмет \hfill}\parbox {7 cm}{Програмування \hfill}\parbox {6cm}{\hfill}\par}\vspace{-7pt} \end{small}}{\center Екзаменаційний білет \No \textbf 
{ 369 }\par}}
{
\begin {large}
1. Що буде виведено за виконання?
code
def f(n):
    print('f', end='')
    return n

try:
    print(f(6) and f(0))
except:
    print('ERR')
code

2. Що буде виведено за виконання?
code
a = 5
b = 8
c = 4

def f(a):
    a += 5
    b = 3
    c = 1
    return a + b + c

a = 0
b = 3
c = 7

try:
    print(f(a), a, b, c, sep=';')
except:
    print('ERR', a, b, c, sep=';')
code

3. Що буде виведено за виконання?
code
try:
    j = 5
    while j<9:
        j += 2
    print(j, end=';')
except:
    print('ERR')
code

4. Що буде виведено за виконання?
code
try:
    for e in range(1, -4, -2):
        if e>=7:
            continue
        print(e, end=';')
        if e<8:
            break
    else:
        print(e,end=';')
    print(e)
except:
    print('ERR')
code

5. Що буде виведено за виконання?\par (Дійсні значення, якщо наявні, в цьому завданні будуть виводитись у форматі з фіксованою крапкою та, можливо, з доволі великою кількістю десяткових розрядів. У відповіді для дійсних значень у ЦЬОМУ завданні достатньо вказати один десятковий розряд після крапки, відкинувши решту розрядів без округлення. Наприклад, замість 13.66666666666667 достатньо вказати 13.6, замість 25.23 достатньо вказати 25.2)
code
try:
    print(7, end=';')
    print(2j==4j, end=';')
    print(9, end=';')
except TypeError:
    print(3, end=';')
except KeyboardInterrupt:
    print(3, end=';')
except:
    print('ERR', end=';')
else:
    print(4, end=';')
finally:
    print(5)
code

6. Що буде виведено за виконання?
code
def f(n):
    if n>9:
        f(n//10)
    print(n%10, end='')
 
f(123456)
code

7. Що буде виведено за виконання?
code
class A:
    b = 1

    def g1(self): print(2, end=';')
    def _g2(self): print(3, end=';')

    def main(self):
        try: print(self.b, end=';')
        except: print('ERR', end=';')

        try: self.g1()
        except: print('ERR', end=';')

        try: self._g2()
        except: print('ERR', end=';')


class B(A):
    b = 4

    def g1(self): print(5, end=';')
    def _g2(self): print(6, end=';')

try: B().main()
except: print('ERR', end=';')
code

8. Що буде виведено за виконання?
code
def f(arg, n):
    arg[32] = 7
    n = 6
    arg = list(arg.keys())

try:
    n = 4
    t = {88:6, 63:3, 24:9, 63:5}
    f(t, n)
    print(n, end=';')
    for x in t.items():
        print(*x, sep=';', end=';')
except:
    print('ERR')
code

9. Що буде виведено за виконання?
code
class A:
    a = 1
    b = 2

    def __init__(self):
        a = 3

class B(A):
    b = 4

    def __init__(self):
        super().__init__()
        self.c = 7

obj = B()
obj.c = obj.a + 10
B.c = 13

try: print(obj.c, end= ';')
except: print('ERR', end=';')

try: print(A.c, end= ';')
except: print('ERR', end=';')

try: print(B.c, end= ';')
except: print('ERR', end=';')
code

10. 
Для графа на рисунку з вершини 4 запущено алгоритм пошуку в глибину, що будує кістякове дерево. Списки суміжності вершин переглядаються алгоритмом за зростанням міток вершин.\par
\begin{center}
	\adjustimage{max size={0.9\linewidth}{0.9\paperheight}}
	{../10/Traversals/210.png}
\end{center}
\par Які з наведених ребер входять до складу отриманого кістякового дерева?\par
1) (3, 4)\par
2) (3, 6)\par
3) (0, 1)\par
4) (0, 2)\par
5) (2, 5)\par
6) (0, 4)\par
{\vskip 15pt} У формі вибрати номери відповідних ребер.\par
%answer:2 3 5 6
\end {large}}
%end
{\smallskip \begin {small} Затверджено на засіданні кафедри теоретичної кібернетики, протокол \No 12 від   19.05.2020 р.\par\smallskip\smallskip
{\parbox {6.5 cm}{Зав. кафедри \dotfill Ю.В.~Крак}}\hspace {0.7cm}{\hbox to 6.5 cm{ Екзаменатор \dotfill Т.О.~Карнаух}} \end{small}}
%begin
{{\begin {small}\par
{ \center  Київський національний університет імені Тараса Шевченка \par}
{\parbox {5 cm}{Освітня програма \hfill}{\parbox {7 cm}{Прикладна математика, Системний аналіз \hfill}}\parbox {6 cm}{\hfill Семестр 2}}\par
{\parbox {5 cm} {Навчальний предмет \hfill}\parbox {7 cm}{Програмування \hfill}\parbox {6cm}{\hfill}\par}\vspace{-7pt} \end{small}}{\center Екзаменаційний білет \No \textbf 
{ 370 }\par}}
{
\begin {large}
1. Що буде виведено за виконання?
code
try:
    res = False>"1j"
    if isinstance(res, bool):
        print(f'bool;{res}')
    elif isinstance(res, int):
        print(f'int;{res}')
    elif isinstance(res, float):
        print(f'float;{res:.2f}')
    else:
        print('???')
except:
    print('ERR')
code

2. Що буде виведено за виконання?
code
a = 0
b = 9
c = 4

def g(b):
    a = 4
    b = 5
    c = 1
    return a + b + c

a = 2
b = 6
c = 5

try:
    print(g(b), a, b, c, sep=';')
except:
    print('ERR', a, b, c, sep=';')
code

3. Що буде виведено за виконання?
code
try:
    i = 10
    while i>9:
        i += -3
    print(i, end=';')
except:
    print('ERR')
code

4. Що буде виведено за виконання?
code
try:
    for f in range(-14, 11, 2):
        if f<8:
            continue
        print(f, end=';');f = -2
    else:
        print('end',end=';')
    print(f)
except:
    print('ERR')
code

5. Що буде виведено за виконання?\par (Дійсні значення, якщо наявні, в цьому завданні будуть виводитись у форматі з фіксованою крапкою та, можливо, з доволі великою кількістю десяткових розрядів. У відповіді для дійсних значень у ЦЬОМУ завданні достатньо вказати один десятковий розряд після крапки, відкинувши решту розрядів без округлення. Наприклад, замість 13.66666666666667 достатньо вказати 13.6, замість 25.23 достатньо вказати 25.2)
code
try:
    print(8, end=';')
    print(int('7'), end=';')
    print(2, end=';')
except ZeroDivisionError:
    print(6, end=';')
except ValueError:
    print(0, end=';')
except:
    print('ERR', end=';')
else:
    print(9, end=';')
finally:
    print(1)
code

6. Що буде виведено за виконання?
code
def f(n):
    if n>9:
        f(n//10)
    else:
        print(n%10, end='')

f(543216)
code

7. Що буде виведено за виконання?
code
class A:
    c = 1

    def _f1(self): print(2, end=';')
    def _f2(self): print(3, end=';')

    def main(self):
        try: print(self.c, end=';')
        except: print('ERR', end=';')

        try: self._f1()
        except: print('ERR', end=';')

        try: self._f2()
        except: print('ERR', end=';')


class B(A):
    c = 4

    def _f1(self): print(5, end=';')
    def _f2(self): print(6, end=';')

try: B().main()
except: print('ERR', end=';')
code

8. Що буде виведено за виконання?
code
def f(arg, n):
    arg[35] = 2
    n = 6
    arg = set(arg.values())

try:
    n = 0
    d = {49:5, 35:1, 57:8, 35:2}
    f(d, n)
    print(n, end=';')
    for x, y in d.items():
        print(x, y, sep=';', end=';')
except:
    print('ERR')
code

9. Що буде виведено за виконання?
code
class A:
    a = 1
    c = 2

    def __init__(self):
        c = 3

class B(A):
    a = 4

    def __init__(self):
        super().__init__()
        self.b = 7

obj = B()
obj.c = obj.b + 10
B.c = 13

try: print(obj.c, end= ';')
except: print('ERR', end=';')

try: print(A.c, end= ';')
except: print('ERR', end=';')

try: print(B.c, end= ';')
except: print('ERR', end=';')
code

10. 
Для графа на рисунку з вершини 1 запущено алгоритм пошуку в глибину, що будує кістякове дерево. Списки суміжності вершин переглядаються алгоритмом за зростанням міток вершин.\par
\begin{center}
	\adjustimage{max size={0.9\linewidth}{0.9\paperheight}}
	{../10/Traversals/55.png}
\end{center}
\par Які з наведених ребер входять до складу отриманого кістякового дерева?\par
1) (2, 6)\par
2) (0, 2)\par
3) (4, 6)\par
4) (3, 6)\par
5) (2, 5)\par
6) (2, 4)\par
{\vskip 15pt} У формі вибрати номери відповідних ребер.\par
%answer:2 3 6
\end {large}}
%end
{\smallskip \begin {small} Затверджено на засіданні кафедри теоретичної кібернетики, протокол \No 12 від   19.05.2020 р.\par\smallskip\smallskip
{\parbox {6.5 cm}{Зав. кафедри \dotfill Ю.В.~Крак}}\hspace {0.7cm}{\hbox to 6.5 cm{ Екзаменатор \dotfill Т.О.~Карнаух}} \end{small}}
%begin
{{\begin {small}\par
{ \center  Київський національний університет імені Тараса Шевченка \par}
{\parbox {5 cm}{Освітня програма \hfill}{\parbox {7 cm}{Прикладна математика, Системний аналіз \hfill}}\parbox {6 cm}{\hfill Семестр 2}}\par
{\parbox {5 cm} {Навчальний предмет \hfill}\parbox {7 cm}{Програмування \hfill}\parbox {6cm}{\hfill}\par}\vspace{-7pt} \end{small}}{\center Екзаменаційний білет \No \textbf 
{ 371 }\par}}
{
\begin {large}
1. Що буде виведено за виконання?
code
try:
    res = 4**4+2
    if isinstance(res, bool):
        print(f'bool;{res}')
    elif isinstance(res, int):
        print(f'int;{res}')
    elif isinstance(res, float):
        print(f'float;{res:.2f}')
    else:
        print('???')
except:
    print('ERR')
code

2. Що буде виведено за виконання?
code
a = 0
b = 3
c = 1

def h(a):
    global c
    a = 4
    b = 3
    c = 5
    return a + b + c

a = 2
b = 8
c = 9

try:
    print(h(a), a, b, c, sep=';')
except:
    print('ERR', a, b, c, sep=';')
code

3. Що буде виведено за виконання?
code
try:
    n = 8
    while n>1:
        n += -2
    print(n, end=';')
except:
    print('ERR')
code

4. Що буде виведено за виконання?
code
try:
    for d in range(-3, -14, 2):
        if d<=4:
            break
        print(d, end=';')
    else:
        print(d,end=';')
    print(d)
except:
    print('ERR')
code

5. Що буде виведено за виконання?\par (Дійсні значення, якщо наявні, в цьому завданні будуть виводитись у форматі з фіксованою крапкою та, можливо, з доволі великою кількістю десяткових розрядів. У відповіді для дійсних значень у ЦЬОМУ завданні достатньо вказати один десятковий розряд після крапки, відкинувши решту розрядів без округлення. Наприклад, замість 13.66666666666667 достатньо вказати 13.6, замість 25.23 достатньо вказати 25.2)
code
try:
    print(6, end=';')
    print(8//0.0, end=';')
    print(4, end=';')
except ZeroDivisionError:
    print(1, end=';')
except KeyboardInterrupt:
    print(2, end=';')
except:
    print('ERR', end=';')
else:
    print(7, end=';')
finally:
    print(0)
code

6. Що буде виведено за виконання?
code
def f(n):
    if n>9:
        f(n//100)
        print(n%10, end='')
 
f(123456)
code

7. Що буде виведено за виконання?
code
class A:
    d = 1

    def h1(self): print(2, end=';')
    def _h2(self): print(3, end=';')

    def main(self):
        try: print(self.d, end=';')
        except: print('ERR', end=';')

        try: self.h1()
        except: print('ERR', end=';')

        try: self._h2()
        except: print('ERR', end=';')


class B(A):
    d = 4

    def h1(self): print(5, end=';')
    def _h2(self): print(6, end=';')

try: B().main()
except: print('ERR', end=';')
code

8. Що буде виведено за виконання?
code
def f(arg, n):
    arg[50] = 5
    n *= 2

try:
    n = 7
    d = {30:6, 72:6, 32:4, 81:0}
    f(d, n)
    print(n, end=';')
    for x in d.items():
        print(*x, sep=';', end=';')
except:
    print('ERR')
code

9. Що буде виведено за виконання?
code
class A:
    c = 1
    a = 2

    def __init__(self):
        c = 3

class B(A):
    a = 4

    def __init__(self):
        super().__init__()
        self.b = 7

obj = B()
obj.a = obj.a + 10
B.a = 13

try: print(obj.a, end= ';')
except: print('ERR', end=';')

try: print(A.a, end= ';')
except: print('ERR', end=';')

try: print(B.a, end= ';')
except: print('ERR', end=';')
code

10. 
Для графа на рисунку з вершини 5 запущено алгоритм пошуку в глибину, що будує кістякове дерево. Списки суміжності вершин переглядаються алгоритмом за зростанням міток вершин.\par
\begin{center}
	\adjustimage{max size={0.9\linewidth}{0.9\paperheight}}
	{../10/Traversals/196.png}
\end{center}
\par Які з наведених ребер входять до складу отриманого кістякового дерева?\par
1) (1, 5)\par
2) (0, 2)\par
3) (4, 5)\par
4) (3, 6)\par
5) (2, 5)\par
6) (0, 6)\par
{\vskip 15pt} У формі вибрати номери відповідних ребер.\par
%answer:2 4 6
\end {large}}
%end
{\smallskip \begin {small} Затверджено на засіданні кафедри теоретичної кібернетики, протокол \No 12 від   19.05.2020 р.\par\smallskip\smallskip
{\parbox {6.5 cm}{Зав. кафедри \dotfill Ю.В.~Крак}}\hspace {0.7cm}{\hbox to 6.5 cm{ Екзаменатор \dotfill Т.О.~Карнаух}} \end{small}}
%begin
{{\begin {small}\par
{ \center  Київський національний університет імені Тараса Шевченка \par}
{\parbox {5 cm}{Освітня програма \hfill}{\parbox {7 cm}{Прикладна математика, Системний аналіз \hfill}}\parbox {6 cm}{\hfill Семестр 2}}\par
{\parbox {5 cm} {Навчальний предмет \hfill}\parbox {7 cm}{Програмування \hfill}\parbox {6cm}{\hfill}\par}\vspace{-7pt} \end{small}}{\center Екзаменаційний білет \No \textbf 
{ 372 }\par}}
{
\begin {large}
1. Що буде виведено за виконання?
code
try:
    res = 6j!=2.0j
    if isinstance(res, bool):
        print(f'bool;{res}')
    elif isinstance(res, int):
        print(f'int;{res}')
    elif isinstance(res, float):
        print(f'float;{res:.2f}')
    else:
        print('???')
except:
    print('ERR')
code

2. Що буде виведено за виконання?
code
a = 9
b = 1
c = 9

def h(b):
    global c
    a = 2
    b -= 4
    c = 2
    return a + b + c

a = 7
b = 0
c = 6

try:
    print(h(a), a, b, c, sep=';')
except:
    print('ERR', a, b, c, sep=';')
code

3. Що буде виведено за виконання?
code
try:
    n = 8
    while n<=13:
        n += 1
    print(n, end=';')
except:
    print('ERR')
code

4. Що буде виведено за виконання?
code
try:
    for c in range(-8, 0, -3):
        if c>3:
            continue
        print(c, end=';')
    else:
        print(c,end=';')
    print(c)
except:
    print('ERR')
code

5. Що буде виведено за виконання?\par (Дійсні значення, якщо наявні, в цьому завданні будуть виводитись у форматі з фіксованою крапкою та, можливо, з доволі великою кількістю десяткових розрядів. У відповіді для дійсних значень у ЦЬОМУ завданні достатньо вказати один десятковий розряд після крапки, відкинувши решту розрядів без округлення. Наприклад, замість 13.66666666666667 достатньо вказати 13.6, замість 25.23 достатньо вказати 25.2)
code
try:
    print(3, end=';')
    print(int('5'), end=';')
    print(9, end=';')
except ValueError:
    print(6, end=';')
except TypeError:
    print(8, end=';')
except:
    print('ERR', end=';')
else:
    print(3, end=';')
finally:
    print(0)
code

6. Що буде виведено за виконання?
code
def f(n):
    print(n%10, end='')
    if n>9:
        f(n//100)
 
f(987654)
code

7. Що буде виведено за виконання?
code
class A:
    a = 1

    def _g1(self): print(2, end=';')
    def _g2(self): print(3, end=';')

    def main(self):
        try: print(self.a, end=';')
        except: print('ERR', end=';')

        try: self._g1()
        except: print('ERR', end=';')

        try: self._g2()
        except: print('ERR', end=';')


class B(A):
    a = 4

    def _g1(self): print(5, end=';')
    def _g2(self): print(6, end=';')

try: B().main()
except: print('ERR', end=';')
code

8. Що буде виведено за виконання?
code
def f(arg, n):
    n = 5
    tmp = arg
    tmp[-3:-9] = 'ac'
    return len(tmp)>3

try:
    f = [10,11,12,13,14]
    n = 8
    f(f, n)
    print(n, end=';')
    for el in f:
        print(el, end=';')
except:
    print('ERR')
code

9. Що буде виведено за виконання?
code
class A:
    b = 1
    a = 2

    def __init__(self):
        a = 3

class B(A):
    b = 4

    def __init__(self):
        super().__init__()
        self.c = 7

obj = B()
obj.c = obj.b + 10
B.c = 13

try: print(obj.c, end= ';')
except: print('ERR', end=';')

try: print(A.c, end= ';')
except: print('ERR', end=';')

try: print(B.c, end= ';')
except: print('ERR', end=';')
code

10. 
Для графа на рисунку з вершини 1 запущено алгоритм пошуку в глибину, що будує кістякове дерево. Списки суміжності вершин переглядаються алгоритмом за зростанням міток вершин.\par
\begin{center}
	\adjustimage{max size={0.9\linewidth}{0.9\paperheight}}
	{../10/Traversals/212.png}
\end{center}
\par Які з наведених ребер входять до складу отриманого кістякового дерева?\par
1) (2, 3)\par
2) (0, 1)\par
3) (4, 6)\par
4) (3, 4)\par
5) (3, 6)\par
6) (4, 5)\par
{\vskip 15pt} У формі вибрати номери відповідних ребер.\par
%answer:1 2 4 6
\end {large}}
%end
{\smallskip \begin {small} Затверджено на засіданні кафедри теоретичної кібернетики, протокол \No 12 від   19.05.2020 р.\par\smallskip\smallskip
{\parbox {6.5 cm}{Зав. кафедри \dotfill Ю.В.~Крак}}\hspace {0.7cm}{\hbox to 6.5 cm{ Екзаменатор \dotfill Т.О.~Карнаух}} \end{small}}
%begin
{{\begin {small}\par
{ \center  Київський національний університет імені Тараса Шевченка \par}
{\parbox {5 cm}{Освітня програма \hfill}{\parbox {7 cm}{Прикладна математика, Системний аналіз \hfill}}\parbox {6 cm}{\hfill Семестр 2}}\par
{\parbox {5 cm} {Навчальний предмет \hfill}\parbox {7 cm}{Програмування \hfill}\parbox {6cm}{\hfill}\par}\vspace{-7pt} \end{small}}{\center Екзаменаційний білет \No \textbf 
{ 373 }\par}}
{
\begin {large}
1. Що буде виведено за виконання?
code
try:
    res = 4/5%3*5
    if isinstance(res, bool):
        print(f'bool;{res}')
    elif isinstance(res, int):
        print(f'int;{res}')
    elif isinstance(res, float):
        print(f'float;{res:.2f}')
    else:
        print('???')
except:
    print('ERR')
code

2. Що буде виведено за виконання?
code
a = 2
b = 8
c = 1

def f(b):
    a *= 1
    b = 3
    c = 5
    return a + b + c

a = 3
b = 4
c = 5

try:
    print(f(b), a, b, c, sep=';')
except:
    print('ERR', a, b, c, sep=';')
code

3. Що буде виведено за виконання?
code
try:
    m = 15
    while m>=4:
        m += -3
    print(m, end=';')
except:
    print('ERR')
code

4. Що буде виведено за виконання?
code
try:
    for f in range(-9, -6, -1):
        if f>=5:
            continue
        print(f, end=';')
    else:
        print(f,end=';')
    print(f)
except:
    print('ERR')
code

5. Що буде виведено за виконання?\par (Дійсні значення, якщо наявні, в цьому завданні будуть виводитись у форматі з фіксованою крапкою та, можливо, з доволі великою кількістю десяткових розрядів. У відповіді для дійсних значень у ЦЬОМУ завданні достатньо вказати один десятковий розряд після крапки, відкинувши решту розрядів без округлення. Наприклад, замість 13.66666666666667 достатньо вказати 13.6, замість 25.23 достатньо вказати 25.2)
code
try:
    print(5, end=';')
    print(9%2, end=';')
    print(7, end=';')
except Exception:
    print(1, end=';')
except ZeroDivisionError:
    print(2, end=';')
except:
    print('ERR', end=';')
else:
    print(4, end=';')
finally:
    print(2)
code

6. Що буде виведено за виконання?
code
def f(n):
    if n>9:
        f(n//100)
    else:
        print(n%10,end='')

f(987654)
code

7. Що буде виведено за виконання?
code
class A:
    d = 1

    def _h1(self): print(2, end=';')
    def h2(self): print(3, end=';')

    def main(self):
        try: print(self.d, end=';')
        except: print('ERR', end=';')

        try: self._h1()
        except: print('ERR', end=';')

        try: self.h2()
        except: print('ERR', end=';')


class B(A):
    d = 4

    def _h1(self): print(5, end=';')
    def h2(self): print(6, end=';')

try: B().main()
except: print('ERR', end=';')
code

8. Що буде виведено за виконання?
code
def f(arg, n):
    arg[75] = 0
    n += 3

try:
    n = 7
    t = {10:0, 75:6, 10:5}
    f(t, n)
    print(n, end=';')
    for x in t.values():
        print(x, sep=';', end=';')
except:
    print('ERR')
code

9. Що буде виведено за виконання?
code
class A:
    c = 1
    b = 2

    def __init__(self):
        b = 3

class B(A):
    b = 4

    def __init__(self):
        super().__init__()
        self.a = 7

obj = B()
obj.a = obj.b + 10
B.a = 13

try: print(obj.a, end= ';')
except: print('ERR', end=';')

try: print(A.a, end= ';')
except: print('ERR', end=';')

try: print(B.a, end= ';')
except: print('ERR', end=';')
code

10. 
Для графа на рисунку з вершини 4 запущено алгоритм пошуку в глибину, що будує кістякове дерево. Списки суміжності вершин переглядаються алгоритмом за зростанням міток вершин.\par
\begin{center}
	\adjustimage{max size={0.9\linewidth}{0.9\paperheight}}
	{../10/Traversals/162.png}
\end{center}
\par Які з наведених ребер входять до складу отриманого кістякового дерева?\par
1) (3, 6)\par
2) (1, 3)\par
3) (5, 6)\par
4) (1, 4)\par
5) (2, 4)\par
6) (0, 3)\par
{\vskip 15pt} У формі вибрати номери відповідних ребер.\par
%answer:2 3 6
\end {large}}
%end
{\smallskip \begin {small} Затверджено на засіданні кафедри теоретичної кібернетики, протокол \No 12 від   19.05.2020 р.\par\smallskip\smallskip
{\parbox {6.5 cm}{Зав. кафедри \dotfill Ю.В.~Крак}}\hspace {0.7cm}{\hbox to 6.5 cm{ Екзаменатор \dotfill Т.О.~Карнаух}} \end{small}}
%begin
{{\begin {small}\par
{ \center  Київський національний університет імені Тараса Шевченка \par}
{\parbox {5 cm}{Освітня програма \hfill}{\parbox {7 cm}{Прикладна математика, Системний аналіз \hfill}}\parbox {6 cm}{\hfill Семестр 2}}\par
{\parbox {5 cm} {Навчальний предмет \hfill}\parbox {7 cm}{Програмування \hfill}\parbox {6cm}{\hfill}\par}\vspace{-7pt} \end{small}}{\center Екзаменаційний білет \No \textbf 
{ 374 }\par}}
{
\begin {large}
1. Що буде виведено за виконання?
code
try:
    res = 6<=2 or not True>7 or 3.0>=8.0
    if isinstance(res, bool):
        print(f'bool;{res}')
    elif isinstance(res, int):
        print(f'int;{res}')
    elif isinstance(res, float):
        print(f'float;{res:.2f}')
    else:
        print('???')
except:
    print('ERR')
code

2. Що буде виведено за виконання?
code
a = 4
b = 0
c = 5

def g(a):
    global c
    a += 3
    b = 5
    c = 1
    return a + b + c

a = 2
b = 3
c = 7

try:
    print(g(a), a, b, c, sep=';')
except:
    print('ERR', a, b, c, sep=';')
code

3. Що буде виведено за виконання?
code
try:
    i = 3
    while i<6:
        i += 2
    print(i, end=';')
except:
    print('ERR')
code

4. Що буде виведено за виконання?
code
try:
    for c in range(0, 6, -2):
        if c<4:
            continue
        print(c, end=';')
    else:
        print('end',end=';')
    print(c)
except:
    print('ERR')
code

5. Що буде виведено за виконання?\par (Дійсні значення, якщо наявні, в цьому завданні будуть виводитись у форматі з фіксованою крапкою та, можливо, з доволі великою кількістю десяткових розрядів. У відповіді для дійсних значень у ЦЬОМУ завданні достатньо вказати один десятковий розряд після крапки, відкинувши решту розрядів без округлення. Наприклад, замість 13.66666666666667 достатньо вказати 13.6, замість 25.23 достатньо вказати 25.2)
code
try:
    print(8, end=';')
    print(int('c0'), end=';')
    print(5, end=';')
except ValueError:
    print(1, end=';')
except Exception:
    print(9, end=';')
except:
    print('ERR', end=';')
else:
    print(7, end=';')
finally:
    print(3)
code

6. Що буде виведено за виконання?
code
def f(n):
    print(n%10, end='')
    if n>9:
        f(n//10)
 
f(123456)
code

7. Що буде виведено за виконання?
code
class A:
    a = 1

    def _g1(self): print(2, end=';')
    def __g2(self): print(3, end=';')

    def main(self):
        try: print(self.a, end=';')
        except: print('ERR', end=';')

        try: self._g1()
        except: print('ERR', end=';')

        try: self.__g2()
        except: print('ERR', end=';')


class B(A):
    a = 4

    def _g1(self): print(5, end=';')
    def __g2(self): print(6, end=';')

try: B().main()
except: print('ERR', end=';')
code

8. Що буде виведено за виконання?
code
def f(arg, n):
    arg[25] = 2
    n *= -2

try:
    n = 3
    t = {72:2, 25:4, 25:1}
    f(t, n)
    print(n, end=';')
    for x, y in t.items():
        print(x, y, sep=';', end=';')
except:
    print('ERR')
code

9. Що буде виведено за виконання?
code
class A:
    a = 1
    c = 2

    def __init__(self):
        a = 3

class B(A):
    a = 4

    def __init__(self):
        super().__init__()
        self.b = 7

obj = B()
obj.c = obj.b + 10
B.c = 13

try: print(obj.c, end= ';')
except: print('ERR', end=';')

try: print(A.c, end= ';')
except: print('ERR', end=';')

try: print(B.c, end= ';')
except: print('ERR', end=';')
code

10. 
Для графа на рисунку з вершини 6 запущено алгоритм пошуку в глибину, що будує кістякове дерево. Списки суміжності вершин переглядаються алгоритмом за зростанням міток вершин.\par
\begin{center}
	\adjustimage{max size={0.9\linewidth}{0.9\paperheight}}
	{../10/Traversals/26.png}
\end{center}
\par Які з наведених ребер входять до складу отриманого кістякового дерева?\par
1) (3, 6)\par
2) (0, 6)\par
3) (4, 5)\par
4) (0, 1)\par
5) (5, 6)\par
6) (1, 6)\par
{\vskip 15pt} У формі вибрати номери відповідних ребер.\par
%answer:2 3 4
\end {large}}
%end
{\smallskip \begin {small} Затверджено на засіданні кафедри теоретичної кібернетики, протокол \No 12 від   19.05.2020 р.\par\smallskip\smallskip
{\parbox {6.5 cm}{Зав. кафедри \dotfill Ю.В.~Крак}}\hspace {0.7cm}{\hbox to 6.5 cm{ Екзаменатор \dotfill Т.О.~Карнаух}} \end{small}}
%begin
{{\begin {small}\par
{ \center  Київський національний університет імені Тараса Шевченка \par}
{\parbox {5 cm}{Освітня програма \hfill}{\parbox {7 cm}{Прикладна математика, Системний аналіз \hfill}}\parbox {6 cm}{\hfill Семестр 2}}\par
{\parbox {5 cm} {Навчальний предмет \hfill}\parbox {7 cm}{Програмування \hfill}\parbox {6cm}{\hfill}\par}\vspace{-7pt} \end{small}}{\center Екзаменаційний білет \No \textbf 
{ 375 }\par}}
{
\begin {large}
1. Що буде виведено за виконання?
code
try:
    res = -2**2.0+-2
    if isinstance(res, bool):
        print(f'bool;{res}')
    elif isinstance(res, int):
        print(f'int;{res}')
    elif isinstance(res, float):
        print(f'float;{res:.2f}')
    else:
        print('???')
except:
    print('ERR')
code

2. Що буде виведено за виконання?
code
a = 4
b = 7
c = 3

def f(b):
    a = 2
    b = 5
    c = 2
    return a + b + c

a = 1
b = 8
c = 5

try:
    print(f(a), a, b, c, sep=';')
except:
    print('ERR', a, b, c, sep=';')
code

3. Що буде виведено за виконання?
code
try:
    i = 10
    while i>7:
        i += -3
    print(i, end=';')
except:
    print('ERR')
code

4. Що буде виведено за виконання?
code
try:
    for e in range(12, 0, 3):
        if e>8:
            break
        print(e, end=';');e = -4
    else:
        print(e,end=';')
    print(e)
except:
    print('ERR')
code

5. Що буде виведено за виконання?\par (Дійсні значення, якщо наявні, в цьому завданні будуть виводитись у форматі з фіксованою крапкою та, можливо, з доволі великою кількістю десяткових розрядів. У відповіді для дійсних значень у ЦЬОМУ завданні достатньо вказати один десятковий розряд після крапки, відкинувши решту розрядів без округлення. Наприклад, замість 13.66666666666667 достатньо вказати 13.6, замість 25.23 достатньо вказати 25.2)
code
try:
    print(6, end=';')
    print(4j<=9j, end=';')
    print(9, end=';')
except KeyboardInterrupt:
    print(5, end=';')
except TypeError:
    print(1, end=';')
except:
    print('ERR', end=';')
else:
    print(6, end=';')
finally:
    print(7)
code

6. Що буде виведено за виконання?
code
def f(n):
    if n>2:
        f(n//2)
    print(n%2, end='')
 
f(16)
code

7. Що буде виведено за виконання?
code
class A:
    c = 1

    def __f1(self): print(2, end=';')
    def _f2(self): print(3, end=';')

    def main(self):
        try: print(self.c, end=';')
        except: print('ERR', end=';')

        try: self.__f1()
        except: print('ERR', end=';')

        try: self._f2()
        except: print('ERR', end=';')


class B(A):
    c = 4

    def __f1(self): print(5, end=';')
    def _f2(self): print(6, end=';')

try: B().main()
except: print('ERR', end=';')
code

8. Що буде виведено за виконання?
code
def f(arg, n):
    n = 4
    tmp = arg
    tmp[5:] = 'bd'
    return len(tmp)>3

try:
    d = ['0','1','2','3','4']
    n = 1
    f(d, n)
    print(n, end=';')
    for el in d:
        print(el, end=';')
except:
    print('ERR')
code

9. Що буде виведено за виконання?
code
class A:
    a = 1
    b = 2

    def __init__(self):
        a = 3

class B(A):
    a = 4

    def __init__(self):
        super().__init__()
        self.c = 7

obj = B()
obj.c = obj.a + 10
B.c = 13

try: print(obj.c, end= ';')
except: print('ERR', end=';')

try: print(A.c, end= ';')
except: print('ERR', end=';')

try: print(B.c, end= ';')
except: print('ERR', end=';')
code

10. 
Для графа на рисунку з вершини 1 запущено алгоритм пошуку в глибину, що будує кістякове дерево. Списки суміжності вершин переглядаються алгоритмом за зростанням міток вершин.\par
\begin{center}
	\adjustimage{max size={0.9\linewidth}{0.9\paperheight}}
	{../10/Traversals/285.png}
\end{center}
\par Які з наведених ребер входять до складу отриманого кістякового дерева?\par
1) (2, 5)\par
2) (3, 4)\par
3) (1, 6)\par
4) (1, 3)\par
5) (1, 5)\par
6) (1, 2)\par
{\vskip 15pt} У формі вибрати номери відповідних ребер.\par
%answer:1 2
\end {large}}
%end
{\smallskip \begin {small} Затверджено на засіданні кафедри теоретичної кібернетики, протокол \No 12 від   19.05.2020 р.\par\smallskip\smallskip
{\parbox {6.5 cm}{Зав. кафедри \dotfill Ю.В.~Крак}}\hspace {0.7cm}{\hbox to 6.5 cm{ Екзаменатор \dotfill Т.О.~Карнаух}} \end{small}}
%begin
{{\begin {small}\par
{ \center  Київський національний університет імені Тараса Шевченка \par}
{\parbox {5 cm}{Освітня програма \hfill}{\parbox {7 cm}{Прикладна математика, Системний аналіз \hfill}}\parbox {6 cm}{\hfill Семестр 2}}\par
{\parbox {5 cm} {Навчальний предмет \hfill}\parbox {7 cm}{Програмування \hfill}\parbox {6cm}{\hfill}\par}\vspace{-7pt} \end{small}}{\center Екзаменаційний білет \No \textbf 
{ 376 }\par}}
{
\begin {large}
1. Що буде виведено за виконання?
code
try:
    res = 5//3/7*4
    if isinstance(res, bool):
        print(f'bool;{res}')
    elif isinstance(res, int):
        print(f'int;{res}')
    elif isinstance(res, float):
        print(f'float;{res:.2f}')
    else:
        print('???')
except:
    print('ERR')
code

2. Що буде виведено за виконання?
code
a = 6
b = 8
c = 9

def h(b):
    global c
    a = 2
    b += 4
    c = 5
    return a + b + c

a = 1
b = 1
c = 0

try:
    print(h(a), a, b, c, sep=';')
except:
    print('ERR', a, b, c, sep=';')
code

3. Що буде виведено за виконання?
code
try:
    i = 6
    while i<7:
        i += 2
    print(i, end=';')
except:
    print('ERR')
code

4. Що буде виведено за виконання?
code
try:
    for a in range(14, -1, 3):
        if a<=6:
            continue
        print(a, end=';');a = 2
    else:
        print(a,end=';')
    print(a)
except:
    print('ERR')
code

5. Що буде виведено за виконання?\par (Дійсні значення, якщо наявні, в цьому завданні будуть виводитись у форматі з фіксованою крапкою та, можливо, з доволі великою кількістю десяткових розрядів. У відповіді для дійсних значень у ЦЬОМУ завданні достатньо вказати один десятковий розряд після крапки, відкинувши решту розрядів без округлення. Наприклад, замість 13.66666666666667 достатньо вказати 13.6, замість 25.23 достатньо вказати 25.2)
code
try:
    print(4, end=';')
    print(2//0, end=';')
    print(8, end=';')
except Exception:
    print(3, end=';')
except ValueError:
    print(0, end=';')
except:
    print('ERR', end=';')
else:
    print(0, end=';')
finally:
    print(7)
code

6. Що буде виведено за виконання?
code
def f(n):
    if n<=9:
        print(n%10, end='')
    else:
        f(n//10)

f(123456)
code

7. Що буде виведено за виконання?
code
class A:
    b = 1

    def _g1(self): print(2, end=';')
    def _g2(self): print(3, end=';')

    def main(self):
        try: print(self.b, end=';')
        except: print('ERR', end=';')

        try: self._g1()
        except: print('ERR', end=';')

        try: self._g2()
        except: print('ERR', end=';')


class B(A):
    b = 4

    def _g1(self): print(5, end=';')
    def _g2(self): print(6, end=';')

try: B().main()
except: print('ERR', end=';')
code

8. Що буде виведено за виконання?
code
def f(arg, n):
    arg[85] = 1
    n -= -3

try:
    n = 5
    d = {31:1, 61:9, 85:5, 25:6}
    f(d, n)
    print(n, end=';')
    for x in d.items():
        print(*x, sep=';', end=';')
except:
    print('ERR')
code

9. Що буде виведено за виконання?
code
class A:
    b = 1
    c = 2

    def __init__(self):
        c = 3

class B(A):
    c = 4

    def __init__(self):
        super().__init__()
        self.a = 7

obj = B()
obj.b = obj.a + 10
B.b = 13

try: print(obj.b, end= ';')
except: print('ERR', end=';')

try: print(A.b, end= ';')
except: print('ERR', end=';')

try: print(B.b, end= ';')
except: print('ERR', end=';')
code

10. 
Для графа на рисунку з вершини 1 запущено алгоритм пошуку в глибину, що будує кістякове дерево. Списки суміжності вершин переглядаються алгоритмом за зростанням міток вершин.\par
\begin{center}
	\adjustimage{max size={0.9\linewidth}{0.9\paperheight}}
	{../10/Traversals/142.png}
\end{center}
\par Які з наведених ребер входять до складу отриманого кістякового дерева?\par
1) (3, 6)\par
2) (2, 5)\par
3) (4, 5)\par
4) (0, 6)\par
5) (0, 5)\par
6) (0, 3)\par
{\vskip 15pt} У формі вибрати номери відповідних ребер.\par
%answer:1 2 5 6
\end {large}}
%end
{\smallskip \begin {small} Затверджено на засіданні кафедри теоретичної кібернетики, протокол \No 12 від   19.05.2020 р.\par\smallskip\smallskip
{\parbox {6.5 cm}{Зав. кафедри \dotfill Ю.В.~Крак}}\hspace {0.7cm}{\hbox to 6.5 cm{ Екзаменатор \dotfill Т.О.~Карнаух}} \end{small}}
%begin
{{\begin {small}\par
{ \center  Київський національний університет імені Тараса Шевченка \par}
{\parbox {5 cm}{Освітня програма \hfill}{\parbox {7 cm}{Прикладна математика, Системний аналіз \hfill}}\parbox {6 cm}{\hfill Семестр 2}}\par
{\parbox {5 cm} {Навчальний предмет \hfill}\parbox {7 cm}{Програмування \hfill}\parbox {6cm}{\hfill}\par}\vspace{-7pt} \end{small}}{\center Екзаменаційний білет \No \textbf 
{ 377 }\par}}
{
\begin {large}
1. Що буде виведено за виконання?
code
try:
    res = "5.0"<=True
    if isinstance(res, bool):
        print(f'bool;{res}')
    elif isinstance(res, int):
        print(f'int;{res}')
    elif isinstance(res, float):
        print(f'float;{res:.2f}')
    else:
        print('???')
except:
    print('ERR')
code

2. Що буде виведено за виконання?
code
a = 6
b = 4
c = 2

def f(a):
    a -= 4
    b = 2
    c = 1
    return a + b + c

a = 7
b = 9
c = 3

try:
    print(f(a), a, b, c, sep=';')
except:
    print('ERR', a, b, c, sep=';')
code

3. Що буде виведено за виконання?
code
try:
    k = 9
    while k>=0:
        k += -2
    print(k, end=';')
except:
    print('ERR')
code

4. Що буде виведено за виконання?
code
try:
    for b in range(-13, 7, -3):
        if b>=7:
            break
        print(b, end=';')
    else:
        print(b,end=';')
    print(b)
except:
    print('ERR')
code

5. Що буде виведено за виконання?\par (Дійсні значення, якщо наявні, в цьому завданні будуть виводитись у форматі з фіксованою крапкою та, можливо, з доволі великою кількістю десяткових розрядів. У відповіді для дійсних значень у ЦЬОМУ завданні достатньо вказати один десятковий розряд після крапки, відкинувши решту розрядів без округлення. Наприклад, замість 13.66666666666667 достатньо вказати 13.6, замість 25.23 достатньо вказати 25.2)
code
try:
    print(1, end=';')
    print(int('9'), end=';')
    print(5, end=';')
except KeyboardInterrupt:
    print(8, end=';')
except TypeError:
    print(3, end=';')
except:
    print('ERR', end=';')
else:
    print(2, end=';')
finally:
    print(6)
code

6. Що буде виведено за виконання?
code
def f(n):
    if n<=9:
        print(n%10,end='')
    else:
        f(n//10)

f(543216)
code

7. Що буде виведено за виконання?
code
class A:
    a = 1

    def __f1(self): print(2, end=';')
    def _f2(self): print(3, end=';')

    def main(self):
        try: print(self.a, end=';')
        except: print('ERR', end=';')

        try: self.__f1()
        except: print('ERR', end=';')

        try: self._f2()
        except: print('ERR', end=';')


class B(A):
    a = 4

    def __f1(self): print(5, end=';')
    def _f2(self): print(6, end=';')

try: B().main()
except: print('ERR', end=';')
code

8. Що буде виведено за виконання?
code
def f(arg, n):
    n = 6
    tmp = arg
    del tmp[-5:-4]
    return len(tmp)>3

try:
    a = ['0','1','2','3','4']
    n = 9
    f(a, n)
    print(n, end=';')
    for el in a:
        print(el, end=';')
except:
    print('ERR')
code

9. Що буде виведено за виконання?
code
class A:
    c = 1
    a = 2

    def __init__(self):
        c = 3

class B(A):
    a = 4

    def __init__(self):
        super().__init__()
        self.b = 7

obj = B()
obj.c = obj.b + 10
B.c = 13

try: print(obj.c, end= ';')
except: print('ERR', end=';')

try: print(A.c, end= ';')
except: print('ERR', end=';')

try: print(B.c, end= ';')
except: print('ERR', end=';')
code

10. 
Для графа на рисунку з вершини 3 запущено алгоритм пошуку в глибину, що будує кістякове дерево. Списки суміжності вершин переглядаються алгоритмом за зростанням міток вершин.\par
\begin{center}
	\adjustimage{max size={0.9\linewidth}{0.9\paperheight}}
	{../10/Traversals/30.png}
\end{center}
\par Які з наведених ребер входять до складу отриманого кістякового дерева?\par
1) (1, 4)\par
2) (2, 5)\par
3) (5, 6)\par
4) (1, 6)\par
5) (3, 6)\par
6) (0, 5)\par
{\vskip 15pt} У формі вибрати номери відповідних ребер.\par
%answer:1 2 4 6
\end {large}}
%end
{\smallskip \begin {small} Затверджено на засіданні кафедри теоретичної кібернетики, протокол \No 12 від   19.05.2020 р.\par\smallskip\smallskip
{\parbox {6.5 cm}{Зав. кафедри \dotfill Ю.В.~Крак}}\hspace {0.7cm}{\hbox to 6.5 cm{ Екзаменатор \dotfill Т.О.~Карнаух}} \end{small}}
%begin
{{\begin {small}\par
{ \center  Київський національний університет імені Тараса Шевченка \par}
{\parbox {5 cm}{Освітня програма \hfill}{\parbox {7 cm}{Прикладна математика, Системний аналіз \hfill}}\parbox {6 cm}{\hfill Семестр 2}}\par
{\parbox {5 cm} {Навчальний предмет \hfill}\parbox {7 cm}{Програмування \hfill}\parbox {6cm}{\hfill}\par}\vspace{-7pt} \end{small}}{\center Екзаменаційний білет \No \textbf 
{ 378 }\par}}
{
\begin {large}
1. Що буде виведено за виконання?
code
def f(n):
    print('f', end='')
    return n

try:
    print(f(-6) and f(4))
except:
    print('ERR')
code

2. Що буде виведено за виконання?
code
a = 5
b = 8
c = 1

def g(b):
    a *= 3
    b = 4
    c = 2
    return a + b + c

a = 6
b = 8
c = 2

try:
    print(g(a), a, b, c, sep=';')
except:
    print('ERR', a, b, c, sep=';')
code

3. Що буде виведено за виконання?
code
try:
    t = 11
    while t>2:
        t += -2
    print(t, end=';')
except:
    print('ERR')
code

4. Що буде виведено за виконання?
code
try:
    for d in range(-8, 1, -1):
        if d<=3:
            continue
        print(d, end=';')
    else:
        print(d,end=';')
    print(d)
except:
    print('ERR')
code

5. Що буде виведено за виконання?\par (Дійсні значення, якщо наявні, в цьому завданні будуть виводитись у форматі з фіксованою крапкою та, можливо, з доволі великою кількістю десяткових розрядів. У відповіді для дійсних значень у ЦЬОМУ завданні достатньо вказати один десятковий розряд після крапки, відкинувши решту розрядів без округлення. Наприклад, замість 13.66666666666667 достатньо вказати 13.6, замість 25.23 достатньо вказати 25.2)
code
try:
    print(4, end=';')
    print(int('d1'), end=';')
    print(0, end=';')
except ZeroDivisionError:
    print(8, end=';')
except ValueError:
    print(4, end=';')
except:
    print('ERR', end=';')
else:
    print(5, end=';')
finally:
    print(7)
code

6. Що буде виведено за виконання?
code
def f(n):
    if n>2:
        f(n//2)
    else:
        print(n%2, end='')

f(16)
code

7. Що буде виведено за виконання?
code
class A:
    c = 1

    def _h1(self): print(2, end=';')
    def __h2(self): print(3, end=';')

    def main(self):
        try: print(self.c, end=';')
        except: print('ERR', end=';')

        try: self._h1()
        except: print('ERR', end=';')

        try: self.__h2()
        except: print('ERR', end=';')


class B(A):
    c = 4

    def _h1(self): print(5, end=';')
    def __h2(self): print(6, end=';')

try: B().main()
except: print('ERR', end=';')
code

8. Що буде виведено за виконання?
code
def f(arg, n):
    n = 2
    tmp = arg
    tmp[:5] = 'bc'
    return len(tmp)>3

try:
    c = [10,11,12,13,14]
    n = 7
    f(c, n)
    print(n, end=';')
    for el in c:
        print(el, end=';')
except:
    print('ERR')
code

9. Що буде виведено за виконання?
code
class A:
    b = 1
    a = 2

    def __init__(self):
        a = 3

class B(A):
    b = 4

    def __init__(self):
        super().__init__()
        self.c = 7

obj = B()
obj.a = obj.c + 10
B.a = 13

try: print(obj.a, end= ';')
except: print('ERR', end=';')

try: print(A.a, end= ';')
except: print('ERR', end=';')

try: print(B.a, end= ';')
except: print('ERR', end=';')
code

10. 
Для графа на рисунку з вершини 2 запущено алгоритм пошуку в глибину, що будує кістякове дерево. Списки суміжності вершин переглядаються алгоритмом за зростанням міток вершин.\par
\begin{center}
	\adjustimage{max size={0.9\linewidth}{0.9\paperheight}}
	{../10/Traversals/266.png}
\end{center}
\par Які з наведених ребер входять до складу отриманого кістякового дерева?\par
1) (0, 6)\par
2) (0, 2)\par
3) (3, 6)\par
4) (5, 6)\par
5) (2, 6)\par
6) (2, 3)\par
{\vskip 15pt} У формі вибрати номери відповідних ребер.\par
%answer:2 3
\end {large}}
%end
{\smallskip \begin {small} Затверджено на засіданні кафедри теоретичної кібернетики, протокол \No 12 від   19.05.2020 р.\par\smallskip\smallskip
{\parbox {6.5 cm}{Зав. кафедри \dotfill Ю.В.~Крак}}\hspace {0.7cm}{\hbox to 6.5 cm{ Екзаменатор \dotfill Т.О.~Карнаух}} \end{small}}
%begin
{{\begin {small}\par
{ \center  Київський національний університет імені Тараса Шевченка \par}
{\parbox {5 cm}{Освітня програма \hfill}{\parbox {7 cm}{Прикладна математика, Системний аналіз \hfill}}\parbox {6 cm}{\hfill Семестр 2}}\par
{\parbox {5 cm} {Навчальний предмет \hfill}\parbox {7 cm}{Програмування \hfill}\parbox {6cm}{\hfill}\par}\vspace{-7pt} \end{small}}{\center Екзаменаційний білет \No \textbf 
{ 379 }\par}}
{
\begin {large}
1. Що буде виведено за виконання?
code
def f(n):
    print('f', end='')
    return n<0

try:
    print(f(1) or f(-5))
except:
    print('ERR')
code

2. Що буде виведено за виконання?
code
a = 4
b = 7
c = 9

def g(a):
    global c
    a += 1
    b = 3
    c = 5
    return a + b + c

a = 3
b = 5
c = 0

try:
    print(g(b), a, b, c, sep=';')
except:
    print('ERR', a, b, c, sep=';')
code

3. Що буде виведено за виконання?
code
try:
    j = 14
    while j>=9:
        j += -3
    print(j, end=';')
except:
    print('ERR')
code

4. Що буде виведено за виконання?
code
try:
    for a in range(11, -11, 2):
        if a<4:
            continue
        print(a, end=';')
    else:
        print(a,end=';')
    print(a)
except:
    print('ERR')
code

5. Що буде виведено за виконання?\par (Дійсні значення, якщо наявні, в цьому завданні будуть виводитись у форматі з фіксованою крапкою та, можливо, з доволі великою кількістю десяткових розрядів. У відповіді для дійсних значень у ЦЬОМУ завданні достатньо вказати один десятковий розряд після крапки, відкинувши решту розрядів без округлення. Наприклад, замість 13.66666666666667 достатньо вказати 13.6, замість 25.23 достатньо вказати 25.2)
code
try:
    print(6, end=';')
    print(3%1, end=';')
    print(2, end=';')
except Exception:
    print(9, end=';')
except TypeError:
    print(8, end=';')
except:
    print('ERR', end=';')
else:
    print(6, end=';')
finally:
    print(3)
code

6. Що буде виведено за виконання?
code
def f(n):
    if n<=9:
        print(n%10, end='')
    else:
        f(n//10)

f(987651)
code

7. Що буде виведено за виконання?
code
class A:
    b = 1

    def _h1(self): print(2, end=';')
    def h2(self): print(3, end=';')

    def main(self):
        try: print(self.b, end=';')
        except: print('ERR', end=';')

        try: self._h1()
        except: print('ERR', end=';')

        try: self.h2()
        except: print('ERR', end=';')


class B(A):
    b = 4

    def _h1(self): print(5, end=';')
    def h2(self): print(6, end=';')

try: B().main()
except: print('ERR', end=';')
code

8. Що буде виведено за виконання?
code
def f(arg, n):
    arg[76] = 7
    n = 5
    arg = list(arg.keys())

try:
    n = 2
    t = {18:6, 34:2, 76:9, 76:6}
    f(t, n)
    print(n, end=';')
    for x in t.items():
        print(x, sep=';', end=';')
except:
    print('ERR')
code

9. Що буде виведено за виконання?
code
class A:
    a = 1
    c = 2

    def __init__(self):
        a = 3

class B(A):
    a = 4

    def __init__(self):
        super().__init__()
        self.b = 7

obj = B()
obj.b = obj.c + 10
B.b = 13

try: print(obj.b, end= ';')
except: print('ERR', end=';')

try: print(A.b, end= ';')
except: print('ERR', end=';')

try: print(B.b, end= ';')
except: print('ERR', end=';')
code

10. 
Для графа на рисунку з вершини 3 запущено алгоритм пошуку в глибину, що будує кістякове дерево. Списки суміжності вершин переглядаються алгоритмом за зростанням міток вершин.\par
\begin{center}
	\adjustimage{max size={0.9\linewidth}{0.9\paperheight}}
	{../10/Traversals/82.png}
\end{center}
\par Які з наведених ребер входять до складу отриманого кістякового дерева?\par
1) (0, 1)\par
2) (0, 3)\par
3) (1, 5)\par
4) (2, 3)\par
5) (0, 4)\par
6) (2, 6)\par
{\vskip 15pt} У формі вибрати номери відповідних ребер.\par
%answer:1 2 3 5 6
\end {large}}
%end
{\smallskip \begin {small} Затверджено на засіданні кафедри теоретичної кібернетики, протокол \No 12 від   19.05.2020 р.\par\smallskip\smallskip
{\parbox {6.5 cm}{Зав. кафедри \dotfill Ю.В.~Крак}}\hspace {0.7cm}{\hbox to 6.5 cm{ Екзаменатор \dotfill Т.О.~Карнаух}} \end{small}}
%begin
{{\begin {small}\par
{ \center  Київський національний університет імені Тараса Шевченка \par}
{\parbox {5 cm}{Освітня програма \hfill}{\parbox {7 cm}{Прикладна математика, Системний аналіз \hfill}}\parbox {6 cm}{\hfill Семестр 2}}\par
{\parbox {5 cm} {Навчальний предмет \hfill}\parbox {7 cm}{Програмування \hfill}\parbox {6cm}{\hfill}\par}\vspace{-7pt} \end{small}}{\center Екзаменаційний білет \No \textbf 
{ 380 }\par}}
{
\begin {large}
1. Що буде виведено за виконання?
code
try:
    res = 4**0.5*2
    if isinstance(res, bool):
        print(f'bool;{res}')
    elif isinstance(res, int):
        print(f'int;{res}')
    elif isinstance(res, float):
        print(f'float;{res:.2f}')
    else:
        print('???')
except:
    print('ERR')
code

2. Що буде виведено за виконання?
code
a = 9
b = 6
c = 4

def h(a):
    global c
    a = 3
    b = 4
    c = 1
    return a + b + c

a = 2
b = 0
c = 7

try:
    print(h(b), a, b, c, sep=';')
except:
    print('ERR', a, b, c, sep=';')
code

3. Що буде виведено за виконання?
code
try:
    m = 7
    while m>=7:
        m += -3
    print(m, end=';')
except:
    print('ERR')
code

4. Що буде виведено за виконання?
code
try:
    for c in range(-5, -13, -2):
        if c>5:
            continue
        print(c, end=';')
        if c<=7:
            break
    else:
        print(c,end=';')
    print(c)
except:
    print('ERR')
code

5. Що буде виведено за виконання?\par (Дійсні значення, якщо наявні, в цьому завданні будуть виводитись у форматі з фіксованою крапкою та, можливо, з доволі великою кількістю десяткових розрядів. У відповіді для дійсних значень у ЦЬОМУ завданні достатньо вказати один десятковий розряд після крапки, відкинувши решту розрядів без округлення. Наприклад, замість 13.66666666666667 достатньо вказати 13.6, замість 25.23 достатньо вказати 25.2)
code
try:
    print(1, end=';')
    print(4j!=0j, end=';')
    print(0, end=';')
except ZeroDivisionError:
    print(2, end=';')
except KeyboardInterrupt:
    print(9, end=';')
except:
    print('ERR', end=';')
else:
    print(5, end=';')
finally:
    print(7)
code

6. Що буде виведено за виконання?
code
def f(s):
    s1 = s
    for i in range(len(s)):
        s1[i] += 1
        return s1

s = [1, 2, 3]
f(s)
print(s)
code

7. Що буде виведено за виконання?
code
class A:
    d = 1

    def g1(self): print(2, end=';')
    def _g2(self): print(3, end=';')

    def main(self):
        try: print(self.d, end=';')
        except: print('ERR', end=';')

        try: self.g1()
        except: print('ERR', end=';')

        try: self._g2()
        except: print('ERR', end=';')


class B(A):
    d = 4

    def g1(self): print(5, end=';')
    def _g2(self): print(6, end=';')

try: B().main()
except: print('ERR', end=';')
code

8. Що буде виведено за виконання?
code
def f(arg, n):
    n = 0
    tmp = arg
    del tmp[0:]
    return len(tmp)>3

try:
    e = [10,11,12,13,14]
    n = 3
    f(e, n)
    print(n, end=';')
    for el in e:
        print(el, end=';')
except:
    print('ERR')
code

9. Що буде виведено за виконання?
code
class A:
    a = 1
    b = 2

    def __init__(self):
        b = 3

class B(A):
    b = 4

    def __init__(self):
        super().__init__()
        self.c = 7

obj = B()
obj.a = obj.b + 10
B.a = 13

try: print(obj.a, end= ';')
except: print('ERR', end=';')

try: print(A.a, end= ';')
except: print('ERR', end=';')

try: print(B.a, end= ';')
except: print('ERR', end=';')
code

10. 
Для графа на рисунку з вершини 6 запущено алгоритм пошуку в глибину, що будує кістякове дерево. Списки суміжності вершин переглядаються алгоритмом за зростанням міток вершин.\par
\begin{center}
	\adjustimage{max size={0.9\linewidth}{0.9\paperheight}}
	{../10/Traversals/256.png}
\end{center}
\par Які з наведених ребер входять до складу отриманого кістякового дерева?\par
1) (0, 3)\par
2) (3, 4)\par
3) (1, 6)\par
4) (1, 3)\par
5) (4, 6)\par
6) (0, 6)\par
{\vskip 15pt} У формі вибрати номери відповідних ребер.\par
%answer:1 4 6
\end {large}}
%end
{\smallskip \begin {small} Затверджено на засіданні кафедри теоретичної кібернетики, протокол \No 12 від   19.05.2020 р.\par\smallskip\smallskip
{\parbox {6.5 cm}{Зав. кафедри \dotfill Ю.В.~Крак}}\hspace {0.7cm}{\hbox to 6.5 cm{ Екзаменатор \dotfill Т.О.~Карнаух}} \end{small}}
%begin
{{\begin {small}\par
{ \center  Київський національний університет імені Тараса Шевченка \par}
{\parbox {5 cm}{Освітня програма \hfill}{\parbox {7 cm}{Прикладна математика, Системний аналіз \hfill}}\parbox {6 cm}{\hfill Семестр 2}}\par
{\parbox {5 cm} {Навчальний предмет \hfill}\parbox {7 cm}{Програмування \hfill}\parbox {6cm}{\hfill}\par}\vspace{-7pt} \end{small}}{\center Екзаменаційний білет \No \textbf 
{ 381 }\par}}
{
\begin {large}
1. Що буде виведено за виконання?
code
try:
    res = 2>=4.0 and not 6.0<=4 and True>3
    if isinstance(res, bool):
        print(f'bool;{res}')
    elif isinstance(res, int):
        print(f'int;{res}')
    elif isinstance(res, float):
        print(f'float;{res:.2f}')
    else:
        print('???')
except:
    print('ERR')
code

2. Що буде виведено за виконання?
code
a = 9
b = 6
c = 0

def f(a):
    global c
    a = 5
    b -= 2
    c = 1
    return a + b + c

a = 3
b = 4
c = 2

try:
    print(f(b), a, b, c, sep=';')
except:
    print('ERR', a, b, c, sep=';')
code

3. Що буде виведено за виконання?
code
try:
    i = 5
    while i>5:
        i += -2
    print(i, end=';')
except:
    print('ERR')
code

4. Що буде виведено за виконання?
code
try:
    for b in range(2, 12, 2):
        if b<=6:
            break
        print(b, end=';');b = -10
    else:
        print(b,end=';')
    print(b)
except:
    print('ERR')
code

5. Що буде виведено за виконання?\par (Дійсні значення, якщо наявні, в цьому завданні будуть виводитись у форматі з фіксованою крапкою та, можливо, з доволі великою кількістю десяткових розрядів. У відповіді для дійсних значень у ЦЬОМУ завданні достатньо вказати один десятковий розряд після крапки, відкинувши решту розрядів без округлення. Наприклад, замість 13.66666666666667 достатньо вказати 13.6, замість 25.23 достатньо вказати 25.2)
code
try:
    print(7, end=';')
    print(int('2'), end=';')
    print(1, end=';')
except ValueError:
    print(4, end=';')
except TypeError:
    print(0, end=';')
except:
    print('ERR', end=';')
else:
    print(6, end=';')
finally:
    print(3)
code

6. Що буде виведено за виконання?
code
def f(n):
    if n>9:
        f(n//10)
    print(n%10, end='')
 
f(543216)
code

7. Що буде виведено за виконання?
code
class A:
    a = 1

    def _f1(self): print(2, end=';')
    def __f2(self): print(3, end=';')

    def main(self):
        try: print(self.a, end=';')
        except: print('ERR', end=';')

        try: self._f1()
        except: print('ERR', end=';')

        try: self.__f2()
        except: print('ERR', end=';')


class B(A):
    a = 4

    def _f1(self): print(5, end=';')
    def __f2(self): print(6, end=';')

try: B().main()
except: print('ERR', end=';')
code

8. Що буде виведено за виконання?
code
def f(arg, n):
    arg[82] = 8
    n += 2

try:
    n = 4
    s = {31:6, 12:8, 31:0}
    f(s, n)
    print(n, end=';')
    for x in s.keys():
        print(x, sep=';', end=';')
except:
    print('ERR')
code

9. Що буде виведено за виконання?
code
class A:
    c = 1
    b = 2

    def __init__(self):
        b = 3

class B(A):
    c = 4

    def __init__(self):
        super().__init__()
        self.a = 7

obj = B()
obj.a = obj.c + 10
B.a = 13

try: print(obj.a, end= ';')
except: print('ERR', end=';')

try: print(A.a, end= ';')
except: print('ERR', end=';')

try: print(B.a, end= ';')
except: print('ERR', end=';')
code

10. 
Для графа на рисунку з вершини 1 запущено алгоритм пошуку в глибину, що будує кістякове дерево. Списки суміжності вершин переглядаються алгоритмом за зростанням міток вершин.\par
\begin{center}
	\adjustimage{max size={0.9\linewidth}{0.9\paperheight}}
	{../10/Traversals/271.png}
\end{center}
\par Які з наведених ребер входять до складу отриманого кістякового дерева?\par
1) (0, 1)\par
2) (2, 3)\par
3) (1, 5)\par
4) (5, 6)\par
5) (1, 6)\par
6) (0, 4)\par
{\vskip 15pt} У формі вибрати номери відповідних ребер.\par
%answer:1 2 4 6
\end {large}}
%end
{\smallskip \begin {small} Затверджено на засіданні кафедри теоретичної кібернетики, протокол \No 12 від   19.05.2020 р.\par\smallskip\smallskip
{\parbox {6.5 cm}{Зав. кафедри \dotfill Ю.В.~Крак}}\hspace {0.7cm}{\hbox to 6.5 cm{ Екзаменатор \dotfill Т.О.~Карнаух}} \end{small}}
%begin
{{\begin {small}\par
{ \center  Київський національний університет імені Тараса Шевченка \par}
{\parbox {5 cm}{Освітня програма \hfill}{\parbox {7 cm}{Прикладна математика, Системний аналіз \hfill}}\parbox {6 cm}{\hfill Семестр 2}}\par
{\parbox {5 cm} {Навчальний предмет \hfill}\parbox {7 cm}{Програмування \hfill}\parbox {6cm}{\hfill}\par}\vspace{-7pt} \end{small}}{\center Екзаменаційний білет \No \textbf 
{ 382 }\par}}
{
\begin {large}
1. Що буде виведено за виконання?
code
try:
    res = not 5!=8.0 or False==3.0 or 9<7
    if isinstance(res, bool):
        print(f'bool;{res}')
    elif isinstance(res, int):
        print(f'int;{res}')
    elif isinstance(res, float):
        print(f'float;{res:.2f}')
    else:
        print('???')
except:
    print('ERR')
code

2. Що буде виведено за виконання?
code
a = 9
b = 2
c = 1

def g(b):
    global c
    a = 2
    b = 5
    c = 4
    return a + b + c

a = 7
b = 3
c = 0

try:
    print(g(a), a, b, c, sep=';')
except:
    print('ERR', a, b, c, sep=';')
code

3. Що буде виведено за виконання?
code
try:
    m = 4
    while m<=5:
        m += 1
    print(m, end=';')
except:
    print('ERR')
code

4. Що буде виведено за виконання?
code
try:
    for e in range(4, 9, -2):
        if e<3:
            continue
        print(e, end=';');e = 5
    else:
        print(e,end=';')
    print(e)
except:
    print('ERR')
code

5. Що буде виведено за виконання?\par (Дійсні значення, якщо наявні, в цьому завданні будуть виводитись у форматі з фіксованою крапкою та, можливо, з доволі великою кількістю десяткових розрядів. У відповіді для дійсних значень у ЦЬОМУ завданні достатньо вказати один десятковий розряд після крапки, відкинувши решту розрядів без округлення. Наприклад, замість 13.66666666666667 достатньо вказати 13.6, замість 25.23 достатньо вказати 25.2)
code
try:
    print(8, end=';')
    print(int('a5'), end=';')
    print(9, end=';')
except ZeroDivisionError:
    print(6, end=';')
except Exception:
    print(3, end=';')
except:
    print('ERR', end=';')
else:
    print(5, end=';')
finally:
    print(9)
code

6. Що буде виведено за виконання?
code
def f(s):
    for el in s:
        el += 1
        return s
 
s = [1, 2, 3]
s = f(s)
print(s)
code

7. Що буде виведено за виконання?
code
class A:
    b = 1

    def _g1(self): print(2, end=';')
    def _g2(self): print(3, end=';')

    def main(self):
        try: print(self.b, end=';')
        except: print('ERR', end=';')

        try: self._g1()
        except: print('ERR', end=';')

        try: self._g2()
        except: print('ERR', end=';')


class B(A):
    b = 4

    def _g1(self): print(5, end=';')
    def _g2(self): print(6, end=';')

try: B().main()
except: print('ERR', end=';')
code

8. Що буде виведено за виконання?
code
def f(arg, n):
    arg[59] = 7
    n = 9
    arg = set(arg.items())

try:
    n = 1
    s = {18:0, 25:9, 25:0}
    f(s, n)
    print(n, end=';')
    for x in s.keys():
        print(x, sep=';', end=';')
except:
    print('ERR')
code

9. Що буде виведено за виконання?
code
class A:
    b = 1
    c = 2

    def __init__(self):
        b = 3

class B(A):
    c = 4

    def __init__(self):
        super().__init__()
        self.a = 7

obj = B()
obj.b = obj.c + 10
B.b = 13

try: print(obj.b, end= ';')
except: print('ERR', end=';')

try: print(A.b, end= ';')
except: print('ERR', end=';')

try: print(B.b, end= ';')
except: print('ERR', end=';')
code

10. 
Для графа на рисунку з вершини 2 запущено алгоритм пошуку в глибину, що будує кістякове дерево. Списки суміжності вершин переглядаються алгоритмом за зростанням міток вершин.\par
\begin{center}
	\adjustimage{max size={0.9\linewidth}{0.9\paperheight}}
	{../10/Traversals/229.png}
\end{center}
\par Які з наведених ребер входять до складу отриманого кістякового дерева?\par
1) (0, 3)\par
2) (2, 3)\par
3) (1, 3)\par
4) (0, 2)\par
5) (1, 5)\par
6) (3, 4)\par
{\vskip 15pt} У формі вибрати номери відповідних ребер.\par
%answer:1 3 4 5 6
\end {large}}
%end
{\smallskip \begin {small} Затверджено на засіданні кафедри теоретичної кібернетики, протокол \No 12 від   19.05.2020 р.\par\smallskip\smallskip
{\parbox {6.5 cm}{Зав. кафедри \dotfill Ю.В.~Крак}}\hspace {0.7cm}{\hbox to 6.5 cm{ Екзаменатор \dotfill Т.О.~Карнаух}} \end{small}}
%begin
{{\begin {small}\par
{ \center  Київський національний університет імені Тараса Шевченка \par}
{\parbox {5 cm}{Освітня програма \hfill}{\parbox {7 cm}{Прикладна математика, Системний аналіз \hfill}}\parbox {6 cm}{\hfill Семестр 2}}\par
{\parbox {5 cm} {Навчальний предмет \hfill}\parbox {7 cm}{Програмування \hfill}\parbox {6cm}{\hfill}\par}\vspace{-7pt} \end{small}}{\center Екзаменаційний білет \No \textbf 
{ 383 }\par}}
{
\begin {large}
1. Що буде виведено за виконання?
code
try:
    res = "False">"5"
    if isinstance(res, bool):
        print(f'bool;{res}')
    elif isinstance(res, int):
        print(f'int;{res}')
    elif isinstance(res, float):
        print(f'float;{res:.2f}')
    else:
        print('???')
except:
    print('ERR')
code

2. Що буде виведено за виконання?
code
a = 8
b = 5
c = 1

def h(b):
    a *= 3
    b = 4
    c = 3
    return a + b + c

a = 7
b = 2
c = 7

try:
    print(h(b), a, b, c, sep=';')
except:
    print('ERR', a, b, c, sep=';')
code

3. Що буде виведено за виконання?
code
try:
    j = 2
    while j<12:
        j += 1
    print(j, end=';')
except:
    print('ERR')
code

4. Що буде виведено за виконання?
code
try:
    for d in range(-15, -8, 3):
        if d>7:
            break
        print(d, end=';')
    else:
        print(d,end=';')
    print(d)
except:
    print('ERR')
code

5. Що буде виведено за виконання?\par (Дійсні значення, якщо наявні, в цьому завданні будуть виводитись у форматі з фіксованою крапкою та, можливо, з доволі великою кількістю десяткових розрядів. У відповіді для дійсних значень у ЦЬОМУ завданні достатньо вказати один десятковий розряд після крапки, відкинувши решту розрядів без округлення. Наприклад, замість 13.66666666666667 достатньо вказати 13.6, замість 25.23 достатньо вказати 25.2)
code
try:
    print(0, end=';')
    print(8/3, end=';')
    print(1, end=';')
except KeyboardInterrupt:
    print(4, end=';')
except Exception:
    print(2, end=';')
except:
    print('ERR', end=';')
else:
    print(7, end=';')
finally:
    print(0)
code

6. Що буде виведено за виконання?
code
def f(n):
    print(n%10, end='')
    if n>9:
        f(n//10)
 
f(543216)
code

7. Що буде виведено за виконання?
code
class A:
    d = 1

    def _h1(self): print(2, end=';')
    def h2(self): print(3, end=';')

    def main(self):
        try: print(self.d, end=';')
        except: print('ERR', end=';')

        try: self._h1()
        except: print('ERR', end=';')

        try: self.h2()
        except: print('ERR', end=';')


class B(A):
    d = 4

    def _h1(self): print(5, end=';')
    def h2(self): print(6, end=';')

try: B().main()
except: print('ERR', end=';')
code

8. Що буде виведено за виконання?
code
def f(arg, n):
    n = 9
    tmp = arg
    tmp[2:8] = 'yz'
    return len(tmp)>3

try:
    b = ['0','1','2','3','4']
    n = 6
    f(b, n)
    print(n, end=';')
    for el in b:
        print(el, end=';')
except:
    print('ERR')
code

9. Що буде виведено за виконання?
code
class A:
    c = 1
    a = 2

    def __init__(self):
        c = 3

class B(A):
    c = 4

    def __init__(self):
        super().__init__()
        self.b = 7

obj = B()
obj.a = obj.a + 10
B.a = 13

try: print(obj.a, end= ';')
except: print('ERR', end=';')

try: print(A.a, end= ';')
except: print('ERR', end=';')

try: print(B.a, end= ';')
except: print('ERR', end=';')
code

10. 
Для графа на рисунку з вершини 4 запущено алгоритм пошуку в глибину, що будує кістякове дерево. Списки суміжності вершин переглядаються алгоритмом за зростанням міток вершин.\par
\begin{center}
	\adjustimage{max size={0.9\linewidth}{0.9\paperheight}}
	{../10/Traversals/92.png}
\end{center}
\par Які з наведених ребер входять до складу отриманого кістякового дерева?\par
1) (1, 6)\par
2) (0, 6)\par
3) (4, 5)\par
4) (1, 3)\par
5) (0, 3)\par
6) (0, 5)\par
{\vskip 15pt} У формі вибрати номери відповідних ребер.\par
%answer:1 5 6
\end {large}}
%end
{\smallskip \begin {small} Затверджено на засіданні кафедри теоретичної кібернетики, протокол \No 12 від   19.05.2020 р.\par\smallskip\smallskip
{\parbox {6.5 cm}{Зав. кафедри \dotfill Ю.В.~Крак}}\hspace {0.7cm}{\hbox to 6.5 cm{ Екзаменатор \dotfill Т.О.~Карнаух}} \end{small}}
%begin
{{\begin {small}\par
{ \center  Київський національний університет імені Тараса Шевченка \par}
{\parbox {5 cm}{Освітня програма \hfill}{\parbox {7 cm}{Прикладна математика, Системний аналіз \hfill}}\parbox {6 cm}{\hfill Семестр 2}}\par
{\parbox {5 cm} {Навчальний предмет \hfill}\parbox {7 cm}{Програмування \hfill}\parbox {6cm}{\hfill}\par}\vspace{-7pt} \end{small}}{\center Екзаменаційний білет \No \textbf 
{ 384 }\par}}
{
\begin {large}
1. Що буде виведено за виконання?
code
try:
    res = True==9.0 and not 8<6 or 5.0>4
    if isinstance(res, bool):
        print(f'bool;{res}')
    elif isinstance(res, int):
        print(f'int;{res}')
    elif isinstance(res, float):
        print(f'float;{res:.2f}')
    else:
        print('???')
except:
    print('ERR')
code

2. Що буде виведено за виконання?
code
a = 5
b = 8
c = 4

def f(a):
    a = 3
    b = 1
    c = 1
    return a + b + c

a = 6
b = 3
c = 9

try:
    print(f(b), a, b, c, sep=';')
except:
    print('ERR', a, b, c, sep=';')
code

3. Що буде виведено за виконання?
code
try:
    k = 7
    while k<=13:
        k += 2
    print(k, end=';')
except:
    print('ERR')
code

4. Що буде виведено за виконання?
code
try:
    for f in range(-4, -2, -1):
        if f>=8:
            continue
        print(f, end=';')
    else:
        print('end',end=';')
    print(f)
except:
    print('ERR')
code

5. Що буде виведено за виконання?\par (Дійсні значення, якщо наявні, в цьому завданні будуть виводитись у форматі з фіксованою крапкою та, можливо, з доволі великою кількістю десяткових розрядів. У відповіді для дійсних значень у ЦЬОМУ завданні достатньо вказати один десятковий розряд після крапки, відкинувши решту розрядів без округлення. Наприклад, замість 13.66666666666667 достатньо вказати 13.6, замість 25.23 достатньо вказати 25.2)
code
try:
    print(7, end=';')
    print(3j>=7j, end=';')
    print(8, end=';')
except KeyboardInterrupt:
    print(9, end=';')
except ValueError:
    print(5, end=';')
except:
    print('ERR', end=';')
else:
    print(6, end=';')
finally:
    print(4)
code

6. Що буде виведено за виконання?
code
def f(n):
    print(n%10,end='')
    if n>9:
        f(n//100)
 
f(123456)
code

7. Що буде виведено за виконання?
code
class A:
    c = 1

    def __f1(self): print(2, end=';')
    def _f2(self): print(3, end=';')

    def main(self):
        try: print(self.c, end=';')
        except: print('ERR', end=';')

        try: self.__f1()
        except: print('ERR', end=';')

        try: self._f2()
        except: print('ERR', end=';')


class B(A):
    c = 4

    def __f1(self): print(5, end=';')
    def _f2(self): print(6, end=';')

try: B().main()
except: print('ERR', end=';')
code

8. Що буде виведено за виконання?
code
def f(arg, n):
    n = 7
    tmp = arg
    tmp[-8:] = [1,0]
    return len(tmp)>3

try:
    h = [10,11,12,13,14]
    n = 8
    f(h, n)
    print(n, end=';')
    for el in h:
        print(el, end=';')
except:
    print('ERR')
code

9. Що буде виведено за виконання?
code
class A:
    b = 1
    a = 2

    def __init__(self):
        a = 3

class B(A):
    a = 4

    def __init__(self):
        super().__init__()
        self.c = 7

obj = B()
obj.b = obj.c + 10
B.b = 13

try: print(obj.b, end= ';')
except: print('ERR', end=';')

try: print(A.b, end= ';')
except: print('ERR', end=';')

try: print(B.b, end= ';')
except: print('ERR', end=';')
code

10. 
Для графа на рисунку з вершини 4 запущено алгоритм пошуку в глибину, що будує кістякове дерево. Списки суміжності вершин переглядаються алгоритмом за зростанням міток вершин.\par
\begin{center}
	\adjustimage{max size={0.9\linewidth}{0.9\paperheight}}
	{../10/Traversals/199.png}
\end{center}
\par Які з наведених ребер входять до складу отриманого кістякового дерева?\par
1) (5, 6)\par
2) (1, 3)\par
3) (0, 5)\par
4) (0, 2)\par
5) (3, 5)\par
6) (2, 4)\par
{\vskip 15pt} У формі вибрати номери відповідних ребер.\par
%answer:2 3 4 5 6
\end {large}}
%end
{\smallskip \begin {small} Затверджено на засіданні кафедри теоретичної кібернетики, протокол \No 12 від   19.05.2020 р.\par\smallskip\smallskip
{\parbox {6.5 cm}{Зав. кафедри \dotfill Ю.В.~Крак}}\hspace {0.7cm}{\hbox to 6.5 cm{ Екзаменатор \dotfill Т.О.~Карнаух}} \end{small}}
%begin
{{\begin {small}\par
{ \center  Київський національний університет імені Тараса Шевченка \par}
{\parbox {5 cm}{Освітня програма \hfill}{\parbox {7 cm}{Прикладна математика, Системний аналіз \hfill}}\parbox {6 cm}{\hfill Семестр 2}}\par
{\parbox {5 cm} {Навчальний предмет \hfill}\parbox {7 cm}{Програмування \hfill}\parbox {6cm}{\hfill}\par}\vspace{-7pt} \end{small}}{\center Екзаменаційний білет \No \textbf 
{ 385 }\par}}
{
\begin {large}
1. Що буде виведено за виконання?
code
try:
    res = "True">="1.0"
    if isinstance(res, bool):
        print(f'bool;{res}')
    elif isinstance(res, int):
        print(f'int;{res}')
    elif isinstance(res, float):
        print(f'float;{res:.2f}')
    else:
        print('???')
except:
    print('ERR')
code

2. Що буде виведено за виконання?
code
a = 0
b = 6
c = 5

def h(a):
    global c
    a = 3
    b = 4
    c = 5
    return a + b + c

a = 2
b = 8
c = 4

try:
    print(h(b), a, b, c, sep=';')
except:
    print('ERR', a, b, c, sep=';')
code

3. Що буде виведено за виконання?
code
try:
    j = 8
    while j>=0:
        j += -3
    print(j, end=';')
except:
    print('ERR')
code

4. Що буде виведено за виконання?
code
try:
    for f in range(10, 5, -3):
        if f>5:
            continue
        print(f, end=';');f = -2
    else:
        print(f,end=';')
    print(f)
except:
    print('ERR')
code

5. Що буде виведено за виконання?\par (Дійсні значення, якщо наявні, в цьому завданні будуть виводитись у форматі з фіксованою крапкою та, можливо, з доволі великою кількістю десяткових розрядів. У відповіді для дійсних значень у ЦЬОМУ завданні достатньо вказати один десятковий розряд після крапки, відкинувши решту розрядів без округлення. Наприклад, замість 13.66666666666667 достатньо вказати 13.6, замість 25.23 достатньо вказати 25.2)
code
try:
    print(2, end=';')
    print(1/0, end=';')
    print(7, end=';')
except TypeError:
    print(1, end=';')
except ZeroDivisionError:
    print(8, end=';')
except:
    print('ERR', end=';')
else:
    print(5, end=';')
finally:
    print(9)
code

6. Що буде виведено за виконання?
code
def f(n):
    print(n%2, end='')
    if n>2:
        f(n//2)
 
f(16)
code

7. Що буде виведено за виконання?
code
class A:
    d = 1

    def f1(self): print(2, end=';')
    def _f2(self): print(3, end=';')

    def main(self):
        try: print(self.d, end=';')
        except: print('ERR', end=';')

        try: self.f1()
        except: print('ERR', end=';')

        try: self._f2()
        except: print('ERR', end=';')


class B(A):
    d = 4

    def f1(self): print(5, end=';')
    def _f2(self): print(6, end=';')

try: B().main()
except: print('ERR', end=';')
code

8. Що буде виведено за виконання?
code
def f(arg, n):
    n = 3
    tmp = arg
    del tmp[4:-1]
    return len(tmp)>3

try:
    g = ['0','1','2','3','4']
    n = 2
    f(g, n)
    print(n, end=';')
    for el in g:
        print(el, end=';')
except:
    print('ERR')
code

9. Що буде виведено за виконання?
code
class A:
    c = 1
    b = 2

    def __init__(self):
        c = 3

class B(A):
    c = 4

    def __init__(self):
        super().__init__()
        self.a = 7

obj = B()
obj.a = obj.c + 10
B.a = 13

try: print(obj.a, end= ';')
except: print('ERR', end=';')

try: print(A.a, end= ';')
except: print('ERR', end=';')

try: print(B.a, end= ';')
except: print('ERR', end=';')
code

10. 
Для графа на рисунку з вершини 2 запущено алгоритм пошуку в глибину, що будує кістякове дерево. Списки суміжності вершин переглядаються алгоритмом за зростанням міток вершин.\par
\begin{center}
	\adjustimage{max size={0.9\linewidth}{0.9\paperheight}}
	{../10/Traversals/101.png}
\end{center}
\par Які з наведених ребер входять до складу отриманого кістякового дерева?\par
1) (2, 3)\par
2) (1, 3)\par
3) (2, 4)\par
4) (3, 6)\par
5) (3, 4)\par
6) (1, 4)\par
{\vskip 15pt} У формі вибрати номери відповідних ребер.\par
%answer:1 6
\end {large}}
%end
{\smallskip \begin {small} Затверджено на засіданні кафедри теоретичної кібернетики, протокол \No 12 від   19.05.2020 р.\par\smallskip\smallskip
{\parbox {6.5 cm}{Зав. кафедри \dotfill Ю.В.~Крак}}\hspace {0.7cm}{\hbox to 6.5 cm{ Екзаменатор \dotfill Т.О.~Карнаух}} \end{small}}
%begin
{{\begin {small}\par
{ \center  Київський національний університет імені Тараса Шевченка \par}
{\parbox {5 cm}{Освітня програма \hfill}{\parbox {7 cm}{Прикладна математика, Системний аналіз \hfill}}\parbox {6 cm}{\hfill Семестр 2}}\par
{\parbox {5 cm} {Навчальний предмет \hfill}\parbox {7 cm}{Програмування \hfill}\parbox {6cm}{\hfill}\par}\vspace{-7pt} \end{small}}{\center Екзаменаційний білет \No \textbf 
{ 386 }\par}}
{
\begin {large}
1. Що буде виведено за виконання?
code
try:
    res = 2>=5 or not 3!=2.0 and 7.0<=False
    if isinstance(res, bool):
        print(f'bool;{res}')
    elif isinstance(res, int):
        print(f'int;{res}')
    elif isinstance(res, float):
        print(f'float;{res:.2f}')
    else:
        print('???')
except:
    print('ERR')
code

2. Що буде виведено за виконання?
code
a = 0
b = 8
c = 4

def h(b):
    a = 4
    b += 1
    c = 2
    return a + b + c

a = 6
b = 1
c = 3

try:
    print(h(b), a, b, c, sep=';')
except:
    print('ERR', a, b, c, sep=';')
code

3. Що буде виведено за виконання?
code
try:
    t = 3
    while t<=10:
        t += 1
    print(t, end=';')
except:
    print('ERR')
code

4. Що буде виведено за виконання?
code
try:
    for b in range(-10, -10, -1):
        if b>=3:
            continue
        print(b, end=';')
        if b<=5:
            break
    else:
        print(b,end=';')
    print(b)
except:
    print('ERR')
code

5. Що буде виведено за виконання?\par (Дійсні значення, якщо наявні, в цьому завданні будуть виводитись у форматі з фіксованою крапкою та, можливо, з доволі великою кількістю десяткових розрядів. У відповіді для дійсних значень у ЦЬОМУ завданні достатньо вказати один десятковий розряд після крапки, відкинувши решту розрядів без округлення. Наприклад, замість 13.66666666666667 достатньо вказати 13.6, замість 25.23 достатньо вказати 25.2)
code
try:
    print(4, end=';')
    print(int('b3'), end=';')
    print(6, end=';')
except TypeError:
    print(0, end=';')
except Exception:
    print(2, end=';')
except:
    print('ERR', end=';')
else:
    print(2, end=';')
finally:
    print(8)
code

6. Що буде виведено за виконання?
code
def f(s):
    s1 = s
    for i in range(len(s)):
        s1[i] += 1
    return s1

s = [1, 2, 3]
f(s)
print(s)
code

7. Що буде виведено за виконання?
code
class A:
    b = 1

    def _h1(self): print(2, end=';')
    def h2(self): print(3, end=';')

    def main(self):
        try: print(self.b, end=';')
        except: print('ERR', end=';')

        try: self._h1()
        except: print('ERR', end=';')

        try: self.h2()
        except: print('ERR', end=';')


class B(A):
    b = 4

    def _h1(self): print(5, end=';')
    def h2(self): print(6, end=';')

try: B().main()
except: print('ERR', end=';')
code

8. Що буде виведено за виконання?
code
def f(arg, n):
    arg[90] = 9
    n = 7
    arg = tuple(arg.values())

try:
    n = 3
    d = {28:2, 85:0, 70:6, 70:1}
    f(d, n)
    print(n, end=';')
    for x, y in d.items():
        print(x, y, sep=';', end=';')
except:
    print('ERR')
code

9. Що буде виведено за виконання?
code
class A:
    c = 1
    a = 2

    def __init__(self):
        a = 3

class B(A):
    a = 4

    def __init__(self):
        super().__init__()
        self.b = 7

obj = B()
obj.b = obj.b + 10
B.b = 13

try: print(obj.b, end= ';')
except: print('ERR', end=';')

try: print(A.b, end= ';')
except: print('ERR', end=';')

try: print(B.b, end= ';')
except: print('ERR', end=';')
code

10. 
Для графа на рисунку з вершини 6 запущено алгоритм пошуку в глибину, що будує кістякове дерево. Списки суміжності вершин переглядаються алгоритмом за зростанням міток вершин.\par
\begin{center}
	\adjustimage{max size={0.9\linewidth}{0.9\paperheight}}
	{../10/Traversals/278.png}
\end{center}
\par Які з наведених ребер входять до складу отриманого кістякового дерева?\par
1) (4, 5)\par
2) (0, 5)\par
3) (1, 6)\par
4) (1, 5)\par
5) (1, 2)\par
6) (2, 6)\par
{\vskip 15pt} У формі вибрати номери відповідних ребер.\par
%answer:1 2 3 5
\end {large}}
%end
{\smallskip \begin {small} Затверджено на засіданні кафедри теоретичної кібернетики, протокол \No 12 від   19.05.2020 р.\par\smallskip\smallskip
{\parbox {6.5 cm}{Зав. кафедри \dotfill Ю.В.~Крак}}\hspace {0.7cm}{\hbox to 6.5 cm{ Екзаменатор \dotfill Т.О.~Карнаух}} \end{small}}
%begin
{{\begin {small}\par
{ \center  Київський національний університет імені Тараса Шевченка \par}
{\parbox {5 cm}{Освітня програма \hfill}{\parbox {7 cm}{Прикладна математика, Системний аналіз \hfill}}\parbox {6 cm}{\hfill Семестр 2}}\par
{\parbox {5 cm} {Навчальний предмет \hfill}\parbox {7 cm}{Програмування \hfill}\parbox {6cm}{\hfill}\par}\vspace{-7pt} \end{small}}{\center Екзаменаційний білет \No \textbf 
{ 387 }\par}}
{
\begin {large}
1. Що буде виведено за виконання?
code
try:
    res = 4.0j<"2j"
    if isinstance(res, bool):
        print(f'bool;{res}')
    elif isinstance(res, int):
        print(f'int;{res}')
    elif isinstance(res, float):
        print(f'float;{res:.2f}')
    else:
        print('???')
except:
    print('ERR')
code

2. Що буде виведено за виконання?
code
a = 9
b = 5
c = 5

def f(a):
    global c
    a = 5
    b *= 5
    c = 3
    return a + b + c

a = 7
b = 1
c = 2

try:
    print(f(a), a, b, c, sep=';')
except:
    print('ERR', a, b, c, sep=';')
code

3. Що буде виведено за виконання?
code
try:
    n = 6
    while n<5:
        n += 2
    print(n, end=';')
except:
    print('ERR')
code

4. Що буде виведено за виконання?
code
try:
    for e in range(13, -14, -2):
        if e<=8:
            break
        print(e, end=';')
    else:
        print(e,end=';')
    print(e)
except:
    print('ERR')
code

5. Що буде виведено за виконання?\par (Дійсні значення, якщо наявні, в цьому завданні будуть виводитись у форматі з фіксованою крапкою та, можливо, з доволі великою кількістю десяткових розрядів. У відповіді для дійсних значень у ЦЬОМУ завданні достатньо вказати один десятковий розряд після крапки, відкинувши решту розрядів без округлення. Наприклад, замість 13.66666666666667 достатньо вказати 13.6, замість 25.23 достатньо вказати 25.2)
code
try:
    print(4, end=';')
    print(int('0'), end=';')
    print(3, end=';')
except KeyboardInterrupt:
    print(1, end=';')
except ZeroDivisionError:
    print(5, end=';')
except:
    print('ERR', end=';')
else:
    print(7, end=';')
finally:
    print(9)
code

6. Що буде виведено за виконання?
code
def f(n):
    if n>9:
        return f(n//10)+1
    else:
        return 1

print(f(123456))
code

7. Що буде виведено за виконання?
code
class A:
    c = 1

    def g1(self): print(2, end=';')
    def _g2(self): print(3, end=';')

    def main(self):
        try: print(self.c, end=';')
        except: print('ERR', end=';')

        try: self.g1()
        except: print('ERR', end=';')

        try: self._g2()
        except: print('ERR', end=';')


class B(A):
    c = 4

    def g1(self): print(5, end=';')
    def _g2(self): print(6, end=';')

try: B().main()
except: print('ERR', end=';')
code

8. Що буде виведено за виконання?
code
def f(arg, n):
    n = 0
    tmp = arg
    tmp[-6:6] = [4,5]
    return len(tmp)>3

try:
    e = [10,11,12,13,14]
    n = 1
    f(e, n)
    print(n, end=';')
    for el in e:
        print(el, end=';')
except:
    print('ERR')
code

9. Що буде виведено за виконання?
code
class A:
    b = 1
    a = 2

    def __init__(self):
        b = 3

class B(A):
    b = 4

    def __init__(self):
        super().__init__()
        self.c = 7

obj = B()
obj.c = obj.a + 10
B.c = 13

try: print(obj.c, end= ';')
except: print('ERR', end=';')

try: print(A.c, end= ';')
except: print('ERR', end=';')

try: print(B.c, end= ';')
except: print('ERR', end=';')
code

10. 
Для графа на рисунку з вершини 2 запущено алгоритм пошуку в глибину, що будує кістякове дерево. Списки суміжності вершин переглядаються алгоритмом за зростанням міток вершин.\par
\begin{center}
	\adjustimage{max size={0.9\linewidth}{0.9\paperheight}}
	{../10/Traversals/5.png}
\end{center}
\par Які з наведених ребер входять до складу отриманого кістякового дерева?\par
1) (2, 6)\par
2) (3, 5)\par
3) (0, 4)\par
4) (2, 5)\par
5) (0, 3)\par
6) (0, 5)\par
{\vskip 15pt} У формі вибрати номери відповідних ребер.\par
%answer:2 3 5
\end {large}}
%end
{\smallskip \begin {small} Затверджено на засіданні кафедри теоретичної кібернетики, протокол \No 12 від   19.05.2020 р.\par\smallskip\smallskip
{\parbox {6.5 cm}{Зав. кафедри \dotfill Ю.В.~Крак}}\hspace {0.7cm}{\hbox to 6.5 cm{ Екзаменатор \dotfill Т.О.~Карнаух}} \end{small}}
%begin
{{\begin {small}\par
{ \center  Київський національний університет імені Тараса Шевченка \par}
{\parbox {5 cm}{Освітня програма \hfill}{\parbox {7 cm}{Прикладна математика, Системний аналіз \hfill}}\parbox {6 cm}{\hfill Семестр 2}}\par
{\parbox {5 cm} {Навчальний предмет \hfill}\parbox {7 cm}{Програмування \hfill}\parbox {6cm}{\hfill}\par}\vspace{-7pt} \end{small}}{\center Екзаменаційний білет \No \textbf 
{ 388 }\par}}
{
\begin {large}
1. Що буде виведено за виконання?
code
try:
    res = 7/6//6*9
    if isinstance(res, bool):
        print(f'bool;{res}')
    elif isinstance(res, int):
        print(f'int;{res}')
    elif isinstance(res, float):
        print(f'float;{res:.2f}')
    else:
        print('???')
except:
    print('ERR')
code

2. Що буде виведено за виконання?
code
a = 6
b = 3
c = 0

def g(a):
    global c
    a -= 2
    b = 1
    c = 4
    return a + b + c

a = 9
b = 4
c = 8

try:
    print(g(b), a, b, c, sep=';')
except:
    print('ERR', a, b, c, sep=';')
code

3. Що буде виведено за виконання?
code
try:
    m = 5
    while m<=8:
        m += 2
    print(m, end=';')
except:
    print('ERR')
code

4. Що буде виведено за виконання?
code
try:
    for d in range(7, -7, 2):
        if d<4:
            continue
        print(d, end=';')
    else:
        print('end',end=';')
    print(d)
except:
    print('ERR')
code

5. Що буде виведено за виконання?\par (Дійсні значення, якщо наявні, в цьому завданні будуть виводитись у форматі з фіксованою крапкою та, можливо, з доволі великою кількістю десяткових розрядів. У відповіді для дійсних значень у ЦЬОМУ завданні достатньо вказати один десятковий розряд після крапки, відкинувши решту розрядів без округлення. Наприклад, замість 13.66666666666667 достатньо вказати 13.6, замість 25.23 достатньо вказати 25.2)
code
try:
    print(6, end=';')
    print(9%2, end=';')
    print(3, end=';')
except ValueError:
    print(5, end=';')
except ZeroDivisionError:
    print(6, end=';')
except:
    print('ERR', end=';')
else:
    print(4, end=';')
finally:
    print(1)
code

6. Що буде виведено за виконання?
code
def f(s):
    for i in range(len(s)):
        s[i] += 1
        return
 
s = [1, 2, 3]
f(s)
print(s)
code

7. Що буде виведено за виконання?
code
class A:
    a = 1

    def _g1(self): print(2, end=';')
    def _g2(self): print(3, end=';')

    def main(self):
        try: print(self.a, end=';')
        except: print('ERR', end=';')

        try: self._g1()
        except: print('ERR', end=';')

        try: self._g2()
        except: print('ERR', end=';')


class B(A):
    a = 4

    def _g1(self): print(5, end=';')
    def _g2(self): print(6, end=';')

try: B().main()
except: print('ERR', end=';')
code

8. Що буде виведено за виконання?
code
def f(arg, n):
    n = 4
    tmp = arg
    del tmp[:-8]
    return len(tmp)>3

try:
    a = ['0','1','2','3','4']
    n = 5
    f(a, n)
    print(n, end=';')
    for el in a:
        print(el, end=';')
except:
    print('ERR')
code

9. Що буде виведено за виконання?
code
class A:
    a = 1
    c = 2

    def __init__(self):
        c = 3

class B(A):
    c = 4

    def __init__(self):
        super().__init__()
        self.b = 7

obj = B()
obj.b = obj.c + 10
B.b = 13

try: print(obj.b, end= ';')
except: print('ERR', end=';')

try: print(A.b, end= ';')
except: print('ERR', end=';')

try: print(B.b, end= ';')
except: print('ERR', end=';')
code

10. 
Для графа на рисунку з вершини 6 запущено алгоритм пошуку в глибину, що будує кістякове дерево. Списки суміжності вершин переглядаються алгоритмом за зростанням міток вершин.\par
\begin{center}
	\adjustimage{max size={0.9\linewidth}{0.9\paperheight}}
	{../10/Traversals/127.png}
\end{center}
\par Які з наведених ребер входять до складу отриманого кістякового дерева?\par
1) (2, 3)\par
2) (3, 5)\par
3) (0, 5)\par
4) (1, 6)\par
5) (3, 4)\par
6) (0, 6)\par
{\vskip 15pt} У формі вибрати номери відповідних ребер.\par
%answer:1 2 4 5 6
\end {large}}
%end
{\smallskip \begin {small} Затверджено на засіданні кафедри теоретичної кібернетики, протокол \No 12 від   19.05.2020 р.\par\smallskip\smallskip
{\parbox {6.5 cm}{Зав. кафедри \dotfill Ю.В.~Крак}}\hspace {0.7cm}{\hbox to 6.5 cm{ Екзаменатор \dotfill Т.О.~Карнаух}} \end{small}}
%begin
{{\begin {small}\par
{ \center  Київський національний університет імені Тараса Шевченка \par}
{\parbox {5 cm}{Освітня програма \hfill}{\parbox {7 cm}{Прикладна математика, Системний аналіз \hfill}}\parbox {6 cm}{\hfill Семестр 2}}\par
{\parbox {5 cm} {Навчальний предмет \hfill}\parbox {7 cm}{Програмування \hfill}\parbox {6cm}{\hfill}\par}\vspace{-7pt} \end{small}}{\center Екзаменаційний білет \No \textbf 
{ 389 }\par}}
{
\begin {large}
1. Що буде виведено за виконання?
code
try:
    res = 9//4*5*7
    if isinstance(res, bool):
        print(f'bool;{res}')
    elif isinstance(res, int):
        print(f'int;{res}')
    elif isinstance(res, float):
        print(f'float;{res:.2f}')
    else:
        print('???')
except:
    print('ERR')
code

2. Що буде виведено за виконання?
code
a = 7
b = 1
c = 2

def f(b):
    a = 2
    b = 3
    c = 2
    return a + b + c

a = 0
b = 8
c = 4

try:
    print(f(b), a, b, c, sep=';')
except:
    print('ERR', a, b, c, sep=';')
code

3. Що буде виведено за виконання?
code
try:
    n = 13
    while n>8:
        n += -2
    print(n, end=';')
except:
    print('ERR')
code

4. Що буде виведено за виконання?
code
try:
    for c in range(5, 11, -3):
        if c>=7:
            continue
        print(c, end=';');c = 3
    else:
        print(c,end=';')
    print(c)
except:
    print('ERR')
code

5. Що буде виведено за виконання?\par (Дійсні значення, якщо наявні, в цьому завданні будуть виводитись у форматі з фіксованою крапкою та, можливо, з доволі великою кількістю десяткових розрядів. У відповіді для дійсних значень у ЦЬОМУ завданні достатньо вказати один десятковий розряд після крапки, відкинувши решту розрядів без округлення. Наприклад, замість 13.66666666666667 достатньо вказати 13.6, замість 25.23 достатньо вказати 25.2)
code
try:
    print(8, end=';')
    print(2j<6j, end=';')
    print(0, end=';')
except KeyboardInterrupt:
    print(7, end=';')
except ValueError:
    print(5, end=';')
except:
    print('ERR', end=';')
else:
    print(9, end=';')
finally:
    print(8)
code

6. Що буде виведено за виконання?
code
def f(s1):
    s = s1[:]
    for i in range(len(s)):
        s[i] += 1
 
s = [1, 2, 3]
f(s)
print(s)
code

7. Що буде виведено за виконання?
code
class A:
    a = 1

    def __f1(self): print(2, end=';')
    def _f2(self): print(3, end=';')

    def main(self):
        try: print(self.a, end=';')
        except: print('ERR', end=';')

        try: self.__f1()
        except: print('ERR', end=';')

        try: self._f2()
        except: print('ERR', end=';')


class B(A):
    a = 4

    def __f1(self): print(5, end=';')
    def _f2(self): print(6, end=';')

try: B().main()
except: print('ERR', end=';')
code

8. Що буде виведено за виконання?
code
def f(arg, n):
    arg[40] = 2
    n -= 3

try:
    n = 6
    d = {47:8, 17:5, 60:6, 60:8}
    f(d, n)
    print(n, end=';')
    for x in d.values():
        print(x, sep=';', end=';')
except:
    print('ERR')
code

9. Що буде виведено за виконання?
code
class A:
    a = 1
    b = 2

    def __init__(self):
        b = 3

class B(A):
    b = 4

    def __init__(self):
        super().__init__()
        self.c = 7

obj = B()
obj.a = obj.b + 10
B.a = 13

try: print(obj.a, end= ';')
except: print('ERR', end=';')

try: print(A.a, end= ';')
except: print('ERR', end=';')

try: print(B.a, end= ';')
except: print('ERR', end=';')
code

10. 
Для графа на рисунку з вершини 5 запущено алгоритм пошуку в глибину, що будує кістякове дерево. Списки суміжності вершин переглядаються алгоритмом за зростанням міток вершин.\par
\begin{center}
	\adjustimage{max size={0.9\linewidth}{0.9\paperheight}}
	{../10/Traversals/106.png}
\end{center}
\par Які з наведених ребер входять до складу отриманого кістякового дерева?\par
1) (1, 5)\par
2) (3, 5)\par
3) (4, 6)\par
4) (0, 4)\par
5) (1, 2)\par
6) (3, 6)\par
{\vskip 15pt} У формі вибрати номери відповідних ребер.\par
%answer:1 3 4 5 6
\end {large}}
%end
{\smallskip \begin {small} Затверджено на засіданні кафедри теоретичної кібернетики, протокол \No 12 від   19.05.2020 р.\par\smallskip\smallskip
{\parbox {6.5 cm}{Зав. кафедри \dotfill Ю.В.~Крак}}\hspace {0.7cm}{\hbox to 6.5 cm{ Екзаменатор \dotfill Т.О.~Карнаух}} \end{small}}
%begin
{{\begin {small}\par
{ \center  Київський національний університет імені Тараса Шевченка \par}
{\parbox {5 cm}{Освітня програма \hfill}{\parbox {7 cm}{Прикладна математика, Системний аналіз \hfill}}\parbox {6 cm}{\hfill Семестр 2}}\par
{\parbox {5 cm} {Навчальний предмет \hfill}\parbox {7 cm}{Програмування \hfill}\parbox {6cm}{\hfill}\par}\vspace{-7pt} \end{small}}{\center Екзаменаційний білет \No \textbf 
{ 390 }\par}}
{
\begin {large}
1. Що буде виведено за виконання?
code
try:
    res = not 8.0<=8 and 7>2.0 or True>=6
    if isinstance(res, bool):
        print(f'bool;{res}')
    elif isinstance(res, int):
        print(f'int;{res}')
    elif isinstance(res, float):
        print(f'float;{res:.2f}')
    else:
        print('???')
except:
    print('ERR')
code

2. Що буде виведено за виконання?
code
a = 6
b = 5
c = 3

def g(b):
    global c
    a = 1
    b = 4
    c = 5
    return a + b + c

a = 7
b = 1
c = 9

try:
    print(g(a), a, b, c, sep=';')
except:
    print('ERR', a, b, c, sep=';')
code

3. Що буде виведено за виконання?
code
try:
    n = 9
    while n<11:
        n += 1
    print(n, end=';')
except:
    print('ERR')
code

4. Що буде виведено за виконання?
code
try:
    for a in range(-3, -12, 3):
        if a>6:
            break
        print(a, end=';')
    else:
        print(a,end=';')
    print(a)
except:
    print('ERR')
code

5. Що буде виведено за виконання?\par (Дійсні значення, якщо наявні, в цьому завданні будуть виводитись у форматі з фіксованою крапкою та, можливо, з доволі великою кількістю десяткових розрядів. У відповіді для дійсних значень у ЦЬОМУ завданні достатньо вказати один десятковий розряд після крапки, відкинувши решту розрядів без округлення. Наприклад, замість 13.66666666666667 достатньо вказати 13.6, замість 25.23 достатньо вказати 25.2)
code
try:
    print(2, end=';')
    print(1//0.0, end=';')
    print(0, end=';')
except Exception:
    print(4, end=';')
except TypeError:
    print(7, end=';')
except:
    print('ERR', end=';')
else:
    print(6, end=';')
finally:
    print(3)
code

6. Що буде виведено за виконання?
code
def f(n):
    if n>9: f(n//10)
    else:
        print(n%10, end='')
 
f(123456)
code

7. Що буде виведено за виконання?
code
class A:
    c = 1

    def _h1(self): print(2, end=';')
    def __h2(self): print(3, end=';')

    def main(self):
        try: print(self.c, end=';')
        except: print('ERR', end=';')

        try: self._h1()
        except: print('ERR', end=';')

        try: self.__h2()
        except: print('ERR', end=';')


class B(A):
    c = 4

    def _h1(self): print(5, end=';')
    def __h2(self): print(6, end=';')

try: B().main()
except: print('ERR', end=';')
code

8. Що буде виведено за виконання?
code
def f(arg, n):
    n = 7
    tmp = arg
    tmp[-9:3] = 'xy'
    return len(tmp)>3

try:
    h = [10,11,12,13,14]
    n = 0
    f(h, n)
    print(n, end=';')
    for el in h:
        print(el, end=';')
except:
    print('ERR')
code

9. Що буде виведено за виконання?
code
class A:
    b = 1
    c = 2

    def __init__(self):
        b = 3

class B(A):
    b = 4

    def __init__(self):
        super().__init__()
        self.a = 7

obj = B()
obj.c = obj.a + 10
B.c = 13

try: print(obj.c, end= ';')
except: print('ERR', end=';')

try: print(A.c, end= ';')
except: print('ERR', end=';')

try: print(B.c, end= ';')
except: print('ERR', end=';')
code

10. 
Для графа на рисунку з вершини 0 запущено алгоритм пошуку в глибину, що будує кістякове дерево. Списки суміжності вершин переглядаються алгоритмом за зростанням міток вершин.\par
\begin{center}
	\adjustimage{max size={0.9\linewidth}{0.9\paperheight}}
	{../10/Traversals/205.png}
\end{center}
\par Які з наведених ребер входять до складу отриманого кістякового дерева?\par
1) (5, 6)\par
2) (0, 5)\par
3) (0, 4)\par
4) (1, 6)\par
5) (3, 6)\par
6) (0, 6)\par
{\vskip 15pt} У формі вибрати номери відповідних ребер.\par
%answer:1 3 4
\end {large}}
%end
{\smallskip \begin {small} Затверджено на засіданні кафедри теоретичної кібернетики, протокол \No 12 від   19.05.2020 р.\par\smallskip\smallskip
{\parbox {6.5 cm}{Зав. кафедри \dotfill Ю.В.~Крак}}\hspace {0.7cm}{\hbox to 6.5 cm{ Екзаменатор \dotfill Т.О.~Карнаух}} \end{small}}
%begin
{{\begin {small}\par
{ \center  Київський національний університет імені Тараса Шевченка \par}
{\parbox {5 cm}{Освітня програма \hfill}{\parbox {7 cm}{Прикладна математика, Системний аналіз \hfill}}\parbox {6 cm}{\hfill Семестр 2}}\par
{\parbox {5 cm} {Навчальний предмет \hfill}\parbox {7 cm}{Програмування \hfill}\parbox {6cm}{\hfill}\par}\vspace{-7pt} \end{small}}{\center Екзаменаційний білет \No \textbf 
{ 391 }\par}}
{
\begin {large}
1. Що буде виведено за виконання?
code
try:
    res = 4**-3--2
    if isinstance(res, bool):
        print(f'bool;{res}')
    elif isinstance(res, int):
        print(f'int;{res}')
    elif isinstance(res, float):
        print(f'float;{res:.2f}')
    else:
        print('???')
except:
    print('ERR')
code

2. Що буде виведено за виконання?
code
a = 9
b = 3
c = 6

def f(b):
    a = 2
    b *= 3
    c = 5
    return a + b + c

a = 8
b = 4
c = 7

try:
    print(f(a), a, b, c, sep=';')
except:
    print('ERR', a, b, c, sep=';')
code

3. Що буде виведено за виконання?
code
try:
    t = 6
    while t>3:
        t += -3
    print(t, end=';')
except:
    print('ERR')
code

4. Що буде виведено за виконання?
code
try:
    for b in range(-12, 15, 2):
        if b<5:
            continue
        print(b, end=';');b = -5
    else:
        print('end',end=';')
    print(b)
except:
    print('ERR')
code

5. Що буде виведено за виконання?\par (Дійсні значення, якщо наявні, в цьому завданні будуть виводитись у форматі з фіксованою крапкою та, можливо, з доволі великою кількістю десяткових розрядів. У відповіді для дійсних значень у ЦЬОМУ завданні достатньо вказати один десятковий розряд після крапки, відкинувши решту розрядів без округлення. Наприклад, замість 13.66666666666667 достатньо вказати 13.6, замість 25.23 достатньо вказати 25.2)
code
try:
    print(6, end=';')
    print(4%1, end=';')
    print(7, end=';')
except TypeError:
    print(8, end=';')
except ValueError:
    print(0, end=';')
except:
    print('ERR', end=';')
else:
    print(3, end=';')
finally:
    print(5)
code

6. Що буде виведено за виконання?
code
def f(s):
    for i in range(len(s)):
        s[i] += 1
    return
 
s = [1, 2, 3]
f(s)
print(s)
code

7. Що буде виведено за виконання?
code
class A:
    b = 1

    def _h1(self): print(2, end=';')
    def _h2(self): print(3, end=';')

    def main(self):
        try: print(self.b, end=';')
        except: print('ERR', end=';')

        try: self._h1()
        except: print('ERR', end=';')

        try: self._h2()
        except: print('ERR', end=';')


class B(A):
    b = 4

    def _h1(self): print(5, end=';')
    def _h2(self): print(6, end=';')

try: B().main()
except: print('ERR', end=';')
code

8. Що буде виведено за виконання?
code
def f(arg, n):
    arg[13] = 8
    n = 8
    arg = set(arg.keys())

try:
    n = 4
    t = {22:4, 87:8, 50:4, 89:3}
    f(t, n)
    print(n, end=';')
    for x in t.items():
        print(*x, sep=';', end=';')
except:
    print('ERR')
code

9. Що буде виведено за виконання?
code
class A:
    b = 1
    a = 2

    def __init__(self):
        a = 3

class B(A):
    b = 4

    def __init__(self):
        super().__init__()
        self.c = 7

obj = B()
obj.c = obj.b + 10
B.c = 13

try: print(obj.c, end= ';')
except: print('ERR', end=';')

try: print(A.c, end= ';')
except: print('ERR', end=';')

try: print(B.c, end= ';')
except: print('ERR', end=';')
code

10. 
Для графа на рисунку з вершини 5 запущено алгоритм пошуку в глибину, що будує кістякове дерево. Списки суміжності вершин переглядаються алгоритмом за зростанням міток вершин.\par
\begin{center}
	\adjustimage{max size={0.9\linewidth}{0.9\paperheight}}
	{../10/Traversals/158.png}
\end{center}
\par Які з наведених ребер входять до складу отриманого кістякового дерева?\par
1) (3, 4)\par
2) (2, 6)\par
3) (2, 4)\par
4) (2, 5)\par
5) (5, 6)\par
6) (0, 4)\par
{\vskip 15pt} У формі вибрати номери відповідних ребер.\par
%answer:4 6
\end {large}}
%end
{\smallskip \begin {small} Затверджено на засіданні кафедри теоретичної кібернетики, протокол \No 12 від   19.05.2020 р.\par\smallskip\smallskip
{\parbox {6.5 cm}{Зав. кафедри \dotfill Ю.В.~Крак}}\hspace {0.7cm}{\hbox to 6.5 cm{ Екзаменатор \dotfill Т.О.~Карнаух}} \end{small}}
%begin
{{\begin {small}\par
{ \center  Київський національний університет імені Тараса Шевченка \par}
{\parbox {5 cm}{Освітня програма \hfill}{\parbox {7 cm}{Прикладна математика, Системний аналіз \hfill}}\parbox {6 cm}{\hfill Семестр 2}}\par
{\parbox {5 cm} {Навчальний предмет \hfill}\parbox {7 cm}{Програмування \hfill}\parbox {6cm}{\hfill}\par}\vspace{-7pt} \end{small}}{\center Екзаменаційний білет \No \textbf 
{ 392 }\par}}
{
\begin {large}
1. Що буде виведено за виконання?
code
try:
    res = 9==False
    if isinstance(res, bool):
        print(f'bool;{res}')
    elif isinstance(res, int):
        print(f'int;{res}')
    elif isinstance(res, float):
        print(f'float;{res:.2f}')
    else:
        print('???')
except:
    print('ERR')
code

2. Що буде виведено за виконання?
code
a = 3
b = 1
c = 6

def g(a):
    a = 1
    b = 4
    c = 5
    return a + b + c

a = 2
b = 9
c = 8

try:
    print(g(b), a, b, c, sep=';')
except:
    print('ERR', a, b, c, sep=';')
code

3. Що буде виведено за виконання?
code
try:
    k = 12
    while k>=6:
        k += -2
    print(k, end=';')
except:
    print('ERR')
code

4. Що буде виведено за виконання?
code
try:
    for c in range(-14, 2, -1):
        if c<=4:
            break
        print(c, end=';')
    else:
        print(c,end=';')
    print(c)
except:
    print('ERR')
code

5. Що буде виведено за виконання?\par (Дійсні значення, якщо наявні, в цьому завданні будуть виводитись у форматі з фіксованою крапкою та, можливо, з доволі великою кількістю десяткових розрядів. У відповіді для дійсних значень у ЦЬОМУ завданні достатньо вказати один десятковий розряд після крапки, відкинувши решту розрядів без округлення. Наприклад, замість 13.66666666666667 достатньо вказати 13.6, замість 25.23 достатньо вказати 25.2)
code
try:
    print(2, end=';')
    print(int('d1'), end=';')
    print(9, end=';')
except ZeroDivisionError:
    print(4, end=';')
except Exception:
    print(3, end=';')
except:
    print('ERR', end=';')
else:
    print(9, end=';')
finally:
    print(6)
code

6. Що буде виведено за виконання?
code
def f(n):
    if n>9:
        f(n//100)
    else:
        print(n%10,end='')

f(123456)
code

7. Що буде виведено за виконання?
code
class A:
    d = 1

    def __f1(self): print(2, end=';')
    def _f2(self): print(3, end=';')

    def main(self):
        try: print(self.d, end=';')
        except: print('ERR', end=';')

        try: self.__f1()
        except: print('ERR', end=';')

        try: self._f2()
        except: print('ERR', end=';')


class B(A):
    d = 4

    def __f1(self): print(5, end=';')
    def _f2(self): print(6, end=';')

try: B().main()
except: print('ERR', end=';')
code

8. Що буде виведено за виконання?
code
def f(arg, n):
    arg[39] = 3
    n = 7
    arg = list(arg.items())

try:
    n = 1
    s = {72:3, 17:2, 39:6, 17:2}
    f(s, n)
    print(n, end=';')
    for x in s.values():
        print(x, sep=';', end=';')
except:
    print('ERR')
code

9. Що буде виведено за виконання?
code
class A:
    a = 1
    c = 2

    def __init__(self):
        a = 3

class B(A):
    c = 4

    def __init__(self):
        super().__init__()
        self.b = 7

obj = B()
obj.a = obj.a + 10
B.a = 13

try: print(obj.a, end= ';')
except: print('ERR', end=';')

try: print(A.a, end= ';')
except: print('ERR', end=';')

try: print(B.a, end= ';')
except: print('ERR', end=';')
code

10. 
Для графа на рисунку з вершини 1 запущено алгоритм пошуку в глибину, що будує кістякове дерево. Списки суміжності вершин переглядаються алгоритмом за зростанням міток вершин.\par
\begin{center}
	\adjustimage{max size={0.9\linewidth}{0.9\paperheight}}
	{../10/Traversals/57.png}
\end{center}
\par Які з наведених ребер входять до складу отриманого кістякового дерева?\par
1) (1, 6)\par
2) (2, 5)\par
3) (4, 5)\par
4) (1, 3)\par
5) (1, 5)\par
6) (0, 4)\par
{\vskip 15pt} У формі вибрати номери відповідних ребер.\par
%answer:3 6
\end {large}}
%end
{\smallskip \begin {small} Затверджено на засіданні кафедри теоретичної кібернетики, протокол \No 12 від   19.05.2020 р.\par\smallskip\smallskip
{\parbox {6.5 cm}{Зав. кафедри \dotfill Ю.В.~Крак}}\hspace {0.7cm}{\hbox to 6.5 cm{ Екзаменатор \dotfill Т.О.~Карнаух}} \end{small}}
%begin
{{\begin {small}\par
{ \center  Київський національний університет імені Тараса Шевченка \par}
{\parbox {5 cm}{Освітня програма \hfill}{\parbox {7 cm}{Прикладна математика, Системний аналіз \hfill}}\parbox {6 cm}{\hfill Семестр 2}}\par
{\parbox {5 cm} {Навчальний предмет \hfill}\parbox {7 cm}{Програмування \hfill}\parbox {6cm}{\hfill}\par}\vspace{-7pt} \end{small}}{\center Екзаменаційний білет \No \textbf 
{ 393 }\par}}
{
\begin {large}
1. Що буде виведено за виконання?
code
try:
    res = 6/6%3*4
    if isinstance(res, bool):
        print(f'bool;{res}')
    elif isinstance(res, int):
        print(f'int;{res}')
    elif isinstance(res, float):
        print(f'float;{res:.2f}')
    else:
        print('???')
except:
    print('ERR')
code

2. Що буде виведено за виконання?
code
a = 4
b = 5
c = 0

def f(b):
    global c
    a = 3
    b = 2
    c = 4
    return a + b + c

a = 7
b = 7
c = 1

try:
    print(f(a), a, b, c, sep=';')
except:
    print('ERR', a, b, c, sep=';')
code

3. Що буде виведено за виконання?
code
try:
    t = 11
    while t>4:
        t += -3
    print(t, end=';')
except:
    print('ERR')
code

4. Що буде виведено за виконання?
code
try:
    for f in range(15, 3, -2):
        if f<8:
            continue
        print(f, end=';')
    else:
        print(f,end=';')
    print(f)
except:
    print('ERR')
code

5. Що буде виведено за виконання?\par (Дійсні значення, якщо наявні, в цьому завданні будуть виводитись у форматі з фіксованою крапкою та, можливо, з доволі великою кількістю десяткових розрядів. У відповіді для дійсних значень у ЦЬОМУ завданні достатньо вказати один десятковий розряд після крапки, відкинувши решту розрядів без округлення. Наприклад, замість 13.66666666666667 достатньо вказати 13.6, замість 25.23 достатньо вказати 25.2)
code
try:
    print(5, end=';')
    print(7j==2j, end=';')
    print(0, end=';')
except KeyboardInterrupt:
    print(8, end=';')
except ZeroDivisionError:
    print(1, end=';')
except:
    print('ERR', end=';')
else:
    print(2, end=';')
finally:
    print(2)
code

6. Що буде виведено за виконання?
code
def f(s):
    for el in s:
        el += 1
        return s
 
s = [1, 2, 3]
s = f(s)
print(s)
code

7. Що буде виведено за виконання?
code
class A:
    a = 1

    def g1(self): print(2, end=';')
    def _g2(self): print(3, end=';')

    def main(self):
        try: print(self.a, end=';')
        except: print('ERR', end=';')

        try: self.g1()
        except: print('ERR', end=';')

        try: self._g2()
        except: print('ERR', end=';')


class B(A):
    a = 4

    def g1(self): print(5, end=';')
    def _g2(self): print(6, end=';')

try: B().main()
except: print('ERR', end=';')
code

8. Що буде виведено за виконання?
code
def f(arg, n):
    arg[75] = 1
    n = 9
    arg = tuple(arg.values())

try:
    n = 2
    d = {22:9, 19:0, 67:2, 67:8}
    f(d, n)
    print(n, end=';')
    for x in d.keys():
        print(x, sep=';', end=';')
except:
    print('ERR')
code

9. Що буде виведено за виконання?
code
class A:
    b = 1
    c = 2

    def __init__(self):
        b = 3

class B(A):
    b = 4

    def __init__(self):
        super().__init__()
        self.a = 7

obj = B()
obj.b = obj.c + 10
B.b = 13

try: print(obj.b, end= ';')
except: print('ERR', end=';')

try: print(A.b, end= ';')
except: print('ERR', end=';')

try: print(B.b, end= ';')
except: print('ERR', end=';')
code

10. 
Для графа на рисунку з вершини 1 запущено алгоритм пошуку в глибину, що будує кістякове дерево. Списки суміжності вершин переглядаються алгоритмом за зростанням міток вершин.\par
\begin{center}
	\adjustimage{max size={0.9\linewidth}{0.9\paperheight}}
	{../10/Traversals/185.png}
\end{center}
\par Які з наведених ребер входять до складу отриманого кістякового дерева?\par
1) (0, 4)\par
2) (4, 6)\par
3) (1, 2)\par
4) (1, 5)\par
5) (0, 3)\par
6) (1, 6)\par
{\vskip 15pt} У формі вибрати номери відповідних ребер.\par
%answer:2 3 5
\end {large}}
%end
{\smallskip \begin {small} Затверджено на засіданні кафедри теоретичної кібернетики, протокол \No 12 від   19.05.2020 р.\par\smallskip\smallskip
{\parbox {6.5 cm}{Зав. кафедри \dotfill Ю.В.~Крак}}\hspace {0.7cm}{\hbox to 6.5 cm{ Екзаменатор \dotfill Т.О.~Карнаух}} \end{small}}
%begin
{{\begin {small}\par
{ \center  Київський національний університет імені Тараса Шевченка \par}
{\parbox {5 cm}{Освітня програма \hfill}{\parbox {7 cm}{Прикладна математика, Системний аналіз \hfill}}\parbox {6 cm}{\hfill Семестр 2}}\par
{\parbox {5 cm} {Навчальний предмет \hfill}\parbox {7 cm}{Програмування \hfill}\parbox {6cm}{\hfill}\par}\vspace{-7pt} \end{small}}{\center Екзаменаційний білет \No \textbf 
{ 394 }\par}}
{
\begin {large}
1. Що буде виведено за виконання?
code
try:
    res = 4/8//9*8
    if isinstance(res, bool):
        print(f'bool;{res}')
    elif isinstance(res, int):
        print(f'int;{res}')
    elif isinstance(res, float):
        print(f'float;{res:.2f}')
    else:
        print('???')
except:
    print('ERR')
code

2. Що буде виведено за виконання?
code
a = 2
b = 0
c = 5

def g(b):
    global c
    a -= 1
    b = 4
    c = 1
    return a + b + c

a = 1
b = 0
c = 3

try:
    print(g(a), a, b, c, sep=';')
except:
    print('ERR', a, b, c, sep=';')
code

3. Що буде виведено за виконання?
code
try:
    k = 8
    while k>=1:
        k += -2
    print(k, end=';')
except:
    print('ERR')
code

4. Що буде виведено за виконання?
code
try:
    for e in range(8, -15, -3):
        if e>=6:
            continue
        print(e, end=';')
    else:
        print(e,end=';')
    print(e)
except:
    print('ERR')
code

5. Що буде виведено за виконання?\par (Дійсні значення, якщо наявні, в цьому завданні будуть виводитись у форматі з фіксованою крапкою та, можливо, з доволі великою кількістю десяткових розрядів. У відповіді для дійсних значень у ЦЬОМУ завданні достатньо вказати один десятковий розряд після крапки, відкинувши решту розрядів без округлення. Наприклад, замість 13.66666666666667 достатньо вказати 13.6, замість 25.23 достатньо вказати 25.2)
code
try:
    print(8, end=';')
    print(9//0, end=';')
    print(0, end=';')
except Exception:
    print(1, end=';')
except ValueError:
    print(3, end=';')
except:
    print('ERR', end=';')
else:
    print(5, end=';')
finally:
    print(6)
code

6. Що буде виведено за виконання?
code
def f(s):
    for i in range(len(s)):
        s[i] += 1
    return
 
s = [1, 2, 3]
f(s)
print(s)
code

7. Що буде виведено за виконання?
code
class A:
    b = 1

    def _g1(self): print(2, end=';')
    def __g2(self): print(3, end=';')

    def main(self):
        try: print(self.b, end=';')
        except: print('ERR', end=';')

        try: self._g1()
        except: print('ERR', end=';')

        try: self.__g2()
        except: print('ERR', end=';')


class B(A):
    b = 4

    def _g1(self): print(5, end=';')
    def __g2(self): print(6, end=';')

try: B().main()
except: print('ERR', end=';')
code

8. Що буде виведено за виконання?
code
def f(arg, n):
    arg[22] = 1
    n += -3

try:
    n = 7
    s = {59:3, 47:0, 69:7, 59:1}
    f(s, n)
    print(n, end=';')
    for x in s.items():
        print(x, sep=';', end=';')
except:
    print('ERR')
code

9. Що буде виведено за виконання?
code
class A:
    c = 1
    a = 2

    def __init__(self):
        a = 3

class B(A):
    a = 4

    def __init__(self):
        super().__init__()
        self.b = 7

obj = B()
obj.c = obj.b + 10
B.c = 13

try: print(obj.c, end= ';')
except: print('ERR', end=';')

try: print(A.c, end= ';')
except: print('ERR', end=';')

try: print(B.c, end= ';')
except: print('ERR', end=';')
code

10. 
Для графа на рисунку з вершини 3 запущено алгоритм пошуку в глибину, що будує кістякове дерево. Списки суміжності вершин переглядаються алгоритмом за зростанням міток вершин.\par
\begin{center}
	\adjustimage{max size={0.9\linewidth}{0.9\paperheight}}
	{../10/Traversals/296.png}
\end{center}
\par Які з наведених ребер входять до складу отриманого кістякового дерева?\par
1) (1, 4)\par
2) (4, 6)\par
3) (2, 5)\par
4) (2, 4)\par
5) (0, 1)\par
6) (1, 3)\par
{\vskip 15pt} У формі вибрати номери відповідних ребер.\par
%answer:1 2 3 4 5
\end {large}}
%end
{\smallskip \begin {small} Затверджено на засіданні кафедри теоретичної кібернетики, протокол \No 12 від   19.05.2020 р.\par\smallskip\smallskip
{\parbox {6.5 cm}{Зав. кафедри \dotfill Ю.В.~Крак}}\hspace {0.7cm}{\hbox to 6.5 cm{ Екзаменатор \dotfill Т.О.~Карнаух}} \end{small}}
%begin
{{\begin {small}\par
{ \center  Київський національний університет імені Тараса Шевченка \par}
{\parbox {5 cm}{Освітня програма \hfill}{\parbox {7 cm}{Прикладна математика, Системний аналіз \hfill}}\parbox {6 cm}{\hfill Семестр 2}}\par
{\parbox {5 cm} {Навчальний предмет \hfill}\parbox {7 cm}{Програмування \hfill}\parbox {6cm}{\hfill}\par}\vspace{-7pt} \end{small}}{\center Екзаменаційний білет \No \textbf 
{ 395 }\par}}
{
\begin {large}
1. Що буде виведено за виконання?
code
try:
    res = 5//5/7*7
    if isinstance(res, bool):
        print(f'bool;{res}')
    elif isinstance(res, int):
        print(f'int;{res}')
    elif isinstance(res, float):
        print(f'float;{res:.2f}')
    else:
        print('???')
except:
    print('ERR')
code

2. Що буде виведено за виконання?
code
a = 4
b = 8
c = 7

def h(a):
    a += 2
    b = 4
    c = 5
    return a + b + c

a = 9
b = 5
c = 6

try:
    print(h(b), a, b, c, sep=';')
except:
    print('ERR', a, b, c, sep=';')
code

3. Що буде виведено за виконання?
code
try:
    k = 0
    while k<=12:
        k += 1
    print(k, end=';')
except:
    print('ERR')
code

4. Що буде виведено за виконання?
code
try:
    for a in range(-6, -9, 3):
        if a>3:
            continue
        print(a, end=';');a = 9
    else:
        print(a,end=';')
    print(a)
except:
    print('ERR')
code

5. Що буде виведено за виконання?\par (Дійсні значення, якщо наявні, в цьому завданні будуть виводитись у форматі з фіксованою крапкою та, можливо, з доволі великою кількістю десяткових розрядів. У відповіді для дійсних значень у ЦЬОМУ завданні достатньо вказати один десятковий розряд після крапки, відкинувши решту розрядів без округлення. Наприклад, замість 13.66666666666667 достатньо вказати 13.6, замість 25.23 достатньо вказати 25.2)
code
try:
    print(4, end=';')
    print(int('7'), end=';')
    print(5, end=';')
except TypeError:
    print(4, end=';')
except KeyboardInterrupt:
    print(8, end=';')
except:
    print('ERR', end=';')
else:
    print(0, end=';')
finally:
    print(1)
code

6. Що буде виведено за виконання?
code
def f(n):
    if n>9:
        f(n//10)
    print(n%10, end='')
 
f(543216)
code

7. Що буде виведено за виконання?
code
class A:
    c = 1

    def _f1(self): print(2, end=';')
    def f2(self): print(3, end=';')

    def main(self):
        try: print(self.c, end=';')
        except: print('ERR', end=';')

        try: self._f1()
        except: print('ERR', end=';')

        try: self.f2()
        except: print('ERR', end=';')


class B(A):
    c = 4

    def _f1(self): print(5, end=';')
    def f2(self): print(6, end=';')

try: B().main()
except: print('ERR', end=';')
code

8. Що буде виведено за виконання?
code
def f(arg, n):
    arg[61] = 5
    n *= -2

try:
    n = 5
    t = {61:1, 86:5, 86:9}
    f(t, n)
    print(n, end=';')
    for x in t.values():
        print(x, sep=';', end=';')
except:
    print('ERR')
code

9. Що буде виведено за виконання?
code
class A:
    c = 1
    b = 2

    def __init__(self):
        b = 3

class B(A):
    b = 4

    def __init__(self):
        super().__init__()
        self.a = 7

obj = B()
obj.a = obj.b + 10
B.a = 13

try: print(obj.a, end= ';')
except: print('ERR', end=';')

try: print(A.a, end= ';')
except: print('ERR', end=';')

try: print(B.a, end= ';')
except: print('ERR', end=';')
code

10. 
Для графа на рисунку з вершини 3 запущено алгоритм пошуку в глибину, що будує кістякове дерево. Списки суміжності вершин переглядаються алгоритмом за зростанням міток вершин.\par
\begin{center}
	\adjustimage{max size={0.9\linewidth}{0.9\paperheight}}
	{../10/Traversals/118.png}
\end{center}
\par Які з наведених ребер входять до складу отриманого кістякового дерева?\par
1) (1, 3)\par
2) (0, 4)\par
3) (3, 5)\par
4) (0, 3)\par
5) (1, 4)\par
6) (5, 6)\par
{\vskip 15pt} У формі вибрати номери відповідних ребер.\par
%answer:4 5 6
\end {large}}
%end
{\smallskip \begin {small} Затверджено на засіданні кафедри теоретичної кібернетики, протокол \No 12 від   19.05.2020 р.\par\smallskip\smallskip
{\parbox {6.5 cm}{Зав. кафедри \dotfill Ю.В.~Крак}}\hspace {0.7cm}{\hbox to 6.5 cm{ Екзаменатор \dotfill Т.О.~Карнаух}} \end{small}}
%begin
{{\begin {small}\par
{ \center  Київський національний університет імені Тараса Шевченка \par}
{\parbox {5 cm}{Освітня програма \hfill}{\parbox {7 cm}{Прикладна математика, Системний аналіз \hfill}}\parbox {6 cm}{\hfill Семестр 2}}\par
{\parbox {5 cm} {Навчальний предмет \hfill}\parbox {7 cm}{Програмування \hfill}\parbox {6cm}{\hfill}\par}\vspace{-7pt} \end{small}}{\center Екзаменаційний білет \No \textbf 
{ 396 }\par}}
{
\begin {large}
1. Що буде виведено за виконання?
code
try:
    res = 5.0<4.0 or not 9==7 and 5!=False
    if isinstance(res, bool):
        print(f'bool;{res}')
    elif isinstance(res, int):
        print(f'int;{res}')
    elif isinstance(res, float):
        print(f'float;{res:.2f}')
    else:
        print('???')
except:
    print('ERR')
code

2. Що буде виведено за виконання?
code
a = 4
b = 0
c = 9

def h(a):
    a = 3
    b = 1
    c = 2
    return a + b + c

a = 6
b = 2
c = 5

try:
    print(h(a), a, b, c, sep=';')
except:
    print('ERR', a, b, c, sep=';')
code

3. Що буде виведено за виконання?
code
try:
    i = 15
    while i>7:
        i += -2
    print(i, end=';')
except:
    print('ERR')
code

4. Що буде виведено за виконання?
code
try:
    for d in range(6, 4, 3):
        if d<=7:
            break
        print(d, end=';')
    else:
        print(d,end=';')
    print(d)
except:
    print('ERR')
code

5. Що буде виведено за виконання?\par (Дійсні значення, якщо наявні, в цьому завданні будуть виводитись у форматі з фіксованою крапкою та, можливо, з доволі великою кількістю десяткових розрядів. У відповіді для дійсних значень у ЦЬОМУ завданні достатньо вказати один десятковий розряд після крапки, відкинувши решту розрядів без округлення. Наприклад, замість 13.66666666666667 достатньо вказати 13.6, замість 25.23 достатньо вказати 25.2)
code
try:
    print(3, end=';')
    print(2/3, end=';')
    print(6, end=';')
except Exception:
    print(7, end=';')
except ZeroDivisionError:
    print(9, end=';')
except:
    print('ERR', end=';')
else:
    print(0, end=';')
finally:
    print(1)
code

6. Що буде виведено за виконання?
code
def f(n):
    if n>9: f(n//10)
    else:
        print(n%10, end='')
 
f(123456)
code

7. Що буде виведено за виконання?
code
class A:
    d = 1

    def h1(self): print(2, end=';')
    def _h2(self): print(3, end=';')

    def main(self):
        try: print(self.d, end=';')
        except: print('ERR', end=';')

        try: self.h1()
        except: print('ERR', end=';')

        try: self._h2()
        except: print('ERR', end=';')


class B(A):
    d = 4

    def h1(self): print(5, end=';')
    def _h2(self): print(6, end=';')

try: B().main()
except: print('ERR', end=';')
code

8. Що буде виведено за виконання?
code
def f(arg, n):
    arg[16] = 4
    n = 5
    arg = list(arg.keys())

try:
    n = 4
    s = {60:5, 16:1, 76:8, 16:6}
    f(s, n)
    print(n, end=';')
    for x in s.values():
        print(x, sep=';', end=';')
except:
    print('ERR')
code

9. Що буде виведено за виконання?
code
class A:
    a = 1
    b = 2

    def __init__(self):
        a = 3

class B(A):
    a = 4

    def __init__(self):
        super().__init__()
        self.c = 7

obj = B()
obj.a = obj.c + 10
B.a = 13

try: print(obj.a, end= ';')
except: print('ERR', end=';')

try: print(A.a, end= ';')
except: print('ERR', end=';')

try: print(B.a, end= ';')
except: print('ERR', end=';')
code

10. 
Для графа на рисунку з вершини 1 запущено алгоритм пошуку в глибину, що будує кістякове дерево. Списки суміжності вершин переглядаються алгоритмом за зростанням міток вершин.\par
\begin{center}
	\adjustimage{max size={0.9\linewidth}{0.9\paperheight}}
	{../10/Traversals/254.png}
\end{center}
\par Які з наведених ребер входять до складу отриманого кістякового дерева?\par
1) (0, 2)\par
2) (0, 1)\par
3) (2, 4)\par
4) (0, 5)\par
5) (0, 3)\par
6) (1, 5)\par
{\vskip 15pt} У формі вибрати номери відповідних ребер.\par
%answer:1 2 3 4 5
\end {large}}
%end
{\smallskip \begin {small} Затверджено на засіданні кафедри теоретичної кібернетики, протокол \No 12 від   19.05.2020 р.\par\smallskip\smallskip
{\parbox {6.5 cm}{Зав. кафедри \dotfill Ю.В.~Крак}}\hspace {0.7cm}{\hbox to 6.5 cm{ Екзаменатор \dotfill Т.О.~Карнаух}} \end{small}}
%begin
{{\begin {small}\par
{ \center  Київський національний університет імені Тараса Шевченка \par}
{\parbox {5 cm}{Освітня програма \hfill}{\parbox {7 cm}{Прикладна математика, Системний аналіз \hfill}}\parbox {6 cm}{\hfill Семестр 2}}\par
{\parbox {5 cm} {Навчальний предмет \hfill}\parbox {7 cm}{Програмування \hfill}\parbox {6cm}{\hfill}\par}\vspace{-7pt} \end{small}}{\center Екзаменаційний білет \No \textbf 
{ 397 }\par}}
{
\begin {large}
1. Що буде виведено за виконання?
code
try:
    res = "6.0j"==False
    if isinstance(res, bool):
        print(f'bool;{res}')
    elif isinstance(res, int):
        print(f'int;{res}')
    elif isinstance(res, float):
        print(f'float;{res:.2f}')
    else:
        print('???')
except:
    print('ERR')
code

2. Що буде виведено за виконання?
code
a = 1
b = 2
c = 0

def g(a):
    a = 3
    b *= 5
    c = 4
    return a + b + c

a = 4
b = 9
c = 5

try:
    print(g(a), a, b, c, sep=';')
except:
    print('ERR', a, b, c, sep=';')
code

3. Що буде виведено за виконання?
code
try:
    j = 13
    while j>=4:
        j += -3
    print(j, end=';')
except:
    print('ERR')
code

4. Що буде виведено за виконання?
code
try:
    for b in range(3, -5, -3):
        if b<5:
            continue
        print(b, end=';')
    else:
        print(b,end=';')
    print(b)
except:
    print('ERR')
code

5. Що буде виведено за виконання?\par (Дійсні значення, якщо наявні, в цьому завданні будуть виводитись у форматі з фіксованою крапкою та, можливо, з доволі великою кількістю десяткових розрядів. У відповіді для дійсних значень у ЦЬОМУ завданні достатньо вказати один десятковий розряд після крапки, відкинувши решту розрядів без округлення. Наприклад, замість 13.66666666666667 достатньо вказати 13.6, замість 25.23 достатньо вказати 25.2)
code
try:
    print(7, end=';')
    print(int('9'), end=';')
    print(3, end=';')
except ValueError:
    print(2, end=';')
except TypeError:
    print(6, end=';')
except:
    print('ERR', end=';')
else:
    print(8, end=';')
finally:
    print(5)
code

6. Що буде виведено за виконання?
code
def f(s):
    s1 = s
    for i in range(len(s)):
        s1[i] += 1
        return s1

s = [1, 2, 3]
f(s)
print(s)
code

7. Що буде виведено за виконання?
code
class A:
    b = 1

    def _g1(self): print(2, end=';')
    def g2(self): print(3, end=';')

    def main(self):
        try: print(self.b, end=';')
        except: print('ERR', end=';')

        try: self._g1()
        except: print('ERR', end=';')

        try: self.g2()
        except: print('ERR', end=';')


class B(A):
    b = 4

    def _g1(self): print(5, end=';')
    def g2(self): print(6, end=';')

try: B().main()
except: print('ERR', end=';')
code

8. Що буде виведено за виконання?
code
def f(arg, n):
    arg[14] = 6
    n -= 3

try:
    n = 4
    s = {88:6, 26:2, 85:0, 38:5}
    f(s, n)
    print(n, end=';')
    for x in s :
        print(x, sep=';', end=';')
except:
    print('ERR')
code

9. Що буде виведено за виконання?
code
class A:
    b = 1
    c = 2

    def __init__(self):
        c = 3

class B(A):
    b = 4

    def __init__(self):
        super().__init__()
        self.a = 7

obj = B()
obj.b = obj.c + 10
B.b = 13

try: print(obj.b, end= ';')
except: print('ERR', end=';')

try: print(A.b, end= ';')
except: print('ERR', end=';')

try: print(B.b, end= ';')
except: print('ERR', end=';')
code

10. 
Для графа на рисунку з вершини 6 запущено алгоритм пошуку в глибину, що будує кістякове дерево. Списки суміжності вершин переглядаються алгоритмом за зростанням міток вершин.\par
\begin{center}
	\adjustimage{max size={0.9\linewidth}{0.9\paperheight}}
	{../10/Traversals/99.png}
\end{center}
\par Які з наведених ребер входять до складу отриманого кістякового дерева?\par
1) (1, 4)\par
2) (4, 6)\par
3) (3, 4)\par
4) (2, 4)\par
5) (4, 5)\par
6) (2, 6)\par
{\vskip 15pt} У формі вибрати номери відповідних ребер.\par
%answer:3 5
\end {large}}
%end
{\smallskip \begin {small} Затверджено на засіданні кафедри теоретичної кібернетики, протокол \No 12 від   19.05.2020 р.\par\smallskip\smallskip
{\parbox {6.5 cm}{Зав. кафедри \dotfill Ю.В.~Крак}}\hspace {0.7cm}{\hbox to 6.5 cm{ Екзаменатор \dotfill Т.О.~Карнаух}} \end{small}}
%begin
{{\begin {small}\par
{ \center  Київський національний університет імені Тараса Шевченка \par}
{\parbox {5 cm}{Освітня програма \hfill}{\parbox {7 cm}{Прикладна математика, Системний аналіз \hfill}}\parbox {6 cm}{\hfill Семестр 2}}\par
{\parbox {5 cm} {Навчальний предмет \hfill}\parbox {7 cm}{Програмування \hfill}\parbox {6cm}{\hfill}\par}\vspace{-7pt} \end{small}}{\center Екзаменаційний білет \No \textbf 
{ 398 }\par}}
{
\begin {large}
1. Що буде виведено за виконання?
code
try:
    res = 3.0<=6 and not 2<True or 7.0!=8
    if isinstance(res, bool):
        print(f'bool;{res}')
    elif isinstance(res, int):
        print(f'int;{res}')
    elif isinstance(res, float):
        print(f'float;{res:.2f}')
    else:
        print('???')
except:
    print('ERR')
code

2. Що буде виведено за виконання?
code
a = 8
b = 3
c = 8

def f(a):
    global c
    a = 5
    b = 4
    c = 2
    return a + b + c

a = 3
b = 5
c = 6

try:
    print(f(a), a, b, c, sep=';')
except:
    print('ERR', a, b, c, sep=';')
code

3. Що буде виведено за виконання?
code
try:
    m = 7
    while m>0:
        m += -2
    print(m, end=';')
except:
    print('ERR')
code

4. Що буде виведено за виконання?
code
try:
    for e in range(-7, 13, -2):
        if e<=8:
            continue
        print(e, end=';')
    else:
        print(e,end=';')
    print(e)
except:
    print('ERR')
code

5. Що буде виведено за виконання?\par (Дійсні значення, якщо наявні, в цьому завданні будуть виводитись у форматі з фіксованою крапкою та, можливо, з доволі великою кількістю десяткових розрядів. У відповіді для дійсних значень у ЦЬОМУ завданні достатньо вказати один десятковий розряд після крапки, відкинувши решту розрядів без округлення. Наприклад, замість 13.66666666666667 достатньо вказати 13.6, замість 25.23 достатньо вказати 25.2)
code
try:
    print(4, end=';')
    print(0j>1j, end=';')
    print(5, end=';')
except KeyboardInterrupt:
    print(7, end=';')
except KeyboardInterrupt:
    print(8, end=';')
except:
    print('ERR', end=';')
else:
    print(3, end=';')
finally:
    print(9)
code

6. Що буде виведено за виконання?
code
def f(n):
    if n>9:
        f(n//10)
    else:
        print(n%10, end='')

f(543216)
code

7. Що буде виведено за виконання?
code
class A:
    d = 1

    def _h1(self): print(2, end=';')
    def __h2(self): print(3, end=';')

    def main(self):
        try: print(self.d, end=';')
        except: print('ERR', end=';')

        try: self._h1()
        except: print('ERR', end=';')

        try: self.__h2()
        except: print('ERR', end=';')


class B(A):
    d = 4

    def _h1(self): print(5, end=';')
    def __h2(self): print(6, end=';')

try: B().main()
except: print('ERR', end=';')
code

8. Що буде виведено за виконання?
code
def f(arg, n):
    n = 1
    tmp = arg
    tmp[:1] = ()
    return len(tmp)>3

try:
    f = ['0','1','2','3','4']
    n = 3
    f(f, n)
    print(n, end=';')
    for el in f:
        print(el, end=';')
except:
    print('ERR')
code

9. Що буде виведено за виконання?
code
class A:
    b = 1
    a = 2

    def __init__(self):
        b = 3

class B(A):
    a = 4

    def __init__(self):
        super().__init__()
        self.c = 7

obj = B()
obj.a = obj.c + 10
B.a = 13

try: print(obj.a, end= ';')
except: print('ERR', end=';')

try: print(A.a, end= ';')
except: print('ERR', end=';')

try: print(B.a, end= ';')
except: print('ERR', end=';')
code

10. 
Для графа на рисунку з вершини 4 запущено алгоритм пошуку в глибину, що будує кістякове дерево. Списки суміжності вершин переглядаються алгоритмом за зростанням міток вершин.\par
\begin{center}
	\adjustimage{max size={0.9\linewidth}{0.9\paperheight}}
	{../10/Traversals/284.png}
\end{center}
\par Які з наведених ребер входять до складу отриманого кістякового дерева?\par
1) (1, 2)\par
2) (2, 5)\par
3) (3, 4)\par
4) (4, 6)\par
5) (3, 5)\par
6) (2, 4)\par
{\vskip 15pt} У формі вибрати номери відповідних ребер.\par
%answer:5 6
\end {large}}
%end
{\smallskip \begin {small} Затверджено на засіданні кафедри теоретичної кібернетики, протокол \No 12 від   19.05.2020 р.\par\smallskip\smallskip
{\parbox {6.5 cm}{Зав. кафедри \dotfill Ю.В.~Крак}}\hspace {0.7cm}{\hbox to 6.5 cm{ Екзаменатор \dotfill Т.О.~Карнаух}} \end{small}}
%begin
{{\begin {small}\par
{ \center  Київський національний університет імені Тараса Шевченка \par}
{\parbox {5 cm}{Освітня програма \hfill}{\parbox {7 cm}{Прикладна математика, Системний аналіз \hfill}}\parbox {6 cm}{\hfill Семестр 2}}\par
{\parbox {5 cm} {Навчальний предмет \hfill}\parbox {7 cm}{Програмування \hfill}\parbox {6cm}{\hfill}\par}\vspace{-7pt} \end{small}}{\center Екзаменаційний білет \No \textbf 
{ 399 }\par}}
{
\begin {large}
1. Що буде виведено за виконання?
code
try:
    res = not False>=9.0 or 9>6.0 and 4==3
    if isinstance(res, bool):
        print(f'bool;{res}')
    elif isinstance(res, int):
        print(f'int;{res}')
    elif isinstance(res, float):
        print(f'float;{res:.2f}')
    else:
        print('???')
except:
    print('ERR')
code

2. Що буде виведено за виконання?
code
a = 1
b = 9
c = 7

def h(b):
    a = 3
    b = 5
    c = 1
    return a + b + c

a = 4
b = 0
c = 2

try:
    print(h(b), a, b, c, sep=';')
except:
    print('ERR', a, b, c, sep=';')
code

3. Що буде виведено за виконання?
code
try:
    t = 1
    while t<9:
        t += 2
    print(t, end=';')
except:
    print('ERR')
code

4. Що буде виведено за виконання?
code
try:
    for d in range(-2, 8, -1):
        if d>=3:
            break
        print(d, end=';');d = 6
        if d>=8:
            break
    else:
        print('end',end=';')
    print(d)
except:
    print('ERR')
code

5. Що буде виведено за виконання?\par (Дійсні значення, якщо наявні, в цьому завданні будуть виводитись у форматі з фіксованою крапкою та, можливо, з доволі великою кількістю десяткових розрядів. У відповіді для дійсних значень у ЦЬОМУ завданні достатньо вказати один десятковий розряд після крапки, відкинувши решту розрядів без округлення. Наприклад, замість 13.66666666666667 достатньо вказати 13.6, замість 25.23 достатньо вказати 25.2)
code
try:
    print(4, end=';')
    print(2%0.0, end=';')
    print(6, end=';')
except Exception:
    print(1, end=';')
except ValueError:
    print(1, end=';')
except:
    print('ERR', end=';')
else:
    print(6, end=';')
finally:
    print(5)
code

6. Що буде виведено за виконання?
code
def f(n):
    if n<=9:
        print(n%10,end='')
    else:
        f(n//10)

f(543216)
code

7. Що буде виведено за виконання?
code
class A:
    a = 1

    def _f1(self): print(2, end=';')
    def _f2(self): print(3, end=';')

    def main(self):
        try: print(self.a, end=';')
        except: print('ERR', end=';')

        try: self._f1()
        except: print('ERR', end=';')

        try: self._f2()
        except: print('ERR', end=';')


class B(A):
    a = 4

    def _f1(self): print(5, end=';')
    def _f2(self): print(6, end=';')

try: B().main()
except: print('ERR', end=';')
code

8. Що буде виведено за виконання?
code
def f(arg, n):
    n = 5
    tmp = arg
    del tmp[8:4]
    return len(tmp)>3

try:
    g = [10,11,12,13,14]
    n = 2
    f(g, n)
    print(n, end=';')
    for el in g:
        print(el, end=';')
except:
    print('ERR')
code

9. Що буде виведено за виконання?
code
class A:
    a = 1
    c = 2

    def __init__(self):
        a = 3

class B(A):
    c = 4

    def __init__(self):
        super().__init__()
        self.b = 7

obj = B()
obj.a = obj.b + 10
B.a = 13

try: print(obj.a, end= ';')
except: print('ERR', end=';')

try: print(A.a, end= ';')
except: print('ERR', end=';')

try: print(B.a, end= ';')
except: print('ERR', end=';')
code

10. 
Для графа на рисунку з вершини 6 запущено алгоритм пошуку в глибину, що будує кістякове дерево. Списки суміжності вершин переглядаються алгоритмом за зростанням міток вершин.\par
\begin{center}
	\adjustimage{max size={0.9\linewidth}{0.9\paperheight}}
	{../10/Traversals/86.png}
\end{center}
\par Які з наведених ребер входять до складу отриманого кістякового дерева?\par
1) (2, 3)\par
2) (0, 6)\par
3) (4, 5)\par
4) (0, 5)\par
5) (5, 6)\par
6) (2, 6)\par
{\vskip 15pt} У формі вибрати номери відповідних ребер.\par
%answer:1 2 4
\end {large}}
%end
{\smallskip \begin {small} Затверджено на засіданні кафедри теоретичної кібернетики, протокол \No 12 від   19.05.2020 р.\par\smallskip\smallskip
{\parbox {6.5 cm}{Зав. кафедри \dotfill Ю.В.~Крак}}\hspace {0.7cm}{\hbox to 6.5 cm{ Екзаменатор \dotfill Т.О.~Карнаух}} \end{small}}
%begin
{{\begin {small}\par
{ \center  Київський національний університет імені Тараса Шевченка \par}
{\parbox {5 cm}{Освітня програма \hfill}{\parbox {7 cm}{Прикладна математика, Системний аналіз \hfill}}\parbox {6 cm}{\hfill Семестр 2}}\par
{\parbox {5 cm} {Навчальний предмет \hfill}\parbox {7 cm}{Програмування \hfill}\parbox {6cm}{\hfill}\par}\vspace{-7pt} \end{small}}{\center Екзаменаційний білет \No \textbf 
{ 400 }\par}}
{
\begin {large}
1. Що буде виведено за виконання?
code
try:
    res = "3"!="True"
    if isinstance(res, bool):
        print(f'bool;{res}')
    elif isinstance(res, int):
        print(f'int;{res}')
    elif isinstance(res, float):
        print(f'float;{res:.2f}')
    else:
        print('???')
except:
    print('ERR')
code

2. Що буде виведено за виконання?
code
a = 1
b = 3
c = 6

def f(b):
    global c
    a -= 1
    b = 3
    c = 2
    return a + b + c

a = 7
b = 2
c = 8

try:
    print(f(a), a, b, c, sep=';')
except:
    print('ERR', a, b, c, sep=';')
code

3. Що буде виведено за виконання?
code
try:
    n = 6
    while n>=9:
        n += -3
    print(n, end=';')
except:
    print('ERR')
code

4. Що буде виведено за виконання?
code
try:
    for a in range(1, 10, 2):
        if a>7:
            continue
        print(a, end=';')
    else:
        print(a,end=';')
    print(a)
except:
    print('ERR')
code

5. Що буде виведено за виконання?\par (Дійсні значення, якщо наявні, в цьому завданні будуть виводитись у форматі з фіксованою крапкою та, можливо, з доволі великою кількістю десяткових розрядів. У відповіді для дійсних значень у ЦЬОМУ завданні достатньо вказати один десятковий розряд після крапки, відкинувши решту розрядів без округлення. Наприклад, замість 13.66666666666667 достатньо вказати 13.6, замість 25.23 достатньо вказати 25.2)
code
try:
    print(8, end=';')
    print(int('c3'), end=';')
    print(9, end=';')
except TypeError:
    print(0, end=';')
except ZeroDivisionError:
    print(7, end=';')
except:
    print('ERR', end=';')
else:
    print(4, end=';')
finally:
    print(2)
code

6. Що буде виведено за виконання?
code
def f(n):
    if n>9:
        f(n//100)
    print(n%10,end='')
 
f(987654)
code

7. Що буде виведено за виконання?
code
class A:
    c = 1

    def __h1(self): print(2, end=';')
    def _h2(self): print(3, end=';')

    def main(self):
        try: print(self.c, end=';')
        except: print('ERR', end=';')

        try: self.__h1()
        except: print('ERR', end=';')

        try: self._h2()
        except: print('ERR', end=';')


class B(A):
    c = 4

    def __h1(self): print(5, end=';')
    def _h2(self): print(6, end=';')

try: B().main()
except: print('ERR', end=';')
code

8. Що буде виведено за виконання?
code
def f(arg, n):
    arg[16] = 5
    n = 8
    arg = tuple(arg.values())

try:
    n = 0
    t = {46:3, 21:7, 16:5, 21:8}
    f(t, n)
    print(n, end=';')
    for x in t.items():
        print(*x, sep=';', end=';')
except:
    print('ERR')
code

9. Що буде виведено за виконання?
code
class A:
    c = 1
    b = 2

    def __init__(self):
        b = 3

class B(A):
    c = 4

    def __init__(self):
        super().__init__()
        self.a = 7

obj = B()
obj.b = obj.a + 10
B.b = 13

try: print(obj.b, end= ';')
except: print('ERR', end=';')

try: print(A.b, end= ';')
except: print('ERR', end=';')

try: print(B.b, end= ';')
except: print('ERR', end=';')
code

10. 
Для графа на рисунку з вершини 1 запущено алгоритм пошуку в глибину, що будує кістякове дерево. Списки суміжності вершин переглядаються алгоритмом за зростанням міток вершин.\par
\begin{center}
	\adjustimage{max size={0.9\linewidth}{0.9\paperheight}}
	{../10/Traversals/115.png}
\end{center}
\par Які з наведених ребер входять до складу отриманого кістякового дерева?\par
1) (3, 6)\par
2) (1, 2)\par
3) (4, 5)\par
4) (0, 5)\par
5) (5, 6)\par
6) (1, 4)\par
{\vskip 15pt} У формі вибрати номери відповідних ребер.\par
%answer:2 4 5
\end {large}}
%end
{\smallskip \begin {small} Затверджено на засіданні кафедри теоретичної кібернетики, протокол \No 12 від   19.05.2020 р.\par\smallskip\smallskip
{\parbox {6.5 cm}{Зав. кафедри \dotfill Ю.В.~Крак}}\hspace {0.7cm}{\hbox to 6.5 cm{ Екзаменатор \dotfill Т.О.~Карнаух}} \end{small}}
%begin
{{\begin {small}\par
{ \center  Київський національний університет імені Тараса Шевченка \par}
{\parbox {5 cm}{Освітня програма \hfill}{\parbox {7 cm}{Прикладна математика, Системний аналіз \hfill}}\parbox {6 cm}{\hfill Семестр 2}}\par
{\parbox {5 cm} {Навчальний предмет \hfill}\parbox {7 cm}{Програмування \hfill}\parbox {6cm}{\hfill}\par}\vspace{-7pt} \end{small}}{\center Екзаменаційний білет \No \textbf 
{ 401 }\par}}
{
\begin {large}
1. Що буде виведено за виконання?
code
def f(n):
    print('f', end='')
    return n

try:
    print(f(-7) and f(8))
except:
    print('ERR')
code

2. Що буде виведено за виконання?
code
a = 3
b = 9
c = 8

def h(b):
    global c
    a += 1
    b = 4
    c = 2
    return a + b + c

a = 5
b = 0
c = 1

try:
    print(h(a), a, b, c, sep=';')
except:
    print('ERR', a, b, c, sep=';')
code

3. Що буде виведено за виконання?
code
try:
    i = 7
    while i<=7:
        i += 1
    print(i, end=';')
except:
    print('ERR')
code

4. Що буде виведено за виконання?
code
try:
    for c in range(-1, -4, 2):
        if c>6:
            break
        print(c, end=';');c = 0
    else:
        print(c,end=';')
    print(c)
except:
    print('ERR')
code

5. Що буде виведено за виконання?\par (Дійсні значення, якщо наявні, в цьому завданні будуть виводитись у форматі з фіксованою крапкою та, можливо, з доволі великою кількістю десяткових розрядів. У відповіді для дійсних значень у ЦЬОМУ завданні достатньо вказати один десятковий розряд після крапки, відкинувши решту розрядів без округлення. Наприклад, замість 13.66666666666667 достатньо вказати 13.6, замість 25.23 достатньо вказати 25.2)
code
try:
    print(0, end=';')
    print(int('8'), end=';')
    print(4, end=';')
except ValueError:
    print(5, end=';')
except TypeError:
    print(6, end=';')
except:
    print('ERR', end=';')
else:
    print(7, end=';')
finally:
    print(3)
code

6. Що буде виведено за виконання?
code
def f(n):
    print(n%10, end='')
    if n>9:
        f(n//100)
 
f(987654)
code

7. Що буде виведено за виконання?
code
class A:
    b = 1

    def __g1(self): print(2, end=';')
    def _g2(self): print(3, end=';')

    def main(self):
        try: print(self.b, end=';')
        except: print('ERR', end=';')

        try: self.__g1()
        except: print('ERR', end=';')

        try: self._g2()
        except: print('ERR', end=';')


class B(A):
    b = 4

    def __g1(self): print(5, end=';')
    def _g2(self): print(6, end=';')

try: B().main()
except: print('ERR', end=';')
code

8. Що буде виведено за виконання?
code
def f(arg, n):
    arg[48] = 0
    n = 6
    arg = set(arg.items())

try:
    n = 3
    d = {52:7, 77:0, 68:0, 57:7}
    f(d, n)
    print(n, end=';')
    for x in d.keys():
        print(x, sep=';', end=';')
except:
    print('ERR')
code

9. Що буде виведено за виконання?
code
class A:
    a = 1
    b = 2

    def __init__(self):
        b = 3

class B(A):
    a = 4

    def __init__(self):
        super().__init__()
        self.c = 7

obj = B()
obj.c = obj.b + 10
B.c = 13

try: print(obj.c, end= ';')
except: print('ERR', end=';')

try: print(A.c, end= ';')
except: print('ERR', end=';')

try: print(B.c, end= ';')
except: print('ERR', end=';')
code

10. 
Для графа на рисунку з вершини 2 запущено алгоритм пошуку в глибину, що будує кістякове дерево. Списки суміжності вершин переглядаються алгоритмом за зростанням міток вершин.\par
\begin{center}
	\adjustimage{max size={0.9\linewidth}{0.9\paperheight}}
	{../10/Traversals/151.png}
\end{center}
\par Які з наведених ребер входять до складу отриманого кістякового дерева?\par
1) (1, 6)\par
2) (0, 6)\par
3) (2, 5)\par
4) (0, 5)\par
5) (1, 2)\par
6) (0, 1)\par
{\vskip 15pt} У формі вибрати номери відповідних ребер.\par
%answer:5 6
\end {large}}
%end
{\smallskip \begin {small} Затверджено на засіданні кафедри теоретичної кібернетики, протокол \No 12 від   19.05.2020 р.\par\smallskip\smallskip
{\parbox {6.5 cm}{Зав. кафедри \dotfill Ю.В.~Крак}}\hspace {0.7cm}{\hbox to 6.5 cm{ Екзаменатор \dotfill Т.О.~Карнаух}} \end{small}}
%begin
{{\begin {small}\par
{ \center  Київський національний університет імені Тараса Шевченка \par}
{\parbox {5 cm}{Освітня програма \hfill}{\parbox {7 cm}{Прикладна математика, Системний аналіз \hfill}}\parbox {6 cm}{\hfill Семестр 2}}\par
{\parbox {5 cm} {Навчальний предмет \hfill}\parbox {7 cm}{Програмування \hfill}\parbox {6cm}{\hfill}\par}\vspace{-7pt} \end{small}}{\center Екзаменаційний білет \No \textbf 
{ 402 }\par}}
{
\begin {large}
1. Що буде виведено за виконання?
code
try:
    res = 7//3/8*3
    if isinstance(res, bool):
        print(f'bool;{res}')
    elif isinstance(res, int):
        print(f'int;{res}')
    elif isinstance(res, float):
        print(f'float;{res:.2f}')
    else:
        print('???')
except:
    print('ERR')
code

2. Що буде виведено за виконання?
code
a = 4
b = 8
c = 6

def g(a):
    a = 1
    b = 2
    c = 3
    return a + b + c

a = 2
b = 3
c = 7

try:
    print(g(a), a, b, c, sep=';')
except:
    print('ERR', a, b, c, sep=';')
code

3. Що буде виведено за виконання?
code
try:
    j = 9
    while j>3:
        j += -3
    print(j, end=';')
except:
    print('ERR')
code

4. Що буде виведено за виконання?
code
try:
    for f in range(9, 14, -3):
        if f<4:
            continue
        print(f, end=';')
    else:
        print(f,end=';')
    print(f)
except:
    print('ERR')
code

5. Що буде виведено за виконання?\par (Дійсні значення, якщо наявні, в цьому завданні будуть виводитись у форматі з фіксованою крапкою та, можливо, з доволі великою кількістю десяткових розрядів. У відповіді для дійсних значень у ЦЬОМУ завданні достатньо вказати один десятковий розряд після крапки, відкинувши решту розрядів без округлення. Наприклад, замість 13.66666666666667 достатньо вказати 13.6, замість 25.23 достатньо вказати 25.2)
code
try:
    print(2, end=';')
    print(9j<4j, end=';')
    print(1, end=';')
except KeyboardInterrupt:
    print(4, end=';')
except ZeroDivisionError:
    print(3, end=';')
except:
    print('ERR', end=';')
else:
    print(2, end=';')
finally:
    print(9)
code

6. Що буде виведено за виконання?
code
def f(s):
    s1 = s
    for i in range(len(s)):
        s1[i] += 1
    return s1

s = [1, 2, 3]
f(s)
print(s)
code

7. Що буде виведено за виконання?
code
class A:
    c = 1

    def _f1(self): print(2, end=';')
    def _f2(self): print(3, end=';')

    def main(self):
        try: print(self.c, end=';')
        except: print('ERR', end=';')

        try: self._f1()
        except: print('ERR', end=';')

        try: self._f2()
        except: print('ERR', end=';')


class B(A):
    c = 4

    def _f1(self): print(5, end=';')
    def _f2(self): print(6, end=';')

try: B().main()
except: print('ERR', end=';')
code

8. Що буде виведено за виконання?
code
def f(arg, n):
    arg[26] = 7
    n *= -3

try:
    n = 6
    d = {66:5, 51:0, 43:2}
    f(d, n)
    print(n, end=';')
    for x in d :
        print(x, sep=';', end=';')
except:
    print('ERR')
code

9. Що буде виведено за виконання?
code
class A:
    c = 1
    a = 2

    def __init__(self):
        c = 3

class B(A):
    a = 4

    def __init__(self):
        super().__init__()
        self.b = 7

obj = B()
obj.c = obj.a + 10
B.c = 13

try: print(obj.c, end= ';')
except: print('ERR', end=';')

try: print(A.c, end= ';')
except: print('ERR', end=';')

try: print(B.c, end= ';')
except: print('ERR', end=';')
code

10. 
Для графа на рисунку з вершини 3 запущено алгоритм пошуку в глибину, що будує кістякове дерево. Списки суміжності вершин переглядаються алгоритмом за зростанням міток вершин.\par
\begin{center}
	\adjustimage{max size={0.9\linewidth}{0.9\paperheight}}
	{../10/Traversals/145.png}
\end{center}
\par Які з наведених ребер входять до складу отриманого кістякового дерева?\par
1) (0, 3)\par
2) (0, 1)\par
3) (1, 2)\par
4) (2, 3)\par
5) (4, 6)\par
6) (1, 4)\par
{\vskip 15pt} У формі вибрати номери відповідних ребер.\par
%answer:1 2 3 5 6
\end {large}}
%end
{\smallskip \begin {small} Затверджено на засіданні кафедри теоретичної кібернетики, протокол \No 12 від   19.05.2020 р.\par\smallskip\smallskip
{\parbox {6.5 cm}{Зав. кафедри \dotfill Ю.В.~Крак}}\hspace {0.7cm}{\hbox to 6.5 cm{ Екзаменатор \dotfill Т.О.~Карнаух}} \end{small}}
%begin
{{\begin {small}\par
{ \center  Київський національний університет імені Тараса Шевченка \par}
{\parbox {5 cm}{Освітня програма \hfill}{\parbox {7 cm}{Прикладна математика, Системний аналіз \hfill}}\parbox {6 cm}{\hfill Семестр 2}}\par
{\parbox {5 cm} {Навчальний предмет \hfill}\parbox {7 cm}{Програмування \hfill}\parbox {6cm}{\hfill}\par}\vspace{-7pt} \end{small}}{\center Екзаменаційний білет \No \textbf 
{ 403 }\par}}
{
\begin {large}
1. Що буде виведено за виконання?
code
def f(n):
    print('f', end='')
    return n==2

try:
    print(f(9) or f(-2))
except:
    print('ERR')
code

2. Що буде виведено за виконання?
code
a = 2
b = 6
c = 4

def f(a):
    a -= 3
    b = 5
    c = 1
    return a + b + c

a = 7
b = 5
c = 9

try:
    print(f(b), a, b, c, sep=';')
except:
    print('ERR', a, b, c, sep=';')
code

3. Що буде виведено за виконання?
code
try:
    j = 2
    while j<6:
        j += 2
    print(j, end=';')
except:
    print('ERR')
code

4. Що буде виведено за виконання?
code
try:
    for a in range(-11, -3, 3):
        if a<=6:
            break
        print(a, end=';')
    else:
        print(a,end=';')
    print(a)
except:
    print('ERR')
code

5. Що буде виведено за виконання?\par (Дійсні значення, якщо наявні, в цьому завданні будуть виводитись у форматі з фіксованою крапкою та, можливо, з доволі великою кількістю десяткових розрядів. У відповіді для дійсних значень у ЦЬОМУ завданні достатньо вказати один десятковий розряд після крапки, відкинувши решту розрядів без округлення. Наприклад, замість 13.66666666666667 достатньо вказати 13.6, замість 25.23 достатньо вказати 25.2)
code
try:
    print(5, end=';')
    print(0/2, end=';')
    print(6, end=';')
except Exception:
    print(8, end=';')
except ValueError:
    print(7, end=';')
except:
    print('ERR', end=';')
else:
    print(1, end=';')
finally:
    print(4)
code

6. Що буде виведено за виконання?
code
def f(n):
    if n>9:
        return f(n//100)+1
    else:
        return 1

print(f(123456))
code

7. Що буде виведено за виконання?
code
class A:
    d = 1

    def _g1(self): print(2, end=';')
    def __g2(self): print(3, end=';')

    def main(self):
        try: print(self.d, end=';')
        except: print('ERR', end=';')

        try: self._g1()
        except: print('ERR', end=';')

        try: self.__g2()
        except: print('ERR', end=';')


class B(A):
    d = 4

    def _g1(self): print(5, end=';')
    def __g2(self): print(6, end=';')

try: B().main()
except: print('ERR', end=';')
code

8. Що буде виведено за виконання?
code
def f(arg, n):
    arg[54] = 1
    n += 2

try:
    n = 3
    t = {84:3, 85:0, 86:3}
    f(t, n)
    print(n, end=';')
    for x in t.keys():
        print(x, sep=';', end=';')
except:
    print('ERR')
code

9. Що буде виведено за виконання?
code
class A:
    a = 1
    c = 2

    def __init__(self):
        a = 3

class B(A):
    a = 4

    def __init__(self):
        super().__init__()
        self.b = 7

obj = B()
obj.a = obj.c + 10
B.a = 13

try: print(obj.a, end= ';')
except: print('ERR', end=';')

try: print(A.a, end= ';')
except: print('ERR', end=';')

try: print(B.a, end= ';')
except: print('ERR', end=';')
code

10. 
Для графа на рисунку з вершини 4 запущено алгоритм пошуку в глибину, що будує кістякове дерево. Списки суміжності вершин переглядаються алгоритмом за зростанням міток вершин.\par
\begin{center}
	\adjustimage{max size={0.9\linewidth}{0.9\paperheight}}
	{../10/Traversals/72.png}
\end{center}
\par Які з наведених ребер входять до складу отриманого кістякового дерева?\par
1) (1, 4)\par
2) (1, 6)\par
3) (4, 5)\par
4) (1, 3)\par
5) (0, 4)\par
6) (3, 4)\par
{\vskip 15pt} У формі вибрати номери відповідних ребер.\par
%answer:2 4 5
\end {large}}
%end
{\smallskip \begin {small} Затверджено на засіданні кафедри теоретичної кібернетики, протокол \No 12 від   19.05.2020 р.\par\smallskip\smallskip
{\parbox {6.5 cm}{Зав. кафедри \dotfill Ю.В.~Крак}}\hspace {0.7cm}{\hbox to 6.5 cm{ Екзаменатор \dotfill Т.О.~Карнаух}} \end{small}}
%begin
{{\begin {small}\par
{ \center  Київський національний університет імені Тараса Шевченка \par}
{\parbox {5 cm}{Освітня програма \hfill}{\parbox {7 cm}{Прикладна математика, Системний аналіз \hfill}}\parbox {6 cm}{\hfill Семестр 2}}\par
{\parbox {5 cm} {Навчальний предмет \hfill}\parbox {7 cm}{Програмування \hfill}\parbox {6cm}{\hfill}\par}\vspace{-7pt} \end{small}}{\center Екзаменаційний білет \No \textbf 
{ 404 }\par}}
{
\begin {large}
1. Що буде виведено за виконання?
code
try:
    res = 6.0>7 and not 4>=2.0 or False<6
    if isinstance(res, bool):
        print(f'bool;{res}')
    elif isinstance(res, int):
        print(f'int;{res}')
    elif isinstance(res, float):
        print(f'float;{res:.2f}')
    else:
        print('???')
except:
    print('ERR')
code

2. Що буде виведено за виконання?
code
a = 5
b = 9
c = 0

def f(b):
    global c
    a = 4
    b = 5
    c = 5
    return a + b + c

a = 1
b = 6
c = 1

try:
    print(f(a), a, b, c, sep=';')
except:
    print('ERR', a, b, c, sep=';')
code

3. Що буде виведено за виконання?
code
try:
    t = 8
    while t<=7:
        t += 1
    print(t, end=';')
except:
    print('ERR')
code

4. Що буде виведено за виконання?
code
try:
    for d in range(12, 11, -2):
        if d>=8:
            continue
        print(d, end=';');d = 7
    else:
        print(d,end=';')
    print(d)
except:
    print('ERR')
code

5. Що буде виведено за виконання?\par (Дійсні значення, якщо наявні, в цьому завданні будуть виводитись у форматі з фіксованою крапкою та, можливо, з доволі великою кількістю десяткових розрядів. У відповіді для дійсних значень у ЦЬОМУ завданні достатньо вказати один десятковий розряд після крапки, відкинувши решту розрядів без округлення. Наприклад, замість 13.66666666666667 достатньо вказати 13.6, замість 25.23 достатньо вказати 25.2)
code
try:
    print(1, end=';')
    print(3//0, end=';')
    print(0, end=';')
except Exception:
    print(2, end=';')
except TypeError:
    print(5, end=';')
except:
    print('ERR', end=';')
else:
    print(8, end=';')
finally:
    print(9)
code

6. Що буде виведено за виконання?
code
def f(s1):
    s = s1[:]
    for i in range(len(s)):
        s[i] += 1
 
s = [1, 2, 3]
f(s)
print(s)
code

7. Що буде виведено за виконання?
code
class A:
    a = 1

    def f1(self): print(2, end=';')
    def _f2(self): print(3, end=';')

    def main(self):
        try: print(self.a, end=';')
        except: print('ERR', end=';')

        try: self.f1()
        except: print('ERR', end=';')

        try: self._f2()
        except: print('ERR', end=';')


class B(A):
    a = 4

    def f1(self): print(5, end=';')
    def _f2(self): print(6, end=';')

try: B().main()
except: print('ERR', end=';')
code

8. Що буде виведено за виконання?
code
def f(arg, n):
    n = 4
    tmp = arg
    tmp[-7:-2] = (14)
    return len(tmp)>3

try:
    c = ['0','1','2','3','4']
    n = 8
    f(c, n)
    print(n, end=';')
    for el in c:
        print(el, end=';')
except:
    print('ERR')
code

9. Що буде виведено за виконання?
code
class A:
    a = 1
    b = 2

    def __init__(self):
        b = 3

class B(A):
    b = 4

    def __init__(self):
        super().__init__()
        self.c = 7

obj = B()
obj.b = obj.c + 10
B.b = 13

try: print(obj.b, end= ';')
except: print('ERR', end=';')

try: print(A.b, end= ';')
except: print('ERR', end=';')

try: print(B.b, end= ';')
except: print('ERR', end=';')
code

10. 
Для графа на рисунку з вершини 6 запущено алгоритм пошуку в глибину, що будує кістякове дерево. Списки суміжності вершин переглядаються алгоритмом за зростанням міток вершин.\par
\begin{center}
	\adjustimage{max size={0.9\linewidth}{0.9\paperheight}}
	{../10/Traversals/260.png}
\end{center}
\par Які з наведених ребер входять до складу отриманого кістякового дерева?\par
1) (3, 5)\par
2) (1, 3)\par
3) (0, 6)\par
4) (0, 5)\par
5) (0, 1)\par
6) (2, 5)\par
{\vskip 15pt} У формі вибрати номери відповідних ребер.\par
%answer:1 2 3 5 6
\end {large}}
%end
{\smallskip \begin {small} Затверджено на засіданні кафедри теоретичної кібернетики, протокол \No 12 від   19.05.2020 р.\par\smallskip\smallskip
{\parbox {6.5 cm}{Зав. кафедри \dotfill Ю.В.~Крак}}\hspace {0.7cm}{\hbox to 6.5 cm{ Екзаменатор \dotfill Т.О.~Карнаух}} \end{small}}
%begin
{{\begin {small}\par
{ \center  Київський національний університет імені Тараса Шевченка \par}
{\parbox {5 cm}{Освітня програма \hfill}{\parbox {7 cm}{Прикладна математика, Системний аналіз \hfill}}\parbox {6 cm}{\hfill Семестр 2}}\par
{\parbox {5 cm} {Навчальний предмет \hfill}\parbox {7 cm}{Програмування \hfill}\parbox {6cm}{\hfill}\par}\vspace{-7pt} \end{small}}{\center Екзаменаційний білет \No \textbf 
{ 405 }\par}}
{
\begin {large}
1. Що буде виведено за виконання?
code
try:
    res = not 5.0!=9 or 8==2 and 3.0<=True
    if isinstance(res, bool):
        print(f'bool;{res}')
    elif isinstance(res, int):
        print(f'int;{res}')
    elif isinstance(res, float):
        print(f'float;{res:.2f}')
    else:
        print('???')
except:
    print('ERR')
code

2. Що буде виведено за виконання?
code
a = 1
b = 6
c = 8

def h(a):
    global c
    a = 2
    b += 3
    c = 4
    return a + b + c

a = 3
b = 4
c = 7

try:
    print(h(b), a, b, c, sep=';')
except:
    print('ERR', a, b, c, sep=';')
code

3. Що буде виведено за виконання?
code
try:
    n = 4
    while n<8:
        n += 2
    print(n, end=';')
except:
    print('ERR')
code

4. Що буде виведено за виконання?
code
try:
    for c in range(8, 4, -1):
        if c>5:
            continue
        print(c, end=';')
    else:
        print(c,end=';')
    print(c)
except:
    print('ERR')
code

5. Що буде виведено за виконання?\par (Дійсні значення, якщо наявні, в цьому завданні будуть виводитись у форматі з фіксованою крапкою та, можливо, з доволі великою кількістю десяткових розрядів. У відповіді для дійсних значень у ЦЬОМУ завданні достатньо вказати один десятковий розряд після крапки, відкинувши решту розрядів без округлення. Наприклад, замість 13.66666666666667 достатньо вказати 13.6, замість 25.23 достатньо вказати 25.2)
code
try:
    print(6, end=';')
    print(int('a7'), end=';')
    print(1, end=';')
except ZeroDivisionError:
    print(4, end=';')
except KeyboardInterrupt:
    print(9, end=';')
except:
    print('ERR', end=';')
else:
    print(8, end=';')
finally:
    print(5)
code

6. Що буде виведено за виконання?
code
def f(n):
    if n>9:
        f(n//100)
    else:
        print(n%10,end='')

f(987654)
code

7. Що буде виведено за виконання?
code
class A:
    c = 1

    def _h1(self): print(2, end=';')
    def h2(self): print(3, end=';')

    def main(self):
        try: print(self.c, end=';')
        except: print('ERR', end=';')

        try: self._h1()
        except: print('ERR', end=';')

        try: self.h2()
        except: print('ERR', end=';')


class B(A):
    c = 4

    def _h1(self): print(5, end=';')
    def h2(self): print(6, end=';')

try: B().main()
except: print('ERR', end=';')
code

8. Що буде виведено за виконання?
code
def f(arg, n):
    arg[65] = 1
    n += -2

try:
    n = 3
    d = {65:1, 26:3, 65:2}
    f(d, n)
    print(n, end=';')
    for x in d.values():
        print(x, sep=';', end=';')
except:
    print('ERR')
code

9. Що буде виведено за виконання?
code
class A:
    b = 1
    a = 2

    def __init__(self):
        a = 3

class B(A):
    a = 4

    def __init__(self):
        super().__init__()
        self.c = 7

obj = B()
obj.a = obj.b + 10
B.a = 13

try: print(obj.a, end= ';')
except: print('ERR', end=';')

try: print(A.a, end= ';')
except: print('ERR', end=';')

try: print(B.a, end= ';')
except: print('ERR', end=';')
code

10. 
Для графа на рисунку з вершини 1 запущено алгоритм пошуку в глибину, що будує кістякове дерево. Списки суміжності вершин переглядаються алгоритмом за зростанням міток вершин.\par
\begin{center}
	\adjustimage{max size={0.9\linewidth}{0.9\paperheight}}
	{../10/Traversals/217.png}
\end{center}
\par Які з наведених ребер входять до складу отриманого кістякового дерева?\par
1) (1, 3)\par
2) (2, 5)\par
3) (0, 3)\par
4) (0, 5)\par
5) (1, 2)\par
6) (3, 6)\par
{\vskip 15pt} У формі вибрати номери відповідних ребер.\par
%answer:2 3 4 5 6
\end {large}}
%end
{\smallskip \begin {small} Затверджено на засіданні кафедри теоретичної кібернетики, протокол \No 12 від   19.05.2020 р.\par\smallskip\smallskip
{\parbox {6.5 cm}{Зав. кафедри \dotfill Ю.В.~Крак}}\hspace {0.7cm}{\hbox to 6.5 cm{ Екзаменатор \dotfill Т.О.~Карнаух}} \end{small}}
%begin
{{\begin {small}\par
{ \center  Київський національний університет імені Тараса Шевченка \par}
{\parbox {5 cm}{Освітня програма \hfill}{\parbox {7 cm}{Прикладна математика, Системний аналіз \hfill}}\parbox {6 cm}{\hfill Семестр 2}}\par
{\parbox {5 cm} {Навчальний предмет \hfill}\parbox {7 cm}{Програмування \hfill}\parbox {6cm}{\hfill}\par}\vspace{-7pt} \end{small}}{\center Екзаменаційний білет \No \textbf 
{ 406 }\par}}
{
\begin {large}
1. Що буде виведено за виконання?
code
try:
    res = 8j<=9.0
    if isinstance(res, bool):
        print(f'bool;{res}')
    elif isinstance(res, int):
        print(f'int;{res}')
    elif isinstance(res, float):
        print(f'float;{res:.2f}')
    else:
        print('???')
except:
    print('ERR')
code

2. Що буде виведено за виконання?
code
a = 5
b = 4
c = 9

def g(b):
    global c
    a = 4
    b = 2
    c = 1
    return a + b + c

a = 2
b = 0
c = 7

try:
    print(g(b), a, b, c, sep=';')
except:
    print('ERR', a, b, c, sep=';')
code

3. Що буде виведено за виконання?
code
try:
    m = 2
    while m<10:
        m += 2
    print(m, end=';')
except:
    print('ERR')
code

4. Що буде виведено за виконання?
code
try:
    for f in range(-14, 9, 2):
        if f<=7:
            continue
        print(f, end=';');f = 1
        if f>6:
            break
    else:
        print('end',end=';')
    print(f)
except:
    print('ERR')
code

5. Що буде виведено за виконання?\par (Дійсні значення, якщо наявні, в цьому завданні будуть виводитись у форматі з фіксованою крапкою та, можливо, з доволі великою кількістю десяткових розрядів. У відповіді для дійсних значень у ЦЬОМУ завданні достатньо вказати один десятковий розряд після крапки, відкинувши решту розрядів без округлення. Наприклад, замість 13.66666666666667 достатньо вказати 13.6, замість 25.23 достатньо вказати 25.2)
code
try:
    print(6, end=';')
    print(int('d7'), end=';')
    print(3, end=';')
except Exception:
    print(0, end=';')
except ZeroDivisionError:
    print(2, end=';')
except:
    print('ERR', end=';')
else:
    print(6, end=';')
finally:
    print(9)
code

6. Що буде виведено за виконання?
code
def f(n):
    if n>9:
        return f(n//10)+1
    else:
        return 1

print(f(123456))
code

7. Що буде виведено за виконання?
code
class A:
    d = 1

    def _f1(self): print(2, end=';')
    def f2(self): print(3, end=';')

    def main(self):
        try: print(self.d, end=';')
        except: print('ERR', end=';')

        try: self._f1()
        except: print('ERR', end=';')

        try: self.f2()
        except: print('ERR', end=';')


class B(A):
    d = 4

    def _f1(self): print(5, end=';')
    def f2(self): print(6, end=';')

try: B().main()
except: print('ERR', end=';')
code

8. Що буде виведено за виконання?
code
def f(arg, n):
    arg[50] = 2
    n *= -3

try:
    n = 7
    s = {50:7, 67:0, 70:6, 50:1}
    f(s, n)
    print(n, end=';')
    for x in s.items():
        print(x, sep=';', end=';')
except:
    print('ERR')
code

9. Що буде виведено за виконання?
code
class A:
    c = 1
    a = 2

    def __init__(self):
        c = 3

class B(A):
    c = 4

    def __init__(self):
        super().__init__()
        self.b = 7

obj = B()
obj.c = obj.a + 10
B.c = 13

try: print(obj.c, end= ';')
except: print('ERR', end=';')

try: print(A.c, end= ';')
except: print('ERR', end=';')

try: print(B.c, end= ';')
except: print('ERR', end=';')
code

10. 
Для графа на рисунку з вершини 2 запущено алгоритм пошуку в глибину, що будує кістякове дерево. Списки суміжності вершин переглядаються алгоритмом за зростанням міток вершин.\par
\begin{center}
	\adjustimage{max size={0.9\linewidth}{0.9\paperheight}}
	{../10/Traversals/152.png}
\end{center}
\par Які з наведених ребер входять до складу отриманого кістякового дерева?\par
1) (2, 4)\par
2) (0, 2)\par
3) (3, 4)\par
4) (0, 1)\par
5) (0, 3)\par
6) (0, 6)\par
{\vskip 15pt} У формі вибрати номери відповідних ребер.\par
%answer:2 3 4 5 6
\end {large}}
%end
{\smallskip \begin {small} Затверджено на засіданні кафедри теоретичної кібернетики, протокол \No 12 від   19.05.2020 р.\par\smallskip\smallskip
{\parbox {6.5 cm}{Зав. кафедри \dotfill Ю.В.~Крак}}\hspace {0.7cm}{\hbox to 6.5 cm{ Екзаменатор \dotfill Т.О.~Карнаух}} \end{small}}
%begin
{{\begin {small}\par
{ \center  Київський національний університет імені Тараса Шевченка \par}
{\parbox {5 cm}{Освітня програма \hfill}{\parbox {7 cm}{Прикладна математика, Системний аналіз \hfill}}\parbox {6 cm}{\hfill Семестр 2}}\par
{\parbox {5 cm} {Навчальний предмет \hfill}\parbox {7 cm}{Програмування \hfill}\parbox {6cm}{\hfill}\par}\vspace{-7pt} \end{small}}{\center Екзаменаційний білет \No \textbf 
{ 407 }\par}}
{
\begin {large}
1. Що буде виведено за виконання?
code
try:
    res = 9/4//6*6
    if isinstance(res, bool):
        print(f'bool;{res}')
    elif isinstance(res, int):
        print(f'int;{res}')
    elif isinstance(res, float):
        print(f'float;{res:.2f}')
    else:
        print('???')
except:
    print('ERR')
code

2. Що буде виведено за виконання?
code
a = 3
b = 8
c = 2

def h(a):
    a = 3
    b = 5
    c = 3
    return a + b + c

a = 5
b = 9
c = 3

try:
    print(h(b), a, b, c, sep=';')
except:
    print('ERR', a, b, c, sep=';')
code

3. Що буде виведено за виконання?
code
try:
    j = 6
    while j<=11:
        j += 1
    print(j, end=';')
except:
    print('ERR')
code

4. Що буде виведено за виконання?
code
try:
    for b in range(4, 15, -2):
        if b<4:
            break
        print(b, end=';')
    else:
        print(b,end=';')
    print(b)
except:
    print('ERR')
code

5. Що буде виведено за виконання?\par (Дійсні значення, якщо наявні, в цьому завданні будуть виводитись у форматі з фіксованою крапкою та, можливо, з доволі великою кількістю десяткових розрядів. У відповіді для дійсних значень у ЦЬОМУ завданні достатньо вказати один десятковий розряд після крапки, відкинувши решту розрядів без округлення. Наприклад, замість 13.66666666666667 достатньо вказати 13.6, замість 25.23 достатньо вказати 25.2)
code
try:
    print(8, end=';')
    print(7/1, end=';')
    print(2, end=';')
except KeyboardInterrupt:
    print(1, end=';')
except ValueError:
    print(0, end=';')
except:
    print('ERR', end=';')
else:
    print(3, end=';')
finally:
    print(4)
code

6. Що буде виведено за виконання?
code
def f(n):
    if n<=9:
        print(n%10, end='')
    else:
        f(n//10)

f(987651)
code

7. Що буде виведено за виконання?
code
class A:
    b = 1

    def _h1(self): print(2, end=';')
    def _h2(self): print(3, end=';')

    def main(self):
        try: print(self.b, end=';')
        except: print('ERR', end=';')

        try: self._h1()
        except: print('ERR', end=';')

        try: self._h2()
        except: print('ERR', end=';')


class B(A):
    b = 4

    def _h1(self): print(5, end=';')
    def _h2(self): print(6, end=';')

try: B().main()
except: print('ERR', end=';')
code

8. Що буде виведено за виконання?
code
def f(arg, n):
    arg[10] = 1
    n = 8
    arg = tuple(arg.items())

try:
    n = 1
    s = {49:7, 48:1, 44:1, 44:7}
    f(s, n)
    print(n, end=';')
    for x in s.values():
        print(x, sep=';', end=';')
except:
    print('ERR')
code

9. Що буде виведено за виконання?
code
class A:
    c = 1
    b = 2

    def __init__(self):
        c = 3

class B(A):
    c = 4

    def __init__(self):
        super().__init__()
        self.a = 7

obj = B()
obj.b = obj.b + 10
B.b = 13

try: print(obj.b, end= ';')
except: print('ERR', end=';')

try: print(A.b, end= ';')
except: print('ERR', end=';')

try: print(B.b, end= ';')
except: print('ERR', end=';')
code

10. 
Для графа на рисунку з вершини 6 запущено алгоритм пошуку в глибину, що будує кістякове дерево. Списки суміжності вершин переглядаються алгоритмом за зростанням міток вершин.\par
\begin{center}
	\adjustimage{max size={0.9\linewidth}{0.9\paperheight}}
	{../10/Traversals/264.png}
\end{center}
\par Які з наведених ребер входять до складу отриманого кістякового дерева?\par
1) (0, 1)\par
2) (4, 6)\par
3) (2, 4)\par
4) (0, 2)\par
5) (2, 3)\par
6) (1, 6)\par
{\vskip 15pt} У формі вибрати номери відповідних ребер.\par
%answer:1 3 4 5 6
\end {large}}
%end
{\smallskip \begin {small} Затверджено на засіданні кафедри теоретичної кібернетики, протокол \No 12 від   19.05.2020 р.\par\smallskip\smallskip
{\parbox {6.5 cm}{Зав. кафедри \dotfill Ю.В.~Крак}}\hspace {0.7cm}{\hbox to 6.5 cm{ Екзаменатор \dotfill Т.О.~Карнаух}} \end{small}}
%begin
{{\begin {small}\par
{ \center  Київський національний університет імені Тараса Шевченка \par}
{\parbox {5 cm}{Освітня програма \hfill}{\parbox {7 cm}{Прикладна математика, Системний аналіз \hfill}}\parbox {6 cm}{\hfill Семестр 2}}\par
{\parbox {5 cm} {Навчальний предмет \hfill}\parbox {7 cm}{Програмування \hfill}\parbox {6cm}{\hfill}\par}\vspace{-7pt} \end{small}}{\center Екзаменаційний білет \No \textbf 
{ 408 }\par}}
{
\begin {large}
1. Що буде виведено за виконання?
code
def f(n):
    print('f', end='')
    return n

try:
    print(f(5) and f(7))
except:
    print('ERR')
code

2. Що буде виведено за виконання?
code
a = 0
b = 2
c = 8

def g(b):
    a = 5
    b *= 3
    c = 1
    return a + b + c

a = 0
b = 4
c = 7

try:
    print(g(b), a, b, c, sep=';')
except:
    print('ERR', a, b, c, sep=';')
code

3. Що буде виведено за виконання?
code
try:
    i = 9
    while i<13:
        i += 1
    print(i, end=';')
except:
    print('ERR')
code

4. Що буде виведено за виконання?
code
try:
    for e in range(-4, 12, -1):
        if e>=3:
            continue
        print(e, end=';')
    else:
        print(e,end=';')
    print(e)
except:
    print('ERR')
code

5. Що буде виведено за виконання?\par (Дійсні значення, якщо наявні, в цьому завданні будуть виводитись у форматі з фіксованою крапкою та, можливо, з доволі великою кількістю десяткових розрядів. У відповіді для дійсних значень у ЦЬОМУ завданні достатньо вказати один десятковий розряд після крапки, відкинувши решту розрядів без округлення. Наприклад, замість 13.66666666666667 достатньо вказати 13.6, замість 25.23 достатньо вказати 25.2)
code
try:
    print(5, end=';')
    print(4//0.0, end=';')
    print(3, end=';')
except TypeError:
    print(2, end=';')
except ValueError:
    print(6, end=';')
except:
    print('ERR', end=';')
else:
    print(1, end=';')
finally:
    print(9)
code

6. Що буде виведено за виконання?
code
def f(n):
    if n>2:
        f(n//2)
    print(n%2, end='')
 
f(16)
code

7. Що буде виведено за виконання?
code
class A:
    a = 1

    def _g1(self): print(2, end=';')
    def __g2(self): print(3, end=';')

    def main(self):
        try: print(self.a, end=';')
        except: print('ERR', end=';')

        try: self._g1()
        except: print('ERR', end=';')

        try: self.__g2()
        except: print('ERR', end=';')


class B(A):
    a = 4

    def _g1(self): print(5, end=';')
    def __g2(self): print(6, end=';')

try: B().main()
except: print('ERR', end=';')
code

8. Що буде виведено за виконання?
code
def f(arg, n):
    n = 6
    tmp = arg
    tmp[-4:] = [2,3]
    return len(tmp)>3

try:
    d = ['0','1','2','3','4']
    n = 9
    f(d, n)
    print(n, end=';')
    for el in d:
        print(el, end=';')
except:
    print('ERR')
code

9. Що буде виведено за виконання?
code
class A:
    b = 1
    c = 2

    def __init__(self):
        c = 3

class B(A):
    c = 4

    def __init__(self):
        super().__init__()
        self.a = 7

obj = B()
obj.a = obj.c + 10
B.a = 13

try: print(obj.a, end= ';')
except: print('ERR', end=';')

try: print(A.a, end= ';')
except: print('ERR', end=';')

try: print(B.a, end= ';')
except: print('ERR', end=';')
code

10. 
Для графа на рисунку з вершини 0 запущено алгоритм пошуку в глибину, що будує кістякове дерево. Списки суміжності вершин переглядаються алгоритмом за зростанням міток вершин.\par
\begin{center}
	\adjustimage{max size={0.9\linewidth}{0.9\paperheight}}
	{../10/Traversals/272.png}
\end{center}
\par Які з наведених ребер входять до складу отриманого кістякового дерева?\par
1) (3, 5)\par
2) (0, 2)\par
3) (3, 6)\par
4) (1, 3)\par
5) (4, 5)\par
6) (0, 1)\par
{\vskip 15pt} У формі вибрати номери відповідних ребер.\par
%answer:1 4 5 6
\end {large}}
%end
{\smallskip \begin {small} Затверджено на засіданні кафедри теоретичної кібернетики, протокол \No 12 від   19.05.2020 р.\par\smallskip\smallskip
{\parbox {6.5 cm}{Зав. кафедри \dotfill Ю.В.~Крак}}\hspace {0.7cm}{\hbox to 6.5 cm{ Екзаменатор \dotfill Т.О.~Карнаух}} \end{small}}
%begin
{{\begin {small}\par
{ \center  Київський національний університет імені Тараса Шевченка \par}
{\parbox {5 cm}{Освітня програма \hfill}{\parbox {7 cm}{Прикладна математика, Системний аналіз \hfill}}\parbox {6 cm}{\hfill Семестр 2}}\par
{\parbox {5 cm} {Навчальний предмет \hfill}\parbox {7 cm}{Програмування \hfill}\parbox {6cm}{\hfill}\par}\vspace{-7pt} \end{small}}{\center Екзаменаційний білет \No \textbf 
{ 409 }\par}}
{
\begin {large}
1. Що буде виведено за виконання?
code
try:
    res = not 5<=7.0 and 8.0>3 and 7>=False
    if isinstance(res, bool):
        print(f'bool;{res}')
    elif isinstance(res, int):
        print(f'int;{res}')
    elif isinstance(res, float):
        print(f'float;{res:.2f}')
    else:
        print('???')
except:
    print('ERR')
code

2. Що буде виведено за виконання?
code
a = 2
b = 3
c = 1

def h(a):
    a += 5
    b = 4
    c = 2
    return a + b + c

a = 6
b = 9
c = 5

try:
    print(h(b), a, b, c, sep=';')
except:
    print('ERR', a, b, c, sep=';')
code

3. Що буде виведено за виконання?
code
try:
    k = 4
    while k<=5:
        k += 2
    print(k, end=';')
except:
    print('ERR')
code

4. Що буде виведено за виконання?
code
try:
    for f in range(-8, 13, -3):
        if f<8:
            continue
        print(f, end=';')
    else:
        print(f,end=';')
    print(f)
except:
    print('ERR')
code

5. Що буде виведено за виконання?\par (Дійсні значення, якщо наявні, в цьому завданні будуть виводитись у форматі з фіксованою крапкою та, можливо, з доволі великою кількістю десяткових розрядів. У відповіді для дійсних значень у ЦЬОМУ завданні достатньо вказати один десятковий розряд після крапки, відкинувши решту розрядів без округлення. Наприклад, замість 13.66666666666667 достатньо вказати 13.6, замість 25.23 достатньо вказати 25.2)
code
try:
    print(7, end=';')
    print(int('8'), end=';')
    print(5, end=';')
except Exception:
    print(0, end=';')
except ZeroDivisionError:
    print(0, end=';')
except:
    print('ERR', end=';')
else:
    print(5, end=';')
finally:
    print(8)
code

6. Що буде виведено за виконання?
code
def f(n):
    if n>9:
        f(n//100)
    else:
        print(n%10,end='')

f(123456)
code

7. Що буде виведено за виконання?
code
class A:
    c = 1

    def __g1(self): print(2, end=';')
    def _g2(self): print(3, end=';')

    def main(self):
        try: print(self.c, end=';')
        except: print('ERR', end=';')

        try: self.__g1()
        except: print('ERR', end=';')

        try: self._g2()
        except: print('ERR', end=';')


class B(A):
    c = 4

    def __g1(self): print(5, end=';')
    def _g2(self): print(6, end=';')

try: B().main()
except: print('ERR', end=';')
code

8. Що буде виведено за виконання?
code
def f(arg, n):
    arg[75] = 1
    n = 5
    arg = set(arg.values())

try:
    n = 0
    t = {12:4, 77:4, 75:2, 29:0}
    f(t, n)
    print(n, end=';')
    for x in t.values():
        print(x, sep=';', end=';')
except:
    print('ERR')
code

9. Що буде виведено за виконання?
code
class A:
    b = 1
    c = 2

    def __init__(self):
        c = 3

class B(A):
    c = 4

    def __init__(self):
        super().__init__()
        self.a = 7

obj = B()
obj.c = obj.a + 10
B.c = 13

try: print(obj.c, end= ';')
except: print('ERR', end=';')

try: print(A.c, end= ';')
except: print('ERR', end=';')

try: print(B.c, end= ';')
except: print('ERR', end=';')
code

10. 
Для графа на рисунку з вершини 5 запущено алгоритм пошуку в глибину, що будує кістякове дерево. Списки суміжності вершин переглядаються алгоритмом за зростанням міток вершин.\par
\begin{center}
	\adjustimage{max size={0.9\linewidth}{0.9\paperheight}}
	{../10/Traversals/235.png}
\end{center}
\par Які з наведених ребер входять до складу отриманого кістякового дерева?\par
1) (3, 4)\par
2) (1, 2)\par
3) (0, 1)\par
4) (0, 5)\par
5) (1, 4)\par
6) (1, 6)\par
{\vskip 15pt} У формі вибрати номери відповідних ребер.\par
%answer:1 2 3 4
\end {large}}
%end
{\smallskip \begin {small} Затверджено на засіданні кафедри теоретичної кібернетики, протокол \No 12 від   19.05.2020 р.\par\smallskip\smallskip
{\parbox {6.5 cm}{Зав. кафедри \dotfill Ю.В.~Крак}}\hspace {0.7cm}{\hbox to 6.5 cm{ Екзаменатор \dotfill Т.О.~Карнаух}} \end{small}}
%begin
{{\begin {small}\par
{ \center  Київський національний університет імені Тараса Шевченка \par}
{\parbox {5 cm}{Освітня програма \hfill}{\parbox {7 cm}{Прикладна математика, Системний аналіз \hfill}}\parbox {6 cm}{\hfill Семестр 2}}\par
{\parbox {5 cm} {Навчальний предмет \hfill}\parbox {7 cm}{Програмування \hfill}\parbox {6cm}{\hfill}\par}\vspace{-7pt} \end{small}}{\center Екзаменаційний білет \No \textbf 
{ 410 }\par}}
{
\begin {large}
1. Що буде виведено за виконання?
code
try:
    res = "False">"7j"
    if isinstance(res, bool):
        print(f'bool;{res}')
    elif isinstance(res, int):
        print(f'int;{res}')
    elif isinstance(res, float):
        print(f'float;{res:.2f}')
    else:
        print('???')
except:
    print('ERR')
code

2. Що буде виведено за виконання?
code
a = 8
b = 4
c = 5

def f(b):
    global c
    a -= 1
    b = 5
    c = 4
    return a + b + c

a = 7
b = 2
c = 6

try:
    print(f(b), a, b, c, sep=';')
except:
    print('ERR', a, b, c, sep=';')
code

3. Що буде виведено за виконання?
code
try:
    k = 10
    while k>=2:
        k += -2
    print(k, end=';')
except:
    print('ERR')
code

4. Що буде виведено за виконання?
code
try:
    for b in range(-5, -11, 3):
        if b>=6:
            continue
        print(b, end=';');b = -9
    else:
        print(b,end=';')
    print(b)
except:
    print('ERR')
code

5. Що буде виведено за виконання?\par (Дійсні значення, якщо наявні, в цьому завданні будуть виводитись у форматі з фіксованою крапкою та, можливо, з доволі великою кількістю десяткових розрядів. У відповіді для дійсних значень у ЦЬОМУ завданні достатньо вказати один десятковий розряд після крапки, відкинувши решту розрядів без округлення. Наприклад, замість 13.66666666666667 достатньо вказати 13.6, замість 25.23 достатньо вказати 25.2)
code
try:
    print(1, end=';')
    print(2j<=8j, end=';')
    print(3, end=';')
except TypeError:
    print(4, end=';')
except KeyboardInterrupt:
    print(9, end=';')
except:
    print('ERR', end=';')
else:
    print(6, end=';')
finally:
    print(7)
code

6. Що буде виведено за виконання?
code
def f(n):
    print(n%2, end='')
    if n>2:
        f(n//2)
 
f(16)
code

7. Що буде виведено за виконання?
code
class A:
    d = 1

    def h1(self): print(2, end=';')
    def _h2(self): print(3, end=';')

    def main(self):
        try: print(self.d, end=';')
        except: print('ERR', end=';')

        try: self.h1()
        except: print('ERR', end=';')

        try: self._h2()
        except: print('ERR', end=';')


class B(A):
    d = 4

    def h1(self): print(5, end=';')
    def _h2(self): print(6, end=';')

try: B().main()
except: print('ERR', end=';')
code

8. Що буде виведено за виконання?
code
def f(arg, n):
    n = 6
    tmp = arg
    tmp[:-6] = 'ab'
    return len(tmp)>3

try:
    b = [10,11,12,13,14]
    n = 9
    f(b, n)
    print(n, end=';')
    for el in b:
        print(el, end=';')
except:
    print('ERR')
code

9. Що буде виведено за виконання?
code
class A:
    a = 1
    b = 2

    def __init__(self):
        a = 3

class B(A):
    a = 4

    def __init__(self):
        super().__init__()
        self.c = 7

obj = B()
obj.b = obj.a + 10
B.b = 13

try: print(obj.b, end= ';')
except: print('ERR', end=';')

try: print(A.b, end= ';')
except: print('ERR', end=';')

try: print(B.b, end= ';')
except: print('ERR', end=';')
code

10. 
Для графа на рисунку з вершини 5 запущено алгоритм пошуку в глибину, що будує кістякове дерево. Списки суміжності вершин переглядаються алгоритмом за зростанням міток вершин.\par
\begin{center}
	\adjustimage{max size={0.9\linewidth}{0.9\paperheight}}
	{../10/Traversals/277.png}
\end{center}
\par Які з наведених ребер входять до складу отриманого кістякового дерева?\par
1) (5, 6)\par
2) (3, 6)\par
3) (3, 5)\par
4) (2, 3)\par
5) (1, 3)\par
6) (4, 6)\par
{\vskip 15pt} У формі вибрати номери відповідних ребер.\par
%answer:3 5 6
\end {large}}
%end
{\smallskip \begin {small} Затверджено на засіданні кафедри теоретичної кібернетики, протокол \No 12 від   19.05.2020 р.\par\smallskip\smallskip
{\parbox {6.5 cm}{Зав. кафедри \dotfill Ю.В.~Крак}}\hspace {0.7cm}{\hbox to 6.5 cm{ Екзаменатор \dotfill Т.О.~Карнаух}} \end{small}}
%begin
{{\begin {small}\par
{ \center  Київський національний університет імені Тараса Шевченка \par}
{\parbox {5 cm}{Освітня програма \hfill}{\parbox {7 cm}{Прикладна математика, Системний аналіз \hfill}}\parbox {6 cm}{\hfill Семестр 2}}\par
{\parbox {5 cm} {Навчальний предмет \hfill}\parbox {7 cm}{Програмування \hfill}\parbox {6cm}{\hfill}\par}\vspace{-7pt} \end{small}}{\center Екзаменаційний білет \No \textbf 
{ 411 }\par}}
{
\begin {large}
1. Що буде виведено за виконання?
code
try:
    res = 3/9%5*5
    if isinstance(res, bool):
        print(f'bool;{res}')
    elif isinstance(res, int):
        print(f'int;{res}')
    elif isinstance(res, float):
        print(f'float;{res:.2f}')
    else:
        print('???')
except:
    print('ERR')
code

2. Що буде виведено за виконання?
code
a = 0
b = 1
c = 6

def f(a):
    a = 2
    b = 4
    c = 1
    return a + b + c

a = 8
b = 7
c = 4

try:
    print(f(b), a, b, c, sep=';')
except:
    print('ERR', a, b, c, sep=';')
code

3. Що буде виведено за виконання?
code
try:
    i = 12
    while i>6:
        i += -2
    print(i, end=';')
except:
    print('ERR')
code

4. Що буде виведено за виконання?
code
try:
    for e in range(-2, 8, 2):
        if e<=4:
            break
        print(e, end=';')
    else:
        print(e,end=';')
    print(e)
except:
    print('ERR')
code

5. Що буде виведено за виконання?\par (Дійсні значення, якщо наявні, в цьому завданні будуть виводитись у форматі з фіксованою крапкою та, можливо, з доволі великою кількістю десяткових розрядів. У відповіді для дійсних значень у ЦЬОМУ завданні достатньо вказати один десятковий розряд після крапки, відкинувши решту розрядів без округлення. Наприклад, замість 13.66666666666667 достатньо вказати 13.6, замість 25.23 достатньо вказати 25.2)
code
try:
    print(6, end=';')
    print(int('2'), end=';')
    print(3, end=';')
except KeyboardInterrupt:
    print(7, end=';')
except ValueError:
    print(8, end=';')
except:
    print('ERR', end=';')
else:
    print(4, end=';')
finally:
    print(0)
code

6. Що буде виведено за виконання?
code
def f(n):
    if n>9:
        f(n//100)
        print(n%10, end='')
 
f(123456)
code

7. Що буде виведено за виконання?
code
class A:
    a = 1

    def _f1(self): print(2, end=';')
    def __f2(self): print(3, end=';')

    def main(self):
        try: print(self.a, end=';')
        except: print('ERR', end=';')

        try: self._f1()
        except: print('ERR', end=';')

        try: self.__f2()
        except: print('ERR', end=';')


class B(A):
    a = 4

    def _f1(self): print(5, end=';')
    def __f2(self): print(6, end=';')

try: B().main()
except: print('ERR', end=';')
code

8. Що буде виведено за виконання?
code
def f(arg, n):
    arg[30] = 2
    n -= -2

try:
    n = 5
    t = {79:4, 30:9, 87:7, 13:4}
    f(t, n)
    print(n, end=';')
    for x in t.keys():
        print(x, sep=';', end=';')
except:
    print('ERR')
code

9. Що буде виведено за виконання?
code
class A:
    c = 1
    b = 2

    def __init__(self):
        c = 3

class B(A):
    b = 4

    def __init__(self):
        super().__init__()
        self.a = 7

obj = B()
obj.b = obj.c + 10
B.b = 13

try: print(obj.b, end= ';')
except: print('ERR', end=';')

try: print(A.b, end= ';')
except: print('ERR', end=';')

try: print(B.b, end= ';')
except: print('ERR', end=';')
code

10. 
Для графа на рисунку з вершини 1 запущено алгоритм пошуку в глибину, що будує кістякове дерево. Списки суміжності вершин переглядаються алгоритмом за зростанням міток вершин.\par
\begin{center}
	\adjustimage{max size={0.9\linewidth}{0.9\paperheight}}
	{../10/Traversals/233.png}
\end{center}
\par Які з наведених ребер входять до складу отриманого кістякового дерева?\par
1) (0, 5)\par
2) (1, 4)\par
3) (0, 3)\par
4) (2, 4)\par
5) (1, 6)\par
6) (1, 3)\par
{\vskip 15pt} У формі вибрати номери відповідних ребер.\par
%answer:1 3 4
\end {large}}
%end
{\smallskip \begin {small} Затверджено на засіданні кафедри теоретичної кібернетики, протокол \No 12 від   19.05.2020 р.\par\smallskip\smallskip
{\parbox {6.5 cm}{Зав. кафедри \dotfill Ю.В.~Крак}}\hspace {0.7cm}{\hbox to 6.5 cm{ Екзаменатор \dotfill Т.О.~Карнаух}} \end{small}}
%begin
{{\begin {small}\par
{ \center  Київський національний університет імені Тараса Шевченка \par}
{\parbox {5 cm}{Освітня програма \hfill}{\parbox {7 cm}{Прикладна математика, Системний аналіз \hfill}}\parbox {6 cm}{\hfill Семестр 2}}\par
{\parbox {5 cm} {Навчальний предмет \hfill}\parbox {7 cm}{Програмування \hfill}\parbox {6cm}{\hfill}\par}\vspace{-7pt} \end{small}}{\center Екзаменаційний білет \No \textbf 
{ 412 }\par}}
{
\begin {large}
1. Що буде виведено за виконання?
code
try:
    res = 1.0<9
    if isinstance(res, bool):
        print(f'bool;{res}')
    elif isinstance(res, int):
        print(f'int;{res}')
    elif isinstance(res, float):
        print(f'float;{res:.2f}')
    else:
        print('???')
except:
    print('ERR')
code

2. Що буде виведено за виконання?
code
a = 9
b = 3
c = 7

def g(b):
    global c
    a = 3
    b = 4
    c = 1
    return a + b + c

a = 0
b = 5
c = 4

try:
    print(g(a), a, b, c, sep=';')
except:
    print('ERR', a, b, c, sep=';')
code

3. Що буде виведено за виконання?
code
try:
    m = 5
    while m>=1:
        m += -3
    print(m, end=';')
except:
    print('ERR')
code

4. Що буде виведено за виконання?
code
try:
    for c in range(15, 1, 3):
        if c>5:
            break
        print(c, end=';')
        if c<3:
            break
    else:
        print('end',end=';')
    print(c)
except:
    print('ERR')
code

5. Що буде виведено за виконання?\par (Дійсні значення, якщо наявні, в цьому завданні будуть виводитись у форматі з фіксованою крапкою та, можливо, з доволі великою кількістю десяткових розрядів. У відповіді для дійсних значень у ЦЬОМУ завданні достатньо вказати один десятковий розряд після крапки, відкинувши решту розрядів без округлення. Наприклад, замість 13.66666666666667 достатньо вказати 13.6, замість 25.23 достатньо вказати 25.2)
code
try:
    print(5, end=';')
    print(1%0, end=';')
    print(9, end=';')
except Exception:
    print(3, end=';')
except TypeError:
    print(1, end=';')
except:
    print('ERR', end=';')
else:
    print(4, end=';')
finally:
    print(7)
code

6. Що буде виведено за виконання?
code
def f(n):
    if n>2:
        f(n//2)
    else:
        print(n%2, end='')

f(16)
code

7. Що буде виведено за виконання?
code
class A:
    b = 1

    def h1(self): print(2, end=';')
    def _h2(self): print(3, end=';')

    def main(self):
        try: print(self.b, end=';')
        except: print('ERR', end=';')

        try: self.h1()
        except: print('ERR', end=';')

        try: self._h2()
        except: print('ERR', end=';')


class B(A):
    b = 4

    def h1(self): print(5, end=';')
    def _h2(self): print(6, end=';')

try: B().main()
except: print('ERR', end=';')
code

8. Що буде виведено за виконання?
code
def f(arg, n):
    arg[52] = 5
    n = 9
    arg = list(arg.keys())

try:
    n = 2
    d = {52:8, 24:7, 52:1}
    f(d, n)
    print(n, end=';')
    for x in d.keys():
        print(x, sep=';', end=';')
except:
    print('ERR')
code

9. Що буде виведено за виконання?
code
class A:
    c = 1
    a = 2

    def __init__(self):
        a = 3

class B(A):
    c = 4

    def __init__(self):
        super().__init__()
        self.b = 7

obj = B()
obj.a = obj.b + 10
B.a = 13

try: print(obj.a, end= ';')
except: print('ERR', end=';')

try: print(A.a, end= ';')
except: print('ERR', end=';')

try: print(B.a, end= ';')
except: print('ERR', end=';')
code

10. 
Для графа на рисунку з вершини 4 запущено алгоритм пошуку в глибину, що будує кістякове дерево. Списки суміжності вершин переглядаються алгоритмом за зростанням міток вершин.\par
\begin{center}
	\adjustimage{max size={0.9\linewidth}{0.9\paperheight}}
	{../10/Traversals/299.png}
\end{center}
\par Які з наведених ребер входять до складу отриманого кістякового дерева?\par
1) (2, 5)\par
2) (4, 6)\par
3) (1, 4)\par
4) (1, 2)\par
5) (2, 4)\par
6) (4, 5)\par
{\vskip 15pt} У формі вибрати номери відповідних ребер.\par
%answer:1 4
\end {large}}
%end
{\smallskip \begin {small} Затверджено на засіданні кафедри теоретичної кібернетики, протокол \No 12 від   19.05.2020 р.\par\smallskip\smallskip
{\parbox {6.5 cm}{Зав. кафедри \dotfill Ю.В.~Крак}}\hspace {0.7cm}{\hbox to 6.5 cm{ Екзаменатор \dotfill Т.О.~Карнаух}} \end{small}}
%begin
{{\begin {small}\par
{ \center  Київський національний університет імені Тараса Шевченка \par}
{\parbox {5 cm}{Освітня програма \hfill}{\parbox {7 cm}{Прикладна математика, Системний аналіз \hfill}}\parbox {6 cm}{\hfill Семестр 2}}\par
{\parbox {5 cm} {Навчальний предмет \hfill}\parbox {7 cm}{Програмування \hfill}\parbox {6cm}{\hfill}\par}\vspace{-7pt} \end{small}}{\center Екзаменаційний білет \No \textbf 
{ 413 }\par}}
{
\begin {large}
1. Що буде виведено за виконання?
code
try:
    res = 9**-1.0*1
    if isinstance(res, bool):
        print(f'bool;{res}')
    elif isinstance(res, int):
        print(f'int;{res}')
    elif isinstance(res, float):
        print(f'float;{res:.2f}')
    else:
        print('???')
except:
    print('ERR')
code

2. Що буде виведено за виконання?
code
a = 3
b = 9
c = 0

def g(a):
    a *= 2
    b = 3
    c = 3
    return a + b + c

a = 1
b = 7
c = 6

try:
    print(g(a), a, b, c, sep=';')
except:
    print('ERR', a, b, c, sep=';')
code

3. Що буде виведено за виконання?
code
try:
    m = 0
    while m<=9:
        m += 1
    print(m, end=';')
except:
    print('ERR')
code

4. Що буде виведено за виконання?
code
try:
    for d in range(-13, 14, -1):
        if d<=3:
            continue
        print(d, end=';');d = 10
    else:
        print(d,end=';')
    print(d)
except:
    print('ERR')
code

5. Що буде виведено за виконання?\par (Дійсні значення, якщо наявні, в цьому завданні будуть виводитись у форматі з фіксованою крапкою та, можливо, з доволі великою кількістю десяткових розрядів. У відповіді для дійсних значень у ЦЬОМУ завданні достатньо вказати один десятковий розряд після крапки, відкинувши решту розрядів без округлення. Наприклад, замість 13.66666666666667 достатньо вказати 13.6, замість 25.23 достатньо вказати 25.2)
code
try:
    print(5, end=';')
    print(6/2, end=';')
    print(0, end=';')
except ZeroDivisionError:
    print(8, end=';')
except KeyboardInterrupt:
    print(9, end=';')
except:
    print('ERR', end=';')
else:
    print(2, end=';')
finally:
    print(8)
code

6. Що буде виведено за виконання?
code
def f(s):
    for el in s:
        el += 1
    return s
 
s = [1, 2, 3]
s = f(s)
print(s)
code

7. Що буде виведено за виконання?
code
class A:
    d = 1

    def __g1(self): print(2, end=';')
    def _g2(self): print(3, end=';')

    def main(self):
        try: print(self.d, end=';')
        except: print('ERR', end=';')

        try: self.__g1()
        except: print('ERR', end=';')

        try: self._g2()
        except: print('ERR', end=';')


class B(A):
    d = 4

    def __g1(self): print(5, end=';')
    def _g2(self): print(6, end=';')

try: B().main()
except: print('ERR', end=';')
code

8. Що буде виведено за виконання?
code
def f(arg, n):
    n = 2
    tmp = arg
    del tmp[6:-5]
    return len(tmp)>3

try:
    g = ['0','1','2','3','4']
    n = 8
    f(g, n)
    print(n, end=';')
    for el in g:
        print(el, end=';')
except:
    print('ERR')
code

9. Що буде виведено за виконання?
code
class A:
    a = 1
    c = 2

    def __init__(self):
        a = 3

class B(A):
    c = 4

    def __init__(self):
        super().__init__()
        self.b = 7

obj = B()
obj.c = obj.b + 10
B.c = 13

try: print(obj.c, end= ';')
except: print('ERR', end=';')

try: print(A.c, end= ';')
except: print('ERR', end=';')

try: print(B.c, end= ';')
except: print('ERR', end=';')
code

10. 
Для графа на рисунку з вершини 1 запущено алгоритм пошуку в глибину, що будує кістякове дерево. Списки суміжності вершин переглядаються алгоритмом за зростанням міток вершин.\par
\begin{center}
	\adjustimage{max size={0.9\linewidth}{0.9\paperheight}}
	{../10/Traversals/14.png}
\end{center}
\par Які з наведених ребер входять до складу отриманого кістякового дерева?\par
1) (0, 6)\par
2) (4, 5)\par
3) (3, 4)\par
4) (1, 2)\par
5) (2, 3)\par
6) (3, 5)\par
{\vskip 15pt} У формі вибрати номери відповідних ребер.\par
%answer:1 3 5 6
\end {large}}
%end
{\smallskip \begin {small} Затверджено на засіданні кафедри теоретичної кібернетики, протокол \No 12 від   19.05.2020 р.\par\smallskip\smallskip
{\parbox {6.5 cm}{Зав. кафедри \dotfill Ю.В.~Крак}}\hspace {0.7cm}{\hbox to 6.5 cm{ Екзаменатор \dotfill Т.О.~Карнаух}} \end{small}}
%begin
{{\begin {small}\par
{ \center  Київський національний університет імені Тараса Шевченка \par}
{\parbox {5 cm}{Освітня програма \hfill}{\parbox {7 cm}{Прикладна математика, Системний аналіз \hfill}}\parbox {6 cm}{\hfill Семестр 2}}\par
{\parbox {5 cm} {Навчальний предмет \hfill}\parbox {7 cm}{Програмування \hfill}\parbox {6cm}{\hfill}\par}\vspace{-7pt} \end{small}}{\center Екзаменаційний білет \No \textbf 
{ 414 }\par}}
{
\begin {large}
1. Що буде виведено за виконання?
code
try:
    res = True<4 or not 5==6 or 9.0!=4.0
    if isinstance(res, bool):
        print(f'bool;{res}')
    elif isinstance(res, int):
        print(f'int;{res}')
    elif isinstance(res, float):
        print(f'float;{res:.2f}')
    else:
        print('???')
except:
    print('ERR')
code

2. Що буде виведено за виконання?
code
a = 4
b = 9
c = 5

def f(b):
    global c
    a *= 5
    b = 4
    c = 2
    return a + b + c

a = 2
b = 0
c = 3

try:
    print(f(b), a, b, c, sep=';')
except:
    print('ERR', a, b, c, sep=';')
code

3. Що буде виведено за виконання?
code
try:
    k = 5
    while k<6:
        k += 2
    print(k, end=';')
except:
    print('ERR')
code

4. Що буде виведено за виконання?
code
try:
    for a in range(-12, -3, -3):
        if a>=7:
            continue
        print(a, end=';')
    else:
        print(a,end=';')
    print(a)
except:
    print('ERR')
code

5. Що буде виведено за виконання?\par (Дійсні значення, якщо наявні, в цьому завданні будуть виводитись у форматі з фіксованою крапкою та, можливо, з доволі великою кількістю десяткових розрядів. У відповіді для дійсних значень у ЦЬОМУ завданні достатньо вказати один десятковий розряд після крапки, відкинувши решту розрядів без округлення. Наприклад, замість 13.66666666666667 достатньо вказати 13.6, замість 25.23 достатньо вказати 25.2)
code
try:
    print(6, end=';')
    print(4j!=9j, end=';')
    print(1, end=';')
except TypeError:
    print(9, end=';')
except Exception:
    print(7, end=';')
except:
    print('ERR', end=';')
else:
    print(3, end=';')
finally:
    print(5)
code

6. Що буде виведено за виконання?
code
def f(n):
    if n>9:
        f(n//10)
    print(n%10, end='')
 
f(123456)
code

7. Що буде виведено за виконання?
code
class A:
    b = 1

    def _f1(self): print(2, end=';')
    def _f2(self): print(3, end=';')

    def main(self):
        try: print(self.b, end=';')
        except: print('ERR', end=';')

        try: self._f1()
        except: print('ERR', end=';')

        try: self._f2()
        except: print('ERR', end=';')


class B(A):
    b = 4

    def _f1(self): print(5, end=';')
    def _f2(self): print(6, end=';')

try: B().main()
except: print('ERR', end=';')
code

8. Що буде виведено за виконання?
code
def f(arg, n):
    arg[66] = 1
    n -= 2

try:
    n = 6
    d = {29:0, 66:2, 29:1}
    f(d, n)
    print(n, end=';')
    for x in d.keys():
        print(x, sep=';', end=';')
except:
    print('ERR')
code

9. Що буде виведено за виконання?
code
class A:
    b = 1
    a = 2

    def __init__(self):
        a = 3

class B(A):
    b = 4

    def __init__(self):
        super().__init__()
        self.c = 7

obj = B()
obj.c = obj.a + 10
B.c = 13

try: print(obj.c, end= ';')
except: print('ERR', end=';')

try: print(A.c, end= ';')
except: print('ERR', end=';')

try: print(B.c, end= ';')
except: print('ERR', end=';')
code

10. 
Для графа на рисунку з вершини 6 запущено алгоритм пошуку в глибину, що будує кістякове дерево. Списки суміжності вершин переглядаються алгоритмом за зростанням міток вершин.\par
\begin{center}
	\adjustimage{max size={0.9\linewidth}{0.9\paperheight}}
	{../10/Traversals/280.png}
\end{center}
\par Які з наведених ребер входять до складу отриманого кістякового дерева?\par
1) (0, 6)\par
2) (2, 5)\par
3) (1, 4)\par
4) (0, 1)\par
5) (1, 3)\par
6) (0, 3)\par
{\vskip 15pt} У формі вибрати номери відповідних ребер.\par
%answer:1 2 3 4 5
\end {large}}
%end
{\smallskip \begin {small} Затверджено на засіданні кафедри теоретичної кібернетики, протокол \No 12 від   19.05.2020 р.\par\smallskip\smallskip
{\parbox {6.5 cm}{Зав. кафедри \dotfill Ю.В.~Крак}}\hspace {0.7cm}{\hbox to 6.5 cm{ Екзаменатор \dotfill Т.О.~Карнаух}} \end{small}}
%begin
{{\begin {small}\par
{ \center  Київський національний університет імені Тараса Шевченка \par}
{\parbox {5 cm}{Освітня програма \hfill}{\parbox {7 cm}{Прикладна математика, Системний аналіз \hfill}}\parbox {6 cm}{\hfill Семестр 2}}\par
{\parbox {5 cm} {Навчальний предмет \hfill}\parbox {7 cm}{Програмування \hfill}\parbox {6cm}{\hfill}\par}\vspace{-7pt} \end{small}}{\center Екзаменаційний білет \No \textbf 
{ 415 }\par}}
{
\begin {large}
1. Що буде виведено за виконання?
code
try:
    res = 8//7*4*9
    if isinstance(res, bool):
        print(f'bool;{res}')
    elif isinstance(res, int):
        print(f'int;{res}')
    elif isinstance(res, float):
        print(f'float;{res:.2f}')
    else:
        print('???')
except:
    print('ERR')
code

2. Що буде виведено за виконання?
code
a = 2
b = 6
c = 8

def h(b):
    a = 2
    b = 5
    c = 3
    return a + b + c

a = 1
b = 2
c = 9

try:
    print(h(a), a, b, c, sep=';')
except:
    print('ERR', a, b, c, sep=';')
code

3. Що буде виведено за виконання?
code
try:
    j = 1
    while j<=12:
        j += 2
    print(j, end=';')
except:
    print('ERR')
code

4. Що буде виведено за виконання?
code
try:
    for c in range(0, 0, -2):
        if c<6:
            break
        print(c, end=';');c = -8
    else:
        print(c,end=';')
    print(c)
except:
    print('ERR')
code

5. Що буде виведено за виконання?\par (Дійсні значення, якщо наявні, в цьому завданні будуть виводитись у форматі з фіксованою крапкою та, можливо, з доволі великою кількістю десяткових розрядів. У відповіді для дійсних значень у ЦЬОМУ завданні достатньо вказати один десятковий розряд після крапки, відкинувши решту розрядів без округлення. Наприклад, замість 13.66666666666667 достатньо вказати 13.6, замість 25.23 достатньо вказати 25.2)
code
try:
    print(2, end=';')
    print(int('c0'), end=';')
    print(6, end=';')
except ValueError:
    print(3, end=';')
except ZeroDivisionError:
    print(9, end=';')
except:
    print('ERR', end=';')
else:
    print(5, end=';')
finally:
    print(1)
code

6. Що буде виведено за виконання?
code
def f(n):
    print(n%10, end='')
    if n>9:
        f(n//10)
 
f(123456)
code

7. Що буде виведено за виконання?
code
class A:
    a = 1

    def _h1(self): print(2, end=';')
    def h2(self): print(3, end=';')

    def main(self):
        try: print(self.a, end=';')
        except: print('ERR', end=';')

        try: self._h1()
        except: print('ERR', end=';')

        try: self.h2()
        except: print('ERR', end=';')


class B(A):
    a = 4

    def _h1(self): print(5, end=';')
    def h2(self): print(6, end=';')

try: B().main()
except: print('ERR', end=';')
code

8. Що буде виведено за виконання?
code
def f(arg, n):
    n = 5
    tmp = arg
    tmp[1:7] = (12)
    return len(tmp)>3

try:
    c = [10,11,12,13,14]
    n = 3
    f(c, n)
    print(n, end=';')
    for el in c:
        print(el, end=';')
except:
    print('ERR')
code

9. Що буде виведено за виконання?
code
class A:
    c = 1
    a = 2

    def __init__(self):
        a = 3

class B(A):
    a = 4

    def __init__(self):
        super().__init__()
        self.b = 7

obj = B()
obj.a = obj.c + 10
B.a = 13

try: print(obj.a, end= ';')
except: print('ERR', end=';')

try: print(A.a, end= ';')
except: print('ERR', end=';')

try: print(B.a, end= ';')
except: print('ERR', end=';')
code

10. 
Для графа на рисунку з вершини 3 запущено алгоритм пошуку в глибину, що будує кістякове дерево. Списки суміжності вершин переглядаються алгоритмом за зростанням міток вершин.\par
\begin{center}
	\adjustimage{max size={0.9\linewidth}{0.9\paperheight}}
	{../10/Traversals/93.png}
\end{center}
\par Які з наведених ребер входять до складу отриманого кістякового дерева?\par
1) (1, 6)\par
2) (4, 5)\par
3) (0, 2)\par
4) (0, 5)\par
5) (1, 5)\par
6) (1, 3)\par
{\vskip 15pt} У формі вибрати номери відповідних ребер.\par
%answer:2 3 4 5 6
\end {large}}
%end
{\smallskip \begin {small} Затверджено на засіданні кафедри теоретичної кібернетики, протокол \No 12 від   19.05.2020 р.\par\smallskip\smallskip
{\parbox {6.5 cm}{Зав. кафедри \dotfill Ю.В.~Крак}}\hspace {0.7cm}{\hbox to 6.5 cm{ Екзаменатор \dotfill Т.О.~Карнаух}} \end{small}}
%begin
{{\begin {small}\par
{ \center  Київський національний університет імені Тараса Шевченка \par}
{\parbox {5 cm}{Освітня програма \hfill}{\parbox {7 cm}{Прикладна математика, Системний аналіз \hfill}}\parbox {6 cm}{\hfill Семестр 2}}\par
{\parbox {5 cm} {Навчальний предмет \hfill}\parbox {7 cm}{Програмування \hfill}\parbox {6cm}{\hfill}\par}\vspace{-7pt} \end{small}}{\center Екзаменаційний білет \No \textbf 
{ 416 }\par}}
{
\begin {large}
1. Що буде виведено за виконання?
code
try:
    res = 2**4+-1
    if isinstance(res, bool):
        print(f'bool;{res}')
    elif isinstance(res, int):
        print(f'int;{res}')
    elif isinstance(res, float):
        print(f'float;{res:.2f}')
    else:
        print('???')
except:
    print('ERR')
code

2. Що буде виведено за виконання?
code
a = 7
b = 3
c = 5

def g(a):
    global c
    a = 1
    b = 2
    c = 5
    return a + b + c

a = 0
b = 1
c = 8

try:
    print(g(b), a, b, c, sep=';')
except:
    print('ERR', a, b, c, sep=';')
code

3. Що буде виведено за виконання?
code
try:
    t = 14
    while t>=8:
        t += -2
    print(t, end=';')
except:
    print('ERR')
code

4. Що буде виведено за виконання?
code
try:
    for d in range(13, -14, 3):
        if d>3:
            continue
        print(d, end=';')
    else:
        print(d,end=';')
    print(d)
except:
    print('ERR')
code

5. Що буде виведено за виконання?\par (Дійсні значення, якщо наявні, в цьому завданні будуть виводитись у форматі з фіксованою крапкою та, можливо, з доволі великою кількістю десяткових розрядів. У відповіді для дійсних значень у ЦЬОМУ завданні достатньо вказати один десятковий розряд після крапки, відкинувши решту розрядів без округлення. Наприклад, замість 13.66666666666667 достатньо вказати 13.6, замість 25.23 достатньо вказати 25.2)
code
try:
    print(4, end=';')
    print(8//3, end=';')
    print(7, end=';')
except TypeError:
    print(0, end=';')
except Exception:
    print(2, end=';')
except:
    print('ERR', end=';')
else:
    print(4, end=';')
finally:
    print(0)
code

6. Що буде виведено за виконання?
code
def f(n):
    print(n%10, end='')
    if n>9:
        f(n//10)
 
f(543216)
code

7. Що буде виведено за виконання?
code
class A:
    c = 1

    def _f1(self): print(2, end=';')
    def __f2(self): print(3, end=';')

    def main(self):
        try: print(self.c, end=';')
        except: print('ERR', end=';')

        try: self._f1()
        except: print('ERR', end=';')

        try: self.__f2()
        except: print('ERR', end=';')


class B(A):
    c = 4

    def _f1(self): print(5, end=';')
    def __f2(self): print(6, end=';')

try: B().main()
except: print('ERR', end=';')
code

8. Що буде виведено за виконання?
code
def f(arg, n):
    n = 4
    tmp = arg
    tmp[3:5] = 'xy'
    return len(tmp)>3

try:
    e = [10,11,12,13,14]
    n = 7
    f(e, n)
    print(n, end=';')
    for el in e:
        print(el, end=';')
except:
    print('ERR')
code

9. Що буде виведено за виконання?
code
class A:
    b = 1
    c = 2

    def __init__(self):
        b = 3

class B(A):
    b = 4

    def __init__(self):
        super().__init__()
        self.a = 7

obj = B()
obj.b = obj.a + 10
B.b = 13

try: print(obj.b, end= ';')
except: print('ERR', end=';')

try: print(A.b, end= ';')
except: print('ERR', end=';')

try: print(B.b, end= ';')
except: print('ERR', end=';')
code

10. 
Для графа на рисунку з вершини 5 запущено алгоритм пошуку в глибину, що будує кістякове дерево. Списки суміжності вершин переглядаються алгоритмом за зростанням міток вершин.\par
\begin{center}
	\adjustimage{max size={0.9\linewidth}{0.9\paperheight}}
	{../10/Traversals/170.png}
\end{center}
\par Які з наведених ребер входять до складу отриманого кістякового дерева?\par
1) (1, 5)\par
2) (0, 3)\par
3) (0, 4)\par
4) (4, 6)\par
5) (1, 2)\par
6) (1, 4)\par
{\vskip 15pt} У формі вибрати номери відповідних ребер.\par
%answer:1 2 3 4 5
\end {large}}
%end
{\smallskip \begin {small} Затверджено на засіданні кафедри теоретичної кібернетики, протокол \No 12 від   19.05.2020 р.\par\smallskip\smallskip
{\parbox {6.5 cm}{Зав. кафедри \dotfill Ю.В.~Крак}}\hspace {0.7cm}{\hbox to 6.5 cm{ Екзаменатор \dotfill Т.О.~Карнаух}} \end{small}}
%begin
{{\begin {small}\par
{ \center  Київський національний університет імені Тараса Шевченка \par}
{\parbox {5 cm}{Освітня програма \hfill}{\parbox {7 cm}{Прикладна математика, Системний аналіз \hfill}}\parbox {6 cm}{\hfill Семестр 2}}\par
{\parbox {5 cm} {Навчальний предмет \hfill}\parbox {7 cm}{Програмування \hfill}\parbox {6cm}{\hfill}\par}\vspace{-7pt} \end{small}}{\center Екзаменаційний білет \No \textbf 
{ 417 }\par}}
{
\begin {large}
1. Що буде виведено за виконання?
code
try:
    res = True>=2.0 or not 8.0>3 or 9<8
    if isinstance(res, bool):
        print(f'bool;{res}')
    elif isinstance(res, int):
        print(f'int;{res}')
    elif isinstance(res, float):
        print(f'float;{res:.2f}')
    else:
        print('???')
except:
    print('ERR')
code

2. Що буде виведено за виконання?
code
a = 6
b = 4
c = 6

def f(a):
    global c
    a = 4
    b = 1
    c = 2
    return a + b + c

a = 7
b = 0
c = 5

try:
    print(f(b), a, b, c, sep=';')
except:
    print('ERR', a, b, c, sep=';')
code

3. Що буде виведено за виконання?
code
try:
    n = 13
    while n>5:
        n += -3
    print(n, end=';')
except:
    print('ERR')
code

4. Що буде виведено за виконання?
code
try:
    for a in range(9, 7, -3):
        if a>4:
            continue
        print(a, end=';')
        if a<4:
            break
    else:
        print('end',end=';')
    print(a)
except:
    print('ERR')
code

5. Що буде виведено за виконання?\par (Дійсні значення, якщо наявні, в цьому завданні будуть виводитись у форматі з фіксованою крапкою та, можливо, з доволі великою кількістю десяткових розрядів. У відповіді для дійсних значень у ЦЬОМУ завданні достатньо вказати один десятковий розряд після крапки, відкинувши решту розрядів без округлення. Наприклад, замість 13.66666666666667 достатньо вказати 13.6, замість 25.23 достатньо вказати 25.2)
code
try:
    print(6, end=';')
    print(1j>=7j, end=';')
    print(5, end=';')
except ZeroDivisionError:
    print(3, end=';')
except ValueError:
    print(2, end=';')
except:
    print('ERR', end=';')
else:
    print(7, end=';')
finally:
    print(8)
code

6. Що буде виведено за виконання?
code
def f(n):
    print(n%10,end='')
    if n>9:
        f(n//100)
 
f(123456)
code

7. Що буде виведено за виконання?
code
class A:
    c = 1

    def _g1(self): print(2, end=';')
    def _g2(self): print(3, end=';')

    def main(self):
        try: print(self.c, end=';')
        except: print('ERR', end=';')

        try: self._g1()
        except: print('ERR', end=';')

        try: self._g2()
        except: print('ERR', end=';')


class B(A):
    c = 4

    def _g1(self): print(5, end=';')
    def _g2(self): print(6, end=';')

try: B().main()
except: print('ERR', end=';')
code

8. Що буде виведено за виконання?
code
def f(arg, n):
    n = 1
    tmp = arg
    del tmp[-9:1]
    return len(tmp)>3

try:
    b = ['0','1','2','3','4']
    n = 0
    f(b, n)
    print(n, end=';')
    for el in b:
        print(el, end=';')
except:
    print('ERR')
code

9. Що буде виведено за виконання?
code
class A:
    a = 1
    c = 2

    def __init__(self):
        a = 3

class B(A):
    c = 4

    def __init__(self):
        super().__init__()
        self.b = 7

obj = B()
obj.b = obj.c + 10
B.b = 13

try: print(obj.b, end= ';')
except: print('ERR', end=';')

try: print(A.b, end= ';')
except: print('ERR', end=';')

try: print(B.b, end= ';')
except: print('ERR', end=';')
code

10. 
Для графа на рисунку з вершини 5 запущено алгоритм пошуку в глибину, що будує кістякове дерево. Списки суміжності вершин переглядаються алгоритмом за зростанням міток вершин.\par
\begin{center}
	\adjustimage{max size={0.9\linewidth}{0.9\paperheight}}
	{../10/Traversals/36.png}
\end{center}
\par Які з наведених ребер входять до складу отриманого кістякового дерева?\par
1) (0, 2)\par
2) (2, 5)\par
3) (1, 4)\par
4) (0, 3)\par
5) (2, 6)\par
6) (1, 5)\par
{\vskip 15pt} У формі вибрати номери відповідних ребер.\par
%answer:1 5
\end {large}}
%end
{\smallskip \begin {small} Затверджено на засіданні кафедри теоретичної кібернетики, протокол \No 12 від   19.05.2020 р.\par\smallskip\smallskip
{\parbox {6.5 cm}{Зав. кафедри \dotfill Ю.В.~Крак}}\hspace {0.7cm}{\hbox to 6.5 cm{ Екзаменатор \dotfill Т.О.~Карнаух}} \end{small}}
%begin
{{\begin {small}\par
{ \center  Київський національний університет імені Тараса Шевченка \par}
{\parbox {5 cm}{Освітня програма \hfill}{\parbox {7 cm}{Прикладна математика, Системний аналіз \hfill}}\parbox {6 cm}{\hfill Семестр 2}}\par
{\parbox {5 cm} {Навчальний предмет \hfill}\parbox {7 cm}{Програмування \hfill}\parbox {6cm}{\hfill}\par}\vspace{-7pt} \end{small}}{\center Екзаменаційний білет \No \textbf 
{ 418 }\par}}
{
\begin {large}
1. Що буде виведено за виконання?
code
try:
    res = 2<=9.0 and not 6==5 and 4.0!=False
    if isinstance(res, bool):
        print(f'bool;{res}')
    elif isinstance(res, int):
        print(f'int;{res}')
    elif isinstance(res, float):
        print(f'float;{res:.2f}')
    else:
        print('???')
except:
    print('ERR')
code

2. Що буде виведено за виконання?
code
a = 3
b = 8
c = 9

def h(b):
    a = 5
    b = 3
    c = 4
    return a + b + c

a = 2
b = 1
c = 4

try:
    print(h(a), a, b, c, sep=';')
except:
    print('ERR', a, b, c, sep=';')
code

3. Що буде виведено за виконання?
code
try:
    t = 8
    while t<11:
        t += 1
    print(t, end=';')
except:
    print('ERR')
code

4. Що буде виведено за виконання?
code
try:
    for e in range(-10, 2, 2):
        if e<8:
            break
        print(e, end=';')
    else:
        print(e,end=';')
    print(e)
except:
    print('ERR')
code

5. Що буде виведено за виконання?\par (Дійсні значення, якщо наявні, в цьому завданні будуть виводитись у форматі з фіксованою крапкою та, можливо, з доволі великою кількістю десяткових розрядів. У відповіді для дійсних значень у ЦЬОМУ завданні достатньо вказати один десятковий розряд після крапки, відкинувши решту розрядів без округлення. Наприклад, замість 13.66666666666667 достатньо вказати 13.6, замість 25.23 достатньо вказати 25.2)
code
try:
    print(9, end=';')
    print(int('b3'), end=';')
    print(2, end=';')
except KeyboardInterrupt:
    print(8, end=';')
except TypeError:
    print(4, end=';')
except:
    print('ERR', end=';')
else:
    print(9, end=';')
finally:
    print(0)
code

6. Що буде виведено за виконання?
code
def f(n):
    if n<=9:
        print(n%10, end='')
    else:
        f(n//10)

f(123456)
code

7. Що буде виведено за виконання?
code
class A:
    b = 1

    def __g1(self): print(2, end=';')
    def _g2(self): print(3, end=';')

    def main(self):
        try: print(self.b, end=';')
        except: print('ERR', end=';')

        try: self.__g1()
        except: print('ERR', end=';')

        try: self._g2()
        except: print('ERR', end=';')


class B(A):
    b = 4

    def __g1(self): print(5, end=';')
    def _g2(self): print(6, end=';')

try: B().main()
except: print('ERR', end=';')
code

8. Що буде виведено за виконання?
code
def f(arg, n):
    arg[76] = 6
    n *= 3

try:
    n = 4
    s = {73:4, 38:8, 73:4}
    f(s, n)
    print(n, end=';')
    for x, y in s.items():
        print(x, y, sep=';', end=';')
except:
    print('ERR')
code

9. Що буде виведено за виконання?
code
class A:
    a = 1
    b = 2

    def __init__(self):
        b = 3

class B(A):
    a = 4

    def __init__(self):
        super().__init__()
        self.c = 7

obj = B()
obj.a = obj.b + 10
B.a = 13

try: print(obj.a, end= ';')
except: print('ERR', end=';')

try: print(A.a, end= ';')
except: print('ERR', end=';')

try: print(B.a, end= ';')
except: print('ERR', end=';')
code

10. 
Для графа на рисунку з вершини 0 запущено алгоритм пошуку в глибину, що будує кістякове дерево. Списки суміжності вершин переглядаються алгоритмом за зростанням міток вершин.\par
\begin{center}
	\adjustimage{max size={0.9\linewidth}{0.9\paperheight}}
	{../10/Traversals/265.png}
\end{center}
\par Які з наведених ребер входять до складу отриманого кістякового дерева?\par
1) (2, 6)\par
2) (1, 4)\par
3) (0, 6)\par
4) (1, 2)\par
5) (4, 6)\par
6) (4, 5)\par
{\vskip 15pt} У формі вибрати номери відповідних ребер.\par
%answer:4 5 6
\end {large}}
%end
{\smallskip \begin {small} Затверджено на засіданні кафедри теоретичної кібернетики, протокол \No 12 від   19.05.2020 р.\par\smallskip\smallskip
{\parbox {6.5 cm}{Зав. кафедри \dotfill Ю.В.~Крак}}\hspace {0.7cm}{\hbox to 6.5 cm{ Екзаменатор \dotfill Т.О.~Карнаух}} \end{small}}
%begin
{{\begin {small}\par
{ \center  Київський національний університет імені Тараса Шевченка \par}
{\parbox {5 cm}{Освітня програма \hfill}{\parbox {7 cm}{Прикладна математика, Системний аналіз \hfill}}\parbox {6 cm}{\hfill Семестр 2}}\par
{\parbox {5 cm} {Навчальний предмет \hfill}\parbox {7 cm}{Програмування \hfill}\parbox {6cm}{\hfill}\par}\vspace{-7pt} \end{small}}{\center Екзаменаційний білет \No \textbf 
{ 419 }\par}}
{
\begin {large}
1. Що буде виведено за виконання?
code
def f(n):
    print('f', end='')
    return n>-1

try:
    print(f(-8) or f(-3))
except:
    print('ERR')
code

2. Що буде виведено за виконання?
code
a = 8
b = 1
c = 7

def h(b):
    a = 1
    b += 2
    c = 1
    return a + b + c

a = 6
b = 1
c = 8

try:
    print(h(a), a, b, c, sep=';')
except:
    print('ERR', a, b, c, sep=';')
code

3. Що буде виведено за виконання?
code
try:
    n = 3
    while n<=6:
        n += 2
    print(n, end=';')
except:
    print('ERR')
code

4. Що буде виведено за виконання?
code
try:
    for b in range(14, -4, -2):
        if b<=7:
            continue
        print(b, end=';');b = -7
    else:
        print(b,end=';')
    print(b)
except:
    print('ERR')
code

5. Що буде виведено за виконання?\par (Дійсні значення, якщо наявні, в цьому завданні будуть виводитись у форматі з фіксованою крапкою та, можливо, з доволі великою кількістю десяткових розрядів. У відповіді для дійсних значень у ЦЬОМУ завданні достатньо вказати один десятковий розряд після крапки, відкинувши решту розрядів без округлення. Наприклад, замість 13.66666666666667 достатньо вказати 13.6, замість 25.23 достатньо вказати 25.2)
code
try:
    print(6, end=';')
    print(7%0.0, end=';')
    print(5, end=';')
except Exception:
    print(1, end=';')
except ValueError:
    print(6, end=';')
except:
    print('ERR', end=';')
else:
    print(5, end=';')
finally:
    print(0)
code

6. Що буде виведено за виконання?
code
def f(n):
    if n>9:
        f(n//100)
    print(n%10, end='')
 
f(123456)
code

7. Що буде виведено за виконання?
code
class A:
    a = 1

    def h1(self): print(2, end=';')
    def _h2(self): print(3, end=';')

    def main(self):
        try: print(self.a, end=';')
        except: print('ERR', end=';')

        try: self.h1()
        except: print('ERR', end=';')

        try: self._h2()
        except: print('ERR', end=';')


class B(A):
    a = 4

    def h1(self): print(5, end=';')
    def _h2(self): print(6, end=';')

try: B().main()
except: print('ERR', end=';')
code

8. Що буде виведено за виконання?
code
def f(arg, n):
    arg[29] = 8
    n = 7
    arg = set(arg.keys())

try:
    n = 3
    d = {63:9, 56:3, 63:3}
    f(d, n)
    print(n, end=';')
    for x in d.items():
        print(x, sep=';', end=';')
except:
    print('ERR')
code

9. Що буде виведено за виконання?
code
class A:
    c = 1
    b = 2

    def __init__(self):
        c = 3

class B(A):
    b = 4

    def __init__(self):
        super().__init__()
        self.a = 7

obj = B()
obj.c = obj.a + 10
B.c = 13

try: print(obj.c, end= ';')
except: print('ERR', end=';')

try: print(A.c, end= ';')
except: print('ERR', end=';')

try: print(B.c, end= ';')
except: print('ERR', end=';')
code

10. 
Для графа на рисунку з вершини 1 запущено алгоритм пошуку в глибину, що будує кістякове дерево. Списки суміжності вершин переглядаються алгоритмом за зростанням міток вершин.\par
\begin{center}
	\adjustimage{max size={0.9\linewidth}{0.9\paperheight}}
	{../10/Traversals/193.png}
\end{center}
\par Які з наведених ребер входять до складу отриманого кістякового дерева?\par
1) (0, 5)\par
2) (0, 2)\par
3) (3, 5)\par
4) (0, 1)\par
5) (1, 2)\par
6) (2, 4)\par
{\vskip 15pt} У формі вибрати номери відповідних ребер.\par
%answer:2 3 4 6
\end {large}}
%end
{\smallskip \begin {small} Затверджено на засіданні кафедри теоретичної кібернетики, протокол \No 12 від   19.05.2020 р.\par\smallskip\smallskip
{\parbox {6.5 cm}{Зав. кафедри \dotfill Ю.В.~Крак}}\hspace {0.7cm}{\hbox to 6.5 cm{ Екзаменатор \dotfill Т.О.~Карнаух}} \end{small}}
%begin
{{\begin {small}\par
{ \center  Київський національний університет імені Тараса Шевченка \par}
{\parbox {5 cm}{Освітня програма \hfill}{\parbox {7 cm}{Прикладна математика, Системний аналіз \hfill}}\parbox {6 cm}{\hfill Семестр 2}}\par
{\parbox {5 cm} {Навчальний предмет \hfill}\parbox {7 cm}{Програмування \hfill}\parbox {6cm}{\hfill}\par}\vspace{-7pt} \end{small}}{\center Екзаменаційний білет \No \textbf 
{ 420 }\par}}
{
\begin {large}
1. Що буде виведено за виконання?
code
try:
    res = -3**-2-False
    if isinstance(res, bool):
        print(f'bool;{res}')
    elif isinstance(res, int):
        print(f'int;{res}')
    elif isinstance(res, float):
        print(f'float;{res:.2f}')
    else:
        print('???')
except:
    print('ERR')
code

2. Що буде виведено за виконання?
code
a = 0
b = 5
c = 3

def g(a):
    global c
    a -= 4
    b = 5
    c = 3
    return a + b + c

a = 4
b = 2
c = 9

try:
    print(g(a), a, b, c, sep=';')
except:
    print('ERR', a, b, c, sep=';')
code

3. Що буде виведено за виконання?
code
try:
    i = 7
    while i<8:
        i += 1
    print(i, end=';')
except:
    print('ERR')
code

4. Що буде виведено за виконання?
code
try:
    for f in range(-6, -5, -1):
        if f>=5:
            continue
        print(f, end=';')
    else:
        print(f,end=';')
    print(f)
except:
    print('ERR')
code

5. Що буде виведено за виконання?\par (Дійсні значення, якщо наявні, в цьому завданні будуть виводитись у форматі з фіксованою крапкою та, можливо, з доволі великою кількістю десяткових розрядів. У відповіді для дійсних значень у ЦЬОМУ завданні достатньо вказати один десятковий розряд після крапки, відкинувши решту розрядів без округлення. Наприклад, замість 13.66666666666667 достатньо вказати 13.6, замість 25.23 достатньо вказати 25.2)
code
try:
    print(3, end=';')
    print(int('7'), end=';')
    print(1, end=';')
except ZeroDivisionError:
    print(9, end=';')
except KeyboardInterrupt:
    print(4, end=';')
except:
    print('ERR', end=';')
else:
    print(2, end=';')
finally:
    print(8)
code

6. Що буде виведено за виконання?
code
def f(s):
    for i in range(len(s)):
        s[i] += 1
        return
 
s = [1, 2, 3]
f(s)
print(s)
code

7. Що буде виведено за виконання?
code
class A:
    d = 1

    def _f1(self): print(2, end=';')
    def f2(self): print(3, end=';')

    def main(self):
        try: print(self.d, end=';')
        except: print('ERR', end=';')

        try: self._f1()
        except: print('ERR', end=';')

        try: self.f2()
        except: print('ERR', end=';')


class B(A):
    d = 4

    def _f1(self): print(5, end=';')
    def f2(self): print(6, end=';')

try: B().main()
except: print('ERR', end=';')
code

8. Що буде виведено за виконання?
code
def f(arg, n):
    arg[29] = 8
    n = 6
    arg = tuple(arg.values())

try:
    n = 4
    t = {76:0, 29:1, 85:3, 76:9}
    f(t, n)
    print(n, end=';')
    for x, y in t.items():
        print(x, y, sep=';', end=';')
except:
    print('ERR')
code

9. Що буде виведено за виконання?
code
class A:
    b = 1
    a = 2

    def __init__(self):
        a = 3

class B(A):
    b = 4

    def __init__(self):
        super().__init__()
        self.c = 7

obj = B()
obj.b = obj.c + 10
B.b = 13

try: print(obj.b, end= ';')
except: print('ERR', end=';')

try: print(A.b, end= ';')
except: print('ERR', end=';')

try: print(B.b, end= ';')
except: print('ERR', end=';')
code

10. 
Для графа на рисунку з вершини 4 запущено алгоритм пошуку в глибину, що будує кістякове дерево. Списки суміжності вершин переглядаються алгоритмом за зростанням міток вершин.\par
\begin{center}
	\adjustimage{max size={0.9\linewidth}{0.9\paperheight}}
	{../10/Traversals/197.png}
\end{center}
\par Які з наведених ребер входять до складу отриманого кістякового дерева?\par
1) (2, 4)\par
2) (5, 6)\par
3) (4, 5)\par
4) (4, 6)\par
5) (2, 5)\par
6) (0, 4)\par
{\vskip 15pt} У формі вибрати номери відповідних ребер.\par
%answer:2 5 6
\end {large}}
%end
{\smallskip \begin {small} Затверджено на засіданні кафедри теоретичної кібернетики, протокол \No 12 від   19.05.2020 р.\par\smallskip\smallskip
{\parbox {6.5 cm}{Зав. кафедри \dotfill Ю.В.~Крак}}\hspace {0.7cm}{\hbox to 6.5 cm{ Екзаменатор \dotfill Т.О.~Карнаух}} \end{small}}
%begin
{{\begin {small}\par
{ \center  Київський національний університет імені Тараса Шевченка \par}
{\parbox {5 cm}{Освітня програма \hfill}{\parbox {7 cm}{Прикладна математика, Системний аналіз \hfill}}\parbox {6 cm}{\hfill Семестр 2}}\par
{\parbox {5 cm} {Навчальний предмет \hfill}\parbox {7 cm}{Програмування \hfill}\parbox {6cm}{\hfill}\par}\vspace{-7pt} \end{small}}{\center Екзаменаційний білет \No \textbf 
{ 421 }\par}}
{
\begin {large}
1. Що буде виведено за виконання?
code
try:
    res = True>="2.0j"
    if isinstance(res, bool):
        print(f'bool;{res}')
    elif isinstance(res, int):
        print(f'int;{res}')
    elif isinstance(res, float):
        print(f'float;{res:.2f}')
    else:
        print('???')
except:
    print('ERR')
code

2. Що буде виведено за виконання?
code
a = 5
b = 8
c = 9

def f(a):
    global c
    a += 1
    b = 4
    c = 3
    return a + b + c

a = 2
b = 6
c = 1

try:
    print(f(b), a, b, c, sep=';')
except:
    print('ERR', a, b, c, sep=';')
code

3. Що буде виведено за виконання?
code
try:
    n = 5
    while n<=10:
        n += 1
    print(n, end=';')
except:
    print('ERR')
code

4. Що буде виведено за виконання?
code
try:
    for b in range(10, -2, -3):
        if b>=7:
            break
        print(b, end=';')
    else:
        print('end',end=';')
    print(b)
except:
    print('ERR')
code

5. Що буде виведено за виконання?\par (Дійсні значення, якщо наявні, в цьому завданні будуть виводитись у форматі з фіксованою крапкою та, можливо, з доволі великою кількістю десяткових розрядів. У відповіді для дійсних значень у ЦЬОМУ завданні достатньо вказати один десятковий розряд після крапки, відкинувши решту розрядів без округлення. Наприклад, замість 13.66666666666667 достатньо вказати 13.6, замість 25.23 достатньо вказати 25.2)
code
try:
    print(5, end=';')
    print(int('4'), end=';')
    print(9, end=';')
except Exception:
    print(0, end=';')
except TypeError:
    print(1, end=';')
except:
    print('ERR', end=';')
else:
    print(3, end=';')
finally:
    print(6)
code

6. Що буде виведено за виконання?
code
def f(n):
    if n>9: f(n//10)
    else:
        print(n%10, end='')
 
f(123456)
code

7. Що буде виведено за виконання?
code
class A:
    d = 1

    def __f1(self): print(2, end=';')
    def _f2(self): print(3, end=';')

    def main(self):
        try: print(self.d, end=';')
        except: print('ERR', end=';')

        try: self.__f1()
        except: print('ERR', end=';')

        try: self._f2()
        except: print('ERR', end=';')


class B(A):
    d = 4

    def __f1(self): print(5, end=';')
    def _f2(self): print(6, end=';')

try: B().main()
except: print('ERR', end=';')
code

8. Що буде виведено за виконання?
code
def f(arg, n):
    arg[17] = 8
    n = 5
    arg = list(arg.items())

try:
    n = 1
    s = {58:7, 24:9, 58:4}
    f(s, n)
    print(n, end=';')
    for x in s.values():
        print(x, sep=';', end=';')
except:
    print('ERR')
code

9. Що буде виведено за виконання?
code
class A:
    c = 1
    a = 2

    def __init__(self):
        c = 3

class B(A):
    a = 4

    def __init__(self):
        super().__init__()
        self.b = 7

obj = B()
obj.a = obj.c + 10
B.a = 13

try: print(obj.a, end= ';')
except: print('ERR', end=';')

try: print(A.a, end= ';')
except: print('ERR', end=';')

try: print(B.a, end= ';')
except: print('ERR', end=';')
code

10. 
Для графа на рисунку з вершини 5 запущено алгоритм пошуку в глибину, що будує кістякове дерево. Списки суміжності вершин переглядаються алгоритмом за зростанням міток вершин.\par
\begin{center}
	\adjustimage{max size={0.9\linewidth}{0.9\paperheight}}
	{../10/Traversals/60.png}
\end{center}
\par Які з наведених ребер входять до складу отриманого кістякового дерева?\par
1) (1, 5)\par
2) (3, 6)\par
3) (2, 5)\par
4) (0, 2)\par
5) (1, 6)\par
6) (1, 2)\par
{\vskip 15pt} У формі вибрати номери відповідних ребер.\par
%answer:1 2 4 6
\end {large}}
%end
{\smallskip \begin {small} Затверджено на засіданні кафедри теоретичної кібернетики, протокол \No 12 від   19.05.2020 р.\par\smallskip\smallskip
{\parbox {6.5 cm}{Зав. кафедри \dotfill Ю.В.~Крак}}\hspace {0.7cm}{\hbox to 6.5 cm{ Екзаменатор \dotfill Т.О.~Карнаух}} \end{small}}
%begin
{{\begin {small}\par
{ \center  Київський національний університет імені Тараса Шевченка \par}
{\parbox {5 cm}{Освітня програма \hfill}{\parbox {7 cm}{Прикладна математика, Системний аналіз \hfill}}\parbox {6 cm}{\hfill Семестр 2}}\par
{\parbox {5 cm} {Навчальний предмет \hfill}\parbox {7 cm}{Програмування \hfill}\parbox {6cm}{\hfill}\par}\vspace{-7pt} \end{small}}{\center Екзаменаційний білет \No \textbf 
{ 422 }\par}}
{
\begin {large}
1. Що буде виведено за виконання?
code
try:
    res = 3/7%6*8
    if isinstance(res, bool):
        print(f'bool;{res}')
    elif isinstance(res, int):
        print(f'int;{res}')
    elif isinstance(res, float):
        print(f'float;{res:.2f}')
    else:
        print('???')
except:
    print('ERR')
code

2. Що буде виведено за виконання?
code
a = 4
b = 6
c = 1

def h(a):
    a = 1
    b = 4
    c = 2
    return a + b + c

a = 3
b = 2
c = 7

try:
    print(h(a), a, b, c, sep=';')
except:
    print('ERR', a, b, c, sep=';')
code

3. Що буде виведено за виконання?
code
try:
    n = 10
    while n>2:
        n += -2
    print(n, end=';')
except:
    print('ERR')
code

4. Що буде виведено за виконання?
code
try:
    for d in range(11, -7, -1):
        if d<=8:
            continue
        print(d, end=';')
    else:
        print(d,end=';')
    print(d)
except:
    print('ERR')
code

5. Що буде виведено за виконання?\par (Дійсні значення, якщо наявні, в цьому завданні будуть виводитись у форматі з фіксованою крапкою та, можливо, з доволі великою кількістю десяткових розрядів. У відповіді для дійсних значень у ЦЬОМУ завданні достатньо вказати один десятковий розряд після крапки, відкинувши решту розрядів без округлення. Наприклад, замість 13.66666666666667 достатньо вказати 13.6, замість 25.23 достатньо вказати 25.2)
code
try:
    print(7, end=';')
    print(int('b8'), end=';')
    print(2, end=';')
except ZeroDivisionError:
    print(8, end=';')
except KeyboardInterrupt:
    print(0, end=';')
except:
    print('ERR', end=';')
else:
    print(7, end=';')
finally:
    print(3)
code

6. Що буде виведено за виконання?
code
def f(n):
    print(n%10, end='')
    if n>9:
        f(n//10)
 
f(543216)
code

7. Що буде виведено за виконання?
code
class A:
    c = 1

    def _g1(self): print(2, end=';')
    def _g2(self): print(3, end=';')

    def main(self):
        try: print(self.c, end=';')
        except: print('ERR', end=';')

        try: self._g1()
        except: print('ERR', end=';')

        try: self._g2()
        except: print('ERR', end=';')


class B(A):
    c = 4

    def _g1(self): print(5, end=';')
    def _g2(self): print(6, end=';')

try: B().main()
except: print('ERR', end=';')
code

8. Що буде виведено за виконання?
code
def f(arg, n):
    n = 6
    tmp = arg
    tmp[7:] = [2,3]
    return len(tmp)>3

try:
    d = ['0','1','2','3','4']
    n = 0
    f(d, n)
    print(n, end=';')
    for el in d:
        print(el, end=';')
except:
    print('ERR')
code

9. Що буде виведено за виконання?
code
class A:
    b = 1
    a = 2

    def __init__(self):
        a = 3

class B(A):
    b = 4

    def __init__(self):
        super().__init__()
        self.c = 7

obj = B()
obj.b = obj.c + 10
B.b = 13

try: print(obj.b, end= ';')
except: print('ERR', end=';')

try: print(A.b, end= ';')
except: print('ERR', end=';')

try: print(B.b, end= ';')
except: print('ERR', end=';')
code

10. 
Для графа на рисунку з вершини 5 запущено алгоритм пошуку в глибину, що будує кістякове дерево. Списки суміжності вершин переглядаються алгоритмом за зростанням міток вершин.\par
\begin{center}
	\adjustimage{max size={0.9\linewidth}{0.9\paperheight}}
	{../10/Traversals/81.png}
\end{center}
\par Які з наведених ребер входять до складу отриманого кістякового дерева?\par
1) (0, 1)\par
2) (2, 5)\par
3) (1, 2)\par
4) (0, 6)\par
5) (5, 6)\par
6) (3, 5)\par
{\vskip 15pt} У формі вибрати номери відповідних ребер.\par
%answer:1 2
\end {large}}
%end
{\smallskip \begin {small} Затверджено на засіданні кафедри теоретичної кібернетики, протокол \No 12 від   19.05.2020 р.\par\smallskip\smallskip
{\parbox {6.5 cm}{Зав. кафедри \dotfill Ю.В.~Крак}}\hspace {0.7cm}{\hbox to 6.5 cm{ Екзаменатор \dotfill Т.О.~Карнаух}} \end{small}}
%begin
{{\begin {small}\par
{ \center  Київський національний університет імені Тараса Шевченка \par}
{\parbox {5 cm}{Освітня програма \hfill}{\parbox {7 cm}{Прикладна математика, Системний аналіз \hfill}}\parbox {6 cm}{\hfill Семестр 2}}\par
{\parbox {5 cm} {Навчальний предмет \hfill}\parbox {7 cm}{Програмування \hfill}\parbox {6cm}{\hfill}\par}\vspace{-7pt} \end{small}}{\center Екзаменаційний білет \No \textbf 
{ 423 }\par}}
{
\begin {large}
1. Що буде виведено за виконання?
code
try:
    res = 7==8 or not 3.0>=4 and True<7.0
    if isinstance(res, bool):
        print(f'bool;{res}')
    elif isinstance(res, int):
        print(f'int;{res}')
    elif isinstance(res, float):
        print(f'float;{res:.2f}')
    else:
        print('???')
except:
    print('ERR')
code

2. Що буде виведено за виконання?
code
a = 7
b = 4
c = 0

def h(b):
    a -= 2
    b = 5
    c = 4
    return a + b + c

a = 3
b = 3
c = 4

try:
    print(h(a), a, b, c, sep=';')
except:
    print('ERR', a, b, c, sep=';')
code

3. Що буде виведено за виконання?
code
try:
    j = 0
    while j<9:
        j += 2
    print(j, end=';')
except:
    print('ERR')
code

4. Що буде виведено за виконання?
code
try:
    for f in range(-11, -8, 3):
        if f<6:
            continue
        print(f, end=';');f = -1
    else:
        print(f,end=';')
    print(f)
except:
    print('ERR')
code

5. Що буде виведено за виконання?\par (Дійсні значення, якщо наявні, в цьому завданні будуть виводитись у форматі з фіксованою крапкою та, можливо, з доволі великою кількістю десяткових розрядів. У відповіді для дійсних значень у ЦЬОМУ завданні достатньо вказати один десятковий розряд після крапки, відкинувши решту розрядів без округлення. Наприклад, замість 13.66666666666667 достатньо вказати 13.6, замість 25.23 достатньо вказати 25.2)
code
try:
    print(6, end=';')
    print(1%2, end=';')
    print(9, end=';')
except ValueError:
    print(5, end=';')
except ValueError:
    print(4, end=';')
except:
    print('ERR', end=';')
else:
    print(2, end=';')
finally:
    print(4)
code

6. Що буде виведено за виконання?
code
def f(n):
    if n>9:
        f(n//10)
    else:
        print(n%10, end='')

f(543216)
code

7. Що буде виведено за виконання?
code
class A:
    a = 1

    def h1(self): print(2, end=';')
    def _h2(self): print(3, end=';')

    def main(self):
        try: print(self.a, end=';')
        except: print('ERR', end=';')

        try: self.h1()
        except: print('ERR', end=';')

        try: self._h2()
        except: print('ERR', end=';')


class B(A):
    a = 4

    def h1(self): print(5, end=';')
    def _h2(self): print(6, end=';')

try: B().main()
except: print('ERR', end=';')
code

8. Що буде виведено за виконання?
code
def f(arg, n):
    n = 9
    tmp = arg
    tmp[-4:7] = [1,2]
    return len(tmp)>3

try:
    f = [10,11,12,13,14]
    n = 5
    f(f, n)
    print(n, end=';')
    for el in f:
        print(el, end=';')
except:
    print('ERR')
code

9. Що буде виведено за виконання?
code
class A:
    a = 1
    b = 2

    def __init__(self):
        a = 3

class B(A):
    b = 4

    def __init__(self):
        super().__init__()
        self.c = 7

obj = B()
obj.b = obj.a + 10
B.b = 13

try: print(obj.b, end= ';')
except: print('ERR', end=';')

try: print(A.b, end= ';')
except: print('ERR', end=';')

try: print(B.b, end= ';')
except: print('ERR', end=';')
code

10. 
Для графа на рисунку з вершини 1 запущено алгоритм пошуку в глибину, що будує кістякове дерево. Списки суміжності вершин переглядаються алгоритмом за зростанням міток вершин.\par
\begin{center}
	\adjustimage{max size={0.9\linewidth}{0.9\paperheight}}
	{../10/Traversals/214.png}
\end{center}
\par Які з наведених ребер входять до складу отриманого кістякового дерева?\par
1) (1, 6)\par
2) (2, 4)\par
3) (0, 1)\par
4) (1, 3)\par
5) (0, 3)\par
6) (0, 6)\par
{\vskip 15pt} У формі вибрати номери відповідних ребер.\par
%answer:2 3 5
\end {large}}
%end
{\smallskip \begin {small} Затверджено на засіданні кафедри теоретичної кібернетики, протокол \No 12 від   19.05.2020 р.\par\smallskip\smallskip
{\parbox {6.5 cm}{Зав. кафедри \dotfill Ю.В.~Крак}}\hspace {0.7cm}{\hbox to 6.5 cm{ Екзаменатор \dotfill Т.О.~Карнаух}} \end{small}}
%begin
{{\begin {small}\par
{ \center  Київський національний університет імені Тараса Шевченка \par}
{\parbox {5 cm}{Освітня програма \hfill}{\parbox {7 cm}{Прикладна математика, Системний аналіз \hfill}}\parbox {6 cm}{\hfill Семестр 2}}\par
{\parbox {5 cm} {Навчальний предмет \hfill}\parbox {7 cm}{Програмування \hfill}\parbox {6cm}{\hfill}\par}\vspace{-7pt} \end{small}}{\center Екзаменаційний білет \No \textbf 
{ 424 }\par}}
{
\begin {large}
1. Що буде виведено за виконання?
code
def f(n):
    print('f', end='')
    return n

try:
    print(f(-4) and f(-1))
except:
    print('ERR')
code

2. Що буде виведено за виконання?
code
a = 0
b = 8
c = 9

def g(b):
    global c
    a = 3
    b = 5
    c = 1
    return a + b + c

a = 5
b = 0
c = 1

try:
    print(g(b), a, b, c, sep=';')
except:
    print('ERR', a, b, c, sep=';')
code

3. Що буде виведено за виконання?
code
try:
    t = 15
    while t>=4:
        t += -3
    print(t, end=';')
except:
    print('ERR')
code

4. Що буде виведено за виконання?
code
try:
    for e in range(7, -10, 2):
        if e>5:
            continue
        print(e, end=';')
    else:
        print(e,end=';')
    print(e)
except:
    print('ERR')
code

5. Що буде виведено за виконання?\par (Дійсні значення, якщо наявні, в цьому завданні будуть виводитись у форматі з фіксованою крапкою та, можливо, з доволі великою кількістю десяткових розрядів. У відповіді для дійсних значень у ЦЬОМУ завданні достатньо вказати один десятковий розряд після крапки, відкинувши решту розрядів без округлення. Наприклад, замість 13.66666666666667 достатньо вказати 13.6, замість 25.23 достатньо вказати 25.2)
code
try:
    print(2, end=';')
    print(1/0, end=';')
    print(3, end=';')
except TypeError:
    print(7, end=';')
except ZeroDivisionError:
    print(0, end=';')
except:
    print('ERR', end=';')
else:
    print(9, end=';')
finally:
    print(5)
code

6. Що буде виведено за виконання?
code
def f(s):
    for el in s:
        el += 1
    return s
 
s = [1, 2, 3]
s = f(s)
print(s)
code

7. Що буде виведено за виконання?
code
class A:
    b = 1

    def _h1(self): print(2, end=';')
    def __h2(self): print(3, end=';')

    def main(self):
        try: print(self.b, end=';')
        except: print('ERR', end=';')

        try: self._h1()
        except: print('ERR', end=';')

        try: self.__h2()
        except: print('ERR', end=';')


class B(A):
    b = 4

    def _h1(self): print(5, end=';')
    def __h2(self): print(6, end=';')

try: B().main()
except: print('ERR', end=';')
code

8. Що буде виведено за виконання?
code
def f(arg, n):
    arg[20] = 2
    n += 2

try:
    n = 6
    t = {56:2, 65:5, 56:9}
    f(t, n)
    print(n, end=';')
    for x in t.items():
        print(*x, sep=';', end=';')
except:
    print('ERR')
code

9. Що буде виведено за виконання?
code
class A:
    a = 1
    c = 2

    def __init__(self):
        c = 3

class B(A):
    a = 4

    def __init__(self):
        super().__init__()
        self.b = 7

obj = B()
obj.c = obj.b + 10
B.c = 13

try: print(obj.c, end= ';')
except: print('ERR', end=';')

try: print(A.c, end= ';')
except: print('ERR', end=';')

try: print(B.c, end= ';')
except: print('ERR', end=';')
code

10. 
Для графа на рисунку з вершини 4 запущено алгоритм пошуку в глибину, що будує кістякове дерево. Списки суміжності вершин переглядаються алгоритмом за зростанням міток вершин.\par
\begin{center}
	\adjustimage{max size={0.9\linewidth}{0.9\paperheight}}
	{../10/Traversals/167.png}
\end{center}
\par Які з наведених ребер входять до складу отриманого кістякового дерева?\par
1) (1, 4)\par
2) (0, 3)\par
3) (1, 5)\par
4) (0, 2)\par
5) (1, 3)\par
6) (0, 1)\par
{\vskip 15pt} У формі вибрати номери відповідних ребер.\par
%answer:1 3 4 6
\end {large}}
%end
{\smallskip \begin {small} Затверджено на засіданні кафедри теоретичної кібернетики, протокол \No 12 від   19.05.2020 р.\par\smallskip\smallskip
{\parbox {6.5 cm}{Зав. кафедри \dotfill Ю.В.~Крак}}\hspace {0.7cm}{\hbox to 6.5 cm{ Екзаменатор \dotfill Т.О.~Карнаух}} \end{small}}
%begin
{{\begin {small}\par
{ \center  Київський національний університет імені Тараса Шевченка \par}
{\parbox {5 cm}{Освітня програма \hfill}{\parbox {7 cm}{Прикладна математика, Системний аналіз \hfill}}\parbox {6 cm}{\hfill Семестр 2}}\par
{\parbox {5 cm} {Навчальний предмет \hfill}\parbox {7 cm}{Програмування \hfill}\parbox {6cm}{\hfill}\par}\vspace{-7pt} \end{small}}{\center Екзаменаційний білет \No \textbf 
{ 425 }\par}}
{
\begin {large}
1. Що буде виведено за виконання?
code
try:
    res = 1**1.0*True
    if isinstance(res, bool):
        print(f'bool;{res}')
    elif isinstance(res, int):
        print(f'int;{res}')
    elif isinstance(res, float):
        print(f'float;{res:.2f}')
    else:
        print('???')
except:
    print('ERR')
code

2. Що буде виведено за виконання?
code
a = 7
b = 5
c = 2

def g(b):
    a = 2
    b *= 5
    c = 3
    return a + b + c

a = 1
b = 8
c = 9

try:
    print(g(a), a, b, c, sep=';')
except:
    print('ERR', a, b, c, sep=';')
code

3. Що буде виведено за виконання?
code
try:
    t = 2
    while t<=5:
        t += 1
    print(t, end=';')
except:
    print('ERR')
code

4. Що буде виведено за виконання?
code
try:
    for c in range(-7, 5, -2):
        if c>=3:
            continue
        print(c, end=';');c = 2
    else:
        print(c,end=';')
    print(c)
except:
    print('ERR')
code

5. Що буде виведено за виконання?\par (Дійсні значення, якщо наявні, в цьому завданні будуть виводитись у форматі з фіксованою крапкою та, можливо, з доволі великою кількістю десяткових розрядів. У відповіді для дійсних значень у ЦЬОМУ завданні достатньо вказати один десятковий розряд після крапки, відкинувши решту розрядів без округлення. Наприклад, замість 13.66666666666667 достатньо вказати 13.6, замість 25.23 достатньо вказати 25.2)
code
try:
    print(8, end=';')
    print(6j<=0j, end=';')
    print(9, end=';')
except Exception:
    print(2, end=';')
except KeyboardInterrupt:
    print(7, end=';')
except:
    print('ERR', end=';')
else:
    print(6, end=';')
finally:
    print(3)
code

6. Що буде виведено за виконання?
code
def f(n):
    if n>9:
        f(n//100)
    print(n%10, end='')
 
f(123456)
code

7. Що буде виведено за виконання?
code
class A:
    b = 1

    def _f1(self): print(2, end=';')
    def f2(self): print(3, end=';')

    def main(self):
        try: print(self.b, end=';')
        except: print('ERR', end=';')

        try: self._f1()
        except: print('ERR', end=';')

        try: self.f2()
        except: print('ERR', end=';')


class B(A):
    b = 4

    def _f1(self): print(5, end=';')
    def f2(self): print(6, end=';')

try: B().main()
except: print('ERR', end=';')
code

8. Що буде виведено за виконання?
code
def f(arg, n):
    arg[10] = 6
    n = 6
    arg = set(arg.keys())

try:
    n = 2
    s = {65:2, 38:6, 60:3}
    f(s, n)
    print(n, end=';')
    for x in s.items():
        print(*x, sep=';', end=';')
except:
    print('ERR')
code

9. Що буде виведено за виконання?
code
class A:
    c = 1
    b = 2

    def __init__(self):
        b = 3

class B(A):
    b = 4

    def __init__(self):
        super().__init__()
        self.a = 7

obj = B()
obj.a = obj.b + 10
B.a = 13

try: print(obj.a, end= ';')
except: print('ERR', end=';')

try: print(A.a, end= ';')
except: print('ERR', end=';')

try: print(B.a, end= ';')
except: print('ERR', end=';')
code

10. 
Для графа на рисунку з вершини 1 запущено алгоритм пошуку в глибину, що будує кістякове дерево. Списки суміжності вершин переглядаються алгоритмом за зростанням міток вершин.\par
\begin{center}
	\adjustimage{max size={0.9\linewidth}{0.9\paperheight}}
	{../10/Traversals/282.png}
\end{center}
\par Які з наведених ребер входять до складу отриманого кістякового дерева?\par
1) (0, 3)\par
2) (0, 5)\par
3) (1, 2)\par
4) (0, 4)\par
5) (4, 5)\par
6) (3, 5)\par
{\vskip 15pt} У формі вибрати номери відповідних ребер.\par
%answer:1 3 5 6
\end {large}}
%end
{\smallskip \begin {small} Затверджено на засіданні кафедри теоретичної кібернетики, протокол \No 12 від   19.05.2020 р.\par\smallskip\smallskip
{\parbox {6.5 cm}{Зав. кафедри \dotfill Ю.В.~Крак}}\hspace {0.7cm}{\hbox to 6.5 cm{ Екзаменатор \dotfill Т.О.~Карнаух}} \end{small}}
%begin
{{\begin {small}\par
{ \center  Київський національний університет імені Тараса Шевченка \par}
{\parbox {5 cm}{Освітня програма \hfill}{\parbox {7 cm}{Прикладна математика, Системний аналіз \hfill}}\parbox {6 cm}{\hfill Семестр 2}}\par
{\parbox {5 cm} {Навчальний предмет \hfill}\parbox {7 cm}{Програмування \hfill}\parbox {6cm}{\hfill}\par}\vspace{-7pt} \end{small}}{\center Екзаменаційний білет \No \textbf 
{ 426 }\par}}
{
\begin {large}
1. Що буде виведено за виконання?
code
try:
    res = 1**0.5*1
    if isinstance(res, bool):
        print(f'bool;{res}')
    elif isinstance(res, int):
        print(f'int;{res}')
    elif isinstance(res, float):
        print(f'float;{res:.2f}')
    else:
        print('???')
except:
    print('ERR')
code

2. Що буде виведено за виконання?
code
a = 6
b = 8
c = 3

def f(b):
    a = 5
    b = 4
    c = 3
    return a + b + c

a = 9
b = 4
c = 5

try:
    print(f(a), a, b, c, sep=';')
except:
    print('ERR', a, b, c, sep=';')
code

3. Що буде виведено за виконання?
code
try:
    m = 7
    while m>=9:
        m += -3
    print(m, end=';')
except:
    print('ERR')
code

4. Що буде виведено за виконання?
code
try:
    for a in range(-9, 3, -3):
        if a>4:
            break
        print(a, end=';')
    else:
        print(a,end=';')
    print(a)
except:
    print('ERR')
code

5. Що буде виведено за виконання?\par (Дійсні значення, якщо наявні, в цьому завданні будуть виводитись у форматі з фіксованою крапкою та, можливо, з доволі великою кількістю десяткових розрядів. У відповіді для дійсних значень у ЦЬОМУ завданні достатньо вказати один десятковий розряд після крапки, відкинувши решту розрядів без округлення. Наприклад, замість 13.66666666666667 достатньо вказати 13.6, замість 25.23 достатньо вказати 25.2)
code
try:
    print(8, end=';')
    print(int('5'), end=';')
    print(4, end=';')
except Exception:
    print(1, end=';')
except ZeroDivisionError:
    print(0, end=';')
except:
    print('ERR', end=';')
else:
    print(6, end=';')
finally:
    print(3)
code

6. Що буде виведено за виконання?
code
def f(n):
    if n>9:
        f(n//100)
        print(n%10, end='')
 
f(123456)
code

7. Що буде виведено за виконання?
code
class A:
    a = 1

    def g1(self): print(2, end=';')
    def _g2(self): print(3, end=';')

    def main(self):
        try: print(self.a, end=';')
        except: print('ERR', end=';')

        try: self.g1()
        except: print('ERR', end=';')

        try: self._g2()
        except: print('ERR', end=';')


class B(A):
    a = 4

    def g1(self): print(5, end=';')
    def _g2(self): print(6, end=';')

try: B().main()
except: print('ERR', end=';')
code

8. Що буде виведено за виконання?
code
def f(arg, n):
    arg[21] = 6
    n -= -3

try:
    n = 3
    s = {84:2, 21:5, 58:1, 76:4}
    f(s, n)
    print(n, end=';')
    for x in s :
        print(x, sep=';', end=';')
except:
    print('ERR')
code

9. Що буде виведено за виконання?
code
class A:
    b = 1
    c = 2

    def __init__(self):
        b = 3

class B(A):
    b = 4

    def __init__(self):
        super().__init__()
        self.a = 7

obj = B()
obj.a = obj.c + 10
B.a = 13

try: print(obj.a, end= ';')
except: print('ERR', end=';')

try: print(A.a, end= ';')
except: print('ERR', end=';')

try: print(B.a, end= ';')
except: print('ERR', end=';')
code

10. 
Для графа на рисунку з вершини 6 запущено алгоритм пошуку в глибину, що будує кістякове дерево. Списки суміжності вершин переглядаються алгоритмом за зростанням міток вершин.\par
\begin{center}
	\adjustimage{max size={0.9\linewidth}{0.9\paperheight}}
	{../10/Traversals/40.png}
\end{center}
\par Які з наведених ребер входять до складу отриманого кістякового дерева?\par
1) (0, 1)\par
2) (2, 3)\par
3) (2, 4)\par
4) (0, 4)\par
5) (1, 4)\par
6) (4, 5)\par
{\vskip 15pt} У формі вибрати номери відповідних ребер.\par
%answer:1 2 3 5 6
\end {large}}
%end
{\smallskip \begin {small} Затверджено на засіданні кафедри теоретичної кібернетики, протокол \No 12 від   19.05.2020 р.\par\smallskip\smallskip
{\parbox {6.5 cm}{Зав. кафедри \dotfill Ю.В.~Крак}}\hspace {0.7cm}{\hbox to 6.5 cm{ Екзаменатор \dotfill Т.О.~Карнаух}} \end{small}}
%begin
{{\begin {small}\par
{ \center  Київський національний університет імені Тараса Шевченка \par}
{\parbox {5 cm}{Освітня програма \hfill}{\parbox {7 cm}{Прикладна математика, Системний аналіз \hfill}}\parbox {6 cm}{\hfill Семестр 2}}\par
{\parbox {5 cm} {Навчальний предмет \hfill}\parbox {7 cm}{Програмування \hfill}\parbox {6cm}{\hfill}\par}\vspace{-7pt} \end{small}}{\center Екзаменаційний білет \No \textbf 
{ 427 }\par}}
{
\begin {large}
1. Що буде виведено за виконання?
code
try:
    res = 2**-2+-2
    if isinstance(res, bool):
        print(f'bool;{res}')
    elif isinstance(res, int):
        print(f'int;{res}')
    elif isinstance(res, float):
        print(f'float;{res:.2f}')
    else:
        print('???')
except:
    print('ERR')
code

2. Що буде виведено за виконання?
code
a = 0
b = 6
c = 8

def h(a):
    global c
    a = 1
    b += 3
    c = 1
    return a + b + c

a = 7
b = 6
c = 9

try:
    print(h(a), a, b, c, sep=';')
except:
    print('ERR', a, b, c, sep=';')
code

3. Що буде виведено за виконання?
code
try:
    i = 8
    while i>5:
        i += -2
    print(i, end=';')
except:
    print('ERR')
code

4. Що буде виведено за виконання?
code
try:
    for d in range(1, 6, 3):
        if d<=3:
            break
        print(d, end=';');d = -3
    else:
        print(d,end=';')
    print(d)
except:
    print('ERR')
code

5. Що буде виведено за виконання?\par (Дійсні значення, якщо наявні, в цьому завданні будуть виводитись у форматі з фіксованою крапкою та, можливо, з доволі великою кількістю десяткових розрядів. У відповіді для дійсних значень у ЦЬОМУ завданні достатньо вказати один десятковий розряд після крапки, відкинувши решту розрядів без округлення. Наприклад, замість 13.66666666666667 достатньо вказати 13.6, замість 25.23 достатньо вказати 25.2)
code
try:
    print(9, end=';')
    print(7j==2j, end=';')
    print(1, end=';')
except KeyboardInterrupt:
    print(0, end=';')
except TypeError:
    print(5, end=';')
except:
    print('ERR', end=';')
else:
    print(2, end=';')
finally:
    print(4)
code

6. Що буде виведено за виконання?
code
def f(s):
    s1 = s
    for i in range(len(s)):
        s1[i] += 1
        return s1

s = [1, 2, 3]
f(s)
print(s)
code

7. Що буде виведено за виконання?
code
class A:
    d = 1

    def _h1(self): print(2, end=';')
    def _h2(self): print(3, end=';')

    def main(self):
        try: print(self.d, end=';')
        except: print('ERR', end=';')

        try: self._h1()
        except: print('ERR', end=';')

        try: self._h2()
        except: print('ERR', end=';')


class B(A):
    d = 4

    def _h1(self): print(5, end=';')
    def _h2(self): print(6, end=';')

try: B().main()
except: print('ERR', end=';')
code

8. Що буде виведено за виконання?
code
def f(arg, n):
    arg[41] = 3
    n = 7
    arg = tuple(arg.values())

try:
    n = 3
    t = {63:1, 41:9, 81:1, 81:8}
    f(t, n)
    print(n, end=';')
    for x in t.keys():
        print(x, sep=';', end=';')
except:
    print('ERR')
code

9. Що буде виведено за виконання?
code
class A:
    a = 1
    c = 2

    def __init__(self):
        a = 3

class B(A):
    c = 4

    def __init__(self):
        super().__init__()
        self.b = 7

obj = B()
obj.a = obj.c + 10
B.a = 13

try: print(obj.a, end= ';')
except: print('ERR', end=';')

try: print(A.a, end= ';')
except: print('ERR', end=';')

try: print(B.a, end= ';')
except: print('ERR', end=';')
code

10. 
Для графа на рисунку з вершини 5 запущено алгоритм пошуку в глибину, що будує кістякове дерево. Списки суміжності вершин переглядаються алгоритмом за зростанням міток вершин.\par
\begin{center}
	\adjustimage{max size={0.9\linewidth}{0.9\paperheight}}
	{../10/Traversals/90.png}
\end{center}
\par Які з наведених ребер входять до складу отриманого кістякового дерева?\par
1) (2, 5)\par
2) (0, 1)\par
3) (0, 5)\par
4) (3, 4)\par
5) (1, 4)\par
6) (3, 5)\par
{\vskip 15pt} У формі вибрати номери відповідних ребер.\par
%answer:2 3 5
\end {large}}
%end
{\smallskip \begin {small} Затверджено на засіданні кафедри теоретичної кібернетики, протокол \No 12 від   19.05.2020 р.\par\smallskip\smallskip
{\parbox {6.5 cm}{Зав. кафедри \dotfill Ю.В.~Крак}}\hspace {0.7cm}{\hbox to 6.5 cm{ Екзаменатор \dotfill Т.О.~Карнаух}} \end{small}}
%begin
{{\begin {small}\par
{ \center  Київський національний університет імені Тараса Шевченка \par}
{\parbox {5 cm}{Освітня програма \hfill}{\parbox {7 cm}{Прикладна математика, Системний аналіз \hfill}}\parbox {6 cm}{\hfill Семестр 2}}\par
{\parbox {5 cm} {Навчальний предмет \hfill}\parbox {7 cm}{Програмування \hfill}\parbox {6cm}{\hfill}\par}\vspace{-7pt} \end{small}}{\center Екзаменаційний білет \No \textbf 
{ 428 }\par}}
{
\begin {large}
1. Що буде виведено за виконання?
code
try:
    res = "2j">True
    if isinstance(res, bool):
        print(f'bool;{res}')
    elif isinstance(res, int):
        print(f'int;{res}')
    elif isinstance(res, float):
        print(f'float;{res:.2f}')
    else:
        print('???')
except:
    print('ERR')
code

2. Що буде виведено за виконання?
code
a = 3
b = 5
c = 2

def g(b):
    global c
    a *= 4
    b = 2
    c = 5
    return a + b + c

a = 1
b = 0
c = 4

try:
    print(g(b), a, b, c, sep=';')
except:
    print('ERR', a, b, c, sep=';')
code

3. Що буде виведено за виконання?
code
try:
    k = 9
    while k<7:
        k += 2
    print(k, end=';')
except:
    print('ERR')
code

4. Що буде виведено за виконання?
code
try:
    for a in range(5, -15, -1):
        if a<8:
            continue
        print(a, end=';')
    else:
        print(a,end=';')
    print(a)
except:
    print('ERR')
code

5. Що буде виведено за виконання?\par (Дійсні значення, якщо наявні, в цьому завданні будуть виводитись у форматі з фіксованою крапкою та, можливо, з доволі великою кількістю десяткових розрядів. У відповіді для дійсних значень у ЦЬОМУ завданні достатньо вказати один десятковий розряд після крапки, відкинувши решту розрядів без округлення. Наприклад, замість 13.66666666666667 достатньо вказати 13.6, замість 25.23 достатньо вказати 25.2)
code
try:
    print(8, end=';')
    print(5//0.0, end=';')
    print(1, end=';')
except ValueError:
    print(2, end=';')
except TypeError:
    print(8, end=';')
except:
    print('ERR', end=';')
else:
    print(9, end=';')
finally:
    print(3)
code

6. Що буде виведено за виконання?
code
def f(n):
    if n>9:
        f(n//100)
    else:
        print(n%10,end='')

f(987654)
code

7. Що буде виведено за виконання?
code
class A:
    c = 1

    def _g1(self): print(2, end=';')
    def __g2(self): print(3, end=';')

    def main(self):
        try: print(self.c, end=';')
        except: print('ERR', end=';')

        try: self._g1()
        except: print('ERR', end=';')

        try: self.__g2()
        except: print('ERR', end=';')


class B(A):
    c = 4

    def _g1(self): print(5, end=';')
    def __g2(self): print(6, end=';')

try: B().main()
except: print('ERR', end=';')
code

8. Що буде виведено за виконання?
code
def f(arg, n):
    n = 3
    tmp = arg
    del tmp[:-2]
    return len(tmp)>3

try:
    a = [10,11,12,13,14]
    n = 2
    f(a, n)
    print(n, end=';')
    for el in a:
        print(el, end=';')
except:
    print('ERR')
code

9. Що буде виведено за виконання?
code
class A:
    c = 1
    b = 2

    def __init__(self):
        b = 3

class B(A):
    c = 4

    def __init__(self):
        super().__init__()
        self.a = 7

obj = B()
obj.b = obj.a + 10
B.b = 13

try: print(obj.b, end= ';')
except: print('ERR', end=';')

try: print(A.b, end= ';')
except: print('ERR', end=';')

try: print(B.b, end= ';')
except: print('ERR', end=';')
code

10. 
Для графа на рисунку з вершини 0 запущено алгоритм пошуку в глибину, що будує кістякове дерево. Списки суміжності вершин переглядаються алгоритмом за зростанням міток вершин.\par
\begin{center}
	\adjustimage{max size={0.9\linewidth}{0.9\paperheight}}
	{../10/Traversals/77.png}
\end{center}
\par Які з наведених ребер входять до складу отриманого кістякового дерева?\par
1) (0, 1)\par
2) (2, 6)\par
3) (2, 5)\par
4) (0, 3)\par
5) (1, 5)\par
6) (1, 6)\par
{\vskip 15pt} У формі вибрати номери відповідних ребер.\par
%answer:1 2 3
\end {large}}
%end
{\smallskip \begin {small} Затверджено на засіданні кафедри теоретичної кібернетики, протокол \No 12 від   19.05.2020 р.\par\smallskip\smallskip
{\parbox {6.5 cm}{Зав. кафедри \dotfill Ю.В.~Крак}}\hspace {0.7cm}{\hbox to 6.5 cm{ Екзаменатор \dotfill Т.О.~Карнаух}} \end{small}}
%begin
{{\begin {small}\par
{ \center  Київський національний університет імені Тараса Шевченка \par}
{\parbox {5 cm}{Освітня програма \hfill}{\parbox {7 cm}{Прикладна математика, Системний аналіз \hfill}}\parbox {6 cm}{\hfill Семестр 2}}\par
{\parbox {5 cm} {Навчальний предмет \hfill}\parbox {7 cm}{Програмування \hfill}\parbox {6cm}{\hfill}\par}\vspace{-7pt} \end{small}}{\center Екзаменаційний білет \No \textbf 
{ 429 }\par}}
{
\begin {large}
1. Що буде виведено за виконання?
code
try:
    res = 6//8/8*9
    if isinstance(res, bool):
        print(f'bool;{res}')
    elif isinstance(res, int):
        print(f'int;{res}')
    elif isinstance(res, float):
        print(f'float;{res:.2f}')
    else:
        print('???')
except:
    print('ERR')
code

2. Що буде виведено за виконання?
code
a = 7
b = 2
c = 6

def f(a):
    global c
    a = 2
    b = 4
    c = 1
    return a + b + c

a = 5
b = 7
c = 1

try:
    print(f(b), a, b, c, sep=';')
except:
    print('ERR', a, b, c, sep=';')
code

3. Що буде виведено за виконання?
code
try:
    j = 12
    while j>6:
        j += -3
    print(j, end=';')
except:
    print('ERR')
code

4. Що буде виведено за виконання?
code
try:
    for b in range(-15, 10, 2):
        if b<=5:
            continue
        print(b, end=';')
    else:
        print('end',end=';')
    print(b)
except:
    print('ERR')
code

5. Що буде виведено за виконання?\par (Дійсні значення, якщо наявні, в цьому завданні будуть виводитись у форматі з фіксованою крапкою та, можливо, з доволі великою кількістю десяткових розрядів. У відповіді для дійсних значень у ЦЬОМУ завданні достатньо вказати один десятковий розряд після крапки, відкинувши решту розрядів без округлення. Наприклад, замість 13.66666666666667 достатньо вказати 13.6, замість 25.23 достатньо вказати 25.2)
code
try:
    print(4, end=';')
    print(0/3, end=';')
    print(6, end=';')
except ValueError:
    print(7, end=';')
except KeyboardInterrupt:
    print(5, end=';')
except:
    print('ERR', end=';')
else:
    print(2, end=';')
finally:
    print(1)
code

6. Що буде виведено за виконання?
code
def f(n):
    print(n%2, end='')
    if n>2:
        f(n//2)
 
f(16)
code

7. Що буде виведено за виконання?
code
class A:
    c = 1

    def _f1(self): print(2, end=';')
    def f2(self): print(3, end=';')

    def main(self):
        try: print(self.c, end=';')
        except: print('ERR', end=';')

        try: self._f1()
        except: print('ERR', end=';')

        try: self.f2()
        except: print('ERR', end=';')


class B(A):
    c = 4

    def _f1(self): print(5, end=';')
    def f2(self): print(6, end=';')

try: B().main()
except: print('ERR', end=';')
code

8. Що буде виведено за виконання?
code
def f(arg, n):
    arg[37] = 4
    n = 8
    arg = list(arg.items())

try:
    n = 0
    d = {60:4, 59:7, 59:6}
    f(d, n)
    print(n, end=';')
    for x in d.values():
        print(x, sep=';', end=';')
except:
    print('ERR')
code

9. Що буде виведено за виконання?
code
class A:
    b = 1
    c = 2

    def __init__(self):
        b = 3

class B(A):
    c = 4

    def __init__(self):
        super().__init__()
        self.a = 7

obj = B()
obj.c = obj.b + 10
B.c = 13

try: print(obj.c, end= ';')
except: print('ERR', end=';')

try: print(A.c, end= ';')
except: print('ERR', end=';')

try: print(B.c, end= ';')
except: print('ERR', end=';')
code

10. 
Для графа на рисунку з вершини 5 запущено алгоритм пошуку в глибину, що будує кістякове дерево. Списки суміжності вершин переглядаються алгоритмом за зростанням міток вершин.\par
\begin{center}
	\adjustimage{max size={0.9\linewidth}{0.9\paperheight}}
	{../10/Traversals/78.png}
\end{center}
\par Які з наведених ребер входять до складу отриманого кістякового дерева?\par
1) (1, 5)\par
2) (1, 6)\par
3) (1, 2)\par
4) (0, 1)\par
5) (2, 5)\par
6) (5, 6)\par
{\vskip 15pt} У формі вибрати номери відповідних ребер.\par
%answer:1 4
\end {large}}
%end
{\smallskip \begin {small} Затверджено на засіданні кафедри теоретичної кібернетики, протокол \No 12 від   19.05.2020 р.\par\smallskip\smallskip
{\parbox {6.5 cm}{Зав. кафедри \dotfill Ю.В.~Крак}}\hspace {0.7cm}{\hbox to 6.5 cm{ Екзаменатор \dotfill Т.О.~Карнаух}} \end{small}}
%begin
{{\begin {small}\par
{ \center  Київський національний університет імені Тараса Шевченка \par}
{\parbox {5 cm}{Освітня програма \hfill}{\parbox {7 cm}{Прикладна математика, Системний аналіз \hfill}}\parbox {6 cm}{\hfill Семестр 2}}\par
{\parbox {5 cm} {Навчальний предмет \hfill}\parbox {7 cm}{Програмування \hfill}\parbox {6cm}{\hfill}\par}\vspace{-7pt} \end{small}}{\center Екзаменаційний білет \No \textbf 
{ 430 }\par}}
{
\begin {large}
1. Що буде виведено за виконання?
code
try:
    res = 9.0j<"False"
    if isinstance(res, bool):
        print(f'bool;{res}')
    elif isinstance(res, int):
        print(f'int;{res}')
    elif isinstance(res, float):
        print(f'float;{res:.2f}')
    else:
        print('???')
except:
    print('ERR')
code

2. Що буде виведено за виконання?
code
a = 0
b = 9
c = 7

def f(a):
    a = 5
    b -= 4
    c = 1
    return a + b + c

a = 5
b = 4
c = 8

try:
    print(f(b), a, b, c, sep=';')
except:
    print('ERR', a, b, c, sep=';')
code

3. Що буде виведено за виконання?
code
try:
    k = 5
    while k>=8:
        k += -2
    print(k, end=';')
except:
    print('ERR')
code

4. Що буде виведено за виконання?
code
try:
    for f in range(2, -6, -2):
        if f>6:
            continue
        print(f, end=';');f = 8
        if f>=7:
            break
    else:
        print(f,end=';')
    print(f)
except:
    print('ERR')
code

5. Що буде виведено за виконання?\par (Дійсні значення, якщо наявні, в цьому завданні будуть виводитись у форматі з фіксованою крапкою та, можливо, з доволі великою кількістю десяткових розрядів. У відповіді для дійсних значень у ЦЬОМУ завданні достатньо вказати один десятковий розряд після крапки, відкинувши решту розрядів без округлення. Наприклад, замість 13.66666666666667 достатньо вказати 13.6, замість 25.23 достатньо вказати 25.2)
code
try:
    print(3, end=';')
    print(int('d7'), end=';')
    print(8, end=';')
except Exception:
    print(4, end=';')
except ZeroDivisionError:
    print(9, end=';')
except:
    print('ERR', end=';')
else:
    print(0, end=';')
finally:
    print(6)
code

6. Що буде виведено за виконання?
code
def f(s):
    s1 = s
    for i in range(len(s)):
        s1[i] += 1
    return s1

s = [1, 2, 3]
f(s)
print(s)
code

7. Що буде виведено за виконання?
code
class A:
    d = 1

    def __f1(self): print(2, end=';')
    def _f2(self): print(3, end=';')

    def main(self):
        try: print(self.d, end=';')
        except: print('ERR', end=';')

        try: self.__f1()
        except: print('ERR', end=';')

        try: self._f2()
        except: print('ERR', end=';')


class B(A):
    d = 4

    def __f1(self): print(5, end=';')
    def _f2(self): print(6, end=';')

try: B().main()
except: print('ERR', end=';')
code

8. Що буде виведено за виконання?
code
def f(arg, n):
    arg[70] = 2
    n *= -2

try:
    n = 4
    t = {11:5, 47:6, 47:1}
    f(t, n)
    print(n, end=';')
    for x in t :
        print(x, sep=';', end=';')
except:
    print('ERR')
code

9. Що буде виведено за виконання?
code
class A:
    b = 1
    a = 2

    def __init__(self):
        a = 3

class B(A):
    b = 4

    def __init__(self):
        super().__init__()
        self.c = 7

obj = B()
obj.a = obj.b + 10
B.a = 13

try: print(obj.a, end= ';')
except: print('ERR', end=';')

try: print(A.a, end= ';')
except: print('ERR', end=';')

try: print(B.a, end= ';')
except: print('ERR', end=';')
code

10. 
Для графа на рисунку з вершини 4 запущено алгоритм пошуку в глибину, що будує кістякове дерево. Списки суміжності вершин переглядаються алгоритмом за зростанням міток вершин.\par
\begin{center}
	\adjustimage{max size={0.9\linewidth}{0.9\paperheight}}
	{../10/Traversals/171.png}
\end{center}
\par Які з наведених ребер входять до складу отриманого кістякового дерева?\par
1) (4, 6)\par
2) (0, 6)\par
3) (0, 5)\par
4) (2, 4)\par
5) (1, 3)\par
6) (3, 5)\par
{\vskip 15pt} У формі вибрати номери відповідних ребер.\par
%answer:2 3 4 5 6
\end {large}}
%end
{\smallskip \begin {small} Затверджено на засіданні кафедри теоретичної кібернетики, протокол \No 12 від   19.05.2020 р.\par\smallskip\smallskip
{\parbox {6.5 cm}{Зав. кафедри \dotfill Ю.В.~Крак}}\hspace {0.7cm}{\hbox to 6.5 cm{ Екзаменатор \dotfill Т.О.~Карнаух}} \end{small}}
%begin
{{\begin {small}\par
{ \center  Київський національний університет імені Тараса Шевченка \par}
{\parbox {5 cm}{Освітня програма \hfill}{\parbox {7 cm}{Прикладна математика, Системний аналіз \hfill}}\parbox {6 cm}{\hfill Семестр 2}}\par
{\parbox {5 cm} {Навчальний предмет \hfill}\parbox {7 cm}{Програмування \hfill}\parbox {6cm}{\hfill}\par}\vspace{-7pt} \end{small}}{\center Екзаменаційний білет \No \textbf 
{ 431 }\par}}
{
\begin {large}
1. Що буде виведено за виконання?
code
def f(n):
    print('f', end='')
    return n!=3

try:
    print(f(2) or f(5))
except:
    print('ERR')
code

2. Що буде виведено за виконання?
code
a = 0
b = 9
c = 8

def h(b):
    global c
    a = 2
    b = 5
    c = 3
    return a + b + c

a = 3
b = 4
c = 2

try:
    print(h(b), a, b, c, sep=';')
except:
    print('ERR', a, b, c, sep=';')
code

3. Що буде виведено за виконання?
code
try:
    i = 1
    while i<=12:
        i += 2
    print(i, end=';')
except:
    print('ERR')
code

4. Що буде виведено за виконання?
code
try:
    for e in range(3, -12, -2):
        if e<7:
            continue
        print(e, end=';')
    else:
        print(e,end=';')
    print(e)
except:
    print('ERR')
code

5. Що буде виведено за виконання?\par (Дійсні значення, якщо наявні, в цьому завданні будуть виводитись у форматі з фіксованою крапкою та, можливо, з доволі великою кількістю десяткових розрядів. У відповіді для дійсних значень у ЦЬОМУ завданні достатньо вказати один десятковий розряд після крапки, відкинувши решту розрядів без округлення. Наприклад, замість 13.66666666666667 достатньо вказати 13.6, замість 25.23 достатньо вказати 25.2)
code
try:
    print(1, end=';')
    print(5//0, end=';')
    print(8, end=';')
except Exception:
    print(2, end=';')
except ZeroDivisionError:
    print(9, end=';')
except:
    print('ERR', end=';')
else:
    print(0, end=';')
finally:
    print(4)
code

6. Що буде виведено за виконання?
code
def f(n):
    if n<=9:
        print(n%10, end='')
    else:
        f(n//10)

f(123456)
code

7. Що буде виведено за виконання?
code
class A:
    a = 1

    def _h1(self): print(2, end=';')
    def __h2(self): print(3, end=';')

    def main(self):
        try: print(self.a, end=';')
        except: print('ERR', end=';')

        try: self._h1()
        except: print('ERR', end=';')

        try: self.__h2()
        except: print('ERR', end=';')


class B(A):
    a = 4

    def _h1(self): print(5, end=';')
    def __h2(self): print(6, end=';')

try: B().main()
except: print('ERR', end=';')
code

8. Що буде виведено за виконання?
code
def f(arg, n):
    n = 4
    tmp = arg
    tmp[-5:] = 'ac'
    return len(tmp)>3

try:
    h = ['0','1','2','3','4']
    n = 1
    f(h, n)
    print(n, end=';')
    for el in h:
        print(el, end=';')
except:
    print('ERR')
code

9. Що буде виведено за виконання?
code
class A:
    a = 1
    b = 2

    def __init__(self):
        b = 3

class B(A):
    a = 4

    def __init__(self):
        super().__init__()
        self.c = 7

obj = B()
obj.c = obj.a + 10
B.c = 13

try: print(obj.c, end= ';')
except: print('ERR', end=';')

try: print(A.c, end= ';')
except: print('ERR', end=';')

try: print(B.c, end= ';')
except: print('ERR', end=';')
code

10. 
Для графа на рисунку з вершини 5 запущено алгоритм пошуку в глибину, що будує кістякове дерево. Списки суміжності вершин переглядаються алгоритмом за зростанням міток вершин.\par
\begin{center}
	\adjustimage{max size={0.9\linewidth}{0.9\paperheight}}
	{../10/Traversals/98.png}
\end{center}
\par Які з наведених ребер входять до складу отриманого кістякового дерева?\par
1) (3, 4)\par
2) (0, 3)\par
3) (1, 6)\par
4) (0, 1)\par
5) (4, 6)\par
6) (5, 6)\par
{\vskip 15pt} У формі вибрати номери відповідних ребер.\par
%answer:1 2 3 4 6
\end {large}}
%end
{\smallskip \begin {small} Затверджено на засіданні кафедри теоретичної кібернетики, протокол \No 12 від   19.05.2020 р.\par\smallskip\smallskip
{\parbox {6.5 cm}{Зав. кафедри \dotfill Ю.В.~Крак}}\hspace {0.7cm}{\hbox to 6.5 cm{ Екзаменатор \dotfill Т.О.~Карнаух}} \end{small}}
%begin
{{\begin {small}\par
{ \center  Київський національний університет імені Тараса Шевченка \par}
{\parbox {5 cm}{Освітня програма \hfill}{\parbox {7 cm}{Прикладна математика, Системний аналіз \hfill}}\parbox {6 cm}{\hfill Семестр 2}}\par
{\parbox {5 cm} {Навчальний предмет \hfill}\parbox {7 cm}{Програмування \hfill}\parbox {6cm}{\hfill}\par}\vspace{-7pt} \end{small}}{\center Екзаменаційний білет \No \textbf 
{ 432 }\par}}
{
\begin {large}
1. Що буде виведено за виконання?
code
def f(n):
    print('f', end='')
    return n

try:
    print(f(6) or f(-4))
except:
    print('ERR')
code

2. Що буде виведено за виконання?
code
a = 1
b = 3
c = 6

def g(b):
    global c
    a = 3
    b += 2
    c = 3
    return a + b + c

a = 2
b = 3
c = 4

try:
    print(g(b), a, b, c, sep=';')
except:
    print('ERR', a, b, c, sep=';')
code

3. Що буде виведено за виконання?
code
try:
    m = 3
    while m<13:
        m += 1
    print(m, end=';')
except:
    print('ERR')
code

4. Що буде виведено за виконання?
code
try:
    for c in range(6, -9, -1):
        if c>=4:
            break
        print(c, end=';')
    else:
        print(c,end=';')
    print(c)
except:
    print('ERR')
code

5. Що буде виведено за виконання?\par (Дійсні значення, якщо наявні, в цьому завданні будуть виводитись у форматі з фіксованою крапкою та, можливо, з доволі великою кількістю десяткових розрядів. У відповіді для дійсних значень у ЦЬОМУ завданні достатньо вказати один десятковий розряд після крапки, відкинувши решту розрядів без округлення. Наприклад, замість 13.66666666666667 достатньо вказати 13.6, замість 25.23 достатньо вказати 25.2)
code
try:
    print(7, end=';')
    print(int('c3'), end=';')
    print(6, end=';')
except KeyboardInterrupt:
    print(1, end=';')
except ValueError:
    print(9, end=';')
except:
    print('ERR', end=';')
else:
    print(8, end=';')
finally:
    print(6)
code

6. Що буде виведено за виконання?
code
def f(s):
    for i in range(len(s)):
        s[i] += 1
    return
 
s = [1, 2, 3]
f(s)
print(s)
code

7. Що буде виведено за виконання?
code
class A:
    b = 1

    def _g1(self): print(2, end=';')
    def _g2(self): print(3, end=';')

    def main(self):
        try: print(self.b, end=';')
        except: print('ERR', end=';')

        try: self._g1()
        except: print('ERR', end=';')

        try: self._g2()
        except: print('ERR', end=';')


class B(A):
    b = 4

    def _g1(self): print(5, end=';')
    def _g2(self): print(6, end=';')

try: B().main()
except: print('ERR', end=';')
code

8. Що буде виведено за виконання?
code
def f(arg, n):
    arg[82] = 8
    n += 3

try:
    n = 7
    d = {25:0, 59:7, 38:0, 38:5}
    f(d, n)
    print(n, end=';')
    for x in d.values():
        print(x, sep=';', end=';')
except:
    print('ERR')
code

9. Що буде виведено за виконання?
code
class A:
    c = 1
    a = 2

    def __init__(self):
        c = 3

class B(A):
    a = 4

    def __init__(self):
        super().__init__()
        self.b = 7

obj = B()
obj.c = obj.b + 10
B.c = 13

try: print(obj.c, end= ';')
except: print('ERR', end=';')

try: print(A.c, end= ';')
except: print('ERR', end=';')

try: print(B.c, end= ';')
except: print('ERR', end=';')
code

10. 
Для графа на рисунку з вершини 0 запущено алгоритм пошуку в глибину, що будує кістякове дерево. Списки суміжності вершин переглядаються алгоритмом за зростанням міток вершин.\par
\begin{center}
	\adjustimage{max size={0.9\linewidth}{0.9\paperheight}}
	{../10/Traversals/184.png}
\end{center}
\par Які з наведених ребер входять до складу отриманого кістякового дерева?\par
1) (0, 5)\par
2) (4, 5)\par
3) (2, 5)\par
4) (4, 6)\par
5) (2, 6)\par
6) (0, 1)\par
{\vskip 15pt} У формі вибрати номери відповідних ребер.\par
%answer:1 3 4 6
\end {large}}
%end
{\smallskip \begin {small} Затверджено на засіданні кафедри теоретичної кібернетики, протокол \No 12 від   19.05.2020 р.\par\smallskip\smallskip
{\parbox {6.5 cm}{Зав. кафедри \dotfill Ю.В.~Крак}}\hspace {0.7cm}{\hbox to 6.5 cm{ Екзаменатор \dotfill Т.О.~Карнаух}} \end{small}}
%begin
{{\begin {small}\par
{ \center  Київський національний університет імені Тараса Шевченка \par}
{\parbox {5 cm}{Освітня програма \hfill}{\parbox {7 cm}{Прикладна математика, Системний аналіз \hfill}}\parbox {6 cm}{\hfill Семестр 2}}\par
{\parbox {5 cm} {Навчальний предмет \hfill}\parbox {7 cm}{Програмування \hfill}\parbox {6cm}{\hfill}\par}\vspace{-7pt} \end{small}}{\center Екзаменаційний білет \No \textbf 
{ 433 }\par}}
{
\begin {large}
1. Що буде виведено за виконання?
code
try:
    res = 3>2 and not 5.0!=False or 9<=6.0
    if isinstance(res, bool):
        print(f'bool;{res}')
    elif isinstance(res, int):
        print(f'int;{res}')
    elif isinstance(res, float):
        print(f'float;{res:.2f}')
    else:
        print('???')
except:
    print('ERR')
code

2. Що буде виведено за виконання?
code
a = 5
b = 7
c = 3

def g(a):
    a = 1
    b = 2
    c = 4
    return a + b + c

a = 1
b = 4
c = 2

try:
    print(g(a), a, b, c, sep=';')
except:
    print('ERR', a, b, c, sep=';')
code

3. Що буде виведено за виконання?
code
try:
    i = 6
    while i<=10:
        i += 1
    print(i, end=';')
except:
    print('ERR')
code

4. Що буде виведено за виконання?
code
try:
    for f in range(-1, -13, 2):
        if f<=3:
            continue
        print(f, end=';')
    else:
        print(f,end=';')
    print(f)
except:
    print('ERR')
code

5. Що буде виведено за виконання?\par (Дійсні значення, якщо наявні, в цьому завданні будуть виводитись у форматі з фіксованою крапкою та, можливо, з доволі великою кількістю десяткових розрядів. У відповіді для дійсних значень у ЦЬОМУ завданні достатньо вказати один десятковий розряд після крапки, відкинувши решту розрядів без округлення. Наприклад, замість 13.66666666666667 достатньо вказати 13.6, замість 25.23 достатньо вказати 25.2)
code
try:
    print(2, end=';')
    print(int('5'), end=';')
    print(7, end=';')
except TypeError:
    print(4, end=';')
except KeyboardInterrupt:
    print(3, end=';')
except:
    print('ERR', end=';')
else:
    print(0, end=';')
finally:
    print(7)
code

6. Що буде виведено за виконання?
code
def f(n):
    if n<=9:
        print(n%10,end='')
    else:
        f(n//10)

f(543216)
code

7. Що буде виведено за виконання?
code
class A:
    c = 1

    def g1(self): print(2, end=';')
    def _g2(self): print(3, end=';')

    def main(self):
        try: print(self.c, end=';')
        except: print('ERR', end=';')

        try: self.g1()
        except: print('ERR', end=';')

        try: self._g2()
        except: print('ERR', end=';')


class B(A):
    c = 4

    def g1(self): print(5, end=';')
    def _g2(self): print(6, end=';')

try: B().main()
except: print('ERR', end=';')
code

8. Що буде виведено за виконання?
code
def f(arg, n):
    arg[88] = 1
    n = 9
    arg = list(arg.keys())

try:
    n = 4
    s = {88:6, 77:2, 88:9}
    f(s, n)
    print(n, end=';')
    for x in s.items():
        print(x, sep=';', end=';')
except:
    print('ERR')
code

9. Що буде виведено за виконання?
code
class A:
    c = 1
    a = 2

    def __init__(self):
        c = 3

class B(A):
    a = 4

    def __init__(self):
        super().__init__()
        self.b = 7

obj = B()
obj.c = obj.b + 10
B.c = 13

try: print(obj.c, end= ';')
except: print('ERR', end=';')

try: print(A.c, end= ';')
except: print('ERR', end=';')

try: print(B.c, end= ';')
except: print('ERR', end=';')
code

10. 
Для графа на рисунку з вершини 4 запущено алгоритм пошуку в глибину, що будує кістякове дерево. Списки суміжності вершин переглядаються алгоритмом за зростанням міток вершин.\par
\begin{center}
	\adjustimage{max size={0.9\linewidth}{0.9\paperheight}}
	{../10/Traversals/39.png}
\end{center}
\par Які з наведених ребер входять до складу отриманого кістякового дерева?\par
1) (2, 6)\par
2) (1, 5)\par
3) (3, 5)\par
4) (3, 6)\par
5) (1, 4)\par
6) (1, 2)\par
{\vskip 15pt} У формі вибрати номери відповідних ребер.\par
%answer:1 3 4 5 6
\end {large}}
%end
{\smallskip \begin {small} Затверджено на засіданні кафедри теоретичної кібернетики, протокол \No 12 від   19.05.2020 р.\par\smallskip\smallskip
{\parbox {6.5 cm}{Зав. кафедри \dotfill Ю.В.~Крак}}\hspace {0.7cm}{\hbox to 6.5 cm{ Екзаменатор \dotfill Т.О.~Карнаух}} \end{small}}
%begin
{{\begin {small}\par
{ \center  Київський національний університет імені Тараса Шевченка \par}
{\parbox {5 cm}{Освітня програма \hfill}{\parbox {7 cm}{Прикладна математика, Системний аналіз \hfill}}\parbox {6 cm}{\hfill Семестр 2}}\par
{\parbox {5 cm} {Навчальний предмет \hfill}\parbox {7 cm}{Програмування \hfill}\parbox {6cm}{\hfill}\par}\vspace{-7pt} \end{small}}{\center Екзаменаційний білет \No \textbf 
{ 434 }\par}}
{
\begin {large}
1. Що буде виведено за виконання?
code
try:
    res = 4//4//5*5
    if isinstance(res, bool):
        print(f'bool;{res}')
    elif isinstance(res, int):
        print(f'int;{res}')
    elif isinstance(res, float):
        print(f'float;{res:.2f}')
    else:
        print('???')
except:
    print('ERR')
code

2. Що буде виведено за виконання?
code
a = 1
b = 9
c = 7

def h(a):
    a *= 4
    b = 2
    c = 1
    return a + b + c

a = 5
b = 0
c = 8

try:
    print(h(a), a, b, c, sep=';')
except:
    print('ERR', a, b, c, sep=';')
code

3. Що буде виведено за виконання?
code
try:
    j = 6
    while j>=1:
        j += -2
    print(j, end=';')
except:
    print('ERR')
code

4. Що буде виведено за виконання?
code
try:
    for b in range(-3, -1, -3):
        if b<6:
            break
        print(b, end=';')
    else:
        print(b,end=';')
    print(b)
except:
    print('ERR')
code

5. Що буде виведено за виконання?\par (Дійсні значення, якщо наявні, в цьому завданні будуть виводитись у форматі з фіксованою крапкою та, можливо, з доволі великою кількістю десяткових розрядів. У відповіді для дійсних значень у ЦЬОМУ завданні достатньо вказати один десятковий розряд після крапки, відкинувши решту розрядів без округлення. Наприклад, замість 13.66666666666667 достатньо вказати 13.6, замість 25.23 достатньо вказати 25.2)
code
try:
    print(6, end=';')
    print(1j>=1j, end=';')
    print(8, end=';')
except ValueError:
    print(4, end=';')
except TypeError:
    print(0, end=';')
except:
    print('ERR', end=';')
else:
    print(9, end=';')
finally:
    print(5)
code

6. Що буде виведено за виконання?
code
def f(n):
    if n>2:
        f(n//2)
    else:
        print(n%2, end='')

f(16)
code

7. Що буде виведено за виконання?
code
class A:
    b = 1

    def _f1(self): print(2, end=';')
    def f2(self): print(3, end=';')

    def main(self):
        try: print(self.b, end=';')
        except: print('ERR', end=';')

        try: self._f1()
        except: print('ERR', end=';')

        try: self.f2()
        except: print('ERR', end=';')


class B(A):
    b = 4

    def _f1(self): print(5, end=';')
    def f2(self): print(6, end=';')

try: B().main()
except: print('ERR', end=';')
code

8. Що буде виведено за виконання?
code
def f(arg, n):
    n = 7
    tmp = arg
    tmp[2:-3] = ()
    return len(tmp)>3

try:
    h = ['0','1','2','3','4']
    n = 8
    f(h, n)
    print(n, end=';')
    for el in h:
        print(el, end=';')
except:
    print('ERR')
code

9. Що буде виведено за виконання?
code
class A:
    b = 1
    c = 2

    def __init__(self):
        c = 3

class B(A):
    b = 4

    def __init__(self):
        super().__init__()
        self.a = 7

obj = B()
obj.a = obj.a + 10
B.a = 13

try: print(obj.a, end= ';')
except: print('ERR', end=';')

try: print(A.a, end= ';')
except: print('ERR', end=';')

try: print(B.a, end= ';')
except: print('ERR', end=';')
code

10. 
Для графа на рисунку з вершини 1 запущено алгоритм пошуку в глибину, що будує кістякове дерево. Списки суміжності вершин переглядаються алгоритмом за зростанням міток вершин.\par
\begin{center}
	\adjustimage{max size={0.9\linewidth}{0.9\paperheight}}
	{../10/Traversals/246.png}
\end{center}
\par Які з наведених ребер входять до складу отриманого кістякового дерева?\par
1) (3, 4)\par
2) (2, 6)\par
3) (3, 5)\par
4) (1, 6)\par
5) (2, 5)\par
6) (1, 4)\par
{\vskip 15pt} У формі вибрати номери відповідних ребер.\par
%answer:2 5
\end {large}}
%end
{\smallskip \begin {small} Затверджено на засіданні кафедри теоретичної кібернетики, протокол \No 12 від   19.05.2020 р.\par\smallskip\smallskip
{\parbox {6.5 cm}{Зав. кафедри \dotfill Ю.В.~Крак}}\hspace {0.7cm}{\hbox to 6.5 cm{ Екзаменатор \dotfill Т.О.~Карнаух}} \end{small}}
%begin
{{\begin {small}\par
{ \center  Київський національний університет імені Тараса Шевченка \par}
{\parbox {5 cm}{Освітня програма \hfill}{\parbox {7 cm}{Прикладна математика, Системний аналіз \hfill}}\parbox {6 cm}{\hfill Семестр 2}}\par
{\parbox {5 cm} {Навчальний предмет \hfill}\parbox {7 cm}{Програмування \hfill}\parbox {6cm}{\hfill}\par}\vspace{-7pt} \end{small}}{\center Екзаменаційний білет \No \textbf 
{ 435 }\par}}
{
\begin {large}
1. Що буде виведено за виконання?
code
try:
    res = -1**1.0-True
    if isinstance(res, bool):
        print(f'bool;{res}')
    elif isinstance(res, int):
        print(f'int;{res}')
    elif isinstance(res, float):
        print(f'float;{res:.2f}')
    else:
        print('???')
except:
    print('ERR')
code

2. Що буде виведено за виконання?
code
a = 6
b = 9
c = 8

def g(b):
    global c
    a = 5
    b = 3
    c = 1
    return a + b + c

a = 0
b = 0
c = 8

try:
    print(g(a), a, b, c, sep=';')
except:
    print('ERR', a, b, c, sep=';')
code

3. Що буде виведено за виконання?
code
try:
    k = 8
    while k<13:
        k += 2
    print(k, end=';')
except:
    print('ERR')
code

4. Що буде виведено за виконання?
code
try:
    for a in range(8, 5, 3):
        if a>5:
            continue
        print(a, end=';');a = -6
        if a>3:
            break
    else:
        print(a,end=';')
    print(a)
except:
    print('ERR')
code

5. Що буде виведено за виконання?\par (Дійсні значення, якщо наявні, в цьому завданні будуть виводитись у форматі з фіксованою крапкою та, можливо, з доволі великою кількістю десяткових розрядів. У відповіді для дійсних значень у ЦЬОМУ завданні достатньо вказати один десятковий розряд після крапки, відкинувши решту розрядів без округлення. Наприклад, замість 13.66666666666667 достатньо вказати 13.6, замість 25.23 достатньо вказати 25.2)
code
try:
    print(2, end=';')
    print(3%1, end=';')
    print(7, end=';')
except ZeroDivisionError:
    print(0, end=';')
except Exception:
    print(4, end=';')
except:
    print('ERR', end=';')
else:
    print(6, end=';')
finally:
    print(3)
code

6. Що буде виведено за виконання?
code
def f(n):
    print(n%10,end='')
    if n>9:
        f(n//100)
 
f(123456)
code

7. Що буде виведено за виконання?
code
class A:
    a = 1

    def __h1(self): print(2, end=';')
    def _h2(self): print(3, end=';')

    def main(self):
        try: print(self.a, end=';')
        except: print('ERR', end=';')

        try: self.__h1()
        except: print('ERR', end=';')

        try: self._h2()
        except: print('ERR', end=';')


class B(A):
    a = 4

    def __h1(self): print(5, end=';')
    def _h2(self): print(6, end=';')

try: B().main()
except: print('ERR', end=';')
code

8. Що буде виведено за виконання?
code
def f(arg, n):
    n = 9
    tmp = arg
    tmp[1:-4] = 'bd'
    return len(tmp)>3

try:
    b = [10,11,12,13,14]
    n = 7
    f(b, n)
    print(n, end=';')
    for el in b:
        print(el, end=';')
except:
    print('ERR')
code

9. Що буде виведено за виконання?
code
class A:
    a = 1
    b = 2

    def __init__(self):
        a = 3

class B(A):
    a = 4

    def __init__(self):
        super().__init__()
        self.c = 7

obj = B()
obj.c = obj.b + 10
B.c = 13

try: print(obj.c, end= ';')
except: print('ERR', end=';')

try: print(A.c, end= ';')
except: print('ERR', end=';')

try: print(B.c, end= ';')
except: print('ERR', end=';')
code

10. 
Для графа на рисунку з вершини 4 запущено алгоритм пошуку в глибину, що будує кістякове дерево. Списки суміжності вершин переглядаються алгоритмом за зростанням міток вершин.\par
\begin{center}
	\adjustimage{max size={0.9\linewidth}{0.9\paperheight}}
	{../10/Traversals/238.png}
\end{center}
\par Які з наведених ребер входять до складу отриманого кістякового дерева?\par
1) (0, 6)\par
2) (0, 3)\par
3) (4, 5)\par
4) (0, 4)\par
5) (3, 4)\par
6) (2, 6)\par
{\vskip 15pt} У формі вибрати номери відповідних ребер.\par
%answer:2 4 6
\end {large}}
%end
{\smallskip \begin {small} Затверджено на засіданні кафедри теоретичної кібернетики, протокол \No 12 від   19.05.2020 р.\par\smallskip\smallskip
{\parbox {6.5 cm}{Зав. кафедри \dotfill Ю.В.~Крак}}\hspace {0.7cm}{\hbox to 6.5 cm{ Екзаменатор \dotfill Т.О.~Карнаух}} \end{small}}
%begin
{{\begin {small}\par
{ \center  Київський національний університет імені Тараса Шевченка \par}
{\parbox {5 cm}{Освітня програма \hfill}{\parbox {7 cm}{Прикладна математика, Системний аналіз \hfill}}\parbox {6 cm}{\hfill Семестр 2}}\par
{\parbox {5 cm} {Навчальний предмет \hfill}\parbox {7 cm}{Програмування \hfill}\parbox {6cm}{\hfill}\par}\vspace{-7pt} \end{small}}{\center Екзаменаційний білет \No \textbf 
{ 436 }\par}}
{
\begin {large}
1. Що буде виведено за виконання?
code
try:
    res = "3.0"!=5
    if isinstance(res, bool):
        print(f'bool;{res}')
    elif isinstance(res, int):
        print(f'int;{res}')
    elif isinstance(res, float):
        print(f'float;{res:.2f}')
    else:
        print('???')
except:
    print('ERR')
code

2. Що буде виведено за виконання?
code
a = 6
b = 2
c = 7

def f(a):
    a -= 5
    b = 4
    c = 3
    return a + b + c

a = 9
b = 4
c = 3

try:
    print(f(a), a, b, c, sep=';')
except:
    print('ERR', a, b, c, sep=';')
code

3. Що буде виведено за виконання?
code
try:
    n = 4
    while n<=11:
        n += 2
    print(n, end=';')
except:
    print('ERR')
code

4. Що буде виведено за виконання?
code
try:
    for c in range(-6, 1, -2):
        if c>=7:
            continue
        print(c, end=';')
    else:
        print(c,end=';')
    print(c)
except:
    print('ERR')
code

5. Що буде виведено за виконання?\par (Дійсні значення, якщо наявні, в цьому завданні будуть виводитись у форматі з фіксованою крапкою та, можливо, з доволі великою кількістю десяткових розрядів. У відповіді для дійсних значень у ЦЬОМУ завданні достатньо вказати один десятковий розряд після крапки, відкинувши решту розрядів без округлення. Наприклад, замість 13.66666666666667 достатньо вказати 13.6, замість 25.23 достатньо вказати 25.2)
code
try:
    print(8, end=';')
    print(2%2, end=';')
    print(1, end=';')
except Exception:
    print(5, end=';')
except TypeError:
    print(9, end=';')
except:
    print('ERR', end=';')
else:
    print(5, end=';')
finally:
    print(3)
code

6. Що буде виведено за виконання?
code
def f(n):
    if n>9:
        return f(n//100)+1
    else:
        return 1

print(f(123456))
code

7. Що буде виведено за виконання?
code
class A:
    d = 1

    def _g1(self): print(2, end=';')
    def _g2(self): print(3, end=';')

    def main(self):
        try: print(self.d, end=';')
        except: print('ERR', end=';')

        try: self._g1()
        except: print('ERR', end=';')

        try: self._g2()
        except: print('ERR', end=';')


class B(A):
    d = 4

    def _g1(self): print(5, end=';')
    def _g2(self): print(6, end=';')

try: B().main()
except: print('ERR', end=';')
code

8. Що буде виведено за виконання?
code
def f(arg, n):
    n = 0
    tmp = arg
    del tmp[:2]
    return len(tmp)>3

try:
    e = [10,11,12,13,14]
    n = 8
    f(e, n)
    print(n, end=';')
    for el in e:
        print(el, end=';')
except:
    print('ERR')
code

9. Що буде виведено за виконання?
code
class A:
    b = 1
    a = 2

    def __init__(self):
        a = 3

class B(A):
    a = 4

    def __init__(self):
        super().__init__()
        self.c = 7

obj = B()
obj.a = obj.c + 10
B.a = 13

try: print(obj.a, end= ';')
except: print('ERR', end=';')

try: print(A.a, end= ';')
except: print('ERR', end=';')

try: print(B.a, end= ';')
except: print('ERR', end=';')
code

10. 
Для графа на рисунку з вершини 1 запущено алгоритм пошуку в глибину, що будує кістякове дерево. Списки суміжності вершин переглядаються алгоритмом за зростанням міток вершин.\par
\begin{center}
	\adjustimage{max size={0.9\linewidth}{0.9\paperheight}}
	{../10/Traversals/136.png}
\end{center}
\par Які з наведених ребер входять до складу отриманого кістякового дерева?\par
1) (1, 4)\par
2) (0, 5)\par
3) (0, 6)\par
4) (2, 5)\par
5) (1, 3)\par
6) (0, 4)\par
{\vskip 15pt} У формі вибрати номери відповідних ребер.\par
%answer:2 4 6
\end {large}}
%end
{\smallskip \begin {small} Затверджено на засіданні кафедри теоретичної кібернетики, протокол \No 12 від   19.05.2020 р.\par\smallskip\smallskip
{\parbox {6.5 cm}{Зав. кафедри \dotfill Ю.В.~Крак}}\hspace {0.7cm}{\hbox to 6.5 cm{ Екзаменатор \dotfill Т.О.~Карнаух}} \end{small}}
%begin
{{\begin {small}\par
{ \center  Київський національний університет імені Тараса Шевченка \par}
{\parbox {5 cm}{Освітня програма \hfill}{\parbox {7 cm}{Прикладна математика, Системний аналіз \hfill}}\parbox {6 cm}{\hfill Семестр 2}}\par
{\parbox {5 cm} {Навчальний предмет \hfill}\parbox {7 cm}{Програмування \hfill}\parbox {6cm}{\hfill}\par}\vspace{-7pt} \end{small}}{\center Екзаменаційний білет \No \textbf 
{ 437 }\par}}
{
\begin {large}
1. Що буде виведено за виконання?
code
def f(n):
    print('f', end='')
    return n>-3

try:
    print(f(7) and f(-5))
except:
    print('ERR')
code

2. Що буде виведено за виконання?
code
a = 4
b = 9
c = 5

def h(a):
    a = 3
    b = 2
    c = 5
    return a + b + c

a = 2
b = 6
c = 3

try:
    print(h(a), a, b, c, sep=';')
except:
    print('ERR', a, b, c, sep=';')
code

3. Що буде виведено за виконання?
code
try:
    t = 7
    while t<6:
        t += 1
    print(t, end=';')
except:
    print('ERR')
code

4. Що буде виведено за виконання?
code
try:
    for e in range(-4, -9, 2):
        if e>=4:
            continue
        print(e, end=';');e = -4
    else:
        print(e,end=';')
    print(e)
except:
    print('ERR')
code

5. Що буде виведено за виконання?\par (Дійсні значення, якщо наявні, в цьому завданні будуть виводитись у форматі з фіксованою крапкою та, можливо, з доволі великою кількістю десяткових розрядів. У відповіді для дійсних значень у ЦЬОМУ завданні достатньо вказати один десятковий розряд після крапки, відкинувши решту розрядів без округлення. Наприклад, замість 13.66666666666667 достатньо вказати 13.6, замість 25.23 достатньо вказати 25.2)
code
try:
    print(2, end=';')
    print(0/0.0, end=';')
    print(7, end=';')
except KeyboardInterrupt:
    print(4, end=';')
except ValueError:
    print(6, end=';')
except:
    print('ERR', end=';')
else:
    print(1, end=';')
finally:
    print(8)
code

6. Що буде виведено за виконання?
code
def f(s):
    for el in s:
        el += 1
        return s
 
s = [1, 2, 3]
s = f(s)
print(s)
code

7. Що буде виведено за виконання?
code
class A:
    a = 1

    def __h1(self): print(2, end=';')
    def _h2(self): print(3, end=';')

    def main(self):
        try: print(self.a, end=';')
        except: print('ERR', end=';')

        try: self.__h1()
        except: print('ERR', end=';')

        try: self._h2()
        except: print('ERR', end=';')


class B(A):
    a = 4

    def __h1(self): print(5, end=';')
    def _h2(self): print(6, end=';')

try: B().main()
except: print('ERR', end=';')
code

8. Що буде виведено за виконання?
code
def f(arg, n):
    arg[59] = 4
    n += 3

try:
    n = 5
    t = {68:5, 59:7, 37:7, 37:4}
    f(t, n)
    print(n, end=';')
    for x in t.values():
        print(x, sep=';', end=';')
except:
    print('ERR')
code

9. Що буде виведено за виконання?
code
class A:
    c = 1
    b = 2

    def __init__(self):
        c = 3

class B(A):
    c = 4

    def __init__(self):
        super().__init__()
        self.a = 7

obj = B()
obj.b = obj.a + 10
B.b = 13

try: print(obj.b, end= ';')
except: print('ERR', end=';')

try: print(A.b, end= ';')
except: print('ERR', end=';')

try: print(B.b, end= ';')
except: print('ERR', end=';')
code

10. 
Для графа на рисунку з вершини 1 запущено алгоритм пошуку в глибину, що будує кістякове дерево. Списки суміжності вершин переглядаються алгоритмом за зростанням міток вершин.\par
\begin{center}
	\adjustimage{max size={0.9\linewidth}{0.9\paperheight}}
	{../10/Traversals/69.png}
\end{center}
\par Які з наведених ребер входять до складу отриманого кістякового дерева?\par
1) (1, 6)\par
2) (3, 4)\par
3) (3, 6)\par
4) (0, 2)\par
5) (2, 6)\par
6) (0, 1)\par
{\vskip 15pt} У формі вибрати номери відповідних ребер.\par
%answer:2 3 4 5 6
\end {large}}
%end
{\smallskip \begin {small} Затверджено на засіданні кафедри теоретичної кібернетики, протокол \No 12 від   19.05.2020 р.\par\smallskip\smallskip
{\parbox {6.5 cm}{Зав. кафедри \dotfill Ю.В.~Крак}}\hspace {0.7cm}{\hbox to 6.5 cm{ Екзаменатор \dotfill Т.О.~Карнаух}} \end{small}}
%begin
{{\begin {small}\par
{ \center  Київський національний університет імені Тараса Шевченка \par}
{\parbox {5 cm}{Освітня програма \hfill}{\parbox {7 cm}{Прикладна математика, Системний аналіз \hfill}}\parbox {6 cm}{\hfill Семестр 2}}\par
{\parbox {5 cm} {Навчальний предмет \hfill}\parbox {7 cm}{Програмування \hfill}\parbox {6cm}{\hfill}\par}\vspace{-7pt} \end{small}}{\center Екзаменаційний білет \No \textbf 
{ 438 }\par}}
{
\begin {large}
1. Що буде виведено за виконання?
code
try:
    res = "1j"==False
    if isinstance(res, bool):
        print(f'bool;{res}')
    elif isinstance(res, int):
        print(f'int;{res}')
    elif isinstance(res, float):
        print(f'float;{res:.2f}')
    else:
        print('???')
except:
    print('ERR')
code

2. Що буде виведено за виконання?
code
a = 7
b = 1
c = 3

def f(a):
    global c
    a = 4
    b = 4
    c = 5
    return a + b + c

a = 1
b = 2
c = 8

try:
    print(f(b), a, b, c, sep=';')
except:
    print('ERR', a, b, c, sep=';')
code

3. Що буде виведено за виконання?
code
try:
    i = 14
    while i>7:
        i += -3
    print(i, end=';')
except:
    print('ERR')
code

4. Що буде виведено за виконання?
code
try:
    for d in range(9, 14, 3):
        if d>8:
            break
        print(d, end=';')
    else:
        print('end',end=';')
    print(d)
except:
    print('ERR')
code

5. Що буде виведено за виконання?\par (Дійсні значення, якщо наявні, в цьому завданні будуть виводитись у форматі з фіксованою крапкою та, можливо, з доволі великою кількістю десяткових розрядів. У відповіді для дійсних значень у ЦЬОМУ завданні достатньо вказати один десятковий розряд після крапки, відкинувши решту розрядів без округлення. Наприклад, замість 13.66666666666667 достатньо вказати 13.6, замість 25.23 достатньо вказати 25.2)
code
try:
    print(9, end=';')
    print(int('5'), end=';')
    print(0, end=';')
except ZeroDivisionError:
    print(1, end=';')
except KeyboardInterrupt:
    print(7, end=';')
except:
    print('ERR', end=';')
else:
    print(2, end=';')
finally:
    print(6)
code

6. Що буде виведено за виконання?
code
def f(n):
    print(n%10, end='')
    if n>9:
        f(n//10)
 
f(123456)
code

7. Що буде виведено за виконання?
code
class A:
    c = 1

    def f1(self): print(2, end=';')
    def _f2(self): print(3, end=';')

    def main(self):
        try: print(self.c, end=';')
        except: print('ERR', end=';')

        try: self.f1()
        except: print('ERR', end=';')

        try: self._f2()
        except: print('ERR', end=';')


class B(A):
    c = 4

    def f1(self): print(5, end=';')
    def _f2(self): print(6, end=';')

try: B().main()
except: print('ERR', end=';')
code

8. Що буде виведено за виконання?
code
def f(arg, n):
    arg[55] = 1
    n = 7
    arg = tuple(arg.items())

try:
    n = 4
    t = {55:5, 13:4, 22:1, 22:6}
    f(t, n)
    print(n, end=';')
    for x, y in t.items():
        print(x, y, sep=';', end=';')
except:
    print('ERR')
code

9. Що буде виведено за виконання?
code
class A:
    a = 1
    c = 2

    def __init__(self):
        c = 3

class B(A):
    c = 4

    def __init__(self):
        super().__init__()
        self.b = 7

obj = B()
obj.c = obj.b + 10
B.c = 13

try: print(obj.c, end= ';')
except: print('ERR', end=';')

try: print(A.c, end= ';')
except: print('ERR', end=';')

try: print(B.c, end= ';')
except: print('ERR', end=';')
code

10. 
Для графа на рисунку з вершини 0 запущено алгоритм пошуку в глибину, що будує кістякове дерево. Списки суміжності вершин переглядаються алгоритмом за зростанням міток вершин.\par
\begin{center}
	\adjustimage{max size={0.9\linewidth}{0.9\paperheight}}
	{../10/Traversals/273.png}
\end{center}
\par Які з наведених ребер входять до складу отриманого кістякового дерева?\par
1) (1, 3)\par
2) (2, 3)\par
3) (1, 6)\par
4) (2, 6)\par
5) (1, 2)\par
6) (2, 5)\par
{\vskip 15pt} У формі вибрати номери відповідних ребер.\par
%answer:1 5
\end {large}}
%end
{\smallskip \begin {small} Затверджено на засіданні кафедри теоретичної кібернетики, протокол \No 12 від   19.05.2020 р.\par\smallskip\smallskip
{\parbox {6.5 cm}{Зав. кафедри \dotfill Ю.В.~Крак}}\hspace {0.7cm}{\hbox to 6.5 cm{ Екзаменатор \dotfill Т.О.~Карнаух}} \end{small}}
%begin
{{\begin {small}\par
{ \center  Київський національний університет імені Тараса Шевченка \par}
{\parbox {5 cm}{Освітня програма \hfill}{\parbox {7 cm}{Прикладна математика, Системний аналіз \hfill}}\parbox {6 cm}{\hfill Семестр 2}}\par
{\parbox {5 cm} {Навчальний предмет \hfill}\parbox {7 cm}{Програмування \hfill}\parbox {6cm}{\hfill}\par}\vspace{-7pt} \end{small}}{\center Екзаменаційний білет \No \textbf 
{ 439 }\par}}
{
\begin {large}
1. Що буде виведено за виконання?
code
try:
    res = 9/6*4*4
    if isinstance(res, bool):
        print(f'bool;{res}')
    elif isinstance(res, int):
        print(f'int;{res}')
    elif isinstance(res, float):
        print(f'float;{res:.2f}')
    else:
        print('???')
except:
    print('ERR')
code

2. Що буде виведено за виконання?
code
a = 6
b = 8
c = 0

def h(b):
    global c
    a *= 5
    b = 1
    c = 2
    return a + b + c

a = 2
b = 1
c = 5

try:
    print(h(b), a, b, c, sep=';')
except:
    print('ERR', a, b, c, sep=';')
code

3. Що буде виведено за виконання?
code
try:
    j = 7
    while j<12:
        j += 1
    print(j, end=';')
except:
    print('ERR')
code

4. Що буде виведено за виконання?
code
try:
    for b in range(4, -4, -3):
        if b<=7:
            break
        print(b, end=';')
    else:
        print('end',end=';')
    print(b)
except:
    print('ERR')
code

5. Що буде виведено за виконання?\par (Дійсні значення, якщо наявні, в цьому завданні будуть виводитись у форматі з фіксованою крапкою та, можливо, з доволі великою кількістю десяткових розрядів. У відповіді для дійсних значень у ЦЬОМУ завданні достатньо вказати один десятковий розряд після крапки, відкинувши решту розрядів без округлення. Наприклад, замість 13.66666666666667 достатньо вказати 13.6, замість 25.23 достатньо вказати 25.2)
code
try:
    print(4, end=';')
    print(int('a9'), end=';')
    print(8, end=';')
except ZeroDivisionError:
    print(3, end=';')
except TypeError:
    print(0, end=';')
except:
    print('ERR', end=';')
else:
    print(5, end=';')
finally:
    print(3)
code

6. Що буде виведено за виконання?
code
def f(s1):
    s = s1[:]
    for i in range(len(s)):
        s[i] += 1
 
s = [1, 2, 3]
f(s)
print(s)
code

7. Що буде виведено за виконання?
code
class A:
    b = 1

    def _f1(self): print(2, end=';')
    def __f2(self): print(3, end=';')

    def main(self):
        try: print(self.b, end=';')
        except: print('ERR', end=';')

        try: self._f1()
        except: print('ERR', end=';')

        try: self.__f2()
        except: print('ERR', end=';')


class B(A):
    b = 4

    def _f1(self): print(5, end=';')
    def __f2(self): print(6, end=';')

try: B().main()
except: print('ERR', end=';')
code

8. Що буде виведено за виконання?
code
def f(arg, n):
    n = 6
    tmp = arg
    tmp[0:3] = 'yz'
    return len(tmp)>3

try:
    a = ['0','1','2','3','4']
    n = 3
    f(a, n)
    print(n, end=';')
    for el in a:
        print(el, end=';')
except:
    print('ERR')
code

9. Що буде виведено за виконання?
code
class A:
    c = 1
    a = 2

    def __init__(self):
        a = 3

class B(A):
    c = 4

    def __init__(self):
        super().__init__()
        self.b = 7

obj = B()
obj.a = obj.c + 10
B.a = 13

try: print(obj.a, end= ';')
except: print('ERR', end=';')

try: print(A.a, end= ';')
except: print('ERR', end=';')

try: print(B.a, end= ';')
except: print('ERR', end=';')
code

10. 
Для графа на рисунку з вершини 5 запущено алгоритм пошуку в глибину, що будує кістякове дерево. Списки суміжності вершин переглядаються алгоритмом за зростанням міток вершин.\par
\begin{center}
	\adjustimage{max size={0.9\linewidth}{0.9\paperheight}}
	{../10/Traversals/38.png}
\end{center}
\par Які з наведених ребер входять до складу отриманого кістякового дерева?\par
1) (1, 4)\par
2) (3, 6)\par
3) (0, 1)\par
4) (5, 6)\par
5) (1, 6)\par
6) (0, 2)\par
{\vskip 15pt} У формі вибрати номери відповідних ребер.\par
%answer:1 3 4
\end {large}}
%end
{\smallskip \begin {small} Затверджено на засіданні кафедри теоретичної кібернетики, протокол \No 12 від   19.05.2020 р.\par\smallskip\smallskip
{\parbox {6.5 cm}{Зав. кафедри \dotfill Ю.В.~Крак}}\hspace {0.7cm}{\hbox to 6.5 cm{ Екзаменатор \dotfill Т.О.~Карнаух}} \end{small}}
%begin
{{\begin {small}\par
{ \center  Київський національний університет імені Тараса Шевченка \par}
{\parbox {5 cm}{Освітня програма \hfill}{\parbox {7 cm}{Прикладна математика, Системний аналіз \hfill}}\parbox {6 cm}{\hfill Семестр 2}}\par
{\parbox {5 cm} {Навчальний предмет \hfill}\parbox {7 cm}{Програмування \hfill}\parbox {6cm}{\hfill}\par}\vspace{-7pt} \end{small}}{\center Екзаменаційний білет \No \textbf 
{ 440 }\par}}
{
\begin {large}
1. Що буде виведено за виконання?
code
def f(n):
    print('f', end='')
    return n==1

try:
    print(f(9) or f(-7))
except:
    print('ERR')
code

2. Що буде виведено за виконання?
code
a = 6
b = 7
c = 5

def h(b):
    a = 1
    b = 3
    c = 2
    return a + b + c

a = 4
b = 9
c = 0

try:
    print(h(b), a, b, c, sep=';')
except:
    print('ERR', a, b, c, sep=';')
code

3. Що буде виведено за виконання?
code
try:
    t = 9
    while t>3:
        t += -2
    print(t, end=';')
except:
    print('ERR')
code

4. Що буде виведено за виконання?
code
try:
    for f in range(-12, 7, -1):
        if f<4:
            continue
        print(f, end=';')
    else:
        print(f,end=';')
    print(f)
except:
    print('ERR')
code

5. Що буде виведено за виконання?\par (Дійсні значення, якщо наявні, в цьому завданні будуть виводитись у форматі з фіксованою крапкою та, можливо, з доволі великою кількістю десяткових розрядів. У відповіді для дійсних значень у ЦЬОМУ завданні достатньо вказати один десятковий розряд після крапки, відкинувши решту розрядів без округлення. Наприклад, замість 13.66666666666667 достатньо вказати 13.6, замість 25.23 достатньо вказати 25.2)
code
try:
    print(4, end=';')
    print(6j<3j, end=';')
    print(2, end=';')
except ValueError:
    print(7, end=';')
except Exception:
    print(1, end=';')
except:
    print('ERR', end=';')
else:
    print(8, end=';')
finally:
    print(9)
code

6. Що буде виведено за виконання?
code
def f(n):
    if n>9:
        f(n//10)
    print(n%10, end='')
 
f(123456)
code

7. Що буде виведено за виконання?
code
class A:
    d = 1

    def _g1(self): print(2, end=';')
    def g2(self): print(3, end=';')

    def main(self):
        try: print(self.d, end=';')
        except: print('ERR', end=';')

        try: self._g1()
        except: print('ERR', end=';')

        try: self.g2()
        except: print('ERR', end=';')


class B(A):
    d = 4

    def _g1(self): print(5, end=';')
    def g2(self): print(6, end=';')

try: B().main()
except: print('ERR', end=';')
code

8. Що буде виведено за виконання?
code
def f(arg, n):
    n = 4
    tmp = arg
    tmp[4:] = [1,0]
    return len(tmp)>3

try:
    d = [10,11,12,13,14]
    n = 5
    f(d, n)
    print(n, end=';')
    for el in d:
        print(el, end=';')
except:
    print('ERR')
code

9. Що буде виведено за виконання?
code
class A:
    a = 1
    c = 2

    def __init__(self):
        a = 3

class B(A):
    c = 4

    def __init__(self):
        super().__init__()
        self.b = 7

obj = B()
obj.b = obj.c + 10
B.b = 13

try: print(obj.b, end= ';')
except: print('ERR', end=';')

try: print(A.b, end= ';')
except: print('ERR', end=';')

try: print(B.b, end= ';')
except: print('ERR', end=';')
code

10. 
Для графа на рисунку з вершини 3 запущено алгоритм пошуку в глибину, що будує кістякове дерево. Списки суміжності вершин переглядаються алгоритмом за зростанням міток вершин.\par
\begin{center}
	\adjustimage{max size={0.9\linewidth}{0.9\paperheight}}
	{../10/Traversals/17.png}
\end{center}
\par Які з наведених ребер входять до складу отриманого кістякового дерева?\par
1) (0, 4)\par
2) (1, 6)\par
3) (0, 1)\par
4) (2, 4)\par
5) (0, 6)\par
6) (1, 4)\par
{\vskip 15pt} У формі вибрати номери відповідних ребер.\par
%answer:1 3 4 5
\end {large}}
%end
{\smallskip \begin {small} Затверджено на засіданні кафедри теоретичної кібернетики, протокол \No 12 від   19.05.2020 р.\par\smallskip\smallskip
{\parbox {6.5 cm}{Зав. кафедри \dotfill Ю.В.~Крак}}\hspace {0.7cm}{\hbox to 6.5 cm{ Екзаменатор \dotfill Т.О.~Карнаух}} \end{small}}
%begin
{{\begin {small}\par
{ \center  Київський національний університет імені Тараса Шевченка \par}
{\parbox {5 cm}{Освітня програма \hfill}{\parbox {7 cm}{Прикладна математика, Системний аналіз \hfill}}\parbox {6 cm}{\hfill Семестр 2}}\par
{\parbox {5 cm} {Навчальний предмет \hfill}\parbox {7 cm}{Програмування \hfill}\parbox {6cm}{\hfill}\par}\vspace{-7pt} \end{small}}{\center Екзаменаційний білет \No \textbf 
{ 441 }\par}}
{
\begin {large}
1. Що буде виведено за виконання?
code
try:
    res = 7//9/9*3
    if isinstance(res, bool):
        print(f'bool;{res}')
    elif isinstance(res, int):
        print(f'int;{res}')
    elif isinstance(res, float):
        print(f'float;{res:.2f}')
    else:
        print('???')
except:
    print('ERR')
code

2. Що буде виведено за виконання?
code
a = 3
b = 1
c = 5

def g(b):
    global c
    a = 3
    b = 5
    c = 1
    return a + b + c

a = 8
b = 0
c = 4

try:
    print(g(a), a, b, c, sep=';')
except:
    print('ERR', a, b, c, sep=';')
code

3. Що буде виведено за виконання?
code
try:
    m = 8
    while m<=5:
        m += 2
    print(m, end=';')
except:
    print('ERR')
code

4. Що буде виведено за виконання?
code
try:
    for e in range(0, -10, 2):
        if e>=8:
            continue
        print(e, end=';');e = 4
    else:
        print(e,end=';')
    print(e)
except:
    print('ERR')
code

5. Що буде виведено за виконання?\par (Дійсні значення, якщо наявні, в цьому завданні будуть виводитись у форматі з фіксованою крапкою та, можливо, з доволі великою кількістю десяткових розрядів. У відповіді для дійсних значень у ЦЬОМУ завданні достатньо вказати один десятковий розряд після крапки, відкинувши решту розрядів без округлення. Наприклад, замість 13.66666666666667 достатньо вказати 13.6, замість 25.23 достатньо вказати 25.2)
code
try:
    print(3, end=';')
    print(int('9'), end=';')
    print(8, end=';')
except ValueError:
    print(5, end=';')
except Exception:
    print(6, end=';')
except:
    print('ERR', end=';')
else:
    print(2, end=';')
finally:
    print(1)
code

6. Що буде виведено за виконання?
code
def f(n):
    if n<=9:
        print(n%10, end='')
    else:
        f(n//10)

f(987651)
code

7. Що буде виведено за виконання?
code
class A:
    b = 1

    def _h1(self): print(2, end=';')
    def _h2(self): print(3, end=';')

    def main(self):
        try: print(self.b, end=';')
        except: print('ERR', end=';')

        try: self._h1()
        except: print('ERR', end=';')

        try: self._h2()
        except: print('ERR', end=';')


class B(A):
    b = 4

    def _h1(self): print(5, end=';')
    def _h2(self): print(6, end=';')

try: B().main()
except: print('ERR', end=';')
code

8. Що буде виведено за виконання?
code
def f(arg, n):
    arg[90] = 2
    n *= 2

try:
    n = 4
    d = {88:8, 16:0, 61:7, 88:7}
    f(d, n)
    print(n, end=';')
    for x in d :
        print(x, sep=';', end=';')
except:
    print('ERR')
code

9. Що буде виведено за виконання?
code
class A:
    b = 1
    c = 2

    def __init__(self):
        b = 3

class B(A):
    b = 4

    def __init__(self):
        super().__init__()
        self.a = 7

obj = B()
obj.a = obj.b + 10
B.a = 13

try: print(obj.a, end= ';')
except: print('ERR', end=';')

try: print(A.a, end= ';')
except: print('ERR', end=';')

try: print(B.a, end= ';')
except: print('ERR', end=';')
code

10. 
Для графа на рисунку з вершини 5 запущено алгоритм пошуку в глибину, що будує кістякове дерево. Списки суміжності вершин переглядаються алгоритмом за зростанням міток вершин.\par
\begin{center}
	\adjustimage{max size={0.9\linewidth}{0.9\paperheight}}
	{../10/Traversals/230.png}
\end{center}
\par Які з наведених ребер входять до складу отриманого кістякового дерева?\par
1) (0, 2)\par
2) (1, 6)\par
3) (0, 4)\par
4) (2, 5)\par
5) (5, 6)\par
6) (0, 1)\par
{\vskip 15pt} У формі вибрати номери відповідних ребер.\par
%answer:1 2 3 4 6
\end {large}}
%end
{\smallskip \begin {small} Затверджено на засіданні кафедри теоретичної кібернетики, протокол \No 12 від   19.05.2020 р.\par\smallskip\smallskip
{\parbox {6.5 cm}{Зав. кафедри \dotfill Ю.В.~Крак}}\hspace {0.7cm}{\hbox to 6.5 cm{ Екзаменатор \dotfill Т.О.~Карнаух}} \end{small}}
%begin
{{\begin {small}\par
{ \center  Київський національний університет імені Тараса Шевченка \par}
{\parbox {5 cm}{Освітня програма \hfill}{\parbox {7 cm}{Прикладна математика, Системний аналіз \hfill}}\parbox {6 cm}{\hfill Семестр 2}}\par
{\parbox {5 cm} {Навчальний предмет \hfill}\parbox {7 cm}{Програмування \hfill}\parbox {6cm}{\hfill}\par}\vspace{-7pt} \end{small}}{\center Екзаменаційний білет \No \textbf 
{ 442 }\par}}
{
\begin {large}
1. Що буде виведено за виконання?
code
try:
    res = 9**-0.5*False
    if isinstance(res, bool):
        print(f'bool;{res}')
    elif isinstance(res, int):
        print(f'int;{res}')
    elif isinstance(res, float):
        print(f'float;{res:.2f}')
    else:
        print('???')
except:
    print('ERR')
code

2. Що буде виведено за виконання?
code
a = 7
b = 6
c = 0

def f(b):
    global c
    a += 1
    b = 2
    c = 3
    return a + b + c

a = 4
b = 1
c = 9

try:
    print(f(b), a, b, c, sep=';')
except:
    print('ERR', a, b, c, sep=';')
code

3. Що буде виведено за виконання?
code
try:
    i = 6
    while i<9:
        i += 2
    print(i, end=';')
except:
    print('ERR')
code

4. Що буде виведено за виконання?
code
try:
    for c in range(14, 4, -1):
        if c<5:
            continue
        print(c, end=';')
    else:
        print(c,end=';')
    print(c)
except:
    print('ERR')
code

5. Що буде виведено за виконання?\par (Дійсні значення, якщо наявні, в цьому завданні будуть виводитись у форматі з фіксованою крапкою та, можливо, з доволі великою кількістю десяткових розрядів. У відповіді для дійсних значень у ЦЬОМУ завданні достатньо вказати один десятковий розряд після крапки, відкинувши решту розрядів без округлення. Наприклад, замість 13.66666666666667 достатньо вказати 13.6, замість 25.23 достатньо вказати 25.2)
code
try:
    print(4, end=';')
    print(0//0, end=';')
    print(7, end=';')
except KeyboardInterrupt:
    print(5, end=';')
except ZeroDivisionError:
    print(1, end=';')
except:
    print('ERR', end=';')
else:
    print(0, end=';')
finally:
    print(7)
code

6. Що буде виведено за виконання?
code
def f(n):
    print(n%10, end='')
    if n>9:
        f(n//100)
 
f(987654)
code

7. Що буде виведено за виконання?
code
class A:
    c = 1

    def _f1(self): print(2, end=';')
    def __f2(self): print(3, end=';')

    def main(self):
        try: print(self.c, end=';')
        except: print('ERR', end=';')

        try: self._f1()
        except: print('ERR', end=';')

        try: self.__f2()
        except: print('ERR', end=';')


class B(A):
    c = 4

    def _f1(self): print(5, end=';')
    def __f2(self): print(6, end=';')

try: B().main()
except: print('ERR', end=';')
code

8. Що буде виведено за виконання?
code
def f(arg, n):
    arg[48] = 5
    n -= -2

try:
    n = 6
    s = {24:3, 48:7, 24:8}
    f(s, n)
    print(n, end=';')
    for x in s.items():
        print(*x, sep=';', end=';')
except:
    print('ERR')
code

9. Що буде виведено за виконання?
code
class A:
    b = 1
    a = 2

    def __init__(self):
        a = 3

class B(A):
    a = 4

    def __init__(self):
        super().__init__()
        self.c = 7

obj = B()
obj.b = obj.c + 10
B.b = 13

try: print(obj.b, end= ';')
except: print('ERR', end=';')

try: print(A.b, end= ';')
except: print('ERR', end=';')

try: print(B.b, end= ';')
except: print('ERR', end=';')
code

10. 
Для графа на рисунку з вершини 5 запущено алгоритм пошуку в глибину, що будує кістякове дерево. Списки суміжності вершин переглядаються алгоритмом за зростанням міток вершин.\par
\begin{center}
	\adjustimage{max size={0.9\linewidth}{0.9\paperheight}}
	{../10/Traversals/65.png}
\end{center}
\par Які з наведених ребер входять до складу отриманого кістякового дерева?\par
1) (3, 5)\par
2) (0, 1)\par
3) (2, 6)\par
4) (3, 4)\par
5) (0, 5)\par
6) (0, 2)\par
{\vskip 15pt} У формі вибрати номери відповідних ребер.\par
%answer:1 2 3 4 5
\end {large}}
%end
{\smallskip \begin {small} Затверджено на засіданні кафедри теоретичної кібернетики, протокол \No 12 від   19.05.2020 р.\par\smallskip\smallskip
{\parbox {6.5 cm}{Зав. кафедри \dotfill Ю.В.~Крак}}\hspace {0.7cm}{\hbox to 6.5 cm{ Екзаменатор \dotfill Т.О.~Карнаух}} \end{small}}
%begin
{{\begin {small}\par
{ \center  Київський національний університет імені Тараса Шевченка \par}
{\parbox {5 cm}{Освітня програма \hfill}{\parbox {7 cm}{Прикладна математика, Системний аналіз \hfill}}\parbox {6 cm}{\hfill Семестр 2}}\par
{\parbox {5 cm} {Навчальний предмет \hfill}\parbox {7 cm}{Програмування \hfill}\parbox {6cm}{\hfill}\par}\vspace{-7pt} \end{small}}{\center Екзаменаційний білет \No \textbf 
{ 443 }\par}}
{
\begin {large}
1. Що буде виведено за виконання?
code
try:
    res = -2**-4-2
    if isinstance(res, bool):
        print(f'bool;{res}')
    elif isinstance(res, int):
        print(f'int;{res}')
    elif isinstance(res, float):
        print(f'float;{res:.2f}')
    else:
        print('???')
except:
    print('ERR')
code

2. Що буде виведено за виконання?
code
a = 2
b = 8
c = 5

def g(a):
    a -= 5
    b = 4
    c = 2
    return a + b + c

a = 3
b = 7
c = 9

try:
    print(g(a), a, b, c, sep=';')
except:
    print('ERR', a, b, c, sep=';')
code

3. Що буде виведено за виконання?
code
try:
    k = 11
    while k>=0:
        k += -3
    print(k, end=';')
except:
    print('ERR')
code

4. Що буде виведено за виконання?
code
try:
    for d in range(-7, -8, 3):
        if d>3:
            break
        print(d, end=';');d = 9
        if d<=7:
            break
    else:
        print(d,end=';')
    print(d)
except:
    print('ERR')
code

5. Що буде виведено за виконання?\par (Дійсні значення, якщо наявні, в цьому завданні будуть виводитись у форматі з фіксованою крапкою та, можливо, з доволі великою кількістю десяткових розрядів. У відповіді для дійсних значень у ЦЬОМУ завданні достатньо вказати один десятковий розряд після крапки, відкинувши решту розрядів без округлення. Наприклад, замість 13.66666666666667 достатньо вказати 13.6, замість 25.23 достатньо вказати 25.2)
code
try:
    print(9, end=';')
    print(4/1, end=';')
    print(6, end=';')
except TypeError:
    print(2, end=';')
except ValueError:
    print(3, end=';')
except:
    print('ERR', end=';')
else:
    print(8, end=';')
finally:
    print(9)
code

6. Що буде виведено за виконання?
code
def f(n):
    if n>9:
        f(n//10)
    print(n%10, end='')
 
f(543216)
code

7. Що буде виведено за виконання?
code
class A:
    a = 1

    def __h1(self): print(2, end=';')
    def _h2(self): print(3, end=';')

    def main(self):
        try: print(self.a, end=';')
        except: print('ERR', end=';')

        try: self.__h1()
        except: print('ERR', end=';')

        try: self._h2()
        except: print('ERR', end=';')


class B(A):
    a = 4

    def __h1(self): print(5, end=';')
    def _h2(self): print(6, end=';')

try: B().main()
except: print('ERR', end=';')
code

8. Що буде виведено за виконання?
code
def f(arg, n):
    arg[84] = 2
    n = 8
    arg = set(arg.values())

try:
    n = 3
    d = {84:7, 33:0, 22:5, 33:4}
    f(d, n)
    print(n, end=';')
    for x in d.keys():
        print(x, sep=';', end=';')
except:
    print('ERR')
code

9. Що буде виведено за виконання?
code
class A:
    c = 1
    b = 2

    def __init__(self):
        b = 3

class B(A):
    c = 4

    def __init__(self):
        super().__init__()
        self.a = 7

obj = B()
obj.a = obj.c + 10
B.a = 13

try: print(obj.a, end= ';')
except: print('ERR', end=';')

try: print(A.a, end= ';')
except: print('ERR', end=';')

try: print(B.a, end= ';')
except: print('ERR', end=';')
code

10. 
Для графа на рисунку з вершини 4 запущено алгоритм пошуку в глибину, що будує кістякове дерево. Списки суміжності вершин переглядаються алгоритмом за зростанням міток вершин.\par
\begin{center}
	\adjustimage{max size={0.9\linewidth}{0.9\paperheight}}
	{../10/Traversals/75.png}
\end{center}
\par Які з наведених ребер входять до складу отриманого кістякового дерева?\par
1) (2, 5)\par
2) (0, 3)\par
3) (0, 1)\par
4) (5, 6)\par
5) (1, 4)\par
6) (0, 4)\par
{\vskip 15pt} У формі вибрати номери відповідних ребер.\par
%answer:1 2 3 4 6
\end {large}}
%end
{\smallskip \begin {small} Затверджено на засіданні кафедри теоретичної кібернетики, протокол \No 12 від   19.05.2020 р.\par\smallskip\smallskip
{\parbox {6.5 cm}{Зав. кафедри \dotfill Ю.В.~Крак}}\hspace {0.7cm}{\hbox to 6.5 cm{ Екзаменатор \dotfill Т.О.~Карнаух}} \end{small}}
%begin
{{\begin {small}\par
{ \center  Київський національний університет імені Тараса Шевченка \par}
{\parbox {5 cm}{Освітня програма \hfill}{\parbox {7 cm}{Прикладна математика, Системний аналіз \hfill}}\parbox {6 cm}{\hfill Семестр 2}}\par
{\parbox {5 cm} {Навчальний предмет \hfill}\parbox {7 cm}{Програмування \hfill}\parbox {6cm}{\hfill}\par}\vspace{-7pt} \end{small}}{\center Екзаменаційний білет \No \textbf 
{ 444 }\par}}
{
\begin {large}
1. Що буде виведено за виконання?
code
try:
    res = not False>=7 or 5.0!=5 or 7.0>8
    if isinstance(res, bool):
        print(f'bool;{res}')
    elif isinstance(res, int):
        print(f'int;{res}')
    elif isinstance(res, float):
        print(f'float;{res:.2f}')
    else:
        print('???')
except:
    print('ERR')
code

2. Що буде виведено за виконання?
code
a = 7
b = 6
c = 2

def f(a):
    a = 4
    b = 2
    c = 4
    return a + b + c

a = 9
b = 2
c = 5

try:
    print(f(b), a, b, c, sep=';')
except:
    print('ERR', a, b, c, sep=';')
code

3. Що буде виведено за виконання?
code
try:
    t = 4
    while t<=7:
        t += 1
    print(t, end=';')
except:
    print('ERR')
code

4. Що буде виведено за виконання?
code
try:
    for a in range(-5, -13, -2):
        if a<=6:
            continue
        print(a, end=';')
    else:
        print(a,end=';')
    print(a)
except:
    print('ERR')
code

5. Що буде виведено за виконання?\par (Дійсні значення, якщо наявні, в цьому завданні будуть виводитись у форматі з фіксованою крапкою та, можливо, з доволі великою кількістю десяткових розрядів. У відповіді для дійсних значень у ЦЬОМУ завданні достатньо вказати один десятковий розряд після крапки, відкинувши решту розрядів без округлення. Наприклад, замість 13.66666666666667 достатньо вказати 13.6, замість 25.23 достатньо вказати 25.2)
code
try:
    print(5, end=';')
    print(3j!=5j, end=';')
    print(1, end=';')
except TypeError:
    print(4, end=';')
except ZeroDivisionError:
    print(0, end=';')
except:
    print('ERR', end=';')
else:
    print(8, end=';')
finally:
    print(7)
code

6. Що буде виведено за виконання?
code
def f(n):
    if n>2:
        f(n//2)
    print(n%2, end='')
 
f(16)
code

7. Що буде виведено за виконання?
code
class A:
    d = 1

    def g1(self): print(2, end=';')
    def _g2(self): print(3, end=';')

    def main(self):
        try: print(self.d, end=';')
        except: print('ERR', end=';')

        try: self.g1()
        except: print('ERR', end=';')

        try: self._g2()
        except: print('ERR', end=';')


class B(A):
    d = 4

    def g1(self): print(5, end=';')
    def _g2(self): print(6, end=';')

try: B().main()
except: print('ERR', end=';')
code

8. Що буде виведено за виконання?
code
def f(arg, n):
    arg[15] = 6
    n = 5
    arg = tuple(arg.items())

try:
    n = 1
    s = {15:2, 82:6, 89:8, 82:6}
    f(s, n)
    print(n, end=';')
    for x in s.items():
        print(*x, sep=';', end=';')
except:
    print('ERR')
code

9. Що буде виведено за виконання?
code
class A:
    a = 1
    b = 2

    def __init__(self):
        a = 3

class B(A):
    b = 4

    def __init__(self):
        super().__init__()
        self.c = 7

obj = B()
obj.a = obj.b + 10
B.a = 13

try: print(obj.a, end= ';')
except: print('ERR', end=';')

try: print(A.a, end= ';')
except: print('ERR', end=';')

try: print(B.a, end= ';')
except: print('ERR', end=';')
code

10. 
Для графа на рисунку з вершини 6 запущено алгоритм пошуку в глибину, що будує кістякове дерево. Списки суміжності вершин переглядаються алгоритмом за зростанням міток вершин.\par
\begin{center}
	\adjustimage{max size={0.9\linewidth}{0.9\paperheight}}
	{../10/Traversals/207.png}
\end{center}
\par Які з наведених ребер входять до складу отриманого кістякового дерева?\par
1) (0, 3)\par
2) (5, 6)\par
3) (1, 4)\par
4) (4, 5)\par
5) (0, 1)\par
6) (2, 4)\par
{\vskip 15pt} У формі вибрати номери відповідних ребер.\par
%answer:1 3 4 5 6
\end {large}}
%end
{\smallskip \begin {small} Затверджено на засіданні кафедри теоретичної кібернетики, протокол \No 12 від   19.05.2020 р.\par\smallskip\smallskip
{\parbox {6.5 cm}{Зав. кафедри \dotfill Ю.В.~Крак}}\hspace {0.7cm}{\hbox to 6.5 cm{ Екзаменатор \dotfill Т.О.~Карнаух}} \end{small}}
%begin
{{\begin {small}\par
{ \center  Київський національний університет імені Тараса Шевченка \par}
{\parbox {5 cm}{Освітня програма \hfill}{\parbox {7 cm}{Прикладна математика, Системний аналіз \hfill}}\parbox {6 cm}{\hfill Семестр 2}}\par
{\parbox {5 cm} {Навчальний предмет \hfill}\parbox {7 cm}{Програмування \hfill}\parbox {6cm}{\hfill}\par}\vspace{-7pt} \end{small}}{\center Екзаменаційний білет \No \textbf 
{ 445 }\par}}
{
\begin {large}
1. Що буде виведено за виконання?
code
try:
    res = not 9==3 and 6.0<=True and 6<3.0
    if isinstance(res, bool):
        print(f'bool;{res}')
    elif isinstance(res, int):
        print(f'int;{res}')
    elif isinstance(res, float):
        print(f'float;{res:.2f}')
    else:
        print('???')
except:
    print('ERR')
code

2. Що буде виведено за виконання?
code
a = 1
b = 6
c = 4

def g(a):
    a -= 5
    b = 3
    c = 1
    return a + b + c

a = 0
b = 8
c = 3

try:
    print(g(b), a, b, c, sep=';')
except:
    print('ERR', a, b, c, sep=';')
code

3. Що буде виведено за виконання?
code
try:
    m = 0
    while m<8:
        m += 2
    print(m, end=';')
except:
    print('ERR')
code

4. Що буде виведено за виконання?
code
try:
    for e in range(-9, -1, -3):
        if e<=7:
            continue
        print(e, end=';');e = 10
    else:
        print(e,end=';')
    print(e)
except:
    print('ERR')
code

5. Що буде виведено за виконання?\par (Дійсні значення, якщо наявні, в цьому завданні будуть виводитись у форматі з фіксованою крапкою та, можливо, з доволі великою кількістю десяткових розрядів. У відповіді для дійсних значень у ЦЬОМУ завданні достатньо вказати один десятковий розряд після крапки, відкинувши решту розрядів без округлення. Наприклад, замість 13.66666666666667 достатньо вказати 13.6, замість 25.23 достатньо вказати 25.2)
code
try:
    print(2, end=';')
    print(int('c6'), end=';')
    print(6, end=';')
except Exception:
    print(9, end=';')
except KeyboardInterrupt:
    print(4, end=';')
except:
    print('ERR', end=';')
else:
    print(0, end=';')
finally:
    print(1)
code

6. Що буде виведено за виконання?
code
def f(n):
    if n>9:
        f(n//100)
    else:
        print(n%10,end='')

f(123456)
code

7. Що буде виведено за виконання?
code
class A:
    d = 1

    def _g1(self): print(2, end=';')
    def g2(self): print(3, end=';')

    def main(self):
        try: print(self.d, end=';')
        except: print('ERR', end=';')

        try: self._g1()
        except: print('ERR', end=';')

        try: self.g2()
        except: print('ERR', end=';')


class B(A):
    d = 4

    def _g1(self): print(5, end=';')
    def g2(self): print(6, end=';')

try: B().main()
except: print('ERR', end=';')
code

8. Що буде виведено за виконання?
code
def f(arg, n):
    arg[81] = 6
    n -= -3

try:
    n = 5
    s = {18:2, 57:4, 21:7, 52:1}
    f(s, n)
    print(n, end=';')
    for x, y in s.items():
        print(x, y, sep=';', end=';')
except:
    print('ERR')
code

9. Що буде виведено за виконання?
code
class A:
    b = 1
    c = 2

    def __init__(self):
        c = 3

class B(A):
    c = 4

    def __init__(self):
        super().__init__()
        self.a = 7

obj = B()
obj.a = obj.b + 10
B.a = 13

try: print(obj.a, end= ';')
except: print('ERR', end=';')

try: print(A.a, end= ';')
except: print('ERR', end=';')

try: print(B.a, end= ';')
except: print('ERR', end=';')
code

10. 
Для графа на рисунку з вершини 5 запущено алгоритм пошуку в глибину, що будує кістякове дерево. Списки суміжності вершин переглядаються алгоритмом за зростанням міток вершин.\par
\begin{center}
	\adjustimage{max size={0.9\linewidth}{0.9\paperheight}}
	{../10/Traversals/64.png}
\end{center}
\par Які з наведених ребер входять до складу отриманого кістякового дерева?\par
1) (0, 4)\par
2) (2, 6)\par
3) (4, 6)\par
4) (5, 6)\par
5) (0, 6)\par
6) (1, 6)\par
{\vskip 15pt} У формі вибрати номери відповідних ребер.\par
%answer:5 6
\end {large}}
%end
{\smallskip \begin {small} Затверджено на засіданні кафедри теоретичної кібернетики, протокол \No 12 від   19.05.2020 р.\par\smallskip\smallskip
{\parbox {6.5 cm}{Зав. кафедри \dotfill Ю.В.~Крак}}\hspace {0.7cm}{\hbox to 6.5 cm{ Екзаменатор \dotfill Т.О.~Карнаух}} \end{small}}
%begin
{{\begin {small}\par
{ \center  Київський національний університет імені Тараса Шевченка \par}
{\parbox {5 cm}{Освітня програма \hfill}{\parbox {7 cm}{Прикладна математика, Системний аналіз \hfill}}\parbox {6 cm}{\hfill Семестр 2}}\par
{\parbox {5 cm} {Навчальний предмет \hfill}\parbox {7 cm}{Програмування \hfill}\parbox {6cm}{\hfill}\par}\vspace{-7pt} \end{small}}{\center Екзаменаційний білет \No \textbf 
{ 446 }\par}}
{
\begin {large}
1. Що буде виведено за виконання?
code
try:
    res = 4**-0.5*-1
    if isinstance(res, bool):
        print(f'bool;{res}')
    elif isinstance(res, int):
        print(f'int;{res}')
    elif isinstance(res, float):
        print(f'float;{res:.2f}')
    else:
        print('???')
except:
    print('ERR')
code

2. Що буде виведено за виконання?
code
a = 5
b = 2
c = 4

def h(b):
    global c
    a += 4
    b = 1
    c = 3
    return a + b + c

a = 7
b = 8
c = 3

try:
    print(h(b), a, b, c, sep=';')
except:
    print('ERR', a, b, c, sep=';')
code

3. Що буде виведено за виконання?
code
try:
    m = 6
    while m>1:
        m += -3
    print(m, end=';')
except:
    print('ERR')
code

4. Що буде виведено за виконання?
code
try:
    for f in range(10, 0, -3):
        if f<6:
            continue
        print(f, end=';')
    else:
        print(f,end=';')
    print(f)
except:
    print('ERR')
code

5. Що буде виведено за виконання?\par (Дійсні значення, якщо наявні, в цьому завданні будуть виводитись у форматі з фіксованою крапкою та, можливо, з доволі великою кількістю десяткових розрядів. У відповіді для дійсних значень у ЦЬОМУ завданні достатньо вказати один десятковий розряд після крапки, відкинувши решту розрядів без округлення. Наприклад, замість 13.66666666666667 достатньо вказати 13.6, замість 25.23 достатньо вказати 25.2)
code
try:
    print(8, end=';')
    print(int('3'), end=';')
    print(5, end=';')
except ZeroDivisionError:
    print(7, end=';')
except Exception:
    print(2, end=';')
except:
    print('ERR', end=';')
else:
    print(6, end=';')
finally:
    print(2)
code

6. Що буде виведено за виконання?
code
def f(n):
    if n>9:
        f(n//100)
    print(n%10,end='')
 
f(987654)
code

7. Що буде виведено за виконання?
code
class A:
    b = 1

    def h1(self): print(2, end=';')
    def _h2(self): print(3, end=';')

    def main(self):
        try: print(self.b, end=';')
        except: print('ERR', end=';')

        try: self.h1()
        except: print('ERR', end=';')

        try: self._h2()
        except: print('ERR', end=';')


class B(A):
    b = 4

    def h1(self): print(5, end=';')
    def _h2(self): print(6, end=';')

try: B().main()
except: print('ERR', end=';')
code

8. Що буде виведено за виконання?
code
def f(arg, n):
    arg[11] = 6
    n += 3

try:
    n = 7
    t = {35:1, 26:0, 66:4}
    f(t, n)
    print(n, end=';')
    for x in t.values():
        print(x, sep=';', end=';')
except:
    print('ERR')
code

9. Що буде виведено за виконання?
code
class A:
    b = 1
    a = 2

    def __init__(self):
        b = 3

class B(A):
    b = 4

    def __init__(self):
        super().__init__()
        self.c = 7

obj = B()
obj.c = obj.c + 10
B.c = 13

try: print(obj.c, end= ';')
except: print('ERR', end=';')

try: print(A.c, end= ';')
except: print('ERR', end=';')

try: print(B.c, end= ';')
except: print('ERR', end=';')
code

10. 
Для графа на рисунку з вершини 1 запущено алгоритм пошуку в глибину, що будує кістякове дерево. Списки суміжності вершин переглядаються алгоритмом за зростанням міток вершин.\par
\begin{center}
	\adjustimage{max size={0.9\linewidth}{0.9\paperheight}}
	{../10/Traversals/198.png}
\end{center}
\par Які з наведених ребер входять до складу отриманого кістякового дерева?\par
1) (0, 6)\par
2) (1, 6)\par
3) (3, 5)\par
4) (2, 5)\par
5) (1, 3)\par
6) (1, 5)\par
{\vskip 15pt} У формі вибрати номери відповідних ребер.\par
%answer:1 3 5
\end {large}}
%end
{\smallskip \begin {small} Затверджено на засіданні кафедри теоретичної кібернетики, протокол \No 12 від   19.05.2020 р.\par\smallskip\smallskip
{\parbox {6.5 cm}{Зав. кафедри \dotfill Ю.В.~Крак}}\hspace {0.7cm}{\hbox to 6.5 cm{ Екзаменатор \dotfill Т.О.~Карнаух}} \end{small}}
%begin
{{\begin {small}\par
{ \center  Київський національний університет імені Тараса Шевченка \par}
{\parbox {5 cm}{Освітня програма \hfill}{\parbox {7 cm}{Прикладна математика, Системний аналіз \hfill}}\parbox {6 cm}{\hfill Семестр 2}}\par
{\parbox {5 cm} {Навчальний предмет \hfill}\parbox {7 cm}{Програмування \hfill}\parbox {6cm}{\hfill}\par}\vspace{-7pt} \end{small}}{\center Екзаменаційний білет \No \textbf 
{ 447 }\par}}
{
\begin {large}
1. Що буде виведено за виконання?
code
try:
    res = 6<=8.0
    if isinstance(res, bool):
        print(f'bool;{res}')
    elif isinstance(res, int):
        print(f'int;{res}')
    elif isinstance(res, float):
        print(f'float;{res:.2f}')
    else:
        print('???')
except:
    print('ERR')
code

2. Що буде виведено за виконання?
code
a = 2
b = 9
c = 5

def f(b):
    a *= 4
    b = 2
    c = 5
    return a + b + c

a = 0
b = 6
c = 1

try:
    print(f(a), a, b, c, sep=';')
except:
    print('ERR', a, b, c, sep=';')
code

3. Що буде виведено за виконання?
code
try:
    k = 9
    while k<=12:
        k += 1
    print(k, end=';')
except:
    print('ERR')
code

4. Що буде виведено за виконання?
code
try:
    for d in range(6, -15, -1):
        if d>=5:
            break
        print(d, end=';')
    else:
        print(d,end=';')
    print(d)
except:
    print('ERR')
code

5. Що буде виведено за виконання?\par (Дійсні значення, якщо наявні, в цьому завданні будуть виводитись у форматі з фіксованою крапкою та, можливо, з доволі великою кількістю десяткових розрядів. У відповіді для дійсних значень у ЦЬОМУ завданні достатньо вказати один десятковий розряд після крапки, відкинувши решту розрядів без округлення. Наприклад, замість 13.66666666666667 достатньо вказати 13.6, замість 25.23 достатньо вказати 25.2)
code
try:
    print(7, end=';')
    print(3//0.0, end=';')
    print(5, end=';')
except TypeError:
    print(4, end=';')
except KeyboardInterrupt:
    print(1, end=';')
except:
    print('ERR', end=';')
else:
    print(8, end=';')
finally:
    print(0)
code

6. Що буде виведено за виконання?
code
def f(s):
    for i in range(len(s)):
        s[i] += 1
        return
 
s = [1, 2, 3]
f(s)
print(s)
code

7. Що буде виведено за виконання?
code
class A:
    c = 1

    def _f1(self): print(2, end=';')
    def __f2(self): print(3, end=';')

    def main(self):
        try: print(self.c, end=';')
        except: print('ERR', end=';')

        try: self._f1()
        except: print('ERR', end=';')

        try: self.__f2()
        except: print('ERR', end=';')


class B(A):
    c = 4

    def _f1(self): print(5, end=';')
    def __f2(self): print(6, end=';')

try: B().main()
except: print('ERR', end=';')
code

8. Що буде виведено за виконання?
code
def f(arg, n):
    arg[90] = 5
    n = 6
    arg = set(arg.values())

try:
    n = 0
    d = {59:7, 65:2, 85:8, 59:8}
    f(d, n)
    print(n, end=';')
    for x in d.keys():
        print(x, sep=';', end=';')
except:
    print('ERR')
code

9. Що буде виведено за виконання?
code
class A:
    a = 1
    b = 2

    def __init__(self):
        b = 3

class B(A):
    a = 4

    def __init__(self):
        super().__init__()
        self.c = 7

obj = B()
obj.b = obj.a + 10
B.b = 13

try: print(obj.b, end= ';')
except: print('ERR', end=';')

try: print(A.b, end= ';')
except: print('ERR', end=';')

try: print(B.b, end= ';')
except: print('ERR', end=';')
code

10. 
Для графа на рисунку з вершини 2 запущено алгоритм пошуку в глибину, що будує кістякове дерево. Списки суміжності вершин переглядаються алгоритмом за зростанням міток вершин.\par
\begin{center}
	\adjustimage{max size={0.9\linewidth}{0.9\paperheight}}
	{../10/Traversals/263.png}
\end{center}
\par Які з наведених ребер входять до складу отриманого кістякового дерева?\par
1) (0, 3)\par
2) (1, 3)\par
3) (3, 4)\par
4) (2, 6)\par
5) (1, 4)\par
6) (1, 2)\par
{\vskip 15pt} У формі вибрати номери відповідних ребер.\par
%answer:1 2 3 4 6
\end {large}}
%end
{\smallskip \begin {small} Затверджено на засіданні кафедри теоретичної кібернетики, протокол \No 12 від   19.05.2020 р.\par\smallskip\smallskip
{\parbox {6.5 cm}{Зав. кафедри \dotfill Ю.В.~Крак}}\hspace {0.7cm}{\hbox to 6.5 cm{ Екзаменатор \dotfill Т.О.~Карнаух}} \end{small}}
%begin
{{\begin {small}\par
{ \center  Київський національний університет імені Тараса Шевченка \par}
{\parbox {5 cm}{Освітня програма \hfill}{\parbox {7 cm}{Прикладна математика, Системний аналіз \hfill}}\parbox {6 cm}{\hfill Семестр 2}}\par
{\parbox {5 cm} {Навчальний предмет \hfill}\parbox {7 cm}{Програмування \hfill}\parbox {6cm}{\hfill}\par}\vspace{-7pt} \end{small}}{\center Екзаменаційний білет \No \textbf 
{ 448 }\par}}
{
\begin {large}
1. Що буде виведено за виконання?
code
try:
    res = "True">="4.0j"
    if isinstance(res, bool):
        print(f'bool;{res}')
    elif isinstance(res, int):
        print(f'int;{res}')
    elif isinstance(res, float):
        print(f'float;{res:.2f}')
    else:
        print('???')
except:
    print('ERR')
code

2. Що буде виведено за виконання?
code
a = 0
b = 7
c = 8

def h(a):
    a = 5
    b = 2
    c = 1
    return a + b + c

a = 4
b = 3
c = 9

try:
    print(h(a), a, b, c, sep=';')
except:
    print('ERR', a, b, c, sep=';')
code

3. Що буде виведено за виконання?
code
try:
    n = 1
    while n<=5:
        n += 1
    print(n, end=';')
except:
    print('ERR')
code

4. Що буде виведено за виконання?
code
try:
    for b in range(-2, 11, 3):
        if b>8:
            continue
        print(b, end=';')
    else:
        print(b,end=';')
    print(b)
except:
    print('ERR')
code

5. Що буде виведено за виконання?\par (Дійсні значення, якщо наявні, в цьому завданні будуть виводитись у форматі з фіксованою крапкою та, можливо, з доволі великою кількістю десяткових розрядів. У відповіді для дійсних значень у ЦЬОМУ завданні достатньо вказати один десятковий розряд після крапки, відкинувши решту розрядів без округлення. Наприклад, замість 13.66666666666667 достатньо вказати 13.6, замість 25.23 достатньо вказати 25.2)
code
try:
    print(9, end=';')
    print(int('d2'), end=';')
    print(6, end=';')
except ValueError:
    print(9, end=';')
except ValueError:
    print(5, end=';')
except:
    print('ERR', end=';')
else:
    print(0, end=';')
finally:
    print(3)
code

6. Що буде виведено за виконання?
code
def f(n):
    if n>9:
        return f(n//10)+1
    else:
        return 1

print(f(123456))
code

7. Що буде виведено за виконання?
code
class A:
    a = 1

    def _g1(self): print(2, end=';')
    def _g2(self): print(3, end=';')

    def main(self):
        try: print(self.a, end=';')
        except: print('ERR', end=';')

        try: self._g1()
        except: print('ERR', end=';')

        try: self._g2()
        except: print('ERR', end=';')


class B(A):
    a = 4

    def _g1(self): print(5, end=';')
    def _g2(self): print(6, end=';')

try: B().main()
except: print('ERR', end=';')
code

8. Що буде виведено за виконання?
code
def f(arg, n):
    arg[88] = 5
    n *= -2

try:
    n = 3
    d = {42:2, 87:4, 40:8, 42:0}
    f(d, n)
    print(n, end=';')
    for x in d.items():
        print(x, sep=';', end=';')
except:
    print('ERR')
code

9. Що буде виведено за виконання?
code
class A:
    c = 1
    b = 2

    def __init__(self):
        c = 3

class B(A):
    b = 4

    def __init__(self):
        super().__init__()
        self.a = 7

obj = B()
obj.c = obj.a + 10
B.c = 13

try: print(obj.c, end= ';')
except: print('ERR', end=';')

try: print(A.c, end= ';')
except: print('ERR', end=';')

try: print(B.c, end= ';')
except: print('ERR', end=';')
code

10. 
Для графа на рисунку з вершини 1 запущено алгоритм пошуку в глибину, що будує кістякове дерево. Списки суміжності вершин переглядаються алгоритмом за зростанням міток вершин.\par
\begin{center}
	\adjustimage{max size={0.9\linewidth}{0.9\paperheight}}
	{../10/Traversals/107.png}
\end{center}
\par Які з наведених ребер входять до складу отриманого кістякового дерева?\par
1) (2, 4)\par
2) (0, 5)\par
3) (5, 6)\par
4) (0, 3)\par
5) (4, 5)\par
6) (4, 6)\par
{\vskip 15pt} У формі вибрати номери відповідних ребер.\par
%answer:1 3 4
\end {large}}
%end
{\smallskip \begin {small} Затверджено на засіданні кафедри теоретичної кібернетики, протокол \No 12 від   19.05.2020 р.\par\smallskip\smallskip
{\parbox {6.5 cm}{Зав. кафедри \dotfill Ю.В.~Крак}}\hspace {0.7cm}{\hbox to 6.5 cm{ Екзаменатор \dotfill Т.О.~Карнаух}} \end{small}}
%begin
{{\begin {small}\par
{ \center  Київський національний університет імені Тараса Шевченка \par}
{\parbox {5 cm}{Освітня програма \hfill}{\parbox {7 cm}{Прикладна математика, Системний аналіз \hfill}}\parbox {6 cm}{\hfill Семестр 2}}\par
{\parbox {5 cm} {Навчальний предмет \hfill}\parbox {7 cm}{Програмування \hfill}\parbox {6cm}{\hfill}\par}\vspace{-7pt} \end{small}}{\center Екзаменаційний білет \No \textbf 
{ 449 }\par}}
{
\begin {large}
1. Що буде виведено за виконання?
code
try:
    res = 5/5//3*6
    if isinstance(res, bool):
        print(f'bool;{res}')
    elif isinstance(res, int):
        print(f'int;{res}')
    elif isinstance(res, float):
        print(f'float;{res:.2f}')
    else:
        print('???')
except:
    print('ERR')
code

2. Що буде виведено за виконання?
code
a = 6
b = 1
c = 0

def g(b):
    global c
    a = 3
    b = 1
    c = 5
    return a + b + c

a = 3
b = 2
c = 1

try:
    print(g(b), a, b, c, sep=';')
except:
    print('ERR', a, b, c, sep=';')
code

3. Що буде виведено за виконання?
code
try:
    j = 3
    while j<13:
        j += 2
    print(j, end=';')
except:
    print('ERR')
code

4. Що буде виведено за виконання?
code
try:
    for c in range(-15, 12, 2):
        if c<4:
            continue
        print(c, end=';');c = -1
    else:
        print(c,end=';')
    print(c)
except:
    print('ERR')
code

5. Що буде виведено за виконання?\par (Дійсні значення, якщо наявні, в цьому завданні будуть виводитись у форматі з фіксованою крапкою та, можливо, з доволі великою кількістю десяткових розрядів. У відповіді для дійсних значень у ЦЬОМУ завданні достатньо вказати один десятковий розряд після крапки, відкинувши решту розрядів без округлення. Наприклад, замість 13.66666666666667 достатньо вказати 13.6, замість 25.23 достатньо вказати 25.2)
code
try:
    print(1, end=';')
    print(8%3, end=';')
    print(7, end=';')
except KeyboardInterrupt:
    print(4, end=';')
except ZeroDivisionError:
    print(9, end=';')
except:
    print('ERR', end=';')
else:
    print(0, end=';')
finally:
    print(6)
code

6. Що буде виведено за виконання?
code
def f(s):
    for el in s:
        el += 1
    return s
 
s = [1, 2, 3]
s = f(s)
print(s)
code

7. Що буде виведено за виконання?
code
class A:
    b = 1

    def __h1(self): print(2, end=';')
    def _h2(self): print(3, end=';')

    def main(self):
        try: print(self.b, end=';')
        except: print('ERR', end=';')

        try: self.__h1()
        except: print('ERR', end=';')

        try: self._h2()
        except: print('ERR', end=';')


class B(A):
    b = 4

    def __h1(self): print(5, end=';')
    def _h2(self): print(6, end=';')

try: B().main()
except: print('ERR', end=';')
code

8. Що буде виведено за виконання?
code
def f(arg, n):
    arg[82] = 5
    n = 9
    arg = list(arg.keys())

try:
    n = 2
    t = {82:6, 27:9, 35:1, 27:6}
    f(t, n)
    print(n, end=';')
    for x, y in t.items():
        print(x, y, sep=';', end=';')
except:
    print('ERR')
code

9. Що буде виведено за виконання?
code
class A:
    a = 1
    c = 2

    def __init__(self):
        c = 3

class B(A):
    c = 4

    def __init__(self):
        super().__init__()
        self.b = 7

obj = B()
obj.b = obj.a + 10
B.b = 13

try: print(obj.b, end= ';')
except: print('ERR', end=';')

try: print(A.b, end= ';')
except: print('ERR', end=';')

try: print(B.b, end= ';')
except: print('ERR', end=';')
code

10. 
Для графа на рисунку з вершини 2 запущено алгоритм пошуку в глибину, що будує кістякове дерево. Списки суміжності вершин переглядаються алгоритмом за зростанням міток вершин.\par
\begin{center}
	\adjustimage{max size={0.9\linewidth}{0.9\paperheight}}
	{../10/Traversals/275.png}
\end{center}
\par Які з наведених ребер входять до складу отриманого кістякового дерева?\par
1) (2, 4)\par
2) (2, 3)\par
3) (4, 6)\par
4) (2, 5)\par
5) (1, 2)\par
6) (0, 1)\par
{\vskip 15pt} У формі вибрати номери відповідних ребер.\par
%answer:3 6
\end {large}}
%end
{\smallskip \begin {small} Затверджено на засіданні кафедри теоретичної кібернетики, протокол \No 12 від   19.05.2020 р.\par\smallskip\smallskip
{\parbox {6.5 cm}{Зав. кафедри \dotfill Ю.В.~Крак}}\hspace {0.7cm}{\hbox to 6.5 cm{ Екзаменатор \dotfill Т.О.~Карнаух}} \end{small}}
%begin
{{\begin {small}\par
{ \center  Київський національний університет імені Тараса Шевченка \par}
{\parbox {5 cm}{Освітня програма \hfill}{\parbox {7 cm}{Прикладна математика, Системний аналіз \hfill}}\parbox {6 cm}{\hfill Семестр 2}}\par
{\parbox {5 cm} {Навчальний предмет \hfill}\parbox {7 cm}{Програмування \hfill}\parbox {6cm}{\hfill}\par}\vspace{-7pt} \end{small}}{\center Екзаменаційний білет \No \textbf 
{ 450 }\par}}
{
\begin {large}
1. Що буде виведено за виконання?
code
try:
    res = not 2<4 or 4.0!=5 and True>8.0
    if isinstance(res, bool):
        print(f'bool;{res}')
    elif isinstance(res, int):
        print(f'int;{res}')
    elif isinstance(res, float):
        print(f'float;{res:.2f}')
    else:
        print('???')
except:
    print('ERR')
code

2. Що буде виведено за виконання?
code
a = 4
b = 9
c = 8

def f(a):
    a = 2
    b = 4
    c = 3
    return a + b + c

a = 5
b = 7
c = 6

try:
    print(f(a), a, b, c, sep=';')
except:
    print('ERR', a, b, c, sep=';')
code

3. Що буде виведено за виконання?
code
try:
    k = 5
    while k<=7:
        k += 1
    print(k, end=';')
except:
    print('ERR')
code

4. Що буде виведено за виконання?
code
try:
    for a in range(-10, -11, -2):
        if a<=3:
            break
        print(a, end=';')
    else:
        print('end',end=';')
    print(a)
except:
    print('ERR')
code

5. Що буде виведено за виконання?\par (Дійсні значення, якщо наявні, в цьому завданні будуть виводитись у форматі з фіксованою крапкою та, можливо, з доволі великою кількістю десяткових розрядів. У відповіді для дійсних значень у ЦЬОМУ завданні достатньо вказати один десятковий розряд після крапки, відкинувши решту розрядів без округлення. Наприклад, замість 13.66666666666667 достатньо вказати 13.6, замість 25.23 достатньо вказати 25.2)
code
try:
    print(7, end=';')
    print(4j>6j, end=';')
    print(8, end=';')
except Exception:
    print(5, end=';')
except TypeError:
    print(3, end=';')
except:
    print('ERR', end=';')
else:
    print(1, end=';')
finally:
    print(2)
code

6. Що буде виведено за виконання?
code
def f(n):
    print(n%10,end='')
    if n>9:
        f(n//100)
 
f(123456)
code

7. Що буде виведено за виконання?
code
class A:
    c = 1

    def _f1(self): print(2, end=';')
    def f2(self): print(3, end=';')

    def main(self):
        try: print(self.c, end=';')
        except: print('ERR', end=';')

        try: self._f1()
        except: print('ERR', end=';')

        try: self.f2()
        except: print('ERR', end=';')


class B(A):
    c = 4

    def _f1(self): print(5, end=';')
    def f2(self): print(6, end=';')

try: B().main()
except: print('ERR', end=';')
code

8. Що буде виведено за виконання?
code
def f(arg, n):
    n = 1
    tmp = arg
    del tmp[-1:4]
    return len(tmp)>3

try:
    f = ['0','1','2','3','4']
    n = 2
    f(f, n)
    print(n, end=';')
    for el in f:
        print(el, end=';')
except:
    print('ERR')
code

9. Що буде виведено за виконання?
code
class A:
    c = 1
    a = 2

    def __init__(self):
        c = 3

class B(A):
    c = 4

    def __init__(self):
        super().__init__()
        self.b = 7

obj = B()
obj.b = obj.c + 10
B.b = 13

try: print(obj.b, end= ';')
except: print('ERR', end=';')

try: print(A.b, end= ';')
except: print('ERR', end=';')

try: print(B.b, end= ';')
except: print('ERR', end=';')
code

10. 
Для графа на рисунку з вершини 2 запущено алгоритм пошуку в глибину, що будує кістякове дерево. Списки суміжності вершин переглядаються алгоритмом за зростанням міток вершин.\par
\begin{center}
	\adjustimage{max size={0.9\linewidth}{0.9\paperheight}}
	{../10/Traversals/161.png}
\end{center}
\par Які з наведених ребер входять до складу отриманого кістякового дерева?\par
1) (2, 5)\par
2) (2, 4)\par
3) (4, 5)\par
4) (1, 3)\par
5) (0, 5)\par
6) (0, 1)\par
{\vskip 15pt} У формі вибрати номери відповідних ребер.\par
%answer:2 3 4 5 6
\end {large}}
%end
{\smallskip \begin {small} Затверджено на засіданні кафедри теоретичної кібернетики, протокол \No 12 від   19.05.2020 р.\par\smallskip\smallskip
{\parbox {6.5 cm}{Зав. кафедри \dotfill Ю.В.~Крак}}\hspace {0.7cm}{\hbox to 6.5 cm{ Екзаменатор \dotfill Т.О.~Карнаух}} \end{small}}
%begin
{{\begin {small}\par
{ \center  Київський національний університет імені Тараса Шевченка \par}
{\parbox {5 cm}{Освітня програма \hfill}{\parbox {7 cm}{Прикладна математика, Системний аналіз \hfill}}\parbox {6 cm}{\hfill Семестр 2}}\par
{\parbox {5 cm} {Навчальний предмет \hfill}\parbox {7 cm}{Програмування \hfill}\parbox {6cm}{\hfill}\par}\vspace{-7pt} \end{small}}{\center Екзаменаційний білет \No \textbf 
{ 451 }\par}}
{
\begin {large}
1. Що буде виведено за виконання?
code
try:
    res = not 4==9.0 and 2.0>=3 or False<=2
    if isinstance(res, bool):
        print(f'bool;{res}')
    elif isinstance(res, int):
        print(f'int;{res}')
    elif isinstance(res, float):
        print(f'float;{res:.2f}')
    else:
        print('???')
except:
    print('ERR')
code

2. Що буде виведено за виконання?
code
a = 7
b = 5
c = 6

def h(a):
    global c
    a -= 2
    b = 4
    c = 3
    return a + b + c

a = 2
b = 4
c = 8

try:
    print(h(a), a, b, c, sep=';')
except:
    print('ERR', a, b, c, sep=';')
code

3. Що буде виведено за виконання?
code
try:
    j = 2
    while j<10:
        j += 2
    print(j, end=';')
except:
    print('ERR')
code

4. Що буде виведено за виконання?
code
try:
    for c in range(15, -14, -2):
        if c>7:
            break
        print(c, end=';');c = 0
    else:
        print(c,end=';')
    print(c)
except:
    print('ERR')
code

5. Що буде виведено за виконання?\par (Дійсні значення, якщо наявні, в цьому завданні будуть виводитись у форматі з фіксованою крапкою та, можливо, з доволі великою кількістю десяткових розрядів. У відповіді для дійсних значень у ЦЬОМУ завданні достатньо вказати один десятковий розряд після крапки, відкинувши решту розрядів без округлення. Наприклад, замість 13.66666666666667 достатньо вказати 13.6, замість 25.23 достатньо вказати 25.2)
code
try:
    print(7, end=';')
    print(int('a1'), end=';')
    print(3, end=';')
except ValueError:
    print(5, end=';')
except TypeError:
    print(2, end=';')
except:
    print('ERR', end=';')
else:
    print(0, end=';')
finally:
    print(4)
code

6. Що буде виведено за виконання?
code
def f(n):
    if n>9:
        f(n//100)
        print(n%10, end='')
 
f(123456)
code

7. Що буде виведено за виконання?
code
class A:
    a = 1

    def __h1(self): print(2, end=';')
    def _h2(self): print(3, end=';')

    def main(self):
        try: print(self.a, end=';')
        except: print('ERR', end=';')

        try: self.__h1()
        except: print('ERR', end=';')

        try: self._h2()
        except: print('ERR', end=';')


class B(A):
    a = 4

    def __h1(self): print(5, end=';')
    def _h2(self): print(6, end=';')

try: B().main()
except: print('ERR', end=';')
code

8. Що буде виведено за виконання?
code
def f(arg, n):
    arg[42] = 6
    n *= 2

try:
    n = 7
    t = {38:1, 53:4, 72:7, 38:2}
    f(t, n)
    print(n, end=';')
    for x in t.keys():
        print(x, sep=';', end=';')
except:
    print('ERR')
code

9. Що буде виведено за виконання?
code
class A:
    c = 1
    a = 2

    def __init__(self):
        c = 3

class B(A):
    c = 4

    def __init__(self):
        super().__init__()
        self.b = 7

obj = B()
obj.c = obj.a + 10
B.c = 13

try: print(obj.c, end= ';')
except: print('ERR', end=';')

try: print(A.c, end= ';')
except: print('ERR', end=';')

try: print(B.c, end= ';')
except: print('ERR', end=';')
code

10. 
Для графа на рисунку з вершини 6 запущено алгоритм пошуку в глибину, що будує кістякове дерево. Списки суміжності вершин переглядаються алгоритмом за зростанням міток вершин.\par
\begin{center}
	\adjustimage{max size={0.9\linewidth}{0.9\paperheight}}
	{../10/Traversals/9.png}
\end{center}
\par Які з наведених ребер входять до складу отриманого кістякового дерева?\par
1) (2, 6)\par
2) (0, 2)\par
3) (5, 6)\par
4) (3, 4)\par
5) (0, 5)\par
6) (3, 5)\par
{\vskip 15pt} У формі вибрати номери відповідних ребер.\par
%answer:1 2 6
\end {large}}
%end
{\smallskip \begin {small} Затверджено на засіданні кафедри теоретичної кібернетики, протокол \No 12 від   19.05.2020 р.\par\smallskip\smallskip
{\parbox {6.5 cm}{Зав. кафедри \dotfill Ю.В.~Крак}}\hspace {0.7cm}{\hbox to 6.5 cm{ Екзаменатор \dotfill Т.О.~Карнаух}} \end{small}}
%begin
{{\begin {small}\par
{ \center  Київський національний університет імені Тараса Шевченка \par}
{\parbox {5 cm}{Освітня програма \hfill}{\parbox {7 cm}{Прикладна математика, Системний аналіз \hfill}}\parbox {6 cm}{\hfill Семестр 2}}\par
{\parbox {5 cm} {Навчальний предмет \hfill}\parbox {7 cm}{Програмування \hfill}\parbox {6cm}{\hfill}\par}\vspace{-7pt} \end{small}}{\center Екзаменаційний білет \No \textbf 
{ 452 }\par}}
{
\begin {large}
1. Що буде виведено за виконання?
code
try:
    res = "8">=6.0
    if isinstance(res, bool):
        print(f'bool;{res}')
    elif isinstance(res, int):
        print(f'int;{res}')
    elif isinstance(res, float):
        print(f'float;{res:.2f}')
    else:
        print('???')
except:
    print('ERR')
code

2. Що буде виведено за виконання?
code
a = 7
b = 2
c = 1

def f(b):
    global c
    a = 2
    b = 4
    c = 5
    return a + b + c

a = 0
b = 4
c = 9

try:
    print(f(b), a, b, c, sep=';')
except:
    print('ERR', a, b, c, sep=';')
code

3. Що буде виведено за виконання?
code
try:
    m = 8
    while m<9:
        m += 2
    print(m, end=';')
except:
    print('ERR')
code

4. Що буде виведено за виконання?
code
try:
    for f in range(12, -5, -1):
        if f>=5:
            continue
        print(f, end=';')
    else:
        print(f,end=';')
    print(f)
except:
    print('ERR')
code

5. Що буде виведено за виконання?\par (Дійсні значення, якщо наявні, в цьому завданні будуть виводитись у форматі з фіксованою крапкою та, можливо, з доволі великою кількістю десяткових розрядів. У відповіді для дійсних значень у ЦЬОМУ завданні достатньо вказати один десятковий розряд після крапки, відкинувши решту розрядів без округлення. Наприклад, замість 13.66666666666667 достатньо вказати 13.6, замість 25.23 достатньо вказати 25.2)
code
try:
    print(9, end=';')
    print(6/3, end=';')
    print(8, end=';')
except ZeroDivisionError:
    print(5, end=';')
except Exception:
    print(2, end=';')
except:
    print('ERR', end=';')
else:
    print(9, end=';')
finally:
    print(4)
code

6. Що буде виведено за виконання?
code
def f(n):
    print(n%10, end='')
    if n>9:
        f(n//10)
 
f(543216)
code

7. Що буде виведено за виконання?
code
class A:
    d = 1

    def _g1(self): print(2, end=';')
    def _g2(self): print(3, end=';')

    def main(self):
        try: print(self.d, end=';')
        except: print('ERR', end=';')

        try: self._g1()
        except: print('ERR', end=';')

        try: self._g2()
        except: print('ERR', end=';')


class B(A):
    d = 4

    def _g1(self): print(5, end=';')
    def _g2(self): print(6, end=';')

try: B().main()
except: print('ERR', end=';')
code

8. Що буде виведено за виконання?
code
def f(arg, n):
    arg[71] = 6
    n = 5
    arg = list(arg.values())

try:
    n = 1
    d = {71:2, 47:1, 33:6, 33:8}
    f(d, n)
    print(n, end=';')
    for x in d.keys():
        print(x, sep=';', end=';')
except:
    print('ERR')
code

9. Що буде виведено за виконання?
code
class A:
    a = 1
    b = 2

    def __init__(self):
        b = 3

class B(A):
    b = 4

    def __init__(self):
        super().__init__()
        self.c = 7

obj = B()
obj.b = obj.c + 10
B.b = 13

try: print(obj.b, end= ';')
except: print('ERR', end=';')

try: print(A.b, end= ';')
except: print('ERR', end=';')

try: print(B.b, end= ';')
except: print('ERR', end=';')
code

10. 
Для графа на рисунку з вершини 5 запущено алгоритм пошуку в глибину, що будує кістякове дерево. Списки суміжності вершин переглядаються алгоритмом за зростанням міток вершин.\par
\begin{center}
	\adjustimage{max size={0.9\linewidth}{0.9\paperheight}}
	{../10/Traversals/140.png}
\end{center}
\par Які з наведених ребер входять до складу отриманого кістякового дерева?\par
1) (1, 6)\par
2) (0, 5)\par
3) (3, 5)\par
4) (3, 4)\par
5) (0, 4)\par
6) (4, 6)\par
{\vskip 15pt} У формі вибрати номери відповідних ребер.\par
%answer:1 2 4 5
\end {large}}
%end
{\smallskip \begin {small} Затверджено на засіданні кафедри теоретичної кібернетики, протокол \No 12 від   19.05.2020 р.\par\smallskip\smallskip
{\parbox {6.5 cm}{Зав. кафедри \dotfill Ю.В.~Крак}}\hspace {0.7cm}{\hbox to 6.5 cm{ Екзаменатор \dotfill Т.О.~Карнаух}} \end{small}}
%begin
{{\begin {small}\par
{ \center  Київський національний університет імені Тараса Шевченка \par}
{\parbox {5 cm}{Освітня програма \hfill}{\parbox {7 cm}{Прикладна математика, Системний аналіз \hfill}}\parbox {6 cm}{\hfill Семестр 2}}\par
{\parbox {5 cm} {Навчальний предмет \hfill}\parbox {7 cm}{Програмування \hfill}\parbox {6cm}{\hfill}\par}\vspace{-7pt} \end{small}}{\center Екзаменаційний білет \No \textbf 
{ 453 }\par}}
{
\begin {large}
1. Що буде виведено за виконання?
code
try:
    res = -3**2+True
    if isinstance(res, bool):
        print(f'bool;{res}')
    elif isinstance(res, int):
        print(f'int;{res}')
    elif isinstance(res, float):
        print(f'float;{res:.2f}')
    else:
        print('???')
except:
    print('ERR')
code

2. Що буде виведено за виконання?
code
a = 6
b = 3
c = 5

def h(b):
    global c
    a = 1
    b = 3
    c = 3
    return a + b + c

a = 8
b = 8
c = 3

try:
    print(h(a), a, b, c, sep=';')
except:
    print('ERR', a, b, c, sep=';')
code

3. Що буде виведено за виконання?
code
try:
    t = 6
    while t<=11:
        t += 1
    print(t, end=';')
except:
    print('ERR')
code

4. Що буде виведено за виконання?
code
try:
    for b in range(1, 6, -3):
        if b<3:
            continue
        print(b, end=';')
    else:
        print('end',end=';')
    print(b)
except:
    print('ERR')
code

5. Що буде виведено за виконання?\par (Дійсні значення, якщо наявні, в цьому завданні будуть виводитись у форматі з фіксованою крапкою та, можливо, з доволі великою кількістю десяткових розрядів. У відповіді для дійсних значень у ЦЬОМУ завданні достатньо вказати один десятковий розряд після крапки, відкинувши решту розрядів без округлення. Наприклад, замість 13.66666666666667 достатньо вказати 13.6, замість 25.23 достатньо вказати 25.2)
code
try:
    print(3, end=';')
    print(1//0.0, end=';')
    print(6, end=';')
except KeyboardInterrupt:
    print(8, end=';')
except ValueError:
    print(7, end=';')
except:
    print('ERR', end=';')
else:
    print(0, end=';')
finally:
    print(5)
code

6. Що буде виведено за виконання?
code
def f(n):
    if n>9:
        f(n//10)
    print(n%10, end='')
 
f(543216)
code

7. Що буде виведено за виконання?
code
class A:
    a = 1

    def _f1(self): print(2, end=';')
    def __f2(self): print(3, end=';')

    def main(self):
        try: print(self.a, end=';')
        except: print('ERR', end=';')

        try: self._f1()
        except: print('ERR', end=';')

        try: self.__f2()
        except: print('ERR', end=';')


class B(A):
    a = 4

    def _f1(self): print(5, end=';')
    def __f2(self): print(6, end=';')

try: B().main()
except: print('ERR', end=';')
code

8. Що буде виведено за виконання?
code
def f(arg, n):
    arg[72] = 9
    n = 9
    arg = set(arg.keys())

try:
    n = 3
    t = {16:1, 63:3, 16:5}
    f(t, n)
    print(n, end=';')
    for x in t.values():
        print(x, sep=';', end=';')
except:
    print('ERR')
code

9. Що буде виведено за виконання?
code
class A:
    a = 1
    c = 2

    def __init__(self):
        c = 3

class B(A):
    a = 4

    def __init__(self):
        super().__init__()
        self.b = 7

obj = B()
obj.a = obj.b + 10
B.a = 13

try: print(obj.a, end= ';')
except: print('ERR', end=';')

try: print(A.a, end= ';')
except: print('ERR', end=';')

try: print(B.a, end= ';')
except: print('ERR', end=';')
code

10. 
Для графа на рисунку з вершини 2 запущено алгоритм пошуку в глибину, що будує кістякове дерево. Списки суміжності вершин переглядаються алгоритмом за зростанням міток вершин.\par
\begin{center}
	\adjustimage{max size={0.9\linewidth}{0.9\paperheight}}
	{../10/Traversals/128.png}
\end{center}
\par Які з наведених ребер входять до складу отриманого кістякового дерева?\par
1) (2, 5)\par
2) (0, 2)\par
3) (3, 6)\par
4) (3, 5)\par
5) (0, 4)\par
6) (1, 3)\par
{\vskip 15pt} У формі вибрати номери відповідних ребер.\par
%answer:2 3 4 5 6
\end {large}}
%end
{\smallskip \begin {small} Затверджено на засіданні кафедри теоретичної кібернетики, протокол \No 12 від   19.05.2020 р.\par\smallskip\smallskip
{\parbox {6.5 cm}{Зав. кафедри \dotfill Ю.В.~Крак}}\hspace {0.7cm}{\hbox to 6.5 cm{ Екзаменатор \dotfill Т.О.~Карнаух}} \end{small}}
%begin
{{\begin {small}\par
{ \center  Київський національний університет імені Тараса Шевченка \par}
{\parbox {5 cm}{Освітня програма \hfill}{\parbox {7 cm}{Прикладна математика, Системний аналіз \hfill}}\parbox {6 cm}{\hfill Семестр 2}}\par
{\parbox {5 cm} {Навчальний предмет \hfill}\parbox {7 cm}{Програмування \hfill}\parbox {6cm}{\hfill}\par}\vspace{-7pt} \end{small}}{\center Екзаменаційний білет \No \textbf 
{ 454 }\par}}
{
\begin {large}
1. Що буде виведено за виконання?
code
try:
    res = 7!=True or not 9.0>=6 and 8<=5.0
    if isinstance(res, bool):
        print(f'bool;{res}')
    elif isinstance(res, int):
        print(f'int;{res}')
    elif isinstance(res, float):
        print(f'float;{res:.2f}')
    else:
        print('???')
except:
    print('ERR')
code

2. Що буде виведено за виконання?
code
a = 0
b = 6
c = 5

def g(a):
    a = 5
    b = 2
    c = 4
    return a + b + c

a = 9
b = 4
c = 1

try:
    print(g(b), a, b, c, sep=';')
except:
    print('ERR', a, b, c, sep=';')
code

3. Що буде виведено за виконання?
code
try:
    n = 15
    while n>=4:
        n += -2
    print(n, end=';')
except:
    print('ERR')
code

4. Що буде виведено за виконання?
code
try:
    for a in range(13, -3, 3):
        if a>=6:
            continue
        print(a, end=';')
    else:
        print(a,end=';')
    print(a)
except:
    print('ERR')
code

5. Що буде виведено за виконання?\par (Дійсні значення, якщо наявні, в цьому завданні будуть виводитись у форматі з фіксованою крапкою та, можливо, з доволі великою кількістю десяткових розрядів. У відповіді для дійсних значень у ЦЬОМУ завданні достатньо вказати один десятковий розряд після крапки, відкинувши решту розрядів без округлення. Наприклад, замість 13.66666666666667 достатньо вказати 13.6, замість 25.23 достатньо вказати 25.2)
code
try:
    print(1, end=';')
    print(int('2'), end=';')
    print(4, end=';')
except ZeroDivisionError:
    print(9, end=';')
except TypeError:
    print(6, end=';')
except:
    print('ERR', end=';')
else:
    print(3, end=';')
finally:
    print(7)
code

6. Що буде виведено за виконання?
code
def f(n):
    if n>9:
        f(n//100)
    print(n%10, end='')
 
f(123456)
code

7. Що буде виведено за виконання?
code
class A:
    d = 1

    def _f1(self): print(2, end=';')
    def f2(self): print(3, end=';')

    def main(self):
        try: print(self.d, end=';')
        except: print('ERR', end=';')

        try: self._f1()
        except: print('ERR', end=';')

        try: self.f2()
        except: print('ERR', end=';')


class B(A):
    d = 4

    def _f1(self): print(5, end=';')
    def f2(self): print(6, end=';')

try: B().main()
except: print('ERR', end=';')
code

8. Що буде виведено за виконання?
code
def f(arg, n):
    n = 5
    tmp = arg
    del tmp[:-6]
    return len(tmp)>3

try:
    g = ['0','1','2','3','4']
    n = 3
    f(g, n)
    print(n, end=';')
    for el in g:
        print(el, end=';')
except:
    print('ERR')
code

9. Що буде виведено за виконання?
code
class A:
    c = 1
    b = 2

    def __init__(self):
        c = 3

class B(A):
    b = 4

    def __init__(self):
        super().__init__()
        self.a = 7

obj = B()
obj.c = obj.b + 10
B.c = 13

try: print(obj.c, end= ';')
except: print('ERR', end=';')

try: print(A.c, end= ';')
except: print('ERR', end=';')

try: print(B.c, end= ';')
except: print('ERR', end=';')
code

10. 
Для графа на рисунку з вершини 3 запущено алгоритм пошуку в глибину, що будує кістякове дерево. Списки суміжності вершин переглядаються алгоритмом за зростанням міток вершин.\par
\begin{center}
	\adjustimage{max size={0.9\linewidth}{0.9\paperheight}}
	{../10/Traversals/126.png}
\end{center}
\par Які з наведених ребер входять до складу отриманого кістякового дерева?\par
1) (3, 4)\par
2) (3, 6)\par
3) (3, 5)\par
4) (1, 3)\par
5) (1, 2)\par
6) (0, 5)\par
{\vskip 15pt} У формі вибрати номери відповідних ребер.\par
%answer:5 6
\end {large}}
%end
{\smallskip \begin {small} Затверджено на засіданні кафедри теоретичної кібернетики, протокол \No 12 від   19.05.2020 р.\par\smallskip\smallskip
{\parbox {6.5 cm}{Зав. кафедри \dotfill Ю.В.~Крак}}\hspace {0.7cm}{\hbox to 6.5 cm{ Екзаменатор \dotfill Т.О.~Карнаух}} \end{small}}
%begin
{{\begin {small}\par
{ \center  Київський національний університет імені Тараса Шевченка \par}
{\parbox {5 cm}{Освітня програма \hfill}{\parbox {7 cm}{Прикладна математика, Системний аналіз \hfill}}\parbox {6 cm}{\hfill Семестр 2}}\par
{\parbox {5 cm} {Навчальний предмет \hfill}\parbox {7 cm}{Програмування \hfill}\parbox {6cm}{\hfill}\par}\vspace{-7pt} \end{small}}{\center Екзаменаційний білет \No \textbf 
{ 455 }\par}}
{
\begin {large}
1. Що буде виведено за виконання?
code
try:
    res = False==9 and not 3.0>4.0 or 4<8
    if isinstance(res, bool):
        print(f'bool;{res}')
    elif isinstance(res, int):
        print(f'int;{res}')
    elif isinstance(res, float):
        print(f'float;{res:.2f}')
    else:
        print('???')
except:
    print('ERR')
code

2. Що буде виведено за виконання?
code
a = 0
b = 9
c = 1

def f(a):
    global c
    a = 5
    b += 1
    c = 2
    return a + b + c

a = 3
b = 3
c = 1

try:
    print(f(b), a, b, c, sep=';')
except:
    print('ERR', a, b, c, sep=';')
code

3. Що буде виведено за виконання?
code
try:
    i = 1
    while i<6:
        i += 2
    print(i, end=';')
except:
    print('ERR')
code

4. Що буде виведено за виконання?
code
try:
    for d in range(-8, 15, 2):
        if d<=8:
            continue
        print(d, end=';')
    else:
        print(d,end=';')
    print(d)
except:
    print('ERR')
code

5. Що буде виведено за виконання?\par (Дійсні значення, якщо наявні, в цьому завданні будуть виводитись у форматі з фіксованою крапкою та, можливо, з доволі великою кількістю десяткових розрядів. У відповіді для дійсних значень у ЦЬОМУ завданні достатньо вказати один десятковий розряд після крапки, відкинувши решту розрядів без округлення. Наприклад, замість 13.66666666666667 достатньо вказати 13.6, замість 25.23 достатньо вказати 25.2)
code
try:
    print(8, end=';')
    print(0j<3j, end=';')
    print(7, end=';')
except KeyboardInterrupt:
    print(1, end=';')
except Exception:
    print(4, end=';')
except:
    print('ERR', end=';')
else:
    print(0, end=';')
finally:
    print(5)
code

6. Що буде виведено за виконання?
code
def f(s1):
    s = s1[:]
    for i in range(len(s)):
        s[i] += 1
 
s = [1, 2, 3]
f(s)
print(s)
code

7. Що буде виведено за виконання?
code
class A:
    b = 1

    def g1(self): print(2, end=';')
    def _g2(self): print(3, end=';')

    def main(self):
        try: print(self.b, end=';')
        except: print('ERR', end=';')

        try: self.g1()
        except: print('ERR', end=';')

        try: self._g2()
        except: print('ERR', end=';')


class B(A):
    b = 4

    def g1(self): print(5, end=';')
    def _g2(self): print(6, end=';')

try: B().main()
except: print('ERR', end=';')
code

8. Що буде виведено за виконання?
code
def f(arg, n):
    arg[38] = 5
    n = 8
    arg = tuple(arg.items())

try:
    n = 0
    s = {81:6, 77:2, 38:3, 81:5}
    f(s, n)
    print(n, end=';')
    for x in s.values():
        print(x, sep=';', end=';')
except:
    print('ERR')
code

9. Що буде виведено за виконання?
code
class A:
    b = 1
    a = 2

    def __init__(self):
        b = 3

class B(A):
    a = 4

    def __init__(self):
        super().__init__()
        self.c = 7

obj = B()
obj.a = obj.b + 10
B.a = 13

try: print(obj.a, end= ';')
except: print('ERR', end=';')

try: print(A.a, end= ';')
except: print('ERR', end=';')

try: print(B.a, end= ';')
except: print('ERR', end=';')
code

10. 
Для графа на рисунку з вершини 5 запущено алгоритм пошуку в глибину, що будує кістякове дерево. Списки суміжності вершин переглядаються алгоритмом за зростанням міток вершин.\par
\begin{center}
	\adjustimage{max size={0.9\linewidth}{0.9\paperheight}}
	{../10/Traversals/208.png}
\end{center}
\par Які з наведених ребер входять до складу отриманого кістякового дерева?\par
1) (3, 5)\par
2) (3, 6)\par
3) (0, 3)\par
4) (0, 4)\par
5) (1, 5)\par
6) (3, 4)\par
{\vskip 15pt} У формі вибрати номери відповідних ребер.\par
%answer:2 3 4 5
\end {large}}
%end
{\smallskip \begin {small} Затверджено на засіданні кафедри теоретичної кібернетики, протокол \No 12 від   19.05.2020 р.\par\smallskip\smallskip
{\parbox {6.5 cm}{Зав. кафедри \dotfill Ю.В.~Крак}}\hspace {0.7cm}{\hbox to 6.5 cm{ Екзаменатор \dotfill Т.О.~Карнаух}} \end{small}}
%begin
{{\begin {small}\par
{ \center  Київський національний університет імені Тараса Шевченка \par}
{\parbox {5 cm}{Освітня програма \hfill}{\parbox {7 cm}{Прикладна математика, Системний аналіз \hfill}}\parbox {6 cm}{\hfill Семестр 2}}\par
{\parbox {5 cm} {Навчальний предмет \hfill}\parbox {7 cm}{Програмування \hfill}\parbox {6cm}{\hfill}\par}\vspace{-7pt} \end{small}}{\center Екзаменаційний білет \No \textbf 
{ 456 }\par}}
{
\begin {large}
1. Що буде виведено за виконання?
code
try:
    res = 4j!="False"
    if isinstance(res, bool):
        print(f'bool;{res}')
    elif isinstance(res, int):
        print(f'int;{res}')
    elif isinstance(res, float):
        print(f'float;{res:.2f}')
    else:
        print('???')
except:
    print('ERR')
code

2. Що буде виведено за виконання?
code
a = 6
b = 4
c = 7

def g(b):
    a = 1
    b *= 4
    c = 5
    return a + b + c

a = 0
b = 9
c = 8

try:
    print(g(a), a, b, c, sep=';')
except:
    print('ERR', a, b, c, sep=';')
code

3. Що буде виведено за виконання?
code
try:
    m = 13
    while m>=9:
        m += -2
    print(m, end=';')
except:
    print('ERR')
code

4. Що буде виведено за виконання?
code
try:
    for e in range(7, -12, 2):
        if e>4:
            break
        print(e, end=';');e = 4
        if e<=4:
            break
    else:
        print(e,end=';')
    print(e)
except:
    print('ERR')
code

5. Що буде виведено за виконання?\par (Дійсні значення, якщо наявні, в цьому завданні будуть виводитись у форматі з фіксованою крапкою та, можливо, з доволі великою кількістю десяткових розрядів. У відповіді для дійсних значень у ЦЬОМУ завданні достатньо вказати один десятковий розряд після крапки, відкинувши решту розрядів без округлення. Наприклад, замість 13.66666666666667 достатньо вказати 13.6, замість 25.23 достатньо вказати 25.2)
code
try:
    print(3, end=';')
    print(int('8'), end=';')
    print(2, end=';')
except KeyboardInterrupt:
    print(9, end=';')
except Exception:
    print(6, end=';')
except:
    print('ERR', end=';')
else:
    print(9, end=';')
finally:
    print(4)
code

6. Що буде виведено за виконання?
code
def f(n):
    if n>9:
        f(n//100)
    else:
        print(n%10,end='')

f(987654)
code

7. Що буде виведено за виконання?
code
class A:
    c = 1

    def _h1(self): print(2, end=';')
    def __h2(self): print(3, end=';')

    def main(self):
        try: print(self.c, end=';')
        except: print('ERR', end=';')

        try: self._h1()
        except: print('ERR', end=';')

        try: self.__h2()
        except: print('ERR', end=';')


class B(A):
    c = 4

    def _h1(self): print(5, end=';')
    def __h2(self): print(6, end=';')

try: B().main()
except: print('ERR', end=';')
code

8. Що буде виведено за виконання?
code
def f(arg, n):
    arg[90] = 0
    n += -3

try:
    n = 3
    s = {90:4, 18:1, 18:4}
    f(s, n)
    print(n, end=';')
    for x in s :
        print(x, sep=';', end=';')
except:
    print('ERR')
code

9. Що буде виведено за виконання?
code
class A:
    b = 1
    c = 2

    def __init__(self):
        c = 3

class B(A):
    b = 4

    def __init__(self):
        super().__init__()
        self.a = 7

obj = B()
obj.a = obj.c + 10
B.a = 13

try: print(obj.a, end= ';')
except: print('ERR', end=';')

try: print(A.a, end= ';')
except: print('ERR', end=';')

try: print(B.a, end= ';')
except: print('ERR', end=';')
code

10. 
Для графа на рисунку з вершини 3 запущено алгоритм пошуку в глибину, що будує кістякове дерево. Списки суміжності вершин переглядаються алгоритмом за зростанням міток вершин.\par
\begin{center}
	\adjustimage{max size={0.9\linewidth}{0.9\paperheight}}
	{../10/Traversals/182.png}
\end{center}
\par Які з наведених ребер входять до складу отриманого кістякового дерева?\par
1) (1, 4)\par
2) (3, 4)\par
3) (5, 6)\par
4) (2, 3)\par
5) (1, 6)\par
6) (1, 3)\par
{\vskip 15pt} У формі вибрати номери відповідних ребер.\par
%answer:3 6
\end {large}}
%end
{\smallskip \begin {small} Затверджено на засіданні кафедри теоретичної кібернетики, протокол \No 12 від   19.05.2020 р.\par\smallskip\smallskip
{\parbox {6.5 cm}{Зав. кафедри \dotfill Ю.В.~Крак}}\hspace {0.7cm}{\hbox to 6.5 cm{ Екзаменатор \dotfill Т.О.~Карнаух}} \end{small}}
%begin
{{\begin {small}\par
{ \center  Київський національний університет імені Тараса Шевченка \par}
{\parbox {5 cm}{Освітня програма \hfill}{\parbox {7 cm}{Прикладна математика, Системний аналіз \hfill}}\parbox {6 cm}{\hfill Семестр 2}}\par
{\parbox {5 cm} {Навчальний предмет \hfill}\parbox {7 cm}{Програмування \hfill}\parbox {6cm}{\hfill}\par}\vspace{-7pt} \end{small}}{\center Екзаменаційний білет \No \textbf 
{ 457 }\par}}
{
\begin {large}
1. Що буде виведено за виконання?
code
try:
    res = 8//3%7*7
    if isinstance(res, bool):
        print(f'bool;{res}')
    elif isinstance(res, int):
        print(f'int;{res}')
    elif isinstance(res, float):
        print(f'float;{res:.2f}')
    else:
        print('???')
except:
    print('ERR')
code

2. Що буде виведено за виконання?
code
a = 5
b = 2
c = 2

def f(a):
    global c
    a -= 3
    b = 4
    c = 2
    return a + b + c

a = 8
b = 9
c = 5

try:
    print(f(b), a, b, c, sep=';')
except:
    print('ERR', a, b, c, sep=';')
code

3. Що буде виведено за виконання?
code
try:
    n = 9
    while n>2:
        n += -3
    print(n, end=';')
except:
    print('ERR')
code

4. Що буде виведено за виконання?
code
try:
    for f in range(-11, -6, -1):
        if f<5:
            continue
        print(f, end=';')
    else:
        print('end',end=';')
    print(f)
except:
    print('ERR')
code

5. Що буде виведено за виконання?\par (Дійсні значення, якщо наявні, в цьому завданні будуть виводитись у форматі з фіксованою крапкою та, можливо, з доволі великою кількістю десяткових розрядів. У відповіді для дійсних значень у ЦЬОМУ завданні достатньо вказати один десятковий розряд після крапки, відкинувши решту розрядів без округлення. Наприклад, замість 13.66666666666667 достатньо вказати 13.6, замість 25.23 достатньо вказати 25.2)
code
try:
    print(1, end=';')
    print(7%0, end=';')
    print(0, end=';')
except TypeError:
    print(2, end=';')
except ValueError:
    print(5, end=';')
except:
    print('ERR', end=';')
else:
    print(8, end=';')
finally:
    print(6)
code

6. Що буде виведено за виконання?
code
def f(n):
    if n<=9:
        print(n%10, end='')
    else:
        f(n//10)

f(987651)
code

7. Що буде виведено за виконання?
code
class A:
    c = 1

    def _f1(self): print(2, end=';')
    def f2(self): print(3, end=';')

    def main(self):
        try: print(self.c, end=';')
        except: print('ERR', end=';')

        try: self._f1()
        except: print('ERR', end=';')

        try: self.f2()
        except: print('ERR', end=';')


class B(A):
    c = 4

    def _f1(self): print(5, end=';')
    def f2(self): print(6, end=';')

try: B().main()
except: print('ERR', end=';')
code

8. Що буде виведено за виконання?
code
def f(arg, n):
    n = 1
    tmp = arg
    tmp[-2:8] = [4,5]
    return len(tmp)>3

try:
    c = [10,11,12,13,14]
    n = 2
    f(c, n)
    print(n, end=';')
    for el in c:
        print(el, end=';')
except:
    print('ERR')
code

9. Що буде виведено за виконання?
code
class A:
    c = 1
    a = 2

    def __init__(self):
        a = 3

class B(A):
    a = 4

    def __init__(self):
        super().__init__()
        self.b = 7

obj = B()
obj.c = obj.b + 10
B.c = 13

try: print(obj.c, end= ';')
except: print('ERR', end=';')

try: print(A.c, end= ';')
except: print('ERR', end=';')

try: print(B.c, end= ';')
except: print('ERR', end=';')
code

10. 
Для графа на рисунку з вершини 3 запущено алгоритм пошуку в глибину, що будує кістякове дерево. Списки суміжності вершин переглядаються алгоритмом за зростанням міток вершин.\par
\begin{center}
	\adjustimage{max size={0.9\linewidth}{0.9\paperheight}}
	{../10/Traversals/50.png}
\end{center}
\par Які з наведених ребер входять до складу отриманого кістякового дерева?\par
1) (5, 6)\par
2) (0, 1)\par
3) (2, 3)\par
4) (2, 4)\par
5) (1, 3)\par
6) (4, 6)\par
{\vskip 15pt} У формі вибрати номери відповідних ребер.\par
%answer:1 2 4 5 6
\end {large}}
%end
{\smallskip \begin {small} Затверджено на засіданні кафедри теоретичної кібернетики, протокол \No 12 від   19.05.2020 р.\par\smallskip\smallskip
{\parbox {6.5 cm}{Зав. кафедри \dotfill Ю.В.~Крак}}\hspace {0.7cm}{\hbox to 6.5 cm{ Екзаменатор \dotfill Т.О.~Карнаух}} \end{small}}
%begin
{{\begin {small}\par
{ \center  Київський національний університет імені Тараса Шевченка \par}
{\parbox {5 cm}{Освітня програма \hfill}{\parbox {7 cm}{Прикладна математика, Системний аналіз \hfill}}\parbox {6 cm}{\hfill Семестр 2}}\par
{\parbox {5 cm} {Навчальний предмет \hfill}\parbox {7 cm}{Програмування \hfill}\parbox {6cm}{\hfill}\par}\vspace{-7pt} \end{small}}{\center Екзаменаційний білет \No \textbf 
{ 458 }\par}}
{
\begin {large}
1. Що буде виведено за виконання?
code
def f(n):
    print('f', end='')
    return n

try:
    print(f(4) and f(-6))
except:
    print('ERR')
code

2. Що буде виведено за виконання?
code
a = 6
b = 0
c = 1

def h(b):
    a *= 1
    b = 5
    c = 3
    return a + b + c

a = 4
b = 3
c = 7

try:
    print(h(b), a, b, c, sep=';')
except:
    print('ERR', a, b, c, sep=';')
code

3. Що буде виведено за виконання?
code
try:
    i = 12
    while i>6:
        i += -2
    print(i, end=';')
except:
    print('ERR')
code

4. Що буде виведено за виконання?
code
try:
    for c in range(3, 9, -2):
        if c>=7:
            break
        print(c, end=';')
        if c>5:
            break
    else:
        print(c,end=';')
    print(c)
except:
    print('ERR')
code

5. Що буде виведено за виконання?\par (Дійсні значення, якщо наявні, в цьому завданні будуть виводитись у форматі з фіксованою крапкою та, можливо, з доволі великою кількістю десяткових розрядів. У відповіді для дійсних значень у ЦЬОМУ завданні достатньо вказати один десятковий розряд після крапки, відкинувши решту розрядів без округлення. Наприклад, замість 13.66666666666667 достатньо вказати 13.6, замість 25.23 достатньо вказати 25.2)
code
try:
    print(3, end=';')
    print(4j==7j, end=';')
    print(7, end=';')
except ZeroDivisionError:
    print(6, end=';')
except KeyboardInterrupt:
    print(2, end=';')
except:
    print('ERR', end=';')
else:
    print(1, end=';')
finally:
    print(5)
code

6. Що буде виведено за виконання?
code
def f(n):
    print(n%2, end='')
    if n>2:
        f(n//2)
 
f(16)
code

7. Що буде виведено за виконання?
code
class A:
    d = 1

    def h1(self): print(2, end=';')
    def _h2(self): print(3, end=';')

    def main(self):
        try: print(self.d, end=';')
        except: print('ERR', end=';')

        try: self.h1()
        except: print('ERR', end=';')

        try: self._h2()
        except: print('ERR', end=';')


class B(A):
    d = 4

    def h1(self): print(5, end=';')
    def _h2(self): print(6, end=';')

try: B().main()
except: print('ERR', end=';')
code

8. Що буде виведено за виконання?
code
def f(arg, n):
    arg[17] = 7
    n = 7
    arg = set(arg.items())

try:
    n = 4
    t = {17:9, 11:1, 11:3}
    f(t, n)
    print(n, end=';')
    for x in t.items():
        print(x, sep=';', end=';')
except:
    print('ERR')
code

9. Що буде виведено за виконання?
code
class A:
    a = 1
    b = 2

    def __init__(self):
        a = 3

class B(A):
    a = 4

    def __init__(self):
        super().__init__()
        self.c = 7

obj = B()
obj.a = obj.a + 10
B.a = 13

try: print(obj.a, end= ';')
except: print('ERR', end=';')

try: print(A.a, end= ';')
except: print('ERR', end=';')

try: print(B.a, end= ';')
except: print('ERR', end=';')
code

10. 
Для графа на рисунку з вершини 1 запущено алгоритм пошуку в глибину, що будує кістякове дерево. Списки суміжності вершин переглядаються алгоритмом за зростанням міток вершин.\par
\begin{center}
	\adjustimage{max size={0.9\linewidth}{0.9\paperheight}}
	{../10/Traversals/11.png}
\end{center}
\par Які з наведених ребер входять до складу отриманого кістякового дерева?\par
1) (2, 4)\par
2) (4, 6)\par
3) (5, 6)\par
4) (0, 3)\par
5) (3, 6)\par
6) (0, 5)\par
{\vskip 15pt} У формі вибрати номери відповідних ребер.\par
%answer:3 4
\end {large}}
%end
{\smallskip \begin {small} Затверджено на засіданні кафедри теоретичної кібернетики, протокол \No 12 від   19.05.2020 р.\par\smallskip\smallskip
{\parbox {6.5 cm}{Зав. кафедри \dotfill Ю.В.~Крак}}\hspace {0.7cm}{\hbox to 6.5 cm{ Екзаменатор \dotfill Т.О.~Карнаух}} \end{small}}
%begin
{{\begin {small}\par
{ \center  Київський національний університет імені Тараса Шевченка \par}
{\parbox {5 cm}{Освітня програма \hfill}{\parbox {7 cm}{Прикладна математика, Системний аналіз \hfill}}\parbox {6 cm}{\hfill Семестр 2}}\par
{\parbox {5 cm} {Навчальний предмет \hfill}\parbox {7 cm}{Програмування \hfill}\parbox {6cm}{\hfill}\par}\vspace{-7pt} \end{small}}{\center Екзаменаційний білет \No \textbf 
{ 459 }\par}}
{
\begin {large}
1. Що буде виведено за виконання?
code
try:
    res = 8/5*8*7
    if isinstance(res, bool):
        print(f'bool;{res}')
    elif isinstance(res, int):
        print(f'int;{res}')
    elif isinstance(res, float):
        print(f'float;{res:.2f}')
    else:
        print('???')
except:
    print('ERR')
code

2. Що буде виведено за виконання?
code
a = 2
b = 7
c = 9

def g(a):
    global c
    a = 1
    b = 5
    c = 4
    return a + b + c

a = 7
b = 3
c = 8

try:
    print(g(b), a, b, c, sep=';')
except:
    print('ERR', a, b, c, sep=';')
code

3. Що буде виведено за виконання?
code
try:
    k = 11
    while k>=0:
        k += -3
    print(k, end=';')
except:
    print('ERR')
code

4. Що буде виведено за виконання?
code
try:
    for e in range(-14, 10, 3):
        if e>3:
            continue
        print(e, end=';');e = 1
    else:
        print(e,end=';')
    print(e)
except:
    print('ERR')
code

5. Що буде виведено за виконання?\par (Дійсні значення, якщо наявні, в цьому завданні будуть виводитись у форматі з фіксованою крапкою та, можливо, з доволі великою кількістю десяткових розрядів. У відповіді для дійсних значень у ЦЬОМУ завданні достатньо вказати один десятковий розряд після крапки, відкинувши решту розрядів без округлення. Наприклад, замість 13.66666666666667 достатньо вказати 13.6, замість 25.23 достатньо вказати 25.2)
code
try:
    print(3, end=';')
    print(8%1, end=';')
    print(0, end=';')
except ValueError:
    print(9, end=';')
except ZeroDivisionError:
    print(4, end=';')
except:
    print('ERR', end=';')
else:
    print(6, end=';')
finally:
    print(5)
code

6. Що буде виведено за виконання?
code
def f(n):
    if n>9:
        f(n//10)
    else:
        print(n%10, end='')

f(543216)
code

7. Що буде виведено за виконання?
code
class A:
    b = 1

    def _g1(self): print(2, end=';')
    def _g2(self): print(3, end=';')

    def main(self):
        try: print(self.b, end=';')
        except: print('ERR', end=';')

        try: self._g1()
        except: print('ERR', end=';')

        try: self._g2()
        except: print('ERR', end=';')


class B(A):
    b = 4

    def _g1(self): print(5, end=';')
    def _g2(self): print(6, end=';')

try: B().main()
except: print('ERR', end=';')
code

8. Що буде виведено за виконання?
code
def f(arg, n):
    arg[75] = 0
    n = 6
    arg = tuple(arg.values())

try:
    n = 2
    s = {80:2, 65:4, 47:1, 80:9}
    f(s, n)
    print(n, end=';')
    for x in s.items():
        print(*x, sep=';', end=';')
except:
    print('ERR')
code

9. Що буде виведено за виконання?
code
class A:
    b = 1
    a = 2

    def __init__(self):
        b = 3

class B(A):
    b = 4

    def __init__(self):
        super().__init__()
        self.c = 7

obj = B()
obj.c = obj.b + 10
B.c = 13

try: print(obj.c, end= ';')
except: print('ERR', end=';')

try: print(A.c, end= ';')
except: print('ERR', end=';')

try: print(B.c, end= ';')
except: print('ERR', end=';')
code

10. 
Для графа на рисунку з вершини 4 запущено алгоритм пошуку в глибину, що будує кістякове дерево. Списки суміжності вершин переглядаються алгоритмом за зростанням міток вершин.\par
\begin{center}
	\adjustimage{max size={0.9\linewidth}{0.9\paperheight}}
	{../10/Traversals/8.png}
\end{center}
\par Які з наведених ребер входять до складу отриманого кістякового дерева?\par
1) (2, 5)\par
2) (1, 3)\par
3) (0, 6)\par
4) (2, 6)\par
5) (4, 5)\par
6) (0, 3)\par
{\vskip 15pt} У формі вибрати номери відповідних ребер.\par
%answer:2 3 4 5
\end {large}}
%end
{\smallskip \begin {small} Затверджено на засіданні кафедри теоретичної кібернетики, протокол \No 12 від   19.05.2020 р.\par\smallskip\smallskip
{\parbox {6.5 cm}{Зав. кафедри \dotfill Ю.В.~Крак}}\hspace {0.7cm}{\hbox to 6.5 cm{ Екзаменатор \dotfill Т.О.~Карнаух}} \end{small}}
%begin
{{\begin {small}\par
{ \center  Київський національний університет імені Тараса Шевченка \par}
{\parbox {5 cm}{Освітня програма \hfill}{\parbox {7 cm}{Прикладна математика, Системний аналіз \hfill}}\parbox {6 cm}{\hfill Семестр 2}}\par
{\parbox {5 cm} {Навчальний предмет \hfill}\parbox {7 cm}{Програмування \hfill}\parbox {6cm}{\hfill}\par}\vspace{-7pt} \end{small}}{\center Екзаменаційний білет \No \textbf 
{ 460 }\par}}
{
\begin {large}
1. Що буде виведено за виконання?
code
def f(n):
    print('f', end='')
    return n

try:
    print(f(2) and f(-8))
except:
    print('ERR')
code

2. Що буде виведено за виконання?
code
a = 4
b = 2
c = 0

def f(b):
    a = 2
    b = 1
    c = 3
    return a + b + c

a = 5
b = 6
c = 1

try:
    print(f(a), a, b, c, sep=';')
except:
    print('ERR', a, b, c, sep=';')
code

3. Що буде виведено за виконання?
code
try:
    n = 5
    while n<=8:
        n += 1
    print(n, end=';')
except:
    print('ERR')
code

4. Що буде виведено за виконання?
code
try:
    for a in range(5, -2, -3):
        if a<=6:
            break
        print(a, end=';')
    else:
        print(a,end=';')
    print(a)
except:
    print('ERR')
code

5. Що буде виведено за виконання?\par (Дійсні значення, якщо наявні, в цьому завданні будуть виводитись у форматі з фіксованою крапкою та, можливо, з доволі великою кількістю десяткових розрядів. У відповіді для дійсних значень у ЦЬОМУ завданні достатньо вказати один десятковий розряд після крапки, відкинувши решту розрядів без округлення. Наприклад, замість 13.66666666666667 достатньо вказати 13.6, замість 25.23 достатньо вказати 25.2)
code
try:
    print(9, end=';')
    print(int('b1'), end=';')
    print(0, end=';')
except Exception:
    print(7, end=';')
except TypeError:
    print(3, end=';')
except:
    print('ERR', end=';')
else:
    print(2, end=';')
finally:
    print(8)
code

6. Що буде виведено за виконання?
code
def f(n):
    if n>2:
        f(n//2)
    else:
        print(n%2, end='')

f(16)
code

7. Що буде виведено за виконання?
code
class A:
    a = 1

    def __h1(self): print(2, end=';')
    def _h2(self): print(3, end=';')

    def main(self):
        try: print(self.a, end=';')
        except: print('ERR', end=';')

        try: self.__h1()
        except: print('ERR', end=';')

        try: self._h2()
        except: print('ERR', end=';')


class B(A):
    a = 4

    def __h1(self): print(5, end=';')
    def _h2(self): print(6, end=';')

try: B().main()
except: print('ERR', end=';')
code

8. Що буде виведено за виконання?
code
def f(arg, n):
    arg[11] = 7
    n = 7
    arg = list(arg.keys())

try:
    n = 4
    d = {11:8, 75:4, 18:8, 48:7}
    f(d, n)
    print(n, end=';')
    for x in d.keys():
        print(x, sep=';', end=';')
except:
    print('ERR')
code

9. Що буде виведено за виконання?
code
class A:
    b = 1
    c = 2

    def __init__(self):
        c = 3

class B(A):
    c = 4

    def __init__(self):
        super().__init__()
        self.a = 7

obj = B()
obj.b = obj.c + 10
B.b = 13

try: print(obj.b, end= ';')
except: print('ERR', end=';')

try: print(A.b, end= ';')
except: print('ERR', end=';')

try: print(B.b, end= ';')
except: print('ERR', end=';')
code

10. 
Для графа на рисунку з вершини 4 запущено алгоритм пошуку в глибину, що будує кістякове дерево. Списки суміжності вершин переглядаються алгоритмом за зростанням міток вершин.\par
\begin{center}
	\adjustimage{max size={0.9\linewidth}{0.9\paperheight}}
	{../10/Traversals/133.png}
\end{center}
\par Які з наведених ребер входять до складу отриманого кістякового дерева?\par
1) (2, 6)\par
2) (2, 3)\par
3) (1, 6)\par
4) (2, 4)\par
5) (5, 6)\par
6) (1, 4)\par
{\vskip 15pt} У формі вибрати номери відповідних ребер.\par
%answer:1 2 5 6
\end {large}}
%end
{\smallskip \begin {small} Затверджено на засіданні кафедри теоретичної кібернетики, протокол \No 12 від   19.05.2020 р.\par\smallskip\smallskip
{\parbox {6.5 cm}{Зав. кафедри \dotfill Ю.В.~Крак}}\hspace {0.7cm}{\hbox to 6.5 cm{ Екзаменатор \dotfill Т.О.~Карнаух}} \end{small}}
%begin
{{\begin {small}\par
{ \center  Київський національний університет імені Тараса Шевченка \par}
{\parbox {5 cm}{Освітня програма \hfill}{\parbox {7 cm}{Прикладна математика, Системний аналіз \hfill}}\parbox {6 cm}{\hfill Семестр 2}}\par
{\parbox {5 cm} {Навчальний предмет \hfill}\parbox {7 cm}{Програмування \hfill}\parbox {6cm}{\hfill}\par}\vspace{-7pt} \end{small}}{\center Екзаменаційний білет \No \textbf 
{ 461 }\par}}
{
\begin {large}
1. Що буде виведено за виконання?
code
try:
    res = 3==True or not 9<=7.0 or 8.0!=5
    if isinstance(res, bool):
        print(f'bool;{res}')
    elif isinstance(res, int):
        print(f'int;{res}')
    elif isinstance(res, float):
        print(f'float;{res:.2f}')
    else:
        print('???')
except:
    print('ERR')
code

2. Що буде виведено за виконання?
code
a = 1
b = 8
c = 4

def g(b):
    a += 3
    b = 5
    c = 2
    return a + b + c

a = 3
b = 6
c = 7

try:
    print(g(a), a, b, c, sep=';')
except:
    print('ERR', a, b, c, sep=';')
code

3. Що буде виведено за виконання?
code
try:
    t = 9
    while t<10:
        t += 2
    print(t, end=';')
except:
    print('ERR')
code

4. Що буде виведено за виконання?
code
try:
    for d in range(2, -7, 3):
        if d<8:
            continue
        print(d, end=';');d = 2
    else:
        print(d,end=';')
    print(d)
except:
    print('ERR')
code

5. Що буде виведено за виконання?\par (Дійсні значення, якщо наявні, в цьому завданні будуть виводитись у форматі з фіксованою крапкою та, можливо, з доволі великою кількістю десяткових розрядів. У відповіді для дійсних значень у ЦЬОМУ завданні достатньо вказати один десятковий розряд після крапки, відкинувши решту розрядів без округлення. Наприклад, замість 13.66666666666667 достатньо вказати 13.6, замість 25.23 достатньо вказати 25.2)
code
try:
    print(0, end=';')
    print(4//2, end=';')
    print(5, end=';')
except ValueError:
    print(2, end=';')
except Exception:
    print(8, end=';')
except:
    print('ERR', end=';')
else:
    print(3, end=';')
finally:
    print(6)
code

6. Що буде виведено за виконання?
code
def f(n):
    if n>9:
        f(n//10)
    print(n%10, end='')
 
f(123456)
code

7. Що буде виведено за виконання?
code
class A:
    b = 1

    def _f1(self): print(2, end=';')
    def __f2(self): print(3, end=';')

    def main(self):
        try: print(self.b, end=';')
        except: print('ERR', end=';')

        try: self._f1()
        except: print('ERR', end=';')

        try: self.__f2()
        except: print('ERR', end=';')


class B(A):
    b = 4

    def _f1(self): print(5, end=';')
    def __f2(self): print(6, end=';')

try: B().main()
except: print('ERR', end=';')
code

8. Що буде виведено за виконання?
code
def f(arg, n):
    arg[23] = 0
    n -= -2

try:
    n = 4
    d = {55:8, 39:9, 55:8}
    f(d, n)
    print(n, end=';')
    for x in d.keys():
        print(x, sep=';', end=';')
except:
    print('ERR')
code

9. Що буде виведено за виконання?
code
class A:
    a = 1
    c = 2

    def __init__(self):
        c = 3

class B(A):
    c = 4

    def __init__(self):
        super().__init__()
        self.b = 7

obj = B()
obj.a = obj.c + 10
B.a = 13

try: print(obj.a, end= ';')
except: print('ERR', end=';')

try: print(A.a, end= ';')
except: print('ERR', end=';')

try: print(B.a, end= ';')
except: print('ERR', end=';')
code

10. 
Для графа на рисунку з вершини 0 запущено алгоритм пошуку в глибину, що будує кістякове дерево. Списки суміжності вершин переглядаються алгоритмом за зростанням міток вершин.\par
\begin{center}
	\adjustimage{max size={0.9\linewidth}{0.9\paperheight}}
	{../10/Traversals/165.png}
\end{center}
\par Які з наведених ребер входять до складу отриманого кістякового дерева?\par
1) (4, 6)\par
2) (0, 4)\par
3) (2, 4)\par
4) (2, 3)\par
5) (5, 6)\par
6) (1, 2)\par
{\vskip 15pt} У формі вибрати номери відповідних ребер.\par
%answer:2 3 4 5 6
\end {large}}
%end
{\smallskip \begin {small} Затверджено на засіданні кафедри теоретичної кібернетики, протокол \No 12 від   19.05.2020 р.\par\smallskip\smallskip
{\parbox {6.5 cm}{Зав. кафедри \dotfill Ю.В.~Крак}}\hspace {0.7cm}{\hbox to 6.5 cm{ Екзаменатор \dotfill Т.О.~Карнаух}} \end{small}}
%begin
{{\begin {small}\par
{ \center  Київський національний університет імені Тараса Шевченка \par}
{\parbox {5 cm}{Освітня програма \hfill}{\parbox {7 cm}{Прикладна математика, Системний аналіз \hfill}}\parbox {6 cm}{\hfill Семестр 2}}\par
{\parbox {5 cm} {Навчальний предмет \hfill}\parbox {7 cm}{Програмування \hfill}\parbox {6cm}{\hfill}\par}\vspace{-7pt} \end{small}}{\center Екзаменаційний білет \No \textbf 
{ 462 }\par}}
{
\begin {large}
1. Що буде виведено за виконання?
code
try:
    res = 4//4%3*6
    if isinstance(res, bool):
        print(f'bool;{res}')
    elif isinstance(res, int):
        print(f'int;{res}')
    elif isinstance(res, float):
        print(f'float;{res:.2f}')
    else:
        print('???')
except:
    print('ERR')
code

2. Що буде виведено за виконання?
code
a = 0
b = 2
c = 9

def g(a):
    global c
    a = 4
    b *= 1
    c = 5
    return a + b + c

a = 5
b = 1
c = 2

try:
    print(g(a), a, b, c, sep=';')
except:
    print('ERR', a, b, c, sep=';')
code

3. Що буде виведено за виконання?
code
try:
    j = 8
    while j>=8:
        j += -2
    print(j, end=';')
except:
    print('ERR')
code

4. Що буде виведено за виконання?
code
try:
    for b in range(-1, 8, -3):
        if b<=4:
            continue
        print(b, end=';')
    else:
        print(b,end=';')
    print(b)
except:
    print('ERR')
code

5. Що буде виведено за виконання?\par (Дійсні значення, якщо наявні, в цьому завданні будуть виводитись у форматі з фіксованою крапкою та, можливо, з доволі великою кількістю десяткових розрядів. У відповіді для дійсних значень у ЦЬОМУ завданні достатньо вказати один десятковий розряд після крапки, відкинувши решту розрядів без округлення. Наприклад, замість 13.66666666666667 достатньо вказати 13.6, замість 25.23 достатньо вказати 25.2)
code
try:
    print(9, end=';')
    print(7j>4j, end=';')
    print(1, end=';')
except TypeError:
    print(0, end=';')
except KeyboardInterrupt:
    print(2, end=';')
except:
    print('ERR', end=';')
else:
    print(8, end=';')
finally:
    print(4)
code

6. Що буде виведено за виконання?
code
def f(s):
    for i in range(len(s)):
        s[i] += 1
        return
 
s = [1, 2, 3]
f(s)
print(s)
code

7. Що буде виведено за виконання?
code
class A:
    a = 1

    def g1(self): print(2, end=';')
    def _g2(self): print(3, end=';')

    def main(self):
        try: print(self.a, end=';')
        except: print('ERR', end=';')

        try: self.g1()
        except: print('ERR', end=';')

        try: self._g2()
        except: print('ERR', end=';')


class B(A):
    a = 4

    def g1(self): print(5, end=';')
    def _g2(self): print(6, end=';')

try: B().main()
except: print('ERR', end=';')
code

8. Що буде виведено за виконання?
code
def f(arg, n):
    n = 6
    tmp = arg
    tmp[-6:] = 'ab'
    return len(tmp)>3

try:
    c = [10,11,12,13,14]
    n = 7
    f(c, n)
    print(n, end=';')
    for el in c:
        print(el, end=';')
except:
    print('ERR')
code

9. Що буде виведено за виконання?
code
class A:
    c = 1
    b = 2

    def __init__(self):
        c = 3

class B(A):
    c = 4

    def __init__(self):
        super().__init__()
        self.a = 7

obj = B()
obj.a = obj.b + 10
B.a = 13

try: print(obj.a, end= ';')
except: print('ERR', end=';')

try: print(A.a, end= ';')
except: print('ERR', end=';')

try: print(B.a, end= ';')
except: print('ERR', end=';')
code

10. 
Для графа на рисунку з вершини 1 запущено алгоритм пошуку в глибину, що будує кістякове дерево. Списки суміжності вершин переглядаються алгоритмом за зростанням міток вершин.\par
\begin{center}
	\adjustimage{max size={0.9\linewidth}{0.9\paperheight}}
	{../10/Traversals/51.png}
\end{center}
\par Які з наведених ребер входять до складу отриманого кістякового дерева?\par
1) (2, 5)\par
2) (0, 2)\par
3) (1, 4)\par
4) (0, 5)\par
5) (1, 3)\par
6) (0, 6)\par
{\vskip 15pt} У формі вибрати номери відповідних ребер.\par
%answer:2 3 4
\end {large}}
%end
{\smallskip \begin {small} Затверджено на засіданні кафедри теоретичної кібернетики, протокол \No 12 від   19.05.2020 р.\par\smallskip\smallskip
{\parbox {6.5 cm}{Зав. кафедри \dotfill Ю.В.~Крак}}\hspace {0.7cm}{\hbox to 6.5 cm{ Екзаменатор \dotfill Т.О.~Карнаух}} \end{small}}
%begin
{{\begin {small}\par
{ \center  Київський національний університет імені Тараса Шевченка \par}
{\parbox {5 cm}{Освітня програма \hfill}{\parbox {7 cm}{Прикладна математика, Системний аналіз \hfill}}\parbox {6 cm}{\hfill Семестр 2}}\par
{\parbox {5 cm} {Навчальний предмет \hfill}\parbox {7 cm}{Програмування \hfill}\parbox {6cm}{\hfill}\par}\vspace{-7pt} \end{small}}{\center Екзаменаційний білет \No \textbf 
{ 463 }\par}}
{
\begin {large}
1. Що буде виведено за виконання?
code
try:
    res = True<"7.0j"
    if isinstance(res, bool):
        print(f'bool;{res}')
    elif isinstance(res, int):
        print(f'int;{res}')
    elif isinstance(res, float):
        print(f'float;{res:.2f}')
    else:
        print('???')
except:
    print('ERR')
code

2. Що буде виведено за виконання?
code
a = 1
b = 5
c = 9

def h(b):
    global c
    a = 2
    b = 4
    c = 5
    return a + b + c

a = 4
b = 3
c = 0

try:
    print(h(a), a, b, c, sep=';')
except:
    print('ERR', a, b, c, sep=';')
code

3. Що буде виведено за виконання?
code
try:
    t = 7
    while t>5:
        t += -3
    print(t, end=';')
except:
    print('ERR')
code

4. Що буде виведено за виконання?
code
try:
    for f in range(-13, 2, -2):
        if f>=4:
            continue
        print(f, end=';');f = 5
    else:
        print('end',end=';')
    print(f)
except:
    print('ERR')
code

5. Що буде виведено за виконання?\par (Дійсні значення, якщо наявні, в цьому завданні будуть виводитись у форматі з фіксованою крапкою та, можливо, з доволі великою кількістю десяткових розрядів. У відповіді для дійсних значень у ЦЬОМУ завданні достатньо вказати один десятковий розряд після крапки, відкинувши решту розрядів без округлення. Наприклад, замість 13.66666666666667 достатньо вказати 13.6, замість 25.23 достатньо вказати 25.2)
code
try:
    print(3, end=';')
    print(int('c9'), end=';')
    print(6, end=';')
except ZeroDivisionError:
    print(7, end=';')
except Exception:
    print(5, end=';')
except:
    print('ERR', end=';')
else:
    print(1, end=';')
finally:
    print(6)
code

6. Що буде виведено за виконання?
code
def f(n):
    if n>9:
        f(n//100)
    print(n%10,end='')
 
f(987654)
code

7. Що буде виведено за виконання?
code
class A:
    d = 1

    def _h1(self): print(2, end=';')
    def h2(self): print(3, end=';')

    def main(self):
        try: print(self.d, end=';')
        except: print('ERR', end=';')

        try: self._h1()
        except: print('ERR', end=';')

        try: self.h2()
        except: print('ERR', end=';')


class B(A):
    d = 4

    def _h1(self): print(5, end=';')
    def h2(self): print(6, end=';')

try: B().main()
except: print('ERR', end=';')
code

8. Що буде виведено за виконання?
code
def f(arg, n):
    n = 8
    tmp = arg
    tmp[6:-8] = [15]
    return len(tmp)>3

try:
    g = ['0','1','2','3','4']
    n = 4
    f(g, n)
    print(n, end=';')
    for el in g:
        print(el, end=';')
except:
    print('ERR')
code

9. Що буде виведено за виконання?
code
class A:
    a = 1
    c = 2

    def __init__(self):
        a = 3

class B(A):
    c = 4

    def __init__(self):
        super().__init__()
        self.b = 7

obj = B()
obj.a = obj.b + 10
B.a = 13

try: print(obj.a, end= ';')
except: print('ERR', end=';')

try: print(A.a, end= ';')
except: print('ERR', end=';')

try: print(B.a, end= ';')
except: print('ERR', end=';')
code

10. 
Для графа на рисунку з вершини 6 запущено алгоритм пошуку в глибину, що будує кістякове дерево. Списки суміжності вершин переглядаються алгоритмом за зростанням міток вершин.\par
\begin{center}
	\adjustimage{max size={0.9\linewidth}{0.9\paperheight}}
	{../10/Traversals/87.png}
\end{center}
\par Які з наведених ребер входять до складу отриманого кістякового дерева?\par
1) (4, 6)\par
2) (0, 1)\par
3) (1, 2)\par
4) (2, 6)\par
5) (1, 5)\par
6) (3, 6)\par
{\vskip 15pt} У формі вибрати номери відповідних ребер.\par
%answer:2 3 5
\end {large}}
%end
{\smallskip \begin {small} Затверджено на засіданні кафедри теоретичної кібернетики, протокол \No 12 від   19.05.2020 р.\par\smallskip\smallskip
{\parbox {6.5 cm}{Зав. кафедри \dotfill Ю.В.~Крак}}\hspace {0.7cm}{\hbox to 6.5 cm{ Екзаменатор \dotfill Т.О.~Карнаух}} \end{small}}
%begin
{{\begin {small}\par
{ \center  Київський національний університет імені Тараса Шевченка \par}
{\parbox {5 cm}{Освітня програма \hfill}{\parbox {7 cm}{Прикладна математика, Системний аналіз \hfill}}\parbox {6 cm}{\hfill Семестр 2}}\par
{\parbox {5 cm} {Навчальний предмет \hfill}\parbox {7 cm}{Програмування \hfill}\parbox {6cm}{\hfill}\par}\vspace{-7pt} \end{small}}{\center Екзаменаційний білет \No \textbf 
{ 464 }\par}}
{
\begin {large}
1. Що буде виведено за виконання?
code
def f(n):
    print('f', end='')
    return n<-3

try:
    print(f(-1) or f(-9))
except:
    print('ERR')
code

2. Що буде виведено за виконання?
code
a = 8
b = 6
c = 7

def g(a):
    a = 3
    b = 1
    c = 3
    return a + b + c

a = 2
b = 8
c = 4

try:
    print(g(b), a, b, c, sep=';')
except:
    print('ERR', a, b, c, sep=';')
code

3. Що буде виведено за виконання?
code
try:
    k = 4
    while k<=5:
        k += 1
    print(k, end=';')
except:
    print('ERR')
code

4. Що буде виведено за виконання?
code
try:
    for c in range(-3, 3, 2):
        if c>5:
            continue
        print(c, end=';')
    else:
        print(c,end=';')
    print(c)
except:
    print('ERR')
code

5. Що буде виведено за виконання?\par (Дійсні значення, якщо наявні, в цьому завданні будуть виводитись у форматі з фіксованою крапкою та, можливо, з доволі великою кількістю десяткових розрядів. У відповіді для дійсних значень у ЦЬОМУ завданні достатньо вказати один десятковий розряд після крапки, відкинувши решту розрядів без округлення. Наприклад, замість 13.66666666666667 достатньо вказати 13.6, замість 25.23 достатньо вказати 25.2)
code
try:
    print(9, end=';')
    print(8/0, end=';')
    print(5, end=';')
except TypeError:
    print(2, end=';')
except KeyboardInterrupt:
    print(3, end=';')
except:
    print('ERR', end=';')
else:
    print(4, end=';')
finally:
    print(7)
code

6. Що буде виведено за виконання?
code
def f(n):
    if n>9:
        f(n//100)
    else:
        print(n%10,end='')

f(123456)
code

7. Що буде виведено за виконання?
code
class A:
    c = 1

    def _f1(self): print(2, end=';')
    def _f2(self): print(3, end=';')

    def main(self):
        try: print(self.c, end=';')
        except: print('ERR', end=';')

        try: self._f1()
        except: print('ERR', end=';')

        try: self._f2()
        except: print('ERR', end=';')


class B(A):
    c = 4

    def _f1(self): print(5, end=';')
    def _f2(self): print(6, end=';')

try: B().main()
except: print('ERR', end=';')
code

8. Що буде виведено за виконання?
code
def f(arg, n):
    arg[41] = 0
    n = 6
    arg = list(arg.items())

try:
    n = 1
    s = {80:4, 41:0, 80:9}
    f(s, n)
    print(n, end=';')
    for x in s.values():
        print(x, sep=';', end=';')
except:
    print('ERR')
code

9. Що буде виведено за виконання?
code
class A:
    a = 1
    b = 2

    def __init__(self):
        b = 3

class B(A):
    a = 4

    def __init__(self):
        super().__init__()
        self.c = 7

obj = B()
obj.c = obj.a + 10
B.c = 13

try: print(obj.c, end= ';')
except: print('ERR', end=';')

try: print(A.c, end= ';')
except: print('ERR', end=';')

try: print(B.c, end= ';')
except: print('ERR', end=';')
code

10. 
Для графа на рисунку з вершини 3 запущено алгоритм пошуку в глибину, що будує кістякове дерево. Списки суміжності вершин переглядаються алгоритмом за зростанням міток вершин.\par
\begin{center}
	\adjustimage{max size={0.9\linewidth}{0.9\paperheight}}
	{../10/Traversals/56.png}
\end{center}
\par Які з наведених ребер входять до складу отриманого кістякового дерева?\par
1) (2, 6)\par
2) (1, 2)\par
3) (0, 1)\par
4) (1, 4)\par
5) (0, 6)\par
6) (3, 4)\par
{\vskip 15pt} У формі вибрати номери відповідних ребер.\par
%answer:1 2 3 6
\end {large}}
%end
{\smallskip \begin {small} Затверджено на засіданні кафедри теоретичної кібернетики, протокол \No 12 від   19.05.2020 р.\par\smallskip\smallskip
{\parbox {6.5 cm}{Зав. кафедри \dotfill Ю.В.~Крак}}\hspace {0.7cm}{\hbox to 6.5 cm{ Екзаменатор \dotfill Т.О.~Карнаух}} \end{small}}
%begin
{{\begin {small}\par
{ \center  Київський національний університет імені Тараса Шевченка \par}
{\parbox {5 cm}{Освітня програма \hfill}{\parbox {7 cm}{Прикладна математика, Системний аналіз \hfill}}\parbox {6 cm}{\hfill Семестр 2}}\par
{\parbox {5 cm} {Навчальний предмет \hfill}\parbox {7 cm}{Програмування \hfill}\parbox {6cm}{\hfill}\par}\vspace{-7pt} \end{small}}{\center Екзаменаційний білет \No \textbf 
{ 465 }\par}}
{
\begin {large}
1. Що буде виведено за виконання?
code
try:
    res = 7/7*6*4
    if isinstance(res, bool):
        print(f'bool;{res}')
    elif isinstance(res, int):
        print(f'int;{res}')
    elif isinstance(res, float):
        print(f'float;{res:.2f}')
    else:
        print('???')
except:
    print('ERR')
code

2. Що буде виведено за виконання?
code
a = 3
b = 7
c = 0

def h(a):
    global c
    a += 1
    b = 2
    c = 4
    return a + b + c

a = 4
b = 8
c = 5

try:
    print(h(b), a, b, c, sep=';')
except:
    print('ERR', a, b, c, sep=';')
code

3. Що буде виведено за виконання?
code
try:
    i = 10
    while i>=7:
        i += -3
    print(i, end=';')
except:
    print('ERR')
code

4. Що буде виведено за виконання?
code
try:
    for a in range(11, 13, -1):
        if a>6:
            break
        print(a, end=';')
    else:
        print(a,end=';')
    print(a)
except:
    print('ERR')
code

5. Що буде виведено за виконання?\par (Дійсні значення, якщо наявні, в цьому завданні будуть виводитись у форматі з фіксованою крапкою та, можливо, з доволі великою кількістю десяткових розрядів. У відповіді для дійсних значень у ЦЬОМУ завданні достатньо вказати один десятковий розряд після крапки, відкинувши решту розрядів без округлення. Наприклад, замість 13.66666666666667 достатньо вказати 13.6, замість 25.23 достатньо вказати 25.2)
code
try:
    print(0, end=';')
    print(int('1'), end=';')
    print(6, end=';')
except ZeroDivisionError:
    print(2, end=';')
except ValueError:
    print(4, end=';')
except:
    print('ERR', end=';')
else:
    print(5, end=';')
finally:
    print(8)
code

6. Що буде виведено за виконання?
code
def f(n):
    if n>9:
        return f(n//100)+1
    else:
        return 1

print(f(123456))
code

7. Що буде виведено за виконання?
code
class A:
    b = 1

    def __g1(self): print(2, end=';')
    def _g2(self): print(3, end=';')

    def main(self):
        try: print(self.b, end=';')
        except: print('ERR', end=';')

        try: self.__g1()
        except: print('ERR', end=';')

        try: self._g2()
        except: print('ERR', end=';')


class B(A):
    b = 4

    def __g1(self): print(5, end=';')
    def _g2(self): print(6, end=';')

try: B().main()
except: print('ERR', end=';')
code

8. Що буде виведено за виконання?
code
def f(arg, n):
    arg[15] = 4
    n += 2

try:
    n = 6
    d = {73:4, 55:4, 21:1, 21:7}
    f(d, n)
    print(n, end=';')
    for x in d :
        print(x, sep=';', end=';')
except:
    print('ERR')
code

9. Що буде виведено за виконання?
code
class A:
    b = 1
    a = 2

    def __init__(self):
        a = 3

class B(A):
    a = 4

    def __init__(self):
        super().__init__()
        self.c = 7

obj = B()
obj.c = obj.b + 10
B.c = 13

try: print(obj.c, end= ';')
except: print('ERR', end=';')

try: print(A.c, end= ';')
except: print('ERR', end=';')

try: print(B.c, end= ';')
except: print('ERR', end=';')
code

10. 
Для графа на рисунку з вершини 5 запущено алгоритм пошуку в глибину, що будує кістякове дерево. Списки суміжності вершин переглядаються алгоритмом за зростанням міток вершин.\par
\begin{center}
	\adjustimage{max size={0.9\linewidth}{0.9\paperheight}}
	{../10/Traversals/179.png}
\end{center}
\par Які з наведених ребер входять до складу отриманого кістякового дерева?\par
1) (2, 3)\par
2) (3, 5)\par
3) (2, 5)\par
4) (1, 6)\par
5) (0, 4)\par
6) (0, 1)\par
{\vskip 15pt} У формі вибрати номери відповідних ребер.\par
%answer:1 3 4 5 6
\end {large}}
%end
{\smallskip \begin {small} Затверджено на засіданні кафедри теоретичної кібернетики, протокол \No 12 від   19.05.2020 р.\par\smallskip\smallskip
{\parbox {6.5 cm}{Зав. кафедри \dotfill Ю.В.~Крак}}\hspace {0.7cm}{\hbox to 6.5 cm{ Екзаменатор \dotfill Т.О.~Карнаух}} \end{small}}
%begin
{{\begin {small}\par
{ \center  Київський національний університет імені Тараса Шевченка \par}
{\parbox {5 cm}{Освітня програма \hfill}{\parbox {7 cm}{Прикладна математика, Системний аналіз \hfill}}\parbox {6 cm}{\hfill Семестр 2}}\par
{\parbox {5 cm} {Навчальний предмет \hfill}\parbox {7 cm}{Програмування \hfill}\parbox {6cm}{\hfill}\par}\vspace{-7pt} \end{small}}{\center Екзаменаційний білет \No \textbf 
{ 466 }\par}}
{
\begin {large}
1. Що буде виведено за виконання?
code
def f(n):
    print('f', end='')
    return n>=1

try:
    print(f(1) and f(8))
except:
    print('ERR')
code

2. Що буде виведено за виконання?
code
a = 9
b = 6
c = 5

def f(b):
    a -= 3
    b = 3
    c = 1
    return a + b + c

a = 7
b = 2
c = 6

try:
    print(f(a), a, b, c, sep=';')
except:
    print('ERR', a, b, c, sep=';')
code

3. Що буде виведено за виконання?
code
try:
    n = 5
    while n>3:
        n += -2
    print(n, end=';')
except:
    print('ERR')
code

4. Що буде виведено за виконання?
code
try:
    for b in range(10, 4, -2):
        if b<3:
            continue
        print(b, end=';')
        if b>=8:
            break
    else:
        print(b,end=';')
    print(b)
except:
    print('ERR')
code

5. Що буде виведено за виконання?\par (Дійсні значення, якщо наявні, в цьому завданні будуть виводитись у форматі з фіксованою крапкою та, можливо, з доволі великою кількістю десяткових розрядів. У відповіді для дійсних значень у ЦЬОМУ завданні достатньо вказати один десятковий розряд після крапки, відкинувши решту розрядів без округлення. Наприклад, замість 13.66666666666667 достатньо вказати 13.6, замість 25.23 достатньо вказати 25.2)
code
try:
    print(7, end=';')
    print(0//3, end=';')
    print(9, end=';')
except ZeroDivisionError:
    print(1, end=';')
except ValueError:
    print(3, end=';')
except:
    print('ERR', end=';')
else:
    print(1, end=';')
finally:
    print(8)
code

6. Що буде виведено за виконання?
code
def f(s):
    for i in range(len(s)):
        s[i] += 1
    return
 
s = [1, 2, 3]
f(s)
print(s)
code

7. Що буде виведено за виконання?
code
class A:
    c = 1

    def _f1(self): print(2, end=';')
    def __f2(self): print(3, end=';')

    def main(self):
        try: print(self.c, end=';')
        except: print('ERR', end=';')

        try: self._f1()
        except: print('ERR', end=';')

        try: self.__f2()
        except: print('ERR', end=';')


class B(A):
    c = 4

    def _f1(self): print(5, end=';')
    def __f2(self): print(6, end=';')

try: B().main()
except: print('ERR', end=';')
code

8. Що буде виведено за виконання?
code
def f(arg, n):
    arg[65] = 3
    n = 9
    arg = tuple(arg.values())

try:
    n = 3
    t = {11:8, 79:8, 14:3, 11:0}
    f(t, n)
    print(n, end=';')
    for x in t.values():
        print(x, sep=';', end=';')
except:
    print('ERR')
code

9. Що буде виведено за виконання?
code
class A:
    c = 1
    a = 2

    def __init__(self):
        c = 3

class B(A):
    c = 4

    def __init__(self):
        super().__init__()
        self.b = 7

obj = B()
obj.c = obj.b + 10
B.c = 13

try: print(obj.c, end= ';')
except: print('ERR', end=';')

try: print(A.c, end= ';')
except: print('ERR', end=';')

try: print(B.c, end= ';')
except: print('ERR', end=';')
code

10. 
Для графа на рисунку з вершини 4 запущено алгоритм пошуку в глибину, що будує кістякове дерево. Списки суміжності вершин переглядаються алгоритмом за зростанням міток вершин.\par
\begin{center}
	\adjustimage{max size={0.9\linewidth}{0.9\paperheight}}
	{../10/Traversals/149.png}
\end{center}
\par Які з наведених ребер входять до складу отриманого кістякового дерева?\par
1) (2, 6)\par
2) (3, 6)\par
3) (1, 3)\par
4) (2, 3)\par
5) (1, 4)\par
6) (0, 1)\par
{\vskip 15pt} У формі вибрати номери відповідних ребер.\par
%answer:1 3 4 5 6
\end {large}}
%end
{\smallskip \begin {small} Затверджено на засіданні кафедри теоретичної кібернетики, протокол \No 12 від   19.05.2020 р.\par\smallskip\smallskip
{\parbox {6.5 cm}{Зав. кафедри \dotfill Ю.В.~Крак}}\hspace {0.7cm}{\hbox to 6.5 cm{ Екзаменатор \dotfill Т.О.~Карнаух}} \end{small}}
%begin
{{\begin {small}\par
{ \center  Київський національний університет імені Тараса Шевченка \par}
{\parbox {5 cm}{Освітня програма \hfill}{\parbox {7 cm}{Прикладна математика, Системний аналіз \hfill}}\parbox {6 cm}{\hfill Семестр 2}}\par
{\parbox {5 cm} {Навчальний предмет \hfill}\parbox {7 cm}{Програмування \hfill}\parbox {6cm}{\hfill}\par}\vspace{-7pt} \end{small}}{\center Екзаменаційний білет \No \textbf 
{ 467 }\par}}
{
\begin {large}
1. Що буде виведено за виконання?
code
try:
    res = 6/8/4*5
    if isinstance(res, bool):
        print(f'bool;{res}')
    elif isinstance(res, int):
        print(f'int;{res}')
    elif isinstance(res, float):
        print(f'float;{res:.2f}')
    else:
        print('???')
except:
    print('ERR')
code

2. Що буде виведено за виконання?
code
a = 7
b = 1
c = 6

def f(a):
    a = 2
    b = 4
    c = 1
    return a + b + c

a = 2
b = 0
c = 3

try:
    print(f(a), a, b, c, sep=';')
except:
    print('ERR', a, b, c, sep=';')
code

3. Що буде виведено за виконання?
code
try:
    j = 14
    while j>=0:
        j += -2
    print(j, end=';')
except:
    print('ERR')
code

4. Що буде виведено за виконання?
code
try:
    for d in range(15, 7, 2):
        if d>=7:
            continue
        print(d, end=';')
    else:
        print(d,end=';')
    print(d)
except:
    print('ERR')
code

5. Що буде виведено за виконання?\par (Дійсні значення, якщо наявні, в цьому завданні будуть виводитись у форматі з фіксованою крапкою та, можливо, з доволі великою кількістю десяткових розрядів. У відповіді для дійсних значень у ЦЬОМУ завданні достатньо вказати один десятковий розряд після крапки, відкинувши решту розрядів без округлення. Наприклад, замість 13.66666666666667 достатньо вказати 13.6, замість 25.23 достатньо вказати 25.2)
code
try:
    print(3, end=';')
    print(0j!=6j, end=';')
    print(4, end=';')
except KeyboardInterrupt:
    print(7, end=';')
except Exception:
    print(2, end=';')
except:
    print('ERR', end=';')
else:
    print(6, end=';')
finally:
    print(9)
code

6. Що буде виведено за виконання?
code
def f(n):
    print(n%10, end='')
    if n>9:
        f(n//100)
 
f(987654)
code

7. Що буде виведено за виконання?
code
class A:
    a = 1

    def __g1(self): print(2, end=';')
    def _g2(self): print(3, end=';')

    def main(self):
        try: print(self.a, end=';')
        except: print('ERR', end=';')

        try: self.__g1()
        except: print('ERR', end=';')

        try: self._g2()
        except: print('ERR', end=';')


class B(A):
    a = 4

    def __g1(self): print(5, end=';')
    def _g2(self): print(6, end=';')

try: B().main()
except: print('ERR', end=';')
code

8. Що буде виведено за виконання?
code
def f(arg, n):
    arg[63] = 5
    n = 5
    arg = set(arg.keys())

try:
    n = 0
    d = {75:2, 88:3, 75:7}
    f(d, n)
    print(n, end=';')
    for x, y in d.items():
        print(x, y, sep=';', end=';')
except:
    print('ERR')
code

9. Що буде виведено за виконання?
code
class A:
    b = 1
    c = 2

    def __init__(self):
        b = 3

class B(A):
    b = 4

    def __init__(self):
        super().__init__()
        self.a = 7

obj = B()
obj.a = obj.c + 10
B.a = 13

try: print(obj.a, end= ';')
except: print('ERR', end=';')

try: print(A.a, end= ';')
except: print('ERR', end=';')

try: print(B.a, end= ';')
except: print('ERR', end=';')
code

10. 
Для графа на рисунку з вершини 2 запущено алгоритм пошуку в глибину, що будує кістякове дерево. Списки суміжності вершин переглядаються алгоритмом за зростанням міток вершин.\par
\begin{center}
	\adjustimage{max size={0.9\linewidth}{0.9\paperheight}}
	{../10/Traversals/268.png}
\end{center}
\par Які з наведених ребер входять до складу отриманого кістякового дерева?\par
1) (1, 4)\par
2) (0, 3)\par
3) (2, 3)\par
4) (1, 3)\par
5) (3, 4)\par
6) (4, 5)\par
{\vskip 15pt} У формі вибрати номери відповідних ребер.\par
%answer:1 2 3 4 6
\end {large}}
%end
{\smallskip \begin {small} Затверджено на засіданні кафедри теоретичної кібернетики, протокол \No 12 від   19.05.2020 р.\par\smallskip\smallskip
{\parbox {6.5 cm}{Зав. кафедри \dotfill Ю.В.~Крак}}\hspace {0.7cm}{\hbox to 6.5 cm{ Екзаменатор \dotfill Т.О.~Карнаух}} \end{small}}
%begin
{{\begin {small}\par
{ \center  Київський національний університет імені Тараса Шевченка \par}
{\parbox {5 cm}{Освітня програма \hfill}{\parbox {7 cm}{Прикладна математика, Системний аналіз \hfill}}\parbox {6 cm}{\hfill Семестр 2}}\par
{\parbox {5 cm} {Навчальний предмет \hfill}\parbox {7 cm}{Програмування \hfill}\parbox {6cm}{\hfill}\par}\vspace{-7pt} \end{small}}{\center Екзаменаційний білет \No \textbf 
{ 468 }\par}}
{
\begin {large}
1. Що буде виведено за виконання?
code
try:
    res = 1**1.0*1
    if isinstance(res, bool):
        print(f'bool;{res}')
    elif isinstance(res, int):
        print(f'int;{res}')
    elif isinstance(res, float):
        print(f'float;{res:.2f}')
    else:
        print('???')
except:
    print('ERR')
code

2. Що буде виведено за виконання?
code
a = 8
b = 1
c = 9

def f(b):
    a = 5
    b += 2
    c = 4
    return a + b + c

a = 0
b = 4
c = 3

try:
    print(f(b), a, b, c, sep=';')
except:
    print('ERR', a, b, c, sep=';')
code

3. Що буде виведено за виконання?
code
try:
    n = 0
    while n<8:
        n += 1
    print(n, end=';')
except:
    print('ERR')
code

4. Що буде виведено за виконання?
code
try:
    for e in range(-4, 15, -1):
        if e<=8:
            break
        print(e, end=';');e = 6
    else:
        print(e,end=';')
    print(e)
except:
    print('ERR')
code

5. Що буде виведено за виконання?\par (Дійсні значення, якщо наявні, в цьому завданні будуть виводитись у форматі з фіксованою крапкою та, можливо, з доволі великою кількістю десяткових розрядів. У відповіді для дійсних значень у ЦЬОМУ завданні достатньо вказати один десятковий розряд після крапки, відкинувши решту розрядів без округлення. Наприклад, замість 13.66666666666667 достатньо вказати 13.6, замість 25.23 достатньо вказати 25.2)
code
try:
    print(5, end=';')
    print(int('d2'), end=';')
    print(4, end=';')
except TypeError:
    print(7, end=';')
except KeyboardInterrupt:
    print(9, end=';')
except:
    print('ERR', end=';')
else:
    print(3, end=';')
finally:
    print(0)
code

6. Що буде виведено за виконання?
code
def f(s):
    for el in s:
        el += 1
        return s
 
s = [1, 2, 3]
s = f(s)
print(s)
code

7. Що буде виведено за виконання?
code
class A:
    d = 1

    def _h1(self): print(2, end=';')
    def h2(self): print(3, end=';')

    def main(self):
        try: print(self.d, end=';')
        except: print('ERR', end=';')

        try: self._h1()
        except: print('ERR', end=';')

        try: self.h2()
        except: print('ERR', end=';')


class B(A):
    d = 4

    def _h1(self): print(5, end=';')
    def h2(self): print(6, end=';')

try: B().main()
except: print('ERR', end=';')
code

8. Що буде виведено за виконання?
code
def f(arg, n):
    arg[87] = 5
    n = 8
    arg = list(arg.values())

try:
    n = 2
    d = {87:7, 39:5, 88:1}
    f(d, n)
    print(n, end=';')
    for x in d.items():
        print(x, sep=';', end=';')
except:
    print('ERR')
code

9. Що буде виведено за виконання?
code
class A:
    c = 1
    b = 2

    def __init__(self):
        b = 3

class B(A):
    b = 4

    def __init__(self):
        super().__init__()
        self.a = 7

obj = B()
obj.b = obj.a + 10
B.b = 13

try: print(obj.b, end= ';')
except: print('ERR', end=';')

try: print(A.b, end= ';')
except: print('ERR', end=';')

try: print(B.b, end= ';')
except: print('ERR', end=';')
code

10. 
Для графа на рисунку з вершини 0 запущено алгоритм пошуку в глибину, що будує кістякове дерево. Списки суміжності вершин переглядаються алгоритмом за зростанням міток вершин.\par
\begin{center}
	\adjustimage{max size={0.9\linewidth}{0.9\paperheight}}
	{../10/Traversals/287.png}
\end{center}
\par Які з наведених ребер входять до складу отриманого кістякового дерева?\par
1) (0, 1)\par
2) (0, 2)\par
3) (0, 6)\par
4) (5, 6)\par
5) (4, 5)\par
6) (1, 5)\par
{\vskip 15pt} У формі вибрати номери відповідних ребер.\par
%answer:1 4 5
\end {large}}
%end
{\smallskip \begin {small} Затверджено на засіданні кафедри теоретичної кібернетики, протокол \No 12 від   19.05.2020 р.\par\smallskip\smallskip
{\parbox {6.5 cm}{Зав. кафедри \dotfill Ю.В.~Крак}}\hspace {0.7cm}{\hbox to 6.5 cm{ Екзаменатор \dotfill Т.О.~Карнаух}} \end{small}}
%begin
{{\begin {small}\par
{ \center  Київський національний університет імені Тараса Шевченка \par}
{\parbox {5 cm}{Освітня програма \hfill}{\parbox {7 cm}{Прикладна математика, Системний аналіз \hfill}}\parbox {6 cm}{\hfill Семестр 2}}\par
{\parbox {5 cm} {Навчальний предмет \hfill}\parbox {7 cm}{Програмування \hfill}\parbox {6cm}{\hfill}\par}\vspace{-7pt} \end{small}}{\center Екзаменаційний білет \No \textbf 
{ 469 }\par}}
{
\begin {large}
1. Що буде виведено за виконання?
code
try:
    res = 9**-1.0*-1
    if isinstance(res, bool):
        print(f'bool;{res}')
    elif isinstance(res, int):
        print(f'int;{res}')
    elif isinstance(res, float):
        print(f'float;{res:.2f}')
    else:
        print('???')
except:
    print('ERR')
code

2. Що буде виведено за виконання?
code
a = 9
b = 5
c = 5

def h(b):
    global c
    a = 5
    b = 3
    c = 4
    return a + b + c

a = 4
b = 3
c = 2

try:
    print(h(b), a, b, c, sep=';')
except:
    print('ERR', a, b, c, sep=';')
code

3. Що буде виведено за виконання?
code
try:
    i = 3
    while i<=6:
        i += 2
    print(i, end=';')
except:
    print('ERR')
code

4. Що буде виведено за виконання?
code
try:
    for c in range(-9, 13, 3):
        if c>=7:
            continue
        print(c, end=';');c = -4
    else:
        print(c,end=';')
    print(c)
except:
    print('ERR')
code

5. Що буде виведено за виконання?\par (Дійсні значення, якщо наявні, в цьому завданні будуть виводитись у форматі з фіксованою крапкою та, можливо, з доволі великою кількістю десяткових розрядів. У відповіді для дійсних значень у ЦЬОМУ завданні достатньо вказати один десятковий розряд після крапки, відкинувши решту розрядів без округлення. Наприклад, замість 13.66666666666667 достатньо вказати 13.6, замість 25.23 достатньо вказати 25.2)
code
try:
    print(6, end=';')
    print(8/0.0, end=';')
    print(5, end=';')
except TypeError:
    print(1, end=';')
except ZeroDivisionError:
    print(5, end=';')
except:
    print('ERR', end=';')
else:
    print(1, end=';')
finally:
    print(0)
code

6. Що буде виведено за виконання?
code
def f(n):
    if n>9:
        return f(n//10)+1
    else:
        return 1

print(f(123456))
code

7. Що буде виведено за виконання?
code
class A:
    d = 1

    def f1(self): print(2, end=';')
    def _f2(self): print(3, end=';')

    def main(self):
        try: print(self.d, end=';')
        except: print('ERR', end=';')

        try: self.f1()
        except: print('ERR', end=';')

        try: self._f2()
        except: print('ERR', end=';')


class B(A):
    d = 4

    def f1(self): print(5, end=';')
    def _f2(self): print(6, end=';')

try: B().main()
except: print('ERR', end=';')
code

8. Що буде виведено за виконання?
code
def f(arg, n):
    arg[17] = 0
    n *= -3

try:
    n = 5
    s = {46:4, 59:1, 25:8, 25:2}
    f(s, n)
    print(n, end=';')
    for x in s.keys():
        print(x, sep=';', end=';')
except:
    print('ERR')
code

9. Що буде виведено за виконання?
code
class A:
    a = 1
    c = 2

    def __init__(self):
        a = 3

class B(A):
    a = 4

    def __init__(self):
        super().__init__()
        self.b = 7

obj = B()
obj.b = obj.a + 10
B.b = 13

try: print(obj.b, end= ';')
except: print('ERR', end=';')

try: print(A.b, end= ';')
except: print('ERR', end=';')

try: print(B.b, end= ';')
except: print('ERR', end=';')
code

10. 
Для графа на рисунку з вершини 3 запущено алгоритм пошуку в глибину, що будує кістякове дерево. Списки суміжності вершин переглядаються алгоритмом за зростанням міток вершин.\par
\begin{center}
	\adjustimage{max size={0.9\linewidth}{0.9\paperheight}}
	{../10/Traversals/192.png}
\end{center}
\par Які з наведених ребер входять до складу отриманого кістякового дерева?\par
1) (3, 5)\par
2) (1, 4)\par
3) (1, 2)\par
4) (3, 4)\par
5) (1, 5)\par
6) (0, 1)\par
{\vskip 15pt} У формі вибрати номери відповідних ребер.\par
%answer:2 6
\end {large}}
%end
{\smallskip \begin {small} Затверджено на засіданні кафедри теоретичної кібернетики, протокол \No 12 від   19.05.2020 р.\par\smallskip\smallskip
{\parbox {6.5 cm}{Зав. кафедри \dotfill Ю.В.~Крак}}\hspace {0.7cm}{\hbox to 6.5 cm{ Екзаменатор \dotfill Т.О.~Карнаух}} \end{small}}
%begin
{{\begin {small}\par
{ \center  Київський національний університет імені Тараса Шевченка \par}
{\parbox {5 cm}{Освітня програма \hfill}{\parbox {7 cm}{Прикладна математика, Системний аналіз \hfill}}\parbox {6 cm}{\hfill Семестр 2}}\par
{\parbox {5 cm} {Навчальний предмет \hfill}\parbox {7 cm}{Програмування \hfill}\parbox {6cm}{\hfill}\par}\vspace{-7pt} \end{small}}{\center Екзаменаційний білет \No \textbf 
{ 470 }\par}}
{
\begin {large}
1. Що буде виведено за виконання?
code
def f(n):
    print('f', end='')
    return n

try:
    print(f(0) or f(-3))
except:
    print('ERR')
code

2. Що буде виведено за виконання?
code
a = 1
b = 8
c = 0

def h(b):
    global c
    a = 2
    b = 1
    c = 5
    return a + b + c

a = 6
b = 7
c = 9

try:
    print(h(b), a, b, c, sep=';')
except:
    print('ERR', a, b, c, sep=';')
code

3. Що буде виведено за виконання?
code
try:
    t = 11
    while t>9:
        t += -3
    print(t, end=';')
except:
    print('ERR')
code

4. Що буде виведено за виконання?
code
try:
    for f in range(-13, 2, -3):
        if f<3:
            continue
        print(f, end=';')
    else:
        print('end',end=';')
    print(f)
except:
    print('ERR')
code

5. Що буде виведено за виконання?\par (Дійсні значення, якщо наявні, в цьому завданні будуть виводитись у форматі з фіксованою крапкою та, можливо, з доволі великою кількістю десяткових розрядів. У відповіді для дійсних значень у ЦЬОМУ завданні достатньо вказати один десятковий розряд після крапки, відкинувши решту розрядів без округлення. Наприклад, замість 13.66666666666667 достатньо вказати 13.6, замість 25.23 достатньо вказати 25.2)
code
try:
    print(6, end=';')
    print(int('3'), end=';')
    print(8, end=';')
except ValueError:
    print(7, end=';')
except Exception:
    print(2, end=';')
except:
    print('ERR', end=';')
else:
    print(9, end=';')
finally:
    print(4)
code

6. Що буде виведено за виконання?
code
def f(n):
    print(n%10, end='')
    if n>9:
        f(n//10)
 
f(123456)
code

7. Що буде виведено за виконання?
code
class A:
    c = 1

    def _g1(self): print(2, end=';')
    def _g2(self): print(3, end=';')

    def main(self):
        try: print(self.c, end=';')
        except: print('ERR', end=';')

        try: self._g1()
        except: print('ERR', end=';')

        try: self._g2()
        except: print('ERR', end=';')


class B(A):
    c = 4

    def _g1(self): print(5, end=';')
    def _g2(self): print(6, end=';')

try: B().main()
except: print('ERR', end=';')
code

8. Що буде виведено за виконання?
code
def f(arg, n):
    arg[23] = 0
    n = 7
    arg = tuple(arg.items())

try:
    n = 4
    s = {27:2, 55:3, 82:7, 82:1}
    f(s, n)
    print(n, end=';')
    for x in s.keys():
        print(x, sep=';', end=';')
except:
    print('ERR')
code

9. Що буде виведено за виконання?
code
class A:
    b = 1
    c = 2

    def __init__(self):
        c = 3

class B(A):
    c = 4

    def __init__(self):
        super().__init__()
        self.a = 7

obj = B()
obj.c = obj.a + 10
B.c = 13

try: print(obj.c, end= ';')
except: print('ERR', end=';')

try: print(A.c, end= ';')
except: print('ERR', end=';')

try: print(B.c, end= ';')
except: print('ERR', end=';')
code

10. 
Для графа на рисунку з вершини 4 запущено алгоритм пошуку в глибину, що будує кістякове дерево. Списки суміжності вершин переглядаються алгоритмом за зростанням міток вершин.\par
\begin{center}
	\adjustimage{max size={0.9\linewidth}{0.9\paperheight}}
	{../10/Traversals/83.png}
\end{center}
\par Які з наведених ребер входять до складу отриманого кістякового дерева?\par
1) (1, 4)\par
2) (4, 5)\par
3) (5, 6)\par
4) (2, 4)\par
5) (1, 6)\par
6) (3, 4)\par
{\vskip 15pt} У формі вибрати номери відповідних ребер.\par
%answer:3 5
\end {large}}
%end
{\smallskip \begin {small} Затверджено на засіданні кафедри теоретичної кібернетики, протокол \No 12 від   19.05.2020 р.\par\smallskip\smallskip
{\parbox {6.5 cm}{Зав. кафедри \dotfill Ю.В.~Крак}}\hspace {0.7cm}{\hbox to 6.5 cm{ Екзаменатор \dotfill Т.О.~Карнаух}} \end{small}}
%begin
{{\begin {small}\par
{ \center  Київський національний університет імені Тараса Шевченка \par}
{\parbox {5 cm}{Освітня програма \hfill}{\parbox {7 cm}{Прикладна математика, Системний аналіз \hfill}}\parbox {6 cm}{\hfill Семестр 2}}\par
{\parbox {5 cm} {Навчальний предмет \hfill}\parbox {7 cm}{Програмування \hfill}\parbox {6cm}{\hfill}\par}\vspace{-7pt} \end{small}}{\center Екзаменаційний білет \No \textbf 
{ 471 }\par}}
{
\begin {large}
1. Що буде виведено за виконання?
code
try:
    res = 9//6//5*8
    if isinstance(res, bool):
        print(f'bool;{res}')
    elif isinstance(res, int):
        print(f'int;{res}')
    elif isinstance(res, float):
        print(f'float;{res:.2f}')
    else:
        print('???')
except:
    print('ERR')
code

2. Що буде виведено за виконання?
code
a = 2
b = 3
c = 4

def h(a):
    global c
    a -= 3
    b = 4
    c = 2
    return a + b + c

a = 1
b = 5
c = 9

try:
    print(h(b), a, b, c, sep=';')
except:
    print('ERR', a, b, c, sep=';')
code

3. Що буде виведено за виконання?
code
try:
    k = 5
    while k>7:
        k += -2
    print(k, end=';')
except:
    print('ERR')
code

4. Що буде виведено за виконання?
code
try:
    for b in range(-12, 6, -3):
        if b>8:
            break
        print(b, end=';')
    else:
        print(b,end=';')
    print(b)
except:
    print('ERR')
code

5. Що буде виведено за виконання?\par (Дійсні значення, якщо наявні, в цьому завданні будуть виводитись у форматі з фіксованою крапкою та, можливо, з доволі великою кількістю десяткових розрядів. У відповіді для дійсних значень у ЦЬОМУ завданні достатньо вказати один десятковий розряд після крапки, відкинувши решту розрядів без округлення. Наприклад, замість 13.66666666666667 достатньо вказати 13.6, замість 25.23 достатньо вказати 25.2)
code
try:
    print(8, end=';')
    print(int('a6'), end=';')
    print(4, end=';')
except TypeError:
    print(3, end=';')
except ValueError:
    print(2, end=';')
except:
    print('ERR', end=';')
else:
    print(7, end=';')
finally:
    print(0)
code

6. Що буде виведено за виконання?
code
def f(n):
    if n>9: f(n//10)
    else:
        print(n%10, end='')
 
f(123456)
code

7. Що буде виведено за виконання?
code
class A:
    b = 1

    def h1(self): print(2, end=';')
    def _h2(self): print(3, end=';')

    def main(self):
        try: print(self.b, end=';')
        except: print('ERR', end=';')

        try: self.h1()
        except: print('ERR', end=';')

        try: self._h2()
        except: print('ERR', end=';')


class B(A):
    b = 4

    def h1(self): print(5, end=';')
    def _h2(self): print(6, end=';')

try: B().main()
except: print('ERR', end=';')
code

8. Що буде виведено за виконання?
code
def f(arg, n):
    arg[59] = 2
    n -= 3

try:
    n = 4
    t = {39:2, 27:1, 66:5, 27:5}
    f(t, n)
    print(n, end=';')
    for x in t.values():
        print(x, sep=';', end=';')
except:
    print('ERR')
code

9. Що буде виведено за виконання?
code
class A:
    b = 1
    a = 2

    def __init__(self):
        a = 3

class B(A):
    a = 4

    def __init__(self):
        super().__init__()
        self.c = 7

obj = B()
obj.b = obj.c + 10
B.b = 13

try: print(obj.b, end= ';')
except: print('ERR', end=';')

try: print(A.b, end= ';')
except: print('ERR', end=';')

try: print(B.b, end= ';')
except: print('ERR', end=';')
code

10. 
Для графа на рисунку з вершини 0 запущено алгоритм пошуку в глибину, що будує кістякове дерево. Списки суміжності вершин переглядаються алгоритмом за зростанням міток вершин.\par
\begin{center}
	\adjustimage{max size={0.9\linewidth}{0.9\paperheight}}
	{../10/Traversals/240.png}
\end{center}
\par Які з наведених ребер входять до складу отриманого кістякового дерева?\par
1) (0, 6)\par
2) (2, 3)\par
3) (4, 5)\par
4) (2, 4)\par
5) (5, 6)\par
6) (1, 4)\par
{\vskip 15pt} У формі вибрати номери відповідних ребер.\par
%answer:2 3 4 5 6
\end {large}}
%end
{\smallskip \begin {small} Затверджено на засіданні кафедри теоретичної кібернетики, протокол \No 12 від   19.05.2020 р.\par\smallskip\smallskip
{\parbox {6.5 cm}{Зав. кафедри \dotfill Ю.В.~Крак}}\hspace {0.7cm}{\hbox to 6.5 cm{ Екзаменатор \dotfill Т.О.~Карнаух}} \end{small}}
%begin
{{\begin {small}\par
{ \center  Київський національний університет імені Тараса Шевченка \par}
{\parbox {5 cm}{Освітня програма \hfill}{\parbox {7 cm}{Прикладна математика, Системний аналіз \hfill}}\parbox {6 cm}{\hfill Семестр 2}}\par
{\parbox {5 cm} {Навчальний предмет \hfill}\parbox {7 cm}{Програмування \hfill}\parbox {6cm}{\hfill}\par}\vspace{-7pt} \end{small}}{\center Екзаменаційний білет \No \textbf 
{ 472 }\par}}
{
\begin {large}
1. Що буде виведено за виконання?
code
try:
    res = 5//3%7*9
    if isinstance(res, bool):
        print(f'bool;{res}')
    elif isinstance(res, int):
        print(f'int;{res}')
    elif isinstance(res, float):
        print(f'float;{res:.2f}')
    else:
        print('???')
except:
    print('ERR')
code

2. Що буде виведено за виконання?
code
a = 8
b = 0
c = 6

def g(a):
    global c
    a *= 1
    b = 5
    c = 2
    return a + b + c

a = 7
b = 0
c = 4

try:
    print(g(a), a, b, c, sep=';')
except:
    print('ERR', a, b, c, sep=';')
code

3. Що буде виведено за виконання?
code
try:
    m = 12
    while m>=1:
        m += -3
    print(m, end=';')
except:
    print('ERR')
code

4. Що буде виведено за виконання?
code
try:
    for a in range(-15, 1, -2):
        if a<=4:
            continue
        print(a, end=';');a = 3
    else:
        print(a,end=';')
    print(a)
except:
    print('ERR')
code

5. Що буде виведено за виконання?\par (Дійсні значення, якщо наявні, в цьому завданні будуть виводитись у форматі з фіксованою крапкою та, можливо, з доволі великою кількістю десяткових розрядів. У відповіді для дійсних значень у ЦЬОМУ завданні достатньо вказати один десятковий розряд після крапки, відкинувши решту розрядів без округлення. Наприклад, замість 13.66666666666667 достатньо вказати 13.6, замість 25.23 достатньо вказати 25.2)
code
try:
    print(9, end=';')
    print(1%0.0, end=';')
    print(5, end=';')
except Exception:
    print(3, end=';')
except ZeroDivisionError:
    print(7, end=';')
except:
    print('ERR', end=';')
else:
    print(5, end=';')
finally:
    print(2)
code

6. Що буде виведено за виконання?
code
def f(n):
    if n<=9:
        print(n%10, end='')
    else:
        f(n//10)

f(123456)
code

7. Що буде виведено за виконання?
code
class A:
    a = 1

    def _f1(self): print(2, end=';')
    def __f2(self): print(3, end=';')

    def main(self):
        try: print(self.a, end=';')
        except: print('ERR', end=';')

        try: self._f1()
        except: print('ERR', end=';')

        try: self.__f2()
        except: print('ERR', end=';')


class B(A):
    a = 4

    def _f1(self): print(5, end=';')
    def __f2(self): print(6, end=';')

try: B().main()
except: print('ERR', end=';')
code

8. Що буде виведено за виконання?
code
def f(arg, n):
    arg[33] = 0
    n += -2

try:
    n = 5
    t = {73:5, 33:5, 64:8, 64:9}
    f(t, n)
    print(n, end=';')
    for x in t.values():
        print(x, sep=';', end=';')
except:
    print('ERR')
code

9. Що буде виведено за виконання?
code
class A:
    c = 1
    a = 2

    def __init__(self):
        c = 3

class B(A):
    c = 4

    def __init__(self):
        super().__init__()
        self.b = 7

obj = B()
obj.a = obj.b + 10
B.a = 13

try: print(obj.a, end= ';')
except: print('ERR', end=';')

try: print(A.a, end= ';')
except: print('ERR', end=';')

try: print(B.a, end= ';')
except: print('ERR', end=';')
code

10. 
Для графа на рисунку з вершини 6 запущено алгоритм пошуку в глибину, що будує кістякове дерево. Списки суміжності вершин переглядаються алгоритмом за зростанням міток вершин.\par
\begin{center}
	\adjustimage{max size={0.9\linewidth}{0.9\paperheight}}
	{../10/Traversals/89.png}
\end{center}
\par Які з наведених ребер входять до складу отриманого кістякового дерева?\par
1) (0, 5)\par
2) (4, 6)\par
3) (3, 5)\par
4) (2, 4)\par
5) (0, 4)\par
6) (1, 2)\par
{\vskip 15pt} У формі вибрати номери відповідних ребер.\par
%answer:3 5
\end {large}}
%end
{\smallskip \begin {small} Затверджено на засіданні кафедри теоретичної кібернетики, протокол \No 12 від   19.05.2020 р.\par\smallskip\smallskip
{\parbox {6.5 cm}{Зав. кафедри \dotfill Ю.В.~Крак}}\hspace {0.7cm}{\hbox to 6.5 cm{ Екзаменатор \dotfill Т.О.~Карнаух}} \end{small}}
%begin
{{\begin {small}\par
{ \center  Київський національний університет імені Тараса Шевченка \par}
{\parbox {5 cm}{Освітня програма \hfill}{\parbox {7 cm}{Прикладна математика, Системний аналіз \hfill}}\parbox {6 cm}{\hfill Семестр 2}}\par
{\parbox {5 cm} {Навчальний предмет \hfill}\parbox {7 cm}{Програмування \hfill}\parbox {6cm}{\hfill}\par}\vspace{-7pt} \end{small}}{\center Екзаменаційний білет \No \textbf 
{ 473 }\par}}
{
\begin {large}
1. Що буде виведено за виконання?
code
try:
    res = 4**0.5*2
    if isinstance(res, bool):
        print(f'bool;{res}')
    elif isinstance(res, int):
        print(f'int;{res}')
    elif isinstance(res, float):
        print(f'float;{res:.2f}')
    else:
        print('???')
except:
    print('ERR')
code

2. Що буде виведено за виконання?
code
a = 3
b = 7
c = 2

def f(b):
    a *= 1
    b = 5
    c = 4
    return a + b + c

a = 1
b = 6
c = 8

try:
    print(f(a), a, b, c, sep=';')
except:
    print('ERR', a, b, c, sep=';')
code

3. Що буде виведено за виконання?
code
try:
    m = 2
    while m<=12:
        m += 2
    print(m, end=';')
except:
    print('ERR')
code

4. Що буде виведено за виконання?
code
try:
    for d in range(12, 10, 3):
        if d>6:
            break
        print(d, end=';')
    else:
        print(d,end=';')
    print(d)
except:
    print('ERR')
code

5. Що буде виведено за виконання?\par (Дійсні значення, якщо наявні, в цьому завданні будуть виводитись у форматі з фіксованою крапкою та, можливо, з доволі великою кількістю десяткових розрядів. У відповіді для дійсних значень у ЦЬОМУ завданні достатньо вказати один десятковий розряд після крапки, відкинувши решту розрядів без округлення. Наприклад, замість 13.66666666666667 достатньо вказати 13.6, замість 25.23 достатньо вказати 25.2)
code
try:
    print(1, end=';')
    print(int('6'), end=';')
    print(0, end=';')
except KeyboardInterrupt:
    print(4, end=';')
except Exception:
    print(9, end=';')
except:
    print('ERR', end=';')
else:
    print(8, end=';')
finally:
    print(1)
code

6. Що буде виведено за виконання?
code
def f(n):
    if n<=9:
        print(n%10,end='')
    else:
        f(n//10)

f(543216)
code

7. Що буде виведено за виконання?
code
class A:
    b = 1

    def _g1(self): print(2, end=';')
    def _g2(self): print(3, end=';')

    def main(self):
        try: print(self.b, end=';')
        except: print('ERR', end=';')

        try: self._g1()
        except: print('ERR', end=';')

        try: self._g2()
        except: print('ERR', end=';')


class B(A):
    b = 4

    def _g1(self): print(5, end=';')
    def _g2(self): print(6, end=';')

try: B().main()
except: print('ERR', end=';')
code

8. Що буде виведено за виконання?
code
def f(arg, n):
    arg[50] = 3
    n -= -3

try:
    n = 6
    d = {85:4, 41:9, 36:8, 41:5}
    f(d, n)
    print(n, end=';')
    for x, y in d.items():
        print(x, y, sep=';', end=';')
except:
    print('ERR')
code

9. Що буде виведено за виконання?
code
class A:
    c = 1
    b = 2

    def __init__(self):
        c = 3

class B(A):
    b = 4

    def __init__(self):
        super().__init__()
        self.a = 7

obj = B()
obj.c = obj.a + 10
B.c = 13

try: print(obj.c, end= ';')
except: print('ERR', end=';')

try: print(A.c, end= ';')
except: print('ERR', end=';')

try: print(B.c, end= ';')
except: print('ERR', end=';')
code

10. 
Для графа на рисунку з вершини 3 запущено алгоритм пошуку в глибину, що будує кістякове дерево. Списки суміжності вершин переглядаються алгоритмом за зростанням міток вершин.\par
\begin{center}
	\adjustimage{max size={0.9\linewidth}{0.9\paperheight}}
	{../10/Traversals/123.png}
\end{center}
\par Які з наведених ребер входять до складу отриманого кістякового дерева?\par
1) (3, 4)\par
2) (1, 3)\par
3) (1, 5)\par
4) (2, 5)\par
5) (0, 6)\par
6) (1, 2)\par
{\vskip 15pt} У формі вибрати номери відповідних ребер.\par
%answer:1 2 4 5 6
\end {large}}
%end
{\smallskip \begin {small} Затверджено на засіданні кафедри теоретичної кібернетики, протокол \No 12 від   19.05.2020 р.\par\smallskip\smallskip
{\parbox {6.5 cm}{Зав. кафедри \dotfill Ю.В.~Крак}}\hspace {0.7cm}{\hbox to 6.5 cm{ Екзаменатор \dotfill Т.О.~Карнаух}} \end{small}}
%begin
{{\begin {small}\par
{ \center  Київський національний університет імені Тараса Шевченка \par}
{\parbox {5 cm}{Освітня програма \hfill}{\parbox {7 cm}{Прикладна математика, Системний аналіз \hfill}}\parbox {6 cm}{\hfill Семестр 2}}\par
{\parbox {5 cm} {Навчальний предмет \hfill}\parbox {7 cm}{Програмування \hfill}\parbox {6cm}{\hfill}\par}\vspace{-7pt} \end{small}}{\center Екзаменаційний білет \No \textbf 
{ 474 }\par}}
{
\begin {large}
1. Що буде виведено за виконання?
code
try:
    res = 3**1.0+-2
    if isinstance(res, bool):
        print(f'bool;{res}')
    elif isinstance(res, int):
        print(f'int;{res}')
    elif isinstance(res, float):
        print(f'float;{res:.2f}')
    else:
        print('???')
except:
    print('ERR')
code

2. Що буде виведено за виконання?
code
a = 5
b = 3
c = 7

def f(a):
    a = 4
    b = 2
    c = 1
    return a + b + c

a = 1
b = 6
c = 0

try:
    print(f(b), a, b, c, sep=';')
except:
    print('ERR', a, b, c, sep=';')
code

3. Що буде виведено за виконання?
code
try:
    j = 9
    while j>=2:
        j += -2
    print(j, end=';')
except:
    print('ERR')
code

4. Що буде виведено за виконання?
code
try:
    for e in range(-11, -13, 2):
        if e>=5:
            continue
        print(e, end=';')
    else:
        print(e,end=';')
    print(e)
except:
    print('ERR')
code

5. Що буде виведено за виконання?\par (Дійсні значення, якщо наявні, в цьому завданні будуть виводитись у форматі з фіксованою крапкою та, можливо, з доволі великою кількістю десяткових розрядів. У відповіді для дійсних значень у ЦЬОМУ завданні достатньо вказати один десятковий розряд після крапки, відкинувши решту розрядів без округлення. Наприклад, замість 13.66666666666667 достатньо вказати 13.6, замість 25.23 достатньо вказати 25.2)
code
try:
    print(8, end=';')
    print(6/2, end=';')
    print(5, end=';')
except ValueError:
    print(2, end=';')
except ZeroDivisionError:
    print(0, end=';')
except:
    print('ERR', end=';')
else:
    print(7, end=';')
finally:
    print(4)
code

6. Що буде виведено за виконання?
code
def f(n):
    if n>2:
        f(n//2)
    print(n%2, end='')
 
f(16)
code

7. Що буде виведено за виконання?
code
class A:
    c = 1

    def _h1(self): print(2, end=';')
    def h2(self): print(3, end=';')

    def main(self):
        try: print(self.c, end=';')
        except: print('ERR', end=';')

        try: self._h1()
        except: print('ERR', end=';')

        try: self.h2()
        except: print('ERR', end=';')


class B(A):
    c = 4

    def _h1(self): print(5, end=';')
    def h2(self): print(6, end=';')

try: B().main()
except: print('ERR', end=';')
code

8. Що буде виведено за виконання?
code
def f(arg, n):
    arg[69] = 2
    n = 6
    arg = set(arg.keys())

try:
    n = 2
    t = {81:3, 69:4, 81:3}
    f(t, n)
    print(n, end=';')
    for x in t.values():
        print(x, sep=';', end=';')
except:
    print('ERR')
code

9. Що буде виведено за виконання?
code
class A:
    a = 1
    b = 2

    def __init__(self):
        b = 3

class B(A):
    a = 4

    def __init__(self):
        super().__init__()
        self.c = 7

obj = B()
obj.c = obj.b + 10
B.c = 13

try: print(obj.c, end= ';')
except: print('ERR', end=';')

try: print(A.c, end= ';')
except: print('ERR', end=';')

try: print(B.c, end= ';')
except: print('ERR', end=';')
code

10. 
Для графа на рисунку з вершини 2 запущено алгоритм пошуку в глибину, що будує кістякове дерево. Списки суміжності вершин переглядаються алгоритмом за зростанням міток вершин.\par
\begin{center}
	\adjustimage{max size={0.9\linewidth}{0.9\paperheight}}
	{../10/Traversals/141.png}
\end{center}
\par Які з наведених ребер входять до складу отриманого кістякового дерева?\par
1) (4, 6)\par
2) (0, 4)\par
3) (0, 5)\par
4) (0, 1)\par
5) (4, 5)\par
6) (2, 4)\par
{\vskip 15pt} У формі вибрати номери відповідних ребер.\par
%answer:1 4 5
\end {large}}
%end
{\smallskip \begin {small} Затверджено на засіданні кафедри теоретичної кібернетики, протокол \No 12 від   19.05.2020 р.\par\smallskip\smallskip
{\parbox {6.5 cm}{Зав. кафедри \dotfill Ю.В.~Крак}}\hspace {0.7cm}{\hbox to 6.5 cm{ Екзаменатор \dotfill Т.О.~Карнаух}} \end{small}}
%begin
{{\begin {small}\par
{ \center  Київський національний університет імені Тараса Шевченка \par}
{\parbox {5 cm}{Освітня програма \hfill}{\parbox {7 cm}{Прикладна математика, Системний аналіз \hfill}}\parbox {6 cm}{\hfill Семестр 2}}\par
{\parbox {5 cm} {Навчальний предмет \hfill}\parbox {7 cm}{Програмування \hfill}\parbox {6cm}{\hfill}\par}\vspace{-7pt} \end{small}}{\center Екзаменаційний білет \No \textbf 
{ 475 }\par}}
{
\begin {large}
1. Що буде виведено за виконання?
code
try:
    res = 3j<="5.0"
    if isinstance(res, bool):
        print(f'bool;{res}')
    elif isinstance(res, int):
        print(f'int;{res}')
    elif isinstance(res, float):
        print(f'float;{res:.2f}')
    else:
        print('???')
except:
    print('ERR')
code

2. Що буде виведено за виконання?
code
a = 8
b = 9
c = 4

def g(a):
    global c
    a = 5
    b = 3
    c = 2
    return a + b + c

a = 2
b = 0
c = 1

try:
    print(g(a), a, b, c, sep=';')
except:
    print('ERR', a, b, c, sep=';')
code

3. Що буде виведено за виконання?
code
try:
    j = 7
    while j<7:
        j += 1
    print(j, end=';')
except:
    print('ERR')
code

4. Що буде виведено за виконання?
code
try:
    for d in range(-2, 14, -1):
        if d<7:
            break
        print(d, end=';')
        if d<6:
            break
    else:
        print(d,end=';')
    print(d)
except:
    print('ERR')
code

5. Що буде виведено за виконання?\par (Дійсні значення, якщо наявні, в цьому завданні будуть виводитись у форматі з фіксованою крапкою та, можливо, з доволі великою кількістю десяткових розрядів. У відповіді для дійсних значень у ЦЬОМУ завданні достатньо вказати один десятковий розряд після крапки, відкинувши решту розрядів без округлення. Наприклад, замість 13.66666666666667 достатньо вказати 13.6, замість 25.23 достатньо вказати 25.2)
code
try:
    print(3, end=';')
    print(9j>=5j, end=';')
    print(9, end=';')
except TypeError:
    print(8, end=';')
except KeyboardInterrupt:
    print(0, end=';')
except:
    print('ERR', end=';')
else:
    print(3, end=';')
finally:
    print(7)
code

6. Що буде виведено за виконання?
code
def f(s):
    s1 = s
    for i in range(len(s)):
        s1[i] += 1
    return s1

s = [1, 2, 3]
f(s)
print(s)
code

7. Що буде виведено за виконання?
code
class A:
    d = 1

    def __h1(self): print(2, end=';')
    def _h2(self): print(3, end=';')

    def main(self):
        try: print(self.d, end=';')
        except: print('ERR', end=';')

        try: self.__h1()
        except: print('ERR', end=';')

        try: self._h2()
        except: print('ERR', end=';')


class B(A):
    d = 4

    def __h1(self): print(5, end=';')
    def _h2(self): print(6, end=';')

try: B().main()
except: print('ERR', end=';')
code

8. Що буде виведено за виконання?
code
def f(arg, n):
    arg[80] = 5
    n *= 2

try:
    n = 7
    s = {76:1, 80:2, 76:7}
    f(s, n)
    print(n, end=';')
    for x in s.items():
        print(x, sep=';', end=';')
except:
    print('ERR')
code

9. Що буде виведено за виконання?
code
class A:
    c = 1
    a = 2

    def __init__(self):
        c = 3

class B(A):
    c = 4

    def __init__(self):
        super().__init__()
        self.b = 7

obj = B()
obj.a = obj.b + 10
B.a = 13

try: print(obj.a, end= ';')
except: print('ERR', end=';')

try: print(A.a, end= ';')
except: print('ERR', end=';')

try: print(B.a, end= ';')
except: print('ERR', end=';')
code

10. 
Для графа на рисунку з вершини 0 запущено алгоритм пошуку в глибину, що будує кістякове дерево. Списки суміжності вершин переглядаються алгоритмом за зростанням міток вершин.\par
\begin{center}
	\adjustimage{max size={0.9\linewidth}{0.9\paperheight}}
	{../10/Traversals/66.png}
\end{center}
\par Які з наведених ребер входять до складу отриманого кістякового дерева?\par
1) (0, 2)\par
2) (3, 4)\par
3) (3, 6)\par
4) (0, 3)\par
5) (4, 5)\par
6) (4, 6)\par
{\vskip 15pt} У формі вибрати номери відповідних ребер.\par
%answer:1 3 5
\end {large}}
%end
{\smallskip \begin {small} Затверджено на засіданні кафедри теоретичної кібернетики, протокол \No 12 від   19.05.2020 р.\par\smallskip\smallskip
{\parbox {6.5 cm}{Зав. кафедри \dotfill Ю.В.~Крак}}\hspace {0.7cm}{\hbox to 6.5 cm{ Екзаменатор \dotfill Т.О.~Карнаух}} \end{small}}
%begin
{{\begin {small}\par
{ \center  Київський національний університет імені Тараса Шевченка \par}
{\parbox {5 cm}{Освітня програма \hfill}{\parbox {7 cm}{Прикладна математика, Системний аналіз \hfill}}\parbox {6 cm}{\hfill Семестр 2}}\par
{\parbox {5 cm} {Навчальний предмет \hfill}\parbox {7 cm}{Програмування \hfill}\parbox {6cm}{\hfill}\par}\vspace{-7pt} \end{small}}{\center Екзаменаційний білет \No \textbf 
{ 476 }\par}}
{
\begin {large}
1. Що буде виведено за виконання?
code
try:
    res = 3/9/9*3
    if isinstance(res, bool):
        print(f'bool;{res}')
    elif isinstance(res, int):
        print(f'int;{res}')
    elif isinstance(res, float):
        print(f'float;{res:.2f}')
    else:
        print('???')
except:
    print('ERR')
code

2. Що буде виведено за виконання?
code
a = 3
b = 4
c = 9

def f(b):
    a = 4
    b = 5
    c = 3
    return a + b + c

a = 7
b = 5
c = 6

try:
    print(f(a), a, b, c, sep=';')
except:
    print('ERR', a, b, c, sep=';')
code

3. Що буде виведено за виконання?
code
try:
    t = 13
    while t>5:
        t += -3
    print(t, end=';')
except:
    print('ERR')
code

4. Що буде виведено за виконання?
code
try:
    for a in range(-6, -14, 2):
        if a<=8:
            continue
        print(a, end=';')
    else:
        print(a,end=';')
    print(a)
except:
    print('ERR')
code

5. Що буде виведено за виконання?\par (Дійсні значення, якщо наявні, в цьому завданні будуть виводитись у форматі з фіксованою крапкою та, можливо, з доволі великою кількістю десяткових розрядів. У відповіді для дійсних значень у ЦЬОМУ завданні достатньо вказати один десятковий розряд після крапки, відкинувши решту розрядів без округлення. Наприклад, замість 13.66666666666667 достатньо вказати 13.6, замість 25.23 достатньо вказати 25.2)
code
try:
    print(4, end=';')
    print(int('b5'), end=';')
    print(2, end=';')
except ValueError:
    print(1, end=';')
except Exception:
    print(6, end=';')
except:
    print('ERR', end=';')
else:
    print(1, end=';')
finally:
    print(4)
code

6. Що буде виведено за виконання?
code
def f(s):
    s1 = s
    for i in range(len(s)):
        s1[i] += 1
        return s1

s = [1, 2, 3]
f(s)
print(s)
code

7. Що буде виведено за виконання?
code
class A:
    a = 1

    def _f1(self): print(2, end=';')
    def _f2(self): print(3, end=';')

    def main(self):
        try: print(self.a, end=';')
        except: print('ERR', end=';')

        try: self._f1()
        except: print('ERR', end=';')

        try: self._f2()
        except: print('ERR', end=';')


class B(A):
    a = 4

    def _f1(self): print(5, end=';')
    def _f2(self): print(6, end=';')

try: B().main()
except: print('ERR', end=';')
code

8. Що буде виведено за виконання?
code
def f(arg, n):
    n = 0
    tmp = arg
    del tmp[-3:-5]
    return len(tmp)>3

try:
    a = ['0','1','2','3','4']
    n = 9
    f(a, n)
    print(n, end=';')
    for el in a:
        print(el, end=';')
except:
    print('ERR')
code

9. Що буде виведено за виконання?
code
class A:
    a = 1
    c = 2

    def __init__(self):
        c = 3

class B(A):
    c = 4

    def __init__(self):
        super().__init__()
        self.b = 7

obj = B()
obj.c = obj.b + 10
B.c = 13

try: print(obj.c, end= ';')
except: print('ERR', end=';')

try: print(A.c, end= ';')
except: print('ERR', end=';')

try: print(B.c, end= ';')
except: print('ERR', end=';')
code

10. 
Для графа на рисунку з вершини 6 запущено алгоритм пошуку в глибину, що будує кістякове дерево. Списки суміжності вершин переглядаються алгоритмом за зростанням міток вершин.\par
\begin{center}
	\adjustimage{max size={0.9\linewidth}{0.9\paperheight}}
	{../10/Traversals/213.png}
\end{center}
\par Які з наведених ребер входять до складу отриманого кістякового дерева?\par
1) (0, 3)\par
2) (0, 1)\par
3) (1, 2)\par
4) (1, 6)\par
5) (3, 6)\par
6) (2, 5)\par
{\vskip 15pt} У формі вибрати номери відповідних ребер.\par
%answer:2 3
\end {large}}
%end
{\smallskip \begin {small} Затверджено на засіданні кафедри теоретичної кібернетики, протокол \No 12 від   19.05.2020 р.\par\smallskip\smallskip
{\parbox {6.5 cm}{Зав. кафедри \dotfill Ю.В.~Крак}}\hspace {0.7cm}{\hbox to 6.5 cm{ Екзаменатор \dotfill Т.О.~Карнаух}} \end{small}}
%begin
{{\begin {small}\par
{ \center  Київський національний університет імені Тараса Шевченка \par}
{\parbox {5 cm}{Освітня програма \hfill}{\parbox {7 cm}{Прикладна математика, Системний аналіз \hfill}}\parbox {6 cm}{\hfill Семестр 2}}\par
{\parbox {5 cm} {Навчальний предмет \hfill}\parbox {7 cm}{Програмування \hfill}\parbox {6cm}{\hfill}\par}\vspace{-7pt} \end{small}}{\center Екзаменаційний білет \No \textbf 
{ 477 }\par}}
{
\begin {large}
1. Що буде виведено за виконання?
code
try:
    res = 9/6*5*5
    if isinstance(res, bool):
        print(f'bool;{res}')
    elif isinstance(res, int):
        print(f'int;{res}')
    elif isinstance(res, float):
        print(f'float;{res:.2f}')
    else:
        print('???')
except:
    print('ERR')
code

2. Що буде виведено за виконання?
code
a = 8
b = 2
c = 5

def h(a):
    global c
    a = 1
    b = 2
    c = 5
    return a + b + c

a = 2
b = 0
c = 7

try:
    print(h(b), a, b, c, sep=';')
except:
    print('ERR', a, b, c, sep=';')
code

3. Що буде виведено за виконання?
code
try:
    i = 15
    while i>=4:
        i += -3
    print(i, end=';')
except:
    print('ERR')
code

4. Що буде виведено за виконання?
code
try:
    for c in range(-8, -6, 3):
        if c<=6:
            continue
        print(c, end=';');c = -10
    else:
        print('end',end=';')
    print(c)
except:
    print('ERR')
code

5. Що буде виведено за виконання?\par (Дійсні значення, якщо наявні, в цьому завданні будуть виводитись у форматі з фіксованою крапкою та, можливо, з доволі великою кількістю десяткових розрядів. У відповіді для дійсних значень у ЦЬОМУ завданні достатньо вказати один десятковий розряд після крапки, відкинувши решту розрядів без округлення. Наприклад, замість 13.66666666666667 достатньо вказати 13.6, замість 25.23 достатньо вказати 25.2)
code
try:
    print(5, end=';')
    print(int('3'), end=';')
    print(0, end=';')
except TypeError:
    print(9, end=';')
except ZeroDivisionError:
    print(6, end=';')
except:
    print('ERR', end=';')
else:
    print(2, end=';')
finally:
    print(7)
code

6. Що буде виведено за виконання?
code
def f(s):
    s1 = s
    for i in range(len(s)):
        s1[i] += 1
    return s1

s = [1, 2, 3]
f(s)
print(s)
code

7. Що буде виведено за виконання?
code
class A:
    c = 1

    def _g1(self): print(2, end=';')
    def __g2(self): print(3, end=';')

    def main(self):
        try: print(self.c, end=';')
        except: print('ERR', end=';')

        try: self._g1()
        except: print('ERR', end=';')

        try: self.__g2()
        except: print('ERR', end=';')


class B(A):
    c = 4

    def _g1(self): print(5, end=';')
    def __g2(self): print(6, end=';')

try: B().main()
except: print('ERR', end=';')
code

8. Що буде виведено за виконання?
code
def f(arg, n):
    arg[81] = 7
    n *= 3

try:
    n = 3
    d = {35:4, 20:3, 87:0, 20:6}
    f(d, n)
    print(n, end=';')
    for x in d.keys():
        print(x, sep=';', end=';')
except:
    print('ERR')
code

9. Що буде виведено за виконання?
code
class A:
    c = 1
    b = 2

    def __init__(self):
        b = 3

class B(A):
    c = 4

    def __init__(self):
        super().__init__()
        self.a = 7

obj = B()
obj.a = obj.c + 10
B.a = 13

try: print(obj.a, end= ';')
except: print('ERR', end=';')

try: print(A.a, end= ';')
except: print('ERR', end=';')

try: print(B.a, end= ';')
except: print('ERR', end=';')
code

10. 
Для графа на рисунку з вершини 5 запущено алгоритм пошуку в глибину, що будує кістякове дерево. Списки суміжності вершин переглядаються алгоритмом за зростанням міток вершин.\par
\begin{center}
	\adjustimage{max size={0.9\linewidth}{0.9\paperheight}}
	{../10/Traversals/166.png}
\end{center}
\par Які з наведених ребер входять до складу отриманого кістякового дерева?\par
1) (4, 6)\par
2) (3, 6)\par
3) (1, 3)\par
4) (0, 5)\par
5) (0, 2)\par
6) (1, 5)\par
{\vskip 15pt} У формі вибрати номери відповідних ребер.\par
%answer:1 2 3 4 5
\end {large}}
%end
{\smallskip \begin {small} Затверджено на засіданні кафедри теоретичної кібернетики, протокол \No 12 від   19.05.2020 р.\par\smallskip\smallskip
{\parbox {6.5 cm}{Зав. кафедри \dotfill Ю.В.~Крак}}\hspace {0.7cm}{\hbox to 6.5 cm{ Екзаменатор \dotfill Т.О.~Карнаух}} \end{small}}
%begin
{{\begin {small}\par
{ \center  Київський національний університет імені Тараса Шевченка \par}
{\parbox {5 cm}{Освітня програма \hfill}{\parbox {7 cm}{Прикладна математика, Системний аналіз \hfill}}\parbox {6 cm}{\hfill Семестр 2}}\par
{\parbox {5 cm} {Навчальний предмет \hfill}\parbox {7 cm}{Програмування \hfill}\parbox {6cm}{\hfill}\par}\vspace{-7pt} \end{small}}{\center Екзаменаційний білет \No \textbf 
{ 478 }\par}}
{
\begin {large}
1. Що буде виведено за виконання?
code
try:
    res = 2.0<6 and not 6.0>=2 and 7>False
    if isinstance(res, bool):
        print(f'bool;{res}')
    elif isinstance(res, int):
        print(f'int;{res}')
    elif isinstance(res, float):
        print(f'float;{res:.2f}')
    else:
        print('???')
except:
    print('ERR')
code

2. Що буде виведено за виконання?
code
a = 9
b = 5
c = 8

def h(b):
    a += 3
    b = 4
    c = 1
    return a + b + c

a = 2
b = 5
c = 4

try:
    print(h(b), a, b, c, sep=';')
except:
    print('ERR', a, b, c, sep=';')
code

3. Що буде виведено за виконання?
code
try:
    n = 7
    while n>3:
        n += -2
    print(n, end=';')
except:
    print('ERR')
code

4. Що буде виведено за виконання?
code
try:
    for b in range(-7, -10, -1):
        if b<5:
            break
        print(b, end=';')
    else:
        print(b,end=';')
    print(b)
except:
    print('ERR')
code

5. Що буде виведено за виконання?\par (Дійсні значення, якщо наявні, в цьому завданні будуть виводитись у форматі з фіксованою крапкою та, можливо, з доволі великою кількістю десяткових розрядів. У відповіді для дійсних значень у ЦЬОМУ завданні достатньо вказати один десятковий розряд після крапки, відкинувши решту розрядів без округлення. Наприклад, замість 13.66666666666667 достатньо вказати 13.6, замість 25.23 достатньо вказати 25.2)
code
try:
    print(8, end=';')
    print(1//0, end=';')
    print(9, end=';')
except KeyboardInterrupt:
    print(0, end=';')
except ValueError:
    print(7, end=';')
except:
    print('ERR', end=';')
else:
    print(8, end=';')
finally:
    print(5)
code

6. Що буде виведено за виконання?
code
def f(s):
    for i in range(len(s)):
        s[i] += 1
    return
 
s = [1, 2, 3]
f(s)
print(s)
code

7. Що буде виведено за виконання?
code
class A:
    b = 1

    def __h1(self): print(2, end=';')
    def _h2(self): print(3, end=';')

    def main(self):
        try: print(self.b, end=';')
        except: print('ERR', end=';')

        try: self.__h1()
        except: print('ERR', end=';')

        try: self._h2()
        except: print('ERR', end=';')


class B(A):
    b = 4

    def __h1(self): print(5, end=';')
    def _h2(self): print(6, end=';')

try: B().main()
except: print('ERR', end=';')
code

8. Що буде виведено за виконання?
code
def f(arg, n):
    n = 2
    tmp = arg
    tmp[:6] = 'bc'
    return len(tmp)>3

try:
    d = [10,11,12,13,14]
    n = 3
    f(d, n)
    print(n, end=';')
    for el in d:
        print(el, end=';')
except:
    print('ERR')
code

9. Що буде виведено за виконання?
code
class A:
    b = 1
    c = 2

    def __init__(self):
        b = 3

class B(A):
    c = 4

    def __init__(self):
        super().__init__()
        self.a = 7

obj = B()
obj.c = obj.b + 10
B.c = 13

try: print(obj.c, end= ';')
except: print('ERR', end=';')

try: print(A.c, end= ';')
except: print('ERR', end=';')

try: print(B.c, end= ';')
except: print('ERR', end=';')
code

10. 
Для графа на рисунку з вершини 0 запущено алгоритм пошуку в глибину, що будує кістякове дерево. Списки суміжності вершин переглядаються алгоритмом за зростанням міток вершин.\par
\begin{center}
	\adjustimage{max size={0.9\linewidth}{0.9\paperheight}}
	{../10/Traversals/1.png}
\end{center}
\par Які з наведених ребер входять до складу отриманого кістякового дерева?\par
1) (2, 5)\par
2) (2, 3)\par
3) (1, 5)\par
4) (0, 4)\par
5) (3, 5)\par
6) (0, 5)\par
{\vskip 15pt} У формі вибрати номери відповідних ребер.\par
%answer:2 3 5
\end {large}}
%end
{\smallskip \begin {small} Затверджено на засіданні кафедри теоретичної кібернетики, протокол \No 12 від   19.05.2020 р.\par\smallskip\smallskip
{\parbox {6.5 cm}{Зав. кафедри \dotfill Ю.В.~Крак}}\hspace {0.7cm}{\hbox to 6.5 cm{ Екзаменатор \dotfill Т.О.~Карнаух}} \end{small}}
%begin
{{\begin {small}\par
{ \center  Київський національний університет імені Тараса Шевченка \par}
{\parbox {5 cm}{Освітня програма \hfill}{\parbox {7 cm}{Прикладна математика, Системний аналіз \hfill}}\parbox {6 cm}{\hfill Семестр 2}}\par
{\parbox {5 cm} {Навчальний предмет \hfill}\parbox {7 cm}{Програмування \hfill}\parbox {6cm}{\hfill}\par}\vspace{-7pt} \end{small}}{\center Екзаменаційний білет \No \textbf 
{ 479 }\par}}
{
\begin {large}
1. Що буде виведено за виконання?
code
try:
    res = 3.0<3 or not 8>=False and 5<=6.0
    if isinstance(res, bool):
        print(f'bool;{res}')
    elif isinstance(res, int):
        print(f'int;{res}')
    elif isinstance(res, float):
        print(f'float;{res:.2f}')
    else:
        print('???')
except:
    print('ERR')
code

2. Що буде виведено за виконання?
code
a = 3
b = 9
c = 8

def g(b):
    a = 1
    b = 3
    c = 4
    return a + b + c

a = 6
b = 4
c = 1

try:
    print(g(a), a, b, c, sep=';')
except:
    print('ERR', a, b, c, sep=';')
code

3. Що буде виведено за виконання?
code
try:
    k = 8
    while k>8:
        k += -2
    print(k, end=';')
except:
    print('ERR')
code

4. Що буде виведено за виконання?
code
try:
    for e in range(4, 0, -3):
        if e>=3:
            continue
        print(e, end=';');e = -5
    else:
        print(e,end=';')
    print(e)
except:
    print('ERR')
code

5. Що буде виведено за виконання?\par (Дійсні значення, якщо наявні, в цьому завданні будуть виводитись у форматі з фіксованою крапкою та, можливо, з доволі великою кількістю десяткових розрядів. У відповіді для дійсних значень у ЦЬОМУ завданні достатньо вказати один десятковий розряд після крапки, відкинувши решту розрядів без округлення. Наприклад, замість 13.66666666666667 достатньо вказати 13.6, замість 25.23 достатньо вказати 25.2)
code
try:
    print(6, end=';')
    print(2j<=2j, end=';')
    print(4, end=';')
except TypeError:
    print(3, end=';')
except ZeroDivisionError:
    print(9, end=';')
except:
    print('ERR', end=';')
else:
    print(1, end=';')
finally:
    print(7)
code

6. Що буде виведено за виконання?
code
def f(n):
    if n>9:
        f(n//10)
    print(n%10, end='')
 
f(123456)
code

7. Що буде виведено за виконання?
code
class A:
    d = 1

    def _f1(self): print(2, end=';')
    def f2(self): print(3, end=';')

    def main(self):
        try: print(self.d, end=';')
        except: print('ERR', end=';')

        try: self._f1()
        except: print('ERR', end=';')

        try: self.f2()
        except: print('ERR', end=';')


class B(A):
    d = 4

    def _f1(self): print(5, end=';')
    def f2(self): print(6, end=';')

try: B().main()
except: print('ERR', end=';')
code

8. Що буде виведено за виконання?
code
def f(arg, n):
    n = 6
    tmp = arg
    del tmp[-7:]
    return len(tmp)>3

try:
    e = ['0','1','2','3','4']
    n = 5
    f(e, n)
    print(n, end=';')
    for el in e:
        print(el, end=';')
except:
    print('ERR')
code

9. Що буде виведено за виконання?
code
class A:
    b = 1
    a = 2

    def __init__(self):
        a = 3

class B(A):
    b = 4

    def __init__(self):
        super().__init__()
        self.c = 7

obj = B()
obj.a = obj.b + 10
B.a = 13

try: print(obj.a, end= ';')
except: print('ERR', end=';')

try: print(A.a, end= ';')
except: print('ERR', end=';')

try: print(B.a, end= ';')
except: print('ERR', end=';')
code

10. 
Для графа на рисунку з вершини 5 запущено алгоритм пошуку в глибину, що будує кістякове дерево. Списки суміжності вершин переглядаються алгоритмом за зростанням міток вершин.\par
\begin{center}
	\adjustimage{max size={0.9\linewidth}{0.9\paperheight}}
	{../10/Traversals/176.png}
\end{center}
\par Які з наведених ребер входять до складу отриманого кістякового дерева?\par
1) (3, 5)\par
2) (0, 3)\par
3) (5, 6)\par
4) (4, 6)\par
5) (0, 1)\par
6) (1, 2)\par
{\vskip 15pt} У формі вибрати номери відповідних ребер.\par
%answer:1 2 4 5 6
\end {large}}
%end
{\smallskip \begin {small} Затверджено на засіданні кафедри теоретичної кібернетики, протокол \No 12 від   19.05.2020 р.\par\smallskip\smallskip
{\parbox {6.5 cm}{Зав. кафедри \dotfill Ю.В.~Крак}}\hspace {0.7cm}{\hbox to 6.5 cm{ Екзаменатор \dotfill Т.О.~Карнаух}} \end{small}}
%begin
{{\begin {small}\par
{ \center  Київський національний університет імені Тараса Шевченка \par}
{\parbox {5 cm}{Освітня програма \hfill}{\parbox {7 cm}{Прикладна математика, Системний аналіз \hfill}}\parbox {6 cm}{\hfill Семестр 2}}\par
{\parbox {5 cm} {Навчальний предмет \hfill}\parbox {7 cm}{Програмування \hfill}\parbox {6cm}{\hfill}\par}\vspace{-7pt} \end{small}}{\center Екзаменаційний білет \No \textbf 
{ 480 }\par}}
{
\begin {large}
1. Що буде виведено за виконання?
code
try:
    res = 3**-1-False
    if isinstance(res, bool):
        print(f'bool;{res}')
    elif isinstance(res, int):
        print(f'int;{res}')
    elif isinstance(res, float):
        print(f'float;{res:.2f}')
    else:
        print('???')
except:
    print('ERR')
code

2. Що буде виведено за виконання?
code
a = 1
b = 3
c = 6

def g(a):
    global c
    a = 2
    b -= 5
    c = 3
    return a + b + c

a = 9
b = 7
c = 0

try:
    print(g(a), a, b, c, sep=';')
except:
    print('ERR', a, b, c, sep=';')
code

3. Що буде виведено за виконання?
code
try:
    m = 6
    while m>=6:
        m += -3
    print(m, end=';')
except:
    print('ERR')
code

4. Що буде виведено за виконання?
code
try:
    for f in range(-14, 11, -2):
        if f>4:
            continue
        print(f, end=';')
    else:
        print(f,end=';')
    print(f)
except:
    print('ERR')
code

5. Що буде виведено за виконання?\par (Дійсні значення, якщо наявні, в цьому завданні будуть виводитись у форматі з фіксованою крапкою та, можливо, з доволі великою кількістю десяткових розрядів. У відповіді для дійсних значень у ЦЬОМУ завданні достатньо вказати один десятковий розряд після крапки, відкинувши решту розрядів без округлення. Наприклад, замість 13.66666666666667 достатньо вказати 13.6, замість 25.23 достатньо вказати 25.2)
code
try:
    print(5, end=';')
    print(2%1, end=';')
    print(0, end=';')
except KeyboardInterrupt:
    print(6, end=';')
except Exception:
    print(4, end=';')
except:
    print('ERR', end=';')
else:
    print(3, end=';')
finally:
    print(8)
code

6. Що буде виведено за виконання?
code
def f(n):
    print(n%2, end='')
    if n>2:
        f(n//2)
 
f(16)
code

7. Що буде виведено за виконання?
code
class A:
    a = 1

    def g1(self): print(2, end=';')
    def _g2(self): print(3, end=';')

    def main(self):
        try: print(self.a, end=';')
        except: print('ERR', end=';')

        try: self.g1()
        except: print('ERR', end=';')

        try: self._g2()
        except: print('ERR', end=';')


class B(A):
    a = 4

    def g1(self): print(5, end=';')
    def _g2(self): print(6, end=';')

try: B().main()
except: print('ERR', end=';')
code

8. Що буде виведено за виконання?
code
def f(arg, n):
    arg[17] = 9
    n -= -3

try:
    n = 7
    t = {27:3, 75:9, 33:9, 33:8}
    f(t, n)
    print(n, end=';')
    for x in t :
        print(x, sep=';', end=';')
except:
    print('ERR')
code

9. Що буде виведено за виконання?
code
class A:
    a = 1
    b = 2

    def __init__(self):
        a = 3

class B(A):
    b = 4

    def __init__(self):
        super().__init__()
        self.c = 7

obj = B()
obj.c = obj.a + 10
B.c = 13

try: print(obj.c, end= ';')
except: print('ERR', end=';')

try: print(A.c, end= ';')
except: print('ERR', end=';')

try: print(B.c, end= ';')
except: print('ERR', end=';')
code

10. 
Для графа на рисунку з вершини 3 запущено алгоритм пошуку в глибину, що будує кістякове дерево. Списки суміжності вершин переглядаються алгоритмом за зростанням міток вершин.\par
\begin{center}
	\adjustimage{max size={0.9\linewidth}{0.9\paperheight}}
	{../10/Traversals/186.png}
\end{center}
\par Які з наведених ребер входять до складу отриманого кістякового дерева?\par
1) (0, 4)\par
2) (2, 3)\par
3) (3, 6)\par
4) (2, 6)\par
5) (4, 5)\par
6) (4, 6)\par
{\vskip 15pt} У формі вибрати номери відповідних ребер.\par
%answer:1 2 5 6
\end {large}}
%end
{\smallskip \begin {small} Затверджено на засіданні кафедри теоретичної кібернетики, протокол \No 12 від   19.05.2020 р.\par\smallskip\smallskip
{\parbox {6.5 cm}{Зав. кафедри \dotfill Ю.В.~Крак}}\hspace {0.7cm}{\hbox to 6.5 cm{ Екзаменатор \dotfill Т.О.~Карнаух}} \end{small}}
%begin
{{\begin {small}\par
{ \center  Київський національний університет імені Тараса Шевченка \par}
{\parbox {5 cm}{Освітня програма \hfill}{\parbox {7 cm}{Прикладна математика, Системний аналіз \hfill}}\parbox {6 cm}{\hfill Семестр 2}}\par
{\parbox {5 cm} {Навчальний предмет \hfill}\parbox {7 cm}{Програмування \hfill}\parbox {6cm}{\hfill}\par}\vspace{-7pt} \end{small}}{\center Екзаменаційний білет \No \textbf 
{ 481 }\par}}
{
\begin {large}
1. Що буде виведено за виконання?
code
try:
    res = 4//4//3*7
    if isinstance(res, bool):
        print(f'bool;{res}')
    elif isinstance(res, int):
        print(f'int;{res}')
    elif isinstance(res, float):
        print(f'float;{res:.2f}')
    else:
        print('???')
except:
    print('ERR')
code

2. Що буде виведено за виконання?
code
a = 1
b = 2
c = 6

def h(a):
    global c
    a *= 5
    b = 4
    c = 3
    return a + b + c

a = 4
b = 8
c = 5

try:
    print(h(a), a, b, c, sep=';')
except:
    print('ERR', a, b, c, sep=';')
code

3. Що буде виведено за виконання?
code
try:
    m = 8
    while m<=9:
        m += 2
    print(m, end=';')
except:
    print('ERR')
code

4. Що буде виведено за виконання?
code
try:
    for c in range(-3, -11, 2):
        if c<6:
            continue
        print(c, end=';')
    else:
        print(c,end=';')
    print(c)
except:
    print('ERR')
code

5. Що буде виведено за виконання?\par (Дійсні значення, якщо наявні, в цьому завданні будуть виводитись у форматі з фіксованою крапкою та, можливо, з доволі великою кількістю десяткових розрядів. У відповіді для дійсних значень у ЦЬОМУ завданні достатньо вказати один десятковий розряд після крапки, відкинувши решту розрядів без округлення. Наприклад, замість 13.66666666666667 достатньо вказати 13.6, замість 25.23 достатньо вказати 25.2)
code
try:
    print(4, end=';')
    print(int('5'), end=';')
    print(6, end=';')
except KeyboardInterrupt:
    print(9, end=';')
except ZeroDivisionError:
    print(7, end=';')
except:
    print('ERR', end=';')
else:
    print(2, end=';')
finally:
    print(8)
code

6. Що буде виведено за виконання?
code
def f(s):
    for el in s:
        el += 1
    return s
 
s = [1, 2, 3]
s = f(s)
print(s)
code

7. Що буде виведено за виконання?
code
class A:
    c = 1

    def __h1(self): print(2, end=';')
    def _h2(self): print(3, end=';')

    def main(self):
        try: print(self.c, end=';')
        except: print('ERR', end=';')

        try: self.__h1()
        except: print('ERR', end=';')

        try: self._h2()
        except: print('ERR', end=';')


class B(A):
    c = 4

    def __h1(self): print(5, end=';')
    def _h2(self): print(6, end=';')

try: B().main()
except: print('ERR', end=';')
code

8. Що буде виведено за виконання?
code
def f(arg, n):
    arg[71] = 7
    n += 2

try:
    n = 4
    s = {31:5, 79:3, 51:1, 31:6}
    f(s, n)
    print(n, end=';')
    for x in s.items():
        print(*x, sep=';', end=';')
except:
    print('ERR')
code

9. Що буде виведено за виконання?
code
class A:
    c = 1
    a = 2

    def __init__(self):
        c = 3

class B(A):
    c = 4

    def __init__(self):
        super().__init__()
        self.b = 7

obj = B()
obj.b = obj.a + 10
B.b = 13

try: print(obj.b, end= ';')
except: print('ERR', end=';')

try: print(A.b, end= ';')
except: print('ERR', end=';')

try: print(B.b, end= ';')
except: print('ERR', end=';')
code

10. 
Для графа на рисунку з вершини 5 запущено алгоритм пошуку в глибину, що будує кістякове дерево. Списки суміжності вершин переглядаються алгоритмом за зростанням міток вершин.\par
\begin{center}
	\adjustimage{max size={0.9\linewidth}{0.9\paperheight}}
	{../10/Traversals/29.png}
\end{center}
\par Які з наведених ребер входять до складу отриманого кістякового дерева?\par
1) (1, 4)\par
2) (1, 6)\par
3) (4, 6)\par
4) (1, 2)\par
5) (2, 3)\par
6) (5, 6)\par
{\vskip 15pt} У формі вибрати номери відповідних ребер.\par
%answer:3 4 5
\end {large}}
%end
{\smallskip \begin {small} Затверджено на засіданні кафедри теоретичної кібернетики, протокол \No 12 від   19.05.2020 р.\par\smallskip\smallskip
{\parbox {6.5 cm}{Зав. кафедри \dotfill Ю.В.~Крак}}\hspace {0.7cm}{\hbox to 6.5 cm{ Екзаменатор \dotfill Т.О.~Карнаух}} \end{small}}
%begin
{{\begin {small}\par
{ \center  Київський національний університет імені Тараса Шевченка \par}
{\parbox {5 cm}{Освітня програма \hfill}{\parbox {7 cm}{Прикладна математика, Системний аналіз \hfill}}\parbox {6 cm}{\hfill Семестр 2}}\par
{\parbox {5 cm} {Навчальний предмет \hfill}\parbox {7 cm}{Програмування \hfill}\parbox {6cm}{\hfill}\par}\vspace{-7pt} \end{small}}{\center Екзаменаційний білет \No \textbf 
{ 482 }\par}}
{
\begin {large}
1. Що буде виведено за виконання?
code
try:
    res = True==7 and not 9!=6 or 2.0>4.0
    if isinstance(res, bool):
        print(f'bool;{res}')
    elif isinstance(res, int):
        print(f'int;{res}')
    elif isinstance(res, float):
        print(f'float;{res:.2f}')
    else:
        print('???')
except:
    print('ERR')
code

2. Що буде виведено за виконання?
code
a = 9
b = 7
c = 0

def f(b):
    a = 2
    b -= 1
    c = 3
    return a + b + c

a = 3
b = 2
c = 3

try:
    print(f(b), a, b, c, sep=';')
except:
    print('ERR', a, b, c, sep=';')
code

3. Що буде виведено за виконання?
code
try:
    k = 14
    while k>=3:
        k += -2
    print(k, end=';')
except:
    print('ERR')
code

4. Що буде виведено за виконання?
code
try:
    for e in range(-5, -9, -3):
        if e<=4:
            break
        print(e, end=';')
    else:
        print(e,end=';')
    print(e)
except:
    print('ERR')
code

5. Що буде виведено за виконання?\par (Дійсні значення, якщо наявні, в цьому завданні будуть виводитись у форматі з фіксованою крапкою та, можливо, з доволі великою кількістю десяткових розрядів. У відповіді для дійсних значень у ЦЬОМУ завданні достатньо вказати один десятковий розряд після крапки, відкинувши решту розрядів без округлення. Наприклад, замість 13.66666666666667 достатньо вказати 13.6, замість 25.23 достатньо вказати 25.2)
code
try:
    print(0, end=';')
    print(1//0, end=';')
    print(3, end=';')
except Exception:
    print(3, end=';')
except ValueError:
    print(0, end=';')
except:
    print('ERR', end=';')
else:
    print(4, end=';')
finally:
    print(8)
code

6. Що буде виведено за виконання?
code
def f(n):
    print(n%10,end='')
    if n>9:
        f(n//100)
 
f(123456)
code

7. Що буде виведено за виконання?
code
class A:
    b = 1

    def _f1(self): print(2, end=';')
    def __f2(self): print(3, end=';')

    def main(self):
        try: print(self.b, end=';')
        except: print('ERR', end=';')

        try: self._f1()
        except: print('ERR', end=';')

        try: self.__f2()
        except: print('ERR', end=';')


class B(A):
    b = 4

    def _f1(self): print(5, end=';')
    def __f2(self): print(6, end=';')

try: B().main()
except: print('ERR', end=';')
code

8. Що буде виведено за виконання?
code
def f(arg, n):
    arg[36] = 1
    n = 9
    arg = set(arg.values())

try:
    n = 1
    t = {51:1, 39:8, 36:3}
    f(t, n)
    print(n, end=';')
    for x in t.items():
        print(*x, sep=';', end=';')
except:
    print('ERR')
code

9. Що буде виведено за виконання?
code
class A:
    a = 1
    c = 2

    def __init__(self):
        c = 3

class B(A):
    c = 4

    def __init__(self):
        super().__init__()
        self.b = 7

obj = B()
obj.c = obj.c + 10
B.c = 13

try: print(obj.c, end= ';')
except: print('ERR', end=';')

try: print(A.c, end= ';')
except: print('ERR', end=';')

try: print(B.c, end= ';')
except: print('ERR', end=';')
code

10. 
Для графа на рисунку з вершини 0 запущено алгоритм пошуку в глибину, що будує кістякове дерево. Списки суміжності вершин переглядаються алгоритмом за зростанням міток вершин.\par
\begin{center}
	\adjustimage{max size={0.9\linewidth}{0.9\paperheight}}
	{../10/Traversals/74.png}
\end{center}
\par Які з наведених ребер входять до складу отриманого кістякового дерева?\par
1) (0, 1)\par
2) (1, 3)\par
3) (1, 2)\par
4) (3, 5)\par
5) (0, 6)\par
6) (1, 4)\par
{\vskip 15pt} У формі вибрати номери відповідних ребер.\par
%answer:1 3 4 5 6
\end {large}}
%end
{\smallskip \begin {small} Затверджено на засіданні кафедри теоретичної кібернетики, протокол \No 12 від   19.05.2020 р.\par\smallskip\smallskip
{\parbox {6.5 cm}{Зав. кафедри \dotfill Ю.В.~Крак}}\hspace {0.7cm}{\hbox to 6.5 cm{ Екзаменатор \dotfill Т.О.~Карнаух}} \end{small}}
%begin
{{\begin {small}\par
{ \center  Київський національний університет імені Тараса Шевченка \par}
{\parbox {5 cm}{Освітня програма \hfill}{\parbox {7 cm}{Прикладна математика, Системний аналіз \hfill}}\parbox {6 cm}{\hfill Семестр 2}}\par
{\parbox {5 cm} {Навчальний предмет \hfill}\parbox {7 cm}{Програмування \hfill}\parbox {6cm}{\hfill}\par}\vspace{-7pt} \end{small}}{\center Екзаменаційний білет \No \textbf 
{ 483 }\par}}
{
\begin {large}
1. Що буде виведено за виконання?
code
try:
    res = not 2>=True or 4==8.0 or 8<9.0
    if isinstance(res, bool):
        print(f'bool;{res}')
    elif isinstance(res, int):
        print(f'int;{res}')
    elif isinstance(res, float):
        print(f'float;{res:.2f}')
    else:
        print('???')
except:
    print('ERR')
code

2. Що буде виведено за виконання?
code
a = 7
b = 4
c = 6

def g(b):
    global c
    a += 4
    b = 5
    c = 2
    return a + b + c

a = 8
b = 0
c = 1

try:
    print(g(b), a, b, c, sep=';')
except:
    print('ERR', a, b, c, sep=';')
code

3. Що буде виведено за виконання?
code
try:
    j = 6
    while j<13:
        j += 1
    print(j, end=';')
except:
    print('ERR')
code

4. Що буде виведено за виконання?
code
try:
    for b in range(14, 9, -2):
        if b>=3:
            continue
        print(b, end=';');b = -7
    else:
        print('end',end=';')
    print(b)
except:
    print('ERR')
code

5. Що буде виведено за виконання?\par (Дійсні значення, якщо наявні, в цьому завданні будуть виводитись у форматі з фіксованою крапкою та, можливо, з доволі великою кількістю десяткових розрядів. У відповіді для дійсних значень у ЦЬОМУ завданні достатньо вказати один десятковий розряд після крапки, відкинувши решту розрядів без округлення. Наприклад, замість 13.66666666666667 достатньо вказати 13.6, замість 25.23 достатньо вказати 25.2)
code
try:
    print(7, end=';')
    print(6/1, end=';')
    print(2, end=';')
except TypeError:
    print(9, end=';')
except ZeroDivisionError:
    print(1, end=';')
except:
    print('ERR', end=';')
else:
    print(5, end=';')
finally:
    print(9)
code

6. Що буде виведено за виконання?
code
def f(n):
    if n>9:
        f(n//100)
    else:
        print(n%10,end='')

f(123456)
code

7. Що буде виведено за виконання?
code
class A:
    a = 1

    def _g1(self): print(2, end=';')
    def g2(self): print(3, end=';')

    def main(self):
        try: print(self.a, end=';')
        except: print('ERR', end=';')

        try: self._g1()
        except: print('ERR', end=';')

        try: self.g2()
        except: print('ERR', end=';')


class B(A):
    a = 4

    def _g1(self): print(5, end=';')
    def g2(self): print(6, end=';')

try: B().main()
except: print('ERR', end=';')
code

8. Що буде виведено за виконання?
code
def f(arg, n):
    n = 1
    tmp = arg
    tmp[8:0] = [13]
    return len(tmp)>3

try:
    h = [10,11,12,13,14]
    n = 7
    f(h, n)
    print(n, end=';')
    for el in h:
        print(el, end=';')
except:
    print('ERR')
code

9. Що буде виведено за виконання?
code
class A:
    b = 1
    a = 2

    def __init__(self):
        b = 3

class B(A):
    b = 4

    def __init__(self):
        super().__init__()
        self.c = 7

obj = B()
obj.a = obj.b + 10
B.a = 13

try: print(obj.a, end= ';')
except: print('ERR', end=';')

try: print(A.a, end= ';')
except: print('ERR', end=';')

try: print(B.a, end= ';')
except: print('ERR', end=';')
code

10. 
Для графа на рисунку з вершини 1 запущено алгоритм пошуку в глибину, що будує кістякове дерево. Списки суміжності вершин переглядаються алгоритмом за зростанням міток вершин.\par
\begin{center}
	\adjustimage{max size={0.9\linewidth}{0.9\paperheight}}
	{../10/Traversals/20.png}
\end{center}
\par Які з наведених ребер входять до складу отриманого кістякового дерева?\par
1) (2, 6)\par
2) (4, 6)\par
3) (0, 5)\par
4) (5, 6)\par
5) (0, 4)\par
6) (1, 5)\par
{\vskip 15pt} У формі вибрати номери відповідних ребер.\par
%answer:1 2 3 5 6
\end {large}}
%end
{\smallskip \begin {small} Затверджено на засіданні кафедри теоретичної кібернетики, протокол \No 12 від   19.05.2020 р.\par\smallskip\smallskip
{\parbox {6.5 cm}{Зав. кафедри \dotfill Ю.В.~Крак}}\hspace {0.7cm}{\hbox to 6.5 cm{ Екзаменатор \dotfill Т.О.~Карнаух}} \end{small}}
%begin
{{\begin {small}\par
{ \center  Київський національний університет імені Тараса Шевченка \par}
{\parbox {5 cm}{Освітня програма \hfill}{\parbox {7 cm}{Прикладна математика, Системний аналіз \hfill}}\parbox {6 cm}{\hfill Семестр 2}}\par
{\parbox {5 cm} {Навчальний предмет \hfill}\parbox {7 cm}{Програмування \hfill}\parbox {6cm}{\hfill}\par}\vspace{-7pt} \end{small}}{\center Екзаменаційний білет \No \textbf 
{ 484 }\par}}
{
\begin {large}
1. Що буде виведено за виконання?
code
try:
    res = "9">"False"
    if isinstance(res, bool):
        print(f'bool;{res}')
    elif isinstance(res, int):
        print(f'int;{res}')
    elif isinstance(res, float):
        print(f'float;{res:.2f}')
    else:
        print('???')
except:
    print('ERR')
code

2. Що буде виведено за виконання?
code
a = 7
b = 4
c = 5

def f(a):
    a = 5
    b = 2
    c = 1
    return a + b + c

a = 9
b = 2
c = 0

try:
    print(f(a), a, b, c, sep=';')
except:
    print('ERR', a, b, c, sep=';')
code

3. Що буде виведено за виконання?
code
try:
    i = 9
    while i<11:
        i += 1
    print(i, end=';')
except:
    print('ERR')
code

4. Що буде виведено за виконання?
code
try:
    for f in range(8, -15, -1):
        if f>7:
            continue
        print(f, end=';');f = -9
    else:
        print(f,end=';')
    print(f)
except:
    print('ERR')
code

5. Що буде виведено за виконання?\par (Дійсні значення, якщо наявні, в цьому завданні будуть виводитись у форматі з фіксованою крапкою та, можливо, з доволі великою кількістю десяткових розрядів. У відповіді для дійсних значень у ЦЬОМУ завданні достатньо вказати один десятковий розряд після крапки, відкинувши решту розрядів без округлення. Наприклад, замість 13.66666666666667 достатньо вказати 13.6, замість 25.23 достатньо вказати 25.2)
code
try:
    print(3, end=';')
    print(8j>=8j, end=';')
    print(2, end=';')
except KeyboardInterrupt:
    print(7, end=';')
except ValueError:
    print(0, end=';')
except:
    print('ERR', end=';')
else:
    print(4, end=';')
finally:
    print(1)
code

6. Що буде виведено за виконання?
code
def f(n):
    if n>9:
        f(n//100)
        print(n%10, end='')
 
f(123456)
code

7. Що буде виведено за виконання?
code
class A:
    d = 1

    def _g1(self): print(2, end=';')
    def _g2(self): print(3, end=';')

    def main(self):
        try: print(self.d, end=';')
        except: print('ERR', end=';')

        try: self._g1()
        except: print('ERR', end=';')

        try: self._g2()
        except: print('ERR', end=';')


class B(A):
    d = 4

    def _g1(self): print(5, end=';')
    def _g2(self): print(6, end=';')

try: B().main()
except: print('ERR', end=';')
code

8. Що буде виведено за виконання?
code
def f(arg, n):
    n = 4
    tmp = arg
    tmp[-8:] = [1,0]
    return len(tmp)>3

try:
    b = ['0','1','2','3','4']
    n = 9
    f(b, n)
    print(n, end=';')
    for el in b:
        print(el, end=';')
except:
    print('ERR')
code

9. Що буде виведено за виконання?
code
class A:
    c = 1
    b = 2

    def __init__(self):
        b = 3

class B(A):
    b = 4

    def __init__(self):
        super().__init__()
        self.a = 7

obj = B()
obj.b = obj.a + 10
B.b = 13

try: print(obj.b, end= ';')
except: print('ERR', end=';')

try: print(A.b, end= ';')
except: print('ERR', end=';')

try: print(B.b, end= ';')
except: print('ERR', end=';')
code

10. 
Для графа на рисунку з вершини 3 запущено алгоритм пошуку в глибину, що будує кістякове дерево. Списки суміжності вершин переглядаються алгоритмом за зростанням міток вершин.\par
\begin{center}
	\adjustimage{max size={0.9\linewidth}{0.9\paperheight}}
	{../10/Traversals/163.png}
\end{center}
\par Які з наведених ребер входять до складу отриманого кістякового дерева?\par
1) (1, 5)\par
2) (1, 2)\par
3) (0, 3)\par
4) (1, 4)\par
5) (0, 5)\par
6) (0, 6)\par
{\vskip 15pt} У формі вибрати номери відповідних ребер.\par
%answer:1 2 3 4 6
\end {large}}
%end
{\smallskip \begin {small} Затверджено на засіданні кафедри теоретичної кібернетики, протокол \No 12 від   19.05.2020 р.\par\smallskip\smallskip
{\parbox {6.5 cm}{Зав. кафедри \dotfill Ю.В.~Крак}}\hspace {0.7cm}{\hbox to 6.5 cm{ Екзаменатор \dotfill Т.О.~Карнаух}} \end{small}}
%begin
{{\begin {small}\par
{ \center  Київський національний університет імені Тараса Шевченка \par}
{\parbox {5 cm}{Освітня програма \hfill}{\parbox {7 cm}{Прикладна математика, Системний аналіз \hfill}}\parbox {6 cm}{\hfill Семестр 2}}\par
{\parbox {5 cm} {Навчальний предмет \hfill}\parbox {7 cm}{Програмування \hfill}\parbox {6cm}{\hfill}\par}\vspace{-7pt} \end{small}}{\center Екзаменаційний білет \No \textbf 
{ 485 }\par}}
{
\begin {large}
1. Що буде виведено за виконання?
code
try:
    res = not 2>7.0 and 3<=4 and False!=5.0
    if isinstance(res, bool):
        print(f'bool;{res}')
    elif isinstance(res, int):
        print(f'int;{res}')
    elif isinstance(res, float):
        print(f'float;{res:.2f}')
    else:
        print('???')
except:
    print('ERR')
code

2. Що буде виведено за виконання?
code
a = 9
b = 5
c = 1

def f(a):
    a = 1
    b -= 1
    c = 2
    return a + b + c

a = 6
b = 7
c = 0

try:
    print(f(a), a, b, c, sep=';')
except:
    print('ERR', a, b, c, sep=';')
code

3. Що буде виведено за виконання?
code
try:
    t = 3
    while t<=13:
        t += 2
    print(t, end=';')
except:
    print('ERR')
code

4. Що буде виведено за виконання?
code
try:
    for d in range(-1, 5, 3):
        if d>=8:
            continue
        print(d, end=';')
    else:
        print(d,end=';')
    print(d)
except:
    print('ERR')
code

5. Що буде виведено за виконання?\par (Дійсні значення, якщо наявні, в цьому завданні будуть виводитись у форматі з фіксованою крапкою та, можливо, з доволі великою кількістю десяткових розрядів. У відповіді для дійсних значень у ЦЬОМУ завданні достатньо вказати один десятковий розряд після крапки, відкинувши решту розрядів без округлення. Наприклад, замість 13.66666666666667 достатньо вказати 13.6, замість 25.23 достатньо вказати 25.2)
code
try:
    print(5, end=';')
    print(int('a6'), end=';')
    print(0, end=';')
except Exception:
    print(1, end=';')
except TypeError:
    print(4, end=';')
except:
    print('ERR', end=';')
else:
    print(7, end=';')
finally:
    print(6)
code

6. Що буде виведено за виконання?
code
def f(n):
    print(n%10, end='')
    if n>9:
        f(n//100)
 
f(987654)
code

7. Що буде виведено за виконання?
code
class A:
    c = 1

    def h1(self): print(2, end=';')
    def _h2(self): print(3, end=';')

    def main(self):
        try: print(self.c, end=';')
        except: print('ERR', end=';')

        try: self.h1()
        except: print('ERR', end=';')

        try: self._h2()
        except: print('ERR', end=';')


class B(A):
    c = 4

    def h1(self): print(5, end=';')
    def _h2(self): print(6, end=';')

try: B().main()
except: print('ERR', end=';')
code

8. Що буде виведено за виконання?
code
def f(arg, n):
    arg[36] = 4
    n += -2

try:
    n = 6
    d = {17:4, 81:2, 17:6}
    f(d, n)
    print(n, end=';')
    for x in d.keys():
        print(x, sep=';', end=';')
except:
    print('ERR')
code

9. Що буде виведено за виконання?
code
class A:
    a = 1
    b = 2

    def __init__(self):
        b = 3

class B(A):
    b = 4

    def __init__(self):
        super().__init__()
        self.c = 7

obj = B()
obj.c = obj.a + 10
B.c = 13

try: print(obj.c, end= ';')
except: print('ERR', end=';')

try: print(A.c, end= ';')
except: print('ERR', end=';')

try: print(B.c, end= ';')
except: print('ERR', end=';')
code

10. 
Для графа на рисунку з вершини 0 запущено алгоритм пошуку в глибину, що будує кістякове дерево. Списки суміжності вершин переглядаються алгоритмом за зростанням міток вершин.\par
\begin{center}
	\adjustimage{max size={0.9\linewidth}{0.9\paperheight}}
	{../10/Traversals/215.png}
\end{center}
\par Які з наведених ребер входять до складу отриманого кістякового дерева?\par
1) (0, 1)\par
2) (1, 6)\par
3) (2, 4)\par
4) (2, 3)\par
5) (0, 3)\par
6) (3, 5)\par
{\vskip 15pt} У формі вибрати номери відповідних ребер.\par
%answer:1 2 3 4 6
\end {large}}
%end
{\smallskip \begin {small} Затверджено на засіданні кафедри теоретичної кібернетики, протокол \No 12 від   19.05.2020 р.\par\smallskip\smallskip
{\parbox {6.5 cm}{Зав. кафедри \dotfill Ю.В.~Крак}}\hspace {0.7cm}{\hbox to 6.5 cm{ Екзаменатор \dotfill Т.О.~Карнаух}} \end{small}}
%begin
{{\begin {small}\par
{ \center  Київський національний університет імені Тараса Шевченка \par}
{\parbox {5 cm}{Освітня програма \hfill}{\parbox {7 cm}{Прикладна математика, Системний аналіз \hfill}}\parbox {6 cm}{\hfill Семестр 2}}\par
{\parbox {5 cm} {Навчальний предмет \hfill}\parbox {7 cm}{Програмування \hfill}\parbox {6cm}{\hfill}\par}\vspace{-7pt} \end{small}}{\center Екзаменаційний білет \No \textbf 
{ 486 }\par}}
{
\begin {large}
1. Що буде виведено за виконання?
code
try:
    res = 4.0j==True
    if isinstance(res, bool):
        print(f'bool;{res}')
    elif isinstance(res, int):
        print(f'int;{res}')
    elif isinstance(res, float):
        print(f'float;{res:.2f}')
    else:
        print('???')
except:
    print('ERR')
code

2. Що буде виведено за виконання?
code
a = 8
b = 6
c = 1

def g(b):
    global c
    a = 4
    b = 3
    c = 5
    return a + b + c

a = 3
b = 0
c = 3

try:
    print(g(b), a, b, c, sep=';')
except:
    print('ERR', a, b, c, sep=';')
code

3. Що буде виведено за виконання?
code
try:
    t = 10
    while t>9:
        t += -3
    print(t, end=';')
except:
    print('ERR')
code

4. Що буде виведено за виконання?
code
try:
    for a in range(5, -12, -3):
        if a<5:
            break
        print(a, end=';')
    else:
        print(a,end=';')
    print(a)
except:
    print('ERR')
code

5. Що буде виведено за виконання?\par (Дійсні значення, якщо наявні, в цьому завданні будуть виводитись у форматі з фіксованою крапкою та, можливо, з доволі великою кількістю десяткових розрядів. У відповіді для дійсних значень у ЦЬОМУ завданні достатньо вказати один десятковий розряд після крапки, відкинувши решту розрядів без округлення. Наприклад, замість 13.66666666666667 достатньо вказати 13.6, замість 25.23 достатньо вказати 25.2)
code
try:
    print(2, end=';')
    print(5j==1j, end=';')
    print(8, end=';')
except ZeroDivisionError:
    print(3, end=';')
except Exception:
    print(9, end=';')
except:
    print('ERR', end=';')
else:
    print(0, end=';')
finally:
    print(9)
code

6. Що буде виведено за виконання?
code
def f(n):
    if n>9:
        return f(n//10)+1
    else:
        return 1

print(f(123456))
code

7. Що буде виведено за виконання?
code
class A:
    d = 1

    def f1(self): print(2, end=';')
    def _f2(self): print(3, end=';')

    def main(self):
        try: print(self.d, end=';')
        except: print('ERR', end=';')

        try: self.f1()
        except: print('ERR', end=';')

        try: self._f2()
        except: print('ERR', end=';')


class B(A):
    d = 4

    def f1(self): print(5, end=';')
    def _f2(self): print(6, end=';')

try: B().main()
except: print('ERR', end=';')
code

8. Що буде виведено за виконання?
code
def f(arg, n):
    n = 8
    tmp = arg
    tmp[3:-9] = [4,5]
    return len(tmp)>3

try:
    f = [10,11,12,13,14]
    n = 0
    f(f, n)
    print(n, end=';')
    for el in f:
        print(el, end=';')
except:
    print('ERR')
code

9. Що буде виведено за виконання?
code
class A:
    b = 1
    c = 2

    def __init__(self):
        b = 3

class B(A):
    b = 4

    def __init__(self):
        super().__init__()
        self.a = 7

obj = B()
obj.b = obj.c + 10
B.b = 13

try: print(obj.b, end= ';')
except: print('ERR', end=';')

try: print(A.b, end= ';')
except: print('ERR', end=';')

try: print(B.b, end= ';')
except: print('ERR', end=';')
code

10. 
Для графа на рисунку з вершини 6 запущено алгоритм пошуку в глибину, що будує кістякове дерево. Списки суміжності вершин переглядаються алгоритмом за зростанням міток вершин.\par
\begin{center}
	\adjustimage{max size={0.9\linewidth}{0.9\paperheight}}
	{../10/Traversals/227.png}
\end{center}
\par Які з наведених ребер входять до складу отриманого кістякового дерева?\par
1) (2, 3)\par
2) (1, 2)\par
3) (1, 4)\par
4) (4, 6)\par
5) (4, 5)\par
6) (1, 5)\par
{\vskip 15pt} У формі вибрати номери відповідних ребер.\par
%answer:1 5
\end {large}}
%end
{\smallskip \begin {small} Затверджено на засіданні кафедри теоретичної кібернетики, протокол \No 12 від   19.05.2020 р.\par\smallskip\smallskip
{\parbox {6.5 cm}{Зав. кафедри \dotfill Ю.В.~Крак}}\hspace {0.7cm}{\hbox to 6.5 cm{ Екзаменатор \dotfill Т.О.~Карнаух}} \end{small}}
%begin
{{\begin {small}\par
{ \center  Київський національний університет імені Тараса Шевченка \par}
{\parbox {5 cm}{Освітня програма \hfill}{\parbox {7 cm}{Прикладна математика, Системний аналіз \hfill}}\parbox {6 cm}{\hfill Семестр 2}}\par
{\parbox {5 cm} {Навчальний предмет \hfill}\parbox {7 cm}{Програмування \hfill}\parbox {6cm}{\hfill}\par}\vspace{-7pt} \end{small}}{\center Екзаменаційний білет \No \textbf 
{ 487 }\par}}
{
\begin {large}
1. Що буде виведено за виконання?
code
try:
    res = "True"<"2.0j"
    if isinstance(res, bool):
        print(f'bool;{res}')
    elif isinstance(res, int):
        print(f'int;{res}')
    elif isinstance(res, float):
        print(f'float;{res:.2f}')
    else:
        print('???')
except:
    print('ERR')
code

2. Що буде виведено за виконання?
code
a = 5
b = 1
c = 7

def h(b):
    a = 1
    b = 2
    c = 4
    return a + b + c

a = 2
b = 4
c = 9

try:
    print(h(b), a, b, c, sep=';')
except:
    print('ERR', a, b, c, sep=';')
code

3. Що буде виведено за виконання?
code
try:
    k = 5
    while k<=9:
        k += 1
    print(k, end=';')
except:
    print('ERR')
code

4. Що буде виведено за виконання?
code
try:
    for f in range(13, 8, -1):
        if f>7:
            continue
        print(f, end=';')
        if f<8:
            break
    else:
        print(f,end=';')
    print(f)
except:
    print('ERR')
code

5. Що буде виведено за виконання?\par (Дійсні значення, якщо наявні, в цьому завданні будуть виводитись у форматі з фіксованою крапкою та, можливо, з доволі великою кількістю десяткових розрядів. У відповіді для дійсних значень у ЦЬОМУ завданні достатньо вказати один десятковий розряд після крапки, відкинувши решту розрядів без округлення. Наприклад, замість 13.66666666666667 достатньо вказати 13.6, замість 25.23 достатньо вказати 25.2)
code
try:
    print(4, end=';')
    print(int('6'), end=';')
    print(7, end=';')
except KeyboardInterrupt:
    print(2, end=';')
except TypeError:
    print(1, end=';')
except:
    print('ERR', end=';')
else:
    print(8, end=';')
finally:
    print(3)
code

6. Що буде виведено за виконання?
code
def f(n):
    if n>9:
        f(n//100)
    print(n%10,end='')
 
f(987654)
code

7. Що буде виведено за виконання?
code
class A:
    a = 1

    def _f1(self): print(2, end=';')
    def f2(self): print(3, end=';')

    def main(self):
        try: print(self.a, end=';')
        except: print('ERR', end=';')

        try: self._f1()
        except: print('ERR', end=';')

        try: self.f2()
        except: print('ERR', end=';')


class B(A):
    a = 4

    def _f1(self): print(5, end=';')
    def f2(self): print(6, end=';')

try: B().main()
except: print('ERR', end=';')
code

8. Що буде виведено за виконання?
code
def f(arg, n):
    arg[78] = 4
    n *= 3

try:
    n = 5
    s = {84:8, 30:5, 78:3, 78:8}
    f(s, n)
    print(n, end=';')
    for x in s.items():
        print(x, sep=';', end=';')
except:
    print('ERR')
code

9. Що буде виведено за виконання?
code
class A:
    b = 1
    a = 2

    def __init__(self):
        b = 3

class B(A):
    b = 4

    def __init__(self):
        super().__init__()
        self.c = 7

obj = B()
obj.a = obj.b + 10
B.a = 13

try: print(obj.a, end= ';')
except: print('ERR', end=';')

try: print(A.a, end= ';')
except: print('ERR', end=';')

try: print(B.a, end= ';')
except: print('ERR', end=';')
code

10. 
Для графа на рисунку з вершини 2 запущено алгоритм пошуку в глибину, що будує кістякове дерево. Списки суміжності вершин переглядаються алгоритмом за зростанням міток вершин.\par
\begin{center}
	\adjustimage{max size={0.9\linewidth}{0.9\paperheight}}
	{../10/Traversals/104.png}
\end{center}
\par Які з наведених ребер входять до складу отриманого кістякового дерева?\par
1) (0, 4)\par
2) (3, 6)\par
3) (0, 5)\par
4) (1, 2)\par
5) (2, 6)\par
6) (0, 2)\par
{\vskip 15pt} У формі вибрати номери відповідних ребер.\par
%answer:1 2 3 6
\end {large}}
%end
{\smallskip \begin {small} Затверджено на засіданні кафедри теоретичної кібернетики, протокол \No 12 від   19.05.2020 р.\par\smallskip\smallskip
{\parbox {6.5 cm}{Зав. кафедри \dotfill Ю.В.~Крак}}\hspace {0.7cm}{\hbox to 6.5 cm{ Екзаменатор \dotfill Т.О.~Карнаух}} \end{small}}
%begin
{{\begin {small}\par
{ \center  Київський національний університет імені Тараса Шевченка \par}
{\parbox {5 cm}{Освітня програма \hfill}{\parbox {7 cm}{Прикладна математика, Системний аналіз \hfill}}\parbox {6 cm}{\hfill Семестр 2}}\par
{\parbox {5 cm} {Навчальний предмет \hfill}\parbox {7 cm}{Програмування \hfill}\parbox {6cm}{\hfill}\par}\vspace{-7pt} \end{small}}{\center Екзаменаційний білет \No \textbf 
{ 488 }\par}}
{
\begin {large}
1. Що буде виведено за виконання?
code
def f(n):
    print('f', end='')
    return n

try:
    print(f(-2) and f(3))
except:
    print('ERR')
code

2. Що буде виведено за виконання?
code
a = 3
b = 9
c = 5

def h(b):
    global c
    a = 4
    b += 5
    c = 3
    return a + b + c

a = 4
b = 2
c = 8

try:
    print(h(a), a, b, c, sep=';')
except:
    print('ERR', a, b, c, sep=';')
code

3. Що буде виведено за виконання?
code
try:
    n = 11
    while n>0:
        n += -3
    print(n, end=';')
except:
    print('ERR')
code

4. Що буде виведено за виконання?
code
try:
    for a in range(-10, -7, 2):
        if a<=5:
            continue
        print(a, end=';')
    else:
        print('end',end=';')
    print(a)
except:
    print('ERR')
code

5. Що буде виведено за виконання?\par (Дійсні значення, якщо наявні, в цьому завданні будуть виводитись у форматі з фіксованою крапкою та, можливо, з доволі великою кількістю десяткових розрядів. У відповіді для дійсних значень у ЦЬОМУ завданні достатньо вказати один десятковий розряд після крапки, відкинувши решту розрядів без округлення. Наприклад, замість 13.66666666666667 достатньо вказати 13.6, замість 25.23 достатньо вказати 25.2)
code
try:
    print(5, end=';')
    print(int('c4'), end=';')
    print(5, end=';')
except ValueError:
    print(3, end=';')
except ValueError:
    print(6, end=';')
except:
    print('ERR', end=';')
else:
    print(8, end=';')
finally:
    print(7)
code

6. Що буде виведено за виконання?
code
def f(n):
    print(n%10, end='')
    if n>9:
        f(n//10)
 
f(543216)
code

7. Що буде виведено за виконання?
code
class A:
    b = 1

    def _h1(self): print(2, end=';')
    def _h2(self): print(3, end=';')

    def main(self):
        try: print(self.b, end=';')
        except: print('ERR', end=';')

        try: self._h1()
        except: print('ERR', end=';')

        try: self._h2()
        except: print('ERR', end=';')


class B(A):
    b = 4

    def _h1(self): print(5, end=';')
    def _h2(self): print(6, end=';')

try: B().main()
except: print('ERR', end=';')
code

8. Що буде виведено за виконання?
code
def f(arg, n):
    n = 0
    tmp = arg
    tmp[:-7] = (14)
    return len(tmp)>3

try:
    e = [10,11,12,13,14]
    n = 4
    f(e, n)
    print(n, end=';')
    for el in e:
        print(el, end=';')
except:
    print('ERR')
code

9. Що буде виведено за виконання?
code
class A:
    c = 1
    a = 2

    def __init__(self):
        a = 3

class B(A):
    a = 4

    def __init__(self):
        super().__init__()
        self.b = 7

obj = B()
obj.c = obj.c + 10
B.c = 13

try: print(obj.c, end= ';')
except: print('ERR', end=';')

try: print(A.c, end= ';')
except: print('ERR', end=';')

try: print(B.c, end= ';')
except: print('ERR', end=';')
code

10. 
Для графа на рисунку з вершини 0 запущено алгоритм пошуку в глибину, що будує кістякове дерево. Списки суміжності вершин переглядаються алгоритмом за зростанням міток вершин.\par
\begin{center}
	\adjustimage{max size={0.9\linewidth}{0.9\paperheight}}
	{../10/Traversals/0.png}
\end{center}
\par Які з наведених ребер входять до складу отриманого кістякового дерева?\par
1) (0, 5)\par
2) (1, 5)\par
3) (1, 4)\par
4) (4, 5)\par
5) (0, 6)\par
6) (1, 3)\par
{\vskip 15pt} У формі вибрати номери відповідних ребер.\par
%answer:4 6
\end {large}}
%end
{\smallskip \begin {small} Затверджено на засіданні кафедри теоретичної кібернетики, протокол \No 12 від   19.05.2020 р.\par\smallskip\smallskip
{\parbox {6.5 cm}{Зав. кафедри \dotfill Ю.В.~Крак}}\hspace {0.7cm}{\hbox to 6.5 cm{ Екзаменатор \dotfill Т.О.~Карнаух}} \end{small}}
%begin
{{\begin {small}\par
{ \center  Київський національний університет імені Тараса Шевченка \par}
{\parbox {5 cm}{Освітня програма \hfill}{\parbox {7 cm}{Прикладна математика, Системний аналіз \hfill}}\parbox {6 cm}{\hfill Семестр 2}}\par
{\parbox {5 cm} {Навчальний предмет \hfill}\parbox {7 cm}{Програмування \hfill}\parbox {6cm}{\hfill}\par}\vspace{-7pt} \end{small}}{\center Екзаменаційний білет \No \textbf 
{ 489 }\par}}
{
\begin {large}
1. Що буде виведено за виконання?
code
try:
    res = "5"<=False
    if isinstance(res, bool):
        print(f'bool;{res}')
    elif isinstance(res, int):
        print(f'int;{res}')
    elif isinstance(res, float):
        print(f'float;{res:.2f}')
    else:
        print('???')
except:
    print('ERR')
code

2. Що буде виведено за виконання?
code
a = 6
b = 8
c = 1

def f(a):
    global c
    a = 3
    b = 5
    c = 1
    return a + b + c

a = 6
b = 3
c = 0

try:
    print(f(a), a, b, c, sep=';')
except:
    print('ERR', a, b, c, sep=';')
code

3. Що буде виведено за виконання?
code
try:
    m = 10
    while m>=8:
        m += -2
    print(m, end=';')
except:
    print('ERR')
code

4. Що буде виведено за виконання?
code
try:
    for c in range(0, -1, 3):
        if c>8:
            break
        print(c, end=';');c = -3
    else:
        print(c,end=';')
    print(c)
except:
    print('ERR')
code

5. Що буде виведено за виконання?\par (Дійсні значення, якщо наявні, в цьому завданні будуть виводитись у форматі з фіксованою крапкою та, можливо, з доволі великою кількістю десяткових розрядів. У відповіді для дійсних значень у ЦЬОМУ завданні достатньо вказати один десятковий розряд після крапки, відкинувши решту розрядів без округлення. Наприклад, замість 13.66666666666667 достатньо вказати 13.6, замість 25.23 достатньо вказати 25.2)
code
try:
    print(2, end=';')
    print(9%2, end=';')
    print(1, end=';')
except ZeroDivisionError:
    print(0, end=';')
except Exception:
    print(3, end=';')
except:
    print('ERR', end=';')
else:
    print(8, end=';')
finally:
    print(5)
code

6. Що буде виведено за виконання?
code
def f(n):
    if n>9: f(n//10)
    else:
        print(n%10, end='')
 
f(123456)
code

7. Що буде виведено за виконання?
code
class A:
    d = 1

    def __g1(self): print(2, end=';')
    def _g2(self): print(3, end=';')

    def main(self):
        try: print(self.d, end=';')
        except: print('ERR', end=';')

        try: self.__g1()
        except: print('ERR', end=';')

        try: self._g2()
        except: print('ERR', end=';')


class B(A):
    d = 4

    def __g1(self): print(5, end=';')
    def _g2(self): print(6, end=';')

try: B().main()
except: print('ERR', end=';')
code

8. Що буде виведено за виконання?
code
def f(arg, n):
    n = 2
    tmp = arg
    del tmp[5:-1]
    return len(tmp)>3

try:
    c = ['0','1','2','3','4']
    n = 6
    f(c, n)
    print(n, end=';')
    for el in c:
        print(el, end=';')
except:
    print('ERR')
code

9. Що буде виведено за виконання?
code
class A:
    a = 1
    c = 2

    def __init__(self):
        c = 3

class B(A):
    a = 4

    def __init__(self):
        super().__init__()
        self.b = 7

obj = B()
obj.a = obj.b + 10
B.a = 13

try: print(obj.a, end= ';')
except: print('ERR', end=';')

try: print(A.a, end= ';')
except: print('ERR', end=';')

try: print(B.a, end= ';')
except: print('ERR', end=';')
code

10. 
Для графа на рисунку з вершини 2 запущено алгоритм пошуку в глибину, що будує кістякове дерево. Списки суміжності вершин переглядаються алгоритмом за зростанням міток вершин.\par
\begin{center}
	\adjustimage{max size={0.9\linewidth}{0.9\paperheight}}
	{../10/Traversals/105.png}
\end{center}
\par Які з наведених ребер входять до складу отриманого кістякового дерева?\par
1) (0, 1)\par
2) (5, 6)\par
3) (4, 5)\par
4) (3, 6)\par
5) (1, 4)\par
6) (4, 6)\par
{\vskip 15pt} У формі вибрати номери відповідних ребер.\par
%answer:1 2
\end {large}}
%end
{\smallskip \begin {small} Затверджено на засіданні кафедри теоретичної кібернетики, протокол \No 12 від   19.05.2020 р.\par\smallskip\smallskip
{\parbox {6.5 cm}{Зав. кафедри \dotfill Ю.В.~Крак}}\hspace {0.7cm}{\hbox to 6.5 cm{ Екзаменатор \dotfill Т.О.~Карнаух}} \end{small}}
%begin
{{\begin {small}\par
{ \center  Київський національний університет імені Тараса Шевченка \par}
{\parbox {5 cm}{Освітня програма \hfill}{\parbox {7 cm}{Прикладна математика, Системний аналіз \hfill}}\parbox {6 cm}{\hfill Семестр 2}}\par
{\parbox {5 cm} {Навчальний предмет \hfill}\parbox {7 cm}{Програмування \hfill}\parbox {6cm}{\hfill}\par}\vspace{-7pt} \end{small}}{\center Екзаменаційний білет \No \textbf 
{ 490 }\par}}
{
\begin {large}
1. Що буде виведено за виконання?
code
try:
    res = 7.0>2j
    if isinstance(res, bool):
        print(f'bool;{res}')
    elif isinstance(res, int):
        print(f'int;{res}')
    elif isinstance(res, float):
        print(f'float;{res:.2f}')
    else:
        print('???')
except:
    print('ERR')
code

2. Що буде виведено за виконання?
code
a = 4
b = 7
c = 9

def h(b):
    global c
    a = 4
    b = 3
    c = 2
    return a + b + c

a = 2
b = 5
c = 8

try:
    print(h(a), a, b, c, sep=';')
except:
    print('ERR', a, b, c, sep=';')
code

3. Що буде виведено за виконання?
code
try:
    n = 1
    while n<11:
        n += 2
    print(n, end=';')
except:
    print('ERR')
code

4. Що буде виведено за виконання?
code
try:
    for b in range(7, -2, -2):
        if b<=3:
            break
        print(b, end=';');b = 7
    else:
        print(b,end=';')
    print(b)
except:
    print('ERR')
code

5. Що буде виведено за виконання?\par (Дійсні значення, якщо наявні, в цьому завданні будуть виводитись у форматі з фіксованою крапкою та, можливо, з доволі великою кількістю десяткових розрядів. У відповіді для дійсних значень у ЦЬОМУ завданні достатньо вказати один десятковий розряд після крапки, відкинувши решту розрядів без округлення. Наприклад, замість 13.66666666666667 достатньо вказати 13.6, замість 25.23 достатньо вказати 25.2)
code
try:
    print(9, end=';')
    print(7/0.0, end=';')
    print(4, end=';')
except TypeError:
    print(6, end=';')
except KeyboardInterrupt:
    print(2, end=';')
except:
    print('ERR', end=';')
else:
    print(1, end=';')
finally:
    print(0)
code

6. Що буде виведено за виконання?
code
def f(n):
    if n>9:
        return f(n//100)+1
    else:
        return 1

print(f(123456))
code

7. Що буде виведено за виконання?
code
class A:
    b = 1

    def _f1(self): print(2, end=';')
    def __f2(self): print(3, end=';')

    def main(self):
        try: print(self.b, end=';')
        except: print('ERR', end=';')

        try: self._f1()
        except: print('ERR', end=';')

        try: self.__f2()
        except: print('ERR', end=';')


class B(A):
    b = 4

    def _f1(self): print(5, end=';')
    def __f2(self): print(6, end=';')

try: B().main()
except: print('ERR', end=';')
code

8. Що буде виведено за виконання?
code
def f(arg, n):
    arg[12] = 4
    n = 5
    arg = list(arg.keys())

try:
    n = 0
    s = {12:7, 73:9, 73:1}
    f(s, n)
    print(n, end=';')
    for x in s.keys():
        print(x, sep=';', end=';')
except:
    print('ERR')
code

9. Що буде виведено за виконання?
code
class A:
    c = 1
    b = 2

    def __init__(self):
        c = 3

class B(A):
    b = 4

    def __init__(self):
        super().__init__()
        self.a = 7

obj = B()
obj.a = obj.b + 10
B.a = 13

try: print(obj.a, end= ';')
except: print('ERR', end=';')

try: print(A.a, end= ';')
except: print('ERR', end=';')

try: print(B.a, end= ';')
except: print('ERR', end=';')
code

10. 
Для графа на рисунку з вершини 2 запущено алгоритм пошуку в глибину, що будує кістякове дерево. Списки суміжності вершин переглядаються алгоритмом за зростанням міток вершин.\par
\begin{center}
	\adjustimage{max size={0.9\linewidth}{0.9\paperheight}}
	{../10/Traversals/195.png}
\end{center}
\par Які з наведених ребер входять до складу отриманого кістякового дерева?\par
1) (0, 2)\par
2) (1, 6)\par
3) (1, 5)\par
4) (1, 3)\par
5) (2, 5)\par
6) (2, 6)\par
{\vskip 15pt} У формі вибрати номери відповідних ребер.\par
%answer:1 3 4
\end {large}}
%end
{\smallskip \begin {small} Затверджено на засіданні кафедри теоретичної кібернетики, протокол \No 12 від   19.05.2020 р.\par\smallskip\smallskip
{\parbox {6.5 cm}{Зав. кафедри \dotfill Ю.В.~Крак}}\hspace {0.7cm}{\hbox to 6.5 cm{ Екзаменатор \dotfill Т.О.~Карнаух}} \end{small}}
%begin
{{\begin {small}\par
{ \center  Київський національний університет імені Тараса Шевченка \par}
{\parbox {5 cm}{Освітня програма \hfill}{\parbox {7 cm}{Прикладна математика, Системний аналіз \hfill}}\parbox {6 cm}{\hfill Семестр 2}}\par
{\parbox {5 cm} {Навчальний предмет \hfill}\parbox {7 cm}{Програмування \hfill}\parbox {6cm}{\hfill}\par}\vspace{-7pt} \end{small}}{\center Екзаменаційний білет \No \textbf 
{ 491 }\par}}
{
\begin {large}
1. Що буде виведено за виконання?
code
def f(n):
    print('f', end='')
    return n<=-2

try:
    print(f(-4) or f(4))
except:
    print('ERR')
code

2. Що буде виведено за виконання?
code
a = 6
b = 9
c = 7

def g(a):
    a *= 2
    b = 3
    c = 4
    return a + b + c

a = 0
b = 2
c = 8

try:
    print(g(b), a, b, c, sep=';')
except:
    print('ERR', a, b, c, sep=';')
code

3. Що буде виведено за виконання?
code
try:
    i = 15
    while i>=7:
        i += -3
    print(i, end=';')
except:
    print('ERR')
code

4. Що буде виведено за виконання?
code
try:
    for e in range(11, -3, 3):
        if e<4:
            continue
        print(e, end=';')
    else:
        print(e,end=';')
    print(e)
except:
    print('ERR')
code

5. Що буде виведено за виконання?\par (Дійсні значення, якщо наявні, в цьому завданні будуть виводитись у форматі з фіксованою крапкою та, можливо, з доволі великою кількістю десяткових розрядів. У відповіді для дійсних значень у ЦЬОМУ завданні достатньо вказати один десятковий розряд після крапки, відкинувши решту розрядів без округлення. Наприклад, замість 13.66666666666667 достатньо вказати 13.6, замість 25.23 достатньо вказати 25.2)
code
try:
    print(8, end=';')
    print(int('d4'), end=';')
    print(5, end=';')
except TypeError:
    print(3, end=';')
except ValueError:
    print(2, end=';')
except:
    print('ERR', end=';')
else:
    print(6, end=';')
finally:
    print(1)
code

6. Що буде виведено за виконання?
code
def f(n):
    if n<=9:
        print(n%10, end='')
    else:
        f(n//10)

f(123456)
code

7. Що буде виведено за виконання?
code
class A:
    a = 1

    def _g1(self): print(2, end=';')
    def _g2(self): print(3, end=';')

    def main(self):
        try: print(self.a, end=';')
        except: print('ERR', end=';')

        try: self._g1()
        except: print('ERR', end=';')

        try: self._g2()
        except: print('ERR', end=';')


class B(A):
    a = 4

    def _g1(self): print(5, end=';')
    def _g2(self): print(6, end=';')

try: B().main()
except: print('ERR', end=';')
code

8. Що буде виведено за виконання?
code
def f(arg, n):
    n = 8
    tmp = arg
    tmp[0:] = 'xy'
    return len(tmp)>3

try:
    f = [10,11,12,13,14]
    n = 9
    f(f, n)
    print(n, end=';')
    for el in f:
        print(el, end=';')
except:
    print('ERR')
code

9. Що буде виведено за виконання?
code
class A:
    a = 1
    b = 2

    def __init__(self):
        a = 3

class B(A):
    a = 4

    def __init__(self):
        super().__init__()
        self.c = 7

obj = B()
obj.c = obj.c + 10
B.c = 13

try: print(obj.c, end= ';')
except: print('ERR', end=';')

try: print(A.c, end= ';')
except: print('ERR', end=';')

try: print(B.c, end= ';')
except: print('ERR', end=';')
code

10. 
Для графа на рисунку з вершини 0 запущено алгоритм пошуку в глибину, що будує кістякове дерево. Списки суміжності вершин переглядаються алгоритмом за зростанням міток вершин.\par
\begin{center}
	\adjustimage{max size={0.9\linewidth}{0.9\paperheight}}
	{../10/Traversals/150.png}
\end{center}
\par Які з наведених ребер входять до складу отриманого кістякового дерева?\par
1) (0, 5)\par
2) (1, 2)\par
3) (3, 5)\par
4) (5, 6)\par
5) (2, 6)\par
6) (0, 4)\par
{\vskip 15pt} У формі вибрати номери відповідних ребер.\par
%answer:2 4 5
\end {large}}
%end
{\smallskip \begin {small} Затверджено на засіданні кафедри теоретичної кібернетики, протокол \No 12 від   19.05.2020 р.\par\smallskip\smallskip
{\parbox {6.5 cm}{Зав. кафедри \dotfill Ю.В.~Крак}}\hspace {0.7cm}{\hbox to 6.5 cm{ Екзаменатор \dotfill Т.О.~Карнаух}} \end{small}}
%begin
{{\begin {small}\par
{ \center  Київський національний університет імені Тараса Шевченка \par}
{\parbox {5 cm}{Освітня програма \hfill}{\parbox {7 cm}{Прикладна математика, Системний аналіз \hfill}}\parbox {6 cm}{\hfill Семестр 2}}\par
{\parbox {5 cm} {Навчальний предмет \hfill}\parbox {7 cm}{Програмування \hfill}\parbox {6cm}{\hfill}\par}\vspace{-7pt} \end{small}}{\center Екзаменаційний білет \No \textbf 
{ 492 }\par}}
{
\begin {large}
1. Що буде виведено за виконання?
code
def f(n):
    print('f', end='')
    return n!=3

try:
    print(f(1) and f(9))
except:
    print('ERR')
code

2. Що буде виведено за виконання?
code
a = 5
b = 0
c = 8

def g(a):
    a = 4
    b = 3
    c = 1
    return a + b + c

a = 2
b = 7
c = 9

try:
    print(g(b), a, b, c, sep=';')
except:
    print('ERR', a, b, c, sep=';')
code

3. Що буде виведено за виконання?
code
try:
    n = 7
    while n<=7:
        n += 1
    print(n, end=';')
except:
    print('ERR')
code

4. Що буде виведено за виконання?
code
try:
    for d in range(2, 3, -1):
        if d>=6:
            continue
        print(d, end=';')
    else:
        print(d,end=';')
    print(d)
except:
    print('ERR')
code

5. Що буде виведено за виконання?\par (Дійсні значення, якщо наявні, в цьому завданні будуть виводитись у форматі з фіксованою крапкою та, можливо, з доволі великою кількістю десяткових розрядів. У відповіді для дійсних значень у ЦЬОМУ завданні достатньо вказати один десятковий розряд після крапки, відкинувши решту розрядів без округлення. Наприклад, замість 13.66666666666667 достатньо вказати 13.6, замість 25.23 достатньо вказати 25.2)
code
try:
    print(7, end=';')
    print(0j<=9j, end=';')
    print(9, end=';')
except KeyboardInterrupt:
    print(1, end=';')
except ZeroDivisionError:
    print(0, end=';')
except:
    print('ERR', end=';')
else:
    print(3, end=';')
finally:
    print(2)
code

6. Що буде виведено за виконання?
code
def f(n):
    if n>2:
        f(n//2)
    else:
        print(n%2, end='')

f(16)
code

7. Що буде виведено за виконання?
code
class A:
    c = 1

    def __h1(self): print(2, end=';')
    def _h2(self): print(3, end=';')

    def main(self):
        try: print(self.c, end=';')
        except: print('ERR', end=';')

        try: self.__h1()
        except: print('ERR', end=';')

        try: self._h2()
        except: print('ERR', end=';')


class B(A):
    c = 4

    def __h1(self): print(5, end=';')
    def _h2(self): print(6, end=';')

try: B().main()
except: print('ERR', end=';')
code

8. Що буде виведено за виконання?
code
def f(arg, n):
    arg[84] = 7
    n -= 3

try:
    n = 3
    t = {87:7, 26:9, 26:0}
    f(t, n)
    print(n, end=';')
    for x in t :
        print(x, sep=';', end=';')
except:
    print('ERR')
code

9. Що буде виведено за виконання?
code
class A:
    b = 1
    c = 2

    def __init__(self):
        c = 3

class B(A):
    c = 4

    def __init__(self):
        super().__init__()
        self.a = 7

obj = B()
obj.a = obj.b + 10
B.a = 13

try: print(obj.a, end= ';')
except: print('ERR', end=';')

try: print(A.a, end= ';')
except: print('ERR', end=';')

try: print(B.a, end= ';')
except: print('ERR', end=';')
code

10. 
Для графа на рисунку з вершини 6 запущено алгоритм пошуку в глибину, що будує кістякове дерево. Списки суміжності вершин переглядаються алгоритмом за зростанням міток вершин.\par
\begin{center}
	\adjustimage{max size={0.9\linewidth}{0.9\paperheight}}
	{../10/Traversals/290.png}
\end{center}
\par Які з наведених ребер входять до складу отриманого кістякового дерева?\par
1) (2, 3)\par
2) (0, 3)\par
3) (1, 3)\par
4) (4, 6)\par
5) (2, 4)\par
6) (0, 6)\par
{\vskip 15pt} У формі вибрати номери відповідних ребер.\par
%answer:1 2 3 5 6
\end {large}}
%end
{\smallskip \begin {small} Затверджено на засіданні кафедри теоретичної кібернетики, протокол \No 12 від   19.05.2020 р.\par\smallskip\smallskip
{\parbox {6.5 cm}{Зав. кафедри \dotfill Ю.В.~Крак}}\hspace {0.7cm}{\hbox to 6.5 cm{ Екзаменатор \dotfill Т.О.~Карнаух}} \end{small}}
%begin
{{\begin {small}\par
{ \center  Київський національний університет імені Тараса Шевченка \par}
{\parbox {5 cm}{Освітня програма \hfill}{\parbox {7 cm}{Прикладна математика, Системний аналіз \hfill}}\parbox {6 cm}{\hfill Семестр 2}}\par
{\parbox {5 cm} {Навчальний предмет \hfill}\parbox {7 cm}{Програмування \hfill}\parbox {6cm}{\hfill}\par}\vspace{-7pt} \end{small}}{\center Екзаменаційний білет \No \textbf 
{ 493 }\par}}
{
\begin {large}
1. Що буде виведено за виконання?
code
try:
    res = 3**-2+1
    if isinstance(res, bool):
        print(f'bool;{res}')
    elif isinstance(res, int):
        print(f'int;{res}')
    elif isinstance(res, float):
        print(f'float;{res:.2f}')
    else:
        print('???')
except:
    print('ERR')
code

2. Що буде виведено за виконання?
code
a = 4
b = 3
c = 1

def h(a):
    a = 2
    b = 5
    c = 2
    return a + b + c

a = 6
b = 9
c = 3

try:
    print(h(a), a, b, c, sep=';')
except:
    print('ERR', a, b, c, sep=';')
code

3. Що буде виведено за виконання?
code
try:
    i = 4
    while i<10:
        i += 2
    print(i, end=';')
except:
    print('ERR')
code

4. Що буде виведено за виконання?
code
try:
    for d in range(6, -4, -2):
        if d<5:
            break
        print(d, end=';');d = -8
    else:
        print(d,end=';')
    print(d)
except:
    print('ERR')
code

5. Що буде виведено за виконання?\par (Дійсні значення, якщо наявні, в цьому завданні будуть виводитись у форматі з фіксованою крапкою та, можливо, з доволі великою кількістю десяткових розрядів. У відповіді для дійсних значень у ЦЬОМУ завданні достатньо вказати один десятковий розряд після крапки, відкинувши решту розрядів без округлення. Наприклад, замість 13.66666666666667 достатньо вказати 13.6, замість 25.23 достатньо вказати 25.2)
code
try:
    print(8, end=';')
    print(4%3, end=';')
    print(9, end=';')
except Exception:
    print(5, end=';')
except ValueError:
    print(7, end=';')
except:
    print('ERR', end=';')
else:
    print(6, end=';')
finally:
    print(6)
code

6. Що буде виведено за виконання?
code
def f(s):
    for el in s:
        el += 1
        return s
 
s = [1, 2, 3]
s = f(s)
print(s)
code

7. Що буде виведено за виконання?
code
class A:
    c = 1

    def _h1(self): print(2, end=';')
    def h2(self): print(3, end=';')

    def main(self):
        try: print(self.c, end=';')
        except: print('ERR', end=';')

        try: self._h1()
        except: print('ERR', end=';')

        try: self.h2()
        except: print('ERR', end=';')


class B(A):
    c = 4

    def _h1(self): print(5, end=';')
    def h2(self): print(6, end=';')

try: B().main()
except: print('ERR', end=';')
code

8. Що буде виведено за виконання?
code
def f(arg, n):
    n = 5
    tmp = arg
    tmp[-4:1] = 'bd'
    return len(tmp)>3

try:
    h = ['0','1','2','3','4']
    n = 7
    f(h, n)
    print(n, end=';')
    for el in h:
        print(el, end=';')
except:
    print('ERR')
code

9. Що буде виведено за виконання?
code
class A:
    a = 1
    c = 2

    def __init__(self):
        c = 3

class B(A):
    c = 4

    def __init__(self):
        super().__init__()
        self.b = 7

obj = B()
obj.b = obj.a + 10
B.b = 13

try: print(obj.b, end= ';')
except: print('ERR', end=';')

try: print(A.b, end= ';')
except: print('ERR', end=';')

try: print(B.b, end= ';')
except: print('ERR', end=';')
code

10. 
Для графа на рисунку з вершини 4 запущено алгоритм пошуку в глибину, що будує кістякове дерево. Списки суміжності вершин переглядаються алгоритмом за зростанням міток вершин.\par
\begin{center}
	\adjustimage{max size={0.9\linewidth}{0.9\paperheight}}
	{../10/Traversals/80.png}
\end{center}
\par Які з наведених ребер входять до складу отриманого кістякового дерева?\par
1) (2, 3)\par
2) (5, 6)\par
3) (0, 1)\par
4) (1, 5)\par
5) (2, 5)\par
6) (3, 4)\par
{\vskip 15pt} У формі вибрати номери відповідних ребер.\par
%answer:1 2 3 4 6
\end {large}}
%end
{\smallskip \begin {small} Затверджено на засіданні кафедри теоретичної кібернетики, протокол \No 12 від   19.05.2020 р.\par\smallskip\smallskip
{\parbox {6.5 cm}{Зав. кафедри \dotfill Ю.В.~Крак}}\hspace {0.7cm}{\hbox to 6.5 cm{ Екзаменатор \dotfill Т.О.~Карнаух}} \end{small}}
%begin
{{\begin {small}\par
{ \center  Київський національний університет імені Тараса Шевченка \par}
{\parbox {5 cm}{Освітня програма \hfill}{\parbox {7 cm}{Прикладна математика, Системний аналіз \hfill}}\parbox {6 cm}{\hfill Семестр 2}}\par
{\parbox {5 cm} {Навчальний предмет \hfill}\parbox {7 cm}{Програмування \hfill}\parbox {6cm}{\hfill}\par}\vspace{-7pt} \end{small}}{\center Екзаменаційний білет \No \textbf 
{ 494 }\par}}
{
\begin {large}
1. Що буде виведено за виконання?
code
def f(n):
    print('f', end='')
    return n

try:
    print(f(-1) or f(8))
except:
    print('ERR')
code

2. Що буде виведено за виконання?
code
a = 1
b = 5
c = 3

def h(b):
    global c
    a = 1
    b *= 5
    c = 4
    return a + b + c

a = 4
b = 2
c = 3

try:
    print(h(a), a, b, c, sep=';')
except:
    print('ERR', a, b, c, sep=';')
code

3. Що буде виведено за виконання?
code
try:
    k = 0
    while k<8:
        k += 2
    print(k, end=';')
except:
    print('ERR')
code

4. Що буде виведено за виконання?
code
try:
    for e in range(1, -5, -3):
        if e>3:
            continue
        print(e, end=';')
    else:
        print(e,end=';')
    print(e)
except:
    print('ERR')
code

5. Що буде виведено за виконання?\par (Дійсні значення, якщо наявні, в цьому завданні будуть виводитись у форматі з фіксованою крапкою та, можливо, з доволі великою кількістю десяткових розрядів. У відповіді для дійсних значень у ЦЬОМУ завданні достатньо вказати один десятковий розряд після крапки, відкинувши решту розрядів без округлення. Наприклад, замість 13.66666666666667 достатньо вказати 13.6, замість 25.23 достатньо вказати 25.2)
code
try:
    print(9, end=';')
    print(1//0, end=';')
    print(5, end=';')
except Exception:
    print(8, end=';')
except TypeError:
    print(4, end=';')
except:
    print('ERR', end=';')
else:
    print(3, end=';')
finally:
    print(7)
code

6. Що буде виведено за виконання?
code
def f(n):
    if n>9:
        f(n//10)
    else:
        print(n%10, end='')

f(543216)
code

7. Що буде виведено за виконання?
code
class A:
    a = 1

    def f1(self): print(2, end=';')
    def _f2(self): print(3, end=';')

    def main(self):
        try: print(self.a, end=';')
        except: print('ERR', end=';')

        try: self.f1()
        except: print('ERR', end=';')

        try: self._f2()
        except: print('ERR', end=';')


class B(A):
    a = 4

    def f1(self): print(5, end=';')
    def _f2(self): print(6, end=';')

try: B().main()
except: print('ERR', end=';')
code

8. Що буде виведено за виконання?
code
def f(arg, n):
    arg[43] = 9
    n = 8
    arg = tuple(arg.items())

try:
    n = 3
    d = {43:7, 62:6, 88:8, 88:7}
    f(d, n)
    print(n, end=';')
    for x in d.items():
        print(x, sep=';', end=';')
except:
    print('ERR')
code

9. Що буде виведено за виконання?
code
class A:
    c = 1
    b = 2

    def __init__(self):
        c = 3

class B(A):
    c = 4

    def __init__(self):
        super().__init__()
        self.a = 7

obj = B()
obj.c = obj.a + 10
B.c = 13

try: print(obj.c, end= ';')
except: print('ERR', end=';')

try: print(A.c, end= ';')
except: print('ERR', end=';')

try: print(B.c, end= ';')
except: print('ERR', end=';')
code

10. 
Для графа на рисунку з вершини 6 запущено алгоритм пошуку в глибину, що будує кістякове дерево. Списки суміжності вершин переглядаються алгоритмом за зростанням міток вершин.\par
\begin{center}
	\adjustimage{max size={0.9\linewidth}{0.9\paperheight}}
	{../10/Traversals/173.png}
\end{center}
\par Які з наведених ребер входять до складу отриманого кістякового дерева?\par
1) (3, 4)\par
2) (5, 6)\par
3) (0, 4)\par
4) (1, 2)\par
5) (2, 3)\par
6) (1, 3)\par
{\vskip 15pt} У формі вибрати номери відповідних ребер.\par
%answer:1 4 6
\end {large}}
%end
{\smallskip \begin {small} Затверджено на засіданні кафедри теоретичної кібернетики, протокол \No 12 від   19.05.2020 р.\par\smallskip\smallskip
{\parbox {6.5 cm}{Зав. кафедри \dotfill Ю.В.~Крак}}\hspace {0.7cm}{\hbox to 6.5 cm{ Екзаменатор \dotfill Т.О.~Карнаух}} \end{small}}
%begin
{{\begin {small}\par
{ \center  Київський національний університет імені Тараса Шевченка \par}
{\parbox {5 cm}{Освітня програма \hfill}{\parbox {7 cm}{Прикладна математика, Системний аналіз \hfill}}\parbox {6 cm}{\hfill Семестр 2}}\par
{\parbox {5 cm} {Навчальний предмет \hfill}\parbox {7 cm}{Програмування \hfill}\parbox {6cm}{\hfill}\par}\vspace{-7pt} \end{small}}{\center Екзаменаційний білет \No \textbf 
{ 495 }\par}}
{
\begin {large}
1. Що буде виведено за виконання?
code
def f(n):
    print('f', end='')
    return n

try:
    print(f(-7) or f(-8))
except:
    print('ERR')
code

2. Що буде виведено за виконання?
code
a = 5
b = 4
c = 6

def g(a):
    a += 5
    b = 2
    c = 3
    return a + b + c

a = 1
b = 7
c = 8

try:
    print(g(b), a, b, c, sep=';')
except:
    print('ERR', a, b, c, sep=';')
code

3. Що буде виведено за виконання?
code
try:
    t = 2
    while t<=5:
        t += 1
    print(t, end=';')
except:
    print('ERR')
code

4. Що буде виведено за виконання?
code
try:
    for b in range(9, -8, 2):
        if b<=6:
            continue
        print(b, end=';')
        if b>=4:
            break
    else:
        print('end',end=';')
    print(b)
except:
    print('ERR')
code

5. Що буде виведено за виконання?\par (Дійсні значення, якщо наявні, в цьому завданні будуть виводитись у форматі з фіксованою крапкою та, можливо, з доволі великою кількістю десяткових розрядів. У відповіді для дійсних значень у ЦЬОМУ завданні достатньо вказати один десятковий розряд після крапки, відкинувши решту розрядів без округлення. Наприклад, замість 13.66666666666667 достатньо вказати 13.6, замість 25.23 достатньо вказати 25.2)
code
try:
    print(2, end=';')
    print(int('0'), end=';')
    print(4, end=';')
except KeyboardInterrupt:
    print(2, end=';')
except ZeroDivisionError:
    print(8, end=';')
except:
    print('ERR', end=';')
else:
    print(7, end=';')
finally:
    print(6)
code

6. Що буде виведено за виконання?
code
def f(n):
    if n<=9:
        print(n%10,end='')
    else:
        f(n//10)

f(543216)
code

7. Що буде виведено за виконання?
code
class A:
    d = 1

    def _g1(self): print(2, end=';')
    def __g2(self): print(3, end=';')

    def main(self):
        try: print(self.d, end=';')
        except: print('ERR', end=';')

        try: self._g1()
        except: print('ERR', end=';')

        try: self.__g2()
        except: print('ERR', end=';')


class B(A):
    d = 4

    def _g1(self): print(5, end=';')
    def __g2(self): print(6, end=';')

try: B().main()
except: print('ERR', end=';')
code

8. Що буде виведено за виконання?
code
def f(arg, n):
    n = 1
    tmp = arg
    del tmp[-3:2]
    return len(tmp)>3

try:
    b = ['0','1','2','3','4']
    n = 3
    f(b, n)
    print(n, end=';')
    for el in b:
        print(el, end=';')
except:
    print('ERR')
code

9. Що буде виведено за виконання?
code
class A:
    c = 1
    a = 2

    def __init__(self):
        c = 3

class B(A):
    a = 4

    def __init__(self):
        super().__init__()
        self.b = 7

obj = B()
obj.b = obj.c + 10
B.b = 13

try: print(obj.b, end= ';')
except: print('ERR', end=';')

try: print(A.b, end= ';')
except: print('ERR', end=';')

try: print(B.b, end= ';')
except: print('ERR', end=';')
code

10. 
Для графа на рисунку з вершини 0 запущено алгоритм пошуку в глибину, що будує кістякове дерево. Списки суміжності вершин переглядаються алгоритмом за зростанням міток вершин.\par
\begin{center}
	\adjustimage{max size={0.9\linewidth}{0.9\paperheight}}
	{../10/Traversals/134.png}
\end{center}
\par Які з наведених ребер входять до складу отриманого кістякового дерева?\par
1) (1, 3)\par
2) (2, 6)\par
3) (0, 6)\par
4) (2, 5)\par
5) (2, 4)\par
6) (0, 3)\par
{\vskip 15pt} У формі вибрати номери відповідних ребер.\par
%answer:1 2 4 5 6
\end {large}}
%end
{\smallskip \begin {small} Затверджено на засіданні кафедри теоретичної кібернетики, протокол \No 12 від   19.05.2020 р.\par\smallskip\smallskip
{\parbox {6.5 cm}{Зав. кафедри \dotfill Ю.В.~Крак}}\hspace {0.7cm}{\hbox to 6.5 cm{ Екзаменатор \dotfill Т.О.~Карнаух}} \end{small}}
%begin
{{\begin {small}\par
{ \center  Київський національний університет імені Тараса Шевченка \par}
{\parbox {5 cm}{Освітня програма \hfill}{\parbox {7 cm}{Прикладна математика, Системний аналіз \hfill}}\parbox {6 cm}{\hfill Семестр 2}}\par
{\parbox {5 cm} {Навчальний предмет \hfill}\parbox {7 cm}{Програмування \hfill}\parbox {6cm}{\hfill}\par}\vspace{-7pt} \end{small}}{\center Екзаменаційний білет \No \textbf 
{ 496 }\par}}
{
\begin {large}
1. Що буде виведено за виконання?
code
try:
    res = 4**-1.0*2
    if isinstance(res, bool):
        print(f'bool;{res}')
    elif isinstance(res, int):
        print(f'int;{res}')
    elif isinstance(res, float):
        print(f'float;{res:.2f}')
    else:
        print('???')
except:
    print('ERR')
code

2. Що буде виведено за виконання?
code
a = 5
b = 7
c = 2

def f(b):
    global c
    a = 3
    b = 1
    c = 4
    return a + b + c

a = 0
b = 8
c = 6

try:
    print(f(b), a, b, c, sep=';')
except:
    print('ERR', a, b, c, sep=';')
code

3. Що буде виведено за виконання?
code
try:
    j = 1
    while j<=12:
        j += 1
    print(j, end=';')
except:
    print('ERR')
code

4. Що буде виведено за виконання?
code
try:
    for a in range(3, 12, -1):
        if a>=4:
            continue
        print(a, end=';')
    else:
        print(a,end=';')
    print(a)
except:
    print('ERR')
code

5. Що буде виведено за виконання?\par (Дійсні значення, якщо наявні, в цьому завданні будуть виводитись у форматі з фіксованою крапкою та, можливо, з доволі великою кількістю десяткових розрядів. У відповіді для дійсних значень у ЦЬОМУ завданні достатньо вказати один десятковий розряд після крапки, відкинувши решту розрядів без округлення. Наприклад, замість 13.66666666666667 достатньо вказати 13.6, замість 25.23 достатньо вказати 25.2)
code
try:
    print(9, end=';')
    print(0//0.0, end=';')
    print(1, end=';')
except KeyboardInterrupt:
    print(5, end=';')
except ValueError:
    print(3, end=';')
except:
    print('ERR', end=';')
else:
    print(4, end=';')
finally:
    print(5)
code

6. Що буде виведено за виконання?
code
def f(n):
    if n>9:
        f(n//10)
    print(n%10, end='')
 
f(543216)
code

7. Що буде виведено за виконання?
code
class A:
    b = 1

    def _g1(self): print(2, end=';')
    def g2(self): print(3, end=';')

    def main(self):
        try: print(self.b, end=';')
        except: print('ERR', end=';')

        try: self._g1()
        except: print('ERR', end=';')

        try: self.g2()
        except: print('ERR', end=';')


class B(A):
    b = 4

    def _g1(self): print(5, end=';')
    def g2(self): print(6, end=';')

try: B().main()
except: print('ERR', end=';')
code

8. Що буде виведено за виконання?
code
def f(arg, n):
    arg[60] = 1
    n = 7
    arg = list(arg.values())

try:
    n = 2
    s = {74:5, 49:8, 49:5}
    f(s, n)
    print(n, end=';')
    for x in s.keys():
        print(x, sep=';', end=';')
except:
    print('ERR')
code

9. Що буде виведено за виконання?
code
class A:
    a = 1
    b = 2

    def __init__(self):
        b = 3

class B(A):
    a = 4

    def __init__(self):
        super().__init__()
        self.c = 7

obj = B()
obj.c = obj.a + 10
B.c = 13

try: print(obj.c, end= ';')
except: print('ERR', end=';')

try: print(A.c, end= ';')
except: print('ERR', end=';')

try: print(B.c, end= ';')
except: print('ERR', end=';')
code

10. 
Для графа на рисунку з вершини 6 запущено алгоритм пошуку в глибину, що будує кістякове дерево. Списки суміжності вершин переглядаються алгоритмом за зростанням міток вершин.\par
\begin{center}
	\adjustimage{max size={0.9\linewidth}{0.9\paperheight}}
	{../10/Traversals/85.png}
\end{center}
\par Які з наведених ребер входять до складу отриманого кістякового дерева?\par
1) (0, 5)\par
2) (0, 6)\par
3) (1, 6)\par
4) (5, 6)\par
5) (1, 3)\par
6) (0, 2)\par
{\vskip 15pt} У формі вибрати номери відповідних ребер.\par
%answer:2 5 6
\end {large}}
%end
{\smallskip \begin {small} Затверджено на засіданні кафедри теоретичної кібернетики, протокол \No 12 від   19.05.2020 р.\par\smallskip\smallskip
{\parbox {6.5 cm}{Зав. кафедри \dotfill Ю.В.~Крак}}\hspace {0.7cm}{\hbox to 6.5 cm{ Екзаменатор \dotfill Т.О.~Карнаух}} \end{small}}
%begin
{{\begin {small}\par
{ \center  Київський національний університет імені Тараса Шевченка \par}
{\parbox {5 cm}{Освітня програма \hfill}{\parbox {7 cm}{Прикладна математика, Системний аналіз \hfill}}\parbox {6 cm}{\hfill Семестр 2}}\par
{\parbox {5 cm} {Навчальний предмет \hfill}\parbox {7 cm}{Програмування \hfill}\parbox {6cm}{\hfill}\par}\vspace{-7pt} \end{small}}{\center Екзаменаційний білет \No \textbf 
{ 497 }\par}}
{
\begin {large}
1. Що буде виведено за виконання?
code
try:
    res = -3**2--1
    if isinstance(res, bool):
        print(f'bool;{res}')
    elif isinstance(res, int):
        print(f'int;{res}')
    elif isinstance(res, float):
        print(f'float;{res:.2f}')
    else:
        print('???')
except:
    print('ERR')
code

2. Що буде виведено за виконання?
code
a = 9
b = 0
c = 2

def f(a):
    a -= 1
    b = 1
    c = 2
    return a + b + c

a = 5
b = 9
c = 4

try:
    print(f(a), a, b, c, sep=';')
except:
    print('ERR', a, b, c, sep=';')
code

3. Що буде виведено за виконання?
code
try:
    j = 8
    while j>4:
        j += -2
    print(j, end=';')
except:
    print('ERR')
code

4. Що буде виведено за виконання?
code
try:
    for f in range(-4, 11, 3):
        if f>=7:
            break
        print(f, end=';')
    else:
        print(f,end=';')
    print(f)
except:
    print('ERR')
code

5. Що буде виведено за виконання?\par (Дійсні значення, якщо наявні, в цьому завданні будуть виводитись у форматі з фіксованою крапкою та, можливо, з доволі великою кількістю десяткових розрядів. У відповіді для дійсних значень у ЦЬОМУ завданні достатньо вказати один десятковий розряд після крапки, відкинувши решту розрядів без округлення. Наприклад, замість 13.66666666666667 достатньо вказати 13.6, замість 25.23 достатньо вказати 25.2)
code
try:
    print(6, end=';')
    print(int('b9'), end=';')
    print(3, end=';')
except Exception:
    print(1, end=';')
except ZeroDivisionError:
    print(2, end=';')
except:
    print('ERR', end=';')
else:
    print(7, end=';')
finally:
    print(8)
code

6. Що буде виведено за виконання?
code
def f(n):
    if n>9:
        f(n//100)
    else:
        print(n%10,end='')

f(987654)
code

7. Що буде виведено за виконання?
code
class A:
    d = 1

    def _h1(self): print(2, end=';')
    def __h2(self): print(3, end=';')

    def main(self):
        try: print(self.d, end=';')
        except: print('ERR', end=';')

        try: self._h1()
        except: print('ERR', end=';')

        try: self.__h2()
        except: print('ERR', end=';')


class B(A):
    d = 4

    def _h1(self): print(5, end=';')
    def __h2(self): print(6, end=';')

try: B().main()
except: print('ERR', end=';')
code

8. Що буде виведено за виконання?
code
def f(arg, n):
    arg[67] = 0
    n *= -3

try:
    n = 6
    d = {52:9, 25:8, 25:5}
    f(d, n)
    print(n, end=';')
    for x in d :
        print(x, sep=';', end=';')
except:
    print('ERR')
code

9. Що буде виведено за виконання?
code
class A:
    b = 1
    a = 2

    def __init__(self):
        b = 3

class B(A):
    b = 4

    def __init__(self):
        super().__init__()
        self.c = 7

obj = B()
obj.b = obj.b + 10
B.b = 13

try: print(obj.b, end= ';')
except: print('ERR', end=';')

try: print(A.b, end= ';')
except: print('ERR', end=';')

try: print(B.b, end= ';')
except: print('ERR', end=';')
code

10. 
Для графа на рисунку з вершини 1 запущено алгоритм пошуку в глибину, що будує кістякове дерево. Списки суміжності вершин переглядаються алгоритмом за зростанням міток вершин.\par
\begin{center}
	\adjustimage{max size={0.9\linewidth}{0.9\paperheight}}
	{../10/Traversals/62.png}
\end{center}
\par Які з наведених ребер входять до складу отриманого кістякового дерева?\par
1) (4, 5)\par
2) (1, 5)\par
3) (0, 1)\par
4) (0, 5)\par
5) (1, 6)\par
6) (0, 4)\par
{\vskip 15pt} У формі вибрати номери відповідних ребер.\par
%answer:1 3 6
\end {large}}
%end
{\smallskip \begin {small} Затверджено на засіданні кафедри теоретичної кібернетики, протокол \No 12 від   19.05.2020 р.\par\smallskip\smallskip
{\parbox {6.5 cm}{Зав. кафедри \dotfill Ю.В.~Крак}}\hspace {0.7cm}{\hbox to 6.5 cm{ Екзаменатор \dotfill Т.О.~Карнаух}} \end{small}}
%begin
{{\begin {small}\par
{ \center  Київський національний університет імені Тараса Шевченка \par}
{\parbox {5 cm}{Освітня програма \hfill}{\parbox {7 cm}{Прикладна математика, Системний аналіз \hfill}}\parbox {6 cm}{\hfill Семестр 2}}\par
{\parbox {5 cm} {Навчальний предмет \hfill}\parbox {7 cm}{Програмування \hfill}\parbox {6cm}{\hfill}\par}\vspace{-7pt} \end{small}}{\center Екзаменаційний білет \No \textbf 
{ 498 }\par}}
{
\begin {large}
1. Що буде виведено за виконання?
code
try:
    res = 9**-0.5*False
    if isinstance(res, bool):
        print(f'bool;{res}')
    elif isinstance(res, int):
        print(f'int;{res}')
    elif isinstance(res, float):
        print(f'float;{res:.2f}')
    else:
        print('???')
except:
    print('ERR')
code

2. Що буде виведено за виконання?
code
a = 3
b = 1
c = 6

def f(b):
    global c
    a = 3
    b += 5
    c = 4
    return a + b + c

a = 8
b = 7
c = 0

try:
    print(f(b), a, b, c, sep=';')
except:
    print('ERR', a, b, c, sep=';')
code

3. Що буде виведено за виконання?
code
try:
    n = 9
    while n>=6:
        n += -2
    print(n, end=';')
except:
    print('ERR')
code

4. Що буде виведено за виконання?
code
try:
    for c in range(-5, -11, -2):
        if c<=8:
            continue
        print(c, end=';');c = 8
    else:
        print(c,end=';')
    print(c)
except:
    print('ERR')
code

5. Що буде виведено за виконання?\par (Дійсні значення, якщо наявні, в цьому завданні будуть виводитись у форматі з фіксованою крапкою та, можливо, з доволі великою кількістю десяткових розрядів. У відповіді для дійсних значень у ЦЬОМУ завданні достатньо вказати один десятковий розряд після крапки, відкинувши решту розрядів без округлення. Наприклад, замість 13.66666666666667 достатньо вказати 13.6, замість 25.23 достатньо вказати 25.2)
code
try:
    print(0, end=';')
    print(int('2'), end=';')
    print(6, end=';')
except TypeError:
    print(3, end=';')
except KeyboardInterrupt:
    print(0, end=';')
except:
    print('ERR', end=';')
else:
    print(1, end=';')
finally:
    print(4)
code

6. Що буде виведено за виконання?
code
def f(n):
    print(n%10, end='')
    if n>9:
        f(n//10)
 
f(123456)
code

7. Що буде виведено за виконання?
code
class A:
    b = 1

    def __f1(self): print(2, end=';')
    def _f2(self): print(3, end=';')

    def main(self):
        try: print(self.b, end=';')
        except: print('ERR', end=';')

        try: self.__f1()
        except: print('ERR', end=';')

        try: self._f2()
        except: print('ERR', end=';')


class B(A):
    b = 4

    def __f1(self): print(5, end=';')
    def _f2(self): print(6, end=';')

try: B().main()
except: print('ERR', end=';')
code

8. Що буде виведено за виконання?
code
def f(arg, n):
    arg[65] = 6
    n += -2

try:
    n = 3
    s = {65:1, 48:9, 74:2, 48:6}
    f(s, n)
    print(n, end=';')
    for x in s.values():
        print(x, sep=';', end=';')
except:
    print('ERR')
code

9. Що буде виведено за виконання?
code
class A:
    b = 1
    c = 2

    def __init__(self):
        c = 3

class B(A):
    c = 4

    def __init__(self):
        super().__init__()
        self.a = 7

obj = B()
obj.c = obj.a + 10
B.c = 13

try: print(obj.c, end= ';')
except: print('ERR', end=';')

try: print(A.c, end= ';')
except: print('ERR', end=';')

try: print(B.c, end= ';')
except: print('ERR', end=';')
code

10. 
Для графа на рисунку з вершини 2 запущено алгоритм пошуку в глибину, що будує кістякове дерево. Списки суміжності вершин переглядаються алгоритмом за зростанням міток вершин.\par
\begin{center}
	\adjustimage{max size={0.9\linewidth}{0.9\paperheight}}
	{../10/Traversals/135.png}
\end{center}
\par Які з наведених ребер входять до складу отриманого кістякового дерева?\par
1) (5, 6)\par
2) (4, 5)\par
3) (1, 5)\par
4) (3, 4)\par
5) (3, 6)\par
6) (1, 3)\par
{\vskip 15pt} У формі вибрати номери відповідних ребер.\par
%answer:3 5 6
\end {large}}
%end
{\smallskip \begin {small} Затверджено на засіданні кафедри теоретичної кібернетики, протокол \No 12 від   19.05.2020 р.\par\smallskip\smallskip
{\parbox {6.5 cm}{Зав. кафедри \dotfill Ю.В.~Крак}}\hspace {0.7cm}{\hbox to 6.5 cm{ Екзаменатор \dotfill Т.О.~Карнаух}} \end{small}}
%begin
{{\begin {small}\par
{ \center  Київський національний університет імені Тараса Шевченка \par}
{\parbox {5 cm}{Освітня програма \hfill}{\parbox {7 cm}{Прикладна математика, Системний аналіз \hfill}}\parbox {6 cm}{\hfill Семестр 2}}\par
{\parbox {5 cm} {Навчальний предмет \hfill}\parbox {7 cm}{Програмування \hfill}\parbox {6cm}{\hfill}\par}\vspace{-7pt} \end{small}}{\center Екзаменаційний білет \No \textbf 
{ 499 }\par}}
{
\begin {large}
1. Що буде виведено за виконання?
code
try:
    res = 8.0j>=True
    if isinstance(res, bool):
        print(f'bool;{res}')
    elif isinstance(res, int):
        print(f'int;{res}')
    elif isinstance(res, float):
        print(f'float;{res:.2f}')
    else:
        print('???')
except:
    print('ERR')
code

2. Що буде виведено за виконання?
code
a = 1
b = 4
c = 4

def g(a):
    a = 5
    b = 3
    c = 2
    return a + b + c

a = 8
b = 9
c = 6

try:
    print(g(b), a, b, c, sep=';')
except:
    print('ERR', a, b, c, sep=';')
code

3. Що буде виведено за виконання?
code
try:
    m = 7
    while m>2:
        m += -3
    print(m, end=';')
except:
    print('ERR')
code

4. Що буде виведено за виконання?
code
try:
    for d in range(15, -12, -3):
        if d>4:
            continue
        print(d, end=';');d = -2
    else:
        print('end',end=';')
    print(d)
except:
    print('ERR')
code

5. Що буде виведено за виконання?\par (Дійсні значення, якщо наявні, в цьому завданні будуть виводитись у форматі з фіксованою крапкою та, можливо, з доволі великою кількістю десяткових розрядів. У відповіді для дійсних значень у ЦЬОМУ завданні достатньо вказати один десятковий розряд після крапки, відкинувши решту розрядів без округлення. Наприклад, замість 13.66666666666667 достатньо вказати 13.6, замість 25.23 достатньо вказати 25.2)
code
try:
    print(7, end=';')
    print(5/2, end=';')
    print(9, end=';')
except Exception:
    print(8, end=';')
except ZeroDivisionError:
    print(5, end=';')
except:
    print('ERR', end=';')
else:
    print(9, end=';')
finally:
    print(0)
code

6. Що буде виведено за виконання?
code
def f(s1):
    s = s1[:]
    for i in range(len(s)):
        s[i] += 1
 
s = [1, 2, 3]
f(s)
print(s)
code

7. Що буде виведено за виконання?
code
class A:
    a = 1

    def _h1(self): print(2, end=';')
    def _h2(self): print(3, end=';')

    def main(self):
        try: print(self.a, end=';')
        except: print('ERR', end=';')

        try: self._h1()
        except: print('ERR', end=';')

        try: self._h2()
        except: print('ERR', end=';')


class B(A):
    a = 4

    def _h1(self): print(5, end=';')
    def _h2(self): print(6, end=';')

try: B().main()
except: print('ERR', end=';')
code

8. Що буде виведено за виконання?
code
def f(arg, n):
    arg[89] = 8
    n -= 2

try:
    n = 7
    t = {78:4, 90:2, 89:8, 78:4}
    f(t, n)
    print(n, end=';')
    for x in t.items():
        print(*x, sep=';', end=';')
except:
    print('ERR')
code

9. Що буде виведено за виконання?
code
class A:
    b = 1
    c = 2

    def __init__(self):
        c = 3

class B(A):
    c = 4

    def __init__(self):
        super().__init__()
        self.a = 7

obj = B()
obj.a = obj.b + 10
B.a = 13

try: print(obj.a, end= ';')
except: print('ERR', end=';')

try: print(A.a, end= ';')
except: print('ERR', end=';')

try: print(B.a, end= ';')
except: print('ERR', end=';')
code

10. 
Для графа на рисунку з вершини 0 запущено алгоритм пошуку в глибину, що будує кістякове дерево. Списки суміжності вершин переглядаються алгоритмом за зростанням міток вершин.\par
\begin{center}
	\adjustimage{max size={0.9\linewidth}{0.9\paperheight}}
	{../10/Traversals/228.png}
\end{center}
\par Які з наведених ребер входять до складу отриманого кістякового дерева?\par
1) (0, 1)\par
2) (0, 6)\par
3) (2, 5)\par
4) (2, 6)\par
5) (1, 2)\par
6) (4, 5)\par
{\vskip 15pt} У формі вибрати номери відповідних ребер.\par
%answer:1 3 5 6
\end {large}}
%end
{\smallskip \begin {small} Затверджено на засіданні кафедри теоретичної кібернетики, протокол \No 12 від   19.05.2020 р.\par\smallskip\smallskip
{\parbox {6.5 cm}{Зав. кафедри \dotfill Ю.В.~Крак}}\hspace {0.7cm}{\hbox to 6.5 cm{ Екзаменатор \dotfill Т.О.~Карнаух}} \end{small}}
%begin
{{\begin {small}\par
{ \center  Київський національний університет імені Тараса Шевченка \par}
{\parbox {5 cm}{Освітня програма \hfill}{\parbox {7 cm}{Прикладна математика, Системний аналіз \hfill}}\parbox {6 cm}{\hfill Семестр 2}}\par
{\parbox {5 cm} {Навчальний предмет \hfill}\parbox {7 cm}{Програмування \hfill}\parbox {6cm}{\hfill}\par}\vspace{-7pt} \end{small}}{\center Екзаменаційний білет \No \textbf 
{ 500 }\par}}
{
\begin {large}
1. Що буде виведено за виконання?
code
def f(n):
    print('f', end='')
    return n==-4

try:
    print(f(-6) and f(6))
except:
    print('ERR')
code

2. Що буде виведено за виконання?
code
a = 5
b = 3
c = 7

def h(b):
    global c
    a = 1
    b = 4
    c = 5
    return a + b + c

a = 2
b = 0
c = 1

try:
    print(h(a), a, b, c, sep=';')
except:
    print('ERR', a, b, c, sep=';')
code

3. Що буде виведено за виконання?
code
try:
    k = 5
    while k>=1:
        k += -2
    print(k, end=';')
except:
    print('ERR')
code

4. Що буде виведено за виконання?
code
try:
    for a in range(0, -1, 2):
        if a<8:
            continue
        print(a, end=';')
    else:
        print(a,end=';')
    print(a)
except:
    print('ERR')
code

5. Що буде виведено за виконання?\par (Дійсні значення, якщо наявні, в цьому завданні будуть виводитись у форматі з фіксованою крапкою та, можливо, з доволі великою кількістю десяткових розрядів. У відповіді для дійсних значень у ЦЬОМУ завданні достатньо вказати один десятковий розряд після крапки, відкинувши решту розрядів без округлення. Наприклад, замість 13.66666666666667 достатньо вказати 13.6, замість 25.23 достатньо вказати 25.2)
code
try:
    print(7, end=';')
    print(8j<0j, end=';')
    print(1, end=';')
except TypeError:
    print(6, end=';')
except ValueError:
    print(3, end=';')
except:
    print('ERR', end=';')
else:
    print(2, end=';')
finally:
    print(4)
code

6. Що буде виведено за виконання?
code
def f(n):
    if n<=9:
        print(n%10, end='')
    else:
        f(n//10)

f(987651)
code

7. Що буде виведено за виконання?
code
class A:
    c = 1

    def f1(self): print(2, end=';')
    def _f2(self): print(3, end=';')

    def main(self):
        try: print(self.c, end=';')
        except: print('ERR', end=';')

        try: self.f1()
        except: print('ERR', end=';')

        try: self._f2()
        except: print('ERR', end=';')


class B(A):
    c = 4

    def f1(self): print(5, end=';')
    def _f2(self): print(6, end=';')

try: B().main()
except: print('ERR', end=';')
code

8. Що буде виведено за виконання?
code
def f(arg, n):
    arg[15] = 6
    n = 9
    arg = set(arg.keys())

try:
    n = 3
    t = {31:9, 15:3, 86:5, 31:3}
    f(t, n)
    print(n, end=';')
    for x, y in t.items():
        print(x, y, sep=';', end=';')
except:
    print('ERR')
code

9. Що буде виведено за виконання?
code
class A:
    c = 1
    b = 2

    def __init__(self):
        c = 3

class B(A):
    c = 4

    def __init__(self):
        super().__init__()
        self.a = 7

obj = B()
obj.c = obj.c + 10
B.c = 13

try: print(obj.c, end= ';')
except: print('ERR', end=';')

try: print(A.c, end= ';')
except: print('ERR', end=';')

try: print(B.c, end= ';')
except: print('ERR', end=';')
code

10. 
Для графа на рисунку з вершини 4 запущено алгоритм пошуку в глибину, що будує кістякове дерево. Списки суміжності вершин переглядаються алгоритмом за зростанням міток вершин.\par
\begin{center}
	\adjustimage{max size={0.9\linewidth}{0.9\paperheight}}
	{../10/Traversals/4.png}
\end{center}
\par Які з наведених ребер входять до складу отриманого кістякового дерева?\par
1) (0, 6)\par
2) (2, 3)\par
3) (1, 6)\par
4) (1, 3)\par
5) (0, 2)\par
6) (0, 3)\par
{\vskip 15pt} У формі вибрати номери відповідних ребер.\par
%answer:3 4 5 6
\end {large}}
%end
{\smallskip \begin {small} Затверджено на засіданні кафедри теоретичної кібернетики, протокол \No 12 від   19.05.2020 р.\par\smallskip\smallskip
{\parbox {6.5 cm}{Зав. кафедри \dotfill Ю.В.~Крак}}\hspace {0.7cm}{\hbox to 6.5 cm{ Екзаменатор \dotfill Т.О.~Карнаух}} \end{small}}
\end{document}
